\documentclass[11pt]{article}

\usepackage[margin=1in]{geometry}
\usepackage{amsmath,amssymb,amsthm,mathtools,mathrsfs,bm}
\usepackage{enumitem}
\usepackage{booktabs,longtable,tabularx,array}
\usepackage{xcolor}
\usepackage{microtype}
\usepackage[colorlinks=true,linkcolor=blue,citecolor=blue,urlcolor=blue]{hyperref}
\usepackage[nameinlink,capitalise]{cleveref}
\setlength{\emergencystretch}{4em}
\raggedbottom
\hbadness=10000

\newtheorem{theorem}{Theorem}[section]
\newtheorem{proposition}[theorem]{Proposition}
\newtheorem{lemma}[theorem]{Lemma}
\newtheorem{corollary}[theorem]{Corollary}
\newtheorem{countertheorem}[theorem]{Countertheorem}
\newtheorem{openproblem}[theorem]{Open Problem}
\theoremstyle{definition}
\newtheorem{definition}[theorem]{Definition}
\newtheorem{assumption}[theorem]{Assumption}
\newtheorem{programme}[theorem]{Programme}
\newtheorem{construction}[theorem]{Construction}
\theoremstyle{remark}
\newtheorem{remark}[theorem]{Remark}
\newtheorem{assessment}[theorem]{Assessment}
\newtheorem{firewall}[theorem]{Firewall}

\newcommand{\R}{\mathbb R}
\newcommand{\C}{\mathbb C}
\newcommand{\Torus}{\mathbb T}
\newcommand{\N}{\mathbb N}
\newcommand{\dd}{\mathrm d}
\newcommand{\PP}{\mathbb P}
\newcommand{\BNS}{\mathcal B}
\newcommand{\Ker}{\operatorname{Ker}}
\newcommand{\Ran}{\operatorname{Ran}}
\newcommand{\Span}{\operatorname{span}}
\newcommand{\supp}{\operatorname{supp}}
\newcommand{\dist}{\operatorname{dist}}
\newcommand{\esssup}{\operatorname*{ess\,sup}}
\newcommand{\ip}[2]{\left\langle #1,#2\right\rangle}
\newcommand{\norm}[1]{\left\lVert #1\right\rVert}
\newcommand{\abs}[1]{\left\lvert #1\right\rvert}
\newcommand{\one}{\mathbf 1}

\definecolor{deepgreen}{RGB}{25,110,65}
\definecolor{deepblue}{RGB}{35,75,145}
\definecolor{deepred}{RGB}{150,45,45}
\definecolor{deeporange}{RGB}{170,95,15}
\newcommand{\PROVED}{\textcolor{deepgreen}{\textbf{PROVED}}}
\newcommand{\INHERITED}{\textcolor{deepblue}{\textbf{INHERITED}}}
\newcommand{\CONDITIONAL}{\textcolor{deeporange}{\textbf{CONDITIONAL}}}
\newcommand{\OPEN}{\textcolor{deepred}{\textbf{OPEN}}}
\newcommand{\OBSTRUCTION}{\textcolor{deepred}{\textbf{OBSTRUCTION}}}

\title{\textbf{Renewal NS Terminal Source Concentration and\\
Singular-Thread Closure Programme}\\[0.45em]
\large NS--A, Version V1: Matched Source-Modulus Saturation,\\
the Cubic Normalization Audit, and the Corrected Thread Alternative}
\author{A.~P\'elissier}
\date{2 September 2026}

\begin{document}
\maketitle

\begin{abstract}
This note starts the append-only NS--A programme.  It does not reopen the
source-identification, pressure/Hodge, caloric-envelope, Dini, localization,
or donor calculations already completed in the dedicated Navier--Stokes
series.  The inherited endpoint is the following narrow chain:
\[
 \boxed{
 \begin{gathered}
 \text{terminal source modulus}
 \longrightarrow
 \text{singular-set concentration}\\
 \longrightarrow
 \text{thread or microdelocalized survivor}
 \end{gathered}}
\]
The first task is to audit the normalization at the interface between the
last two arrows.

For the signed longitudinal increment measure used by the inherited
backward-caloric contact, a dyadic parabolic cell of physical scale
\(\ell_k\) becomes a fixed cylinder under the unit-viscosity rescaling
\[
 U_k(s,y)=\frac{\ell_k}{\nu}
 u\!\left(b_k+\frac{\ell_k^2}{\nu}s,x_k+\ell_k y\right).
\]
Version V1 proves the exact scale factor which was missing from the previous
heavy-thread handoff: the normalized cell contact is bounded by
\[
 C\frac{\nu^2}{r_k}
 \int_{Q^\sharp}|U_k|^3,
 \qquad
 r_k=\|A^{1/4}u(b_k)\|_2\longrightarrow\infty.
\]
Consequently a fixed coherent payment in one cell forces its dimensionless
local cubic packet to grow at least like \(r_k/\nu^2\).  The usual uniform
local-energy and local-dissipation packet therefore cannot hold on that same
thread.  In particular, strong local \(L^3\) compactness of the
unit-viscosity rescalings cannot carry the fixed \(r_k^{-1}\)-normalized
contact to a nonzero ancient Navier--Stokes profile.  The compactness theorem
itself remains valid when its hypotheses hold; what fails is the asserted
nontriviality transfer from this particular scalar.

The contact-preserving amplitude normalization is instead
\(
ho_k=(r_k/\nu^2)^{1/3}\), \(W_k=U_k/\rho_k\).  It converts the rescaled
equation into a strong-coupling equation, or after the fast-time change into
a vanishing-viscosity equation.  Thus the normalized failure object is an
Euler/Euler--Reynolds candidate, or a further concentration defect, not
without additional input a compact unit-viscosity ancient solution.

On a source-aligned heavy thread, the same cell also forces
\(M_k(\tau_k)\gtrsim\sqrt{\tau_k}\) at its matched terminal age.  Hence it
realizes the critical square-root obstruction to the inherited Dini
criterion; failure of this conclusion is an explicit local
source-incidence/router residual.  On the microdelocalized branch, the sum of
the dimensionless cubic capacities of all near-singular cells is at least
\(c r_k/\nu^2\).  This yields a precise alternative between a cubic-heavy
thread and genuine capacity spreading through at least order \(r_k\) local
cells, together with a conditional singular-set packing gate.  These are
new localization consequences for the already occurring scalar, not new
source prices.  No global three-dimensional regularity conclusion is
claimed.
\end{abstract}

\tableofcontents

% =============================================================================
\section{Context, frozen inputs, and append-only protocol}
\label{sec:context}
% =============================================================================

\subsection{Why this fresh series starts after V35}

The attached programme architecture separates the remaining
Navier--Stokes work into two logically different tasks.  NS--A concerns only
terminal concentration on the singular set.  NS--B concerns the already
reduced fixed-datum Reynolds covariance and paid-stress packet.  The present
file is NS--A.  Its purpose is not to develop another ambient source theory,
but to decide what the terminal source-paid scalar can actually produce
after the V35 localization theorem.

The following conclusions are treated as frozen theorem layers.
\begin{enumerate}[label=\textup{(F\arabic*)},leftmargin=3.8em]
\item The unforced periodic nonlinear occurrence is the projected divergence
      of \(u\otimes u\).  Pressure and the equivalent Lamb, vorticity,
      stretching, and work descriptions are followers on the complete
      history, not additional source births.
\item On a unit critical-norm first-passage chronology, the energy-neutral
      backward-caloric contact has a fixed lower margin while the native
      source length
      \(
       \int\|A^{-1/4}\BNS(u)\|_2\,\dd t
      \)
      tends to zero packetwise.
\item The source-attached gain has the sharp age envelope
      \(
       C[\nu(b_k-t)]^{-1/2}
      \), and the terminal source modulus
      \(
       M_k(\tau)=\int_{b_k-\tau}^{b_k}
                  \|A^{-1/4}\BNS(u)(t)\|_2\,\dd t
      \)
      yields an exact gain-tail estimate.
\item The inherited Dini condition
      \[
       \lim_{\varepsilon\downarrow0}\limsup_{k\to\infty}
       \left[
        \varepsilon^{-1/2}M_k(\varepsilon)
        +\int_0^\varepsilon M_k(\tau)\tau^{-3/2}\,\dd\tau
       \right]=0
      \]
      makes the selected contact vanish.  A uniform power improvement over
      \(\sqrt\tau\) suffices, whereas an abstract critical
      \(O(\sqrt\tau)\) modulus is sharply insufficient.
\item The absolute longitudinal carrier vanishes on every compact terminally
      regular set.  Every cluster of the coherent contact-centre law is
      supported on the terminal singular set \(\Sigma_T\).
\item The surviving spatial law has the exact trichotomy: singular-distance/
      scale separation, one heavy near-singular parabolic cell, or
      near-singular microdelocalization.
\item The fixed-datum donor programme has already compressed the low-output
      residue to the five-dimensional path
      \(
       (\theta_{\rm c}\mathbb K[u])^\circ\in\operatorname{Sym}_0(3)
      \)
      together with a positive Reynolds-covariance square function.  That
      packet is opened only after NS--A selects a genuine surviving thread or
      a source-relative microdelocalized Gram.
\end{enumerate}

\begin{firewall}[No duplicated theorem layer]
\label{fire:no-duplication}
This file does not re-prove the V35 uniform-integrability theorem, the
terminal Dini tail estimate, regular-region vanishing, singular-set support,
the measure-theoretic thread/microdelocalization trichotomy, the abstract
Aubin--Lions compactness implication, the rank-at-most-two \(2\mathrm D3\mathrm C\)
regularity handoff, or the V13 donor table.  It begins with a normalization
audit of the physical scalar on the V33--V35 parabolic cells.  Any statement
below which uses an inherited row names that row explicitly.
\end{firewall}

\subsection{Frozen source map}

The first version was prepared against the following exact attached files.
The hashes are recorded only to freeze the reading base; they are not
mathematical hypotheses.
\begin{center}
\begingroup
\scriptsize
\begin{tabularx}{\textwidth}{>{\raggedright\arraybackslash}p{0.48\textwidth}X}
\toprule
File & SHA--256 \\
\midrule
\path{Pasted markdown(20260902-235348).md}
& \texttt{44ccab676db008dbd9460f1518d432ef0}\newline\texttt{eb0ea08a4599cf1e3bac5155f72170b}\\[0.35em]
\path{Pasted text (2)(20260902-235711).txt}
& \texttt{ada139e74cc525a13a5dce3d7ca26f3}\newline\texttt{ce2c9286a922d1df9be739c726bc3f989}\\[0.35em]
\path{Renewal_Dedicated_Navier_Stokes_Source_Relative_Insertion_Carrier_Motion_and_Regularity_Closure_Programme_V35(3).tex}
& \texttt{1a6ac99acdff7055ca046166a90cabaf0}\newline\texttt{1180e89ca90dcd9e046410c88d427ea}\\[0.35em]
\path{Renewal_NS_Rank_Three_Quadratic_Donor_and_Critical_Dipole_Programme_V13(3).tex}
& \texttt{7796896b084bada8e30abdccaae81860}\newline\texttt{1342aed1a26ea5301c0eb382e6e6ea6e}\\
\bottomrule
\end{tabularx}
\endgroup
\end{center}

\subsection{Append-only rule}

This file is the complete Version~V1 source.  On every later iteration the
complete source is to be preserved verbatim, the unique final
\verb|\end{document}| removed, and the new material appended immediately
before a restored final \verb|\end{document}|.  The filename is then
incremented to V2, V3, and so on.  A row declared closed in the proof-status
ledger may be reopened only by a theorem invalidating one of its displayed
hypotheses.

\begin{programme}[Termination semantics]
\label{prog:termination}
Every phase of NS--A must end in one of three forms:
\[
 \boxed{
 \text{PASS}
 \quad\big|\quad
 \text{TYPED OBSTRUCTION}
 \quad\big|\quad
 \text{STOP AND EXPORT}.}
\]
Here PASS means that the selected scalar vanishes and the inherited critical
continuation theorem is invoked.  A typed obstruction is one of the exact
source-modulus, local-energy, dissipation, incidence, packing, or
strong-coupling objects constructed below.  STOP AND EXPORT means that one
rank-three fixed-datum thread has been selected and is sent to NS--B.  There
is no fourth outcome called additional general compiler machinery.
\end{programme}

% =============================================================================
\section{Frozen caloric contact and the canonical parabolic cells}
\label{sec:frozen-contact}
% =============================================================================

Let \(u\) be a smooth mean-zero divergence-free solution of the unforced
periodic Navier--Stokes equation on \([0,T)\times\Torus^3\):
\begin{equation}
 \partial_tu+\BNS(u)+\nu Au=0,
 \qquad
 A=-\PP\Delta,
 \qquad
 \BNS(u)=\PP\operatorname{div}(u\otimes u),
 \label{eq:NS}
\end{equation}
where \(\nu>0\) is fixed.  Retain the inherited disjoint unit
critical-norm first-passage intervals
\(
 I_k=(a_k,b_k]
\)
and put
\begin{equation}
 r_k:=\|A^{1/4}u(b_k)\|_2,
 \qquad r_k\longrightarrow\infty.
 \label{eq:rk}
\end{equation}
For \(t<b_k\), set
\begin{equation}
 a_k(t):=2\nu(b_k-t).
 \label{eq:heat-age}
\end{equation}

Let \(L_a\) be the mean-zero periodic convolution kernel of
\(\Lambda e^{-aA}\), with \(\Lambda=A^{1/2}\).  The inherited longitudinal
increment identity defines the signed measure
\begin{equation}
 \boxed{
 \begin{aligned}
 \eta_k(\dd t,\dd x,\dd z)
 :={}&-\frac1{4r_k}
 \nabla L_{a_k(t)}(z)\cdot\delta_zu(x,t)\\
 &\qquad\qquad\times|\delta_zu(x,t)|^2
 \,\dd z\,\dd x\,\dd t,
 \end{aligned}}
 \label{eq:eta}
\end{equation}
where \(\delta_zu(x)=u(x+z)-u(x)\).  Its total mass is the inherited
held-out backward-caloric scalar.  We work on a subsequence on which
\begin{equation}
 h_k:=\eta_k(\Omega_k)>0,
 \qquad
 h_k\ge h_*>0.
 \label{eq:fixed-contact}
\end{equation}
Changing the global sign if necessary causes no loss.

Write \(\eta_k=\eta_{k,+}-\eta_{k,-}\), put
\begin{equation}
 P_k:=\eta_{k,+}(\Omega_k),
 \qquad
 \kappa_k:=\frac{h_k}{P_k}\eta_{k,+},
 \qquad
 \pi_k:=\frac{\kappa_k}{h_k}=\frac{\eta_{k,+}}{P_k}.
 \label{eq:coherent-replay}
\end{equation}
Since \(P_k-\eta_{k,-}(\Omega_k)=h_k\), one has \(P_k\ge h_k\), and hence
\begin{equation}
 \boxed{
 0\le\kappa_k\le\eta_{k,+}\le|\eta_k|.}
 \label{eq:payment-dominated}
\end{equation}
This elementary domination is the only property of the coherent replay
needed for the normalization audit.

The inherited critical kernel estimate is
\begin{equation}
 \boxed{
 |\nabla L_a(z)|\le C_0
 \left[
  \one_{\{0<d(z,0)\le1\}}d(z,0)^{-5}
  +\one_{\{d(z,0)>1\}}
 \right]}
 \label{eq:kernel-bound}
\end{equation}
for all \(a\ge0\).  We use it as a frozen input.

Fix the inherited parabolic annulus constants
\begin{equation}
 0<\alpha<\beta<\infty.
 \label{eq:annulus-constants}
\end{equation}
For an integer \(n\), write \(\ell=2^{-n}\), let \(Q\) be a periodic dyadic
cube of side \(\ell\), and define
\begin{equation}
 \boxed{
 \begin{aligned}
 \mathcal C_{k,n,Q}:=\bigl\{(t,x,z):{}&
 \ell^2/4<a_k(t)\le\ell^2,\quad x\in Q,\\
 &\alpha\sqrt{a_k(t)}
 \le d(z,0)\le
 \beta\sqrt{a_k(t)}\bigr\}. 
 \end{aligned}}
 \label{eq:parabolic-cell}
\end{equation}
For late small cells choose a lift of \(Q\) to \(\R^3\) and a point
\(x_{k,n,Q}\) in that lift.  Define the unit-viscosity rescaling
\begin{equation}
 \boxed{
 U_{k,n,Q}(s,y)
 :=\frac{\ell}{\nu}
 u\!\left(
  b_k+\frac{\ell^2}{\nu}s,
  x_{k,n,Q}+\ell y
 \right).}
 \label{eq:unit-viscosity-rescaling}
\end{equation}
The rescaled pressure is
\(
 \Pi_{k,n,Q}=\ell^2p/\nu^2
\), and \((U_{k,n,Q},\Pi_{k,n,Q})\) solves the
unit-viscosity equation on its rescaled domain.

\begin{lemma}[Fixed geometry of one canonical cell]
\label{lem:fixed-geometry}
On \(\mathcal C_{k,n,Q}\), the variables
\begin{equation}
 s=\frac{\nu(t-b_k)}{\ell^2},
 \qquad
 y=\frac{x-x_{k,n,Q}}{\ell},
 \qquad
 \zeta=\frac z\ell
 \label{eq:scaled-variables}
\end{equation}
satisfy
\begin{equation}
 -\frac12\le s<-\frac18,
 \qquad
 y\in\widehat Q,
 \qquad
 \frac\alpha2<|\zeta|\le\beta,
 \label{eq:fixed-cell-geometry}
\end{equation}
where \(\widehat Q\) is a unit cube.  Moreover
\begin{equation}
 \dd t\,\dd x\,\dd z
 =\frac{\ell^8}{\nu}\,\dd s\,\dd y\,\dd\zeta,
 \qquad
 \delta_zu=\frac\nu\ell\delta_\zeta U_{k,n,Q}.
 \label{eq:jacobian}
\end{equation}
All constants are independent of \(k,n,Q\).
\end{lemma}

\begin{proof}
The age condition in \eqref{eq:parabolic-cell} is
\(
 \ell^2/4<-2\ell^2s\le\ell^2
\), because
\(
 a_k(t)=2\nu(b_k-t)=-2\ell^2s
\).  This is exactly the displayed interval for \(s\).  Hence
\(
 \sqrt{a_k(t)}/\ell=\sqrt{-2s}\in(1/2,1]
\), and the separation condition gives the displayed fixed annulus for
\(\zeta\).  The cube rescales to a unit cube.  The Jacobian and increment
identities follow directly from \eqref{eq:scaled-variables} and
\eqref{eq:unit-viscosity-rescaling}.
\end{proof}

\begin{definition}[Dimensionless local cubic packets]
\label{def:cubic-packets}
For a cell \(\gamma=(k,n,Q)\), let \(D_\gamma\) be its image in
\((s,y,\zeta)\)-coordinates and define
\begin{align}
 \mathscr I_\gamma
 &:={}
 \int_{D_\gamma}
 |\delta_\zeta U_\gamma(s,y)|^3
 \,\dd s\,\dd y\,\dd\zeta,
 \label{eq:increment-packet}\\
 \mathscr X_\gamma
 &:={}
 \int_{-1/2}^{-1/8}
 \int_{\widehat Q+B_\beta}
 |U_\gamma(s,y)|^3\,\dd y\,\dd s,
 \label{eq:velocity-packet}
\end{align}
where \(B_\beta\) is the Euclidean ball of radius \(\beta\).  These are
scale-invariant quantities on the unit-viscosity cell.  They are not new
source actions.
\end{definition}

% =============================================================================
\section{The exact cubic normalization audit}
\label{sec:cubic-audit}
% =============================================================================

\begin{proposition}[Exact principal-part scale factor]
\label{prop:principal-scale-factor}
Let \(\mathscr K\) be the Euclidean vector kernel in the inherited periodic
principal-part decomposition
\begin{equation}
 \nabla L_a(z)=a^{-5/2}\mathscr K(z/\sqrt a)+\mathscr R_a(z).
 \label{eq:principal-decomposition}
\end{equation}
Let \(\eta_{k}^{\rm prin}\) denote the part of \eqref{eq:eta} obtained by
retaining the first term in \eqref{eq:principal-decomposition}.  Then on
one canonical cell \(\gamma=(k,n,Q)\),
\begin{equation}
 \boxed{
 \eta_{k}^{\rm prin}(\mathcal C_{k,n,Q})
 =\frac{\nu^2}{r_k}\,\mathfrak C(U_\gamma),}
 \label{eq:exact-principal-scaling}
\end{equation}
where
\begin{equation}
 \boxed{
 \begin{aligned}
 \mathfrak C(U):={}&-\frac14
 \int_{D_\gamma}
 (-2s)^{-5/2}
 \mathscr K\!\left(\frac\zeta{\sqrt{-2s}}\right)\\
 &\qquad\qquad\cdot\delta_\zeta U(s,y)
 |\delta_\zeta U(s,y)|^2
 \,\dd s\,\dd y\,\dd\zeta.
 \end{aligned}}
 \label{eq:dimensionless-contact}
\end{equation}
In particular, the inherited scalar is not invariant under the
unit-viscosity blow-up: it carries the explicit prefactor \(\nu^2/r_k\).
\end{proposition}

\begin{proof}
On the cell, \(a_k(t)=-2\ell^2s\) and \(z=\ell\zeta\), so
\[
 a_k(t)^{-5/2}\mathscr K(z/\sqrt{a_k(t)})
 =\ell^{-5}(-2s)^{-5/2}
  \mathscr K\!\left(\frac\zeta{\sqrt{-2s}}\right).
\]
By \eqref{eq:jacobian}, the cubic increment contributes
\(\nu^3\ell^{-3}\), while the full Jacobian contributes
\(\ell^8/\nu\).  Hence
\[
 \ell^{-5}\cdot\nu^3\ell^{-3}\cdot\frac{\ell^8}{\nu}
 =\nu^2.
\]
Multiplication by the normalization \((4r_k)^{-1}\) gives
\eqref{eq:exact-principal-scaling} and \eqref{eq:dimensionless-contact}.
\end{proof}

\begin{theorem}[Full periodic cell variation bound]
\label{thm:full-cell-bound}
There is a constant \(C_{\alpha,\beta}\), depending only on the periodic
cell and the fixed annulus constants, such that every sufficiently small
canonical cell satisfies
\begin{equation}
 \boxed{
 |\eta_k|(\mathcal C_{k,n,Q})
 \le
 C_{\alpha,\beta}\frac{\nu^2}{r_k}\mathscr I_{k,n,Q}.}
 \label{eq:full-cell-increment-bound}
\end{equation}
Moreover,
\begin{equation}
 \boxed{
 \mathscr I_{k,n,Q}
 \le C_\beta\mathscr X_{k,n,Q},}
 \label{eq:increment-to-velocity}
\end{equation}
and consequently
\begin{equation}
 \boxed{
 |\eta_k|(\mathcal C_{k,n,Q})
 \le C\frac{\nu^2}{r_k}\mathscr X_{k,n,Q}.}
 \label{eq:full-cell-velocity-bound}
\end{equation}
No periodic-remainder or signed-counterflow hypothesis is needed for these
upper bounds.
\end{theorem}

\begin{proof}
On a sufficiently small cell, \(d(z,0)\le\beta\ell\le1\), while
\(d(z,0)\ge(\alpha/2)\ell\) by \Cref{lem:fixed-geometry}.  Therefore the
full periodic estimate \eqref{eq:kernel-bound} gives
\[
 |\nabla L_{a_k(t)}(z)|\le C_{\alpha,\beta}\ell^{-5}.
\]
Using \eqref{eq:eta}, dropping the longitudinal alignment factor, and then
using \eqref{eq:jacobian},
\[
 \begin{aligned}
 |\eta_k|(\mathcal C_{k,n,Q})
 &\le\frac{C}{r_k}\ell^{-5}
   \int_{\mathcal C_{k,n,Q}}|\delta_zu|^3
   \,\dd t\,\dd x\,\dd z\\
 &=\frac{C}{r_k}\ell^{-5}
   \left(\frac{\nu^3}{\ell^3}\right)
   \left(\frac{\ell^8}{\nu}\right)
   \int_{D_\gamma}|\delta_\zeta U_\gamma|^3
   \,\dd s\,\dd y\,\dd\zeta\\
 &=C\frac{\nu^2}{r_k}\mathscr I_\gamma.
 \end{aligned}
\]
This proves \eqref{eq:full-cell-increment-bound}.

The elementary inequality
\(
 |v-w|^3\le4(|v|^3+|w|^3)
\)
gives, for every admissible \(\zeta\),
\[
 \int_{\widehat Q}|\delta_\zeta U(s,y)|^3\,\dd y
 \le8\int_{\widehat Q+B_\beta}|U(s,y)|^3\,\dd y.
\]
The \(\zeta\)-domain has uniformly bounded volume and the \(s\)-interval is
fixed.  Integration proves \eqref{eq:increment-to-velocity}, and the last
bound follows.
\end{proof}

\begin{theorem}[A heavy coherent cell forces cubic divergence]
\label{thm:heavy-cell-cubic-divergence}
Assume that canonical cells \(\gamma_k=(k,n_k,Q_k)\) satisfy
\begin{equation}
 \pi_k(\mathcal C_{\gamma_k})\ge p_*>0.
 \label{eq:heavy-payment-cell}
\end{equation}
Then
\begin{equation}
 \boxed{
 \mathscr I_{\gamma_k}
 \ge c\frac{h_*p_*}{\nu^2}r_k,
 \qquad
 \mathscr X_{\gamma_k}
 \ge c\frac{h_*p_*}{\nu^2}r_k.}
 \label{eq:heavy-cubic-lower}
\end{equation}
In particular,
\begin{equation}
 \mathscr X_{\gamma_k}\longrightarrow\infty.
 \label{eq:local-L3-divergence}
\end{equation}
This conclusion uses the actual coherent payment and the full periodic
kernel.  It does not assume that the local signed counterflow vanishes.
\end{theorem}

\begin{proof}
By \eqref{eq:coherent-replay} and \eqref{eq:heavy-payment-cell},
\[
 \kappa_k(\mathcal C_{\gamma_k})
 =h_k\pi_k(\mathcal C_{\gamma_k})
 \ge h_*p_*.
\]
The domination \eqref{eq:payment-dominated} gives
\(
 h_*p_*\le|\eta_k|(\mathcal C_{\gamma_k})
\).  Apply \eqref{eq:full-cell-increment-bound} and
\eqref{eq:increment-to-velocity} and rearrange.
\end{proof}

\begin{assessment}[The missing normalization factor]
\label{ass:missing-factor}
The fixed cell probability is a probability under the coherent payment law,
not a scale-invariant lower bound for the unit-viscosity velocity itself.
The signed measure from which that law is formed contains the factor
\(r_k^{-1}\).  Under the parabolic Navier--Stokes rescaling all powers of
\(\ell_k\) cancel, but \(r_k^{-1}\) remains.  Therefore a fixed cell payment
requires the rescaled cubic packet to grow like \(r_k\).  Omitting this
factor turns a concentration sequence into an apparent compact nonzero
profile.
\end{assessment}

% =============================================================================
\section{Local energy and dissipation escalation}
\label{sec:energy-escalation}
% =============================================================================

For a fixed unit cube \(\widehat Q\), choose bounded Lipschitz domains
\[
 \widehat Q+B_\beta\Subset G_0\Subset G_1
\]
inside the rescaled periodic cover.  For a rescaled field \(U\), define
\begin{align}
 \mathscr E[U]
 &:={}
 \esssup_{-1/2<s<-1/8}
 \|U(s)\|_{L^2(G_1)}^2,
 \label{eq:local-energy}\\
 \mathscr D[U]
 &:={}
 \int_{-1/2}^{-1/8}
 \left(
  \|\nabla U(s)\|_{L^2(G_1)}^2
  +\|U(s)\|_{L^2(G_1)}^2
 \right)\,\dd s.
 \label{eq:local-dissipation}
\end{align}
The lower-order term is included only to make the local Sobolev estimate
independent of boundary conditions.

\begin{lemma}[Scale-invariant local cubic interpolation]
\label{lem:local-cubic-interpolation}
There is a constant \(C_{G_0,G_1}\) such that
\begin{equation}
 \boxed{
 \int_{-1/2}^{-1/8}\int_{G_0}|U|^3\,\dd y\,\dd s
 \le
 C_{G_0,G_1}
 \mathscr E[U]^{3/4}\mathscr D[U]^{3/4}.}
 \label{eq:local-cubic-interpolation}
\end{equation}
\end{lemma}

\begin{proof}
Choose \(\chi\in C_c^\infty(G_1)\) with \(\chi=1\) on \(G_0\).  The
Sobolev inequality on \(\R^3\) gives
\[
 \|U(s)\|_{L^6(G_0)}
 \le\|\chi U(s)\|_6
 \le C\left(
  \|\nabla U(s)\|_{L^2(G_1)}
  +\|U(s)\|_{L^2(G_1)}
 \right).
\]
Interpolation between \(L^2\) and \(L^6\) yields
\[
 \|U(s)\|_{L^3(G_0)}^3
 \le
 \|U(s)\|_{L^2(G_0)}^{3/2}
 \|U(s)\|_{L^6(G_0)}^{3/2}.
\]
Take the \(L^\infty_sL^2_y\) factor outside the time integral and apply
H\"older with exponents \(4/3\) and \(4\) to the remaining
\(L^6\)-factor.  Since the time interval has fixed length, this gives
\eqref{eq:local-cubic-interpolation}.
\end{proof}

\begin{theorem}[Quantitative local-energy/dissipation escalation]
\label{thm:energy-dissipation-escalation}
Under the heavy-cell hypothesis of
\Cref{thm:heavy-cell-cubic-divergence}, let
\(
 U_k=U_{\gamma_k}
\).  Then
\begin{equation}
 \boxed{
 \mathscr E[U_k]\mathscr D[U_k]
 \ge
 c\left(\frac{h_*p_*}{\nu^2}r_k\right)^{4/3}.}
 \label{eq:ED-product-lower}
\end{equation}
Consequently
\begin{equation}
 \boxed{
 \max\{\mathscr E[U_k],\mathscr D[U_k]\}
 \ge
 c\left(\frac{h_*p_*}{\nu^2}r_k\right)^{2/3}
 \longrightarrow\infty.}
 \label{eq:ED-max-lower}
\end{equation}
Thus every heavy coherent thread carries a quantitative failure of the
standard scale-invariant local compactness packet.
\end{theorem}

\begin{proof}
Choose \(G_0\) containing \(\widehat Q+B_\beta\).  The lower bound
\eqref{eq:heavy-cubic-lower} and
\Cref{lem:local-cubic-interpolation} give
\[
 c\frac{h_*p_*}{\nu^2}r_k
 \le C\mathscr E[U_k]^{3/4}\mathscr D[U_k]^{3/4}.
\]
Raise both sides to the power \(4/3\) to obtain
\eqref{eq:ED-product-lower}.  Since
\(
 \max\{a,b\}\ge\sqrt{ab}
\)
for nonnegative \(a,b\), \eqref{eq:ED-max-lower} follows.
\end{proof}

\begin{remark}[Original-variable meaning]
\label{rem:original-variable-meaning}
Up to fixed enlargement constants, the two dimensionless quantities are
\begin{align}
 \mathscr E[U_k]
 &\asymp
 \esssup_{t\in I_{k,n_k}}
 \frac1{\nu^2\ell_k}
 \int_{B_{C\ell_k}(x_k)}|u(t,x)|^2\,\dd x,
 \label{eq:original-energy}\\
 \int|\nabla U_k|^2
 &\asymp
 \frac1{\nu\ell_k}
 \int_{I_{k,n_k}}
 \int_{B_{C\ell_k}(x_k)}|\nabla u|^2\,\dd x\,\dd t.
 \label{eq:original-dissipation}
\end{align}
Hence \Cref{thm:energy-dissipation-escalation} is a terminal rate statement
for the usual scale-invariant local kinetic-energy and dissipation packets.
It does not introduce a new norm or a new source clock.
\end{remark}

\begin{countertheorem}[The compact ancient-profile handoff cannot carry this fixed contact]
\label{cth:compact-handoff}
There is no sequence of canonical cells for which all three properties hold:
\begin{enumerate}[label=\textup{(C\arabic*)},leftmargin=3.4em]
\item \(r_k\to\infty\) and the total coherent scalar satisfies
      \(h_k\ge h_*>0\);
\item the cell has fixed coherent probability
      \(\pi_k(\mathcal C_{\gamma_k})\ge p_*>0\);
\item the unit-viscosity rescalings obey a uniform local packet
      \[
       \sup_k\left[
        \|U_k\|_{L^\infty_sL^2_y(G_1)}
        +\|\nabla U_k\|_{L^2_{s,y}(G_1)}
       \right]<\infty.
      \]
\end{enumerate}
In particular, a subsequence satisfying the displayed uniform packet cannot
converge strongly in local \(L^3\) while retaining a nonzero amount of the
\(r_k^{-1}\)-normalized coherent contact in the selected cell.
\end{countertheorem}

\begin{proof}
Property \textup{(C3)} makes both \(\mathscr E[U_k]\) and
\(\mathscr D[U_k]\) uniformly bounded.  This contradicts
\eqref{eq:ED-max-lower}.  Equivalently,
\Cref{lem:local-cubic-interpolation} makes \(\mathscr X_{\gamma_k}\)
uniformly bounded, while \eqref{eq:full-cell-velocity-bound} then gives
\[
 \kappa_k(\mathcal C_{\gamma_k})
 \le|\eta_k|(\mathcal C_{\gamma_k})
 \le C\frac{\nu^2}{r_k}\mathscr X_{\gamma_k}
 \longrightarrow0,
\]
contradicting
\(
 \kappa_k(\mathcal C_{\gamma_k})\ge h_*p_*
\).
\end{proof}

\begin{assessment}[Precise correction to the inherited V32--V35 thread conclusion]
\label{ass:thread-correction}
The inherited compactness implication remains mathematically correct as an
implication from a uniform local energy, gradient, and pressure packet: such
a packet can yield an ancient suitable limit.  The correction is narrower
and decisive.  The fixed coherent cell mass generated by the
\(r_k^{-1}\)-normalized caloric contact is incompatible with that uniform
packet.  Strong local \(L^3\) convergence cannot ``pass the fixed contact''
because the contact carries the vanishing prefactor \(\nu^2/r_k\).

Therefore this scalar cannot by itself be used as the nontriviality marker
for a compact unit-viscosity ancient Navier--Stokes profile.  A nonzero
ancient profile may still be obtained from an independent scale-invariant
nontriviality row, such as an appropriate local epsilon-regularity
contrapositive.  That row must be stated and paid separately; it is not
contained in the coherent contact probability.
\end{assessment}

% =============================================================================
\section{Matched source-modulus saturation on a heavy thread}
\label{sec:source-modulus-thread}
% =============================================================================

The cubic audit applies without source-side localization.  To connect a
heavy cell to the terminal source modulus one must additionally verify that
the coherent cell payment is represented by the localized source/writer
pair, rather than by a localization router or local counterflow.

Put
\begin{equation}
 c(t):=A^{-1/4}\BNS(u)(t),
 \qquad
 \dd\mu_k^{\rm src}(t):=\|c(t)\|_2\,\dd t,
 \label{eq:source-c}
\end{equation}
and let \(q_k(t)\) be the inherited normalized backward-caloric writer.  Its
frozen age bound is
\begin{equation}
 \boxed{
 \|q_k(t)\|_2
 \le\frac1{2\sqrt{e\nu(b_k-t)}}.}
 \label{eq:writer-age-bound}
\end{equation}
The terminal source modulus is
\begin{equation}
 M_k(\tau)
 :=\int_{b_k-\tau}^{b_k}\|c(t)\|_2\,\dd t.
 \label{eq:source-modulus}
\end{equation}

Let \(E_{k,n,Q}(t)\) be the contraction on the source Hilbert space supplied
by the already assembled quadratic spatial localization for the cell.  It
may be zero outside the age slab.  Define the localized source-pair scalar
\begin{equation}
 s_{k,n,Q}^{\rm src}
 :=-\operatorname{Re}
 \int_{I_{k,n}}
 \ip{E_{k,n,Q}(t)c(t)}{E_{k,n,Q}(t)q_k(t)}\,\dd t.
 \label{eq:localized-source-scalar}
\end{equation}
For a selected cell \(\gamma=(k,n,Q)\), define the local incidence residual
\begin{equation}
 \boxed{
 \Delta_\gamma^{\rm inc}
 :=\left|
  \kappa_k(\mathcal C_{k,n,Q})
  -s_{k,n,Q}^{\rm src}
 \right|.}
 \label{eq:local-incidence-residual}
\end{equation}
This single residual includes the increment-to-caloric localization router,
local signed counterflow, and any failure of the declared source-side cell
to represent the coherent payment.  It is not silently set to zero.

\begin{lemma}[Matched-age source bound]
\label{lem:matched-source-bound}
For \(\ell=2^{-n}\),
\begin{equation}
 \boxed{
 |s_{k,n,Q}^{\rm src}|
 \le
 \frac{C}{\ell}
 M_k\!\left(\frac{\ell^2}{2\nu}\right),}
 \label{eq:matched-source-bound}
\end{equation}
where \(C\) is absolute.
\end{lemma}

\begin{proof}
On the age slab \(I_{k,n}\),
\[
 \frac{\ell^2}{8\nu}<b_k-t\le\frac{\ell^2}{2\nu}.
\]
Hence \eqref{eq:writer-age-bound} gives
\(
 \|q_k(t)\|_2\le C/\ell
\).  Since \(E_{k,n,Q}(t)\) is a contraction,
\[
 \begin{aligned}
 |s_{k,n,Q}^{\rm src}|
 &\le\int_{I_{k,n}}
      \|E_{k,n,Q}c(t)\|_2
      \|E_{k,n,Q}q_k(t)\|_2\,\dd t\\
 &\le\frac C\ell
      \int_{I_{k,n}}\|c(t)\|_2\,\dd t\\
 &\le\frac C\ell
      M_k\!\left(\frac{\ell^2}{2\nu}\right).
 \end{aligned}
\]
\end{proof}

\begin{theorem}[Heavy thread: local incidence obstruction or critical source saturation]
\label{thm:heavy-thread-source-alternative}
Assume \eqref{eq:heavy-payment-cell} on cells of scales
\(\ell_k\to0\).  After passage to a subsequence, at least one of the
following holds.
\begin{enumerate}[label=\textup{(S\arabic*)},leftmargin=3.4em]
\item \emph{Local source-incidence/router survival:}
      \begin{equation}
       \boxed{
       \limsup_{k\to\infty}
       \Delta_{\gamma_k}^{\rm inc}>0.}
       \label{eq:incidence-survival}
      \end{equation}
\item \emph{Critical matched source-modulus saturation:} with
      \(
       \tau_k=\ell_k^2/(2\nu)
      \),
      \begin{equation}
       \boxed{
       M_k(\tau_k)\ge c h_*p_*\ell_k,
       \qquad
       \frac{M_k(\tau_k)}{\sqrt{\tau_k}}
       \ge c h_*p_*\sqrt\nu.}
       \label{eq:critical-source-saturation}
      \end{equation}
      Consequently, for every fixed \(\varepsilon>0\),
      \begin{equation}
       \boxed{
       \limsup_{k\to\infty}
       \left[
        \varepsilon^{-1/2}M_k(\varepsilon)
        +\int_0^\varepsilon
          M_k(\tau)\tau^{-3/2}\,\dd\tau
       \right]
       \ge c h_*p_*\sqrt\nu.}
       \label{eq:Dini-defect-lower}
      \end{equation}
\end{enumerate}
Thus a source-aligned heavy singular thread does not merely coexist with
failure of the Dini criterion: it realizes the sharp square-root obstruction
at the physical parabolic age of that thread.
\end{theorem}

\begin{proof}
If item \textup{(S1)} fails, pass to a subsequence on which
\(
 \Delta_{\gamma_k}^{\rm inc}\to0
\).  Since
\(
 \kappa_k(\mathcal C_{\gamma_k})
 \ge h_*p_*
\), eventually
\(
 |s_{\gamma_k}^{\rm src}|\ge h_*p_*/2
\).  Lemma~\ref{lem:matched-source-bound} gives
\[
 M_k(\tau_k)\ge c h_*p_*\ell_k.
\]
Since \(\sqrt{\tau_k}=\ell_k/\sqrt{2\nu}\), the second inequality in
\eqref{eq:critical-source-saturation} follows.

Fix \(\varepsilon>0\).  For late \(k\), \(4\tau_k\le\varepsilon\).  By
monotonicity of \(M_k\),
\[
 \begin{aligned}
 \int_0^\varepsilon M_k(\tau)\tau^{-3/2}\,\dd\tau
 &\ge
 \int_{\tau_k}^{4\tau_k}
 M_k(\tau)\tau^{-3/2}\,\dd\tau\\
 &\ge
 M_k(\tau_k)
 \int_{\tau_k}^{4\tau_k}\tau^{-3/2}\,\dd\tau\\
 &=\frac{M_k(\tau_k)}{\sqrt{\tau_k}}.
 \end{aligned}
\]
Use \eqref{eq:critical-source-saturation} and take the limsup.
\end{proof}

\begin{corollary}[No uniform supercritical modulus on a source-aligned heavy thread]
\label{cor:no-supercritical-heavy-thread}
On the second branch of
\Cref{thm:heavy-thread-source-alternative}, there do not exist
\(\epsilon>0\), \(C<\infty\), and \(\tau_0>0\) such that
\begin{equation}
 M_k(\tau)\le C\tau^{1/2+\epsilon}
 \label{eq:supercritical-modulus}
\end{equation}
for all late \(k\) and all \(0<\tau<\tau_0\).
\end{corollary}

\begin{proof}
Evaluate \eqref{eq:supercritical-modulus} at \(\tau_k\).  Dividing by
\(\sqrt{\tau_k}\) gives an upper bound \(C\tau_k^\epsilon\to0\), which
contradicts \eqref{eq:critical-source-saturation}.
\end{proof}

\begin{assessment}[What has and has not been obtained from the source modulus]
\label{ass:source-modulus-meaning}
The theorem localizes the critical Dini failure to the same age scale as a
source-aligned heavy cell.  It does not derive a supercritical source modulus
from Leray energy.  Rather, it proves that every heavy thread has an exact
early fork: either the physical source-side localization fails to represent
the coherent payment, or the terminal source modulus saturates the critical
square-root scale there.  This is the requested source-modulus-to-thread
interface.
\end{assessment}

% =============================================================================
\section{The contact-preserving amplitude normalization}
\label{sec:amplitude-normalization}
% =============================================================================

The unit-viscosity rescaling preserves the equation but not the normalized
contact.  The unique scalar amplitude which removes the factor in
\eqref{eq:exact-principal-scaling} is
\begin{equation}
 \boxed{
 \rho_k:=\left(\frac{r_k}{\nu^2}\right)^{1/3},
 \qquad
 W_k:=\rho_k^{-1}U_k.}
 \label{eq:contact-amplitude}
\end{equation}
Since \(\nu\) is fixed and \(r_k\to\infty\), one has
\(\rho_k\to\infty\).

\begin{theorem}[Strong-coupling and vanishing-viscosity equations]
\label{thm:strong-coupling}
Let \(U_k\) solve the unit-viscosity projected equation
\begin{equation}
 \partial_sU_k+\PP\operatorname{div}(U_k\otimes U_k)=\Delta U_k.
 \label{eq:U-equation}
\end{equation}
Then \(W_k\) from \eqref{eq:contact-amplitude} solves
\begin{equation}
 \boxed{
 \partial_sW_k
 +\rho_k\PP\operatorname{div}(W_k\otimes W_k)
 =\Delta W_k.}
 \label{eq:strong-coupling-equation}
\end{equation}
Moreover, for any reference time \(s_k^0\), the fast-time field
\begin{equation}
 Z_k(\tau,y):=W_k\!\left(s_k^0+\frac\tau{\rho_k},y\right)
 \label{eq:fast-time-field}
\end{equation}
solves
\begin{equation}
 \boxed{
 \partial_\tau Z_k
 +\PP\operatorname{div}(Z_k\otimes Z_k)
 =\rho_k^{-1}\Delta Z_k.}
 \label{eq:vanishing-viscosity-equation}
\end{equation}
Finally, the principal contact is amplitude normalized:
\begin{equation}
 \boxed{
 \frac{\nu^2}{r_k}\mathfrak C(U_k)
 =\mathfrak C(W_k).}
 \label{eq:contact-normalized}
\end{equation}
\end{theorem}

\begin{proof}
Substitute \(U_k=\rho_kW_k\) into \eqref{eq:U-equation} and divide by
\(\rho_k\).  This gives \eqref{eq:strong-coupling-equation}.  The chain rule
in \eqref{eq:fast-time-field} gives
\(
 \partial_\tau Z_k=\rho_k^{-1}\partial_sW_k
\), so division of \eqref{eq:strong-coupling-equation} by \(\rho_k\) gives
\eqref{eq:vanishing-viscosity-equation}.  The functional
\(\mathfrak C\) is cubic, and \(\rho_k^3=r_k/\nu^2\); hence
\eqref{eq:contact-normalized}.
\end{proof}

\begin{proposition}[Conditional Euler--Reynolds limit]
\label{prop:Euler-Reynolds-limit}
Let \(K\Subset\R\times\R^3\) be a compact cylinder on which
\eqref{eq:vanishing-viscosity-equation} is defined.  Suppose, after passage
to a subsequence, that
\begin{align}
 Z_k&\rightharpoonup Z
 &&\text{in }L^2(K),
 \label{eq:Z-weak}\\
 Z_k\otimes Z_k&\rightharpoonup Q
 &&\text{in }L^1(K),
 \label{eq:Q-weak}
\end{align}
and that \((Z_k)\) is locally bounded in \(L^1\), so that the viscous term
vanishes against smooth tests after multiplication by \(\rho_k^{-1}\).
Then
\begin{equation}
 \boxed{
 \partial_\tau Z+\PP\operatorname{div}Q=0}
 \label{eq:Q-limit-equation}
\end{equation}
in distributions.  Writing
\(
 \mathcal R:=Q-Z\otimes Z
\), this is the Euler--Reynolds equation
\begin{equation}
 \boxed{
 \partial_\tau Z
 +\PP\operatorname{div}(Z\otimes Z+\mathcal R)=0.}
 \label{eq:Euler-Reynolds}
\end{equation}
If \(Z_k\to Z\) strongly in local \(L^2\), then \(\mathcal R=0\) and the
limit solves Euler.
\end{proposition}

\begin{proof}
Test \eqref{eq:vanishing-viscosity-equation} against a compactly supported
smooth divergence-free field \(\varphi\).  The weak convergences pass the
linear and quadratic terms.  The viscous term is
\(
 \rho_k^{-1}\int Z_k\cdot\Delta\varphi
\), which tends to zero by local
\(L^1\)-boundedness and \(\rho_k\to\infty\).  This gives
\eqref{eq:Q-limit-equation}.  The decomposition of \(Q\) is a definition,
and strong local \(L^2\) convergence gives
\(Z_k\otimes Z_k\to Z\otimes Z\) in local \(L^1\).
\end{proof}

\begin{firewall}[The normalized survivor is not automatically ancient Navier--Stokes]
\label{fire:not-ancient-NS}
The amplitude normalization which retains a nonzero cubic contact changes
the equation.  In the original rescaled time it produces coupling
\(\rho_k\to\infty\); in the fast time it produces viscosity
\(\rho_k^{-1}\to0\).  A compact limit, if independently obtained, is
therefore an Euler or Euler--Reynolds object.  Failure of compactness is a
further concentration defect.  Neither branch is a unit-viscosity ancient
Navier--Stokes profile without an additional normalization or a separate
nontriviality theorem.
\end{firewall}

% =============================================================================
\section{Near-singular microdelocalization as cubic-capacity spreading}
\label{sec:microdelocalization}
% =============================================================================

Let \(\mathscr S_k(R)\) be the inherited family of canonical cells whose
spatial cubes lie within \(R\) cell scales of \(\Sigma_T\).  Let
\begin{equation}
 \eta_k(R):=\pi_k\!\left(
  \bigcup_{\gamma\in\mathscr S_k(R)}\mathcal C_\gamma
 \right).
 \label{eq:near-singular-mass}
\end{equation}
The canonical cells are disjoint after their null boundaries are fixed.

\begin{theorem}[Aggregate near-singular cubic-capacity lower bound]
\label{thm:aggregate-cubic-capacity}
If for some fixed \(R<\infty\) and \(\eta_*>0\),
\begin{equation}
 \eta_k(R)\ge\eta_*
 \label{eq:near-singular-fixed-mass}
\end{equation}
for all late \(k\), then
\begin{equation}
 \boxed{
 \sum_{\gamma\in\mathscr S_k(R)}\mathscr I_\gamma
 \ge c\frac{h_*\eta_*}{\nu^2}r_k,}
 \label{eq:aggregate-I-lower}
\end{equation}
and
\begin{equation}
 \boxed{
 \sum_{\gamma\in\mathscr S_k(R)}\mathscr X_\gamma
 \ge c\frac{h_*\eta_*}{\nu^2}r_k.}
 \label{eq:aggregate-X-lower}
\end{equation}
Thus the microdelocalized coherent law carries a linearly diverging sum of
scale-invariant local cubic capacities, even though no single cell has a
fixed payment probability.
\end{theorem}

\begin{proof}
By disjointness and \eqref{eq:payment-dominated},
\[
 \begin{aligned}
 h_*\eta_*
 &\le
 \kappa_k\!\left(
  \bigcup_{\gamma\in\mathscr S_k(R)}\mathcal C_\gamma
 \right)\\
 &\le
 \sum_{\gamma\in\mathscr S_k(R)}
 |\eta_k|(\mathcal C_\gamma).
 \end{aligned}
\]
Apply \eqref{eq:full-cell-increment-bound} to every cell and sum.  The
constant is uniform because all cells use the same fixed parabolic annulus.
This proves \eqref{eq:aggregate-I-lower}.  Apply
\eqref{eq:increment-to-velocity} cellwise to obtain
\eqref{eq:aggregate-X-lower}.
\end{proof}

\begin{definition}[Local energy--dissipation capacity]
\label{def:ED-capacity}
For each cell choose the fixed enlargements used in
\Cref{lem:local-cubic-interpolation}, and put
\begin{equation}
 \boxed{
 \mathscr K_\gamma^{\rm ED}
 :=\mathscr E[U_\gamma]^{3/4}
   \mathscr D[U_\gamma]^{3/4}.}
 \label{eq:ED-capacity}
\end{equation}
\end{definition}

\begin{corollary}[Capacity-spreading alternative]
\label{cor:capacity-spreading}
Under the hypotheses of \Cref{thm:aggregate-cubic-capacity},
\begin{equation}
 \boxed{
 \sum_{\gamma\in\mathscr S_k(R)}
 \mathscr K_\gamma^{\rm ED}
 \ge c\frac{h_*\eta_*}{\nu^2}r_k.}
 \label{eq:aggregate-ED-capacity}
\end{equation}
Consequently, for every threshold \(K_k>0\), at least one of the following
holds:
\begin{enumerate}[label=\textup{(M\arabic*)},leftmargin=3.4em]
\item there is a near-singular cell with
      \(\mathscr K_\gamma^{\rm ED}>K_k\);
\item the number of near-singular cells carrying nonzero capacity is at least
      \begin{equation}
       \boxed{
       c\frac{h_*\eta_*}{\nu^2}
       \frac{r_k}{K_k}.}
       \label{eq:capacity-cell-count}
      \end{equation}
\end{enumerate}
If \(K_k\equiv K_*<\infty\), a fixed near-singular payment with uniformly
bounded local compactness capacity must therefore spread through at least
order \(r_k\) parabolic cells.
\end{corollary}

\begin{proof}
Apply \Cref{lem:local-cubic-interpolation} cellwise in
\eqref{eq:aggregate-X-lower}.  This gives
\eqref{eq:aggregate-ED-capacity}.  If item \textup{(M1)} fails, every
nonzero summand is at most \(K_k\).  The number of nonzero summands is then
at least the lower bound for their sum divided by \(K_k\).
\end{proof}

\begin{theorem}[Conditional singular-set packing gate]
\label{thm:packing-gate}
Let \(\mathcal A_k\subseteq\mathscr S_k(R)\) be a finite family of canonical
cells, and assume
\begin{equation}
 \sup_{\gamma\in\mathcal A_k}
 \mathscr K_\gamma^{\rm ED}\le K_k.
 \label{eq:local-capacity-upper}
\end{equation}
Then
\begin{equation}
 \boxed{
 \kappa_k\!\left(
  \bigcup_{\gamma\in\mathcal A_k}\mathcal C_\gamma
 \right)
 \le
 C\frac{\nu^2}{r_k}K_k|\mathcal A_k|.}
 \label{eq:packing-payment-upper}
\end{equation}
In particular, if
\begin{equation}
 K_k|\mathcal A_k|=o(r_k/\nu^2),
 \label{eq:packing-small}
\end{equation}
then the coherent payment of that family tends to zero.

At one comparable spatial scale \(\ell_k\), suppose the number of dyadic
cubes meeting the \(R\ell_k\)-neighbourhood of \(\Sigma_T\) is at most
\(C_\Sigma\ell_k^{-d}\).  If the local capacities are uniformly bounded,
a fixed payment on that scale requires
\begin{equation}
 \boxed{
 r_k\ell_k^d\le C.}
 \label{eq:Minkowski-gate}
\end{equation}
Thus the branch closes whenever \(r_k\ell_k^d\to\infty\), unless a local
energy--dissipation capacity diverges.
\end{theorem}

\begin{proof}
For each cell, combine \eqref{eq:payment-dominated},
\eqref{eq:full-cell-velocity-bound}, and
\Cref{lem:local-cubic-interpolation}:
\[
 \kappa_k(\mathcal C_\gamma)
 \le C\frac{\nu^2}{r_k}\mathscr K_\gamma^{\rm ED}.
\]
Sum over the disjoint family \(\mathcal A_k\) and use
\eqref{eq:local-capacity-upper}.  This proves
\eqref{eq:packing-payment-upper} and its little-oh consequence.

For the final statement, take \(\mathcal A_k\) to be the comparable-scale
near-singular cells.  Its cardinality is at most
\(C_\Sigma\ell_k^{-d}\).  If its coherent payment is bounded below and
\(K_k\le K_*\), \eqref{eq:packing-payment-upper} rearranges to
\(r_k\ell_k^d\le C\).
\end{proof}

\begin{firewall}[Packing is geometry, not source multiplicity]
\label{fire:packing-not-births}
The cardinality in \eqref{eq:capacity-cell-count} counts the spatial/
parabolic support needed to carry one already assembled scalar under a
uniform local-capacity bound.  It is not the dimension of the nonlinear
source carrier and may not be charged as multiple source birth.  Likewise,
\eqref{eq:Minkowski-gate} is a conditional geometric closure row.  V1 does
not assert a new bound for the Minkowski dimension of \(\Sigma_T\).
\end{firewall}

% =============================================================================
\section{Corrected singular-thread and microdelocalization alternative}
\label{sec:corrected-alternative}
% =============================================================================

\begin{theorem}[Version V1 terminal concentration alternative]
\label{thm:V1-integrated}
Assume the inherited unforced periodic unit critical-norm chronology, the
fixed held-out contact margin \eqref{eq:fixed-contact}, and the V35 removal
of represented return, cutoff, upper-tail, forcing, boundary, gluing,
carrier, and physical-input scalars.  Retain the V35 terminal singular-set
support and spatial trichotomy.  Then the following sharpened conclusions
hold.
\begin{enumerate}[label=\textup{(V1.\arabic*)},leftmargin=5.2em]
\item If the inherited terminal source Dini condition holds, the selected
      scalar vanishes and the chronology stops by the inherited critical
      continuation theorem.
\item On a heavy near-singular parabolic thread, the dimensionless local
      cubic packet satisfies \eqref{eq:heavy-cubic-lower}, and the local
      energy--dissipation product satisfies \eqref{eq:ED-product-lower}.
      Hence the uniform unit-viscosity compactness packet fails.
\item On that heavy thread, either the local source-incidence/router residual
      \eqref{eq:incidence-survival} survives, or the terminal source modulus
      saturates the critical square-root rate at the matched physical age,
      with the nonzero Dini defect \eqref{eq:Dini-defect-lower}.
\item The amplitude which preserves the normalized cubic contact is
      \(\rho_k=(r_k/\nu^2)^{1/3}\).  It produces the strong-coupling equation
      \eqref{eq:strong-coupling-equation}, or after fast time the
      vanishing-viscosity equation \eqref{eq:vanishing-viscosity-equation}.
      A compact failure object is therefore Euler/Euler--Reynolds unless an
      independent Navier--Stokes nontriviality row is supplied.
\item On the near-singular microdelocalized branch, the aggregate local cubic
      and energy--dissipation capacities satisfy
      \eqref{eq:aggregate-X-lower} and
      \eqref{eq:aggregate-ED-capacity}.  Thus either a cubic-heavy subthread
      emerges or the packet genuinely spreads through order \(r_k\) cells
      under a uniform local-capacity bound.
\item A predeclared singular-set packing bound closes any finite or
      comparable-scale cohort for which
      \(K_k|\mathcal A_k|=o(r_k/\nu^2)\).  Failure returns an explicit local
      capacity blow-up, packing growth, scale spreading, or
      singular-distance/scale separation.
\end{enumerate}
No ancient suitable profile, donor budget, or global regularity conclusion
is inferred merely from the coherent cell probability.
\end{theorem}

\begin{proof}
Item \textup{(V1.1)} is the frozen V35 Dini continuation implication.
Items \textup{(V1.2)--(V1.3)} are
\Cref{thm:heavy-cell-cubic-divergence,thm:energy-dissipation-escalation,cth:compact-handoff,thm:heavy-thread-source-alternative}.
Item \textup{(V1.4)} is \Cref{thm:strong-coupling} and
\Cref{prop:Euler-Reynolds-limit}.  Item \textup{(V1.5)} is
\Cref{thm:aggregate-cubic-capacity,cor:capacity-spreading}.  The final item
is \Cref{thm:packing-gate}.  The inherited singular-distance/scale branch is
untouched and remains an explicit alternative.
\end{proof}

\begin{assessment}[Correct NS--A interface to NS--B]
\label{ass:NSA-NSB-interface}
The donor programme is not opened merely because V35 selected a heavy
payment cell.  V1 first removes a normalization ambiguity:
\[
 \boxed{
 \begin{gathered}
 \text{heavy coherent cell}
 \Longrightarrow
 \text{local cubic and energy/dissipation escalation},\\
 \text{not automatically a compact nonzero unit-viscosity ancient profile}.
 \end{gathered}}
\]
The correct export to NS--B is one of the following fixed-datum objects:
\begin{enumerate}[label=\textup{(E\arabic*)},leftmargin=3.4em]
\item a further localized subthread selected by an independent
      scale-invariant Navier--Stokes nontriviality marker;
\item a contact-normalized strong-coupling/vanishing-viscosity thread with a
      specified quadratic defect;
\item a near-singular microdelocalized family carrying the aggregate capacity
      lower bound \eqref{eq:aggregate-ED-capacity}.
\end{enumerate}
Only after rank-at-most-two descendants and represented returns are removed
from one of these objects should the V13 path
\((\theta_{\rm c}\mathbb K[u])^\circ\) and the Reynolds-covariance square
function be tested.
\end{assessment}

% =============================================================================
\section{Proof-status ledger}
\label{sec:ledger}
% =============================================================================

\begingroup
\small
\setlength{\LTpre}{0.4em}
\setlength{\LTpost}{0.4em}
\begin{longtable}{p{0.23\textwidth}p{0.19\textwidth}p{0.50\textwidth}}
\caption{NS--A Version V1 proof-status ledger.}
\label{tab:ledger}\\
\toprule
Row & Status & Exact content \\
\midrule
\endfirsthead
\toprule
Row & Status & Exact content \\
\midrule
\endhead
V35 terminal source modulus and Dini criterion
& \INHERITED
& Exact gain-tail reduction and continuation under the uniform terminal Dini condition; not re-proved.\\[0.3em]
V35 regular-region vanishing and singular-set support
& \INHERITED
& Every coherent centre-law cluster is supported on \(\Sigma_T\); not re-proved.\\[0.3em]
V35 spatial trichotomy
& \INHERITED
& Singular-distance/scale separation, one heavy near-singular cell, or near-singular microdelocalization.\\[0.3em]
Principal contact rescaling
& \PROVED
& Exact prefactor \(\nu^2/r_k\) in \eqref{eq:exact-principal-scaling}.\\[0.3em]
Full periodic cell bound
& \PROVED
& \(|\eta_k|(\mathcal C)\le C(\nu^2/r_k)\mathscr I_\mathcal C\le C(\nu^2/r_k)\mathscr X_\mathcal C\).\\[0.3em]
Heavy cell local cubic packet
& \PROVED
& Fixed coherent cell probability forces \(\mathscr X_\mathcal C\gtrsim r_k/\nu^2\).\\[0.3em]
Local energy/dissipation rate
& \PROVED
& Product and maximum lower bounds \eqref{eq:ED-product-lower}--\eqref{eq:ED-max-lower}.\\[0.3em]
Uniform unit-viscosity compactness on the heavy contact thread
& \OBSTRUCTION
& Incompatible with the fixed \(r_k^{-1}\)-normalized coherent contact.  The prior nontriviality transfer omitted the surviving \(r_k^{-1}\) factor.\\[0.3em]
Abstract ancient compactness implication
& \INHERITED
& Still valid when its uniform hypotheses hold, but this contact cannot supply its nonzero limit.\\[0.3em]
Matched source-modulus localization
& \PROVED / \CONDITIONAL
& Passing local source incidence forces critical \(\sqrt\tau\) saturation and a fixed Dini defect; otherwise the incidence/router residual survives.\\[0.3em]
Contact-preserving normalization
& \PROVED
& \(W_k=(\nu^2/r_k)^{1/3}U_k\) gives strong coupling; fast time gives vanishing viscosity.\\[0.3em]
Euler--Reynolds failure object
& \CONDITIONAL
& Exact distributional limit under the displayed compactness of \(Z_k\) and its quadratic products.\\[0.3em]
Microdelocalized aggregate capacity
& \PROVED
& Near-singular payment mass \(\eta_*>0\) forces total cubic and ED capacity at least \(c r_k/\nu^2\).\\[0.3em]
Singular-set packing gate
& \PROVED / \CONDITIONAL INPUT
& Exact implication from a supplied packing/local-capacity bound; no new dimension estimate for \(\Sigma_T\) is asserted.\\[0.3em]
Singular-distance/scale separation
& \OPEN
& Not changed by V1; remains a distinct localization obstruction.\\[0.3em]
Rank-three donor/paid-stress estimate
& \OPEN / NOT YET EXPORTED
& V13 is opened only on a corrected fixed-datum survivor.\\[0.3em]
Global three-dimensional regularity
& \OPEN
& No global regularity conclusion is claimed.\\
\bottomrule
\end{longtable}
\endgroup

% =============================================================================
\section{Exact V1 endpoint and V2 attack order}
\label{sec:V2-programme}
% =============================================================================

The first fresh obstruction after V1 is
\begin{center}
\fbox{\begin{minipage}{0.91\textwidth}
\centering
\textbf{Critical source-modulus saturation coupled to a noncompact local
energy--dissipation packet.}
\par\medskip
A heavy source-aligned cell saturates the terminal \(\sqrt\tau\) modulus at
its matched parabolic age.  Independently of source alignment, the same
coherent payment forces the unit-viscosity local cubic packet to grow like
\(r_k\), and hence forces scale-invariant local energy or dissipation to grow
like at least \(r_k^{2/3}\).  Compact unit-viscosity ancient-profile
extraction therefore cannot be the immediate next arrow.  The contact-
normalized object is strong-coupling/vanishing-viscosity; the alternative is
a further subscale or microdelocalized capacity defect.
\end{minipage}}
\end{center}

\begin{programme}[Version V2 attack order]
\label{prog:V2}
Preserve Version~V1 verbatim.  Continue in the following order.
\begin{enumerate}[label=\textup{(V2.\arabic*)},leftmargin=5.2em]
\item Freeze the normalization correction.  Do not again infer nonzero
      unit-viscosity ancient mass from a fixed coherent cell probability.
\item On a heavy source-aligned thread, insert the localized local-energy
      inequality and harmonic-pressure decomposition into
      \eqref{eq:ED-product-lower}.  Determine whether the divergence lies in
      local kinetic energy, local dissipation, harmonic pressure, or the
      source-incidence residual.
\item Seek an independent scale-invariant nontriviality marker on the same
      cell.  If one exists, perform the ancient suitable-profile extraction
      using that marker, not the \(r_k^{-1}\)-normalized caloric contact.
\item On the contact-normalized branch, study the fast-time
      vanishing-viscosity sequence.  Prove compact Euler landing, produce the
      exact Euler--Reynolds defect, or return a further concentration scale.
\item On near-singular microdelocalization, apply the aggregate capacity
      theorem before counting cells.  Test a physical singular-set packing,
      local-energy capacity, or finite-measure density theorem against
      \eqref{eq:aggregate-ED-capacity}.
\item Keep singular-distance/scale separation distinct.  Test whether the
      terminal singular set, the heat length, and the source-age law yield a
      return to a matched cell; otherwise retain the ratio escape as the
      typed obstruction.
\item Open the V13 donor packet only after one rank-three fixed-datum survivor
      has been selected by items \textup{(V2.2)--(V2.5)}.  Then test the
      five-dimensional paid stress and positive Reynolds covariance on that
      same thread, charging no cell or scale twice.
\item Invoke the inherited critical continuation theorem immediately if the
      terminal Dini condition or every corrected singular-contact branch
      closes.
\end{enumerate}
A return to uniform cell action, weak-\(L^1\) gain, raw participation,
source-dimension counting, or the uncorrected compact-contact handoff would
repeat a closed loop.
\end{programme}

\begin{firewall}[No global theorem in V1]
\label{fire:no-global}
Version~V1 proves a normalization correction, a matched source-modulus
alternative, quantitative local energy/dissipation escalation, a
strong-coupling/vanishing-viscosity reclassification, and an aggregate
microdelocalized capacity lower bound.  It does not prove a supercritical
terminal source modulus, exclude every Euler/Euler--Reynolds or rank-three
survivor, control the singular-distance/scale-separated branch, close the
V13 paid-stress variation, or establish global three-dimensional
Navier--Stokes regularity.
\end{firewall}

% =============================================================================
\section*{Version V1 history}
\addcontentsline{toc}{section}{Version V1 history}
% =============================================================================

\begin{description}[leftmargin=3.5cm,style=nextline]
\item[Version V1 -- 2 September 2026.]
Started the dedicated NS--A append-only series at the exact V35 terminal
frontier.  Preserved the terminal Dini criterion, singular-set support,
spatial trichotomy, and V13 donor packet as inherited layers.  Computed the
exact \(\nu^2/r_k\) factor for the caloric cell under unit-viscosity
rescaling and proved the full periodic total-variation bound.  Deducted that
a heavy coherent cell forces dimensionless local \(L^3\) mass of order
\(r_k\), together with quantitative local energy/dissipation escalation.
Corrected the previous nonzero ancient-profile handoff: its uniform local
compactness packet is incompatible with this fixed normalized contact.
Identified the contact-preserving strong-coupling/vanishing-viscosity
normalization and its conditional Euler--Reynolds limit.  Proved that a
source-aligned heavy cell saturates the critical square-root terminal source
modulus at its matched age, while failure returns a local incidence/router
residual.  Upgraded near-singular microdelocalization to an aggregate cubic
and energy--dissipation capacity lower bound and a conditional singular-set
packing gate.  No global regularity conclusion was claimed.
\end{description}

For Version~V2, preserve this complete source verbatim, remove only the final
\verb|\end{document}|, append new mathematical material immediately before
it, restore the final command, and use the filename
\begin{center}
\scriptsize
\path{Renewal_NS_Terminal_Source_Concentration_and_Singular_Thread_Closure_Programme_V2.tex}.
\end{center}


% =============================================================================
% BEGIN VERSION V2 MATERIAL -- append-only continuation of NS--A V1
% =============================================================================
\clearpage
\hypersetup{pdftitle={NS-A Terminal Source Concentration Programme: Cumulative V2},pdfauthor={A. Pelissier}}
\makeatletter
\addtocontents{toc}{\protect\begingroup\protect\makeatletter}
\addtocontents{toc}{\protect\def\protect\l@section{\protect\@dottedtocline{1}{0em}{3.5em}}}
\addtocontents{toc}{\protect\makeatother}
\makeatother
\renewcommand{\thesection}{V2.\arabic{section}}
\renewcommand{\thesubsection}{V2.\arabic{section}.\arabic{subsection}}
\renewcommand{\theequation}{V2.\arabic{equation}}
\setcounter{section}{0}
\setcounter{equation}{0}
\renewcommand{\theHsection}{V2.\arabic{section}}
\renewcommand{\theHsubsection}{V2.\arabic{section}.\arabic{subsection}}
\renewcommand{\theHtheorem}{V2.\arabic{section}.\arabic{theorem}}
\renewcommand{\theHequation}{V2.\arabic{equation}}

\begin{center}
{\LARGE\bfseries NS--A: Version V2}\\[0.5em]
{\large Local-Energy Absorption, the Caloric Primitive,\\
and the Long Fast-Time Contact}\\[0.5em]
{\normalsize 2 September 2026}
\end{center}
\addcontentsline{toc}{part}{Version V2: negative-primitive dissipation and long fast-time contact}

\begin{abstract}
The V1 normalization audit is retained.  This continuation inserts the
actual localized Navier--Stokes balance into the contact-normalized thread,
rather than asking again for an unnormalized compact ancient profile.  The
result is an exact identity expressing the selected signed contact as one
complete nonlinear localization current, a quadratic stock divided by the
large coupling, and a signed test of a positive increment-dissipation
measure.  A separate replay discrepancy records the difference between a
positive coherent payment and that signed local contact.

The caloric primitive in this identity is evaluated explicitly.  In three
dimensions the Euclidean kernel of $\Lambda e^{-aA}$ has one radial zero,
at $|z|=R_\Lambda\sqrt a$, with the certified enclosure
\[
 3.0039505<R_\Lambda<3.0039506.
\]
It is positive inside that radius and negative outside.  Therefore a
positive contact whose complete localization and replay discrepancies
vanish cannot persist entirely inside the positive-primitive region unless
its normalized kinetic stock grows.  Otherwise it forces a nonzero positive
increment-dissipation measure tested on the negative-primitive region.
This is a conditional exclusion using the actual kernel, not a claim that
local interface currents vanish automatically.

On the dissipation branch the unit-viscosity local dissipation grows at
least like $r_k/\nu^2$, strengthening V1's $r_k^{2/3}$ energy/dissipation
alternative after the stated interface hypotheses pass.  The original
energy inequality then gives $\sum_k r_k\ell_k<\infty$ on disjoint selected
passages.  The complementary kinetic branch with nonvanishing
$r_k\ell_k$ produces an atomic terminal kinetic-energy concentration
measure.  The local energy equation further identifies the surviving
averaged dissipation with a nonvanishing incoming kinetic--pressure energy
current.

Finally, a fixed slow-time cell occupies a fast-time interval of length
$O(\rho_k)$, not a fixed interval.  Its contact is a $\rho_k^{-1}$-weighted
integral on that growing interval.  Uniformly bounded fixed fast-time
cubic windows consequently carry vanishing payment.  Under explicit local
velocity and pressure bounds, the slow-time moments satisfy a spatial
Euler--Reynolds equilibrium constraint and an energy-flux balance with a
positive averaged dissipation measure.  Neither a dynamical Euler limit on
one bounded fast window nor a nonzero Reynolds covariance alone identifies
this long-window contact.  The local energy equation then absorbs the complete viscous contact term
into two actual signed fluxes and a spacetime kinetic stock.  If the fluxes
and localization discrepancy vanish, the unnormalized cubic packet must
grow at least like $(r_k/\nu^2)^{3/2}$, stronger than the V1 linear rate.
No terminal Dini improvement or global three-dimensional regularity theorem
is asserted.
\end{abstract}

% =============================================================================
\section{The selected thread and the two local scalar meanings}
\label{sec:v2-thread}
% =============================================================================

Keep the V1 first-passage chronology, viscosity, coherent replay, and
canonical cells.  Select a heavy cell $\gamma_k=(k,n_k,Q_k)$, let
$\ell_k=2^{-n_k}$, and choose the rescaling centre at a fixed corner of its
Euclidean lift.  Thus the rescaled cube is the same unit cube $Q_0$ for all
$k$.  Write
\begin{equation}
 \rho_k=\left(\frac{r_k}{\nu^2}\right)^{1/3},\qquad
 \varepsilon_k=\rho_k^{-1},\qquad U_k=\rho_k W_k.
 \label{eq:v2-normalization}
\end{equation}
The pressure used below is
\begin{equation}
 P_k(s,y)=\frac{\ell_k^2}{\nu^2\rho_k^2}
 p\!\left(b_k+\frac{\ell_k^2}{\nu}s,x_k+\ell_k y\right).
 \label{eq:v2-pressure-normalization}
\end{equation}
In particular the unprojected equation, including its pressure coefficient,
is
\begin{equation}
 \boxed{\partial_sW_k+
 \rho_k\bigl((W_k\cdot\nabla)W_k+\nabla P_k\bigr)=\Delta W_k,
 \qquad \operatorname{div}W_k=0.}
 \label{eq:v2-equation}
\end{equation}
This is V1's equation with the pressure normalization made explicit.

Let $J=[-1/2,-1/8]$.  If the canonical age slab is truncated by the original
first-passage interval, use its actual rescaled subinterval
$J_k=[s_k^-,s_k^+]\subset J$.  All integrations on the selected passage are
over $J_k$; no extension of its time support is used in the summability
argument.  Smoothness before $T$ supplies both endpoint traces.

Fix, independently of the field, $0<\alpha_0<\alpha<\beta<\beta_0$ and
$\zeta_0\in C_c^\infty((\alpha_0,\beta_0))$ with
$0\le\zeta_0\le1$ and $\zeta_0=1$ on $[\alpha,\beta]$.  Put
\begin{equation}
 \zeta(s,z)=\zeta_0\!\left(\frac{|z|}{\sqrt{-2s}}\right).
 \label{eq:v2-separation-cutoff}
\end{equation}
Choose $\chi\in C_c^\infty(G_0)$, $0\le\chi\le1$, with $\chi=1$ on
$Q_0$.  Fixed larger domains $G_1,G_2$ contain all translated supports
$\supp\chi+B_{2\beta_0}$ and the pressure-localization collars.  Every
constant below may depend on these fixed cutoffs and the original torus,
but not on $k,\ell_k,r_k$, or $\rho_k$.

\begin{definition}[Exact rescaled periodic primitive and signed window]
\label{def:v2-primitive-window}
Use the full periodic kernel from \eqref{eq:eta} and define
\begin{equation}
 F_k(s,z)=\ell_k^4 L_{-2\ell_k^2s}(\ell_k z),\qquad
 \Psi_k(s,z)=F_k(s,z)\zeta(s,z).
 \label{eq:v2-full-primitive}
\end{equation}
The signed window contact is
\begin{equation}
 \boxed{\mathcal C_k=-\frac14\int_{J_k}\int\!\int
 \chi(y)\zeta(s,z)\nabla_zF_k(s,z)\cdot\delta_zW_k(s,y)
 |\delta_zW_k(s,y)|^2\,\dd y\,\dd z\,\dd s.}
 \label{eq:v2-signed-window}
\end{equation}
Let $\mathcal P_k$ be the integral of the same nonnegative window
$\chi\zeta$ against the V1 coherent payment $\kappa_k$, after the physical
change of variables.  Define the signed replay discrepancy
\begin{equation}
 \boxed{\mathcal R_k^{\rm rep}:=\mathcal P_k-\mathcal C_k.}
 \label{eq:v2-replay-discrepancy}
\end{equation}
\end{definition}

The exact scaling in V1 gives \eqref{eq:v2-signed-window} for the full
periodic measure, not just its Euclidean principal part: the factor
$\nu^2\rho_k^3/r_k$ is exactly one.  Since the cutoffs equal one on the
original cell and the replay is positive,
\begin{equation}
 \boxed{\mathcal P_k\ge\kappa_k(\mathcal C_{\gamma_k})
 \ge h_*p_*=:p_0>0.}
 \label{eq:v2-window-payment}
\end{equation}

\begin{assessment}[Positive replay is not a signed local equation]
\label{ass:v2-replay-scope}
The raw formula \eqref{eq:eta} is the literal caloric increment carrier.
A represented-source short, a change of localization, or cancellation of
positive and negative increments is not performed by that formula itself.
Consequently the local signed quantity used in a PDE identity must be
\eqref{eq:v2-signed-window}.  Any discrepancy from the coherent payment or
from a proposed locally held-out contact remains in
$\mathcal R_k^{\rm rep}$.  V1's full-variation and cubic-divergence bounds
are unchanged; they used domination by total variation, not equality of
these two local scalars.  No proof below infers
$\mathcal C_k\ge p_0$ solely from a heavy coherent cell.
\end{assessment}

% =============================================================================
\section{Evaluation of the caloric primitive and its unique radial sign change}
\label{sec:v2-kernel}
% =============================================================================

The new coefficient needed by the local balance is the primitive itself,
not only the inherited estimate for its gradient.  Use the Euclidean
Fourier convention
\begin{equation}
 \mathcal L_a(z)=(2\pi)^{-3}\int_{\R^3}
 e^{\mathrm i\xi\cdot z}|\xi|e^{-a|\xi|^2}\,\dd\xi.
 \label{eq:v2-euclidean-kernel}
\end{equation}
The sign conclusions below are independent of a positive Fourier-volume
normalization.

\begin{lemma}[Uniform primitive and time-derivative control]
\label{lem:v2-primitive-bounds}
On $J$ and every fixed bounded $z$-set,
\begin{equation}
 \boxed{F_k(s,z)=\mathcal L_{-2s}(z)+O(\ell_k^4),\qquad
 \sup_k\bigl(\|F_k\|_\infty+\|\nabla_zF_k\|_\infty
 +\|\partial_sF_k\|_\infty\bigr)<\infty.}
 \label{eq:v2-primitive-bounds}
\end{equation}
The convergence also holds after one spatial or time derivative.  In
particular $\Psi_k$ and $\partial_s\Psi_k$ are uniformly bounded on the
fixed separation supports.
\end{lemma}

\begin{proof}
For the Euclidean heat kernel $G_a$, spectral integration gives
\[
 \mathcal L_a(z)=\frac1{\sqrt\pi}\int_0^\infty
 q^{-1/2}(-\Delta)G_{a+q}(z)\,\dd q.
\]
Indeed, its Fourier multiplier is
$|\xi|^2e^{-a|\xi|^2}\pi^{-1/2}
\int_0^\infty q^{-1/2}e^{-q|\xi|^2}\dd q
=|\xi|e^{-a|\xi|^2}$.
For $a$ in a compact subinterval of $(0,\infty)$, the Gaussian derivative
bounds and the substitution $(a+q)/|z|^2=v$ give
\[
 |\partial_a^j\nabla_z^m\mathcal L_a(z)|
 \le C_{j,m}(1+|z|)^{-4-2j-m},\qquad j,m\in\{0,1\}.
\]
For completeness, on $q\ge a$ the integrand is bounded by
$Cq^{-3-j-m/2}\exp(-c|z|^2/q)$, whose integral has the displayed power;
on $q<a$ it has exponential decay in $|z|$.  On bounded $z$-sets the
integral is uniformly finite because $a$ is bounded away from zero.

Let $\Gamma$ be the original spatial period lattice.  Absolute convergence
and Fourier coefficients show that $L_a$ is the periodization of
$\mathcal L_a$; its zero coefficient vanishes.  Homogeneity yields
\[
 F_k(s,z)=\sum_{\gamma\in\Gamma}
 \mathcal L_{-2s}(z+\gamma/\ell_k).
\]
The zero image is the asserted limit.  Summing the remaining bounds gives
$O(\ell_k^4)$ for the kernel, $O(\ell_k^5)$ for its spatial derivative,
and $O(\ell_k^6)$ for its time derivative.  The lattice sums converge in
three dimensions.  This proves every assertion.
\end{proof}

\begin{theorem}[Explicit primitive and the universal sign radius]
\label{thm:v2-kernel-sign}
Let
\[
 D(x)=e^{-x^2}\int_0^x e^{t^2}\,\dd t,\qquad
 N(x)=x+(1-2x^2)D(x).
\]
For $r=|z|>0$,
\begin{equation}
 \boxed{\mathcal L_a(z)=a^{-2}\frac{N(x)}{8\pi^2x},
 \qquad x=\frac{r}{2\sqrt a}.}
 \label{eq:v2-Dawson-kernel}
\end{equation}
The value at $r=0$ is $a^{-2}/(4\pi^2)$.  There is exactly one
$R_\Lambda>0$ such that
\begin{equation}
 \boxed{
 \mathcal L_a(z)>0\ \text{if }|z|<R_\Lambda\sqrt a,\qquad
 \mathcal L_a(z)<0\ \text{if }|z|>R_\Lambda\sqrt a.}
 \label{eq:v2-kernel-sign}
\end{equation}
Its enclosure is
\begin{equation}
 \boxed{3.0039505<R_\Lambda<3.0039506.}
 \label{eq:v2-kernel-root-enclosure}
\end{equation}
\end{theorem}

\begin{proof}
Spherical integration in \eqref{eq:v2-euclidean-kernel} gives
\[
 \mathcal L_1(r)=\frac1{2\pi^2r}
 \int_0^\infty q^2e^{-q^2}\sin(rq)\,\dd q.
\]
The function $I(x)=\int_0^\infty e^{-q^2}\sin(2xq)\dd q$ satisfies
$I(0)=0$ and $I'=1-2xI$, by integration by parts.  The same initial-value
problem is satisfied by $D$, so $I=D$.  Differentiating twice under the
Gaussian integral and using
$D''=(4x^2-2)D-2x$ proves \eqref{eq:v2-Dawson-kernel} at $a=1$.
Scaling gives general $a$, and $D(x)=x+O(x^3)$ gives the value at zero.

For $0<x\le1/\sqrt2$, $N(x)>0$.  At any positive zero beyond that point,
$D(x)=x/(2x^2-1)$, and differentiation gives the exact value
\begin{equation}
 \left.N'(x)\right|_{N(x)=0}=-\frac2{2x^2-1}<0.
 \label{eq:v2-zero-crossing}
\end{equation}
Thus every zero is a strict crossing from positive to negative; there can
be at most one.  The rational sign certificate below gives a positive
value at $x=3.0039505/2$ and a negative value at $x=3.0039506/2$.
The intermediate value theorem proves existence, uniqueness, and the
claimed sign pattern.

Here are full details of the certificate, so the enclosure does not rely
on floating-point root finding.  For rational $q=x^2>1/2$, set
\[
 E_N(q)=\sum_{n=0}^N\frac{q^n}{n!},\qquad
 S_N(q)=\sum_{n=0}^N\frac{q^n}{n!(2n+1)},\qquad
 T_N(q)=\frac{q^{N+1}}{(N+1)!}\frac1{1-q/(N+2)}.
\]
When $q<N+2$, the omitted exponential tail is at most $T_N$, and the
omitted $S$-tail is at most $T_N/(2N+3)$.  Since
\[
 \frac{e^{x^2}N(x)}x=e^q-(2q-1)\sum_{n=0}^\infty
 \frac{q^n}{n!(2n+1)},
\]
this expression belongs to the rational interval
\begin{equation}
 \left[E_N-(2q-1)\left(S_N+\frac{T_N}{2N+3}\right),\quad
 E_N+T_N-(2q-1)S_N\right].
 \label{eq:v2-rational-sign-interval}
\end{equation}
Substitution of $N=32$ gives the following outward rational enclosures:
\[
\begin{array}{c|c}
 2x & e^{x^2}N(x)/x\\ \hline
 3.0039505 & (661,662)\,10^{-10}\\
 3.0039506 & (-1149,-1148)\,10^{-10}.
\end{array}
\]
They follow by finite rational additions and comparisons in
\eqref{eq:v2-rational-sign-interval}; the exact reproducible verifier is
included in Section~\ref{sec:v2-verification}.  The intervals have opposite
strict signs, completing the proof.
\end{proof}

\begin{corollary}[The positive region survives periodic rescaling]
\label{cor:v2-periodic-positive}
For every fixed $\delta>0$, on any fixed window with
$|z|/\sqrt{-2s}\le R_\Lambda-\delta$, the full periodic $F_k$ is positive
for all sufficiently large $k$.  Every negative point of $F_k$ on a fixed
window eventually satisfies
$|z|/\sqrt{-2s}>R_\Lambda-\delta$.
\end{corollary}

\begin{proof}
The Euclidean primitive has a positive minimum on the indicated compact
set by \Cref{thm:v2-kernel-sign}, including at $z=0$.  Uniform convergence
in \Cref{lem:v2-primitive-bounds} preserves that minimum up to a factor of
two.  The second assertion is its contrapositive.
\end{proof}

% =============================================================================
\section{The local balance after contact normalization}
\label{sec:v2-local-balance}
% =============================================================================

The inherited localized K\'arm\'an--Howarth--Monin identity is used, not
re-proved as a new theorem.  For the equation \eqref{eq:v2-equation}, its
coefficients and the required cutoff terms are as follows.  Define
\begin{align}
 S_k(s,z)&=\frac14\int\chi(y)|\delta_zW_k(s,y)|^2\,\dd y,
 &Y_k(s,z)&=\frac14\int\chi(y)\delta_zW_k|\delta_zW_k|^2\,\dd y,
 \label{eq:v2-SY}\\
 D_k(s,z)&=\frac12\int\chi(y)|\nabla_y\delta_zW_k|^2\,\dd y,
 &B_k^{\rm lin}(s,z)&=\frac14\int\Delta\chi(y)|\delta_zW_k|^2\,\dd y,
 \label{eq:v2-DBlin}
\end{align}
with the single nonlinear centre current
\begin{equation}
 B_k^{\rm nl}(s,z)=\int
 \left(\frac14 W_k(y)|\delta_zW_k|^2
       +\frac12\delta_zP_k\,\delta_zW_k\right)
 \cdot\nabla\chi(y)\,\dd y.
 \label{eq:v2-Bnl}
\end{equation}
The correctly normalized inherited identity is
\begin{equation}
 \boxed{\partial_sS_k+\rho_k\operatorname{div}_zY_k+D_k
 =B_k^{\rm lin}+\rho_k B_k^{\rm nl}.}
 \label{eq:v2-normalized-KHM}
\end{equation}
The coefficient $\rho_k$ multiplies both convection and pressure; dropping
it from the pressure current would change the identity.

Introduce the following three quantities on the same selected window:
\begin{align}
 \mathcal F_k&=\int_{J_k}\int
 \left(\Psi_k B_k^{\rm nl}
       +F_k\nabla_z\zeta\cdot Y_k\right)\,\dd z\,\dd s,
 \label{eq:v2-complete-current}\\
 \mathcal Q_k&=\int\Psi_k(s_k^-,z)S_k(s_k^-,z)\,\dd z
             -\int\Psi_k(s_k^+,z)S_k(s_k^+,z)\,\dd z\notag\\
 &\quad+\int_{J_k}\int
       \bigl(\partial_s\Psi_k\,S_k+\Psi_k B_k^{\rm lin}\bigr)
       \,\dd z\,\dd s,
 \label{eq:v2-quadratic-stock}\\
 \mathcal V_k&=\int_{J_k}\int\Psi_k(s,z)D_k(s,z)\,\dd z\,\dd s.
 \label{eq:v2-signed-dissipation}
\end{align}
The two summands in $\mathcal F_k$ are centre transport (including pressure)
and the separation-cutoff current.  They are assembled before any sign or
bound is imposed.  The moving heat annulus is accounted for by
$\partial_s\Psi_k$ in $\mathcal Q_k$.

\begin{theorem}[Exact normalized contact--current--dissipation identity]
\label{thm:v2-exact-contact-balance}
For every selected smooth Navier--Stokes window,
\begin{equation}
 \boxed{\mathcal C_k=\mathcal F_k+
 \varepsilon_k(\mathcal Q_k-\mathcal V_k),\qquad
 \mathcal P_k=\mathcal R_k^{\rm rep}+\mathcal F_k+
 \varepsilon_k(\mathcal Q_k-\mathcal V_k).}
 \label{eq:v2-master-balance}
\end{equation}
Let
\begin{equation}
 \mathcal E_k=\sup_{s\in J_k}\int_{G_1}|W_k(s,y)|^2\,\dd y,
 \qquad
 \mathcal D_k=\int_{J_k}\int_{G_1}|\nabla W_k|^2\,\dd y\,\dd s.
 \label{eq:v2-ED-W}
\end{equation}
Then
\begin{equation}
 \boxed{|\mathcal Q_k|\le C\mathcal E_k,\qquad
 |\mathcal V_k|\le C\mathcal D_k.}
 \label{eq:v2-quadratic-bounds}
\end{equation}
\end{theorem}

\begin{proof}
Multiply \eqref{eq:v2-normalized-KHM} by $\Psi_k$, integrate on
$J_k\times\R^3_z$, and integrate the $z$-divergence by parts.  Since
$\Psi_k$ is compactly supported in $z$,
\[
 -\int\nabla_z\Psi_k\cdot Y_k
 =\int\Psi_k B_k^{\rm nl}
 +\varepsilon_k\int\Psi_k
       (B_k^{\rm lin}-\partial_sS_k-D_k).
\]
All displayed integrals here include time.  The product rule gives
$\zeta\nabla F_k=\nabla\Psi_k-F_k\nabla\zeta$, which supplies the
second summand in \eqref{eq:v2-complete-current}.  Integrating
$-\Psi_k\partial_sS_k$ in time gives exactly the endpoints and
$\partial_s\Psi_k$ term in \eqref{eq:v2-quadratic-stock}.  This proves
the first identity; \eqref{eq:v2-replay-discrepancy} proves the second.

The kernel, cutoff, and time-derivative bounds of
\Cref{lem:v2-primitive-bounds}, together with
$|\delta_zW|^2\le2(|W(y+z)|^2+|W(y)|^2)$, bound every term of
$\mathcal Q_k$ by $C\mathcal E_k$.  The time length is at most $3/8$.
Likewise
$|\nabla\delta_zW|^2\le2(|\nabla W(y+z)|^2+|\nabla W(y)|^2)$,
and all translated points lie in $G_1$.  The separation set has uniformly
bounded volume.  This proves the second bound without any trace estimate
for $\nabla W$ on a sharp spatial boundary.
\end{proof}

\begin{remark}[Why the smooth window has a replay discrepancy]
The spatial smoothing above avoids a viscous boundary trace which is not
controlled by a local $H^1$ norm.  It does not silently replace the heavy
cell by a different signed scalar: $\mathcal P_k$ still dominates its
actual coherent payment, while the entire replacement discrepancy is
$\mathcal R_k^{\rm rep}$.  A passing local-source incidence certificate
must also control this discrepancy before a signed local lower bound is
used.  It is not a second nonlinear occurrence.
\end{remark}

% =============================================================================
\section{The negative-primitive gate and the improved local rate}
\label{sec:v2-gate}
% =============================================================================

Define a positive two-point dissipation measure on the fixed window by
\begin{equation}
 \dd\mathfrak d_k(s,z)=\varepsilon_kD_k(s,z)\,\dd s\,\dd z,
 \qquad
 \mathcal B_k:=\mathcal R_k^{\rm rep}+\mathcal F_k.
 \label{eq:v2-positive-increment-measure}
\end{equation}
Only $\mathcal B_k$ is the complete scalar correction in the following
bounds; its constituent currents are never independently charged.

\begin{theorem}[Negative-primitive dissipation is necessary after the local corrections pass]
\label{thm:v2-negative-gate}
Every selected window satisfies
\begin{equation}
 \boxed{\mathcal P_k\le |\mathcal B_k|+
 C\varepsilon_k\mathcal E_k+
 \int(\Psi_k)_-\,\dd\mathfrak d_k.}
 \label{eq:v2-negative-gate}
\end{equation}
Consequently, if $\mathcal P_k\ge p_0$ and
\begin{equation}
 |\mathcal B_k|+C\varepsilon_k\mathcal E_k\le p_0/2,
 \label{eq:v2-gate-smallness}
\end{equation}
then
\begin{equation}
 \boxed{\int(\Psi_k)_-\,\dd\mathfrak d_k\ge p_0/2,
 \qquad \varepsilon_k\mathcal D_k\ge c p_0.}
 \label{eq:v2-dissipation-lower}
\end{equation}
For each fixed $\delta>0$, the first lower bound is carried, eventually,
only at relative separations
$|z|/\sqrt{-2s}>R_\Lambda-\delta$.
\end{theorem}

\begin{proof}
By \eqref{eq:v2-master-balance},
\[
 \mathcal P_k=\mathcal B_k+\varepsilon_k\mathcal Q_k
             -\int\Psi_k\,\dd\mathfrak d_k.
\]
The measure $\mathfrak d_k$ is positive, so
$-\int\Psi_k\dd\mathfrak d_k\le\int(\Psi_k)_-\dd\mathfrak d_k$.
Use \eqref{eq:v2-quadratic-bounds} to obtain
\eqref{eq:v2-negative-gate}.  The assumed smallness gives the first lower
bound.  Since $\zeta\ge0$ and $F_k$ is uniformly bounded,
\[
 \int(\Psi_k)_-\dd\mathfrak d_k
 \le C\varepsilon_k\int_{J_k}\int_{\supp_z\zeta}D_k\,\dd z\,\dd s
 \le C\varepsilon_k\mathcal D_k.
\]
This gives the second lower bound.  The support assertion is
\Cref{cor:v2-periodic-positive}.
\end{proof}

\begin{corollary}[Exclusion of a positive-primitive contact window]
\label{cor:v2-inner-exclusion}
Suppose the separation cutoff is supported in
$|z|/\sqrt{-2s}\le R_\Lambda-\delta$ for one fixed $\delta>0$.
If $\mathcal B_k\to0$ and
$\varepsilon_k\mathcal E_k\to0$, then
\begin{equation}
 \boxed{\limsup_{k\to\infty}\mathcal P_k\le0.}
 \label{eq:v2-inner-exclusion}
\end{equation}
Thus a fixed positive coherent payment in that window is impossible.
No upper bound for $\mathcal D_k$ is needed in this corollary.
\end{corollary}

\begin{proof}
For large $k$, $F_k\ge0$ throughout the support of $\zeta$.  Hence
$(\Psi_k)_-=0$.  Apply \eqref{eq:v2-negative-gate}.
\end{proof}

\begin{theorem}[Boundary-corrected quadratic rate alternative]
\label{thm:v2-rate-alternative}
Write the corresponding unnormalized local quantities as
\[
 E_k^U=\sup_{J_k}\int_{G_1}|U_k|^2=\rho_k^2\mathcal E_k,
 \qquad D_k^U=\int_{J_k}\int_{G_1}|\nabla U_k|^2
             =\rho_k^2\mathcal D_k.
\]
Then
\begin{equation}
 \boxed{\mathcal P_k\le |\mathcal B_k|
 +C\frac{\nu^2}{r_k}(E_k^U+D_k^U).}
 \label{eq:v2-rate-upper}
\end{equation}
On a heavy window with $\mathcal B_k\to0$,
\begin{equation}
 \boxed{E_k^U+D_k^U\ge c p_0\frac{r_k}{\nu^2}.}
 \label{eq:v2-improved-quadratic-rate}
\end{equation}
After subsequence selection, either $E_k^U\ge c p_0r_k/\nu^2$, or the
negative-primitive dissipation conclusion holds and
$D_k^U\ge c p_0r_k/\nu^2$.  In particular, when
$E_k^U=o(r_k/\nu^2)$ the latter branch is forced.
\end{theorem}

\begin{proof}
Use \eqref{eq:v2-master-balance} and
\eqref{eq:v2-quadratic-bounds}, and substitute
$\varepsilon_k\rho_k^{-2}=\rho_k^{-3}=\nu^2/r_k$.
This proves \eqref{eq:v2-rate-upper}; the heavy lower bound gives
\eqref{eq:v2-improved-quadratic-rate}.  Choose the kinetic threshold small
enough that its complement implies
$C\varepsilon_k\mathcal E_k\le p_0/4$.  For late $k$,
$|\mathcal B_k|\le p_0/4$, so
\Cref{thm:v2-negative-gate} applies on that complement.
\end{proof}

\begin{assessment}[What improves over V1]
V1 used an absolute cubic upper bound and obtained a local energy/dissipation
maximum of order $r_k^{2/3}$.  The rate in
\eqref{eq:v2-improved-quadratic-rate} is of order $r_k$ because the
Navier--Stokes evolution converts the signed cubic scalar into quadratic
stock and dissipation, after the complete local scalar correction passes.
This stronger hypothesis is explicit: a heavy positive replay alone still
only licenses V1's bound.  The positive-primitive exclusion and the
negative-primitive dissipation test evaluate the actual caloric coefficient;
they do not introduce a replacement ambient action.
\end{assessment}

% =============================================================================
\section{Local energy influx and the harmonic-pressure row}
\label{sec:v2-influx}
% =============================================================================

The pressure decomposition is imported from the main programme's
Theorem V34, ``Local pressure equals controlled near field plus harmonic
far field.''  Applying its homogeneous quadratic normalization gives
\begin{equation}
 P_k=P_k^{\rm loc}+P_k^{\rm harm}+c_k(s),\qquad
 P_k^{\rm loc}=\mathcal R_i\mathcal R_j(\vartheta W_{k,i}W_{k,j}),
 \label{eq:v2-local-pressure}
\end{equation}
where $\vartheta=1$ on the relevant inner collar.  The harmonic part is
harmonic there.  The same spatial Calder\'on--Zygmund estimate gives
$\|P_k^{\rm loc}(s)\|_{3/2}\le C\|W_k(s)\|_{L^3(G_2)}^2$.
No pressure compactness is asserted without control of the harmonic
oscillation.

\begin{proposition}[Physical upper bounds for the complete local current]
\label{prop:v2-pressure-current}
Let
\[
 X_k=\int_{J_k}\int_{G_2}|W_k|^3,\qquad
 H_k^{\rm harm}=\int_{J_k}\int_{G_1}
 |P_k^{\rm harm}-(P_k^{\rm harm})_{G_1}(s)|^{3/2},
\]
with the fixed domains enlarged as necessary to contain all pressure and
translation supports.  Then
\begin{equation}
 \boxed{|\mathcal F_k|\le
 C\left(X_k+(H_k^{\rm harm})^{2/3}X_k^{1/3}\right).}
 \label{eq:v2-current-upper}
\end{equation}
This is a bound, not a vanishing theorem.  In particular, bounded $X_k$ and
bounded harmonic-pressure oscillation do not make $\mathcal F_k$ small.
\end{proposition}

\begin{proof}
All kernel and cutoff coefficients in \eqref{eq:v2-complete-current} are
bounded.  The separation-edge term and the convective centre term are
bounded by $C\int|W_k|^3$, using translation of the fixed spatial supports
and H\"older's inequality.  For the pressure term,
\[
 \int|\delta_zP_k|\,|\delta_zW_k|
 \le C\|P_k-c_k^\sharp(s)\|_{L^{3/2}_{s,y}}
       \|W_k\|_{L^3_{s,y}},
\]
where the time-dependent pressure gauge is chosen to subtract the harmonic
mean.  Integrating over the bounded $z$-set changes only the constant.
The local pressure contribution is at most $CX_k$, by
\eqref{eq:v2-local-pressure}; the harmonic contribution is at most
$C(H_k^{\rm harm})^{2/3}X_k^{1/3}$.  Adding the estimates proves the
claim.  The estimate never assigns separate positive costs to the two
currents comprising $\mathcal F_k$.
\end{proof}

Choose a fixed $\psi\in C_c^\infty(G_2)$, $0\le\psi\le1$, with
$\psi=1$ on $G_1$.  Put
\begin{equation}
 e_{k,\psi}(s)=\frac12\int\psi|W_k|^2,\qquad
 \mathcal I_{k,\psi}=\int_{J_k}\int
 \left(\frac{|W_k|^2}{2}+P_k\right)W_k\cdot\nabla\psi.
 \label{eq:v2-energy-influx}
\end{equation}
The sign convention is such that positive $\mathcal I_{k,\psi}$ supplies
energy to the region where $\psi$ is large.  A time-dependent constant
added to $P_k$ does not change this quantity, since
$\int W_k\cdot\nabla\psi=0$.

\begin{theorem}[A persistent averaged dissipation requires incoming energy work]
\label{thm:v2-influx}
The selected smooth solution obeys the exact identity
\begin{equation}
 \boxed{
 \begin{aligned}
 \mathcal I_{k,\psi}=\varepsilon_k\bigg[
 &e_{k,\psi}(s_k^+)-e_{k,\psi}(s_k^-)\\
 &+\int_{J_k}\int\psi|\nabla W_k|^2
 -\frac12\int_{J_k}\int |W_k|^2\Delta\psi\bigg].
 \end{aligned}}
 \label{eq:v2-local-energy-identity}
\end{equation}
If the dissipation lower bound in \eqref{eq:v2-dissipation-lower} holds
and
\begin{equation}
 \varepsilon_k\sup_{s\in J_k}\int_{G_2}|W_k(s)|^2\longrightarrow0,
 \label{eq:v2-large-collar-energy-smallness}
\end{equation}
then
\begin{equation}
 \boxed{\liminf_{k\to\infty}\mathcal I_{k,\psi}\ge c p_0>0.}
 \label{eq:v2-influx-lower}
\end{equation}
Thus this dissipation branch cannot coexist with vanishing incoming local
kinetic--pressure energy work and the displayed endpoint/cutoff control.
\end{theorem}

\begin{proof}
Take the scalar product of \eqref{eq:v2-equation} with $W_k$ and use
incompressibility:
\[
 \partial_s\frac{|W_k|^2}{2}
 +\rho_k\operatorname{div}\!\left[
     \left(\frac{|W_k|^2}{2}+P_k\right)W_k\right]
 =\Delta\frac{|W_k|^2}{2}-|\nabla W_k|^2.
\]
Multiply by $\psi$ and integrate in space and time.  Integration by parts
and division by $\rho_k$ give
\eqref{eq:v2-local-energy-identity}.  The non-dissipative bracketed terms
have absolute value at most
$C\sup_{J_k}\int_{G_2}|W_k|^2$.  Since $\psi=1$ on $G_1$, the positive
dissipation term is at least $\mathcal D_k$.  Use
\eqref{eq:v2-dissipation-lower} and
\eqref{eq:v2-large-collar-energy-smallness}.
\end{proof}

\begin{remark}[Two interfaces, one physical energy history]
The two-point current $\mathcal F_k$ in the contact identity is not the
one-point energy influx $\mathcal I_{k,\psi}$.  Vanishing of the former
must not be substituted for vanishing of the latter.  Theorem
\ref{thm:v2-influx} instead proves the incidence: the dissipation forced by
the former balance must be supplied by the latter balance, or by a surviving
endpoint/cutoff kinetic term.  They are two tests of the same energy
history, not two additive source payments.  Formula
\eqref{eq:v2-local-pressure} identifies the unbounded pressure part which
must remain explicit in either test.
\end{remark}

% =============================================================================
\section{Fixed-datum consequences in the original physical variables}
\label{sec:v2-physical-consequences}
% =============================================================================

The following consequences use one fixed original solution.  No
varying-data family is substituted for the terminal chronology.

\begin{theorem}[The dissipation branch has summable $r_k\ell_k$]
\label{thm:v2-weighted-length}
Let $\mathcal K$ be a selected subsequence of the disjoint first-passage
intervals.  Suppose its heavy windows satisfy
\eqref{eq:v2-gate-smallness}, with one uniform $p_0>0$.  Then
\begin{equation}
 \boxed{\nu\int_{I_k}\int_{x_k+\ell_kG_1}|\nabla u|^2
 \ge c p_0 r_k\ell_k\qquad(k\in\mathcal K),}
 \label{eq:v2-physical-dissipation-lower}
\end{equation}
and therefore
\begin{equation}
 \boxed{\sum_{k\in\mathcal K}r_k\ell_k
 \le C_{p_0}\nu\int_0^T\|\nabla u(t)\|_2^2\,\dd t<\infty.}
 \label{eq:v2-weighted-length-sum}
\end{equation}
In particular $r_k\ell_k\to0$ along this branch.  If the unit norm records
are enumerated by $r_k=r_0+k$, this is also a weighted summability statement
$\sum_{k\in\mathcal K}k\ell_k<\infty$.
\end{theorem}

\begin{proof}
For the actual age subinterval corresponding to $J_k$, direct change of
variables gives
\[
 \nu\int\!\int_{x_k+\ell_kG_1}|\nabla u|^2
 =\nu^2\ell_k\rho_k^2\mathcal D_k
 =r_k\ell_k\,\varepsilon_k\mathcal D_k.
\]
The last equality uses $r_k=\nu^2\rho_k^3$.  Apply
\eqref{eq:v2-dissipation-lower} and enlarge the time integral to $I_k$.
The spatial domain is an embedded local chart for late $k$.  Since the
$I_k$ are disjoint, summation counts the original positive dissipation at
most once.  The Leray energy bound proves
\eqref{eq:v2-weighted-length-sum}.  Finitely many earlier packets can be
omitted without changing the conclusion.
\end{proof}

\begin{corollary}[Physical normalization of the incoming work]
\label{cor:v2-physical-influx}
Let $\psi_k(x)=\psi((x-x_k)/\ell_k)$.  On the actual age subinterval,
\begin{equation}
 \boxed{\int\!\int
 \left(\frac{|u|^2}{2}+p\right)u\cdot\nabla\psi_k\,\dd x\,\dd t
 =r_k\ell_k\,\mathcal I_{k,\psi}.}
 \label{eq:v2-physical-influx}
\end{equation}
The forced incoming work and the dissipation in
\eqref{eq:v2-physical-dissipation-lower} are therefore on the same
$r_k\ell_k$ physical scale.
\end{corollary}

\begin{proof}
The energy current contributes $\nu^3\rho_k^3\ell_k^{-3}$, the cutoff
gradient contributes $\ell_k^{-1}$, and $\dd t\dd x$ contributes
$\ell_k^5/\nu$.  Their product is
$\nu^2\rho_k^3\ell_k=r_k\ell_k$.
\end{proof}

\begin{lemma}[Weak terminal trace and terminal kinetic-energy defects]
\label{lem:v2-terminal-trace}
The fixed energy-bounded smooth solution has a weak $L^2$ terminal trace
$u_T$.  For every sequence $t_j\uparrow T$, a subsequence of the measures
$|u(t_j)|^2\dd x$ converges weakly-* on the torus to a positive finite
measure $\lambda$, and
\begin{equation}
 \boxed{\lambda=|u_T|^2\dd x+\lambda^{\rm conc},\qquad
 \lambda^{\rm conc}\ge0,\qquad
 \supp\lambda^{\rm conc}\subseteq\Sigma_T.}
 \label{eq:v2-terminal-energy-defect}
\end{equation}
No uniqueness, non-atomicity, or vanishing of $\lambda^{\rm conc}$ is
asserted.
\end{lemma}

\begin{proof}
For every smooth divergence-free test field $\phi$,
\[
 \left|\frac{\dd}{\dd t}\langle u(t),\phi\rangle\right|
 \le\nu\|u(t)\|_2\|\Delta\phi\|_2
      +\|u(t)\|_2^2\|\nabla\phi\|_\infty.
\]
The right side is uniform.  The scalar has a limit at $T$; boundedness in
$L^2$ and density of smooth solenoidal tests give one weak terminal trace.
For general tests use the solenoidal projection, since $u$ is solenoidal.
Bounded total energy gives weak-* subsequential compactness of the kinetic
measures.  For any continuous nonnegative $f$, weak lower semicontinuity
of $\|\sqrt f\,u\|_2^2$ yields
$\int f|u_T|^2\le\int f\dd\lambda$.  This gives the positive difference.
On every compact terminally regular region the solution converges smoothly
to its terminal trace, so the two measures agree there.  Exhaust the open
regular set by compact subsets to obtain the support statement.
\end{proof}

\begin{theorem}[The nonshrinking kinetic branch creates a terminal energy atom]
\label{thm:v2-kinetic-atom}
Suppose selected heavy near-singular cells are on the kinetic branch
\begin{equation}
 E_k^U\ge c_0r_k/\nu^2,\qquad
 \liminf_{k\to\infty}r_k\ell_k\ge a_0>0.
 \label{eq:v2-kinetic-branch}
\end{equation}
Then, after taking a subsequence, there exist $t_k\in I_k$ with $t_k\to T$,
a centre $x_*\in\Sigma_T$, and a terminal measure in
\eqref{eq:v2-terminal-energy-defect} for which
\begin{equation}
 \boxed{\lambda^{\rm conc}(\{x_*\})\ge c a_0>0.}
 \label{eq:v2-energy-atom}
\end{equation}
Consequently, if all terminal kinetic-energy cluster measures are
non-atomic, then $r_k\ell_k\to0$ also on the kinetic branch.  Strong
$L^2$ convergence to $u_T$ is one sufficient stronger hypothesis for this
non-atomicity, but is not supplied here.
\end{theorem}

\begin{proof}
The original kinetic mass is
\[
 \int_{x_k+\ell_kG_1}|u(t,x)|^2\,\dd x
 =\nu^2\ell_k\int_{G_1}|U_k(s,y)|^2\,\dd y.
\]
Choose $t_k$ approximating the supremum defining $E_k^U$ within a factor
of two.  Then the local kinetic mass is at least
$(c_0/2)r_k\ell_k\ge c a_0$.  The times lie in the selected first-passage
intervals and tend to $T$.  The near-singular condition and compactness of
the torus give $x_k\to x_*\in\Sigma_T$ along a subsequence.  Let
$\lambda$ be a cluster measure along those same times, as in
\Cref{lem:v2-terminal-trace}.  For every fixed $r>0$, the shrinking local
sets eventually lie inside $\overline B_r(x_*)$.  The closed-set
Portmanteau inequality therefore gives
$\lambda(\overline B_r(x_*))\ge c a_0$.  Let $r\downarrow0$.
The absolutely continuous part $|u_T|^2\dd x$ has no atoms, proving
\eqref{eq:v2-energy-atom}.  For the final assertion, any subsequence with
$r_k\ell_k$ bounded below would produce such an atom.  Strong $L^2$
convergence implies convergence of the squared densities in $L^1$ and
hence gives $\lambda^{\rm conc}=0$.
\end{proof}

\begin{assessment}[The physical-scale conclusion is not yet a contradiction]
The dissipation branch yields the actual finite-datum stock
\eqref{eq:v2-weighted-length-sum}; the kinetic branch returns a terminal
energy atom unless $r_k\ell_k$ vanishes.  Small $r_k\ell_k$ is compatible
with finite energy and is not excluded by V1's source-modulus saturation.
Thus neither a finite total dissipation measure nor singular-set support
alone closes this rate.  In particular no claim about the Hausdorff or
Minkowski dimension of $\Sigma_T$ has been used to erase an atom or to
infer a summable unweighted contact.
\end{assessment}

% =============================================================================
\section{The contact lives on a growing fast-time interval}
\label{sec:v2-fast-time}
% =============================================================================

The V1 fast-time equation is retained.  Its time-domain effect must also
be retained when a nonzero contact is used.  For $s_k^0\in J$, put
\[
 Z_k(\tau,y)=W_k(s_k^0+\tau/\rho_k,y),\qquad
 \widehat J_k=\rho_k(J_k-s_k^0).
\]
Then $Z_k$ solves Navier--Stokes with viscosity $\varepsilon_k$.

\begin{theorem}[Exact long-window contact transformation]
\label{thm:v2-long-window}
The selected contact is exactly
\begin{equation}
 \boxed{
 \begin{aligned}
 \mathcal C_k=-\frac1{4\rho_k}\int_{\widehat J_k}\int\!\int
 &\chi(y)\zeta(s_k^0+\tau/\rho_k,z)
 \nabla_zF_k(s_k^0+\tau/\rho_k,z)\cdot\delta_zZ_k(\tau,y)\\
 &\hspace{3em}\times|\delta_zZ_k(\tau,y)|^2
 \,\dd y\,\dd z\,\dd\tau.
 \end{aligned}}
 \label{eq:v2-fast-contact}
\end{equation}
The fast-time interval has length $\rho_k|J_k|$.  Thus a nondegenerate
slow-time cell is an interval of order $\rho_k$ in fast time.
\end{theorem}

\begin{proof}
Substitute $s=s_k^0+\tau/\rho_k$ in
\eqref{eq:v2-signed-window}.  The only time Jacobian is
$\dd s=\rho_k^{-1}\dd\tau$; the field and kernel arguments transform
as displayed.  No normalization of the contact removes this time factor.
\end{proof}

Let $\widehat\kappa_k$ denote the nonnegative coherent window measure
pushed forward to fast time, including its spatial and separation
variables.  Its mass is $\mathcal P_k$.

\begin{theorem}[A bounded fast window cannot retain the heavy payment under bounded cubic capacity]
\label{thm:v2-fast-window-exclusion}
For every fast-time interval $I\subset\widehat J_k$,
\begin{equation}
 \boxed{\widehat\kappa_k\{\tau\in I\}
 \le\frac C{\rho_k}\int_I\int_{G_1}|Z_k(\tau,y)|^3
 \,\dd y\,\dd\tau.}
 \label{eq:v2-fast-window-bound}
\end{equation}
Consequently, if for one fixed $L>0$
\begin{equation}
 \sup_k\sup_{\tau_0}\int_{\widehat J_k\cap(\tau_0-L,\tau_0+L)}
 \int_{G_1}|Z_k|^3\le M_L<\infty,
 \label{eq:v2-fixed-fast-capacity}
\end{equation}
then every moving fast interval of length $2L$ has coherent payment
$O(\rho_k^{-1})$.  If $\mathcal P_k\ge p_0$, the number of intervals of
length at most $2L$ required to cover its time support is at least
$c p_0\rho_k/M_L$.  In particular $\rho_k|J_k|$ must grow at least
linearly in $\rho_k$ under these hypotheses.

Conversely, if one bounded moving fast window carries a fixed payment,
its local fast-time cubic integral must grow at least like $\rho_k$.
\end{theorem}

\begin{proof}
Coherent replay is dominated by the total variation of the raw increment
measure.  In the rescaled coordinates that variation has density
$\frac14|\nabla F_k\cdot\delta W_k|\,|\delta W_k|^2$.
The uniform kernel bound, the bounded separation support, and translation
of the spatial domain bound it after integration by
$C\int_{G_1}|W_k|^3\dd y\dd s$.  Transforming to fast time proves
\eqref{eq:v2-fast-window-bound}.  Apply
\eqref{eq:v2-fixed-fast-capacity} and sum over a covering by such intervals
to obtain the count.  Covering $\widehat J_k$ by at most
$1+|\widehat J_k|/(2L)$ intervals gives the stated length consequence.
The last assertion follows by rearranging
\eqref{eq:v2-fast-window-bound}.
\end{proof}

\begin{firewall}[A bounded fast-time Euler limit is not the long-window contact]
The conditional Euler--Reynolds implication of V1 concerns a fixed compact
fast-time cylinder.  Theorem~\ref{thm:v2-long-window} shows that the
normalized contact uses a different object: a $\rho_k^{-1}$-weighted
integral on a growing fast-time interval.  Under
\eqref{eq:v2-fixed-fast-capacity}, every such compact cylinder has vanishing
payment, including cylinders whose time origins move with $k$.  Thus a
nonzero limit field on one compact fast-time window does not establish
nonzero contact in that limit.  The alternative is a long-time occupation
of those windows or a new cubic concentration within one window, not an
automatic compact-contact handoff.
\end{firewall}

% =============================================================================
\section{The slow-time stress and positive averaged dissipation limit}
\label{sec:v2-slow-limit}
% =============================================================================

Only the moments needed by the actual equation and its local energy balance
are retained here.  No path law, invariant probability measure, or complete
profile family is introduced.  Consider $W_k,P_k$ on an open cylinder
$\mathcal O$ containing the closure of the selected relative-time and
spatial window.  They remain the rescalings of the same original solution.

\begin{assumption}[Local moment bounds on the contact-normalized thread]
\label{ass:v2-local-moments}
For every compact subcylinder $K\Subset\mathcal O$, pressure gauges are
chosen so that
\begin{equation}
 \sup_k\left(\|W_k\|_{L^3(K)}+
                    \|P_k\|_{L^{3/2}(K)}\right)<\infty.
 \label{eq:v2-local-moment-bounds}
\end{equation}
These are additional bounds.  V1's lower cubic estimate does not provide
them.  By the inherited pressure decomposition, bounded local cubic mass
controls the near-pressure part, but harmonic-pressure oscillation remains
an independent requirement.
\end{assumption}

\begin{theorem}[Spatial equilibrium of the slow-time quadratic moments]
\label{thm:v2-slow-momentum}
Under \Cref{ass:v2-local-moments}, a subsequence has local weak limits
\[
 W_k\rightharpoonup W\quad\hbox{in }L^3,\qquad
 W_k\otimes W_k\rightharpoonup Q\quad\hbox{in }L^{3/2},\qquad
 P_k\rightharpoonup\overline P\quad\hbox{in }L^{3/2}.
\]
They satisfy
\begin{equation}
 \boxed{\operatorname{div}_yW=0,\qquad
 \operatorname{div}_yQ+\nabla_y\overline P=0.}
 \label{eq:v2-slow-equilibrium}
\end{equation}
The quadratic defect
\begin{equation}
 \boxed{\mathbb R=Q-W\otimes W\succeq0\quad\hbox{a.e.}}
 \label{eq:v2-positive-Reynolds}
\end{equation}
is a positive semidefinite symmetric tensor.  In particular
\[
 \operatorname{div}_y(W\otimes W+\mathbb R)+\nabla_y\overline P=0.
\]
This is a spatial Euler--Reynolds equilibrium constraint at almost every
slow time.  It is not the time-dependent Euler equation in the variable
$s$.  No equation for $\partial_sW$ follows from this argument.
\end{theorem}

\begin{proof}
Local weak compactness in the reflexive spaces and a diagonal extraction
give the displayed limits.  Divide \eqref{eq:v2-equation} by $\rho_k$:
\[
 \varepsilon_k\partial_sW_k+
 \operatorname{div}_y(W_k\otimes W_k)+\nabla_yP_k
 =\varepsilon_k\Delta_yW_k.
\]
Against a fixed compactly supported test field, the first and last terms
are bounded by $C\varepsilon_k\|W_k\|_{L^1(K)}$ and hence vanish.
The weak product and pressure limits give
\eqref{eq:v2-slow-equilibrium}; divergence passes linearly.  Since there
are no time derivatives in that constraint, testing first with products of
a temporal scalar and a spatial test gives the almost-every-time statement.

For nonnegative smooth $f$ and a fixed real vector $a$, weak $L^2$
convergence on its compact support gives
\[
 \int f\,a^{\mathsf T}Qa
 =\lim_k\int f|a\cdot W_k|^2
 \ge\int f|a\cdot W|^2.
\]
Using a countable dense set of vectors $a$ proves
\eqref{eq:v2-positive-Reynolds}.  The last limitation follows directly
from the vanishing coefficient of $\partial_sW_k$; it is not legitimate
to restore that coefficient after taking the limit.
\end{proof}

\begin{theorem}[Positive averaged dissipation and its energy-current incidence]
\label{thm:v2-slow-energy}
Under \Cref{ass:v2-local-moments}, the measures
\begin{equation}
 \mu_k:=\varepsilon_k|\nabla W_k|^2\,\dd s\,\dd y,
 \qquad
 J_k^{E}:=\left(\frac{|W_k|^2}{2}+P_k\right)W_k\,\dd s\,\dd y
 \label{eq:v2-energy-measures}
\end{equation}
are locally bounded in total variation; $\mu_k$ is nonnegative.
After a further subsequence they converge locally weakly-* to $\mu\ge0$
and a vector-valued Radon measure $J^E$.  The limiting equation is
\begin{equation}
 \boxed{\operatorname{div}_yJ^E+\mu=0.}
 \label{eq:v2-limit-energy}
\end{equation}
If the negative-primitive gate holds on selected windows in a common
compact subset of $\mathcal O$, then $\mu$ is nonzero on a slightly
larger compact subset.

If additionally $W_k\to W$ strongly in local $L^3$, then
$\mathbb R=0$ and
\begin{equation}
 \boxed{J^E=\left(\frac{|W|^2}{2}+\overline P\right)W\,\dd s\,\dd y.}
 \label{eq:v2-strong-energy-current}
\end{equation}
Neither strong $L^3$ convergence nor $\mathbb R=0$ is used here to assert
$\mu=0$; energy-flux conservation would be a further assertion.
\end{theorem}

\begin{proof}
Let $e_k=|W_k|^2/2$.  The smooth local energy identity used in
\Cref{thm:v2-influx}, divided by $\rho_k$, is
\begin{equation}
 \varepsilon_k\partial_se_k+
 \operatorname{div}_y\left[(e_k+P_k)W_k\right]
 =\varepsilon_k\Delta_ye_k-\varepsilon_k|\nabla W_k|^2.
 \label{eq:v2-slow-energy-equation}
\end{equation}
For a nonnegative $f\in C_c^\infty(\mathcal O)$, integration gives
\[
 \int f\,\dd\mu_k
 =\varepsilon_k\int e_k(\partial_sf+\Delta_yf)
   +\int(e_k+P_k)W_k\cdot\nabla_yf.
\]
The right side is uniformly bounded by
$C_f(\varepsilon_k\|W_k\|_2^2+
\|W_k\|_3^3+\|P_k\|_{3/2}\|W_k\|_3)$.
Choosing $f=1$ on any prescribed compact subset proves local boundedness
of the positive measures $\mu_k$.  H\"older gives local $L^1$ bounds for
the energy currents.  Weak-* compactness of Radon measures and a diagonal
subsequence give both limits.  The terms containing
$\varepsilon_ke_k$ vanish against every compactly supported test, which
proves \eqref{eq:v2-limit-energy}.

The gate gives $\mu_k(J_k\times G_1)\ge c p_0$.  Choose one nonnegative
smooth cutoff equal to one on a compact set containing all those windows
and supported in $\mathcal O$.  Its limiting integral is at least $c p_0$,
so $\mu\ne0$.  This argument retains possible concentration at the
selected time endpoints; it does not pass a moving discontinuous cutoff
through a weak limit.

Under strong local $L^3$ convergence the quadratic products converge
strongly in $L^{3/2}$ and the kinetic cubic currents in $L^1$.  The
pressure products converge in distributions because
$P_k\rightharpoonup\overline P$ in $L^{3/2}$ and $W_k\to W$ strongly
in $L^3$.  Hence \eqref{eq:v2-strong-energy-current} holds, with no cubic
concentration measure remaining.  This passage does not prove that the
spatial divergence of the limiting energy current is zero; its divergence
is exactly $-\mu$ by \eqref{eq:v2-limit-energy}.
\end{proof}

\begin{corollary}[No-influx exclusion on a complete spatial domain]
\label{cor:v2-no-influx}
Suppose the limit in \eqref{eq:v2-limit-energy} is defined on a complete
spatial domain and admits nonnegative spatial cutoffs $\psi_R\to1$ such
that for every compactly supported nonnegative temporal test $\theta$,
\begin{equation}
 \int\theta(s)\nabla\psi_R(y)\cdot\dd J^E(s,y)\longrightarrow0.
 \label{eq:v2-no-influx-condition}
\end{equation}
Then $\mu=0$.  In particular this holds on a fixed spatial torus, or on
$\R^3$ with finite total energy-current variation on each compact time
slab and standard cutoffs $|\nabla\psi_R|\le C/R$.  It contradicts the
nonzero averaged dissipation conclusion of \Cref{thm:v2-slow-energy}.
A local contact cylinder alone does not supply
\eqref{eq:v2-no-influx-condition}.
\end{corollary}

\begin{proof}
Test \eqref{eq:v2-limit-energy} against $\theta\psi_R$:
\[
 \int\theta\psi_R\,\dd\mu
 =\int\theta\nabla\psi_R\cdot\dd J^E.
\]
Positivity and Fatou's lemma, together with
\eqref{eq:v2-no-influx-condition}, give
$\int\theta\dd\mu=0$.  On a fixed torus use $\psi_R=1$ exactly.
Under the stated finite-variation hypothesis on $\R^3$, the right side
is bounded by $C\|J^E\|_{\rm TV}/R$.  No such bound is inferred merely
from local moment estimates on a growing rescaled torus.
\end{proof}

\begin{assessment}[The precise limit exported by the normalized thread]
\label{ass:v2-export}
The new finite-datum alternative is a kinetic trace concentration, a
surviving complete local scalar correction, or a negative-primitive
increment-dissipation packet.  On the last branch with the stated local
moment bounds, the export includes the positive measure $\mu$ and its
incoming energy-current incidence, together with $Q$ and its possible
Reynolds defect $\mathbb R$.  The two positive objects are different:
$\mathbb R$ measures loss in the quadratic velocity limit, whereas
$\mu$ is the surviving $\rho_k^{-1}$-weighted gradient dissipation.
The argument does not require $\mathbb R\ne0$.

The local energy-defect formulation is related to the classical
energy-flux analysis of Duchon--Robert and Eyink
\cite{v2-DR,v2-Eyink}; those works are background references, not a source
of the new $r_k$, $\ell_k$, and $\rho_k$ bounds proved above.  This
continuation uses the inherited KHM and pressure identities and derives
the displayed normalization-dependent conclusions directly.
\end{assessment}


% =============================================================================
\section{Absorbing the viscous term into the actual local energy flux}
\label{sec:v2-absorption}
% =============================================================================

The positive dissipation classifier above does not make dissipation an
independent source.  The next result takes one further upstream step:
ordinary local energy absorbs its entire contribution to the selected
contact.  It also removes a representational separation-edge current by
using a compact radial primitive of the selected vector writer itself.

\begin{lemma}[Compact primitive of an annular vector writer]
\label{lem:v2-annular-primitive}
Let $J^\circ\Subset(-\infty,0)$ contain $J$ in its interior, and let
$\vartheta(s,r)\ge0$ be smooth, compactly supported in
$J^\circ\times(\alpha_1,\beta_1)$, with $0<\alpha_1<\beta_1<\infty$.
Writing $l_s(r)=\mathcal L_{-2s}(z)$ for $|z|=r$, put
\begin{equation}
 \boxed{\Phi(s,z)=-\int_{|z|}^{\infty}
 \vartheta(s,r)\,\partial_rl_s(r)\,\dd r.}
 \label{eq:v2-annular-primitive}
\end{equation}
Then $\Phi$ is smooth, compactly supported in $J^\circ\times B_{\beta_1}$,
constant in $z$ near zero, and
\begin{equation}
 \boxed{\nabla_z\Phi(s,z)=\vartheta(s,|z|)
                         \nabla_z\mathcal L_{-2s}(z).}
 \label{eq:v2-primitive-vector-equality}
\end{equation}
All coefficient bounds needed below are uniform on the fixed window.
\end{lemma}

\begin{proof}
Differentiation in $r=|z|$ gives
$\partial_r\Phi=\vartheta\partial_rl_s$.  For $r<\alpha_1$ this
primitive is constant; for $r>\beta_1$ it is zero.  It is therefore smooth
at the origin as well as at the annular cutoffs.  The age is bounded away
from zero, so differentiating under the integral gives every displayed
coefficient bound and the compact time support.
\end{proof}

\begin{remark}[The sign radius is not a primitive-gauge-independent restriction]
The value $R_\Lambda$ is an exact coefficient of the full zero-mode-normalized
caloric primitive.  It is not a support restriction on every equivalent
representation of an annular vector writer.  The primitive
\eqref{eq:v2-annular-primitive} has the same selected gradient without the
term $\mathcal L\nabla\vartheta$ generated by a product cutoff, but has
a different scalar sign pattern.  Corresponding quadratic and boundary
terms change with it.  Thus the sign gate of
\Cref{thm:v2-negative-gate} is valid precisely for its displayed complete
balance; a primitive's sign may not be used after its compensating boundary
terms are omitted.  The absorption estimate below does not depend on the
sign of $\Phi$.
\end{remark}

For one solution $(W_k,P_k)$, choose a fixed smooth spatial cutoff $\chi$,
and write the pointwise increment quantities
\[
 q_k=\tfrac14|\delta_zW_k|^2,\quad
 Y_k^0=\tfrac14\delta_zW_k|\delta_zW_k|^2,\quad
 V_k^0=\tfrac14W_k|\delta_zW_k|^2
                  +\tfrac12\delta_zP_k\,\delta_zW_k,\quad
 d_k^0=\tfrac12|\nabla_y\delta_zW_k|^2.
\]
All following integrals have compact support supplied by their writers.
Define the selected principal contact and its centre flux by
\begin{align}
 \mathcal H_k^\Phi&=-\int\chi(y)\nabla_z\Phi(s,z)\cdot Y_k^0,
 \label{eq:v2-primitive-contact}\\
 \mathcal F_k^\Phi&=\int\Phi(s,z)V_k^0\cdot\nabla_y\chi(y).
 \label{eq:v2-primitive-centre-flux}
\end{align}
Choose one smooth nonnegative space--time cutoff $\omega(s,y)$ equal to
one on the spatial projections, and all bounded translations, of
$\supp(\chi\Phi)$.  It has compact support in a larger fixed cylinder
$\mathcal Q$.  Put
\begin{equation}
 \mathcal F_k^{E,\omega}=\int
 \left(\frac{|W_k|^2}{2}+P_k\right)W_k\cdot\nabla_y\omega,
 \qquad E_{2,k}=\int_{\mathcal Q}|W_k|^2.
 \label{eq:v2-energy-flux-compact}
\end{equation}
The kinetic and pressure parts are integrated together before taking the
positive part of this single scalar.

\begin{theorem}[Local-energy absorption of the complete viscous contact]
\label{thm:v2-absorbed-contact}
There is a constant depending only on the fixed cutoffs and coefficient
bounds such that
\begin{equation}
 \boxed{|\mathcal H_k^\Phi|
 \le |\mathcal F_k^\Phi|
      +C(\mathcal F_k^{E,\omega})_+
      +C\varepsilon_k E_{2,k}.}
 \label{eq:v2-absorbed-contact}
\end{equation}
No upper bound on $\int|\nabla W_k|^2$ is an assumption of this theorem.
In particular, vanishing of the two displayed flux terms and
$E_{2,k}=o(\rho_k)$ imply $\mathcal H_k^\Phi\to0$.
\end{theorem}

\begin{proof}
The inherited pointwise increment identity in the normalized variables is
\[
 \partial_sq_k+\rho_k\operatorname{div}_yV_k^0
       +\rho_k\operatorname{div}_zY_k^0
 =\Delta_yq_k-d_k^0.
\]
Test against $\chi\Phi$ and integrate by parts.  Compact time support
removes the endpoint terms.  The exact result is
\[
 \mathcal H_k^\Phi=\mathcal F_k^\Phi
 +\varepsilon_k\left[
  \int q_k(\partial_s+\Delta_y)(\chi\Phi)
  -\int\chi\Phi d_k^0\right].
\]
The first bracketed term is bounded in absolute value by $CE_{2,k}$.
The second is bounded by
$C\int\omega|\nabla W_k|^2$, using the increment inequality for gradients
and the fixed bounded separation volume.  The ordinary local energy
identity \eqref{eq:v2-slow-energy-equation}, tested with the nonnegative
$\omega$, gives exactly
\[
 \varepsilon_k\int\omega|\nabla W_k|^2
 =\varepsilon_k\int\frac{|W_k|^2}{2}
                 (\partial_s\omega+\Delta_y\omega)
   +\mathcal F_k^{E,\omega}
 \le C\varepsilon_k E_{2,k}+(\mathcal F_k^{E,\omega})_+.
\]
The left side is nonnegative.  Substitution proves
\eqref{eq:v2-absorbed-contact}, retaining the full viscous term and both
complete signed fluxes.  No gradient stock was assumed finite uniformly
in $k$.
\end{proof}

To compare this smooth principal contact with a literal coherent cell,
choose $\vartheta$ and $\chi$ equal to one on that cell in its relative
variables.  Let $\widetilde{\mathcal P}_k\ge p_0$ be the corresponding
nonnegative coherent window mass, and let
$\widetilde{\mathcal C}_k$ be the full periodic signed window integrated
on the smooth time support.  Extend the raw signed density by its literal
formula on that support, where the original smooth solution exists.
Any contribution outside the selected first-passage time support is kept
in
\begin{equation}
 \mathcal R_k^{\rm loc}:=
 \widetilde{\mathcal P}_k-\widetilde{\mathcal C}_k.
 \label{eq:v2-smooth-localization-residual}
\end{equation}
This residual includes replay, smoothing, and temporal-support incidence;
no claim of its vanishing follows from a heavy positive cell.

\begin{lemma}[The periodic principal remainder is negligible on the fixed-datum chronology]
\label{lem:v2-fixed-datum-remainder}
For the preceding fixed smooth window,
\begin{equation}
 \boxed{\widetilde{\mathcal C}_k-\mathcal H_k^\Phi\longrightarrow0.}
 \label{eq:v2-principal-error}
\end{equation}
This conclusion does not require a uniform upper bound for the normalized
local cubic packet.
\end{lemma}

\begin{proof}
The inherited periodic principal decomposition has a gradient remainder
bounded independently of the heat age on a fixed injectivity-radius ball.
All physical separations in this window are $O(\ell_k)$, so it applies for
late $k$.  In original variables the absolute remainder is bounded by
\[
 \frac{C\ell_k^3}{r_k}
 \int_{b_k-C_1\ell_k^2/\nu}^{b_k-c_1\ell_k^2/\nu}
       \|u(t)\|_3^3\,\dd t,
\]
with harmless adjustment of positive fixed constants $c_1,C_1$.
The factor $\ell_k^3$ is the physical separation volume; translation and
$|\delta_zu|^3\le4(|u(x+z)|^3+|u(x)|^3)$ give the estimate.
The fixed-datum Leray bounds imply
\[
 \|u(t)\|_3^3\le
 C\|u(t)\|_2^{3/2}\|\nabla u(t)\|_2^{3/2}
 \in L^1(0,T).
\]
The intervals lie in $(0,T)$ for late $k$ and shrink to $T$.  Absolute
continuity of this integral proves the limit, independently of how the
rescaled cubic mass behaves.  Finally
\eqref{eq:v2-primitive-vector-equality} identifies its Euclidean principal
term with $\mathcal H_k^\Phi$ exactly.
\end{proof}

\begin{theorem}[Zero-flux contact forces superlinear cubic concentration]
\label{thm:v2-superlinear-cubic}
Suppose $\widetilde{\mathcal P}_k\ge p_0>0$,
$\mathcal R_k^{\rm loc}\to0$, and
\begin{equation}
 |\mathcal F_k^\Phi|+(\mathcal F_k^{E,\omega})_+\longrightarrow0.
 \label{eq:v2-zero-two-flux}
\end{equation}
Then, on the fixed enlarged smooth cylinder $\mathcal Q$,
\begin{equation}
 \boxed{\int_{\mathcal Q}|W_k|^2\ge c\rho_k,\qquad
 \int_{\mathcal Q}|W_k|^3\ge c\rho_k^{3/2},\qquad
 \int_{\mathcal Q}|U_k|^3
 \ge c\left(\frac{r_k}{\nu^2}\right)^{3/2}.}
 \label{eq:v2-superlinear-cubic}
\end{equation}
In particular, a heavy coherent contact with bounded normalized local
cubic mass and vanishing localization discrepancy must retain at least
one of the two complete physical fluxes in
\eqref{eq:v2-zero-two-flux}.
\end{theorem}

\begin{proof}
By \eqref{eq:v2-smooth-localization-residual} and
\Cref{lem:v2-fixed-datum-remainder},
$|\mathcal H_k^\Phi|\ge p_0/2$ eventually.  The absorbed estimate and
\eqref{eq:v2-zero-two-flux} therefore give
$E_{2,k}\ge c\varepsilon_k^{-1}=c\rho_k$.
On the fixed finite-volume cylinder, H\"older's inequality gives
\[
 E_{2,k}\le |\mathcal Q|^{1/3}
             \left(\int_{\mathcal Q}|W_k|^3\right)^{2/3}.
\]
This proves the second lower bound.  Since $U_k=\rho_kW_k$, its cubic
integral is $\rho_k^3$ times the latter integral.  Thus it is at least
$c\rho_k^{9/2}=c(r_k/\nu^2)^{3/2}$.  The final assertion is the
contrapositive, and does not assume an independent dissipation budget.
\end{proof}

\begin{corollary}[Same-scale fusion without a cell-count debit]
\label{cor:v2-fused-contact}
Consider any finite cohort at the same age scale and endpoint, using the
same radial writer $\Phi$ and one common rescaled field.  Let its centre
cutoffs be $\chi_j$, and put $\chi_* =\sum_j\chi_j$.  Suppose
$\chi_*$ and an enclosing nonnegative energy cutoff $\omega_*$ have
uniform coefficient/derivative bounds of the type used in
\Cref{thm:v2-absorbed-contact}.  Then
\begin{equation}
 \boxed{
 \left|\sum_j\mathcal H_k^{\Phi,\chi_j}\right|
 \le |\mathcal F_k^{\Phi,\chi_*}|
 +C(\mathcal F_k^{E,\omega_*})_+
 +C\varepsilon_k\int_{\mathcal Q_*}|W_k|^2,}
 \label{eq:v2-fused-bound}
\end{equation}
with no factor equal to the number of cells.  The scalar corrections must
also be fused before this estimate is used.  Thus the same direct
vanishing test applies to a microdelocalized comparable-scale cohort once
its two actual fused fluxes and kinetic term are controlled.
\end{corollary}

\begin{proof}
The contact and increment-flux identities are linear in the centre cutoff.
Sum them before estimating: the centre derivative is
$\nabla\chi_* =\sum_j\nabla\chi_j$, so internal signed interfaces
cancel exactly wherever the fused cutoff is constant.  Apply
\Cref{thm:v2-absorbed-contact} once to $\chi_*$ and $\omega_*$.
Its constants depend on their derivative bounds and on the fixed radial
writer, not the number of summands or the spatial volume.  The remaining
kinetic integral is over the actual enlarged union.  It is not multiplied
by the cell count, and no claim is made that it is uniformly bounded.
\end{proof}

\begin{assessment}[Upstream endpoint after absorption]
Positive averaged dissipation is a useful exact classifier of the original
local balance, but is not the last unaccounted source.  The local energy
equation supplies its earlier incidence with incoming kinetic--pressure
work.  At bounded contact-normalized cubic mass the surviving scalar must
therefore reside in actual boundary flux or in its localization/replay
discrepancy.  If those rows vanish, the stronger
$(r_k/\nu^2)^{3/2}$ cubic concentration is forced.  This is the next
physical frontier; a fresh unweighted dissipation action or another
bounded-fast-window profile theorem would not evaluate it.
\end{assessment}

% =============================================================================
\section{Integrated endpoint, status ledger, and the next physical test}
\label{sec:v2-endpoint}
% =============================================================================

\begin{theorem}[Version V2 fixed-datum terminal alternative]
\label{thm:v2-integrated}
On any V1 heavy parabolic thread, retain the full periodic window and its
coherent payment.  After subsequence selection, the following statements
hold at their explicit hypotheses.
\begin{enumerate}[label=\textup{(V2.\arabic*)},leftmargin=5.2em]
\item A surviving replay/localization correction is represented by the single
      scalar $\mathcal B_k=\mathcal R_k^{\rm rep}+\mathcal F_k$.
      No signed local lower bound is inferred before this row is controlled.
\item When $\mathcal B_k\to0$, the selected quadratic local stock satisfies
      $E_k^U+D_k^U\gtrsim r_k/\nu^2$.  Either the kinetic term is of that
      order, or the negative-primitive dissipation gate holds.
\item If the kinetic term is $o(r_k/\nu^2)$, a window lying strictly inside
      $R_\Lambda$ heat lengths is excluded.  Any survivor forces
      $D_k^U\gtrsim r_k/\nu^2$ and the positive increment-dissipation
      measure is tested on the negative region of the actual primitive.
\item The dissipation branch satisfies $\sum r_k\ell_k<\infty$ on the
      disjoint original intervals.  Under the larger-collar kinetic control,
      it also forces incoming kinetic--pressure energy work of size
      $r_k\ell_k$.  A kinetic branch with $r_k\ell_k$ bounded below
      instead yields an atomic terminal kinetic-energy defect.
\item The normalized scalar is a $\rho_k^{-1}$-weighted growing-fast-time
      integral.  A fixed fast-time cylinder either has vanishing coherent
      payment or has local cubic integral at least of order $\rho_k$.
\item Under the stated local cubic and pressure bounds, the slow-time
      quadratic moments satisfy spatial Euler--Reynolds equilibrium and
      their energy current satisfies
      $\operatorname{div}_yJ^E+\mu=0$ with $\mu\ge0$.  The
      dissipation branch makes $\mu\ne0$.  A passing complete-domain
      no-influx condition excludes this branch.
\end{enumerate}
The inherited terminal source-modulus saturation remains available on a
source-aligned thread.  In addition, \Cref{thm:v2-superlinear-cubic,cor:v2-fused-contact} gives
the stronger zero-flux cubic-concentration alternative and its same-scale
fused version.  No improvement of that modulus, exclusion of all
local energy-influx configurations, or global regularity conclusion is
included in this theorem.
\end{theorem}

\begin{proof}
The exact local identity and replay distinction are
\Cref{thm:v2-exact-contact-balance}.  The quadratic rate and sign gate are
\Cref{thm:v2-rate-alternative,thm:v2-negative-gate,cor:v2-inner-exclusion}.
The original-variable conclusions are
\Cref{thm:v2-weighted-length,thm:v2-influx,cor:v2-physical-influx,thm:v2-kinetic-atom}.
The fast-time assertion is
\Cref{thm:v2-long-window,thm:v2-fast-window-exclusion}.  The limiting
assertions are
\Cref{thm:v2-slow-momentum,thm:v2-slow-energy,cor:v2-no-influx}.
The absorbed zero-flux alternative and its cohort version are
\Cref{thm:v2-absorbed-contact,thm:v2-superlinear-cubic,cor:v2-fused-contact}.
The final limitation preserves the hypotheses of the inherited V35 Dini
criterion and V1 source-incidence alternative.
\end{proof}

\begingroup\small
\begin{longtable}{p{0.25\textwidth}p{0.17\textwidth}p{0.50\textwidth}}
\caption{NS--A Version V2 proof-status ledger.}\label{tab:v2-status}\\
\toprule Row & Status & Exact content\\\midrule
\endfirsthead
\toprule Row & Status & Exact content\\\midrule
\endhead
V1 source and normalization & Preserved & Complete preterminal source retained; no earlier equation changed.\\[0.25em]
Periodic caloric primitive & Proved & Uniform primitive/time-derivative control; exact Dawson expression and certified unique radial zero.\\[0.25em]
Local contact balance & Proved & One complete nonlinear current plus $\rho_k^{-1}$ quadratic stock and signed positive-measure dissipation test.\\[0.25em]
Coherent versus signed contact & Explicit residual & $\mathcal R_k^{\rm rep}$ is retained; heavy positive replay is not a signed PDE lower bound.\\[0.25em]
Inner positive-primitive window & Conditional exclusion & No persistent positive payment when the complete correction and scaled kinetic stock vanish; no dissipation upper bound required.\\[0.25em]
Improved local quadratic rate & Proved implication & After local corrections pass, $E_k^U+D_k^U\gtrsim r_k/\nu^2$.\\[0.25em]
Dissipation chronology & Proved implication & On its selected branch $\sum r_k\ell_k<\infty$ by the actual disjoint Leray dissipation.\\[0.25em]
Kinetic terminal alternative & Proved implication & Nonshrinking $r_k\ell_k$ and kinetic escalation force a terminal energy atom; absence of atoms is not proved.\\[0.25em]
Incoming energy work & Proved implication & Positive averaged dissipation requires local kinetic--pressure influx or surviving kinetic endpoint/cutoff terms.\\[0.25em]
Fixed fast-time contact & Conditional exclusion & Bounded cubic capacity in every fixed moving fast window forces its coherent payment to vanish.\\[0.25em]
Slow-time moment landing & Conditional & Spatial equilibrium of the stress and positive averaged dissipation/energy-current balance under local cubic and pressure bounds.\\[0.25em]
Global no-influx test & Conditional & A complete-domain vanishing boundary-current test forces the averaged dissipation measure to vanish.\\[0.25em]
Viscous-term absorption & Proved & The local energy equation gives the two-flux plus spacetime kinetic-stock bound without a uniform dissipation assumption.\\[0.25em]
Zero-flux cubic growth & Proved implication & Passing localization and both flux tests force $\int|U_k|^3\gtrsim(r_k/\nu^2)^{3/2}$.\\[0.25em]
Same-scale microdelocalized cohort & Proved implication & Fuse signed contacts and currents first; the absorbed estimate has no cell-count factor.\\[0.25em]
Terminal Dini improvement & Open & Not derived from Leray bounds or from the new sign radius.\\[0.25em]
Microdelocalized and distance-separated branches & Open & V1 capacity bounds retained; no uniform bound for the complete cell family is acquired here.\\[0.25em]
NS--B donor estimate & Not claimed & Test only on the same selected rank-three finite-datum fields; $\mu\ne0$ does not prove $\mathbb R\ne0$.\\[0.25em]
Global regularity & Open & No global three-dimensional Navier--Stokes regularity claim.\\
\bottomrule
\end{longtable}
\endgroup

\begin{programme}[Version V3: evaluate the surviving physical current]
\label{prog:v2-next}
Preserve Versions V1--V2 verbatim.  The absorbed estimate of \Cref{thm:v2-absorbed-contact} is the first
continuation test: evaluate its two actual fused fluxes before treating
viscous action as an independent obstruction.  The next attack is the actual
kinetic--pressure influx on a source-aligned negative-primitive thread,
not another abstract compactness theorem.  First evaluate the complete
signed local correction $\mathcal B_k$, including the separation edge;
a nonvanishing value is an interface survivor, not dissipation.  On its
vanishing branch, distinguish the terminal kinetic atom from the
$D_k^U\gtrsim r_k/\nu^2$-scale dissipation alternative, with the notation
$D_k^U$ of \Cref{thm:v2-rate-alternative}.  Then test whether the same
fixed-datum incoming energy current is controlled by the terminal source
modulus, local harmonic-pressure decay, or an actual enclosing no-influx
balance.

Any fast-time extraction must retain the growing time horizon in
\eqref{eq:v2-fast-contact}.  Compactness on one bounded fast window is
insufficient unless its cubic mass grows enough to retain payment.  A
rank-three donor test may be opened on the selected finite-k thread after
the inherited dynamic rank-at-most-two and represented-return audits pass;
it must retain the distinction between the Reynolds tensor and the
positive averaged dissipation.  V1's microdelocalized and
singular-distance/scale-separated branches remain separate until their
own scalar corrections or source-modulus estimates are evaluated.
\end{programme}

% =============================================================================
\section{Exact verification record and append-only history}
\label{sec:v2-verification}
% =============================================================================

The following standalone Python code checks the rational enclosure used in
\Cref{thm:v2-kernel-sign}.  All sign decisions use fractions of integers;
no special-function implementation or floating-point root finder enters
the certificate.
\begingroup\footnotesize
\begin{verbatim}
from fractions import Fraction as F
from math import factorial

def kernel_sign_interval(radius: F, n: int = 32):
    q = radius * radius / 4
    assert F(1, 2) < q < n + 2
    e = sum((q**j / factorial(j) for j in range(n + 1)), F(0))
    s = sum((q**j / (factorial(j) * (2*j + 1))
             for j in range(n + 1)), F(0))
    tail = q**(n + 1) / factorial(n + 1) / (1 - q / (n + 2))
    lo = e - (2*q - 1) * (s + tail / (2*n + 3))
    hi = e + tail - (2*q - 1) * s
    return lo, hi

lo, hi = kernel_sign_interval(F(30039505, 10**7))
assert F(661, 10**10) < lo <= hi < F(662, 10**10)
lo, hi = kernel_sign_interval(F(30039506, 10**7))
assert F(-1149, 10**10) < lo <= hi < F(-1148, 10**10)
print("Certified: 3.0039505 < R_Lambda < 3.0039506")
\end{verbatim}
\endgroup

\paragraph{Append-only integrity.}
The complete V1 preterminal byte string has SHA--256
\begin{center}\scriptsize\ttfamily
338b1f56e968962d0e00789e1cbe4bc8baf78f4cecaf007b8570360258da8d4b
\end{center}
and is preserved exactly before the V2 material begins.  The only removed
active command is the final document terminator, which is restored after
this continuation.  The V1 title and historical status statements above
are therefore left unchanged.  The symbol printed as $ho_k$ in its
abstract is a typographical rendering of $\rho_k$, whose definition is
correctly displayed in \eqref{eq:contact-amplitude} and retained in
\eqref{eq:v2-normalization}.

\paragraph{Version V2 history.}
Evaluated the radial caloric primitive and its unique sign change with an
exact rational enclosure.  Inserted the correctly pressure-normalized KHM
balance into the contact-normalized thread and kept the coherent/signed
replay discrepancy explicit.  Proved the negative-primitive dissipation
gate and the conditional inner-window exclusion.  Strengthened the
boundary-corrected local quadratic rate to order $r_k$, derived the
fixed-datum weighted length sum and the terminal kinetic-atom alternative,
and identified the required incoming kinetic--pressure energy work.
Separated a fixed fast-time Euler limit from the growing-window contact.
Proved the conditional slow-time stress-equilibrium and positive averaged
dissipation limits, with their precise no-influx exclusion.  Absorbed the full viscous term with the ordinary local energy equation,
removed a representational annular-edge term by an exact compact radial
primitive, and obtained the stronger zero-flux superlinear cubic bound and
a count-free same-scale fused estimate.  The terminal
Dini estimate, unclosed interfaces, and unselected microdelocalized branches
remain open.

\begin{thebibliography}{9}
\bibitem{v2-DR}
J.~Duchon and R.~Robert,
\emph{Dissipation d'\'energie pour des solutions faibles des \'equations
d'Euler et Navier--Stokes incompressibles},
S\'eminaire Goulaouic--Schwartz (1999--2000), expos\'e 13, 1--10.
\url{https://www.numdam.org/item/SEDP_1999-2000____A13_0/}.

\bibitem{v2-Eyink}
G.~L.~Eyink,
\emph{Energy dissipation without viscosity in ideal hydrodynamics I.
Fourier analysis and local energy transfer},
Physica D \textbf{78} (1994), 222--240.
\href{https://doi.org/10.1016/0167-2789(94)90117-1}
{doi:10.1016/0167-2789(94)90117-1}.
\end{thebibliography}

For Version V3, preserve this complete source verbatim, remove only the
sole final active document terminator, append new material immediately
before it, and restore it.  The next filename is
\begin{center}\scriptsize
\path{Renewal_NS_Terminal_Source_Concentration_and_Singular_Thread_Closure_Programme_V3.tex}.
\end{center}

% =============================================================================
\addtocontents{toc}{\protect\endgroup}
% END VERSION V2 MATERIAL
% =============================================================================


% =============================================================================
% BEGIN VERSION V3 MATERIAL -- append-only continuation of NS--A V1--V2
% =============================================================================
\clearpage
\hypersetup{pdftitle={NS-A Terminal Source Concentration Programme: Cumulative V3},pdfauthor={A. Pelissier}}
\makeatletter
\addtocontents{toc}{\protect\begingroup\protect\makeatletter}
\addtocontents{toc}{\protect\def\protect\l@section{\protect\@dottedtocline{1}{0em}{3.5em}}}
\addtocontents{toc}{\protect\makeatother}
\makeatother
\renewcommand{\thesection}{V3.\arabic{section}}
\renewcommand{\thesubsection}{V3.\arabic{section}.\arabic{subsection}}
\renewcommand{\theequation}{V3.\arabic{equation}}
\setcounter{section}{0}
\setcounter{equation}{0}
\renewcommand{\theHsection}{V3.\arabic{section}}
\renewcommand{\theHsubsection}{V3.\arabic{section}.\arabic{subsection}}
\renewcommand{\theHtheorem}{V3.\arabic{section}.\arabic{theorem}}
\renewcommand{\theHequation}{V3.\arabic{equation}}
\newcommand{\symzero}{\operatorname{sym}_0}
\newcommand{\HSnorm}[1]{\lVert#1\rVert_{\mathrm{HS}}}

\begin{center}
{\LARGE\bfseries NS--A: Version V3}\\[0.5em]
{\large Mean-Centred Energy, the Exact Pressure-Curvature Card,\\
and Shrinking-Neighbourhood Replenishment}\\[0.5em]
{\normalsize 2 September 2026}
\end{center}
\addcontentsline{toc}{part}{Version V3: affine-null pressure flux and shrinking-neighbourhood replenishment}

\begin{abstract}
Version V2 reduced the selected caloric contact to two complete signed
boundary fluxes and a coupling-divided kinetic stock.  The present
continuation evaluates the pressure part of those fluxes.  A local
weighted-mean subtraction gives an exact relative energy identity in
which the derivative of the mean cancels.  The resulting energy flux and
the original two-point pressure flux both annihilate every affine spatial
pressure.  This is an annihilator of the two tests, not an enlargement of
the physical pressure gauge.  Consequently the unresolved harmonic
pressure enters through its Hessian rather than its value or gradient.

The far-pressure Hessian is calculated from the actual kinetic tensor.
Its Euclidean kernel is the fourth derivative of the Newtonian potential.
On the five-dimensional trace-free symmetric tensor space, its angular
operator has the exact eigenvalues $36,-24,-24,6,6$, multiplied by
$(4\pi |z|^5)^{-1}$.  The two physical pressure currents are paired with
this same Hessian through two explicit local tensor moments; a controlled
third-order Taylor remainder records all higher jets.  The operator is
signed, so no positive pressure price is inferred from its coefficients.

More importantly, the affine cancellation gives a pressure-tail estimate
using only the fixed solution's total kinetic energy.  For the normalized
global energy cap
$\mathsf E_k=E_*/(\nu^2\rho_k^2\ell_k)$, the sum of the two far-pressure
current magnitudes is at most $C\mathsf E_k^{3/2}R^{-5}$.  Taking
$R_k=R_0(1+\mathsf E_k)^{3/10}\log(e+\ell_k^{-1})$ makes that tail vanish,
while $\ell_kR_k\to0$.  No local pressure, cubic, or dissipation upper bound
is assumed.  A bounded normalized local cubic packet improves the exponent
$3/10$ to $1/5$.  Thus the selected fluxes can be evaluated, up to a
vanishing error, with pressure generated only inside one shrinking
physical neighbourhood of the same terminal point.

An exact transported-mean formulation also removes uniform sweeping from
both relative balances.  Its physical centre travels only $O(\sqrt{\ell_k})$
on a fixed relative-time window under the global energy bound.  Capture of
the original coherent payment by that moving tube remains a stated
incidence condition.  The remaining fixed-datum target is local signed
kinetic/near-pressure replenishment, mean-centred kinetic concentration,
or the original localization discrepancy.  The Dini improvement and
unconditional exclusion of those survivors are not claimed.
\end{abstract}

% =============================================================================
\section{One fixed-datum thread and the precise new acquisition}
\label{sec:v3-setting}
% =============================================================================

Keep the complete V1--V2 source, chronology, and normalization.  In
particular
\begin{equation}
 \rho_k=(r_k/\nu^2)^{1/3},\quad \varepsilon_k=\rho_k^{-1},\quad
 W_k(s,y)=\frac{\ell_k}{\nu\rho_k}
 u\!\left(b_k+\frac{\ell_k^2}{\nu}s,x_k+\ell_ky\right)
 \label{eq:v3-normalization}
\end{equation}
solves \eqref{eq:v2-equation}, with the pressure normalization
\eqref{eq:v2-pressure-normalization}.  All fields below are obtained from
this one solution, not from varying initial data.  Fix a compact
relative-time interval inside $(-\infty,0)$ containing the support of the
V2 smooth writer $\Phi$.  For late $k$ its physical image lies inside
$(0,T)$.  The distinction between that smooth window and the original
first-passage time support remains in $\mathcal R_k^{\rm loc}$ of
\eqref{eq:v2-smooth-localization-residual}.

Use the compact annular primitive $\Phi(s,z)$ of
\Cref{lem:v2-annular-primitive} and a fixed centre cutoff $\chi(y)$.
Choose $\psi\ge0$ smooth and compactly supported, with $\psi=1$ on a
neighbourhood of all points $y,y+z$ used by $\chi\Phi$.  Choose
$\theta\ge0$ smooth and compactly supported in negative relative time,
with $\theta=1$ near the time support of $\Phi$.  A fixed cylinder
$\mathcal Q=J^\circ\times B_L$ contains all supports and translations,
including those of $\psi$, $\nabla\psi$, and $\Delta\psi$.  Enlarging $L$
by a fixed factor below is harmless.  Cutoffs and $L$ do not depend on $k$.

Write $V_\psi=\int\psi>0$ and define the actual local mean and fluctuation
\begin{equation}
 b_k(s)=V_\psi^{-1}\int\psi(y)W_k(s,y)\,\dd y,
 \qquad v_k(s,y)=W_k(s,y)-b_k(s).
 \label{eq:v3-local-mean}
\end{equation}
Thus $\int\psi v_k=0$ at every time.  Let
\begin{equation}
 E_*:=\sup_{0<t<T}\|u(t)\|_2^2<\infty,\qquad
 \mathsf E_k:=\frac{E_*}{\nu^2\rho_k^2\ell_k}.
 \label{eq:v3-global-energy-cap}
\end{equation}
On the rescaled torus $\Torus_k^3=\ell_k^{-1}\Torus^3$,
\begin{equation}
 \int_{\Torus_k^3}|W_k(s)|^2\le\mathsf E_k.
 \label{eq:v3-rescaled-energy}
\end{equation}
This identity follows directly by spatial change of variables.  It holds
throughout the chosen relative-time window, irrespective of local
concentration.

\begin{assessment}[What is inherited and what is being evaluated]
The local pressure decomposition, the caloric primitive, the KHM identity,
and V2's viscous absorption are inherited.  The new work is the common
affine-pressure cancellation of the two flux tests, the actual signed
five-coordinate curvature kernel, and its fixed-datum shrinking-radius
estimate.  No uniform cell action, global inviscid compactness, terminal
non-atomicity, or NS--B paid-stress variation is postulated.  The
five-coordinate curvature tensor below is not identified with the donor
programme's five-coordinate paid kinetic-stress path.
\end{assessment}

% =============================================================================
\section{Mean-centred energy and the common affine-pressure annihilator}
\label{sec:v3-relative-energy}
% =============================================================================

Keep the selected contact $\mathcal H_k^\Phi$ and the complete two-point
centre flux $\mathcal F_k^\Phi$ of
\eqref{eq:v2-primitive-contact}--\eqref{eq:v2-primitive-centre-flux}.
Replace only the auxiliary one-point energy test by
\begin{equation}
 \mathcal I_k^{\rm rel}
 :=\int\theta(s)
 \left(\frac{|v_k|^2}{2}W_k+P_kv_k\right)
 \cdot\nabla\psi\,\dd y\,\dd s,
 \qquad
 K_k^{\rm rel}:=\int_{\mathcal Q}|v_k|^2.
 \label{eq:v3-relative-flux}
\end{equation}
Convection still uses the actual velocity $W_k$; subtracting its local
mean is not permission to delete sweeping across a fixed boundary.

\begin{theorem}[Exact local mean-centred energy identity]
\label{thm:v3-relative-energy}
For every finite $k$,
\begin{equation}
 \boxed{
 \varepsilon_k\int\theta\psi|\nabla W_k|^2
 =\mathcal I_k^{\rm rel}
 +\varepsilon_k\int\frac{|v_k|^2}{2}
       (\theta'\psi+\theta\Delta\psi).}
 \label{eq:v3-relative-energy}
\end{equation}
In particular,
\begin{equation}
 \varepsilon_k\int\theta\psi|\nabla W_k|^2
 \le (\mathcal I_k^{\rm rel})_++C\varepsilon_kK_k^{\rm rel}.
 \label{eq:v3-relative-dissipation}
\end{equation}
There is no term involving $|b_k'|$, and no bound on that derivative is
an assumption.
\end{theorem}

\begin{proof}
For readability omit $k$.  Since $b$ is spatially constant,
$\operatorname{div}v=0$ and $\nabla v=\nabla W$.  Taking the scalar
product of \eqref{eq:v2-equation} with $v$ gives
\[
 \partial_s\frac{|v|^2}{2}
 +\rho\operatorname{div}_y
       \left(\frac{|v|^2}{2}W+Pv\right)
 =\Delta_y\frac{|v|^2}{2}-|\nabla W|^2-b'(s)\cdot v.
\]
Multiply by $\theta\psi$, integrate, and divide by $\rho$.
The last term vanishes exactly because
$\int\psi v=0$ at every time.  Compact support of $\theta$ removes
endpoints, and integration by parts gives
\eqref{eq:v3-relative-energy}.  The quadratic term is bounded by
$CK^{\rm rel}$; the dissipation is nonnegative.  These observations give
\eqref{eq:v3-relative-dissipation}.  All calculations take place on smooth
finite-$k$ fields, so differentiating the local mean is legitimate even
when its derivative has no uniform estimate.
\end{proof}

Separate pressure only for the purpose of identifying its linear test:
\begin{align}
 \mathcal L_k^\Phi[p]
 &=\frac12\int\Phi(s,z)\delta_zp(s,y)\,
                    \delta_zW_k(s,y)\cdot\nabla\chi(y)
                    \,\dd y\,\dd z\,\dd s,
 \label{eq:v3-pressure-test-two}\\
 \mathcal L_k^E[p]
 &=\int\theta(s)p(s,y)v_k(s,y)\cdot\nabla\psi(y)
                    \,\dd y\,\dd s.
 \label{eq:v3-pressure-test-one}
\end{align}
The full nonlinear fluxes, not their separately positive parts, remain the
quantities used in the contact bound.

\begin{theorem}[Both pressure tests annihilate affine spatial fields]
\label{thm:v3-affine-null}
For arbitrary smooth functions $a(s)\in\R^3$ and $c(s)\in\R$,
\begin{equation}
 \boxed{\mathcal L_k^\Phi[c+a\cdot y]=0,
        \qquad\mathcal L_k^E[c+a\cdot y]=0.}
 \label{eq:v3-affine-null}
\end{equation}
Hence any affine Taylor jet of a harmonic pressure component may be removed
from both tests exactly, with no compensating estimate.
\end{theorem}

\begin{proof}
For the first test, $\delta_z(c+a\cdot y)=a\cdot z$ is independent of
$y$.  For each $s,z$, incompressibility gives
\[
 \int\delta_zW_k\cdot\nabla\chi
 =-\int\chi\operatorname{div}_y\delta_zW_k=0.
\]
For the second test, integration by parts gives
\[
 \int(c+a\cdot y)v_k\cdot\nabla\psi
 =-a\cdot\int\psi v_k=0.
\]
Integration over the remaining variables proves both assertions.  The
identities are local, so it suffices for the displayed affine field to be
defined on a neighbourhood of the test supports.
\end{proof}

\begin{firewall}[A test annihilator is not an enlarged pressure gauge]
The Navier--Stokes pressure is determined, on the periodic solution, only
up to a function of time.  An affine pressure with nonzero gradient is a
real force in the equation.  Theorem~\ref{thm:v3-affine-null} says that
these two specific integrated tests annihilate that force's affine
potential after the energy mean has been centred.  It neither changes the
solution nor removes the force from its momentum equation.  The uncentred
one-point energy influx of V2 does not in general annihilate an affine
pressure: its affine contribution is $-a\cdot\int\psi W_k$.
\end{firewall}

\begin{theorem}[Relative-energy absorption on the original selected contact]
\label{thm:v3-relative-absorption}
With constants independent of the mean and its derivative,
\begin{equation}
 \boxed{|
 \mathcal H_k^\Phi|
 \le |\mathcal F_k^\Phi|
   +C(\mathcal I_k^{\rm rel})_+
   +C\varepsilon_kK_k^{\rm rel}.}
 \label{eq:v3-relative-absorption}
\end{equation}
All pressure contributions in the two complete fluxes in this inequality
have the common affine annihilator \eqref{eq:v3-affine-null}.
\end{theorem}

\begin{proof}
Use V2's exact compact-writer identity, not a new KHM calculation:
\[
 \mathcal H_k^\Phi=\mathcal F_k^\Phi+
 \varepsilon_k\left[
  \int\tfrac14|\delta_zW_k|^2
                  (\partial_s+\Delta_y)(\chi\Phi)
 -\int\tfrac12\chi\Phi|\nabla_y\delta_zW_k|^2\right].
\]
At each time $\delta_zW_k=\delta_zv_k$.  The first bracketed term is
therefore bounded by $CK_k^{\rm rel}$, with translated supports contained
in $\mathcal Q$.  The absolute value of the last term is at most
$C\int\theta\psi|\nabla W_k|^2$ by the same increment estimate and
bounded separation volume used in V2.  Apply
\eqref{eq:v3-relative-dissipation}.  The two affine cancellations are
\Cref{thm:v3-affine-null}.
\end{proof}

\begin{lemma}[Energy and cubic bounds for the centred tests]
\label{lem:v3-mean-bounds}
Uniformly on the fixed window,
\begin{equation}
 |b_k(s)|\le C\mathsf E_k^{1/2},\qquad
 \int_{B_L}|v_k(s)|^2\le C\mathsf E_k,
 \qquad
 K_k^{\rm rel}\le C\mathsf E_k.
 \label{eq:v3-mean-energy-bounds}
\end{equation}
If $X_k:=\int_{\mathcal Q}|W_k|^3$, then
\begin{equation}
 \int_{\mathcal Q}|v_k|^3\le CX_k,
 \qquad
 \int_{\mathcal Q}|v_k|\le CX_k^{1/3}.
 \label{eq:v3-mean-cubic-bounds}
\end{equation}
Here $\mathcal Q$ includes the support of $\psi$ and a fixed time interval;
no local cubic upper bound is asserted by these formulas.
\end{lemma}

\begin{proof}
Weighted Cauchy--Schwarz in \eqref{eq:v3-local-mean} bounds
$|b_k|$ by $C\|W_k\|_{L^2(\supp\psi)}$.  Use
$|W_k-b_k|^2\le2|W_k|^2+2|b_k|^2$ and
\eqref{eq:v3-rescaled-energy} on the fixed-volume ball.  For the cubic
bounds, Jensen gives
$|b_k|^3\le V_\psi^{-1}\int\psi|W_k|^3$.
The elementary cubic difference inequality and then H\"older on the fixed
cylinder give \eqref{eq:v3-mean-cubic-bounds}.
\end{proof}

% =============================================================================
\section{The actual far-pressure source and its first surviving jet}
\label{sec:v3-pressure-source}
% =============================================================================

The earlier near-field plus harmonic-far-field decomposition is retained.
Here it is made source explicit for the two affine-null tests.  Let
$G_k$ be the mean-zero Green function of $-\Delta$ on $\Torus_k^3$,
so that $-\Delta G_k=\delta_0-|\Torus_k^3|^{-1}$.
Taking the divergence of \eqref{eq:v2-equation} gives
\begin{equation}
 -\Delta P_k=\partial_i\partial_j(W_{k,i}W_{k,j}),\qquad
 P_k=\partial_i\partial_jG_k*(W_{k,i}W_{k,j})+c_k(s).
 \label{eq:v3-pressure-poisson}
\end{equation}
Repeated indices are summed from one to three.  This is the actual
periodic pressure, including its periodic Green kernel; a separately
chosen harmonic field is not substituted.

For $R\ge4L$ and $2R$ below a fixed fraction of the rescaled injectivity
radius, choose one radial cutoff $\eta_R$, equal to one on $B_R$ and zero
outside $B_{2R}$, with $0\le\eta_R\le1$.  Split
\begin{align}
 P_{k,R}^{\rm near}
 &=\partial_i\partial_jG_k*(\eta_R W_{k,i}W_{k,j}),
 \label{eq:v3-near-pressure}\\
 P_{k,R}^{\rm far}(s,y)
 &=\int_{\Torus_k^3}(1-\eta_R(z))
       \partial_i\partial_jG_k(y-z)W_{k,i}(s,z)W_{k,j}(s,z)\,\dd z.
 \label{eq:v3-far-pressure}
\end{align}
The second integral has no diagonal singularity for $y\in B_L$ and is
harmonic there.  The first is interpreted by the usual periodic
Calder\'on--Zygmund distribution.  Their sum is $P_k-c_k(s)$ exactly.
This is an observation-side splitting of the one tensor $W_k\otimes W_k$;
it does not truncate or recompute the Navier--Stokes evolution.

\begin{lemma}[Uniform periodic Green derivative estimates]
\label{lem:v3-Green-derivatives}
For $m=4,5$ and $z\ne0$,
\begin{equation}
 |\nabla^mG_k(z)|\le C_m d_k(z,0)^{-m-1},
 \label{eq:v3-Green-derivatives}
\end{equation}
where $d_k$ is the rescaled torus distance.  On the centred Voronoi fundamental
cell, away from zero,
\begin{equation}
 \nabla^4G_k(z)=\nabla^4\frac1{4\pi|z|}+\mathcal E_k^{(4)}(z),
 \qquad |\mathcal E_k^{(4)}(z)|\le C\ell_k^5.
 \label{eq:v3-periodic-Green-remainder}
\end{equation}
The fourth-order tensor is trace free in either contracted index pair
away from the diagonal.
\end{lemma}

\begin{proof}
On the fixed original flat torus its Green function equals
$(4\pi|z|)^{-1}$ plus a smooth function in a ball about zero.  This follows
by subtracting a cut-off Newtonian potential from its distributional
Poisson equation; the remainder solves an equation with smooth right-hand
side there.  Away from zero the Green function is smooth on a compact set.
Consequently $|\nabla^mG(z)|\le C_m d(z,0)^{-m-1}$ on the fixed torus.
Since $G_k(z)=\ell_kG(\ell_kz)$, differentiation gives
\eqref{eq:v3-Green-derivatives}.

For the sharper uniform remainder, the fourth derivative can be
periodized directly:
\[
 \nabla^4G_k(z)=\sum_{\gamma\in\ell_k^{-1}\Gamma}
                 \nabla^4\frac1{4\pi|z+\gamma|},\qquad z\ne0.
\]
One justification is $G_k=\int_0^\infty(H_{k,t}-|\Torus_k^3|^{-1})\dd t$:
after four derivatives, periodize the heat kernel and interchange sum and
integral.  The absolute integral of each Gaussian derivative is bounded
by $C|z+\gamma|^{-5}$, a summable lattice bound away from the diagonal.
On a centred Voronoi fundamental cell all nonzero images are at least a
fixed multiple of $\ell_k^{-1}$ away, with the corresponding lattice
multiplicity bound.  Their sum is $O(\ell_k^5)$, proving
\eqref{eq:v3-periodic-Green-remainder}.
Finally $\Delta G_k=|\Torus_k^3|^{-1}$ away from zero; two further
spatial derivatives of that constant vanish.  This proves both pairwise
trace identities.
\end{proof}

\begin{theorem}[Affine-subtracted far pressure is an inverse-fifth-power tail]
\label{thm:v3-far-tail}
Set
\begin{equation}
 a_{k,R}(s)=\nabla P_{k,R}^{\rm far}(s,0),\qquad
 q_{k,R}(s,y)=P_{k,R}^{\rm far}(s,y)
 -P_{k,R}^{\rm far}(s,0)-a_{k,R}(s)\cdot y.
 \label{eq:v3-affine-subtracted-pressure}
\end{equation}
Then
\begin{equation}
 \boxed{\|q_{k,R}(s)\|_{L^\infty(B_L)}
       \le C_L R^{-5}\mathsf E_k.}
 \label{eq:v3-affine-tail}
\end{equation}
In particular,
\begin{equation}
 \boxed{|
 \mathcal L_k^\Phi[P_{k,R}^{\rm far}]|
 +|\mathcal L_k^E[P_{k,R}^{\rm far}]|
 \le C R^{-5}\mathsf E_k^{3/2}.}
 \label{eq:v3-energy-only-pressure-tail}
\end{equation}
There is also the sharper local-cubic form
\begin{equation}
 \boxed{|
 \mathcal L_k^\Phi[P_{k,R}^{\rm far}]|
 +|\mathcal L_k^E[P_{k,R}^{\rm far}]|
 \le C R^{-5}\mathsf E_k X_k^{1/3}.}
 \label{eq:v3-cubic-pressure-tail}
\end{equation}
No local pressure estimate or gradient upper bound is a premise.
\end{theorem}

\begin{proof}
For $y\in B_L$ and any source point in \eqref{eq:v3-far-pressure},
$d_k(y,z)\ge R-L\ge3R/4$.  Differentiating that nonsingular integral twice
and applying \eqref{eq:v3-Green-derivatives} gives
\[
 \sup_{B_L}|\nabla_y^2P_{k,R}^{\rm far}(s)|
 \le C R^{-5}\int|W_k(s,z)|^2\dd z
 \le C R^{-5}\mathsf E_k.
\]
Taylor's integral formula proves \eqref{eq:v3-affine-tail}.
By \Cref{thm:v3-affine-null}, replace the far pressure in both tests by
$q_{k,R}$.  The pressure increments in the first test are at most twice
its sup norm.  Also $\delta_zW_k=\delta_zv_k$, so the bounded $z$-volume
and translated supports give
\[
 |\mathcal L_k^\Phi[P_{k,R}^{\rm far}]|
 +|\mathcal L_k^E[P_{k,R}^{\rm far}]|
 \le C R^{-5}\mathsf E_k\int_{\mathcal Q}|v_k|.
\]
Use the energy bounds of \Cref{lem:v3-mean-bounds} and Cauchy--Schwarz on
the fixed cylinder for \eqref{eq:v3-energy-only-pressure-tail}.  Use
\eqref{eq:v3-mean-cubic-bounds} for \eqref{eq:v3-cubic-pressure-tail}.
All pressure jets were removed by exact annihilation, not by assigning
positive costs to them.
\end{proof}

% =============================================================================
\section{Evaluation of the five-coordinate pressure-curvature card}
\label{sec:v3-tidal-card}
% =============================================================================

For a matrix $A$, put
$\symzero A=(A+A^{\mathsf T})/2-(\operatorname{tr}A)I/3$.
For symmetric $A$ also write $A^\circ=\symzero A$.
Let $N(z)=(4\pi|z|)^{-1}$ and let $n=z/|z|$.

\begin{theorem}[Exact signed angular operator of the far-pressure Hessian]
\label{thm:v3-tidal-symbol}
For every symmetric trace-free matrix $S$ define
\begin{equation}
 \boxed{\mathcal T_nS
 =6S-30\bigl(n\otimes Sn+Sn\otimes n\bigr)
 +(105n\otimes n-15I)(n^{\mathsf T}Sn).}
 \label{eq:v3-tidal-operator}
\end{equation}
Then
\begin{equation}
 \partial_a\partial_b\partial_i\partial_jN(z)A_{ij}
 =\frac1{4\pi|z|^5}(\mathcal T_n A^\circ)_{ab}
 \label{eq:v3-Newtonian-card}
\end{equation}
for every symmetric $A$.  The map $\mathcal T_n$ is self-adjoint on
$\operatorname{Sym}_0(3)$ and has the complete orthogonal decomposition
\begin{equation}
 \boxed{\mathcal T_n=36\Pi_{0,n}-24\Pi_{1,n}+6\Pi_{2,n},
 \qquad(\operatorname{rank}\Pi_{0,n},\operatorname{rank}\Pi_{1,n},
 \operatorname{rank}\Pi_{2,n})=(1,2,2).}
 \label{eq:v3-tidal-spectrum}
\end{equation}
In particular its Hilbert--Schmidt operator norm is $36$, and
\begin{equation}
 \langle S,\mathcal T_nS\rangle_{\rm HS}
 =36\|\Pi_{0,n}S\|_{\rm HS}^2
  -24\|\Pi_{1,n}S\|_{\rm HS}^2
  +6\|\Pi_{2,n}S\|_{\rm HS}^2.
 \label{eq:v3-signed-tidal-form}
\end{equation}
\end{theorem}

\begin{proof}
Four differentiations of $|z|^{-1}$ give
\begin{align*}
 4\pi\partial_{abij}N(z)
 ={}&105\frac{z_az_bz_iz_j}{|z|^9}
   +3\frac{\delta_{ab}\delta_{ij}+\delta_{ai}\delta_{bj}
                         +\delta_{aj}\delta_{bi}}{|z|^5}\\
 &-15|z|^{-7}\bigl(
 \delta_{ab}z_iz_j+\delta_{ai}z_bz_j+\delta_{aj}z_bz_i
 +\delta_{bi}z_az_j+\delta_{bj}z_az_i+\delta_{ij}z_az_b\bigr).
\end{align*}
Contraction with a trace-free $S$ is exactly
\eqref{eq:v3-tidal-operator}.  Contracting $i,j$ with $I$ gives
$\partial_{ab}\Delta N=0$ away from zero, so the isotropic part of $A$
vanishes.  Symmetry of the fourth derivative under interchange of the
pairs $(a,b)$ and $(i,j)$ proves self-adjointness.

Here are explicit projectors, which also verify the spectrum without any
numerical eigenvalue calculation.  Write $a=n^{\mathsf T}Sn$ and
$w=(I-n\otimes n)Sn$.  Then
\[
 \Pi_{0,n}S=\frac{3a}{2}(n\otimes n-I/3),\qquad
 \Pi_{1,n}S=n\otimes w+w\otimes n,\qquad
 \Pi_{2,n}S=S-\Pi_{0,n}S-\Pi_{1,n}S.
\]
Their ranges are, respectively, the axisymmetric trace-free line, the two
normal--tangential mixed directions, and the two tangential trace-free
directions.  They are mutually orthogonal in the Hilbert--Schmidt product;
for the last range $Sn=0$.  Substitution in
\eqref{eq:v3-tidal-operator} gives $36S,-24S,6S$ on these three ranges.
Their dimensions sum to five, so they exhaust the space.  The norm and
quadratic-form identities follow.
\end{proof}

The actual curvature card is
\begin{equation}
 H_{k,R}(s):=\nabla_y^2P_{k,R}^{\rm far}(s,0)
 \in\operatorname{Sym}_0(3).
 \label{eq:v3-actual-Hessian}
\end{equation}
Let $\mathcal V_k$ be the centred Voronoi fundamental cell of the rescaled torus.

\begin{corollary}[Literal kinetic-tensor acquisition and periodic error]
\label{cor:v3-actual-tidal-card}
The same finite-$k$ velocity determines
\begin{equation}
 \boxed{
 H_{k,R}(s)=\frac1{4\pi}\int_{\mathcal V_k}
 \frac{1-\eta_R(z)}{|z|^5}\,
 \mathcal T_{z/|z|}\bigl((W_k\otimes W_k)^\circ(s,z)\bigr)\,\dd z
 +\mathfrak E_{k,R}(s),}
 \label{eq:v3-actual-card}
\end{equation}
where
\begin{equation}
 \|\mathfrak E_{k,R}(s)\|_{\rm HS}\le C\ell_k^5\mathsf E_k,
 \qquad
 \|H_{k,R}(s)\|_{\rm HS}\le C R^{-5}\mathsf E_k.
 \label{eq:v3-tidal-card-errors}
\end{equation}
The periodic remainder also annihilates isotropic source tensors.  No
source covariance is fitted and no pressure field is freely selected.
\end{corollary}

\begin{proof}
Differentiate \eqref{eq:v3-far-pressure} twice at zero.  The fourth
Green derivative has zero trace in the source index pair by
\Cref{lem:v3-Green-derivatives}.  Insert
\eqref{eq:v3-periodic-Green-remainder} and then
\eqref{eq:v3-Newtonian-card}; its angular kernel is even in $z$, so the
argument $-z$ has the same value.  Integrating the image remainder against
$|W_k|^2$ gives its displayed bound.  The full-kernel estimate from the
proof of \Cref{thm:v3-far-tail} gives the last bound.  Both the full kernel
and its Euclidean term annihilate isotropic tensors, hence so does their
difference.
\end{proof}

\begin{assessment}[A signed response, not a new positive source price]
The local tensor $(W_k\otimes W_k)^\circ$ is a deterministic part of the
already occurring quadratic source.  The coefficients in
\eqref{eq:v3-tidal-spectrum} are genuinely signed.  Positivity of kinetic
energy does not make the curvature operator positive or determine the
sign of its two scalar boundary responses.  In particular the number five
is the dimension of this trace-free pressure jet, not a count of new
nonlinear occurrences.  Equality of that dimension with
$\dim\operatorname{Sym}_0(3)$ in the NS--B paid-stress reduction does not
identify the two histories or their time derivatives.
\end{assessment}

Define two local response tensors using the actual finite-$k$ field:
\begin{align}
 M_k^\Phi(s)&=-\frac12\symzero\int\!\int
   \chi(y)\Phi(s,z)\,\delta_zW_k(s,y)\otimes z\,\dd y\,\dd z,
 \label{eq:v3-two-point-moment}\\
 M_k^E(s)&=-\theta(s)\symzero\int
                 \psi(y)v_k(s,y)\otimes y\,\dd y.
 \label{eq:v3-energy-moment}
\end{align}

\begin{theorem}[Both far-pressure fluxes use the same curvature card]
\label{thm:v3-common-curvature-pairing}
For $\bullet\in\{\Phi,E\}$,
\begin{equation}
 \boxed{\mathcal L_k^\bullet[P_{k,R}^{\rm far}]
 =\int H_{k,R}(s):M_k^\bullet(s)\,\dd s
   +\mathfrak r_k^\bullet(R),}
 \label{eq:v3-common-tidal-pairing}
\end{equation}
with the controlled higher-jet remainder
\begin{equation}
 |\mathfrak r_k^\Phi(R)|+|\mathfrak r_k^E(R)|
 \le C R^{-6}\mathsf E_k^{3/2}.
 \label{eq:v3-higher-jet-remainder}
\end{equation}
One may replace the factor $\mathsf E_k^{3/2}$ on the right by
$\mathsf E_kX_k^{1/3}$.  Thus the approximation retains both actual mixed
pressure/velocity pairings, rather than replacing them by the norm of
$H_{k,R}$ alone.
\end{theorem}

\begin{proof}
By \eqref{eq:v3-Green-derivatives} with $m=5$,
$\|\nabla_y^3P_{k,R}^{\rm far}(s)\|_{L^\infty(B_L)}
\le C R^{-6}\mathsf E_k$.
Taylor expansion about zero therefore writes the pressure as its affine
jet, $\tfrac12y^{\mathsf T}H_{k,R}y$, and a remainder bounded by
$C_LR^{-6}\mathsf E_k$ on the whole test domain.  The affine jet is
annihilated by \Cref{thm:v3-affine-null}.  For the quadratic term,
\[
 \nabla_y\delta_z\bigl(\tfrac12y^{\mathsf T}Hy\bigr)=Hz.
\]
Integrating the $y$-derivative off $\chi$ in
\eqref{eq:v3-pressure-test-two}, using
$\operatorname{div}\delta_zW_k=0$, gives
$-\tfrac12\int\chi\Phi\,\delta_zW_k\cdot Hz$.
Likewise the one-point pressure term becomes
$-\int\theta\psi v_k\cdot Hy$.
Since $H$ is symmetric and trace free, these are precisely the pairings
with \eqref{eq:v3-two-point-moment} and \eqref{eq:v3-energy-moment}.
The pressure remainder is estimated exactly as in
\Cref{thm:v3-far-tail}, now with $R^{-6}$ in place of $R^{-5}$.
\end{proof}

\begin{remark}[Angular cancellation precedes annular estimation]
Formula \eqref{eq:v3-actual-card} can be partitioned into physical dyadic
source annuli.  At each time the exact sum of the signed angular tensors
is formed first.  An optional upper bound is
\[
 C\sum_{j\ge0}(2^jR)^{-5}
 \int_{\{2^jR\le|z|<2^{j+1}R\}}|W_k(s,z)|^2\dd z
 +C\ell_k^5\mathsf E_k,
\]
with the cutoff factor included and the last annulus truncated by
$\mathcal V_k$.  These are disjoint source regions of one observation,
not independently chargeable pressure events.  Nor does a nonzero norm in
one annulus ensure a nonzero assembled scalar after the signed angular
and temporal pairings have been made.
\end{remark}

% =============================================================================
\section{An unconditional shrinking physical radius for pressure replenishment}
\label{sec:v3-shrinking-radius}
% =============================================================================

The radius $R$ so far is measured in the original cell's rescaled
coordinates.  The next choice is made from the actual first-passage
parameters and the one fixed energy bound, before any far-pressure value
is inspected.

\begin{theorem}[Energy-only removal of pressure sources outside a shrinking neighbourhood]
\label{thm:v3-shrinking-pressure}
Let $\ell_k\to0$, $\rho_k\to\infty$, and retain the original fixed
solution.  Set
\begin{equation}
 L_k^{\log}=\log(e+\ell_k^{-1}),\qquad
 R_k=R_0(1+\mathsf E_k)^{3/10}L_k^{\log},\qquad
 d_k=\ell_kR_k,
 \label{eq:v3-energy-radius}
\end{equation}
where the fixed $R_0$ is sufficiently large compared with $L$.  Then
\begin{equation}
 \boxed{R_k\to\infty,\qquad d_k\to0,}
 \label{eq:v3-radius-limits}
\end{equation}
and the two actual far-pressure currents satisfy
\begin{equation}
 \boxed{|
 \mathcal L_k^\Phi[P_{k,R_k}^{\rm far}]|
 +|\mathcal L_k^E[P_{k,R_k}^{\rm far}]|
 \le C(L_k^{\log})^{-5}\longrightarrow0.}
 \label{eq:v3-vanishing-far-flux}
\end{equation}
Thus both pressure tests may be evaluated using only the tensor source
inside $B_{2d_k}(x_k)$ in physical variables, up to a vanishing error.
No local cubic, pressure, enstrophy, or dissipation upper bound is required.
\end{theorem}

\begin{proof}
The logarithm diverges, so $R_k\to\infty$.  Since $0<3/10<1$,
\[
 d_k\le C L_k^{\log}
 \left(\ell_k+
 E_*^{3/10}\nu^{-3/5}\rho_k^{-3/5}\ell_k^{7/10}\right)
 \longrightarrow0.
\]
Here $\nu,E_*$ are fixed, $\rho_k\ge1$ eventually, and every positive
power of $\ell_k$ dominates its logarithm.  Hence for large $k$ the
source cutoff fits in the original injectivity-radius ball, as required
for the decomposition.  By \eqref{eq:v3-energy-only-pressure-tail},
\[
 \text{sum of far-current magnitudes}
 \le C\frac{\mathsf E_k^{3/2}}
 {(1+\mathsf E_k)^{3/2}(L_k^{\log})^5}
 \le C(L_k^{\log})^{-5}.
\]
The near pressure uses $\eta_{R_k}W_k\otimes W_k$, supported within
$2R_k$ rescaled lengths.  Multiplication by $\ell_k$ gives the stated
physical source neighbourhood.  This bound holds on every fixed smooth
relative-time window; it does not invoke disjointness or repeated
summation of the original energy.
\end{proof}

\begin{corollary}[Improved radius on a bounded normalized cubic branch]
\label{cor:v3-cubic-radius}
If additionally $\sup_kX_k<\infty$, one may instead take
\begin{equation}
 \widehat R_k=R_0(1+\mathsf E_k)^{1/5}L_k^{\log},\qquad
 \widehat d_k=\ell_k\widehat R_k.
 \label{eq:v3-cubic-radius}
\end{equation}
The same tail conclusion holds, and
\begin{equation}
 \widehat d_k\le C L_k^{\log}
 \left(\ell_k+E_*^{1/5}\nu^{-2/5}
                  \rho_k^{-2/5}\ell_k^{4/5}\right)\longrightarrow0.
 \label{eq:v3-cubic-physical-radius}
\end{equation}
This sharper exponent is conditional on the local cubic upper bound;
the energy-only theorem is not.
\end{corollary}

\begin{proof}
Insert \eqref{eq:v3-cubic-radius} in
\eqref{eq:v3-cubic-pressure-tail}.  The quotient
$\mathsf E_k/(1+\mathsf E_k)$ is at most one.  The physical-radius
estimate follows by raising \eqref{eq:v3-global-energy-cap} to the
power $1/5$ and using the same power--logarithm comparison.
\end{proof}

\begin{corollary}[Macroscopic harmonic pressure cannot be the surviving flux source]
\label{cor:v3-macroscopic-pressure}
On a heavy near-singular thread with $x_k\to x_*\in\Sigma_T$, all the
nonvanishing pressure contribution to either test is generated, up to
$o(1)$, inside a physical neighbourhood shrinking to $x_*$.  In
particular a pressure source remaining at a fixed positive physical
distance from $x_k$ has vanishing contribution to both tests.
\end{corollary}

\begin{proof}
Use \Cref{thm:v3-shrinking-pressure} and $d_k\to0$.
Alternatively fix a positive physical distance $d$ below injectivity
radius and take $R=d/\ell_k$.  The tail bound becomes
\[
 C\mathsf E_k^{3/2}\ell_k^5d^{-5}
 =C E_*^{3/2}\nu^{-3}\rho_k^{-3}\ell_k^{7/2}d^{-5}
 \longrightarrow0.
\]
The centre convergence then gives the same terminal point.  The statement
concerns pressure source locality for the selected tests, not local
regularity of the velocity on the shrinking neighbourhood.
\end{proof}

\begin{firewall}[The pressure radius is not a bounded scale-ratio certificate]
The radii above tend to zero in physical units but satisfy
$R_k=d_k/\ell_k\to\infty$.  Thus they eliminate a genuinely remote
harmonic-pressure input; they do not confine the remaining donor to a
fixed number of cell scales.  A mesoscopic source between $\ell_k$ and
$d_k$ is still permitted.  On the singular-distance/scale-separated
branch, absolute approach of a centre to $\Sigma_T$ likewise does not
bound its distance measured in either $\ell_k$ or $d_k$.  Those geometric
comparisons remain explicit and are not inferred from pressure locality.
\end{firewall}

% =============================================================================
\section{The reduced same-scalar replenishment test}
\label{sec:v3-near-flux}
% =============================================================================

For the radius in \Cref{thm:v3-shrinking-pressure}, define two full signed
near-source fluxes by replacing pressure, and pressure only:
\begin{align}
 \mathcal F_{k,\rm near}^\Phi
 &=\int\Phi(s,z)\left(
   \tfrac14W_k(y)|\delta_zW_k|^2
  +\tfrac12\delta_zP_{k,R_k}^{\rm near}\,\delta_zW_k\right)
     \cdot\nabla\chi(y),
 \label{eq:v3-near-two-flux}\\
 \mathcal I_{k,\rm near}^{\rm rel}
 &=\int\theta(s)\left(
  \tfrac12|v_k|^2W_k+P_{k,R_k}^{\rm near}v_k\right)
     \cdot\nabla\psi(y).
 \label{eq:v3-near-energy-flux}
\end{align}
All missing integration variables have the measures in
\eqref{eq:v3-pressure-test-two}--\eqref{eq:v3-pressure-test-one}.
These fluxes retain the signed kinetic--pressure cancellation.  No
positive part of a kinetic or pressure constituent is taken separately.

\begin{theorem}[Fixed-datum near-pressure reduction of the selected contact]
\label{thm:v3-local-replenishment}
Every selected smooth contact satisfies
\begin{equation}
 \boxed{|
 \mathcal H_k^\Phi|
 \le |\mathcal F_{k,\rm near}^\Phi|
 +C(\mathcal I_{k,\rm near}^{\rm rel})_+
 +C\varepsilon_kK_k^{\rm rel}
 +C(L_k^{\log})^{-5}.}
 \label{eq:v3-local-replenishment}
\end{equation}
If its actual coherent window mass obeys
$\widetilde{\mathcal P}_k\ge p_0>0$, and
$\mathcal R_k^{\rm loc}\to0$, then the following conclusions hold.
If $K_k^{\rm rel}=o(\rho_k)$,
\begin{equation}
 \liminf_k\left(
 |\mathcal F_{k,\rm near}^\Phi|
 +C(\mathcal I_{k,\rm near}^{\rm rel})_+\right)\ge p_0/2.
 \label{eq:v3-near-flux-required}
\end{equation}
If instead these two flux terms tend to zero, then
\begin{equation}
 K_k^{\rm rel}\ge c\rho_k,\qquad
 \int_{\mathcal Q}|v_k|^3\ge c\rho_k^{3/2}.
 \label{eq:v3-relative-concentration}
\end{equation}
In particular the latter growth cannot consist solely of a spatially
uniform mean velocity on the selected cylinder.
\end{theorem}

\begin{proof}
The differences between the complete fluxes in
\eqref{eq:v3-relative-absorption} and the corresponding near-source
fluxes are exactly
$\mathcal L_k^\Phi[P_{k,R_k}^{\rm far}]$ and
$\mathcal L_k^E[P_{k,R_k}^{\rm far}]$; the time-only pressure gauge is
annihilated.  Use \eqref{eq:v3-vanishing-far-flux}, the triangle
inequality, and the one-Lipschitz property of $x\mapsto x_+$ to prove
\eqref{eq:v3-local-replenishment}.

V2's fixed-datum periodic-remainder theorem and its explicit localization
residual give
$\widetilde{\mathcal P}_k=\mathcal H_k^\Phi+
\mathcal R_k^{\rm loc}+o(1)$.
Thus $|\mathcal H_k^\Phi|\ge p_0-o(1)$ on the passing localization
branch.  The first conclusion follows immediately if
$\varepsilon_kK_k^{\rm rel}\to0$.  If both near flux terms vanish,
the same inequality instead gives $K_k^{\rm rel}\ge c\rho_k$.
H\"older on the fixed-volume cylinder gives the cubic lower bound for
$v_k$.  By construction $v_k$ has zero $\psi$-mean, proving the final
qualification.  It remains true that
$\int|U_k|^3\gtrsim\rho_k^{9/2}$, by
\eqref{eq:v3-mean-cubic-bounds}; this is V2's rate now forced in local
fluctuations as well.
\end{proof}

\begin{assessment}[Exact scope of the acquisition]
V2 left an unspecified harmonic-pressure contribution in each flux.  The
preceding theorem proves that its genuinely remote part is negligible on
every fixed-datum thread, without an assumed pressure compactness bound.
The remaining pressure is the response to the literal tensor inside
$B_{2d_k}(x_k)$, and its separated far portion has the evaluated curvature
card \eqref{eq:v3-actual-card}.  The theorem does not evaluate the sign of
the remaining near-source current or show it vanishes.  A near-pressure
flux is not an independent forcing: it is a nonlocal response of the same
quadratic velocity history within a shrinking region.
\end{assessment}

\begin{proposition}[Signed fusion of the pressure-card tests]
\label{prop:v3-fusion}
For a finite family of centre cutoffs $\chi_j$ using the same $\Phi$,
common field, and common near/far pressure split,
\[
 \sum_j\mathcal L_k^{\Phi,\chi_j}[p]
 =\mathcal L_k^{\Phi,\sum_j\chi_j}[p],\qquad
 \sum_jM_k^{\Phi,\chi_j}=M_k^{\Phi,\sum_j\chi_j}.
\]
Consequently internal interfaces cancel before the common curvature card
is paired or estimated.  The mean-centred energy test of a fused cohort
uses one mean computed from its single enclosing $\psi$, not the sum of
independently centred one-cell energy fluxes.
\end{proposition}

\begin{proof}
Both the two-point pressure functional and its tensor moment are linear
in $\chi$.  Sum before integration by parts to obtain the two identities.
The one-point relative energy proof uses
$\int\psi(W-b)=0$ with the mean computed from that same $\psi$; separate
means need not agree, so their fluxes are not substituted for it.
\end{proof}

\begin{firewall}[Diameter is retained on a delocalized cohort]
The pressure-tail constants in this version refer to an explicitly fixed
relative spatial test domain $B_L$.  A microdelocalized cohort with
unbounded rescaled diameter must retain the corresponding Taylor powers,
cutoff derivative bounds, and spatial-volume factors.  Linear fusion
removes a spurious cell-count multiplier but does not give a uniform
fixed-diameter estimate on an arbitrarily large union.  V1's aggregate
capacity theorem and V2's correctly assembled cohort identity remain in
force at their own hypotheses; no complete microdelocalization exclusion
is inferred here.
\end{firewall}

% =============================================================================
\section{Uniform sweeping can be removed on a literally transported tube}
\label{sec:v3-material-tube}
% =============================================================================

The preceding theorem uses the original fixed centre.  This final
calculation identifies exactly which additional incidence is needed when
uniform sweeping is to be removed as well.  It does not assume that the
old heavy cell is already captured by a material tube.

On the rescaled torus, solve the smooth centre equation
\begin{equation}
 \xi_k'(s)=\rho_k\beta_k(s),\qquad
 \beta_k(s)=V_\psi^{-1}\int\psi(y)
                  W_k(s,y+\xi_k(s))\,\dd y,
 \qquad \xi_k(s_0)=0,
 \label{eq:v3-centre-ODE}
\end{equation}
for a fixed $s_0\in J^\circ$.  The convolution defining its vector field
is smooth for each finite $k$, so a unique solution exists on the compact
time interval; use a continuous lift of the torus path.  Put
\begin{equation}
 V_k(s,y)=W_k(s,y+\xi_k(s))-\beta_k(s),\qquad
 \Pi_k(s,y)=P_k(s,y+\xi_k(s)).
 \label{eq:v3-material-fields}
\end{equation}
Then $\int\psi V_k=0$ at every time.

\begin{theorem}[Mean-transported contact and energy balances]
\label{thm:v3-material-absorption}
The transported fields satisfy the exact equation
\begin{equation}
 \partial_sV_k+\rho_k\bigl((V_k\cdot\nabla)V_k+\nabla\Pi_k\bigr)
 =\Delta V_k-\beta_k'(s),\qquad\operatorname{div}V_k=0.
 \label{eq:v3-material-equation}
\end{equation}
The spatially constant acceleration cancels from the increment equation
and from the $\theta\psi$-integrated relative energy equation.  Hence,
with the contact and fluxes evaluated in the fixed moving coordinates,
\begin{equation}
 \boxed{|\mathcal H_{k,\rm mov}^\Phi|
 \le |\mathcal F_{k,\rm mov}^\Phi|
 +C(\mathcal I_{k,\rm mov}^{\rm rel})_+
 +C\varepsilon_k\int_{\mathcal Q}|V_k|^2.}
 \label{eq:v3-material-absorption}
\end{equation}
Here the convective part of the two-point flux is
$\tfrac14V_k|\delta_zV_k|^2$, and that of the energy flux is
$\tfrac12|V_k|^2V_k$.  Thus neither contains uniform sweeping.
Both pressure tests still annihilate affine fields.  The energy-only
shrinking-pressure theorem applies along this path, using the global
energy of the translated original $W_k$.
\end{theorem}

\begin{proof}
The chain rule and \eqref{eq:v3-centre-ODE} give
\[
 \partial_sV_k=(\partial_sW_k)(s,y+\xi_k)
       +\rho_k\beta_k\cdot\nabla W_k(s,y+\xi_k)-\beta_k'.
\]
Substitution of \eqref{eq:v2-equation} proves
\eqref{eq:v3-material-equation}.  Subtracting its values at $y+z$ and
$y$ removes $-\beta_k'$ exactly.  Taking its scalar product with $V_k$
and integrating against $\theta\psi$ also removes that term because
$\int\psi V_k=0$.  The proof of
\Cref{thm:v3-relative-absorption} now applies with convective velocity
$V_k$; there is no separate moving-cutoff estimate.  Integration by parts
proves the same affine-pressure cancellations.

Translation preserves the actual global energy cap for $W_k$.
The local mean bound implies
$\int_{B_L}|V_k|^2\le C\mathsf E_k$, uniformly in the centre.
The pressure Poisson source can be kept as the translated actual tensor
$W_k\otimes W_k$; no altered initial datum or truncated evolution is
used.  Repeating the pressure-tail proof therefore gives
\eqref{eq:v3-vanishing-far-flux} along the path.  The time dependence of
the source split is harmless here: its coefficients are never
differentiated in time in either pressure test.
\end{proof}

\begin{theorem}[The transported centre stays in a shrinking physical neighbourhood]
\label{thm:v3-travel}
On the fixed relative-time interval,
\begin{equation}
 \boxed{\sup_{s\in J^\circ}
 \ell_k|\xi_k(s)-\xi_k(s_0)|
 \le C\nu^{-1}E_*^{1/2}\ell_k^{1/2}\longrightarrow0.}
 \label{eq:v3-physical-travel}
\end{equation}
Thus when $x_k\to x_*\in\Sigma_T$, the entire physical moving tube and
its energy-only pressure-source neighbourhood shrink to $x_*$.  Their
diameter need not be comparable with $\ell_k$.
\end{theorem}

\begin{proof}
Weighted Cauchy--Schwarz gives
$|\beta_k(s)|\le C\mathsf E_k^{1/2}$, uniformly in its translated
centre.  Integrating \eqref{eq:v3-centre-ODE} over the fixed time interval
bounds the lifted displacement by $C\rho_k\mathsf E_k^{1/2}$.
Multiply by $\ell_k$ and use
\eqref{eq:v3-global-energy-cap}:
\[
 \ell_k\rho_k\mathsf E_k^{1/2}
 =\nu^{-1}E_*^{1/2}\ell_k^{1/2}.
\]
This proves the estimate.  Adding the $O(\ell_k)$ test radius and the
radius $d_k\to0$ of \Cref{thm:v3-shrinking-pressure} proves the last
claim.  The estimate concerns a lifted path's total possible displacement,
so periodic wrapping introduces no ambiguity.
\end{proof}

\begin{definition}[Material capture of the selected coherent scalar]
Let $\mathcal P_{k,\rm mov}$ be the coherent mass tested with the actual
moving spatial cutoff and the same relative separation/time writer.
Define its signed localization residual against the full periodic moving
contact exactly as in \eqref{eq:v2-smooth-localization-residual}.
A material capture certificate consists of
$\mathcal P_{k,\rm mov}\ge p_0>0$ and vanishing of that residual.
\end{definition}

\begin{corollary}[Intrinsic local replenishment on a captured material thread]
\label{cor:v3-material-survivor}
On a material capture branch, if
$\int_{\mathcal Q}|V_k|^2=o(\rho_k)$, then a fixed amount of at least one
complete near-pressure boundary flux survives in
\eqref{eq:v3-material-absorption}.  Its convective part uses fluctuations
only, its pressure contribution has no affine jet, and its pressure source
is confined, up to $o(1)$, to a shrinking physical neighbourhood of the
same terminal point.  More generally, on a subsequence with
$\mathcal P_{k,\rm mov}\ge p_0$ and vanishing material-capture residual,
vanishing of both near-pressure flux terms forces
$\int_{\mathcal Q}|V_k|^2\ge c\rho_k$ for all sufficiently large $k$.
Without the capture certificate no positive lower bound for the moving
contact is asserted.
\end{corollary}

\begin{proof}
V2's periodic-remainder estimate is uniform in the spatial centre, since
its proof bounds the remainder by the original global cubic integral.
It therefore also applies to the moving smooth cutoff.  Combine material
capture with \Cref{thm:v3-material-absorption} and the moving version of
\Cref{thm:v3-shrinking-pressure}.  The same rearrangement used in
\Cref{thm:v3-local-replenishment} proves the alternatives.
\end{proof}

\begin{firewall}[Physical travel does not prove payment capture]
The bound \eqref{eq:v3-physical-travel} allows travel through an unbounded
number of $\ell_k$-cells.  It does not imply that a heavy fixed Eulerian
cell becomes a heavy material tube, or that every microdelocalized law has
one captured trajectory.  Such a conclusion would require control of the
actual coherent measure under this transport.  Failure is retained as a
material-capture/localization discrepancy; no source price is assigned to
$\beta_k'$, to the uniform sweep, or to the number of crossed cells.
\end{firewall}

% =============================================================================
\section{Proof-status, fixed-datum export, and the next estimate}
\label{sec:v3-endpoint}
% =============================================================================

\begin{theorem}[Version V3 terminal replenishment reduction]
\label{thm:v3-integrated}
On every fixed-datum V1--V2 thread with the displayed fixed relative
writers, the following new assertions are proved.
The energy absorption can be centred by the actual local mean without
estimating its time derivative.  Both resulting pressure tests annihilate
affine spatial pressure, so their far component begins with the actual
trace-free Hessian \eqref{eq:v3-actual-card}.  Its Euclidean angular
operator is exactly \eqref{eq:v3-tidal-spectrum}, and the two tests have
the common-card pairings \eqref{eq:v3-common-tidal-pairing}.

The fixed global energy bound alone supplies a radius $d_k\to0$ for
which both pressure tails outside $B_{2d_k}(x_k)$ vanish.  Therefore the
original contact is controlled by the two signed near-pressure fluxes,
the mean-centred kinetic stock divided by $\rho_k$, and the original
localization/replay residual, as in
\Cref{thm:v3-local-replenishment}.  On a separately certified material
capture branch, uniform sweeping also disappears and the entire tube
remains in a shrinking physical neighbourhood of the same terminal point.
No conclusion asserts that the remaining near-pressure fluxes vanish,
that material capture always occurs, or that the terminal Dini criterion
has been proved.
\end{theorem}

\begin{proof}
The first paragraph is
\Cref{thm:v3-relative-energy,thm:v3-affine-null,thm:v3-relative-absorption,thm:v3-tidal-symbol,cor:v3-actual-tidal-card,thm:v3-common-curvature-pairing}.
The pressure-radius statement is \Cref{thm:v3-shrinking-pressure}; its
insertion into the same selected scalar is
\Cref{thm:v3-local-replenishment}.  The transported conclusions are
\Cref{thm:v3-material-absorption,thm:v3-travel,cor:v3-material-survivor}.
Every remaining qualification is a hypothesis explicitly retained in those
statements.
\end{proof}

\begingroup\small
\begin{longtable}{p{0.26\textwidth}p{0.15\textwidth}p{0.51\textwidth}}
\caption{NS--A Version V3 proof-status ledger.}\label{tab:v3-status}\\
\toprule Row & Status & Exact conclusion\\\midrule
\endfirsthead
\toprule Row & Status & Exact conclusion\\\midrule
\endhead
V1--V2 source & Preserved & Complete preterminal byte string retained; earlier proofs not rewritten.\\[0.25em]
Weighted-mean relative energy & Proved & The time derivative of the mean cancels exactly in the physical energy test.\\[0.25em]
Affine pressure & Proved test null & Both pressure functionals annihilate the affine jet; this is not a new pressure gauge.\\[0.25em]
Mean-centred viscous absorption & Proved & Same selected contact; no independent dissipation or mean-acceleration upper bound.\\[0.25em]
Actual far-pressure card & Proved & One trace-free Hessian from the actual exterior kinetic tensor and the periodic Green kernel.\\[0.25em]
Angular coefficients & Evaluated & $36,-24,-24,6,6$ on the five trace-free directions, with multiplier $(4\pi|z|^5)^{-1}$.\\[0.25em]
Two common-card responses & Proved & Explicit two-point and relative-energy moment tensors, with an $R^{-6}$ higher-jet error.\\[0.25em]
Shrinking pressure-source radius & Proved & Energy-only radius $\ell_k(1+\mathsf E_k)^{3/10}\log(e+\ell_k^{-1})$; both far fluxes vanish.\\[0.25em]
Improved radius & Conditional & Exponent $1/5$ replaces $3/10$ under a bounded normalized local cubic packet.\\[0.25em]
Near-source scalar test & Proved & Exact bound with two complete signed near-pressure fluxes, mean-free kinetic stock, and stated residual.\\[0.25em]
Mean-transported tube & Proved identity & Uniform sweeping and uniform acceleration disappear from the appropriate relative balances.\\[0.25em]
Physical tube travel & Proved & At most $C\nu^{-1}E_*^{1/2}\sqrt{\ell_k}$ on a fixed relative-time window.\\[0.25em]
Capture of original payment & Open input & Small physical travel does not guarantee a heavy transported tube.\\[0.25em]
Mesoscopic replenishment & Open & Pressure sources between $\ell_k$ and the shrinking outer radius may retain the signed flux.\\[0.25em]
Terminal Dini improvement & Open & No decay rate for the actual terminal source modulus is inferred from pressure locality alone.\\[0.25em]
Full microdelocalization and distance separation & Open & Domain diameter and relative distances remain explicit; no uniform packing theorem supplied.\\[0.25em]
NS--B paid stress and Reynolds budget & Not re-proved & The new curvature card is not that paid-stress path or its variation estimate.\\[0.25em]
Global regularity & Not claimed & The remaining same-history local flux or concentration must still be controlled.\\
\bottomrule
\end{longtable}
\endgroup

\paragraph{Fixed-datum interface to NS--B.}
The export is now more local and more concrete: the original finite-$k$
source history, its selected centre and time window, the tensor source in
$B_{2d_k}(x_k)$, the signed near-source fluxes, and---when a separated
pressure portion is retained---the common Hessian and the two moment
matrices in \eqref{eq:v3-common-tidal-pairing}.  A rank-three or Reynolds
classification still requires the inherited carrier and actual covariance
rows.  No rank-three claim follows merely from the three-dimensional
spatial kernel, the dimension five of its Hessian, or nonzero pressure.
The donor programme should test this same finite-$k$ survivor after its
represented and dynamically closed rank-at-most-two components have been
removed.  It should not replace it by an unrelated inviscid limit.

\begin{programme}[Version V4: the remaining near-source kinetic--pressure cancellation]
\label{prog:v3-next}
Preserve Versions V1--V3 verbatim.  On one surviving source-aligned thread,
evaluate the two complete near-pressure fluxes of
\eqref{eq:v3-near-two-flux}--\eqref{eq:v3-near-energy-flux}; do not reopen
macroscopic harmonic pressure, affine pressure, or mean acceleration as
independent costs.  The first unresolved estimate is their joint
same-scalar upper bound by the actual terminal source/dissipation history,
with the local kinetic term retained.  Distinguish a bounded cell-ratio
near source from a mesoscopic source whose ratio to $\ell_k$ diverges
inside $d_k$.  The pressure-curvature card gives its signed angular
coefficients but no lower positivity margin.

A transported calculation may be used only after the coherent payment is
shown to be captured by the actual tube.  On such a branch the intrinsic
flux uses $V_k$ rather than a uniform sweeping velocity.  If the
near-pressure fluxes and the declared residual vanish, invoke the existing
same-scalar contradiction when the mean-free kinetic term is small;
otherwise retain its quantitative concentration.  For diffuse cohorts,
first assemble the signed currents and then retain the actual diameter
and relative-source-scale bounds.  The terminal Dini theorem is invoked
only when its displayed source-modulus hypothesis is genuinely obtained.
\end{programme}

\paragraph{Literature boundary.}
The distinction between inviscid energy dissipation and the compactness
needed for an Euler landing is part of the classical literature; for
example Drivas--Eyink \cite{v3-DE} give viscosity-uniform regularity
conditions and dissipation estimates on a fixed torus.  Those hypotheses
are not supplied by this programme's expanding rescaled domain.
Huang's recent preprint \cite{v3-Huang} studies the pressure concentration
forced by endpoint energy atoms.  It is relevant to V2's kinetic-atom
branch, but no theorem of that preprint is used to assert non-atomicity
or a vanishing pressure current here.  The mean-centred identities,
coefficients, and shrinking-radius bounds of this version have complete
proofs above; these references provide outside context only.

% =============================================================================
\section*{Version V3 verification and history}
\addcontentsline{toc}{section}{Version V3 verification and history}
% =============================================================================

The operator certificate can be checked exactly on the orthogonal basis
\[
 \operatorname{diag}(-1,-1,2),\quad E_{13}+E_{31},\quad
 E_{23}+E_{32},\quad \operatorname{diag}(1,-1,0),\quad E_{12}+E_{21}.
\]
At $n=e_3$, direct fourth differentiation of $|z|^{-1}$ gives,
respectively, $36,-24,-24,6,6$ times those matrices; the identity matrix
is mapped to zero.  Rotational covariance transports this calculation to
every unit vector.  This finite check supplements the explicit projector
proof of \Cref{thm:v3-tidal-symbol}; no numerical eigenvalue or positivity
test substitutes for that proof.

\paragraph{Append-only integrity.}
The full V2 preterminal byte string is retained without alteration.
Its SHA--256 digest is
\begin{center}\scriptsize\ttfamily
ea4c6f870524a30480a085fb84450e947bda717f6f5ba338381e9da22966d48a
\end{center}
Only its final active document terminator was removed, and a single active
terminator is restored after this continuation.  The original title and
historical proof-status statements remain part of the preserved prefix.

\paragraph{Version V3 history.}
Derived the local mean-centred energy balance and the common affine-pressure
annihilator.  Evaluated the actual far-pressure Hessian, its signed angular
spectrum, its two local response moments, and its higher-jet remainder.
Proved an energy-only shrinking pressure-source radius and its sharper
bounded-cubic counterpart.  Inserted the resulting near-pressure fluxes
into the original selected scalar, retaining the localization discrepancy
and mean-free concentration alternative.  Established the transported-mean
balances and shrinking physical travel while leaving material payment
capture explicit.  Local kinetic--pressure cancellation, the terminal
Dini estimate, and full delocalized closure remain unresolved.

\begin{thebibliography}{9}
\bibitem{v3-DE}
T.~D.~Drivas and G.~L.~Eyink,
\emph{An Onsager singularity theorem for Leray solutions of incompressible
Navier--Stokes}, Nonlinearity \textbf{32} (2019), 4465--4482.
\href{https://doi.org/10.1088/1361-6544/ab2f42}{doi:10.1088/1361-6544/ab2f42}.
\url{https://arxiv.org/abs/1710.05205}.

\bibitem{v3-Huang}
H.~Huang,
\emph{Endpoint Energy Atoms Force Local Pressure Concentration in
Three-Dimensional Navier--Stokes Flow}, preprint, 31 August 2026.
\url{https://arxiv.org/abs/2608.30715v1}.
\end{thebibliography}

For Version V4, preserve this complete source verbatim, remove only its
final active document terminator, append the new proofs immediately before
it, and restore the terminator.  The next filename is
\begin{center}\scriptsize
\path{Renewal_NS_Terminal_Source_Concentration_and_Singular_Thread_Closure_Programme_V4.tex}.
\end{center}

\addtocontents{toc}{\protect\endgroup}
% =============================================================================
% END VERSION V3 MATERIAL
% =============================================================================


% =============================================================================
% BEGIN VERSION V4 MATERIAL -- append-only continuation of NS--A V1--V3
% =============================================================================
\clearpage
\hypersetup{pdftitle={NS-A Terminal Source Concentration Programme: Cumulative V4},pdfauthor={A. Pelissier}}
\makeatletter
\addtocontents{toc}{\protect\begingroup\protect\makeatletter}
\addtocontents{toc}{\protect\def\protect\l@section{\protect\@dottedtocline{1}{0em}{3.5em}}}
\addtocontents{toc}{\protect\makeatother}
\makeatother
\renewcommand{\thesection}{V4.\arabic{section}}
\renewcommand{\thesubsection}{V4.\arabic{section}.\arabic{subsection}}
\renewcommand{\theequation}{V4.\arabic{equation}}
\setcounter{section}{0}
\setcounter{equation}{0}
\renewcommand{\theHsection}{V4.\arabic{section}}
\renewcommand{\theHsubsection}{V4.\arabic{section}.\arabic{subsection}}
\renewcommand{\theHtheorem}{V4.\arabic{section}.\arabic{theorem}}
\renewcommand{\theHequation}{V4.\arabic{equation}}

\begin{center}
{\LARGE\bfseries NS--A: Version V4}\\[0.5em]
{\large Forward Record-to-Record Replenishment,\\
the Canonical Energy Cutoff, and the No-Silent-Abort Alternative}\\[0.5em]
{\normalsize 2 September 2026}
\end{center}
\addcontentsline{toc}{part}{Version V4: forward record-to-record replenishment and singular stress landing}

\begin{abstract}
The complete V1--V3 source is retained.  Two auxiliary terminal-concentration
calculations supplied after V3 change the order of the next attack.  They show
that a time-marginal improvement of the coherent source law is not the missing
local lemma: logarithmically spread ridge packets can have
$o(\sqrt\tau)$ marginal constants and still retain a fixed caloric payment.
They also isolate an ``aborted virtual collapse'' at a regular record time
$\beta_n$ and identify the unused datum: after $\beta_n$ the critical norm
must reach the next doubled record.  V4 therefore moves forward in time from
one record to the next instead of seeking another static upper bound for the
V3 near-pressure flux.

Let $\mathcal R_n$ be a doubling record for
$\|A^{1/4}u\|_2$, let $\beta_n$ be its first hitting time, and let
$E_*=\sup_{t<T}\|u(t)\|_2^2$.  The canonical energy cutoff
\[
 K_n=\mathcal R_n^2/E_*
\]
has the following exact property.  At the next record, the critical state has
high-frequency norm at least $\sqrt3\,\mathcal R_n$ above $K_n$, whereas the
entire old state has norm only $\mathcal R_n$.  Duhamel's formula therefore
forces an interval-fresh nonlinear-source response of norm at least
$(\sqrt3-1)\mathcal R_n$ above $K_n$.  More generally, at every prescribed
frequency front, either the fresh response below the front or the fresh
response above it has norm at least $(\sqrt2-1)\mathcal R_n$.  A lower-frequency
renewal spends at least $c\mathcal R_n/\Lambda_n$ of the native
$L_t^1H_x^{-1/2}$ source length.  A higher-frequency renewal has an explicit
terminal age core obtained by inverting the exact high-pass heat envelope.
For the canonical cutoff the corresponding physical heat lengths are
summable.

The terminal core has an exact pressure-free quadratic realization.  Its
Riesz direction produces a backward caloric divergence-free test field
$\phi_n$, and the complete response is
\[
 \int (u\otimes u)^\circ:S(\phi_n).
\]
Thus every doubled record yields a literal quadratic stress--strain card,
not merely an abstract norm increment.  The coherent positive replay of this
assembled scalar has vanishing mass on every compact terminally regular
region.  Its spatial centre laws are consequently supported on the terminal
singular set and admit, at the summable canonical response scale, the exact
trichotomy: singular-distance/scale separation, one heavy coherent stress
replay, or near-singular stress microdelocalization.  The heavy replay is then
shorted against the actual signed stress on the same ball.  It yields either a
normalized counterflow witness or a signed local stress thread; only the latter
is a correctly normalized fixed-datum card eligible for the inherited
rank-three donor and paid-stress audit.  The diffuse branch and the ancestry of
the interval-fresh response remain explicit.

V4 does not rule out an infinite chain of scale resets or fresh terminal
responses.  It proves that an aborted cascade cannot disappear silently: its
successor is a lower-frequency renewal, a same-or-finer renewal, a spatial
relocation/delocalization, or an explicit source-incidence/ancestry defect.
No terminal Dini improvement or global three-dimensional regularity theorem is
claimed.
\end{abstract}

% =============================================================================
\section{Auxiliary-note audit and the upstream change of direction}
\label{sec:v4-audit}
% =============================================================================

The two files
\begin{center}\small
\path{terminal_concentration_report(1).tex},\qquad
\path{terminal_concentration_report(4).tex}
\end{center}
are byte-for-byte identical; their common SHA--256 digest is
\begin{center}\scriptsize\ttfamily
bc2374ee79ba3e390387536047c092e76fb959c485ba085340401428ec306130
\end{center}
The companion first-pass file
\begin{center}\small
\path{NS_Terminal_Programme_V1_first_pass.tex}
\end{center}
has digest
\begin{center}\scriptsize\ttfamily
8236c202ee03616cb3e679990aa6bae3f430957d0ec5e662d9c7c3cf0c303e1a.
\end{center}
These hashes only freeze the reviewed inputs.

The auxiliary report gives three conclusions relevant after V3.  First, the
caloric payment is a net parabolic-band energy deposition and has an exact
ridge-weighted source representation.  Second, a logarithmically spread
ridge measure can keep a fixed payment even when its unweighted time marginal
obeys an arbitrarily small $o(\sqrt\tau)$ coefficient.  Third, an unconditional
proof of the terminal Dini condition on every hypothetical terminal chronology
would itself be a global regularity theorem.  The first-pass calculation then
sharpens the surviving spread object to a virtual Type-I collapse aimed at a
regular record time $\beta_n$.  Every estimate used there stops at $\beta_n$;
the calculation explicitly recommends using the next record
$\beta_{n+1}$ or a genuinely backward rigidity theorem.

\begin{assessment}[What is imported and what is new in V4]
\label{ass:v4-import}
V4 does not repeat the auxiliary flux identities, Abel--$L^4$ record law,
log-age trichotomy, log-spread counterexample, or virtual Type-I estimate.
Their role is diagnostic: another estimate using only the time before one
record would repeat the identified loop.  The new acquisition is a theorem
about the interval after that record.  It uses only the exact critical-state
equation, the global kinetic-energy cap, spectral orthogonality, and the
already inherited finite native source length.
\end{assessment}

\begin{firewall}[No return to a marginal-Dini surrogate]
\label{fire:v4-no-marginal-loop}
The source-age estimates below identify a terminal core of an actual
record-to-record vector response.  They are not used to assert that a bound
$M_n(\tau)=o(\sqrt\tau)$ implies the V35 Dini condition.  Logarithmic spreading
can make the Dini integral nonzero while every individual marginal coefficient
is small.  Any later closure must retain the joint age--frequency law, the
quadratic stress card, or the inter-record ancestry; it may not replace them by
a marginal rate already shown insufficient.
\end{firewall}

% =============================================================================
\section{Doubling records and the canonical energy cutoff}
\label{sec:v4-records}
% =============================================================================

Retain the fixed smooth unforced periodic solution of \eqref{eq:NS}, and put
\begin{equation}
 e(t)=A^{1/4}u(t),\qquad c(t)=A^{-1/4}\BNS(u)(t).
 \label{eq:v4-ec}
\end{equation}
Thus
\begin{equation}
 \partial_te+\nu Ae=-\Lambda c,
 \qquad \Lambda=A^{1/2}.
 \label{eq:v4-critical-equation}
\end{equation}
The inherited native estimate gives
\begin{equation}
 \mathfrak M_c:=\int_0^T\|c(t)\|_2\,\dd t<\infty.
 \label{eq:v4-native-source-total}
\end{equation}
The energy cap $E_*$ is the quantity already fixed in
\eqref{eq:v3-global-energy-cap}.  If $E_*=0$, the solution is zero and no
terminal chronology exists, so assume $E_*>0$.

Choose $\mathcal R_*>\|e(0)\|_2$ and put
\begin{equation}
 \mathcal R_n=2^n\mathcal R_*,
 \qquad
 \beta_n=\inf\{t<T:\|e(t)\|_2=\mathcal R_n\},
 \qquad
 J_n^+=(\beta_n,\beta_{n+1}].
 \label{eq:v4-doubling-records}
\end{equation}
On the hypothetical finite maximal-time branch these times are defined for all
large $n$, increase to $T$, and the intervals $J_n^+$ are disjoint.  Write
\begin{equation}
 \Delta_n=\beta_{n+1}-\beta_n,
 \qquad
 m_n=\int_{J_n^+}\|c(t)\|_2\,\dd t.
 \label{eq:v4-gap-source-mass}
\end{equation}
Then
\begin{equation}
 \sum_nm_n\le\mathfrak M_c,
 \qquad m_n\longrightarrow0,
 \qquad \Delta_n\longrightarrow0.
 \label{eq:v4-source-gap-summability}
\end{equation}
For $K>0$ let
\begin{equation}
 \mathsf P_{\le K}=\one_{[0,K]}(\Lambda),
 \qquad
 \mathsf P_{>K}=\one_{(K,\infty)}(\Lambda).
 \label{eq:v4-spectral-projections}
\end{equation}

\begin{lemma}[The kinetic-energy head bound]
\label{lem:v4-energy-head}
For every $t<T$ and every $K>0$,
\begin{equation}
 \boxed{\|\mathsf P_{\le K}e(t)\|_2^2\le K E_*.}
 \label{eq:v4-energy-head}
\end{equation}
Consequently, at the canonical cutoff
\begin{equation}
 \boxed{K_n:=\frac{\mathcal R_n^2}{E_*},
 \qquad \ell_n^{\rm out}:=K_n^{-1}=\frac{E_*}{\mathcal R_n^2},}
 \label{eq:v4-canonical-cutoff}
\end{equation}
one has
\begin{equation}
 \|\mathsf P_{>K_n}e(\beta_{n+1})\|_2
 \ge\sqrt3\,\mathcal R_n.
 \label{eq:v4-next-high-tail}
\end{equation}
\end{lemma}

\begin{proof}
In the spectral decomposition of $\Lambda$,
\[
 \|\mathsf P_{\le K}e(t)\|_2^2
 =\sum_{\kappa\le K}\kappa\|E_\kappa u(t)\|_2^2
 \le K\|u(t)\|_2^2\le K E_*.
\]
At $t=\beta_{n+1}$ the total squared critical norm is
$4\mathcal R_n^2$.  For $K=K_n$, the low part has squared norm at most
$\mathcal R_n^2$.  Orthogonality gives
\eqref{eq:v4-next-high-tail}.
\end{proof}

Define the interval-fresh high-frequency response
\begin{equation}
 \boxed{
 F_n^{\rm can}:=-\int_{\beta_n}^{\beta_{n+1}}
 \mathsf P_{>K_n}\Lambda
 e^{-\nu(\beta_{n+1}-t)A}c(t)\,\dd t.}
 \label{eq:v4-canonical-response}
\end{equation}
The adjective ``interval-fresh'' means only that this Duhamel term is
integrated over $J_n^+$.  It does not assert a new source direction after the
ancestry short.

\begin{theorem}[Canonical record-to-record high-frequency replenishment]
\label{thm:v4-canonical-replenishment}
For every sufficiently late doubled record,
\begin{equation}
 \boxed{\|F_n^{\rm can}\|_2
 \ge\gamma_*\mathcal R_n,
 \qquad \gamma_*:=\sqrt3-1.}
 \label{eq:v4-canonical-response-lower}
\end{equation}
Thus the next record cannot be obtained solely by transporting the state
present at $\beta_n$: above the energy-determined cutoff $K_n$, a fixed
fraction of the preceding record height is generated by the nonlinear source
on the interval after $\beta_n$.
\end{theorem}

\begin{proof}
Duhamel's formula for \eqref{eq:v4-critical-equation}, followed by the spectral
projection, is
\begin{equation}
 \mathsf P_{>K_n}e(\beta_{n+1})
 =e^{-\nu\Delta_nA}\mathsf P_{>K_n}e(\beta_n)
 +F_n^{\rm can}.
 \label{eq:v4-canonical-Duhamel}
\end{equation}
The semigroup is contractive and
$\|e(\beta_n)\|_2=\mathcal R_n$, so the old term has norm at most
$\mathcal R_n$.  The left side has norm at least
$\sqrt3\mathcal R_n$ by \Cref{lem:v4-energy-head}.  The reverse triangle
inequality gives \eqref{eq:v4-canonical-response-lower}.
\end{proof}

\begin{corollary}[Divergent response leverage on a finite native source stock]
\label{cor:v4-response-leverage}
Let
\begin{equation}
 H_n:=\|F_n^{\rm can}\|_2,
 \qquad L_n:=\frac{H_n}{m_n}.
 \label{eq:v4-HL}
\end{equation}
Then $m_n>0$ and
\begin{equation}
 \boxed{L_n\ge\gamma_*\frac{\mathcal R_n}{m_n}
 \longrightarrow\infty.}
 \label{eq:v4-leverage-divergence}
\end{equation}
In particular, every doubled record is followed by a response whose gain per
unit native source length diverges.
\end{corollary}

\begin{proof}
A zero $m_n$ would make \eqref{eq:v4-canonical-response} vanish.  The lower
bound is \eqref{eq:v4-canonical-response-lower}.  Since
$\mathcal R_n\to\infty$ and $m_n\to0$, the quotient diverges.
\end{proof}

% =============================================================================
\section{Every prescribed front has a lower- or higher-frequency renewal}
\label{sec:v4-no-silent-abort}
% =============================================================================

The canonical cutoff is independent of any previously selected virtual
cascade front.  The next theorem treats an arbitrary front and supplies the
precise forward-in-time successor requested by the auxiliary first pass.
For a cutoff $\Lambda_n>0$, define
\begin{align}
 F_{n,\le}^{\Lambda}
 &:=-\int_{J_n^+}\mathsf P_{\le\Lambda_n}\Lambda
 e^{-\nu(\beta_{n+1}-t)A}c(t)\,\dd t,
 \label{eq:v4-low-response}\\
 F_{n,>}^{\Lambda}
 &:=-\int_{J_n^+}\mathsf P_{>\Lambda_n}\Lambda
 e^{-\nu(\beta_{n+1}-t)A}c(t)\,\dd t.
 \label{eq:v4-high-response}
\end{align}

\begin{theorem}[No-silent-abort frequency alternative]
\label{thm:v4-no-silent-abort}
Put $\gamma_0=\sqrt2-1$.  For every cutoff sequence $\Lambda_n>0$ and every
late $n$, at least one of the following holds:
\begin{enumerate}[label=\textup{(A\arabic*)},leftmargin=3.6em]
\item \emph{Lower-frequency renewal:}
      \begin{equation}
       \boxed{\|F_{n,\le}^{\Lambda}\|_2
       \ge\gamma_0\mathcal R_n,
       \qquad
       m_n\ge\gamma_0\frac{\mathcal R_n}{\Lambda_n}.}
       \label{eq:v4-low-renewal}
      \end{equation}
\item \emph{Same-or-finer renewal:}
      \begin{equation}
       \boxed{\|F_{n,>}^{\Lambda}\|_2
       \ge\gamma_0\mathcal R_n.}
       \label{eq:v4-high-renewal}
      \end{equation}
      More precisely, if
      $\|\mathsf P_{>\Lambda_n}e(\beta_{n+1})\|_2\ge\sqrt2\mathcal R_n$,
      then
      \begin{equation}
       \|F_{n,>}^{\Lambda}\|_2
       \ge
       \left(\sqrt2-e^{-\nu\Lambda_n^2\Delta_n}\right)\mathcal R_n.
       \label{eq:v4-high-renewal-damped}
      \end{equation}
\end{enumerate}
Both alternatives may occur.  On the lower branch,
\begin{equation}
 \boxed{\sum_{n\in\mathcal L}
 \frac{\mathcal R_n}{\Lambda_n}
 \le\gamma_0^{-1}\mathfrak M_c}
 \label{eq:v4-low-renewal-sum}
\end{equation}
for every set $\mathcal L$ of indices on which \textup{(A1)} is selected.
\end{theorem}

\begin{proof}
At $\beta_{n+1}$ the low and high components are orthogonal and their squared
norms sum to $4\mathcal R_n^2$.  Hence at least one component has norm at least
$\sqrt2\mathcal R_n$.  Apply Duhamel's formula to that component.  Its old
part is the contraction of the corresponding component of $e(\beta_n)$ and
therefore has norm at most $\mathcal R_n$.  The fresh part consequently has
norm at least $(\sqrt2-1)\mathcal R_n$.

For the low component,
\[
 \|F_{n,\le}^{\Lambda}\|_2
 \le\int_{J_n^+}
 \|\mathsf P_{\le\Lambda_n}\Lambda
 e^{-\nu(\beta_{n+1}-t)A}\|_{2\to2}\|c(t)\|_2\,\dd t
 \le\Lambda_nm_n,
\]
which proves the source-length debit in \eqref{eq:v4-low-renewal}.  Summation
uses the disjointness of $J_n^+$ and \eqref{eq:v4-native-source-total}.

For the high component, spectral support gives
\[
 \|e^{-\nu\Delta_nA}\mathsf P_{>\Lambda_n}e(\beta_n)\|_2
 \le e^{-\nu\Lambda_n^2\Delta_n}\mathcal R_n.
\]
The reverse triangle inequality proves
\eqref{eq:v4-high-renewal-damped}, and its weaker form is
\eqref{eq:v4-high-renewal}.
\end{proof}

\begin{corollary}[Successor of a virtual cascade front]
\label{cor:v4-virtual-front}
Suppose an auxiliary log-age or band argument selects a virtual front
$\Lambda_n^{\rm vir}$ at the regular record time $\beta_n$.  At the next
record, the front has one of two exact successors:
\begin{equation}
 \boxed{
 \begin{gathered}
 \text{an interval-fresh response of size }\gamma_0\mathcal R_n
 \text{ below }\Lambda_n^{\rm vir},\\
 \text{or an interval-fresh response of size }\gamma_0\mathcal R_n
 \text{ above }\Lambda_n^{\rm vir}.
 \end{gathered}}
 \label{eq:v4-virtual-successor}
\end{equation}
The first is a lower-frequency reorganization and pays the debit
\eqref{eq:v4-low-renewal}; the second is a same-or-finer replenishment.  Thus
an aborted virtual collapse can remain regular at $\beta_n$, but it cannot
simply disappear from the chronology without one of these two source
responses.
\end{corollary}

\begin{proof}
Apply \Cref{thm:v4-no-silent-abort} with
$\Lambda_n=\Lambda_n^{\rm vir}$.  No estimate from the time before $\beta_n$
is used.
\end{proof}

\begin{corollary}[The energy cutoff forces the higher-frequency branch]
\label{cor:v4-canonical-no-low}
If $\Lambda_n\le K_n=\mathcal R_n^2/E_*$, then
\begin{equation}
 \|\mathsf P_{\le\Lambda_n}e(\beta_{n+1})\|_2\le\mathcal R_n.
 \label{eq:v4-low-head-canonical}
\end{equation}
Hence the endpoint high component is at least $\sqrt3\mathcal R_n$, and
\begin{equation}
 \boxed{\|F_{n,>}^{\Lambda}\|_2
 \ge\left(\sqrt3-e^{-\nu\Lambda_n^2\Delta_n}\right)\mathcal R_n
 \ge\gamma_*\mathcal R_n.}
 \label{eq:v4-subcanonical-high-renewal}
\end{equation}
In particular, every virtual front at or below the canonical energy frequency
must be continued by a same-or-finer interval-fresh response.
\end{corollary}

\begin{proof}
Use \Cref{lem:v4-energy-head} and orthogonality at the next record, then repeat
the high-component argument of \Cref{thm:v4-no-silent-abort}.
\end{proof}

\begin{firewall}[Interval freshness is not orthogonal source birth]
\label{fire:v4-fresh-not-birth}
The responses in \eqref{eq:v4-low-response}--\eqref{eq:v4-high-response}
are generated after $\beta_n$, but they may be deterministic descendants of
an earlier source, represented return, or continuing carrier.  V4 proves a
Duhamel time-incidence statement, not an orthogonal source-first-birth theorem.
The inherited ancestry and return shorts must be applied before the response
is assigned a new source price.
\end{firewall}

% =============================================================================
\section{Exact terminal age core of a higher-frequency renewal}
\label{sec:v4-age-core}
% =============================================================================

The higher-frequency alternative can exploit arbitrarily large leverage near
the right endpoint.  The exact envelope determines how close to that endpoint
a fixed vector response must be represented.

\begin{lemma}[Exact high-pass backward-heat envelope]
\label{lem:v4-highpass-envelope}
For $K>0$ and $s>0$, put
\begin{equation}
 G_K(s):=\sup_{\kappa\ge K}\kappa e^{-\nu s\kappa^2}.
 \label{eq:v4-GK-def}
\end{equation}
Then $G_K$ is decreasing and
\begin{equation}
 \boxed{
 G_K(s)=
 \begin{cases}
 (2e\nu s)^{-1/2},&0<s\le(2\nu K^2)^{-1},\\[0.25em]
 K e^{-\nu sK^2},&s\ge(2\nu K^2)^{-1}.
 \end{cases}}
 \label{eq:v4-GK-formula}
\end{equation}
For $L>0$, define
\begin{equation}
 \boxed{
 \sigma(K,L)=
 \begin{cases}
 \dfrac{2}{e\nu L^2},&L\ge 2K/\sqrt e,\\[0.8em]
 \dfrac{1}{\nu K^2}\log\dfrac{2K}{L},&0<L<2K/\sqrt e.
 \end{cases}}
 \label{eq:v4-sigma-KL}
\end{equation}
Then $G_K(\sigma(K,L))=L/2$; the two formulas agree at the
interface.
\end{lemma}

\begin{proof}
The derivative of $\kappa e^{-\nu s\kappa^2}$ vanishes at
$\kappa=(2\nu s)^{-1/2}$.  This point lies in $[K,\infty)$ exactly in the
first time range.  Otherwise the function decreases on $[K,\infty)$ and its
maximum is the endpoint value.  Both branches are decreasing and meet at
$K e^{-1/2}$.  Solving $G_K(\sigma)=L/2$ on the appropriate branch gives
\eqref{eq:v4-sigma-KL}; the branch condition is exactly the condition that the
solution lie on that side of $(2\nu K^2)^{-1}$.
\end{proof}

\begin{theorem}[Terminal vector core and Abel source record]
\label{thm:v4-terminal-vector-core}
Let $a<b<T$, $\Delta=b-a$, $K>0$, and
\begin{equation}
 F=-\int_a^b\mathsf P_{>K}\Lambda
 e^{-\nu(b-t)A}c(t)\,\dd t,
 \qquad H=\|F\|_2>0,
 \qquad m=\int_a^b\|c(t)\|_2\,\dd t.
 \label{eq:v4-general-high-response}
\end{equation}
Put $L=H/m$, $\widehat\sigma=\min\{\Delta,\sigma(K,L)\}$, and
\begin{equation}
 F^{\rm core}=-\int_{b-\widehat\sigma}^{b}
 \mathsf P_{>K}\Lambda e^{-\nu(b-t)A}c(t)\,\dd t.
 \label{eq:v4-core-response}
\end{equation}
Then
\begin{equation}
 \boxed{\|F^{\rm core}\|_2\ge\frac H2.}
 \label{eq:v4-core-half}
\end{equation}
Moreover, with $f(s)=\|c(b-s)\|_2$ and
$M(\tau)=\int_0^{\min\{\tau,\Delta\}}f(s)\,\dd s$,
\begin{equation}
 \boxed{
 \int_0^\Delta\frac{f(s)}{\sqrt s}\,\dd s
 =\Delta^{-1/2}M(\Delta)
 +\frac12\int_0^\Delta M(\tau)\tau^{-3/2}\,\dd\tau
 \ge\sqrt{2e\nu}\,H.}
 \label{eq:v4-Abel-source-record}
\end{equation}
Thus the core is a vector-valued terminal source packet, while the final
formula is the exact Abel action of its actual source-time marginal.  No
vanishing Dini conclusion is inferred.
\end{theorem}

\begin{proof}
If $\sigma(K,L)\ge\Delta$, then $F^{\rm core}=F$.  Otherwise, by
\Cref{lem:v4-highpass-envelope},
\[
 \|F-F^{\rm core}\|_2
 \le\int_a^{b-\sigma(K,L)}G_K(b-t)\|c(t)\|_2\,\dd t
 \le mG_K(\sigma(K,L))=\frac H2.
\]
The reverse triangle inequality proves \eqref{eq:v4-core-half}.

The universal envelope
$G_K(s)\le(2e\nu s)^{-1/2}$ gives
\[
 H\le\frac1{\sqrt{2e\nu}}
 \int_0^\Delta f(s)s^{-1/2}\,\dd s.
\]
Since $M$ is absolutely continuous and $M(\tau)=O(\tau)$ near zero for each
fixed smooth interval, Stieltjes integration by parts gives
\[
 \int_{(0,\Delta]}\tau^{-1/2}\,\dd M(\tau)
 =\Delta^{-1/2}M(\Delta)
 +\frac12\int_0^\Delta M(\tau)\tau^{-3/2}\,\dd\tau.
\]
The left side is the weighted source integral, proving
\eqref{eq:v4-Abel-source-record}.
\end{proof}

Apply the theorem to \eqref{eq:v4-canonical-response}; denote the resulting
quantities by
\begin{equation}
 \widehat\sigma_n,
 \qquad F_n^{\rm core},
 \qquad H_n^{\rm core}=\|F_n^{\rm core}\|_2,
 \qquad d_n^{\rm age}=\sqrt{\nu\widehat\sigma_n}.
 \label{eq:v4-canonical-core-data}
\end{equation}
Then
\begin{equation}
 H_n^{\rm core}\ge\frac{\gamma_*}{2}\mathcal R_n.
 \label{eq:v4-core-record-lower}
\end{equation}

\begin{theorem}[Summable physical heat lengths of the canonical future cores]
\label{thm:v4-summable-core-lengths}
The canonical terminal cores satisfy
\begin{equation}
 \boxed{\sum_n d_n^{\rm age}<\infty,
 \qquad
 \sum_n\ell_n^{\rm out}<\infty.}
 \label{eq:v4-two-scale-summability}
\end{equation}
More precisely, if $L_n=H_n/m_n$, then on the branch
$L_n\ge2K_n/\sqrt e$,
\begin{equation}
 d_n^{\rm age}
 \le\sqrt{\frac2e}\,L_n^{-1}
 \le C\frac{m_n}{\mathcal R_n},
 \qquad
 d_n^{\rm age}\le\frac{\ell_n^{\rm out}}{\sqrt2},
 \label{eq:v4-leverage-core-scale}
\end{equation}
while on the complementary branch
\begin{equation}
 d_n^{\rm age}
 \le\ell_n^{\rm out}
 \sqrt{C+\log\mathcal R_n}.
 \label{eq:v4-cutoff-core-scale}
\end{equation}
Consequently the combined response scale
\begin{equation}
 \boxed{\delta_n^{\rm rec}:=
 \ell_n^{\rm out}+d_n^{\rm age}}
 \label{eq:v4-response-scale}
\end{equation}
is summable and tends to zero.
\end{theorem}

\begin{proof}
The geometric series for
$\ell_n^{\rm out}=E_*/\mathcal R_n^2$ is summable.  On the first branch,
\eqref{eq:v4-sigma-KL} and \eqref{eq:v4-canonical-core-data} give
\[
 d_n^{\rm age}\le\sqrt{\frac2e}\,L_n^{-1}
 \le\sqrt{\frac2e}\,\frac{m_n}{\gamma_*\mathcal R_n}.
\]
The branch condition also places $\widehat\sigma_n$ below
$(2\nu K_n^2)^{-1}$, proving the second assertion in
\eqref{eq:v4-leverage-core-scale}.  Summation is bounded by a constant times
$\sum_nm_n/\mathcal R_n$, which is finite by
\eqref{eq:v4-source-gap-summability}.

On the second branch,
\[
 d_n^{\rm age}
 \le K_n^{-1}\sqrt{\log(2K_n/L_n)}.
\]
Using $H_n\ge\gamma_*\mathcal R_n$ and $m_n\le\mathfrak M_c$,
\[
 \frac{2K_n}{L_n}
 \le\frac{2m_n\mathcal R_n}{\gamma_*E_*}
 \le C\mathcal R_n.
\]
The logarithm is positive on this branch and is at most
$C+\log\mathcal R_n$.  Since $\mathcal R_n=2^n\mathcal R_*$,
\[
 \sum_n\frac{\sqrt{C+\log\mathcal R_n}}{\mathcal R_n^2}<\infty,
\]
which proves \eqref{eq:v4-cutoff-core-scale} and the theorem.
\end{proof}

\begin{corollary}[A forward Abel record growing with the doubled norm]
\label{cor:v4-forward-Abel-record}
For the canonical intervals,
\begin{equation}
 \boxed{
 \int_{\beta_n}^{\beta_{n+1}}
 \frac{\|c(t)\|_2}{\sqrt{\beta_{n+1}-t}}\,\dd t
 \ge\sqrt{2e\nu}\,\gamma_*\mathcal R_n.}
 \label{eq:v4-forward-Abel-lower}
\end{equation}
The lower bound diverges geometrically with the record.  It is compatible
with the finite unweighted stock \eqref{eq:v4-native-source-total} only through
increasing terminal leverage and does not exclude logarithmic spreading over
many age scales.
\end{corollary}

\begin{proof}
Use \eqref{eq:v4-Abel-source-record} with
$H=H_n\ge\gamma_*\mathcal R_n$.
\end{proof}

% =============================================================================
\section{The terminal core is an exact quadratic stress--strain card}
\label{sec:v4-stress-card}
% =============================================================================

The preceding results are vector statements in the critical state space.  The
next identity returns them to the literal quadratic Navier--Stokes occurrence.
Put
\begin{equation}
 \omega_n:=\frac{F_n^{\rm core}}{H_n^{\rm core}},
 \qquad \|\omega_n\|_2=1,
 \qquad \mathsf P_{>K_n}\omega_n=\omega_n,
 \label{eq:v4-core-direction}
\end{equation}
and, for
$t\in(\beta_{n+1}-\widehat\sigma_n,\beta_{n+1})$, define
\begin{align}
 q_n(t)
 &:=\mathsf P_{>K_n}\Lambda
 e^{-\nu(\beta_{n+1}-t)A}\omega_n,
 \label{eq:v4-core-q}\\
 \phi_n(t)
 &:=A^{-1/4}q_n(t)
 =A^{1/4}\mathsf P_{>K_n}
 e^{-\nu(\beta_{n+1}-t)A}\omega_n.
 \label{eq:v4-core-phi}
\end{align}
The field $\phi_n$ is divergence-free and satisfies
$-\partial_t\phi_n+\nu A\phi_n=0$ on the open interval.
Write
$S(\phi)=\frac12(\nabla\phi+\nabla\phi^{\mathsf T})$ and
$M^\circ=M-\frac13(\operatorname{tr}M)I$.

\begin{theorem}[Pressure-free quadratic realization of the future core]
\label{thm:v4-stress-realization}
The complete core response has the exact representations
\begin{equation}
 \boxed{
 \begin{aligned}
 H_n^{\rm core}
 &=-\int_{\beta_{n+1}-\widehat\sigma_n}^{\beta_{n+1}}
   \langle c(t),q_n(t)\rangle\,\dd t\\
 &=\int_{\beta_{n+1}-\widehat\sigma_n}^{\beta_{n+1}}
   \int_{\Torus^3}
   (u\otimes u)^\circ(t,x):S(\phi_n)(t,x)\,\dd x\,\dd t.
 \end{aligned}}
 \label{eq:v4-stress-realization}
\end{equation}
In particular,
\begin{equation}
 \frac1{\mathcal R_n}
 \int\!\!\int (u\otimes u)^\circ:S(\phi_n)
 \ge\frac{\gamma_*}{2}.
 \label{eq:v4-normalized-stress-margin}
\end{equation}
No pressure term occurs: it is removed by the divergence-free test before any
spatial positive replay is formed.
\end{theorem}

\begin{proof}
By the definition of $\omega_n$ and $F_n^{\rm core}$,
\[
 H_n^{\rm core}
 =\langle F_n^{\rm core},\omega_n\rangle
 =-\int\langle c(t),q_n(t)\rangle\,\dd t.
\]
Since $q_n=A^{1/4}\phi_n$ and
$c=A^{-1/4}\BNS(u)$, self-adjoint spectral calculus gives
$\langle c,q_n\rangle=\langle\BNS(u),\phi_n\rangle$.
The Leray projection may be dropped against the divergence-free field
$\phi_n$, and periodic integration by parts gives
\[
 -\langle\BNS(u),\phi_n\rangle
 =-\langle\operatorname{div}(u\otimes u),\phi_n\rangle
 =\int(u\otimes u):\nabla\phi_n.
\]
The stress is symmetric, so only $S(\phi_n)$ contributes.  Its trace is
$\operatorname{div}\phi_n=0$, so the isotropic part of $u\otimes u$ also
vanishes.  This proves \eqref{eq:v4-stress-realization}; the margin is
\eqref{eq:v4-core-record-lower}.
\end{proof}

\begin{definition}[Coherent positive replay of the assembled future stress card]
\label{def:v4-stress-law}
On the core spacetime slab let $\zeta_n$ be the finite signed measure with
density
\begin{equation}
 \dd\zeta_n(t,x)=
 (u\otimes u)^\circ(t,x):S(\phi_n)(t,x)\,\dd x\,\dd t.
 \label{eq:v4-zeta-density}
\end{equation}
Put
\begin{equation}
 Z_n^+=\zeta_{n,+}(\Omega_n),
 \qquad
 \dd\kappa_n^{\rm str}:=
 \frac{H_n^{\rm core}}{Z_n^+}\,\dd\zeta_{n,+},
 \qquad
 \dd\pi_n^{\rm str}:=
 \frac{\dd\kappa_n^{\rm str}}{H_n^{\rm core}}
 =\frac{\dd\zeta_{n,+}}{Z_n^+}.
 \label{eq:v4-stress-probability}
\end{equation}
Let $\rho_n^{\rm str}$ be the spatial marginal of
$\pi_n^{\rm str}$.  By \eqref{eq:v4-stress-realization},
\begin{equation}
 \boxed{Z_n^+\ge H_n^{\rm core}
 \ge\frac{\gamma_*}{2}\mathcal R_n,\qquad
 0\le\kappa_n^{\rm str}\le\zeta_{n,+}\le|\zeta_n|,\qquad
 \kappa_n^{\rm str}(\Omega_n)=H_n^{\rm core}.}
 \label{eq:v4-positive-stress-mass}
\end{equation}
The coherent replay is formed only after the complete signed stress scalar.
It reallocates that one resultant through its positive carrier; it does not
assert that a positive replay cell equals the signed local stress before the
counterflow test below passes.
\end{definition}

\begin{theorem}[Terminally regular regions cannot carry the future stress law]
\label{thm:v4-stress-singular-support}
Define the inherited terminal regular and singular sets by
\begin{equation}
 \begin{aligned}
 \mathcal R_T:=\bigl\{x_0\in\Torus^3:\ &\text{there is }r>0\text{ such that }u\\
 &\text{extends smoothly to }(T-r^2,T]\times B_{2r}(x_0)\bigr\},\\
 \Sigma_T:=\Torus^3\setminus\mathcal R_T.
 \end{aligned}
 \label{eq:v4-terminal-sets}
\end{equation}
Let $K\Subset\mathcal R_T$ be compact.  Then
\begin{equation}
 \boxed{\rho_n^{\rm str}(K)\longrightarrow0.}
 \label{eq:v4-regular-stress-vanishing}
\end{equation}
Consequently every weak-* cluster point of the spatial laws
$\rho_n^{\rm str}$ is supported on $\Sigma_T$.
\end{theorem}

\begin{proof}
Terminal regularity gives a compact neighbourhood $K'$ with $K\Subset K'\Subset\mathcal R_T$, a time $t_K<T$,
and a constant $C_K$ such that $|u(t,x)|\le C_K$ on
$(t_K,T]\times K'$.  For large $n$ the core interval lies in this time slab.
For $s=\beta_{n+1}-t>0$, spectral calculus gives
\begin{equation}
 \|S(\phi_n)(t)\|_2
 \le C\|A^{3/4}e^{-\nu sA}\omega_n\|_2
 \le C(\nu s)^{-3/4}.
 \label{eq:v4-test-strain-bound}
\end{equation}
Indeed,
$\sup_{\kappa>0}\kappa^{3/2}e^{-\nu s\kappa^2}
=C(\nu s)^{-3/4}$.
Therefore
\[
 |\zeta_n|\bigl((\beta_{n+1}-\widehat\sigma_n,\beta_{n+1})\times K\bigr)
 \le C_K\int_0^{\widehat\sigma_n}s^{-3/4}\,\dd s
 \le C_K\widehat\sigma_n^{1/4}.
\]
By \eqref{eq:v4-positive-stress-mass}, the coherent replay is dominated by
$|\zeta_n|$, and division by $H_n^{\rm core}\ge(\gamma_*/2)\mathcal R_n$
shows that $\rho_n^{\rm str}(K)\to0$.  If a subsequence converges weakly-*
to $\rho$, then every continuous cutoff supported compactly in
$\mathcal R_T$ has integral tending to zero and hence has zero $\rho$-integral.
Such cutoffs exhaust the open regular set, so $\supp\rho\subseteq\Sigma_T$.
\end{proof}

\begin{remark}[Why this support theorem is different from the V35 cubic one]
The V35 support theorem concerned the centre law of the original longitudinal
caloric contact.  The present theorem concerns a new record-to-record
backward-heat stress test selected by the future Duhamel response.  Its proof
uses the integrable $s^{-3/4}$ strain singularity and terminal smoothness on a
regular compact set.  It does not identify the two positive replay laws or
transfer a cell probability from one to the other.
\end{remark}

% =============================================================================
\section{Canonical spatial trichotomy for the future stress law}
\label{sec:v4-spatial-trichotomy}
% =============================================================================

The two physical scales of the future core are the inverse output cutoff
$\ell_n^{\rm out}$ and the terminal heat length $d_n^{\rm age}$.  Their sum
$\delta_n^{\rm rec}$ is the summable scale in
\eqref{eq:v4-response-scale}.  For $R<\infty$, put
\begin{equation}
 \mathcal N_n(R)=
 \{x\in\Torus^3:\operatorname{dist}(x,\Sigma_T)
       \le R\delta_n^{\rm rec}\},
 \qquad
 \eta_n(R)=\rho_n^{\rm str}(\mathcal N_n(R)).
 \label{eq:v4-near-singular-stress-mass}
\end{equation}

\begin{theorem}[Canonical record-response spatial trichotomy]
\label{thm:v4-spatial-trichotomy}
After passage to a subsequence, exactly one of the following occurs.
\begin{enumerate}[label=\textup{(S\arabic*)},leftmargin=3.5em]
\item \emph{Singular-distance/response-scale separation.}  For every fixed
      $R<\infty$,
      \begin{equation}
       \boxed{\eta_n(R)\longrightarrow0.}
       \label{eq:v4-distance-separation}
      \end{equation}
      The stress centres approach $\Sigma_T$ in absolute probability by
      \Cref{thm:v4-stress-singular-support}, but their distance in units of
      $\delta_n^{\rm rec}$ diverges.
\item \emph{One heavy singular stress thread.}  There are finite $R,A$, positive
      $\eta,p$, centres $x_n$, and points $y_n\in\Sigma_T$ such that
      \begin{equation}
       \boxed{
       \rho_n^{\rm str}
       \bigl(B_{A\delta_n^{\rm rec}}(x_n)\cap\mathcal N_n(R)\bigr)
       \ge p,
       \qquad
       d(x_n,y_n)\le(A+R)\delta_n^{\rm rec}.}
       \label{eq:v4-heavy-stress-thread}
      \end{equation}
      After a further subsequence, $x_n,y_n\to x_*\in\Sigma_T$.
\item \emph{Near-singular stress microdelocalization.}  There are finite $R$ and
      $\eta>0$ such that $\eta_n(R)\ge\eta$, while for every fixed
      $A<\infty$,
      \begin{equation}
       \boxed{
       \sup_{x\in\Torus^3}
       \rho_n^{\rm str}
       \bigl(B_{A\delta_n^{\rm rec}}(x)\cap\mathcal N_n(R)\bigr)
       \longrightarrow0.}
       \label{eq:v4-stress-microdelocalization}
      \end{equation}
\end{enumerate}
The trichotomy concerns one coherent replay of one assembled stress scalar; it
is not a source-count alternative.
\end{theorem}

\begin{proof}
Pass diagonally so that $\eta_n(R)$ has a limit for every positive integer
$R$.  If all limits vanish, monotonicity in $R$ gives \textup{(S1)}.  Absolute
approach to $\Sigma_T$ follows from
\Cref{thm:v4-stress-singular-support}; combining this with
$\delta_n^{\rm rec}\to0$ gives the stated ratio interpretation.

Otherwise choose an integer $R$ and $\eta>0$ for which
$\eta_n(R)\ge\eta$ eventually.  Restrict $\rho_n^{\rm str}$ to
$\mathcal N_n(R)$.  For every positive integer $A$, pass diagonally again so
that
\[
 c_n(A)=\sup_x\rho_n^{\rm str}
       (B_{A\delta_n^{\rm rec}}(x)\cap\mathcal N_n(R))
\]
has a limit.  If one such limit is positive, choose almost maximizing balls.
Their restricted positive mass gives the first inequality in
\eqref{eq:v4-heavy-stress-thread}.  Such a ball intersects
$\mathcal N_n(R)$, so its centre lies within $(A+R)\delta_n^{\rm rec}$ of
some $y_n\in\Sigma_T$.  Compactness of $\Sigma_T$ and the vanishing scale give
the common cluster point.  This is \textup{(S2)}.

If every $c_n(A)$ tends to zero, \textup{(S3)} holds.  The alternatives are
mutually exclusive after the selected subsequence and exhaust the two diagonal
choices.
\end{proof}

\begin{corollary}[Heavy replay gives a counterflow witness or a signed local stress thread]
\label{cor:v4-heavy-counterflow}
Assume branch \textup{(S2)} of \Cref{thm:v4-spatial-trichotomy}, and write
\begin{align}
 B_n&:=B_{A\delta_n^{\rm rec}}(x_n),
 \label{eq:v4-heavy-ball}\\
 \Delta_n^{\rm str}&:=
 \left|\kappa_n^{\rm str}(I_n^{\rm core}\times B_n)
       -\zeta_n(I_n^{\rm core}\times B_n)\right|,
 \qquad
 I_n^{\rm core}:=(\beta_{n+1}-\widehat\sigma_n,\beta_{n+1}).
 \label{eq:v4-local-stress-counterflow}
\end{align}
After passage to a subsequence, at least one of the following holds:
\begin{enumerate}[label=\textup{(C\arabic*)},leftmargin=3.5em]
\item \emph{Local stress counterflow or replay mismatch:}
      \begin{equation}
       \boxed{\limsup_{n\to\infty}
       \frac{\Delta_n^{\rm str}}{\mathcal R_n}>0.}
       \label{eq:v4-stress-counterflow-survival}
      \end{equation}
\item \emph{Signed local stress thread:}
      \begin{equation}
       \boxed{
       \zeta_n
       \bigl((\beta_{n+1}-\widehat\sigma_n,\beta_{n+1})\times B_n\bigr)
       \ge \frac{p\gamma_*}{4}\mathcal R_n}
       \label{eq:v4-signed-local-stress-margin}
      \end{equation}
      for every sufficiently large $n$.
\end{enumerate}
The first branch is retained as one signed counterflow/interface obstruction;
it is not charged in addition to the full stress resultant.
\end{corollary}

\begin{proof}
The heavy-ball conclusion and the definition of the coherent replay give
\[
 \kappa_n^{\rm str}
 \bigl((\beta_{n+1}-\widehat\sigma_n,\beta_{n+1})\times B_n\bigr)
 =H_n^{\rm core}\rho_n^{\rm str}(B_n)
 \ge pH_n^{\rm core}
 \ge\frac{p\gamma_*}{2}\mathcal R_n.
\]
If the first branch fails, pass to a subsequence on which
$\Delta_n^{\rm str}/\mathcal R_n\to0$.  It is eventually at most
$(p\gamma_*/4)\mathcal R_n$.  Subtracting this discrepancy from the preceding
coherent lower bound proves \eqref{eq:v4-signed-local-stress-margin}.
\end{proof}

\begin{corollary}[A correctly normalized fixed-datum export card]
\label{cor:v4-export-card}
On branch \textup{(S2)}, and on the signed-thread branch of
\Cref{cor:v4-heavy-counterflow}, the packet
\begin{equation}
 \boxed{
 \begin{aligned}
 \mathfrak C_n^{\rm rec}=\bigl(&J_n^+,K_n,\widehat\sigma_n,x_n,
 \delta_n^{\rm rec},B_n,\\
 &\omega_n,\phi_n,(u\otimes u)^\circ,
 \kappa_n^{\rm str},\Delta_n^{\rm str}\bigr).
 \end{aligned}}
 \label{eq:v4-export-card}
\end{equation}
has all of the following properties:
\begin{enumerate}[label=\textup{(E\arabic*)},leftmargin=3.5em]
\item its source time lies in a terminal core of physical heat length
      $d_n^{\rm age}$ and its output test lies above $K_n$;
\item both $d_n^{\rm age}$ and $K_n^{-1}$ are summable;
\item its normalized quadratic stress--strain resultant is bounded below by
      $\gamma_*/2$;
\item a fixed fraction of its coherent replay is carried in a ball of radius
      $A\delta_n^{\rm rec}$ whose distance from $\Sigma_T$ is
      $O(\delta_n^{\rm rec})$;
\item after the local counterflow short, the actual signed stress in that ball
      has the fixed margin \eqref{eq:v4-signed-local-stress-margin};
\item no compact terminally regular region carries a positive fraction.
\end{enumerate}
After the inherited represented-return, source-ancestry, and dynamically
closed rank-at-most-two audits pass, this is a literal fixed-datum packet on
which the NS--B Fourier donor and paid-stress calculation may be opened.
\end{corollary}

\begin{proof}
The first two properties are
\Cref{thm:v4-terminal-vector-core,thm:v4-summable-core-lengths}; the third is
\eqref{eq:v4-normalized-stress-margin}; the fourth is
\eqref{eq:v4-heavy-stress-thread}; the fifth is
\Cref{cor:v4-heavy-counterflow}; and the sixth is
\Cref{thm:v4-stress-singular-support}.  The final sentence is the declared
NS--A/NS--B interface: V4 supplies the physical packet but does not perform the
separate donor and ancestry audits.
\end{proof}

\begin{firewall}[The five stress coordinates are not yet the paid-stress path]
\label{fire:v4-not-paid-stress}
The pointwise trace-free stress in \eqref{eq:v4-stress-realization} takes values
in a five-dimensional matrix space.  This dimensional fact does not identify
it with the NS--B path $(\theta_{\rm c}\mathbb K[u])^\circ$, prove bounded
variation of that path, or establish a positive Reynolds-covariance square
function.  Those conclusions require the inherited selector, low-output,
interband, and ancestry reductions on the same packet.
\end{firewall}

% =============================================================================
\section{Integrated V4 theorem and exact remaining alternatives}
\label{sec:v4-integrated}
% =============================================================================

\begin{theorem}[Version V4 forward terminal-renewal theorem]
\label{thm:v4-integrated}
Preserving every V1--V3 theorem at its stated hypotheses, V4 proves the
following for any hypothetical finite maximal time.
\begin{enumerate}[label=\textup{(V4.\arabic*)},leftmargin=5.3em]
\item On every interval between consecutive doubled critical-norm records, the
      canonical high-frequency nonlinear-source response satisfies
      \eqref{eq:v4-canonical-response-lower} above the energy cutoff
      $K_n=\mathcal R_n^2/E_*$.
\item Relative to any predeclared virtual front, the next record has a fresh
      lower-frequency or fresh same-or-finer response of size at least
      $(\sqrt2-1)\mathcal R_n$.  Lower renewals obey the finite native-source
      debit \eqref{eq:v4-low-renewal-sum}.
\item Every canonical high response contains a terminal vector core retaining
      at least half its norm.  The associated output and heat lengths are
      summable, and the actual post-record source satisfies the divergent Abel
      lower bound \eqref{eq:v4-forward-Abel-lower}.
\item The core is exactly the pressure-free quadratic stress--strain card
      \eqref{eq:v4-stress-realization}; no positive pressure, dissipation, or
      source component is separately charged.
\item The coherent replay of this assembled card vanishes on compact
      terminally regular regions and has the spatial alternatives
      \eqref{eq:v4-distance-separation},
      \eqref{eq:v4-heavy-stress-thread}, or
      \eqref{eq:v4-stress-microdelocalization} at the summable response scale.
\item A heavy coherent replay supplies either the local counterflow witness
      \eqref{eq:v4-stress-counterflow-survival} or the signed local stress margin
      \eqref{eq:v4-signed-local-stress-margin}.  Only the latter supplies the
      fixed-datum export card \eqref{eq:v4-export-card}.  The distance-separated
      and microdelocalized branches remain NS--A localization obstructions; a
      failed ancestry or localization premise remains a typed source/transport
      obstruction.
\end{enumerate}
Thus the auxiliary ``aborted virtual collapse'' is not excluded, but its later
chronology is no longer invisible: every abort is followed by a quantitatively
large lower-scale renewal, same-or-finer renewal, or spatially noncompact
successor packet.
\end{theorem}

\begin{proof}
Items \textup{(V4.1)--(V4.2)} are
\Cref{thm:v4-canonical-replenishment,thm:v4-no-silent-abort,cor:v4-virtual-front}.
Item \textup{(V4.3)} is
\Cref{thm:v4-terminal-vector-core,thm:v4-summable-core-lengths,cor:v4-forward-Abel-record}.
Item \textup{(V4.4)} is \Cref{thm:v4-stress-realization}.
Items \textup{(V4.5)--(V4.6)} are
\Cref{thm:v4-stress-singular-support,thm:v4-spatial-trichotomy,cor:v4-heavy-counterflow,cor:v4-export-card}.
The final interpretation is precisely the exhaustive low/high frequency split
and the subsequent spatial trichotomy.
\end{proof}

\begin{assessment}[How V4 changes the V3 attack order]
\label{ass:v4-direction}
V3 correctly reduced its selected caloric scalar to local near-pressure
replenishment, kinetic concentration, or localization failure.  The auxiliary
notes show that trying to make every such pre-record current small by a
marginal source modulus would repeat a sharp no-go.  V4 therefore inserts the
next record before continuing that estimate.  The resulting future packet is
forced unconditionally by the norm chronology and has a literal quadratic
stress writer.  A V3 material or pressure-localized thread is compared with
this packet only through an actual common-history incidence map.  Agreement
produces a cofinal local replenishment chain; disagreement is spatial
relocation, scale reset, represented ancestry, or a localization residual.
The two scalars are not silently identified.
\end{assessment}

\begingroup\small
\begin{longtable}{p{0.27\textwidth}p{0.16\textwidth}p{0.49\textwidth}}
\caption{NS--A Version V4 proof-status ledger.}\label{tab:v4-status}\\
\toprule Row & Status & Exact conclusion\\\midrule
\endfirsthead
\toprule Row & Status & Exact conclusion\\\midrule
\endhead
V1--V3 source & \PROVED{} preserved & Complete preterminal source retained byte-for-byte; no earlier theorem rewritten.\\[0.25em]
Auxiliary marginal-Dini and virtual-collapse findings & Reviewed / inherited & Used to change the attack order; their proofs are not copied.\\[0.25em]
Canonical energy head & \PROVED & $\|P_{\le K}e\|^2\le KE_*$ and $K_n=\mathcal R_n^2/E_*$ force a $\sqrt3\mathcal R_n$ next-record high tail.\\[0.25em]
Canonical post-record response & \PROVED & Interval-fresh high-frequency Duhamel term has norm at least $(\sqrt3-1)\mathcal R_n$.\\[0.25em]
Arbitrary virtual front & \PROVED & Lower- or higher-frequency fresh response of size $(\sqrt2-1)\mathcal R_n$; both may occur.\\[0.25em]
Lower-frequency renewal stock & \PROVED & $m_n\ge(\sqrt2-1)\mathcal R_n/\Lambda_n$ and the weighted debit is summable.\\[0.25em]
High-pass terminal age core & \PROVED & Exact inverse envelope; final age window retains at least half the vector response.\\[0.25em]
Canonical response lengths & \PROVED & Inverse output lengths and physical heat lengths are summable.\\[0.25em]
Forward Abel source record & \PROVED & Weighted source action on $J_n^+$ is at least $c\sqrt\nu\,\mathcal R_n$.\\[0.25em]
Quadratic stress realization & \PROVED & Pressure-free trace-free stress paired with one backward caloric strain.\\[0.25em]
Regular-region stress replay & \PROVED & Absolute stress carrier is $O(\widehat\sigma_n^{1/4})$ on each regular compact; normalized coherent mass vanishes.\\[0.25em]
Future stress singular-set support & \PROVED & Every spatial cluster law is supported on $\Sigma_T$.\\[0.25em]
Canonical spatial trichotomy & \PROVED & Relative separation, one heavy coherent stress replay, or near-singular microdelocalization.\\[0.25em]
Local signed-stress short & \PROVED & A heavy replay gives either a normalized counterflow/interface witness or an actual signed local stress margin of order $\mathcal R_n$.\\[0.25em]
NS--B export & {\scriptsize\CONDITIONAL} & Ready after counterflow, return, ancestry, and rank-at-most-two audits.\\[0.25em]
Identification with the V3 caloric thread & \OPEN & Requires a common-history localization/response incidence; not inferred from matching endpoints.\\[0.25em]
Exclusion of all renewal chains & \OPEN & Infinite lower resets, same-or-finer renewals, relocation, and microdelocalization are not ruled out.\\[0.25em]
Terminal Dini improvement & \OPEN & No marginal or ridge-weighted vanishing theorem is obtained.\\[0.25em]
Global three-dimensional regularity & \OPEN & Not claimed.\\
\bottomrule
\end{longtable}
\endgroup

\begin{programme}[Version V5: cofinal ancestry and spatial landing of the future packets]
\label{prog:v4-next}
Preserve Versions V1--V4 verbatim.  Work on the exact cards
\eqref{eq:v4-export-card}, not on a new ambient source norm.  First compare the
future response at record $n$ with the represented tail, local V3
near-replenishment card, and response at record $n+1$ through one common
history.  Short old transported ancestry before declaring the interval-fresh
response a new source.  The first cofinal alternatives are:
\begin{enumerate}[label=\textup{(V5.\arabic*)},leftmargin=5.2em]
\item a summable lower-frequency reset debit from
      \eqref{eq:v4-low-renewal-sum};
\item a same-or-finer terminal core whose output/age scales are transported to
      the next record;
\item spatial relocation or singular-distance/scale separation;
\item one heavy coherent stress replay, followed by its local signed-counterflow
      short and, on the passing branch, a fixed-datum donor audit;
\item stress microdelocalization requiring an aggregate source-relative
      capacity theorem.
\end{enumerate}
On the heavy branch, first apply \Cref{cor:v4-heavy-counterflow}, then the
inherited rank-at-most-two smooth handoff and represented-return short.  Only a surviving rank-three card is
exported to NS--B, where the paid kinetic-stress path and positive Reynolds
covariance are evaluated on this same finite-$n$ history.  On the diffuse
branch, assemble signed cells before taking positive parts and retain the
actual diameter and both scales $K_n^{-1}$ and $d_n^{\rm age}$.

Do not return to an $o(\sqrt\tau)$ marginal source estimate, a bounded
fast-time window, an uncorrected compact ancient-profile argument, or separate
positive prices for pressure, dissipation, low/high response, and return.  The
terminal Dini theorem is invoked only if its joint ridge-weighted hypothesis is
actually established.  If the cofinal cards are all represented, relocated
into a closed smooth carrier, or made scalar-null, invoke the inherited
critical continuation theorem immediately.
\end{programme}

\begin{firewall}[No global theorem in V4]
\label{fire:v4-no-global}
V4 proves a forward replenishment theorem, exact low/high scale alternatives,
a terminal vector core with summable response lengths, a pressure-free
quadratic stress card, regular-region vanishing, a canonical singular-set
spatial trichotomy, and an exact coherent-to-signed local counterflow short.  It does not prove that repeated replenishment is
impossible, identify every interval-fresh response as an orthogonal source
birth, control the NS--B paid-stress variation, eliminate the log-spread ridge,
or establish global regularity.
\end{firewall}

% =============================================================================
\section*{Version V4 verification and history}
\addcontentsline{toc}{section}{Version V4 verification and history}
% =============================================================================

\paragraph{Append-only integrity.}
The complete V3 preterminal byte string has SHA--256
\begin{center}\scriptsize\ttfamily
1952fc268d8a5abbec227ed8c29266cd4085d52b8a413c0615f412fca2706a57
\end{center}
and is preserved exactly before this V4 material begins.  Only the sole final
active document terminator was removed; one active terminator is restored
below.

\paragraph{Version V4 history.}
Reviewed the identical terminal-concentration reports and their first-pass
companion.  Retained their no-go for marginal Dini surrogates and adopted their
recommendation to use the time after a regular record.  Proved the canonical
energy-cutoff high-frequency replenishment theorem and the arbitrary-front
lower/higher renewal alternative.  Quantified the finite native source debit
of lower renewals.  Inverted the high-pass heat envelope to construct a
terminal vector core, proved a forward Abel source record, and obtained
summable canonical output and heat lengths.  Rewrote the core exactly as a
pressure-free trace-free kinetic-stress/backward-strain scalar.  Proved
absolute regular-region vanishing for its coherent replay, singular-set
support, and the response-scale trichotomy into relative separation, one heavy
coherent stress replay, or microdelocalization.  Shorted the heavy replay
against its actual signed local stress and retained local counterflow as an
explicit obstruction.  Identified only the signed passing card as the first
forward-chronology packet eligible for the inherited fixed-datum donor audit.  No global regularity conclusion was claimed.

For Version V5, preserve this complete source verbatim, remove only its sole
final active document terminator, append new mathematical material immediately
before it, and restore the terminator.  The next filename is
\begin{center}\scriptsize
\path{Renewal_NS_Terminal_Source_Concentration_and_Singular_Thread_Closure_Programme_V5.tex}.
\end{center}

\addtocontents{toc}{\protect\endgroup}
% =============================================================================
% END VERSION V4 MATERIAL
% =============================================================================


% =============================================================================
% BEGIN VERSION V5 MATERIAL -- append-only continuation of NS--A V1--V4
% =============================================================================
\clearpage
\hypersetup{pdftitle={NS-A Terminal Source Concentration Programme: Cumulative V5},pdfauthor={A. Pelissier}}
\makeatletter
\addtocontents{toc}{\protect\begingroup\protect\makeatletter}
\addtocontents{toc}{\protect\def\protect\l@section{\protect\@dottedtocline{1}{0em}{3.5em}}}
\addtocontents{toc}{\protect\makeatother}
\makeatother
\renewcommand{\thesection}{V5.\arabic{section}}
\renewcommand{\thesubsection}{V5.\arabic{section}.\arabic{subsection}}
\renewcommand{\theequation}{V5.\arabic{equation}}
\setcounter{section}{0}
\setcounter{equation}{0}
\renewcommand{\theHsection}{V5.\arabic{section}}
\renewcommand{\theHsubsection}{V5.\arabic{section}.\arabic{subsection}}
\renewcommand{\theHtheorem}{V5.\arabic{section}.\arabic{theorem}}
\renewcommand{\theHequation}{V5.\arabic{equation}}

\begin{center}
{\LARGE\bfseries NS--A: Version V5}\\[0.5em]
{\large Annular Future-Response Resolution, Parabolic Source--Output Transport,\\
and the First Two-Record Ancestry Short}\\[0.5em]
{\normalsize 2 September 2026}
\end{center}
\addcontentsline{toc}{part}{Version V5: annular future response, parabolic transport, and two-record ancestry}

\begin{abstract}
Versions V1--V4 are retained verbatim.  V4 proved that every doubled critical
record is followed by a canonical high-frequency nonlinear-source response,
extracted a terminal vector core with summable output and heat lengths, and
rewrote that core as a pressure-free quadratic stress--strain scalar.  Its
remaining alternatives were frequency ancestry, spatial relocation,
source-relative counterflow, one heavy singular stress thread, or
microdelocalization.  The present continuation resolves the first three rows
without returning to the time-marginal Dini route ruled out by the attached
terminal-concentration report.

The canonical core is first decomposed by a smooth dyadic Parseval partition
above the V4 energy cutoff.  This gives an exact positive response-energy law.
After passage to a subsequence, either that law is logarithmically spread over
infinitely many output annuli, in which case no uniformly finite band head
captures the response, or one annulus carries a fixed fraction of the core.
On the latter branch a source-first Riesz test produces one literal annular
scalar of size comparable with the doubled record.  The same scalar has two
exact representations: a pairing with the actual
$A^{-1/4}\mathcal B(u)$ source and a pairing of the trace-free kinetic stress
with one annular backward-caloric strain.  Its scalar source action is bounded
below by response squared divided by native source length; no shell is charged
as a separate source occurrence.

The annular strain has a Schwartz kernel on the physical scale
$\Gamma_n^{-1}$.  If its stress leverage
$\Gamma_n^{3/2}\int\|u\|_4^2/H_n^{\rm ann}$ is bounded, exponential age
damping and off-diagonal kernel decay force a fixed part of the signed scalar
into a parabolic near-diagonal region:
\[
 \beta_{n+1}-t\lesssim(\nu\Gamma_n^2)^{-1},
 \qquad d(x,y)\lesssim\Gamma_n^{-1}.
\]
If that leverage is unbounded, the exact escaping ratio is retained.  The
near-diagonal signed measure has both input and output spatial marginals
supported, in the limit, on the terminal singular set.  It admits the refined
trichotomy: singular-distance/annular-scale separation, one heavy joint
source--output cell, or near-singular microdelocalization.  A heavy positive
cell is shorted against its signed value before it becomes a physical local
margin.

For two selected annular responses, V5 transports the earlier vector to the
later endpoint and proves one exact Pythagoras.  The mismatch splits into an
orthogonal response-ancestry residual and an aligned amplitude-replenishment
residual.  If both are small, the annuli overlap and the inter-record time is
parabolically short unless the response-height ratio escapes.  On a heavy
spatial branch the annular radii are summable; consequently the centres either
form one Cauchy thread ending at a point of the singular set or have an
explicit relative relocation defect.  Only a signed heavy cell which survives
the literal represented dictionary and the inherited rank-at-most-two audit is
exported to the fixed-datum Reynolds/paid-stress programme.  Frequency spread,
stress leverage, counterflow, microdelocalization, relocation, and represented
reuse remain typed alternatives.  No terminal Dini improvement or global
three-dimensional regularity theorem is claimed.
\end{abstract}

% =============================================================================
\section{Scope after V4 and the auxiliary no-go results}
\label{sec:v5-scope}
% =============================================================================

The two identical terminal-concentration reports and their companion
first-pass programme were already audited in V4.  Their load-bearing
conclusions remain frozen: the ridge-weighted Dini functional is the exact
coherent scalar, a small coefficient in an unweighted $O(\!\sqrt\tau)$
time marginal does not remove logarithmically spread ridge payment, and an
argument using only times before one regular record cannot distinguish a true
terminal collapse from a collapse which aborts at that record.  V4 used the
missing future interval and forced the response
$F_n^{\rm can}$ in \eqref{eq:v4-canonical-response}.

\begin{assessment}[The exact V5 target]
\label{ass:v5-target}
V5 does not seek another upper bound for the original source marginal.  It
asks what the already forced vector $F_n^{\rm core}$ does after three pieces
of information are retained simultaneously:
\[
 \boxed{
 \text{output frequency}\quad+\quad
 \text{physical source/stress location}\quad+\quad
 \text{transport to the next record}.}
\]
The first output is a finite annular card or an all-shell response law.  The
second is a matched parabolic source--output packet or an explicit leverage,
counterflow, or delocalization defect.  The third is an old-response reuse,
aligned replenishment, or orthogonal ancestry residual.  No additional
compiler layer is introduced.
\end{assessment}

\begin{firewall}[Rows not reopened in V5]
\label{fire:v5-no-reopen}
The V35 Dini estimate, the V1 normalization correction, the V2 local-energy
absorption, the V3 affine-pressure annihilation and shrinking far-pressure
radius, the V4 record-to-record response lower bound, terminal vector core,
Abel source record, stress realization, regular-region vanishing, and
coherent spatial trichotomy are not re-proved.  Likewise V5 does not repeat
the donor programme's Fourier alias table, positive Reynolds-covariance
factorization, paid-stress reduction, or rank-at-most-two smooth theorem.
Those rows are invoked only after one literal V5 card reaches their declared
interface.
\end{firewall}

% =============================================================================
\section{A positive annular law for the canonical future response}
\label{sec:v5-annular-law}
% =============================================================================

Fix a nonnegative function $\psi\in C_c^\infty((1/2,2))$ satisfying
\begin{equation}
 \sum_{j\in\mathbb Z}\psi(2^{-j}s)^2=1,
 \qquad s>0,
 \qquad 0\le\psi\le1.
 \label{eq:v5-dyadic-partition}
\end{equation}
For each V4 record put
\begin{equation}
 \Psi_{n,j}:=\psi\!\left(\frac{\Lambda}{2^jK_n}\right),
 \qquad
 \Gamma_{n,j}:=2^jK_n,
 \qquad
 G_{n,j}:=\Psi_{n,j}F_n^{\rm core}.
 \label{eq:v5-annular-components}
\end{equation}
Since $F_n^{\rm core}=\mathsf P_{>K_n}F_n^{\rm core}$, only indices with
$\Gamma_{n,j}>K_n/2$ can contribute.  Define
\begin{equation}
 p_{n,j}:=\frac{\|G_{n,j}\|_2^2}{(H_n^{\rm core})^2},
 \qquad
 \mathfrak N_n^{\rm freq}:=
 \left(\sum_jp_{n,j}^2\right)^{-1}.
 \label{eq:v5-frequency-law}
\end{equation}

\begin{theorem}[Exact annular response-energy resolution]
\label{thm:v5-annular-resolution}
For every $n$,
\begin{equation}
 \boxed{
 \sum_j\|G_{n,j}\|_2^2=(H_n^{\rm core})^2,
 \qquad \sum_jp_{n,j}=1.}
 \label{eq:v5-annular-Parseval}
\end{equation}
After passage to a subsequence, exactly one of the following holds.
\begin{enumerate}[label=\textup{(F\arabic*)},leftmargin=3.5em]
\item \emph{Annular response concentration.}  There are $\eta>0$ and indices
      $j_n$ such that
      \begin{equation}
       \boxed{p_{n,j_n}\ge\eta.}
       \label{eq:v5-heavy-annulus}
      \end{equation}
\item \emph{Log-frequency response spreading.}
      \begin{equation}
       \boxed{\sup_jp_{n,j}\longrightarrow0,
       \qquad \mathfrak N_n^{\rm freq}\longrightarrow\infty.}
       \label{eq:v5-frequency-spread}
      \end{equation}
      For every fixed $M<\infty$ and every set $J_n\subset\mathbb Z$ with
      $|J_n|\le M$,
      \begin{equation}
       \boxed{
       \frac{\sum_{j\in J_n}\|G_{n,j}\|_2^2}
            {(H_n^{\rm core})^2}\longrightarrow0.}
       \label{eq:v5-no-finite-band-head}
      \end{equation}
\end{enumerate}
The second branch is an output-response statement.  It is not identified with
the auxiliary report's coherent-payment log-spread law without a same-history
source incidence.
\end{theorem}

\begin{proof}
The spectral multipliers commute and
$\sum_j\Psi_{n,j}^2=I$ on the positive spectrum.  Hence
\[
 \sum_j\|G_{n,j}\|_2^2
 =\sum_j\langle F_n^{\rm core},\Psi_{n,j}^2F_n^{\rm core}\rangle
 =\|F_n^{\rm core}\|_2^2,
\]
which proves \eqref{eq:v5-annular-Parseval}.  Along a subsequence either
$\sup_jp_{n,j}$ has positive lower limit, in which case choose almost
maximizing $j_n$, or it tends to zero.  On the latter branch
$\sum_jp_{n,j}^2\le\sup_jp_{n,j}$, so participation diverges.  Finally,
\[
 \sum_{j\in J_n}p_{n,j}
 \le M\sup_jp_{n,j}\longrightarrow0,
\]
proving \eqref{eq:v5-no-finite-band-head}.
\end{proof}

\begin{corollary}[Positive all-shell response stock on the spread branch]
\label{cor:v5-all-shell-stock}
On branch \textup{(F2)}, the positive trace-class operator
\begin{equation}
 \mathbb C_n^{\rm out}:=
 \sum_j G_{n,j}\otimes G_{n,j}
 \label{eq:v5-output-covariance}
\end{equation}
satisfies
\begin{equation}
 \boxed{
 \operatorname{tr}\mathbb C_n^{\rm out}
 =(H_n^{\rm core})^2
 \ge\frac{\gamma_*^2}{4}\mathcal R_n^2.}
 \label{eq:v5-output-covariance-trace}
\end{equation}
No uniformly finite collection of its annular blocks carries a positive
fraction of the trace.  This is a positive output square function of the V4
response; it is not yet the V13 positive Reynolds covariance.
\end{corollary}

\begin{proof}
Each rank-one summand is positive and its trace is
$\|G_{n,j}\|_2^2$.  Sum and use
\eqref{eq:v4-core-record-lower}.  The finite-head statement is
\eqref{eq:v5-no-finite-band-head}.
\end{proof}

% =============================================================================
\section{The exact annular source and kinetic-stress card}
\label{sec:v5-annular-card}
% =============================================================================

Work now on branch \textup{(F1)}.  Write
\begin{equation}
 \Psi_n=\Psi_{n,j_n},
 \qquad \Gamma_n=\Gamma_{n,j_n},
 \qquad G_n=G_{n,j_n},
 \qquad H_n^{\rm ann}=\|G_n\|_2.
 \label{eq:v5-selected-annulus}
\end{equation}
Then
\begin{equation}
 \Gamma_n>\frac{K_n}{2},
 \qquad
 H_n^{\rm ann}\ge\sqrt\eta H_n^{\rm core}
 \ge h_*\mathcal R_n,
 \qquad h_*:=\frac{\sqrt\eta\,\gamma_*}{2}.
 \label{eq:v5-annular-record-margin}
\end{equation}
Define the source-first annular test direction
\begin{equation}
 \boxed{
 \varpi_n:=\frac{\Psi_nG_n}{H_n^{\rm ann}}
 =\frac{\Psi_n^2F_n^{\rm core}}{H_n^{\rm ann}}.}
 \label{eq:v5-annular-test-direction}
\end{equation}
It obeys $\|\varpi_n\|_2\le1$ and has spectral support in
$[\Gamma_n/2,2\Gamma_n]$.  On the V4 core interval
$I_n^{\rm core}$ define
\begin{align}
 q_n^{\rm ann}(t)
 &:=\Lambda e^{-\nu(\beta_{n+1}-t)A}\varpi_n,
 \label{eq:v5-annular-q}\\
 \phi_n^{\rm ann}(t)
 &:=A^{-1/4}q_n^{\rm ann}(t)
 =A^{1/4}e^{-\nu(\beta_{n+1}-t)A}\varpi_n.
 \label{eq:v5-annular-phi}
\end{align}

\begin{theorem}[One annular response, two exact physical representations]
\label{thm:v5-annular-two-representations}
The selected annular response has the exact identities
\begin{equation}
 \boxed{
 \begin{aligned}
 H_n^{\rm ann}
 &=-\int_{I_n^{\rm core}}
   \langle c(t),q_n^{\rm ann}(t)\rangle\,\dd t\\
 &=\int_{I_n^{\rm core}}\int_{\Torus^3}
   (u\otimes u)^\circ:S(\phi_n^{\rm ann})\,\dd x\,\dd t.
 \end{aligned}}
 \label{eq:v5-annular-two-representations}
\end{equation}
Moreover,
\begin{equation}
 \boxed{
 \|q_n^{\rm ann}(t)\|_2
 \le C\Gamma_ne^{-c\nu(\beta_{n+1}-t)\Gamma_n^2},
 \qquad
 \|S(\phi_n^{\rm ann})(t)\|_2
 \le C\Gamma_n^{3/2}
 e^{-c\nu(\beta_{n+1}-t)\Gamma_n^2}.}
 \label{eq:v5-annular-operator-bounds}
\end{equation}
The source and stress formulas describe the same scalar and may not be charged
as two occurrences.
\end{theorem}

\begin{proof}
Since $G_n=\Psi_nF_n^{\rm core}$ and $\Psi_n$ is self-adjoint,
\[
 \langle F_n^{\rm core},\varpi_n\rangle
 =\frac{\langle F_n^{\rm core},\Psi_n^2F_n^{\rm core}\rangle}
       {H_n^{\rm ann}}
 =H_n^{\rm ann}.
\]
Insert the Duhamel definition \eqref{eq:v4-core-response}.  The support of
$\varpi_n$ lies above $K_n$, so the V4 high-pass projection acts as the
identity.  This gives the first line of
\eqref{eq:v5-annular-two-representations}.  The second follows exactly as in
\Cref{thm:v4-stress-realization}: move $A^{-1/4}$ by spectral calculus, drop
the Leray projection against the divergence-free test, integrate by parts,
and remove the isotropic stress against the trace-free strain.

On the support of $\varpi_n$, the spectral variable lies between
$\Gamma_n/2$ and $2\Gamma_n$.  Hence
\[
 \kappa e^{-\nu s\kappa^2}
 \le2\Gamma_ne^{-\nu s\Gamma_n^2/4},
 \qquad
 \kappa^{3/2}e^{-\nu s\kappa^2}
 \le C\Gamma_n^{3/2}e^{-\nu s\Gamma_n^2/4}.
\]
The first is the norm of $q_n^{\rm ann}$; one derivative together with
$A^{1/4}$ gives the second.
\end{proof}

Put
\begin{equation}
 \widehat m_n:=\int_{I_n^{\rm core}}\|c(t)\|_2\,\dd t,
 \qquad
 \widehat s_n:=\int_{I_n^{\rm core}}\|u(t)\|_4^2\,\dd t.
 \label{eq:v5-two-source-lengths}
\end{equation}
Both are lengths of the same fixed solution on the selected core.  The first
is the native $H^{-1/2}$ nonlinear-source length; the second is only the
$L^2$ length of the kinetic stress needed for the spatial kernel estimate.

\begin{theorem}[Source-visible annular action and finite scale debits]
\label{thm:v5-annular-source-action}
Let
\begin{equation}
 \dd\mu_n(t)=\|c(t)\|_2\,\dd t,
 \qquad
 \widehat c(t)=\frac{c(t)}{\|c(t)\|_2}
 \quad\text{on }\{c(t)\ne0\},
 \label{eq:v5-source-polar}
\end{equation}
and define
\begin{equation}
 a_n(t):=-\operatorname{Re}
 \langle\widehat c(t),q_n^{\rm ann}(t)\rangle,
 \qquad
 \mathcal A_n^{\rm src}:=\int|a_n(t)|^2\,\dd\mu_n(t).
 \label{eq:v5-annular-source-action}
\end{equation}
Then
\begin{equation}
 \boxed{
 \mathcal A_n^{\rm src}
 \ge\frac{(H_n^{\rm ann})^2}{\widehat m_n},
 \qquad
 H_n^{\rm ann}\le C\Gamma_n\widehat m_n.}
 \label{eq:v5-source-action-lower}
\end{equation}
Also
\begin{equation}
 \boxed{H_n^{\rm ann}\le C\Gamma_n^{3/2}\widehat s_n.}
 \label{eq:v5-stress-length-bound}
\end{equation}
Consequently, on every annular-concentration subsequence,
\begin{equation}
 \boxed{
 \sum_n\frac{\mathcal R_n}{\Gamma_n}<\infty,
 \qquad
 \sum_n\frac{\mathcal R_n}{\Gamma_n^{3/2}}<\infty.}
 \label{eq:v5-annular-scale-debits}
\end{equation}
These are scale constraints on one response chronology, not counts of source
births.
\end{theorem}

\begin{proof}
The first representation in
\eqref{eq:v5-annular-two-representations} gives
$H_n^{\rm ann}=\int a_n\,\dd\mu_n$.  Cauchy--Schwarz yields
\[
 (H_n^{\rm ann})^2
 \le\widehat m_n\mathcal A_n^{\rm src}.
\]
The first operator bound in
\eqref{eq:v5-annular-operator-bounds}, with the exponential discarded, gives
$H_n^{\rm ann}\le C\Gamma_n\widehat m_n$.

Let $T_u=(u\otimes u)^\circ$.  Since
$\|T_u\|_2\le\|u\otimes u\|_2=\|u\|_4^2$, the stress representation and the
second operator bound give
\[
 H_n^{\rm ann}
 \le\int_{I_n^{\rm core}}\|T_u(t)\|_2
       \|S(\phi_n^{\rm ann})(t)\|_2\,\dd t
 \le C\Gamma_n^{3/2}\widehat s_n.
\]
The intervals $I_n^{\rm core}\subset J_n^+$ are disjoint.  Hence
$\sum_n\widehat m_n\le\mathfrak M_c$.  Also
\[
 \|u(t)\|_4^{8/3}
 \le\|u(t)\|_2^{2/3}\|u(t)\|_6^2
 \le C E_*^{1/3}\|\nabla u(t)\|_2^2,
\]
so $\|u\|_4^2\in L^{4/3}(0,T)\subset L^1(0,T)$ and
$\sum_n\widehat s_n<\infty$.  Divide
\eqref{eq:v5-source-action-lower} and
\eqref{eq:v5-stress-length-bound} by the lower margin
\eqref{eq:v5-annular-record-margin}, then sum.
\end{proof}

\begin{assessment}[What the annular action does not prove]
The lower bound for $\mathcal A_n^{\rm src}$ identifies the amount of
source-visible scalar gain forced on the selected annulus.  It is not an upper
bound, a Carleson theorem, or the V13 donor action.  The scale sums in
\eqref{eq:v5-annular-scale-debits} use one occurrence on disjoint time
intervals.  They do not license summation of the positive parts of every
frequency representation of that occurrence.
\end{assessment}

% =============================================================================
\section{Annular kernel localization and parabolic source--output transport}
\label{sec:v5-parabolic-transport}
% =============================================================================

Choose $\widetilde\psi\in C_c^\infty((1/4,4))$ with
$\widetilde\psi=1$ on $[1/2,2]$.  For a frequency $\Gamma\ge1$ and age
$s\ge0$, let $\mathcal K_{\Gamma,s}(x,y)$ be the tensor-valued periodic kernel
of
\begin{equation}
 \mathcal S_{\Gamma,s}
 :=S A^{1/4}e^{-\nu sA}
   \widetilde\psi(\Lambda/\Gamma),
 \label{eq:v5-annular-strain-operator}
\end{equation}
acting from divergence-free vectors to symmetric trace-free tensors.

\begin{lemma}[Uniform periodic annular kernel bound]
\label{lem:v5-annular-kernel}
For every $C_0<\infty$ and every integer $N\ge0$, there is
$C_{N,C_0}$ such that, whenever
$0\le\nu s\Gamma^2\le C_0$,
\begin{equation}
 \boxed{
 |\mathcal K_{\Gamma,s}(x,y)|
 \le C_{N,C_0}\Gamma^{9/2}
       (1+\Gamma d(x,y))^{-N}.}
 \label{eq:v5-kernel-pointwise}
\end{equation}
Consequently, for every $M>0$ and $R\ge1$, the part of the kernel on
$d(x,y)>R/\Gamma$ defines an $L^2$ operator of norm at most
\begin{equation}
 \boxed{
 C_{M,C_0}\Gamma^{3/2}(1+R)^{-M}.}
 \label{eq:v5-kernel-offdiagonal}
\end{equation}
\end{lemma}

\begin{proof}
On divergence-free mean-zero fields the Stokes operator is $-\Delta$ and its
fractional powers are scalar Fourier multipliers.  After writing
$\tau=\nu s\Gamma^2$, the Euclidean symbol of
\eqref{eq:v5-annular-strain-operator} has the form
\[
 \Gamma^{3/2}\,M_\tau(\xi/\Gamma),
\]
where $M_\tau$ is smooth, supported in the fixed annulus
$1/4<|\xi|<4$, and has Schwartz seminorms bounded uniformly for
$0\le\tau\le C_0$.  Its Euclidean kernel is therefore
$\Gamma^{9/2}k_\tau(\Gamma z)$ with
$|k_\tau(z)|\le C_{N,C_0}(1+|z|)^{-N}$.  Periodize over the fixed torus
lattice.  The zero image gives the displayed term and the nonzero lattice
images satisfy the same bound after increasing the constant, proving
\eqref{eq:v5-kernel-pointwise}.

Integrating the right-hand side on
$d(x,y)>R/\Gamma$ gives
\[
 C_N\Gamma^{9/2}\Gamma^{-3}
 \int_R^\infty(1+r)^{-N}r^2\,\dd r
 \le C_M\Gamma^{3/2}(1+R)^{-M}
\]
when $N$ is chosen larger than $M+3$.  The same bound holds for both Schur
integrals in $x$ and $y$, so the Schur test gives
\eqref{eq:v5-kernel-offdiagonal}.
\end{proof}

Define the dimensionless annular stress leverage
\begin{equation}
 \boxed{
 \lambda_n^{\rm str}:=
 \frac{\Gamma_n^{3/2}\widehat s_n}{H_n^{\rm ann}}.}
 \label{eq:v5-stress-leverage}
\end{equation}
By \eqref{eq:v5-stress-length-bound}, it is bounded below by a positive
constant, but it need not be bounded above.

\begin{theorem}[Stress-leverage escape or matched parabolic transport]
\label{thm:v5-parabolic-transport}
After passage to a subsequence, exactly one of the following occurs.
\begin{enumerate}[label=\textup{(P\arabic*)},leftmargin=3.5em]
\item \emph{Stress-leverage escape:}
      \begin{equation}
       \boxed{\lambda_n^{\rm str}\longrightarrow\infty.}
       \label{eq:v5-stress-leverage-escape}
      \end{equation}
\item \emph{Matched parabolic source--output transport:} there is
      $L_*<\infty$ with $\lambda_n^{\rm str}\le L_*$.  There are constants
      $R_t,R_x<\infty$, depending only on $L_*$ and the fixed annular
      multipliers, for which the following holds.  Put
      \begin{equation}
       \tau_n^{\rm ann}
       :=\min\left\{\widehat\sigma_n,
                    \frac{R_t}{\nu\Gamma_n^2}\right\},
       \qquad
       I_n^{\rm ann}:=
       (\beta_{n+1}-\tau_n^{\rm ann},\beta_{n+1}).
       \label{eq:v5-matched-age}
      \end{equation}
      Define the signed measure on
      $I_n^{\rm ann}\times\Torus^3_x\times\Torus^3_y$ by
      \begin{equation}
       \boxed{
       \begin{aligned}
       \dd\Xi_n(t,x,y):={}&
       (u\otimes u)^\circ(t,x):
       \mathcal K_{\Gamma_n,\beta_{n+1}-t}(x,y)
       \varpi_n(y)\\
       &\times\one_{\{d(x,y)\le R_x/\Gamma_n\}}
       \,\dd y\,\dd x\,\dd t.
       \end{aligned}}
       \label{eq:v5-transport-measure}
      \end{equation}
      Then its signed resultant satisfies
      \begin{equation}
       \boxed{
       \Xi_n(\Omega_n^{\rm tr})
       \ge\frac12H_n^{\rm ann}
       \ge\frac{h_*}{2}\mathcal R_n.}
       \label{eq:v5-near-transport-margin}
      \end{equation}
\end{enumerate}
Thus bounded stress leverage forces a fixed part of the actual quadratic
scalar into the same terminal age and the same spatial scale as its selected
output annulus.
\end{theorem}

\begin{proof}
Pass to a subsequence on which $\lambda_n^{\rm str}$ either tends to infinity
or is bounded by $L_*$.  Work on the bounded branch.  By
\eqref{eq:v5-annular-operator-bounds}, the part of the stress scalar at ages
$s=\beta_{n+1}-t\ge R_t/(\nu\Gamma_n^2)$ has absolute value at most
\[
 C\Gamma_n^{3/2}e^{-cR_t}\widehat s_n
 =Ce^{-cR_t}\lambda_n^{\rm str}H_n^{\rm ann}.
\]
Choose $R_t$ so this is at most $H_n^{\rm ann}/4$.  If the full core is
shorter, that old-age part is empty.

On the remaining interval one has
$\nu s\Gamma_n^2\le R_t$.  Use
\Cref{lem:v5-annular-kernel} and split the kernel at distance
$R_x/\Gamma_n$.  The far part of the scalar is bounded by
\[
 C_M\Gamma_n^{3/2}(1+R_x)^{-M}
 \int_{I_n^{\rm ann}}\|u(t)\|_4^2\,\dd t
 \le C_M(1+R_x)^{-M}L_*H_n^{\rm ann}.
\]
Choose $M$ and then $R_x$ so this is at most
$H_n^{\rm ann}/4$.  The full scalar is $H_n^{\rm ann}>0$ by
\Cref{thm:v5-annular-two-representations}.  Removing an old-age contribution
and a spatially far contribution of total absolute value at most
$H_n^{\rm ann}/2$ leaves the near signed resultant
\eqref{eq:v5-near-transport-margin}.
\end{proof}

\begin{firewall}[Bounded leverage is not an energy estimate]
\label{fire:v5-leverage-scope}
The second branch of \Cref{thm:v5-parabolic-transport} is an exact implication
from bounded \emph{relative} stress leverage.  V5 does not derive that bound
from kinetic energy, dissipation, or the V4 Abel record.  On the first branch
the escaping ratio \eqref{eq:v5-stress-leverage-escape} is retained as the
physical obstruction; it may not be replaced by a shell count or by a new
positive source price.
\end{firewall}

% =============================================================================
\section{Singular-set landing of the near-diagonal transport law}
\label{sec:v5-transport-landing}
% =============================================================================

On branch \textup{(P2)}, write
\begin{equation}
 H_n^{\rm tr}:=\Xi_n(\Omega_n^{\rm tr})>0,
 \qquad
 Z_n^{\rm tr,+}:=\Xi_{n,+}(\Omega_n^{\rm tr}),
 \label{eq:v5-transport-resultant}
\end{equation}
and define the coherent replay
\begin{equation}
 \dd\kappa_n^{\rm tr}:=
 \frac{H_n^{\rm tr}}{Z_n^{\rm tr,+}}\,\dd\Xi_{n,+},
 \qquad
 \dd\pi_n^{\rm tr}:=
 \frac{\dd\Xi_{n,+}}{Z_n^{\rm tr,+}}.
 \label{eq:v5-transport-replay}
\end{equation}
Let $\rho_n^x$ and $\rho_n^y$ be the input and output spatial marginals of
$\pi_n^{\rm tr}$.  By construction,
\begin{equation}
 \boxed{d(x,y)\le R_x\Gamma_n^{-1}}
 \label{eq:v5-near-diagonal-support}
\end{equation}
throughout the transport measure.

\begin{theorem}[Both transport marginals land on the terminal singular set]
\label{thm:v5-transport-singular-support}
For every compact $K\Subset\mathcal R_T$,
\begin{equation}
 \boxed{\rho_n^x(K)+\rho_n^y(K)\longrightarrow0.}
 \label{eq:v5-two-marginal-regular-vanishing}
\end{equation}
Consequently every weak-* cluster point of either marginal is supported on
$\Sigma_T$.  The input and output cluster points have the same support in the
following weak sense: their distance in the Prokhorov metric is at most
$R_x/\Gamma_n$ and therefore tends to zero.
\end{theorem}

\begin{proof}
Choose $K'$ compactly contained in $\mathcal R_T$ with a positive collar
around $K$.  Terminal regularity gives $C_K$ and $t_K<T$ such that
$\|u(t)\|_{L^4(K')}^2\le C_K$ for $t_K<t<T$.  The annular intervals lie in
this time slab for late $n$.  The $L^1$ norm in the output variable of the
absolute kernel in \eqref{eq:v5-kernel-pointwise} is at most
$C\Gamma_n^{3/2}$.  Therefore Young's inequality gives
\[
 |\Xi_n|\{x\in K\}
 \le C_K\Gamma_n^{3/2}|I_n^{\rm ann}|
 \le\frac{C_KR_t}{\nu\Gamma_n^{1/2}}
 \longrightarrow0.
\]
If $y\in K$, then \eqref{eq:v5-near-diagonal-support} places $x$ in $K'$ for
late $n$, so the same estimate gives
$|\Xi_n|\{y\in K\}\to0$.  Since
$0\le\kappa_n^{\rm tr}\le\Xi_{n,+}\le|\Xi_n|$ and
$H_n^{\rm tr}\ge(h_*/2)\mathcal R_n$, division by the coherent mass proves
\eqref{eq:v5-two-marginal-regular-vanishing}.

For the final assertion, every coupling $\pi_n^{\rm tr}$ of its two marginals
is supported on pairs at distance at most $R_x/\Gamma_n$.  The standard
coupling characterization of the Prokhorov distance gives the stated bound.
\end{proof}

Put $r_n^{\rm ann}:=\Gamma_n^{-1}$ and, for $R<\infty$, define
\begin{equation}
 \mathcal N_n^{\rm ann}(R)
 :=\{x:\dist(x,\Sigma_T)\le Rr_n^{\rm ann}\},
 \qquad
 \eta_n^{\rm ann}(R):=\rho_n^x(\mathcal N_n^{\rm ann}(R)).
 \label{eq:v5-annular-singular-neighbourhood}
\end{equation}

\begin{theorem}[Joint annular transport trichotomy]
\label{thm:v5-joint-spatial-trichotomy}
After passage to a subsequence, exactly one of the following occurs.
\begin{enumerate}[label=\textup{(J\arabic*)},leftmargin=3.5em]
\item \emph{Singular-distance/annular-scale separation.}  For every fixed
      $R<\infty$,
      \begin{equation}
       \boxed{\eta_n^{\rm ann}(R)\longrightarrow0.}
       \label{eq:v5-annular-distance-separation}
      \end{equation}
\item \emph{One heavy near-diagonal singular cell.}  There are finite
      $R,A$, positive $p$, centres $x_n$, and points $z_n\in\Sigma_T$ such
      that
      \begin{equation}
       \boxed{
       \begin{aligned}
       &\pi_n^{\rm tr}\{x\in B_{Ar_n^{\rm ann}}(x_n),
          \ y\in B_{(A+R_x)r_n^{\rm ann}}(x_n)\}\ge p,\\
       &d(x_n,z_n)\le(A+R)r_n^{\rm ann}.
       \end{aligned}}
       \label{eq:v5-heavy-joint-cell}
      \end{equation}
      After a further subsequence, $x_n,z_n\to x_*\in\Sigma_T$.
\item \emph{Near-singular transport microdelocalization.}  There are finite
      $R$ and $\eta>0$ with $\eta_n^{\rm ann}(R)\ge\eta$, while for every
      fixed $A<\infty$,
      \begin{equation}
       \boxed{
       \sup_x\rho_n^x
       (B_{Ar_n^{\rm ann}}(x)\cap\mathcal N_n^{\rm ann}(R))
       \longrightarrow0.}
       \label{eq:v5-transport-microdelocalization}
      \end{equation}
\end{enumerate}
The same alternatives hold for the output marginal after enlarging every
radius by the fixed amount $R_xr_n^{\rm ann}$.
\end{theorem}

\begin{proof}
Apply the diagonal concentration argument of
\Cref{thm:v4-spatial-trichotomy} to the probability laws $\rho_n^x$ and the
scales $r_n^{\rm ann}$.  If every relative singular neighbourhood has
vanishing mass, item \textup{(J1)} holds.  Otherwise fix one neighbourhood
carrying positive mass.  If some ball of a fixed multiple of
$r_n^{\rm ann}$ carries a fixed fraction, choose it.  Its intersection with
$\mathcal N_n^{\rm ann}(R)$ supplies $z_n\in\Sigma_T$ at the displayed
distance.  The near-diagonal support
\eqref{eq:v5-near-diagonal-support} places the corresponding output coordinate
in the enlarged ball, proving \eqref{eq:v5-heavy-joint-cell}.  Compactness of
$\Sigma_T$ gives the common cluster point.  If no such ball exists, item
\textup{(J3)} holds.  The output statements follow from the same deterministic
distance bound.
\end{proof}

\begin{corollary}[A heavy transport replay must pass a signed local short]
\label{cor:v5-heavy-transport-short}
On branch \textup{(J2)}, let $\mathcal C_n^{\rm tr}$ denote the joint cell in
\eqref{eq:v5-heavy-joint-cell} together with the matched age interval, and put
\begin{equation}
 \Delta_n^{\rm tr}:=
 \left|\kappa_n^{\rm tr}(\mathcal C_n^{\rm tr})
       -\Xi_n(\mathcal C_n^{\rm tr})\right|.
 \label{eq:v5-transport-counterflow}
\end{equation}
After passage to a subsequence, at least one of the following holds.
\begin{enumerate}[label=\textup{(C\arabic*)},leftmargin=3.5em]
\item \emph{Transport counterflow or replay mismatch:}
      \begin{equation}
       \boxed{\limsup_n
       \frac{\Delta_n^{\rm tr}}{\mathcal R_n}>0.}
       \label{eq:v5-transport-counterflow-survival}
      \end{equation}
\item \emph{Signed local source--output transport:}
      \begin{equation}
       \boxed{
       \Xi_n(\mathcal C_n^{\rm tr})
       \ge\frac{ph_*}{4}\mathcal R_n}
       \label{eq:v5-signed-local-transport}
      \end{equation}
      for all sufficiently large $n$.
\end{enumerate}
The second branch is a signed scalar on one matched parabolic card.  The first
is retained as one complete local counterflow/interface obstruction.
\end{corollary}

\begin{proof}
By \eqref{eq:v5-transport-replay},
\[
 \kappa_n^{\rm tr}(\mathcal C_n^{\rm tr})
 =H_n^{\rm tr}\pi_n^{\rm tr}(\mathcal C_n^{\rm tr})
 \ge pH_n^{\rm tr}
 \ge\frac{ph_*}{2}\mathcal R_n.
\]
If the first branch fails, pass to a subsequence on which
$\Delta_n^{\rm tr}/\mathcal R_n\to0$ and subtract it from the preceding lower
bound.
\end{proof}

% =============================================================================
\section{The first literal two-record response-ancestry short}
\label{sec:v5-two-record-ancestry}
% =============================================================================

The annular law is defined at the endpoint $b_n:=\beta_{n+1}$.  Let
$\mathcal A\subset\mathbb N$ be a cofinal set on which annular concentration
has been selected.  For two consecutive elements $m<n$ of $\mathcal A$, put
\begin{equation}
 \tau_{m,n}:=b_n-b_m,
 \qquad
 V_{m\to n}:=e^{-\nu\tau_{m,n}A}G_m.
 \label{eq:v5-transported-old-annulus}
\end{equation}
Let $P_{m\to n}$ be the orthogonal projection onto
$\Span\{V_{m\to n}\}$, with $P_{m\to n}=0$ if the vector vanishes, and set
\begin{equation}
 R_{m,n}^{\perp}:=(I-P_{m\to n})G_n,
 \qquad
 R_{m,n}^{\parallel}:=P_{m\to n}G_n-V_{m\to n}.
 \label{eq:v5-two-ancestry-residuals}
\end{equation}

\begin{theorem}[Exact old-tail, aligned-replenishment, and orthogonal-residual Pythagoras]
\label{thm:v5-two-record-Pythagoras}
For every consecutive selected pair,
\begin{equation}
 \boxed{
 \|G_n-V_{m\to n}\|_2^2
 =\|R_{m,n}^{\perp}\|_2^2
  +\|R_{m,n}^{\parallel}\|_2^2.}
 \label{eq:v5-two-record-Pythagoras}
\end{equation}
Fix $0<\epsilon<1/(2\sqrt2)$.  At least one of the following holds.
\begin{enumerate}[label=\textup{(A\arabic*)},leftmargin=3.5em]
\item \emph{Orthogonal response ancestry:}
      \begin{equation}
       \boxed{\|R_{m,n}^{\perp}\|_2\ge\epsilon H_n^{\rm ann}.}
       \label{eq:v5-orthogonal-ancestry}
      \end{equation}
\item \emph{Aligned amplitude replenishment:}
      \begin{equation}
       \boxed{\|R_{m,n}^{\parallel}\|_2
       \ge\epsilon H_n^{\rm ann}.}
       \label{eq:v5-aligned-replenishment}
      \end{equation}
\item \emph{Actual immediate old-annulus reuse:} both preceding inequalities
      fail.  Then the two annular supports intersect,
      \begin{equation}
       \boxed{\frac14\le\frac{\Gamma_n}{\Gamma_m}\le4,}
       \label{eq:v5-comparable-annuli}
      \end{equation}
      and
      \begin{equation}
       \boxed{
       e^{-\nu\Gamma_m^2\tau_{m,n}/4}H_m^{\rm ann}
       \ge(1-\sqrt2\epsilon)H_n^{\rm ann}.}
       \label{eq:v5-reuse-damping-gate}
      \end{equation}
      In particular, if $H_m^{\rm ann}/H_n^{\rm ann}\le C_*$, then
      \begin{equation}
       \boxed{
       \nu\Gamma_m^2\tau_{m,n}
       \le4\log\frac{C_*}{1-\sqrt2\epsilon}.}
       \label{eq:v5-parabolic-reuse-time}
      \end{equation}
\end{enumerate}
The first residual is a candidate new response direction.  The second is a
new payment along an old response direction and does not enlarge source rank.
The third is the only branch on which the selected annular vector itself is
reused with its transported coefficient.
\end{theorem}

\begin{proof}
The first term in \eqref{eq:v5-two-ancestry-residuals} is orthogonal to
$\Ran P_{m\to n}$, while the second belongs to that one-dimensional range.
Their sum is $G_n-V_{m\to n}$, proving
\eqref{eq:v5-two-record-Pythagoras}.

Suppose both residual bounds fail.  Then
\[
 \|G_n-V_{m\to n}\|_2
 <\sqrt2\epsilon H_n^{\rm ann},
\]
so
\begin{equation}
 \|V_{m\to n}\|_2
 \ge(1-\sqrt2\epsilon)H_n^{\rm ann}>0.
 \label{eq:v5-reuse-norm-lower}
\end{equation}
The semigroup preserves Fourier support.  If the annuli
$[\Gamma_m/2,2\Gamma_m]$ and
$[\Gamma_n/2,2\Gamma_n]$ were disjoint, then $G_n$ and
$V_{m\to n}$ would be orthogonal, so
$\|G_n-V_{m\to n}\|_2\ge H_n^{\rm ann}$, a contradiction.  Their
intersection is equivalent to \eqref{eq:v5-comparable-annuli}.
Finally, the spectral support of $G_m$ gives
\[
 \|V_{m\to n}\|_2
 \le e^{-\nu\Gamma_m^2\tau_{m,n}/4}H_m^{\rm ann}.
\]
Combine this with \eqref{eq:v5-reuse-norm-lower}; if the height ratio is at
most $C_*$, take logarithms.
\end{proof}

\begin{corollary}[Exact short against the literal represented dictionary]
\label{cor:v5-represented-dictionary}
At endpoint $b_n$, let $\mathcal J_n^{\rm rep}$ be the closed span of a
\emph{finite literal same-history dictionary}: the transported preceding
annular response, every V3 near-replenishment vector which has an actual map
to $b_n$, and every declared represented return vector.  Let
$P_n^{\rm rep}$ be its orthogonal projection.  Then
\begin{equation}
 \boxed{
 (H_n^{\rm ann})^2
 =\|P_n^{\rm rep}G_n\|_2^2
  +\|(I-P_n^{\rm rep})G_n\|_2^2.}
 \label{eq:v5-represented-Pythagoras}
\end{equation}
A lower bound
\begin{equation}
 \|(I-P_n^{\rm rep})G_n\|_2
 \ge\theta H_n^{\rm ann}
 \label{eq:v5-heldout-response}
\end{equation}
produces a literal held-out response of size
$\theta h_*\mathcal R_n$.  If this residual tends to zero relatively, the
annular output is represented by the declared dictionary and receives no new
source price.  The theorem does not assert that a V3 or return vector belongs
to the dictionary before its common-endpoint transport has been physically
specified.
\end{corollary}

\begin{proof}
This is orthogonal Pythagoras.  The lower bound uses
\eqref{eq:v5-annular-record-margin}.  The final statements are the definitions
of the represented and held-out branches.
\end{proof}

\begin{firewall}[Time support is not source first birth]
\label{fire:v5-time-not-birth}
The V4 response is integrated over a new time interval, but interval
disjointness does not by itself make its nonlinear source an orthogonal
physical first birth.  Equations \eqref{eq:v5-two-record-Pythagoras} and
\eqref{eq:v5-represented-Pythagoras} short the response after actual endpoint
transport.  A surviving aligned mismatch is replenishment of an old direction;
only the held-out component can enlarge the response-source atlas, and even
that component must pass the inherited source-ancestry audit.
\end{firewall}

% =============================================================================
\section{Summable annular radii and the cofinal centre chain}
\label{sec:v5-centre-chain}
% =============================================================================

On every selected annular branch,
$\Gamma_n>K_n/2$, so
\begin{equation}
 \boxed{
 \sum_n r_n^{\rm ann}
 =\sum_n\Gamma_n^{-1}
 \le2\sum_nK_n^{-1}<\infty.}
 \label{eq:v5-summable-annular-radii}
\end{equation}
Assume now the heavy joint-cell branch \textup{(J2)} and retain its centres
$x_n$ and singular points $z_n$.

\begin{theorem}[Cofinal singular thread or relative spatial relocation]
\label{thm:v5-centre-chain}
After passage to a subsequence, exactly one of the following occurs.
\begin{enumerate}[label=\textup{(S\arabic*)},leftmargin=3.5em]
\item \emph{Summable centre transport.}  There is $C_*<\infty$ such that for
      consecutive selected indices $m<n$,
      \begin{equation}
       \boxed{
       d(x_m,x_n)\le C_*(r_m^{\rm ann}+r_n^{\rm ann}).}
       \label{eq:v5-centre-step}
      \end{equation}
      The centres form a Cauchy sequence and converge to one point
      $x_*\in\Sigma_T$.  The sum of their successive travel distances is
      finite.
\item \emph{Relative relocation.}  Along the selected pairs,
      \begin{equation}
       \boxed{
       \frac{d(x_m,x_n)}{r_m^{\rm ann}+r_n^{\rm ann}}
       \longrightarrow\infty.}
       \label{eq:v5-relative-relocation}
      \end{equation}
\end{enumerate}
Thus a heavy packet either lands on one cofinal singular thread at its own
annular scale or retains an explicit changing-centre obstruction.
\end{theorem}

\begin{proof}
Apply subsequence selection to the nonnegative relative step ratios.  They are
either bounded or tend to infinity.  On the bounded branch,
\[
 \sum d(x_m,x_n)
 \le2C_*\sum_nr_n^{\rm ann}<\infty
\]
by \eqref{eq:v5-summable-annular-radii}.  Hence $x_n$ is Cauchy.  From
\eqref{eq:v5-heavy-joint-cell},
$d(x_n,z_n)\le C r_n^{\rm ann}\to0$ with $z_n\in\Sigma_T$; closedness of
$\Sigma_T$ places the common limit in $\Sigma_T$.
\end{proof}

\begin{corollary}[Material capture is the exact interface between ancestry and centre transport]
\label{cor:v5-material-capture-interface}
Suppose an independently supplied same-history material/response incidence
has the following property: whenever the immediate old-annulus reuse branch
of \Cref{thm:v5-two-record-Pythagoras} passes, the transported packet is
captured inside a fixed enlargement of the two heavy joint cells.  Then
relative relocation \eqref{eq:v5-relative-relocation} is incompatible with
old-annulus reuse.  Hence either an ancestry residual survives, or the centres
obey \eqref{eq:v5-centre-step} and form one singular thread.
\end{corollary}

\begin{proof}
A fixed enlargement of cells of radii
$r_m^{\rm ann}$ and $r_n^{\rm ann}$ can intersect only when their centre
distance is bounded by a fixed multiple of the sum of those radii.  This is
\eqref{eq:v5-centre-step}.  Its failure is therefore a failure of the stated
capture incidence or of immediate old-annulus reuse.
\end{proof}

\begin{firewall}[Spatial cluster is weaker than material ancestry]
\label{fire:v5-cluster-not-capture}
Compactness of the torus always gives subsequential centre convergence.  The
content of \Cref{thm:v5-centre-chain} is the much stronger summable travel
bound at the physical annular radii.  V5 does not derive material capture from
weak spatial convergence, matching singular endpoints, or overlap of scalar
marginals.  The common-history incidence in
\Cref{cor:v5-material-capture-interface} remains a separate physical row.
\end{firewall}

% =============================================================================
\section{Integrated V5 alternative and the first eligible donor export}
\label{sec:v5-integrated}
% =============================================================================

\begin{theorem}[Version V5 annular transport and ancestry theorem]
\label{thm:v5-integrated}
Retain the V4 canonical record response and terminal core.  After passage to a
subsequence, the following exhaustive ordered alternatives hold.
\begin{enumerate}[label=\textup{(V5.\arabic*)},leftmargin=5.2em]
\item The positive output-response law is logarithmically spread as in
      \eqref{eq:v5-frequency-spread}.  It carries the all-shell trace lower
      bound \eqref{eq:v5-output-covariance-trace}, while every uniformly finite
      band head is negligible.
\item One annulus carries a fixed fraction of the response.  It has the exact
      source/stress representations
      \eqref{eq:v5-annular-two-representations}, source-action lower bound
      \eqref{eq:v5-source-action-lower}, and finite scale debits
      \eqref{eq:v5-annular-scale-debits}.
\item On the annular branch, either stress leverage escapes through
      \eqref{eq:v5-stress-leverage-escape}, or a fixed scalar margin lies in
      the matched parabolic near-diagonal region
      \eqref{eq:v5-matched-age}--\eqref{eq:v5-near-transport-margin}.
\item The near-diagonal transport law has both marginals supported on
      $\Sigma_T$ in the limit and satisfies the relative-distance, heavy-cell,
      or microdelocalized alternatives of
      \Cref{thm:v5-joint-spatial-trichotomy}.
\item A heavy cell carries either the local counterflow/replay defect
      \eqref{eq:v5-transport-counterflow-survival} or the signed local
      source--output scalar \eqref{eq:v5-signed-local-transport}.
\item Across two records, the selected response has the exact ancestry split
      \eqref{eq:v5-two-record-Pythagoras}.  It yields an orthogonal response
      residual, an aligned amplitude replenishment, or immediate old-annulus
      reuse.  On the reuse branch the scales are comparable and the bridge is
      parabolically short unless the response-height ratio escapes.
\item Against the complete literal represented dictionary, a fixed held-out
      component satisfies \eqref{eq:v5-heldout-response}, or the output is
      represented and receives no new source price.
\item On heavy cells the centres either form one summably transported singular
      thread or satisfy the relative relocation law
      \eqref{eq:v5-relative-relocation}.  Turning the first alternative into
      material ancestry requires the explicit capture incidence of
      \Cref{cor:v5-material-capture-interface}.
\end{enumerate}
No positive part is taken before the corresponding signed scalar is assembled,
and no shell, cell, stress representation, or return is counted as a second
nonlinear occurrence.
\end{theorem}

\begin{proof}
Items \textup{(V5.1)--(V5.2)} are
\Cref{thm:v5-annular-resolution,cor:v5-all-shell-stock,thm:v5-annular-two-representations,thm:v5-annular-source-action}.
Item \textup{(V5.3)} is \Cref{thm:v5-parabolic-transport}.  Items
\textup{(V5.4)--(V5.5)} are
\Cref{thm:v5-transport-singular-support,thm:v5-joint-spatial-trichotomy,cor:v5-heavy-transport-short}.
Items \textup{(V5.6)--(V5.7)} are
\Cref{thm:v5-two-record-Pythagoras,cor:v5-represented-dictionary}.  The final
item is
\Cref{thm:v5-centre-chain,cor:v5-material-capture-interface}.  Each theorem
retains its failed premise as the displayed preceding branch, so the ordered
list is exhaustive after the successive subsequence choices.
\end{proof}

\begin{corollary}[First V5 packet eligible for NS--B]
\label{cor:v5-NSB-export}
Suppose a cofinal subsequence satisfies all of the following:
\begin{enumerate}[label=\textup{(E\arabic*)},leftmargin=3.5em]
\item annular concentration and bounded stress leverage;
\item the heavy joint-cell and signed local transport branches;
\item summable centre transport and the material-capture incidence;
\item a fixed held-out response margin
      \eqref{eq:v5-heldout-response};
\item the inherited represented-return and dynamically closed
      rank-at-most-two audits pass.
\end{enumerate}
Then the card
\begin{equation}
 \boxed{
 \begin{aligned}
 \mathfrak C_n^{\rm ann}=\bigl(&J_n^+,I_n^{\rm ann},K_n,\Gamma_n,x_n,
 r_n^{\rm ann},H_n^{\rm ann},\\
 &G_n,\varpi_n,q_n^{\rm ann},\phi_n^{\rm ann},
 \Xi_n,\kappa_n^{\rm tr},\Delta_n^{\rm tr},
 (I-P_n^{\rm rep})G_n\bigr)
 \end{aligned}}
 \label{eq:v5-annular-export-card}
\end{equation}
is a literal fixed-datum rank-three source--stress--response thread.  The
NS--A programme stops on this subsequence and exports the card to NS--B for
the V13 paid kinetic-stress variation and fused positive Reynolds-covariance
square-function test.
\end{corollary}

\begin{proof}
The first two hypotheses give the matched age, frequency, near-diagonal
spatial card, and signed lower margin.  The third places those cards on one
cofinal singular thread.  The fourth removes the declared response dictionary.
The fifth removes represented return and every isolated rank-at-most-two
Fourier carrier, leaving precisely a rank-three held-out fixed-datum response.
All entries in \eqref{eq:v5-annular-export-card} were constructed from the
same V4 Duhamel source and its exact stress representation.  The export is
therefore licensed without identifying the output square function
\eqref{eq:v5-output-covariance} with the V13 Reynolds covariance in advance.
\end{proof}

\begin{assessment}[What V5 acquires and where it stops]
\label{ass:v5-endpoint}
V4 supplied a future stress card but left its output frequency unrestricted.
V5 makes the first decisive split.  Frequency concentration yields one
matched annular source--stress--response card; frequency spreading yields a
positive all-shell response law which no finite band bank captures.  On the
annular branch, bounded stress leverage forces simultaneous age and spatial
locality, and the first two-record short distinguishes old-vector reuse from
aligned replenishment and orthogonal response ancestry.  Summable annular
radii then turn bounded relative centre motion into one singular thread.

What remains is no longer an unnamed compactness problem.  It is one of:
\[
 \boxed{
 \begin{gathered}
 \text{all-shell output spread or stress-leverage escape},\\
 \text{relative singular-distance escape or transport counterflow},\\
 \text{spatial microdelocalization or relative relocation},\\
 \text{represented old-tail reuse or aligned replenishment},\\
 \text{or one held-out rank-three annular thread exported to NS--B}.
 \end{gathered}}
\]
The donor programme is opened only in the final branch.  In particular, the
five-dimensional trace-free stress target and the positive output covariance
constructed here are not relabelled as the V13 paid-stress path and Reynolds
covariance before their same-history incidence is proved.
\end{assessment}

% =============================================================================
\section{Version V5 proof-status ledger and next attack}
\label{sec:v5-ledger}
% =============================================================================

\begingroup\small
\begin{longtable}{p{0.24\textwidth}p{0.19\textwidth}p{0.49\textwidth}}
\caption{NS--A Version V5 proof-status ledger.}\label{tab:v5-ledger}\\
\toprule Row & Status & Exact conclusion\\\midrule
\endfirsthead
\toprule Row & Status & Exact conclusion\\\midrule
\endhead
V1--V4 prefix & \PROVED{} preserved & Exact current V4 preterminal source retained byte-for-byte; no earlier theorem was rewritten.\\[0.25em]
Auxiliary Dini/log-spread and virtual-collapse results & \INHERITED & Used only to keep the attack after the record; no marginal-Dini theorem is repeated.\\[0.25em]
Future response annular resolution & \PROVED & Exact Parseval law \eqref{eq:v5-annular-Parseval} and concentration versus log-frequency spreading.\\[0.25em]
Finite annular head on the spread branch & \OBSTRUCTION & Every fixed number of bands carries vanishing response fraction; positive all-shell trace remains.\\[0.25em]
Annular source/stress identity & \PROVED & One scalar has exact native-source and pressure-free trace-free stress representations.\\[0.25em]
Source-visible annular action & \PROVED & $\mathcal A_n^{\rm src}\ge(H_n^{\rm ann})^2/\widehat m_n$; no shell is charged as a second source.\\[0.25em]
Annular scale chronologies & \PROVED & $\sum\mathcal R_n/\Gamma_n$ and $\sum\mathcal R_n/\Gamma_n^{3/2}$ are finite on the concentrated branch.\\[0.25em]
Annular kernel localization & \PROVED & Uniform periodic Schwartz kernel and off-diagonal $L^2$ bounds at scale $\Gamma_n^{-1}$.\\[0.25em]
Matched parabolic transport & \PROVED / \CONDITIONAL & Bounded stress leverage forces age $O((\nu\Gamma_n^2)^{-1})$ and distance $O(\Gamma_n^{-1})$; otherwise leverage escapes.\\[0.25em]
Transport singular support & \PROVED & Both source-location and output-location marginals land on $\Sigma_T$.\\[0.25em]
Joint spatial trichotomy & \PROVED & Relative separation, one heavy near-diagonal cell, or microdelocalization at annular scale.\\[0.25em]
Heavy-cell signed short & \PROVED & Local replay yields counterflow/interface survival or a signed source--output margin of order $\mathcal R_n$.\\[0.25em]
Immediate two-record ancestry & \PROVED & Exact orthogonal/aligned Pythagoras; true reuse forces overlapping annuli and a parabolic time gate modulo height-ratio escape.\\[0.25em]
Complete represented dictionary & \PROVED{} identity / \OPEN{} estimate & Held-out orthogonal response versus represented V3/return/old-tail head after literal common-endpoint transport.\\[0.25em]
Cofinal centre chain & \PROVED & Summable radii give one Cauchy singular thread or relative relocation.\\[0.25em]
Material capture & \OPEN & Required to identify response reuse with transport inside successive heavy cells.\\[0.25em]
NS--B export & \CONDITIONAL & Licensed only on the signed, captured, held-out, rank-three annular branch.\\[0.25em]
Terminal Dini improvement & \OPEN & Not inferred from annular concentration, scale sums, or leverage.\\[0.25em]
Global regularity & \OPEN & No surviving branch is unconditionally excluded.\\
\bottomrule
\end{longtable}
\endgroup

\begin{programme}[Version V6: source ancestry and the Reynolds/paid-stress incidence]
\label{prog:v5-next}
Preserve Versions V1--V5 verbatim.  Continue in the following order.
\begin{enumerate}[label=\textup{(V6.\arabic*)},leftmargin=5.2em]
\item On the annular held-out branch, pull the V13 finite-lag Reynolds covariance
      and the paid kinetic-stress path onto the exact card
      \eqref{eq:v5-annular-export-card}.  Prove their same-history incidence or
      return the precise source/router residual.
\item Apply source first-birth and final-head promotion to the complete
      dictionary in \eqref{eq:v5-represented-Pythagoras}.  Do not charge the
      aligned replenishment residual as a new source direction.
\item On immediate old-annulus reuse, retain the comparable-frequency and
      parabolic-time gates.  Test whether repeated reuse has a finite
      source-paid return stock or leaves an undamped recurrent high-frequency
      carrier.
\item On the all-shell branch, compare the positive response operator
      \eqref{eq:v5-output-covariance} with the fused Reynolds-covariance output
      law.  A failed comparison is an incidence defect, not evidence that the
      two covariances coincide.
\item On stress-leverage escape, disintegrate the kinetic-stress length in age
      and space before taking positive parts.  Seek one heavier annular card or
      return a typed high-stress/low-resultant counterflow sequence.
\item On microdelocalization, fuse the joint signed transport cells before
      variation and retain their diameter.  Test a source-relative capacity or
      Reynolds square-function upper bound without multiplying by the number
      of cells.
\item On relative relocation, test the physical material-capture map.  Either
      transport enters the summable centre chain or a literal changing-centre
      source/return packet remains.
\item Stop and export immediately when one same-history rank-three card has a
      held-out source response, paid-stress path, and positive Reynolds
      covariance on the same record/time/output carrier.
\item Invoke the inherited continuation theorem only if every branch of
      \Cref{thm:v5-integrated} is made scalar-null or the exact ridge-weighted
      terminal Dini condition is independently proved.
\end{enumerate}
No new general response graph or ambient source norm is admissible in place of
these physical acquisitions.
\end{programme}

\begin{firewall}[No global theorem in V5]
\label{fire:v5-no-global}
Version V5 proves a positive annular response resolution, exact source/stress
representations, source-visible action and scale debits, a periodic annular
kernel theorem, a leverage-or-parabolic-transport alternative, singular-set
landing of the joint transport law, the first two-record response Pythagoras,
and a summable-centre chain alternative.  It does not prove bounded stress
leverage, material capture, finite return stock, equality with the V13
Reynolds covariance, paid-stress variation, exclusion of all-shell spreading,
the terminal Dini improvement, or global three-dimensional Navier--Stokes
regularity.
\end{firewall}

% =============================================================================
\section*{Version V5 verification and history}
\addcontentsline{toc}{section}{Version V5 verification and history}
% =============================================================================

\paragraph{Append-only integrity.}
The complete current V4 preterminal byte string has SHA--256 digest
\begin{center}\scriptsize\ttfamily
ee1f420884ef1b3f725c9400a84f99720c1f0f3837b96a4dc3716c711c8c6da0
\end{center}
and is preserved exactly before the V5 material begins.  This digest
identifies the sole V4 prefix used for the continuation.  Only the final active
document terminator of that source was removed; one active terminator is
restored below.

\paragraph{Version V5 history.}
Kept the auxiliary report's lossless-Dini and log-spread no-go results frozen
and continued from the exact V4 future core.  Constructed the positive dyadic
annular response law and separated a fixed annular response from genuine
all-shell spreading.  On the concentrated branch, produced one source-first
annular test with exact native-source and pressure-free kinetic-stress
representations, proved the source-action lower bound and finite scale debits,
and evaluated the periodic annular strain kernel.  Derived the bounded-stress-
leverage implication to matched terminal age and near-diagonal source--output
transport.  Proved terminal singular-set support for both transport marginals,
the refined relative-distance/heavy-cell/microdelocalized trichotomy, and the
local coherent-to-signed counterflow short.  Transported consecutive annular
responses to one endpoint and proved the exact orthogonal-ancestry/aligned-
replenishment Pythagoras, with comparable-frequency and parabolic-time gates on
actual reuse.  Finally used summability of the annular radii to separate one
Cauchy singular thread from relative spatial relocation and stated the exact
conditional interface to the V13 fixed-datum Reynolds/paid-stress packet.  No
global regularity conclusion was claimed.

For Version V6, preserve this complete source verbatim, remove only its final
active document terminator, append new mathematical material immediately
before it, and restore the terminator.  The next filename is
\begin{center}\scriptsize
\path{Renewal_NS_Terminal_Source_Concentration_and_Singular_Thread_Closure_Programme_V6.tex}.
\end{center}

\addtocontents{toc}{\protect\endgroup}
% =============================================================================
% END VERSION V5 MATERIAL
% =============================================================================


% =============================================================================
% BEGIN VERSION V6 MATERIAL -- append-only continuation of NS--A V1--V5
% =============================================================================
\clearpage
\hypersetup{pdftitle={NS-A Terminal Source Concentration Programme: Cumulative V6},pdfauthor={A. Pelissier}}
\makeatletter
\addtocontents{toc}{\protect\begingroup\protect\makeatletter}
\addtocontents{toc}{\protect\def\protect\l@section{\protect\@dottedtocline{1}{0em}{3.5em}}}
\addtocontents{toc}{\protect\makeatother}
\makeatother
\renewcommand{\thesection}{V6.\arabic{section}}
\renewcommand{\thesubsection}{V6.\arabic{section}.\arabic{subsection}}
\renewcommand{\theequation}{V6.\arabic{equation}}
\setcounter{section}{0}
\setcounter{equation}{0}
\renewcommand{\theHsection}{V6.\arabic{section}}
\renewcommand{\theHsubsection}{V6.\arabic{section}.\arabic{subsection}}
\renewcommand{\theHtheorem}{V6.\arabic{section}.\arabic{theorem}}
\renewcommand{\theHequation}{V6.\arabic{equation}}

\begin{center}
{\LARGE\bfseries NS--A: Version V6}\\[0.5em]
{\large The Heat--Reynolds Incidence of the Held-Out Annular Card,\\
Dynamic Final-Head Reflection, and the Active-Future Alternative}\\[0.5em]
{\normalsize 2 September 2026}
\end{center}
\addcontentsline{toc}{part}{Version V6: heat--Reynolds incidence and reflected final-head transport}

\begin{abstract}
Versions V1--V5 are retained verbatim.  The newly supplied fixed-datum
Reynolds programme, shared cofinal-survivor programme, and complete terminal
session report change the correct order of the next attack.  The fixed-datum
programme has already reduced its logarithmic-age branch to an exact octave
transport defect and its compact-age branch to a local covariance--writer
Riesz short with an explicit collar.  The shared programme expressly forbids
importing a positive Reynolds operator into Navier--Stokes before one
same-history bridge ties it to the corrected terminal chronology.  The
session report proves that every further one-sided renewal or restart
estimate remains abort blind; its first genuinely two-sided quantity is the
reflected record pairing, for which viscosity cancels exactly.  V6 therefore
does not repeat the donor table, the octave compiler, or another Dini
majorant.  It constructs the missing bridge on the exact V5 annular card and
then applies the reflected pairing after the literal represented-head short.

First, the heat-smoothed native source in the V5 annular scalar is decomposed
exactly as
\[
 e^{-\nu sA}A^{-1/4}\mathcal B(u)
 =A^{-1/4}\mathbb P\operatorname{div}\mathbb C_s[u]
  +A^{-1/4}\mathbb P\operatorname{div}(v_s\otimes v_s),
\]
where $v_s=e^{-\nu sA}u$ and
$\mathbb C_s[u]=P_{\nu s}(u\otimes u)-v_s\otimes v_s\succeq0$ is the
literal heat Reynolds covariance.  Hence the matched-age annular payment is
the sum of one positive-covariance source--writer scalar and one resolved
Euler re-entry scalar.  One of them carries a fixed fraction of the V5 record
margin.  On the covariance branch an exact history-space Pythagoras gives a
source action of order $\nu\mathcal R_n^2$ on the matched parabolic window.
This is the first positive quadratic carrier proved to be incident with the
same endpoint, age, and output annulus as the NS--A failure; it is not inferred
from a marginal covariance law.

The final-head promotion is then performed before this source is priced.
Applying the literal V5 represented projection to the annular Duhamel
response and extracting its own terminal age core yields a held-out vector of
fixed record size.  Its Riesz writer again splits exactly into Reynolds and
resolved parts.  Either the logarithmic source-age depth escapes, the resolved
part survives, or the held-out covariance action is bounded below by
$\nu\mathcal R_n^2$ divided by that explicit depth.  Localizing this card by a
quadratic spatial partition produces one exact covariance--strain term plus
the inherited commutator collar.  Failure to place it in the V5 singular cell
is retained as the common-history material-capture residual required by the
fixed-datum programme.

Finally, V6 compresses the represented dictionary inside a sharp annular
spectral head and forms the reflected held-out pairing across
$b_n=\beta_{n+1}$.  A static head need not be dynamically reducing: the exact
derivative contains the commutator $\nu[A,R_n]$ in addition to the past and
future nonlinear source works.  If the band-router and endpoint-incidence
residuals pass, the reflected pairing begins with mass comparable with
$\mathcal R_n^2$.  At one fixed parabolic lag there are four exhaustive
outcomes: the next record gap is too short, dynamic-head motion survives, a
past/future reflected nonlinear work has record size, or the pairing remains
large.  On the last branch, a small semigroup-head defect forces an active
future held-out Duhamel response of size comparable with $\mathcal R_n$.
That future response again splits into a positive heat-Reynolds source and a
resolved source, and its native source mass pays $\mathcal R_n/\Gamma_n$ on
the actual post-record interval.  Thus a passive future is not silently used:
reflected retention either produces a named dynamic defect or a literal
same-annulus replenishment after the record.

The paid radial-stress path is not identified with this covariance carrier.
Its own same-history action density remains a separate finite acquisition, as
the supplied NS--B note requires.  No Liouville theorem, terminal Dini
improvement, exclusion of the all-shell branch, or global three-dimensional
regularity theorem is claimed.
\end{abstract}

% =============================================================================
\section{Scope after the new attachment audit}
\label{sec:v6-scope}
% =============================================================================

The V5 endpoint asked for the V13 Reynolds and paid-stress objects to be
pulled onto the annular export card.  The newly supplied material determines
which part of that request is mathematically licensed.

\begin{enumerate}[label=\textup{(N6.\arabic*)},leftmargin=5.0em]
\item The fixed-datum programme already proves the common-lag heat covariance
      $\mathbb C_s[u]\succeq0$, its source
      $A^{-1/4}\mathbb P\operatorname{div}\mathbb C_s[u]$, the local Riesz
      short, and the complete localization collar.  It also proves an exact
      adjacent-age alternative on the diffuse front.  These rows are frozen
      and are not copied here.
\item The shared cofinal-survivor programme explicitly leaves Navier--Stokes
      measure valued until NS--A and NS--B construct one literal positive
      quadratic carrier on the corrected terminal chronology.  A covariance
      existing on a different record is not enough.
\item The complete session report proves that one-sided restart, lifetime, and
      returned-memory inequalities do not exclude aborted cascades.  The
      reflected record pairing is the first available two-sided quantity and
      has an exact cancellation of viscosity.
\item Accordingly, V6 performs two card-specific acquisitions only: the
      heat--Reynolds factorization of the V5 annular source, and the reflected
      dynamic-head test of that same annular chronology.
\end{enumerate}

\begin{firewall}[Rows not reopened in V6]
\label{fire:v6-no-reopening}
No theorem below re-proves the terminal Dini no-go, the log-spread
counterexample, the NS--B octave recurrence, the five-dimensional donor table,
the radial-selector obstruction, the V5 annular kernel estimate, the V5
spatial trichotomy, or the abstract reflected-record identity.  The latter is
specialized below only because the represented projection and its dynamic
commutator are new data.  An all-shell response law, stress-leverage escape,
relative relocation, and microdelocalization remain exact V5 alternatives
unless their own hypotheses are explicitly invoked.
\end{firewall}

Retain a cofinal V5 annular-concentration subsequence.  Thus
\begin{equation}
 b_n:=\beta_{n+1},\qquad
 s_n(t):=b_n-t,\qquad
 H_n^{\rm ann}\ge h_*\mathcal R_n,
 \label{eq:v6-standing-annular}
\end{equation}
with $\varpi_n$, $q_n^{\rm ann}$, $I_n^{\rm core}$, $\Gamma_n$, and
$P_n^{\rm rep}$ as in V5.  On the bounded-stress-leverage branch retain the
matched interval
\begin{equation}
 I_n^{\rm ann}=(b_n-\tau_n^{\rm ann},b_n),
 \qquad
 \tau_n^{\rm ann}\le\frac{R_t}{\nu\Gamma_n^2}
 \label{eq:v6-matched-interval}
\end{equation}
from \eqref{eq:v5-matched-age}.  All scalar products are real unless a real
part is displayed.

% =============================================================================
\section{The exact heat--Reynolds factorization of the V5 annular scalar}
\label{sec:v6-Reynolds-factorization}
% =============================================================================

For $s\ge0$, let
\begin{equation}
 E_s:=e^{-\nu sA}
 \label{eq:v6-vector-heat}
\end{equation}
act on divergence-free vectors and let $P_{\nu s}=e^{\nu s\Delta}$ act
componentwise on tensors.  For $t\in I_n^{\rm core}$ define
\begin{align}
 v_{n,t}&:=E_{s_n(t)}u(t),
 \label{eq:v6-resolved-field}\\
 \mathbb C_{n,t}&:=P_{\nu s_n(t)}(u\otimes u)(t)
                 -v_{n,t}\otimes v_{n,t},
 \label{eq:v6-heat-covariance}\\
 g_{n,t}^{\rm cov}&:=A^{-1/4}\mathbb P\operatorname{div}
                    \mathbb C_{n,t}^{\circ},
 \label{eq:v6-cov-source}\\
 g_{n,t}^{\rm res}&:=A^{-1/4}\mathbb P\operatorname{div}
                    (v_{n,t}\otimes v_{n,t})^{\circ}.
 \label{eq:v6-res-source}
\end{align}
The trace-free signs may be removed without changing either projected
divergence.

\begin{proposition}[Positive heat covariance and exact source factorization]
\label{prop:v6-heat-Reynolds-factorization}
For every $n$ and every $t\in I_n^{\rm core}$,
\begin{equation}
 \boxed{\mathbb C_{n,t}(x)\succeq0\quad\text{for every }x\in\Torus^3,}
 \label{eq:v6-covariance-positive}
\end{equation}
and
\begin{equation}
 \boxed{E_{s_n(t)}c(t)=g_{n,t}^{\rm cov}+g_{n,t}^{\rm res}.}
 \label{eq:v6-source-factorization}
\end{equation}
If
\begin{equation}
 Z_n:=\Lambda\varpi_n,
 \qquad
 M_n:=S(A^{1/4}\varpi_n),
 \label{eq:v6-annular-fixed-writers}
\end{equation}
then
\begin{align}
 -\langle g_{n,t}^{\rm cov},Z_n\rangle
 &=\int_{\Torus^3}\mathbb C_{n,t}^{\circ}:M_n\,\dd x,
 \label{eq:v6-covariance-tensor-pairing}\\
 -\langle g_{n,t}^{\rm res},Z_n\rangle
 &=\int_{\Torus^3}(v_{n,t}\otimes v_{n,t})^{\circ}:M_n\,\dd x.
 \label{eq:v6-resolved-tensor-pairing}
\end{align}
Thus the V5 native source scalar has a same-history Reynolds/resolved split,
not two independently declared source occurrences.
\end{proposition}

\begin{proof}
For $a\in\R^3$, positivity of the heat kernel and Jensen's inequality give
\[
 a^{\mathsf T}\mathbb C_{n,t}(x)a
 =P_{\nu s_n(t)}[(a\cdot u(t))^2](x)
  -\bigl(P_{\nu s_n(t)}[a\cdot u(t)](x)\bigr)^2\ge0.
\]
This proves \eqref{eq:v6-covariance-positive}.  All operators in
$A^{-1/4}\mathbb P\operatorname{div}$ are Fourier multipliers, so they
commute with the componentwise heat semigroup.  Consequently
\[
 \begin{aligned}
 E_sc
 &=A^{-1/4}\mathbb P\operatorname{div}
       P_{\nu s}(u\otimes u)\\
 &=A^{-1/4}\mathbb P\operatorname{div}\mathbb C_s[u]
   +A^{-1/4}\mathbb P\operatorname{div}(v_s\otimes v_s),
 \end{aligned}
\]
which is \eqref{eq:v6-source-factorization}; isotropic tensors contribute
only gradients and are killed by $\mathbb P$.

Let $T=A^{-1/4}\mathbb P\operatorname{div}$ on symmetric tensors.  Its
adjoint on divergence-free vectors is
$T^*z=-S(A^{-1/4}z)$.  Since
$A^{-1/4}Z_n=A^{1/4}\varpi_n$, one has
\[
 -\langle T\mathbb C_{n,t}^{\circ},Z_n\rangle
 =\langle\mathbb C_{n,t}^{\circ},S(A^{1/4}\varpi_n)\rangle.
\]
The same calculation with $v_{n,t}\otimes v_{n,t}$ proves the two tensor
formulas.
\end{proof}

Define the time-restricted annular scalar
\begin{equation}
 \mathfrak H_n^{\rm age}
 :=-\int_{I_n^{\rm ann}}
       \langle c(t),q_n^{\rm ann}(t)\rangle\,\dd t
 \label{eq:v6-age-scalar}
\end{equation}
and its two components
\begin{align}
 \mathfrak H_n^{\rm cov}
 &:=-\int_{I_n^{\rm ann}}
       \langle g_{n,t}^{\rm cov},Z_n\rangle\,\dd t,
 \label{eq:v6-age-cov-scalar}\\
 \mathfrak H_n^{\rm res}
 &:=-\int_{I_n^{\rm ann}}
       \langle g_{n,t}^{\rm res},Z_n\rangle\,\dd t.
 \label{eq:v6-age-res-scalar}
\end{align}

\begin{theorem}[Matched-age covariance or resolved-Euler incidence]
\label{thm:v6-matched-Reynolds-incidence}
On the bounded-stress-leverage branch of V5,
\begin{equation}
 \boxed{
 \mathfrak H_n^{\rm age}
 =\mathfrak H_n^{\rm cov}+\mathfrak H_n^{\rm res}
 \ge\frac34H_n^{\rm ann}.}
 \label{eq:v6-age-split-lower}
\end{equation}
Consequently, after the fixed priority convention ``covariance first'', at
least one of the following holds:
\begin{enumerate}[label=\textup{(R6.\arabic*)},leftmargin=5.0em]
\item \emph{Positive-covariance source--writer incidence:}
      \begin{equation}
       \boxed{\mathfrak H_n^{\rm cov}
       \ge\frac38H_n^{\rm ann}
       \ge\frac{3h_*}{8}\mathcal R_n.}
       \label{eq:v6-cov-record-margin}
      \end{equation}
\item \emph{Resolved-Euler re-entry:}
      \begin{equation}
       \boxed{\mathfrak H_n^{\rm res}
       \ge\frac38H_n^{\rm ann}
       \ge\frac{3h_*}{8}\mathcal R_n.}
       \label{eq:v6-res-record-margin}
      \end{equation}
\end{enumerate}
The first branch is a literal positive quadratic carrier on the V5
record/age/output card.  The second is not relabelled as Reynolds work.
\end{theorem}

\begin{proof}
The first V5 representation and self-adjointness of $E_s$ give, for every
subinterval $I$ of the core,
\[
 -\int_I\langle c,\Lambda E_{s_n(t)}\varpi_n\rangle\,\dd t
 =-\int_I\langle E_{s_n(t)}c,Z_n\rangle\,\dd t.
\]
Insert \eqref{eq:v6-source-factorization} to obtain the exact equality in
\eqref{eq:v6-age-split-lower}.

In the proof of the V5 matched-age theorem, $R_t$ was chosen so that the
absolute contribution from
$s_n(t)\ge R_t/(\nu\Gamma_n^2)$ is at most
$H_n^{\rm ann}/4$.  If the full core is shorter, that contribution is empty.
Therefore removing the old-age part from the positive total scalar
$H_n^{\rm ann}$ leaves
$\mathfrak H_n^{\rm age}\ge3H_n^{\rm ann}/4$.  If both components were
strictly smaller than $3H_n^{\rm ann}/8$, their sum would be smaller than
$3H_n^{\rm ann}/4$.  This proves the alternative and the record-scale lower
bounds.
\end{proof}

\begin{theorem}[Exact Riesz short and record-scale covariance action]
\label{thm:v6-matched-Riesz-action}
For $\bullet\in\{\mathrm{cov},\mathrm{res}\}$ put
\begin{equation}
 \mathcal A_n^{\bullet}
 :=\int_{I_n^{\rm ann}}\|g_{n,t}^{\bullet}\|_2^2\,\dd t,
 \qquad
 \mathcal D_n^{\rm ann}
 :=\tau_n^{\rm ann}\|Z_n\|_2^2.
 \label{eq:v6-matched-actions}
\end{equation}
Whenever $\mathfrak H_n^{\bullet}>0$,
\begin{equation}
 \boxed{
 \mathcal A_n^{\bullet}
 =\frac{(\mathfrak H_n^{\bullet})^2}{\mathcal D_n^{\rm ann}}
 +\int_{I_n^{\rm ann}}
 \left\|g_{n,t}^{\bullet}
 +\frac{\mathfrak H_n^{\bullet}}{\mathcal D_n^{\rm ann}}Z_n
 \right\|_2^2\,\dd t.}
 \label{eq:v6-matched-Riesz-Pythagoras}
\end{equation}
Moreover,
\begin{equation}
 \boxed{\mathcal D_n^{\rm ann}\le\frac{4R_t}{\nu}.}
 \label{eq:v6-writer-stock-bound}
\end{equation}
On either record-margin branch of
\Cref{thm:v6-matched-Reynolds-incidence},
\begin{equation}
 \boxed{
 \mathcal A_n^{\bullet}
 \ge\frac{9h_*^2}{256R_t}\,\nu\mathcal R_n^2.}
 \label{eq:v6-record-action-lower}
\end{equation}
For $\bullet=\mathrm{cov}$ this is an action of the actual Reynolds source
$A^{-1/4}\mathbb P\operatorname{div}\mathbb C_{n,t}$; it is not an action
assigned to each covariance tensor component separately.
\end{theorem}

\begin{proof}
In the Hilbert space
$L^2(I_n^{\rm ann};L^2_\sigma)$, regard $Z_n$ as the constant-in-time
writer.  Since
$\mathfrak H_n^{\bullet}=-\int\langle g_{n,t}^{\bullet},Z_n\rangle\,\dd t$,
expanding the final square in
\eqref{eq:v6-matched-Riesz-Pythagoras} proves the identity.  The spectral
support of $\varpi_n$ gives $\|Z_n\|_2\le2\Gamma_n$, while
\eqref{eq:v6-matched-interval} gives
\[
 \mathcal D_n^{\rm ann}
 \le\frac{R_t}{\nu\Gamma_n^2}\,4\Gamma_n^2
 =\frac{4R_t}{\nu}.
\]
Drop the nonnegative residual term and use either
\eqref{eq:v6-cov-record-margin} or
\eqref{eq:v6-res-record-margin}.
\end{proof}

\begin{assessment}[The shared cofinal programme now has its missing NS carrier]
\label{ass:v6-positive-carrier}
On branch \eqref{eq:v6-cov-record-margin}, the packet
\begin{equation}
 \boxed{
 \mathfrak Q_n^{\rm cov}
 =\bigl(b_n,I_n^{\rm ann},\Gamma_n,\varpi_n,
        \mathbb C_{n,t},g_{n,t}^{\rm cov},Z_n,M_n,
        \mathfrak H_n^{\rm cov},\mathcal A_n^{\rm cov}\bigr)}
 \label{eq:v6-positive-covariance-card}
\end{equation}
is a literal positive quadratic carrier tied to the corrected NS--A
chronology.  This closes the \emph{time/age/output} incidence which the shared
cofinal-survivor note kept open.  It does not yet identify the covariance card
with the V5 selected spatial cell, nor does it identify the covariance action
with the NS--B radial paid-stress action.  Those are later mixed rows.
\end{assessment}

% =============================================================================
\section{Final-head promotion before the covariance source is priced}
\label{sec:v6-heldout-source}
% =============================================================================

Assume the V5 held-out branch
\begin{equation}
 d_n:=(I-P_n^{\rm rep})G_n,
 \qquad
 D_n:=\|d_n\|_2\ge\theta H_n^{\rm ann}
 \label{eq:v6-static-heldout}
\end{equation}
for one fixed $\theta>0$.  Put
\begin{equation}
 B_n:=(I-P_n^{\rm rep})\Psi_n\mathsf P_{>K_n}.
 \label{eq:v6-heldout-synthesis}
\end{equation}
Then $d_n=B_nF_n^{\rm core}$ and $\|B_n\|\le1$.  The order in
\eqref{eq:v6-heldout-synthesis} is physical: the response is first restricted
to the V5 annulus and only then shorted against the common-endpoint
dictionary.

\begin{lemma}[Annular envelope for the static held-out synthesis]
\label{lem:v6-heldout-envelope}
For every $s\ge0$,
\begin{equation}
 \boxed{
 \|B_n\Lambda e^{-\nu sA}\|_{2\to2}
 \le2\Gamma_ne^{-\nu s\Gamma_n^2/4}.}
 \label{eq:v6-heldout-envelope}
\end{equation}
Consequently
\begin{equation}
 D_n\le2\Gamma_n\widehat m_n.
 \label{eq:v6-heldout-source-upper}
\end{equation}
\end{lemma}

\begin{proof}
The left factors in $B_n$ are contractions.  The multiplier $\Psi_n$ is
supported where $\Gamma_n/2\le\Lambda\le2\Gamma_n$.  On this set
$\Lambda e^{-\nu sA}\le2\Gamma_ne^{-\nu s\Gamma_n^2/4}$ in operator norm.
Integrate the resulting estimate over $I_n^{\rm core}$.
\end{proof}

Define the dimensionless held-out source-age depth
\begin{equation}
 \boxed{
 L_n^{\rm ho}:=\log\frac{4\Gamma_n\widehat m_n}{D_n}\ge\log2,}
 \label{eq:v6-heldout-depth}
\end{equation}
where the inequality is \eqref{eq:v6-heldout-source-upper}, and put
\begin{equation}
 \sigma_n^{\rm ho}
 :=\min\left\{\widehat\sigma_n,
       \frac{4L_n^{\rm ho}}{\nu\Gamma_n^2}\right\},
 \qquad
 I_n^{\rm ho}:=(b_n-\sigma_n^{\rm ho},b_n].
 \label{eq:v6-heldout-age-core}
\end{equation}
Let
\begin{equation}
 d_n^{\rm ho}:=-\int_{I_n^{\rm ho}}
 B_n\Lambda e^{-\nu(b_n-t)A}c(t)\,\dd t,
 \qquad
 H_n^{\rm ho}:=\|d_n^{\rm ho}\|_2.
 \label{eq:v6-heldout-vector-core}
\end{equation}

\begin{theorem}[Terminal core of the final-head residual]
\label{thm:v6-heldout-core}
The held-out terminal core satisfies
\begin{equation}
 \boxed{
 H_n^{\rm ho}\ge\frac12D_n
 \ge\frac{\theta}{2}H_n^{\rm ann}
 \ge\frac{\theta h_*}{2}\mathcal R_n.}
 \label{eq:v6-heldout-core-lower}
\end{equation}
It remains orthogonal to the V5 represented dictionary:
\begin{equation}
 \boxed{P_n^{\rm rep}d_n^{\rm ho}=0.}
 \label{eq:v6-heldout-core-orthogonal}
\end{equation}
Thus final-head promotion has been performed before any Reynolds or resolved
source action is attached to the vector.
\end{theorem}

\begin{proof}
If $\sigma_n^{\rm ho}=\widehat\sigma_n$, then
$d_n^{\rm ho}=d_n$.  Otherwise the omitted part has norm at most
\[
 2\Gamma_ne^{-\nu\sigma_n^{\rm ho}\Gamma_n^2/4}\widehat m_n
 =2\Gamma_ne^{-L_n^{\rm ho}}\widehat m_n
 =\frac{D_n}{2}
\]
by \eqref{eq:v6-heldout-depth}.  The reverse triangle inequality gives the
first bound.  The remaining lower bounds are
\eqref{eq:v6-static-heldout} and
\eqref{eq:v5-annular-record-margin}.  Every integrand in
\eqref{eq:v6-heldout-vector-core} lies in
$\Ran(I-P_n^{\rm rep})$, proving
\eqref{eq:v6-heldout-core-orthogonal}.
\end{proof}

Set
\begin{equation}
 \omega_n^{\rm ho}:=\frac{d_n^{\rm ho}}{H_n^{\rm ho}},
 \qquad
 a_n^{\rm ho}:=B_n^*\omega_n^{\rm ho},
 \qquad
 Z_n^{\rm ho}:=\Lambda a_n^{\rm ho}.
 \label{eq:v6-heldout-writer}
\end{equation}
Then $\|a_n^{\rm ho}\|_2\le1$, it is supported in
$(K_n,\infty)\cap[\Gamma_n/2,2\Gamma_n]$, and
$\|Z_n^{\rm ho}\|_2\le2\Gamma_n$.  Define on $I_n^{\rm ho}$
\begin{align}
 \mathfrak H_n^{{\rm ho},{\rm cov}}
 &:=-\int\langle g_{n,t}^{\rm cov},Z_n^{\rm ho}\rangle\,\dd t,
 \label{eq:v6-heldout-cov-scalar}\\
 \mathfrak H_n^{{\rm ho},{\rm res}}
 &:=-\int\langle g_{n,t}^{\rm res},Z_n^{\rm ho}\rangle\,\dd t.
 \label{eq:v6-heldout-res-scalar}
\end{align}

\begin{theorem}[Held-out Reynolds/resolved source alternative]
\label{thm:v6-heldout-Reynolds-alternative}
The final-head core has the exact same-history representations
\begin{equation}
 \boxed{
 \begin{aligned}
 H_n^{\rm ho}
 &=-\int_{I_n^{\rm ho}}
   \langle c(t),\Lambda e^{-\nu(b_n-t)A}a_n^{\rm ho}\rangle\,\dd t\\
 &=\mathfrak H_n^{{\rm ho},{\rm cov}}
   +\mathfrak H_n^{{\rm ho},{\rm res}}.
 \end{aligned}}
 \label{eq:v6-heldout-source-split}
\end{equation}
At least one of the two real scalars is at least $H_n^{\rm ho}/2$.  On a
component $\bullet\in\{\mathrm{cov},\mathrm{res}\}$ satisfying that lower
bound, put
\begin{equation}
 \mathcal A_n^{{\rm ho},\bullet}
 :=\int_{I_n^{\rm ho}}\|g_{n,t}^{\bullet}\|_2^2\,\dd t,
 \qquad
 \mathcal D_n^{\rm ho}:=\sigma_n^{\rm ho}\|Z_n^{\rm ho}\|_2^2.
 \label{eq:v6-heldout-component-actions}
\end{equation}
Then
\begin{equation}
 \boxed{
 \mathcal A_n^{{\rm ho},\bullet}
 =\frac{(\mathfrak H_n^{{\rm ho},\bullet})^2}
        {\mathcal D_n^{\rm ho}}
 +\int_{I_n^{\rm ho}}
 \left\|g_{n,t}^{\bullet}
 +\frac{\mathfrak H_n^{{\rm ho},\bullet}}
        {\mathcal D_n^{\rm ho}}Z_n^{\rm ho}\right\|_2^2\,\dd t,}
 \label{eq:v6-heldout-Riesz-Pythagoras}
\end{equation}
with
\begin{equation}
 \boxed{
 \mathcal D_n^{\rm ho}\le\frac{16L_n^{\rm ho}}{\nu},
 \qquad
 \mathcal A_n^{{\rm ho},\bullet}
 \ge\frac{\theta^2h_*^2}{256L_n^{\rm ho}}
       \nu\mathcal R_n^2.}
 \label{eq:v6-heldout-action-lower}
\end{equation}
Hence a cofinal held-out branch has the exact alternative
\begin{equation}
 \boxed{
 L_n^{\rm ho}\longrightarrow\infty
 \quad\big|\quad
 \text{record-scale resolved action}
 \quad\big|\quad
 \text{record-scale positive-covariance action}.}
 \label{eq:v6-heldout-three-way}
\end{equation}
The first line is a source-age-depth escape, not a new frequency count.
\end{theorem}

\begin{proof}
Taking the scalar product of
\eqref{eq:v6-heldout-vector-core} with
$\omega_n^{\rm ho}$ and moving $B_n$ to the other side gives the first line
of \eqref{eq:v6-heldout-source-split}.  Move the heat semigroup from the
writer to $c$ and use \eqref{eq:v6-source-factorization} for the second line.
Since the sum is the positive number $H_n^{\rm ho}$, one component is at
least half of it.

The Pythagoras identity is the same one-dimensional Hilbert-space projection
calculation as in \Cref{thm:v6-matched-Riesz-action}.  By
\eqref{eq:v6-heldout-age-core} and the annular support of
$a_n^{\rm ho}$,
\[
 \mathcal D_n^{\rm ho}
 \le\frac{4L_n^{\rm ho}}{\nu\Gamma_n^2}\,4\Gamma_n^2
 =\frac{16L_n^{\rm ho}}{\nu}.
\]
Use
$\mathfrak H_n^{{\rm ho},\bullet}\ge H_n^{\rm ho}/2$ and
\eqref{eq:v6-heldout-core-lower}; dropping the nonnegative residual yields
\eqref{eq:v6-heldout-action-lower}.  Along a subsequence the depth is either
bounded or tends to infinity.  On the bounded branch choose the component by
the stated priority convention.
\end{proof}

\begin{remark}[Tensor form of the held-out covariance scalar]
\label{rem:v6-heldout-tensor}
On the covariance branch,
\begin{equation}
 \boxed{
 \mathfrak H_n^{{\rm ho},{\rm cov}}
 =\int_{I_n^{\rm ho}}\int_{\Torus^3}
   \mathbb C_{n,t}^{\circ}:
   S(A^{1/4}a_n^{\rm ho})\,\dd x\,\dd t.}
 \label{eq:v6-heldout-cov-tensor}
\end{equation}
Thus the positive tensor, its projected source, the annular writer, and the
held-out response are all deterministic functions of the same original
solution and the same V5 represented dictionary.  This is the required
source/time/output incidence.  It does not yet supply a spatial-cell Read or
the paid radial-stress path.
\end{remark}

% =============================================================================
\section{The spatial common-history row and the inherited covariance collar}
\label{sec:v6-spatial-incidence}
% =============================================================================

Let $\{\chi_{n,\alpha}\}_\alpha$ be the V5 quadratic spatial partition at
scale comparable with $\Gamma_n^{-1}$, transported on the same physical
record, so that
\begin{equation}
 \sum_\alpha\chi_{n,\alpha}^2=1.
 \label{eq:v6-quadratic-partition}
\end{equation}
For a covariance-dominant held-out card define
\begin{equation}
 h_{n,\alpha}^{\rm cov}
 :=-\int_{I_n^{\rm ho}}
 \langle\chi_{n,\alpha}g_{n,t}^{\rm cov},
        \chi_{n,\alpha}Z_n^{\rm ho}\rangle\,\dd t.
 \label{eq:v6-local-cov-source-scalar}
\end{equation}
Then
\begin{equation}
 \boxed{\sum_\alpha h_{n,\alpha}^{\rm cov}
 =\mathfrak H_n^{{\rm ho},{\rm cov}}.}
 \label{eq:v6-local-cov-sum}
\end{equation}
This is an exact source-side localization after the common-lag covariance has
been formed.

\begin{proposition}[Exact local covariance--strain and collar identity]
\label{prop:v6-local-covariance-collar}
Let $T=A^{-1/4}\mathbb P\operatorname{div}$ and
\begin{equation}
 M_n^{\rm ho}:=S(A^{1/4}a_n^{\rm ho}).
 \label{eq:v6-heldout-strain-writer}
\end{equation}
For every cell,
\begin{equation}
 \boxed{
 \begin{aligned}
 h_{n,\alpha}^{\rm cov}
 ={}&\int_{I_n^{\rm ho}}\int
       \chi_{n,\alpha}^2\mathbb C_{n,t}^{\circ}:
       M_n^{\rm ho}\,\dd x\,\dd t\\
 &-\int_{I_n^{\rm ho}}
   \langle\mathbb C_{n,t}^{\circ},
          [T^*,\chi_{n,\alpha}^2]Z_n^{\rm ho}\rangle\,\dd t.
 \end{aligned}}
 \label{eq:v6-local-collar-identity}
\end{equation}
The second line is the complete localization collar.  It is not a second
Reynolds occurrence and may not be dropped when the first line is used as a
local paid stress.
\end{proposition}

\begin{proof}
By definition and self-adjoint multiplication,
\[
 h_{n,\alpha}^{\rm cov}
 =-\int\langle T\mathbb C_{n,t}^{\circ},
                  \chi_{n,\alpha}^2Z_n^{\rm ho}\rangle\,\dd t
 =-\int\langle\mathbb C_{n,t}^{\circ},
                  T^*(\chi_{n,\alpha}^2Z_n^{\rm ho})\rangle\,\dd t.
\]
Insert
$T^*(\chi^2Z)=\chi^2T^*Z+[T^*,\chi^2]Z$ and
$-T^*Z=S(A^{-1/4}Z)=M_n^{\rm ho}$.
\end{proof}

For a covariance-dominant card let
\begin{equation}
 P_n^{\rm cov,+}:=\sum_\alpha(h_{n,\alpha}^{\rm cov})_+,
 \qquad
 \pi_n^{\rm cov}(\alpha)
 :=\frac{(h_{n,\alpha}^{\rm cov})_+}{P_n^{\rm cov,+}}.
 \label{eq:v6-cov-cell-law}
\end{equation}
Since the signed sum is positive,
$P_n^{\rm cov,+}\ge\mathfrak H_n^{{\rm ho},{\rm cov}}>0$.
Let $\alpha_n^{\rm V5}$ denote the already selected V5 singular cell when
that branch is present.

\begin{corollary}[Exact material-capture/collar interface to NS--B]
\label{cor:v6-covariance-cell-interface}
After passage to a subsequence, one of the following holds.
\begin{enumerate}[label=\textup{(C6.\arabic*)},leftmargin=5.0em]
\item \emph{Covariance-cell capture fails:}
      \begin{equation}
       \boxed{\pi_n^{\rm cov}(\alpha_n^{\rm V5})\longrightarrow0.}
       \label{eq:v6-cov-cell-capture-fails}
      \end{equation}
      The positive covariance action and the V5 stress replay have compatible
      scales but do not occupy the same selected spatial record.
\item \emph{A covariance cell is captured:} there is $p>0$ with
      $\pi_n^{\rm cov}(\alpha_n^{\rm V5})\ge p$.  Put
      \begin{equation}
       \Delta_n^{\rm cov,rep}
       :=\left|
       \frac{\mathfrak H_n^{{\rm ho},{\rm cov}}}
            {P_n^{\rm cov,+}}
       (h_{n,\alpha_n^{\rm V5}}^{\rm cov})_+
       -h_{n,\alpha_n^{\rm V5}}^{\rm cov}\right|.
       \label{eq:v6-cov-replay-defect}
      \end{equation}
      Then either
      \begin{equation}
       \limsup_n
       \frac{\Delta_n^{\rm cov,rep}}
            {\mathfrak H_n^{{\rm ho},{\rm cov}}}>0,
       \label{eq:v6-cov-counterflow}
      \end{equation}
      or the signed local source scalar is at least
      \begin{equation}
       \boxed{
       h_{n,\alpha_n^{\rm V5}}^{\rm cov}
       \ge\frac p2\mathfrak H_n^{{\rm ho},{\rm cov}}.}
       \label{eq:v6-local-cov-margin}
      \end{equation}
      In the latter case, \eqref{eq:v6-local-collar-identity} gives either a
      local covariance--strain margin of the same order or a collar of that
      order.
\end{enumerate}
Thus the time/age/output covariance incidence is unconditional on its branch;
the spatial identification demanded by NS--B is exactly capture plus the
signed replay and collar tests above.
\end{corollary}

\begin{proof}
Along a subsequence the nonnegative selected-cell probabilities either tend
to zero or have a positive lower bound.  Work on the second branch.  The
coherent replay of the positive cell law assigns the selected cell mass at
least $p\mathfrak H_n^{{\rm ho},{\rm cov}}$.  If the normalized discrepancy
in \eqref{eq:v6-cov-replay-defect} tends to zero, the signed cell value is at
least half that amount for late $n$.  Finally, in
\eqref{eq:v6-local-collar-identity} a signed sum of the physical local tensor
pairing and the negative collar has the displayed lower bound; one of their
absolute values is therefore at least half of it.
\end{proof}

\begin{firewall}[Paid stress is not manufactured by covariance incidence]
\label{fire:v6-paid-stress-separate}
The supplied NS--B V6 source distinguishes the positive common-lag Reynolds
covariance from the radial paid-stress path and from its selector-motion
collar.  Equations \eqref{eq:v6-positive-covariance-card} and
\eqref{eq:v6-local-collar-identity} populate the covariance source and its
local writer on the NS--A chronology.  They do not prove that the NS--B radial
path, its BV action, or its Fisher-height density occurs on this same cell.
A direct mixed action density, or its singular Radon--Nikodym residual, remains
the exact paid-stress incidence row.
\end{firewall}

% =============================================================================
\section{Band-compatible final heads and their reflected dynamics}
\label{sec:v6-reflected-head}
% =============================================================================

The reflected pairing in the session report is exact for a spectral band.
The V5 represented dictionary need not itself be spectrally reducing.  We
therefore separate static head fitting from dynamic head invariance.

Let
\begin{equation}
 Q_n:=\one_{[\Gamma_n/4,4\Gamma_n]}(\Lambda),
 \label{eq:v6-sharp-annular-head}
\end{equation}
so that $Q_nG_n=G_n$.  Let $\mathcal J_n^{\rm band}$ be the span of
$Q_nv$ with $v\in\mathcal J_n^{\rm rep}$, let $P_n^{\rm band}$ be its
orthogonal projection, and put
\begin{equation}
 R_n:=Q_n-P_n^{\rm band}.
 \label{eq:v6-dynamic-heldout-projection}
\end{equation}
Since $\Ran P_n^{\rm band}\subseteq\Ran Q_n$, $R_n$ is an orthogonal
projection and $R_n\le Q_n$.

Define the band-router and endpoint-incidence residuals
\begin{align}
 \delta_n^{\rm band}
 &:=\|(P_n^{\rm band}-P_n^{\rm rep})G_n\|_2,
 \label{eq:v6-band-router-residual}\\
 \delta_n^{\rm end}
 &:=\|R_n(e(b_n)-G_n)\|_2.
 \label{eq:v6-endpoint-incidence-residual}
\end{align}

\begin{lemma}[Static held-out response to a band-compatible endpoint state]
\label{lem:v6-band-endpoint-lower}
Under \eqref{eq:v6-static-heldout},
\begin{equation}
 \boxed{
 \|R_nG_n\|_2
 \ge\theta H_n^{\rm ann}-\delta_n^{\rm band}.}
 \label{eq:v6-band-heldout-lower}
\end{equation}
If
\begin{equation}
 \delta_n^{\rm band}\le\frac\theta2H_n^{\rm ann},
 \qquad
 \delta_n^{\rm end}\le\frac\theta4H_n^{\rm ann},
 \label{eq:v6-router-endpoint-pass}
\end{equation}
then, with
\begin{equation}
 X_n:=\|R_ne(b_n)\|_2,
 \qquad c_0:=\frac{\theta h_*}{4},
 \label{eq:v6-X-c0}
\end{equation}
one has
\begin{equation}
 \boxed{X_n\ge c_0\mathcal R_n.}
 \label{eq:v6-endpoint-heldout-lower}
\end{equation}
\end{lemma}

\begin{proof}
Since $G_n=Q_nG_n$,
\[
 R_nG_n=(I-P_n^{\rm band})G_n
 =(I-P_n^{\rm rep})G_n
  +(P_n^{\rm rep}-P_n^{\rm band})G_n.
\]
The reverse triangle inequality proves
\eqref{eq:v6-band-heldout-lower}.  Also
$R_ne(b_n)=R_nG_n+R_n(e(b_n)-G_n)$.  Apply the same inequality, then use
\eqref{eq:v5-annular-record-margin}.
\end{proof}

For an admissible lag
\begin{equation}
 0\le\tau<\min\{b_n,T-b_n\},
 \label{eq:v6-admissible-reflection-lag}
\end{equation}
write $e_n^{\pm}(\tau)=e(b_n\pm\tau)$ and
$c_n^{\pm}(\tau)=c(b_n\pm\tau)$.  Define the reflected dynamic-head pairing
\begin{equation}
 \Theta_n(\tau)
 :=\langle R_ne_n^-(\tau),R_ne_n^+(\tau)\rangle.
 \label{eq:v6-reflected-heldout-pairing}
\end{equation}

\begin{theorem}[Exact reflected identity after final-head promotion]
\label{thm:v6-reflected-final-head}
For every admissible lag,
\begin{equation}
 \boxed{
 \begin{aligned}
 \Theta_n'(\tau)
 ={}&\nu\langle e_n^-(\tau),[A,R_n]e_n^+(\tau)\rangle\\
 &+\langle R_n\Lambda c_n^-(\tau),R_ne_n^+(\tau)\rangle
  -\langle R_ne_n^-(\tau),R_n\Lambda c_n^+(\tau)\rangle.
 \end{aligned}}
 \label{eq:v6-reflected-head-identity}
\end{equation}
Moreover,
\begin{equation}
 \Theta_n(0)=X_n^2,
 \qquad \Theta_n'(0)=0.
 \label{eq:v6-reflected-initial-data}
\end{equation}
The viscosity term vanishes identically for every pair of states in the
occupied band when the held-out head reduces $A$.  Conversely, vanishing of
the bilinear commutator form on the whole occupied band is equivalent to that
reducing property.  Along one trajectory accidental cancellation can occur;
otherwise the displayed commutator is the dynamic-head promotion residual and
is not a nonlinear source birth.
\end{theorem}

\begin{proof}
The critical variable satisfies
$\partial_te=-\nu Ae-\Lambda c$.  Therefore
\[
 \partial_\tau e_n^-=\nu Ae_n^-+\Lambda c_n^-,
 \qquad
 \partial_\tau e_n^+=-\nu Ae_n^+-\Lambda c_n^+.
\]
Differentiate \eqref{eq:v6-reflected-heldout-pairing}.  The two viscous terms
are
\[
 \nu\langle Ae_n^-,R_ne_n^+\rangle
 -\nu\langle R_ne_n^-,R_nAe_n^+\rangle
 =\nu\langle e_n^-,(AR_n-R_nA)e_n^+\rangle,
\]
using $R_n^2=R_n$ and self-adjointness.  This gives
\eqref{eq:v6-reflected-head-identity}.  The value at zero is immediate.  At
zero the two source terms are equal, and
$\langle e,[A,R_n]e\rangle=0$ because the commutator of two self-adjoint
operators is skew symmetric on the real Hilbert space.  Hence the derivative
vanishes.
\end{proof}

\begin{lemma}[Exact semigroup defect of a nonreducing final head]
\label{lem:v6-semigroup-head-defect}
Let $E_\tau=e^{-\nu\tau A}$.  Then
\begin{equation}
 \boxed{
 [R_n,E_\tau]
 =\nu\int_0^\tau
 E_{\tau-s}[A,R_n]E_s\,\dd s.}
 \label{eq:v6-semigroup-commutator}
\end{equation}
In particular, with
\begin{equation}
 \mathfrak s_n(\tau):=\|[R_n,E_\tau]e(b_n)\|_2,
 \label{eq:v6-semigroup-defect}
\end{equation}
one has
\begin{equation}
 \boxed{
 \|R_nE_\tau e(b_n)\|_2
 \le e^{-\nu\tau\Gamma_n^2/16}X_n
      +\mathfrak s_n(\tau).}
 \label{eq:v6-old-head-decay}
\end{equation}
\end{lemma}

\begin{proof}
Differentiate
$E_{\tau-s}R_nE_s$ in $s$ and integrate from $0$ to $\tau$:
\[
 R_nE_\tau-E_\tau R_n
 =\nu\int_0^\tau E_{\tau-s}[A,R_n]E_s\,\dd s.
\]
Since $R_ne(b_n)\in\Ran Q_n$ and every spectral value on this range is at
least $\Gamma_n/4$,
\[
 \|E_\tau R_ne(b_n)\|_2
 \le e^{-\nu\tau\Gamma_n^2/16}X_n.
\]
Add the commutator defect.
\end{proof}

% =============================================================================
\section{Reflection loss, dynamic-head survival, or active future replenishment}
\label{sec:v6-reflected-alternative}
% =============================================================================

Fix once and for all $0<\eta<1/2$ and, under
\eqref{eq:v6-router-endpoint-pass}, put
\begin{equation}
 a_0:=\frac{(1-\eta)c_0^2}{2},
 \qquad
 L_0:=16\log\frac8{a_0},
 \qquad
 \tau_n^*:=\frac{L_0}{\nu\Gamma_n^2}.
 \label{eq:v6-reflection-constants}
\end{equation}
The constants depend only on the fixed V5 annular and held-out margins and on
$\eta$.

For $0<\tau\le\beta_{n+2}-b_n$, define the absolute reflected carriers
\begin{align}
 \mathcal K_n(\tau)
 &:=\nu\int_0^\tau
 \left|\langle e_n^-(s),[A,R_n]e_n^+(s)\rangle\right|\,\dd s,
 \label{eq:v6-reflected-K}\\
 \mathcal W_n^-(\tau)
 &:=\int_0^\tau
 \left|\langle R_n\Lambda c_n^-(s),R_ne_n^+(s)\rangle\right|\,\dd s,
 \label{eq:v6-reflected-past-work}\\
 \mathcal W_n^+(\tau)
 &:=\int_0^\tau
 \left|\langle R_ne_n^-(s),R_n\Lambda c_n^+(s)\rangle\right|\,\dd s.
 \label{eq:v6-reflected-future-work}
\end{align}

\begin{theorem}[Two-sided reflected final-head alternative]
\label{thm:v6-two-sided-alternative}
Assume the V5 annular and static held-out branches together with
\eqref{eq:v6-router-endpoint-pass}.  After passage to a subsequence, at least
one of the following holds.
\begin{enumerate}[label=\textup{(T6.\arabic*)},leftmargin=5.0em]
\item \emph{Parabolic future-gap compression:}
      \begin{equation}
       \boxed{
       \nu\Gamma_n^2(\beta_{n+2}-b_n)<L_0.}
       \label{eq:v6-future-gap-compression}
      \end{equation}
\item \emph{Semigroup/dynamic-head survival:}
      \begin{equation}
       \boxed{
       \mathfrak s_n(\tau_n^*)>\frac{a_0}{4}\mathcal R_n.}
       \label{eq:v6-semigroup-head-survival}
      \end{equation}
\item \emph{Reflected loss with a physical carrier:}
      \begin{equation}
       \boxed{
       \Theta_n(\tau_n^*)\le(1-\eta)X_n^2,}
       \label{eq:v6-reflection-loss}
      \end{equation}
      and
      \begin{equation}
       \boxed{
       \mathcal K_n(\tau_n^*)
       +\mathcal W_n^-(\tau_n^*)
       +\mathcal W_n^+(\tau_n^*)
       \ge\eta c_0^2\mathcal R_n^2.}
       \label{eq:v6-reflected-carrier-lower}
      \end{equation}
      Hence at least one dynamic-head, past-source, or future-source carrier
      has size at least $\eta c_0^2\mathcal R_n^2/3$.
\item \emph{Reflected retention and active future replenishment:}
      \begin{equation}
       \boxed{
       \Theta_n(\tau_n^*)>(1-\eta)X_n^2,}
       \label{eq:v6-reflection-retention}
      \end{equation}
      and the post-record held-out response
      \begin{equation}
       Y_n^+:=-\int_0^{\tau_n^*}
       R_n\Lambda e^{-\nu(\tau_n^*-s)A}c(b_n+s)\,\dd s
       \label{eq:v6-active-future-response}
      \end{equation}
      satisfies
      \begin{equation}
       \boxed{
       \|Y_n^+\|_2\ge\frac{a_0}{2}\mathcal R_n,
       \qquad
       m_n^+:=\int_{b_n}^{b_n+\tau_n^*}\|c(t)\|_2\,\dd t
       \ge\frac{a_0}{8}\frac{\mathcal R_n}{\Gamma_n}.}
       \label{eq:v6-active-future-lower}
      \end{equation}
\end{enumerate}
The four alternatives use the actual solution on both sides of $b_n$.  In
particular, the last branch proves that a retained reflected head cannot be a
passive continuation once the dynamic-head defect is small and a full
parabolic future lag is available.
\end{theorem}

\begin{proof}
Pass first to the branch on which the future gap contains $\tau_n^*$; otherwise
item \textup{(T6.1)} holds.  If
\eqref{eq:v6-semigroup-head-survival} holds, take item \textup{(T6.2)}.
Assume its failure.

If \eqref{eq:v6-reflection-loss} holds, integrate
\eqref{eq:v6-reflected-head-identity} from $0$ to $\tau_n^*$ and use
\eqref{eq:v6-reflected-initial-data}.  The absolute values of the three
integrated terms dominate
$X_n^2-\Theta_n(\tau_n^*)\ge\eta X_n^2$.  The endpoint lower bound
\eqref{eq:v6-endpoint-heldout-lower} proves
\eqref{eq:v6-reflected-carrier-lower}.

It remains to treat \eqref{eq:v6-reflection-retention}.  Since
$b_n=\beta_{n+1}$ is the first time at which the critical norm equals
$\mathcal R_{n+1}=2\mathcal R_n$, every earlier time has critical norm at
most $2\mathcal R_n$.  Cauchy--Schwarz therefore gives
\[
 \|R_ne(b_n+\tau_n^*)\|_2
 \ge\frac{\Theta_n(\tau_n^*)}
          {\|R_ne(b_n-\tau_n^*)\|_2}
 >\frac{(1-\eta)c_0^2\mathcal R_n^2}{2\mathcal R_n}
 =a_0\mathcal R_n.
\]
Duhamel's formula after $b_n$ is
\[
 R_ne(b_n+\tau_n^*)=R_nE_{\tau_n^*}e(b_n)+Y_n^+.
\]
By \Cref{lem:v6-semigroup-head-defect}, the failure of
\eqref{eq:v6-semigroup-head-survival}, $X_n\le2\mathcal R_n$, and the
choice of $L_0$ give
\[
 \begin{aligned}
 \|R_nE_{\tau_n^*}e(b_n)\|_2
 &\le2e^{-L_0/16}\mathcal R_n+\frac{a_0}{4}\mathcal R_n\\
 &\le\frac{a_0}{2}\mathcal R_n.
 \end{aligned}
\]
The reverse triangle inequality yields the response lower bound.  Since
$R_n\le Q_n$ and the spectrum of $Q_n$ lies below $4\Gamma_n$,
\[
 \|Y_n^+\|_2
 \le4\Gamma_n\int_{b_n}^{b_n+\tau_n^*}\|c(t)\|_2\,\dd t.
\]
Rearrange to obtain the source-mass lower bound.
\end{proof}

\begin{corollary}[Cofinal source debit of the retained reflected branch]
\label{cor:v6-reflected-source-debit}
On every subsequence satisfying item \textup{(T6.4)},
\begin{equation}
 \boxed{
 \sum_n\frac{\mathcal R_n}{\Gamma_n}
 \le\frac8{a_0}\int_0^T\|c(t)\|_2\,\dd t<\infty.}
 \label{eq:v6-reflected-scale-sum}
\end{equation}
The intervals used in the sum are contained in the disjoint record intervals
$(\beta_{n+1},\beta_{n+2}]$.  This is one payment by the native source stock,
not a second charge in addition to the V5 scale constraint.
\end{corollary}

\begin{proof}
Sum the second inequality of \eqref{eq:v6-active-future-lower}.  The post-record
intervals are pairwise disjoint and lie in $(0,T)$.
\end{proof}

The active future response itself has the same Reynolds/resolved factorization.
Put
\begin{equation}
 y_n^+:=\frac{Y_n^+}{\|Y_n^+\|_2}.
 \label{eq:v6-future-direction}
\end{equation}
Since $Y_n^+\in\Ran R_n$, one has $R_ny_n^+=y_n^+$ and
$\supp y_n^+\subset[\Gamma_n/4,4\Gamma_n]$.
For $0\le s\le\tau_n^*$ set
\begin{align}
 \mathbb C_{n,s}^+
 &:=P_{\nu(\tau_n^*-s)}(u\otimes u)(b_n+s)
   -v_{n,s}^+\otimes v_{n,s}^+,
 \label{eq:v6-future-covariance}\\
 v_{n,s}^+
 &:=e^{-\nu(\tau_n^*-s)A}u(b_n+s),
 \label{eq:v6-future-resolved-field}
\end{align}
and define $g_{n,s}^{+,\rm cov}$ and $g_{n,s}^{+,\rm res}$ by the analogues
of \eqref{eq:v6-cov-source}--\eqref{eq:v6-res-source}.

\begin{corollary}[The active future is covariance work or resolved re-entry]
\label{cor:v6-future-Reynolds-split}
On item \textup{(T6.4)},
\begin{equation}
 \boxed{
 \begin{aligned}
 \|Y_n^+\|_2
 ={}&-\int_0^{\tau_n^*}
   \langle g_{n,s}^{+,\rm cov},\Lambda y_n^+\rangle\,\dd s\\
 &-\int_0^{\tau_n^*}
   \langle g_{n,s}^{+,\rm res},\Lambda y_n^+\rangle\,\dd s.
 \end{aligned}}
 \label{eq:v6-future-Reynolds-split}
\end{equation}
At least one term on the right is at least $\|Y_n^+\|_2/2$.  On the
covariance branch the positive tensor $\mathbb C_{n,s}^+\succeq0$ has source
action
\begin{equation}
 \boxed{
 \int_0^{\tau_n^*}\|g_{n,s}^{+,\rm cov}\|_2^2\,\dd s
 \ge\frac{a_0^2}{256L_0}\,\nu\mathcal R_n^2.}
 \label{eq:v6-future-cov-action}
\end{equation}
The same lower bound holds for the resolved source action on its branch.
\end{corollary}

\begin{proof}
Take the scalar product of \eqref{eq:v6-active-future-response} with
$y_n^+$.  The adjoint of
$R_n\Lambda e^{-\nu(\tau_n^*-s)A}$ is
$e^{-\nu(\tau_n^*-s)A}\Lambda R_n$.  Since $R_ny_n^+=y_n^+$, moving the
heat semigroup to $c(b_n+s)$ gives
\eqref{eq:v6-future-Reynolds-split} by
\Cref{prop:v6-heat-Reynolds-factorization}.  One real component is at least
half the positive sum.

The constant writer has norm
$\|\Lambda y_n^+\|_2\le4\Gamma_n$ and duration
$\tau_n^*=L_0/(\nu\Gamma_n^2)$, so its squared spacetime stock is at most
$16L_0/\nu$.  The one-dimensional Riesz Pythagoras therefore gives
\[
 \int\|g\|_2^2
 \ge\frac{(\|Y_n^+\|_2/2)^2}{16L_0/\nu}.
\]
Use \eqref{eq:v6-active-future-lower}.
\end{proof}

\begin{assessment}[What the reflected theorem adds to the restart ledger]
\label{ass:v6-reflection-versus-restart}
The restart-cost theorem controls how soon the next critical record can be
reached from a given $L^\infty$ state, but its fixed-fraction cost is
summable and cannot rule out infinitely many relaunches.  Theorem
\ref{thm:v6-two-sided-alternative} has a different content.  It compares the
same represented held-out band on both sides of an actual record.  Failure of
retention returns an antisymmetric two-sided source work or dynamic-head
motion; retention over one parabolic lag forces a concrete post-record
held-out response and then a positive-covariance/resolved split.  This is the
first V1--V6 statement whose positive branch uses the solution after the
record and whose target has already survived the static final-head short.
\end{assessment}

% =============================================================================
\section{Source first birth, future succession, and the exact stopping line}
\label{sec:v6-source-first-birth}
% =============================================================================

The future response $Y_n^+$ belongs to $\Ran R_n$, hence is orthogonal to the
band-compressed represented dictionary at the old endpoint.  It can still be
represented at its own endpoint by a later return or promoted head.  Let
$\mathcal J_{n,+}^{\rm rep}$ be the finite literal dictionary at
$b_n+\tau_n^*$ after all actual transports to that endpoint, and let
$P_{n,+}^{\rm rep}$ be its orthogonal projection.

\begin{proposition}[Exact successor promotion of the active future response]
\label{prop:v6-future-promotion}
Every response from item \textup{(T6.4)} satisfies
\begin{equation}
 \boxed{
 \|Y_n^+\|_2^2
 =\|P_{n,+}^{\rm rep}Y_n^+\|_2^2
  +\|(I-P_{n,+}^{\rm rep})Y_n^+\|_2^2.}
 \label{eq:v6-future-promotion}
\end{equation}
Therefore either
\begin{equation}
 \|(I-P_{n,+}^{\rm rep})Y_n^+\|_2
 \ge\vartheta_+\|Y_n^+\|_2
 \label{eq:v6-future-heldout-successor}
\end{equation}
for one fixed $\vartheta_+>0$ on a subsequence, producing a new held-out
successor of size $\gtrsim\mathcal R_n$, or the future response is represented
by the declared dictionary and receives no new source-rank price.  The latter
case is recurrent replenishment, even though the native source mass in
\eqref{eq:v6-active-future-lower} is paid on a new time interval.
\end{proposition}

\begin{proof}
This is orthogonal Pythagoras at the actual future endpoint.  Along a
subsequence the relative held-out norm either has a positive lower bound or
tends to zero.  The size statement follows from
\eqref{eq:v6-active-future-lower}.  Time-interval freshness does not alter the
source-rank interpretation after the endpoint dictionary has been assembled.
\end{proof}

\begin{definition}[V6 two-sided covariance survivor card]
\label{def:v6-two-sided-card}
A V5 annular branch is called a \emph{V6 two-sided covariance survivor} when:
\begin{enumerate}[label=\textup{(S\arabic*)},leftmargin=3.5em]
\item the V5 static held-out and rank-three audits pass;
\item the band-router and endpoint-incidence residuals in
      \eqref{eq:v6-router-endpoint-pass} pass;
\item the pre-record held-out factorization
      \eqref{eq:v6-heldout-source-split} is covariance dominant;
      an active-future covariance branch from
      \eqref{eq:v6-future-Reynolds-split} remains a successor card until its
      own spatial capture and collar have been tested;
\item the common-history spatial capture, signed replay, and collar tests of
      \Cref{cor:v6-covariance-cell-interface} pass on one cofinal singular
      thread;
\item the paid-stress mixed action density is supplied independently, or its
      absence is retained as the terminal NS--B incidence residual.
\end{enumerate}
The card consists of the V5 entries together with
\begin{equation}
 \boxed{
 \bigl(B_n,d_n^{\rm ho},a_n^{\rm ho},
       \mathbb C_{n,t},g_{n,t}^{\rm cov},
       R_n,\Theta_n,[A,R_n]\bigr),}
 \label{eq:v6-two-sided-card}
\end{equation}
and, on item~\textup{(T6.4)}, also the successor entries
\(Y_n^+\) and \(\mathbb C_{n,s}^+\).  These future entries do not enter the
export until their own common-history spatial incidence has passed.
\end{definition}

\begin{theorem}[Version V6 heat--Reynolds and reflected-survivor theorem]
\label{thm:v6-integrated}
Retain the ordered V5 alternatives.  On every annular-concentrated,
bounded-leverage, statically held-out branch, V6 proves the following
exhaustive refinement after subsequence selection.
\begin{enumerate}[label=\textup{(V6.\arabic*)},leftmargin=5.2em]
\item The matched-age V5 scalar is either a positive heat-Reynolds
      covariance--writer scalar or a resolved-Euler re-entry scalar of fixed
      record size.  On either branch the corresponding source action is at
      least $c\nu\mathcal R_n^2$.
\item After the literal V5 final-head projection, the held-out response has a
      terminal vector core of fixed record size.  It has source-age depth
      $L_n^{\rm ho}$ and an exact covariance/resolved Riesz Pythagoras.  Either
      $L_n^{\rm ho}\to\infty$, resolved re-entry survives, or the positive
      covariance action is at least
      $c\nu\mathcal R_n^2/L_n^{\rm ho}$.
\item The covariance branch is localized by one exact quadratic partition.
      The only difference between the local source pairing and the physical
      covariance--strain is the complete commutator collar
      \eqref{eq:v6-local-collar-identity}.  Failure to occupy the V5 singular
      cell is the explicit material-capture residual
      \eqref{eq:v6-cov-cell-capture-fails}.
\item Compressing the represented dictionary to a sharp annular head returns
      either the band-router residual, the endpoint response/state incidence
      residual, or a noncollapsed held-out endpoint state of size
      $\gtrsim\mathcal R_n$.
\item On that state, the reflected pairing obeys the exact dynamic-head
      identity \eqref{eq:v6-reflected-head-identity}.  At the fixed parabolic
      lag, it returns future-gap compression, semigroup-head motion, a
      record-sized reflected dynamic/past/future carrier, or reflected
      retention.
\item Reflected retention with a small dynamic-head defect forces the active
      future response \eqref{eq:v6-active-future-response}, whose norm is
      $\gtrsim\mathcal R_n$ and whose native source mass is
      $\gtrsim\mathcal R_n/\Gamma_n$.  That response is again either positive
      heat-Reynolds covariance work or resolved re-entry, with source action
      $\gtrsim\nu\mathcal R_n^2$.
\item The future response is shorted once more at its own endpoint.  Its
      represented part is recurrent replenishment of the first source; only
      the final held-out component may enlarge the source atlas.
\item A pre-record covariance branch satisfying all spatial, ancestry, and
      independently supplied paid-stress incidence rows is the V6 two-sided
      survivor of \Cref{def:v6-two-sided-card} and is exported to NS--B.  An
      active-future covariance successor is exported only after its analogous
      spatial capture and collar pass.  If the paid-stress mixed row is absent,
      that absence is the exact stopping obstruction.
\end{enumerate}
No item makes the positive output operator
$\mathbb C_n^{\rm out}$ equal to the heat Reynolds covariance, and no item
excludes the V5 all-shell, stress-leverage, microdelocalized, or relative-
relocation branches without their own hypotheses.
\end{theorem}

\begin{proof}
Item \textup{(V6.1)} is
\Cref{thm:v6-matched-Reynolds-incidence,thm:v6-matched-Riesz-action}.
Item \textup{(V6.2)} is
\Cref{thm:v6-heldout-core,thm:v6-heldout-Reynolds-alternative}.
Item \textup{(V6.3)} is
\Cref{prop:v6-local-covariance-collar,cor:v6-covariance-cell-interface}.
Item \textup{(V6.4)} is \Cref{lem:v6-band-endpoint-lower}.
Item \textup{(V6.5)} is
\Cref{thm:v6-reflected-final-head,lem:v6-semigroup-head-defect,thm:v6-two-sided-alternative}.
Item \textup{(V6.6)} is
\Cref{thm:v6-two-sided-alternative,cor:v6-future-Reynolds-split}; its source
stock is \Cref{cor:v6-reflected-source-debit}.  Item \textup{(V6.7)} is
\Cref{prop:v6-future-promotion}.  The final item is the definition of the
card together with the paid-stress firewall.  Every alternative omitted from
the hypotheses is retained from \Cref{thm:v5-integrated}.
\end{proof}

\begin{firewall}[No reflected Liouville theorem in V6]
\label{fire:v6-no-Liouville}
The two-sided identity supplies a physical test functional and identifies the
exact branches on which it can fail to retain a represented held-out head.
It does not prove a sign for the reflected nonlinear work, monotonicity of
$\Theta_n$, a positive lower bound for all lags, compactness of the reflected
cards, or a Liouville theorem for the bounded ancient mild or strong-coupling
Euler profiles isolated in the session report.  Those are genuine PDE inputs.
A passive-future estimate would merely reproduce the one-sided Duhamel bound
and is not used as closure.
\end{firewall}

% =============================================================================
\section{Version V6 proof-status ledger and V7 attack order}
\label{sec:v6-ledger}
% =============================================================================

\begingroup\small
\begin{longtable}{p{0.24\textwidth}p{0.19\textwidth}p{0.49\textwidth}}
\caption{NS--A Version V6 proof-status ledger.}\label{tab:v6-ledger}\\
\toprule Row & Status & Exact conclusion\\\midrule
\endfirsthead
\toprule Row & Status & Exact conclusion\\\midrule
\endhead
V1--V5 prefix & \PROVED{} preserved & Complete V5 preterminal source retained byte-for-byte; no prior theorem was rewritten.\\[0.25em]
New NS--B octave/front calculation & \INHERITED{} / not recopied & Its age transport and local Riesz/collar rows determine the interface; NS--A supplies the missing terminal chronology and spatial card.\\[0.25em]
Shared cofinal NS carrier & \PROVED{} on covariance branch & The positive heat covariance, its projected source, annular writer, record margin, and source action occur on the same V5 time/age/output card.\\[0.25em]
Matched Reynolds/resolved split & \PROVED{} exact & The V5 scalar is covariance work plus resolved-Euler re-entry; one carries at least $3H_n^{\rm ann}/8$.\\[0.25em]
Matched covariance source action & \PROVED{} lower packet & The Riesz action is at least $c\nu\mathcal R_n^2$ on the bounded-leverage parabolic window.\\[0.25em]
Final-head terminal core & \PROVED{} exact & The V5 held-out vector has a later-age core of at least half its norm and remains orthogonal to the represented dictionary.\\[0.25em]
Held-out source-age depth & \PROVED{} alternative & Bounded $L_n^{\rm ho}$ gives covariance or resolved action $\gtrsim\nu\mathcal R_n^2$; otherwise the explicit depth escapes.\\[0.25em]
Spatial covariance incidence & \PROVED{} identity / \OPEN{} capture & Quadratic localization is exact and the only tensor conversion term is the commutator collar; capture by the V5 cell is a separate mixed row.\\[0.25em]
Paid radial-stress incidence & \OPEN{} finite mixed row & Positive covariance does not identify the NS--B radial path, BV action, selector height, or localized collar density.\\[0.25em]
Band-compatible final head & \PROVED{} exact & Static held-out response gives a noncollapsed spectral held-out state unless the band-router or endpoint-incidence residual survives.\\[0.25em]
Reflected promoted-head identity & \PROVED{} exact & Viscosity cancels modulo the dynamic-head commutator $[A,R_n]$; past and future nonlinear works remain signed.\\[0.25em]
Two-sided parabolic alternative & \PROVED{} exact & Future-gap compression, dynamic-head survival, reflected source-work loss, or retained reflection with active future response.\\[0.25em]
Active-future source debit & \PROVED{} exact & Retention forces response $\gtrsim\mathcal R_n$, mass $\gtrsim\mathcal R_n/\Gamma_n$, and a future covariance/resolved action $\gtrsim\nu\mathcal R_n^2$.\\[0.25em]
Future final-head promotion & \PROVED{} identity / \OPEN{} estimate & Represented recurrent replenishment is charged to the first source; only the future held-out residual may be new.\\[0.25em]
All-shell output comparison & \OPEN{} incidence & $\mathbb C_n^{\rm out}$ remains an output response operator and is not equated with the heat Reynolds covariance.\\[0.25em]
Reflected Liouville / OS record positivity & \OPEN{} PDE row & No sign, monotonicity, compactness, or ancient-profile rigidity is proved.\\[0.25em]
Global regularity & \OPEN & No surviving terminal branch is unconditionally excluded.\\
\bottomrule
\end{longtable}
\endgroup

\begin{programme}[Version V7: reflected spatial landing or an explicit Liouville input]
\label{prog:v6-next}
Preserve Versions V1--V6 verbatim.  Continue only on one branch already
selected above.
\begin{enumerate}[label=\textup{(V7.\arabic*)},leftmargin=5.2em]
\item On a covariance-dominant V6 card, apply the already proved NS--B local
      Riesz/collar theorem to the exact cell law
      \eqref{eq:v6-local-cov-source-scalar}.  Do not repeat its octave or
      radial-selector calculations.
\item Test the common-history material capture of the covariance card and the
      V5 heavy singular transport cell.  Either the two cell laws couple at
      scale $\Gamma_n^{-1}$, or return the literal capture/relocation measure.
\item On the reflected-loss branch, localize the particular surviving carrier
      in \eqref{eq:v6-reflected-carrier-lower} by the same quadratic partition.
      Keep dynamic-head, past-source, and future-source carriers separate.
\item On reflected retention, spatially localize the selected future
      covariance source using the quadratic partition and the exact
      covariance--strain/collar identity.  Then short $Y_n^+$ against the
      next actual endpoint dictionary and iterate only if the response is
      represented.  Test whether recurrent reuse has a finite source-paid
      return stock or yields a cofinal undamped high-frequency reflected
      carrier.
\item If compactness of one spatially captured reflected card is available,
      pass its exact covariance, source ancestry, and $\Theta$ pairing to the
      corresponding ancient limit.  The target is a sign/rigidity theorem for
      that specified limit, not an unnamed ancient profile.
\item On the all-shell branch, no finite annular head may be inserted.  A
      comparison with NS--B is licensed only through one populated joint
      response/covariance Gram or a direct common-history Read.
\item Stop and export as soon as one card carries, on the same spatial and
      temporal record, the held-out response, positive heat covariance,
      local collar, and independently occurring paid-stress path.  If the
      next missing statement is reflected Liouville, terminal trace,
      backward uniqueness, epsilon regularity, or paid-stress BV, name it and
      stop rather than introducing another compiler.
\item Invoke the inherited critical continuation theorem only if every V5--V6
      scalar branch is null or the exact ridge-weighted Dini condition is
      independently proved.
\end{enumerate}
No new one-sided restart estimate, unweighted time marginal, shell count,
source norm, or generic response graph is admissible.
\end{programme}

\begin{firewall}[No global theorem in V6]
\label{fire:v6-no-global}
Version V6 constructs the exact heat--Reynolds incidence of the V5 annular
card, performs final-head promotion before source action, gives the local
covariance/collar identity, and derives the dynamic reflected-head and active-
future alternatives.  It does not control the paid radial path, prove
material capture, bound the source-age depth, eliminate the resolved branch,
exclude all-shell spreading, establish OS positivity or reflected Liouville,
prove the terminal Dini condition, or establish global three-dimensional
Navier--Stokes regularity.
\end{firewall}

% =============================================================================
\section*{Version V6 verification and history}
\addcontentsline{toc}{section}{Version V6 verification and history}
% =============================================================================

\paragraph{Append-only integrity.}
The complete V5 preterminal byte string has SHA--256 digest
\begin{center}\scriptsize\ttfamily
3ae63560de22e25f674465ecdd56fedac89152a5565f619acc5c15e34e573d39
\end{center}
and is preserved exactly before the V6 material begins.  Only the final
active document terminator was removed; one active terminator is restored
below.

\paragraph{Version V6 history.}
Audited the supplied NS--B V6, shared cofinal-survivor V3, and complete NS
session report.  Froze the NS--B octave/front and radial-selector calculations,
accepted the shared programme's prohibition on an analogy-based Reynolds
import, and moved the attack to the session report's first two-sided
quantity.  Factored the V5 heat-smoothed annular source exactly into the
positive heat Reynolds covariance source and resolved Euler re-entry, proved
the matched-age record margin and Riesz action, and thereby populated the
first literal positive quadratic carrier on the corrected NS--A chronology.
Performed the V5 final-head short before pricing the source, extracted a
held-out terminal age core, and proved its covariance/resolved action
alternative with an explicit logarithmic source-age depth.  Localized the
covariance source by a quadratic partition and retained the exact commutator
collar and common-history capture residual.  Compressed the represented
head inside a sharp annulus, derived the reflected promoted-head identity and
semigroup commutator, and proved the four-way parabolic alternative: short
future gap, dynamic-head survival, reflected source-work loss, or retained
reflection.  On the retained branch, forced an active post-record held-out
response, its native source debit, its Reynolds/resolved split, and its own
endpoint promotion.  Kept the paid radial-stress path, spatial capture,
all-shell comparison, and reflected Liouville theorem as independent open
rows.  No global regularity conclusion was claimed.

For Version V7, preserve this complete source verbatim, remove only the final
active document terminator, append new mathematical material immediately
before it, and restore the terminator.  The next filename is
\begin{center}\scriptsize
\path{Renewal_NS_Terminal_Source_Concentration_and_Singular_Thread_Closure_Programme_V7.tex}.
\end{center}

\addtocontents{toc}{\protect\endgroup}
% =============================================================================
% END VERSION V6 MATERIAL
% =============================================================================


% =============================================================================
% BEGIN VERSION V7 MATERIAL -- complete V1--V6 preterminal source preserved
% =============================================================================
\clearpage
\hypersetup{pdftitle={NS-A Terminal Source Concentration Programme: Cumulative V7},pdfauthor={A. Pelissier}}
\makeatletter
\addtocontents{toc}{\protect\begingroup\protect\makeatletter}
\addtocontents{toc}{\protect\def\protect\l@section{\protect\@dottedtocline{1}{0em}{3.5em}}}
\addtocontents{toc}{\protect\def\protect\l@subsection{\protect\@dottedtocline{2}{1.5em}{4.2em}}}
\addtocontents{toc}{\protect\makeatother}
\makeatother
\renewcommand{\thesection}{V7.\arabic{section}}
\renewcommand{\thesubsection}{V7.\arabic{section}.\arabic{subsection}}
\renewcommand{\theequation}{V7.\arabic{equation}}
\setcounter{section}{0}
\setcounter{equation}{0}
\renewcommand{\theHsection}{V7.\arabic{section}}
\renewcommand{\theHsubsection}{V7.\arabic{section}.\arabic{subsection}}
\renewcommand{\theHtheorem}{V7.\arabic{section}.\arabic{theorem}}
\renewcommand{\theHequation}{V7.\arabic{equation}}

\begin{center}
{\LARGE\bfseries NS--A: Version V7}\\[0.5em]
{\large An Upstream Audit of Reflected Record Positivity:\\
Smooth Navier--Stokes Counterexamples and Transport-Corrected Critical Geometry}\\[0.5em]
{\normalsize 2 September 2026}
\end{center}
\addcontentsline{toc}{part}{Version V7: reflected-record counterexamples and the transport correction}

\begin{abstract}
The V6 reflected successor alternatives are exact, but iterating them does not
supply a nonsummable physical budget.  Before extending that iteration, this
version tests the proposed sign mechanism on actual Navier--Stokes solutions.
The test fails.  For every fixed viscosity, every prescribed finite number of
doubling records, and every positive strain tolerance, one smooth mean-zero
periodic solution has a negative reflected critical pairing at each of those
first-passage records, at a lag shorter than the next record gap, while the
integrated velocity gradient on each reflected window is below the tolerance.
The explicit solutions are globally smooth two-dimensional/three-component
shear flows.  Arbitrarily small genuinely rank-three perturbations preserve
the conclusion on the complete finite record interval.

The proof uses a one-dimensional linear parabolic equation, a uniform scaled
Sobolev approximation, and an explicit square-root multiplier estimate.  It
also disproves the general cross-time strain estimate used in the supplied
session report.  The mechanism is rapid advection of a short-wave phase, not
large deformation or a singularity.  The actual material pullback removes
this loss on the same examples.  A periodic fractional-energy calculation
then proves that material transport changes the critical metric by at most
an exponential of the integrated strain, while it need not be close to the
identity as an operator on states.  These results retire unrestricted
Eulerian record positivity as a closing premise.  They do not disprove any
V6 identity, a terminal-only estimate with additional hypotheses, or
three-dimensional regularity.  The new direction is to evaluate the
transport-corrected signed critical work, not to repeat the successor
classification.
\end{abstract}

% =============================================================================
\section{The first upstream issue is the proposed reflected sign}
\label{sec:v7-upstream}
% =============================================================================

The consolidation mandate asks for one decisive physical acquisition rather
than another general compiler.  The original NS--A objective remains terminal
source concentration and exclusion of its source-visible singular survivor.
NS--B's paid stress and Reynolds square function remain a separate interface.
The exact V6 heat--Reynolds decomposition and the derivative formula
\eqref{eq:v6-reflected-head-identity} are retained without re-proof.

The step audited here is different: can regular record times, the
Navier--Stokes equation, energy-neutral quadratic forcing, and positivity of
the heat covariance force a positive Eulerian pairing across a record?  The
session report asks whether, for one fixed positive constant, the pairing
\begin{equation}
 \Psi_b(\tau)
 :=\langle A^{1/4}u(b-\tau),A^{1/4}u(b+\tau)\rangle
 \label{eq:v7-Psi}
\end{equation}
stays above that constant times the critical record square whenever the lag
precedes the next doubled record on a hypothetical fixed-datum terminal
tail.  The exact terminal-tail question is not disproved here.  V7 tests,
and disproves, its extension to arbitrary regular record chronologies and
the idea that the finite-record properties alone supply the sign.  The
report's proof of the two-sided BKM-type bound also invokes a general
estimate for
\begin{equation}
 \Omega(a,b):=\langle\mathcal B(a),\Lambda b\rangle
             -\langle\Lambda a,\mathcal B(b)\rangle.
 \label{eq:v7-Omega}
\end{equation}
Both issues can be tested before any terminal compactness or donor export.
The precise audited locations are the reflected-pairing theorem, the displayed
pointwise estimate preceding (W.2), and the question labelled
\texttt{op:OS} in \path{NS_session_report(1).tex}.

\begin{assessment}[What is corrected, and what is not]
\label{ass:v7-correction-scope}
The unrestricted positive-record premise is false; a complete Navier--Stokes
counterexample is proved below.  The report's general cross-field estimate
\eqref{eq:v7-false-cross-estimate} is also false, including its use as a
pointwise estimate on all smooth time-shifted trajectories.  Consequently
that estimate and the proof which invokes it are not imported as established
results.  This does not by itself refute the report's more specially
conditioned conclusion involving a slow future norm product; that conclusion
would need an independent proof.

V6's reflected derivative and its four-way conditional alternative do not
use the false cross-field estimate and remain unchanged.  A negative
Eulerian reflected value is now explicitly known to occur through regular
transport, so it is not evidence of a singular source.  No claim that all
one-sided PDE estimates are insufficient is inferred from the failure of the
particular estimates tried so far.
\end{assessment}

All calculations in V7 use normalized Haar measure on
\(\Torus^3=(\R/2\pi\mathbb Z)^3\).  Thus
\(\Lambda=(-\Delta)^{1/2}\) on divergence-free, mean-zero fields and
\[
 \|u\|_{\dot H^{1/2}}^2=\langle u,\Lambda u\rangle.
\]
The viscosity \(\nu>0\) is fixed throughout the construction.  No forcing is
introduced.

% =============================================================================
\section{A finite Fourier witness against the cross-field strain estimate}
\label{sec:v7-finite-witness}
% =============================================================================

\begin{countertheorem}[Critical cross work is not controlled by strain and critical norms]
\label{cth:v7-cross-BKM}
There is no universal constant \(C\) such that every pair of smooth mean-zero
divergence-free periodic fields satisfies
\begin{equation}
 \boxed{
 |\Omega(a,b)|\le C\bigl(\|\nabla a\|_\infty+\|\nabla b\|_\infty\bigr)
 \bigl(\|a\|_{\dot H^{1/2}}^2+\|b\|_{\dot H^{1/2}}^2\bigr).}
 \label{eq:v7-false-cross-estimate}
\end{equation}
In fact, the ratio of the two sides with \(C=1\) is at least
\(\sqrt N/20\) on an explicit Fourier family.
\end{countertheorem}

\begin{proof}
For an integer \(N\ge1\), set
\begin{equation}
 a_N=(\sin y,0,N^{-1/2}\sin Nx),\qquad
 b_N=(\sin y,0,N^{-1/2}\sin y\cos Nx).
 \label{eq:v7-finite-fields}
\end{equation}
Both fields have zero mean and zero divergence.  Their nonlinearities are
already divergence-free, so the Leray projector changes nothing:
\[
 \mathcal B(a_N)=(0,0,\sqrt N\sin y\cos Nx),\qquad
 \mathcal B(b_N)=(0,0,-\sqrt N\sin^2y\sin Nx).
\]
The vertical part of \(a_N\) has frequency modulus \(N\), and that of
\(b_N\) has frequency modulus \(\sqrt{N^2+1}\).  Orthogonality gives
\begin{equation}
 \boxed{
 \Omega(a_N,b_N)=\frac{N+\sqrt{N^2+1}}4\ge\frac N2,}
 \label{eq:v7-finite-cross-work}
\end{equation}
and
\[
 \|a_N\|_{\dot H^{1/2}}^2=1,\qquad
 \|b_N\|_{\dot H^{1/2}}^2
 =\frac12+\frac{\sqrt{N^2+1}}{4N}.
\]
Their sum is less than two.  Direct differentiation, with the matrix
operator norm bounded by the sum of absolute entries, gives
\[
 \|\nabla a_N\|_\infty+\|\nabla b_N\|_\infty
 \le2+2\sqrt N+N^{-1/2}\le5\sqrt N.
\]
The ratio is therefore at least \(\sqrt N/20\), proving the assertion.
\end{proof}

\begin{remark}[Why a diagonal commutator bound cannot be polarized this way]
The cancellation in a diagonal transport energy pairs the same transported
field with itself.  In \eqref{eq:v7-Omega}, the two transported fields are
different.  The leading advection terms can add rather than cancel.  The
example does not yet assert that \(a_N,b_N\) are two states of one solution;
the next sections supply that stronger, dynamical counterexample.  No donor
alias table is being recomputed here: \eqref{eq:v7-finite-fields} is a direct
witness to one failed estimate.
\end{remark}

% =============================================================================
\section{An exact smooth Navier--Stokes family}
\label{sec:v7-shear-family}
% =============================================================================

Let \(N\ge2\) be an integer and define
\begin{equation}
 a_N:=\nu N^4,\qquad c=a_Nt.
 \label{eq:v7-amplitude-time}
\end{equation}
The scalar \(a_N\) in this section is an amplitude, not the vector field
\(a_N\) used only in \eqref{eq:v7-finite-fields}.  Let the complex periodic
function \(f_N(c,y)\) solve
\begin{equation}
 \left\{
 \begin{aligned}
 \partial_cf_N+iNe^{-c/N^4}\sin y\,f_N
   &=N^{-4}\partial_y^2f_N-N^{-2}f_N,\\
 f_N(0,y)&=1.
 \end{aligned}\right.
 \label{eq:v7-scalar-PDE}
\end{equation}
Set
\begin{equation}
 \boxed{
 u_N(t,x,y,z)=
 \left(a_Ne^{-\nu t}\sin y,\ 0,\
       \frac{2a_N}{\sqrt N}
       \operatorname{Im}\bigl(e^{iNx}f_N(a_Nt,y)\bigr)\right).}
 \label{eq:v7-exact-NS-family}
\end{equation}

\begin{proposition}[Exact, unforced, globally smooth realization]
\label{prop:v7-exact-NS}
For each \(N\) and the same fixed \(\nu\),
\eqref{eq:v7-exact-NS-family} is a globally smooth mean-zero solution of the
three-dimensional periodic Navier--Stokes equations.  Its pressure is
constant and may be set to zero.  Its horizontal shear and vertical component
satisfy
\begin{equation}
 U_t=\nu U_{yy},\qquad
 w_t+U(t,y)w_x=\nu(w_{xx}+w_{yy}).
 \label{eq:v7-shear-scalar-system}
\end{equation}
Moreover, for \(c\ge0\),
\begin{equation}
 \|f_N(c)\|_\infty\le1,\qquad
 \|\partial_yf_N(c)\|_\infty\le Nc.
 \label{eq:v7-scalar-sup-bounds}
\end{equation}
On \(0\le a_Nt\le M\),
\begin{equation}
 \boxed{
 \|\nabla u_N(t)\|_\infty
 \le a_N\bigl[1+2(1+M)\sqrt N\bigr].}
 \label{eq:v7-gradient-bound}
\end{equation}
\end{proposition}

\begin{proof}
For a field \((U(y,t),0,w(x,y,t))\) independent of \(z\), incompressibility
is automatic, and its convective derivative is \((0,0,Uw_x)\).  This vector
is divergence-free.  Hence the momentum equations reduce exactly to
\eqref{eq:v7-shear-scalar-system}.  Substitution of the complex Fourier mode
in \(x\), and division by \(a_N\), give
\eqref{eq:v7-scalar-PDE}.  Both components have zero spatial mean.

The scalar equation is uniformly parabolic for each fixed \(N\), with smooth
coefficients on every finite time interval.  Its differentiated energy
estimates give finite bounds at every Sobolev order on every such interval:
the imaginary multiplication term contributes only lower-order commutators,
and diffusion is nonpositive in each energy identity.  Fourier Galerkin
approximation and these estimates construct the solution, and the same
bounds continue it through every finite time.  This proves global smoothness
without appealing to three-dimensional global regularity.

The equation for \(|f_N|^2\) has nonpositive zeroth-order dissipation; its
maximum principle gives \(\|f_N(c)\|_\infty\le e^{-c/N^2}\le1\).
Differentiating in \(y\) adds the source
\(-iNe^{-c/N^4}\cos y\,f_N\).  The same maximum principle, or its
inhomogeneous Duhamel form, bounds its supremum by
\(N\int_0^c e^{-s/N^4}\,ds\le Nc\).
Finally, the vertical \(x\)-derivative is bounded by \(2a_N\sqrt N\),
the vertical \(y\)-derivative by \(2a_NM\sqrt N\), and the shear derivative
by \(a_N\).  Their sum proves \eqref{eq:v7-gradient-bound}.
\end{proof}

\subsection{Uniform approximation on the full record interval}

For an integer \(m\ge0\), write
\[
 \|g\|_{H_N^m}^2:=\sum_{j=0}^mN^{-2j}\|\partial_y^jg\|_{L^2(\Torus)}^2,
 \qquad F_N(c,y):=e^{-iNc\sin y}.
\]
These norms measure transverse derivatives in units of the longitudinal
frequency \(N\).

\begin{lemma}[Scaled Sobolev approximation, with fixed viscosity]
\label{lem:v7-scaled-approx}
For every finite \(M\) and integer \(m\ge0\), there is a finite constant
\(C_{M,m}\), independent of \(N\) and \(\nu\), such that
\begin{equation}
 \boxed{
 \sup_{0\le c\le M}\|f_N(c)-F_N(c)\|_{H_N^m}
 \le C_{M,m}N^{-2}.}
 \label{eq:v7-scaled-approx}
\end{equation}
In particular, the approximation is uniform through every fixed finite
collection of the critical records constructed below, not merely on the
short reflected windows.
\end{lemma}

\begin{proof}
The scaled derivatives of \(F_N\) are bounded uniformly for \(c\in[0,M]\):
\[
 \sup_{0\le c\le M}\|F_N(c)\|_{H_N^r}\le C_{M,r}
 \quad(r\ge0).
\]
This follows by successive differentiation of the displayed exponential;
each factor \(N\) is cancelled by a scaled derivative.
For \(g_N=f_N-F_N\), the equation has zero initial datum and forcing
\begin{equation}
 R_N=N^{-4}\partial_y^2F_N-N^{-2}F_N
       +iN(1-e^{-c/N^4})\sin y\,F_N.
 \label{eq:v7-approx-forcing}
\end{equation}
The first two terms have \(H_N^m\) norm at most \(C_{M,m}N^{-2}\).
Since \(N(1-e^{-c/N^4})\le MN^{-3}\), the third is smaller.

For completeness, the only potentially large differentiated coefficient is
\(iNe^{-c/N^4}\sin y\).  Its commutator with a scaled derivative is
\[
 [N^{-j}\partial_y^j,iNe^{-c/N^4}\sin y]g
 =ie^{-c/N^4}\sum_{\ell=1}^j\binom j\ell
 N^{1-\ell}(\partial_y^\ell\sin y)
 (N^{-(j-\ell)}\partial_y^{j-\ell}g).
\]
All coefficients are uniformly bounded.  The undifferentiated imaginary
multiplication cancels in the real energy, and both diffusion and
\(-N^{-2}\) are dissipative.  Summing differentiated energies therefore gives
\[
 \frac{d}{dc}\|g_N\|_{H_N^m}
 \le C_m\|g_N\|_{H_N^m}+C_{M,m}N^{-2}
\]
in the upper-derivative sense at zero norm.  Integration from zero proves
\eqref{eq:v7-scaled-approx}.
\end{proof}

% =============================================================================
\section{The critical symbol and its reflected limit}
\label{sec:v7-symbol}
% =============================================================================

Let
\[
 D_y=-i\partial_y,\qquad
 J_N=(1+N^{-2}D_y^2)^{1/2},\qquad
 U_{N,c}g=e^{-iNc\sin y}g,\qquad
 w_c(y)=\sqrt{1+c^2\cos^2y}.
\]
The operator \(J_N\) is the exact three-dimensional critical multiplier on
one longitudinal \(x\)-mode, after division by \(N\).

\begin{lemma}[Uniform square-root symbol limit]
\label{lem:v7-square-root}
For \(c\in[0,M]\) and every \(b\in H^2(\Torus)\),
\begin{equation}
 \boxed{
 \|U_{N,c}^*J_NU_{N,c}b-w_cb\|_2
 \le C_MN^{-1}\|b\|_{H^2}.}
 \label{eq:v7-square-root}
\end{equation}
The estimate is uniform on compact \(c\)-intervals.
\end{lemma}

\begin{proof}
Conjugation gives
\[
 U_{N,c}^*(D_y/N)U_{N,c}=D_y/N-c\cos y.
\]
Put \(g=-c\cos y\),
\[
 A_{N,c}=1+(D_y/N+g)^2,\qquad B_c=1+g^2.
\]
The first is self-adjoint with domain \(H^2\) and bounded below by one;
the second is multiplication by a smooth function bounded below by one.
On \(H^2\),
\[
 A_{N,c}-B_c=N^{-1}(D_yg+gD_y)+N^{-2}D_y^2.
\]
The resolvent representation of the square root yields, on \(b\in H^2\),
\begin{equation}
 (A_{N,c}^{1/2}-B_c^{1/2})b
 =\frac1\pi\int_0^\infty
 \lambda^{1/2}(A_{N,c}+\lambda)^{-1}
 (A_{N,c}-B_c)(B_c+\lambda)^{-1}b\,d\lambda.
 \label{eq:v7-root-resolvent}
\end{equation}
One may first integrate over a compact \(\lambda\)-interval; the estimate
below proves convergence and hence the identity on the stated domain.
The derivatives through order two of \((B_c+\lambda)^{-1}\) are bounded
by \(C_M(1+\lambda)^{-1}\).  Consequently
\[
 \|(A_{N,c}-B_c)(B_c+\lambda)^{-1}b\|_2
 \le \frac{C_M}{N(1+\lambda)}\|b\|_{H^2}.
\]
Also \(\|(A_{N,c}+\lambda)^{-1}\|\le(1+\lambda)^{-1}\), and
\(\int_0^\infty\lambda^{1/2}(1+\lambda)^{-2}\,d\lambda<\infty\).
Equation \eqref{eq:v7-root-resolvent} proves the bound, because
\(A_{N,c}^{1/2}=U_{N,c}^*J_NU_{N,c}\) and \(B_c^{1/2}=w_c\).
\end{proof}

Define the deterministic functions
\begin{align}
 \mathcal E(c)&:=\frac12+2\int_{\Torus}w_c(y)\,dy,
 \label{eq:v7-E-limit}\\
 \mathcal S(c,h)&:=\frac12+2\int_{\Torus}
                   w_c(y)\cos(2h\sin y)\,dy.
 \label{eq:v7-S-limit}
\end{align}
Here and below the one-dimensional torus integrals also use normalized
Haar measure.

\begin{theorem}[Critical-energy and two-sided pairing limits]
\label{thm:v7-symbol-limits}
Let \(r_N(t)=\|u_N(t)\|_{\dot H^{1/2}}\).  Uniformly for \(c\in[0,M]\),
\begin{equation}
 \boxed{
 a_N^{-2}r_N(c/a_N)^2=\mathcal E(c)+O_M(N^{-1}).}
 \label{eq:v7-energy-limit}
\end{equation}
For a fixed \(h>0\), uniformly when \(h/N\le c\le M-h/N\),
\begin{equation}
 \boxed{
 \frac{\Psi_{c/a_N}\bigl(h/(a_NN)\bigr)}{a_N^2}
 =\mathcal S(c,h)+O_{M,h}(N^{-1}).}
 \label{eq:v7-reflected-limit}
\end{equation}
Moreover, \(\mathcal E(0)=5/2\), \(\mathcal E\) is strictly increasing
on \((0,\infty)\), and \(\mathcal E(c)\to\infty\).
\end{theorem}

\begin{proof}
Orthogonality in vector components and in the longitudinal Fourier variable
gives the exact expression
\begin{equation}
 \frac{r_N(c/a_N)^2}{a_N^2}
 =\frac12e^{-2c/N^4}+2\langle f_N(c),J_Nf_N(c)\rangle.
 \label{eq:v7-exact-energy}
\end{equation}
The complex scalar product is conjugate-linear in its first argument; the
value in this formula is real.  The factor two is
\((2/\sqrt N)^2\cdot(N/2)\).
Similarly, the exact reflected expression is
\begin{equation}
 \frac{\Psi_{c/a_N}(h/(a_NN))}{a_N^2}
 =\frac12e^{-2c/N^4}
   +2\operatorname{Re}\langle f_N(c-h/N),J_Nf_N(c+h/N)\rangle.
 \label{eq:v7-exact-reflection}
\end{equation}
Since \(\|J_Ng\|_2=\|g\|_{H_N^1}\),
Lemma~\ref{lem:v7-scaled-approx} changes either expression by
\(O_{M,h}(N^{-2})\) when \(f_N\) is replaced by \(F_N\).
For the energy, apply Lemma~\ref{lem:v7-square-root} with \(b=1\).
For the cross term, use
\[
 F_N(c-h/N)=U_{N,c}e^{ih\sin y},\qquad
 F_N(c+h/N)=U_{N,c}e^{-ih\sin y}.
\]
The same lemma, now with \(b=e^{-ih\sin y}\), gives the real integral
\(\int w_c\cos(2h\sin y)\).  This proves both limits.
Finally,
\[
 \mathcal E'(c)=2\int_{\Torus}
 \frac{c\cos^2y}{\sqrt{1+c^2\cos^2y}}\,dy>0\quad(c>0),
 \qquad
 \mathcal E(c)\ge\frac12+\frac{4c}{\pi}.
\]
The remaining claims follow.
\end{proof}

\begin{lemma}[An exact negative sign, with a rational margin]
\label{lem:v7-negative-sign}
For every \(c\ge100\),
\begin{equation}
 \boxed{\mathcal S(c,2)<-\frac18\mathcal E(c).}
 \label{eq:v7-negative-limit}
\end{equation}
No numerical integration is needed for this inequality.
\end{lemma}

\begin{proof}
Pointwise,
\(0\le w_c-c|\cos y|\le1\).  Substitution \(v=\sin y\) on the four
quarter-circles gives
\[
 \int_{\Torus}|\cos y|\cos(4\sin y)\,dy=\frac{\sin4}{2\pi}.
\]
Thus
\[
 \mathcal S(c,2)\le\frac52+\frac{c\sin4}{\pi},\qquad
 \mathcal E(c)\le\frac52+\frac{4c}{\pi}.
\]
The elementary bounds \(3<\pi<22/7\) and \(\sin4<-3/4\) imply
\begin{equation}
 \mathcal S(c,2)+\frac18\mathcal E(c)
 \le\frac{45}{16}-\frac{19c}{264}
 \le-\frac{2315}{528}<0.
 \label{eq:v7-rational-margin}
\end{equation}
For an explicit rational certificate for the sine bound, let
\[
 P_{21}=\sum_{j=0}^{10}\frac{(-1)^j4^{2j+1}}{(2j+1)!}.
\]
Taylor's formula through degree 22 gives
\[
 \sin4\le P_{21}+\frac{4^{23}}{23!}
 =-\frac{12439009754071268}{16436269594119375}<-\frac34.
\]
This completes the sign proof.
\end{proof}

% =============================================================================
\section{Negative reflection at actual first-passage doubling records}
\label{sec:v7-record-theorem}
% =============================================================================

Fix \(c_*=100\).  For every integer \(j\ge0\), let \(c_j\) be the unique
solution of
\begin{equation}
 \mathcal E(c_j)=4^j\mathcal E(c_*).
 \label{eq:v7-record-cj}
\end{equation}
In particular, \(c_0=100\) and \(c_{j+1}>c_j\).  Set
\begin{equation}
 R_{N,j}:=2^ja_N\sqrt{\mathcal E(100)},\qquad
 \beta_{N,j}:=\inf\{t\ge0:r_N(t)=R_{N,j}\}.
 \label{eq:v7-record-definition}
\end{equation}
The next theorem proves that the indicated times exist on every prescribed
finite level range for all sufficiently large \(N\).

\begin{theorem}[One smooth datum realizes any prescribed finite counterexample chronology]
\label{thm:v7-record-counterexample}
Fix \(\nu>0\), an integer \(J\ge0\), and \(\varepsilon>0\).
For all sufficiently large \(N\), the single initial datum
\eqref{eq:v7-exact-NS-family} has the following properties.
\begin{enumerate}[label=\textup{(\roman*)},leftmargin=3em]
\item The records \(\beta_{N,0},\ldots,\beta_{N,J+1}\) all exist, are
      first-passage records of the same solution, and satisfy
      \begin{equation}
       a_N\beta_{N,j}\longrightarrow c_j,
       \qquad R_{N,0}>2r_N(0).
       \label{eq:v7-record-location}
      \end{equation}
\item The single lag
      \begin{equation}
       \tau_N:=\frac{2}{a_NN}=\frac{2}{\nu N^5}
       \label{eq:v7-lag}
      \end{equation}
      is admissible at every \(0\le j\le J\) and precedes the next
      doubled record:
      \[
       0<\tau_N<\min\{\beta_{N,j},\beta_{N,j+1}-\beta_{N,j}\}.
      \]
\item At every one of those records,
      \begin{equation}
       \boxed{
       \Psi_{\beta_{N,j}}(\tau_N)<-\frac18R_{N,j}^2.}
       \label{eq:v7-negative-record}
      \end{equation}
\item On each reflected window,
      \begin{equation}
       \boxed{
       \int_{\beta_{N,j}-\tau_N}^{\beta_{N,j}+\tau_N}
          \|\nabla u_N(t)\|_\infty\,dt<\varepsilon.}
       \label{eq:v7-small-strain-record}
      \end{equation}
\end{enumerate}
The solution is globally smooth, has no forcing, and has positive
heat-Reynolds covariance at every heat lag.  The source identity
\(\langle A^{-1/4}\mathcal B(u_N),A^{1/4}u_N\rangle=0\) holds throughout.
\end{theorem}

\begin{proof}
Choose a finite \(M>c_{J+1}+1\).  Uniform convergence in
\eqref{eq:v7-energy-limit}, strict increase of \(\mathcal E\), and the
intermediate value theorem show that, for each fixed \(j\le J+1\), the
first crossing of \(4^j\mathcal E(100)\) by
\(a_N^{-2}r_N(c/a_N)^2\) lies arbitrarily close to \(c_j\).  More
explicitly, for a small fixed \(\delta>0\), every
\(c\le c_j-\delta\) has limiting energy strictly below the level, while
\(c_j+\delta\) has limiting energy strictly above it.  Uniform convergence
preserves these strict inequalities.  This proves existence, first-passage
location, and ordering for the complete finite set of levels.  No monotonicity
of the approximate energy is assumed.  Also
\[
 a_N^{-2}r_N(0)^2=\frac52,\qquad
 \mathcal E(100)\ge\frac12+\frac{400}{\pi}>10,
\]
which proves the last assertion of \eqref{eq:v7-record-location}.

In the \(c\)-coordinate the reflected lag is \(2/N\), whereas the record
centres converge to distinct positive numbers.  Thus the lag is smaller than
the initial distance and every subsequent record gap for large \(N\).
Applying \eqref{eq:v7-reflected-limit} at
\(c=a_N\beta_{N,j}\), and using \eqref{eq:v7-negative-limit}, gives
\[
 a_N^{-2}\Psi_{\beta_{N,j}}(\tau_N)
 \longrightarrow\mathcal S(c_j,2)
 <-\tfrac18\mathcal E(c_j)
 =-\tfrac18a_N^{-2}R_{N,j}^2.
\]
The strict margin is uniform over the finite level set, proving
\eqref{eq:v7-negative-record}.

All reflected windows lie in \(0\le a_Nt\le M\).  Their length is
\(4/(a_NN)\).  By \eqref{eq:v7-gradient-bound},
\begin{equation}
 \int_{\beta_{N,j}-\tau_N}^{\beta_{N,j}+\tau_N}
       \|\nabla u_N(t)\|_\infty\,dt
 \le\frac4N+\frac{8(1+M)}{\sqrt N}\longrightarrow0.
 \label{eq:v7-explicit-strain-small}
\end{equation}
This proves the fourth item.  The final statements follow respectively from
Proposition~\ref{prop:v7-exact-NS}, the heat-kernel covariance identity already
proved in V6, and the ordinary incompressible kinetic-energy cancellation.
They are properties of this exact solution, not extra assumptions.
\end{proof}

\begin{corollary}[Failure of unrestricted OS-type record positivity]
\label{cor:v7-no-record-positivity}
There is no positive lower factor, even one chosen for each datum separately,
which is guaranteed solely from the properties in
Theorem~\ref{thm:v7-record-counterexample} to satisfy
\[
 \Psi_{\beta_j}(\tau)\ge\theta r(\beta_j)^2
 \quad\text{for every admissible }\tau<\beta_{j+1}-\beta_j
\]
on all the specified records.  There is also no universal estimate of the
form \eqref{eq:v7-false-cross-estimate} on all pairs
\(u(b-s),u(b+s)\) of an actual smooth Navier--Stokes trajectory.
\end{corollary}

\begin{proof}
A strictly negative value rules out every positive factor at that record.
For the second assertion, suppose a universal constant \(C\) worked along
all smooth trajectories.  Integrating the exact identity
\(\Psi_b'(s)=\Omega(u(b-s),u(b+s))\) would give
\begin{align*}
 |\Psi_b(\tau)-r(b)^2|
 &\le C\sup_{|t-b|\le\tau}\bigl(2r(t)^2\bigr)
       \int_{b-\tau}^{b+\tau}\|\nabla u(t)\|_\infty\,dt.
\end{align*}
For the records above, the supremum is at most \(4R_{N,j}^2\) for large
\(N\), by the uniform energy limit and the vanishing \(c\)-length of the
window.  The left side exceeds \(9R_{N,j}^2/8\), whereas the normalized
right side tends to zero by \eqref{eq:v7-explicit-strain-small}.  This is a
contradiction.
\end{proof}

\begin{firewall}[Finite chronologies are not a blow-up construction]
\label{fire:v7-finite-not-terminal}
The theorem says that, for every prescribed finite number of levels, one
datum realizes the complete sequence.  It does not assert the existence of
one datum with infinitely many terminal records.
The amplitude and energy vary across the sequence of examples.  For each
fixed example the displayed solution is globally smooth.  Thus the result
refutes an unrestricted finite-record sign theorem and the stated universal
cross-time estimate.  It does not refute a theorem restricted to an actual
fixed-datum terminal singular chronology with additional source incidence,
scale, or compactness assumptions.  The lag \eqref{eq:v7-lag} is short
compared with the viscous time of the occupied frequencies; this is not a
counterexample to V6's exact fixed-parabolic-lag alternative.
\end{firewall}

% =============================================================================
\section{The obstruction persists on genuinely rank-three data}
\label{sec:v7-rank-three}
% =============================================================================

The explicit carrier above has a rank-two Fourier subgroup.  Its global
smoothness is therefore consistent with the inherited rank-at-most-two
exclusion.  The sign obstruction is nevertheless not confined to that exact
subgroup.

\begin{lemma}[Persistence near a smooth reference on a fixed finite interval]
\label{lem:v7-smooth-stability}
Let \(U\) be a smooth periodic Navier--Stokes solution on \([0,T_0]\),
with fixed \(\nu>0\), and let \(m\ge4\).  For every \(\delta>0\)
there is \(\delta_0>0\) such that every smooth divergence-free datum
\(V_0\) with \(\|V_0-U(0)\|_{H^m}<\delta_0\) generates a smooth
solution through \(T_0\) satisfying
\[
 \sup_{0\le t\le T_0}\|V(t)-U(t)\|_{H^m}<\delta.
\]
No uniformity in the reference solution is asserted.
\end{lemma}

\begin{proof}
Write \(w=V-U\).  Its equation is
\[
 w_t+\nu Aw+
 \mathbb P(U\cdot\nabla w+w\cdot\nabla U+w\cdot\nabla w)=0.
\]
Integer differentiation, the divergence-free cancellation in the highest
transport derivatives, and the Sobolev product estimates for \(m\ge4\)
give
\[
 \frac12\frac d{dt}\|w\|_{H^m}^2+\nu\|\nabla w\|_{H^m}^2
 \le C_m\bigl(1+\|U\|_{H^{m+1}}+\|w\|_{H^m}\bigr)
            \|w\|_{H^m}^2.
\]
For clarity, each commutator contains either a derivative of \(U\) times
an order-at-most-\(m\) derivative of \(w\), or two derivatives of \(w\)
with total order at most \(m+1\); the usual \(H^m\) algebra and
\(H^m\hookrightarrow W^{1,\infty}\) bounds give exactly the displayed
terms after the top divergence-free term is cancelled.
On the bootstrap \(\|w\|_{H^m}\le\min\{1,\delta\}\), Gronwall has a
finite exponent because \(U\) is smooth through \(T_0\).  Taking
\(\delta_0\) small closes the bootstrap strictly.  Local smooth existence,
obtained by the same Galerkin estimates, then continues to \(T_0\): bounded
\(H^m\) norm prevents loss of the local continuation interval.  This proves
the claim.
\end{proof}

\begin{corollary}[Rank-three finite record counterexamples]
\label{cor:v7-rank-three-counterexample}
Fix \(\nu>0\), a finite \(J\), and \(\varepsilon>0\).  There is a
smooth mean-zero datum whose Fourier subgroup has lattice rank three and
whose unforced solution is smooth through first-passage records
\(\widetilde\beta_0,\ldots,\widetilde\beta_{J+1}\), with levels
\(\widetilde R_j=2^j\widetilde R_0\) and
\(\widetilde R_0>2\|u(0)\|_{\dot H^{1/2}}\), such that one admissible
lag \(\tau\) satisfies, for every \(j\le J\),
\begin{equation}
 \boxed{
 \Psi_{\widetilde\beta_j}(\tau)<-\frac1{10}\widetilde R_j^2,
 \qquad
 \int_{\widetilde\beta_j-\tau}^{\widetilde\beta_j+\tau}
       \|\nabla u(t)\|_\infty\,dt<\varepsilon.}
 \label{eq:v7-rank-three-result}
\end{equation}
The lag precedes the next doubled record.  The solution beyond the stated
finite interval is not addressed by this perturbation argument.
\end{corollary}

\begin{proof}
Use the reference family above with a sufficiently large fixed \(N\), and
choose a time \(T_0=M/a_N\) after the last required crossing.  Perturb its
datum by
\[
 \delta_N(0,\sin z,0),\qquad \delta_N\ne0.
\]
This is mean-zero and divergence-free.  The initial Fourier support now
contains nonzero multiples of the three independent directions
\((N,0,0),(0,1,0),(0,0,1)\), so its generated subgroup has rank three.

Here is why first-passage times, not only values at fixed times, are retained.
Choose disjoint small neighbourhoods of the limiting \(c_j\)'s.  The
uniform energy convergence makes the reference energy strictly below each
level before its neighbourhood and strictly above that level at a point
inside its right half.  These are strict finite-interval inequalities and
persist under a small uniform \(H^m\) perturbation.  Every perturbed first
crossing therefore lies in the corresponding neighbourhood.  The uniform
reflected limit and its strict rational sign margin, after the neighbourhoods
are reduced if necessary, make the reference reflected value less than
\(-\widetilde R_j^2/9\) at every possible centre in those neighbourhoods.
Small uniform perturbations of the two states preserve the weaker margin
\(-\widetilde R_j^2/10\).  The record levels can be kept equal to
\(R_{N,j}\); the strict initial inequality also persists.

Take the reference strain integral below \(\varepsilon/2\), uniformly for
these possible centres, using \eqref{eq:v7-gradient-bound}.  The
\(H^m\hookrightarrow W^{1,\infty}\) embedding preserves the required
\(\varepsilon\) bound for small enough perturbation.  The neighbourhoods
remain separated by more than twice the fixed lag.  Finally apply
Lemma~\ref{lem:v7-smooth-stability}, choosing \(\delta_N\) small enough
for all these finitely many strict conditions.  This proves the complete
same-datum finite chronology claim.
\end{proof}

\begin{remark}[A rank test is not an estimate for a selected terminal source]
The perturbation removes the exact rank-two carrier objection, but does not
claim to satisfy the additional held-out, source-age, material-capture, or
terminal-singular hypotheses of the V5--V6 export card.  The conclusion is
precisely that Fourier rank three alone does not restore reflected positivity
or the invalid cross-time strain estimate.
\end{remark}

% =============================================================================
\section{The actual transport removes the false reflected loss}
\label{sec:v7-material-correction}
% =============================================================================

The counterexample also identifies what the Eulerian pairing forgot.
For \(t_\pm=b\pm\tau\), define the horizontal transport of the displayed
shear by
\begin{equation}
 X_{-,+}(x,y,z)
 =\left(x+\sigma_{-,+}\sin y,y,z\right),\qquad
 \sigma_{-,+}:=\int_{t_-}^{t_+}a_Ne^{-\nu t}\,dt.
 \label{eq:v7-shear-transport}
\end{equation}
This is a volume-preserving periodic diffeomorphism.  It is the exact
horizontal part of the Lagrangian flow of \(u_N\).  The full flow also
moves \(z\); since \(u_N\) is independent of \(z\), this extra movement
has no effect on the following pullback.

\begin{proposition}[Transport-corrected pairing on the same records]
\label{prop:v7-material-pairing}
Define
\begin{equation}
 \Psi_b^{\rm mat}(\tau)
 :=\langle u_N(t_-),\Lambda\bigl(u_N(t_+)\circ X_{-,+}\bigr)\rangle.
 \label{eq:v7-material-Psi}
\end{equation}
For \(b=c/a_N\), \(\tau=h/(a_NN)\), and compact \(c\)-intervals,
\begin{equation}
 \boxed{
 a_N^{-2}\Psi_b^{\rm mat}(\tau)=\mathcal E(c)+o(1).}
 \label{eq:v7-material-limit}
\end{equation}
Consequently, at the records of
Theorem~\ref{thm:v7-record-counterexample},
\[
 \frac{\Psi_{\beta_{N,j}}^{\rm mat}(\tau_N)}{R_{N,j}^2}\longrightarrow1,
 \qquad
 \frac{\Psi_{\beta_{N,j}}(\tau_N)}{R_{N,j}^2}<-\frac18.
\]
The correction is determined by the actual velocity, not fitted to the
observed negative scalar.
\end{proposition}

\begin{proof}
The shear pullback multiplies the future complex vertical amplitude by
\(e^{iN\sigma_{-,+}\sin y}\).  In the \(c\)-coordinate,
\[
 N\sigma_{-,+}
 =N\int_{c-h/N}^{c+h/N}e^{-s/N^4}\,ds
 =2h+O_{M,h}(N^{-4}).
\]
Thus the pulled-back inviscid reference at \(c+h/N\) differs from
\(F_N(c-h/N)\) by \(O_{M,h}(N^{-4})\) in \(H_N^1\).
Multiplication by \(e^{iN\sigma_{-,+}\sin y}\) is uniformly bounded on
\(H_N^1\), because \(N\sigma_{-,+}\) stays bounded.  Applying
Lemma~\ref{lem:v7-scaled-approx} to both actual amplitudes proves
\[
 \|e^{iN\sigma_{-,+}\sin y}f_N(c+h/N)-f_N(c-h/N)\|_{H_N^1}
 =O_{M,h}(N^{-2}).
\]
The exact pairing formula analogous to
\eqref{eq:v7-exact-reflection} now differs from the past diagonal energy
by \(o(a_N^2)\).  The shear contribution is
\(a_N^2e^{-2\nu b}/2\), and its difference from the past shear energy is
also \(o(a_N^2)\).  Equation \eqref{eq:v7-energy-limit} finishes the proof.
\end{proof}

\subsection{A positive critical-metric estimate for arbitrary material transport}

The following estimate explains why small strain controls the transported
metric but not an untransported cross-time overlap.  In this subsection,
$\Lambda$ acts componentwise as $(-\Delta)^{1/2}$, also on fields which
are not divergence-free.  Composition need not preserve the solenoidal
subspace; this is a metric comparison, not a projected momentum identity.  It is a direct periodic
fractional-energy calculation.  The fractional seminorm/Fourier
correspondence is classical; a reference is Di Nezza--Palatucci--Valdinoci,
\emph{Hitchhiker's guide to the fractional Sobolev spaces},
\emph{Bull. Sci. Math.} \textbf{136} (2012), 521--573,
arXiv:1104.4345.  The periodic calculation needed here is included in full.

\begin{lemma}[Exact periodic difference-kernel representation]
\label{lem:v7-periodic-energy}
There is a fixed positive constant \(c_{\rm per}\) such that every smooth
periodic vector field satisfies
\begin{equation}
 \|f\|_{\dot H^{1/2}}^2
 =c_{\rm per}\int_{\Torus^3}\int_{\Torus^3}
 |f(x)-f(y)|^2\mathcal K(x-y)\,dx\,dy,
 \qquad
 \mathcal K(z)=\sum_{m\in\mathbb Z^3}|z+2\pi m|^{-4}.
 \label{eq:v7-periodic-kernel}
\end{equation}
The constant includes the chosen Haar normalization.  Polarization gives
the corresponding bilinear identity.
\end{lemma}

\begin{proof}
For a Fourier polynomial, Parseval gives
\[
 \int_{\Torus^3}|f(x)-f(x-z)|^2\,dx
 =2\sum_{k\in\mathbb Z^3}|\widehat f(k)|^2(1-\cos(k\cdot z)).
\]
Periodizing the kernel turns its \(z\)-integral into a Euclidean integral.
The latter is finite: near zero the numerator is \(O(|z|^2)\), and at
infinity \(|z|^{-4}\) is integrable in dimension three.  Rotation and
scaling give
\[
 \int_{\R^3}\frac{1-\cos(k\cdot z)}{|z|^4}\,dz
 =|k|\int_{\R^3}\frac{1-\cos z_1}{|z|^4}\,dz.
\]
The final integral is finite and strictly positive.  Choosing
\(c_{\rm per}\) to absorb it gives the exact multiplier \(|k|\).
Smooth approximation proves \eqref{eq:v7-periodic-kernel}; polarization
proves its bilinear form.
\end{proof}

\begin{theorem}[Material deformation controls the critical metric, not state alignment]
\label{thm:v7-material-metric}
Let \(X\) be the flow map from time \(s\) to time \(t\) of a smooth
periodic divergence-free velocity \(v\), and set
\[
 L_{s,t}:=\int_s^t\|S(v(r))\|_{{\rm op},\infty}\,dr,
 \qquad S(v)=\tfrac12(\nabla v+\nabla v^{\mathsf T}).
\]
For every smooth periodic vector field \(f\),
\begin{equation}
 \boxed{
 e^{-4L_{s,t}}\|f\|_{\dot H^{1/2}}^2
 \le\|f\circ X\|_{\dot H^{1/2}}^2
 \le e^{4L_{s,t}}\|f\|_{\dot H^{1/2}}^2.}
 \label{eq:v7-critical-metric-bound}
\end{equation}
For two such fields,
\begin{equation}
 \boxed{
 \left|\langle f\circ X,\Lambda(g\circ X)\rangle
        -\langle f,\Lambda g\rangle\right|
 \le(e^{4L_{s,t}}-1)
       \|f\|_{\dot H^{1/2}}\|g\|_{\dot H^{1/2}}.}
 \label{eq:v7-bilinear-metric-bound}
\end{equation}
These estimates do not imply that \(f\circ X\) is close to \(f\), or
that \(\langle f,\Lambda(f\circ X)\rangle\) is positive.
\end{theorem}

\begin{proof}
Lift the flow to \(\R^3\).  For two starting points \(a,b\), differentiating
the squared distance of their trajectories and integrating \(\nabla v\)
along the joining line give
\[
 -2\|S(v)\|_{{\rm op},\infty}|X(a)-X(b)|^2
 \le\frac d{dr}|X(a)-X(b)|^2
 \le2\|S(v)\|_{{\rm op},\infty}|X(a)-X(b)|^2.
\]
The skew part drops out of the scalar product.  Hence the lift and its inverse
are bi-Lipschitz with factors \(e^{\pm L_{s,t}}\).  Periodicity gives
\(X(a+2\pi m)=X(a)+2\pi m\), so for every lattice vector the same distance
comparison applies to \(a,b-2\pi m\).  Consequently
\[
 e^{-4L_{s,t}}\mathcal K(a-b)
 \le\mathcal K(X(a)-X(b))
 \le e^{4L_{s,t}}\mathcal K(a-b),
\]
and likewise for the inverse flow.

Incompressibility gives Jacobian one.  Change variables using \(X\) in
\eqref{eq:v7-periodic-kernel}, compare the inverse-flow kernel, and obtain
\eqref{eq:v7-critical-metric-bound}.  In the polarized formula, the absolute
kernel difference is bounded by
\((e^{4L_{s,t}}-1)\mathcal K\).  Cauchy--Schwarz with respect to the
positive difference-kernel measure gives
\eqref{eq:v7-bilinear-metric-bound}.

The final warning is intrinsic: a translation is an exact isometry of the
critical seminorm but can send a Fourier mode to its negative.  More
importantly for the present mean-zero problem, the actual shear transports
in Proposition~\ref{prop:v7-material-pairing} have vanishing integrated
strain on the reflected windows, while the untransported pairing is
negative.  Near-isometry of the metric does not give alignment of states.
\end{proof}

\begin{assessment}[A constructive replacement, not an automatic regularity theorem]
\label{ass:v7-replacement}
The difference between \eqref{eq:v7-critical-metric-bound} and the false
estimate \eqref{eq:v7-false-cross-estimate} is the actual material transport.
On the explicit family, its correction restores the complete critical
pairing.  On a general Navier--Stokes trajectory, the pulled-back velocity is
not constant: pressure and viscosity still act, and the critical metric is
transported rather than fixed.  Thus the theorem supplies a legitimate
geometric comparison, not a claim that material reflection is automatically
positive for every solution.

The original longitudinal increment/contact representation and V3's
mean-centred local energy calculation already retain deformation information.
They should be used as the existing starting point for the remaining signed
PDE estimate.  No new source is assigned to a phase lost merely because the
same physical field was observed in different spatial coordinates.
\end{assessment}

% =============================================================================
\section{What V7 closes and the revised next calculation}
\label{sec:v7-endpoint}
% =============================================================================

\begin{theorem}[Upstream conclusion of Version V7]
\label{thm:v7-summary}
Preserving the full V1--V6 source and all its correctly conditioned
identities, V7 establishes the following.
\begin{enumerate}[label=\textup{(\roman*)},leftmargin=3em]
\item The session report's universal cross-field critical strain estimate is
      false, with an explicit finite Fourier witness; it is also false as a
      pointwise bound on all time-shifted smooth Navier--Stokes trajectories.
\item Unrestricted positive reflection at regular first-passage critical
      records is false.  One unforced, globally smooth mean-zero datum can
      realize any prescribed finite number of counterexample doubling records.
\item The counterexample persists for genuinely rank-three data on the
      complete finite interval containing the prescribed records.
\item Its physical mechanism is rapid phase advection.  The actual material
      pullback restores the normalized pairing, and a general periodic
      fractional-energy estimate controls the material critical metric by
      integrated strain.
\end{enumerate}
These conclusions close the unrestricted reflection-sign route, not the
terminal Navier--Stokes problem.  The first-passage records and the pressure-
free scalar transport are actual solution data, not an abstract consistent
sequence of record inequalities.
\end{theorem}

\begin{proof}
The first assertion is Countertheorem~\ref{cth:v7-cross-BKM} and
Corollary~\ref{cor:v7-no-record-positivity}; the second is
Theorem~\ref{thm:v7-record-counterexample}; the third is
Corollary~\ref{cor:v7-rank-three-counterexample}; the fourth is
Proposition~\ref{prop:v7-material-pairing} and
Theorem~\ref{thm:v7-material-metric}.  The quantifier restriction is
Firewall~\ref{fire:v7-finite-not-terminal}.
\end{proof}

\begin{programme}[The next calculation must distinguish deformation from phase transport]
\label{prog:v7-next}
Preserve V1--V7.  Do not iterate the V6 successor alternatives merely to
rename an unresolved replenishment.  Do not seek a universal positive
Eulerian record pairing: V7 disproves that premise.

On one already selected source-visible terminal cell, evaluate the existing
signed critical increment/contact after the actual common material transport
has been accounted for.  The next substantive result must be either a
quantitative upper bound for the remaining deformation/pressure work in the
original source normalization, or a counterexample to that particular bound
which retains its exact fixed-datum terminal hypotheses.  A coordinate-phase
loss, a positive Reynolds tensor by itself, or another response Pythagoras is
not a substitute.

The NS--B collar and paid-stress inputs remain available where their declared
common-history hypotheses genuinely hold.  A new terminal-only reflected
claim is admissible only if it states an additional property absent from the
smooth record family above and proves why the selected singular chronology
has that property.  No global conclusion is made until the original terminal
scalar is controlled.
\end{programme}

\paragraph{Verification scope.}
The rational inequality \eqref{eq:v7-rational-margin} and its sine certificate
are checked using exact rational arithmetic in the accompanying working
calculation.  Fourier evaluation of the inviscid reference checks the
square-root symbol and reflected limit; it is not a Navier--Stokes simulation
or a replacement for Lemmas~\ref{lem:v7-scaled-approx} and
\ref{lem:v7-square-root}.  The asymptotic proof gives existence of a
sufficiently large \(N\), not an asserted explicit smallest \(N\).
Compilation and reference checks are document validation, not formal proof
verification.

\paragraph{Append-only history.}
V7 retains the complete V6 preterminal byte string.  The new material begins
at Section~\ref{sec:v7-upstream}.  An explicit upstream correction is appended
rather than silently rewriting an earlier claim.  The V6 conditional
reflected identities are retained, while the unrestricted sign premise and
the session report's cross-field estimate are rejected by proved examples.

\addtocontents{toc}{\protect\endgroup}
% =============================================================================
% END VERSION V7 MATERIAL -- append later versions before final terminator
% =============================================================================


\clearpage
% =============================================================================
% VERSION V8: NEW MATERIAL.  The complete V1--V7 source above is preserved
% verbatim; only its sole final document terminator was removed before append.
% =============================================================================
\hypersetup{pdftitle={NS-A Terminal Source Concentration Programme: Cumulative V8},pdfauthor={A. Pelissier}}
\makeatletter
\addtocontents{toc}{\protect\begingroup\protect\makeatletter}
\addtocontents{toc}{\protect\def\protect\l@section{\protect\@dottedtocline{1}{0em}{4.2em}}}
\addtocontents{toc}{\protect\def\protect\l@subsection{\protect\@dottedtocline{2}{1.5em}{4.8em}}}
\addtocontents{toc}{\protect\makeatother}
\makeatother
\renewcommand{\thesection}{V8.\arabic{section}}
\renewcommand{\thesubsection}{V8.\arabic{section}.\arabic{subsection}}
\renewcommand{\theequation}{V8.\arabic{equation}}
\setcounter{section}{0}
\setcounter{equation}{0}
\renewcommand{\theHsection}{V8.\arabic{section}}
\renewcommand{\theHsubsection}{V8.\arabic{section}.\arabic{subsection}}
\renewcommand{\theHtheorem}{V8.\arabic{section}.\arabic{theorem}}
\renewcommand{\theHequation}{V8.\arabic{equation}}

\begin{center}
{\LARGE\bfseries NS--A: Version V8}\\[0.5em]
{\large Stochastic Material Ancestry, Shrinking Preimage Cells,\\
and the Material-Sweeping Alternative}\\[0.5em]
{\normalsize 3 September 2026}
\end{center}
\addcontentsline{toc}{part}{Version V8: stochastic material ancestry and the material-sweeping alternative}

\begin{abstract}
The newly supplied NS--B continuation constructs the positive backward
material--caloric scalar router and removes the time, advection, viscous-cutoff,
heat-localization, and tame pressure-boundary collars on its declared branch.
This version does not repeat that calculation.  It resolves the next upstream
question: what full history object ties the already occurring signed annular
source to its terminal cell?  At positive viscosity the answer cannot be one
deterministic parcel.  The transition kernel is non-atomic.  The canonical
object is instead a signed source-marked path cylinder of the actual
drift--diffusion.

The path lift has exact total variation, polar sign, source and endpoint
marginals, terminal partition, and resolved/Reynolds decomposition.  A maximal
path estimate controls the full diffusion trajectory relative to the
deterministic material flow.  On a stronger diameter-compatible tame branch,
one mesoscopic radius simultaneously removes the imported collars and makes
the deterministic preimage cells shrink.  Bounded annular stress leverage
then promotes routed scalar capture to path-level source ancestry on moving
cells of volume $O(\rho_n^3)$.  Their centre either sweeps
super-parabolically, forcing divergent Lagrangian kinetic action, or the whole
path law collapses to the constant path at one point of the terminal singular
set.  On an efficient routed Reynolds branch, the same material tube carries
positive source action $\gtrsim\nu\mathcal R_n^2$ and an explicit diverging
action-density lower bound.  The note does not supply the corresponding local
upper stock, compactness, terminal trace, reflected positivity, Liouville
rigidity, or global regularity.
\end{abstract}

\section{Version V8: stochastic material ancestry, shrinking preimage cells, and the material-sweeping alternative}
\label{sec:v8-entry}

The newly supplied fixed-datum Reynolds/stress continuation has already
constructed the scalar localization which was only conditional at the end of
Version~V7.  Namely, a terminal weight is propagated by the backward
material--caloric equation
\begin{equation}
 (\partial_t+u\cdot\nabla+\nu\Delta)\omega=0,
 \qquad \omega(b)=\omega_b,
 \label{eq:v8-imported-router}
\end{equation}
and the resulting positive bistochastic operator cancels the time,
advective, and viscous cutoff collars in the localized kinetic-stress
identity.  That continuation also gives the exact source/output heat
commutator, the routed resolved/Reynolds split, and the routed
capture-or-escape alternative.  None of those results is repeated here.

The other newly supplied sector notes reinforce an order-of-operations point
which is used here only as proof discipline: a normalized response or
covariance cannot repair an earlier occurrence or source-incidence defect.
In particular, the Yang--Mills continuation moves upstream to an
unnormalized partition-curvature germ, while the Einstein--Standard-Model
continuation shows that thermal dressing cannot repair a wrong source
incidence.  No Yang--Mills or Standard-Model theorem is imported into the
Navier--Stokes argument.

The remaining NS--A issue is more specific.  Equation
\eqref{eq:v8-imported-router} gives a scalar weight on each earlier time
slice, but the phrase ``material capture'' still needs an actual history
object tying a source atom to its terminal cell.  At positive viscosity that
object cannot be a deterministic parcel map: molecular diffusion gives a
non-atomic transition law.  Version~V8 therefore constructs the canonical
\emph{signed path cylinder} of the already occurring annular scalar and then
executes it on the V5 heavy-thread card.

The new conclusions are as follows.
\begin{enumerate}[label=\textup{(U8.\arabic*)},leftmargin=5em]
\item Every integrable signed source density has a canonical lift to the
      path law of the actual drift-diffusion.  Its total variation, source
      marginal, terminal marginal, sign, terminal-cell split, and
      resolved/Reynolds split are exact.
\item A deterministic one-parcel representation is impossible for every
      positive time lag when \(\nu>0\).  The correct viscous ancestry is the
      stochastic path cylinder together with its deterministic material-flow
      shadow.
\item The full diffusion path stays, in mean-square supremum, within
      \(C\nu(b-t)e^{2\mathfrak B}\) of the deterministic material trajectory.
      This strengthens the endpoint estimate needed by the imported scalar
      router.
\item A stronger, diameter-compatible tame fork is isolated.  Either strain
      amplifies the shrinking source, cubic, or annular stock after a second
      exponential weight, or one may choose a mesoscopic radius for which
      both the imported collars and the deterministic material-cell diameter
      vanish.
\item On the bounded-leverage captured annular branch, the normalized signed
      path cylinder is asymptotically supported on the deterministic preimage
      of a shrinking terminal ball.  Those moving cells have volume
      \(O(\rho_n^3)\), without a cell-count factor.
\item The material-cell centre either sweeps super-parabolically, forcing a
      divergent one-trajectory kinetic action, or the complete path cylinder
      collapses to the constant terminal singular path.
\item On a captured Reynolds source--writer branch, either the unsigned
      source--writer turnover is unbounded relative to its signed resultant,
      or a positive source action of order \(\nu\mathcal R_n^2\) is carried by
      the same shrinking material tube.  On the nonsweeping branch its
      normalized source-action law lands at one point of \(\Sigma_T\).
\end{enumerate}

\begin{firewall}[A path representation is not a stochastic forcing]
\label{fire:v8-path-not-source}
The Brownian path below is the Markov representation of the molecular
viscosity already present in the deterministic Navier--Stokes equation.  It
is not an added random forcing, a new physical source, or a replacement of
the fixed datum.  The signed path cylinder merely disintegrates an already
assembled scalar against the transition law of
\eqref{eq:v8-imported-router}.  Its source price remains the price of the
original source atom.  Positive and negative parts are taken only after the
complete scalar and its terminal route have been assembled.
\end{firewall}

% =============================================================================
\section{The canonical signed path cylinder of one source scalar}
\label{sec:v8-path-cylinder}
% =============================================================================

Fix a smooth interval \(I=(a,b)\Subset(0,T_*)\).  Let
\(\mathbf P^u_{t,x}\) be the law on continuous periodic paths of the strong
solution of
\begin{equation}
 \dd X_r=u(r,X_r)\,\dd r+\sqrt{2\nu}\,\dd B_r,
 \qquad X_t=x,
 \qquad t\le r\le b.
 \label{eq:v8-SDE}
\end{equation}
Before time \(t\), extend the path constantly by \(X_r=x\).  Thus all laws
may be regarded on the fixed Polish space
\(
 \mathscr P_{a,b}:=I\times\Torus^3\times C([a,b];\Torus^3)
\), with the starting time and starting point retained as marks.  Denote the
transition density from \((t,x)\) to \((b,y)\) by
\(K^u_{t,b}(x,y)\), and write
\begin{equation}
 \mathcal U^u_{t,b}f(x)
 :=\int_{\Torus^3}K^u_{t,b}(x,y)f(y)\,\dd y
 =\mathbf E^u_{t,x}f(X_b).
 \label{eq:v8-router-kernel}
\end{equation}
For smooth divergence-free \(u\), this is exactly the imported router
\eqref{eq:v8-imported-router}.

Let \(j\in L^1(I\times\Torus^3;\R)\) be an already assembled signed scalar
source density.  For a terminal weight
\(0\le\omega_b\le1\), define the signed path measure
\(\boldsymbol\Xi_b^{\omega_b}[j]\) by
\begin{equation}
 \boxed{
 \begin{aligned}
 \int_{\mathscr P_{a,b}}F\,\dd\boldsymbol\Xi_b^{\omega_b}[j]
 :={}&\int_a^b\int_{\Torus^3}j(t,x)\,
 \mathbf E^u_{t,x}\!\left[
   \omega_b(X_b)F(t,x,X_{[a,b]})
 \right]\dd x\,\dd t
 \end{aligned}}
 \label{eq:v8-path-lift}
\end{equation}
for every bounded Borel function \(F\).  The unlocalized lift is
\(\boldsymbol\Xi_b[j]:=\boldsymbol\Xi_b^1[j]\).

\begin{theorem}[Exact path-cylinder incidence and total variation]
\label{thm:v8-path-incidence}
The measure in \eqref{eq:v8-path-lift} has the following exact properties.
\begin{enumerate}[label=\textup{(P\arabic*)},leftmargin=3.4em]
\item Its signed mass and total variation are
      \begin{equation}
       \boxed{
       \begin{aligned}
       \boldsymbol\Xi_b^{\omega_b}[j](\mathscr P_{a,b})
        &=\int_a^b\int j(t,x)
          \mathcal U^u_{t,b}\omega_b(x)\,\dd x\,\dd t,\\
       \|\boldsymbol\Xi_b^{\omega_b}[j]\|_{\rm TV}
        &=\int_a^b\int |j(t,x)|
          \mathcal U^u_{t,b}\omega_b(x)\,\dd x\,\dd t.
       \end{aligned}}
       \label{eq:v8-path-mass-TV}
      \end{equation}
      Its polar sign is \(\operatorname{sgn}j(t,x)\); no additional sign is
      created by path transport.
\item The marked source marginal is
      \begin{equation}
       \boxed{
       (\pi_{t,x})_*\boldsymbol\Xi_b^{\omega_b}[j]
       =j(t,x)\mathcal U^u_{t,b}\omega_b(x)\,\dd t\,\dd x.}
       \label{eq:v8-source-marginal}
      \end{equation}
\item The terminal endpoint marginal is
      \begin{equation}
       \boxed{
       (\pi_b)_*\boldsymbol\Xi_b^{\omega_b}[j](\dd y)
       =\omega_b(y)
        \left[\int_a^b\int_{\Torus^3}
          j(t,x)K^u_{t,b}(x,y)\,\dd x\,\dd t\right]\dd y.}
       \label{eq:v8-endpoint-marginal}
      \end{equation}
\item The terminal cell and its complement split one measure exactly:
      \begin{equation}
       \boxed{
       \boldsymbol\Xi_b[j]
       =\boldsymbol\Xi_b^{\omega_b}[j]
        +\boldsymbol\Xi_b^{1-\omega_b}[j].}
       \label{eq:v8-path-cell-split}
      \end{equation}
      More generally, every finite nonnegative terminal partition of unity
      \(\sum_Q\omega_{b,Q}=1\) gives
      \begin{equation}
       \boxed{
       \boldsymbol\Xi_b[j]
       =\sum_Q\boldsymbol\Xi_b^{\omega_{b,Q}}[j].}
       \label{eq:v8-path-partition}
      \end{equation}
      A cohort is obtained by summing its terminal weights before any
      variation or positive part is taken.
\item If \(j=\sum_{\alpha=1}^m j_\alpha\) is an exact signed source,
      resolved, covariance, return, or pressure decomposition formed before
      positivity, then
      \begin{equation}
       \boxed{
       \boldsymbol\Xi_b^{\omega_b}[j]
       =\sum_{\alpha=1}^m
        \boldsymbol\Xi_b^{\omega_b}[j_\alpha].}
       \label{eq:v8-path-component-split}
      \end{equation}
\end{enumerate}
Hence the imported routed capture/escape scalar is the endpoint projection
of one literal signed path cylinder, not an independent material occurrence.
\end{theorem}

\begin{proof}
Let
\[
 \dd\mathbf Q(t,x,\gamma)
 :=\dd t\,\dd x\,\dd\mathbf P^u_{t,x}(\gamma).
\]
The measure \(\boldsymbol\Xi_b^{\omega_b}[j]\) has density
\(j(t,x)\omega_b(\gamma(b))\) relative to the positive measure
\(\mathbf Q\).  Its absolute value therefore has density
\(|j(t,x)|\omega_b(\gamma(b))\), proving both identities in
\eqref{eq:v8-path-mass-TV} and the polar-sign assertion after conditional
expectation.  Testing against functions of \((t,x)\) gives
\eqref{eq:v8-source-marginal}.  Testing against a function of \(\gamma(b)\),
then inserting the transition density, gives
\eqref{eq:v8-endpoint-marginal}.  Equations
\eqref{eq:v8-path-cell-split}--\eqref{eq:v8-path-component-split} are
linearity before total variation or positive parts.
\end{proof}

\begin{corollary}[No double charge in a common path lift]
\label{cor:v8-path-no-double-charge}
Suppose the same signed scalar has a native-source representation, a
pressure-free stress representation, and a resolved/Reynolds representation.
After these representations have been identified on one spacetime history,
all three induce the same signed mass in
\eqref{eq:v8-path-mass-TV}.  Decomposing its path cylinder by terminal cell,
source type, or endpoint response does not create a second source price.
Only a nonzero residual between the corresponding signed path measures is a
new incidence obstruction.
\end{corollary}

\begin{proof}
The common scalar equality fixes the total signed functional.  Apply
\Cref{thm:v8-path-incidence} to each exact component decomposition and use
linearity.  If two proposed representations give different path measures,
their difference has zero or nonzero signed tests according to the actual
mixed incidence; it cannot be removed by repricing either marginal.
\end{proof}

% =============================================================================
\section{Positive viscosity forbids deterministic one-parcel ancestry}
\label{sec:v8-no-deterministic}
% =============================================================================

\begin{countertheorem}[The material--caloric router is not a composition operator]
\label{cth:v8-no-deterministic-router}
Assume \(\nu>0\), \(t<b\), and \(u\) is smooth.  There is no measurable map
\(Y_{t,b}:\Torus^3\to\Torus^3\) such that
\begin{equation}
 \mathcal U^u_{t,b}f(x)=f(Y_{t,b}(x))
 \label{eq:v8-false-deterministic}
\end{equation}
for every bounded continuous \(f\), even after modification on a null set.
Equivalently, the conditional endpoint law of a source point is genuinely
non-atomic.  A deterministic material trajectory may be retained as the
centre of the diffusion, but not as the complete viscous ancestry law.
\end{countertheorem}

\begin{proof}
Uniform parabolicity and the strong maximum principle give a smooth strictly
positive transition density
\(K^u_{t,b}(x,y)>0\) for every \(x,y\) when \(b-t>0\).  Choose a nonconstant
real continuous function \(f\).  For every \(x\), strict Jensen inequality
for the probability density \(K^u_{t,b}(x,\cdot)\) gives
\begin{equation}
 \bigl(\mathcal U^u_{t,b}f(x)\bigr)^2
 <\mathcal U^u_{t,b}(f^2)(x).
 \label{eq:v8-strict-Jensen}
\end{equation}
If \eqref{eq:v8-false-deterministic} held, both sides would equal
\(f(Y_{t,b}(x))^2\), a contradiction.
\end{proof}

\begin{assessment}[Correct meaning of full material incidence]
\label{ass:v8-material-meaning}
For \(\nu=0\), a scalar terminal cell is pulled back by the deterministic
volume-preserving flow.  For \(\nu>0\), the exact replacement is the path
kernel in \eqref{eq:v8-path-lift}.  Thus a request that every source atom lie
on one deterministic parcel is stronger than the deterministic PDE and is
false at every positive lag.  The valid NS--A target is instead:
\[
 \boxed{
 \begin{gathered}
 \text{source atom}
 \longrightarrow
 \text{actual drift--diffusion path law}\\[-0.1em]
 \longrightarrow
 \text{shrinking deterministic material-flow shadow}
 \end{gathered}}
\]
Represented return vectors and source dictionaries still require their own
linear ancestry maps; the theorem does not identify them merely from the
scalar path cylinder.
\end{assessment}

% =============================================================================
\section{A maximal material-diffusion estimate}
\label{sec:v8-maximal-spread}
% =============================================================================

For \((t,x)\), let \(A_r^{t,x}\) be the deterministic material trajectory
on the periodic cover,
\begin{equation}
 \frac{\dd}{\dd r}A_r^{t,x}=u(r,A_r^{t,x}),
 \qquad A_t^{t,x}=x.
 \label{eq:v8-deterministic-flow}
\end{equation}
Use the same initial lift for \(X_r^{t,x}\) in \eqref{eq:v8-SDE}.  Put
\begin{equation}
 \mathfrak B(t,b)
 :=\int_t^b\|\nabla u(r)\|_\infty\,\dd r.
 \label{eq:v8-BKM-window}
\end{equation}

\begin{lemma}[Maximal stochastic deviation from the material trajectory]
\label{lem:v8-maximal-deviation}
For every \(t<b\),
\begin{equation}
 \boxed{
 \mathbf E^u_{t,x}\!
 \left[\sup_{t\le r\le b}
  d_{\Torus^3}(X_r^{t,x},A_r^{t,x})^2\right]
 \le24\nu(b-t)e^{2\mathfrak B(t,b)}.}
 \label{eq:v8-maximal-deviation}
\end{equation}
The bound is uniform in \(x\).
\end{lemma}

\begin{proof}
On compatible lifts set \(Y_r=X_r^{t,x}-A_r^{t,x}\) and
\(L(r)=\|\nabla u(r)\|_\infty\).  Then
\[
 Y_r=\int_t^r\bigl(u(q,X_q)-u(q,A_q)\bigr)\,\dd q
       +\sqrt{2\nu}(B_r-B_t).
\]
With
\(R_r=\sup_{t\le q\le r}|Y_q|\) and
\(M_r=\sup_{t\le q\le r}|B_q-B_t|\),
\[
 R_r\le\int_t^rL(q)R_q\,\dd q+\sqrt{2\nu}M_r.
\]
Since \(M_r\) is nondecreasing, Gronwall's inequality gives
\(
 R_b\le\sqrt{2\nu}e^{\mathfrak B(t,b)}M_b
\).
For three-dimensional Brownian motion,
\[
 \mathbf E M_b^2
 \le\sum_{i=1}^3
   \mathbf E\sup_{t\le r\le b}|B_r^i-B_t^i|^2
 \le12(b-t)
\]
by Doob's \(L^2\) inequality.  The torus distance is no larger than the
lifted Euclidean distance, proving \eqref{eq:v8-maximal-deviation}.
\end{proof}

\begin{corollary}[Tube estimate for a signed captured path cylinder]
\label{cor:v8-weighted-tube}
Let \(j\in L^1(I\times\Torus^3)\), let
\(0\le\omega_b\le1\), and define
\begin{equation}
 V[j]:=\int_I\int|j|,
 \qquad
 V_{\omega_b}[j]
 :=\int_I\int|j|\mathcal U^u_{t,b}\omega_b.
 \label{eq:v8-two-variations}
\end{equation}
Assume \(V_{\omega_b}[j]>0\) and
\(V[j]\le C_0V_{\omega_b}[j]\).  Let
\(
 \Pi_b^{\omega_b}[j]
 :=|\boldsymbol\Xi_b^{\omega_b}[j]|/V_{\omega_b}[j]
\).
If \(|I|\le\tau\) and
\(
 \sup_{t\in I}\mathfrak B(t,b)\le\mathfrak B_I
\), then for every \(r>0\),
\begin{equation}
 \boxed{
 \Pi_b^{\omega_b}[j]
 \left\{\sup_{t\le q\le b}
 d(X_q,A_q)>r\right\}
 \le\frac{24C_0\nu\tau e^{2\mathfrak B_I}}{r^2}.}
 \label{eq:v8-weighted-tube}
\end{equation}
\end{corollary}

\begin{proof}
The numerator is at most
\[
 \int_I\int|j(t,x)|\,
 \mathbf P^u_{t,x}\!\left\{\sup_{t\le q\le b}d(X_q,A_q)>r\right\}
 \dd x\,\dd t.
\]
Use Markov's inequality and \Cref{lem:v8-maximal-deviation}, then divide by
\(V_{\omega_b}[j]\) and apply the assumed variation ratio.
\end{proof}

% =============================================================================
\section{A diameter-compatible tame fork on the terminal annular cards}
\label{sec:v8-diameter-tame}
% =============================================================================

Return to the V5 annular branch with terminal endpoint
\(b_n=\beta_{n+1}\), output frequency \(\Gamma_n\to\infty\), matched age
interval \(I_n^{\rm ann}\), and
\begin{equation}
 |I_n^{\rm ann}|
 \le\frac{R_t}{\nu\Gamma_n^2}.
 \label{eq:v8-annular-time-cap}
\end{equation}
Use the shrinking stocks already isolated in the supplied NS--B
continuation:
\begin{equation}
 \mathfrak m_n
 :=\int_{I_n^{\rm ann}}\|c[u(t)]\|_2\,\dd t,
 \qquad
 \mathfrak l_n
 :=\int_{I_n^{\rm ann}}\|u(t)\|_3^3\,\dd t,
 \qquad
 \mathfrak B_n
 :=\int_{I_n^{\rm ann}}\|\nabla u(t)\|_\infty\,\dd t.
 \label{eq:v8-card-stocks}
\end{equation}
For the fixed datum,
\(
 \mathfrak m_n\to0
\) and
\(
 \mathfrak l_n\to0
\).
After fixing the torus, viscosity, and source units once, put
\begin{equation}
 d_n:=\max\{\mathfrak m_n,\mathfrak l_n,\Gamma_n^{-1}\}.
 \label{eq:v8-dn}
\end{equation}
Equivalently, fixed positive reference units may be inserted in each entry;
they only change the constants below.

\begin{theorem}[Material-diameter amplification or a doubly tame radius]
\label{thm:v8-doubly-tame}
After passage to a subsequence, exactly one priority branch occurs.
\begin{enumerate}[label=\textup{(D\arabic*)},leftmargin=3.4em]
\item For some \(\varepsilon>0\), one of
      \begin{equation}
       e^{2\mathfrak B_n}\mathfrak m_n,
       \qquad
       e^{2\mathfrak B_n}\mathfrak l_n,
       \qquad
       \frac{e^{2\mathfrak B_n}}{\Gamma_n}
       \label{eq:v8-double-amplification-list}
      \end{equation}
      is at least \(\varepsilon\) on the selected subsequence.  In the three
      cases, respectively,
      \begin{equation}
       \boxed{
       \mathfrak B_n\ge\frac12\log\frac{\varepsilon}{\mathfrak m_n},
       \quad
       \mathfrak B_n\ge\frac12\log\frac{\varepsilon}{\mathfrak l_n},
       \quad
       \mathfrak B_n\ge\frac12\log(\varepsilon\Gamma_n).}
       \label{eq:v8-double-amplification}
      \end{equation}
      Thus the strain window diverges at a rate sufficient to prevent a
      shrinking deterministic material cell.
\item All three quantities in \eqref{eq:v8-double-amplification-list} tend to
      zero.  With
      \begin{equation}
       \boxed{\rho_n:=\sqrt{d_n},}
       \label{eq:v8-rho-choice}
      \end{equation}
      one has
      \begin{equation}
       \boxed{
       \begin{gathered}
       \rho_n\to0,
       \qquad \Gamma_n\rho_n\to\infty,
       \qquad e^{\mathfrak B_n}\rho_n\to0,\\
       \frac{e^{\mathfrak B_n}\mathfrak m_n}{\rho_n}\to0,
       \qquad
       \frac{e^{\mathfrak B_n}\mathfrak l_n}{\rho_n}\to0,
       \qquad
       \frac{e^{\mathfrak B_n}}{\Gamma_n\rho_n}\to0.
       \end{gathered}}
       \label{eq:v8-doubly-tame-properties}
      \end{equation}
      This radius simultaneously makes the imported heat and pressure
      collars vanish and makes every deterministic material preimage of a
      \(\rho_n\)-ball have vanishing diameter.
\end{enumerate}
\end{theorem}

\begin{proof}
Pass to a subsequence on which
\(e^{2\mathfrak B_n}d_n\) either tends to zero or has a positive lower
bound.  On the latter branch, one of the three entries defining \(d_n\)
attains the maximum along a further subsequence, which gives
\eqref{eq:v8-double-amplification-list}; taking logarithms proves
\eqref{eq:v8-double-amplification}.

On the zero branch, \(d_n\to0\), so \(\rho_n\to0\).  Since
\(d_n\ge\Gamma_n^{-1}\),
\(
 \Gamma_n\rho_n\ge\sqrt{\Gamma_n}\to\infty
\).
Moreover,
\[
 e^{\mathfrak B_n}\rho_n
 =\sqrt{e^{2\mathfrak B_n}d_n}\longrightarrow0.
\]
Each of \(\mathfrak m_n,\mathfrak l_n,\Gamma_n^{-1}\) is at most \(d_n\),
so every remaining ratio in
\eqref{eq:v8-doubly-tame-properties} is bounded by
\(e^{\mathfrak B_n}\sqrt{d_n}\to0\).
\end{proof}

\begin{assessment}[Relation to the supplied router theorem]
\label{ass:v8-router-refinement}
The supplied NS--B theorem needs the last three ratios in
\eqref{eq:v8-doubly-tame-properties}; it does not require
\(e^{\mathfrak B_n}\rho_n\to0\).  The additional factor is exactly what is
needed to shrink the \emph{deterministic preimage diameter}, not merely the
stochastic spread around one material trajectory.  Thus
\Cref{thm:v8-doubly-tame} is a genuine upstream refinement of the spatial
incidence row, rather than another Dini or age law.
\end{assessment}

% =============================================================================
\section{Execution on the bounded-leverage annular scalar}
\label{sec:v8-annular-execution}
% =============================================================================

Retain the V5 annular vector \(\varpi_n\), writer
\begin{equation}
 M_n(t):=S\!\left(A^{1/4}
  e^{-\nu(b_n-t)A}\varpi_n\right),
 \label{eq:v8-annular-writer}
\end{equation}
and the already assembled stress density
\begin{equation}
 j_n^{\rm ann}(t,x)
 :=(u\otimes u)^{\circ}(t,x):M_n(t,x).
 \label{eq:v8-annular-density}
\end{equation}
Its signed mass is
\begin{equation}
 H_n^{\rm ann}
 =\int_{I_n^{\rm ann}}\int j_n^{\rm ann}
 \ge h_*\mathcal R_n.
 \label{eq:v8-annular-margin}
\end{equation}
Work on the V5 bounded stress-leverage branch
\begin{equation}
 \Gamma_n^{3/2}
 \int_{I_n^{\rm ann}}\|u(t)\|_4^2\,\dd t
 \le L_*H_n^{\rm ann}.
 \label{eq:v8-bounded-leverage}
\end{equation}
The stress-leverage escape branch remains an earlier typed output.

\begin{lemma}[Bounded leverage controls the complete signed turnover]
\label{lem:v8-annular-variation}
There is a constant depending only on the fixed annular multipliers and
\(R_t\) such that
\begin{equation}
 \boxed{
 V_n^{\rm ann}
 :=\int_{I_n^{\rm ann}}\int|j_n^{\rm ann}|
 \le C L_*H_n^{\rm ann}.}
 \label{eq:v8-annular-variation}
\end{equation}
\end{lemma}

\begin{proof}
The annular writer estimate already used in V5 gives
\(
 \|M_n(t)\|_2\le C\Gamma_n^{3/2}
 e^{-c\nu(b_n-t)\Gamma_n^2}
\).
Therefore
\[
 V_n^{\rm ann}
 \le\int_{I_n^{\rm ann}}
   \|u(t)\|_4^2\|M_n(t)\|_2\,\dd t
 \le C\Gamma_n^{3/2}
   \int_{I_n^{\rm ann}}\|u(t)\|_4^2\,\dd t.
\]
Apply \eqref{eq:v8-bounded-leverage}.
\end{proof}

On the doubly tame branch choose a terminal cutoff
\(0\le\omega_{n,b_n}\le1\) with
\begin{equation}
 \omega_{n,b_n}=1\text{ on }B(x_n,\rho_n),
 \qquad
 \supp\omega_{n,b_n}\subset B(x_n,2\rho_n),
 \qquad
 \|\nabla\omega_{n,b_n}\|_\infty\le C\rho_n^{-1},
 \label{eq:v8-terminal-weight}
\end{equation}
and let
\(
 \omega_n(t)=\mathcal U^u_{t,b_n}\omega_{n,b_n}
\).
The centres are the V5 heavy-thread centres, so after subsequence
\begin{equation}
 x_n\longrightarrow x_*\in\Sigma_T.
 \label{eq:v8-centre-singular}
\end{equation}
The attached NS--B route theorem gives an exact escape/capture split.  Work
now on its captured annular-stress branch: for some \(\eta>0\),
\begin{equation}
 \boxed{
 h_n^{\rm mc}
 :=\int_{I_n^{\rm ann}}\int
   \omega_n(t,x)j_n^{\rm ann}(t,x)\,\dd x\,\dd t,
 \qquad
 |h_n^{\rm mc}|\ge\eta H_n^{\rm ann}.}
 \label{eq:v8-captured-annular}
\end{equation}
If this branch fails, the routed outside scalar, resolved re-entry,
paid-reader-null, dispersion, selector-off, or collar packet supplied by
NS--B is retained without alteration.

Define
\begin{equation}
 V_n^{\rm mc}
 :=\int_{I_n^{\rm ann}}\int
   \omega_n|j_n^{\rm ann}|.
 \label{eq:v8-captured-variation}
\end{equation}
Then
\begin{equation}
 \boxed{
 |h_n^{\rm mc}|\le V_n^{\rm mc}\le V_n^{\rm ann},
 \qquad
 \frac{V_n^{\rm ann}}{V_n^{\rm mc}}
 \le\frac{CL_*}{\eta}.}
 \label{eq:v8-variation-ratio}
\end{equation}

Let \(\Phi_{t,r}\) denote the deterministic volume-preserving flow of
\(u\), and define the material preimage cell
\begin{equation}
 \boxed{
 \mathcal M_n(t)
 :=\Phi_{t,b_n}^{-1}
   \bigl(B(x_n,3\rho_n)\bigr).}
 \label{eq:v8-material-cell}
\end{equation}

\begin{theorem}[Path-level ancestry and shrinking material preimage]
\label{thm:v8-annular-path-capture}
Let
\begin{equation}
 \Pi_n^{\rm path}
 :=\frac{
   |\boldsymbol\Xi_{b_n}^{\omega_{n,b_n}}[j_n^{\rm ann}]|
  }{V_n^{\rm mc}}
 \label{eq:v8-normalized-path-law}
\end{equation}
be the normalized total-variation law of the captured signed annular
cylinder.  Then:
\begin{enumerate}[label=\textup{(A\arabic*)},leftmargin=3.4em]
\item the terminal endpoint satisfies
      \begin{equation}
       \Pi_n^{\rm path}
       \{X_{b_n}\in B(x_n,2\rho_n)\}=1;
       \label{eq:v8-exact-endpoint-cell}
      \end{equation}
\item the complete path follows its deterministic material trajectory:
      \begin{equation}
       \boxed{
       \Pi_n^{\rm path}
       \left\{\sup_{t\le r\le b_n}
       d(X_r,A_r^{t,x})>\rho_n\right\}
       \longrightarrow0;}
       \label{eq:v8-path-follows-material}
      \end{equation}
\item the marked source marginal is asymptotically carried by the moving
      cell \(\mathcal M_n(t)\):
      \begin{equation}
       \boxed{
       \Pi_n^{\rm path}
       \{x\notin\mathcal M_n(t)\}\longrightarrow0;}
       \label{eq:v8-source-in-material-cell}
      \end{equation}
\item for every \(t\in I_n^{\rm ann}\),
      \begin{equation}
       \boxed{
       |\mathcal M_n(t)|=|B(x_n,3\rho_n)|\le C\rho_n^3,}
       \label{eq:v8-material-volume}
      \end{equation}
      and
      \begin{equation}
       \boxed{
       \operatorname{diam}\mathcal M_n(t)
       \le6e^{\mathfrak B_n}\rho_n\longrightarrow0.}
       \label{eq:v8-material-diameter}
      \end{equation}
\end{enumerate}
Thus the scalar material-capture row is promoted from one routed number to a
full source-marked path cylinder whose deterministic shadow is a shrinking,
volume-preserving material cell.
\end{theorem}

\begin{proof}
Item~\textup{(A1)} follows from the factor
\(\omega_{n,b_n}(X_{b_n})\) in the path measure and the support condition in
\eqref{eq:v8-terminal-weight}.  Apply
\Cref{cor:v8-weighted-tube} with
\(j=j_n^{\rm ann}\), \(r=\rho_n\), the ratio
\eqref{eq:v8-variation-ratio}, and
\eqref{eq:v8-annular-time-cap}.  This gives
\[
 \Pi_n^{\rm path}
 \left\{\sup d(X_r,A_r)>\rho_n\right\}
 \le C_{L_*,\eta,R_t}
 \left(\frac{e^{\mathfrak B_n}}
            {\Gamma_n\rho_n}\right)^2
 \longrightarrow0,
\]
which is item~\textup{(A2)}.

If \(x\notin\mathcal M_n(t)\), then
\(A_{b_n}^{t,x}\notin B(x_n,3\rho_n)\).  On the support of the captured path
one has \(X_{b_n}\in B(x_n,2\rho_n)\); hence
\(d(X_{b_n},A_{b_n}^{t,x})>\rho_n\).  Item~\textup{(A3)} follows from
item~\textup{(A2)}.

The flow \(\Phi_{t,b_n}\) preserves Haar measure because
\(\operatorname{div}u=0\), proving \eqref{eq:v8-material-volume}.  Its inverse
is \(e^{\mathfrak B_n}\)-Lipschitz.  The preimage of a ball of diameter
\(6\rho_n\) therefore has diameter at most
\(6e^{\mathfrak B_n}\rho_n\), and the last limit is
\eqref{eq:v8-doubly-tame-properties}.
\end{proof}

\begin{firewall}[The moving cell is not the original fixed cell]
\label{fire:v8-moving-not-fixed}
The set \(\mathcal M_n(t)\) is determined by the actual flow and the terminal
cell.  It need not coincide with the fixed Eulerian source cell selected in
V5, even though both have the same terminal centre.  Their signed difference
is precisely a material-capture or relocation residual; it is not erased by
\Cref{thm:v8-annular-path-capture}.  The theorem proves path incidence and
shrinking volume, not absence of rapid sweeping.
\end{firewall}

% =============================================================================
\section{Material-centre sweeping or collapse to one singular path}
\label{sec:v8-sweeping}
% =============================================================================

Let \(z_n(t)\) be the unique deterministic material trajectory ending at the
terminal centre:
\begin{equation}
 \Phi_{t,b_n}(z_n(t))=x_n,
 \qquad
 z_n'(t)=u(t,z_n(t)),
 \qquad
 z_n(b_n)=x_n.
 \label{eq:v8-material-centre}
\end{equation}
Put
\begin{equation}
 \sigma_n
 :=\sup_{t\in I_n^{\rm ann}}d(z_n(t),x_n).
 \label{eq:v8-sweeping-radius}
\end{equation}

\begin{theorem}[Super-parabolic sweeping or a constant-path singular thread]
\label{thm:v8-sweeping-alternative}
After passage to a subsequence, exactly one of the following occurs.
\begin{enumerate}[label=\textup{(S\arabic*)},leftmargin=3.4em]
\item \emph{Super-parabolic material sweeping:}
      \begin{equation}
       \boxed{\Gamma_n\sigma_n\longrightarrow\infty.}
       \label{eq:v8-super-sweeping}
      \end{equation}
      The actual centre trajectory then has divergent kinetic action:
      \begin{equation}
       \boxed{
       \int_{I_n^{\rm ann}}|u(t,z_n(t))|^2\,\dd t
       \ge\frac{\sigma_n^2}{|I_n^{\rm ann}|}
       \ge\frac{\nu}{R_t}(\Gamma_n\sigma_n)^2
       \longrightarrow\infty.}
       \label{eq:v8-sweeping-action}
      \end{equation}
\item \emph{Path-concentrated singular thread:}
      \begin{equation}
       \boxed{\sup_n\Gamma_n\sigma_n<\infty.}
       \label{eq:v8-bounded-sweeping}
      \end{equation}
      Define
      \begin{equation}
       \varepsilon_n
       :=\sigma_n+3e^{\mathfrak B_n}\rho_n+\rho_n.
       \label{eq:v8-path-radius}
      \end{equation}
      Then \(\varepsilon_n\to0\) and
      \begin{equation}
       \boxed{
       \Pi_n^{\rm path}
       \left\{\sup_{t\le r\le b_n}
       d(X_r,x_n)>\varepsilon_n\right\}
       \longrightarrow0.}
       \label{eq:v8-path-collapses}
      \end{equation}
      Consequently every fixed-time marginal and the normalized occupation
      measure of the path cylinder converge weakly to
      \(\delta_{x_*}\), and after constant extension to a common time interval
      the path laws converge in probability to the constant path
      \(r\mapsto x_*\).
\end{enumerate}
\end{theorem}

\begin{proof}
Pass to a subsequence on which \(\Gamma_n\sigma_n\) either tends to infinity
or is bounded.  The centre curve is continuous on the closure of
\(I_n^{\rm ann}\), so there is a time \(t_n\) at which its distance from
\(x_n=z_n(b_n)\) equals \(\sigma_n\).  Its length on \([t_n,b_n]\) is at
least that distance.  Therefore
\[
 \sigma_n
 \le \int_{t_n}^{b_n}|z_n'(t)|\,\dd t
 =\int_{t_n}^{b_n}|u(t,z_n(t))|\,\dd t,
\]
and Cauchy--Schwarz gives the exact estimate
\[
 \sigma_n^2
 \le (b_n-t_n)\int_{t_n}^{b_n}|u(t,z_n(t))|^2\,\dd t
 \le |I_n^{\rm ann}|
     \int_{I_n^{\rm ann}}|u(t,z_n(t))|^2\,\dd t.
\]
Together with \eqref{eq:v8-annular-time-cap}, this is
\eqref{eq:v8-sweeping-action}.

Now assume \eqref{eq:v8-bounded-sweeping}.  Then
\(\sigma_n=O(\Gamma_n^{-1})\to0\), while
\(e^{\mathfrak B_n}\rho_n\to0\) and \(\rho_n\to0\), so
\(\varepsilon_n\to0\).

For a starting point \(x\in\mathcal M_n(t)\), the terminal point
\(A_{b_n}^{t,x}\) lies in \(B(x_n,3\rho_n)\).  Compare its deterministic
trajectory at time \(r\) with the centre trajectory \(z_n(r)\).  The inverse
flow from \(b_n\) to \(r\) is \(e^{\mathfrak B_n}\)-Lipschitz, hence
\[
 d(A_r^{t,x},z_n(r))\le3e^{\mathfrak B_n}\rho_n.
\]
Therefore
\(
 d(A_r^{t,x},x_n)
 \le\sigma_n+3e^{\mathfrak B_n}\rho_n
\)
for every \(r\in[t,b_n]\).  Combine this with
\eqref{eq:v8-source-in-material-cell} and
\eqref{eq:v8-path-follows-material} to obtain
\eqref{eq:v8-path-collapses}.  On the same good event, the starting point \(x=A_t^{t,x}\) obeys the same
bound, so the constant extension before the marked start time is also within
\(\varepsilon_n\) of \(x_n\).  Extend each path constantly after \(b_n\) as
well.  The resulting laws live on the common space
\(C([0,T_*];\Torus^3)\).  Finally use \eqref{eq:v8-centre-singular}; bounded
continuous tests on path space and on occupation measures pass by convergence
in probability.
\end{proof}

\begin{remark}[What V7's phase-transport example becomes]
The V7 smooth counterexample shows why branch \textup{(S1)} cannot be removed
by a strain-only estimate: rapid sweeping can occur while the deformation of
the critical metric is small.  The present theorem does not call that phase
motion a singular source.  It records its exact cost along the selected
material centre.  Only branch \textup{(S2)} identifies the source cylinder
with one fixed terminal singular path.
\end{remark}

% =============================================================================
\section{Positive Reynolds source action on the shrinking material tube}
\label{sec:v8-Reynolds-tube}
% =============================================================================

The supplied NS--B continuation gives, on its routed Reynolds branch, a
source vector \(g_n(t,x)\), an annular writer \(W_n(x)\), and a signed scalar
\begin{equation}
 h_n^R
 :=-\int_{I_n^{\rm ann}}\int
   \omega_n(t,x)g_n(t,x)\cdot W_n(x)\,\dd x\,\dd t,
 \qquad
 |h_n^R|\ge h_R\mathcal R_n,
 \label{eq:v8-routed-Reynolds-margin}
\end{equation}
for some \(h_R>0\), together with the writer cap
\begin{equation}
 \int_{I_n^{\rm ann}}\int\omega_n|W_n|^2
 \le\frac{4R_t}{\nu}.
 \label{eq:v8-Reynolds-writer-cap}
\end{equation}
These are imported same-history statements; their covariance construction
and Riesz identity are not re-proved.

To pass from one signed resultant to a path-level source-action statement,
the unsigned turnover must also be retained.  Define
\begin{equation}
 \kappa_n^R
 :=\frac{
  \displaystyle\int_{I_n^{\rm ann}}\int|g_n\cdot W_n|
 }{|h_n^R|}.
 \label{eq:v8-Reynolds-cancellation-ratio}
\end{equation}

\begin{theorem}[Counterflow amplification or material-tube source action]
\label{thm:v8-material-source-action}
After passage to a subsequence, one of the following occurs.
\begin{enumerate}[label=\textup{(R\arabic*)},leftmargin=3.4em]
\item \emph{Unsigned source--writer amplification:}
      \begin{equation}
       \boxed{\kappa_n^R\longrightarrow\infty.}
       \label{eq:v8-Reynolds-counterflow}
      \end{equation}
      This is a same-history counterflow/inefficiency survivor; it is not a
      second Reynolds source.
\item \emph{Material-tube source action:} \(\kappa_n^R\le K_R<\infty\).
      On the doubly tame branch, form the path cylinder
      \(\boldsymbol\Xi_{b_n}^{\omega_{n,b_n}}[-g_n\cdot W_n]\).
      Then its normalized total variation satisfies
      \eqref{eq:v8-path-follows-material}--\eqref{eq:v8-source-in-material-cell}.
      Moreover, with
      \begin{equation}
       h_{n,\mathcal M}^R
       :=-\int_{I_n^{\rm ann}}
          \int_{\mathcal M_n(t)}
          \omega_n g_n\cdot W_n,
       \label{eq:v8-material-Reynolds-scalar}
      \end{equation}
      one has for all late \(n\)
      \begin{equation}
       \boxed{|h_{n,\mathcal M}^R|
       \ge\frac{h_R}{2}\mathcal R_n.}
       \label{eq:v8-material-Reynolds-margin}
      \end{equation}
      Hence the positive source action in the same moving material cell obeys
      \begin{equation}
       \boxed{
       \mathcal A_{n,\mathcal M}^R
       :=\int_{I_n^{\rm ann}}
          \int_{\mathcal M_n(t)}
          \omega_n|g_n|^2
       \ge\frac{\nu h_R^2}{16R_t}\mathcal R_n^2.}
       \label{eq:v8-material-action-lower}
      \end{equation}
      Since
      \begin{equation}
       |\mathscr M_n|
       :=\int_{I_n^{\rm ann}}|\mathcal M_n(t)|\,\dd t
       \le\frac{CR_t}{\nu\Gamma_n^2}\rho_n^3,
       \label{eq:v8-material-spacetime-volume}
      \end{equation}
      the mean source-action density and its essential supremum satisfy
      \begin{equation}
       \boxed{
       \frac{\mathcal A_{n,\mathcal M}^R}{|\mathscr M_n|}
       \ge c\frac{\nu^2h_R^2}{R_t^2}
          \frac{\mathcal R_n^2\Gamma_n^2}{\rho_n^3},
       \qquad
       \esssup_{\mathscr M_n}|g_n|^2
       \ge c\frac{\nu^2h_R^2}{R_t^2}
          \frac{\mathcal R_n^2\Gamma_n^2}{\rho_n^3}.}
       \label{eq:v8-material-density-lower}
      \end{equation}
\end{enumerate}
On branch \textup{(R2)} together with bounded material sweeping
\eqref{eq:v8-bounded-sweeping}, the normalized source-action measures
\begin{equation}
 \dd\alpha_n^R
 :=\frac{
  \one_{\mathcal M_n(t)}\omega_n|g_n|^2\,\dd x\,\dd t
 }{\mathcal A_{n,\mathcal M}^R}
 \label{eq:v8-action-law}
\end{equation}
have every spatial cluster point equal to \(\delta_{x_*}\).  Thus one
positive quadratic source carrier, not merely its signed scalar, lands on the
same material singular thread.
\end{theorem}

\begin{proof}
Pass to a subsequence on which \(\kappa_n^R\) either tends to infinity or is
bounded.  Work on the bounded branch.  In
\Cref{cor:v8-weighted-tube}, take
\(j=-g_n\cdot W_n\).  The unlocalized variation is at most
\(K_R|h_n^R|\), while the captured variation is at least \(|h_n^R|\).
Thus the same tube estimate used in
\Cref{thm:v8-annular-path-capture} applies.

The total variation of the source marginal outside \(\mathcal M_n(t)\) is
\(o(|h_n^R|)\).  Therefore its signed mass is also
\(o(|h_n^R|)\), and
\eqref{eq:v8-routed-Reynolds-margin} gives
\eqref{eq:v8-material-Reynolds-margin} for late \(n\).

Apply Cauchy--Schwarz on the moving spacetime cell:
\[
 |h_{n,\mathcal M}^R|^2
 \le
 \left(\int_{\mathscr M_n}\omega_n|g_n|^2\right)
 \left(\int_{\mathscr M_n}\omega_n|W_n|^2\right).
\]
The second factor is no larger than
\eqref{eq:v8-Reynolds-writer-cap}.  Insert
\eqref{eq:v8-material-Reynolds-margin} to obtain
\eqref{eq:v8-material-action-lower}.  The volume estimate follows from
\eqref{eq:v8-material-volume} and
\eqref{eq:v8-annular-time-cap}.  Divide the action lower bound by that volume;
since \(0\le\omega_n\le1\), the essential-supremum estimate follows as well.

On the bounded-sweeping branch,
\Cref{thm:v8-sweeping-alternative} places every material cell inside a ball
of radius tending to zero about \(x_n\).  The probability measures
\(\alpha_n^R\) are supported there, and
\eqref{eq:v8-centre-singular} proves their spatial convergence to
\(\delta_{x_*}\).
\end{proof}

\begin{remark}[The density blow-up is a concentration statement, not closure]
The right side of \eqref{eq:v8-material-density-lower} diverges on every
cofinal record branch because \(\mathcal R_n,\Gamma_n\to\infty\) and
\(\rho_n\to0\).  This is a positive source-action concentration on a genuine
material tube.  It is not yet an upper-bound contradiction: the inherited
Leray stock does not control this critical local action.  The next analytic
input must be a same-tube compactness, epsilon-regularity, pressure/product,
or critical source-action estimate.
\end{remark}

% =============================================================================
\section{Integrated V8 alternative and the exact next stop}
\label{sec:v8-integrated}
% =============================================================================

\begin{theorem}[Version V8 material-ancestry theorem]
\label{thm:v8-integrated}
Preserving every theorem and limitation of Versions~V1--V7, and importing the
new NS--B scalar-router theorem only at its declared common-history
interface, Version~V8 proves the following ordered conclusion on a selected
bounded-leverage annular terminal branch.
\begin{enumerate}[label=\textup{(V8.\arabic*)},leftmargin=5em]
\item The routed scalar has a canonical signed lift to the actual
      drift--diffusion path cylinder.  Its polar sign, total variation,
      source marginal, endpoint marginal, terminal partition, and source
      decomposition are exact.
\item A deterministic single-parcel lift is impossible for \(\nu>0\).
      Stochastic path incidence is therefore the minimum correct viscous
      material-history object.
\item Either one of the doubled strain-amplification quantities in
      \eqref{eq:v8-double-amplification-list} survives, or the radius
      \(\rho_n=\sqrt{d_n}\) simultaneously makes the heat collar, pressure
      collar, stochastic spread, and deterministic preimage diameter vanish.
\item On the captured annular-stress branch, bounded leverage supplies the
      total-variation efficiency needed to lift the complete signed scalar.
      Its source marginal is asymptotically supported on the volume-preserving
      material preimage \(\mathcal M_n(t)\), whose volume is
      \(O(\rho_n^3)\).
\item The material centre either has the super-parabolic sweeping action
      \eqref{eq:v8-sweeping-action}, or the full captured path law converges
      to the constant path at one point \(x_*\in\Sigma_T\).
\item On the routed Reynolds branch, either unsigned source--writer turnover
      diverges relative to its signed resultant, or the same moving tube
      carries source action at least
      \(c\nu\mathcal R_n^2\) and the density lower bound
      \eqref{eq:v8-material-density-lower}.  On the nonsweeping branch its
      normalized source-action law converges spatially to \(\delta_{x_*}\).
\item If the independently required routed paid-stress margin also passes,
      the existing NS--B covariance/paid-stress packet may be augmented by
      this path cylinder and material-tube source-action law.  A null
      paid-stress reader leaves the Reynolds-only packet intact.
\end{enumerate}
No item proves a local energy upper bound, strong product compactness,
terminal trace, reflected positivity, backward uniqueness, either ancient
Liouville theorem, or global three-dimensional regularity.
\end{theorem}

\begin{proof}
Items~\textup{(V8.1)--(V8.2)} are
\Cref{thm:v8-path-incidence,cth:v8-no-deterministic-router}.  Item
\textup{(V8.3)} is
\Cref{lem:v8-maximal-deviation,thm:v8-doubly-tame}.  Item
\textup{(V8.4)} is
\Cref{lem:v8-annular-variation,thm:v8-annular-path-capture}.  Item
\textup{(V8.5)} is
\Cref{thm:v8-sweeping-alternative}.  Item \textup{(V8.6)} is
\Cref{thm:v8-material-source-action}.  The last item is the declared
same-history conjunction with the supplied routed paid-stress alternative;
no additional estimate is inferred from that conjunction.
\end{proof}

\begingroup\small
\begin{longtable}{p{0.25\textwidth}p{0.19\textwidth}p{0.48\textwidth}}
\caption{NS--A Version V8 proof-status ledger.}\label{tab:v8-ledger}\\
\toprule Row & Status & Exact conclusion\\\midrule
\endfirsthead
\toprule Row & Status & Exact conclusion\\\midrule
\endhead
V1--V7 prefix & \PROVED{} preserved
 &The complete V7 preterminal source is retained byte-for-byte; no prior
 theorem is rewritten.\\[0.3em]
New NS--B router layer & \INHERITED{} at one interface
 &The positive scalar router, heat/pressure collar cancellations, and routed
 resolved/Reynolds/paid-stress alternatives are not recomputed.\\[0.3em]
Signed path cylinder & \PROVED{} exact
 &Source and endpoint marginals, total variation, sign, partitions, and
 component splits are the disintegration identities
 \eqref{eq:v8-path-mass-TV}--\eqref{eq:v8-path-component-split}.\\[0.3em]
Deterministic parcel map & \OBSTRUCTION{} exact
 &Strict parabolic Jensen inequality forbids a deterministic composition
 representation at every positive viscous lag.\\[0.3em]
Maximal material spread & \PROVED{} quantitative
 &The complete stochastic path has mean-square supremum distance at most
 \(24\nu(b-t)e^{2\mathfrak B}\) from the deterministic material path.\\[0.3em]
Diameter-compatible tame radius & \PROVED{} alternative
 &Either doubled strain amplification survives, or one radius makes all
 imported collars and the deterministic preimage diameter vanish.\\[0.3em]
Annular turnover & \PROVED{} bounded
 &Bounded stress leverage controls the total variation of the complete
 annular scalar by its signed record height.\\[0.3em]
Material preimage cell & \PROVED{} path-level capture
 &The captured scalar is carried by a source-marked path cylinder and a
 moving volume \(O(\rho_n^3)\) material cell.\\[0.3em]
Material-centre motion & \PROVED{} exact fork
 &Super-parabolic sweeping gives divergent one-trajectory kinetic action;
 otherwise the complete path law collapses to the singular constant path.\\[0.3em]
Reynolds counterflow & \OBSTRUCTION{} typed
 &Unbounded unsigned source--writer turnover is retained before any positive
 source action is promoted.\\[0.3em]
Material-tube Reynolds action & \PROVED{} on efficient branch
 &The same moving cell carries action \(\gtrsim\nu\mathcal R_n^2\) and the
 explicit density lower bound \eqref{eq:v8-material-density-lower}.\\[0.3em]
Path-level singular landing & \PROVED{} on nonsweeping branch
 &The path, source-position, occupation, and positive source-action laws land
 at one \(x_*\in\Sigma_T\).\\[0.3em]
Represented return vectors & \OPEN{} separate ancestry
 &The scalar path cylinder does not by itself transport every old response or
 return vector in the represented dictionary.\\[0.3em]
Strong compactness / epsilon regularity & \OPEN{} terminal PDE input
 &No upper stock currently contradicts the concentrated critical action or
 passes every nonlinear product.\\[0.3em]
Reflected positivity / Liouville & \OPEN
 &These are tested only after the path-concentrated common-history packet has
 passed the local PDE rows.\\[0.3em]
Global regularity & \OPEN
 &No surviving branch is eliminated for every fixed datum.\\
\bottomrule
\end{longtable}
\endgroup

\begin{programme}[Version V9: evaluate the material-tube PDE packet or stop]
\label{prog:v8-next}
Preserve Versions~V1--V8 verbatim.  Do not add another path metric, age law,
response graph, or material selector.

Proceed only on the path-concentrated branch of
\Cref{thm:v8-material-source-action}, and in the following order.
\begin{enumerate}[label=\textup{(V9.\arabic*)},leftmargin=5.2em]
\item Use the exact moving cells \(\mathcal M_n(t)\) and the already occurring
      source, writer, covariance, and paid-stress fields.  Evaluate the local
      energy and pressure/product packet on this same tube; do not replace it
      by a fixed fitted ball.
\item Compare the positive lower action
      \eqref{eq:v8-material-action-lower} with an actual same-tube upper
      stock.  A failed upper estimate is retained as a critical source-action,
      pressure, or product-concentration survivor.
\item If the normalized tube fields are compact at one physical scale, pass
      the signed source--writer scalar and positive covariance to the limit
      and export the resulting nonzero profile immediately.  If compactness
      fails, return the explicit concentration defect rather than another
      general concentration trichotomy.
\item Transport every represented old-response or return vector through its
      declared common-history map before calling the residual source fresh.
      The path cylinder closes scalar source incidence, not vector ancestry.
\item Stop at the first unavailable epsilon-regularity, terminal-trace,
      reflected, backward-uniqueness, or Liouville row.  No further compiler
      is admitted beyond that named PDE input.
\end{enumerate}
The permitted outputs are
\[
 \boxed{
 \begin{gathered}
 \text{doubled strain amplification}
 \quad\big|\quad
 \text{material route escape or source--writer counterflow}\\
 \big|\quad
 \text{super-parabolic material sweeping}
 \quad\big|\quad
 \text{one path-concentrated positive source-action thread}\\
 \big|\quad
 \text{a named local compactness, pressure, trace, or Liouville stop}.
 \end{gathered}}
\]
\end{programme}

\begin{assessment}[Append-only integrity]
\label{ass:v8-integrity}
The complete Version~V7 source preceding its sole final document terminator is
preserved verbatim.  The present material begins at
Section~\ref{sec:v8-entry}; the final command below is again unique.
\end{assessment}

\addtocontents{toc}{\protect\endgroup}
% =============================================================================
% END VERSION V8 MATERIAL -- append later versions before final terminator
% =============================================================================


\clearpage
% =============================================================================
% VERSION V9: NEW MATERIAL.  The complete V1--V8 source above is preserved
% verbatim; only its sole final document terminator was removed before append.
% =============================================================================
\hypersetup{pdftitle={NS-A Terminal Source Concentration Programme: Cumulative V9},pdfauthor={A. Pelissier}}
\makeatletter
\addtocontents{toc}{\protect\begingroup\protect\makeatletter}
\addtocontents{toc}{\protect\def\protect\l@section{\protect\@dottedtocline{1}{0em}{4.2em}}}
\addtocontents{toc}{\protect\def\protect\l@subsection{\protect\@dottedtocline{2}{1.5em}{4.8em}}}
\addtocontents{toc}{\protect\makeatother}
\makeatother
\renewcommand{\thesection}{V9.\arabic{section}}
\renewcommand{\thesubsection}{V9.\arabic{section}.\arabic{subsection}}
\renewcommand{\theequation}{V9.\arabic{equation}}
\setcounter{section}{0}
\setcounter{equation}{0}
\renewcommand{\theHsection}{V9.\arabic{section}}
\renewcommand{\theHsubsection}{V9.\arabic{section}.\arabic{subsection}}
\renewcommand{\theHtheorem}{V9.\arabic{section}.\arabic{theorem}}
\renewcommand{\theHequation}{V9.\arabic{equation}}

\begin{center}
{\LARGE\bfseries NS--A: Version V9}\\[0.5em]
{\large The Material--Caloric Energy Trace, the Parabolic Aspect-Ratio Stop,\\
and the Critical Product-Concentration Survivor}\\[0.5em]
{\normalsize 3 September 2026}
\end{center}
\addcontentsline{toc}{part}{Version V9: material--caloric energy, aspect ratio, and critical product concentration}

\begin{abstract}
Version~V8 promoted the routed annular scalar to a signed drift--diffusion
path cylinder and, on its efficient branch, placed positive Reynolds source
action of order $\nu\mathcal R_n^2$ on a shrinking material tube.  The
present continuation does not add another path law or concentration grammar.
It evaluates the PDE packet which that tube actually carries.

First, the backward material--caloric weight gives an exact local kinetic
energy identity.  On the doubly tame terminal cards, both its pressure
boundary current and its integrated viscous dissipation vanish, so the
weighted kinetic energy changes by $o(1)$.  The trace of the routed kinetic
stress equation is exactly twice this scalar identity.  The local
pressure--strain tensor is trace free and remains completely invisible to it.
Consequently scalar local energy cannot control the paid anisotropic stress
without an independent pressure/product row.

Second, the tube has an intrinsic two-scale mismatch.  If
$\ell_n=\Gamma_n^{-1}$ is the annular length and $\rho_n$ the material radius,
then $\rho_n/\ell_n\to\infty$, while
$\nu|I_n^{\rm ann}|/\rho_n^2\to0$.  At the annular scale the spatial tube
contains an unbounded number of unit cells; at the tube scale its time window
has vanishing parabolic depth.  Thus the V8 packet is not a one-scale
parabolic compactness cylinder.

Third, the same captured signed stress forces the scale-invariant local stock
\[
 \Gamma_n\int_{I_n^{\rm ann}}\int_{\mathcal M_n(t)}
       \omega_n|u|^4\gtrsim\nu\mathcal R_n^2.
\]
An annular multiplier commutator then yields either a concrete
route--dissipation collar or a positive output-annular Reynolds source action
of the same order.  The latter forces both a global annular $L^4$ record and
an enstrophy-square pulse.  Independently, localizing the covariance source
through the square root of the router gives an exact alternative between a
fractional source-localization collar and a weighted
$\dot H^{1/2}$ covariance-product concentration.

Finally, the terminal weighted kinetic energy either produces an atom in a
kinetic-energy concentration measure at the selected singular point, or the
thread is energy-null while its critical $L^4$ and source-action stocks
diverge.  A scaling counterexample proves that absolute local-energy and
integrated-dissipation smallness cannot bound this critical $L^4$ stock by
interpolation.  The first remaining PDE row is therefore one-scale
anisotropic product/pressure compactness, or an explicit annular-scale
microdelocalization defect.  No terminal trace, Liouville theorem, or global
regularity conclusion is asserted.
\end{abstract}

% =============================================================================
\section{Version V9 scope: test the tube, not another path representation}
\label{sec:v9-entry}
% =============================================================================

Retain the path-concentrated, bounded-sweeping, efficient routed branch of
Version~V8.  Write
\begin{equation}
 I_n:=I_n^{\rm ann}=(a_n,b_n),
 \qquad \ell_n:=\Gamma_n^{-1},
 \qquad
 |I_n|\le \frac{R_t\ell_n^2}{\nu},
 \label{eq:v9-standing-window}
\end{equation}
with
\begin{equation}
 \Gamma_n\to\infty,
 \qquad \rho_n\to0,
 \qquad \Gamma_n\rho_n\to\infty,
 \qquad e^{\mathfrak B_n}\rho_n\to0.
 \label{eq:v9-standing-scales}
\end{equation}
The material--caloric weight $\omega_n$ satisfies
\begin{equation}
 (\partial_t+u\cdot\nabla+\nu\Delta)\omega_n=0,
 \qquad 0\le\omega_n\le1,
 \qquad
 \|\nabla\omega_n(t)\|_\infty
 \le C\frac{e^{\mathfrak B_n}}{\rho_n},
 \label{eq:v9-router}
\end{equation}
and its terminal value is supported in
$B(x_n,2\rho_n)$, where $x_n\to x_*\in\Sigma_T$.
The deterministic material preimage is
\begin{equation}
 \mathcal M_n(t)=\Phi_{t,b_n}^{-1}(B(x_n,3\rho_n)).
 \label{eq:v9-material-cell}
\end{equation}
Define its spacetime trace by
\begin{equation}
 \mathfrak T_n:=\{(t,x):t\in I_n,\ x\in\mathcal M_n(t)\}.
 \label{eq:v9-spacetime-tube}
\end{equation}

We use three already occurring lower packets.  The captured annular stress
has
\begin{equation}
 h_n^{\rm mc}
 =\int_{I_n}\int\omega_n
   (u\otimes u)^\circ:M_n,
 \qquad
 |h_n^{\rm mc}|\ge\eta H_n^{\rm ann},
 \qquad
 H_n^{\rm ann}\ge h_*\mathcal R_n,
 \label{eq:v9-captured-stress}
\end{equation}
where
\begin{equation}
 M_n(t)=S\!\left(A^{1/4}
 e^{-\nu(b_n-t)A}\varpi_n\right).
 \label{eq:v9-stress-writer}
\end{equation}
The routed Reynolds source and writer satisfy
\begin{equation}
 h_n^R=-\int_{I_n}\int\omega_n g_n\cdot W_n,
 \qquad |h_n^R|\ge h_R\mathcal R_n,
 \qquad W_n=\Lambda\varpi_n,
 \label{eq:v9-Reynolds-scalar}
\end{equation}
and
\begin{equation}
 \mathcal A_n^{R,\omega}:=
 \int_{I_n}\int\omega_n|g_n|^2
 \ge c_R\nu\mathcal R_n^2,
 \qquad
 c_R:=\frac{h_R^2}{4R_t}.
 \label{eq:v9-imported-action}
\end{equation}
Replacing $c_R$ by a smaller fixed constant if necessary accommodates the
material-cell restriction of \eqref{eq:v8-material-action-lower}.  No new
source is declared in \eqref{eq:v9-imported-action}; it is the same positive
heat-covariance source already assembled in Versions~V6--V8.

Finally, put
\begin{equation}
 \delta_n^E:=\int_{I_n}\|\nabla u(t)\|_2^2\,\dd t.
 \label{eq:v9-tail-dissipation}
\end{equation}
Since $a_n\uparrow T$ and the Leray energy stock is finite,
\begin{equation}
 \boxed{\delta_n^E\longrightarrow0.}
 \label{eq:v9-tail-dissipation-vanish}
\end{equation}
The first calculation is the exact information carried by this vanishing
stock on the routed tube.

\begin{firewall}[What is and is not reopened]
\label{fire:v9-no-reopen}
The heat covariance, source factorization, scalar router, path-cylinder
marginals, material capture, and paid-stress alternatives of Versions~V6--V8
are frozen.  Version~V9 does not revisit the unweighted terminal Dini
functional, the reflected-sign shortcut rejected in Version~V7, or the
NS--B logarithmic-age compiler.  Its only question is whether the already
selected material tube supplies a one-scale local PDE upper packet.
\end{firewall}

% =============================================================================
\section{The exact material--caloric local energy trace}
\label{sec:v9-local-energy}
% =============================================================================

For a smooth material--caloric weight $\omega$ on $[a,b]$, define
\begin{align}
 \mathscr E_\omega(t)
 &:=\frac12\int_{\Torus^3}\omega(t,x)|u(t,x)|^2\,\dd x,
 \label{eq:v9-local-energy}\\
 \mathscr D_\omega[a,b]
 &:=\int_a^b\int_{\Torus^3}
       \omega|\nabla u|^2\,\dd x\,\dd t,
 \label{eq:v9-local-dissipation}\\
 \mathscr P_\omega[a,b]
 &:=\int_a^b\int_{\Torus^3}
       p\,u\cdot\nabla\omega\,\dd x\,\dd t.
 \label{eq:v9-local-pressure-current}
\end{align}

\begin{theorem}[Exact local energy identity on the material--caloric route]
\label{thm:v9-routed-energy}
Let $u$ be a smooth unforced periodic Navier--Stokes solution and let
\begin{equation}
 (\partial_t+u\cdot\nabla+\nu\Delta)\omega=0.
 \label{eq:v9-router-general}
\end{equation}
Then
\begin{equation}
 \boxed{
 \mathscr E_\omega(b)-\mathscr E_\omega(a)
 +\nu\mathscr D_\omega[a,b]
 =\mathscr P_\omega[a,b].}
 \label{eq:v9-routed-energy-identity}
\end{equation}
No time-cutoff, advective-cutoff, or viscous-cutoff term remains.
\end{theorem}

\begin{proof}
Set $e=|u|^2/2$.  The smooth local energy equation is
\begin{equation}
 \partial_te+u\cdot\nabla e+\operatorname{div}(pu)
 =\nu\Delta e-\nu|\nabla u|^2.
 \label{eq:v9-point-energy}
\end{equation}
Differentiate $\int\omega e$.  The two transport terms assemble as
\[
 -\int e\,u\cdot\nabla\omega
 -\int\omega u\cdot\nabla e
 =-\int u\cdot\nabla(\omega e)=0
\]
by incompressibility and periodicity.  The two Laplacian terms cancel by
self-adjointness:
\[
 -\nu\int e\Delta\omega+\nu\int\omega\Delta e=0.
\]
Finally,
\(
 -\int\omega\operatorname{div}(pu)=\int pu\cdot\nabla\omega
\).
Thus
\[
 \frac{\dd}{\dd t}\mathscr E_\omega(t)
 =-\nu\int\omega|\nabla u|^2
   +\int pu\cdot\nabla\omega.
\]
Integrating proves \eqref{eq:v9-routed-energy-identity}.
\end{proof}

\begin{proposition}[The routed pressure boundary current is tame]
\label{prop:v9-energy-pressure}
For every smooth weight,
\begin{equation}
 \left|\int_{\Torus^3}p\,u\cdot\nabla\omega\right|
 \le C\|\nabla\omega\|_\infty\|u\|_3^3.
 \label{eq:v9-energy-pressure-bound}
\end{equation}
Consequently, on the Version~V8 doubly tame cards,
\begin{equation}
 \boxed{
 |\mathscr P_{\omega_n}[a_n,b_n]|
 \le C\frac{e^{\mathfrak B_n}\mathfrak l_n}{\rho_n}
 \longrightarrow0.}
 \label{eq:v9-energy-pressure-vanish}
\end{equation}
\end{proposition}

\begin{proof}
The periodic pressure equation is
\(
 -\Delta p=\partial_i\partial_j(u_i u_j)
\), with mean zero.  Calder\'on--Zygmund boundedness on $L^{3/2}$ gives
\[
 \|p\|_{3/2}\le C\|u\otimes u\|_{3/2}=C\|u\|_3^2.
\]
H\"older's inequality proves \eqref{eq:v9-energy-pressure-bound}.  Use
\eqref{eq:v9-router} and the definition of $\mathfrak l_n$ in
\eqref{eq:v8-card-stocks}; the final limit is one of
\eqref{eq:v8-doubly-tame-properties}.
\end{proof}

\begin{corollary}[Vanishing routed dissipation and asymptotic energy conservation]
\label{cor:v9-energy-conservation}
Put
\begin{equation}
 \mathscr E_n^-:=\mathscr E_{\omega_n}(a_n),
 \qquad
 \mathscr E_n^+:=\mathscr E_{\omega_n}(b_n),
 \qquad
 \mathscr D_n:=\mathscr D_{\omega_n}[a_n,b_n].
 \label{eq:v9-energy-endpoints}
\end{equation}
Then
\begin{equation}
 \boxed{
 \mathscr D_n\le\delta_n^E\longrightarrow0,
 \qquad
 \mathscr E_n^+-\mathscr E_n^-\longrightarrow0.}
 \label{eq:v9-energy-small-change}
\end{equation}
\end{corollary}

\begin{proof}
The first assertion follows from $0\le\omega_n\le1$ and
\eqref{eq:v9-tail-dissipation-vanish}.  Insert it and
\eqref{eq:v9-energy-pressure-vanish} into
\eqref{eq:v9-routed-energy-identity}.
\end{proof}

The preceding identity is also the complete scalar trace of the routed
kinetic-stress equation.  This fact identifies exactly what local energy
cannot see.

\begin{proposition}[Energy is the trace; pressure--strain is trace free]
\label{prop:v9-trace-blindness}
For the routed tensors of the supplied NS--B continuation,
\begin{equation}
 \operatorname{tr}\mathbb K_\omega=2\mathscr E_\omega,
 \qquad
 \operatorname{tr}\mathbb M_{\nabla,\omega}
 =\int\omega|\nabla u|^2,
 \qquad
 \boxed{\operatorname{tr}\Pi_\omega=0,}
 \label{eq:v9-traces}
\end{equation}
and
\begin{equation}
 \operatorname{tr}\mathbb B_{p,\omega}
 =2\int p\,u\cdot\nabla\omega.
 \label{eq:v9-pressure-current-trace}
\end{equation}
The inherited routed stress equation is
\begin{equation}
 \frac{\dd}{\dd t}\mathbb K_\omega
 =-2\nu\mathbb M_{\nabla,\omega}+2\Pi_\omega+\mathbb B_{p,\omega}.
 \label{eq:v9-imported-routed-stress}
\end{equation}
Taking its trace gives exactly twice
\eqref{eq:v9-routed-energy-identity}.  Its trace-free part remains
\begin{equation}
 \boxed{
 \frac{\dd}{\dd t}\mathbb K_\omega^\circ
 =-2\nu\mathbb M_{\nabla,\omega}^\circ
  +2\Pi_\omega
  +\mathbb B_{p,\omega}^\circ.}
 \label{eq:v9-tracefree-stress}
\end{equation}
Thus the vanishing scalar pressure current in
\eqref{eq:v9-energy-pressure-vanish} does not control the physical local
pressure--strain tensor $\Pi_{\omega_n}$.
\end{proposition}

\begin{proof}
The first two identities are the definitions.  Since
\(\operatorname{tr}S_u=\operatorname{div}u=0\),
\[
 \operatorname{tr}\Pi_\omega
 =\int\omega p\operatorname{tr}S_u=0.
\]
Taking the trace of
\(
 p(u\otimes\nabla\omega+\nabla\omega\otimes u)
\)
gives $2pu\cdot\nabla\omega$.  Tracing
\eqref{eq:v9-imported-routed-stress} now gives the differential form of
\eqref{eq:v9-routed-energy-identity}; subtracting one third of that trace
from the tensor equation gives \eqref{eq:v9-tracefree-stress}.
\end{proof}

\begin{firewall}[A scalar energy estimate cannot close the anisotropic row]
\label{fire:v9-energy-not-stress}
The conclusion is not that $\Pi_{\omega_n}$ is large.  It is that the scalar
energy balance contains no information about it: $\Pi_\omega$ lies entirely
in $\operatorname{Sym}_0(3)$.  Any estimate which controls the local paid
stress through pressure--strain is therefore an additional anisotropic
product theorem, not a corollary of
\eqref{eq:v9-routed-energy-identity}.
\end{firewall}

% =============================================================================
\section{Terminal kinetic-energy atom or an energy-null routed thread}
\label{sec:v9-energy-atom}
% =============================================================================

\begin{theorem}[Energy atom versus energy-null material route]
\label{thm:v9-energy-atom}
After passage to a subsequence, exactly one of the following occurs.
\begin{enumerate}[label=\textup{(E\arabic*)},leftmargin=3.4em]
\item There is $e_*>0$ such that
      \begin{equation}
       \mathscr E_n^+\longrightarrow e_*.
       \label{eq:v9-positive-terminal-energy}
      \end{equation}
      If
      \begin{equation}
       |u(b_n,x)|^2\,\dd x\ \stackrel{*}{\rightharpoonup}\ \mu_T
       \quad\text{in }\mathcal M(\Torus^3),
       \label{eq:v9-energy-measure-limit}
      \end{equation}
      then
      \begin{equation}
       \boxed{\mu_T(\{x_*\})\ge2e_*.}
       \label{eq:v9-energy-atom-lower}
      \end{equation}
\item The routed endpoint energies vanish:
      \begin{equation}
       \boxed{
       \mathscr E_n^+\longrightarrow0,
       \qquad
       \mathscr E_n^-\longrightarrow0.}
       \label{eq:v9-energy-null}
      \end{equation}
      Together with \eqref{eq:v9-energy-small-change}, the complete scalar
      local-energy packet on $I_n$ is then asymptotically null.
\end{enumerate}
\end{theorem}

\begin{proof}
The numbers $\mathscr E_n^+$ lie in the compact interval
$[0,E_*/2]$.  Pass to a subsequence on which they converge to $e_*\ge0$.
If $e_*=0$, \eqref{eq:v9-energy-small-change} gives the second endpoint
limit.

Assume $e_*>0$.  The measures in \eqref{eq:v9-energy-measure-limit} have
uniformly bounded mass, so a weak-* convergent subsequence exists.  At the
terminal time,
\(
 \supp\omega_{n,b_n}\subset B(x_n,2\rho_n)
\)
and $0\le\omega_{n,b_n}\le1$.  Hence
\begin{equation}
 |u(b_n)|^2\,\dd x\bigl(B(x_n,2\rho_n)\bigr)
 \ge 2\mathscr E_n^+.
 \label{eq:v9-ball-energy-lower}
\end{equation}
For every fixed $r>0$, the left ball is contained in
$\overline B(x_*,r)$ for all late $n$.  The closed-set half of the
Portmanteau theorem gives
\[
 2e_*
 \le\limsup_n\int_{\overline B(x_*,r)}|u(b_n)|^2
 \le\mu_T(\overline B(x_*,r)).
\]
Let $r\downarrow0$ and use continuity from above of the finite measure
$\mu_T$.
\end{proof}

\begin{assessment}[What the atom alternative means]
\label{ass:v9-energy-atom-meaning}
The atom is a concentration statement for a subsequential terminal kinetic
energy measure.  It is not by itself a Navier--Stokes singularity certificate.
On the complementary branch the scalar energy carried by the routed cell
vanishes; the next sections show that the critical source and product stocks
nevertheless diverge.  That separation is the relevant intermittency
obstruction.
\end{assessment}

% =============================================================================
\section{The material tube is not one parabolic cylinder}
\label{sec:v9-aspect-ratio}
% =============================================================================

Put
\begin{equation}
 \mathfrak a_n:=\frac{\rho_n}{\ell_n}=\Gamma_n\rho_n.
 \label{eq:v9-aspect}
\end{equation}
By \eqref{eq:v9-standing-scales}, $\mathfrak a_n\to\infty$.

\begin{theorem}[Parabolic aspect-ratio obstruction]
\label{thm:v9-aspect-ratio}
The V8 material cells satisfy, uniformly in $t\in I_n$,
\begin{equation}
 |\mathcal M_n(t)|=|B(x_n,3\rho_n)|\asymp\rho_n^3
 \label{eq:v9-cell-volume}
\end{equation}
and
\begin{equation}
 \boxed{
 \frac{\operatorname{diam}\mathcal M_n(t)}{\ell_n}
 \ge c\mathfrak a_n\longrightarrow\infty.}
 \label{eq:v9-spatial-aspect}
\end{equation}
At the tube radius, however,
\begin{equation}
 \boxed{
 \frac{\nu|I_n|}{\rho_n^2}
 \le\frac{R_t}{\mathfrak a_n^2}
 \longrightarrow0.}
 \label{eq:v9-temporal-aspect}
\end{equation}
Consequently there is no sequence of radii $r_n$ and fixed constants
$c,C>0$ such that simultaneously
\begin{equation}
 \operatorname{diam}\mathcal M_n(t)\le Cr_n
 \quad(t\in I_n),
 \qquad
 \nu|I_n|\ge c r_n^2.
 \label{eq:v9-no-one-scale}
\end{equation}
Thus the selected packet is spatially mesoscopic at the annular time scale
and temporally ultrathin at the material radius.
\end{theorem}

\begin{proof}
The flow $\Phi_{t,b_n}$ preserves Haar measure, proving
\eqref{eq:v9-cell-volume}.  On the fixed compact three-torus there is a
constant $C_0$ such that every measurable set $E$ satisfies
\(
 |E|\le C_0(\operatorname{diam}E)^3
\): if the diameter is below the injectivity radius, contain $E$ in one
geodesic ball; above that radius the estimate follows after increasing
$C_0$.  Hence
\(
 \operatorname{diam}\mathcal M_n(t)\ge c_0\rho_n
\), which proves \eqref{eq:v9-spatial-aspect}.  Formula
\eqref{eq:v9-temporal-aspect} is
\eqref{eq:v9-standing-window} divided by $\rho_n^2/\nu$.

If \eqref{eq:v9-no-one-scale} held, the diameter lower bound would give
$r_n\ge c'\rho_n$, while the time inequality and
\eqref{eq:v9-standing-window} would give $r_n\le C'\ell_n$.  This contradicts
$\rho_n/\ell_n\to\infty$.
\end{proof}

\begin{corollary}[Exact size of the scale mismatch]
\label{cor:v9-aspect-count}
At each fixed time, the material-cell volume measured in annular units is
\begin{equation}
 \boxed{
 \ell_n^{-3}|\mathcal M_n(t)|\asymp\mathfrak a_n^3\longrightarrow\infty,}
 \label{eq:v9-annular-cell-count}
\end{equation}
whereas its available time depth measured in material-radius units is
$O(\mathfrak a_n^{-2})$.  These two powers are geometric identities, not an
estimate obtained by counting the source once per cell.
\end{corollary}

\begin{firewall}[Why one-scale compactness has not yet been reached]
\label{fire:v9-aspect-stop}
At scale $\ell_n$ the interval has the correct parabolic order but the tube
contains $\asymp\mathfrak a_n^3$ spatial cells.  At scale $\rho_n$ the cell
is tight but the observed interval occupies a vanishing fraction of one
parabolic time.  Epsilon regularity, an ancient-profile extraction, or a
one-scale product limit therefore requires either a further annular-cell
capture or a longer same-history tube.  Neither is supplied by path
concentration alone.
\end{firewall}

% =============================================================================
\section{Annular source promotion and the route--dissipation collar}
\label{sec:v9-annular-source}
% =============================================================================

The next calculation asks whether the routed Reynolds scalar can be assigned
to the actual output annulus before its positive action is used.  Let
$\psi\in C_c^\infty((1/4,4))$ be real, with $\psi=1$ on $[1/2,2]$, and set
\begin{equation}
 \Psi_n:=\psi(\Lambda/\Gamma_n).
 \label{eq:v9-annular-projector}
\end{equation}
The V5 support of $\varpi_n$ gives
\begin{equation}
 \Psi_nW_n=W_n,
 \qquad \|W_n\|_2\le2\Gamma_n.
 \label{eq:v9-writer-annular}
\end{equation}

For a smooth divergence-free field $w$, abbreviate
\begin{equation}
 c[w]:=A^{-1/4}\mathbb P\operatorname{div}(w\otimes w).
 \label{eq:v9-source-functional}
\end{equation}

\begin{lemma}[Critical source estimate and Reynolds-source length]
\label{lem:v9-source-length}
For every smooth mean-zero divergence-free $w$,
\begin{equation}
 \boxed{
 \|A^{-1/4}\mathbb P\operatorname{div}(w\otimes w)\|_2
 \le C\|\nabla w\|_2^2.}
 \label{eq:v9-critical-source-bound}
\end{equation}
For the heat-covariance source
\(
 g_n=P_{\nu(b_n-t)}c[u]-c[v_{n,t}]
\)
one consequently has
\begin{equation}
 \boxed{
 \int_{I_n}\|g_n(t)\|_2\,\dd t
 \le\mathfrak m_n+C\delta_n^E.}
 \label{eq:v9-g-length}
\end{equation}
\end{lemma}

\begin{proof}
The multiplier $A^{-1/4}\mathbb P\operatorname{div}$ has order $1/2$.
The fractional product estimate and Sobolev embedding give
\[
 \begin{aligned}
 \|A^{-1/4}\mathbb P\operatorname{div}(w\otimes w)\|_2
 &\le C\|w\otimes w\|_{\dot H^{1/2}}\\
 &\le C\|w\|_6\|\Lambda^{1/2}w\|_3
 \le C\|\nabla w\|_2^2.
 \end{aligned}
\]
Here the last inequality applies
$\dot H^{1/2}\hookrightarrow L^3$ to $\Lambda^{1/2}w$.
The source-commutator identity follows from
\eqref{eq:v6-source-factorization}.  Heat contraction in $\dot H^1$ gives
\(
 \|\nabla v_{n,t}\|_2\le\|\nabla u(t)\|_2
\).
Therefore
\[
 \|g_n(t)\|_2
 \le\|c[u(t)]\|_2+C\|\nabla u(t)\|_2^2.
\]
Integrate.
\end{proof}

\begin{lemma}[Annular multiplier commutator]
\label{lem:v9-annular-commutator}
For every Lipschitz scalar $\omega$ and every vector field $f$,
\begin{equation}
 \boxed{
 \|[\Psi_n,\omega]f\|_2
 \le C\Gamma_n^{-1}\|\nabla\omega\|_\infty\|f\|_2.}
 \label{eq:v9-annular-commutator}
\end{equation}
The constant depends only on the fixed multiplier $\psi$ and the torus.
\end{lemma}

\begin{proof}
The smooth annular multiplier has a periodic convolution kernel $K_n$ whose
first absolute moment satisfies
\begin{equation}
 \int_{\Torus^3}d(y,0)|K_n(y)|\,\dd y\le C\Gamma_n^{-1}.
 \label{eq:v9-kernel-first-moment}
\end{equation}
This follows by periodizing the Euclidean Schwartz kernel
$\Gamma_n^3K(\Gamma_n\cdot)$; the nonzero period images are rapidly
decaying uniformly for $\Gamma_n\ge1$.  Therefore
\[
 [\Psi_n,\omega]f(x)
 =\int K_n(y)[\omega(x-y)-\omega(x)]f(x-y)\,\dd y.
\]
Minkowski's inequality, translation invariance of $L^2$, and the Lipschitz
bound prove \eqref{eq:v9-annular-commutator}.
\end{proof}

Define the dimensionless route--dissipation ratio
\begin{equation}
 \mathfrak C_n^{E\!\to\!R}
 :=\frac{e^{\mathfrak B_n}\delta_n^E}{\rho_n}.
 \label{eq:v9-route-dissipation-ratio}
\end{equation}

\begin{theorem}[Route--dissipation survivor or positive annular source action]
\label{thm:v9-annular-source-action}
After passage to a subsequence, one of the following occurs.
\begin{enumerate}[label=\textup{(A\arabic*)},leftmargin=3.4em]
\item The route--dissipation ratio survives:
      \begin{equation}
       \boxed{\liminf_n\mathfrak C_n^{E\!\to\!R}>0.}
       \label{eq:v9-route-dissipation-survivor}
      \end{equation}
      This is a physical localization/energy collar on the already selected
      tube.
\item The ratio tends to zero.  Put
      \begin{equation}
       h_n^{R,\rm ann}
       :=-\int_{I_n}\int\omega_n(\Psi_ng_n)\cdot W_n.
       \label{eq:v9-annular-Reynolds-scalar}
      \end{equation}
      Then
      \begin{equation}
       \boxed{
       |h_n^{R,\rm ann}|
       \ge\frac{h_R}{2}\mathcal R_n}
       \label{eq:v9-annular-Reynolds-margin}
      \end{equation}
      for all late $n$, and
      \begin{equation}
       \boxed{
       \mathcal A_n^{R,\rm ann}
       :=\int_{I_n}\int\omega_n|\Psi_ng_n|^2
       \ge\frac{\nu h_R^2}{16R_t}\mathcal R_n^2.}
       \label{eq:v9-annular-source-action-lower}
      \end{equation}
      Hence the positive source action can be placed on the literal output
      annulus, not only on the complete fused source.
\end{enumerate}
\end{theorem}

\begin{proof}
Pass to a subsequence on which
$\mathfrak C_n^{E\!\to\!R}$ either tends to zero or has a positive lower
bound.  Work on the zero branch.  Self-adjointness of $\Psi_n$ and
\eqref{eq:v9-writer-annular} give
\[
 \begin{aligned}
 h_n^R-h_n^{R,\rm ann}
 &=-\int_{I_n}\langle[\Psi_n,\omega_n]g_n,W_n\rangle\,\dd t.
 \end{aligned}
\]
By \Cref{lem:v9-annular-commutator}, \eqref{eq:v9-router},
\eqref{eq:v9-writer-annular}, and \Cref{lem:v9-source-length},
\begin{equation}
 |h_n^R-h_n^{R,\rm ann}|
 \le C\frac{e^{\mathfrak B_n}}{\rho_n}
       (\mathfrak m_n+\delta_n^E)
 \longrightarrow0.
 \label{eq:v9-annular-scalar-error}
\end{equation}
The $\mathfrak m_n$ term vanishes by
\eqref{eq:v8-doubly-tame-properties}; the dissipation term is the present
branch assumption.  Combine \eqref{eq:v9-annular-scalar-error} with
\eqref{eq:v9-Reynolds-scalar}.  Since $\mathcal R_n\to\infty$, this proves
\eqref{eq:v9-annular-Reynolds-margin}.

Cauchy--Schwarz gives
\[
 |h_n^{R,\rm ann}|^2
 \le\mathcal A_n^{R,\rm ann}
    \int_{I_n}\int\omega_n|W_n|^2.
\]
The last factor is at most $4R_t/\nu$ by
\eqref{eq:v8-Reynolds-writer-cap}.  Insert
\eqref{eq:v9-annular-Reynolds-margin}.
\end{proof}

\begin{corollary}[Global critical $L^4$ record on the annular window]
\label{cor:v9-global-L4}
On branch \textup{(A2)},
\begin{equation}
 \boxed{
 \Gamma_n\int_{I_n}\|u(t)\|_4^4\,\dd t
 \ge c\nu\mathcal R_n^2.}
 \label{eq:v9-global-L4-lower}
\end{equation}
The constant depends only on the fixed annular multipliers, $R_t$, and the
margin $h_R$.
\end{corollary}

\begin{proof}
Let $T=A^{-1/4}\mathbb P\operatorname{div}$.  On the support of $\Psi_n$,
its symbol has norm at most $C\Gamma_n^{1/2}$, so
\begin{equation}
 \|\Psi_nTF\|_2\le C\Gamma_n^{1/2}\|F\|_2.
 \label{eq:v9-annular-T-bound}
\end{equation}
For the heat covariance
\(
 \mathbb C_{n,t}=P_{\nu(b_n-t)}(u\otimes u)-v_{n,t}\otimes v_{n,t}
\),
heat contraction on $L^2$ and $L^4$ gives
\begin{equation}
 \|\mathbb C_{n,t}\|_2
 \le\|u(t)\|_4^2+\|v_{n,t}\|_4^2
 \le2\|u(t)\|_4^2.
 \label{eq:v9-covariance-L2}
\end{equation}
Since $g_n=T\mathbb C_{n,t}^\circ$ and trace-free projection is contractive,
\[
 \|\Psi_ng_n(t)\|_2^2
 \le C\Gamma_n\|u(t)\|_4^4.
\]
Use $0\le\omega_n\le1$, integrate, and compare with
\eqref{eq:v9-annular-source-action-lower}.
\end{proof}

% =============================================================================
\section{The captured material tube carries a critical local \texorpdfstring{$L^4$}{L4} stock}
\label{sec:v9-material-L4}
% =============================================================================

The preceding corollary is global in space.  The signed annular stress path
from Version~V8 gives the corresponding statement on the actual material
preimage cell.

\begin{lemma}[Square action of the annular stress writer]
\label{lem:v9-writer-square}
There is a fixed $C_M<\infty$ such that
\begin{equation}
 \boxed{
 \int_{I_n}\|M_n(t)\|_2^2\,\dd t
 \le C_M\frac{\Gamma_n}{\nu}.}
 \label{eq:v9-writer-square}
\end{equation}
\end{lemma}

\begin{proof}
The annular support of $\varpi_n$ and
\eqref{eq:v9-stress-writer} give fixed constants $C,c>0$ with
\[
 \|M_n(t)\|_2
 \le C\Gamma_n^{3/2}
       e^{-c\nu(b_n-t)\Gamma_n^2}.
\]
Integrating the square over $0<b_n-t<\infty$ gives
\[
 C^2\Gamma_n^3
 \int_0^\infty e^{-2c\nu s\Gamma_n^2}\,\dd s
 =\frac{C^2}{2c}\frac{\Gamma_n}{\nu}.
\]
Restricting to $I_n$ only decreases the integral.
\end{proof}

\begin{theorem}[Same-tube critical $L^4$ lower bound]
\label{thm:v9-material-L4}
On the captured annular path branch,
\begin{equation}
 \boxed{
 \Gamma_n\int_{I_n}\int_{\mathcal M_n(t)}
       \omega_n(t,x)|u(t,x)|^4\,\dd x\,\dd t
 \ge c_4\nu\mathcal R_n^2}
 \label{eq:v9-material-L4-lower}
\end{equation}
for all late $n$, where $c_4>0$ depends only on
$\eta,h_*$ and the fixed writer constants.
\end{theorem}

\begin{proof}
The source-marginal conclusion
\eqref{eq:v8-source-in-material-cell} for the signed annular path cylinder
means
\begin{equation}
 \frac{1}{V_n^{\rm mc}}
 \int_{I_n}\int_{x\notin\mathcal M_n(t)}
       \omega_n|j_n^{\rm ann}|\longrightarrow0.
 \label{eq:v9-annular-outside-variation}
\end{equation}
By \eqref{eq:v8-variation-ratio}, $V_n^{\rm mc}\le CL_*H_n^{\rm ann}$.
Hence the absolute outside integral is $o(H_n^{\rm ann})$.  Therefore,
using \eqref{eq:v9-captured-stress},
\begin{equation}
 \left|
 \int_{I_n}\int_{\mathcal M_n(t)}
 \omega_nj_n^{\rm ann}
 \right|
 \ge\frac{\eta}{2}H_n^{\rm ann}
 \ge\frac{\eta h_*}{2}\mathcal R_n
 \label{eq:v9-material-stress-margin}
\end{equation}
for all late $n$.

Since trace-free projection is orthogonal in the Frobenius metric,
$|(u\otimes u)^\circ|\le|u|^2$.  Cauchy--Schwarz on the moving spacetime
cell gives
\[
 \begin{aligned}
 \left|\int_{\mathfrak T_n}\omega_n
      (u\otimes u)^\circ:M_n\right|^2
 &\le
 \left(\int_{\mathfrak T_n}\omega_n|u|^4\right)
 \left(\int_{\mathfrak T_n}\omega_n|M_n|^2\right)\\
 &\le
 \left(\int_{\mathfrak T_n}\omega_n|u|^4\right)
 \left(C_M\frac{\Gamma_n}{\nu}\right),
 \end{aligned}
\]
where \Cref{lem:v9-writer-square} was used in the last line.  Insert
\eqref{eq:v9-material-stress-margin} and multiply by $\Gamma_n$.
\end{proof}

\begin{corollary}[Energy-null critical intermittency]
\label{cor:v9-energy-null-intermittency}
On branch \textup{(E2)} of \Cref{thm:v9-energy-atom},
\begin{equation}
 \begin{gathered}
 \mathscr E_n^-+\mathscr E_n^+
 +\nu\mathscr D_n+|\mathscr P_{\omega_n}[a_n,b_n]|
 \longrightarrow0,\\
 \Gamma_n\int_{I_n}\int_{\mathcal M_n(t)}
       \omega_n|u|^4
 \ge c_4\nu\mathcal R_n^2\longrightarrow\infty.
 \end{gathered}
 \label{eq:v9-energy-null-intermittency}
\end{equation}
Thus the routed thread can be null in the complete scalar local-energy packet
while diverging in its scale-invariant quartic product stock.
\end{corollary}

\begin{proof}
The first line is
\Cref{cor:v9-energy-conservation,prop:v9-energy-pressure,thm:v9-energy-atom}.
The second is \Cref{thm:v9-material-L4} and
$\mathcal R_n\to\infty$.
\end{proof}

% =============================================================================
\section{The source-action localization fork is fractional}
\label{sec:v9-fractional-collar}
% =============================================================================

Let
\begin{equation}
 T:=A^{-1/4}\mathbb P\operatorname{div}
 \label{eq:v9-T}
\end{equation}
act on symmetric tensor fields, so that
\(
 g_n=T\mathbb C_{n,t}^\circ
\).
For $t<b_n$, strict positivity of the parabolic transition kernel makes
$\omega_n(t,x)>0$ whenever the terminal cutoff is nonzero.  Put
\begin{equation}
 \zeta_n(t,x):=\sqrt{\omega_n(t,x)}.
 \label{eq:v9-square-root-router}
\end{equation}
The endpoint $t=b_n$ has zero spacetime measure and may be defined by
continuous approximation.

\begin{proposition}[Exact fractional source-localization identity]
\label{prop:v9-fractional-collar}
For every $t\in I_n$,
\begin{equation}
 \boxed{
 \zeta_ng_n
 =T(\zeta_n\mathbb C_{n,t}^\circ)
  +[\zeta_n,T]\mathbb C_{n,t}^\circ.}
 \label{eq:v9-fractional-localization}
\end{equation}
Define
\begin{align}
 \mathcal Q_n^{\rm cov}
 &:=\int_{I_n}
   \|T(\zeta_n\mathbb C_{n,t}^\circ)\|_2^2\,\dd t,
 \label{eq:v9-local-covariance-action}\\
 \mathcal C_n^{\rm frac}
 &:=\int_{I_n}
   \|[\zeta_n,T]\mathbb C_{n,t}^\circ\|_2^2\,\dd t.
 \label{eq:v9-fractional-collar-action}
\end{align}
Then
\begin{equation}
 \boxed{
 \max\{\mathcal Q_n^{\rm cov},\mathcal C_n^{\rm frac}\}
 \ge\frac14\mathcal A_n^{R,\omega}
 \ge\frac{c_R}{4}\nu\mathcal R_n^2.}
 \label{eq:v9-fractional-fork}
\end{equation}
Moreover,
\begin{equation}
 \mathcal Q_n^{\rm cov}
 \le C\int_{I_n}
 \|\zeta_n\mathbb C_{n,t}^\circ\|_{\dot H^{1/2}}^2\,\dd t.
 \label{eq:v9-product-upper}
\end{equation}
Thus the routed source action returns either a fractional localization collar
or a weighted critical covariance-product stock of record size.
\end{proposition}

\begin{proof}
Identity \eqref{eq:v9-fractional-localization} is the definition of the
commutator:
\[
 \zeta_nT\mathbb C
 =T(\zeta_n\mathbb C)+(\zeta_nT-T\zeta_n)\mathbb C.
\]
Take the $L^2(I_n\times\Torus^3)$ norm.  The triangle inequality gives
\[
 (\mathcal A_n^{R,\omega})^{1/2}
 \le(\mathcal Q_n^{\rm cov})^{1/2}
   +(\mathcal C_n^{\rm frac})^{1/2}.
\]
If both terms on the right had square strictly below
$\mathcal A_n^{R,\omega}/4$, their sum would be strictly below the left
side.  This proves the first inequality in
\eqref{eq:v9-fractional-fork}; the second is
\eqref{eq:v9-imported-action}.  Finally, $T$ is a Fourier multiplier of
order $1/2$, so
\(
 \|TF\|_2\le C\|F\|_{\dot H^{1/2}}
\), proving \eqref{eq:v9-product-upper}.
\end{proof}

\begin{assessment}[This is the actual product stop]
\label{ass:v9-product-stop}
The heat and scalar pressure collars made small in the supplied NS--B router
are linear pairings.  The quantity
$\mathcal C_n^{\rm frac}$ is different: it is the square-action cost of
localizing the nonlocal half-order covariance source.  Its vanishing is not
implied by scalar capture.  On its complementary branch,
\eqref{eq:v9-product-upper} identifies the precise critical product norm that
must be compact or uniformly controlled.  Replacing either branch by a local
energy scalar would discard the half derivative carried by the physical
source.
\end{assessment}

% =============================================================================
\section{Enstrophy-square concentration forced by the Reynolds action}
\label{sec:v9-enstrophy-square}
% =============================================================================

Let
\begin{equation}
 D(t):=\|\nabla u(t)\|_2^2.
 \label{eq:v9-enstrophy}
\end{equation}

\begin{theorem}[Record-sized enstrophy-square pulse]
\label{thm:v9-enstrophy-square}
On the material-tube Reynolds-action branch,
\begin{equation}
 \boxed{
 \int_{I_n}D(t)^2\,\dd t
 \ge c_D\nu\mathcal R_n^2}
 \label{eq:v9-enstrophy-square-lower}
\end{equation}
for a fixed $c_D>0$.  Consequently
\begin{equation}
 \boxed{
 \sup_{t\in I_n}D(t)
 \ge c\nu\Gamma_n\mathcal R_n,}
 \label{eq:v9-enstrophy-height-scale}
\end{equation}
and, more sharply in terms of the remaining dissipation mass,
\begin{equation}
 \boxed{
 \sup_{t\in I_n}D(t)
 \ge c_D\frac{\nu\mathcal R_n^2}{\delta_n^E}.}
 \label{eq:v9-enstrophy-height-tail}
\end{equation}
In particular the enstrophy height diverges.  These are fixed-datum temporal
concentration statements; they do not localize the enstrophy to
$\mathcal M_n(t)$ without controlling the fractional collar of
\Cref{prop:v9-fractional-collar}.
\end{theorem}

\begin{proof}
By \Cref{lem:v9-source-length}, pointwise
\[
 \|g_n(t)\|_2
 \le\|c[u(t)]\|_2+\|c[v_{n,t}]\|_2
 \le C D(t),
\]
because heat is contractive in $\dot H^1$.  Hence
\[
 \mathcal A_{n,\mathcal M}^R
 \le\int_{I_n}\|g_n(t)\|_2^2\,\dd t
 \le C\int_{I_n}D(t)^2\,\dd t.
\]
Combine this with \eqref{eq:v8-material-action-lower} to obtain
\eqref{eq:v9-enstrophy-square-lower}.

Since
\(
 \int_{I_n}D^2\le |I_n|(\sup_{I_n}D)^2
\),
\eqref{eq:v9-standing-window} gives
\[
 \sup_{I_n}D
 \ge\left(\frac{c_D\nu\mathcal R_n^2}{|I_n|}\right)^{1/2}
 \ge c\nu\Gamma_n\mathcal R_n.
\]
Also
\(
 \int D^2\le(\sup D)\int D=(\sup D)\delta_n^E
\),
which proves \eqref{eq:v9-enstrophy-height-tail}.  The denominator is
positive for all late $n$, since otherwise $D=0$ almost everywhere on the
window and the source-action lower bound would fail.
\end{proof}

\begin{firewall}[Why this is not a finite-stock contradiction]
\label{fire:v9-enstrophy-no-contradiction}
The Leray class controls $D$ in $L_t^1$, not in $L_t^2$.  The simultaneous
facts
\(
 \int_{I_n}D\to0
\)
and
\(
 \int_{I_n}D^2\to\infty
\)
are compatible through increasingly narrow spikes.  Excluding those spikes
requires a new time-regularity, palinstrophy, one-scale local-energy, or
backward-rigidity estimate.  The theorem identifies the missing stock; it
does not supply it.
\end{firewall}

% =============================================================================
\section{Annular-cell concentration or explicit one-scale microdelocalization}
\label{sec:v9-one-scale}
% =============================================================================

Normalize the positive material-tube quartic measure by
\begin{equation}
 \dd\lambda_n(t,x)
 :=\frac{
  \one_{\mathcal M_n(t)}(x)\omega_n(t,x)|u(t,x)|^4\,\dd x\,\dd t
 }{
  \displaystyle\int_{I_n}\int_{\mathcal M_n(t)}\omega_n|u|^4
 }.
 \label{eq:v9-L4-law}
\end{equation}
For $R>0$, define its annular-scale concentration function
\begin{equation}
 \vartheta_n(R)
 :=\sup_{y\in\Torus^3}
   \lambda_n\bigl(I_n\times B(y,R\ell_n)\bigr).
 \label{eq:v9-concentration-function}
\end{equation}

\begin{theorem}[The explicit one-scale defect]
\label{thm:v9-one-scale-defect}
After passage to a subsequence, exactly one of the following occurs.
\begin{enumerate}[label=\textup{(S\arabic*)},leftmargin=3.4em]
\item There are $R_0<\infty$, $\vartheta_0>0$, and centres $y_n$ such that
      \begin{equation}
       \lambda_n(I_n\times B(y_n,R_0\ell_n))\ge\vartheta_0.
       \label{eq:v9-one-scale-heavy}
      \end{equation}
      Then
      \begin{equation}
       \boxed{
       \Gamma_n\int_{I_n}\int_{B(y_n,R_0\ell_n)}|u|^4
       \ge c\nu\mathcal R_n^2.}
       \label{eq:v9-one-scale-L4}
      \end{equation}
      Under the unit-viscosity rescaling
      \begin{equation}
       U_n(s,z)
       :=\frac{\ell_n}{\nu}
       u\!\left(b_n+\frac{\ell_n^2}{\nu}s,
                y_n+\ell_n z\right),
       \label{eq:v9-unit-rescaling}
      \end{equation}
      the corresponding fixed spatial ball and bounded backward-time strip
      satisfy
      \begin{equation}
       \boxed{
       \int |U_n|^4\,\dd z\,\dd s
       \ge c\frac{\mathcal R_n^2}{\nu^2}
       \longrightarrow\infty.}
       \label{eq:v9-rescaled-L4-divergence}
      \end{equation}
      Hence any passage of the quartic product or Reynolds source action
      requires a concentration defect or an additional amplitude
      normalization; uniform local $L^4$ compactness is impossible on this
      subpacket.
\item For every fixed $R<\infty$,
      \begin{equation}
       \boxed{\vartheta_n(R)\longrightarrow0.}
       \label{eq:v9-one-scale-microdelocalization}
      \end{equation}
      This is the explicit annular-scale spatial microdelocalization defect
      inside one path-concentrated material tube.
\end{enumerate}
\end{theorem}

\begin{proof}
For each integer $m\ge1$, pass diagonally so that
$\vartheta_n(m)$ converges.  If one limit is positive, choose $R_0=m$ and
centres within half of the supremum; this proves
\eqref{eq:v9-one-scale-heavy}.  Otherwise every fixed integer, and hence
every fixed $R$, has zero limit.

On the first branch, multiply the fraction
\eqref{eq:v9-one-scale-heavy} by the total lower bound
\eqref{eq:v9-material-L4-lower}; drop the factors
$\one_{\mathcal M_n}\omega_n\le1$ to obtain
\eqref{eq:v9-one-scale-L4}.  Under \eqref{eq:v9-unit-rescaling},
\[
 \int |u|^4\,\dd x\,\dd t
 =\nu^3\ell_n\int|U_n|^4\,\dd z\,\dd s.
\]
Since $\Gamma_n=\ell_n^{-1}$, division by $\nu^3$ in
\eqref{eq:v9-one-scale-L4} gives
\eqref{eq:v9-rescaled-L4-divergence}.  The rescaled time interval is
contained in $[-R_t,0]$ by \eqref{eq:v9-standing-window}.
\end{proof}

\begin{assessment}[The two spatial notions are now separated]
\label{ass:v9-path-versus-one-scale}
Version~V8 path concentration says that the complete material tube collapses
to the constant singular path in the original coordinates.  Branch
\textup{(S2)} says that, after magnification by the physical annular scale,
its critical quartic stock is nevertheless diffuse through every bounded
spatial head.  There is no contradiction: the tube radius shrinks while
$\rho_n/\ell_n\to\infty$.  This is precisely the aspect-ratio obstruction of
\Cref{thm:v9-aspect-ratio}.
\end{assessment}

% =============================================================================
\section{Absolute energy-class smallness cannot control the critical stock}
\label{sec:v9-interpolation-nogo}
% =============================================================================

The actual terminal branch already separates vanishing scalar energy from a
divergent critical $L^4$ stock.  The following elementary rescaling shows
why this separation cannot be repaired by an abstract interpolation modulus
using only absolute local energy and integrated dissipation.

\begin{countertheorem}[No absolute energy-class modulus for critical $L^4$]
\label{cth:v9-energy-interpolation}
There are smooth compactly supported divergence-free vector fields $V_N$ on
$(-1,0)\times\Torus^3$ and scales $r_N=N^{-1}\downarrow0$ such that
\begin{equation}
 \sup_t\|V_N(t)\|_2^2\longrightarrow0,
 \qquad
 \int\|\nabla V_N\|_2^2\,\dd t\longrightarrow0,
 \label{eq:v9-test-energy-small}
\end{equation}
but
\begin{equation}
 \boxed{
 r_N^{-1}\int|V_N|^4\,\dd x\,\dd t=Q_0>0}
 \label{eq:v9-test-L4-fixed}
\end{equation}
for every $N$.  Consequently no function
$\Omega(a,b)\to0$ as $(a,b)\to(0,0)$ can give a universal estimate
\begin{equation}
 r^{-1}\int_{Q_r}|V|^4
 \le\Omega\!\left(
   \sup_t\int_{B_r}|V|^2,
   \int_{Q_r}|\nabla V|^2\right)
 \label{eq:v9-impossible-modulus}
\end{equation}
for all smooth divergence-free test fields.
\end{countertheorem}

\begin{proof}
Choose a nonzero
$V\in C_c^\infty((-1,0)\times B(0,1/8);\R^3)$ with
$\operatorname{div}V=0$; for example take a compactly supported scalar
potential and its planar curl.  Fix a coordinate ball about $x_0$ in the
torus.  For large integers $N$, define inside that chart
\begin{equation}
 V_N(t,x)=N V(N^2t,N(x-x_0))
 \label{eq:v9-test-rescaling}
\end{equation}
and extend by zero.  Compact support makes the extension smooth.  Direct
changes of variables give
\begin{align}
 \sup_t\|V_N(t)\|_2^2
 &=N^{-1}\sup_s\|V(s)\|_2^2,
 \label{eq:v9-test-L2}\\
 \int\|\nabla V_N\|_2^2\,\dd t
 &=N^{-1}\int\|\nabla V\|_2^2\,\dd s,
 \label{eq:v9-test-H1}\\
 \int|V_N|^4\,\dd x\,\dd t
 &=N^{-1}\int|V|^4\,\dd y\,\dd s.
 \label{eq:v9-test-L4}
\end{align}
Since $r_N=N^{-1}$, the last formula is
\eqref{eq:v9-test-L4-fixed} with
$Q_0=\int|V|^4>0$.  If \eqref{eq:v9-impossible-modulus} held, its right side
would tend to zero by the first two formulas while its left side stayed
$Q_0$.
\end{proof}

\begin{firewall}[This is an interpolation no-go, not a solution counterexample]
\label{fire:v9-test-not-solution}
The fields $V_N$ are not asserted to solve Navier--Stokes.  The theorem shows
only that the two absolute energy-class stocks which vanish in
\Cref{cor:v9-energy-conservation} cannot, by functional interpolation alone,
control a scale-invariant quartic quantity.  Classical epsilon-regularity
uses scale-normalized velocity and pressure packets; it is not contradicted.
The missing one-scale normalization is exactly the issue exposed in
\Cref{thm:v9-aspect-ratio}.
\end{firewall}

% =============================================================================
\section{Integrated Version V9 theorem and stopping boundary}
\label{sec:v9-integrated}
% =============================================================================

\begin{theorem}[Version V9 material-tube PDE theorem]
\label{thm:v9-integrated}
Preserving every theorem and limitation of Versions~V1--V8, Version~V9 proves
the following on the path-concentrated efficient Reynolds branch.
\begin{enumerate}[label=\textup{(V9.\arabic*)},leftmargin=5em]
\item The material--caloric weight gives the exact local energy identity
      \eqref{eq:v9-routed-energy-identity}.  Its pressure boundary current
      and integrated dissipation vanish on the doubly tame terminal cards,
      so the routed kinetic energy changes by $o(1)$.
\item The scalar local-energy identity is exactly the trace of the routed
      kinetic-stress equation.  The physical pressure--strain tensor is
      trace free and remains an independent anisotropic product row even
      after the pressure boundary current vanishes.
\item The terminal routed kinetic energy either yields an atom at
      $x_*\in\Sigma_T$ in a subsequential kinetic-energy concentration
      measure, or both routed endpoint energies vanish.
\item With $\ell_n=\Gamma_n^{-1}$, the material tube has spatial aspect
      $\rho_n/\ell_n\to\infty$ but temporal depth
      $\nu|I_n|/\rho_n^2\to0$.  No radius makes the present packet a
      nondegenerate parabolic cylinder in both variables.
\item Either the route--dissipation ratio
      $e^{\mathfrak B_n}\delta_n^E/\rho_n$ survives, or the Reynolds source
      itself has a positive output-annular action
      $\gtrsim\nu\mathcal R_n^2$ on the same router.
\item The captured annular stress forces the same material tube to carry
      the critical stock \eqref{eq:v9-material-L4-lower}.  On the no-atom
      branch this stock diverges while the complete scalar local-energy
      packet vanishes.
\item Localizing the positive covariance source in square action gives the
      exact fractional-collar/product alternative
      \eqref{eq:v9-fractional-fork}.  The product branch carries a weighted
      $\dot H^{1/2}$ covariance stock of record size.
\item The positive Reynolds source action forces the enstrophy-square pulse
      \eqref{eq:v9-enstrophy-square-lower} and the explicit enstrophy-height
      lower bounds \eqref{eq:v9-enstrophy-height-scale}--
      \eqref{eq:v9-enstrophy-height-tail}.
\item The material-tube quartic law either selects one annular-scale cell,
      on which the unit-viscosity rescaled $L^4$ norm diverges, or returns the
      explicit one-scale microdelocalization
      \eqref{eq:v9-one-scale-microdelocalization}.
\item Absolute local energy and integrated dissipation cannot control this
      critical quartic stock by interpolation; the scale-normalized pressure,
      product, or compactness row is genuinely additional.
\end{enumerate}
No item proves terminal trace, epsilon regularity, strong product
compactness, backward uniqueness, either ancient Liouville theorem, or global
three-dimensional regularity.
\end{theorem}

\begin{proof}
Items~\textup{(V9.1)--(V9.2)} are
\Cref{thm:v9-routed-energy,prop:v9-energy-pressure,cor:v9-energy-conservation,prop:v9-trace-blindness}.
Item~\textup{(V9.3)} is \Cref{thm:v9-energy-atom}.  Item~\textup{(V9.4)} is
\Cref{thm:v9-aspect-ratio}.  Item~\textup{(V9.5)} is
\Cref{thm:v9-annular-source-action}.  Item~\textup{(V9.6)} is
\Cref{thm:v9-material-L4,cor:v9-energy-null-intermittency}.  Item
\textup{(V9.7)} is \Cref{prop:v9-fractional-collar}.  Item
\textup{(V9.8)} is \Cref{thm:v9-enstrophy-square}.  Item
\textup{(V9.9)} is \Cref{thm:v9-one-scale-defect}.  The final item is
\Cref{cth:v9-energy-interpolation}.
\end{proof}

\begingroup\small
\begin{longtable}{p{0.25\textwidth}p{0.19\textwidth}p{0.48\textwidth}}
\caption{NS--A Version V9 proof-status ledger.}\label{tab:v9-ledger}\\
\toprule Row & Status & Exact conclusion\\\midrule
\endfirsthead
\toprule Row & Status & Exact conclusion\\\midrule
\endhead
V1--V8 prefix & \PROVED{} preserved
 &The complete V8 preterminal source is retained byte-for-byte; no preceding
 theorem is rewritten.\\[0.3em]
Material--caloric local energy & \PROVED{} exact
 &Time, transport, and viscous localization cancel; only the physical pressure
 boundary current remains.\\[0.3em]
Terminal pressure boundary current & \PROVED{} vanishing
 &The doubly tame ratio
 $e^{\mathfrak B_n}\mathfrak l_n/\rho_n$ makes the scalar current $o(1)$.\\[0.3em]
Routed dissipation & \PROVED{} vanishing
 &It is bounded by the fixed solution's terminal Leray dissipation tail.\\[0.3em]
Pressure--strain & \OPEN{} independent anisotropic row
 &It is trace free and invisible to scalar energy; the tame boundary-current
 estimate does not control it.\\[0.3em]
Terminal local energy & \PROVED{} exact alternative
 &A nonzero endpoint packet gives an atom at $x_*$; otherwise both routed
 endpoint energies vanish.\\[0.3em]
Parabolic aspect ratio & \OBSTRUCTION{} exact
 &$\rho_n/\ell_n\to\infty$ while
 $\nu|I_n|/\rho_n^2\to0$; the tube is not one parabolic cell.\\[0.3em]
Route--dissipation collar & \OBSTRUCTION{} typed
 &Failure of $e^{\mathfrak B_n}\delta_n^E/\rho_n\to0$ is retained before
 annular source promotion.\\[0.3em]
Output-annular Reynolds action & \PROVED{} on the zero-collar branch
 &The routed scalar and action remain of record size after the literal annular
 projection.\\[0.3em]
Material-tube critical $L^4$ & \PROVED{} lower packet
 &The captured signed stress forces
 $\Gamma_n\int_{\mathfrak T_n}\omega_n|u|^4\gtrsim\nu\mathcal R_n^2$.\\[0.3em]
Energy-null intermittency & \PROVED{} alternative
 &In the no-atom branch, scalar energy and dissipation vanish while the critical
 quartic stock diverges.\\[0.3em]
Fractional source collar & \OBSTRUCTION{} exact
 &Square-action localization of $A^{-1/4}P\operatorname{div}$ returns the
 commutator norm $\mathcal C_n^{\rm frac}$.\\[0.3em]
Covariance product stock & \PROVED{} complementary lower packet
 &If the fractional collar is small, the weighted
 $\dot H^{1/2}$ covariance product is of record size.\\[0.3em]
Enstrophy-square pulse & \PROVED{} fixed datum
 &$\int_{I_n}\|\nabla u\|_2^4\gtrsim\nu\mathcal R_n^2$ and the enstrophy
 height obeys two explicit lower bounds.\\[0.3em]
Annular one-cell capture & \CONDITIONAL{} exact fork
 &One cell has divergent rescaled $L^4$, or every bounded annular head has
 vanishing fraction.\\[0.3em]
Absolute energy interpolation & \OBSTRUCTION{} scaling counterexample
 &Vanishing absolute $L_t^\infty L_x^2$ and $L_t^2\dot H_x^1$ stocks do not
 force the critical $r^{-1}L^4$ stock to vanish.\\[0.3em]
One-scale pressure/product compactness & \OPEN{} first PDE input
 &Required before a suitable or source-action profile can be passed to a
 limit.\\[0.3em]
Terminal trace / Liouville / regularity & \OPEN
 &No ancient class is excluded and no global regularity conclusion follows.\\
\bottomrule
\end{longtable}
\endgroup

\begin{programme}[Version V10: one-scale anisotropic product or stop]
\label{prog:v9-next}
Preserve Versions~V1--V9 verbatim.  Do not introduce another path law,
material radius, age distribution, response graph, or scalar local-energy
estimate.

Proceed only in the following order.
\begin{enumerate}[label=\textup{(V10.\arabic*)},leftmargin=5.4em]
\item Test the actual fractional collar
      $\mathcal C_n^{\rm frac}$ of
      \eqref{eq:v9-fractional-collar-action}.  If it carries a fixed action
      fraction, retain it as the nonlocal source-localization obstruction and
      stop on that subsequence.
\item On the covariance-product branch, use the concentration function
      \eqref{eq:v9-concentration-function}.  If
      \eqref{eq:v9-one-scale-microdelocalization} holds, retain that explicit
      material-tube/annular-scale defect; do not create another abstract
      concentration trichotomy.
\item On one selected annular cell, evaluate the scale-normalized velocity,
      pressure, and covariance-product packet on the same physical history.
      Pass a profile only if those products are compact.  Divergent rescaled
      $L^4$, pressure--strain, or source action is the named defect.
\item Transport every represented response and return vector through its
      declared endpoint map before source novelty is asserted.  The positive
      covariance and the paid stress remain one source occurrence.
\item Stop at the first unavailable epsilon-regularity, terminal-trace,
      backward-uniqueness, or Liouville theorem.  A later scalar or positive
      Gram cannot repair that missing PDE row.
\end{enumerate}
The admissible endpoint is
\[
 \boxed{
 \begin{gathered}
 \text{fractional collar}
 \ \big|\ 
 \text{annular microdelocalization}
 \ \big|\ 
 \text{one noncompact product cell}\\
 \big|\ 
 \text{one compact nonzero profile}
 \ \big|\ 
 \text{named trace/rigidity stop}
 \end{gathered}}
\]
\end{programme}

\begin{assessment}[Append-only integrity]
\label{ass:v9-integrity}
The complete Version~V8 source preceding its sole final document terminator
is preserved verbatim.  The present material begins at
Section~\ref{sec:v9-entry}; the final command below is again unique.
\end{assessment}

\addtocontents{toc}{\protect\endgroup}
% =============================================================================
% END VERSION V9 MATERIAL -- append later versions before final terminator
% =============================================================================


\clearpage
% =============================================================================
% VERSION V10: NEW MATERIAL.  The complete V1--V9 source above is preserved
% verbatim; only its sole final document terminator was removed before append.
% =============================================================================
\hypersetup{pdftitle={NS-A Terminal Source Concentration Programme: Cumulative V10},pdfauthor={A. Pelissier}}
\makeatletter
\addtocontents{toc}{\protect\begingroup\protect\makeatletter}
\addtocontents{toc}{\protect\def\protect\l@section{\protect\@dottedtocline{1}{0em}{4.7em}}}
\addtocontents{toc}{\protect\def\protect\l@subsection{\protect\@dottedtocline{2}{1.5em}{5.3em}}}
\addtocontents{toc}{\protect\makeatother}
\makeatother
\renewcommand{\thesection}{V10.\arabic{section}}
\renewcommand{\thesubsection}{V10.\arabic{section}.\arabic{subsection}}
\renewcommand{\theequation}{V10.\arabic{equation}}
\setcounter{section}{0}
\setcounter{equation}{0}
\renewcommand{\theHsection}{V10.\arabic{section}}
\renewcommand{\theHsubsection}{V10.\arabic{section}.\arabic{subsection}}
\renewcommand{\theHtheorem}{V10.\arabic{section}.\arabic{theorem}}
\renewcommand{\theHequation}{V10.\arabic{equation}}

\begin{center}
{\LARGE\bfseries NS--A: Version V10}\\[0.5em]
{\large The Source--Pressure Hodge Polarization, Target-Native Annular\\
Pseudolocality, and the First Nonzero Covariance-Profile Stop}\\[0.5em]
{\normalsize 3 September 2026}
\end{center}
\addcontentsline{toc}{part}{Version V10: source--pressure polarization, annular pseudolocality, and covariance profile}

\begin{abstract}
Version~V9 reached a genuine one-scale obstruction: on the selected material
thread, scalar energy and dissipation may vanish while an output-annular
Reynolds source action and a critical quartic product diverge.  Its final
fractional localization fork was algebraically correct, but it was formed
from the complete covariance tensor.  The present continuation moves one
arrow upstream and audits which tensor directions the physical source
operator actually sees before any localized product is priced.

For every nonzero Fourier mode, the trace-free symmetric tensor fibre has an
orthogonal decomposition of ranks $2+1+2$.  The rank-two source-visible
summand is the unique minimum tensor preimage of the pressure-shorted source
$A^{-1/4}P\operatorname{div}$; the rank-one summand is the pure gradient
pressure follower; and the remaining rank-two summand has zero divergence.
The exact source norm is one half of the $\dot H^{1/2}$ norm of the
source-visible tensor.  This identifies the Moore--Penrose inverse and shows
that neither a pressure follower nor a divergence-null stress may be charged
as a Reynolds source.

The same calculation exposes a target-native correction to the V9 collar.
A positive covariance may contain arbitrarily large source-null tensor
components which leave the selected annular source unchanged but make both
terms of the unprojected fractional triangle decomposition arbitrarily
large.  The selected scalar must therefore first be projected to its output
annulus and then to the source-visible tensor quotient.  On that quotient an
annular commutator theorem gives a factor
$\|\nabla\omega_n\|_\infty/\Gamma_n$.  The doubly tame material route makes
this factor vanish.  Hence either the unweighted annular source action itself
overdrives the record scale, or the output-localized source has a genuine
source-visible covariance-product action of order $\nu\mathcal R_n^2$ with
no surviving fractional collar.

Finally, the note rescales a selected annular cell by its physical wavelength
and by the record amplitude.  If normalized source action, velocity product,
and pressure fail to be bounded, their precise overdrive is retained.  On
the bounded branch, the scaled interval has nonzero length, the annular
writer is locally compact, and the weighted source converges weakly to a
nonzero vector $G$ of the form
$G=A^{-1/4}P\operatorname{div}F$ for a nonzero source-visible tensor profile
$F$.  Bounded velocity products additionally yield an Euler--Reynolds
constraint $\operatorname{div}Q+\nabla\Pi=0$ and a positive product defect
$Q-V\otimes V\succeq0$.  This is a covariance/source profile, not yet an
ancient unit-viscosity Navier--Stokes solution.  Product compactness,
terminal trace, backward uniqueness, and the relevant Liouville theorem
remain the exact downstream PDE rows.  No global regularity conclusion is
asserted.
\end{abstract}

% =============================================================================
\section{Upstream scope: project to the physical source before localizing}
\label{sec:v10-scope}
% =============================================================================

Retain the path-concentrated, bounded-sweeping, efficient annular branch and
all notation of Version~V9.  In particular,
\begin{equation}
 g_n=T\mathbb C_{n,t}^{\circ},
 \qquad
 T=A^{-1/4}P\operatorname{div},
 \qquad
 g_n^{\rm ann}:=\Psi_ng_n,
 \label{eq:v10-standing-source}
\end{equation}
where $\Psi_n=\psi(\Lambda/\Gamma_n)$ is supported in
$[\Gamma_n/4,4\Gamma_n]$.  The selected annular scalar satisfies
\begin{equation}
 h_n^{R,\rm ann}
 =-\int_{I_n}\!\int_{\Torus^3}
   \omega_ng_n^{\rm ann}\cdot W_n,
 \qquad
 |h_n^{R,\rm ann}|\ge\frac{h_R}{2}\mathcal R_n,
 \label{eq:v10-standing-scalar}
\end{equation}
and
\begin{equation}
 |I_n|\le\frac{R_t}{\nu\Gamma_n^2},
 \qquad
 \|W_n\|_2\le2\Gamma_n,
 \qquad
 \frac{e^{\mathfrak B_n}}{\Gamma_n\rho_n}\longrightarrow0.
 \label{eq:v10-standing-tame}
\end{equation}
The last limit is the final ratio in
\eqref{eq:v8-doubly-tame-properties}.

The first issue is logical rather than technical.  The V9 identity
\eqref{eq:v9-fractional-localization} used the complete tensor
$\mathbb C_{n,t}^{\circ}$.  A large norm of either term on its right need not
belong to the selected output source: multiplication by the router can turn a
source-null tensor into a nonzero source and the two resulting terms then
cancel exactly.  The next countertheorem makes this phenomenon explicit.

\begin{countertheorem}[The complete fractional collar is not target native]
\label{cth:v10-full-collar-not-target}
On $\Torus^3=(\R/2\pi\mathbb Z)^3$, let
\begin{equation}
 D=\operatorname{diag}(0,1,-1),
 \qquad
 H(x)=D\cos x_1,
 \qquad
 \zeta(x)=\frac{2+\cos x_2}{3}.
 \label{eq:v10-null-example}
\end{equation}
Then $H$ is real, symmetric, trace free, and
\begin{equation}
 \boxed{TH=0,\qquad T(\zeta H)\ne0.}
 \label{eq:v10-null-mixing}
\end{equation}
For every $N\ge1$, the matrix field
\begin{equation}
 \mathbb C_N=(N+1)I_3+NH
 \label{eq:v10-positive-null-covariance}
\end{equation}
is pointwise positive definite and has
$\mathbb C_N^{\circ}=NH$.  Consequently
\begin{equation}
 \zeta T\mathbb C_N^{\circ}=0,
 \qquad
 T(\zeta\mathbb C_N^{\circ})=NT(\zeta H),
 \qquad
 [\zeta,T]\mathbb C_N^{\circ}=-NT(\zeta H).
 \label{eq:v10-cancelled-collar}
\end{equation}
Thus the left side of the V9 localization identity vanishes while each term
on the right has squared norm $N^2\|T(\zeta H)\|_2^2$.

More generally, after adding a sufficiently large scalar identity, one may
add $NH$ to any bounded positive tensor packet without changing its physical
source, while making the complete fractional collar arbitrarily large.
\end{countertheorem}

\begin{proof}
The only Fourier modes of $H$ are $k=\pm e_1$.  Since $De_1=0$,
\[
 P_k\widehat H(k)k=0,
\]
and hence $TH=0$.  The product $\zeta H$ contains a nonzero Fourier
coefficient at $k=e_1+e_2$.  With
$n=(e_1+e_2)/\sqrt2$,
\[
 Dn=(0,2^{-1/2},0),
 \qquad
 (I-n\otimes n)Dn
 =\frac{1}{2\sqrt2}(-1,1,0)\ne0.
\]
Therefore $T(\zeta H)$ has a nonzero coefficient at that mode.

The eigenvalues of $H(x)$ are
$0,\cos x_1,-\cos x_1$.  Hence every eigenvalue of
$(N+1)I_3+NH(x)$ is at least one.  Since $H$ is trace free, the trace-free
part of \eqref{eq:v10-positive-null-covariance} is $NH$.  The three
identities in \eqref{eq:v10-cancelled-collar} now follow from $TH=0$ and the
definition $[\zeta,T]=\zeta T-T\zeta$.

For the final statement, if $\mathbb C_0$ is bounded and positive, choose a
scalar $a_N$ larger than the operator norm of the trace-free perturbation
$NH$.  Adding $a_NI_3+NH$ preserves positivity, contributes only $NH$ to the
trace-free part, and leaves $T\mathbb C_0^{\circ}$ unchanged because $TH=0$.
\end{proof}

\begin{firewall}[What is corrected and what remains valid]
\label{fire:v10-v9-collar-scope}
Proposition~\ref{prop:v9-fractional-collar} is an exact triangle identity and
remains valid.  Countertheorem~\ref{cth:v10-full-collar-not-target} shows
that the size of its two right-hand terms is not a source-minimal diagnostic
for the selected annular scalar.  A target-native collar must be formed after
both output-frequency selection and the pressure-shorted source quotient.
No conclusion of Versions~V1--V9 which used only the left-hand source action
is altered.
\end{firewall}

% =============================================================================
\section{Exact source--pressure--null Hodge polarization of a tensor mode}
\label{sec:v10-hodge-polar}
% =============================================================================

Let $\operatorname{Sym}_0(3;\C)$ denote the complex trace-free symmetric
$3\times3$ matrices with the Frobenius inner product.  For a unit vector
$n\in\R^3$, put
\begin{equation}
 P_n=I-n\otimes n,
 \qquad
 a\odot b:=\frac12(a\otimes b+b\otimes a),
 \qquad
 E_n:=\frac32\left(n\otimes n-\frac13I\right).
 \label{eq:v10-mode-notation}
\end{equation}
For $F\in\operatorname{Sym}_0(3;\C)$ define
\begin{align}
 \mathsf P_{\rm sv}(n)F
 &:=2n\odot(P_nFn),
 \label{eq:v10-Psv}\\
 \mathsf P_{\rm gr}(n)F
 &:=(n^{\mathsf T}Fn)E_n,
 \label{eq:v10-Pgr}\\
 \mathsf P_{\rm dn}(n)F
 &:=F-\mathsf P_{\rm sv}(n)F-\mathsf P_{\rm gr}(n)F.
 \label{eq:v10-Pdn}
\end{align}
The labels mean source visible, pure gradient, and divergence null.

\begin{theorem}[Rank $2+1+2$ source--pressure Hodge polarization]
\label{thm:v10-mode-polarization}
For every unit $n$, the three maps in
\eqref{eq:v10-Psv}--\eqref{eq:v10-Pdn} are mutually orthogonal projections
on $\operatorname{Sym}_0(3;\C)$ of ranks $2,1,2$, respectively.  Their
ranges are characterized by
\begin{align}
 \operatorname{Ran}\mathsf P_{\rm sv}(n)
 &=\{2n\odot v:v\perp n\},
 \label{eq:v10-source-visible-range}\\
 \operatorname{Ran}\mathsf P_{\rm gr}(n)
 &=\C E_n,
 \label{eq:v10-gradient-range}\\
 \operatorname{Ran}\mathsf P_{\rm dn}(n)
 &=\{F\in\operatorname{Sym}_0(3;\C):Fn=0\}.
 \label{eq:v10-div-null-range}
\end{align}
Writing
\begin{equation}
 v=P_nFn,
 \qquad q=n^{\mathsf T}Fn,
 \label{eq:v10-vq}
\end{equation}
one has the exact norm decomposition
\begin{equation}
 \boxed{
 |F|^2=2|v|^2+\frac32|q|^2
       +|\mathsf P_{\rm dn}(n)F|^2.}
 \label{eq:v10-mode-Pythagoras}
\end{equation}

For a nonzero Fourier mode $k$, put $n_k=k/|k|$.  Then
\begin{equation}
 \widehat{TF}(k)
 =i|k|^{1/2}P_{n_k}\widehat F(k)n_k,
 \label{eq:v10-T-symbol}
\end{equation}
so
\begin{equation}
 \boxed{T=T\mathsf P_{\rm sv},
 \qquad
 |\widehat{TF}(k)|^2
 =\frac{|k|}{2}
  |\mathsf P_{\rm sv}(n_k)\widehat F(k)|^2.}
 \label{eq:v10-source-polar-norm}
\end{equation}
The gradient component satisfies
\begin{equation}
 \widehat{\operatorname{div}\mathsf P_{\rm gr}F}(k)
 =ik\,q_k,
 \qquad q_k=n_k^{\mathsf T}\widehat F(k)n_k,
 \label{eq:v10-gradient-divergence}
\end{equation}
whereas
\begin{equation}
 \operatorname{div}\mathsf P_{\rm dn}F=0.
 \label{eq:v10-null-divergence}
\end{equation}
Thus $\mathsf P_{\rm gr}$ is exactly the pressure follower removed by the
Leray projection, and $\mathsf P_{\rm dn}$ is invisible even before that
projection.
\end{theorem}

\begin{proof}
Let $v=P_nFn$.  Since $F$ is symmetric,
\[
 F:(n\odot w)=w\cdot Fn=w\cdot v
 \qquad(w\perp n).
\]
Moreover,
\[
 |n\odot w|^2=\frac12|w|^2,
 \qquad
 (2n\odot v)n=v.
\]
It follows that $F\mapsto2n\odot(P_nFn)$ is the orthogonal projection onto
$\{2n\odot w:w\perp n\}$ and has rank two.

The matrix $E_n$ is trace free, satisfies $E_nn=n$, and has eigenvalues
$1,-1/2,-1/2$.  Hence
\[
 |E_n|^2=\frac32,
 \qquad
 F:E_n=\frac32n^{\mathsf T}Fn.
\]
Therefore $F\mapsto(n^{\mathsf T}Fn)E_n$ is the orthogonal projection onto
$\C E_n$.  It is orthogonal to the source-visible range because
$E_n:(n\odot v)=v\cdot E_nn=v\cdot n=0$.

For the remainder $F_{\rm dn}=\mathsf P_{\rm dn}(n)F$,
\[
 F_{\rm dn}n
 =Fn-P_nFn-(n^{\mathsf T}Fn)n=0.
\]
Conversely, every trace-free symmetric $F$ with $Fn=0$ is orthogonal to both
preceding ranges.  Its dimension is two, so the three projections exhaust
the five-dimensional fibre.  Their squared norms are
$2|v|^2$, $(3/2)|q|^2$, and $|F_{\rm dn}|^2$, proving
\eqref{eq:v10-mode-Pythagoras}.

The Fourier symbol of $\operatorname{div}$ is $iF(k)k$.  The symbols of
$A^{-1/4}$ and the Leray projection are $|k|^{-1/2}$ and $P_{n_k}$.
This proves \eqref{eq:v10-T-symbol}; inserting the source-visible norm gives
\eqref{eq:v10-source-polar-norm}.  Finally,
$\mathsf P_{\rm gr}F\,n=q n$ and
$\mathsf P_{\rm dn}F\,n=0$, which proves
\eqref{eq:v10-gradient-divergence}--\eqref{eq:v10-null-divergence}.
\end{proof}

Extend the three mode projections to tensor Fourier multipliers on
$\Torus^3$, assigning value zero at $k=0$; denote them by the same symbols.
They commute with every scalar Fourier multiplier, including $\Lambda$ and
$\Psi_n$.

\begin{corollary}[Exact source norm and Moore--Penrose inverse]
\label{cor:v10-source-MP}
For every mean-zero symmetric trace-free tensor field $F$,
\begin{equation}
 \boxed{
 \|TF\|_2^2
 =\frac12\|\Lambda^{1/2}\mathsf P_{\rm sv}F\|_2^2.}
 \label{eq:v10-global-source-norm}
\end{equation}
On divergence-free vector fields,
\begin{equation}
 \boxed{TT^*=\frac12\Lambda,
 \qquad
 T^\dagger=2T^*\Lambda^{-1},
 \qquad
 T^\dagger T=\mathsf P_{\rm sv}.}
 \label{eq:v10-T-MP}
\end{equation}
Thus $T^\dagger g$ is the unique minimum-$L^2$ tensor preimage of a source
$g\in\operatorname{Ran}T$.  Modewise,
\begin{equation}
 \widehat{T^\dagger g}(k)
 =-2i|k|^{-1/2}n_k\odot\widehat g(k).
 \label{eq:v10-Tdagger-symbol}
\end{equation}

If $p_F$ is the mean-zero pressure follower defined by
\begin{equation}
 -\Delta p_F=\partial_i\partial_jF_{ij},
 \label{eq:v10-pressure-follower}
\end{equation}
then $\widehat p_F(k)=-q_k$ and
\begin{equation}
 \boxed{
 \|p_F\|_{\dot H^{1/2}}^2
 =\frac23\|\Lambda^{1/2}\mathsf P_{\rm gr}F\|_2^2.}
 \label{eq:v10-pressure-gradient-norm}
\end{equation}
The Reynolds source action controls neither this pressure follower nor the
divergence-null tensor component.
\end{corollary}

\begin{proof}
Sum \eqref{eq:v10-source-polar-norm} over Fourier modes.  For a
divergence-free vector $g_k\perp n_k$, the tensor adjoint has symbol
\[
 \widehat{T^*g}(k)=-i|k|^{1/2}n_k\odot g_k.
\]
Since $P_{n_k}(n_k\odot g_k)n_k=g_k/2$, one obtains
$TT^*g=(|k|/2)g$ modewise.  The Moore--Penrose formulas follow, as does
\eqref{eq:v10-Tdagger-symbol}.  Finally, Fourier transformation of
\eqref{eq:v10-pressure-follower} gives $\widehat p_F=-q_k$, and
$|\mathsf P_{\rm gr}F|^2=(3/2)|q_k|^2$ gives
\eqref{eq:v10-pressure-gradient-norm}.
\end{proof}

\begin{assessment}[The exact tensor rows carried by the PDE]
\label{ass:v10-three-tensor-rows}
The decomposition is not another donor table.  It is the polar geometry of
the already assembled continuum source map.  The three rows have different
physical meanings:
\[
 \boxed{
 \begin{array}{c|c|c}
 \mathsf P_{\rm sv}F&\text{pressure-shorted source}&TF\\
 \mathsf P_{\rm gr}F&\text{gradient follower}&p_F\\
 \mathsf P_{\rm dn}F&\text{divergence-null stress}&0
 \end{array}}
\]
Only the first row is priced by the positive Reynolds source action.  A
localized cutoff may mix the rows, but that mixing is a router commutator and
must not be interpreted as a second physical source birth.
\end{assessment}

% =============================================================================
\section{Annular commutators on the source-visible quotient}
\label{sec:v10-annular-commutator}
% =============================================================================

The next estimate uses only the fixed output annulus.  We record a periodic
commutator lemma in the precise form needed below.

\begin{lemma}[Annular multiplier commutator below order one]
\label{lem:v10-annular-calderon}
Let $M(D)$ be a matrix Fourier multiplier on $\Torus^3$ whose symbol is
smooth away from the origin and homogeneous of degree $s<1$.  Let $F$ have
Fourier support in
\begin{equation}
 c_0\Gamma\le|k|\le C_0\Gamma,
 \qquad \Gamma\ge1.
 \label{eq:v10-annular-support}
\end{equation}
For every Lipschitz scalar $a$,
\begin{equation}
 \boxed{
 \|[a,M(D)]F\|_2
 \le C\Gamma^{s-1}\|\nabla a\|_\infty\|F\|_2.}
 \label{eq:v10-general-commutator}
\end{equation}
If $M(D)$ has degree zero, then the stronger output estimate
\begin{equation}
 \boxed{
 \|\Lambda^{1/2}[M(D),a]F\|_2
 \le C\Gamma^{-1/2}\|\nabla a\|_\infty\|F\|_2}
 \label{eq:v10-orderzero-Hhalf-commutator}
\end{equation}
holds.  The constants depend only on finitely many symbol seminorms and on
$c_0,C_0$.
\end{lemma}

\begin{proof}
Choose a periodic Littlewood--Paley partition and an integer $j$ with
$2^j\asymp\Gamma$.  Write
\[
 a=S_{j-4}a+\sum_{q\ge j-4}\Delta_q a.
\]
For the low-frequency factor, the kernel of $M(D)$ restricted to the
$j$-th annulus has first absolute moment at most
$C2^{(s-1)j}$.  The identity
\[
 [a,M]F(x)=\int K_j(y)[a(x)-a(x-y)]F(x-y)\,\dd y
\]
therefore gives the right side of
\eqref{eq:v10-general-commutator} for $S_{j-4}a$.

For $q\ge j-4$, Bernstein's inequality gives
\[
 \|\Delta_qa\|_\infty
 \le C2^{-q}\|\nabla a\|_\infty.
\]
The product $\Delta_qa\,F$ has output frequencies $O(2^q+2^j)$.
The two terms of the commutator are consequently bounded by a summable
series of size
\[
 C\|\nabla a\|_\infty
 \sum_{q\ge j-4}2^{(s-1)q}\|F\|_2
 \le C2^{(s-1)j}\|\nabla a\|_\infty\|F\|_2,
\]
because $s<1$.  This proves \eqref{eq:v10-general-commutator}.
For a degree-zero multiplier, applying $\Lambda^{1/2}$ to the low output
contributes $2^{j/2}$ to the order $-1$ commutator, while the high-frequency
series is
$\sum_{q\ge j-4}2^{-q/2}$.  Both are
$O(2^{-j/2})$, proving
\eqref{eq:v10-orderzero-Hhalf-commutator}.  The periodic low-frequency
remainder is smooth and obeys the same, weaker bound after the constant is
increased.
\end{proof}

For the V9 annulus define the actual source-visible tensor
\begin{equation}
 \boxed{
 \mathbb F_n(t)
 :=\mathsf P_{\rm sv}\Psi_n\mathbb C_{n,t}^{\circ}.}
 \label{eq:v10-visible-annular-tensor}
\end{equation}
Then
\begin{equation}
 \boxed{
 g_n^{\rm ann}=T\mathbb F_n,
 \qquad
 \|g_n^{\rm ann}(t)\|_2^2
 =\frac12\|\Lambda^{1/2}\mathbb F_n(t)\|_2^2.}
 \label{eq:v10-annular-source-polar}
\end{equation}

\begin{corollary}[Relative pseudolocality on the annular source range]
\label{cor:v10-relative-pseudolocality}
For every Lipschitz scalar $a$ and every $t\in I_n$,
\begin{align}
 \|[a,T]\mathbb F_n(t)\|_2
 &\le C\frac{\|\nabla a\|_\infty}{\Gamma_n}
       \|g_n^{\rm ann}(t)\|_2,
 \label{eq:v10-T-relative-commutator}\\
 \|\Lambda^{1/2}[\mathsf P_{\rm sv},a]\mathbb F_n(t)\|_2
 &\le C\frac{\|\nabla a\|_\infty}{\Gamma_n}
       \|g_n^{\rm ann}(t)\|_2.
 \label{eq:v10-Psv-relative-commutator}
\end{align}
Here $[a,T]=aT-Ta$ and
$[\mathsf P_{\rm sv},a]=\mathsf P_{\rm sv}a-a\mathsf P_{\rm sv}$.
\end{corollary}

\begin{proof}
Apply Lemma~\ref{lem:v10-annular-calderon} with $s=1/2$ to $T$ and with
$s=0$ to $\mathsf P_{\rm sv}$.  Since
$\mathbb F_n=\mathsf P_{\rm sv}\mathbb F_n$ and its frequency support lies
in $[\Gamma_n/4,4\Gamma_n]$, the exact polar identity gives
\[
 \|\mathbb F_n\|_2
 \le\sqrt8\,\Gamma_n^{-1/2}\|T\mathbb F_n\|_2.
\]
Insert this bound in the two commutator estimates.
\end{proof}

\begin{proposition}[Exact output-side localization after the Hodge short]
\label{prop:v10-output-localization}
For every scalar $a$,
\begin{equation}
 \boxed{
 a g_n^{\rm ann}
 =T(a\mathbb F_n)+[a,T]\mathbb F_n.}
 \label{eq:v10-output-localization}
\end{equation}
No pressure-follower or divergence-null component occurs in this identity.
By contrast, source-side localization of the complete annular covariance is
\begin{equation}
 T\bigl(a\Psi_n\mathbb C_{n,t}^{\circ}\bigr)
 =T(a\mathbb F_n)
  +T[\mathsf P_{\rm sv},a]
    (I-\mathsf P_{\rm sv})\Psi_n\mathbb C_{n,t}^{\circ}.
 \label{eq:v10-source-side-mixing}
\end{equation}
The last term is the exact source-null-to-visible router mixing.
\end{proposition}

\begin{proof}
The first formula is the definition of $[a,T]$ applied to
$g_n^{\rm ann}=T\mathbb F_n$.  For the second, decompose the annular tensor
as $\mathbb F_n+\mathbb N_n$, where
$\mathbb N_n=(I-\mathsf P_{\rm sv})\Psi_n\mathbb C_{n,t}^{\circ}$ and $T\mathbb N_n=0$.  Then
\[
 T(a\mathbb N_n)
 =T\mathsf P_{\rm sv}(a\mathbb N_n)
 =T[\mathsf P_{\rm sv},a]\mathbb N_n,
\]
which proves \eqref{eq:v10-source-side-mixing}.
\end{proof}

% =============================================================================
\section{The target-native annular product theorem}
\label{sec:v10-target-native-product}
% =============================================================================

The selected scalar contains the router weight on the output side.  This
makes it possible to avoid the square-root localization altogether.  Define
\begin{align}
 \mathcal B_n^{\rm out}
 &:=\int_{I_n}\|\omega_ng_n^{\rm ann}\|_2^2\,\dd t,
 \label{eq:v10-output-weight-action}\\
 \mathcal B_n^{0}
 &:=\int_{I_n}\|g_n^{\rm ann}\|_2^2\,\dd t,
 \label{eq:v10-unweighted-annular-action}\\
 \mathfrak L_n^{\rm ann}
 &:=\frac{\mathcal B_n^0}{\nu\mathcal R_n^2}.
 \label{eq:v10-annular-action-leverage}
\end{align}

\begin{lemma}[The signed annular scalar forces output-weight square action]
\label{lem:v10-output-action-lower}
For every late $n$,
\begin{equation}
 \boxed{
 \mathcal B_n^{\rm out}
 \ge c_{\rm out}\nu\mathcal R_n^2,
 \qquad
 c_{\rm out}:=\frac{h_R^2}{16R_t}.}
 \label{eq:v10-output-action-lower}
\end{equation}
\end{lemma}

\begin{proof}
By \eqref{eq:v10-standing-tame} and the annular writer bound,
\begin{equation}
 \int_{I_n}\|W_n\|_2^2\,\dd t
 \le4\Gamma_n^2|I_n|
 \le\frac{4R_t}{\nu}.
 \label{eq:v10-unweighted-writer-cap}
\end{equation}
Cauchy--Schwarz in spacetime and
\eqref{eq:v10-standing-scalar} give
\[
 \frac{h_R^2}{4}\mathcal R_n^2
 \le |h_n^{R,\rm ann}|^2
 \le\mathcal B_n^{\rm out}
       \int_{I_n}\|W_n\|_2^2\,\dd t.
\]
Insert \eqref{eq:v10-unweighted-writer-cap}.
\end{proof}

Put
\begin{equation}
 \varepsilon_n^{\rm route}
 :=\sup_{t\in I_n}
   \frac{\|\nabla\omega_n(t)\|_\infty}{\Gamma_n}.
 \label{eq:v10-route-epsilon}
\end{equation}
The router estimate and \eqref{eq:v8-doubly-tame-properties} imply
\begin{equation}
 \boxed{
 \varepsilon_n^{\rm route}
 \le C\frac{e^{\mathfrak B_n}}{\Gamma_n\rho_n}
 \longrightarrow0.}
 \label{eq:v10-route-epsilon-zero}
\end{equation}

\begin{theorem}[Annular action overdrive or collar-free source-visible product]
\label{thm:v10-target-native-product}
After passage to a subsequence, exactly one of the following priority
branches occurs.
\begin{enumerate}[label=\textup{(P\arabic*)},leftmargin=3.4em]
\item \emph{Unweighted annular source-action overdrive:}
      \begin{equation}
       \boxed{\mathfrak L_n^{\rm ann}\longrightarrow\infty.}
       \label{eq:v10-annular-action-overdrive}
      \end{equation}
      The selected output source then has action growing faster than
      $\nu\mathcal R_n^2$ before any spatial localization.
\item \emph{Target-native product branch:} for a fixed $L_*<\infty$,
      \begin{equation}
       \mathcal B_n^0\le L_*\nu\mathcal R_n^2.
       \label{eq:v10-annular-action-bounded}
      \end{equation}
      Define
      \begin{align}
       \mathcal C_n^{\rm sv}
       &:=\int_{I_n}
         \|[\omega_n,T]\mathbb F_n\|_2^2\,\dd t,
       \label{eq:v10-visible-collar}\\
       \mathcal Q_n^{\rm sv}
       &:=\int_{I_n}
         \|T(\omega_n\mathbb F_n)\|_2^2\,\dd t.
       \label{eq:v10-visible-product-action}
      \end{align}
      Then
      \begin{equation}
       \boxed{
       \mathcal C_n^{\rm sv}=o(\nu\mathcal R_n^2),
       \qquad
       \mathcal Q_n^{\rm sv}
       \ge\frac{c_{\rm out}}4\nu\mathcal R_n^2}
       \label{eq:v10-visible-product-lower}
      \end{equation}
      for all late $n$.  Moreover,
      \begin{equation}
       \boxed{
       \int_{I_n}
       \|\Lambda^{1/2}(\omega_n\mathbb F_n)\|_2^2\,\dd t
       \ge c_{\rm prod}\nu\mathcal R_n^2}
       \label{eq:v10-preprojected-product-lower}
      \end{equation}
      for a fixed $c_{\rm prod}>0$ depending only on the previous margins and
      the annular cutoffs.
\end{enumerate}
The collar in \eqref{eq:v10-visible-collar} is formed from the actual
source-visible annular tensor and hence cannot be inflated by the source-null
example of Countertheorem~\ref{cth:v10-full-collar-not-target}.
\end{theorem}

\begin{proof}
Apply the ordinary subsequence alternative to the nonnegative numbers
$\mathfrak L_n^{\rm ann}$.  On the bounded branch,
Corollary~\ref{cor:v10-relative-pseudolocality} and
\eqref{eq:v10-route-epsilon-zero} give
\[
 \mathcal C_n^{\rm sv}
 \le C(\varepsilon_n^{\rm route})^2\mathcal B_n^0
 =o(\nu\mathcal R_n^2).
\]
Proposition~\ref{prop:v10-output-localization}, with $a=\omega_n$, gives
\[
 \omega_ng_n^{\rm ann}
 =T(\omega_n\mathbb F_n)+[\omega_n,T]\mathbb F_n.
\]
Hence
\[
 (\mathcal B_n^{\rm out})^{1/2}
 \le(\mathcal Q_n^{\rm sv})^{1/2}
   +(\mathcal C_n^{\rm sv})^{1/2}.
\]
By Lemma~\ref{lem:v10-output-action-lower}, the first term on the left is at
least
$\sqrt{c_{\rm out}\nu}\,\mathcal R_n$.  The collar is lower order, so for
late $n$ its square root is at most one half of that number.  This proves the
second inequality in \eqref{eq:v10-visible-product-lower}.

The exact polar identity gives
\begin{equation}
 \mathcal Q_n^{\rm sv}
 =\frac12\int_{I_n}
   \|\Lambda^{1/2}\mathsf P_{\rm sv}
       (\omega_n\mathbb F_n)\|_2^2\,\dd t.
 \label{eq:v10-local-polar-identity}
\end{equation}
Because $\mathsf P_{\rm sv}\mathbb F_n=\mathbb F_n$,
\begin{equation}
 \mathsf P_{\rm sv}(\omega_n\mathbb F_n)
 =\omega_n\mathbb F_n
  +[\mathsf P_{\rm sv},\omega_n]\mathbb F_n.
 \label{eq:v10-Psv-local-split}
\end{equation}
The second estimate in
Corollary~\ref{cor:v10-relative-pseudolocality} gives
\[
 \int_{I_n}
 \|\Lambda^{1/2}[\mathsf P_{\rm sv},\omega_n]
       \mathbb F_n\|_2^2\,\dd t
 \le C(\varepsilon_n^{\rm route})^2\mathcal B_n^0
 =o(\nu\mathcal R_n^2).
\]
Combining this with \eqref{eq:v10-local-polar-identity},
\eqref{eq:v10-visible-product-lower}, and the reverse triangle inequality
proves \eqref{eq:v10-preprojected-product-lower} after decreasing the fixed
constant.
\end{proof}

\begin{corollary}[The pressure and divergence-null rows remain independent]
\label{cor:v10-null-rows-independent}
Let
\begin{equation}
 \Psi_n\mathbb C_{n,t}^{\circ}
 =\mathbb F_n+\mathbb G_n^{\rm gr}+\mathbb G_n^{\rm dn}
 \label{eq:v10-annular-three-way}
\end{equation}
be the three orthogonal components of
Theorem~\ref{thm:v10-mode-polarization}.  On branch \textup{(P2)},
\eqref{eq:v10-preprojected-product-lower} controls only the first component.
The annular pressure follower and divergence-null stocks are
\begin{align}
 \mathcal P_n^{\rm gr}
 &:=\int_{I_n}
   \|\Lambda^{1/2}(\omega_n\mathbb G_n^{\rm gr})\|_2^2\,\dd t,
 \label{eq:v10-gradient-stock}\\
 \mathcal N_n^{\rm dn}
 &:=\int_{I_n}
   \|\Lambda^{1/2}(\omega_n\mathbb G_n^{\rm dn})\|_2^2\,\dd t.
 \label{eq:v10-divnull-stock}
\end{align}
Neither has an upper bound from the Reynolds source action.  Their boundedness
or concentration is an independent product/pressure row.
\end{corollary}

\begin{proof}
Orthogonality and the source identities are
Theorem~\ref{thm:v10-mode-polarization}.  The physical source annihilates
both displayed components before multiplication by $\omega_n$.  Any
contribution after multiplication is a projection commutator, not a term in
the unlocalized Reynolds source action.  Therefore no estimate for
$\mathcal B_n^0$ or $\mathcal Q_n^{\rm sv}$ bounds
\eqref{eq:v10-gradient-stock}--\eqref{eq:v10-divnull-stock}.
Countertheorem~\ref{cth:v10-full-collar-not-target} gives an explicit family
where a source-null component is arbitrarily large while the source remains
fixed.
\end{proof}

\begin{assessment}[The V9 fractional stop after the upstream audit]
\label{ass:v10-fractional-stop-corrected}
The complete V9 collar remains a valid record of all tensor mixing.  For the
selected annular response, the terminating physical alternative is sharper:
\[
 \boxed{
 \begin{gathered}
 \text{unweighted annular source-action overdrive}\\[0.15em]
 \big|\\[-0.15em]
 \text{collar-free source-visible local product}\\[0.15em]
 \big|\\[-0.15em]
 \text{gradient-pressure or divergence-null tensor stock}.
 \end{gathered}}
\]
This is the source-first Hodge short required before one asks for a velocity
or pressure profile.
\end{assessment}

% =============================================================================
\section{One annular source cell or target-native microdelocalization}
\label{sec:v10-source-cell}
% =============================================================================

On branch \textup{(P2)}, normalize the output-weight action by
\begin{equation}
 \dd\mu_n^{\rm sv}(t,x)
 :=\frac{|\omega_n(t,x)g_n^{\rm ann}(t,x)|^2}
         {\mathcal B_n^{\rm out}}\,\dd x\,\dd t.
 \label{eq:v10-source-action-law}
\end{equation}
This is the actual selected source law after output-frequency and Hodge
projection; it is not the complete covariance-product law of Version~V9.
For $R>0$, set
\begin{equation}
 \vartheta_n^{\rm sv}(R)
 :=\sup_{y\in\Torus^3}
   \mu_n^{\rm sv}(I_n\times B(y,R/\Gamma_n)).
 \label{eq:v10-source-concentration-function}
\end{equation}

\begin{theorem}[Target-native one-scale alternative]
\label{thm:v10-source-cell-alternative}
After passage to a subsequence, exactly one of the following occurs.
\begin{enumerate}[label=\textup{(S\arabic*)},leftmargin=3.4em]
\item \emph{One source-action cell.}  There are fixed
      $R_0<\infty$, $\eta_0>0$, and centres $y_n$ such that
      \begin{equation}
       \boxed{
       \int_{I_n}\int_{B(y_n,R_0/\Gamma_n)}
       |\omega_ng_n^{\rm ann}|^2
       \ge\eta_0c_{\rm out}\nu\mathcal R_n^2.}
       \label{eq:v10-heavy-source-cell}
      \end{equation}
\item \emph{Source-action microdelocalization.}  For every fixed $R$,
      \begin{equation}
       \boxed{\vartheta_n^{\rm sv}(R)\longrightarrow0.}
       \label{eq:v10-source-microdelocalization}
      \end{equation}
      Every bounded collection of annular cells then carries a vanishing
      fraction of the selected source action.
\end{enumerate}
Branch \textup{(S2)} is the target-native refinement of the V9
material-tube/annular-scale defect; no additional concentration grammar is
introduced.
\end{theorem}

\begin{proof}
For every positive integer $m$, pass diagonally so that
$\vartheta_n^{\rm sv}(m)$ converges.  If one limit is positive, choose the
corresponding $R_0=m$, take centres within half the supremum, and use
Lemma~\ref{lem:v10-output-action-lower}.  If every limit is zero, monotonicity
in $R$ gives \eqref{eq:v10-source-microdelocalization} for all fixed $R$.
A union of $M$ fixed-radius cells has mass at most
$M\vartheta_n^{\rm sv}(R)$, proving the final statement.
\end{proof}

The source cell and the V9 quartic cell need not coincide merely because both
lie in the same shrinking material tube.  The next definition records the
minimum same-cell incidence required for a velocity-product profile.

\begin{definition}[Source/product cell-incidence residual]
\label{def:v10-cell-incidence}
On branch \textup{(S1)}, fix a smooth cutoff $\chi_n$ which equals one on
$B(y_n,R_0/\Gamma_n)$, is supported in
$B(y_n,(R_0+1)/\Gamma_n)$, and satisfies
$\|\nabla\chi_n\|_\infty\le C\Gamma_n$.  Put
\begin{equation}
 h_{n,\chi}^{\rm ann}
 :=-\int_{I_n}\!\int\chi_n\omega_ng_n^{\rm ann}\cdot W_n
 \label{eq:v10-local-signed-source}
\end{equation}
and
\begin{equation}
 \Delta_n^{\rm cell}
 :=\frac{1}{\mathcal R_n}
   \left|h_n^{R,\rm ann}-h_{n,\chi}^{\rm ann}\right|.
 \label{eq:v10-cell-incidence-residual}
\end{equation}
A cofinal cell is called signed-source incident when, after enlarging $R_0$
by one fixed amount if necessary,
\begin{equation}
 |h_{n,\chi}^{\rm ann}|\ge h_{\chi}\mathcal R_n
 \label{eq:v10-local-signed-margin}
\end{equation}
for a fixed $h_\chi>0$.
\end{definition}

\begin{assessment}[Why action and signed work are kept distinct]
\label{ass:v10-action-versus-signed-cell}
A cell may carry positive source square action while its signed pairing with
the selected writer is small because of local counterflow, or the record
scalar may be distributed over many cells.  Failure of
\eqref{eq:v10-local-signed-margin} is therefore retained as the existing
signed counterflow/microdelocalization or source/product-incidence defect.
It is not repaired by moving the cell after reading the writer.  The profile
construction below is performed only on the signed-source incident branch.
\end{assessment}

% =============================================================================
\section{Record-amplitude rescaling of a common annular cell}
\label{sec:v10-rescaling}
% =============================================================================

Work on the one-cell, signed-source incident branch and choose the centres
$y_n$ from the common source/product cell.  Put
\begin{equation}
 \tau_n:=\nu\Gamma_n^2(b_n-a_n)\in(0,R_t],
 \qquad
 J_n:=(-\tau_n,0),
 \qquad
 \alpha_n:=\sqrt{\frac{\mathcal R_n}{\nu}}\longrightarrow\infty.
 \label{eq:v10-scaled-time-amplitude}
\end{equation}
The rescaled spatial torus is
$\Torus_n^3:=\Gamma_n(\Torus^3-y_n)$.  For
$s\in J_n$ and $y\in\Torus_n^3$, set
\begin{align}
 V_n(s,y)
 &:=\frac{1}{\Gamma_n\sqrt{\nu\mathcal R_n}}
 u\!\left(b_n+\frac{s}{\nu\Gamma_n^2},
          y_n+\frac{y}{\Gamma_n}\right),
 \label{eq:v10-scaled-velocity}\\
 \Pi_n(s,y)
 &:=\frac{1}{\nu\mathcal R_n\Gamma_n^2}
 p\!\left(b_n+\frac{s}{\nu\Gamma_n^2},
          y_n+\frac{y}{\Gamma_n}\right),
 \label{eq:v10-scaled-pressure}\\
 \Omega_n(s,y)
 &:=\omega_n\!\left(b_n+\frac{s}{\nu\Gamma_n^2},
          y_n+\frac{y}{\Gamma_n}\right).
 \label{eq:v10-scaled-router}
\end{align}
Let
\begin{equation}
 \varphi_n(y)
 :=\Gamma_n^{-3/2}
   \varpi_n\!\left(y_n+\frac{y}{\Gamma_n}\right),
 \qquad
 Z_n:=\Lambda_y\varphi_n.
 \label{eq:v10-scaled-writer}
\end{equation}
Then $\|\varphi_n\|_2=1$, its Fourier support lies in a fixed annulus, and
\begin{equation}
 \frac12\le\|Z_n\|_2\le2.
 \label{eq:v10-scaled-writer-window}
\end{equation}

For $s\le0$, define the heat covariance of $V_n$ at age $-s$ by
\begin{equation}
 \mathbb D_n(s)
 :=e^{-s\Delta_y}(V_n\otimes V_n)(s)
  -\bigl(e^{-s\Delta_y}V_n(s)\bigr)
   \otimes
   \bigl(e^{-s\Delta_y}V_n(s)\bigr).
 \label{eq:v10-scaled-covariance}
\end{equation}
It is pointwise positive semidefinite.  Put
\begin{equation}
 \mathbb F_n^{\rm sc}
 :=\mathsf P_{\rm sv}\psi(\Lambda_y)\mathbb D_n^{\circ},
 \qquad
 G_n:=T_y\mathbb F_n^{\rm sc}.
 \label{eq:v10-scaled-source}
\end{equation}

\begin{proposition}[Exact scaling identities]
\label{prop:v10-scaling-identities}
The scaled velocity solves
\begin{equation}
 \boxed{
 \partial_sV_n
 +\alpha_n\bigl((V_n\cdot\nabla)V_n+\nabla\Pi_n\bigr)
 =\Delta V_n,
 \qquad \operatorname{div}V_n=0.}
 \label{eq:v10-strong-coupling-equation}
\end{equation}
Moreover,
\begin{align}
 \frac{h_n^{R,\rm ann}}{\mathcal R_n}
 &=-\int_{J_n}\!\int_{\Torus_n^3}
     \Omega_nG_n\cdot Z_n\,\dd y\,\dd s,
 \label{eq:v10-scaled-scalar}\\
 \frac{\mathcal B_n^0}{\nu\mathcal R_n^2}
 &=\int_{J_n}\|G_n(s)\|_2^2\,\dd s,
 \label{eq:v10-scaled-source-action}\\
 \frac{\mathcal B_n^{\rm out}}{\nu\mathcal R_n^2}
 &=\int_{J_n}\|\Omega_nG_n(s)\|_2^2\,\dd s,
 \label{eq:v10-scaled-output-action}\\
 \Gamma_n\int_{I_n}\!\int_{B(y_n,R/\Gamma_n)}|u|^4
 &=\nu\mathcal R_n^2
   \int_{J_n}\!\int_{B(0,R)}|V_n|^4.
 \label{eq:v10-scaled-L4}
\end{align}
Finally,
\begin{equation}
 \|\nabla_y\Omega_n\|_\infty
 \le C\frac{e^{\mathfrak B_n}}{\Gamma_n\rho_n}
 \longrightarrow0.
 \label{eq:v10-scaled-router-flat}
\end{equation}
\end{proposition}

\begin{proof}
Substitute
$u=\Gamma_n\sqrt{\nu\mathcal R_n}\,V_n$,
$p=\nu\mathcal R_n\Gamma_n^2\Pi_n$,
$\partial_t=\nu\Gamma_n^2\partial_s$, and
$\nabla_x=\Gamma_n\nabla_y$ into Navier--Stokes and divide by
$\nu\Gamma_n^3\sqrt{\nu\mathcal R_n}$.  This gives
\eqref{eq:v10-strong-coupling-equation}.

The covariance scales as
\[
 \mathbb C_{n,t}
 =\nu\mathcal R_n\Gamma_n^2\mathbb D_n(s),
\]
while the order-$1/2$ source operator gives
\[
 g_n^{\rm ann}(t,x)
 =\nu\mathcal R_n\Gamma_n^{5/2}G_n(s,y).
\]
The writer scales as
$W_n(x)=\Gamma_n^{5/2}Z_n(y)$, and
$\dd x\,\dd t=(\nu\Gamma_n^5)^{-1}\dd y\,\dd s$.
These identities prove
\eqref{eq:v10-scaled-scalar}--\eqref{eq:v10-scaled-output-action}.
The same change of variables and
$u=\Gamma_n\sqrt{\nu\mathcal R_n}V_n$ prove
\eqref{eq:v10-scaled-L4}.  The writer norm window follows from the
frequency support of $\varpi_n$.  Finally,
$\nabla_y=\Gamma_n^{-1}\nabla_x$ and \eqref{eq:v9-router} give
\eqref{eq:v10-scaled-router-flat}.
\end{proof}

\begin{corollary}[A bounded action packet has nonvanishing scaled duration]
\label{cor:v10-time-noncollapse}
Assume \eqref{eq:v10-annular-action-bounded} and the signed-cell margin
\eqref{eq:v10-local-signed-margin}.  Then
\begin{equation}
 \boxed{
 \tau_n\ge\tau_*
 :=\frac{h_\chi^2}{4L_*}>0}
 \label{eq:v10-time-lower}
\end{equation}
for all late $n$.
\end{corollary}

\begin{proof}
Choose the cutoffs in Definition~\ref{def:v10-cell-incidence} as rescalings
of one fixed $\chi\in C_c^\infty(B(0,R_0+1))$.  Then
\eqref{eq:v10-local-signed-margin} and
\eqref{eq:v10-scaled-scalar} give
\[
 h_\chi
 \le\left(\int_{J_n}\!\int\chi^2|G_n|^2\right)^{1/2}
      \left(\int_{J_n}\!\int\chi^2|Z_n|^2\right)^{1/2}.
\]
The first factor is at most $L_*^{1/2}$ by
\eqref{eq:v10-scaled-source-action}; the second is at most
$2\sqrt{\tau_n}$ by \eqref{eq:v10-scaled-writer-window}.  Rearrangement
proves \eqref{eq:v10-time-lower}.
\end{proof}

The one-cell quartic statement is useful only if the source and product use
the same cell.  We therefore strengthen the incidence convention of
Definition~\ref{def:v10-cell-incidence}: on the profile branch the chosen
fixed cutoff also satisfies
\begin{equation}
 \int_{J_n}\!\int_{B(0,R_0+1)}|V_n|^4
 \ge q_*>0.
 \label{eq:v10-common-cell-L4-lower}
\end{equation}
If this fails for every fixed enlargement, the V9 quartic law and the
annular source law occupy different cells and their separation is the
source/product cell-incidence defect.  No cell is reselected after the two
marginals are read.

% =============================================================================
\section{Nonzero source-visible covariance profile or a named compactness defect}
\label{sec:v10-profile}
% =============================================================================

The source-visible tensor profile requires only the bounded annular source
action.  A velocity or pressure profile requires additional product bounds,
which are stated separately.

\begin{theorem}[First nonzero source-visible covariance profile]
\label{thm:v10-covariance-profile}
Assume branch \textup{(P2)}, branch \textup{(S1)}, the signed-cell margin
\eqref{eq:v10-local-signed-margin}, and the common-cell condition
\eqref{eq:v10-common-cell-L4-lower}.  After passage to a subsequence,
\begin{equation}
 \tau_n\longrightarrow\tau_\infty
 \in[\tau_*,R_t].
 \label{eq:v10-time-limit}
\end{equation}
Extend the scaled fields by zero from $J_n$ to
$J:=(-\tau_\infty,0)$.  Then there are
\begin{equation}
 Z\in C^\infty_{\rm loc}(\R^3;\R^3),
 \qquad
 H\in L^2_{\rm loc}(J\times\R^3;\R^3),
 \label{eq:v10-HZ-limits}
\end{equation}
and a tensor
\begin{equation}
 F\in L^2\bigl(J;\dot H^{1/2}_{\rm loc}(\R^3;
          \operatorname{Sym}_0(3))\bigr)
 \label{eq:v10-F-limit-space}
\end{equation}
such that, along the subsequence,
\begin{align}
 Z_n&\longrightarrow Z
       &&\text{strongly in }L^2_{\rm loc}(\R^3),
 \label{eq:v10-writer-strong}\\
 \Omega_nG_n&\rightharpoonup H
       &&\text{weakly in }L^2_{\rm loc}(J\times\R^3),
 \label{eq:v10-source-weak}\\
 \Omega_n\mathbb F_n^{\rm sc}&\rightharpoonup F
       &&\text{weakly in }L^2(J;\dot H^{1/2}_{\rm loc}).
 \label{eq:v10-tensor-weak}
\end{align}
The limits obey
\begin{equation}
 \boxed{H=T_yF\quad\text{in }\mathcal D'(J\times\R^3),}
 \label{eq:v10-limit-source-relation}
\end{equation}
and are nonzero.  More precisely, for the fixed cutoff $\chi$ used in the
cell incidence,
\begin{equation}
 \boxed{
 \left|\int_J\!\int_{\R^3}\chi(y)H(s,y)\cdot Z(y)
       \,\dd y\,\dd s\right|
 \ge h_\chi.}
 \label{eq:v10-limit-nonzero-pairing}
\end{equation}
Thus a bounded target-native branch produces a nonzero source-visible tensor
profile on one fixed annular cell.
\end{theorem}

\begin{proof}
The time limit follows from
Corollary~\ref{cor:v10-time-noncollapse} and the upper bound $R_t$.
The rescaled writers have Fourier support in a fixed compact annulus and
uniform $L^2$ norm.  Bernstein estimates therefore bound every spatial
Sobolev norm on fixed balls.  Rellich compactness and a diagonal argument
give \eqref{eq:v10-writer-strong}, after passing to a subsequence.

Equation \eqref{eq:v10-scaled-source-action} and
\eqref{eq:v10-annular-action-bounded} give a uniform global
$L_s^2L_y^2$ bound for $G_n$.  Hence $\Omega_nG_n$ has a weakly convergent
subsequence, giving \eqref{eq:v10-source-weak}.
The scaled form of \eqref{eq:v10-annular-source-polar} gives
\[
 \int_{J_n}\|\Lambda_y^{1/2}\mathbb F_n^{\rm sc}\|_2^2\,\dd s
 =2\int_{J_n}\|G_n\|_2^2\,\dd s\le2L_*.
\]
Furthermore, the scaled versions of
\eqref{eq:v10-output-localization} and
\eqref{eq:v10-T-relative-commutator}, together with
\eqref{eq:v10-scaled-router-flat}, give
\begin{equation}
 \|\Omega_nG_n-T_y(\Omega_n\mathbb F_n^{\rm sc})
   \|_{L^2(J_n\times\Torus_n^3)}
 \longrightarrow0.
 \label{eq:v10-scaled-collar-zero}
\end{equation}
The same argument using
\eqref{eq:v10-Psv-relative-commutator} shows that
$\Omega_n\mathbb F_n^{\rm sc}$ is uniformly bounded in
$L_s^2\dot H_y^{1/2}$: its source-visible projection is controlled by
$T_y(\Omega_n\mathbb F_n^{\rm sc})$, and its orthogonal component is the
vanishing projection commutator.  This proves
\eqref{eq:v10-tensor-weak}.  Passing to distributions in
\eqref{eq:v10-scaled-collar-zero} gives
\eqref{eq:v10-limit-source-relation}.

Under scaling, the signed-cell margin is
\[
 \left|\int_{J_n}\!\int\chi\Omega_nG_n\cdot Z_n\right|
 \ge h_\chi.
\]
The functions $\chi Z_n\one_{J_n}$ converge strongly in spacetime to
$\chi Z\one_J$, because
$\tau_n\to\tau_\infty$ and \eqref{eq:v10-writer-strong} holds.  Pairing this
strong convergence with \eqref{eq:v10-source-weak} proves
\eqref{eq:v10-limit-nonzero-pairing}.  In particular $H\ne0$, and
\eqref{eq:v10-limit-source-relation} then implies $F\ne0$.
\end{proof}

\begin{firewall}[The profile is source visible, not yet a full covariance]
\label{fire:v10-profile-not-full-product}
The tensor $F$ is the weak limit of the output-annular, pressure-shorted,
material-weighted covariance component.  The theorem does not bound or
identify the gradient-pressure component
$\mathsf P_{\rm gr}\mathbb D_n^{\circ}$, the divergence-null component
$\mathsf P_{\rm dn}\mathbb D_n^{\circ}$, the complete velocity product, or
the Navier--Stokes pressure.  Those rows are exactly what is required to
promote $F$ to a complete local PDE profile.
\end{firewall}

For completeness, the following theorem performs that promotion as far as
the available bounds permit and returns a named defect whenever they do not.

\begin{theorem}[Product/pressure overdrive or a local Euler--Reynolds constraint]
\label{thm:v10-product-pressure-profile}
Retain the hypotheses of
Theorem~\ref{thm:v10-covariance-profile}.  Fix
$R_1>R_0+2$.  After passage to a further subsequence, one of the following
priority branches occurs.
\begin{enumerate}[label=\textup{(E\arabic*)},leftmargin=3.4em]
\item \emph{Normalized product overdrive:}
      \begin{equation}
       \int_{J_n}\int_{B(0,R_1)}|V_n|^4\longrightarrow\infty.
       \label{eq:v10-product-overdrive}
      \end{equation}
\item The products are bounded, but the locally normalized pressure is
      unbounded:
      \begin{equation}
       \inf_{c_n\in L^2(J_n)}
       \int_{J_n}\int_{B(0,R_1)}|\Pi_n(s,y)-c_n(s)|^2
       \longrightarrow\infty.
       \label{eq:v10-pressure-overdrive}
      \end{equation}
\item Both quantities are bounded.  Then, on every smaller ball
      $B(0,R)$, $R<R_1$, there are weak limits
      \begin{align}
       V_n&\rightharpoonup V
       &&\text{in }L^4(J\times B(0,R_1)),
       \label{eq:v10-V-weak}\\
       V_n\otimes V_n&\rightharpoonup Q
       &&\text{in }L^2(J\times B(0,R_1)),
       \label{eq:v10-Q-weak}\\
       \Pi_n-c_n&\rightharpoonup\Pi
       &&\text{in }L^2(J\times B(0,R_1)),
       \label{eq:v10-pressure-weak}
      \end{align}
      satisfying
      \begin{equation}
       \boxed{\operatorname{div}Q+\nabla\Pi=0,
       \qquad \operatorname{div}V=0}
       \label{eq:v10-Euler-Reynolds-constraint}
      \end{equation}
      in distributions on $J\times B(0,R)$.  Moreover,
      \begin{equation}
       \boxed{\mathbb R:=Q-V\otimes V\succeq0}
       \label{eq:v10-positive-product-defect}
      \end{equation}
      as a matrix-valued distribution.
\end{enumerate}
On the third branch, $\mathbb R$ is the exact local product-compactness
defect.  If the products converge strongly to $V\otimes V$, then
$\mathbb R=0$ and \eqref{eq:v10-Euler-Reynolds-constraint} is the stationary
Euler constraint at almost every slow time.  It is not an ancient dynamical
Navier--Stokes equation.
\end{theorem}

\begin{proof}
Pass successively to subsequences on which the two nonnegative quantities in
\eqref{eq:v10-product-overdrive}--\eqref{eq:v10-pressure-overdrive} either
are bounded or diverge.  On the bounded branch, Banach--Alaoglu gives
\eqref{eq:v10-V-weak}--\eqref{eq:v10-pressure-weak}.  The pressure constants
may be chosen as the spatial means over $B(0,R_1)$; they disappear after a
gradient.

Divide \eqref{eq:v10-strong-coupling-equation} by $\alpha_n$:
\begin{equation}
 \alpha_n^{-1}\partial_sV_n
 +\operatorname{div}(V_n\otimes V_n)+\nabla\Pi_n
 =\alpha_n^{-1}\Delta V_n.
 \label{eq:v10-divided-equation}
\end{equation}
For every compactly supported smooth test field, the first and last terms
tend to zero after integration by parts because
$\alpha_n\to\infty$ and $V_n$ is bounded in $L^4$ on the support.  Passing
the remaining terms weakly gives
\eqref{eq:v10-Euler-Reynolds-constraint}.  Divergence-freeness passes
linearly.

To prove positivity, fix $\xi\in\R^3$ and a nonnegative
$\phi\in C_c^\infty(J\times B(0,R_1))$.  The sequence
$\sqrt\phi\,\xi\cdot V_n$ converges weakly in $L^2$ to
$\sqrt\phi\,\xi\cdot V$.  Lower semicontinuity gives
\[
 \int\phi\,|\xi\cdot V|^2
 \le\liminf_n\int\phi\,|\xi\cdot V_n|^2
 =\int\phi\,\xi^{\mathsf T}Q\xi.
\]
This is exactly \eqref{eq:v10-positive-product-defect}.
If the products converge strongly to $V\otimes V$, then $Q=V\otimes V$ and
$\mathbb R=0$; the final statement follows from
\eqref{eq:v10-Euler-Reynolds-constraint}.
\end{proof}

\begin{corollary}[The selected cell remains nontrivial after amplitude normalization]
\label{cor:v10-normalized-cell-nonzero}
On every bounded-product branch satisfying
\eqref{eq:v10-common-cell-L4-lower},
\begin{equation}
 \boxed{
 \liminf_n\int_{J_n}\int_{B(0,R_0+1)}|V_n|^4\ge q_*>0.}
 \label{eq:v10-normalized-L4-nonzero}
\end{equation}
If the weak velocity limit is zero, the mass is carried by the weak quadratic
limit $Q$ and/or by the positive $L^2$ norm defect of the products.  In
particular, $Q=0$ forces a strictly positive product concentration/oscillation
defect rather than a zero packet.
\end{corollary}

\begin{proof}
The first assertion is the common-cell hypothesis.  If $V=0$, then $\mathbb R=Q\succeq0$.  Let $\chi_0$ be a fixed cutoff
which equals one on $B(0,R_0+1)$ and is supported in a slightly larger ball.
After passing to a subsequence, weak lower semicontinuity gives the nonnegative
product norm defect
\[
 \delta_{\rm prod}
 :=\liminf_n\|\chi_0V_n\otimes V_n\|_2^2
   -\|\chi_0Q\|_2^2\ge0.
\]
The lower bound \eqref{eq:v10-normalized-L4-nonzero} implies that if $Q=0$
then $\delta_{\rm prod}\ge q_*>0$.  Thus either the weak quadratic limit or
the finer product concentration/oscillation defect is nonzero.
\end{proof}

\begin{assessment}[Why the limit is not yet the desired ancient profile]
\label{ass:v10-not-ancient}
The amplitude factor $\alpha_n\to\infty$ is forced by the record-sized
source action.  On slow time it removes the time derivative and leaves the
Euler--Reynolds constraint of
\eqref{eq:v10-Euler-Reynolds-constraint}.  On fast time
$\tau=\alpha_ns$, the viscosity becomes $\alpha_n^{-1}$ and the interval
length becomes $\alpha_n\tau_n\to\infty$.  As already audited in Version~V1,
a nonzero slow-time averaged scalar need not survive on one bounded fast-time
window.  A genuine ancient dynamical profile therefore still requires a
terminal occupation/trace compactness theorem or a separate backward
rigidity input.
\end{assessment}

% =============================================================================
\section{Version V10 integrated theorem and exact stopping boundary}
\label{sec:v10-integrated}
% =============================================================================

\begin{theorem}[Version V10 source-visible product and profile theorem]
\label{thm:v10-integrated}
Preserving every theorem and stated limitation of Versions~V1--V9,
Version~V10 proves the following.
\begin{enumerate}[label=\textup{(V10.\arabic*)},leftmargin=5.5em]
\item The complete fractional triangle decomposition of V9 is not target native:
      positive source-null tensors can leave the selected source identically
      unchanged while making both right-hand terms arbitrarily large.
\item Every trace-free symmetric Fourier tensor has the exact orthogonal
      source--pressure--null split of ranks $2+1+2$.  The physical source sees
      only the rank-two component, with norm identity
      \eqref{eq:v10-global-source-norm}; the rank-one component is the pressure
      follower and the final rank-two component has zero divergence.
\item The Moore--Penrose inverse in \eqref{eq:v10-T-MP} is the unique
      source-minimal tensor lift.  Pressure and divergence-null tensors receive
      no Reynolds source price.
\item Annular pseudolocality gains the relative factor
      $\|\nabla\omega_n\|_\infty/\Gamma_n$.  The doubly tame material router
      makes this factor vanish on the selected output annulus.
\item The signed annular scalar itself forces output-weight square action
      \eqref{eq:v10-output-action-lower}.  Hence either the unweighted annular
      source action overdrives $\nu\mathcal R_n^2$, or the target-native
      fractional collar is $o(\nu\mathcal R_n^2)$ and the source-visible
      covariance product satisfies
      \eqref{eq:v10-preprojected-product-lower}.
\item The selected source action either microdelocalizes through every bounded
      annular cell or occupies one cell.  On the latter branch signed work and
      quartic product are required on the same predeclared cell; failure is the
      explicit source/product incidence residual.
\item Record-amplitude and wavelength rescaling gives the exact strong-coupling
      equation \eqref{eq:v10-strong-coupling-equation}.  Bounded normalized
      source action plus a nonzero local signed margin forces a positive lower
      bound for the scaled time depth.
\item On that bounded branch, the material-weighted annular source converges
      weakly to a nonzero vector $H=T_yF$, where $F$ is a nonzero
      source-visible covariance-tensor profile.  No full velocity product or
      pressure compactness is used in this first profile extraction.
\item If the normalized velocity product or pressure is unbounded, the exact
      overdrive is retained.  If both are bounded, every local limit satisfies
      the Euler--Reynolds constraint
      \eqref{eq:v10-Euler-Reynolds-constraint} with positive product defect
      \eqref{eq:v10-positive-product-defect}.
\item The slow-time profile is not an ancient Navier--Stokes solution.  Strong
      product compactness, pressure/source-null control, terminal occupation,
      and a backward rigidity or Liouville theorem remain separate PDE rows.
\end{enumerate}
No item proves a terminal Dini improvement, epsilon regularity, terminal
trace, backward uniqueness, an ancient Liouville theorem, or global
three-dimensional regularity.
\end{theorem}

\begin{proof}
Item~\textup{(V10.1)} is
Countertheorem~\ref{cth:v10-full-collar-not-target}.  Items
\textup{(V10.2)--(V10.3)} are
Theorem~\ref{thm:v10-mode-polarization} and
Corollary~\ref{cor:v10-source-MP}.  Item~\textup{(V10.4)} is
Lemma~\ref{lem:v10-annular-calderon} and
Corollary~\ref{cor:v10-relative-pseudolocality}.  Item~\textup{(V10.5)} is
Lemma~\ref{lem:v10-output-action-lower} and
Theorem~\ref{thm:v10-target-native-product}.  Item~\textup{(V10.6)} is
Theorem~\ref{thm:v10-source-cell-alternative} and
Definition~\ref{def:v10-cell-incidence}.  Item~\textup{(V10.7)} is
Proposition~\ref{prop:v10-scaling-identities} and
Corollary~\ref{cor:v10-time-noncollapse}.  Item~\textup{(V10.8)} is
Theorem~\ref{thm:v10-covariance-profile}.  Item~\textup{(V10.9)} is
Theorem~\ref{thm:v10-product-pressure-profile}.  The final item is
Assessment~\ref{ass:v10-not-ancient}.
\end{proof}

\begingroup\small
\begin{longtable}{p{0.25\textwidth}p{0.19\textwidth}p{0.48\textwidth}}
\caption{NS--A Version V10 proof-status ledger.}\label{tab:v10-ledger}\\
\toprule Row & Status & Exact conclusion\\\midrule
\endfirsthead
\toprule Row & Status & Exact conclusion\\\midrule
\endhead
V1--V9 prefix & \PROVED{} preserved
 &The complete V9 preterminal source is retained byte-for-byte; no previous
 theorem is rewritten.\\[0.3em]
Complete V9 fractional collar & \OBSTRUCTION{} target scope
 &It is an exact identity but can be inflated by positive source-null tensor
 components which leave the selected annular source unchanged.\\[0.3em]
Source--pressure tensor split & \PROVED{} exact
 &Every nonzero mode splits orthogonally into source-visible rank two, pure
 gradient rank one, and divergence-null rank two.\\[0.3em]
Source Moore--Penrose lift & \PROVED{} exact
 &$T^\dagger=2T^*\Lambda^{-1}$ is the unique minimum tensor preimage and
 $\|TF\|_2^2=\frac12\|\Lambda^{1/2}P_{\rm sv}F\|_2^2$.\\[0.3em]
Pressure follower & \PROVED{} separated
 &Its scalar $\dot H^{1/2}$ norm is the exact rank-one tensor norm in
 \eqref{eq:v10-pressure-gradient-norm}; Reynolds source action does not bound it.\\[0.3em]
Annular pseudolocality & \PROVED{} relative
 &On the source-visible annulus, cutoff commutators cost
 $\|\nabla\omega_n\|_\infty/\Gamma_n\to0$.\\[0.3em]
Unweighted annular source action & \OBSTRUCTION{} typed
 &If it grows faster than $\nu\mathcal R_n^2$, the overdrive is retained before
 any local profile argument.\\[0.3em]
Target-native covariance product & \PROVED{} on bounded-action branch
 &The source-visible collar vanishes and
 $\int\|\Lambda^{1/2}(\omega_n\mathbb F_n)\|_2^2
 \gtrsim\nu\mathcal R_n^2$.\\[0.3em]
Gradient/divergence-null tensor rows & \OPEN{} independent
 &Neither is controlled by the Reynolds source action; they are the exact
 pressure and full-product completion rows.\\[0.3em]
Annular source cell & \PROVED{} exact alternative
 &One cell carries fixed source-action mass or every bounded annular cell is
 negligible.\\[0.3em]
Signed source/product cell incidence & \CONDITIONAL{} physical
 &A profile uses one predeclared cell carrying signed source work and quartic
 product; failure is retained rather than repaired by reselection.\\[0.3em]
Record-amplitude rescaling & \PROVED{} exact
 &The normalized equation is strong coupling with
 $\alpha_n=(\mathcal R_n/\nu)^{1/2}\to\infty$ and all source/action scalings are
 explicit.\\[0.3em]
Scaled time depth & \PROVED{} noncollapse
 &Bounded source action and nonzero local signed work imply
 $\tau_n\ge h_\chi^2/(4L_*)$.\\[0.3em]
Source-visible covariance profile & \PROVED{} weak, nonzero
 &A nonzero $F$ and $H=T_yF$ survive on one fixed annular cell.\\[0.3em]
Normalized velocity product & \OBSTRUCTION{} or bounded
 &Divergence gives the explicit product-overdrive branch; boundedness permits a
 weak quadratic limit.\\[0.3em]
Normalized pressure & \OBSTRUCTION{} or bounded
 &Divergence modulo spatial constants is the pressure-concentration branch;
 boundedness permits a weak pressure limit.\\[0.3em]
Euler--Reynolds slow-time constraint & \PROVED{} conditional on bounds
 &$\operatorname{div}Q+\nabla\Pi=0$ and
 $Q-V\otimes V\succeq0$; this is not an ancient dynamical equation.\\[0.3em]
Terminal trace / rigidity / regularity & \OPEN
 &The source-visible profile does not supply product compactness, terminal
 occupation, backward uniqueness, or a Liouville theorem.\\
\bottomrule
\end{longtable}
\endgroup

% =============================================================================
\section{Exact mandate for Version V11}
\label{sec:v10-next}
% =============================================================================

\begin{programme}[Complete the physical tensor row or stop]
\label{prog:v10-next}
Preserve Versions~V1--V10 verbatim.  Do not introduce another path law,
age law, response graph, abstract concentration trichotomy, or ambient source
norm.

Proceed only in the following order.
\begin{enumerate}[label=\textup{(V11.\arabic*)},leftmargin=5.5em]
\item On the unweighted annular action-overdrive branch, compare the actual
      normalized action with the fixed-datum enstrophy-square pulse and the
      inherited critical-acceleration budget.  If no upper stock exists,
      retain action overdrive as the terminal PDE obstruction and stop.
\item On the source-action microdelocalized branch
      \eqref{eq:v10-source-microdelocalization}, use the existing material-tube
      capacity and no-cell-count discipline.  Do not replace it by a third
      probability metric.
\item On one nonzero covariance profile, evaluate the two missing tensor rows
      $\mathsf P_{\rm gr}$ and $\mathsf P_{\rm dn}$ on the same rescaled cell.
      The first is the pressure follower; the second is a divergence-null
      product defect.  A later source estimate may not erase either one.
\item Pass the full velocity product only if its local $L^2$ norm is uniformly
      bounded and its Reynolds defect is shown to vanish.  Pass the physical
      pressure only with the corresponding normalized local pressure bound.
\item Transport every represented response and return through the inherited
      endpoint maps before source novelty is asserted.  The source-visible
      covariance, pressure follower, paid stress, and return are components of
      one occurring quadratic history, not independent births.
\item Stop at the first missing terminal occupation, epsilon-regularity,
      terminal-trace, backward-uniqueness, or Liouville theorem.  A compact
      source-visible tensor alone does not license an ancient Navier--Stokes
      profile.
\end{enumerate}
The admissible endpoint is
\[
 \boxed{
 \begin{gathered}
 \text{annular action overdrive}
 \ \big|\ 
 \text{source microdelocalization}\\[0.15em]
 \text{pressure/null tensor survivor}
 \ \big|\ 
 \text{positive Euler--Reynolds defect}\\[0.15em]
 \text{full compact product profile}
 \ \big|\ 
 \text{named trace/rigidity stop}.
 \end{gathered}}
\]
\end{programme}

\begin{assessment}[Programme position after V10]
\label{ass:v10-position}
The selected annular source now has its exact source-minimal tensor carrier.
The broad fractional collar has been replaced by a target-native
pseudolocality theorem, and a bounded branch yields a nonzero covariance
profile without pretending that pressure or the full velocity product is
compact.  The next unknown is no longer how to represent the Reynolds source.
It is whether the pressure-follower, divergence-null product, and
Euler--Reynolds defect vanish on that same profile, or survive as the named
terminal obstruction.
\end{assessment}

\begin{assessment}[Append-only integrity]
\label{ass:v10-integrity}
The complete Version~V9 source preceding its sole final document terminator
is preserved verbatim.  The present material begins at
Section~\ref{sec:v10-scope}; the final command below is again unique.
\end{assessment}

\addtocontents{toc}{\protect\endgroup}
% =============================================================================
% END VERSION V10 MATERIAL -- append later versions before final terminator
% =============================================================================


\clearpage
% =============================================================================
% VERSION V11: NEW MATERIAL.  The complete V1--V10 source above is preserved
% verbatim; only its sole final document terminator was removed before append.
% =============================================================================
\hypersetup{pdftitle={NS-A Terminal Source Concentration Programme: Cumulative V11},pdfauthor={A. Pelissier}}
\makeatletter
\addtocontents{toc}{\protect\begingroup\protect\makeatletter}
\addtocontents{toc}{\protect\def\protect\l@section{\protect\@dottedtocline{1}{0em}{4.7em}}}
\addtocontents{toc}{\protect\def\protect\l@subsection{\protect\@dottedtocline{2}{1.5em}{5.3em}}}
\addtocontents{toc}{\protect\makeatother}
\makeatother
\renewcommand{\thesection}{V11.\arabic{section}}
\renewcommand{\thesubsection}{V11.\arabic{section}.\arabic{subsection}}
\renewcommand{\theequation}{V11.\arabic{equation}}
\setcounter{section}{0}
\setcounter{equation}{0}
\renewcommand{\theHsection}{V11.\arabic{section}}
\renewcommand{\theHsubsection}{V11.\arabic{section}.\arabic{subsection}}
\renewcommand{\theHtheorem}{V11.\arabic{section}.\arabic{theorem}}
\renewcommand{\theHequation}{V11.\arabic{equation}}

\begin{center}
{\LARGE\bfseries NS--A: Version V11}\\[0.5em]
{\large The Divergence-Quotient Correction, Effective Pressure from Product\\
Compactness, and the Source-Visible Reynolds-Force Stop}\\[0.5em]
{\normalsize 3 September 2026}
\end{center}
\addcontentsline{toc}{part}{Version V11: divergence quotient, effective pressure, and source-visible Reynolds force}

\begin{abstract}
Version~V10 extracted a nonzero source-visible heat-covariance profile on one
annular terminal cell and then listed the pressure follower, the
divergence-null tensor, the full product defect, and physical-pressure
compactness as separate completion rows.  The present continuation moves one
arrow upstream and asks which of those rows are actually required by the
Navier--Stokes momentum target.

The first correction concerns the action-overdrive branch.  The selected
annular Reynolds source is the annular projection of the complete heat-
covariance source and is pointwise bounded by the enstrophy.  Its spacetime
square action is therefore bounded by the enstrophy-square stock.  Hence
annular action overdrive is not a new mechanism: it is an output-resolved
subbranch of the already identified enstrophy-square pulse.  A finite
$L_t^1H_x^{-7/2}$ acceleration stock cannot control this $L_t^2L_x^2$ action;
an explicit moving-frequency pulse family proves the functional mismatch.
The overdrive branch is consequently a named terminal PDE stop rather than a
reason to add another source norm.

The second correction is tensorial.  At product regularity the complete
physical divergence has an exact quotient norm
\[
 \|\Lambda^{-1}\operatorname{div}F\|_2^2
 =\frac12\|\mathsf P_{\rm sv}F\|_2^2
  +\frac23\|\mathsf P_{\rm gr}F\|_2^2,
\]
while its kernel is precisely the divergence-null tensor sector.  Thus the
source-visible part and the scalar pressure follower are the complete
PDE-visible quotient of a trace-free product; the divergence-null component
is not a missing momentum row.  More strongly, only the source-visible part
of a Reynolds defect can obstruct the existence of an Euler pressure.  The
gradient part is absorbed into that pressure and the divergence-null part is
invisible.

This gives a sharper compactness theorem.  A uniform local $L^2$ bound for
the velocity products suffices, without any bound on the displayed physical
pressures, to pass the divided strong-coupling equations against
divergence-free tests.  A local Bogovskii construction then produces an
effective $L^2$ pressure directly from the weak product.  If
$\mathbb R=Q-V\otimes V\succeq0$ is the product defect, the limiting velocity
satisfies the stationary Euler constraint for almost every slow time exactly
when the solenoidal distribution
$P_\sigma\operatorname{div}\mathbb R$ vanishes.  The full tensor
$\mathbb R$ need not vanish.  A sequence of exact stationary Euler shears is
given whose velocities converge weakly to zero while their products converge
to a nonzero constant positive defect with zero divergence.

On the annularly microdelocalized branch, the note performs no new
concentration abstraction.  It first tests whether the positive source action
is captured by the already constructed material tube.  Captured action has an
exact bounded-overlap annular partition, diverging participation, and
cohort fusion with no cell-count factor; failure is the literal square-action
escape from the scalar material route.

The resulting endpoint is narrower than Version~V10: enstrophy-square
overdrive, material-route action escape, source-action microdelocalization,
product overdrive, a source-visible Reynolds force, a stationary Euler
quotient with unresolved terminal occupation, or a physical-pressure
identification defect.  Divergence-null stress and a nonzero but source-null
Reynolds defect are no longer promoted to independent obstructions.  No
terminal trace, ancient Liouville theorem, or global regularity conclusion is
asserted.
\end{abstract}

% =============================================================================
\section{Upstream scope and priority order}
\label{sec:v11-scope}
% =============================================================================

Retain all notation and all hypotheses of the bounded-sweeping, doubly tame,
output-annular branch of Version~V10.  In particular,
\begin{equation}
 g_n^{\rm ann}=\Psi_ng_n,
 \qquad
 \mathcal B_n^0=\int_{I_n}\|g_n^{\rm ann}(t)\|_2^2\,\dd t,
 \qquad
 D(t)=\|\nabla u(t)\|_2^2,
 \label{eq:v11-standing-action}
\end{equation}
and
\begin{equation}
 \delta_n^E:=\int_{I_n}D(t)\,\dd t\longrightarrow0.
 \label{eq:v11-standing-dissipation}
\end{equation}
On the bounded-action one-cell branch, retain the record-amplitude rescalings
$V_n,\Pi_n,\Omega_n$, the interval $J_n\to J$, the weak product limit $Q$,
and the positive defect
\begin{equation}
 \mathbb R=Q-V\otimes V\succeq0
 \label{eq:v11-standing-defect}
\end{equation}
from Version~V10 whenever the local products are bounded.

The order used below is strict.
\begin{enumerate}[label=\textup{(U\arabic*)},leftmargin=3.4em]
\item First absorb annular source-action overdrive into the pre-existing
enstrophy-square branch.
\item Next stop or localize the source-action microdelocalized branch using
the already selected material tube.
\item Only on a bounded-product cell pass to the divergence quotient of the
quadratic tensor.
\item Require vanishing only of the source-visible divergence of
$\mathbb R$, not of the full defect.
\end{enumerate}
This order avoids asking a pressure or divergence-null tensor to repair an
unbounded source action, and avoids demanding full product compactness before
the weaker momentum target has been evaluated.

\begin{firewall}[No new source carrier]
\label{fire:v11-no-new-carrier}
The objects below are deterministic quotients, weak limits, or local pressure
representatives of the exact V10 velocity product.  They do not create a new
Reynolds covariance, source birth, return vector, material path, or terminal
cell.  The V5--V10 represented-response and source-ancestry shorts remain in
force before any surviving source-visible force is called held out.
\end{firewall}

% =============================================================================
\section{Annular action overdrive is enstrophy-square overdrive}
\label{sec:v11-action-enstrophy}
% =============================================================================

For a divergence-free field $w$, write
\begin{equation}
 c[w]:=A^{-1/4}\mathcal B(w)
     =A^{-1/4}P_\sigma\operatorname{div}(w\otimes w).
 \label{eq:v11-cw}
\end{equation}
Let
\begin{equation}
 v_{n,t}:=e^{-\nu(b_n-t)A}u(t).
 \label{eq:v11-heat-field}
\end{equation}

\begin{lemma}[Exact heat-covariance source and enstrophy domination]
\label{lem:v11-source-enstrophy}
For every $t\in I_n$,
\begin{equation}
 \boxed{
 g_n(t)=e^{-\nu(b_n-t)A}c[u(t)]-c[v_{n,t}],
 \qquad
 g_n^{\rm ann}(t)=\Psi_ng_n(t).}
 \label{eq:v11-source-decomposition}
\end{equation}
There is a fixed constant $C$ such that
\begin{equation}
 \boxed{
 \|g_n^{\rm ann}(t)\|_2\le C D(t).}
 \label{eq:v11-source-enstrophy-bound}
\end{equation}
Consequently,
\begin{equation}
 \boxed{
 \mathcal B_n^0\le C^2\int_{I_n}D(t)^2\,\dd t.}
 \label{eq:v11-action-by-enstrophy-square}
\end{equation}
\end{lemma}

\begin{proof}
The heat covariance is
\[
 \mathbb C_{n,t}
 =e^{-\nu(b_n-t)A}(u\otimes u)
   -v_{n,t}\otimes v_{n,t}.
\]
The scalar heat semigroup commutes with $P_\sigma$, $A$, and
$\operatorname{div}$.  Applying
$T=A^{-1/4}P_\sigma\operatorname{div}$ therefore gives the first identity in
\eqref{eq:v11-source-decomposition}.  Since $T=T\mathsf P_{\rm sv}$ and
$T$ commutes with the annular multiplier $\Psi_n$, the V10 definition of the
annular source gives the second identity.

The standard critical product estimate is
\begin{equation}
 \|c[w]\|_2
 \le C\|w\otimes w\|_{\dot H^{1/2}}
 \le C\|w\|_6\|\Lambda^{1/2}w\|_3
 \le C\|\nabla w\|_2^2.
 \label{eq:v11-critical-product}
\end{equation}
Indeed, the middle inequality is the fractional Leibniz rule, while
$\dot H^1\hookrightarrow L^6$ and
$\dot H^{1/2}\hookrightarrow L^3$ applied to $\Lambda^{1/2}w$ give the
last inequality.  Heat contraction in $\dot H^1$ yields
$\|\nabla v_{n,t}\|_2\le\|\nabla u(t)\|_2$.  Hence
\[
 \|g_n(t)\|_2
 \le\|c[u(t)]\|_2+\|c[v_{n,t}]\|_2
 \le C D(t).
\]
The fixed smooth multiplier $\Psi_n$ is an $L^2$ contraction up to a fixed
operator norm, which proves \eqref{eq:v11-source-enstrophy-bound}.
Squaring and integrating gives
\eqref{eq:v11-action-by-enstrophy-square}.
\end{proof}

Define the normalized enstrophy-square stock
\begin{equation}
 \mathfrak E_n^{(2)}
 :=\frac{1}{\nu\mathcal R_n^2}
   \int_{I_n}D(t)^2\,\dd t.
 \label{eq:v11-normalized-enstrophy-square}
\end{equation}

\begin{theorem}[The V10 action branch is subordinate to one PDE stock]
\label{thm:v11-action-subordinate}
The V10 annular leverage satisfies
\begin{equation}
 \boxed{
 \mathfrak L_n^{\rm ann}
 \le C^2\mathfrak E_n^{(2)}.}
 \label{eq:v11-leverage-subordinate}
\end{equation}
Therefore:
\begin{enumerate}[label=\textup{(A\arabic*)},leftmargin=3.4em]
\item if $\mathfrak L_n^{\rm ann}\to\infty$, then
      $\mathfrak E_n^{(2)}\to\infty$ after a subsequence;
\item if $\sup_n\mathfrak E_n^{(2)}<\infty$, then the bounded-action branch
      \eqref{eq:v10-annular-action-bounded} holds automatically;
\item on the overdrive branch,
      \begin{equation}
       \boxed{
       \frac{\delta_n^E}{\nu\mathcal R_n^2}
       \sup_{t\in I_n}D(t)
       \ge\mathfrak E_n^{(2)}\longrightarrow\infty.}
       \label{eq:v11-enstrophy-spike-overdrive}
      \end{equation}
\end{enumerate}
Thus annular source-action overdrive is an output-resolved
\emph{enstrophy-square overdrive}; it is not a third terminal mechanism.
\end{theorem}

\begin{proof}
Divide \eqref{eq:v11-action-by-enstrophy-square} by
$\nu\mathcal R_n^2$ to obtain
\eqref{eq:v11-leverage-subordinate}.  The first two assertions are immediate.
For the third, use
\[
 \int_{I_n}D^2
 \le\left(\sup_{I_n}D\right)\int_{I_n}D
 =\left(\sup_{I_n}D\right)\delta_n^E
\]
and divide by $\nu\mathcal R_n^2$.  The source-action lower bound ensures
$\delta_n^E>0$ for every late $n$.
\end{proof}

The inherited critical-acceleration stock is much weaker in both time
integrability and spatial order.  The following functional counterexample
shows that no direct norm comparison can close
\eqref{eq:v11-enstrophy-spike-overdrive}.

\begin{countertheorem}[A negative-Sobolev $L^1$ stock does not control moving-annulus $L^2$ action]
\label{cth:v11-negative-stock-no-action}
There is a smooth divergence-free field
$f\in C^\infty((0,1)\times\Torus^3)$, written as a sum of disjoint-time
annular packets $f_j$, such that
\begin{equation}
 \boxed{
 \sum_{j=1}^{\infty}
 \int_0^1\|f_j(t)\|_{H^{-7/2}}\,\dd t<\infty,}
 \label{eq:v11-negative-stock-finite}
\end{equation}
but
\begin{equation}
 \boxed{
 \int_0^1\|f_j(t)\|_2^2\,\dd t\longrightarrow\infty.}
 \label{eq:v11-L2-packets-diverge}
\end{equation}
The Fourier support of $f_j$ is contained in the single annulus
$|k|=2^j$.
\end{countertheorem}

\begin{proof}
Use normalized Haar measure on $\Torus^3$ and put
\[
 \phi_j(x)=\sqrt2\sin(2^jx_1)e_2,
 \qquad \|\phi_j\|_2=1,
 \qquad \operatorname{div}\phi_j=0.
\]
Choose a nonzero $\eta\in C_c^\infty(0,1)$ and pairwise disjoint intervals
$J_j=(t_j,t_j+\delta_j)$ with
$\delta_j=2^{-6j}$; their total length is finite.  Define
\[
 f_j(t,x)=2^{4j}\eta\!\left(\frac{t-t_j}{\delta_j}\right)\phi_j(x)
\]
on $J_j$ and zero elsewhere.  All time derivatives vanish to every order at
the endpoints because $\eta$ is compactly supported inside $(0,1)$, so the
sum $f=\sum_jf_j$ is smooth on $(0,1)\times\Torus^3$ after placing the
intervals with their only accumulation at an endpoint.

Since $\phi_j$ has frequency $2^j$,
\[
 \int\|f_j(t)\|_{H^{-7/2}}\,\dd t
 \le C\,2^{4j}2^{-6j}2^{-7j/2}
 =C2^{-11j/2},
\]
which is summable.  On the other hand,
\[
 \int\|f_j(t)\|_2^2\,\dd t
 =\|\eta\|_2^2\,2^{8j}2^{-6j}
 =\|\eta\|_2^2\,2^{2j}\longrightarrow\infty.
\]
\end{proof}

\begin{firewall}[Meaning of the countertheorem]
\label{fire:v11-negative-stock-scope}
Countertheorem~\ref{cth:v11-negative-stock-no-action} is a functional-
analytic obstruction to a norm implication.  It is not a construction of a
Navier--Stokes source history.  It proves that the inherited
$L_t^1H_x^{-7/2}$ acceleration control cannot, by itself, be substituted for
the missing $L_t^2L_x^2$ annular action or enstrophy-square estimate.
Consequently branch \textup{(A1)} is a legitimate terminal PDE stop.
\end{firewall}

% =============================================================================
\section{Material-tube capture or exact annular participation}
\label{sec:v11-microdelocalization}
% =============================================================================

Work now on the bounded-action branch and suppose the source-action law
$\mu_n^{\rm sv}$ of \eqref{eq:v10-source-action-law} satisfies the
microdelocalization alternative
\eqref{eq:v10-source-microdelocalization}.  Retain the V8--V9 material tube
\begin{equation}
 \mathfrak T_n
 :=\{(t,x):t\in I_n,\ x\in\mathcal M_n(t)\}
 \label{eq:v11-material-tube}
\end{equation}
and define its source-action capture fraction
\begin{equation}
 \eta_n^{\rm tube}:=\mu_n^{\rm sv}(\mathfrak T_n).
 \label{eq:v11-tube-capture-fraction}
\end{equation}

\begin{theorem}[Square-action escape or diffuse material-tube participation]
\label{thm:v11-tube-participation}
After passage to a subsequence, exactly one of the following priority
branches occurs.
\begin{enumerate}[label=\textup{(M\arabic*)},leftmargin=3.4em]
\item \emph{Square-action route escape:}
      \begin{equation}
       \boxed{\eta_n^{\rm tube}\longrightarrow0.}
       \label{eq:v11-square-action-escape}
      \end{equation}
      The scalar material route does not capture a positive fraction of the
      selected source square action.
\item \emph{Diffuse material-tube action:} for some $\eta_*>0$,
      $\eta_n^{\rm tube}\ge\eta_*$ for all late $n$.  Put
      \begin{equation}
       \overline\mu_n
       :=(\eta_n^{\rm tube})^{-1}
         \mu_n^{\rm sv}|_{\mathfrak T_n}.
       \label{eq:v11-restricted-action-law}
      \end{equation}
      There is a finite quadratic partition
      $(\chi_{n,j})_{j\in\mathcal J_n}$ with
      \begin{equation}
       \sum_{j\in\mathcal J_n}\chi_{n,j}^2=1,
       \qquad
       \operatorname{diam}\supp\chi_{n,j}\le C\Gamma_n^{-1},
       \label{eq:v11-annular-partition}
      \end{equation}
      and, for
      \begin{equation}
       p_{n,j}
       :=\int_{\mathfrak T_n}\chi_{n,j}(x)^2\,\dd\overline\mu_n(t,x),
       \label{eq:v11-cell-probabilities}
      \end{equation}
      one has
      \begin{equation}
       \boxed{
       \max_jp_{n,j}\longrightarrow0,
       \qquad
       \mathfrak N_n^{\rm tube}
       :=\left(\sum_jp_{n,j}^2\right)^{-1}
       \longrightarrow\infty.}
       \label{eq:v11-participation-diverges}
      \end{equation}
      For every $\theta\in(0,1)$ there are cohorts
      $\mathcal C_n\subseteq\mathcal J_n$ such that
      \begin{equation}
       \sum_{j\in\mathcal C_n}p_{n,j}\longrightarrow\theta.
       \label{eq:v11-cohort-fraction}
      \end{equation}
      With
      $\chi_{n,\mathcal C_n}^2
       :=\sum_{j\in\mathcal C_n}\chi_{n,j}^2$, their source action is the
      exact sum of the cell actions; no factor $|\mathcal C_n|$ occurs.
\end{enumerate}
\end{theorem}

\begin{proof}
First pass to a subsequence on which
$\eta_n^{\rm tube}\to0$ or it has a positive lower bound.  This gives the
first branch or the entrance to the second.

For the partition, take a maximal $c\Gamma_n^{-1}$-separated set on the
torus.  Fixed enlargements of its balls cover the torus and have overlap
bounded only by the dimension.  Choose nonnegative smooth bumps
$\widetilde\chi_{n,j}$ which are one on the smaller cover and supported in
the enlarged balls, and normalize them by
\[
 \chi_{n,j}
 =\frac{\widetilde\chi_{n,j}}
        {(\sum_k\widetilde\chi_{n,k}^2)^{1/2}}.
\]
This gives \eqref{eq:v11-annular-partition}.

Every support is contained in a ball of radius $C\Gamma_n^{-1}$, so the
V10 concentration function gives
\[
 p_{n,j}
 \le\frac{\vartheta_n^{\rm sv}(C)}{\eta_n^{\rm tube}}
 \le\eta_*^{-1}\vartheta_n^{\rm sv}(C)\longrightarrow0.
\]
Because $\sum_jp_{n,j}=1$,
\[
 \sum_jp_{n,j}^2\le\max_jp_{n,j},
\]
which proves \eqref{eq:v11-participation-diverges}.

For a fixed $\theta$, add indices one at a time until their accumulated mass
first reaches $\theta$.  The overshoot is at most $\max_jp_{n,j}$, proving
\eqref{eq:v11-cohort-fraction}.  Finally,
\[
 \int\chi_{n,\mathcal C_n}^2
      |\omega_ng_n^{\rm ann}|^2
 =\sum_{j\in\mathcal C_n}
   \int\chi_{n,j}^2|\omega_ng_n^{\rm ann}|^2
\]
by the definition of the fused quadratic weight.  This is exact and contains
no cell-count multiplier.
\end{proof}

\begin{assessment}[The branch is already terminal at annular resolution]
\label{ass:v11-microdeloc-stop}
The material tube shrinks in the original coordinates, while its annular
participation diverges.  A fused cohort can carry any prescribed positive
fraction of the action, but its relative diameter need not be bounded.  Thus
cohort fusion removes multiplicity from estimates without manufacturing one
compact annular cell.  Excluding this branch requires a genuine local
capacity, epsilon-regularity, or one-scale compactness theorem; another
probability metric would not supply that PDE input.
\end{assessment}

% =============================================================================
\section{The product-regularity divergence quotient}
\label{sec:v11-divergence-quotient}
% =============================================================================

Version~V10 used the order-$1/2$ source norm appropriate to the Reynolds
action.  Weak velocity products live naturally in $L^2$, so the relevant
completion now occurs one derivative lower.  Let
$F\in L_0^2(\Torus^3;\operatorname{Sym}_0(3))$ and retain the Fourier
projections $\mathsf P_{\rm sv}$, $\mathsf P_{\rm gr}$, and
$\mathsf P_{\rm dn}$ from Version~V10.  Define
\begin{equation}
 \mathscr S_F:=\Lambda^{-1}P_\sigma\operatorname{div}F
 \label{eq:v11-product-source}
\end{equation}
and let $p_F$ be the mean-zero scalar satisfying
\begin{equation}
 -\Delta p_F=\partial_i\partial_jF_{ij}.
 \label{eq:v11-product-pressure}
\end{equation}

\begin{theorem}[Exact divergence quotient at product regularity]
\label{thm:v11-divergence-quotient}
For every such $F$,
\begin{equation}
 \boxed{
 \Lambda^{-1}\operatorname{div}F
 =\mathscr S_F-\nabla\Lambda^{-1}p_F.}
 \label{eq:v11-divergence-Helmholtz}
\end{equation}
The two terms are orthogonal in $L^2$, and
\begin{equation}
 \boxed{
 \begin{aligned}
 \|\Lambda^{-1}\operatorname{div}F\|_2^2
 &=\|\mathscr S_F\|_2^2+\|p_F\|_2^2\\
 &=\frac12\|\mathsf P_{\rm sv}F\|_2^2
   +\frac23\|\mathsf P_{\rm gr}F\|_2^2.
 \end{aligned}}
 \label{eq:v11-divergence-quotient-norm}
\end{equation}
Moreover,
\begin{equation}
 \boxed{
 \Ker(\Lambda^{-1}\operatorname{div})
 =\operatorname{Ran}\mathsf P_{\rm dn}.}
 \label{eq:v11-divergence-kernel}
\end{equation}
Consequently the map
\begin{equation}
 [F]\longmapsto(\mathscr S_F,p_F)
 \label{eq:v11-quotient-map}
\end{equation}
is an isometric isomorphism from the quotient by the divergence-null sector,
with the norm in \eqref{eq:v11-divergence-quotient-norm}, onto the product of
the solenoidal vector and scalar pressure-follower spaces.  The unique
minimum-$L^2$ representative of $[F]$ is
\begin{equation}
 \boxed{F_{\min}=\mathsf P_{\rm sv}F+\mathsf P_{\rm gr}F.}
 \label{eq:v11-minimum-divergence-representative}
\end{equation}
\end{theorem}

\begin{proof}
At a nonzero Fourier mode $k$, put $n=k/|k|$ and write, as in Version~V10,
\[
 v=P_n\widehat F(k)n,
 \qquad
 q=n^{\mathsf T}\widehat F(k)n.
\]
Then
\[
 \widehat{\Lambda^{-1}\operatorname{div}F}(k)
 =i(v+qn),
 \qquad
 \widehat{\mathscr S_F}(k)=iv,
 \qquad
 \widehat p_F(k)=-q.
\]
Since
$\widehat{\nabla\Lambda^{-1}p_F}(k)=-iqn$, this proves
\eqref{eq:v11-divergence-Helmholtz}.  The vectors $v$ and $qn$ are
orthogonal.  Furthermore,
\[
 |\mathsf P_{\rm sv}(n)\widehat F(k)|^2=2|v|^2,
 \qquad
 |\mathsf P_{\rm gr}(n)\widehat F(k)|^2=\frac32|q|^2.
\]
Summation proves \eqref{eq:v11-divergence-quotient-norm}.

The divergence vanishes exactly when $v=q=0$, equivalently when
$\widehat F(k)n=0$.  This is precisely the V10 characterization of the
$\mathsf P_{\rm dn}$ range, proving \eqref{eq:v11-divergence-kernel}.
Orthogonality of the three projections makes
$\mathsf P_{\rm sv}F+\mathsf P_{\rm gr}F$ the unique representative
orthogonal to the kernel and hence the unique minimum-$L^2$ representative.

To see surjectivity of \eqref{eq:v11-quotient-map}, take a mean-zero
solenoidal vector $s$ and a mean-zero scalar $p$.  At mode $k$, set
\[
 v_k=-i\widehat s(k),
 \qquad q_k=-\widehat p(k),
 \qquad
 \widehat F(k)=2n_k\odot v_k+q_kE_{n_k}.
\]
This is the minimum representative and maps to $(s,p)$.
\end{proof}

\begin{corollary}[Only one tensor component obstructs an Euler pressure]
\label{cor:v11-only-source-obstructs}
Let $Q\in L^2(\Torus^3;\operatorname{Sym}(3))$ and write
$F=Q^\circ-\int_{\Torus^3}Q^\circ\,\dd x$.  Then
\begin{equation}
 \boxed{
 P_\sigma\operatorname{div}Q
 =P_\sigma\operatorname{div}(\mathsf P_{\rm sv}F),}
 \label{eq:v11-only-source-divergence}
\end{equation}
and
\begin{equation}
 \operatorname{div}Q
 =P_\sigma\operatorname{div}Q
  +\nabla\left(\frac13\operatorname{tr}Q-p_F\right).
 \label{eq:v11-full-product-divergence}
\end{equation}
Thus:
\begin{enumerate}[label=\textup{(i)},leftmargin=2.5em]
\item $\mathsf P_{\rm dn}F$ never enters the momentum equation;
\item $\mathsf P_{\rm gr}F$ changes only the pressure representative;
\item the sole obstruction to writing $\operatorname{div}Q$ as a gradient is
      the source-visible component $\mathsf P_{\rm sv}F$.
\end{enumerate}
\end{corollary}

\begin{proof}
The isotropic part of $Q$ has gradient divergence.  Version~V10 gives
$P_\sigma\operatorname{div}F
 =P_\sigma\operatorname{div}(\mathsf P_{\rm sv}F)$,
$\operatorname{div}\mathsf P_{\rm dn}F=0$, and
$\operatorname{div}\mathsf P_{\rm gr}F=-\nabla p_F$.
Adding the isotropic gradient gives
\eqref{eq:v11-full-product-divergence}.
\end{proof}

\begin{assessment}[Target correction to the V10 tensor ledger]
\label{ass:v11-tensor-target-correction}
The pressure follower remains necessary when the target is convergence of
the \emph{physical pressure itself}.  It is not an obstruction to the
existence of some pressure solving the weak Euler momentum equation.  The
divergence-null sector is not required for either target.  Requiring both
rows to vanish before passing the momentum equation would therefore inflate
the target from a PDE quotient to reconstruction of the complete quadratic
tensor.
\end{assessment}

% =============================================================================
\section{Product compactness gives an effective pressure without physical-pressure compactness}
\label{sec:v11-effective-pressure}
% =============================================================================

We use the following standard right inverse of divergence on a ball.  It is
recorded explicitly because it is the only local pressure tool needed below.

\begin{lemma}[Bogovskii right inverse on a ball]
\label{lem:v11-Bogovskii}
For every ball $B\subset\R^3$ there is a bounded linear map
\begin{equation}
 \mathfrak B_B:L_0^2(B)\longrightarrow H_0^1(B;\R^3)
 \label{eq:v11-Bogovskii-map}
\end{equation}
such that
\begin{equation}
 \boxed{
 \operatorname{div}\mathfrak B_Bq=q,
 \qquad
 \|\nabla\mathfrak B_Bq\|_{L^2(B)}
 \le C_B\|q\|_{L^2(B)}.}
 \label{eq:v11-Bogovskii-bound}
\end{equation}
\end{lemma}

\begin{proof}
After translation and dilation it suffices to work on the unit ball.  Choose
$\rho\in C_c^\infty(B)$ with $\int\rho=1$.  For smooth mean-zero $q$, the
classical Bogovskii formula
\[
 (\mathfrak B_Bq)(x)
 =\int_Bq(y)(x-y)
   \int_1^\infty\rho(y+r(x-y))r^2\,\dd r\,\dd y
\]
has support in $B$ and direct differentiation, using $\int q=0$, gives
$\operatorname{div}\mathfrak B_Bq=q$.  The derivatives of its kernel are
Calder\'on--Zygmund kernels plus integrable remainders; hence the $L^2$
bound in \eqref{eq:v11-Bogovskii-bound} follows from the standard singular-
integral estimate.  Density extends the map to $L_0^2(B)$.  Translation and
dilation preserve the asserted form of the estimate, with a constant
depending only on the ball's shape.
\end{proof}

\begin{theorem}[Pressure-free product passage and automatic effective pressure]
\label{thm:v11-product-only-pressure}
Retain the V10 one-cell branch and suppose that, for some
$R_1>R_0+2$,
\begin{equation}
 \sup_n\int_{J_n}\int_{B(0,R_1)}|V_n|^4<\infty.
 \label{eq:v11-local-product-bound}
\end{equation}
No assumption is made on $\Pi_n$.  After passage to a subsequence, for every
$J'\Subset J$ and $B\Subset B(0,R_1)$ there are weak limits
\begin{equation}
 V_n\rightharpoonup V\quad\hbox{in }L^4(J'\times B),
 \qquad
 V_n\otimes V_n\rightharpoonup Q\quad\hbox{in }L^2(J'\times B),
 \label{eq:v11-product-weak-limits}
\end{equation}
such that
\begin{equation}
 \boxed{
 \int_{J'\times B}Q:\nabla\phi=0}
 \label{eq:v11-pressure-free-constraint}
\end{equation}
for every
$\phi\in C_c^\infty(J'\times B;\R^3)$ with
$\operatorname{div}\phi=0$.

There is a unique zero-spatial-mean
\begin{equation}
 \pi_Q\in L^2(J';L_0^2(B))
 \label{eq:v11-effective-pressure-space}
\end{equation}
such that
\begin{equation}
 \boxed{
 \operatorname{div}Q+\nabla\pi_Q=0
 \quad\hbox{in }\mathcal D'(J'\times B),}
 \label{eq:v11-effective-pressure-equation}
\end{equation}
and
\begin{equation}
 \boxed{
 \|\pi_Q\|_{L^2(J'\times B)}
 \le C_B\|Q\|_{L^2(J'\times B)}.}
 \label{eq:v11-effective-pressure-bound}
\end{equation}
Thus a separate bound for the displayed physical pressures is unnecessary
for existence of a local effective pressure on the weak-product branch.
\end{theorem}

\begin{proof}
The compact local product bound gives the weak convergences in
\eqref{eq:v11-product-weak-limits}.  Divide the strong-coupling equation
\eqref{eq:v10-strong-coupling-equation} by $\alpha_n$ and test it against a
compactly supported divergence-free $\phi$.  The pressure term vanishes.
After integration by parts, the time and viscous terms are bounded by
$C_\phi\alpha_n^{-1}\|V_n\|_{L^1(\supp\phi)}$ and tend to zero because
$\alpha_n\to\infty$ and the $L^4$ norms are bounded.  Weak convergence of
the products gives \eqref{eq:v11-pressure-free-constraint}.

Fix a time $s$ for which $Q(s)\in L^2(B)$ and define a bounded functional on
$L_0^2(B)$ by
\begin{equation}
 \int_B\pi_Q(s)q
 :=-\int_BQ(s):\nabla\mathfrak B_Bq.
 \label{eq:v11-pressure-Riesz-definition}
\end{equation}
Lemma~\ref{lem:v11-Bogovskii} and Riesz representation give a unique
zero-mean $\pi_Q(s)$ with
$\|\pi_Q(s)\|_2\le C_B\|Q(s)\|_2$.  The construction is bounded linear in
$Q(s)$, hence measurable in $s$ and gives
\eqref{eq:v11-effective-pressure-bound}.

For $\varphi\in C_c^\infty(B;\R^3)$ put
$q=\operatorname{div}\varphi$.  Its integral is zero, and
$\varphi-\mathfrak B_Bq$ is divergence free.  By
\eqref{eq:v11-pressure-free-constraint}, first for smooth time-separated
tests and then by density,
\[
 \int_BQ:\nabla\varphi
 =\int_BQ:\nabla\mathfrak B_Bq
 =-\int_B\pi_Qq.
\]
This is exactly
$\operatorname{div}Q+\nabla\pi_Q=0$.  Zero spatial mean gives uniqueness.
\end{proof}

\begin{corollary}[Physical-pressure overdrive is an identification defect]
\label{cor:v11-physical-pressure-identification}
If, in addition, spatially normalized physical pressures satisfy
\begin{equation}
 \Pi_n-c_n\rightharpoonup\Pi
 \quad\hbox{in }L^2(J'\times B)
 \label{eq:v11-physical-pressure-weak}
\end{equation}
for some functions $c_n(s)$, then
$\Pi-\pi_Q$ is spatially constant almost everywhere in time.  If
\eqref{eq:v11-physical-pressure-weak} fails, the weak-product equation
\eqref{eq:v11-effective-pressure-equation} still holds, but identification
of its pressure with the limit of the original Navier--Stokes pressures is
withheld.
\end{corollary}

\begin{proof}
Passing the full divided equation under
\eqref{eq:v11-physical-pressure-weak} gives
$\operatorname{div}Q+\nabla\Pi=0$.  Subtract
\eqref{eq:v11-effective-pressure-equation}.  A distribution with zero
spatial gradient is a function of time only.
\end{proof}

% =============================================================================
\section{The exact Reynolds-defect criterion for a stationary Euler quotient}
\label{sec:v11-Reynolds-force}
% =============================================================================

Retain the product-only limit and put
\begin{equation}
 \mathbb R:=Q-V\otimes V\succeq0.
 \label{eq:v11-Reynolds-defect}
\end{equation}
For a compact cylinder $J'\times B$, define its source-visible force by the
functional
\begin{equation}
 \langle\mathfrak f_{\mathbb R},\phi\rangle
 :=-\int_{J'\times B}\mathbb R:\nabla\phi,
 \qquad
 \operatorname{div}\phi=0.
 \label{eq:v11-source-visible-force}
\end{equation}
It belongs to the dual of
$L^2(J';H_{0,\sigma}^1(B))$.

\begin{theorem}[Source-visible Reynolds force or stationary Euler]
\label{thm:v11-force-or-Euler}
Exactly one of the following occurs on each compact subcylinder.
\begin{enumerate}[label=\textup{(R\arabic*)},leftmargin=3.4em]
\item \emph{Source-visible Reynolds-force survivor:}
      \begin{equation}
       \boxed{\mathfrak f_{\mathbb R}\ne0.}
       \label{eq:v11-Reynolds-force-survives}
      \end{equation}
      This is the part of the weak product defect which cannot be absorbed
      into pressure.
\item \emph{Stationary Euler quotient:}
      $\mathfrak f_{\mathbb R}=0$.  There is
      $\pi_{\mathbb R}\in L^2(J';L_0^2(B))$ such that
      \begin{equation}
       \operatorname{div}\mathbb R+\nabla\pi_{\mathbb R}=0,
       \label{eq:v11-defect-pressure}
      \end{equation}
      and the weak velocity satisfies
      \begin{equation}
       \boxed{
       \operatorname{div}(V\otimes V)
       +\nabla\pi_E=0,
       \qquad
       \pi_E:=\pi_Q-\pi_{\mathbb R}.}
       \label{eq:v11-stationary-Euler}
      \end{equation}
      Equivalently, for almost every slow time $s\in J'$, the spatial
      field $V(s,\cdot)$ is a stationary weak Euler field.  No independence
      of $s$ is asserted.
\end{enumerate}
In particular, neither $\mathbb R=0$ nor convergence of the physical pressure
is necessary for the stationary Euler conclusion.
\end{theorem}

\begin{proof}
If the functional in \eqref{eq:v11-source-visible-force} vanishes, the same
Bogovskii construction used in Theorem~\ref{thm:v11-product-only-pressure},
with $Q$ replaced by $\mathbb R$, gives
\eqref{eq:v11-defect-pressure}.  Subtract this identity from
\eqref{eq:v11-effective-pressure-equation}, using
$Q=V\otimes V+\mathbb R$, to obtain
\eqref{eq:v11-stationary-Euler}.  Conversely, if
\eqref{eq:v11-stationary-Euler} holds for some pressure, subtraction from the
$Q$ equation shows that $\operatorname{div}\mathbb R$ is a gradient and
therefore annihilates every divergence-free test field.  The alternatives
are exhaustive.
\end{proof}

On the periodic whole-space chart, the criterion has an exact tensor norm.

\begin{proposition}[Product-defect source norm on the torus]
\label{prop:v11-defect-source-norm}
Let $\mathbb R\in L^2(\Torus^3;\operatorname{Sym}(3))$.  After removing the
spatial mean of $\mathbb R^\circ$,
\begin{equation}
 \boxed{
 \|P_\sigma\operatorname{div}\mathbb R\|_{\dot H^{-1}}^2
 =\frac12\|\mathsf P_{\rm sv}\mathbb R^\circ\|_2^2.}
 \label{eq:v11-defect-source-norm}
\end{equation}
If the left side vanishes, the pressure absorbed from the defect may be
chosen as
\begin{equation}
 \pi_{\mathbb R}
 =p_{\mathbb R^\circ}-\frac13\operatorname{tr}\mathbb R
 \label{eq:v11-defect-pressure-formula}
\end{equation}
up to a spatial constant.  The component
$\mathsf P_{\rm dn}\mathbb R^\circ$ is absent from both formulas.
\end{proposition}

\begin{proof}
At a nonzero Fourier mode,
$P_\sigma\operatorname{div}\mathbb R=i|k|v_k$, where
$v_k=P_{n_k}\widehat{\mathbb R^\circ}(k)n_k$.  The homogeneous
$H^{-1}$ norm therefore contributes $|v_k|^2$, whereas the source-visible
tensor norm contributes $2|v_k|^2$.  This proves
\eqref{eq:v11-defect-source-norm}.  Formula
\eqref{eq:v11-full-product-divergence} gives
\[
 \operatorname{div}\mathbb R
 =\nabla\left(\frac13\operatorname{tr}\mathbb R
              -p_{\mathbb R^\circ}\right)
\]
on the zero source-visible branch, which is equivalent to
\eqref{eq:v11-defect-pressure-formula}.
\end{proof}

\begin{countertheorem}[A nonzero positive Reynolds defect need not obstruct Euler]
\label{cth:v11-nonzero-defect-harmless}
On $\Torus^3$, define
\begin{equation}
 V_m(x)=\sqrt2\sin(mx_1)e_2.
 \label{eq:v11-shear-sequence}
\end{equation}
Then every $V_m$ is a smooth mean-zero divergence-free stationary Euler
field with pressure zero,
\begin{equation}
 V_m\rightharpoonup0\quad\hbox{in }L^4,
 \qquad
 V_m\otimes V_m\rightharpoonup e_2\otimes e_2
 \quad\hbox{in }L^2.
 \label{eq:v11-shear-limits}
\end{equation}
Thus the limiting defect is
\begin{equation}
 \boxed{\mathbb R=e_2\otimes e_2\ne0,
 \qquad \mathbb R\succeq0,
 \qquad \operatorname{div}\mathbb R=0.}
 \label{eq:v11-harmless-defect}
\end{equation}
Therefore the condition $\mathbb R=0$ is strictly stronger than the
source-visible condition required by
Theorem~\ref{thm:v11-force-or-Euler}.
\end{countertheorem}

\begin{proof}
The field has only a second component and is independent of $x_2$, so
$\operatorname{div}V_m=0$ and
$(V_m\cdot\nabla)V_m=0$.  Hence it is stationary Euler with zero pressure.
The first convergence follows from oscillation.  Since
\[
 V_m\otimes V_m
 =2\sin^2(mx_1)e_2\otimes e_2
 =(1-\cos(2mx_1))e_2\otimes e_2,
\]
the second convergence follows from the Riemann--Lebesgue lemma.  The limit
is constant, positive, nonzero, and divergence free.
\end{proof}

\begin{firewall}[What remains physical]
\label{fire:v11-effective-versus-physical-pressure}
Theorem~\ref{thm:v11-force-or-Euler} produces the minimum pressure needed by
the weak momentum equation.  It does not identify that pressure with the
limit of the original Navier--Stokes pressures, reconstruct the complete
quadratic stress, or erase the paid-stress and return rows.  Those stronger
targets still require their own same-history pressure, selector, and
transport estimates.  The correction is only that they are not premises for
existence of a weak stationary Euler quotient.
\end{firewall}

% =============================================================================
\section{Refined endpoint of the V10 profile branch}
\label{sec:v11-profile-endpoint}
% =============================================================================

Combine Theorems~\ref{thm:v10-covariance-profile},
\ref{thm:v11-product-only-pressure}, and
\ref{thm:v11-force-or-Euler}.  The nonzero source-visible covariance profile
$H=T_yF$ remains as constructed in Version~V10.  The velocity/product side
has the following sharper alternative.

\begin{theorem}[Covariance-only profile, source-visible force, or Euler quotient]
\label{thm:v11-profile-alternative}
On the V10 one-cell bounded-action branch, after the action and
microdelocalization priorities above have been removed, exactly one of the
following occurs.
\begin{enumerate}[label=\textup{(P\arabic*)},leftmargin=3.4em]
\item \emph{Normalized product overdrive:}
      \eqref{eq:v10-product-overdrive} holds.  No pressure or tensor
      completion is attempted.
\item The products are locally bounded and the source-visible Reynolds force
      \eqref{eq:v11-Reynolds-force-survives} is nonzero.
\item The products are locally bounded, the source-visible Reynolds force
      vanishes, and $V\ne0$ in spacetime.  Then $V$ is not identically zero
      and, for almost every slow time, $V(s,\cdot)$ is a stationary weak
      Euler field with effective pressure \eqref{eq:v11-stationary-Euler}.
\item The products are locally bounded, the source-visible Reynolds force
      vanishes, and $V=0$.  Then $V$ is the trivial stationary Euler field,
      while the nonzero V10 tensor $F$ and source $H$ survive as a
      covariance-only oscillation/concentration profile.  A terminal
      occupation or trace theorem is required to promote it to a nonzero
      velocity profile.
\end{enumerate}
A separate physical-pressure overdrive may still obstruct identification of
the original pressure, but it is not a fifth branch for the effective weak
momentum equation.
\end{theorem}

\begin{proof}
Apply the product bounded/unbounded subsequence alternative.  On the bounded
branch, Theorem~\ref{thm:v11-product-only-pressure} constructs $Q$ and its
effective pressure without using $\Pi_n$.  Theorem~\ref{thm:v11-force-or-Euler}
then gives the source-force versus Euler alternative.  Split the Euler branch
according to whether $V$ is zero.  In the zero case, the nonzero pairing
\eqref{eq:v10-limit-nonzero-pairing} still gives $H\ne0$ and $F\ne0$, so the
remaining object is a covariance profile rather than a velocity profile.
\end{proof}

\begin{assessment}[The exact upstream gain]
\label{ass:v11-upstream-gain}
Version~V10 required the full Reynolds defect to vanish before reporting an
Euler profile and listed the divergence-null tensor as an open completion
row.  Version~V11 proves that both requirements are stronger than the
momentum target.  The proper hierarchy is
\[
 \boxed{
 \begin{gathered}
 \text{local product bound}
 \Longrightarrow
 \text{effective pressure and }P_\sigma\operatorname{div}Q=0,\\
 P_\sigma\operatorname{div}\mathbb R=0
 \Longrightarrow
 \text{stationary Euler for }V,\\
 \mathbb R=0
 \Longrightarrow
 \text{stronger full-product compactness}.
 \end{gathered}}
\]
Only the first two arrows belong to the weak momentum target.  The last is a
strictly stronger product-reconstruction demand.
\end{assessment}

% =============================================================================
\section{Version V11 integrated theorem and stopping boundary}
\label{sec:v11-integrated}
% =============================================================================

\begin{theorem}[Version V11 divergence-quotient and Reynolds-force theorem]
\label{thm:v11-integrated}
Preserving every theorem and stated limitation of Versions~V1--V10,
Version~V11 proves the following.
\begin{enumerate}[label=\textup{(V11.\arabic*)},leftmargin=5.8em]
\item The complete output-annular source satisfies
      $\|g_n^{\rm ann}(t)\|_2\le CD(t)$, so its unweighted square action is
      bounded by the enstrophy-square stock.
\item Annular action overdrive forces normalized enstrophy-square overdrive
      and the spike estimate \eqref{eq:v11-enstrophy-spike-overdrive}.  It is
      not an independent source mechanism.
\item A finite $L_t^1H_x^{-7/2}$ stock cannot dominate moving-annulus
      $L_t^2L_x^2$ action, even at the level of smooth divergence-free test
      packets.  The inherited critical-acceleration stock therefore supplies
      no missing upper bound by norm comparison alone.
\item On source microdelocalization, the positive action either escapes the
      existing material tube or has diverging exact annular participation
      inside that tube.  Quadratic cohort fusion carries no cell-count
      factor but does not create a bounded-diameter cell.
\item At product regularity, the tensor quotient by divergence-null stress is
      exactly the pair consisting of the solenoidal source and scalar pressure
      follower, with norm \eqref{eq:v11-divergence-quotient-norm}.
\item The divergence-null tensor is not a missing momentum row.  The gradient
      tensor changes only the pressure representative.  Only the
      source-visible tensor can obstruct a pressure representation.
\item A uniform local $L^2$ product bound passes the strong-coupling equations
      against divergence-free tests without any physical-pressure bound.  A
      Bogovskii short reconstructs an effective $L^2$ pressure with
      \eqref{eq:v11-effective-pressure-bound}.
\item For $\mathbb R=Q-V\otimes V\succeq0$, the exact obstruction is
      $P_\sigma\operatorname{div}\mathbb R$.  Its vanishing gives a stationary
      Euler quotient after pressure absorption, even when $\mathbb R\ne0$.
\item Exact stationary Euler shears have weak product limits with a nonzero
      positive divergence-null defect, proving that full-defect vanishing is
      not necessary.
\item The V10 nonzero covariance profile therefore ends in product overdrive,
      a source-visible Reynolds force, a nonzero stationary Euler quotient,
      or a covariance-only profile requiring terminal occupation.  Physical-
      pressure convergence remains a stronger identification row.
\end{enumerate}
No item proves terminal occupation, epsilon regularity, a terminal trace,
backward uniqueness, an ancient Liouville theorem, or global
three-dimensional regularity.
\end{theorem}

\begin{proof}
Items~\textup{(V11.1)--(V11.3)} are
Lemma~\ref{lem:v11-source-enstrophy},
Theorem~\ref{thm:v11-action-subordinate}, and
Countertheorem~\ref{cth:v11-negative-stock-no-action}.  Item~\textup{(V11.4)}
is Theorem~\ref{thm:v11-tube-participation}.  Items
\textup{(V11.5)--(V11.6)} are
Theorem~\ref{thm:v11-divergence-quotient} and
Corollary~\ref{cor:v11-only-source-obstructs}.  Item~\textup{(V11.7)} is
Theorem~\ref{thm:v11-product-only-pressure}.  Items
\textup{(V11.8)--(V11.9)} are
Theorem~\ref{thm:v11-force-or-Euler},
Proposition~\ref{prop:v11-defect-source-norm}, and
Countertheorem~\ref{cth:v11-nonzero-defect-harmless}.  The final item is
Theorem~\ref{thm:v11-profile-alternative}.
\end{proof}

\begingroup\small
\begin{longtable}{p{0.25\textwidth}p{0.19\textwidth}p{0.48\textwidth}}
\caption{NS--A Version V11 proof-status ledger.}\label{tab:v11-ledger}\\
\toprule Row & Status & Exact conclusion\\\midrule
\endfirsthead
\toprule Row & Status & Exact conclusion\\\midrule
\endhead
V1--V10 prefix & \PROVED{} preserved
 &The complete V10 preterminal source is retained byte-for-byte; no earlier
 theorem is rewritten.\\[0.3em]
Annular source action & \PROVED{} subordinated
 &$\mathcal B_n^0\le C\int D^2$; action overdrive is an annularly resolved
 enstrophy-square overdrive.\\[0.3em]
Critical-acceleration norm shortcut & \OBSTRUCTION{} explicit
 &A finite moving-frequency $L_t^1H_x^{-7/2}$ stock can coexist with divergent
 per-packet $L_t^2L_x^2$ action.\\[0.3em]
Material-tube square-action capture & \PROVED{} exact fork
 &Action escapes the canonical tube or a positive captured fraction has
 diverging annular participation.\\[0.3em]
Microdelocalized cohort fusion & \PROVED{} no count
 &Quadratic partition weights fuse any fixed action fraction exactly; no cell
 count appears, but relative diameter remains uncontrolled.\\[0.3em]
Product divergence quotient & \PROVED{} exact
 &Source-visible and gradient-pressure components are the complete quotient;
 divergence-null stress is its kernel.\\[0.3em]
Divergence-null completion row & \OBSTRUCTION{} removed from target
 &It is unnecessary for the weak momentum equation or existence of an Euler
 pressure.\\[0.3em]
Physical-pressure bound before passage & \OBSTRUCTION{} target inflation
 &Local product compactness alone gives the pressure-free constraint and an
 effective $L^2$ pressure.\\[0.3em]
Physical-pressure identification & \OPEN{} stronger target
 &Convergence of the original pressures is still required to identify the
 effective pressure with the physical limit.\\[0.3em]
Positive Reynolds defect & \PROVED{} refined
 &Only $P_\sigma\operatorname{div}\mathbb R$ obstructs Euler; gradient and
 divergence-null defect components are pressure/gauge followers.\\[0.3em]
Full defect vanishing & \OBSTRUCTION{} not necessary
 &Stationary Euler shears give $\mathbb R\ne0$, $\mathbb R\succeq0$, and
 $\operatorname{div}\mathbb R=0$.\\[0.3em]
Nonzero velocity occupation & \OPEN{} terminal
 &A nonzero covariance profile may coexist with $V=0$; terminal occupation or
 trace is required for a nonzero velocity profile.\\[0.3em]
Stationary Euler quotient & \CONDITIONAL{} exact
 &Bounded product plus zero source-visible Reynolds force gives a slow-time
 family satisfying stationary weak Euler almost everywhere, with an effective
 pressure; nontriviality still requires $V\ne0$.\\[0.3em]
Ancient dynamics / rigidity / regularity & \OPEN
 &Slow-time stationarity is not fast-time occupation, backward uniqueness, or
 an ancient Liouville theorem.\\
\bottomrule
\end{longtable}
\endgroup

% =============================================================================
\section{Exact mandate for Version V12}
\label{sec:v11-next}
% =============================================================================

\begin{programme}[Relate the covariance profile to the product quotient or stop]
\label{prog:v11-next}
Preserve Versions~V1--V11 verbatim.  Do not restore the divergence-null
completion row, demand $\mathbb R=0$ before testing Euler, or introduce a new
path, age, cell, or source metric.

Proceed only in the following order.
\begin{enumerate}[label=\textup{(V12.\arabic*)},leftmargin=5.6em]
\item Stop immediately on enstrophy-square overdrive, material-tube action
      escape, or source-action microdelocalization unless a genuinely new PDE
      upper stock is supplied.
\item On the bounded-product cell, compare the nonzero heat-covariance source
      $H=T_yF$ with the product quotient $(V,Q,\mathbb R)$ on the same history.
      Perform the heat propagation before taking a positive part.  A failure
      of local heat-tail tightness is retained as the exact tail survivor.
\item Determine whether the covariance profile is generated by the limiting
      velocity, by the source-visible Reynolds force, or by an unresolved
      terminal occupation defect.  Do not identify $F$ with
      $\mathbb R$ by analogy.
\item On the zero source-visible-force branch, use the stationary Euler
      quotient only after proving $V\ne0$.  If $V=0$, stop at the covariance-
      only profile and seek terminal occupation/trace rather than another
      tensor compactness theorem.
\item If a nonzero stationary Euler profile is obtained, test the original
      source-paid scalar, terminal trace, and the appropriate backward or
      Liouville statement on that profile.  Physical pressure is required only
      if the target explicitly includes its identification.
\item Preserve every represented return and paid-stress component through its
      existing endpoint map before declaring a new source-visible force.
\end{enumerate}
The admissible endpoint is
\[
 \boxed{
 \begin{gathered}
 \text{enstrophy/action overdrive}
 \ \big|\ 
 \text{material action escape or microdelocalization}\\[0.15em]
 \text{heat-tail or terminal-occupation defect}
 \ \big|\ 
 \text{source-visible Reynolds force}\\[0.15em]
 \text{nonzero stationary Euler quotient}
 \ \big|\ 
 \text{named terminal trace or Liouville stop}.
 \end{gathered}}
\]
\end{programme}

\begin{assessment}[Programme position after V11]
\label{ass:v11-position}
The weak momentum target no longer waits for a complete tensor or physical-
pressure compactness theorem.  Product boundedness gives the effective
pressure automatically, and a nonzero Reynolds defect is harmless precisely
when its solenoidal divergence vanishes.  The remaining upstream bridge is
now exact: relate the already nonzero heat-covariance profile to the limiting
velocity and source-visible product defect on the same history.  If that
bridge fails, its heat-tail, occupation, or source-force residual is the
terminal object; no additional compiler is needed.
\end{assessment}

\begin{assessment}[Append-only integrity]
\label{ass:v11-integrity}
The complete Version~V10 source preceding its sole final document terminator
is preserved verbatim.  The present material begins at
Section~\ref{sec:v11-scope}; the final command below is again unique.
\end{assessment}

\addtocontents{toc}{\protect\endgroup}
% =============================================================================
% END VERSION V11 MATERIAL -- append later versions before final terminator
% =============================================================================


\clearpage
% =============================================================================
% VERSION V12: NEW MATERIAL.  The complete V1--V11 source above is preserved
% verbatim; only its sole final document terminator was removed before append.
% =============================================================================
\hypersetup{pdftitle={NS-A Terminal Source Concentration Programme: Cumulative V12},pdfauthor={A. Pelissier}}
\makeatletter
\addtocontents{toc}{\protect\begingroup\protect\makeatletter}
\addtocontents{toc}{\protect\def\protect\l@section{\protect\@dottedtocline{1}{0em}{4.7em}}}
\addtocontents{toc}{\protect\def\protect\l@subsection{\protect\@dottedtocline{2}{1.5em}{5.3em}}}
\addtocontents{toc}{\protect\makeatother}
\makeatother
\renewcommand{\thesection}{V12.\arabic{section}}
\renewcommand{\thesubsection}{V12.\arabic{section}.\arabic{subsection}}
\renewcommand{\theequation}{V12.\arabic{equation}}
\setcounter{section}{0}
\setcounter{equation}{0}
\renewcommand{\theHsection}{V12.\arabic{section}}
\renewcommand{\theHsubsection}{V12.\arabic{section}.\arabic{subsection}}
\renewcommand{\theHtheorem}{V12.\arabic{section}.\arabic{theorem}}
\renewcommand{\theHequation}{V12.\arabic{equation}}

\begin{center}
{\LARGE\bfseries NS--A: Version V12}\\[0.5em]
{\large The Routed-Profile Correlation Short, Heat--Occupation Defect,\\
and the Nonzero Heat-Unstable Stationary Euler Survivor}\\[0.5em]
{\normalsize 3 September 2026}
\end{center}
\addcontentsline{toc}{part}{Version V12: routed correlation, heat occupation, and the heat-Euler survivor}

\begin{abstract}
Version~V11 showed that local product compactness produces an effective
pressure without compactness of the displayed Navier--Stokes pressure, and
that only the solenoidal divergence of the Reynolds product defect obstructs
a stationary Euler quotient.  It left one upstream bridge: the nonzero
routed heat-covariance profile $H=T_yF$ was not yet related to the weak
velocity/product limit $(V,Q,\mathbb R)$ on the same history.

The present continuation first corrects the formulation of that bridge.  The
profile extracted in Version~V10 is the weak limit of $\Omega_nG_n$, whereas
the product quotient comes from the unweighted fields $V_n$.  Spatial
flatness of the material router does not imply temporal decorrelation.  After
extracting the spatially constant weak limit $\theta(s)$ of $\Omega_n$, the
exact decomposition is
\[
 H=\theta\,\overline H+\mathfrak C_{\rm rt},
 \qquad
 F=\theta\,\overline F+\mathfrak C_{\rm rt}^{F},
 \qquad
 \mathfrak C_{\rm rt}=T_y\mathfrak C_{\rm rt}^{F},
\]
where $(\overline F,\overline H)$ is the unweighted covariance profile and
$\mathfrak C_{\rm rt}$ is a routed time-correlation defect.  A scalar
oscillation example proves that the defect is not removed by spatial
flatness alone.

Second, the note performs the heat propagation before taking any positive
part.  On every fixed interior age and spatial cutoff, the heat-smoothed
velocities have a weak quadratic limit.  Its excess over the square of the
weak heat-smoothed velocity is a positive heat--occupation defect
$\mathbb O_R\succeq0$.  If $Q=V\otimes V+\mathbb R$, the exact truncated
covariance identity is
\[
 \mathbb D_R
 =\mathbb C_R[V]+P_{-s}(\chi_R\mathbb R)-\mathbb O_R,
\]
where $\mathbb C_R[V]\succeq0$ is the intrinsic truncated heat covariance.
The weak limit of the actual source differs from the source of $\mathbb D_R$
by one explicitly named heat-tail distribution.

Combining these two shorts gives an exhaustive target-native alternative.
A nonzero selected scalar is carried by a routed time-correlation defect, a
heat tail, a heat-smoothed source-visible Reynolds defect, a source-visible
heat--occupation defect, or the intrinsic heat covariance of the weak
velocity.  On the final branch the weak velocity is necessarily nonzero.
Thus the former ``covariance-only with $V=0$'' row is resolved into explicit
occupation/tail/force defects; it is not an additional mysterious profile.

On the zero source-visible Reynolds-force branch, global heat tightness
together with vanishing of the full routed, heat-tail, and occupation-source
distributions turns the surviving profile into the exact heat--Euler commutator
\[
 \overline H(s)
 =-\psi(\Lambda_y)c_y[P_{-s}V(s)],
 \qquad
 c_y[w]=(-\Delta_y)^{-1/4}P_\sigma\operatorname{div}(w\otimes w).
\]
The original source-paid scalar therefore survives on a nonzero stationary
Euler field whose heat orbit is not Euler-stationary.  Laplacian eigenfields,
Beltrami eigenfields, and arbitrary unidirectional shears are excluded from
this branch.  A separate countertheorem shows that a slow-time family of
stationary Euler fields may have an arbitrary measurable amplitude, so the
stationary quotient still supplies no ancient dynamics.  The first remaining
row is a fast-time occupation/terminal-trace or Liouville theorem for this
specific heat-unstable survivor.  No global regularity conclusion is claimed.
\end{abstract}

% =============================================================================
\section{Upstream scope: the routed profile is not the unweighted profile}
\label{sec:v12-scope}
% =============================================================================

Retain the bounded-action, one-cell, bounded-product branch of
Versions~V10--V11, and stop on every earlier enstrophy-square overdrive,
material-route action escape, or annular microdelocalization output.  Extend
all scaled fields by zero from $J_n$ to the limiting interval
$J=(-\tau_\infty,0)$ as in Theorem~\ref{thm:v10-covariance-profile}.  Thus
\begin{equation}
 \Omega_nG_n\rightharpoonup H,
 \qquad
 \Omega_n\mathbb F_n^{\rm sc}\rightharpoonup F,
 \qquad
 H=T_yF,
 \label{eq:v12-weighted-profile}
\end{equation}
with the nonzero pairing
\begin{equation}
 \left|\int_J\!\int\chi(y)H(s,y)\cdot Z(y)\,\dd y\,\dd s\right|
 \ge h_\chi>0.
 \label{eq:v12-inherited-pairing}
\end{equation}
The scaled router obeys
\begin{equation}
 0\le\Omega_n\le1,
 \qquad
 \|\nabla_y\Omega_n\|_{L^\infty(J_n\times\Torus_n^3)}
 \longrightarrow0.
 \label{eq:v12-router-flat}
\end{equation}

The local product branch is interpreted on every fixed spatial head.  The
following elementary diagonalization records the exact alternative.

\begin{lemma}[Fixed-radius product overdrive or one global local profile]
\label{lem:v12-global-local-product}
After passage to a subsequence, exactly one of the following occurs.
\begin{enumerate}[label=\textup{(P\arabic*)},leftmargin=3.4em]
\item There is an integer $R\ge1$ such that
      \begin{equation}
       \int_{J_n}\int_{B(0,R)}|V_n|^4\longrightarrow\infty.
       \label{eq:v12-fixed-radius-overdrive}
      \end{equation}
      This is the product-overdrive branch already retained in Version~V10.
\item For every $R<\infty$,
      \begin{equation}
       \sup_n\int_{J_n}\int_{B(0,R)}|V_n|^4<\infty.
       \label{eq:v12-all-radius-product-bound}
      \end{equation}
      On one diagonal subsequence there are fields
      \begin{equation}
       V\in L^4_{\rm loc}(J\times\R^3),
       \qquad
       Q\in L^2_{\rm loc}(J\times\R^3;\operatorname{Sym}(3))
       \label{eq:v12-global-local-VQ}
      \end{equation}
      such that $V_n\rightharpoonup V$ in $L^4_{\rm loc}$ and
      $V_n\otimes V_n\rightharpoonup Q$ in $L^2_{\rm loc}$.
\end{enumerate}
\end{lemma}

\begin{proof}
For each integer $R$, pass to a subsequence on which the nonnegative numbers
in \eqref{eq:v12-fixed-radius-overdrive} either are bounded or tend to
infinity.  If the latter occurs for one $R$, stop at the first such radius.
Otherwise use a diagonal subsequence.  Reflexivity and the bounds at every
integer radius give the weak limits; compatibility on nested balls follows
from uniqueness of weak restrictions.
\end{proof}

We work below on branch \textup{(P2)}.  Put
\begin{equation}
 \mathbb R:=Q-V\otimes V\succeq0,
 \label{eq:v12-Reynolds-defect}
\end{equation}
where positivity is the V10 weak-product theorem.

\begin{lemma}[The selected scalar has an interior-age witness]
\label{lem:v12-interior-age}
There are $\eta\in C_c^\infty(J)$, $0\le\eta\le1$, and
$h_0>0$ such that, with
\begin{equation}
 \Phi(s,y):=\eta(s)\chi(y)Z(y),
 \qquad
 \mathscr J(S):=\int_J\!\int_{\R^3}S\cdot\Phi\,\dd y\,\dd s,
 \label{eq:v12-test-functional}
\end{equation}
we have
\begin{equation}
 \boxed{|\mathscr J(H)|\ge h_0.}
 \label{eq:v12-interior-pairing}
\end{equation}
In particular there are constants $0<\tau_-<\tau_+<\infty$ such that
\begin{equation}
 \supp\eta\subset\{s:-\tau_+\le s\le-\tau_-\}.
 \label{eq:v12-interior-age-window}
\end{equation}
\end{lemma}

\begin{proof}
The integrand in \eqref{eq:v12-inherited-pairing} belongs to $L^1(J)$ by
Cauchy--Schwarz, the local $L^2$ bound for $H$, and smooth compact support of
$\chi Z$.  Choose $\eta_m\in C_c^\infty(J)$, $0\le\eta_m\le1$, increasing
pointwise to one.  Dominated convergence sends the corresponding integrals
to the integral in \eqref{eq:v12-inherited-pairing}.  One of them has
absolute value at least $h_\chi/2$; take this value for $h_0$.
\end{proof}

\begin{firewall}[Why the endpoint age is removed before heat passage]
\label{fire:v12-interior-first}
The heat time is $-s$.  On the support of $\eta$ it lies in the fixed compact
interval $[\tau_-,\tau_+]$.  Thus every heat operator used below is uniformly
smoothing.  Any scalar that survives only through $s=0$ would be a terminal
trace concentration, but the inherited $L^2$ profile pairing already has an
interior witness, so that branch is not silently created here.
\end{firewall}

% =============================================================================
\section{The routed time-correlation short}
\label{sec:v12-router-correlation}
% =============================================================================

For $s\in J$, define the spatial average on the fixed unit ball
\begin{equation}
 \theta_n(s):=\frac{1}{|B(0,1)|}\int_{B(0,1)}\Omega_n(s,y)\,\dd y.
 \label{eq:v12-theta-n}
\end{equation}
After passage to a subsequence,
\begin{equation}
 \theta_n\stackrel{*}{\rightharpoonup}\theta
 \quad\text{in }L^\infty(J),
 \qquad 0\le\theta\le1.
 \label{eq:v12-theta-limit}
\end{equation}

\begin{lemma}[Spatially flat router and unweighted source profile]
\label{lem:v12-unweighted-profile}
For every fixed $R<\infty$,
\begin{equation}
 \sup_{s\in J_n}\sup_{|y|\le R}
 |\Omega_n(s,y)-\theta_n(s)|\longrightarrow0.
 \label{eq:v12-router-spatial-constant}
\end{equation}
On the same diagonal subsequence there are unweighted weak limits
\begin{align}
 \mathbb F_n^{\rm sc}&\rightharpoonup\overline F
 &&\text{in }L^2(J;\dot H^{1/2}_{\rm loc}),
 \label{eq:v12-unweighted-F}\\
 G_n&\rightharpoonup\overline H
 &&\text{in }L^2_{\rm loc}(J\times\R^3),
 \label{eq:v12-unweighted-H}
\end{align}
with
\begin{equation}
 \boxed{\overline H=T_y\overline F.}
 \label{eq:v12-unweighted-source-relation}
\end{equation}
Define
\begin{equation}
 \mathfrak C_{\rm rt}^{F}:=F-\theta\overline F,
 \qquad
 \mathfrak C_{\rm rt}:=H-\theta\overline H.
 \label{eq:v12-router-correlation-defects}
\end{equation}
Then
\begin{equation}
 \boxed{
 H=\theta\overline H+\mathfrak C_{\rm rt},
 \qquad
 \mathfrak C_{\rm rt}=T_y\mathfrak C_{\rm rt}^{F}.}
 \label{eq:v12-router-correlation-short}
\end{equation}
\end{lemma}

\begin{proof}
The mean-value estimate and \eqref{eq:v12-router-flat} give
\[
 |\Omega_n(s,y)-\theta_n(s)|
 \le C_R\|\nabla_y\Omega_n(s)\|_\infty
\]
for $|y|\le R$, proving
\eqref{eq:v12-router-spatial-constant}.  The unweighted source action is
bounded by \eqref{eq:v10-annular-action-bounded}; the polar identity
\eqref{eq:v10-annular-source-polar} gives the corresponding local
$L_s^2\dot H_y^{1/2}$ bound for $\mathbb F_n^{\rm sc}$.  Banach--Alaoglu and
a diagonal argument give \eqref{eq:v12-unweighted-F}--
\eqref{eq:v12-unweighted-H}.  Passing to distributions in
$G_n=T_y\mathbb F_n^{\rm sc}$ proves
\eqref{eq:v12-unweighted-source-relation}.

Equation \eqref{eq:v12-router-correlation-short} is algebraic.  Since
$\theta$ depends only on slow time, it commutes with the spatial multiplier
$T_y$; combine $H=T_yF$ with
$\overline H=T_y\overline F$.
\end{proof}

\begin{definition}[Selected routed time-correlation residual]
\label{def:v12-selected-route-residual}
The scalar seen by the inherited writer is
\begin{equation}
 \boxed{
 \mathfrak r_{\rm rt}:=\mathscr J(\mathfrak C_{\rm rt}).}
 \label{eq:v12-selected-route-residual}
\end{equation}
A nonzero value is a same-history time-correlation between the material
router and the covariance source.  It is neither a new source birth nor a
spatial collar.
\end{definition}

\begin{countertheorem}[Spatial flatness does not remove routed time correlation]
\label{cth:v12-flat-router-time-correlation}
Let $J=(0,2\pi)$ and let $F_0$ be any smooth compactly supported
source-visible tensor with $T_yF_0\ne0$.  Set
\begin{equation}
 \Omega_m(s,y)=\frac{1+\sin(ms)}2,
 \qquad
 F_m(s,y)=\sin(ms)F_0(y).
 \label{eq:v12-time-correlation-example}
\end{equation}
Then $0\le\Omega_m\le1$, $\nabla_y\Omega_m=0$,
\begin{equation}
 \Omega_m\stackrel{*}{\rightharpoonup}\frac12,
 \qquad
 F_m\rightharpoonup0,
 \label{eq:v12-example-marginals}
\end{equation}
but
\begin{equation}
 \boxed{
 \Omega_mF_m\rightharpoonup\frac14F_0,
 \qquad
 \Omega_mT_yF_m\rightharpoonup\frac14T_yF_0.}
 \label{eq:v12-example-correlation}
\end{equation}
Thus the routed profile is not determined by the separate router and
unweighted-source weak limits.
\end{countertheorem}

\begin{proof}
The functions $\sin(ms)$ converge weakly to zero, while
$\sin^2(ms)$ converges weakly to $1/2$.  Expand
\(
 \Omega_mF_m=\frac12\sin(ms)F_0+\frac12\sin^2(ms)F_0
\)
and use that $T_y$ acts only in space.
\end{proof}

\begin{firewall}[The counterexample is functional, not a second PDE branch]
\label{fire:v12-correlation-functional}
Countertheorem~\ref{cth:v12-flat-router-time-correlation} is not asserted to
solve the material-router equation or Navier--Stokes.  It proves the precise
logical point needed here: the spatial-flatness estimate
\eqref{eq:v12-router-flat} alone cannot justify replacing the weak limit of
$\Omega_nG_n$ by $\theta$ times the weak limit of $G_n$.  Vanishing of
$\mathfrak C_{\rm rt}$ is therefore an independent terminal occupation row.
\end{firewall}

% =============================================================================
\section{Heat propagation before positivity and the occupation defect}
\label{sec:v12-heat-occupation}
% =============================================================================

Let $P_\tau^{(n)}=e^{\tau\Delta_y}$ be the heat semigroup on the expanding
rescaled torus and let $P_\tau$ be the Euclidean heat semigroup.  For each
integer $R\ge1$, choose $\chi_R\in C_c^\infty(\R^3)$ with
\begin{equation}
 0\le\chi_R\le1,
 \qquad
 \chi_R=1\text{ on }B(0,R),
 \qquad
 \supp\chi_R\subset B(0,2R).
 \label{eq:v12-heat-cutoff}
\end{equation}
All passages $R\to\infty$ below are taken through the positive integers, and
one diagonal subsequence is fixed simultaneously for those cutoffs.  For
$s\in\supp\eta$, put
\begin{align}
 U_{n,R}(s)&:=P_{-s}^{(n)}(\chi_RV_n(s)),
 \label{eq:v12-heat-velocity-n}\\
 A_{n,R}(s)&:=P_{-s}^{(n)}(\chi_RV_n\otimes V_n)(s),
 \label{eq:v12-heat-product-n}\\
 S_{n,R}(s)&:=U_{n,R}(s)\otimes U_{n,R}(s),
 \label{eq:v12-heat-square-n}\\
 D_{n,R}(s)&:=A_{n,R}(s)-S_{n,R}(s).
 \label{eq:v12-truncated-covariance-n}
\end{align}

\begin{lemma}[Local convergence of the expanding-torus heat kernel]
\label{lem:v12-periodic-heat-limit}
Fix $0<\tau_-<\tau_+<\infty$ and compact sets $K,K'\Subset\R^3$.
For every multi-index $\alpha$, the kernel of $P_\tau^{(n)}$ and its
$\alpha$-derivative converge uniformly on
$[\tau_-,\tau_+]\times K\times K'$ to the Euclidean heat kernel and its
corresponding derivative.
\end{lemma}

\begin{proof}
The spatial period of $\Torus_n^3$ tends to infinity with $\Gamma_n$.
The periodic heat kernel is the periodization
\[
 K_\tau^{(n)}(x-y)
 =\sum_{m\in L_n\mathbb Z^3}
   (4\pi\tau)^{-3/2}
   \exp\!\left(-\frac{|x-y+m|^2}{4\tau}\right),
\]
where $L_n\to\infty$.  On fixed compact sets, every term with $m\ne0$ and
all of its derivatives are bounded by
$C_{\alpha,\tau_-,\tau_+}e^{-cL_n^2}$.  Their sum tends uniformly to zero.
\end{proof}

\begin{theorem}[Exact local heat-product and occupation decomposition]
\label{thm:v12-heat-product-decomposition}
Fix $R$.  After passage to a diagonal subsequence, on
$\supp\eta\times\R^3$ we have
\begin{align}
 U_{n,R}&\rightharpoonup
 U_R:=P_{-s}(\chi_RV)
 &&\text{in }L^4_{\rm loc},
 \label{eq:v12-UR-limit}\\
 A_{n,R}&\rightharpoonup
 A_R:=P_{-s}(\chi_RQ)
 &&\text{in }L^2_{\rm loc},
 \label{eq:v12-AR-limit}\\
 S_{n,R}&\rightharpoonup S_R
 &&\text{in }L^2_{\rm loc}.
 \label{eq:v12-SR-limit}
\end{align}
Define the heat--occupation defect
\begin{equation}
 \boxed{\mathbb O_R:=S_R-U_R\otimes U_R.}
 \label{eq:v12-occupation-defect}
\end{equation}
Then
\begin{equation}
 \boxed{\mathbb O_R\succeq0}
 \label{eq:v12-occupation-positive}
\end{equation}
as a matrix-valued distribution.  Put
\begin{equation}
 \mathbb C_R[V]
 :=P_{-s}(\chi_RV\otimes V)-U_R\otimes U_R.
 \label{eq:v12-truncated-intrinsic-covariance}
\end{equation}
Then
\begin{equation}
 \boxed{\mathbb C_R[V]\succeq0,}
 \label{eq:v12-intrinsic-positive}
\end{equation}
and the weak limit
$\mathbb D_R:=A_R-S_R$ satisfies the exact identity
\begin{equation}
 \boxed{
 \mathbb D_R
 =\mathbb C_R[V]
  +P_{-s}(\chi_R\mathbb R)
  -\mathbb O_R.}
 \label{eq:v12-three-term-covariance}
\end{equation}
Moreover $\mathbb D_R\succeq0$, so
\begin{equation}
 \boxed{
 \mathbb O_R
 \preceq
 \mathbb C_R[V]+P_{-s}(\chi_R\mathbb R).}
 \label{eq:v12-occupation-budget}
\end{equation}
\end{theorem}

\begin{proof}
The all-radius product bound gives local $L^4$ bounds for $V_n$ and local
$L^2$ bounds for $V_n\otimes V_n$.  The heat operators in the fixed age
window are bounded on these spaces.  Weak compactness gives the three limits.
To identify the first two, test against a compactly supported smooth
function, move $P_{-s}^{(n)}$ to the test, use
Lemma~\ref{lem:v12-periodic-heat-limit}, and pass with the weak limits of
$V_n$ and $V_n\otimes V_n$.

For $\xi\in\R^3$ and nonnegative
$\varphi\in C_c^\infty(\supp\eta\times\R^3)$, the scalar sequence
$\sqrt\varphi\,\xi\cdot U_{n,R}$ converges weakly in $L^2$ to
$\sqrt\varphi\,\xi\cdot U_R$.  Lower semicontinuity gives
\[
 \int\varphi|\xi\cdot U_R|^2
 \le\liminf_n\int\varphi|\xi\cdot U_{n,R}|^2
 =\int\varphi\,\xi^{\mathsf T}S_R\xi,
\]
which proves \eqref{eq:v12-occupation-positive}.

For \eqref{eq:v12-intrinsic-positive}, let
$m_R=P_{-s}\chi_R\in[0,1]$.  Pointwise Cauchy--Schwarz with the positive
measure $K_{-s}(y-z)\chi_R(z)\,\dd z$ gives
\[
 |P_{-s}(\chi_R\,\xi\cdot V)|^2
 \le m_RP_{-s}(\chi_R|\xi\cdot V|^2)
 \le P_{-s}(\chi_R|\xi\cdot V|^2).
\]
This is exactly positivity of $\mathbb C_R[V]$.

Since $Q=V\otimes V+\mathbb R$, subtracting
$S_R=U_R\otimes U_R+\mathbb O_R$ from
$A_R=P_{-s}(\chi_RQ)$ gives
\eqref{eq:v12-three-term-covariance}.  Finally, each finite-$n$ tensor
$D_{n,R}$ is positive by the same sub-Markov Jensen argument.  Weak limits
preserve distributional positivity, so $\mathbb D_R\succeq0$; rearrange
\eqref{eq:v12-three-term-covariance} to obtain
\eqref{eq:v12-occupation-budget}.
\end{proof}

\begin{assessment}[What the new positive defect measures]
\label{ass:v12-occupation-meaning}
The defect $\mathbb O_R$ is not the Reynolds defect $\mathbb R$.  The latter
measures weak loss before heat propagation:
$V_n\otimes V_n\rightharpoonup V\otimes V+\mathbb R$.  The former measures
weak loss after the same-time velocity has been heat smoothed:
$U_{n,R}\otimes U_{n,R}\rightharpoonup
 U_R\otimes U_R+\mathbb O_R$.  It is therefore a terminal occupation or
fast-time oscillation defect on the precise heat scale read by the caloric
writer.
\end{assessment}

% =============================================================================
\section{The selected heat-tail identity}
\label{sec:v12-heat-tail}
% =============================================================================

Define the fixed annular source multiplier
\begin{equation}
 \mathscr A S
 :=T_y\mathsf P_{\rm sv}\psi(\Lambda_y)S^\circ.
 \label{eq:v12-source-multiplier}
\end{equation}
Since $T_y=T_y\mathsf P_{\rm sv}$, the explicit projection is retained only
to display the physical quotient.  Put
\begin{equation}
 \mathfrak T_R
 :=\overline H-\mathscr A\mathbb D_R.
 \label{eq:v12-heat-tail-distribution}
\end{equation}
This is the source produced outside the radius-$R$ heat truncation, including
all cross terms created by the square of the full heat-smoothed velocity.

For the target functional of \eqref{eq:v12-test-functional}, define
\begin{align}
 t_R&:=\mathscr J(\theta\mathfrak T_R),
 \label{eq:v12-tail-scalar}\\
 v_R&:=\mathscr J(\theta\mathscr A\mathbb C_R[V]),
 \label{eq:v12-velocity-scalar}\\
 r_R&:=\mathscr J\bigl(\theta\mathscr A
       P_{-s}(\chi_R\mathbb R)\bigr),
 \label{eq:v12-Reynolds-scalar}\\
 o_R&:=\mathscr J(\theta\mathscr A\mathbb O_R).
 \label{eq:v12-occupation-scalar}
\end{align}

\begin{theorem}[Exact five-term selected-scalar decomposition]
\label{thm:v12-five-term-scalar}
For every fixed $R$,
\begin{equation}
 \boxed{
 \mathscr J(H)
 =\mathfrak r_{\rm rt}+t_R+v_R+r_R-o_R.}
 \label{eq:v12-five-term-identity}
\end{equation}
Consequently, after passage to a subsequence, at least one of the following
priority outputs occurs.
\begin{enumerate}[label=\textup{(B\arabic*)},leftmargin=3.4em]
\item $\mathfrak r_{\rm rt}\ne0$: routed time-correlation survivor;
\item $\limsup_{R\to\infty}|t_R|>0$: selected local heat-tail survivor;
\item $\limsup_{R\to\infty}|r_R|>0$: heat-smoothed source-visible Reynolds
      survivor or Reynolds-tail survivor;
\item $\limsup_{R\to\infty}|o_R|>0$: source-visible heat--occupation survivor;
\item all four defects vanish, and
      \begin{equation}
       \boxed{
       \lim_{R\to\infty}v_R=\mathscr J(H)\ne0.}
       \label{eq:v12-intrinsic-scalar-survives}
      \end{equation}
      In this final branch, $V\not\equiv0$ on
      $\supp\eta\times\R^3$.
\end{enumerate}
Thus a nonzero routed covariance profile with zero weak velocity is not a
fifth unexplained object: it necessarily returns one of the first four
physical defects.
\end{theorem}

\begin{proof}
By Lemma~\ref{lem:v12-unweighted-profile},
$H=\theta\overline H+\mathfrak C_{\rm rt}$.  By definition of the heat tail,
$\overline H=\mathfrak T_R+\mathscr A\mathbb D_R$.  Insert the exact tensor
identity \eqref{eq:v12-three-term-covariance} and apply the linear functional
$\mathscr J$.  This gives \eqref{eq:v12-five-term-identity}.

If one of the first four quantities does not vanish in the displayed sense,
retain the first such output.  Otherwise take $R\to\infty$ in
\eqref{eq:v12-five-term-identity};
\eqref{eq:v12-interior-pairing} gives
\eqref{eq:v12-intrinsic-scalar-survives}.  If $V=0$ almost everywhere, then
$\mathbb C_R[V]=0$ for every $R$, hence $v_R=0$, a contradiction.
\end{proof}

\begin{proposition}[Concrete sufficient conditions for the three zero-defect rows]
\label{prop:v12-defect-sufficient}
The following implications hold.
\begin{enumerate}[label=\textup{(\roman*)},leftmargin=2.7em]
\item If $\Omega_n-\theta\to0$ strongly in
      $L^2_{\rm loc}(J\times\R^3)$, then
      $\mathfrak C_{\rm rt}=0$.
\item If
      $\mathbb R\in L^2(\supp\eta\times\R^3)$,
      $P_\sigma\operatorname{div}\mathbb R=0$, and
      $\chi_R\mathbb R\to\mathbb R$ in that global $L^2$ space, then
      \begin{equation}
       r_R\longrightarrow0.
       \label{eq:v12-Reynolds-tail-zero}
      \end{equation}
\item For a fixed $R$, the following local equivalence holds.  If
      $\varphi\in C_c^\infty(\supp\eta\times\R^3)$ is nonnegative and
      $\mathbb O_R$ vanishes on a neighbourhood of $\supp\varphi$, then
      $\sqrt\varphi\,U_{n,R}\to\sqrt\varphi\,U_R$ strongly in $L^2$.
      Conversely, strong $L^2_{\rm loc}$ convergence of $U_{n,R}$ implies
      $\mathbb O_R=0$.  In particular, strong heat-occupation compactness
      removes $o_R$.
\end{enumerate}
\end{proposition}

\begin{proof}
For item~\textup{(i)}, pair
$(\Omega_n-\theta)G_n$ against a compactly supported smooth function and use
Cauchy--Schwarz with the uniform local $L^2$ bound for $G_n$; do the same for
$\mathbb F_n^{\rm sc}$.  Item~\textup{(ii)} follows because
$\mathscr A P_{-s}$ is a bounded $L^2$ multiplier uniformly on the fixed age
window, $\chi_R\mathbb R\to\mathbb R$ in $L^2$, and
\[
 \mathscr A P_{-s}\mathbb R
 =\psi(\Lambda_y)P_{-s}
   (-\Delta_y)^{-1/4}P_\sigma\operatorname{div}\mathbb R=0.
\]
For item~\textup{(iii)}, take the trace in the proof of
Theorem~\ref{thm:v12-heat-product-decomposition}.  For every nonnegative
$\varphi\in C_c^\infty(\supp\eta\times\R^3)$,
\[
 \int\varphi\operatorname{tr}\mathbb O_R
 =\lim_n\int\varphi|U_{n,R}|^2
  -\int\varphi|U_R|^2.
\]
If $\mathbb O_R$ vanishes near $\supp\varphi$, the two localized norms
converge; together with weak convergence this is equivalent to strong
convergence of $\sqrt\varphi\,U_{n,R}$ in $L^2$.  Conversely, strong local
$L^2$ convergence passes the quadratic products in $L^1_{\rm loc}$, hence
forces $S_R=U_R\otimes U_R$ and $\mathbb O_R=0$.
\end{proof}

\begin{firewall}[The selected tail is target native]
\label{fire:v12-tail-target-native}
The tail $t_R$ is not the total mass outside a ball and is not estimated by
multiplying a cell count by a local bound.  It is exactly the part of the
same output-annular, source-visible caloric scalar which is not reconstructed
from the radius-$R$ weak product and heat occupation.  If it survives, the
programme stops at that physical heat-tail defect.
\end{firewall}

% =============================================================================
\section{Nonzero velocity and the heat--Euler commutator}
\label{sec:v12-heat-Euler}
% =============================================================================

Work on branch \textup{(B5)} and on the zero source-visible Reynolds-force
branch of Theorem~\ref{thm:v11-force-or-Euler}.  Then $V\ne0$ by
Theorem~\ref{thm:v12-five-term-scalar}, and there is an effective pressure
$\pi_E$ such that
\begin{equation}
 \boxed{
 \operatorname{div}(V\otimes V)+\nabla\pi_E=0,
 \qquad \operatorname{div}V=0}
 \label{eq:v12-stationary-Euler}
\end{equation}
for almost every slow time.

For a divergence-free field $W$, write
\begin{equation}
 c_y[W]
 :=(-\Delta_y)^{-1/4}P_\sigma\operatorname{div}(W\otimes W).
 \label{eq:v12-critical-source}
\end{equation}

\begin{theorem}[Heat-covariance source equals the heat--Euler commutator]
\label{thm:v12-heat-Euler-commutator}
Assume in addition the global heat-tight branch
\begin{equation}
 V\in L^4(\supp\eta\times\R^3),
 \qquad
 \mathbb R\in L^2(\supp\eta\times\R^3),
 \label{eq:v12-global-heat-tightness}
\end{equation}
and the full profile-level zero-defect conditions
\begin{equation}
 \boxed{
 \begin{gathered}
  \mathfrak C_{\rm rt}=0,\\
  \mathfrak T_R\longrightarrow0,\qquad
  \mathscr A\mathbb O_R\longrightarrow0
  \quad\text{in }\mathcal D'(\supp\eta\times\R^3)
  \quad(R\to\infty).
 \end{gathered}}
 \label{eq:v12-profile-zero-defects}
\end{equation}
Retain also the zero source-visible Reynolds-force branch
\begin{equation}
 P_\sigma\operatorname{div}\mathbb R=0.
 \label{eq:v12-zero-Reynolds-force-global}
\end{equation}
Then
\begin{equation}
 \mathbb C_{-s}[V]
 :=P_{-s}(V\otimes V)
   -(P_{-s}V)\otimes(P_{-s}V)
 \label{eq:v12-full-heat-covariance}
\end{equation}
is well defined on $\supp\eta$, and
\begin{equation}
 \boxed{
 \overline H
 =\psi(\Lambda_y)T_y\mathbb C_{-s}[V]^\circ
 =\psi(\Lambda_y)
   \bigl(P_{-s}c_y[V]-c_y[P_{-s}V]\bigr).}
 \label{eq:v12-heat-commutator-general}
\end{equation}
On the stationary Euler branch, $c_y[V]=0$, hence
\begin{equation}
 \boxed{
 \overline H(s)
 =-\psi(\Lambda_y)c_y[P_{-s}V(s)].}
 \label{eq:v12-stationary-heat-source}
\end{equation}
The original selected scalar becomes
\begin{equation}
 \boxed{
 \left|
 \int_J\!\int
 \theta(s)\eta(s)\chi(y)
 \psi(\Lambda_y)c_y[P_{-s}V(s)]\cdot Z(y)
 \,\dd y\,\dd s
 \right|
 \ge h_0.}
 \label{eq:v12-source-paid-Euler-scalar}
\end{equation}
Thus the profile-level defect-free branch produces a nonzero stationary
Euler field whose heat-smoothed orbit is not Euler-stationary on the
selected annulus.
\end{theorem}

\begin{proof}
The global integrability in \eqref{eq:v12-global-heat-tightness} gives
$\chi_RV\to V$ in $L^4$ and
$\chi_R(V\otimes V)\to V\otimes V$ in $L^2$.  Uniform heat smoothing on the
interior age window therefore sends $\mathbb C_R[V]$ to
\eqref{eq:v12-full-heat-covariance} in the local spaces seen by $\mathscr A$.
Moreover, since $\chi_R\mathbb R\to\mathbb R$ in $L^2$ and
\eqref{eq:v12-zero-Reynolds-force-global} holds,
\[
 \mathscr A P_{-s}(\chi_R\mathbb R)
 \longrightarrow
 \psi(\Lambda_y)P_{-s}
 (-\Delta_y)^{-1/4}P_\sigma\operatorname{div}\mathbb R
 =0
\]
in distributions.  Insert
\eqref{eq:v12-three-term-covariance} into the definition
\eqref{eq:v12-heat-tail-distribution} and let $R\to\infty$ using
\eqref{eq:v12-profile-zero-defects}.  This proves
$\overline H=\mathscr A\mathbb C_{-s}[V]$.

All operators in $T_y$ commute with the heat semigroup.  Therefore
\[
 T_yP_{-s}(V\otimes V)=P_{-s}c_y[V],
 \qquad
 T_y\bigl((P_{-s}V)\otimes(P_{-s}V)\bigr)=c_y[P_{-s}V],
\]
which proves \eqref{eq:v12-heat-commutator-general}.  Equation
\eqref{eq:v12-stationary-Euler} implies
$P_\sigma\operatorname{div}(V\otimes V)=0$, hence $c_y[V]=0$.
This gives \eqref{eq:v12-stationary-heat-source}.  Finally,
$\mathfrak C_{\rm rt}=0$ gives $H=\theta\overline H$; insert
\eqref{eq:v12-stationary-heat-source} into
\eqref{eq:v12-interior-pairing}.
\end{proof}

\begin{proposition}[Carr\'e-du-champ representation of the intrinsic heat covariance]
\label{prop:v12-carre-du-champ}
If, for almost every $s$, the spatial field $V(s)$ belongs to $L^2(\R^3)$,
then for every $\tau>0$
\begin{equation}
 \boxed{
 \mathbb C_\tau[V]_{ij}
 =2\int_0^\tau
 P_{\tau-r}
 \left(\sum_{\ell=1}^3
   \partial_\ell P_rV_i\,
   \partial_\ell P_rV_j\right)\dd r.}
 \label{eq:v12-carre-du-champ}
\end{equation}
The integral is a positive matrix-valued $L^1$ distribution.  If $V$ is
stationary Euler, applying $T_y\psi(\Lambda_y)$ gives
\begin{equation}
 -\psi(\Lambda_y)c_y[P_\tau V]
 =2\int_0^\tau
  \psi(\Lambda_y)P_{\tau-r}T_y
  \left(\sum_{\ell}
  \partial_\ell P_rV\otimes\partial_\ell P_rV\right)\dd r.
 \label{eq:v12-viscous-commutator-source}
\end{equation}
\end{proposition}

\begin{proof}
For smooth $V$, differentiate
\[
 \mathcal F(r)
 :=P_{\tau-r}(P_rV\otimes P_rV).
\]
The product rule gives
\[
 \mathcal F'(r)
 =-2P_{\tau-r}
   \left(\sum_{\ell}
   \partial_\ell P_rV\otimes\partial_\ell P_rV\right).
\]
Since
$\mathcal F(0)=P_\tau(V\otimes V)$ and
$\mathcal F(\tau)=P_\tau V\otimes P_\tau V$, integration proves
\eqref{eq:v12-carre-du-champ}.  For $V\in L^2$, approximate in $L^2$ and use
\(
 \int_0^\tau\|\nabla P_rV\|_2^2\,\dd r
 \le\frac12\|V\|_2^2
\)
to pass in matrix-valued $L^1$.  Equation
\eqref{eq:v12-viscous-commutator-source} is
Theorem~\ref{thm:v12-heat-Euler-commutator} applied to
\eqref{eq:v12-carre-du-champ}.
\end{proof}

\begin{corollary}[Positive caloric deformation stock on the clean branch]
\label{cor:v12-caloric-gradient-stock}
Assume the hypotheses of
Theorem~\ref{thm:v12-heat-Euler-commutator} and, in addition, that
$V(s)\in L^2(\R^3)$ for almost every $s\in\supp\eta$.  Then, with the
left side interpreted in $[0,\infty]$,
\begin{equation}
 \boxed{
 \int_J\eta(s)\theta(s)
  \int_0^{-s}\|\nabla P_rV(s)\|_4^2\,\dd r\,\dd s
 \ge
 \frac{h_0}{C_\psi\|\chi Z\|_2}.}
 \label{eq:v12-caloric-gradient-stock}
\end{equation}
Here $C_\psi<\infty$ depends only on the fixed annular multiplier and the
Fourier normalization.  Thus the clean survivor carries a positive
heat-gradient history before the signed Leray/source read is taken.
\end{corollary}

\begin{proof}
Apply $\psi(\Lambda_y)T_y$ to
\eqref{eq:v12-carre-du-champ}.  Since the multiplier
$\psi(\Lambda_y)T_yP_{\tau-r}$ has fixed compact annular support, it maps
matrix $L^2$ to vector $L^2$ with norm at most $C_\psi$, uniformly for
$0\le r\le\tau\le\tau_+$.  Therefore
\[
 \|\psi(\Lambda_y)c_y[P_\tau V]\|_2
 \le C_\psi\int_0^\tau
 \left\|\sum_\ell
 \partial_\ell P_rV\otimes\partial_\ell P_rV\right\|_2\,\dd r
 \le C_\psi\int_0^\tau\|\nabla P_rV\|_4^2\,\dd r.
\]
Insert this estimate into
\eqref{eq:v12-source-paid-Euler-scalar}, use
$0\le\eta,\theta\le1$, and apply Cauchy--Schwarz in space against
$\chi Z$.  Rearrangement gives
\eqref{eq:v12-caloric-gradient-stock}.
\end{proof}

\begin{corollary}[Simple heat-stable Euler carriers cannot carry the scalar]
\label{cor:v12-heat-stable-excluded}
On the branch of Theorem~\ref{thm:v12-heat-Euler-commutator}, the field
cannot be heat rigid on the portion of slow time read by $\theta\eta$; that
is, it cannot satisfy
\begin{equation}
 P_\sigma\operatorname{div}
 \bigl(P_\tau V(s)\otimes P_\tau V(s)\bigr)=0
 \quad\text{for every relevant heat age }\tau.
 \label{eq:v12-heat-rigid}
\end{equation}
In particular, the following model classes are excluded whenever the stated
heat operations belong to the tightness class of the theorem:
\begin{enumerate}[label=\textup{(\roman*)},leftmargin=2.7em]
\item $P_\tau V(s)=a(s,\tau)V(s)$; in particular a single Laplacian
      eigenspace;
\item a Beltrami Laplacian eigenfield;
\item a unidirectional shear
      $V(s,y)=b(s)f_s(n_s\cdot y)$ with $b(s)\cdot n_s=0$.
\end{enumerate}
Thus the nonzero survivor is genuinely heat unstable; it requires a
multi-directional or nontrivially cancelling stationary Euler structure.
\end{corollary}

\begin{proof}
Condition \eqref{eq:v12-heat-rigid} is precisely
$c_y[P_\tau V]=0$, so it contradicts the nonzero pairing
\eqref{eq:v12-source-paid-Euler-scalar} if it holds almost everywhere on the
portion carrying that pairing.  If $P_\tau V=aV$, then
$c_y[P_\tau V]=a^2c_y[V]=0$ by stationarity.  A Beltrami Laplacian eigenfield
is a special case.  For a unidirectional shear,
$(V\cdot\nabla)V=0$ because $b\cdot n=0$; heat preserves the same form, so
$c_y[P_\tau V]=0$ for every $\tau$.
\end{proof}

\begin{assessment}[What has and has not been promoted]
\label{ass:v12-heat-Euler-meaning}
The profile now has nonzero velocity and carries the original source-paid
scalar, but the surviving quantity is a heat--Euler commutator.  It measures
failure of heat smoothing to preserve the stationary Euler constraint.  It
is not a sign-definite Euler production, a unit-viscosity ancient solution,
or a contradiction to stationary Euler.  The positive object lies one step
before the source, in the carr\'e-du-champ covariance of
\eqref{eq:v12-carre-du-champ}; the Leray divergence and the signed writer may
still cancel.
\end{assessment}

% =============================================================================
\section{Slow-time stationarity carries no chronology}
\label{sec:v12-no-chronology}
% =============================================================================

\begin{countertheorem}[A stationary Euler quotient does not determine an ancient evolution]
\label{cth:v12-stationary-no-time}
Let $W$ be any nonzero smooth divergence-free stationary Euler field with
pressure $p_W$:
\begin{equation}
 \operatorname{div}(W\otimes W)+\nabla p_W=0.
 \label{eq:v12-stationary-W}
\end{equation}
For every bounded measurable scalar function $a:J\to\R$, put
\begin{equation}
 V(s,y)=a(s)W(y),
 \qquad
 \pi(s,y)=a(s)^2p_W(y).
 \label{eq:v12-arbitrary-amplitude}
\end{equation}
Then
\begin{equation}
 \boxed{
 \operatorname{div}(V\otimes V)+\nabla\pi=0,
 \qquad
 \operatorname{div}V=0}
 \label{eq:v12-arbitrary-stationary-family}
\end{equation}
for almost every $s$, regardless of the time regularity of $a$.  In
particular, the stationary quotient permits jumps, rapid oscillations, and
arbitrary slow-time amplitude histories.
\end{countertheorem}

\begin{proof}
Since $a$ is spatially constant,
\[
 \operatorname{div}(V\otimes V)+\nabla\pi
 =a(s)^2
  \bigl(\operatorname{div}(W\otimes W)+\nabla p_W\bigr)=0.
\]
Divergence-freeness is equally immediate.  No time derivative occurs in the
stationary equation.
\end{proof}

\begin{firewall}[No Liouville theorem may be applied to the slow quotient]
\label{fire:v12-no-slow-Liouville}
Countertheorem~\ref{cth:v12-stationary-no-time} does not construct the
specific heat-unstable survivor of
Theorem~\ref{thm:v12-heat-Euler-commutator}.  It proves that the equation
\eqref{eq:v12-stationary-Euler} alone contains no chronology.  Backward
uniqueness, ancient mild-solution Liouville theorems, or reflected-time
arguments may be applied only after a fast-time occupation/terminal-trace
compactness theorem has produced an actual dynamical limit.
\end{firewall}

\begin{firewall}[Represented returns remain shorted first]
\label{fire:v12-represented-first}
Every profile and defect above is read on the held-out annular card inherited
from the V5--V10 represented-head Pythagoras.  If a component of
$\mathfrak C_{\rm rt}$, $\mathfrak T_R$, the Reynolds term, or the occupation
term is generated by a declared old-response or paid-stress return after its
actual endpoint transport, that component remains a follower of the earlier
source and receives no new source price.  Only the final held-out component
is eligible for the terminal survivor classification.
\end{firewall}

% =============================================================================
\section{Version V12 integrated theorem and exact stopping boundary}
\label{sec:v12-integrated}
% =============================================================================

\begin{theorem}[Version V12 routed heat-product and occupation theorem]
\label{thm:v12-integrated}
Preserving every theorem and stated limitation of Versions~V1--V11,
Version~V12 proves the following.
\begin{enumerate}[label=\textup{(V12.\arabic*)},leftmargin=6em]
\item The nonzero V10 profile is routed: its weak limit is that of
      $\Omega_nG_n$, not automatically the weak limit of $G_n$ multiplied by
      a scalar router limit.
\item Spatial router flatness gives a scalar weak router $\theta(s)$ and the
      exact source short
      $H=\theta\overline H+\mathfrak C_{\rm rt}$ with
      $\mathfrak C_{\rm rt}=T_y\mathfrak C_{\rm rt}^{F}$.
\item A bounded oscillatory example shows that the routed time-correlation
      residual can be nonzero even when the router is spatially constant and
      the separate unweighted source converges weakly to zero.
\item At every fixed interior heat age and spatial cutoff, the weak square of
      the heat-smoothed velocity has the positive occupation defect
      $\mathbb O_R$.  The exact covariance identity is
      \eqref{eq:v12-three-term-covariance}, with the positive budget
      \eqref{eq:v12-occupation-budget}.
\item The actual selected scalar has the exact five-term decomposition
      \eqref{eq:v12-five-term-identity}.  Its terminal outputs are routed
      correlation, heat tail, heat-smoothed Reynolds source, heat occupation,
      or intrinsic velocity heat covariance.
\item If the first four defects vanish, the intrinsic velocity term carries
      the original nonzero scalar and the weak velocity satisfies
      $V\ne0$.  Thus a zero-velocity covariance profile necessarily returns a
      typed tail, force, or occupation defect.
\item Global $L^2$ tightness and zero solenoidal divergence of $\mathbb R$
      remove the Reynolds term.  Strong compactness of the heat-smoothed
      velocities is exactly the zero heat-occupation row.
\item On the zero Reynolds-force, globally heat-tight, and full
      profile-level zero-defect branch
      \eqref{eq:v12-profile-zero-defects}, the nonzero stationary Euler
      survivor obeys the heat--Euler commutator identity
      \eqref{eq:v12-stationary-heat-source}, and the inherited source-paid
      scalar becomes \eqref{eq:v12-source-paid-Euler-scalar}.
\item The intrinsic covariance has the positive carr\'e-du-champ formula
      \eqref{eq:v12-carre-du-champ}.  Under the stated $L^2$ entrance it also
      forces the positive caloric deformation stock
      \eqref{eq:v12-caloric-gradient-stock}.  Nevertheless its source/writer
      read is signed.  Heat-rigid eigenshell, Beltrami, and shear carriers
      cannot carry the surviving scalar.
\item The stationary Euler quotient has no time derivative and allows an
      arbitrary measurable slow-time amplitude.  It cannot be promoted to an
      ancient dynamical profile without a separate fast-time occupation or
      terminal-trace theorem.
\end{enumerate}
No item proves heat-tail tightness for every terminal datum, vanishing of the
routed or occupation defects, a fast-time dynamical limit, backward
uniqueness, an ancient Liouville theorem, or global three-dimensional
regularity.
\end{theorem}

\begin{proof}
Items~\textup{(V12.1)--(V12.3)} are
Lemma~\ref{lem:v12-unweighted-profile},
Definition~\ref{def:v12-selected-route-residual}, and
Countertheorem~\ref{cth:v12-flat-router-time-correlation}.  Item
\textup{(V12.4)} is
Theorem~\ref{thm:v12-heat-product-decomposition}.  Items
\textup{(V12.5)--(V12.6)} are
Theorem~\ref{thm:v12-five-term-scalar}.  Item~\textup{(V12.7)} is
Proposition~\ref{prop:v12-defect-sufficient}.  Items
\textup{(V12.8)--(V12.9)} are
Theorem~\ref{thm:v12-heat-Euler-commutator},
Proposition~\ref{prop:v12-carre-du-champ},
Corollary~\ref{cor:v12-caloric-gradient-stock}, and
Corollary~\ref{cor:v12-heat-stable-excluded}.  The final item is
Countertheorem~\ref{cth:v12-stationary-no-time}.
\end{proof}

\begingroup\small
\begin{longtable}{p{0.25\textwidth}p{0.19\textwidth}p{0.48\textwidth}}
\caption{NS--A Version V12 proof-status ledger.}\label{tab:v12-ledger}\\
\toprule Row & Status & Exact conclusion\\\midrule
\endfirsthead
\toprule Row & Status & Exact conclusion\\\midrule
\endhead
V1--V11 prefix & \PROVED{} preserved
 &The complete V11 preterminal source is retained byte-for-byte; no earlier
 theorem is rewritten.\\[0.3em]
Routed versus unweighted profile & \PROVED{} exact short
 &$H=\theta\overline H+\mathfrak C_{\rm rt}$ and the correlation residual is
 itself source visible.\\[0.3em]
Spatial router flatness & \OBSTRUCTION{} insufficient alone
 &Time oscillation can leave a nonzero routed source although both separate
 weak marginals lose it.\\[0.3em]
Interior heat age & \PROVED{} selected
 &The nonzero profile pairing survives on a compact age interval bounded away
 from zero, so endpoint singular heat time is not used.\\[0.3em]
Heat--occupation defect & \PROVED{} positive
 &$\mathbb O_R\succeq0$ is the weak square defect after heat smoothing and is
 distinct from the Reynolds defect before heat smoothing.\\[0.3em]
Heat-product bridge & \PROVED{} exact
 &$\mathbb D_R=\mathbb C_R[V]+P_{-s}(\chi_R\mathbb R)-\mathbb O_R$ before
 source projection or positivity splitting.\\[0.3em]
Local heat tail & \OPEN{} typed survivor
 &$\mathfrak T_R$ is exactly the selected source not reconstructed from the
 radius-$R$ weak product and heat occupation.\\[0.3em]
Zero-velocity covariance branch & \PROVED{} resolved conditionally
 &If routed correlation, heat tail, Reynolds source, and heat occupation all
 vanish, the nonzero scalar forces $V\ne0$.\\[0.3em]
Zero Reynolds force & \CONDITIONAL{} exact sufficient row
 &Global $L^2$ tightness plus $P_\sigma\operatorname{div}\mathbb R=0$ makes
 the heat-smoothed Reynolds scalar vanish.\\[0.3em]
Terminal occupation compactness & \OPEN{} exact criterion
 &$\mathbb O_R=0$ is equivalent to strong local $L^2$ convergence of the
 heat-smoothed velocities.\\[0.3em]
Heat--Euler commutator & \CONDITIONAL{} exact on full zero-defect branch
 &$\overline H=-\psi c_y[P_{-s}V]$ once the routed, distributional heat-tail,
 Reynolds-force, and occupation-source rows vanish.\\[0.3em]
Positive caloric deformation stock & \PROVED{} on the clean $L^2$ branch
 &The original scalar forces \eqref{eq:v12-caloric-gradient-stock} before the
 signed Leray/source read.\\[0.3em]
Simple stationary carriers & \OBSTRUCTION{} excluded
 &One-shell heat rays, Beltrami eigenfields, and unidirectional shears have
 zero heat--Euler source and cannot carry the selected scalar.\\[0.3em]
Slow-time stationary quotient & \OBSTRUCTION{} no chronology
 &Arbitrary measurable amplitudes preserve stationary Euler, so the quotient
 alone is not an ancient evolution.\\[0.3em]
Fast-time occupation / terminal trace & \OPEN{} first PDE row
 &Needed to turn the heat-unstable stationary survivor into an actual dynamic
 profile before backward uniqueness or Liouville.\\[0.3em]
Three-dimensional regularity & \OPEN
 &No terminal defect is eliminated for every fixed datum.\\
\bottomrule
\end{longtable}
\endgroup

% =============================================================================
\section{Exact mandate for Version V13}
\label{sec:v12-next}
% =============================================================================

\begin{programme}[Acquire fast-time occupation or stop at the exact defect]
\label{prog:v12-next}
Preserve Versions~V1--V12 verbatim.  Do not add another path law, age law,
cell metric, product quotient, or pressure completion.

Proceed only in the following order.
\begin{enumerate}[label=\textup{(V13.\arabic*)},leftmargin=5.7em]
\item Stop on nonzero routed time correlation, selected heat tail,
      source-visible Reynolds heat tail, or source-visible heat occupation.
      These are now the exact missing incidence/compactness rows.
\item On the intrinsic branch, retain the nonzero stationary Euler field and
      the precise heat--Euler scalar
      \eqref{eq:v12-source-paid-Euler-scalar}.  Do not replace it by a generic
      stationary Euler norm or by a full Reynolds tensor.
\item Test whether the original fast-time sequence places a fixed fraction of
      that scalar on one bounded fast-time window.  If every bounded window is
      negligible, return fast-time occupation escape rather than averaging it
      into the slow profile.
\item Only after a nonzero fast-time dynamical profile is obtained, evaluate
      terminal trace, reflected pairing, backward uniqueness, or the relevant
      Liouville theorem.  The slow stationary quotient is not an admissible
      substitute.
\item Use the carr\'e-du-champ identity and
      \eqref{eq:v12-caloric-gradient-stock} only before the signed
      Leray/source read.  A positive gradient covariance does not supply a
      sign for the source-paid scalar.
\item Preserve all represented-return and paid-stress components through their
      declared endpoint maps before classifying a heat-tail or occupation
      force as source fresh.
\end{enumerate}
The admissible endpoint is
\[
 \boxed{
 \begin{gathered}
 \text{routed correlation or selected heat tail}
 \ \big|\ 
 \text{Reynolds/occupation source defect}\\[0.15em]
 \text{fast-time occupation escape}
 \ \big|\ 
 \text{one nonzero heat-unstable dynamical profile}\\[0.15em]
 \text{terminal-trace, backward-uniqueness, or Liouville stop}.
 \end{gathered}}
\]
\end{programme}

\begin{assessment}[Programme position after V12]
\label{ass:v12-position}
The covariance/product bridge is no longer an analogy.  Its exact defects are
the routed time correlation, the spatial heat tail, the heat-smoothed
Reynolds source, and the post-heat occupation covariance.  If all four pass,
the weak velocity is nonzero and the original scalar lands on a specific
heat-unstable stationary Euler field.  The remaining gap is temporal rather
than tensorial: one must retain that scalar on a bounded fast-time window or
return occupation escape.  Further slow-time product machinery would be
circular.
\end{assessment}

\paragraph{Version V12 history.}
Moved upstream from the slow Euler--Reynolds quotient to the exact relation
between the routed heat-covariance profile and the unweighted product limit.
Separated router-time correlation, spatial heat tail, heat-smoothed Reynolds
force, and post-heat occupation before identifying the intrinsic velocity
covariance.  On the full profile-level zero-defect branch, derived the
heat--Euler commutator, retained a positive caloric deformation stock, and
proved that the surviving slow profile is nonzero but does not yet carry an
ancient chronology.

\begin{assessment}[Append-only integrity]
\label{ass:v12-integrity}
The complete Version~V11 source preceding its sole final document terminator
is preserved verbatim.  Its preterminal SHA--256 digest is
\begin{center}\scriptsize\ttfamily
0a1bb41c58636abcd6d98b71caa139cd9ddaa131552214feca859dfbc8601837
\end{center}
The present material begins at Section~\ref{sec:v12-scope}; the final command
below is again unique.
\end{assessment}

\addtocontents{toc}{\protect\endgroup}
% =============================================================================
% END VERSION V12 MATERIAL -- append later versions before final terminator
% =============================================================================

\clearpage
% =============================================================================
% VERSION V13: NEW MATERIAL.  The complete V1--V12 source above is preserved
% verbatim; only its sole final document terminator was removed before append.
% =============================================================================
\hypersetup{pdftitle={NS-A Terminal Source Concentration Programme: Cumulative V13},pdfauthor={A. Pelissier}}
\makeatletter
\addtocontents{toc}{\protect\begingroup\protect\makeatletter}
\addtocontents{toc}{\protect\def\protect\l@section{\protect\@dottedtocline{1}{0em}{4.7em}}}
\addtocontents{toc}{\protect\def\protect\l@subsection{\protect\@dottedtocline{2}{1.5em}{5.3em}}}
\addtocontents{toc}{\protect\makeatother}
\makeatother
\renewcommand{\thesection}{V13.\arabic{section}}
\renewcommand{\thesubsection}{V13.\arabic{section}.\arabic{subsection}}
\renewcommand{\theequation}{V13.\arabic{equation}}
\setcounter{section}{0}
\setcounter{equation}{0}
\renewcommand{\theHsection}{V13.\arabic{section}}
\renewcommand{\theHsubsection}{V13.\arabic{section}.\arabic{subsection}}
\renewcommand{\theHtheorem}{V13.\arabic{section}.\arabic{theorem}}
\renewcommand{\theHequation}{V13.\arabic{equation}}

\begin{center}
{\LARGE\bfseries NS--A: Version V13}\\[0.5em]
{\large Fast-Time Vanishing, Extensive Source Occupation,\\
and Closure of the One-Profile Route}\\[0.5em]
{\normalsize 3 September 2026}
\end{center}
\addcontentsline{toc}{part}{Version V13: fast-time vanishing and extensive source occupation}

\begin{abstract}
Version~V12 reached a nonzero, source-paid heat--Euler scalar on a stationary
Euler quotient, but correctly refused to call that quotient an ancient
trajectory.  Its next mandate asked whether the original fast-time sequence
places a fixed part of the scalar on one bounded fast-time window.  The
present continuation evaluates that question on the exact bounded-action
branch rather than introducing another temporal compactness device.

The answer is negative, and the reason is quantitative.  Let
$\alpha_n=(\mathcal R_n/\nu)^{1/2}\to\infty$ be the V10 coupling, let
$q_n(s)$ be the finite-record scalar density on the fixed slow-time window,
and let
\[
 d\Sigma_n(\tau)
 =\alpha_n^{-1}q_n(\tau/\alpha_n)\,d\tau
\]
be the same scalar in fast time.  The target-native source-action bound gives
one uniform $L^2$ bound for $q_n$, while the inherited nonzero pairing gives
one positive lower bound for its integral.  Consequently
\[
 \frac{c}{\alpha_n}
 \le \left\|\frac{d\Sigma_n}{d\tau}\right\|_2^2
 \le \frac{C}{\alpha_n},
 \qquad
 |\Sigma_n(\mathbb R)|\ge c_0>0.
\]
Thus the scalar density converges strongly to zero in every fast-time
$L^p$, $1<p\le2$, while its total signed mass remains nonzero.  More
strongly, every measurable fast-time set of length $o(\alpha_n)$ has
vanishing total variation.  This holds uniformly after arbitrary time
translation.  A fixed fraction of the scalar therefore requires order
$\alpha_n$ bounded windows; no finite or sublinear family of compact
fast-time profiles can carry it.

The only way for one bounded fast-time window to retain a fixed scalar is a
linear source-action overdrive: its normalized source action must be at least
$c\alpha_n$.  That is precisely the enstrophy-square/source-action branch
already separated in Versions~V10--V11.  On the bounded-action branch the
V12 alternative therefore collapses to fast-time occupation escape.  The
slow stationary Euler scalar survives only after the macroscopic dilation
$\tau\mapsto \tau/\alpha_n$; every translated compact fast-time limit loses
its contribution in the original normalization.  Terminal trace, reflected
pairing, backward uniqueness, and one-profile Liouville arguments are hence
downstream of a branch which does not occur here.  The remaining physical
question is an extensive occupation/payment or replenishment estimate on the
same source history.  No global regularity conclusion is claimed.
\end{abstract}

% =============================================================================
\section{Scope: the bounded-action branch decides the fast-time question}
\label{sec:v13-scope}
% =============================================================================

Retain the bounded-action, one-cell, bounded-product branch of
Versions~V10--V12.  Stop on every earlier routed time-correlation, selected
heat-tail, source-visible Reynolds-force, heat-occupation, enstrophy-square
overdrive, material-route escape, or source-action microdelocalization
output.  In particular,
\begin{equation}
 \boxed{
 \int_{J_n}\|G_n(s)\|_2^2\,ds\le L_*}
 \label{eq:v13-action-bound}
\end{equation}
for one fixed $L_*<\infty$, and
\begin{equation}
 \alpha_n=\sqrt{\frac{\mathcal R_n}{\nu}}\longrightarrow\infty.
 \label{eq:v13-alpha}
\end{equation}
The intervals $J_n=(-\tau_n,0)$ converge to
$J=(-\tau_\infty,0)$, with $\tau_n\to\tau_\infty>0$.

Retain the V12 interior cutoff $\eta\in C_c^\infty(J)$, the spatial cutoff
$\chi$, and the annular writers $Z_n\to Z$ strongly on compact sets.  Extend
all fields by zero outside $J_n$.  Define the finite-record selected scalar
density
\begin{equation}
 q_n(s)
 :=\eta(s)\int_{\mathbb R^3}
       \chi(y)\Omega_n(s,y)G_n(s,y)\cdot Z_n(y)\,dy.
 \label{eq:v13-slow-density}
\end{equation}
The integral is read on the expanding torus whenever $n$ is finite; the
fixed compact support of $\chi$ lies in that torus for all late $n$.

\begin{lemma}[Finite-record recovery of the V12 scalar]
\label{lem:v13-finite-scalar}
The functions $q_n$ are bounded in $L^2(J)$ and converge weakly there to
\begin{equation}
 q(s)
 :=\eta(s)\int_{\mathbb R^3}
       \chi(y)H(s,y)\cdot Z(y)\,dy.
 \label{eq:v13-slow-limit-density}
\end{equation}
Moreover,
\begin{equation}
 \boxed{
 \int_Jq_n(s)\,ds\longrightarrow\int_Jq(s)\,ds
 =\mathscr J(H),}
 \label{eq:v13-scalar-convergence}
\end{equation}
and hence, for all late $n$,
\begin{equation}
 \boxed{
 \left|\int_Jq_n(s)\,ds\right|\ge\frac{h_0}{2}.}
 \label{eq:v13-finite-nonzero}
\end{equation}
On the full intrinsic branch of
Theorem~\ref{thm:v12-heat-Euler-commutator},
\begin{equation}
 \boxed{
 q(s)=-\eta(s)\theta(s)
 \int_{\mathbb R^3}\chi(y)
 \psi(\Lambda_y)c_y[P_{-s}V(s)]\cdot Z(y)\,dy.}
 \label{eq:v13-explicit-Euler-density}
\end{equation}
\end{lemma}

\begin{proof}
Since $0\le\Omega_n\le1$ and $\chi Z_n$ is uniformly bounded in $L^2$,
\[
 |q_n(s)|
 \le \|\eta\|_\infty\|\chi Z_n\|_2\|G_n(s)\|_2.
\]
Equation~\eqref{eq:v13-action-bound} gives the uniform $L^2(J)$ bound.
For $\varphi\in L^2(J)$,
\[
 \int_J\varphi q_n
 =\int_J\!\int
  \Omega_nG_n\cdot
  \bigl(\varphi(s)\eta(s)\chi(y)Z_n(y)\bigr)\,dy\,ds.
\]
The bracketed tests converge strongly in $L^2(J\times\mathbb R^3)$ to the
same expression with $Z$, whereas
$\Omega_nG_n\rightharpoonup H$ locally in $L^2$.  This proves weak
convergence to \eqref{eq:v13-slow-limit-density}.  Taking
$\varphi=1$ on $\supp\eta$ proves
\eqref{eq:v13-scalar-convergence}; the V12 interior pairing gives
\eqref{eq:v13-finite-nonzero}.  On the full intrinsic branch,
$H=\theta\overline H$ and
$\overline H=-\psi(\Lambda_y)c_y[P_{-s}V]$ by
\eqref{eq:v12-stationary-heat-source}, giving
\eqref{eq:v13-explicit-Euler-density}.
\end{proof}

\begin{firewall}[This is the original held-out scalar]
\label{fire:v13-heldout-scalar}
The density $q_n$ is formed from the V10--V12 held-out annular source and its
predeclared writer.  It is not a new temporal observable.  Every declared
old-response or paid-stress return remains removed by its existing endpoint
map before $q_n$ is formed.  The passage to fast time below changes only the
time coordinate and never creates a second source price.
\end{firewall}

% =============================================================================
\section{Exact fast-time transformation of the field and scalar}
\label{sec:v13-fast-transform}
% =============================================================================

Put
\begin{equation}
 \widehat J_n:=\alpha_nJ_n=(-\alpha_n\tau_n,0)
 \label{eq:v13-fast-interval}
\end{equation}
and, for $\tau\in\widehat J_n$, define
\begin{equation}
 \begin{gathered}
 \widehat V_n(\tau,y):=V_n(\tau/\alpha_n,y),
 \qquad
 \widehat\Pi_n(\tau,y):=\Pi_n(\tau/\alpha_n,y),\\
 \widehat\Omega_n(\tau,y):=\Omega_n(\tau/\alpha_n,y),
 \qquad
 \widehat G_n(\tau,y):=G_n(\tau/\alpha_n,y).
 \end{gathered}
 \label{eq:v13-fast-fields}
\end{equation}

\begin{proposition}[V10 fast equation and target-native scalar measure]
\label{prop:v13-fast-equation-measure}
The fast velocity satisfies
\begin{equation}
 \boxed{
 \partial_\tau\widehat V_n
 +(\widehat V_n\cdot\nabla)\widehat V_n
 +\nabla\widehat\Pi_n
 =\alpha_n^{-1}\Delta\widehat V_n,
 \qquad \operatorname{div}\widehat V_n=0.}
 \label{eq:v13-fast-equation}
\end{equation}
Define the signed finite measure on $\mathbb R$ by
\begin{equation}
 \boxed{
 d\Sigma_n(\tau)
 :=\frac1{\alpha_n}q_n(\tau/\alpha_n)\,d\tau,}
 \label{eq:v13-fast-measure}
\end{equation}
where $q_n$ is extended by zero.  Then
\begin{equation}
 \boxed{
 \Sigma_n(\mathbb R)=\int_Jq_n(s)\,ds,}
 \label{eq:v13-fast-total}
\end{equation}
while
\begin{equation}
 \supp\Sigma_n
 \subseteq\alpha_n\supp\eta
 \subseteq[-\alpha_n\tau_+,-\alpha_n\tau_-].
 \label{eq:v13-fast-support}
\end{equation}
If $D_n(\tau)=\tau/\alpha_n$, then
\begin{equation}
 \boxed{
 (D_n)_*\Sigma_n=q_n(s)\,ds.}
 \label{eq:v13-dilation-pushforward}
\end{equation}
Consequently,
\begin{equation}
 (D_n)_*\Sigma_n\rightharpoonup q(s)\,ds
 \quad\text{weakly as finite signed measures on }J,
 \label{eq:v13-dilation-limit}
\end{equation}
and the limiting measure has nonzero total mass.
\end{proposition}

\begin{proof}
Substitute $s=\tau/\alpha_n$ in
\eqref{eq:v10-strong-coupling-equation}; the chain rule gives
$\partial_\tau\widehat V_n=\alpha_n^{-1}\partial_sV_n$ and proves
\eqref{eq:v13-fast-equation}.  The change of variables
$\tau=\alpha_ns$ gives
\eqref{eq:v13-fast-total} and
\eqref{eq:v13-dilation-pushforward}.  The support statement follows from
\eqref{eq:v12-interior-age-window}.  Finally,
Lemma~\ref{lem:v13-finite-scalar} gives weak $L^2$, hence weak measure,
convergence of $q_n(s)\,ds$ to $q(s)\,ds$.
\end{proof}

\begin{assessment}[Why the V1--V2 fast-time theorem is not being repeated]
\label{ass:v13-not-v2-repeat}
Versions~V1--V2 already proved that the contact-preserving amplitude
normalization leads to vanishing viscosity and that bounded cubic capacity
on a compact fast window cannot retain the original cubic contact.  The
present section records the exact V10 record-amplitude parameters and the
specific V12 heat--Euler scalar after every later spatial, covariance,
product, pressure, router, and represented-return short has been performed.
The new estimate below is an $L^2$ source-action statement for that literal
held-out scalar; it is not a second cubic-capacity argument.
\end{assessment}

% =============================================================================
\section{The scalar is critical in slow time and vanishing in fast time}
\label{sec:v13-fast-vanishing}
% =============================================================================

Let
\begin{equation}
 M_Z:=\sup_n\|\chi Z_n\|_2<\infty,
 \qquad
 L_\eta:=|\supp\eta|<\infty,
 \qquad
 C_q:=M_Z^2\|\eta\|_\infty^2L_*.
 \label{eq:v13-constants}
\end{equation}

\begin{theorem}[Exact $L^2$ scale of the fast-time scalar density]
\label{thm:v13-fast-L2-scale}
For all late $n$,
\begin{equation}
 \boxed{
 \frac{h_0^2}{4L_\eta}
 \le\|q_n\|_{L^2(J)}^2
 \le C_q.}
 \label{eq:v13-slow-L2-window}
\end{equation}
Writing
\begin{equation}
 \widehat q_n(\tau)
 :=\frac{d\Sigma_n}{d\tau}
 =\alpha_n^{-1}q_n(\tau/\alpha_n),
 \label{eq:v13-fast-density}
\end{equation}
one has the exact scaling identity
\begin{equation}
 \boxed{
 \alpha_n\|\widehat q_n\|_{L^2(\mathbb R)}^2
 =\|q_n\|_{L^2(J)}^2.}
 \label{eq:v13-L2-scaling-identity}
\end{equation}
In particular,
\begin{equation}
 \boxed{
 \frac{h_0^2}{4L_\eta\alpha_n}
 \le\|\widehat q_n\|_2^2
 \le\frac{C_q}{\alpha_n},}
 \label{eq:v13-fast-L2-two-sided}
\end{equation}
so $\widehat q_n\to0$ strongly in $L^2(\mathbb R)$ although
$|\int\widehat q_n|\ge h_0/2$.

More generally, for every $1<p\le2$,
\begin{equation}
 \boxed{
 \|\widehat q_n\|_{L^p(\mathbb R)}
 =\alpha_n^{1/p-1}\|q_n\|_{L^p(J)}
 \longrightarrow0.}
 \label{eq:v13-fast-Lp-vanishing}
\end{equation}
The scalar is therefore nontrivial only at the nonreflexive $L^1$/measure
level in fast time.
\end{theorem}

\begin{proof}
The upper bound in \eqref{eq:v13-slow-L2-window} follows from
\[
 |q_n(s)|\le M_Z\|\eta\|_\infty\|G_n(s)\|_2
\]
and \eqref{eq:v13-action-bound}.  The lower bound follows from
\eqref{eq:v13-finite-nonzero} and Cauchy--Schwarz on $\supp\eta$:
\[
 \frac{h_0}{2}
 \le\left|\int q_n\right|
 \le L_\eta^{1/2}\|q_n\|_2.
\]
The change of variables $\tau=\alpha_ns$ gives
\eqref{eq:v13-L2-scaling-identity} and, for every $p\ge1$,
\[
 \|\widehat q_n\|_p^p
 =\alpha_n^{1-p}\|q_n\|_p^p.
\]
For $1<p\le2$, the fixed support and the uniform $L^2$ bound make
$\|q_n\|_p$ uniformly bounded, proving
\eqref{eq:v13-fast-Lp-vanishing}.  The nonzero integral is
\eqref{eq:v13-fast-total} and \eqref{eq:v13-finite-nonzero}.
\end{proof}

\begin{corollary}[Uniform vanishing on every submacroscopic fast-time set]
\label{cor:v13-submacroscopic-vanishing}
For every measurable $E\subset\mathbb R$,
\begin{equation}
 \boxed{
 |\Sigma_n|(E)
 \le\left(\frac{C_q|E|}{\alpha_n}\right)^{1/2}.}
 \label{eq:v13-set-variation-bound}
\end{equation}
Consequently, if $|E_n|=o(\alpha_n)$, then
\begin{equation}
 \boxed{|\Sigma_n|(E_n)\longrightarrow0.}
 \label{eq:v13-sublinear-set-zero}
\end{equation}
In particular, for every fixed $L<\infty$,
\begin{equation}
 \boxed{
 \sup_{\tau_0\in\mathbb R}
 |\Sigma_n|([\tau_0-L,\tau_0+L])
 \le\left(\frac{2LC_q}{\alpha_n}\right)^{1/2}
 \longrightarrow0.}
 \label{eq:v13-moving-window-zero}
\end{equation}
Thus the vanishing holds even after an arbitrary moving fast-time
translation.
\end{corollary}

\begin{proof}
Apply Cauchy--Schwarz to the density of $\Sigma_n$ on $E$ and use the upper
bound in \eqref{eq:v13-fast-L2-two-sided}.  The two consequences are
immediate.
\end{proof}

\begin{corollary}[Translation-vanishing but dilation-nonzero]
\label{cor:v13-translation-dilation}
Let $T_a(\tau)=\tau-a$.  For every sequence $a_n\in\mathbb R$ and every
compact set $K\Subset\mathbb R$,
\begin{equation}
 \boxed{
 |(T_{a_n})_*\Sigma_n|(K)\longrightarrow0.}
 \label{eq:v13-all-translates-vague-zero}
\end{equation}
Hence every translated subsequence converges vaguely to the zero signed
measure on bounded fast-time sets.  In contrast, the dilated measures
$(D_n)_*\Sigma_n$ converge to the nonzero measure $q(s)\,ds$ of
\eqref{eq:v13-dilation-limit}.  The selected scalar is therefore a
\emph{dilation-only survivor}: no translation extracts a scalar-bearing
compact fast-time packet.
\end{corollary}

\begin{proof}
Apply \eqref{eq:v13-set-variation-bound} to
$E=a_n+K$, whose Lebesgue measure is independent of $n$.  The dilation
statement is Proposition~\ref{prop:v13-fast-equation-measure}.
\end{proof}

% =============================================================================
\section{Extensive participation and the exact bounded-window price}
\label{sec:v13-extensive}
% =============================================================================

The preceding vanishing is stronger than escape from the terminal endpoint:
even a moving compact window inside the occupied fast-time slab carries no
fixed scalar.  The following two statements quantify the resulting
extensive occupation without introducing a new age law.

Fix $L>0$.  Partition the convex hull of
$\alpha_n\supp\eta$ into disjoint half-open intervals
$(I_{n,j})_{j=1}^{N_n}$ of length at most $L$, and put
\begin{equation}
 a_{n,j}:=\Sigma_n(I_{n,j}).
 \label{eq:v13-window-coefficients}
\end{equation}
Zero intervals outside the actual support are harmless.

\begin{theorem}[Linear fast-window participation]
\label{thm:v13-linear-participation}
For all late $n$,
\begin{equation}
 \boxed{
 \sum_{j=1}^{N_n}|a_{n,j}|^2
 \le\frac{LC_q}{\alpha_n}.}
 \label{eq:v13-window-l2}
\end{equation}
Define the signed-window participation number
\begin{equation}
 \mathfrak N_n^{\rm fast}(L)
 :=\frac{\bigl(\sum_j|a_{n,j}|\bigr)^2}
          {\sum_j|a_{n,j}|^2},
 \label{eq:v13-fast-participation}
\end{equation}
with value $+\infty$ if the denominator vanishes.  Then
\begin{equation}
 \boxed{
 \mathfrak N_n^{\rm fast}(L)
 \ge\frac{h_0^2}{4LC_q}\,\alpha_n.}
 \label{eq:v13-linear-participation-lower}
\end{equation}
Moreover, if $\mathcal C_n$ is any cohort of at most $m_n$ partition
intervals, then
\begin{equation}
 \boxed{
 \left|\Sigma_n\!\left(\bigcup_{j\in\mathcal C_n}I_{n,j}\right)\right|
 \le\left(\frac{m_nLC_q}{\alpha_n}\right)^{1/2}.}
 \label{eq:v13-cohort-bound}
\end{equation}
Thus every cohort with $m_n=o(\alpha_n)$ carries a vanishing scalar.  A
fixed fraction $\delta>0$ requires at least
\begin{equation}
 \boxed{
 m_n\ge\frac{\delta^2}{LC_q}\,\alpha_n.}
 \label{eq:v13-linear-window-count}
\end{equation}
\end{theorem}

\begin{proof}
For every partition interval,
\[
 |a_{n,j}|^2
 \le |I_{n,j}|\int_{I_{n,j}}|\widehat q_n|^2
 \le L\int_{I_{n,j}}|\widehat q_n|^2.
\]
Summing and using \eqref{eq:v13-fast-L2-two-sided} proves
\eqref{eq:v13-window-l2}.  Since
$\sum_j|a_{n,j}|\ge|\Sigma_n(\mathbb R)|\ge h_0/2$,
\eqref{eq:v13-linear-participation-lower} follows.  For a cohort,
Cauchy--Schwarz and \eqref{eq:v13-window-l2} give
\[
 \left|\sum_{j\in\mathcal C_n}a_{n,j}\right|
 \le m_n^{1/2}
      \left(\sum_j|a_{n,j}|^2\right)^{1/2},
\]
which is \eqref{eq:v13-cohort-bound}.  Rearrangement gives
\eqref{eq:v13-linear-window-count}.
\end{proof}

For a measurable fast-time set $E$, define its normalized source action by
\begin{equation}
 \mathcal A_n^{\rm fast}(E)
 :=\frac1{\alpha_n}
   \int_E\|\widehat G_n(\tau)\|_2^2\,d\tau.
 \label{eq:v13-fast-action}
\end{equation}
Then
\begin{equation}
 \mathcal A_n^{\rm fast}(\mathbb R)
 =\int_{J_n}\|G_n(s)\|_2^2\,ds\le L_*.
 \label{eq:v13-fast-total-action}
\end{equation}

\begin{theorem}[A scalar-bearing bounded fast window forces linear action overdrive]
\label{thm:v13-window-action-overdrive}
For every measurable $E\subset\mathbb R$ of finite measure,
\begin{equation}
 \boxed{
 |\Sigma_n|(E)^2
 \le\frac{M_Z^2\|\eta\|_\infty^2|E|}{\alpha_n}
       \mathcal A_n^{\rm fast}(E).}
 \label{eq:v13-local-action-price}
\end{equation}
Consequently, if $|E_n|\le L$ and
$|\Sigma_n(E_n)|\ge\delta>0$, then
\begin{equation}
 \boxed{
 \mathcal A_n^{\rm fast}(E_n)
 \ge
 \frac{\delta^2}{M_Z^2\|\eta\|_\infty^2L}\,\alpha_n.}
 \label{eq:v13-linear-action-overdrive}
\end{equation}
This is impossible under \eqref{eq:v13-fast-total-action}.  More generally,
any bounded-window branch carrying a fixed scalar is exactly the
source-action overdrive already returned by V10 and the enstrophy-square
overdrive already returned by V11.
\end{theorem}

\begin{proof}
At fast time $\tau$,
\[
 \left|
 \eta(\tau/\alpha_n)
 \int\chi\widehat\Omega_n\widehat G_n\cdot Z_n
 \right|
 \le M_Z\|\eta\|_\infty\|\widehat G_n(\tau)\|_2.
\]
Integrate its absolute value on $E$, include the factor $\alpha_n^{-1}$
from \eqref{eq:v13-fast-measure}, and apply Cauchy--Schwarz.  Since
$\int_E\|\widehat G_n\|_2^2=\alpha_n
\mathcal A_n^{\rm fast}(E)$, this gives
\eqref{eq:v13-local-action-price}.  Rearrangement proves
\eqref{eq:v13-linear-action-overdrive}; the total-action contradiction is
\eqref{eq:v13-fast-total-action}.  The final identification uses the exact
V10--V11 action/enstrophy-square comparison and adds no new source row.
\end{proof}

\begin{corollary}[Linear fast-time source action]
\label{cor:v13-linear-fast-action}
The fast-time source action on the occupied slab is itself extensive:
\begin{equation}
 \boxed{
 c_A\alpha_n
 \le
 \int_{\alpha_n\supp\eta}\|\widehat G_n(\tau)\|_2^2\,d\tau
 \le L_*\alpha_n,}
 \label{eq:v13-linear-source-action}
\end{equation}
where
\begin{equation}
 c_A
 :=\frac{h_0^2}
 {4M_Z^2\|\eta\|_\infty^2L_\eta}>0.
 \label{eq:v13-cA}
\end{equation}
Thus the occupation escape has a positive mean source-action density; it is
not disappearance of the source.
\end{corollary}

\begin{proof}
The upper bound is \eqref{eq:v13-action-bound} after the fast-time change of
variables.  For the lower bound, apply Cauchy--Schwarz to
\eqref{eq:v13-finite-nonzero}:
\[
 \frac{h_0}{2}
 \le M_Z\|\eta\|_\infty L_\eta^{1/2}
      \left(\int_{\supp\eta}\|G_n(s)\|_2^2\,ds\right)^{1/2}.
\]
Rearrange and multiply by $\alpha_n$.
\end{proof}

% =============================================================================
\section{Sharpness and closure of the compact-profile route}
\label{sec:v13-profile-stop}
% =============================================================================

\begin{countertheorem}[A nonzero dilated scalar need not have any translated atom]
\label{cth:v13-dilation-no-translate}
Let $q\in L^2_c((-2,-1))$ satisfy $\int q\ne0$ and let
$\alpha_n\to\infty$.  Define
\begin{equation}
 d\mu_n(\tau)
 :=\alpha_n^{-1}q(\tau/\alpha_n)\,d\tau.
 \label{eq:v13-scalar-model}
\end{equation}
Then
\begin{equation}
 \int d\mu_n=\int q\ne0,
 \qquad
 \alpha_n\left\|\frac{d\mu_n}{d\tau}\right\|_2^2
 =\|q\|_2^2,
 \label{eq:v13-model-scales}
\end{equation}
while, for every $L<\infty$,
\begin{equation}
 \sup_{a\in\mathbb R}|\mu_n|([a-L,a+L])\longrightarrow0.
 \label{eq:v13-model-window-zero}
\end{equation}
The dilation $\tau\mapsto\tau/\alpha_n$ recovers $q(s)\,ds$, but every
translated compact-window limit is zero.  Hence no abstract measure
compactness theorem can turn a nonzero dilation limit into a scalar-bearing
translated profile without an additional concentration estimate.
\end{countertheorem}

\begin{proof}
The first two identities are changes of variables.  Cauchy--Schwarz gives
\[
 |\mu_n|([a-L,a+L])
 \le(2L)^{1/2}\alpha_n^{-1/2}\|q\|_2,
\]
uniformly in $a$.
\end{proof}

\begin{firewall}[The scalar model is not a Navier--Stokes trajectory]
\label{fire:v13-model-functional}
Countertheorem~\ref{cth:v13-dilation-no-translate} is a sharpness statement
for the temporal compactness implication only.  It is not a solution of the
fast equation.  The actual Navier--Stokes sequence satisfies the stronger
source-action and field identities of the preceding sections.  Those
identities nevertheless produce the same translation-vanishing conclusion
on the bounded-action branch.
\end{firewall}

\begin{theorem}[No scalar-bearing bounded-window dynamical profile on the clean branch]
\label{thm:v13-no-one-profile}
Let $a_n\in\mathbb R$ be arbitrary.  Suppose the translated fast fields
\begin{equation}
 \widehat V_n^{[a]}(\tau,y)
 :=\widehat V_n(a_n+\tau,y)
 \label{eq:v13-translated-fields}
\end{equation}
converge on a fixed compact fast-time cylinder in any topology sufficient to
pass the local heat-covariance source and writer pairing.  Then the
contribution of that cylinder to the original normalized source-paid scalar
is zero.  The same holds for every fixed finite family, and more generally
for every family of $o(\alpha_n)$ bounded windows.

Therefore a nonzero local Euler or Euler--Reynolds limit may exist, but it
cannot inherit the nonzero scalar
\eqref{eq:v12-source-paid-Euler-scalar} in the original normalization.  A
terminal-trace, reflected, backward-uniqueness, or Liouville theorem applied
to one such compact profile does not close the bounded-action branch.
\end{theorem}

\begin{proof}
The scalar contribution of a translated compact time cylinder is the
restriction of $\Sigma_n$ to one translated bounded interval.  It tends to
zero in total variation by
\eqref{eq:v13-moving-window-zero}, independently of the topology used for
the fields.  A finite family is a cohort with bounded cardinality in
Theorem~\ref{thm:v13-linear-participation}; an $o(\alpha_n)$ family is
covered by \eqref{eq:v13-cohort-bound}.  Passing the local field equations
may produce a nonzero dynamical limit, but the original normalized scalar
on that limit is the limit of these vanishing restrictions and is therefore
zero.
\end{proof}

\begin{assessment}[The upstream correction to the V12 endpoint]
\label{ass:v13-upstream-correction}
Version~V12 allowed two temporal outputs on its clean intrinsic branch:
fast-time occupation escape or one scalar-bearing bounded-window dynamical
profile.  Theorem~\ref{thm:v13-window-action-overdrive} shows that the
second output is incompatible with the bounded source-action hypothesis
under which the clean profile was constructed.  It can occur only after
returning to the earlier source-action/enstrophy-square overdrive branch.
Thus, on the actual V12 bounded-action branch, the temporal output is no
longer a dichotomy:
\[
 \boxed{\text{fast-time occupation escape is forced.}}
\]
The nonzero slow-time heat--Euler scalar is a macroscopic occupation average,
not the scalar of one compact fast-time orbit.
\end{assessment}

% =============================================================================
\section{Version V13 integrated theorem and exact stopping boundary}
\label{sec:v13-integrated}
% =============================================================================

\begin{theorem}[Version V13 extensive fast-time occupation theorem]
\label{thm:v13-integrated}
Preserving every theorem and limitation of Versions~V1--V12, Version~V13
proves the following on the full intrinsic bounded-action branch.
\begin{enumerate}[label=\textup{(V13.\arabic*)},leftmargin=6em]
\item The finite-record scalar densities $q_n$ are uniformly bounded and
      nontrivial in slow-time $L^2$.  They converge weakly to the explicit
      heat--Euler density \eqref{eq:v13-explicit-Euler-density}.
\item The exact V10 fast-time change gives the vanishing-viscosity equation
      \eqref{eq:v13-fast-equation} and the target-native signed measure
      $\Sigma_n$ of \eqref{eq:v13-fast-measure}.  Its dilation back to slow
      time recovers $q_n(s)\,ds$ exactly.
\item The fast scalar density has the two-sided critical scale
      $\|d\Sigma_n/d\tau\|_2^2\asymp\alpha_n^{-1}$ and converges strongly to
      zero in every $L^p$, $1<p\le2$, while its total signed mass stays
      bounded away from zero.
\item Every fast-time set of length $o(\alpha_n)$ has vanishing total
      variation.  In particular every moving bounded window and every
      translated compact subsequence is scalar null.
\item A fixed-length partition has scalar participation at least
      $c\alpha_n$.  Any cohort of $o(\alpha_n)$ windows is negligible, and a
      fixed scalar fraction requires linearly many windows.
\item A fixed scalar on one bounded fast window forces normalized source
      action at least $c\alpha_n$.  It is therefore precisely the earlier
      source-action/enstrophy-square overdrive branch and is impossible on
      the present bounded-action subsequence.
\item The total fast-time source action is itself comparable to
      $\alpha_n$.  Occupation escape is thus an extensive positive-action
      history, not loss of the source.
\item No scalar-bearing compact dynamical profile can be extracted by time
      translation on this branch.  The nonzero stationary Euler scalar is
      retained only by the macroscopic dilation, so a one-profile terminal
      trace or Liouville argument is downstream of an absent branch.
\end{enumerate}
No item provides a nonduplicating finite-stock payment bound for the
extensive occupation, a monotone replenishment law, an exclusion of the
source-action overdrive branch, terminal regularity, or global
three-dimensional Navier--Stokes regularity.
\end{theorem}

\begin{proof}
Item~\textup{(V13.1)} is Lemma~\ref{lem:v13-finite-scalar}.  Item
\textup{(V13.2)} is
Proposition~\ref{prop:v13-fast-equation-measure}.  Items
\textup{(V13.3)--(V13.4)} are
Theorem~\ref{thm:v13-fast-L2-scale} and
Corollaries~\ref{cor:v13-submacroscopic-vanishing}--
\ref{cor:v13-translation-dilation}.  Item~\textup{(V13.5)} is
Theorem~\ref{thm:v13-linear-participation}.  Items
\textup{(V13.6)--(V13.7)} are
Theorem~\ref{thm:v13-window-action-overdrive} and
Corollary~\ref{cor:v13-linear-fast-action}.  The final item is
Theorem~\ref{thm:v13-no-one-profile} and
Assessment~\ref{ass:v13-upstream-correction}.
\end{proof}

\begingroup\small
\begin{longtable}{p{0.25\textwidth}p{0.19\textwidth}p{0.48\textwidth}}
\caption{NS--A Version V13 proof-status ledger.}\label{tab:v13-ledger}\\
\toprule Row & Status & Exact conclusion\\\midrule
\endfirsthead
\toprule Row & Status & Exact conclusion\\\midrule
\endhead
V1--V12 prefix & \PROVED{} preserved
 &The complete V12 preterminal source is retained byte-for-byte; no earlier
 theorem is rewritten.\\[0.3em]
Earlier V12 source defects & \INHERITED{} stop first
 &V13 is read only after the four V12 source-incidence and compactness defects
 vanish.\\[0.3em]
Finite-record slow scalar & \PROVED{} nonzero
 &$q_n\rightharpoonup q$ in $L^2$, with nonzero integral and explicit
 heat--Euler density on the intrinsic branch.\\[0.3em]
Fast vanishing-viscosity equation & \PROVED{} exact normalization
 &The V10 record-amplitude sequence satisfies
 \eqref{eq:v13-fast-equation}; no unit-viscosity ancient profile is asserted.\\[0.3em]
Fast scalar measure & \PROVED{} target native
 &$d\Sigma_n=\alpha_n^{-1}q_n(\tau/\alpha_n)d\tau$ is the original held-out
 scalar after time change, not a new age law.\\[0.3em]
Fast $L^p$ density & \PROVED{} vanishing
 &$\|d\Sigma_n/d\tau\|_2^2\asymp\alpha_n^{-1}$ and every
 $1<p\le2$ norm tends to zero while total signed mass remains nonzero.\\[0.3em]
Moving bounded fast window & \OBSTRUCTION{} excluded on bounded action
 &Every translated compact window has vanishing total variation.\\[0.3em]
Submacroscopic fast cohort & \PROVED{} negligible
 &Every time set of length $o(\alpha_n)$ and every cohort of
 $o(\alpha_n)$ bounded windows carries vanishing scalar.\\[0.3em]
Extensive participation & \PROVED{} linear
 &At least $c\alpha_n$ bounded windows are required to carry a fixed scalar
 fraction.\\[0.3em]
Bounded-window scalar atom & \OBSTRUCTION{} action overdrive
 &A fixed atom forces normalized source action $\gtrsim\alpha_n$, returning
 the V10--V11 source-action/enstrophy-square branch.\\[0.3em]
Fast source action & \PROVED{} extensive
 &The total source action on the occupied fast slab is between
 $c\alpha_n$ and $L_*\alpha_n$.\\[0.3em]
Dilation versus translation & \PROVED{} separated
 &Slow dilation recovers the nonzero heat--Euler measure; every translated
 compact measure limit is zero.\\[0.3em]
One-profile Liouville route & \OBSTRUCTION{} downstream here
 &A compact dynamical profile may exist, but it cannot inherit the selected
 scalar in the original normalization on this branch.\\[0.3em]
Extensive occupation payment & \OPEN{} first physical row
 &Requires an actual nonduplicating action-to-stock or monotone replenishment
 estimate on the same source history.\\[0.3em]
Three-dimensional regularity & \OPEN
 &Neither the extensive occupation nor the action-overdrive branch is
 excluded for every fixed datum.\\
\bottomrule
\end{longtable}
\endgroup

% =============================================================================
\section{Exact mandate for Version V14}
\label{sec:v13-next}
% =============================================================================

\begin{programme}[Price the extensive occupation or stop at the replenishment row]
\label{prog:v13-next}
Preserve Versions~V1--V13 verbatim.  Do not seek another bounded fast-time
profile on the bounded-action branch; Theorem~\ref{thm:v13-no-one-profile}
closes that route for the selected scalar.  Do not add another age law,
path law, statistical-solution space, concentration metric, or slow-time
product quotient.

Proceed only in the following order.
\begin{enumerate}[label=\textup{(V14.\arabic*)},leftmargin=5.7em]
\item Retain the exact signed window coefficients $a_{n,j}$ and the original
      source/writer incidence before taking positive parts.  Preserve every
      represented return and paid-stress follower through its endpoint map.
\item Test whether the extensive cohort has a physical, nonduplicating
      action-to-stock incidence: a finite native source payment, a monotone
      replenishment law, or a first-exit stock which bounds the same signed
      scalar across the $\gtrsim\alpha_n$ windows.
\item If such a stock exists, compare its upper budget with
      \eqref{eq:v13-linear-participation-lower} and
      \eqref{eq:v13-linear-source-action}.  Internal cancellation and
      replenishment must be assembled before any Jordan or positive split.
\item If no such incidence is available, stop and export
      \emph{extensive fast-time occupation} as the exact terminal PDE input.
      Do not disguise the missing monotone-replenishment estimate as a
      compactness theorem.
\item Reopen a bounded-window dynamical profile only on the separately typed
      source-action/enstrophy-square overdrive branch, where
      \eqref{eq:v13-linear-action-overdrive} supplies the required action
      atom.  Its normalization and product/pressure compactness must then be
      redone on that branch rather than imported from the bounded-action
      profile.
\end{enumerate}
The admissible endpoint is
\[
 \boxed{
 \begin{gathered}
 \text{finite nonduplicating occupation payment}
 \ \big|\ 
 \text{monotone replenishment survivor}\\[0.15em]
 \text{unpaid extensive fast-time occupation}
 \ \big|\ 
 \text{source-action atom or overdrive}\\[0.15em]
 \text{with a separately normalized profile}.
 \end{gathered}}
\]
\end{programme}

\begin{assessment}[Programme position after V13]
\label{ass:v13-position}
The V12 temporal question has now been answered on its own hypotheses.  The
bounded source-action branch cannot place the source-paid heat--Euler scalar
on a bounded fast-time window; it spreads that scalar over a linearly growing
number of windows with linearly growing total fast-time source action.  The
remaining issue is not profile compactness.  It is whether this extensive
occupation is paid, replenished, or cancelled by a finite physical stock on
the same fixed datum.  That is the monotone-replenishment or critical-action
row already identified upstream; another one-profile extraction would be
circular.
\end{assessment}

\paragraph{Version V13 history.}
Kept every routed, heat-tail, Reynolds-force, occupation, and product short
of V12 fixed.  On the full intrinsic bounded-action branch, formed the exact
finite-record scalar density, carried it to the V10 fast time, and proved its
critical two-sided $L^2$ scaling.  Established uniform total-variation
vanishing on every submacroscopic fast-time set, linear participation across
bounded windows, and linear total fast-time source action.  Proved that any
scalar-bearing bounded window forces the earlier source-action/enstrophy-
square overdrive, so the clean branch necessarily returns extensive
occupation escape.  Separated translation-vanishing compact profiles from
the nonzero dilation limit and moved the next attack to the physical
payment/replenishment row.

\begin{assessment}[Append-only integrity]
\label{ass:v13-integrity}
The complete Version~V12 source preceding its sole final document terminator
is preserved verbatim.  Its preterminal SHA--256 digest is
\begin{center}\scriptsize\ttfamily
72030710c6eef02f9e36bb7ca14629c23e2ddfa6e0d0ca37e672e0fd3b61038d
\end{center}
The present material begins at Section~\ref{sec:v13-scope}; the final command
below is again unique.
\end{assessment}

\addtocontents{toc}{\protect\endgroup}
% =============================================================================
% END VERSION V13 MATERIAL -- append later versions before final terminator
% =============================================================================

\clearpage
% =============================================================================
% VERSION V14: NEW MATERIAL.  The complete V1--V13 source above is preserved
% verbatim; only its sole active \end{document} line was removed before append.
% =============================================================================
\hypersetup{pdftitle={NS-A Terminal Source Concentration Programme: Cumulative V14},pdfauthor={A. Pelissier}}
\makeatletter
\addtocontents{toc}{\protect\begingroup\protect\makeatletter}
\addtocontents{toc}{\protect\def\protect\l@section{\protect\@dottedtocline{1}{0em}{4.7em}}}
\addtocontents{toc}{\protect\def\protect\l@subsection{\protect\@dottedtocline{2}{1.5em}{5.3em}}}
\addtocontents{toc}{\protect\makeatother}
\makeatother
\renewcommand{\thesection}{V14.\arabic{section}}
\renewcommand{\thesubsection}{V14.\arabic{section}.\arabic{subsection}}
\renewcommand{\theequation}{V14.\arabic{equation}}
\setcounter{section}{0}
\setcounter{equation}{0}
\renewcommand{\theHsection}{V14.\arabic{section}}
\renewcommand{\theHsubsection}{V14.\arabic{section}.\arabic{subsection}}
\renewcommand{\theHtheorem}{V14.\arabic{section}.\arabic{theorem}}
\renewcommand{\theHequation}{V14.\arabic{equation}}

\begin{center}
{\LARGE\bfseries NS--A: Version V14}\\[0.5em]
{\large Bounded Scalar Turnover, Canonical Stock Reflection,\\
and the Non-Summable Replenishment Stop}\\[0.5em]
{\normalsize 3 September 2026}
\end{center}
\addcontentsline{toc}{part}{Version V14: bounded scalar turnover and the replenishment stop}

\begin{abstract}
Version~V13 proved that the selected source-paid scalar occupies order
$\alpha_n$ bounded fast-time windows and that its raw fast-time source action
is of order $\alpha_n$.  The first possible reading is that an extensive
number of physical payments has accumulated.  The present continuation tests
that reading before any new compactness or profile construction is admitted.

The first result is an upstream correction.  Although the participation
number diverges linearly, the total variation of the selected scalar is
uniformly bounded:
\[
 \|\Sigma_n\|_{\rm TV}
 =\int_J|q_n(s)|\,\dd s
 \le |\supp\eta|^{1/2}C_q^{1/2}.
\]
After one fixed orientation is selected on a subsequence, the positive Jordan
part is carried by order $\alpha_n$ windows, but its total mass remains
bounded.  The extensive window count therefore does not exhaust a larger and
larger scalar reservoir inside one record.  In fact one bounded abstract stock
can pay every window of one record without replenishment.  Rebuilding that
stock at the next record, however, is itself an unrecorded replenishment.

The second result gives the exact missing calculation.  For any chronologically
ordered signed debit sequence, the minimum external replenishment needed to
keep a nonnegative stock is the one-sided Skorokhod reflection
\[
 R_k^{\min}=\max_{0\le m\le k}(C_m-S_0)_+,
 \qquad C_m=\sum_{i\le m}b_i.
\]
Apply this only after all signed window contributions and represented returns
have been assembled.  If $\lambda_n$ is a predeclared physical calibration
from the normalized V13 scalar to one common stock, then every stock incidence
obeys
\[
 R_N^{\rm ext}
 \ge \frac{h_0}{2}\sum_{n\le N}\lambda_n-S_{\rm in}.
\]
Hence a finite stock and finite total replenishment force
$\sum_n\lambda_n<\infty$.  Under the natural record-amplitude calibration
$\lambda_n=\mathcal R_n$, replenishment grows at least like
$\sum_n\mathcal R_n$ and is exponential in the record number.

The existing finite native source length cannot repair this conclusion.  Any
incidence which bounds the calibrated debit by the native source mass forces
$\sum_n\lambda_n<\infty$, so it pays only after a summable calibration
collapse.  Kinetic energy supplies no nonlinear debit because
$\langle\mathcal B(u),u\rangle=0$.  The exact caloric critical-energy rectangle
does carry the nonlinear contact, but it is an unbounded inventory whose
record-to-record increase is itself supplied by a positive nonlinear
replenishment of order $\mathcal R_n^2$.  Finally, the linear raw fast-time
action is a change-of-time identity and not a stock decrement.

Thus the V13 extensive branch has reached a precise physical stop.  It is
either assigned a summably collapsing stock calibration, accompanied by a
non-summable replenishment law, or left without an action-to-stock incidence.
No finite-stock contradiction, terminal regularity theorem, or global
three-dimensional Navier--Stokes conclusion is claimed.
\end{abstract}

% =============================================================================
\section{The extensive window law has bounded scalar turnover}
\label{sec:v14-turnover}
% =============================================================================

Retain every hypothesis and definition of Version~V13.  In particular,
$q_n$, $\Sigma_n$, the fixed-length fast-time partition
$(I_{n,j})_{j=1}^{N_n}$, and the signed window coefficients
$a_{n,j}=\Sigma_n(I_{n,j})$ are unchanged.  The represented-response,
paid-stress, routed, heat-tail, Reynolds-force, and occupation shorts have
already been performed before these coefficients are formed.

After passage to a subsequence there is one fixed
$\sigma_*\in\{-1,1\}$ such that
\begin{equation}
 \boxed{
 A_n:=\sigma_*\Sigma_n(\mathbb R)
 =\sum_{j=1}^{N_n}b_{n,j}\ge\frac{h_0}{2},
 \qquad b_{n,j}:=\sigma_*a_{n,j}.}
 \label{eq:v14-oriented-window-law}
\end{equation}
The orientation is fixed once on the subsequence.  It is not changed from
window to window or record to record.

Put
\begin{equation}
 \boxed{
 \mathfrak V_*:=L_\eta^{1/2}C_q^{1/2}.}
 \label{eq:v14-variation-cap}
\end{equation}

\begin{theorem}[Bounded scalar variation despite linear participation]
\label{thm:v14-bounded-turnover}
For every late $n$,
\begin{equation}
 \boxed{
 \|\Sigma_n\|_{\rm TV}
 =\int_J|q_n(s)|\,\dd s
 \le\mathfrak V_*,}
 \label{eq:v14-TV-cap}
\end{equation}
and consequently
\begin{equation}
 \boxed{
 \sum_{j=1}^{N_n}|b_{n,j}|
 \le\mathfrak V_*.
 }
 \label{eq:v14-coarse-variation-cap}
\end{equation}
Let
\begin{equation}
 P_n:=\sum_j(b_{n,j})_+,
 \qquad
 C_n:=\sum_j(b_{n,j})_-.
 \label{eq:v14-Jordan-totals}
\end{equation}
Then
\begin{equation}
 \boxed{
 A_n=P_n-C_n,
 \qquad
 P_n+C_n\le\mathfrak V_*,
 \qquad
 P_n\ge\frac{h_0}{2}.}
 \label{eq:v14-Jordan-balance}
\end{equation}
Moreover,
\begin{equation}
 \max_j|b_{n,j}|
 \le\left(\frac{LC_q}{\alpha_n}\right)^{1/2}
 \longrightarrow0.
 \label{eq:v14-max-window-zero}
\end{equation}
If
\(
 N_n^+:=\#\{j:b_{n,j}>0\}
\), then
\begin{equation}
 \boxed{
 N_n^+
 \ge\frac{h_0^2}{4LC_q}\,\alpha_n.}
 \label{eq:v14-positive-window-count}
\end{equation}
Thus the positive Jordan part occupies linearly many windows while its total
mass remains uniformly bounded.
\end{theorem}

\begin{proof}
By the definition of $\Sigma_n$ and the change of variables
$\tau=\alpha_ns$,
\[
 \|\Sigma_n\|_{\rm TV}
 =\int_{\mathbb R}\alpha_n^{-1}|q_n(\tau/\alpha_n)|\,\dd\tau
 =\int_J|q_n(s)|\,\dd s.
\]
Cauchy--Schwarz, the support bound, and
\eqref{eq:v13-slow-L2-window} give
\[
 \int_J|q_n|
 \le L_\eta^{1/2}\|q_n\|_2
 \le L_\eta^{1/2}C_q^{1/2}=\mathfrak V_*.
\]
The variation of a signed measure on a finite measurable partition dominates
the sum of the absolute values of its cell masses, proving
\eqref{eq:v14-coarse-variation-cap}.  Equations
\eqref{eq:v14-oriented-window-law} and
\eqref{eq:v14-Jordan-totals} give the first identity in
\eqref{eq:v14-Jordan-balance}; the other two follow from
\eqref{eq:v14-coarse-variation-cap} and $A_n\ge h_0/2$.

The pointwise bound follows from
\eqref{eq:v13-window-l2}.  Finally,
\[
 P_n^2
 \le N_n^+\sum_j(b_{n,j})_+^2
 \le N_n^+\sum_j|a_{n,j}|^2
 \le N_n^+\frac{LC_q}{\alpha_n}.
\]
Use $P_n\ge h_0/2$ and rearrange.
\end{proof}

\begin{corollary}[One record has a uniformly bounded abstract stock]
\label{cor:v14-one-record-stock}
Order the windows chronologically and put
\begin{equation}
 B_{n,k}:=\sum_{j=1}^{k}b_{n,j},
 \qquad B_{n,0}:=0.
 \label{eq:v14-record-partial-sums}
\end{equation}
Then
\begin{equation}
 \boxed{
 |B_{n,k}|\le\mathfrak V_*
 \quad(0\le k\le N_n).}
 \label{eq:v14-partial-sum-cap}
\end{equation}
Consequently
\begin{equation}
 \mathsf S_{n,k}^{\rm abs}:=
 \mathfrak V_*-B_{n,k}
 \label{eq:v14-abstract-record-stock}
\end{equation}
is nonnegative and satisfies the no-replenishment balance
\begin{equation}
 \boxed{
 \mathsf S_{n,k}^{\rm abs}
 =\mathsf S_{n,k-1}^{\rm abs}-b_{n,k}.}
 \label{eq:v14-abstract-record-balance}
\end{equation}
Thus the linear window count does not force an unbounded scalar reservoir
inside one record.
\end{corollary}

\begin{proof}
Every partial sum is bounded in modulus by
$\sum_j|b_{n,j}|$, and Theorem~\ref{thm:v14-bounded-turnover} bounds the
latter by $\mathfrak V_*$.  The displayed balance is immediate.
\end{proof}

\begin{firewall}[The abstract stock is rebuilt at every record]
\label{fire:v14-abstract-stock-not-physical}
The stock in Corollary~\ref{cor:v14-one-record-stock} is a mathematical
primitive of one finite signed sequence.  It is initialized afresh at value
$\mathfrak V_*$ for each $n$.  It is not a physical Navier--Stokes stock and
cannot be concatenated across records without accounting for that reset as a
replenishment.  The corollary has one purpose: it rules out the inference that
linearly many windows alone create linearly growing scalar turnover.
\end{firewall}

\begin{theorem}[Jordan splitting is the unique minimum-turnover split]
\label{thm:v14-Jordan-minimality}
Fix one finite signed sequence $(b_j)$.  Suppose nonnegative numbers
$p_j,r_j$ satisfy
\begin{equation}
 b_j=p_j-r_j.
 \label{eq:v14-arbitrary-positive-split}
\end{equation}
Then there are unique $d_j\ge0$ such that
\begin{equation}
 \boxed{
 p_j=(b_j)_++d_j,
 \qquad
 r_j=(b_j)_-+d_j.}
 \label{eq:v14-duplicated-turnover}
\end{equation}
Hence
\begin{equation}
 \sum_jp_j=\sum_j(b_j)_++\sum_jd_j,
 \qquad
 \sum_jr_j=\sum_j(b_j)_-+\sum_jd_j.
 \label{eq:v14-positive-split-totals}
\end{equation}
Every positive decomposition other than the Jordan split adds the same
duplicated turnover to both sides.
\end{theorem}

\begin{proof}
If $b_j\ge0$, equation~\eqref{eq:v14-arbitrary-positive-split} gives
$p_j=b_j+r_j$; put $d_j=r_j$.  If $b_j<0$, it gives
$r_j=-b_j+p_j$; put $d_j=p_j$.  These formulas are forced and prove
uniqueness.
\end{proof}

\begin{assessment}[What is and is not extensive]
\label{ass:v14-extensivity}
The number of positive windows and the raw fast-time source action are of
order $\alpha_n$.  The net scalar, its total variation, its minimum Jordan
payment, and the maximum partial-sum stock required inside one record are all
of order one.  Therefore no argument may multiply a per-window scalar price
by the participation number.  A cofinal contradiction can arise only from a
single stock carried across records or from an independently proved physical
price for the positive action.
\end{assessment}

% =============================================================================
\section{Canonical stock reflection and the exact replenishment law}
\label{sec:v14-stock-reflection}
% =============================================================================

The next theorem is applied to the V13 coefficients only after their signed
assembly.  It is included because it is the exact arithmetic content of the
first-exit or monotone-replenishment row; it is not a new response compiler.

\begin{theorem}[One-sided Skorokhod reflection of a signed debit history]
\label{thm:v14-Skorokhod}
Let $d_1,\ldots,d_M\in\mathbb R$ be chronologically ordered signed debits,
put
\begin{equation}
 D_k:=\sum_{i=1}^{k}d_i,
 \qquad D_0:=0,
 \label{eq:v14-cumulative-debit}
\end{equation}
and fix an initial stock $S_0\ge0$.  Define
\begin{equation}
 \boxed{
 R_k^{\min}:=
 \max_{0\le m\le k}(D_m-S_0)_+,
 \qquad
 S_k^{\min}:=S_0-D_k+R_k^{\min}.}
 \label{eq:v14-Skorokhod-formula}
\end{equation}
Then:
\begin{enumerate}[label=\textup{(S\arabic*)},leftmargin=3.5em]
\item $R_k^{\min}$ is nondecreasing, $R_0^{\min}=0$, and
      $S_k^{\min}\ge0$;
\item with $r_k^{\min}=R_k^{\min}-R_{k-1}^{\min}\ge0$,
      \begin{equation}
       \boxed{
       S_k^{\min}=S_{k-1}^{\min}-d_k+r_k^{\min};}
       \label{eq:v14-Skorokhod-balance}
      \end{equation}
\item if $r_k^{\min}>0$, then $S_k^{\min}=0$;
\item for every other nondecreasing $R_k$ with $R_0=0$ and
      $S_k=S_0-D_k+R_k\ge0$, one has
      \begin{equation}
       \boxed{R_k\ge R_k^{\min}\quad(0\le k\le M).}
       \label{eq:v14-Skorokhod-minimality}
      \end{equation}
\end{enumerate}
In particular, zero external replenishment is possible exactly when
\begin{equation}
 \boxed{S_0\ge\max_{0\le k\le M}D_k.}
 \label{eq:v14-no-replenishment-stock}
\end{equation}
Negative debits enter through the signed cumulative sum and replenish the
stock internally; they are not assigned a second positive price.
\end{theorem}

\begin{proof}
The running maximum in \eqref{eq:v14-Skorokhod-formula} is nondecreasing.
For every $k$,
$R_k^{\min}\ge(D_k-S_0)_+$, so
$S_k^{\min}=S_0-D_k+R_k^{\min}\ge0$.  Taking first differences gives
\eqref{eq:v14-Skorokhod-balance}.  If the running maximum increases at $k$,
then $R_k^{\min}=D_k-S_0>0$, and hence $S_k^{\min}=0$.

For any competing replenishment, nonnegativity gives
$R_k\ge D_k-S_0$ for each $k$.  Since $R_k$ is nondecreasing,
$R_k\ge\max_{m\le k}(D_m-S_0)_+=R_k^{\min}$.  The final equivalence is the
case $R_k\equiv0$.
\end{proof}

\begin{countertheorem}[Scalar and action marginals do not determine first-exit stock]
\label{cth:v14-order-matters}
The two debit histories
\begin{equation}
 (1,1,-1),
 \qquad
 (-1,1,1)
 \label{eq:v14-order-counterexample}
\end{equation}
have the same total debit, the same Jordan totals, and the same quadratic
action:
\begin{equation}
 \sum d_j=1,
 \qquad
 \sum(d_j)_+=2,
 \qquad
 \sum(d_j)_-=1,
 \qquad
 \sum d_j^2=3.
 \label{eq:v14-order-same-marginals}
\end{equation}
Nevertheless, with zero initial stock their minimum replenishments are two
and one, respectively.  Thus total scalar, positive/negative variation, and
quadratic action do not determine the first-exit stock; chronology is an
independent physical row.
\end{countertheorem}

\begin{proof}
The cumulative sums are $(1,2,1)$ and $(-1,0,1)$.  Apply
Theorem~\ref{thm:v14-Skorokhod} with $S_0=0$.
\end{proof}

\begin{definition}[Same-history stock incidence for the V13 scalar]
\label{def:v14-stock-incidence}
A stock incidence on a cofinal V13 subsequence consists of:
\begin{enumerate}[label=\textup{(I\arabic*)},leftmargin=3.3em]
\item coefficients $\lambda_n\ge0$ fixed by one physical calibration before
      the window values are inspected;
\item one nonnegative stock carried chronologically through all selected
      windows and all record boundaries, with finite initial value
      $S_{\rm in}$;
\item nonnegative external replenishment increments
      $r_{n,j}^{\rm ext}$;
\item the exact balance
      \begin{equation}
       \boxed{
       S_{n,j}=S_{n,j-1}-\lambda_nb_{n,j}
                    +r_{n,j}^{\rm ext},}
       \label{eq:v14-physical-stock-balance}
      \end{equation}
      with the terminal stock of one record equal to the initial stock of the
      next record.
\end{enumerate}
Every boundary reset, imported packet, represented return, or external inflow
which changes the stock belongs in $r_{n,j}^{\rm ext}$.  A different stock
initialized independently on each record does not satisfy this definition.
\end{definition}

\begin{theorem}[Cofinal debit forces calibration collapse or non-summable replenishment]
\label{thm:v14-cofinal-replenishment}
Assume Definition~\ref{def:v14-stock-incidence} through records
$n_0,\ldots,N$, and put
\begin{equation}
 R_N^{\rm ext}:=
 \sum_{n=n_0}^{N}\sum_{j=1}^{N_n}r_{n,j}^{\rm ext}.
 \label{eq:v14-total-external-replenishment}
\end{equation}
Then
\begin{equation}
 \boxed{
 R_N^{\rm ext}
 \ge
 \sum_{n=n_0}^{N}\lambda_nA_n-S_{\rm in}
 \ge
 \frac{h_0}{2}\sum_{n=n_0}^{N}\lambda_n-S_{\rm in}.}
 \label{eq:v14-replenishment-lower}
\end{equation}
Consequently:
\begin{enumerate}[label=\textup{(R\arabic*)},leftmargin=3.5em]
\item if $\sum_n\lambda_n=\infty$, every finite initial stock requires
      $R_N^{\rm ext}\to\infty$;
\item if the total external replenishment is finite, then necessarily
      $\sum_n\lambda_n<\infty$ and in particular $\lambda_n\to0$;
\item if $\lambda_n\ge\lambda_*>0$ cofinally, then
      \begin{equation}
       \boxed{
       R_N^{\rm ext}
       \ge\frac{h_0\lambda_*}{2}(N-n_0+1)-S_{\rm in}.}
       \label{eq:v14-linear-record-replenishment}
      \end{equation}
\end{enumerate}
Thus a noncollapsed physical price and a finite replenishment budget cannot
coexist on infinitely many selected records.
\end{theorem}

\begin{proof}
Sum \eqref{eq:v14-physical-stock-balance} over all windows and records.  The
stock terms telescope:
\[
 S_{N,N_N}
 =S_{\rm in}-\sum_{n=n_0}^{N}\lambda_nA_n+R_N^{\rm ext}.
\]
The terminal stock is nonnegative, which proves the first inequality in
\eqref{eq:v14-replenishment-lower}.  The second uses
\eqref{eq:v14-oriented-window-law}.  The three consequences are immediate.
Theorem~\ref{thm:v14-Skorokhod} identifies the pointwise minimum
replenishment realizing this endpoint lower bound together with all
intermediate nonnegativity constraints.
\end{proof}

\begin{corollary}[Natural record-amplitude calibration has exponential replenishment]
\label{cor:v14-record-amplitude-replenishment}
The V10 scaling identity \eqref{eq:v10-scaled-scalar} fixes the natural
physical amplitude of the selected interior scalar as
\begin{equation}
 \boxed{
 \mathcal H_n^{\rm phys}:=\mathcal R_nA_n.}
 \label{eq:v14-physical-selected-scalar}
\end{equation}
If this quantity is a debit of one common finite stock, then its external
replenishment satisfies
\begin{equation}
 \boxed{
 R_N^{\rm ext}
 \ge\frac{h_0}{2}\sum_{n=n_0}^{N}\mathcal R_n-S_{\rm in}
 =\frac{h_0\mathcal R_*}{2}
   \bigl(2^{N+1}-2^{n_0}\bigr)-S_{\rm in}.}
 \label{eq:v14-exponential-replenishment}
\end{equation}
Thus the natural same-history scaling returns a non-summable, indeed
record-exponential, replenishment survivor.  It does not produce a finite
stock contradiction.
\end{corollary}

\begin{proof}
Use $\lambda_n=\mathcal R_n=2^n\mathcal R_*$ in
Theorem~\ref{thm:v14-cofinal-replenishment} and sum the geometric series.
\end{proof}

\begin{firewall}[A calibration is part of the physical incidence]
\label{fire:v14-calibration-not-fit}
The coefficients $\lambda_n$ cannot be chosen after the record scalar is
observed merely to make the series summable.  They must be supplied by the
same physical stock and remain compatible under the declared cutoff and
record transports.  Conversely, the normalized number $A_n$ cannot be summed
as a stock decrement without such a calibration.  The two failures are,
respectively, calibration collapse and missing action-to-stock incidence.
\end{firewall}

% =============================================================================
\section{Audit of the existing Navier--Stokes stock candidates}
\label{sec:v14-stock-audit}
% =============================================================================

The preceding theorem is conditional only on an actual balance law.  We now
test the stock candidates already present in the programme rather than
postulating a new one.

\subsection{Finite native source length forces a summable calibration}

Retain the V5 native source masses
$\widehat m_n$ of \eqref{eq:v5-two-source-lengths}.  Their selected intervals
are disjoint and
\begin{equation}
 \sum_n\widehat m_n\le\mathfrak M_c<\infty.
 \label{eq:v14-native-mass-sum}
\end{equation}

\begin{theorem}[Native source payment can only occur through calibration collapse]
\label{thm:v14-native-source-calibration}
Suppose a proposed action-to-stock incidence satisfies, for a fixed
$C_{\rm nat}<\infty$,
\begin{equation}
 \boxed{
 \lambda_nA_n\le C_{\rm nat}\widehat m_n}
 \label{eq:v14-native-payment-bound}
\end{equation}
cofinally.  Then
\begin{equation}
 \boxed{
 \lambda_n\le\frac{2C_{\rm nat}}{h_0}\widehat m_n,
 \qquad
 \sum_n\lambda_n
 \le\frac{2C_{\rm nat}}{h_0}\mathfrak M_c<\infty.}
 \label{eq:v14-native-calibration-sum}
\end{equation}
In particular, the finite native source length cannot pay the V13 scalar with
a uniform positive calibration, and it cannot pay the natural record-amplitude
quantity \eqref{eq:v14-physical-selected-scalar}.
\end{theorem}

\begin{proof}
Use $A_n\ge h_0/2$ in
\eqref{eq:v14-native-payment-bound} and sum, using
\eqref{eq:v14-native-mass-sum}.
\end{proof}

\begin{assessment}[Relation to the inherited source-action theorem]
\label{ass:v14-native-source-action}
Version~V5 already proved the sharp Riesz lower bound
$(H_n^{\rm ann})^2/\widehat m_n$ and the scale debit
$\sum_n\mathcal R_n/\Gamma_n<\infty$.  Theorem~\ref{thm:v14-native-source-calibration}
does not repeat that estimate.  It answers a different question: whether the
finite native length can serve as a common monotone stock for the normalized
extensive occupation.  It can do so only after the stock price itself becomes
summable.
\end{assessment}

\subsection{Kinetic energy has no nonlinear source debit}

\begin{proposition}[Exact kinetic-energy nonincidence]
\label{prop:v14-kinetic-nonincidence}
For the fixed smooth solution,
\begin{equation}
 \boxed{
 \langle c(t),e(t)\rangle
 =\langle\BNS(u(t)),u(t)\rangle=0,}
 \label{eq:v14-energy-neutral-source}
\end{equation}
and
\begin{equation}
 \boxed{
 \frac12\|u(t_2)\|_2^2-\frac12\|u(t_1)\|_2^2
 =-\nu\int_{t_1}^{t_2}\|\nabla u(t)\|_2^2\,\dd t.}
 \label{eq:v14-kinetic-stock-balance}
\end{equation}
Thus the nonlinear critical source makes no direct debit in the finite
kinetic-energy stock.  Assigning the V13 source-paid scalar to that stock
requires an additional mixed incidence or replenishment term; it is not an
identity of Navier--Stokes.
\end{proposition}

\begin{proof}
Spectral calculus gives
$\langle c,e\rangle=\langle A^{-1/4}\BNS(u),A^{1/4}u\rangle
=\langle\BNS(u),u\rangle$.  Incompressibility and periodic integration by
parts make the last expression zero.  Testing Navier--Stokes with $u$ and
integrating proves \eqref{eq:v14-kinetic-stock-balance}.
\end{proof}

\subsection{The caloric critical stock is an unbounded replenished inventory}

For $b<d<T$ and $t\le d$, define
\begin{equation}
 \Phi_d(t):=\frac12
 \left\|e^{-\nu(d-t)A}e(t)\right\|_2^2,
 \qquad
 \mathcal J_d(t):=
 -\left\langle c(t),
 \Lambda e^{-2\nu(d-t)A}e(t)\right\rangle.
 \label{eq:v14-caloric-stock}
\end{equation}

\begin{theorem}[Exact record rectangle for critical-energy replenishment]
\label{thm:v14-caloric-rectangle}
For every $b<d<T$,
\begin{equation}
 \boxed{
 \frac{\dd}{\dd t}\Phi_d(t)=\mathcal J_d(t)}
 \label{eq:v14-caloric-derivative}
\end{equation}
for $t\le d$, and
\begin{equation}
 \boxed{
 \Phi_b(b)-\Phi_d(b)
 =\nu\int_0^{d-b}
 \left\|A^{1/2}e^{-\nu sA}e(b)\right\|_2^2\,\dd s.}
 \label{eq:v14-endpoint-carry-loss}
\end{equation}
Consequently,
\begin{equation}
 \boxed{
 \frac12\bigl(\|e(d)\|_2^2-\|e(b)\|_2^2\bigr)
 =\int_b^d\mathcal J_d(t)\,\dd t
  -\nu\int_0^{d-b}
   \left\|A^{1/2}e^{-\nu sA}e(b)\right\|_2^2\,\dd s.}
 \label{eq:v14-caloric-rectangle}
\end{equation}
At consecutive doubled records,
\begin{equation}
 \boxed{
 \int_{\beta_n}^{\beta_{n+1}}
 \mathcal J_{\beta_{n+1}}(t)\,\dd t
 =\frac32\mathcal R_n^2+\mathcal L_n^{\rm cal},}
 \label{eq:v14-doubled-contact-replenishment}
\end{equation}
where
\begin{equation}
 \mathcal L_n^{\rm cal}:=
 \nu\int_0^{\Delta_n}
 \left\|A^{1/2}e^{-\nu sA}e(\beta_n)\right\|_2^2\,\dd s
 \ge0.
 \label{eq:v14-caloric-loss}
\end{equation}
Thus the nearest exact critical stock is not finite: its record values are
$\frac12\mathcal R_n^2$, and the nonlinear contact is the replenishment which
builds that unbounded inventory while also replacing the caloric carry loss.
\end{theorem}

\begin{proof}
Differentiate
$w(t)=e^{-\nu(d-t)A}e(t)$.  Since
$\partial_te+\nu Ae=-\Lambda c$,
\[
 w'(t)
 =\nu Ae^{-\nu(d-t)A}e(t)
  +e^{-\nu(d-t)A}(-\nu Ae(t)-\Lambda c(t))
 =-e^{-\nu(d-t)A}\Lambda c(t).
\]
Self-adjointness and commutation of the spectral multipliers give
\eqref{eq:v14-caloric-derivative}.  For fixed $e(b)$, differentiation in
$s$ gives
\[
 \frac{\dd}{\dd s}\frac12
 \|e^{-\nu sA}e(b)\|_2^2
 =-\nu\|A^{1/2}e^{-\nu sA}e(b)\|_2^2.
\]
Integrate from zero to $d-b$ to obtain
\eqref{eq:v14-endpoint-carry-loss}.  Now write
\[
 \Phi_d(d)-\Phi_b(b)
 =\bigl(\Phi_d(d)-\Phi_d(b)\bigr)
  +\bigl(\Phi_d(b)-\Phi_b(b)\bigr)
\]
and use the preceding two identities.  At a doubled record,
$\|e(\beta_{n+1})\|_2=2\mathcal R_n$ and
$\|e(\beta_n)\|_2=\mathcal R_n$, proving
\eqref{eq:v14-doubled-contact-replenishment}.
\end{proof}

\begin{definition}[Full-contact incidence residual]
\label{def:v14-full-contact-residual}
For any predeclared calibration $\lambda_n$ from the oriented V13 scalar to
the complete caloric-contact units, define
\begin{equation}
 \boxed{
 \Delta_n^{\rm full}
 :=\int_{\beta_n}^{\beta_{n+1}}
      \mathcal J_{\beta_{n+1}}(t)\,\dd t
   -\lambda_nA_n.}
 \label{eq:v14-full-contact-residual}
\end{equation}
This residual contains the complete signed complement: resolved contact,
unselected Reynolds output, pressure-compatible followers, represented
returns, endpoint promotion, and every other component not retained in the
V13 held-out scalar.  It must be assembled before a positive payment is
assigned to either term.
\end{definition}

\begin{firewall}[The full contact is not the selected annular scalar]
\label{fire:v14-contact-not-selected}
Theorem~\ref{thm:v14-caloric-rectangle} identifies an actual replenished
critical inventory, but it does not identify the V13 held-out annular scalar
with the complete caloric contact.  That identification is exactly the
same-history component incidence measured by
$\Delta_n^{\rm full}$.  A later theorem about the full contact cannot erase a
nonzero residual at this earlier arrow.
\end{firewall}

\subsection{Linear fast-time action is a dilation identity, not a stock law}

Put
\begin{equation}
 \mathfrak Q_n^{\rm raw}:=
 \int_{\alpha_n\supp\eta}
 \|\widehat G_n(\tau)\|_2^2\,\dd\tau.
 \label{eq:v14-raw-fast-action}
\end{equation}

\begin{proposition}[Divergent action-to-scalar ratio]
\label{prop:v14-action-scalar-ratio}
On the V13 bounded-action branch,
\begin{equation}
 \boxed{
 c_A\alpha_n\le\mathfrak Q_n^{\rm raw}\le L_*\alpha_n,}
 \label{eq:v14-raw-action-linear}
\end{equation}
whereas
\begin{equation}
 \frac{h_0}{2}
 \le\|\Sigma_n\|_{\rm TV}
 \le\mathfrak V_*.
 \label{eq:v14-scalar-variation-window}
\end{equation}
Hence
\begin{equation}
 \boxed{
 \frac{\mathfrak Q_n^{\rm raw}}{\|\Sigma_n\|_{\rm TV}}
 \ge\frac{c_A}{\mathfrak V_*}\alpha_n
 \longrightarrow\infty.}
 \label{eq:v14-action-resultant-leverage}
\end{equation}
The raw fast-time action therefore cannot itself be a finite scalar payment.
An action-to-stock incidence must introduce a separate physical conversion
and balance law.
\end{proposition}

\begin{proof}
Equation~\eqref{eq:v14-raw-action-linear} is
\eqref{eq:v13-linear-source-action}.  The lower variation bound is the
absolute value of the total mass, and the upper bound is
Theorem~\ref{thm:v14-bounded-turnover}.  Divide.
\end{proof}

\begin{countertheorem}[Linear raw action can be purely kinematic]
\label{cth:v14-kinematic-action}
Let $J\Subset(-\infty,0)$ be fixed, let $g\in L^2(J)$ have nonzero integral,
let $z$ be a unit vector in a Hilbert space, and let $\alpha_n\to\infty$.
Set
\begin{equation}
 G_n(s)=g(s)z,
 \qquad
 q_n(s)=g(s),
 \qquad
 d\Sigma_n(\tau)=\alpha_n^{-1}g(\tau/\alpha_n)\,\dd\tau.
 \label{eq:v14-kinematic-model}
\end{equation}
Then
\begin{equation}
 \|\Sigma_n\|_{\rm TV}=\|g\|_1,
 \qquad
 \Sigma_n(\mathbb R)=\int_Jg,
 \qquad
 \int_{\alpha_nJ}\|G_n(\tau/\alpha_n)\|^2\,\dd\tau
 =\alpha_n\|g\|_2^2.
 \label{eq:v14-kinematic-model-scales}
\end{equation}
Thus the same bounded scalar history has linearly growing raw action under
nothing more than the fast-time dilation.
\end{countertheorem}

\begin{proof}
All three identities are changes of variables.
\end{proof}

\begin{firewall}[The kinematic model is not a Navier--Stokes solution]
\label{fire:v14-kinematic-model}
Countertheorem~\ref{cth:v14-kinematic-action} is a sharpness statement for the
inference from fast-time action to finite-stock payment.  It does not satisfy
the Navier--Stokes field equation or reproduce the V12 heat-covariance
profile.  The actual solution supplies those additional structures, but no
existing identity turns their positive square action into a decrement of a
finite physical stock.
\end{firewall}

% =============================================================================
\section{The exact terminal alternative after the stock audit}
\label{sec:v14-alternative}
% =============================================================================

\begin{theorem}[Version V14 payment, calibration, and replenishment theorem]
\label{thm:v14-integrated}
Preserving every theorem and limitation of Versions~V1--V13, Version~V14
proves the following on the V13 extensive bounded-action branch.
\begin{enumerate}[label=\textup{(V14.\arabic*)},leftmargin=6em]
\item The linearly growing window participation has uniformly bounded scalar
      total variation.  The positive Jordan part occupies order $\alpha_n$
      windows, but its total mass and every partial-sum stock are uniformly
      bounded on one record.
\item Jordan decomposition after complete signed assembly is the unique
      minimum-turnover positive split.  Every earlier positive splitting adds
      the same duplicated amount to payment and counterflow.
\item The minimum external replenishment for any chronological signed debit
      history is the one-sided reflection
      \eqref{eq:v14-Skorokhod-formula}.  Scalar, variation, and action
      marginals do not determine it without chronology.
\item For any same-history physical stock calibration $\lambda_n$, a finite
      initial stock obeys the cofinal lower bound
      \eqref{eq:v14-replenishment-lower}.  Thus a nonsummable calibration
      forces nonsummable replenishment, while finite replenishment forces
      calibration collapse.
\item Under the record-amplitude calibration dictated by the V10 physical
      scalar, the required replenishment has the exponential lower bound
      \eqref{eq:v14-exponential-replenishment}.
\item Any proposed domination by the finite native source length forces
      $\sum_n\lambda_n<\infty$.  The native source can pay the extensive
      normalized scalar only through a summably vanishing stock conversion.
\item Kinetic energy has no nonlinear source debit.  The caloric critical
      energy has the exact rectangle
      \eqref{eq:v14-caloric-rectangle}, but is an unbounded inventory whose
      record growth is supplied by the nonlinear replenishment
      \eqref{eq:v14-doubled-contact-replenishment}.
\item The raw fast-time source action grows linearly while scalar variation
      remains bounded.  This scaling can arise from time dilation alone and
      is not a finite-stock payment theorem.
\end{enumerate}
Consequently the bounded-action branch terminates in exactly one of
\begin{equation}
 \boxed{
 \begin{gathered}
 \text{summable physical calibration collapse},\\
 \text{non-summable same-history replenishment},\\
 \text{a nonzero full-contact/component incidence residual},\\
 \text{or an unpaid extensive source-action occupation.}
 \end{gathered}}
 \label{eq:v14-final-alternative}
\end{equation}
No finite-stock contradiction or regularity conclusion follows without an
additional physical action-to-stock law.
\end{theorem}

\begin{proof}
Item~\textup{(V14.1)} is
Theorem~\ref{thm:v14-bounded-turnover} and
Corollary~\ref{cor:v14-one-record-stock}.  Item~\textup{(V14.2)} is
Theorem~\ref{thm:v14-Jordan-minimality}.  Item~\textup{(V14.3)} is
Theorem~\ref{thm:v14-Skorokhod} and
Countertheorem~\ref{cth:v14-order-matters}.  Items
\textup{(V14.4)--(V14.5)} are
Theorem~\ref{thm:v14-cofinal-replenishment} and
Corollary~\ref{cor:v14-record-amplitude-replenishment}.  Item
\textup{(V14.6)} is
Theorem~\ref{thm:v14-native-source-calibration}.  Item~\textup{(V14.7)} is
Proposition~\ref{prop:v14-kinetic-nonincidence} and
Theorem~\ref{thm:v14-caloric-rectangle}.  The final item is
Proposition~\ref{prop:v14-action-scalar-ratio} and
Countertheorem~\ref{cth:v14-kinematic-action}.
\end{proof}

\begin{firewall}[This is the payment stop, not a regularity theorem]
\label{fire:v14-no-regularity}
Version~V14 does not prove that every possible physical stock has been
excluded.  It proves that the quantities already present in the NS--A packet
do not supply a finite, noncollapsed, same-history stock balance: native
source length collapses the calibration, kinetic energy shorts the nonlinear
source, critical energy is itself replenished and unbounded, and raw fast
action is kinematic.  A new finite stock may be used only after its physical
occurrence, calibration, and exact balance with the selected scalar are
proved.  No terminal regularity or global three-dimensional Navier--Stokes
claim is made.
\end{firewall}

\begingroup\small
\begin{longtable}{p{0.25\textwidth}p{0.19\textwidth}p{0.48\textwidth}}
\caption{NS--A Version V14 proof-status ledger.}\label{tab:v14-ledger}\\
\toprule Row & Status & Exact conclusion\\\midrule
\endfirsthead
\toprule Row & Status & Exact conclusion\\\midrule
\endhead
V1--V13 prefix & \PROVED{} preserved
 &The complete V13 preterminal source is retained byte-for-byte; no earlier
 theorem is rewritten.\\[0.3em]
Fast-window participation & \INHERITED{} linear
 &Order $\alpha_n$ bounded windows are required for a fixed scalar fraction.\\[0.3em]
Scalar total variation & \PROVED{} uniformly bounded
 &$\|\Sigma_n\|_{\rm TV}\le\mathfrak V_*$; extensive participation is not
 extensive scalar turnover.\\[0.3em]
Positive Jordan windows & \PROVED{} linear count
 &After one fixed orientation, at least $c\alpha_n$ windows have positive
 assembled coefficient, while their total mass remains bounded.\\[0.3em]
One-record stock & \PROVED{} abstract only
 &A stock of size $\mathfrak V_*$ pays every signed window of one record with
 no replenishment; resetting it is not free.\\[0.3em]
Positive split & \PROVED{} unique minimum
 &Every non-Jordan payment/replenishment split duplicates the same turnover
 on both sides.\\[0.3em]
First-exit stock & \PROVED{} exact
 &The one-sided Skorokhod formula is the minimum replenishment compatible with
 nonnegative stock at every window.\\[0.3em]
Chronology from marginals & \OBSTRUCTION{} explicit
 &Equal total, Jordan, and quadratic-action data can require different
 first-exit stocks.\\[0.3em]
Cofinal stock balance & \PROVED{} conditional incidence
 &A common calibrated stock requires replenishment at least
 $(h_0/2)\sum\lambda_n-S_{\rm in}$.\\[0.3em]
Natural record calibration & \PROVED{} divergent replenishment
 &$\lambda_n=\mathcal R_n$ forces an exponential replenishment lower bound.\\[0.3em]
Native source length & \OBSTRUCTION{} calibration collapse
 &Any debit dominated by $\widehat m_n$ has a summable calibration and no
 uniform positive physical price.\\[0.3em]
Kinetic energy stock & \OBSTRUCTION{} source neutral
 &The quadratic nonlinearity is orthogonal to $u$ and contributes no kinetic
 decrement.\\[0.3em]
Caloric critical stock & \PROVED{} exact rectangle
 &The full contact replenishes the unbounded critical inventory and its
 caloric carry loss; it is not a finite stock.\\[0.3em]
Selected/full-contact incidence & \OPEN{} exact residual
 &$\Delta_n^{\rm full}$ is the complete signed complement which must be
 evaluated before the selected scalar is called critical replenishment.\\[0.3em]
Raw fast source action & \OBSTRUCTION{} not payment
 &It grows linearly under the time dilation while scalar variation remains
 bounded.\\[0.3em]
Finite action-to-stock law & \OPEN{} first physical row
 &No existing packet supplies a noncollapsed finite stock and a summable
 replenishment budget on the same history.\\[0.3em]
Three-dimensional regularity & \OPEN
 &Neither replenishment nor unpaid extensive occupation is excluded for every
 fixed datum.\\
\bottomrule
\end{longtable}
\endgroup

% =============================================================================
\section{Exact mandate for Version V15}
\label{sec:v14-next}
% =============================================================================

\begin{programme}[Evaluate the selected scalar inside the complete caloric rectangle]
\label{prog:v14-next}
Preserve Versions~V1--V14 verbatim.  Do not introduce another abstract stock,
window participation law, fast-time compactness space, or positive action
norm.  Version~V14 has closed every inference from those quantities alone.

The sole admissible continuation on the bounded-action branch is the
following upstream same-history calculation.
\begin{enumerate}[label=\textup{(V15.\arabic*)},leftmargin=5.7em]
\item Pull the naturally calibrated selected scalar
      $\mathcal H_n^{\rm phys}=\mathcal R_nA_n$ back to its original physical
      time interval and place it in the complete caloric-contact rectangle
      \eqref{eq:v14-caloric-rectangle}.
\item Assemble, before positivity, the resolved contact, unselected Reynolds
      outputs, pressure-compatible followers, represented returns, endpoint
      promotion, and caloric carry loss.  Evaluate the exact residual
      $\Delta_n^{\rm full}$ of \eqref{eq:v14-full-contact-residual}; do not
      replace it by the difference of marginal bounds.
\item If the complement carries a fixed fraction, return its precise signed
      type as the source/contact incidence obstruction.  If it vanishes in
      the required normalization, identify the selected annular scalar as an
      actual component of the record-to-record critical replenishment.
\item On the latter branch, retain the non-summable critical replenishment as
      the terminal survivor.  It is not a contradiction unless an
      independently occurring finite stock or monotone upper budget is
      supplied.
\item Reopen a bounded fast-time profile only on the separately typed
      source-action/enstrophy-square overdrive branch.  No conclusion from the
      present bounded-action branch may be imported into that normalization.
\end{enumerate}
A successful Version~V15 ends in
\begin{equation}
 \boxed{
 \begin{gathered}
 \text{signed full-contact complement}\\[0.15em]
 \big|\quad
 \text{selected critical-replenishment component}\\[0.15em]
 \big|\quad
 \text{missing finite-stock upper budget}\\[0.15em]
 \big|\quad
 \text{separately normalized action-overdrive profile}.
 \end{gathered}}
 \label{eq:v14-v15-endpoint}
\end{equation}
No downstream Liouville or terminal-trace theorem is applied before this
component incidence is settled.
\end{programme}

\begin{assessment}[Programme position after V14]
\label{ass:v14-position}
The word ``extensive'' has now been disambiguated.  The fast horizon, positive
source action, participation number, and number of positive assembled windows
are extensive.  The signed scalar variation and the stock needed inside one
record are not.  Cofinal obstruction begins only when one physically
calibrated stock is carried through the record sequence.  At that point the
exact reflection law forces either summable calibration collapse or
non-summable replenishment.  The only upstream PDE object already known to
carry the complete nonlinear contact is the caloric critical-energy
rectangle, so the next calculation is its exact signed incidence with the
held-out annular scalar.
\end{assessment}

\paragraph{Version V14 history.}
Preserved the complete extensive fast-time law of Version~V13 and evaluated
its payment interpretation.  Proved uniform scalar total variation, linear
positive-window count, and a bounded one-record stock; established the unique
minimum-turnover Jordan split and the exact one-sided stock-reflection
formula.  Introduced only the physical calibration needed to state a common
cofinal stock balance and proved the calibration-collapse versus
non-summable-replenishment theorem, including the exponential lower bound at
the natural record amplitude.  Audited the existing NS candidates: finite
native source length, kinetic energy, caloric critical energy, and raw
fast-time source action.  The first forces summable calibration, the second
shorts the nonlinearity, the third is an unbounded replenished inventory, and
the fourth is kinematically extensive under dilation.  Isolated the complete
selected/full-contact incidence residual as the next and only upstream
calculation on this branch.

\begin{assessment}[Append-only integrity]
\label{ass:v14-integrity}
The complete Version~V13 source preceding its sole active final document
terminator is preserved verbatim.  Its preterminal SHA--256 digest is
\begin{center}\scriptsize\ttfamily
2143ecdf5ebbd2e3aa2e596a9534e011bc903be2fb16eeb0da723d942bc0a760
\end{center}
The present material begins at Section~\ref{sec:v14-turnover}; the final
command below is again unique.
\end{assessment}

\addtocontents{toc}{\protect\endgroup}
% =============================================================================
% END VERSION V14 MATERIAL -- append later versions before final terminator
% =============================================================================

\clearpage
% =============================================================================
% VERSION V15: NEW MATERIAL.  The complete V1--V14 source above is preserved
% verbatim; only its sole active \end{document} line was removed before append.
% =============================================================================
\hypersetup{pdftitle={NS-A Terminal Source Concentration Programme: Cumulative V15},pdfauthor={A. Pelissier}}
\makeatletter
\addtocontents{toc}{\protect\begingroup\protect\makeatletter}
\addtocontents{toc}{\protect\def\protect\l@section{\protect\@dottedtocline{1}{0em}{4.7em}}}
\addtocontents{toc}{\protect\def\protect\l@subsection{\protect\@dottedtocline{2}{1.5em}{5.3em}}}
\addtocontents{toc}{\protect\makeatother}
\makeatother
\renewcommand{\thesection}{V15.\arabic{section}}
\renewcommand{\thesubsection}{V15.\arabic{section}.\arabic{subsection}}
\renewcommand{\theequation}{V15.\arabic{equation}}
\setcounter{section}{0}
\setcounter{equation}{0}
\renewcommand{\theHsection}{V15.\arabic{section}}
\renewcommand{\theHsubsection}{V15.\arabic{section}.\arabic{subsection}}
\renewcommand{\theHtheorem}{V15.\arabic{section}.\arabic{theorem}}
\renewcommand{\theHequation}{V15.\arabic{equation}}

\begin{center}
{\LARGE\bfseries NS--A: Version V15}\\[0.5em]
{\large Terminal-Writer Calibration, the Exact Caloric-Contact Ledger,\\
and the Forced Full-Complement Alternative}\\[0.5em]
{\normalsize 3 September 2026}
\end{center}
\addcontentsline{toc}{part}{Version V15: terminal-writer calibration and the exact caloric-contact ledger}

\begin{abstract}
Version~V14 asked that the naturally scaled V13 scalar
$\mathcal H_n^{\rm phys}=\mathcal R_nA_n$ be inserted into the complete
record-to-record caloric-contact rectangle.  The first result of the present
continuation is an upstream typing correction.  The quantity
$\mathcal R_nA_n$ is the physical size of one fixed-writer response pairing;
the caloric contact is a critical-energy secant.  Since $A_n$ is uniformly
bounded while the complete contact is at least
$\tfrac32\mathcal R_n^2$, the direct V14 residual obtained by subtracting
$\mathcal R_nA_n$ contains asymptotically the entire caloric contact.  Thus
record amplitude alone is not the missing component incidence.  The audit also
corrects one inherited normalization shorthand: on the selected V5 heavy annulus
one has $\sqrt\eta\le\|\varpi_n\|_2\le1$, rather than an automatic unit norm.
Consequently the V10 rescaled writer has the genuine fixed window
$\sqrt\eta/2\le\|Z_n\|_2\le2$.  Every later argument which used only the
upper edge remains unchanged.

The correct conversion is a same-history mixed writer Gram.  At terminal
age $r=b_n-t$, the full caloric writer is
$Z_{n,t}^{\rm cal}=\Lambda e^{-\nu rA}e(t)$, whereas the selected annular
source is tested against the fixed Riesz writer $W_n=\Lambda\varpi_n$.
After pulling the V13 time, cell, and material-router cutoffs back to the
original variables, project $Z_{n,t}^{\rm cal}$ onto $W_n$ in the actual
weighted history space.  Its coefficient
\[
 \kappa_n=
 \frac{\int\Theta_n Z_{n,t}^{\rm cal}\cdot W_n}
      {\int\Theta_n|W_n|^2}
\]
is the unique least-square terminal-writer calibration.  The physical signed
V13 scalar is $\mathfrak s_n=-\sigma_*\mathcal R_nA_n$, and the only
candidate selected caloric component is therefore
$\mathcal C_n^{\rm sel}=\kappa_n\mathfrak s_n$, not
$\mathcal R_nA_n$ itself.  On the bounded-action branch the normalized
coefficient $\kappa_n/\mathcal R_n$ is uniformly bounded, but neither its
sign nor its lower edge follows from record height or from separate source
and writer norms.

The full contact is then assembled in two exact signed ledgers.  The source
ledger contains the resolved heat--Euler contact, unselected covariance
outputs, material/cell/time escape, the selected mixed-writer component, and
one writer-orthogonal residual.  Gradient and isotropic pressure followers
vanish exactly against the divergence-free caloric writer.  Independently,
the Duhamel endpoint ledger splits the same contact into the part represented
by the literal V5 terminal dictionary and its held-out complement; the
quadratic half-norms are the endpoint-promotion terms.  Subtracting the two
ledgers gives one exact source-to-final-head incidence identity.  The V14
caloric carry loss then enters only when contact is converted to net critical
inventory growth.

Consequently, after a priority subsequence, either a represented endpoint
contact survives, a resolved/unselected/router/writer incidence residual
carries a fixed fraction, or the selected mixed-writer term is asymptotic to
the complete held-out caloric contact.  On the last branch the writer Gram
has record-size coefficient and positive orientation, its projection
efficiency has a fixed floor, the caloric carry loss lies universally in
$[0,\mathcal R_n^2/2]$, and the selected contact is itself non-summable.  This is a
critical-replenishment survivor, not a contradiction: no independently
occurring finite stock or monotone upper budget has been supplied.  The
source-action/enstrophy-square overdrive remains a separately normalized
branch.  No terminal regularity or global three-dimensional Navier--Stokes
conclusion is claimed.
\end{abstract}

% =============================================================================
\section{Scope and the response-amplitude/contact typing correction}
\label{sec:v15-scope}
% =============================================================================

Retain the bounded-action, one-cell, bounded-product branch of
Versions~V10--V14.  In particular,
\begin{equation}
 \frac{h_0}{2}\le A_n\le\mathfrak V_*,
 \qquad
 \mathcal R_n=2^n\mathcal R_*,
 \qquad
 b_n=\beta_{n+1},
 \label{eq:v15-standing-A}
\end{equation}
and the complete caloric contact on the record interval is
\begin{equation}
 \boxed{
 \mathcal C_n^{\rm cal}
 :=\int_{\beta_n}^{b_n}\mathcal J_{b_n}(t)\,\dd t
 =\frac32\mathcal R_n^2+\mathcal L_n^{\rm cal},
 \qquad
 0\le\mathcal L_n^{\rm cal}\le\frac12\mathcal R_n^2,}
 \label{eq:v15-full-contact}
\end{equation}
so in particular
\begin{equation}
 \boxed{
 \frac32\mathcal R_n^2
 \le\mathcal C_n^{\rm cal}
 \le2\mathcal R_n^2.}
 \label{eq:v15-full-contact-window}
\end{equation}
The upper carry bound follows directly from heat contraction:
$\mathcal L_n^{\rm cal}
 =\tfrac12(\|e(\beta_n)\|_2^2-\|S_ne(\beta_n)\|_2^2)$ once $S_n$ is
introduced below.  It is recorded here because it is needed whenever a
relative contact residual is converted to record-amplitude units.
All represented-response, source-ancestry, material-route, cell-incidence,
heat-tail, Reynolds-force, and occupation shorts used to define $A_n$ remain
in force.  No quantity below is formed by reselecting a cell or a writer after
its value has been inspected.

\begin{proposition}[Correction of the annular test-direction normalization]
\label{prop:v15-annular-normalization}
On the selected V5 annular branch
\(
 \|G_n\|_2^2/(H_n^{\rm core})^2\ge\eta
\), the test direction of
\eqref{eq:v5-annular-test-direction} satisfies
\begin{equation}
 \boxed{
 \sqrt\eta\le\|\varpi_n\|_2\le1.}
 \label{eq:v15-varpi-window}
\end{equation}
Under the V10 annular rescaling of
\eqref{eq:v10-scaled-writer}, therefore,
\begin{equation}
 \boxed{
 \sqrt\eta\le\|\varphi_n\|_2\le1,
 \qquad
 \frac{\sqrt\eta}{2}\le\|Z_n\|_2\le2.}
 \label{eq:v15-rescaled-writer-window}
\end{equation}
Thus the unit-norm sentence immediately following
\eqref{eq:v10-scaled-writer} is replaced by
\eqref{eq:v15-rescaled-writer-window}.  All subsequent estimates which used
only \(\|Z_n\|_2\le2\) are unaffected.
\end{proposition}

\begin{proof}
Because
\(
 G_n=\Psi_nF_n^{\rm core}
\),
\(
 H_n^{\rm ann}=\|G_n\|_2
\), and
\(
 \varpi_n=\Psi_n^2F_n^{\rm core}/H_n^{\rm ann}
\), self-adjointness of \(\Psi_n\) gives
\[
 \langle F_n^{\rm core},\varpi_n\rangle
 =\frac{\|\Psi_nF_n^{\rm core}\|_2^2}{H_n^{\rm ann}}
 =H_n^{\rm ann}.
\]
Cauchy--Schwarz and the heavy-annulus inequality yield
\[
 \|\varpi_n\|_2
 \ge\frac{H_n^{\rm ann}}{H_n^{\rm core}}
 \ge\sqrt\eta.
\]
Since \(0\le\Psi_n\le I\),
\[
 \|\Psi_n^2F_n^{\rm core}\|_2
 \le\|\Psi_nF_n^{\rm core}\|_2=H_n^{\rm ann},
\]
which proves the upper bound.  The spatial change of variables in
\eqref{eq:v10-scaled-writer} preserves the \(L^2\) norm, so
\(\|\varphi_n\|_2=\|\varpi_n\|_2\).  Its Fourier support lies in
\(\{1/2\le|\xi|\le2\}\); spectral calculus therefore gives
\[
 \frac12\|\varphi_n\|_2
 \le\|\Lambda_y\varphi_n\|_2
 \le2\|\varphi_n\|_2.
\]
This proves \eqref{eq:v15-rescaled-writer-window}.
\end{proof}

\begin{theorem}[The direct V14 natural-amplitude residual is asymptotically the full contact]
\label{thm:v15-direct-residual-full}
Let
\begin{equation}
 \Delta_{n,\rm nat}^{\rm full}
 :=\mathcal C_n^{\rm cal}-\mathcal R_nA_n
 \label{eq:v15-natural-residual}
\end{equation}
be the residual in \eqref{eq:v14-full-contact-residual} under the V14
record-amplitude choice $\lambda_n=\mathcal R_n$.  Then
\begin{equation}
 \boxed{
 0\le
 \frac{\mathcal R_nA_n}{\mathcal C_n^{\rm cal}}
 \le\frac{2\mathfrak V_*}{3\mathcal R_n}
 \longrightarrow0,}
 \label{eq:v15-natural-negligible}
\end{equation}
and hence
\begin{equation}
 \boxed{
 \frac{\Delta_{n,\rm nat}^{\rm full}}
      {\mathcal C_n^{\rm cal}}
 \longrightarrow1.}
 \label{eq:v15-natural-residual-ratio}
\end{equation}
If the physical sign is restored as in
\Cref{prop:v15-physical-pullback} below, the same conclusion holds with
$\mathcal R_nA_n$ replaced by that signed response scalar.  Therefore a
fixed-writer response amplitude cannot be inserted directly as a
critical-energy contact component.
\end{theorem}

\begin{proof}
Use $A_n\le\mathfrak V_*$ and
$\mathcal C_n^{\rm cal}\ge\tfrac32\mathcal R_n^2$ in the quotient.  This
proves \eqref{eq:v15-natural-negligible}; subtracting it from one proves
\eqref{eq:v15-natural-residual-ratio}.  A signed scalar with the same absolute
value obeys the identical absolute ratio bound.
\end{proof}

\begin{assessment}[What must be supplied before the subtraction becomes typed]
\label{ass:v15-missing-degree}
The V13 scalar is linear in one source and one normalized Riesz writer.  The
caloric rectangle is quadratic in the critical state.  A second
same-history coefficient is therefore required: the component of the actual
caloric writer in the selected Riesz-writer direction.  This is not a new
normalization convention.  It is a mixed physical Gram which can vanish,
change sign, or remain of record size while all separate norms are fixed.
\end{assessment}

\begin{firewall}[Correction of interpretation, not deletion of V14]
\label{fire:v15-v14-correction}
The algebraic definition \eqref{eq:v14-full-contact-residual} remains valid
for any scalar coefficient.  What fails for the choice
$\lambda_n=\mathcal R_n$ is its interpretation as a decomposition into two
quantities of the same physical contact type.  Theorem~\ref{thm:v15-direct-residual-full}
settles the V15 mandate at that direct normalization: the full-complement
branch is forced.  The remainder of V15 constructs the unique least-square
coefficient which does give a contact-typed component.
\end{firewall}

% =============================================================================
\section{The complete caloric contact is a Duhamel secant}
\label{sec:v15-Duhamel-secant}
% =============================================================================

Put
\begin{equation}
 S_n:=e^{-\nu(b_n-\beta_n)A},
 \qquad
 \mathcal F_n:=-\int_{\beta_n}^{b_n}
  \Lambda e^{-\nu(b_n-t)A}c(t)\,\dd t.
 \label{eq:v15-full-Duhamel-response}
\end{equation}
Duhamel's formula gives
\begin{equation}
 e(b_n)=S_ne(\beta_n)+\mathcal F_n.
 \label{eq:v15-full-Duhamel}
\end{equation}

\begin{theorem}[Exact endpoint-secant representations]
\label{thm:v15-Duhamel-secant}
The complete contact has the two equivalent representations
\begin{equation}
 \boxed{
 \begin{aligned}
 \mathcal C_n^{\rm cal}
 &=\frac12\Bigl(
   \|e(b_n)\|_2^2-\|S_ne(\beta_n)\|_2^2\Bigr)\\
 &=\langle S_ne(\beta_n),\mathcal F_n\rangle
   +\frac12\|\mathcal F_n\|_2^2\\
 &=\langle e(b_n),\mathcal F_n\rangle
   -\frac12\|\mathcal F_n\|_2^2.
 \end{aligned}}
 \label{eq:v15-Duhamel-secant}
\end{equation}
Moreover,
\begin{equation}
 \boxed{
 \frac12\Bigl(
 \|e(b_n)\|_2^2-\|e(\beta_n)\|_2^2\Bigr)
 =\mathcal C_n^{\rm cal}-\mathcal L_n^{\rm cal}
 =\frac32\mathcal R_n^2.}
 \label{eq:v15-inventory-secant}
\end{equation}
Thus the linear term in \eqref{eq:v15-Duhamel-secant} and the positive
half-norm are assembled before any contact component is declared; the latter
is the exact endpoint-promotion term.
\end{theorem}

\begin{proof}
The first line is the integrated derivative identity
\eqref{eq:v14-caloric-derivative}.  Insert
\eqref{eq:v15-full-Duhamel} and expand the squared norm to obtain the second
line.  Replacing $S_ne(\beta_n)$ by $e(b_n)-\mathcal F_n$ gives the third.
Equation~\eqref{eq:v14-endpoint-carry-loss} says
\[
 \frac12\|e(\beta_n)\|_2^2
 -\frac12\|S_ne(\beta_n)\|_2^2
 =\mathcal L_n^{\rm cal}.
\]
Subtract this identity from the first line of
\eqref{eq:v15-Duhamel-secant} and use the doubled record values.
\end{proof}

Let $P_n^{\rm rep}$ be the literal finite same-endpoint represented
dictionary projection of Version~V5 and put
\begin{equation}
 Q_n^{\rm rep}:=I-P_n^{\rm rep},
 \qquad
 \mathcal F_n^{\rm rep}:=P_n^{\rm rep}\mathcal F_n,
 \qquad
 \mathcal F_n^{\rm ho}:=Q_n^{\rm rep}\mathcal F_n.
 \label{eq:v15-full-response-heads}
\end{equation}
No commutation of $P_n^{\rm rep}$ with $A$ is assumed.

\begin{proposition}[Exact represented/held-out endpoint ledger]
\label{prop:v15-endpoint-ledger}
Define
\begin{align}
 \mathcal C_n^{\rm rep}
 &:=\langle S_ne(\beta_n),\mathcal F_n^{\rm rep}\rangle
   +\frac12\|\mathcal F_n^{\rm rep}\|_2^2,
 \label{eq:v15-represented-contact}\\
 \mathcal C_n^{\rm ho}
 &:=\langle S_ne(\beta_n),\mathcal F_n^{\rm ho}\rangle
   +\frac12\|\mathcal F_n^{\rm ho}\|_2^2.
 \label{eq:v15-heldout-contact}
\end{align}
Then
\begin{equation}
 \boxed{
 \mathcal C_n^{\rm cal}
 =\mathcal C_n^{\rm rep}+\mathcal C_n^{\rm ho}.}
 \label{eq:v15-endpoint-ledger}
\end{equation}
Equivalently, each term can be written with the terminal state:
\begin{equation}
 \mathcal C_n^{\bullet}
 =\langle e(b_n),\mathcal F_n^{\bullet}\rangle
  -\frac12\|\mathcal F_n^{\bullet}\|_2^2,
 \qquad \bullet\in\{\mathrm{rep},\mathrm{ho}\}.
 \label{eq:v15-terminal-endpoint-ledger}
\end{equation}
The represented response and its endpoint promotion are therefore retained
as one signed term rather than subtracted after positivity.
\end{proposition}

\begin{proof}
The response components are orthogonal, so
$\|\mathcal F_n\|_2^2=
 \|\mathcal F_n^{\rm rep}\|_2^2+
 \|\mathcal F_n^{\rm ho}\|_2^2$; the linear term splits by linearity.  This
proves \eqref{eq:v15-endpoint-ledger}.  Since
$e(b_n)=S_ne(\beta_n)+\mathcal F_n$ and the two response components are
orthogonal,
$\langle e(b_n),\mathcal F_n^\bullet\rangle
 =\langle S_ne(\beta_n),\mathcal F_n^\bullet\rangle
  +\|\mathcal F_n^\bullet\|_2^2$, proving the equivalent formula.
\end{proof}

% =============================================================================
\section{The full forced-Euler source ledger}
\label{sec:v15-source-ledger}
% =============================================================================

For $t\in(\beta_n,b_n)$ put
\begin{equation}
 r_n(t):=b_n-t,
 \qquad
 v_{n,t}:=e^{-\nu r_n(t)A}u(t),
 \qquad
 Z_{n,t}^{\rm cal}:=\Lambda e^{-\nu r_n(t)A}e(t),
 \label{eq:v15-caloric-field-writer}
\end{equation}
and define
\begin{equation}
 g_{n,t}^{\rm cov}
 :=e^{-\nu r_n(t)A}c[u(t)]-c[v_{n,t}],
 \qquad
 c[w]:=A^{-1/4}\BNS(w).
 \label{eq:v15-full-cov-source}
\end{equation}
This is the same heat-covariance source as in
\eqref{eq:v11-source-decomposition}, now retained on the whole record
interval.

\begin{theorem}[Exact resolved/covariance decomposition of the complete contact]
\label{thm:v15-full-forced-Euler}
For every $t\in(\beta_n,b_n)$,
\begin{equation}
 \boxed{
 \mathcal J_{b_n}(t)
 =-\langle c[v_{n,t}],Z_{n,t}^{\rm cal}\rangle
  -\langle g_{n,t}^{\rm cov},Z_{n,t}^{\rm cal}\rangle.}
 \label{eq:v15-pointwise-forced-Euler}
\end{equation}
Consequently, with
\begin{align}
 \mathcal C_n^{\rm res}
 &:=-\int_{\beta_n}^{b_n}
   \langle c[v_{n,t}],Z_{n,t}^{\rm cal}\rangle\,\dd t,
 \label{eq:v15-resolved-contact}\\
 \mathcal C_n^{\rm cov}
 &:=-\int_{\beta_n}^{b_n}
   \langle g_{n,t}^{\rm cov},Z_{n,t}^{\rm cal}\rangle\,\dd t,
 \label{eq:v15-covariance-contact}
\end{align}
one has
\begin{equation}
 \boxed{
 \mathcal C_n^{\rm cal}
 =\mathcal C_n^{\rm res}+\mathcal C_n^{\rm cov}.}
 \label{eq:v15-res-cov-ledger}
\end{equation}
The two terms are signed and are assembled before any positive Reynolds work
is extracted.
\end{theorem}

\begin{proof}
Self-adjointness and commutation of the heat multiplier give
\[
 \mathcal J_{b_n}(t)
 =-\langle e^{-\nu r_nA}c[u(t)],
          \Lambda e^{-\nu r_nA}e(t)\rangle.
\]
Insert
$e^{-\nu r_nA}c[u]=c[v_{n,t}]+g_{n,t}^{\rm cov}$ from
\eqref{eq:v15-full-cov-source}.  This proves the pointwise identity;
integration proves the ledger.
\end{proof}

\begin{proposition}[Pressure-compatible followers have zero caloric contact]
\label{prop:v15-pressure-follower-zero}
Let $a(t,x)$ be any scalar distribution for which the following pairing is
defined.  Then
\begin{equation}
 \boxed{
 \langle\nabla a(t),Z_{n,t}^{\rm cal}\rangle=0.}
 \label{eq:v15-gradient-zero}
\end{equation}
Likewise, adding an isotropic tensor $aI$ to the heat covariance changes
neither $g_{n,t}^{\rm cov}$ nor
\eqref{eq:v15-pointwise-forced-Euler}.  Thus a pressure or isotropic follower
is not an additional term of the caloric-contact complement.
\end{proposition}

\begin{proof}
The critical state $e(t)$ is divergence free and every spectral multiplier
in $Z_{n,t}^{\rm cal}$ preserves divergence freeness.  Periodic integration
by parts gives \eqref{eq:v15-gradient-zero}.  The projected divergence of
$aI$ is $P_\sigma\nabla a=0$.
\end{proof}

% =============================================================================
\section{Pullback of the V13 scalar and restoration of its physical sign}
\label{sec:v15-pullback}
% =============================================================================

Choose the V10 spatial cutoff with $0\le\chi\le1$, and retain the V12 time
cutoff $0\le\eta\le1$ with
\begin{equation}
 \supp\eta\subset[-\tau_+,-\tau_-],
 \qquad 0<\tau_-<\tau_+<\infty.
 \label{eq:v15-interior-ages}
\end{equation}
In the original variables put
\begin{equation}
 \eta_n(t):=\eta\bigl(\nu\Gamma_n^2(t-b_n)\bigr),
 \qquad
 \chi_n(x):=\chi\bigl(\Gamma_n(x-y_n)\bigr),
 \label{eq:v15-physical-cutoffs}
\end{equation}
\begin{equation}
 \Theta_n(t,x):=\eta_n(t)\chi_n(x)\omega_n(t,x),
 \qquad 0\le\Theta_n\le1,
 \label{eq:v15-complete-cutoff}
\end{equation}
where $\omega_n$ is the already selected material--caloric router.  Extend
$\Theta_n$ by zero outside $I_n^{\rm ann}$.  Let
\begin{equation}
 \Psi_n=\psi(\Lambda/\Gamma_n),
 \qquad
 W_n:=\Lambda\varpi_n.
 \label{eq:v15-selected-source-writer}
\end{equation}

\begin{proposition}[Exact original-variable pullback]
\label{prop:v15-physical-pullback}
The signed physical V13 scalar is
\begin{equation}
 \mathfrak s_n
 :=-\int_{\beta_n}^{b_n}\!\int_{\Torus^3}
   \Theta_n\,(\Psi_ng_{n,t}^{\rm cov})\cdot W_n\,\dd x\,\dd t.
 \label{eq:v15-signed-selected-scalar}
\end{equation}
It satisfies the exact identities
\begin{equation}
 \boxed{
 \mathfrak s_n
 =-\mathcal R_n\int_Jq_n(s)\,\dd s
 =-\sigma_*\mathcal R_nA_n.}
 \label{eq:v15-pullback-identity}
\end{equation}
In particular,
\begin{equation}
 \boxed{
 \frac{h_0}{2}\mathcal R_n
 \le|\mathfrak s_n|
 \le\mathfrak V_*\mathcal R_n.}
 \label{eq:v15-selected-size}
\end{equation}
The positive number $\mathcal H_n^{\rm phys}$ of
\eqref{eq:v14-physical-selected-scalar} is therefore
$|\mathfrak s_n|$ with one fixed subsequential orientation restored; it is
not yet a signed caloric-contact component.
\end{proposition}

\begin{proof}
Use the V10 rescaling.  On the selected card,
\[
 \Psi_ng_{n,t}^{\rm cov}(x)
 =\nu\mathcal R_n\Gamma_n^{5/2}G_n(s,y),
 \qquad
 W_n(x)=\Gamma_n^{5/2}Z_n(y),
\]
while
$\dd x\,\dd t=(\nu\Gamma_n^5)^{-1}\dd y\,\dd s$ and
$\Theta_n=\eta\chi\Omega_n$.  Substitution in
\eqref{eq:v15-signed-selected-scalar} gives
$\mathfrak s_n=-\mathcal R_n\int_Jq_n$.  The definition
\eqref{eq:v14-oriented-window-law} gives the second equality.  The size
bounds follow from \eqref{eq:v15-standing-A}.
\end{proof}

\begin{firewall}[Orientation is not contact positivity]
\label{fire:v15-orientation-not-contact}
The sign $\sigma_*$ makes the V13 scalar mass positive after a subsequence.
It does not determine whether the actual contact
$\mathfrak s_n$ assists or opposes the complete caloric writer.  That verdict
is the sign of the mixed coefficient constructed next.  Taking an absolute
value before this mixed row would erase a physical counterflow.
\end{firewall}

% =============================================================================
\section{The terminal-writer mixed Gram and the exact source ledger}
\label{sec:v15-writer-Gram}
% =============================================================================

Use the weighted history inner product
\begin{equation}
 \langle X,Y\rangle_{\Theta_n}
 :=\int_{\beta_n}^{b_n}\!\int_{\Torus^3}
   \Theta_n(t,x)X(t,x)\cdot Y(t,x)\,\dd x\,\dd t.
 \label{eq:v15-weighted-history-product}
\end{equation}
Regard $W_n$ as constant in time.  Put
\begin{equation}
 D_n^W:=\|W_n\|_{\Theta_n}^2.
 \label{eq:v15-writer-denominator}
\end{equation}
Since $\mathfrak s_n\ne0$, one has $D_n^W>0$.  Define
\begin{equation}
 \boxed{
 \kappa_n:=
 \frac{\langle Z_{n,\cdot}^{\rm cal},W_n\rangle_{\Theta_n}}
      {D_n^W},
 \qquad
 Z_{n,t}^{\perp}:=Z_{n,t}^{\rm cal}-\kappa_nW_n.}
 \label{eq:v15-kappa-definition}
\end{equation}
Then
\begin{equation}
 \langle Z_{n,\cdot}^{\perp},W_n\rangle_{\Theta_n}=0.
 \label{eq:v15-writer-orthogonality}
\end{equation}

\begin{theorem}[Least-square terminal-writer calibration and quantitative window]
\label{thm:v15-writer-calibration}
The coefficient $\kappa_n$ is the unique minimizer of
\begin{equation}
 a\longmapsto
 \|Z_{n,\cdot}^{\rm cal}-aW_n\|_{\Theta_n}^2.
 \label{eq:v15-least-square-problem}
\end{equation}
There are constants $d_*>0$ and $K_\kappa<\infty$, depending only on the
fixed V13 interior cutoffs and the bounded-action constants, such that
\begin{equation}
 \boxed{
 D_n^W\ge\frac{d_*}{\nu},
 \qquad
 |\kappa_n|\le K_\kappa\mathcal R_n.}
 \label{eq:v15-kappa-window}
\end{equation}
More explicitly, one may take
\begin{equation}
 d_*:=\frac{h_0^2}{4L_*}
 \label{eq:v15-dstar}
\end{equation}
and
\begin{equation}
 K_\kappa
 :=\frac{2}{h_0}\sqrt{C_\eta L_*},
 \qquad
 C_\eta:=\frac{2(\tau_+-\tau_-)}{e\tau_-}.
 \label{eq:v15-Kkappa}
\end{equation}
Thus $\vartheta_n:=\kappa_n/\mathcal R_n$ is uniformly bounded, but no sign
or positive lower edge is supplied by the separate norm data.
\end{theorem}

\begin{proof}
Expanding the square in \eqref{eq:v15-least-square-problem} gives a strictly
convex quadratic with derivative
$2aD_n^W-2\langle Z^{\rm cal},W_n\rangle_{\Theta_n}$, proving the first
statement and \eqref{eq:v15-writer-orthogonality}.

Let
\[
 \mathcal A_n^\Theta
 :=\int\!\!\int\Theta_n|\Psi_ng_{n,t}^{\rm cov}|^2.
\]
The V10 scaling and \eqref{eq:v13-action-bound} give
\begin{equation}
 \mathcal A_n^\Theta
 \le\nu\mathcal R_n^2
     \int_{J_n}\|G_n(s)\|_2^2\,\dd s
 \le\nu\mathcal R_n^2L_*.
 \label{eq:v15-selected-action-upper}
\end{equation}
Cauchy--Schwarz in the weighted history space and
\eqref{eq:v15-selected-size} yield
\[
 \frac{h_0^2}{4}\mathcal R_n^2
 \le|\mathfrak s_n|^2
 \le\mathcal A_n^\Theta D_n^W,
\]
which proves the lower bound with \eqref{eq:v15-dstar}.

On $\supp\Theta_n$, the interior-age condition gives
\[
 \frac{\tau_-}{\nu\Gamma_n^2}
 \le b_n-t\le
 \frac{\tau_+}{\nu\Gamma_n^2}.
\]
Because $b_n=\beta_{n+1}$ is a first record,
$\|e(t)\|_2\le2\mathcal R_n$ for $t\le b_n$.  The elementary heat
multiplier maximum therefore gives
\[
 \|Z_{n,t}^{\rm cal}\|_2
 \le(2e\nu(b_n-t))^{-1/2}\|e(t)\|_2
 \le\sqrt{\frac{2}{e\tau_-}}\,\Gamma_n\mathcal R_n.
\]
The physical length of the time support is at most
$(\tau_+-\tau_-)/(\nu\Gamma_n^2)$, so
\begin{equation}
 \|Z_{n,\cdot}^{\rm cal}\|_{\Theta_n}^2
 \le\frac{C_\eta}{\nu}\mathcal R_n^2.
 \label{eq:v15-caloric-writer-upper}
\end{equation}
Finally,
$\kappa_n^2D_n^W\le
 \|Z_{n,\cdot}^{\rm cal}\|_{\Theta_n}^2$ by Cauchy--Schwarz in
\eqref{eq:v15-kappa-definition}.  Combine the last two bounds with
$D_n^W\ge d_*/\nu$ to obtain \eqref{eq:v15-Kkappa}.
\end{proof}

Define the actual selected contact component
\begin{equation}
 \boxed{
 \mathcal C_n^{\rm sel}:=\kappa_n\mathfrak s_n
 =-\sigma_*\kappa_n\mathcal R_nA_n.}
 \label{eq:v15-selected-contact}
\end{equation}
It has the correct quadratic record scale when
$\kappa_n\asymp\mathcal R_n$.  The contact calibration from $A_n$ is the
signed coefficient $-\sigma_*\kappa_n\mathcal R_n$, not the V14 stock
coefficient $\mathcal R_n$.

Define the remaining signed source-side terms
\begin{align}
 \mathcal U_n^{\rm out}
 &:=-\int_{\beta_n}^{b_n}
  \langle(I-\Psi_n)g_{n,t}^{\rm cov},
          Z_{n,t}^{\rm cal}\rangle\,\dd t,
 \label{eq:v15-unselected-output}\\
 \mathcal U_n^{\rm route}
 &:=-\int_{\beta_n}^{b_n}\!\int_{
 \Torus^3}
  (1-\Theta_n)(\Psi_ng_{n,t}^{\rm cov})\cdot
  Z_{n,t}^{\rm cal}\,\dd x\,\dd t,
 \label{eq:v15-route-complement}\\
 \mathcal U_n^{\rm wr}
 &:=-\int_{\beta_n}^{b_n}\!\int_{
 \Torus^3}
  \Theta_n(\Psi_ng_{n,t}^{\rm cov})\cdot
  Z_{n,t}^{\perp}\,\dd x\,\dd t.
 \label{eq:v15-writer-orthogonal-contact}
\end{align}
The term \eqref{eq:v15-route-complement} includes the already declared time,
material-router, spatial-cell, and localized replay complement.  The term
\eqref{eq:v15-unselected-output} includes every covariance output not
retained by the selected annular multiplier.  Represented source returns
which were removed before the V13 scalar are therefore never silently
reintroduced as selected work.

\begin{theorem}[Exact source-side caloric-contact ledger]
\label{thm:v15-source-contact-ledger}
The complete contact satisfies
\begin{equation}
 \boxed{
 \mathcal C_n^{\rm cal}
 =\mathcal C_n^{\rm res}
  +\mathcal U_n^{\rm out}
  +\mathcal U_n^{\rm route}
  +\mathcal C_n^{\rm sel}
  +\mathcal U_n^{\rm wr}.}
 \label{eq:v15-source-contact-ledger}
\end{equation}
Equivalently, the net critical-inventory growth is
\begin{equation}
 \boxed{
 \frac32\mathcal R_n^2
 =\mathcal C_n^{\rm res}
  +\mathcal U_n^{\rm out}
  +\mathcal U_n^{\rm route}
  +\mathcal C_n^{\rm sel}
  +\mathcal U_n^{\rm wr}
  -\mathcal L_n^{\rm cal}.}
 \label{eq:v15-inventory-ledger}
\end{equation}
All terms are assembled with their signs before any positive part is taken.
\end{theorem}

\begin{proof}
Starting from \eqref{eq:v15-covariance-contact}, use the exact telescoping
identity
\[
 g_{n,t}^{\rm cov}
 =(I-\Psi_n)g_{n,t}^{\rm cov}
 +(1-\Theta_n)\Psi_ng_{n,t}^{\rm cov}
 +\Theta_n\Psi_ng_{n,t}^{\rm cov}.
\]
In the last term insert
$Z_{n,t}^{\rm cal}=\kappa_nW_n+Z_{n,t}^{\perp}$ and use
\eqref{eq:v15-signed-selected-scalar}.  This proves
\eqref{eq:v15-source-contact-ledger} after adding the resolved term from
\eqref{eq:v15-res-cov-ledger}.  Subtract the carry loss using
\eqref{eq:v15-inventory-secant}.
\end{proof}

\begin{countertheorem}[Record height and marginal norms do not determine terminal-writer incidence]
\label{cth:v15-marginals-no-writer-incidence}
For every $R>0$ there is a one-parameter family in the real Hilbert space
$\R^2$ with fixed source $G=e_1$, selected writer $W=e_1$, and actual writer
\begin{equation}
 Z_\theta=R(\cos\theta\,e_1+\sin\theta\,e_2)
 \label{eq:v15-finite-writer-family}
\end{equation}
for which
\begin{equation}
 \|G\|=\|W\|=1,
 \qquad \|Z_\theta\|=R,
 \qquad -\langle G,W\rangle=-1
 \label{eq:v15-fixed-marginals}
\end{equation}
are independent of $\theta$, while
\begin{equation}
 \boxed{
 \kappa_\theta=R\cos\theta,
 \qquad
 -\langle G,Z_\theta\rangle=-R\cos\theta.}
 \label{eq:v15-varying-incidence}
\end{equation}
Thus a record-size writer norm and a fixed selected source pairing do not
supply the mixed contact coefficient or its sign.
\end{countertheorem}

\begin{proof}
All formulas follow by direct substitution.  The least-square projection of
$Z_\theta$ onto $W$ has coefficient $\langle Z_\theta,W\rangle=R\cos\theta$.
\end{proof}

% =============================================================================
\section{Equality of the source and endpoint ledgers}
\label{sec:v15-two-ledgers}
% =============================================================================

Define the source-to-final-head incidence residual
\begin{equation}
 \boxed{
 \mathcal I_n^{\rm head}
 :=\mathcal C_n^{\rm ho}-\mathcal C_n^{\rm sel}.}
 \label{eq:v15-head-incidence-residual}
\end{equation}

\begin{theorem}[Exact two-ledger incidence identity]
\label{thm:v15-two-ledger-identity}
The residual \eqref{eq:v15-head-incidence-residual} has the exact signed
representation
\begin{equation}
 \boxed{
 \mathcal I_n^{\rm head}
 =\mathcal C_n^{\rm res}
  +\mathcal U_n^{\rm out}
  +\mathcal U_n^{\rm route}
  +\mathcal U_n^{\rm wr}
  -\mathcal C_n^{\rm rep}.}
 \label{eq:v15-two-ledger-identity}
\end{equation}
In particular, equality of the selected source contact with the held-out
terminal contact requires simultaneous assembly of the resolved source,
unselected covariance outputs, material/cell/time escape, writer
misalignment, and represented endpoint return.  None of their separate
marginal bounds can replace this signed identity.
\end{theorem}

\begin{proof}
Subtract $\mathcal C_n^{\rm sel}$ from the endpoint ledger
\eqref{eq:v15-endpoint-ledger} and replace
$\mathcal C_n^{\rm cal}$ by the source ledger
\eqref{eq:v15-source-contact-ledger}.  Rearrangement gives the formula.
\end{proof}

\begin{corollary}[The corrected full-contact residual]
\label{cor:v15-corrected-residual}
The contact-typed complement of the selected scalar is
\begin{equation}
 \boxed{
 \widetilde\Delta_n^{\rm full}
 :=\mathcal C_n^{\rm cal}-\mathcal C_n^{\rm sel}
 =\mathcal C_n^{\rm res}
  +\mathcal U_n^{\rm out}
  +\mathcal U_n^{\rm route}
  +\mathcal U_n^{\rm wr}.}
 \label{eq:v15-corrected-full-residual}
\end{equation}
It differs from the direct V14 residual by the terminal-writer calibration:
\begin{equation}
 \boxed{
 \widetilde\Delta_n^{\rm full}
 -\Delta_{n,\rm nat}^{\rm full}
 =\mathcal R_nA_n-\mathcal C_n^{\rm sel}
 =\mathcal R_nA_n
  +\sigma_*\kappa_n\mathcal R_nA_n.}
 \label{eq:v15-residual-correction}
\end{equation}
The expression is an algebraic comparison of two residual definitions; the
summands on its right are not asserted to have the same stock semantics.
\end{corollary}

\begin{proof}
The first identity is \eqref{eq:v15-source-contact-ledger}.  Substitute
\eqref{eq:v15-natural-residual} and
\eqref{eq:v15-selected-contact} for the second.
\end{proof}

% =============================================================================
\section{The exact fixed-datum alternative}
\label{sec:v15-alternative}
% =============================================================================

The complete contact is positive and obeys
$\mathcal C_n^{\rm cal}\ge\tfrac32\mathcal R_n^2$.  Every ratio below is
therefore well defined.  Apply the following priority order before taking a
positive part.

\begin{theorem}[Represented contact, typed incidence residual, or selected critical replenishment]
\label{thm:v15-main-alternative}
After passage to a subsequence, exactly one of the following priority
branches occurs.
\begin{enumerate}[label=\textup{(V15.\arabic*)},leftmargin=5.7em]
\item \emph{Represented endpoint contact.}  There is $\delta>0$ such that
      \begin{equation}
       \boxed{
       \frac{|\mathcal C_n^{\rm rep}|}
            {\mathcal C_n^{\rm cal}}\ge\delta.}
       \label{eq:v15-represented-survivor}
      \end{equation}
      This is a signed represented-return/endpoint-promotion survivor; it is
      not a fresh source payment.
\item \emph{Source-to-head incidence residual.}  The first branch fails and
      there is $\delta>0$ such that
      \begin{equation}
       \boxed{
       \frac{|\mathcal I_n^{\rm head}|}
            {\mathcal C_n^{\rm cal}}\ge\delta.}
       \label{eq:v15-head-incidence-survivor}
      \end{equation}
      After a further subsequence, at least one of
      \begin{equation}
       \mathcal C_n^{\rm res},\qquad
       \mathcal U_n^{\rm out},\qquad
       \mathcal U_n^{\rm route},\qquad
       \mathcal U_n^{\rm wr}
       \label{eq:v15-four-typed-complements}
      \end{equation}
      has absolute value at least
      $(\delta/8)\mathcal C_n^{\rm cal}$ for all late $n$.
\item \emph{Selected held-out critical replenishment.}  The preceding
      branches fail, so
      \begin{equation}
       \boxed{
       \frac{\mathcal C_n^{\rm rep}}
            {\mathcal C_n^{\rm cal}}\longrightarrow0,
       \qquad
       \frac{\mathcal I_n^{\rm head}}
            {\mathcal C_n^{\rm cal}}\longrightarrow0,}
       \label{eq:v15-clean-incidence}
      \end{equation}
      and consequently
      \begin{equation}
       \boxed{
       \frac{\mathcal C_n^{\rm sel}}
            {\mathcal C_n^{\rm cal}}\longrightarrow1.}
       \label{eq:v15-selected-full-contact}
      \end{equation}
      In this case $\mathcal C_n^{\rm sel}>0$ eventually and
      \begin{equation}
       \boxed{
       -\sigma_*\frac{\kappa_n}{\mathcal R_n}A_n
       =\frac{\mathcal C_n^{\rm cal}}{\mathcal R_n^2}+o(1)
       \ge\frac32+o(1).}
       \label{eq:v15-alignment-lower}
      \end{equation}
\end{enumerate}
The priority order is exhaustive.  The second branch is typed only after the
represented endpoint term has been shown negligible.
\end{theorem}

\begin{proof}
Pass first to a subsequence on which
$|\mathcal C_n^{\rm rep}|/\mathcal C_n^{\rm cal}$ either tends to zero or
has a positive lower bound.  This gives the first branch or its complement.
On the latter, pass again for
$|\mathcal I_n^{\rm head}|/\mathcal C_n^{\rm cal}$.
If it has a positive lower bound, use
\eqref{eq:v15-two-ledger-identity}.  Since
$|\mathcal C_n^{\rm rep}|/\mathcal C_n^{\rm cal}\to0$, for all late $n$
\[
 |\mathcal C_n^{\rm res}+
   \mathcal U_n^{\rm out}+
   \mathcal U_n^{\rm route}+
   \mathcal U_n^{\rm wr}|
 \ge\frac\delta2\mathcal C_n^{\rm cal}.
\]
The triangle inequality shows that one of the four terms has magnitude at
least $\delta\mathcal C_n^{\rm cal}/8$; a finite priority subsequence fixes
its type.

On the remaining branch,
$\mathcal C_n^{\rm ho}/\mathcal C_n^{\rm cal}\to1$ by
\eqref{eq:v15-endpoint-ledger}, and
$\mathcal C_n^{\rm sel}/\mathcal C_n^{\rm cal}\to1$ by the definition of
$\mathcal I_n^{\rm head}$.  Positivity follows from positivity of the full
contact.  Divide
$\mathcal C_n^{\rm sel}=-\sigma_*\kappa_n\mathcal R_nA_n$ by
$\mathcal R_n^2$.  Since
$\mathcal C_n^{\rm sel}-\mathcal C_n^{\rm cal}
 =o(\mathcal C_n^{\rm cal})$ and
$\mathcal C_n^{\rm cal}/\mathcal R_n^2\le2$ by
\eqref{eq:v15-full-contact-window}, the normalized difference is $o(1)$.
This gives \eqref{eq:v15-alignment-lower}.
\end{proof}

\begin{corollary}[Quantitative writer incidence and carry control on the clean branch]
\label{cor:v15-clean-quantitative}
On branch \textup{(V15.3)}, there are constants
$\kappa_*>0$, $\eta_*>0$, and $L_*^{\rm cal}<\infty$ such that, for all
late $n$,
\begin{equation}
 \boxed{
 -\sigma_*\kappa_n>0,
 \qquad
 \frac{|\kappa_n|}{\mathcal R_n}\ge\kappa_*,}
 \label{eq:v15-kappa-lower}
\end{equation}
\begin{equation}
 \boxed{
 \eta_n^{\rm tw}
 :=\frac{\kappa_n^2D_n^W}
 {\|Z_{n,\cdot}^{\rm cal}\|_{\Theta_n}^2}
 \ge\eta_*,}
 \label{eq:v15-writer-efficiency}
\end{equation}
and
\begin{equation}
 \boxed{
 0\le\frac{\mathcal L_n^{\rm cal}}{\mathcal R_n^2}
 \le L_*^{\rm cal}.}
 \label{eq:v15-carry-control}
\end{equation}
One may take, after increasing the starting index,
\begin{equation}
 \kappa_*:=\frac{1}{\mathfrak V_*},
 \qquad
 \eta_*:=\min\left\{\frac12,
          \frac{\kappa_*^2d_*}{2C_\eta}\right\},
 \qquad
 L_*^{\rm cal}:=\frac12.
 \label{eq:v15-explicit-clean-constants}
\end{equation}
Thus the clean branch contains a genuine noncollapsed mixed terminal-writer
Gram; it cannot be produced from the separate marginals of
\Cref{cth:v15-marginals-no-writer-incidence}.
\end{corollary}

\begin{proof}
Equation~\eqref{eq:v15-alignment-lower} and
$A_n\le\mathfrak V_*$ give, for all late $n$,
$|\kappa_n|/\mathcal R_n\ge1/\mathfrak V_*$, with the displayed sign.
Combine this with
$D_n^W\ge d_*/\nu$ and
\eqref{eq:v15-caloric-writer-upper} to obtain
\eqref{eq:v15-writer-efficiency}.

The carry estimate is independent of the clean branch.  Heat contraction in
\eqref{eq:v15-full-contact} gives directly
$0\le\mathcal L_n^{\rm cal}/\mathcal R_n^2\le1/2$, proving
\eqref{eq:v15-carry-control} with the stated constant.
\end{proof}

\begin{corollary}[The clean component is a non-summable replenishment survivor]
\label{cor:v15-nonsummable-contact}
On branch \textup{(V15.3)},
\begin{equation}
 \boxed{
 \mathcal C_n^{\rm sel}
 \ge\frac34\mathcal R_n^2}
 \label{eq:v15-selected-contact-lower}
\end{equation}
for all sufficiently large $n$, and hence
\begin{equation}
 \boxed{
 \sum_n\mathcal C_n^{\rm sel}=+\infty.}
 \label{eq:v15-selected-contact-diverges}
\end{equation}
This is the actual critical-replenishment conclusion licensed by the clean
incidence branch.  It is not a contradiction without an independently
occurring finite stock or a monotone upper budget for the same contact.
\end{corollary}

\begin{proof}
Use \eqref{eq:v15-selected-full-contact} and
$\mathcal C_n^{\rm cal}\ge\tfrac32\mathcal R_n^2$.  The geometric record
sequence makes the series diverge.
\end{proof}

\begin{corollary}[The caloric carry is universally record quadratic]
\label{cor:v15-universal-carry}
On every doubled record, independently of the V15 alternative,
\begin{equation}
 \boxed{
 0\le\frac{\mathcal L_n^{\rm cal}}{\mathcal R_n^2}\le\frac12,
 \qquad
 \frac32\le
 \frac{\mathcal C_n^{\rm cal}}{\mathcal R_n^2}
 \le2.}
 \label{eq:v15-universal-carry}
\end{equation}
Thus there is no separate carry-loss-overdrive branch.  The carry is already
part of the complete record-quadratic contact and must remain inside the
signed ledgers above.
\end{corollary}

\begin{proof}
Use the identity and heat-contraction bound in
\eqref{eq:v15-full-contact}.
\end{proof}

% =============================================================================
\section{Integrated V15 conclusion and exact stopping boundary}
\label{sec:v15-integrated}
% =============================================================================

\begin{theorem}[Version V15 terminal-writer and full-contact theorem]
\label{thm:v15-integrated}
Preserving every theorem and limitation of Versions~V1--V14, Version~V15
proves the following.
\begin{enumerate}[label=\textup{(T15.\arabic*)},leftmargin=5.5em]
\item The direct record-amplitude quantity
      $\mathcal R_nA_n$ is $o(\mathcal C_n^{\rm cal})$; the V14 natural
      residual therefore contains asymptotically the complete caloric
      contact and is not a component test.
\item The complete contact is the exact Duhamel secant
      \eqref{eq:v15-Duhamel-secant}.  The literal represented terminal head
      and its held-out complement have the signed endpoint ledger
      \eqref{eq:v15-endpoint-ledger}; their half-norms are the endpoint
      promotion terms.
\item The same contact has the exact resolved/heat-covariance decomposition
      \eqref{eq:v15-res-cov-ledger}.  Gradient and isotropic pressure
      followers vanish identically against the caloric writer.
\item The original-variable V13 scalar is
      $\mathfrak s_n=-\sigma_*\mathcal R_nA_n$.  Its orientation is restored
      before contact positivity is discussed.
\item The selected V5 test direction has the exact window
      \eqref{eq:v15-varpi-window}, and its V10 rescaled writer has the corrected
      lower edge \eqref{eq:v15-rescaled-writer-window}.  The unique least-square
      mixed writer coefficient is $\kappa_n$ of
      \eqref{eq:v15-kappa-definition}.  Its normalized magnitude is bounded,
      but its sign and lower edge are independent physical incidence data.
\item The source ledger \eqref{eq:v15-source-contact-ledger} and endpoint
      ledger \eqref{eq:v15-endpoint-ledger} are exact.  Their difference is
      the source-to-final-head residual
      \eqref{eq:v15-two-ledger-identity}; the caloric carry loss enters only
      in the net-inventory ledger \eqref{eq:v15-inventory-ledger}.
\item Every bounded-action selected sequence returns a represented endpoint
      contact, a typed resolved/unselected/router/writer complement, or a
      selected held-out caloric component asymptotic to the complete contact.
\item On the clean selected branch the mixed writer Gram has record-size
      coefficient and fixed efficiency, the universal carry lies in the window
      $[0,\mathcal R_n^2/2]$, and the selected caloric component is
      non-summable.
\item A non-summable selected contact is a critical-replenishment survivor,
      not a regularity contradiction.  An independently occurring finite
      stock or monotone upper budget is still absent.
\end{enumerate}
No terminal trace, backward uniqueness, Liouville theorem, or global
three-dimensional regularity conclusion is asserted.
\end{theorem}

\begin{proof}
Items~\textup{(T15.1)--(T15.2)} are
Theorems~\ref{thm:v15-direct-residual-full} and
\ref{thm:v15-Duhamel-secant}, together with
Proposition~\ref{prop:v15-endpoint-ledger}.  Item~\textup{(T15.3)} is
Theorem~\ref{thm:v15-full-forced-Euler} and
Proposition~\ref{prop:v15-pressure-follower-zero}.  Items
\textup{(T15.4)--(T15.5)} are
Proposition~\ref{prop:v15-annular-normalization},
Proposition~\ref{prop:v15-physical-pullback}, and
Theorem~\ref{thm:v15-writer-calibration}.  Item~\textup{(T15.6)} is
Theorems~\ref{thm:v15-source-contact-ledger} and
\ref{thm:v15-two-ledger-identity}.  The final three items are
Theorem~\ref{thm:v15-main-alternative} and
Corollaries~\ref{cor:v15-clean-quantitative},
\ref{cor:v15-nonsummable-contact}, and
\ref{cor:v15-universal-carry}.
\end{proof}

\begin{firewall}[Exact stopping boundary after V15]
\label{fire:v15-stop}
V15 has settled the component-typing question.  The bounded-action V13 scalar
cannot be subtracted directly from the critical contact; it enters only after
the terminal caloric writer is projected onto its fixed Riesz writer.  A
failure of that mixed Gram, of the source/head identity, or of the
represented endpoint short is already the terminal obstruction.  On the
clean branch the programme has an actual non-summable critical replenishment,
but no finite upper stock.  Adding another lower action estimate, temporal
partition, compact profile space, or abstract stock would be downstream of
this stop.
\end{firewall}

\begingroup\small
\begin{longtable}{p{0.25\textwidth}p{0.19\textwidth}p{0.48\textwidth}}
\caption{NS--A Version V15 proof-status ledger.}\label{tab:v15-ledger}\\
\toprule Row & Status & Exact conclusion\\\midrule
\endfirsthead
\toprule Row & Status & Exact conclusion\\\midrule
\endhead
V1--V14 prefix & \PROVED{} preserved
 &The complete V14 preterminal source is retained byte-for-byte; no earlier
 theorem is rewritten.\\[0.3em]
Natural V14 contact subtraction & \OBSTRUCTION{} forced complement
 &$\mathcal R_nA_n/\mathcal C_n^{\rm cal}\to0$; the direct residual is
 asymptotically the full contact.\\[0.3em]
Full Duhamel contact & \PROVED{} exact
 &The contact is a critical-state secant with one linear response term and one
 positive endpoint-promotion half-norm.\\[0.3em]
Represented endpoint head & \PROVED{} exact signed ledger
 &The full response splits into represented and held-out contact without any
 assumption that the represented head reduces $A$.\\[0.3em]
Resolved/Reynolds source split & \PROVED{} exact
 &The complete caloric contact is resolved heat--Euler work plus heat-covariance
 work on the same record.\\[0.3em]
Pressure-compatible followers & \PROVED{} zero
 &Gradient and isotropic followers pair to zero with the divergence-free
 caloric writer and receive no separate contact price.\\[0.3em]
V13 physical pullback & \PROVED{} signed
 &$\mathfrak s_n=-\sigma_*\mathcal R_nA_n$ on the original physical card.\\[0.3em]
Annular test normalization & \PROVED{} corrected
 &$\sqrt\eta\le\|\varpi_n\|_2\le1$ and
 $\sqrt\eta/2\le\|Z_n\|_2\le2$; the earlier unit-norm shorthand is not used.\\[0.3em]
Terminal-writer mixed Gram & \PROVED{} exact acquisition
 &$\kappa_n$ is the unique least-square contact calibration;
 $|\kappa_n|\lesssim\mathcal R_n$, but its sign and floor are not marginal data.\\[0.3em]
Full source ledger & \PROVED{} exact
 &Resolved contact, unselected outputs, route/cell/time escape, selected mixed
 writer, and writer-orthogonal work sum to the complete contact.\\[0.3em]
Source-to-final-head incidence & \PROVED{} exact residual
 &The difference between held-out endpoint contact and selected source contact
 is the assembled signed five-term identity \eqref{eq:v15-two-ledger-identity}.\\[0.3em]
Clean selected branch & \CONDITIONAL{} quantitative
 &Vanishing represented and incidence residuals give full-contact capture,
 record-size writer alignment, positive efficiency, and bounded relative carry loss.\\[0.3em]
Critical replenishment & \PROVED{} survivor on clean branch
 &The selected contact is at least $c\mathcal R_n^2$ and its cofinal sum diverges.\\[0.3em]
Finite stock or monotone upper budget & \OPEN{} physical
 &No occurring finite inventory has been proved to dominate the same calibrated
 contact.\\[0.3em]
Action-overdrive profile & \INHERITED{} separate normalization
 &It may be reopened only on the V10--V11 enstrophy-square/source-action
 overdrive branch.\\[0.3em]
Three-dimensional regularity & \OPEN
 &Neither the typed complement nor the non-summable replenishment is excluded
 for every fixed datum.\\
\bottomrule
\end{longtable}
\endgroup

% =============================================================================
\section{Exact mandate for Version V16}
\label{sec:v15-next}
% =============================================================================

\begin{programme}[Evaluate the mixed terminal writer or stop]
\label{prog:v15-next}
Preserve Versions~V1--V15 verbatim.  Do not introduce another response norm,
source-action lower bound, fast-time window law, stock reflection, or compact
profile space.  The next continuation is restricted to the literal V15
mixed card.

Proceed only in the following order.
\begin{enumerate}[label=\textup{(V16.\arabic*)},leftmargin=5.7em]
\item Evaluate the same-history mixed Gram
      $\langle Z_{n,\cdot}^{\rm cal},W_n\rangle_{\Theta_n}$ and its
      denominator on the already selected fixed-datum cards.  A sign change,
      calibration collapse, or efficiency loss is the terminal-writer
      obstruction.
\item Evaluate the represented endpoint contact
      $\mathcal C_n^{\rm rep}$ before calling any held-out response fresh.
      A fixed signed fraction is returned to its original source ancestry.
\item On the represented-null branch, evaluate the assembled identity
      \eqref{eq:v15-two-ledger-identity}.  Do not bound its five terms
      separately before their signs and common endpoint are fixed.
\item If the residual survives, export exactly one of: resolved heat--Euler
      contact, unselected covariance output, material/cell/time escape, or
      writer-orthogonal work.  Do not rename it a failure of terminal
      compactness.
\item If both represented and incidence residuals vanish, stop NS--A on that
      subsequence and export the actual selected critical-replenishment card
      $(\mathfrak s_n,\kappa_n,\mathcal C_n^{\rm sel})$ to the fixed-datum
      Reynolds/paid-stress programme.  The remaining question there is an
      independent finite-stock upper budget, not another NS--A lower bound.
\item Reopen a compact dynamical profile only on the separately normalized
      source-action/enstrophy-square overdrive branch, preserving its own
      physical scale and source ancestry.
\end{enumerate}
A successful Version~V16 ends in
\begin{equation}
 \boxed{
 \begin{gathered}
 \text{terminal-writer calibration obstruction}\\
 \big|\quad\text{represented endpoint return}\\
 \big|\quad\text{one typed signed full-contact complement}\\
 \big|\quad\text{one exported selected critical-replenishment card}\\
 \big|\quad\text{a separately normalized action-overdrive profile}.
 \end{gathered}}
 \label{eq:v15-v16-endpoint}
\end{equation}
No new Dini modulus, Liouville theorem, terminal trace, or global regularity
claim is admitted before this finite mixed card is evaluated.
\end{programme}

\begin{assessment}[Programme position after V15]
\label{ass:v15-position}
The missing arrow has been reduced to one concrete same-history mixed Gram and
one signed finite ledger.  Record height gives the size of the caloric writer
but not its orientation relative to the selected response writer.  Once that
mixed row is supplied, every other component of the full contact has an exact
name and sign.  The programme is therefore no longer licensed to develop
additional general machinery on the bounded-action branch.
\end{assessment}

\paragraph{Version V15 history.}
Preserved the complete V14 source and audited its mandated insertion into the
caloric rectangle.  Proved that the direct record-amplitude subtraction is
asymptotically the whole contact, derived the full Duhamel secant and the
represented/held-out endpoint ledger, and assembled the full forced-Euler
resolved/covariance ledger.  Pulled the V13 scalar back with its physical
sign, corrected the inherited annular test-direction normalization and its
rescaled writer lower edge, constructed the unique terminal-writer mixed Gram,
proved its record-scale upper window, and gave a finite counterexample showing that
separate source and writer marginals do not determine it.  Derived the exact
source-side contact ledger, the source-to-final-head incidence identity, and
the corrected contact residual.  Finally proved the represented-return,
typed-complement, or selected-critical-replenishment alternative, with
quantitative writer efficiency, carry-loss control, and non-summability on
the clean branch.  No finite-stock upper budget or regularity theorem was
claimed.

\begin{assessment}[Append-only integrity]
\label{ass:v15-integrity}
The complete Version~V14 source preceding its sole final document terminator
is preserved verbatim.  Its preterminal SHA--256 digest is
\begin{center}\scriptsize\ttfamily
f855e0e6ca791a4d595c1e33adc935861fdb23e2e593ae666461f35164fd0a25
\end{center}
The present material begins at Section~\ref{sec:v15-scope}; the final command
below is again unique.
\end{assessment}

\addtocontents{toc}{\protect\endgroup}
% =============================================================================
% END VERSION V15 MATERIAL -- append later versions before final terminator
% =============================================================================

% APPEND ALL LATER VERSIONS IMMEDIATELY ABOVE THE SOLE FINAL LINE BELOW.

\clearpage
% =============================================================================
% VERSION V16: NEW MATERIAL.  The complete V1--V15 source above is preserved
% verbatim; only its sole active \end{document} line was removed before append.
% =============================================================================
\hypersetup{pdftitle={NS-A Terminal Source Concentration Programme: Cumulative V16},pdfauthor={A. Pelissier}}
\makeatletter
\addtocontents{toc}{\protect\begingroup\protect\makeatletter}
\addtocontents{toc}{\protect\def\protect\l@section{\protect\@dottedtocline{1}{0em}{4.7em}}}
\addtocontents{toc}{\protect\def\protect\l@subsection{\protect\@dottedtocline{2}{1.5em}{5.3em}}}
\addtocontents{toc}{\protect\makeatother}
\makeatother
\renewcommand{\thesection}{V16.\arabic{section}}
\renewcommand{\thesubsection}{V16.\arabic{section}.\arabic{subsection}}
\renewcommand{\theequation}{V16.\arabic{equation}}
\setcounter{section}{0}
\setcounter{equation}{0}
\renewcommand{\theHsection}{V16.\arabic{section}}
\renewcommand{\theHsubsection}{V16.\arabic{section}.\arabic{subsection}}
\renewcommand{\theHtheorem}{V16.\arabic{section}.\arabic{theorem}}
\renewcommand{\theHequation}{V16.\arabic{equation}}

\begin{center}
{\LARGE\bfseries NS--A: Version V16}\\[0.5em]
{\large The Dimensionless Terminal-Writer Card, the Three-Column Gram Short,\\
and the Exact Finite Branch Certificate}\\[0.5em]
{\normalsize 3 September 2026}
\end{center}
\addcontentsline{toc}{part}{Version V16: the dimensionless terminal-writer card and finite branch certificate}

\begin{abstract}
Version~V15 reduced the bounded-action terminal branch to one same-history
mixed coefficient,
\[
 \kappa_n=
 \frac{\langle Z_{n,\cdot}^{\rm cal},W_n\rangle_{\Theta_n}}
      {\|W_n\|_{\Theta_n}^2},
\]
one represented endpoint contact, and one assembled source-to-head residual.
The present continuation does not introduce another response norm or profile
space.  It evaluates the mathematical content of that finite card before any
of its signed components are estimated separately.

After the record-amplitude rescaling, the normalized coefficient
$\kappa_n/\mathcal R_n$ is exactly the least-square projection of a uniformly
bounded rescaled caloric writer onto the fixed annular writer.  Its denominator
has a two-sided cutoff-independent window.  The selected covariance source,
the annular writer, and the complete caloric writer therefore determine one
bounded positive three-by-three Gram matrix.  Shorting this Gram by the
annular writer gives an exact two-by-two residual matrix.  Its off-diagonal
entry is precisely the writer-orthogonal contact of Version~V15; its diagonal
entries are the source and terminal-writer Riesz residual actions.  This gives
a sharp determinant bound and shows that a record-sized writer-orthogonal
contact forces an actual source-angle or terminal-writer-angle defect.

The mixed coefficient can be acquired by two positive balanced writer actions,
$\|Z^{\rm cal}\pm\mathcal R_nW_n\|_{\Theta_n}^2$.  The represented endpoint
contact is likewise one difference of positive endpoint norms.  Consequently
the entire V15 priority alternative is decided by five dimensionless scalar
Reads:
\[
 \left(
 \frac{\mathcal C_n^{\rm cal}}{\mathcal R_n^2},
 \frac{\mathcal C_n^{\rm rep}}{\mathcal R_n^2},
 A_n,
 \nu D_n^W,
 \frac{\nu}{\mathcal R_n}
 \langle Z_{n,\cdot}^{\rm cal},W_n\rangle_{\Theta_n}
 \right).
\]
They lie in one compact set, and the selected-contact and incidence-residual
readouts are uniformly Lipschitz on that set.  Thus a finite outward interval
for these five values gives a rigorous branch certificate without first
bounding the resolved, unselected-output, route, and writer-orthogonal terms.
Only after a nonzero aggregate residual has been certified is its signed
four-term splitter opened.

The note also identifies the chronological content of the mixed row.  The
caloric writer is the carried old critical state plus the cumulative Duhamel
response transported to the common endpoint.  Hence the mixed coefficient
contains an ordered first moment of the actual source history.  Endpoint
response, total source length, and marginal norms do not determine this
moment.  On the clean branch the five-scalar card exports the actual selected
critical-replenishment triple
$(\mathfrak s_n,\kappa_n,\mathcal C_n^{\rm sel})$ and NS--A stops on that
subsequence.  The supplied theorem sources contain no populated values of the
five Reads, so no literal terminal trajectory is assigned to one branch in
this version.  No global regularity conclusion is claimed.
\end{abstract}

% =============================================================================
\section{Scope: evaluate the finite mixed card before reopening analysis}
\label{sec:v16-scope}
% =============================================================================

Retain all notation and hypotheses of the bounded-action branch of
Versions~V10--V15.  In particular,
\[
 \mathfrak s_n=-\sigma_*\mathcal R_nA_n,
 \qquad
 \frac{h_0}{2}\le A_n\le\mathfrak V_*,
\]
$W_n=\Lambda\varpi_n$, and
$Z_{n,t}^{\rm cal}=\Lambda e^{-\nu(b_n-t)A}e(t)$.  The complete cutoff is
$0\le\Theta_n\le1$ and is supported where
\[
 \frac{\tau_-}{\nu\Gamma_n^2}
 \le b_n-t\le
 \frac{\tau_+}{\nu\Gamma_n^2}.
\]
The selected source square action satisfies
\begin{equation}
 \mathcal A_n^\Theta
 :=\int\!\!\int
   \Theta_n|\Psi_ng_{n,t}^{\rm cov}|^2
 \le\nu\mathcal R_n^2L_*.
 \label{eq:v16-source-action-window}
\end{equation}

Define three vectors in the single real history space
\[
 \mathcal H_n
 :=L^2\bigl((\beta_n,b_n)\times\Torus^3,
            \dd t\,\dd x;\R^3\bigr)
\]
by
\begin{equation}
 \mathbf g_n:=\Theta_n^{1/2}\Psi_ng_{n,\cdot}^{\rm cov},
 \qquad
 \mathbf w_n:=\Theta_n^{1/2}W_n,
 \qquad
 \mathbf z_n:=\Theta_n^{1/2}Z_{n,\cdot}^{\rm cal}.
 \label{eq:v16-three-history-vectors}
\end{equation}
Multiplication by $\Theta_n^{1/2}$ is understood in the ordinary $L^2$
sense; no differentiability of the square root is used.  Put
\begin{equation}
 \begin{aligned}
 E_n^g&:=\|\mathbf g_n\|_2^2=\mathcal A_n^\Theta,
 &D_n^W&:=\|\mathbf w_n\|_2^2,\\
 E_n^z&:=\|\mathbf z_n\|_2^2,
 &N_n^{zw}&:=\langle\mathbf z_n,\mathbf w_n\rangle,\\
 \mathcal C_n^\Theta
 &:=-\langle\mathbf g_n,\mathbf z_n\rangle.
 \end{aligned}
 \label{eq:v16-five-Gram-scalars}
\end{equation}
Then
\begin{equation}
 \langle\mathbf g_n,\mathbf w_n\rangle=-\mathfrak s_n
 =\sigma_*\mathcal R_nA_n,
 \qquad
 N_n^{zw}=\kappa_nD_n^W,
 \label{eq:v16-known-cross-entries}
\end{equation}
and
\begin{equation}
 \mathcal C_n^\Theta
 =\mathcal C_n^{\rm sel}+\mathcal U_n^{\rm wr}.
 \label{eq:v16-local-cov-contact}
\end{equation}
Thus all three columns occur on the same time, cell, router, and endpoint
history.

\begin{firewall}[No new source is introduced by the Gram]
\label{fire:v16-Gram-not-source}
The vectors in \eqref{eq:v16-three-history-vectors} are the already occurring
selected covariance source, its fixed annular writer, and the complete
caloric writer.  Polarizing them evaluates mixed incidence.  It neither
creates another nonlinear source nor assigns a second price to the selected
response.  The Gram is used only after the source, writer, cutoff, and common
endpoint have been fixed.
\end{firewall}

% =============================================================================
\section{Two-sided denominator control and the fixed-scale mixed coefficient}
\label{sec:v16-fixed-scale}
% =============================================================================

\begin{proposition}[The selected-writer denominator has a two-sided physical window]
\label{prop:v16-denominator-window}
Let
\begin{equation}
 d^*:=4(\tau_+-\tau_-).
 \label{eq:v16-dstar-upper}
\end{equation}
Then
\begin{equation}
 \boxed{
 \frac{d_*}{\nu}\le D_n^W\le\frac{d^*}{\nu}.}
 \label{eq:v16-denominator-window}
\end{equation}
In particular the mixed coefficient is never made large merely by collapse of
its selected-writer denominator on the present branch.
\end{proposition}

\begin{proof}
The lower bound is \eqref{eq:v15-kappa-window}.  For the upper bound, the
spectral support of $\varpi_n$ lies in
$[\Gamma_n/2,2\Gamma_n]$ and
$\|\varpi_n\|_2\le1$, so
\[
 \|W_n\|_2^2=\|\Lambda\varpi_n\|_2^2
 \le4\Gamma_n^2.
\]
The time support of $\Theta_n$ has length at most
$(\tau_+-\tau_-)/(\nu\Gamma_n^2)$ and $0\le\Theta_n\le1$.  Therefore
\[
 D_n^W
 \le\frac{\tau_+-\tau_-}{\nu\Gamma_n^2}
      \|W_n\|_2^2
 \le\frac{4(\tau_+-\tau_-)}{\nu}.
\]
\end{proof}

We now write the mixed coefficient on the fixed V10 scale.  Put
$A_y=-\Delta_y$ and retain
\[
 s=\nu\Gamma_n^2(t-b_n),
 \qquad y=\Gamma_n(x-y_n),
 \qquad
 \alpha_n=\sqrt{\frac{\mathcal R_n}{\nu}}.
\]
Define the normalized caloric writer
\begin{equation}
 \boxed{
 \mathscr Z_n^{\rm cal}(s,y)
 :=\alpha_n^{-1}\Lambda_y e^{sA_y}
   \Lambda_y^{1/2}V_n(s,y),
 \qquad s<0.}
 \label{eq:v16-scaled-caloric-writer}
\end{equation}
The exponential is the forward heat contraction through age $-s>0$.

\begin{theorem}[Exact fixed-scale evaluation of the terminal-writer coefficient]
\label{thm:v16-fixed-scale-kappa}
On the V10 rescaled card,
\begin{equation}
 Z_{n,t}^{\rm cal}(x)
 =\mathcal R_n\Gamma_n^{5/2}
   \mathscr Z_n^{\rm cal}(s,y),
 \qquad
 W_n(x)=\Gamma_n^{5/2}Z_n(y).
 \label{eq:v16-two-writer-scalings}
\end{equation}
Consequently, with
\begin{align}
 \mathfrak d_n
 &:=\nu D_n^W
 =\int_{J_n}\!\!\int_{\Torus_n^3}
   \eta(s)\chi(y)\Omega_n(s,y)|Z_n(y)|^2\,\dd y\,\dd s,
 \label{eq:v16-dimensionless-denominator}\\
 \mathfrak m_n
 &:=\frac{\nu}{\mathcal R_n}N_n^{zw}
 =\int_{J_n}\!\!\int_{\Torus_n^3}
   \eta(s)\chi(y)\Omega_n(s,y)
   \mathscr Z_n^{\rm cal}(s,y)\cdot Z_n(y)
   \,\dd y\,\dd s,
 \label{eq:v16-dimensionless-mixed}
\end{align}
one has the exact formula
\begin{equation}
 \boxed{
 \frac{\kappa_n}{\mathcal R_n}
 =\frac{\mathfrak m_n}{\mathfrak d_n}.}
 \label{eq:v16-kappa-fixed-card}
\end{equation}
Moreover,
\begin{equation}
 \boxed{
 \sup_{s\in\supp\eta}
 \|\mathscr Z_n^{\rm cal}(s)\|_2
 \le\sqrt{\frac{2}{e\tau_-}}.}
 \label{eq:v16-scaled-caloric-bound}
\end{equation}
Thus the unresolved coefficient is one bounded projection on a fixed
rescaled time and frequency card; no diverging ambient norm remains inside
its definition.
\end{theorem}

\begin{proof}
The V10 scaling gives
\[
 u(t,x)=\Gamma_n\sqrt{\nu\mathcal R_n}\,V_n(s,y).
\]
Therefore
\[
 e(t,x)=A_x^{1/4}u(t,x)
 =\Gamma_n^{3/2}\sqrt{\nu\mathcal R_n}\,
   \Lambda_y^{1/2}V_n(s,y).
\]
Since $\nu(b_n-t)\Gamma_n^2=-s$,
\[
 e^{-\nu(b_n-t)A_x}=e^{sA_y}
\]
after the spatial rescaling.  Applying $\Lambda_x=\Gamma_n\Lambda_y$ and
using
$\sqrt{\nu\mathcal R_n}\,\alpha_n=\mathcal R_n$
proves the first identity in \eqref{eq:v16-two-writer-scalings}.  The second
is \eqref{eq:v10-scaled-writer}.  Since
$\dd x\,\dd t=(\nu\Gamma_n^5)^{-1}\dd y\,\dd s$ and
$\Theta_n=\eta\chi\Omega_n$, substitution gives
\eqref{eq:v16-dimensionless-denominator}--
\eqref{eq:v16-dimensionless-mixed}.  Division proves
\eqref{eq:v16-kappa-fixed-card}.

For the uniform bound, the first-record property gives
$\|e(t)\|_2\le2\mathcal R_n$ for $t\le b_n$.  The norm scaling above yields
\[
 \alpha_n^{-1}
 \|\Lambda_y^{1/2}V_n(s)\|_2
 =\frac{\|e(t)\|_2}{\mathcal R_n}\le2.
\]
For $s\le-\tau_-$,
\[
 \|\Lambda_y e^{sA_y}\|_{2\to2}
 \le(2e(-s))^{-1/2}
 \le(2e\tau_-)^{-1/2}.
\]
Multiplication by the preceding factor two proves
\eqref{eq:v16-scaled-caloric-bound}.
\end{proof}

\begin{assessment}[What has and has not been evaluated]
\label{ass:v16-fixed-scale-status}
The theorem removes the changing physical scale from the mixed row.  It does
not assign a numerical value to $\mathfrak m_n$: that value depends on the
actual rescaled critical state on the selected terminal history.  The
supplied theorem documents do not contain samples of
$\mathscr Z_n^{\rm cal}$, $Z_n$, and $\Omega_n$ on one such hypothetical
terminal record.  The remaining task is therefore a finite physical Read or
a new PDE estimate for this displayed integral, not another normalization.
\end{assessment}

% =============================================================================
\section{The complete three-column Gram and its double Riesz short}
\label{sec:v16-three-Gram}
% =============================================================================

Define dimensionless vectors
\begin{equation}
 \widehat{\mathbf g}_n
 :=\frac{\mathbf g_n}{\sqrt\nu\,\mathcal R_n},
 \qquad
 \widehat{\mathbf w}_n:=\sqrt\nu\,\mathbf w_n,
 \qquad
 \widehat{\mathbf z}_n
 :=\frac{\sqrt\nu}{\mathcal R_n}\,\mathbf z_n.
 \label{eq:v16-normalized-three-vectors}
\end{equation}
Put
\begin{equation}
 \mathfrak a_n:=\|\widehat{\mathbf g}_n\|_2^2,
 \qquad
 \mathfrak e_n:=\|\widehat{\mathbf z}_n\|_2^2,
 \qquad
 \mathfrak c_n:=-\langle\widehat{\mathbf g}_n,
                              \widehat{\mathbf z}_n\rangle.
 \label{eq:v16-a-e-c}
\end{equation}
Thus
\begin{equation}
 \mathfrak a_n=\frac{E_n^g}{\nu\mathcal R_n^2},
 \qquad
 \mathfrak e_n=\frac{\nu E_n^z}{\mathcal R_n^2},
 \qquad
 \mathfrak c_n=\frac{\mathcal C_n^\Theta}{\mathcal R_n^2}.
 \label{eq:v16-a-e-c-physical}
\end{equation}

\begin{theorem}[Compact positive three-column terminal Gram]
\label{thm:v16-three-column-Gram}
The matrix
\begin{equation}
 \boxed{
 \mathbb G_n^{(3)}
 :=\begin{pmatrix}
 \mathfrak a_n&\sigma_*A_n&-\mathfrak c_n\\
 \sigma_*A_n&\mathfrak d_n&\mathfrak m_n\\
 -\mathfrak c_n&\mathfrak m_n&\mathfrak e_n
 \end{pmatrix}\succeq0}
 \label{eq:v16-three-Gram}
\end{equation}
is the Gram matrix of
$(\widehat{\mathbf g}_n,
  \widehat{\mathbf w}_n,
  \widehat{\mathbf z}_n)$.
Its entries satisfy the uniform bounds
\begin{equation}
 \boxed{
 \frac{h_0^2}{4d^*}\le\mathfrak a_n\le L_*,
 \qquad
 d_*\le\mathfrak d_n\le d^*,
 \qquad
 0\le\mathfrak e_n\le C_\eta,}
 \label{eq:v16-diagonal-compact-window}
\end{equation}
where $C_\eta$ is defined in \eqref{eq:v15-Kkappa}, and
\begin{equation}
 \boxed{
 |\mathfrak m_n|\le\sqrt{d^*C_\eta},
 \qquad
 |\mathfrak c_n|\le\sqrt{L_*C_\eta}.}
 \label{eq:v16-cross-compact-window}
\end{equation}
Every sequence of these finite cards therefore has a convergent subsequence
inside one compact positive-semidefinite matrix set.
\end{theorem}

\begin{proof}
The cross entries follow from
\eqref{eq:v16-known-cross-entries} and the definitions:
\[
 \langle\widehat{\mathbf g}_n,
        \widehat{\mathbf w}_n\rangle
 =-\frac{\mathfrak s_n}{\mathcal R_n}
 =\sigma_*A_n,
 \qquad
 \langle\widehat{\mathbf z}_n,
        \widehat{\mathbf w}_n\rangle
 =\frac{\nu N_n^{zw}}{\mathcal R_n}
 =\mathfrak m_n.
\]
Hence \eqref{eq:v16-three-Gram} is a Gram matrix and is positive.
The upper bound for $\mathfrak a_n$ is
\eqref{eq:v16-source-action-window}; the upper bound for
$\mathfrak e_n$ is \eqref{eq:v15-caloric-writer-upper}; and the bounds for
$\mathfrak d_n$ are \eqref{eq:v16-denominator-window}.  Since
\[
 A_n^2
 =|\langle\widehat{\mathbf g}_n,
              \widehat{\mathbf w}_n\rangle|^2
 \le\mathfrak a_n\mathfrak d_n,
\]
$A_n\ge h_0/2$ and $\mathfrak d_n\le d^*$ give the lower bound for
$\mathfrak a_n$.  The cross bounds are Cauchy--Schwarz followed by the
diagonal bounds.  Compactness is finite dimensional.
\end{proof}

Short the complete Gram by the middle writer column.  Define
\begin{equation}
 \widehat{\mathbf g}_n^\perp
 :=\widehat{\mathbf g}_n
   -\frac{\sigma_*A_n}{\mathfrak d_n}
    \widehat{\mathbf w}_n,
 \qquad
 \widehat{\mathbf z}_n^\perp
 :=\widehat{\mathbf z}_n
   -\frac{\mathfrak m_n}{\mathfrak d_n}
    \widehat{\mathbf w}_n.
 \label{eq:v16-double-residual-vectors}
\end{equation}
The second residual is exactly
$(\sqrt\nu/\mathcal R_n)\Theta_n^{1/2}Z_{n,\cdot}^{\perp}$.

\begin{theorem}[Double Riesz short and the exact writer-orthogonal determinant]
\label{thm:v16-double-Riesz}
The residual Gram is
\begin{equation}
 \boxed{
 \begin{pmatrix}
 \mathfrak a_n-A_n^2/\mathfrak d_n
 &-\mathcal U_n^{\rm wr}/\mathcal R_n^2\\
 -\mathcal U_n^{\rm wr}/\mathcal R_n^2
 &\mathfrak e_n-\mathfrak m_n^2/\mathfrak d_n
 \end{pmatrix}\succeq0.}
 \label{eq:v16-residual-Gram}
\end{equation}
Equivalently,
\begin{equation}
 \boxed{
 \frac{\mathcal U_n^{\rm wr}}{\mathcal R_n^2}
 =\mathfrak c_n
  +\sigma_*A_n\frac{\mathfrak m_n}{\mathfrak d_n},}
 \label{eq:v16-Uwr-readout}
\end{equation}
and
\begin{equation}
 \boxed{
 \left(\frac{\mathcal U_n^{\rm wr}}{\mathcal R_n^2}\right)^2
 \le
 \left(\mathfrak a_n-\frac{A_n^2}{\mathfrak d_n}\right)
 \left(\mathfrak e_n-\frac{\mathfrak m_n^2}{\mathfrak d_n}\right).}
 \label{eq:v16-Uwr-determinant}
\end{equation}
The determinant of the full three-column Gram is
\begin{equation}
 \boxed{
 \det\mathbb G_n^{(3)}
 =\mathfrak d_n
 \left[
 \left(\mathfrak a_n-\frac{A_n^2}{\mathfrak d_n}\right)
 \left(\mathfrak e_n-\frac{\mathfrak m_n^2}{\mathfrak d_n}\right)
 -\left(\frac{\mathcal U_n^{\rm wr}}{\mathcal R_n^2}\right)^2
 \right].}
 \label{eq:v16-full-Gram-determinant}
\end{equation}
Thus the writer-orthogonal contact is not an additional free scalar: it is
the off-diagonal entry of the positive Schur residual left after both source
and caloric writer are shorted by the same selected writer.
\end{theorem}

\begin{proof}
Both vectors in \eqref{eq:v16-double-residual-vectors} are orthogonal to
$\widehat{\mathbf w}_n$.  Their squared norms are the diagonal entries of
\eqref{eq:v16-residual-Gram}.  Their cross pairing is
\[
 \langle\widehat{\mathbf g}_n^\perp,
        \widehat{\mathbf z}_n^\perp\rangle
 =-\mathfrak c_n
  -\sigma_*A_n\frac{\mathfrak m_n}{\mathfrak d_n}.
\]
On the other hand,
$\mathcal U_n^{\rm wr}=-\langle\mathbf g_n,
\Theta_n^{1/2}Z_n^\perp\rangle$ and the component of $\mathbf g_n$ parallel
to $\mathbf w_n$ is orthogonal to $Z_n^\perp$.  Rescaling therefore makes
the preceding cross pairing equal to
$-\mathcal U_n^{\rm wr}/\mathcal R_n^2$, proving
\eqref{eq:v16-Uwr-readout} and the residual matrix.  Positivity and the
two-by-two determinant criterion give
\eqref{eq:v16-Uwr-determinant}.  The block determinant formula, or direct
Schur complementation of the positive middle entry $\mathfrak d_n$, gives
\eqref{eq:v16-full-Gram-determinant}.
\end{proof}

Define the two Riesz efficiencies
\begin{equation}
 \eta_n^{\rm sw}
 :=\frac{A_n^2}{\mathfrak a_n\mathfrak d_n},
 \qquad
 \eta_n^{\rm tw}
 :=\frac{\mathfrak m_n^2}{\mathfrak d_n\mathfrak e_n}
 \quad(\mathfrak e_n>0).
 \label{eq:v16-two-efficiencies}
\end{equation}
The second is exactly the terminal-writer efficiency of
\eqref{eq:v15-writer-efficiency}.

\begin{corollary}[A uniform source--writer angle and a sharp orthogonal-contact fork]
\label{cor:v16-efficiency-fork}
For every selected card,
\begin{equation}
 \boxed{
 \eta_n^{\rm sw}
 \ge\eta_{\rm sw,*}
 :=\frac{h_0^2}{4L_*d^*}>0.}
 \label{eq:v16-source-writer-efficiency-floor}
\end{equation}
Whenever $\mathfrak e_n>0$,
\begin{equation}
 \boxed{
 \frac{|\mathcal U_n^{\rm wr}|}{\mathcal R_n^2}
 \le
 \sqrt{\mathfrak a_n\mathfrak e_n}
 \sqrt{(1-\eta_n^{\rm sw})(1-\eta_n^{\rm tw})}.}
 \label{eq:v16-Uwr-efficiency}
\end{equation}
Consequently:
\begin{enumerate}[label=\textup{(\roman*)},leftmargin=3em]
\item if $\eta_n^{\rm sw}\to1$ and $\eta_n^{\rm tw}\to1$, then
      $\mathcal U_n^{\rm wr}/\mathcal R_n^2\to0$;
\item if $|\mathcal U_n^{\rm wr}|\ge\delta\mathcal R_n^2$ for some fixed
      $\delta>0$, then at least one of
      \begin{equation}
       1-\eta_n^{\rm sw},
       \qquad 1-\eta_n^{\rm tw}
       \label{eq:v16-angle-defects}
      \end{equation}
      is at least
      $\delta/\sqrt{L_*C_\eta}$.
\end{enumerate}
Thus a writer-orthogonal contact of record size carries a positive angle
defect in one of the two actual Riesz projections.
\end{corollary}

\begin{proof}
Use $A_n\ge h_0/2$, $\mathfrak a_n\le L_*$, and
$\mathfrak d_n\le d^*$ for the first estimate.  Factor the two terms on the
right of \eqref{eq:v16-Uwr-determinant} as
$\mathfrak a_n(1-\eta_n^{\rm sw})$ and
$\mathfrak e_n(1-\eta_n^{\rm tw})$.  This proves
\eqref{eq:v16-Uwr-efficiency}.  The first consequence is immediate.  For the
second, $\mathfrak a_n\mathfrak e_n\le L_*C_\eta$ gives
\[
 (1-\eta_n^{\rm sw})(1-\eta_n^{\rm tw})
 \ge\frac{\delta^2}{L_*C_\eta}.
\]
If both factors were smaller than
$\delta/\sqrt{L_*C_\eta}$, their product would be smaller than the displayed
lower bound.
\end{proof}

% =============================================================================
\section{The one-sided terminal orientation and its exact obstruction}
\label{sec:v16-orientation}
% =============================================================================

The V13 orientation $\sigma_*$ declares the desired terminal-writer ray to be
\begin{equation}
 \mathcal K_n^{\rm des}
 :=\{-\sigma_*\lambda\widehat{\mathbf w}_n:\lambda\ge0\}.
 \label{eq:v16-desired-ray}
\end{equation}

\begin{theorem}[Projection onto the physically oriented writer ray]
\label{thm:v16-oriented-ray}
The unique metric projection of $\widehat{\mathbf z}_n$ onto
$\mathcal K_n^{\rm des}$ has coefficient
\begin{equation}
 \boxed{
 \lambda_n^{\rm des}
 =\left(-\sigma_*\frac{\mathfrak m_n}{\mathfrak d_n}\right)_+.}
 \label{eq:v16-desired-coefficient}
\end{equation}
Its squared residual is
\begin{equation}
 \boxed{
 \inf_{\lambda\ge0}
 \|\widehat{\mathbf z}_n+\sigma_*\lambda
      \widehat{\mathbf w}_n\|_2^2
 =\mathfrak e_n-\frac{\mathfrak m_n^2}{\mathfrak d_n}
  +\frac{(\sigma_*\mathfrak m_n)_+^2}{\mathfrak d_n}.}
 \label{eq:v16-oriented-residual}
\end{equation}
The second positive term is exactly the wrong-sign aligned component; the
first is the writer-orthogonal component.
\end{theorem}

\begin{proof}
Expand
\[
 \|\widehat{\mathbf z}_n+\sigma_*\lambda
      \widehat{\mathbf w}_n\|_2^2
 =\mathfrak e_n+2\sigma_*\lambda\mathfrak m_n
  +\lambda^2\mathfrak d_n.
\]
The unconstrained minimizer is
$-\sigma_*\mathfrak m_n/\mathfrak d_n$; projection onto the half-line
$[0,\infty)$ gives \eqref{eq:v16-desired-coefficient}.  Substitution gives
\eqref{eq:v16-oriented-residual}.
\end{proof}

Define the normalized terminal-writer calibration defect
\begin{equation}
 \boxed{
 \delta_n^{\rm tw}
 :=\frac{\mathcal C_n^{\rm cal}-\mathcal C_n^{\rm sel}}
          {\mathcal R_n^2}
 =\frac{\mathcal C_n^{\rm cal}}{\mathcal R_n^2}
  +\sigma_*A_n\frac{\mathfrak m_n}{\mathfrak d_n}.}
 \label{eq:v16-terminal-calibration-defect}
\end{equation}
This is the complete signed source-side complement, not merely the
least-square residual norm.

\begin{corollary}[Wrong sign or calibration deficit forces a typed full-contact complement]
\label{cor:v16-calibration-obstruction}
If $\sigma_*\mathfrak m_n\ge0$, then
\begin{equation}
 \boxed{
 \mathcal C_n^{\rm sel}\le0,
 \qquad
 \delta_n^{\rm tw}
 \ge\frac{\mathcal C_n^{\rm cal}}{\mathcal R_n^2}
 \ge\frac32.}
 \label{eq:v16-wrong-sign-complement}
\end{equation}
More generally, if
$|\delta_n^{\rm tw}|\ge\delta>0$, then on a further subsequence at least one
of
\begin{equation}
 \mathcal C_n^{\rm res},\qquad
 \mathcal U_n^{\rm out},\qquad
 \mathcal U_n^{\rm route},\qquad
 \mathcal U_n^{\rm wr}
 \label{eq:v16-four-complements}
\end{equation}
has magnitude at least
$(\delta/4)\mathcal R_n^2$.  Thus a sign change or collapse of the mixed
writer does not create a vague compactness failure; it returns one of the
four signed complements already present in the full source ledger.
\end{corollary}

\begin{proof}
By \eqref{eq:v16-kappa-fixed-card},
\[
 \frac{\mathcal C_n^{\rm sel}}{\mathcal R_n^2}
 =-\sigma_*A_n\frac{\mathfrak m_n}{\mathfrak d_n}.
\]
Since $A_n>0$ and $\mathfrak d_n>0$, the first assertion follows.  The source
ledger gives
\[
 \mathcal C_n^{\rm cal}-\mathcal C_n^{\rm sel}
 =\mathcal C_n^{\rm res}
  +\mathcal U_n^{\rm out}
  +\mathcal U_n^{\rm route}
  +\mathcal U_n^{\rm wr}.
\]
If the left side has magnitude at least
$\delta\mathcal R_n^2$, the triangle inequality and a finite priority
subsequence select one of the four terms with magnitude at least one quarter
of that amount.
\end{proof}

% =============================================================================
\section{The mixed row is an ordered Duhamel moment}
\label{sec:v16-chronological-moment}
% =============================================================================

The finite coefficient has a chronological meaning which is lost if only the
endpoint response is retained.  Put
\begin{equation}
 X_n:=S_ne(\beta_n),
 \qquad
 \mathcal F_n(t)
 :=-\int_{\beta_n}^{t}
   \Lambda e^{-\nu(b_n-s)A}c(s)\,\dd s.
 \label{eq:v16-partial-endpoint-response}
\end{equation}
Thus $\mathcal F_n(b_n)=\mathcal F_n$ is the complete endpoint Duhamel
response of Version~V15.

\begin{theorem}[Carried state plus cumulative common-endpoint response]
\label{thm:v16-cumulative-response}
For every $t\in[\beta_n,b_n]$,
\begin{equation}
 \boxed{
 e^{-\nu(b_n-t)A}e(t)=X_n+\mathcal F_n(t).}
 \label{eq:v16-carried-plus-cumulative}
\end{equation}
Consequently the mixed numerator has the exact split
\begin{equation}
 \boxed{
 N_n^{zw}=N_n^{\rm old}+N_n^{\rm mom},}
 \label{eq:v16-mixed-moment-split}
\end{equation}
where
\begin{align}
 N_n^{\rm old}
 &:=\int_{\beta_n}^{b_n}
   \langle\Lambda X_n,\Theta_n(t)W_n\rangle\,\dd t,
 \label{eq:v16-old-writer-moment}\\
 N_n^{\rm mom}
 &:=\int_{\beta_n}^{b_n}
   \langle\Lambda\mathcal F_n(t),
          \Theta_n(t)W_n\rangle\,\dd t.
 \label{eq:v16-new-writer-moment}
\end{align}
The second term is the ordered source moment
\begin{equation}
 \boxed{
 N_n^{\rm mom}
 =-\int_{\beta_n}^{b_n}
   \langle c(s),Q_n^{\rm mom}(s)\rangle\,\dd s,}
 \label{eq:v16-source-moment-pairing}
\end{equation}
with the deterministic follower writer
\begin{equation}
 \boxed{
 Q_n^{\rm mom}(s)
 :=\int_s^{b_n}
   e^{-\nu(b_n-s)A}\Lambda^2
   \bigl(\Theta_n(t)W_n\bigr)\,\dd t.}
 \label{eq:v16-moment-writer}
\end{equation}
No new source occurrence is present: $Q_n^{\rm mom}$ is constructed from the
already selected writer, cutoff, and common endpoint.
\end{theorem}

\begin{proof}
Differentiate
$Y_n(t)=e^{-\nu(b_n-t)A}e(t)$.  Since
$e_t=-\nu Ae-\Lambda c$ and the derivative of the heat factor is
$\nu Ae^{-\nu(b_n-t)A}$,
\[
 Y_n'(t)=-\Lambda e^{-\nu(b_n-t)A}c(t).
\]
Integrating from $\beta_n$ to $t$ gives
\eqref{eq:v16-carried-plus-cumulative}.  Since
$Z_{n,t}^{\rm cal}=\Lambda Y_n(t)$, insertion into
$N_n^{zw}=\int\langle Z_{n,t}^{\rm cal},\Theta_nW_n\rangle\,\dd t$
gives \eqref{eq:v16-mixed-moment-split}.  Finally,
\[
 \Lambda\mathcal F_n(t)
 =-\int_{\beta_n}^{t}
   \Lambda^2e^{-\nu(b_n-s)A}c(s)\,\dd s.
\]
Fubini's theorem on the smooth finite record interval and self-adjointness of
$e^{-\nu(b_n-s)A}\Lambda^2$ give
\eqref{eq:v16-source-moment-pairing}--
\eqref{eq:v16-moment-writer}.
\end{proof}

\begin{countertheorem}[Endpoint response and source length do not determine the ordered moment]
\label{cth:v16-endpoint-no-moment}
Even in a one-dimensional real Hilbert model with unit generator, unit writer,
zero carried state, and unit cutoff, the endpoint response and total source
length do not determine $N^{\rm mom}$.  On $[0,1]$, let
\begin{equation}
 \mathcal F_p(t)=t^p,
 \qquad p\in\{1,2\}.
 \label{eq:v16-scalar-cumulative-family}
\end{equation}
Then
\begin{equation}
 \mathcal F_p(1)=1,
 \qquad
 \int_0^1|\mathcal F_p'(t)|\,\dd t=1,
 \label{eq:v16-same-endpoint-length}
\end{equation}
while
\begin{equation}
 \boxed{
 N_p^{\rm mom}=\int_0^1\mathcal F_p(t)\,\dd t
 =\frac1{p+1}.}
 \label{eq:v16-different-moments}
\end{equation}
Thus the mixed terminal-writer Read contains chronological information not
present in the endpoint vector, its norm, or the unweighted source length.
This is a functional independence statement, not a construction of two
Navier--Stokes trajectories.
\end{countertheorem}

\begin{proof}
All identities follow by direct integration.
\end{proof}

% =============================================================================
\section{Positive polarization of the mixed writer and the represented endpoint}
\label{sec:v16-positive-acquisition}
% =============================================================================

The sign of the mixed row can be read from positive quadratic quantities once
the physical calibration $\mathcal R_n$ has been fixed.  Define
\begin{equation}
 \mathcal E_n^{\pm}
 :=\int_{\beta_n}^{b_n}\!\!\int_{\Torus^3}
   \Theta_n
   |Z_{n,t}^{\rm cal}\pm\mathcal R_nW_n|^2
   \,\dd x\,\dd t.
 \label{eq:v16-balanced-positive-actions}
\end{equation}

\begin{theorem}[Balanced positive acquisition of the terminal-writer cross row]
\label{thm:v16-positive-polarization}
The two positive actions in
\eqref{eq:v16-balanced-positive-actions} satisfy
\begin{align}
 \mathcal E_n^+-\mathcal E_n^-
 &=4\mathcal R_nN_n^{zw},
 \label{eq:v16-positive-difference}\\
 \mathcal E_n^++\mathcal E_n^-
 &=2E_n^z+2\mathcal R_n^2D_n^W.
 \label{eq:v16-positive-sum}
\end{align}
Hence
\begin{equation}
 \boxed{
 \mathfrak m_n
 =\frac{\nu}{4\mathcal R_n^2}
  (\mathcal E_n^+-\mathcal E_n^-),
 \qquad
 \frac{\kappa_n}{\mathcal R_n}
 =\frac{\nu(\mathcal E_n^+-\mathcal E_n^-)}
        {4\mathcal R_n^2\mathfrak d_n}.}
 \label{eq:v16-positive-kappa-read}
\end{equation}
In particular the desired sign $-\sigma_*\kappa_n>0$ is equivalent to
\begin{equation}
 \boxed{
 -\sigma_*(\mathcal E_n^+-\mathcal E_n^-)>0.}
 \label{eq:v16-positive-sign-test}
\end{equation}
The two actions are polarization queries of one already occurring writer
pair; they are not two physical source births.
\end{theorem}

\begin{proof}
Expand the two squares.  The diagonal terms cancel in the difference and add
in the sum.  Substitute the definitions of $\mathfrak m_n$ and
$\mathfrak d_n$.
\end{proof}

For the represented endpoint, recall
$\mathcal F_n^{\rm rep}=P_n^{\rm rep}\mathcal F_n$ and
$X_n=S_ne(\beta_n)$.  Define
\begin{equation}
 \mathcal E_n^{\rm carry}:=\|X_n\|_2^2,
 \qquad
 \mathcal E_n^{\rm rep,+}:=
   \|X_n+\mathcal F_n^{\rm rep}\|_2^2.
 \label{eq:v16-represented-positive-reads}
\end{equation}

\begin{proposition}[Positive endpoint read of the represented contact]
\label{prop:v16-represented-positive-read}
The complete and represented contacts are
\begin{equation}
 \boxed{
 \mathcal C_n^{\rm cal}
 =\frac12\bigl(\|e(b_n)\|_2^2-
                    \mathcal E_n^{\rm carry}\bigr),
 \qquad
 \mathcal C_n^{\rm rep}
 =\frac12\bigl(\mathcal E_n^{\rm rep,+}-
                    \mathcal E_n^{\rm carry}\bigr).}
 \label{eq:v16-two-positive-contact-reads}
\end{equation}
Thus the represented branch is evaluated before freshness by one deterministic
endpoint-head norm difference.  The vector
$X_n+\mathcal F_n^{\rm rep}$ is a represented-head synthesis, not asserted to
be the actual endpoint state when a held-out response remains.
\end{proposition}

\begin{proof}
The first identity is \eqref{eq:v15-Duhamel-secant}.  Expanding the second
norm gives
\[
 \frac12(\|X_n+\mathcal F_n^{\rm rep}\|_2^2-\|X_n\|_2^2)
 =\langle X_n,\mathcal F_n^{\rm rep}\rangle
  +\frac12\|\mathcal F_n^{\rm rep}\|_2^2
 =\mathcal C_n^{\rm rep}.
\]
\end{proof}

The writer-orthogonal cross row can also be obtained from positive actions.
Put
\begin{equation}
 \mathcal P_n^\pm
 :=\|\widehat{\mathbf g}_n
       \pm\widehat{\mathbf z}_n\|_2^2.
 \label{eq:v16-source-caloric-polarizations}
\end{equation}
Then
\begin{equation}
 \boxed{
 \mathfrak c_n=-\frac14(\mathcal P_n^+-\mathcal P_n^-),
 \qquad
 \frac{\mathcal U_n^{\rm wr}}{\mathcal R_n^2}
 =-\frac14(\mathcal P_n^+-\mathcal P_n^-)
  +\sigma_*A_n\frac{\mathfrak m_n}{\mathfrak d_n}.}
 \label{eq:v16-positive-Uwr-read}
\end{equation}
Hence every entry of the three-column Gram can be reconstructed from diagonal
positive actions and balanced polarization queries.

% =============================================================================
\section{The five-scalar decision card}
\label{sec:v16-five-card}
% =============================================================================

Define
\begin{equation}
 \mathfrak c_n^{\rm cal}
 :=\frac{\mathcal C_n^{\rm cal}}{\mathcal R_n^2},
 \qquad
 \mathfrak c_n^{\rm rep}
 :=\frac{\mathcal C_n^{\rm rep}}{\mathcal R_n^2},
 \label{eq:v16-normalized-contact-reads}
\end{equation}
and form the five-scalar packet
\begin{equation}
 \boxed{
 \mathbf K_n^{(16)}
 :=\bigl(
 \mathfrak c_n^{\rm cal},
 \mathfrak c_n^{\rm rep},
 A_n,
 \mathfrak d_n,
 \mathfrak m_n
 \bigr).}
 \label{eq:v16-five-card}
\end{equation}

\begin{theorem}[Exact readout of selected contact and source-to-head incidence]
\label{thm:v16-five-card-readout}
The five values in \eqref{eq:v16-five-card} determine the normalized selected
contact and head-incidence residual exactly:
\begin{align}
 \boxed{
 \mathfrak c_n^{\rm sel}
 :=\frac{\mathcal C_n^{\rm sel}}{\mathcal R_n^2}
 =-\sigma_*A_n\frac{\mathfrak m_n}{\mathfrak d_n},}
 \label{eq:v16-selected-readout}\\
 \boxed{
 \mathfrak i_n^{\rm head}
 :=\frac{\mathcal I_n^{\rm head}}{\mathcal R_n^2}
 =\mathfrak c_n^{\rm cal}-\mathfrak c_n^{\rm rep}
  +\sigma_*A_n\frac{\mathfrak m_n}{\mathfrak d_n}.}
 \label{eq:v16-incidence-readout}
\end{align}
They also determine
\begin{equation}
 \frac{\mathcal C_n^{\rm ho}}{\mathcal R_n^2}
 =\mathfrak c_n^{\rm cal}-\mathfrak c_n^{\rm rep}.
 \label{eq:v16-heldout-readout}
\end{equation}
No resolved, output, route, or writer-orthogonal term is estimated in order to
decide whether the aggregate incidence residual vanishes.
\end{theorem}

\begin{proof}
Use
$\mathfrak s_n=-\sigma_*\mathcal R_nA_n$,
$\kappa_n/\mathcal R_n=\mathfrak m_n/\mathfrak d_n$, and
$\mathcal C_n^{\rm sel}=\kappa_n\mathfrak s_n$ for
\eqref{eq:v16-selected-readout}.  The definition
$\mathcal I_n^{\rm head}=\mathcal C_n^{\rm ho}-
\mathcal C_n^{\rm sel}$ and the endpoint ledger
$\mathcal C_n^{\rm ho}=\mathcal C_n^{\rm cal}-
\mathcal C_n^{\rm rep}$ prove the other formulas.
\end{proof}

\begin{theorem}[Compactness of the five-scalar physical card]
\label{thm:v16-five-card-compact}
The packets \eqref{eq:v16-five-card} lie in the fixed compact set determined
by
\begin{equation}
 \boxed{
 \frac32\le\mathfrak c_n^{\rm cal}\le2,
 \qquad
 |\mathfrak c_n^{\rm rep}|\le\frac{15}{2},
 \qquad
 \frac{h_0}{2}\le A_n\le\mathfrak V_*,}
 \label{eq:v16-five-card-first-bounds}
\end{equation}
\begin{equation}
 \boxed{
 d_*\le\mathfrak d_n\le d^*,
 \qquad
 |\mathfrak m_n|\le M_*:=\sqrt{d^*C_\eta}.}
 \label{eq:v16-five-card-second-bounds}
\end{equation}
Consequently every selected sequence has a subsequence on which all V16
branch readouts converge as ordinary real numbers.
\end{theorem}

\begin{proof}
The full-contact bounds are \eqref{eq:v15-full-contact-window}; the bounds for
$A_n$, $\mathfrak d_n$, and $\mathfrak m_n$ have already been proved.  Let
$X_n=S_ne(\beta_n)$.  Heat contraction gives $\|X_n\|_2\le\mathcal R_n$,
and $\|e(b_n)\|_2=2\mathcal R_n$.  Since
$\mathcal F_n=e(b_n)-X_n$,
\[
 \|\mathcal F_n\|_2\le3\mathcal R_n,
 \qquad
 \|\mathcal F_n^{\rm rep}\|_2\le3\mathcal R_n.
\]
Therefore
\[
 |\mathcal C_n^{\rm rep}|
 \le\|X_n\|_2\|\mathcal F_n^{\rm rep}\|_2
    +\frac12\|\mathcal F_n^{\rm rep}\|_2^2
 \le\frac{15}{2}\mathcal R_n^2.
\]
The result follows from finite-dimensional compactness and continuity of the
readout maps on $\mathfrak d\ge d_*>0$.
\end{proof}

\begin{theorem}[Uniform stability of the finite branch certificate]
\label{thm:v16-card-stability}
Let
$K=(c,r,a,d,m)$ and $K'=(c',r',a',d',m')$ be two packets satisfying the
bounds of \Cref{thm:v16-five-card-compact}.  Put
\begin{equation}
 S(K):=-\sigma_*\frac{am}{d},
 \qquad
 I(K):=c-r-S(K).
 \label{eq:v16-readout-map}
\end{equation}
Then
\begin{equation}
 \boxed{
 |S(K)-S(K')|
 \le
 \frac{M_*}{d_*}|a-a'|
 +\frac{\mathfrak V_*}{d_*}|m-m'|
 +\frac{\mathfrak V_*M_*}{d_*^2}|d-d'|,}
 \label{eq:v16-selected-Lipschitz}
\end{equation}
and
\begin{equation}
 \boxed{
 |I(K)-I(K')|
 \le |c-c'|+|r-r'|+|S(K)-S(K')|.}
 \label{eq:v16-incidence-Lipschitz}
\end{equation}
Thus any outward interval evaluation whose error is smaller than a displayed
branch margin gives the same represented, incidence-residual, or clean verdict
for every exact card inside that interval.
\end{theorem}

\begin{proof}
Insert and subtract $a'm/d$ and $a'm'/d$:
\[
 \left|\frac{am}{d}-\frac{a'm'}{d'}\right|
 \le\frac{|m|}{d}|a-a'|
  +\frac{|a'|}{d}|m-m'|
  +|a'm'|\left|\frac1d-\frac1{d'}\right|.
\]
Use $|m|,|m'|\le M_*$, $|a'|\le\mathfrak V_*$, and
$d,d'\ge d_*$.  This proves \eqref{eq:v16-selected-Lipschitz}; the second
estimate is the triangle inequality.
\end{proof}

% =============================================================================
\section{Direct finite branch theorem and post-verdict typing}
\label{sec:v16-direct-branch}
% =============================================================================

\begin{theorem}[The V15 priority alternative is a five-scalar finite verdict]
\label{thm:v16-direct-verdict}
After passage to a subsequence, exactly one of the following priority branches
occurs.
\begin{enumerate}[label=\textup{(V16.\arabic*)},leftmargin=5.7em]
\item \emph{Represented endpoint return.}  There is $\delta>0$ such that
      \begin{equation}
       \boxed{
       |\mathfrak c_n^{\rm rep}|
       \ge\delta\mathfrak c_n^{\rm cal}.}
       \label{eq:v16-represented-verdict}
      \end{equation}
      The represented endpoint synthesis is returned to its original source
      ancestry before any selected response is called fresh.
\item \emph{Assembled source-to-head residual.}  The first branch fails and
      there is $\delta>0$ such that
      \begin{equation}
       \boxed{
       |\mathfrak i_n^{\rm head}|
       \ge\delta\mathfrak c_n^{\rm cal}.}
       \label{eq:v16-incidence-verdict}
      \end{equation}
      This verdict is obtained directly from
      \eqref{eq:v16-incidence-readout}, before the four signed complement
      terms are opened.
\item \emph{Selected critical-replenishment export.}  The preceding two
      branches fail, so
      \begin{equation}
       \boxed{
       \mathfrak c_n^{\rm rep}\longrightarrow0,
       \qquad
       \mathfrak i_n^{\rm head}\longrightarrow0,}
       \label{eq:v16-clean-five-card}
      \end{equation}
      and
      \begin{equation}
       \boxed{
       -\sigma_*A_n\frac{\mathfrak m_n}{\mathfrak d_n}
       -\mathfrak c_n^{\rm cal}\longrightarrow0.}
       \label{eq:v16-clean-card-relation}
      \end{equation}
      In particular the selected contact is positive and asymptotic to the
      complete contact.
\end{enumerate}
The three branches are decided entirely by the card
$\mathbf K_n^{(16)}$.  No fourth branch is generated by uncertainty in the
individual source-side terms.
\end{theorem}

\begin{proof}
Select a subsequence on which
$|\mathfrak c_n^{\rm rep}|/\mathfrak c_n^{\rm cal}$ either has a positive
lower bound or tends to zero.  In the second case, do the same for
$|\mathfrak i_n^{\rm head}|/\mathfrak c_n^{\rm cal}$.  Since
$3/2\le\mathfrak c_n^{\rm cal}\le2$, failure of both positive-lower-bound
branches gives \eqref{eq:v16-clean-five-card}.  The exact readout
\eqref{eq:v16-incidence-readout} then gives
\eqref{eq:v16-clean-card-relation}.  This is the same exhaustive priority
logic as Version~V15, now expressed on the finite physical card.
\end{proof}

Only after branch \textup{(V16.2)} has been certified do we type its source.
Define the cumulative signed splitter
\begin{equation}
 \begin{aligned}
 B_{n,0}&:=0,\\
 B_{n,1}&:=\mathcal C_n^{\rm res},\\
 B_{n,2}&:=\mathcal C_n^{\rm res}+\mathcal U_n^{\rm out},\\
 B_{n,3}&:=\mathcal C_n^{\rm res}+\mathcal U_n^{\rm out}
                         +\mathcal U_n^{\rm route},\\
 B_{n,4}&:=\mathcal C_n^{\rm res}+\mathcal U_n^{\rm out}
                         +\mathcal U_n^{\rm route}
                         +\mathcal U_n^{\rm wr}.
 \end{aligned}
 \label{eq:v16-cumulative-splitter}
\end{equation}
Then
\begin{equation}
 B_{n,4}=\mathcal I_n^{\rm head}+\mathcal C_n^{\rm rep}.
 \label{eq:v16-splitter-total}
\end{equation}

\begin{corollary}[Post-verdict signed typing without prior absolute estimates]
\label{cor:v16-post-verdict-typing}
On branch \textup{(V16.2)}, for all late $n$,
\begin{equation}
 |B_{n,4}|\ge\frac\delta2\mathcal C_n^{\rm cal}.
 \label{eq:v16-splitter-lower}
\end{equation}
Consequently one increment
$B_{n,j}-B_{n,j-1}$ has magnitude at least
$(\delta/8)\mathcal C_n^{\rm cal}$.  The four increments are, in order,
resolved heat--Euler contact, unselected covariance output,
material/cell/time escape, and writer-orthogonal work.  Their absolute values
are not summed before this aggregate verdict.  If the fourth increment is
selected, \eqref{eq:v16-positive-Uwr-read} and
\eqref{eq:v16-residual-Gram} give its complete positive Gram certificate.
\end{corollary}

\begin{proof}
The represented branch has failed, so
$|\mathcal C_n^{\rm rep}|=o(\mathcal C_n^{\rm cal})$.  Combine this with
\eqref{eq:v16-incidence-verdict} and
\eqref{eq:v16-splitter-total}.  The four-increment conclusion is the triangle
inequality followed by a finite priority subsequence.  The final statement
uses the positive polarization and Schur formulas already proved.
\end{proof}

% =============================================================================
\section{The exact NS--A export card}
\label{sec:v16-export}
% =============================================================================

\begin{definition}[Selected critical-replenishment card]
\label{def:v16-export-card}
On branch \textup{(V16.3)}, retain
\begin{equation}
 \boxed{
 \begin{aligned}
 \mathfrak C_n^{\rm V16}:=\bigl(&b_n,I_n^{\rm ann},\Gamma_n,
 \Theta_n,P_n^{\rm rep},
 \Psi_ng_{n,\cdot}^{\rm cov},W_n,Z_{n,\cdot}^{\rm cal},\\
 &\mathbb G_n^{(3)},
 \mathcal C_n^{\rm cal},\mathcal C_n^{\rm rep},
 \mathfrak s_n,\kappa_n,
 \mathcal C_n^{\rm sel},\mathcal I_n^{\rm head}\bigr).
 \end{aligned}}
 \label{eq:v16-export-card}
\end{equation}
Every entry is attached to the same fixed datum, record interval, selected
annulus, material--caloric route, spatial cutoff, and final endpoint.
\end{definition}

\begin{theorem}[Clean-card export and termination of NS--A on that subsequence]
\label{thm:v16-clean-export}
On branch \textup{(V16.3)},
\begin{equation}
 \boxed{
 \frac{\mathcal C_n^{\rm rep}}{\mathcal R_n^2}\to0,
 \qquad
 \frac{\mathcal I_n^{\rm head}}{\mathcal R_n^2}\to0,
 \qquad
 \frac{\mathcal C_n^{\rm sel}}{\mathcal R_n^2}
 -\frac{\mathcal C_n^{\rm cal}}{\mathcal R_n^2}\to0.}
 \label{eq:v16-clean-export-limits}
\end{equation}
Moreover,
\begin{equation}
 \boxed{
 -\sigma_*\frac{\mathfrak m_n}{\mathfrak d_n}
 =-\sigma_*\frac{\kappa_n}{\mathcal R_n}
 \ge\frac1{\mathfrak V_*}}
 \label{eq:v16-export-alignment}
\end{equation}
for all late $n$, the source--writer efficiency has the fixed floor
\eqref{eq:v16-source-writer-efficiency-floor}, and the terminal-writer
efficiency has the V15 floor \eqref{eq:v15-writer-efficiency}.  Finally,
\begin{equation}
 \boxed{
 \mathcal C_n^{\rm sel}\ge\frac34\mathcal R_n^2,
 \qquad
 \sum_n\mathcal C_n^{\rm sel}=+\infty.}
 \label{eq:v16-export-replenishment}
\end{equation}
The packet \eqref{eq:v16-export-card} is therefore exported to the
fixed-datum Reynolds-covariance and paid-stress programme.  No further
bounded-action NS--A lower estimate is licensed on this subsequence.  The
remaining question is an independently occurring finite-stock or monotone
upper budget for this same calibrated contact.
\end{theorem}

\begin{proof}
The first three limits are \eqref{eq:v16-clean-five-card} and
\eqref{eq:v16-clean-card-relation}.  Since
$\mathfrak c_n^{\rm cal}\ge3/2$ and $A_n\le\mathfrak V_*$,
\eqref{eq:v16-clean-card-relation} gives
$-\sigma_*\mathfrak m_n/\mathfrak d_n\ge1/\mathfrak V_*$ for all late $n$.
The two efficiency statements are
\Cref{cor:v16-efficiency-fork} and
\Cref{cor:v15-clean-quantitative}.  The contact lower bound and divergence
are \Cref{cor:v15-nonsummable-contact}.  All quantities have already been
placed on the same history by their definitions, so the export adds no new
incidence assumption.
\end{proof}

\begin{firewall}[The finite card does not supply its own data]
\label{fire:v16-no-populated-card}
The formulas above are an exact acquisition and decision theorem.  They do
not manufacture the values of the five Reads on a hypothetical singular
Navier--Stokes trajectory.  In particular, compactness of the allowed card
set is not convergence of the underlying fields, and positivity of the
three-column Gram is not a finite-stock upper bound.  Without populated
same-history actions or a new PDE estimate for
\eqref{eq:v16-dimensionless-mixed}, the correct verdict is that the finite
card is unexecuted, not that one of its physical branches has passed.
\end{firewall}

% =============================================================================
\section{Version V16 integrated theorem and exact stopping boundary}
\label{sec:v16-integrated}
% =============================================================================

\begin{theorem}[Version V16 finite-card and terminal-writer theorem]
\label{thm:v16-integrated}
Preserving every theorem and limitation of Versions~V1--V15, Version~V16
proves the following.
\begin{enumerate}[label=\textup{(T16.\arabic*)},leftmargin=5.5em]
\item The selected-writer denominator has the two-sided window
      \eqref{eq:v16-denominator-window}; denominator collapse is excluded on
      the bounded-action selected branch.
\item The mixed coefficient has the fixed-scale representation
      \eqref{eq:v16-kappa-fixed-card}.  Its rescaled caloric-writer column is
      uniformly bounded on the interior age interval.
\item The selected covariance source, annular writer, and complete caloric
      writer form the compact positive Gram \eqref{eq:v16-three-Gram}.
      Shorting by the common writer gives the exact residual Gram
      \eqref{eq:v16-residual-Gram} and determinant
      \eqref{eq:v16-Uwr-determinant}.
\item The source--writer Riesz efficiency has a universal positive floor.
      A record-sized writer-orthogonal contact forces a source-angle or
      terminal-writer-angle defect; it cannot be an independent unpriced
      scalar.
\item Projection onto the physically oriented writer ray separates wrong-sign
      alignment from writer-orthogonal loss.  Wrong sign or any fixed
      calibration deficit forces a typed full-contact complement.
\item The mixed row is the carried old-state term plus the ordered Duhamel
      moment \eqref{eq:v16-source-moment-pairing}.  Endpoint response and
      unweighted source length do not determine this chronological moment.
\item Two positive balanced writer actions reconstruct the mixed cross row,
      and one positive endpoint norm difference reconstructs the represented
      contact.
\item The five-scalar card \eqref{eq:v16-five-card} determines the selected
      contact, held-out contact, and source-to-head residual exactly.  Its
      compact domain gives the quantitative stability estimates
      \eqref{eq:v16-selected-Lipschitz}--
      \eqref{eq:v16-incidence-Lipschitz}.
\item The represented, aggregate-residual, and clean-export alternatives are
      decided before any four-term absolute estimate.  A surviving aggregate
      residual is typed only afterward by the signed cumulative splitter.
\item On the clean branch the complete same-history packet
      \eqref{eq:v16-export-card} carries noncollapsed source--writer and
      terminal-writer incidence and a non-summable selected critical contact.
      NS--A stops and exports it to NS--B on that subsequence.
\end{enumerate}
No populated terminal card, finite-stock upper budget, terminal trace,
backward-uniqueness theorem, Liouville theorem, or global three-dimensional
regularity conclusion is supplied by these finite identities.
\end{theorem}

\begin{proof}
Items~\textup{(T16.1)--(T16.2)} are
\Cref{prop:v16-denominator-window,thm:v16-fixed-scale-kappa}.  Items
\textup{(T16.3)--(T16.4)} are
\Cref{thm:v16-three-column-Gram,thm:v16-double-Riesz,cor:v16-efficiency-fork}.
Item~\textup{(T16.5)} is
\Cref{thm:v16-oriented-ray,cor:v16-calibration-obstruction}.  Item
\textup{(T16.6)} is
\Cref{thm:v16-cumulative-response,cth:v16-endpoint-no-moment}.  Item
\textup{(T16.7)} is
\Cref{thm:v16-positive-polarization,prop:v16-represented-positive-read}.
Items~\textup{(T16.8)--(T16.9)} are
\Cref{thm:v16-five-card-readout,thm:v16-five-card-compact,thm:v16-card-stability,thm:v16-direct-verdict,cor:v16-post-verdict-typing}.
The final item is \Cref{thm:v16-clean-export}.
\end{proof}

\begin{firewall}[Exact stopping boundary after V16]
\label{fire:v16-stop}
On the bounded-action selected branch there is no longer an unnamed mixed
incidence.  It is the finite card \eqref{eq:v16-five-card}, and its branch map
is exact and stable.  A new general source norm, age law, path law, compact
profile space, or lower action estimate would be downstream of this card.
The only licensed continuations are: populate the card; prove a genuinely new
Navier--Stokes estimate for its ordered moment; follow a represented return;
or open the one signed complement already selected by the aggregate residual.
The clean branch has terminated in an NS--B export.
\end{firewall}

\begingroup\small
\begin{longtable}{p{0.25\textwidth}p{0.19\textwidth}p{0.48\textwidth}}
\caption{NS--A Version V16 proof-status ledger.}\label{tab:v16-ledger}\\
\toprule Row & Status & Exact conclusion\\\midrule
\endfirsthead
\toprule Row & Status & Exact conclusion\\\midrule
\endhead
V1--V15 prefix & \PROVED{} preserved
 &The complete V15 preterminal source is retained byte-for-byte; no earlier
 theorem is rewritten.\\[0.3em]
Writer denominator & \PROVED{} two-sided
 &$d_*/\nu\le D_n^W\le d^*/\nu$; the mixed coefficient is not amplified by
 denominator collapse.\\[0.3em]
Fixed-scale mixed coefficient & \PROVED{} exact
 &$\kappa_n/\mathcal R_n$ is the bounded rescaled projection
 $\mathfrak m_n/\mathfrak d_n$.\\[0.3em]
Rescaled caloric writer & \PROVED{} bounded
 &Its interior-age $L^2$ norm is at most
 $\sqrt{2/(e\tau_-)}$.\\[0.3em]
Three-column Gram & \PROVED{} positive compact card
 &The source, selected writer, and caloric writer give
 \eqref{eq:v16-three-Gram} in a fixed compact matrix set.\\[0.3em]
Double Riesz short & \PROVED{} exact
 &The writer-orthogonal contact is the cross entry of the positive residual
 Gram \eqref{eq:v16-residual-Gram}.\\[0.3em]
Source--writer efficiency & \PROVED{} positive floor
 &$\eta_n^{\rm sw}\ge h_0^2/(4L_*d^*)$.\\[0.3em]
Writer-orthogonal contact & \PROVED{} determinant bound
 &A record-sized value forces an actual source-angle or terminal-writer-angle
 defect.\\[0.3em]
Oriented terminal writer & \PROVED{} exact cone projection
 &Wrong-sign aligned mass and orthogonal loss are separated by
 \eqref{eq:v16-oriented-residual}.\\[0.3em]
Chronological mixed moment & \PROVED{} exact
 &The coefficient is carried old state plus one ordered common-endpoint Duhamel
 moment; endpoint marginals do not determine it.\\[0.3em]
Positive mixed acquisition & \PROVED{} exact
 &Two balanced positive actions reconstruct the mixed cross row and its sign.\\[0.3em]
Represented endpoint acquisition & \PROVED{} exact
 &One positive endpoint-head norm difference gives
 $\mathcal C_n^{\rm rep}$ before freshness.\\[0.3em]
Five-scalar card & \PROVED{} sufficient
 &The full, represented, selected, held-out, and incidence reads follow from
 \eqref{eq:v16-selected-readout}--\eqref{eq:v16-incidence-readout}.\\[0.3em]
Finite-card stability & \PROVED{} quantitative
 &The selected and incidence readouts are uniformly Lipschitz on the compact
 physical domain.\\[0.3em]
Aggregate residual typing & \PROVED{} post-verdict
 &The four signed complements are opened only after the direct incidence
 residual has survived.\\[0.3em]
Clean NS--A export & \CONDITIONAL{} exact
 &A populated clean card exports
 $(\mathfrak s_n,\kappa_n,\mathcal C_n^{\rm sel})$ and its full common-history
 Gram to NS--B.\\[0.3em]
Literal five-read population & \OPEN{} finite acquisition
 &No attached terminal trajectory supplies the positive actions and endpoint
 reads needed to choose one branch.\\[0.3em]
Finite-stock upper budget & \OPEN{} NS--B/physical
 &The non-summable selected contact is not dominated by an independently
 occurring finite inventory.\\[0.3em]
Three-dimensional regularity & \OPEN
 &No surviving represented, complement, calibration, or replenishment packet
 is excluded for every datum.\\
\bottomrule
\end{longtable}
\endgroup

% =============================================================================
\section{Exact mandate after Version V16}
\label{sec:v16-next}
% =============================================================================

\begin{programme}[Populate one finite card or follow one already typed branch]
\label{prog:v16-next}
Preserve Versions~V1--V16 verbatim.  There is no automatic general
Version~V17 on the bounded-action clean branch.

A continuation is licensed only in the following order.
\begin{enumerate}[label=\textup{(V17.\arabic*)},leftmargin=5.7em]
\item If the five same-history Reads in \eqref{eq:v16-five-card} are supplied,
      evaluate the stable branch map
      \eqref{eq:v16-selected-readout}--
      \eqref{eq:v16-incidence-readout}.  Do not reopen field compactness before
      this verdict.
\item On a terminal-writer calibration branch, evaluate the two terms
      $N_n^{\rm old}$ and $N_n^{\rm mom}$ of
      \eqref{eq:v16-mixed-moment-split}.  A new PDE estimate must control this
      actual ordered moment; an endpoint norm or another source-action lower
      bound is inadmissible.
\item On a represented branch, return to the already proved comparable-frequency
      and parabolic-time gates for old-annulus reuse and test its actual
      source-paid return stock.  Do not charge aligned replenishment as a new
      source direction.
\item On an aggregate incidence branch, open exactly one increment of
      \eqref{eq:v16-cumulative-splitter}.  Follow only that resolved,
      unselected-output, route, or writer-angle packet.
\item On the clean branch, terminate NS--A and continue only in NS--B with the
      exported card \eqref{eq:v16-export-card} and an independently occurring
      finite-stock or paid-stress upper budget.
\item The separately normalized source-action/enstrophy-square overdrive branch
      remains distinct and may reopen a compact profile only in its own
      normalization.
\end{enumerate}
No new Dini modulus, response graph, path measure, temporal partition,
abstract stock, or slow-time quotient is admitted in place of these finite
physical acquisitions.
\end{programme}

\begin{assessment}[Programme position after V16]
\label{ass:v16-position}
The bounded-action terminal problem has reached a finite executable boundary.
The mixed writer is a compact rescaled projection, its positive polarization
queries are explicit, the represented endpoint is read first, and the
source-to-head residual is computed as one assembled scalar.  The remaining
unknown is not the form of the theorem but the value of one literal card, or
a genuinely new Navier--Stokes estimate for its ordered source moment.  A
clean card has already crossed the single NS--A/NS--B interface and does not
license another NS--A lower-bound layer.
\end{assessment}

\paragraph{Version V16 history.}
Preserved the complete V15 source.  Added the missing upper denominator bound
and rescaled the terminal-writer coefficient to one fixed dimensionless
projection.  Constructed the complete positive source--writer--caloric-writer
Gram, computed its Schur residual and determinant, proved a source--writer
efficiency floor, and converted record-sized writer-orthogonal work into a
quantitative angle defect.  Separated desired writer orientation from
wrong-sign and orthogonal components.  Expressed the mixed row as carried old
state plus an ordered Duhamel moment and proved that endpoint response and
source length do not determine it.  Reconstructed the mixed and represented
rows from positive polarization actions, reduced the full V15 verdict to a
compact five-scalar card, and proved a uniform stability estimate.  Finally
made residual typing post-verdict and stated the exact clean export to NS--B.
No literal card values or finite-stock upper budget were claimed.

\begin{assessment}[Append-only integrity]
\label{ass:v16-integrity}
The complete Version~V15 source preceding its sole final document terminator
is preserved verbatim.  Its preterminal SHA--256 digest is
\begin{center}\scriptsize\ttfamily
26e1514d353886a1c6ff4109c978eb62d3ec63c48039d87356fc6308cdfe4d66
\end{center}
The present material begins at Section~\ref{sec:v16-scope}; the final command
below is again unique.
\end{assessment}

\addtocontents{toc}{\protect\endgroup}
% =============================================================================
% END VERSION V16 MATERIAL -- append later versions before final terminator
% =============================================================================

% APPEND ALL LATER VERSIONS IMMEDIATELY ABOVE THE SOLE FINAL LINE BELOW.

\clearpage
% =============================================================================
% VERSION V17: NEW MATERIAL.  The complete V1--V16 source above is preserved
% verbatim; only its sole active \end{document} line was removed before append.
% =============================================================================
\hypersetup{pdftitle={NS-A Terminal Source Concentration Programme: Cumulative V17},pdfauthor={A. Pelissier}}
\makeatletter
\addtocontents{toc}{\protect\begingroup\protect\makeatletter}
\addtocontents{toc}{\protect\def\protect\l@section{\protect\@dottedtocline{1}{0em}{4.7em}}}
\addtocontents{toc}{\protect\def\protect\l@subsection{\protect\@dottedtocline{2}{1.5em}{5.3em}}}
\addtocontents{toc}{\protect\makeatother}
\makeatother
\renewcommand{\thesection}{V17.\arabic{section}}
\renewcommand{\thesubsection}{V17.\arabic{section}.\arabic{subsection}}
\renewcommand{\theequation}{V17.\arabic{equation}}
\setcounter{section}{0}
\setcounter{equation}{0}
\renewcommand{\theHsection}{V17.\arabic{section}}
\renewcommand{\theHsubsection}{V17.\arabic{section}.\arabic{subsection}}
\renewcommand{\theHtheorem}{V17.\arabic{section}.\arabic{theorem}}
\renewcommand{\theHequation}{V17.\arabic{equation}}

\begin{center}
{\LARGE\bfseries NS--A: Version V17}\\[0.5em]
{\large Native-Source Metric Pullback, the Energy-Neutral Caloric Dipole,\\
and the Upstream Metric-Incidence Stop}\\[0.5em]
{\normalsize 3 September 2026}
\end{center}
\addcontentsline{toc}{part}{Version V17: native-source pullback and the energy-neutral caloric dipole}

\begin{abstract}
Version~V16 reduced the bounded-action branch to a compact five-scalar card.
That reduction is algebraically exact in the half-heat Reynolds metric used by
Versions~V6--V16: the smoothed covariance source is paired with a fixed
annular writer, and the half-caloric state writer is projected onto that same
writer in a weighted history space.  Before the resulting component can be
called the physical critical payment, however, one earlier row must be
checked.  The native Navier--Stokes source is
$c=A^{-1/4}\mathcal B(u)$, its native metric is the ordinary $L^2$ metric,
and it has the exact null direction
$\langle c,A^{1/4}u\rangle=0$.  Heat smoothing changes the source metric and
spatial localization changes its factorization.  A later finite Gram cannot
repair a missing incidence between those two metrics.

Version~V17 therefore goes upstream.  For the complete caloric contact it
constructs the unique energy-neutral native writer
\[
 \Lambda e^{-2\nu rA}e-
 \frac{\langle\Lambda e^{-2\nu rA}e,e\rangle}{\|e\|_2^2}
 e.
\]
Its squared norm is exactly the spectral variance of
$\kappa e^{-2\nu r\kappa^2}$ under the critical-energy distribution.  This
is the finite-age continuation of the inherited frequency-centred critical
dipole.  It yields the sharp source-relative action bound and proves that a
caloric-flat spectral packet cannot carry the contact.

The half-heat factorization is then audited.  At finite cutoff its energy-null
direction is the backward-heat vector $e^{\nu rA}e$, and the physical
centering is orthogonal in the transported native metric, not in the
ordinary half-heat metric.  An explicit two-mode heat-compatible model shows
that the raw V16 card can place the entire contact in its selected component
while the source-native energy-centred projection is zero.  Thus the V16
clean branch is not, by itself, invariant under the exact source-null fibre.
The example is a finite spectral obstruction, not a Navier--Stokes
trajectory.

On the actual annular card, V17 pulls the selected smoothed source back through
the heat multiplier before spatial localization.  The same scalar is then
paired with the source-native writer $E_r(\Theta_nW_n)$.  The inverse heat is
uniformly bounded because the source has already been projected to the fixed
annulus and the selected ages are parabolic.  The resulting source action,
writer denominator, and energy-centred caloric-writer action have uniform
windows.  They form a new compact positive three-column Gram and a corrected
five-scalar card.  Its mixed row is acquired by positive polarization.

The old and new cards differ by one exact metric-incidence residual.  Its
denominator contains two positive debits---smooth localization slack and
heat-metric loss---whereas its numerator contains a signed heat-metric cross
term and the energy-centering transfer.  The source-to-head residual changes
by exactly this metric defect.  Consequently a V16 clean card may be exported
to NS--B only after the metric-incidence residual vanishes, or after the
source-native card itself is populated and passes.  Otherwise the metric
residual is the first upstream obstruction.

Finally, the source-native mixed row is written as carried old state plus an
ordered Duhamel moment with the correct extra heat factors, minus the
energy-centering term.  Its follower writer has a sharp
$O(\Gamma_n/\nu)$ bound, but no sign or noncollapsed lower edge follows from
the current theorem layers.  No new Dini modulus, profile space, path law, or
global regularity conclusion is claimed.
\end{abstract}

% =============================================================================
\section{Scope: the source metric precedes the terminal-writer card}
\label{sec:v17-scope}
% =============================================================================

Retain the bounded-action selected branch and all notation of
Versions~V13--V16.  Thus
\begin{equation}
 b_n=\beta_{n+1},\qquad r_n(t)=b_n-t,
 \qquad E_{n,t}:=e^{-\nu r_n(t)A},
 \label{eq:v17-heat-notation}
\end{equation}
\begin{equation}
 e(t)=A^{1/4}u(t),\qquad c(t)=A^{-1/4}\BNS(u(t)),
 \qquad \langle c(t),e(t)\rangle=0,
 \label{eq:v17-native-source-null}
\end{equation}
Define the physical-time pullback of the fixed V13 slow-time cutoff by
\begin{equation}
 \eta_n(t):=\eta\!\left(\nu\Gamma_n^2(t-b_n)\right).
 \label{eq:v17-physical-time-cutoff}
\end{equation}
The selected interior ages satisfy
\begin{equation}
 \frac{\tau_-}{\nu\Gamma_n^2}
 \le r_n(t)\le
 \frac{\tau_+}{\nu\Gamma_n^2}
 \quad\hbox{on }\supp\eta_n,
 \qquad 0<\tau_-<\tau_+<\infty.
 \label{eq:v17-age-window}
\end{equation}
The next record is the first time at level $2\mathcal R_n$, so
\begin{equation}
 \|e(t)\|_2<2\mathcal R_n
 \quad(\beta_n\le t<b_n),
 \qquad \|e(b_n)\|_2=2\mathcal R_n.
 \label{eq:v17-record-cap}
\end{equation}
The solution cannot vanish at an intermediate time: if $u(t_0)=0$, forward
uniqueness on the smooth interval gives $u\equiv0$ afterward, contradicting
\eqref{eq:v17-record-cap}.  Hence every quotient by the line $\mathbb Re(t)$
below is well defined.

The complete caloric-contact integrand has two equivalent factorizations:
\begin{equation}
 \boxed{
 \mathcal J_{b_n}(t)
 =-\langle E_{n,t}c(t),\Lambda E_{n,t}e(t)\rangle
 =-\langle c(t),\Lambda E_{n,t}^2e(t)\rangle.}
 \label{eq:v17-two-contact-factorizations}
\end{equation}
Version~V16 forms its mixed writer in the first factorization.  The exact
energy-null identity in \eqref{eq:v17-native-source-null}, however, belongs to
the second, native-source factorization.  The present version transports the
selected card to that metric before interpreting its selected component as a
physical payment.

\begin{firewall}[Correction of export type, not of the V16 algebra]
\label{fire:v17-v16-status}
Every V16 Gram and readout remains an exact identity for the displayed
half-heat vectors and the weighted metric chosen there.  The correction is
the physical promotion: a half-heat selected component is not yet a
source-minimal critical payment unless its metric and energy-null incidence
with the native source has been verified.  The V16 clean-export statement is
therefore retained as a smoothed-Reynolds-card conclusion and is strengthened
below by one upstream metric-incidence gate.
\end{firewall}

% =============================================================================
\section{The native energy-neutral caloric dipole}
\label{sec:v17-native-dipole}
% =============================================================================

For $t\in(\beta_n,b_n)$ define the complete native caloric writer
\begin{equation}
 \widehat Z_{n,t}^{\rm cal}
 :=\Lambda E_{n,t}^2e(t).
 \label{eq:v17-native-caloric-writer}
\end{equation}
Put
\begin{equation}
 \lambda_n^{\rm en}(t)
 :=\frac{\langle\widehat Z_{n,t}^{\rm cal},e(t)\rangle}
          {\|e(t)\|_2^2},
 \qquad
 Z_{n,t}^{\rm dip}
 :=\widehat Z_{n,t}^{\rm cal}-\lambda_n^{\rm en}(t)e(t).
 \label{eq:v17-energy-centred-dipole}
\end{equation}

\begin{theorem}[Energy-neutral caloric dipole and native minimum]
\label{thm:v17-native-dipole}
For every smooth time before $b_n$,
\begin{equation}
 \boxed{
 \lambda_n^{\rm en}(t)\ge0,
 \qquad
 \langle Z_{n,t}^{\rm dip},e(t)\rangle=0,}
 \label{eq:v17-centred-properties}
\end{equation}
\begin{equation}
 \boxed{
 \mathcal J_{b_n}(t)
 =-\langle c(t),Z_{n,t}^{\rm dip}\rangle.}
 \label{eq:v17-centred-contact}
\end{equation}
Moreover, for every real $a$,
\begin{equation}
 \boxed{
 \|\widehat Z_{n,t}^{\rm cal}-ae(t)\|_2^2
 =\|Z_{n,t}^{\rm dip}\|_2^2
  +|a-\lambda_n^{\rm en}(t)|^2\|e(t)\|_2^2.}
 \label{eq:v17-native-minimum}
\end{equation}
Thus $Z_{n,t}^{\rm dip}$ is the unique minimum-norm representative of the
caloric writer modulo the exact source-null line $\mathbb Re(t)$.

Let
\begin{equation}
 \mu_t(k):=\frac{|\widehat e(t,k)|^2}{\|e(t)\|_2^2},
 \qquad
 f_r(k):=|k|e^{-2\nu r|k|^2}.
 \label{eq:v17-spectral-law}
\end{equation}
Then $\sum_{k\ne0}\mu_t(k)=1$ and
\begin{equation}
 \boxed{
 \lambda_n^{\rm en}(t)=\sum_{k\ne0}f_{r_n(t)}(k)\mu_t(k),}
 \label{eq:v17-frequency-centroid}
\end{equation}
\begin{equation}
 \boxed{
 \|Z_{n,t}^{\rm dip}\|_2^2
 =\|e(t)\|_2^2
  \sum_{k\ne0}
  |f_{r_n(t)}(k)-\lambda_n^{\rm en}(t)|^2\mu_t(k).}
 \label{eq:v17-caloric-frequency-variance}
\end{equation}
At zero age this is exactly the inherited frequency-centred critical dipole.
\end{theorem}

\begin{proof}
The multiplier $\Lambda E_{n,t}^2$ is nonnegative and self-adjoint, so its
Rayleigh quotient \eqref{eq:v17-energy-centred-dipole} is nonnegative.
Orthogonality is immediate.  Equation~\eqref{eq:v17-centred-contact} follows
from \eqref{eq:v17-native-source-null}.  Orthogonal Pythagoras in the affine
line $\widehat Z_{n,t}^{\rm cal}+\mathbb Re(t)$ gives
\eqref{eq:v17-native-minimum}.

On a nonzero Fourier mode, $\Lambda E_{n,t}^2$ has symbol
$|k|e^{-2\nu r_n(t)|k|^2}$.  Parseval's identity gives
\eqref{eq:v17-frequency-centroid}; subtracting its weighted mean gives
\eqref{eq:v17-caloric-frequency-variance}.  At $r=0$ the symbol is $|k|$,
which is the frequency variable used in the inherited critical dipole.
\end{proof}

\begin{corollary}[Sharp source-relative action and the caloric-flat stop]
\label{cor:v17-native-action}
At every smooth time,
\begin{equation}
 \boxed{
 |\mathcal J_{b_n}(t)|^2
 \le\|c(t)\|_2^2\|Z_{n,t}^{\rm dip}\|_2^2.}
 \label{eq:v17-native-action-bound}
\end{equation}
If $\mathcal J_{b_n}(t)\ne0$, then
\begin{equation}
 \boxed{
 \|c(t)\|_2^2
 \ge\frac{|\mathcal J_{b_n}(t)|^2}
          {\|Z_{n,t}^{\rm dip}\|_2^2}.}
 \label{eq:v17-native-action-lower}
\end{equation}
If the variance in \eqref{eq:v17-caloric-frequency-variance} vanishes, then
$\mathcal J_{b_n}(t)=0$.  Finally,
\begin{equation}
 \boxed{
 0\le\lambda_n^{\rm en}(t)
 \le\frac1{\sqrt{4e\nu r_n(t)}},
 \qquad
 \|Z_{n,t}^{\rm dip}\|_2
 \le\frac{\|e(t)\|_2}{\sqrt{4e\nu r_n(t)}}.}
 \label{eq:v17-centred-age-envelope}
\end{equation}
\end{corollary}

\begin{proof}
The first two statements are Cauchy--Schwarz and its rearrangement applied to
\eqref{eq:v17-centred-contact}.  Zero variance means
$Z_{n,t}^{\rm dip}=0$.  Finally,
\[
 \sup_{\kappa\ge0}\kappa e^{-2\nu r\kappa^2}
 =\frac1{\sqrt{4e\nu r}}.
\]
The Rayleigh quotient and the orthogonal projection are bounded by the norm
of this multiplier.
\end{proof}

\begin{assessment}[What the variance does and does not provide]
\label{ass:v17-variance-scope}
The caloric variance is a sharper target than the uncentred age envelope.  It
vanishes on every spectral packet on which
$\kappa e^{-2\nu r\kappa^2}$ is constant, and at zero age it vanishes on a
single frequency sphere.  No terminal upper estimate for this variance is
proved here.  Accordingly \eqref{eq:v17-native-action-lower} is a precise
source-action alternative, not a Dini improvement or a regularity theorem.
\end{assessment}

% =============================================================================
\section{Why the half-heat Gram is not yet the native source quotient}
\label{sec:v17-half-metric}
% =============================================================================

The metric issue is already visible at a finite Galerkin cutoff, where every
heat inverse is defined.  Suppress $(n,t)$ and write
\begin{equation}
 E=e^{-\nu rA},\qquad h=Ec,\qquad z=\Lambda Ee,
 \qquad \widehat z=Ez=\Lambda E^2e.
 \label{eq:v17-half-symbols}
\end{equation}

\begin{proposition}[Transported energy-null direction in the half-heat space]
\label{prop:v17-half-energy-null}
At finite cutoff put
\begin{equation}
 d:=E^{-1}e,
 \qquad
 \psi^{\rm half}:=z-\lambda^{\rm en}d,
 \qquad
 \lambda^{\rm en}:=\frac{\langle\widehat z,e\rangle}{\|e\|^2}.
 \label{eq:v17-half-centred}
\end{equation}
Then
\begin{equation}
 \boxed{
 \langle h,d\rangle=0,
 \qquad
 E\psi^{\rm half}=\widehat z-\lambda^{\rm en}e,}
 \label{eq:v17-half-null-relations}
\end{equation}
\begin{equation}
 \boxed{
 \langle h,z\rangle=\langle h,\psi^{\rm half}\rangle.}
 \label{eq:v17-half-contact-invariance}
\end{equation}
Moreover, $\psi^{\rm half}$ minimizes
$a\mapsto\|E(z-ad)\|^2$, not in general
$a\mapsto\|z-ad\|^2$.  Thus the native energy quotient is orthogonal in the
transported metric $\langle x,y\rangle_E:=\langle Ex,Ey\rangle$.
\end{proposition}

\begin{proof}
Self-adjointness and invertibility at finite cutoff give
$\langle h,d\rangle=\langle Ec,E^{-1}e\rangle=\langle c,e\rangle=0$.
The second relation in \eqref{eq:v17-half-null-relations} follows by applying
$E$.  Equation~\eqref{eq:v17-half-contact-invariance} then follows from the
first relation.  Finally,
\[
 \|E(z-ad)\|^2=\|\widehat z-ae\|^2,
\]
so Theorem~\ref{thm:v17-native-dipole} gives the stated minimizer.  Ordinary
half-heat orthogonality would instead use
$\langle z,d\rangle/\|d\|^2$, which need not equal $\lambda^{\rm en}$.
\end{proof}

\begin{firewall}[The backward-heat direction is not assumed in the limit]
\label{fire:v17-no-global-backward-heat}
In infinite dimension $E^{-1}e$ need not belong to $L^2$ uniformly, or even
at a specified positive age without an analytic-domain input.  V17 does not
insert such an input.  The source-native construction below pulls back only
the already annular selected source, on which $E^{-1}$ is uniformly bounded,
and keeps the complete caloric writer in the original native source metric.
\end{firewall}

\begin{countertheorem}[A raw clean half-heat card need not be native clean]
\label{cth:v17-raw-clean-counterexample}
There is a two-mode heat-compatible finite spectral model satisfying the
energy-null identity in which the raw half-heat projection places the entire
contact in the selected writer, while the source-native energy-centred mixed
coefficient is zero.

More precisely, let
\begin{equation}
 H=\mathbb R^2,\qquad
 \Lambda=\begin{pmatrix}1&0\\0&2\end{pmatrix},
 \qquad E_r=e^{-r\Lambda^2},
 \qquad e=2^{-1/2}(1,1).
 \label{eq:v17-two-mode-data}
\end{equation}
For every sufficiently small $r>0$ there are vectors $c_r,h_r,g_r,W_r$
with
\begin{equation}
 h_r=E_rc_r,\qquad \langle c_r,e\rangle=0,
 \qquad h_r=g_r+(h_r-g_r),
 \label{eq:v17-two-mode-source}
\end{equation}
for which, with $z_r=\Lambda E_re$,
\begin{equation}
 \boxed{
 -\langle h_r,z_r\rangle
 =\kappa_r^{\rm sm}s_r>0,}
 \label{eq:v17-raw-clean}
\end{equation}
where
\begin{equation}
 s_r=-\langle g_r,W_r\rangle>0,
 \qquad
 \kappa_r^{\rm sm}
 =\frac{\langle z_r,W_r\rangle}{\|W_r\|^2},
 \label{eq:v17-raw-coefficients}
\end{equation}
while both the raw selected writer residual and the raw source complement
have zero contact.  If
\begin{equation}
 Z_r^{\rm dip}
 =\Lambda E_r^2e
 -\frac{\langle\Lambda E_r^2e,e\rangle}{\|e\|^2}e,
 \qquad W_r^{\rm nat}=E_rW_r,
 \label{eq:v17-two-mode-native}
\end{equation}
then
\begin{equation}
 \boxed{
 \langle Z_r^{\rm dip},W_r^{\rm nat}\rangle=0.}
 \label{eq:v17-native-zero}
\end{equation}
\end{countertheorem}

\begin{proof}
At $r=0$ one has
\[
 z_0=2^{-1/2}(1,2),\qquad
 \lambda_0^{\rm en}=\frac32,
 \qquad
 Z_0^{\rm dip}=\frac1{2\sqrt2}(-1,1).
\]
Hence $e\perp Z_0^{\rm dip}$ and
$\langle z_0,e\rangle=3/2>0$.  Choose, continuously for small $r$, a unit
vector $v_r\perp Z_r^{\rm dip}$ with $v_r\to e$, and put
\begin{equation}
 W_r:=E_r^{-1}v_r,
 \qquad g_r:=-W_r.
 \label{eq:v17-two-mode-Wg}
\end{equation}
Then $W_r^{\rm nat}=v_r$, so \eqref{eq:v17-native-zero} holds, whereas
$\langle z_r,W_r\rangle>0$ for all sufficiently small $r$ by continuity.

Choose a continuous unit vector $q_r\perp z_r$ with
$q_0=(2,-1)/\sqrt5$.  Since
$\langle q_0,e\rangle=1/\sqrt{10}\ne0$, one has
$\langle q_r,E_r^{-1}e\rangle\ne0$ for small $r$.  Define
\begin{equation}
 t_r:=-\frac{\langle g_r,E_r^{-1}e\rangle}
             {\langle q_r,E_r^{-1}e\rangle},
 \qquad
 h_r:=g_r+t_rq_r,
 \qquad
 c_r:=E_r^{-1}h_r.
 \label{eq:v17-two-mode-completion}
\end{equation}
Then
$\langle c_r,e\rangle=\langle h_r,E_r^{-1}e\rangle=0$.  Since
$q_r\perp z_r$,
\[
 -\langle h_r,z_r\rangle
 =-\langle g_r,z_r\rangle
 =\langle W_r,z_r\rangle>0.
\]
Also $s_r=\|W_r\|^2$ and therefore
$\kappa_r^{\rm sm}s_r=\langle z_r,W_r\rangle$, proving
\eqref{eq:v17-raw-clean}.  The residual
$z_r-\kappa_r^{\rm sm}W_r$ is orthogonal to $W_r$ and hence to $g_r$;
the source complement $t_rq_r$ is orthogonal to $z_r$.  Thus both raw
complements vanish, while the native centred projection is zero.
\end{proof}

\begin{assessment}[Meaning of the finite counterexample]
\label{ass:v17-counterexample-scope}
The construction respects a genuine heat functional calculus
$E_r=e^{-r\Lambda^2}$ and the exact energy-null identity.  It does not claim
to be a Navier--Stokes trajectory or to satisfy the terminal record
chronology.  It proves the precise logical point needed here: heat
factorization, Gram positivity, separate norm windows, and energy neutrality
do not identify the half-heat selected component with the native
energy-centred payment.  A same-history metric-incidence row is necessary.
\end{assessment}

% =============================================================================
\section{Source-first pullback of the selected annular scalar}
\label{sec:v17-native-pullback}
% =============================================================================

Let
\begin{equation}
 H_{n,t}:=\Psi_ng_{n,t}^{\rm cov}
 \label{eq:v17-smoothed-selected-source}
\end{equation}
be the already selected smoothed covariance source.  It has Fourier support
in $[\Gamma_n/2,2\Gamma_n]$.  Extend the following vectors by zero outside
$\supp\eta_n$ and define
\begin{equation}
 G_{n,t}^{\rm nat}:=E_{n,t}^{-1}H_{n,t},
 \qquad
 W_{n,t}^{\rm nat}:=E_{n,t}\bigl(\Theta_n(t)W_n\bigr).
 \label{eq:v17-native-source-writer}
\end{equation}
The inverse in the first formula is taken only on the fixed annular range.
All spatial localization remains on the writer side; no backward heat is
applied after multiplication by the cell or router cutoff.

\begin{theorem}[Exact source-native pullback and uniform action window]
\label{thm:v17-native-pullback}
The V15 selected scalar has the exact native-source representation
\begin{equation}
 \boxed{
 \mathfrak s_n
 =-\int_{\beta_n}^{b_n}
   \langle G_{n,t}^{\rm nat},W_{n,t}^{\rm nat}\rangle\,\dd t
 =-\sigma_*\mathcal R_nA_n.}
 \label{eq:v17-native-scalar}
\end{equation}
Moreover,
\begin{equation}
 \boxed{
 \int_{\beta_n}^{b_n}\|G_{n,t}^{\rm nat}\|_2^2\,\dd t
 \le e^{8\tau_+}\nu\mathcal R_n^2L_*.}
 \label{eq:v17-native-source-action-upper}
\end{equation}
Put
\begin{equation}
 D_n^{\rm nat}
 :=\int_{\beta_n}^{b_n}\|W_{n,t}^{\rm nat}\|_2^2\,\dd t.
 \label{eq:v17-native-denominator}
\end{equation}
Then
\begin{equation}
 \boxed{
 \frac{d_{\rm nat,*}}{\nu}
 \le D_n^{\rm nat}
 \le\frac{d_{\rm nat}^*}{\nu},}
 \label{eq:v17-native-denominator-window}
\end{equation}
where one may take
\begin{equation}
 d_{\rm nat,*}
 :=\frac{h_0^2}{4e^{8\tau_+}L_*},
 \qquad
 d_{\rm nat}^*:=4(\tau_+-\tau_-).
 \label{eq:v17-native-d-constants}
\end{equation}
\end{theorem}

\begin{proof}
Self-adjointness of $E_{n,t}$ gives
\[
 \langle E_{n,t}^{-1}H_{n,t},
         E_{n,t}(\Theta_nW_n)\rangle
 =\langle H_{n,t},\Theta_nW_n\rangle.
\]
Integration and \eqref{eq:v15-signed-selected-scalar} prove
\eqref{eq:v17-native-scalar}.

On the annular support,
\[
 \|E_{n,t}^{-1}\Psi_n\|_{2\to2}
 \le\exp\!\bigl(\nu r_n(t)(2\Gamma_n)^2\bigr)
 \le e^{4\tau_+}.
\]
The V10 scaling identity and the V13 bounded-action branch give
\[
 \int_{\supp\eta_n}\|H_{n,t}\|_2^2\,\dd t
 \le\nu\mathcal R_n^2L_*.
\]
This proves \eqref{eq:v17-native-source-action-upper}.

Cauchy--Schwarz in \eqref{eq:v17-native-scalar} and
$|\mathfrak s_n|\ge(h_0/2)\mathcal R_n$ give the lower bound for
$D_n^{\rm nat}$.  For the upper bound, heat contraction,
$0\le\Theta_n\le1$, the annular estimate $\|W_n\|_2\le2\Gamma_n$, and
\eqref{eq:v17-age-window} yield
\[
 D_n^{\rm nat}
 \le\int_{\supp\eta_n}\|\Theta_nW_n\|_2^2\,\dd t
 \le4\Gamma_n^2\frac{\tau_+-\tau_-}{\nu\Gamma_n^2}.
\]
\end{proof}

\begin{corollary}[The source-first pullback retains positive Riesz efficiency]
\label{cor:v17-native-efficiency}
Define
\begin{equation}
 \mathfrak a_n^{\rm nat}
 :=\frac1{\nu\mathcal R_n^2}
   \int\|G_{n,t}^{\rm nat}\|_2^2\,\dd t,
 \qquad
 \mathfrak d_n^{\rm nat}:=\nu D_n^{\rm nat}.
 \label{eq:v17-native-a-d}
\end{equation}
Then
\begin{equation}
 \boxed{
 \frac{A_n^2}{\mathfrak a_n^{\rm nat}\mathfrak d_n^{\rm nat}}
 \ge
 \frac{h_0^2}
 {4e^{8\tau_+}L_*d_{\rm nat}^*}>0.}
 \label{eq:v17-native-efficiency-floor}
\end{equation}
\end{corollary}

\begin{proof}
The ratio is the squared cosine of the two vectors in
\eqref{eq:v17-native-scalar}.  Use
$A_n\ge h_0/2$, \eqref{eq:v17-native-source-action-upper}, and the upper
bound in \eqref{eq:v17-native-denominator-window}.
\end{proof}

% =============================================================================
\section{The source-native centred Gram and corrected five-scalar card}
\label{sec:v17-native-card}
% =============================================================================

Use $Z_{n,t}^{\rm dip}$ from \eqref{eq:v17-energy-centred-dipole} and put
\begin{equation}
 \mathcal E_n^{\rm dip}
 :=\int_{\supp\eta_n}\|Z_{n,t}^{\rm dip}\|_2^2\,\dd t.
 \label{eq:v17-dipole-action}
\end{equation}

\begin{lemma}[Uniform native caloric-writer window]
\label{lem:v17-dipole-window}
One has
\begin{equation}
 \boxed{
 \mathcal E_n^{\rm dip}
 \le\frac{\mathcal R_n^2}{e\nu}
       \log\frac{\tau_+}{\tau_-}.}
 \label{eq:v17-dipole-window}
\end{equation}
\end{lemma}

\begin{proof}
Theorem~\ref{thm:v17-native-dipole} and
\eqref{eq:v17-record-cap} give
\[
 \|Z_{n,t}^{\rm dip}\|_2^2
 \le\|\widehat Z_{n,t}^{\rm cal}\|_2^2
 \le\frac{\mathcal R_n^2}{e\nu r_n(t)}.
\]
Integrate in $r=b_n-t$ over the interval in
\eqref{eq:v17-age-window}.
\end{proof}

Define the native mixed coefficient by
\begin{equation}
 N_n^{\rm nat}
 :=\int_{\supp\eta_n}
   \langle Z_{n,t}^{\rm dip},W_{n,t}^{\rm nat}\rangle\,\dd t,
 \qquad
 \kappa_n^{\rm nat}:=\frac{N_n^{\rm nat}}{D_n^{\rm nat}}.
 \label{eq:v17-native-mixed}
\end{equation}
Put
\begin{equation}
 \mathfrak m_n^{\rm nat}:=\frac\nu{\mathcal R_n}N_n^{\rm nat},
 \qquad
 \mathfrak e_n^{\rm nat}:=
 \frac\nu{\mathcal R_n^2}\mathcal E_n^{\rm dip},
 \label{eq:v17-native-m-e}
\end{equation}
so that
\begin{equation}
 \boxed{
 \frac{\kappa_n^{\rm nat}}{\mathcal R_n}
 =\frac{\mathfrak m_n^{\rm nat}}{\mathfrak d_n^{\rm nat}}.}
 \label{eq:v17-native-kappa}
\end{equation}

\begin{theorem}[Compact source-native three-column Gram]
\label{thm:v17-native-three-Gram}
Define
\begin{equation}
 \mathfrak c_n^{\rm nat}
 :=-\frac1{\mathcal R_n^2}
   \int_{\supp\eta_n}
   \langle G_{n,t}^{\rm nat},Z_{n,t}^{\rm dip}\rangle\,\dd t.
 \label{eq:v17-native-source-contact}
\end{equation}
Then
\begin{equation}
 \boxed{
 \mathbb G_n^{\rm nat}
 :=\begin{pmatrix}
 \mathfrak a_n^{\rm nat}&\sigma_*A_n&-\mathfrak c_n^{\rm nat}\\
 \sigma_*A_n&\mathfrak d_n^{\rm nat}&\mathfrak m_n^{\rm nat}\\
 -\mathfrak c_n^{\rm nat}&\mathfrak m_n^{\rm nat}&\mathfrak e_n^{\rm nat}
 \end{pmatrix}\succeq0.}
 \label{eq:v17-native-three-Gram}
\end{equation}
The entries obey the fixed bounds
\begin{equation}
 \boxed{
 0\le\mathfrak a_n^{\rm nat}\le e^{8\tau_+}L_*,
 \qquad
 d_{\rm nat,*}\le\mathfrak d_n^{\rm nat}\le d_{\rm nat}^*,}
 \label{eq:v17-native-Gram-bounds-1}
\end{equation}
\begin{equation}
 \boxed{
 0\le\mathfrak e_n^{\rm nat}
 \le e_{\rm nat}^*:=\frac1e\log\frac{\tau_+}{\tau_-},
 \qquad
 |\mathfrak m_n^{\rm nat}|
 \le\sqrt{d_{\rm nat}^*e_{\rm nat}^*}.}
 \label{eq:v17-native-Gram-bounds-2}
\end{equation}
\end{theorem}

\begin{proof}
The matrix is the Gram of the three history vectors
\[
 \frac{G_n^{\rm nat}}{\sqrt\nu\mathcal R_n},
 \qquad
 \sqrt\nu W_n^{\rm nat},
 \qquad
 \frac{\sqrt\nu}{\mathcal R_n}Z_n^{\rm dip}.
\]
The first cross entry is
$\langle G_n^{\rm nat},W_n^{\rm nat}\rangle/\mathcal R_n
 =-\mathfrak s_n/\mathcal R_n=\sigma_*A_n$.
The other entries are the definitions above.  The bounds are
Theorem~\ref{thm:v17-native-pullback},
Lemma~\ref{lem:v17-dipole-window}, and Cauchy--Schwarz.
\end{proof}

Let
\begin{equation}
 Z_{n,t}^{\rm dip,\perp}
 :=Z_{n,t}^{\rm dip}-\kappa_n^{\rm nat}W_{n,t}^{\rm nat}.
 \label{eq:v17-native-writer-residual}
\end{equation}
Then $Z_n^{\rm dip,\perp}\perp W_n^{\rm nat}$ in the history space.  Put
\begin{equation}
 \mathcal U_n^{\rm nat,wr}
 :=-\int
 \langle G_{n,t}^{\rm nat},Z_{n,t}^{\rm dip,\perp}\rangle\,\dd t.
 \label{eq:v17-native-writer-work}
\end{equation}

\begin{proposition}[Exact source-native selected contact]
\label{prop:v17-native-selected-contact}
The selected native contact is
\begin{equation}
 \boxed{
 \mathcal C_n^{\rm nat,sel}
 :=\kappa_n^{\rm nat}\mathfrak s_n
 =-\sigma_*\mathcal R_n^2A_n
   \frac{\mathfrak m_n^{\rm nat}}{\mathfrak d_n^{\rm nat}}.}
 \label{eq:v17-native-selected-contact}
\end{equation}
The complete contact has the exact source-first ledger
\begin{equation}
 \boxed{
 \mathcal C_n^{\rm cal}
 =\mathcal C_n^{\rm nat,comp}
  +\mathcal C_n^{\rm nat,sel}
  +\mathcal U_n^{\rm nat,wr},}
 \label{eq:v17-native-source-ledger}
\end{equation}
where
\begin{equation}
 \mathcal C_n^{\rm nat,comp}
 :=-\int_{\beta_n}^{b_n}
 \langle c(t)-G_{n,t}^{\rm nat},Z_{n,t}^{\rm dip}\rangle\,\dd t.
 \label{eq:v17-native-complement}
\end{equation}
Here $G_n^{\rm nat}$ is extended by zero outside the selected interior
interval.  The complement includes every nonselected native source and is
opened further only if the aggregate native incidence survives.
\end{proposition}

\begin{proof}
Projection gives
$Z_n^{\rm dip}=\kappa_n^{\rm nat}W_n^{\rm nat}
 +Z_n^{\rm dip,\perp}$.  Pair with $-G_n^{\rm nat}$ and use
\eqref{eq:v17-native-scalar}.  Add the exact source complement
$c-G_n^{\rm nat}$ and use \eqref{eq:v17-centred-contact} integrated over the
record interval.
\end{proof}

\begin{theorem}[Positive acquisition of the native mixed row]
\label{thm:v17-native-positive-read}
Define the two positive history actions
\begin{equation}
 \mathcal E_{n,\rm nat}^{\pm}
 :=\int_{\supp\eta_n}
 \left\|Z_{n,t}^{\rm dip}
       \pm\mathcal R_nW_{n,t}^{\rm nat}\right\|_2^2\,\dd t.
 \label{eq:v17-native-positive-actions}
\end{equation}
Then
\begin{equation}
 \boxed{
 \mathfrak m_n^{\rm nat}
 =\frac\nu{4\mathcal R_n^2}
   \left(\mathcal E_{n,\rm nat}^{+}
        -\mathcal E_{n,\rm nat}^{-}\right).}
 \label{eq:v17-native-positive-polarization}
\end{equation}
Thus the corrected cross row can be acquired by positive polarization after
the native energy-centred writer has been formed.
\end{theorem}

\begin{proof}
Expand the two squares and subtract.  The difference is
$4\mathcal R_nN_n^{\rm nat}$.
\end{proof}

Define the corrected five-scalar card
\begin{equation}
 \boxed{
 \mathbf K_n^{(17)}
 :=\left(
 \mathfrak c_n^{\rm cal},
 \mathfrak c_n^{\rm rep},
 A_n,\mathfrak d_n^{\rm nat},\mathfrak m_n^{\rm nat}
 \right).}
 \label{eq:v17-five-card}
\end{equation}
Its selected and source-to-head readouts are
\begin{gather}
 \boxed{
 S_n^{\rm nat}
 :=-\sigma_*A_n
   \frac{\mathfrak m_n^{\rm nat}}{\mathfrak d_n^{\rm nat}}
   =\frac{\mathcal C_n^{\rm nat,sel}}{\mathcal R_n^2}}
 \label{eq:v17-native-selected-readout}\\
 \boxed{
 I_n^{\rm nat}
 :=\mathfrak c_n^{\rm cal}-\mathfrak c_n^{\rm rep}-S_n^{\rm nat}
   =\frac{\mathcal C_n^{\rm ho}
          -\mathcal C_n^{\rm nat,sel}}{\mathcal R_n^2}}
 \label{eq:v17-native-incidence-readout}
\end{gather}
The compactness and Lipschitz theorem of V16 applies verbatim with
$(\mathfrak d_n,\mathfrak m_n)$ replaced by
$(\mathfrak d_n^{\rm nat},\mathfrak m_n^{\rm nat})$ and with the constants in
\eqref{eq:v17-native-Gram-bounds-1}--\eqref{eq:v17-native-Gram-bounds-2}.

% =============================================================================
\section{Exact comparison with the V16 half-heat card}
\label{sec:v17-metric-defect}
% =============================================================================

For clarity write
\begin{equation}
 z_{n,t}^{\rm half}:=\Lambda E_{n,t}e(t),
 \qquad
 D_n^{\rm sm}:=\int\langle W_n,\Theta_nW_n\rangle\,\dd t,
 \label{eq:v17-smoothed-raw-data}
\end{equation}
so $D_n^{\rm sm}=D_n^W$ and
\begin{equation}
 N_n^{\rm sm}:=
 \int\langle z_{n,t}^{\rm half},\Theta_nW_n\rangle\,\dd t
 \label{eq:v17-smoothed-numerator}
\end{equation}
produces the V16 ratio
$\kappa_n/\mathcal R_n=\mathfrak m_n/\mathfrak d_n$.
Since
$\widehat Z_{n,t}^{\rm cal}=E_{n,t}z_{n,t}^{\rm half}$ and
$W_{n,t}^{\rm nat}=E_{n,t}(\Theta_nW_n)$, define
\begin{equation}
 N_n^{\rm en}
 :=\int\lambda_n^{\rm en}(t)
       \langle e(t),E_{n,t}(\Theta_nW_n)\rangle\,\dd t.
 \label{eq:v17-energy-transfer-row}
\end{equation}

\begin{theorem}[Positive denominator debit and signed numerator transfer]
\label{thm:v17-metric-bridge}
The two writer denominators satisfy the exact positive identity
\begin{equation}
 \boxed{
 \begin{aligned}
 D_n^{\rm sm}-D_n^{\rm nat}
 ={}&\int\langle W_n,(\Theta_n-\Theta_n^2)W_n\rangle\,\dd t\\
 &+\int\langle\Theta_nW_n,
       (I-E_{n,t}^2)\Theta_nW_n\rangle\,\dd t
 \ge0.
 \end{aligned}}
 \label{eq:v17-denominator-debit}
\end{equation}
The native centred numerator satisfies
\begin{equation}
 \boxed{
 N_n^{\rm sm}-N_n^{\rm nat}
 =\int\langle z_{n,t}^{\rm half},
       (I-E_{n,t}^2)(\Theta_nW_n)\rangle\,\dd t
  +N_n^{\rm en}.}
 \label{eq:v17-numerator-transfer}
\end{equation}
The first term is the signed heat-metric cross row; the second is the exact
energy-centering transfer.  Neither has a sign forced by the current
hypotheses.
\end{theorem}

\begin{proof}
Because $E_{n,t}$ is a self-adjoint contraction,
\[
 D_n^{\rm nat}
 =\int\langle\Theta_nW_n,E_{n,t}^2\Theta_nW_n\rangle\,\dd t.
\]
Insert $\|\Theta_n^{1/2}W_n\|^2$ and
$\|\Theta_nW_n\|^2$ between $D_n^{\rm sm}$ and $D_n^{\rm nat}$.
The first difference is the first line of
\eqref{eq:v17-denominator-debit}; the second is the heat-metric term.
Both are nonnegative because $0\le\Theta_n\le1$ and $0\le E_{n,t}^2\le I$.

For the numerator,
\[
 \int\langle\widehat Z_{n,t}^{\rm cal},W_{n,t}^{\rm nat}\rangle\,\dd t
 =\int\langle z_{n,t}^{\rm half},
             E_{n,t}^2(\Theta_nW_n)\rangle\,\dd t.
\]
Subtract the energy component \eqref{eq:v17-energy-transfer-row} to obtain
$N_n^{\rm nat}$, then subtract from \eqref{eq:v17-smoothed-numerator}.
\end{proof}

Define the normalized V16 smoothed readout and the metric-incidence defect by
\begin{equation}
 S_n^{\rm sm}:=-\sigma_*A_n\frac{\mathfrak m_n}{\mathfrak d_n},
 \qquad
 \boxed{\Delta_n^{\rm met}:=S_n^{\rm sm}-S_n^{\rm nat}.}
 \label{eq:v17-metric-defect}
\end{equation}
Let
\begin{equation}
 I_n^{\rm sm}
 :=\mathfrak c_n^{\rm cal}-\mathfrak c_n^{\rm rep}-S_n^{\rm sm}
 \label{eq:v17-smoothed-incidence}
\end{equation}
be the V16 source-to-head readout.

\begin{corollary}[The source-to-head verdict moves by exactly the metric defect]
\label{cor:v17-incidence-transfer}
One has
\begin{equation}
 \boxed{
 I_n^{\rm nat}=I_n^{\rm sm}+\Delta_n^{\rm met}.}
 \label{eq:v17-incidence-transfer}
\end{equation}
In particular, a V16 smoothed clean card is a native clean card if and only if
$\Delta_n^{\rm met}\to0$ in addition to the V16 clean limits.  The finite
counterexample of Countertheorem~\ref{cth:v17-raw-clean-counterexample}
shows that this extra condition is not an algebraic consequence of Gram
positivity or energy neutrality.
\end{corollary}

\begin{proof}
Subtract \eqref{eq:v17-native-incidence-readout} from
\eqref{eq:v17-smoothed-incidence} and use
\eqref{eq:v17-metric-defect}.
\end{proof}

\begin{firewall}[No later positive card repairs a metric-incidence defect]
\label{fire:v17-no-late-repair}
The positive Reynolds covariance remains an actual carrier of the smoothed
source $H_{n,t}$.  The native pullback $G_{n,t}^{\rm nat}$ is a deterministic
source-metric follower of that same carrier, not a second source.  A nonzero
$\Delta_n^{\rm met}$ says that the half-heat projection and the native
energy-centred payment allocate the common scalar differently.  Paid-stress
variation, terminal compactness, or a later finite-stock estimate cannot
retroactively identify those allocations.
\end{firewall}

% =============================================================================
\section{The corrected mixed row is still an ordered Duhamel moment}
\label{sec:v17-native-moment}
% =============================================================================

Retain the common-endpoint carried state and cumulative response of V16:
\begin{equation}
 X_n=S_ne(\beta_n),
 \qquad
 \mathcal F_n(t)
 =-\int_{\beta_n}^{t}
   \Lambda E_{n,s}c(s)\,\dd s,
 \qquad
 E_{n,t}e(t)=X_n+\mathcal F_n(t).
 \label{eq:v17-cumulative-response}
\end{equation}
Define the uncentred native numerator
\begin{equation}
 N_n^{\rm nat,unc}
 :=\int
 \langle\widehat Z_{n,t}^{\rm cal},W_{n,t}^{\rm nat}\rangle\,\dd t.
 \label{eq:v17-native-uncentred}
\end{equation}

\begin{theorem}[Native carried-state, ordered-moment, and energy split]
\label{thm:v17-native-ordered-moment}
The corrected numerator has the exact decomposition
\begin{equation}
 \boxed{
 N_n^{\rm nat}
 =N_n^{\rm nat,old}
  +N_n^{\rm nat,mom}
  -N_n^{\rm en},}
 \label{eq:v17-native-moment-split}
\end{equation}
where
\begin{align}
 N_n^{\rm nat,old}
 &:=\int
 \langle\Lambda E_{n,t}X_n,
         E_{n,t}(\Theta_nW_n)\rangle\,\dd t,
 \label{eq:v17-native-old}\\
 N_n^{\rm nat,mom}
 &:=\int
 \langle\Lambda E_{n,t}\mathcal F_n(t),
         E_{n,t}(\Theta_nW_n)\rangle\,\dd t.
 \label{eq:v17-native-mom}
\end{align}
The ordered term is the original native-source pairing
\begin{equation}
 \boxed{
 N_n^{\rm nat,mom}
 =-\int_{\beta_n}^{b_n}
 \langle c(s),Q_n^{\rm nat}(s)\rangle\,\dd s,}
 \label{eq:v17-native-source-moment}
\end{equation}
with follower writer
\begin{equation}
 \boxed{
 Q_n^{\rm nat}(s)
 :=E_{n,s}\Lambda^2
   \int_s^{b_n}
   E_{n,t}^2(\Theta_n(t)W_n)\,\dd t.}
 \label{eq:v17-native-moment-writer}
\end{equation}
It obeys the uniform pointwise estimate
\begin{equation}
 \boxed{
 \|Q_n^{\rm nat}(s)\|_2
 \le\frac{2\Gamma_n}{e\nu}}
 \label{eq:v17-native-moment-bound}
\end{equation}
whenever $s$ precedes the selected support.  Consequently
\begin{equation}
 \boxed{
 |N_n^{\rm nat,mom}|
 \le\frac{2\Gamma_n}{e\nu}
       \int_{\beta_n}^{b_n}\|c(s)\|_2\,\dd s.}
 \label{eq:v17-native-moment-source-length}
\end{equation}
\end{theorem}

\begin{proof}
Since
$\widehat Z_{n,t}^{\rm cal}=\Lambda E_{n,t}(X_n+\mathcal F_n(t))$,
linearity gives the first two terms in
\eqref{eq:v17-native-moment-split}; subtracting the centred energy row gives
the last.  Insert the Duhamel formula for $\mathcal F_n(t)$ into
\eqref{eq:v17-native-mom}, use Fubini, and move the commuting self-adjoint
multipliers to the writer.  This gives
\eqref{eq:v17-native-source-moment}--\eqref{eq:v17-native-moment-writer}.

Let $r=b_n-s$.  The spectral multiplier bound
\[
 \|E_{n,s}\Lambda^2\|_{2\to2}
 =\sup_{\kappa\ge0}\kappa^2e^{-\nu r\kappa^2}
 =\frac1{e\nu r}
\]
and heat contraction give
\[
 \|Q_n^{\rm nat}(s)\|_2
 \le\frac1{e\nu r}
     \int_s^{b_n}\|\Theta_n(t)W_n\|_2\,\dd t
 \le\frac{\|W_n\|_2}{e\nu}
 \le\frac{2\Gamma_n}{e\nu}.
\]
The source-length estimate is Cauchy--Schwarz in space followed by time
integration.
\end{proof}

\begin{assessment}[Exact remaining PDE estimate]
\label{ass:v17-moment-frontier}
The native ordered writer differs from the V16 follower by the extra
$E_{n,t}^2$ inside the Volterra integral, and the physical coefficient also
contains the energy-centering row $N_n^{\rm en}$.  The estimate
\eqref{eq:v17-native-moment-source-length} recovers the already finite
frequency-discounted source debit; it gives no sign and no contradiction.
A genuinely new terminal theorem must control the signed combination in
\eqref{eq:v17-native-moment-split}, or must show that the metric-incidence
defect itself is a typed surviving packet.
\end{assessment}

% =============================================================================
\section{Corrected branch theorem and NS--A stopping boundary}
\label{sec:v17-branch}
% =============================================================================

\begin{theorem}[Metric-first native terminal-writer alternative]
\label{thm:v17-branch}
After passage to a priority subsequence, exactly one of the following occurs.
\begin{enumerate}[label=\textup{(V17.\arabic*)},leftmargin=5.9em]
\item \emph{Source-metric incidence defect.}  There is $\delta>0$ such that
      \begin{equation}
       \boxed{
       |\Delta_n^{\rm met}|
       \ge\delta\mathfrak c_n^{\rm cal}.}
       \label{eq:v17-metric-survivor}
      \end{equation}
      The half-heat Reynolds allocation and the native energy-centred payment
      differ by a fixed record fraction.  The exact positive denominator
      debits and signed numerator rows are
      \eqref{eq:v17-denominator-debit}--\eqref{eq:v17-numerator-transfer}.
\item \emph{Represented endpoint return.}  The metric defect tends to zero,
      and for some $\delta>0$,
      \begin{equation}
       \boxed{
       |\mathfrak c_n^{\rm rep}|
       \ge\delta\mathfrak c_n^{\rm cal}.}
       \label{eq:v17-native-represented}
      \end{equation}
      The response is returned to its represented source before any new
      payment is assigned.
\item \emph{Native aggregate incidence residual.}  The first two branches
      fail and, for some $\delta>0$,
      \begin{equation}
       \boxed{
       |I_n^{\rm nat}|
       \ge\delta\mathfrak c_n^{\rm cal}.}
       \label{eq:v17-native-incidence-survivor}
      \end{equation}
      Only then is the aggregate native complement
      $\mathcal C_n^{\rm nat,comp}+\mathcal U_n^{\rm nat,wr}$ opened.
\item \emph{Source-native clean export.}  The preceding three branches fail:
      \begin{equation}
       \boxed{
       \Delta_n^{\rm met}\to0,
       \qquad
       \mathfrak c_n^{\rm rep}\to0,
       \qquad
       I_n^{\rm nat}\to0.}
       \label{eq:v17-native-clean}
      \end{equation}
      Then
      \begin{equation}
       \boxed{
       S_n^{\rm nat}-\mathfrak c_n^{\rm cal}\to0,}
       \label{eq:v17-native-clean-relation}
      \end{equation}
      and the source-native selected contact is positive, record quadratic,
      and non-summable.
\end{enumerate}
No branch is chosen from the presently attached theorem layers, because the
new native mixed row and the metric-incidence residual are not populated on a
hypothetical singular trajectory.
\end{theorem}

\begin{proof}
First select according to whether
$|\Delta_n^{\rm met}|/\mathfrak c_n^{\rm cal}$ has a positive lower bound or
tends to zero.  On the second branch, apply the same priority selection to
$|\mathfrak c_n^{\rm rep}|/\mathfrak c_n^{\rm cal}$ and then to
$|I_n^{\rm nat}|/\mathfrak c_n^{\rm cal}$.  Since
$3/2\le\mathfrak c_n^{\rm cal}\le2$, failure of all three positive-lower-bound
branches gives \eqref{eq:v17-native-clean}.  The exact readout
\eqref{eq:v17-native-incidence-readout} then gives
\eqref{eq:v17-native-clean-relation}.
\end{proof}

\begin{corollary}[Quantitative source-native clean branch]
\label{cor:v17-native-clean}
On branch \textup{(V17.4)}, for all late $n$,
\begin{equation}
 \boxed{
 -\sigma_*\frac{\mathfrak m_n^{\rm nat}}
                    {\mathfrak d_n^{\rm nat}}
 \ge\frac1{\mathfrak V_*},}
 \label{eq:v17-native-alignment-floor}
\end{equation}
\begin{equation}
 \boxed{
 \mathcal C_n^{\rm nat,sel}
 \ge\frac34\mathcal R_n^2,
 \qquad
 \sum_n\mathcal C_n^{\rm nat,sel}=+\infty.}
 \label{eq:v17-native-nonsummable}
\end{equation}
The export packet must retain both the positive smoothed covariance source
$H_{n,t}$ and its source-native pullback $G_{n,t}^{\rm nat}$, together with
$E_{n,t}$ and the vanishing metric-incidence residual.  They are two metrics
on one occurring source, not two source births.
\end{corollary}

\begin{proof}
Equation~\eqref{eq:v17-native-clean-relation}, the lower bound
$\mathfrak c_n^{\rm cal}\ge3/2$, and $A_n\le\mathfrak V_*$ give the alignment
floor.  The contact and divergence conclusions follow exactly as in the V16
clean branch, with the native readout replacing the smoothed one.
\end{proof}

\begin{definition}[Source-native V17 export card]
\label{def:v17-export-card}
On branch \textup{(V17.4)} retain
\begin{equation}
 \boxed{
 \begin{aligned}
 \mathfrak C_n^{\rm V17}:=\bigl(&b_n,I_n^{\rm ann},\Gamma_n,
 \Theta_n,E_{n,\cdot},P_n^{\rm rep},\\
 &H_{n,\cdot},G_{n,\cdot}^{\rm nat},
 W_{n,\cdot}^{\rm nat},Z_{n,\cdot}^{\rm dip},
 \mathbb G_n^{\rm nat},\\
 &\mathcal C_n^{\rm cal},\mathcal C_n^{\rm rep},
 \mathfrak s_n,\kappa_n^{\rm nat},
 \mathcal C_n^{\rm nat,sel},I_n^{\rm nat},
 \Delta_n^{\rm met}\bigr).
 \end{aligned}}
 \label{eq:v17-export-card}
\end{equation}
This packet is the source-native refinement of the V16 card.  It preserves
the original positive Reynolds carrier and records the exact metric map by
which it is interpreted in the energy-neutral critical source space.
\end{definition}

% =============================================================================
\section{Version V17 integrated theorem and proof-status ledger}
\label{sec:v17-integrated}
% =============================================================================

\begin{theorem}[Version V17 native-metric and caloric-dipole theorem]
\label{thm:v17-integrated}
Preserving every algebraic theorem and stated limitation of Versions~V1--V16,
Version~V17 proves the following.
\begin{enumerate}[label=\textup{(T17.\arabic*)},leftmargin=5.5em]
\item The complete caloric contact has the unique native energy-neutral
      writer \eqref{eq:v17-energy-centred-dipole}.  Its action is the exact
      caloric frequency variance \eqref{eq:v17-caloric-frequency-variance}.
\item The native source action required by a fixed contact is
      \eqref{eq:v17-native-action-lower}; a caloric-flat packet carries zero
      contact.
\item In the half-heat factorization the source-null direction is transported
      by backward heat and the centering is orthogonal in the native
      transported metric, not automatically in the ordinary half-heat metric.
\item A heat-compatible two-mode model has a raw clean half-heat allocation
      and zero source-native centred mixed coefficient.  Hence the missing
      metric incidence is logically independent of the V16 positive Gram.
\item The selected annular source admits the exact source-first pullback
      \eqref{eq:v17-native-source-writer}.  Annular support makes its inverse
      heat uniformly bounded, and the native source action and writer
      denominator have the fixed windows
      \eqref{eq:v17-native-source-action-upper}--
      \eqref{eq:v17-native-denominator-window}.
\item The native source, selected writer, and energy-centred caloric writer
      form the compact positive Gram \eqref{eq:v17-native-three-Gram}.  Two
      positive actions reconstruct its mixed row.
\item The complete contact has the source-first ledger
      \eqref{eq:v17-native-source-ledger}, and the corrected five-scalar card
      gives the exact native selected and source-to-head readouts.
\item The smoothed and native denominators differ by two positive debits,
      while their numerators differ by the signed heat-metric cross and the
      energy-centering transfer.  The source-to-head verdict moves by exactly
      $\Delta_n^{\rm met}$.
\item The source-native mixed row is carried old state plus the ordered
      Duhamel moment \eqref{eq:v17-native-source-moment}, minus the energy row.
      Its follower has the sharp bound
      \eqref{eq:v17-native-moment-bound}, but no sign is obtained.
\item The physical priority order is metric incidence, represented endpoint,
      native aggregate residual, and only then the source-native clean export
      \eqref{eq:v17-export-card}.
\end{enumerate}
No terminal Dini gain, finite-stock upper budget, backward-uniqueness theorem,
Liouville theorem, or global three-dimensional regularity conclusion follows.
\end{theorem}

\begin{proof}
Items~\textup{(T17.1)--(T17.2)} are
\Cref{thm:v17-native-dipole,cor:v17-native-action}.  Items
\textup{(T17.3)--(T17.4)} are
\Cref{prop:v17-half-energy-null,cth:v17-raw-clean-counterexample}.  Item
\textup{(T17.5)} is
\Cref{thm:v17-native-pullback,cor:v17-native-efficiency}.  Items
\textup{(T17.6)--(T17.7)} are
\Cref{lem:v17-dipole-window,thm:v17-native-three-Gram,prop:v17-native-selected-contact,thm:v17-native-positive-read}.
Item~\textup{(T17.8)} is
\Cref{thm:v17-metric-bridge,cor:v17-incidence-transfer}.  Item
\textup{(T17.9)} is \Cref{thm:v17-native-ordered-moment}.  The final item is
\Cref{thm:v17-branch,cor:v17-native-clean,def:v17-export-card}.
\end{proof}

\begin{firewall}[Exact stopping boundary after V17]
\label{fire:v17-stop}
The V16 finite card was not the end of the upstream audit because its source
metric had not been compared with the native energy-null quotient.  V17
closes that structural gap and returns one explicit residual
$\Delta_n^{\rm met}$.  A later source norm, profile compactness theorem,
paid-stress estimate, or finite-stock argument cannot make this residual
vanish.  The next admissible work is an actual Navier--Stokes estimate for the
signed rows in \eqref{eq:v17-numerator-transfer} or
\eqref{eq:v17-native-moment-split}, a populated native five-card, or the one
already typed metric/represented/native-complement branch.  The clean native
branch has terminated in an NS--B export.
\end{firewall}

\begingroup\small
\begin{longtable}{p{0.25\textwidth}p{0.19\textwidth}p{0.48\textwidth}}
\caption{NS--A Version V17 proof-status ledger.}\label{tab:v17-ledger}\\
\toprule Row & Status & Exact conclusion\\\midrule
\endfirsthead
\toprule Row & Status & Exact conclusion\\\midrule
\endhead
V1--V16 prefix & \PROVED{} preserved
 &The complete V16 preterminal source is retained byte-for-byte; no earlier
 formula is deleted.\\[0.3em]
Native caloric dipole & \PROVED{} exact
 &The complete contact is represented by the unique minimum native writer
 orthogonal to the energy vector.\\[0.3em]
Caloric frequency variance & \PROVED{} exact
 &The dipole action is the variance of
 $\kappa e^{-2\nu r\kappa^2}$ under the critical-energy law.\\[0.3em]
Caloric-flat packet & \PROVED{} zero contact
 &Zero native variance forces the complete caloric contact integrand to
 vanish.\\[0.3em]
Half-heat energy-null direction & \PROVED{} finite-cutoff
 &The null direction is $E_r^{-1}e$ and centering is orthogonal in the
 transported native metric.\\[0.3em]
Raw half-heat clean allocation & \OBSTRUCTION{} explicit
 &A two-mode heat-compatible card can be raw clean while its native centred
 mixed coefficient is zero.\\[0.3em]
Native annular pullback & \PROVED{} exact
 &Inverse heat is applied only before spatial localization and is uniformly
 bounded on the selected parabolic annulus.\\[0.3em]
Native source action & \PROVED{} bounded
 &$\mathfrak a_n^{\rm nat}\le e^{8\tau_+}L_*$.\\[0.3em]
Native writer denominator & \PROVED{} two-sided
 &$d_{\rm nat,*}/\nu\le D_n^{\rm nat}\le d_{\rm nat}^*/\nu$.\\[0.3em]
Native three-column Gram & \PROVED{} positive compact
 &Source, selected writer, and energy-centred caloric writer give
 \eqref{eq:v17-native-three-Gram}.\\[0.3em]
Native mixed acquisition & \PROVED{} positive polarization
 &Two positive actions determine $\mathfrak m_n^{\rm nat}$ and its sign.\\[0.3em]
Metric denominator debit & \PROVED{} positive
 &Smooth localization slack and heat-metric loss add exactly to
 $D_n^{\rm sm}-D_n^{\rm nat}$.\\[0.3em]
Metric numerator transfer & \PROVED{} signed
 &The heat cross and energy-centering row give
 \eqref{eq:v17-numerator-transfer}; no sign is forced.\\[0.3em]
Corrected ordered moment & \PROVED{} exact
 &The native row is old carry plus a heat-weighted Volterra source moment,
 minus the energy-centering transfer.\\[0.3em]
Native five-card population & \OPEN{} finite acquisition
 &No hypothetical singular trajectory in the attachments supplies the
 corrected positive reads.\\[0.3em]
Metric-incidence estimate & \OPEN{} upstream PDE row
 &No current theorem makes $\Delta_n^{\rm met}$ vanish or selects its sign.\\[0.3em]
Clean NS--A export & \CONDITIONAL{} corrected
 &Export requires the source-native card, or the V16 card together with
 vanishing metric incidence.\\[0.3em]
Three-dimensional regularity & \OPEN
 &No surviving metric, represented, native-complement, or replenishment
 packet is excluded for every datum.\\
\bottomrule
\end{longtable}
\endgroup

% =============================================================================
\section{Exact mandate after Version V17}
\label{sec:v17-next}
% =============================================================================

\begin{programme}[Evaluate the metric bridge or stop on its survivor]
\label{prog:v17-next}
Preserve Versions~V1--V17 verbatim.  There is no automatic compiler-style
Version~V18.

A continuation is licensed only in the following order.
\begin{enumerate}[label=\textup{(V18.\arabic*)},leftmargin=5.8em]
\item Evaluate the two positive denominator debits and the two signed
      numerator rows in
      \eqref{eq:v17-denominator-debit}--\eqref{eq:v17-numerator-transfer} on
      the literal selected history.  A nonzero normalized value is the
      source-metric incidence survivor and stops the clean export.
\item Independently test the caloric frequency variance
      \eqref{eq:v17-caloric-frequency-variance}.  Variance collapse with a
      fixed contact is a native source-action overdrive, not a new compact
      profile.
\item If the metric bridge passes, populate the corrected card
      \eqref{eq:v17-five-card} by the positive reads of
      \eqref{eq:v17-native-positive-actions} and the represented endpoint
      read.  Apply the branch theorem before opening any component.
\item On a native aggregate-residual branch, follow only the already selected
      source complement or writer-angle residual.  Do not return to the
      smoothed card to relabel it.
\item On the source-native clean branch, terminate NS--A and continue only in
      NS--B with \eqref{eq:v17-export-card} and an independently occurring
      paid-stress or finite-stock upper budget.
\item A new analytical continuation without populated data must control the
      actual signed combination
      \eqref{eq:v17-native-moment-split}; a bound for its separate absolute
      terms, another endpoint norm, or another source-action lower estimate is
      inadmissible.
\end{enumerate}
No new Dini modulus, response graph, path measure, temporal partition,
abstract stock, slow-time quotient, or uncentred terminal-writer projection
is admitted in place of this metric incidence.
\end{programme}

\begin{assessment}[Programme position after V17]
\label{ass:v17-position}
The bounded-action branch has reached its genuinely upstream finite boundary.
The selected positive Reynolds source and the native critical payment now
live on one card with their heat-metric map displayed.  Energy neutrality is
implemented before component pricing, the exact metric transfer is explicit,
and a raw clean card is known not to imply a native clean card.  The remaining
unknown is one physical signed metric/ordered-moment estimate or one populated
native card.  Continuing with another generic concentration or compactness
layer would be circular.
\end{assessment}

\paragraph{Version V17 history.}
Preserved the complete V16 source.  Reopened only the physical interpretation
of its clean export after identifying an upstream source-metric omission.
Constructed the native energy-neutral finite-age caloric dipole and its exact
frequency-variance action, transported the energy-null direction through the
half-heat factorization, and gave a heat-compatible two-mode counterexample
to raw-card invariance.  Pulled the selected annular source back before
spatial localization, proved uniform native action and writer windows,
constructed the compact native three-column Gram and positive mixed-row
acquisition, and derived the exact positive denominator and signed numerator
metric-transfer formulas.  Rewrote the corrected mixed row as carried state
plus its native ordered Duhamel moment minus the energy row.  Finally inserted
the metric defect ahead of the represented/native-incidence/clean priority
and stated the corrected NS--A to NS--B export.

\begin{assessment}[Append-only integrity]
\label{ass:v17-integrity}
The complete Version~V16 source preceding its sole final document terminator
is preserved verbatim.  Its preterminal SHA--256 digest is
\begin{center}\scriptsize\ttfamily
5512fcd5db7b6c4f9cac9b15ac49233baa2dedf9234a1e757ad4ee45d963708c
\end{center}
The present material begins at Section~\ref{sec:v17-scope}; the final command
below is again unique.
\end{assessment}

\addtocontents{toc}{\protect\endgroup}
% =============================================================================
% END VERSION V17 MATERIAL -- append later versions before final terminator
% =============================================================================

% APPEND ALL LATER VERSIONS IMMEDIATELY ABOVE THE SOLE FINAL LINE BELOW.

\clearpage
% =============================================================================
% VERSION V18: NEW MATERIAL.  The complete V1--V17 source above is preserved
% verbatim; only its sole active \end{document} line was removed before append.
% =============================================================================
\hypersetup{pdftitle={NS-A Terminal Source Concentration Programme: Cumulative V18},pdfauthor={A. Pelissier}}
\makeatletter
\addtocontents{toc}{\protect\begingroup\protect\makeatletter}
\addtocontents{toc}{\protect\def\protect\l@section{\protect\@dottedtocline{1}{0em}{4.7em}}}
\addtocontents{toc}{\protect\def\protect\l@subsection{\protect\@dottedtocline{2}{1.5em}{5.3em}}}
\addtocontents{toc}{\protect\makeatother}
\makeatother
\renewcommand{\thesection}{V18.\arabic{section}}
\renewcommand{\thesubsection}{V18.\arabic{section}.\arabic{subsection}}
\renewcommand{\theequation}{V18.\arabic{equation}}
\setcounter{section}{0}
\setcounter{equation}{0}
\renewcommand{\theHsection}{V18.\arabic{section}}
\renewcommand{\theHsubsection}{V18.\arabic{section}.\arabic{subsection}}
\renewcommand{\theHtheorem}{V18.\arabic{section}.\arabic{theorem}}
\renewcommand{\theHequation}{V18.\arabic{equation}}

\begin{center}
{\LARGE\bfseries NS--A: Version V18}\\[0.5em]
{\large The Energy--Helicity Source-Null Quotient, Helical Caloric Regression,\\
and the Corrected Rank-Three Terminal Gate}\\[0.5em]
{\normalsize 3 September 2026}
\end{center}
\addcontentsline{toc}{part}{Version V18: the energy--helicity quotient and helical rank-three gate}

\begin{abstract}
Version~V17 moved the terminal writer into the native critical-source metric
and removed the exact energy-null line generated by
$e=A^{1/4}u$.  That was a necessary correction, but it was not yet the full
universal quadratic-invariant quotient of the Navier--Stokes nonlinearity.
For every smooth divergence-free periodic field, the projected quadratic
source is also orthogonal to the vorticity.  In critical coordinates,
\[
 c=A^{-1/4}\mathcal B(u),\qquad
 e=A^{1/4}u,\qquad
 h=A^{1/4}\operatorname{curl}u=\operatorname{curl}e,
\]
one has the two exact identities
\[
 \langle c,e\rangle=0,\qquad \langle c,h\rangle=0.
\]
A later metric comparison or positive Reynolds Gram cannot repair omission
of this second source-null direction.

Version~V18 therefore goes one step upstream.  It constructs the
Moore--Penrose projection onto the possibly singular plane
$\operatorname{span}\{e,h\}$ and replaces the V17 energy-centred caloric
writer by its orthogonal complement to this plane.  The complete caloric
contact is unchanged, while the writer action decreases by an exact positive
helicity debit.  The V17 frequency variance is consequently refined to a
linear-regression residual on the signed helical frequency
$X=\sigma|k|$:
\[
 \frac{\|Z^{\rm EH}\|_2^2}{\|e\|_2^2}
 =\inf_{a,b\in\mathbb R}
   \mathbb E\left|F_r-a-bX\right|^2,
 \qquad
 F_r=|X|e^{-2\nu rX^2}.
\]
When $\operatorname{Var}X>0$, this is
\[
 \operatorname{Var}F_r-
 \frac{\operatorname{Cov}(F_r,X)^2}{\operatorname{Var}X}.
\]
Thus the helicity short removes exactly the part of the energy-only caloric
variance explained by signed curl frequency.

This formula gives an exact finite-rank obstruction.  Any field occupying at
most two distinct signed helical frequencies has zero energy--helicity
caloric writer at every fixed age.  A nonzero contact therefore requires at
least three signed helical levels.  For exactly three levels, V18 evaluates
the writer action as the square of one second divided difference divided by
an explicit positive weight.  At zero age, the multiplier is $|X|$; hence a
nonzero instantaneous critical work requires both helicity signs and at
least three distinct signed levels.  An explicit smooth two-frequency
homochiral field has nonzero quadratic nonlinearity and positive V17
energy-only variance, but its complete energy--helicity caloric writer is
identically zero.  This proves that the extra quotient is not cosmetic.

On the selected terminal history, V18 inserts the helicity short into the
V17 three-column Gram.  The old and new mixed rows differ by one explicit
signed helicity-transfer term.  If $\Delta_n^{\rm hel}$ denotes the induced
change of the selected-contact readout, then the complete source-to-head
identity is
\[
 I_n^{\rm EH}
 =I_n^{\rm nat}+\Delta_n^{\rm hel}
 =I_n^{\rm sm}+\Delta_n^{\rm met}+\Delta_n^{\rm hel}.
\]
Accordingly, the correct priority is now: helicity incidence, native metric
incidence, represented endpoint return, invariant aggregate residual, and
only then a clean NS--A export.  On the clean branch the exported contact is
again of order $\mathcal R_n^2$ and is non-summable, but it now carries a
literal three-level signed-helical gate for the NS--B donor audit.  This
helical rank is not the spatial lattice rank of the two-triad donor complex;
the two conditions are complementary and neither is substituted for the
other.  No terminal Dini gain, finite-stock upper bound, Liouville theorem,
or global regularity conclusion is claimed.
\end{abstract}

% =============================================================================
\section{Scope: universal quadratic invariants precede the native metric card}
\label{sec:v18-scope}
% =============================================================================

Retain the smooth fixed-datum record history and all notation of
Versions~V15--V17.  In particular,
\begin{equation}
 e(t)=A^{1/4}u(t),\qquad
 c(t)=A^{-1/4}\BNS(u(t)),\qquad
 E_{n,t}=e^{-\nu(b_n-t)A},
 \label{eq:v18-standing}
\end{equation}
and the complete native caloric writer is
\begin{equation}
 \widehat Z_{n,t}^{\rm cal}=\Lambda E_{n,t}^2e(t).
 \label{eq:v18-caloric-writer}
\end{equation}
Version~V17 formed the energy-centred writer
$Z_{n,t}^{\rm dip}$ by subtracting the projection of
$\widehat Z_{n,t}^{\rm cal}$ onto $\mathbb Re(t)$.  The present version
preserves that identity and tests whether $\mathbb Re(t)$ is the complete
universal source-null carrier relevant to the quadratic nonlinearity.

Put
\begin{equation}
 \omega(t):=\operatorname{curl}u(t),\qquad
 h(t):=A^{1/4}\omega(t)=\operatorname{curl}e(t).
 \label{eq:v18-helicity-vector}
\end{equation}
All fields in this section are smooth and mean zero.  The powers of $A$, the
Leray projection, and curl commute on the periodic divergence-free Fourier
carrier.

\begin{firewall}[Refinement of the V17 minimum, not deletion of it]
\label{fire:v18-v17-status}
Theorem~\ref{thm:v17-native-dipole} remains exact: it gives the unique
minimum modulo the one-dimensional energy line.  V18 proves that the
quadratic source has another universal null direction.  The V17 writer is
therefore an energy-centred representative, but it is not in general the
minimum representative modulo the complete energy--helicity invariant
carrier.  Every V17 algebraic identity is retained; only its physical clean
export acquires one earlier incidence gate.
\end{firewall}

% =============================================================================
\section{The exact energy--helicity source-null carrier}
\label{sec:v18-null-plane}
% =============================================================================

\begin{theorem}[Simultaneous nonlinear energy and helicity cancellation]
\label{thm:v18-energy-helicity-null}
For every smooth mean-zero divergence-free periodic field $u$,
\begin{equation}
 \boxed{
 \langle\BNS(u),u\rangle=0,
 \qquad
 \langle\BNS(u),\operatorname{curl}u\rangle=0.}
 \label{eq:v18-EH-cancellation}
\end{equation}
Equivalently, in the critical source coordinates of
\eqref{eq:v18-standing}--\eqref{eq:v18-helicity-vector},
\begin{equation}
 \boxed{
 \langle c,e\rangle=0,
 \qquad
 \langle c,h\rangle=0.}
 \label{eq:v18-critical-EH-null}
\end{equation}
If $u\ne0$, then $e\ne0$ and $h\ne0$; the two vectors may nevertheless be
linearly dependent.
\end{theorem}

\begin{proof}
Since $u$ is divergence free and periodic,
\[
 \langle\BNS(u),u\rangle
 =\int_{\Torus^3}(u\cdot\nabla)u\cdot u\,\dd x
 =\frac12\int_{\Torus^3}u\cdot\nabla|u|^2\,\dd x=0.
\]
The vorticity $\omega=\operatorname{curl}u$ is divergence free.  The vector
identity
\[
 (u\cdot\nabla)u
 =\nabla\frac{|u|^2}{2}-u\times\omega
\]
gives
\[
 \int_{\Torus^3}(u\cdot\nabla)u\cdot\omega\,\dd x
 =-\int_{\Torus^3}(u\times\omega)\cdot\omega\,\dd x=0;
\]
the gradient term vanishes against the divergence-free field $\omega$.
The Leray projection may be removed in both pairings.  Self-adjointness of
the powers of $A$ then gives
\[
 \langle c,e\rangle
 =\langle A^{-1/4}\BNS(u),A^{1/4}u\rangle
 =\langle\BNS(u),u\rangle,
\]
and similarly
$\langle c,h\rangle=\langle\BNS(u),\omega\rangle$.

If $e=0$, then $u=0$.  If $h=0$, then $\operatorname{curl}u=0$.  On every
nonzero Fourier mode, $k\cdot\widehat u(k)=0$ and
$k\times\widehat u(k)=0$ force $\widehat u(k)=0$; mean zero then gives
$u=0$.  Linear dependence of $e$ and $h$ occurs, for example, on a single
Beltrami eigenshell.
\end{proof}

For a fixed smooth time, define the synthesis
\begin{equation}
 \mathsf N_t:\mathbb R^2\longrightarrow L^2_\sigma(\Torus^3),
 \qquad
 \mathsf N_t(a,b)=ae(t)+bh(t),
 \label{eq:v18-null-synthesis}
\end{equation}
and its positive Gram
\begin{equation}
 \mathsf G_t^{\rm EH}:=\mathsf N_t^*\mathsf N_t
 =\begin{pmatrix}
   \|e(t)\|_2^2&\langle e(t),h(t)\rangle\\
   \langle e(t),h(t)\rangle&\|h(t)\|_2^2
  \end{pmatrix}.
 \label{eq:v18-null-Gram}
\end{equation}
The Moore--Penrose support projection is
\begin{equation}
 \boxed{
 \mathsf P_t^{\rm EH}
 :=\mathsf N_t(\mathsf G_t^{\rm EH})^\dagger\mathsf N_t^*.}
 \label{eq:v18-null-projection}
\end{equation}
It is the orthogonal projection onto
$\mathcal N_t^{\rm EH}:=\operatorname{span}\{e(t),h(t)\}$, including the
rank-one case.

\begin{theorem}[Moore--Penrose energy--helicity caloric quotient]
\label{thm:v18-EH-quotient}
Define
\begin{equation}
 \boxed{
 Z_{n,t}^{\rm EH}
 :=(I-\mathsf P_t^{\rm EH})\widehat Z_{n,t}^{\rm cal}.}
 \label{eq:v18-EH-writer}
\end{equation}
Then
\begin{equation}
 \boxed{
 Z_{n,t}^{\rm EH}\perp e(t),\qquad
 Z_{n,t}^{\rm EH}\perp h(t),}
 \label{eq:v18-EH-orthogonality}
\end{equation}
\begin{equation}
 \boxed{
 \mathcal J_{b_n}(t)
 =-\langle c(t),Z_{n,t}^{\rm EH}\rangle.}
 \label{eq:v18-EH-contact}
\end{equation}
For every $y\in\mathcal N_t^{\rm EH}$,
\begin{equation}
 \boxed{
 \|\widehat Z_{n,t}^{\rm cal}-y\|_2^2
 =\|Z_{n,t}^{\rm EH}\|_2^2
  +\|\mathsf P_t^{\rm EH}\widehat Z_{n,t}^{\rm cal}-y\|_2^2.}
 \label{eq:v18-EH-Pythagoras}
\end{equation}
Hence $Z_{n,t}^{\rm EH}$ is the unique minimum-norm representative of the
caloric writer modulo the universal invariant carrier
$\mathcal N_t^{\rm EH}$.
\end{theorem}

\begin{proof}
The Moore--Penrose formula in \eqref{eq:v18-null-projection} is the
orthogonal projection onto $\operatorname{Ran}\mathsf N_t$.  This proves
\eqref{eq:v18-EH-orthogonality}.  By
Theorem~\ref{thm:v18-energy-helicity-null}, the source $c(t)$ is orthogonal
to the range of this projection.  Subtracting
$\mathsf P_t^{\rm EH}\widehat Z_{n,t}^{\rm cal}$ therefore leaves the
contact unchanged.  Finally,
$\widehat Z=Z^{\rm EH}+\mathsf P^{\rm EH}\widehat Z$ is an orthogonal
sum, and $y$ lies in the second summand.  Orthogonal Pythagoras gives
\eqref{eq:v18-EH-Pythagoras} and the unique minimum.
\end{proof}

\begin{corollary}[Sharp invariant-relative source action]
\label{cor:v18-EH-action}
At every smooth time,
\begin{equation}
 \boxed{
 |\mathcal J_{b_n}(t)|^2
 \le\|c(t)\|_2^2\|Z_{n,t}^{\rm EH}\|_2^2.}
 \label{eq:v18-EH-action-bound}
\end{equation}
If the contact is nonzero, then
\begin{equation}
 \boxed{
 \|c(t)\|_2^2
 \ge
 \frac{|\mathcal J_{b_n}(t)|^2}
      {\|Z_{n,t}^{\rm EH}\|_2^2}.}
 \label{eq:v18-EH-action-lower}
\end{equation}
If $Z_{n,t}^{\rm EH}=0$, then the complete caloric contact integrand
vanishes.
\end{corollary}

\begin{proof}
Apply Cauchy--Schwarz to \eqref{eq:v18-EH-contact}.  The zero statement is
the same identity.
\end{proof}

\begin{firewall}[Universal invariant carrier versus the accidental full nullspace]
\label{fire:v18-universal-not-accidental}
For one fixed source vector $c(t)$, its complete orthogonal complement is
much larger than $\mathcal N_t^{\rm EH}$ and depends on the value of the
source itself.  V18 does not quotient by that accidental nullspace.  It
removes only the two canonical quadratic-invariant directions furnished by
the equation before the terminal source is inspected.  Further accidental
orthogonality is retained in the ordinary source/writer Riesz residual.
\end{firewall}

% =============================================================================
\section{The positive helicity debit inside the V17 writer}
\label{sec:v18-helicity-debit}
% =============================================================================

Let
\begin{equation}
 \mathsf P_t^{E}f
 :=\frac{\langle f,e(t)\rangle}{\|e(t)\|_2^2}e(t),
 \qquad
 h_\perp(t):=(I-\mathsf P_t^{E})h(t).
 \label{eq:v18-helicity-orthogonal}
\end{equation}
When $h_\perp(t)\ne0$, put
\begin{equation}
 \gamma_n^{\rm hel}(t)
 :=\frac{\langle Z_{n,t}^{\rm dip},h_\perp(t)\rangle}
         {\|h_\perp(t)\|_2^2};
 \label{eq:v18-helicity-coefficient}
\end{equation}
when $h_\perp(t)=0$, set $\gamma_n^{\rm hel}(t)=0$.

\begin{theorem}[Sequential energy and helicity short]
\label{thm:v18-helicity-short}
The joint invariant projection factors as
\begin{equation}
 \mathsf P_t^{\rm EH}
 =\mathsf P_t^E+\mathsf P_{h_\perp(t)},
 \label{eq:v18-projection-factor}
\end{equation}
with the second projection understood as zero when $h_\perp=0$.  Therefore
\begin{equation}
 \boxed{
 Z_{n,t}^{\rm EH}
 =Z_{n,t}^{\rm dip}-\gamma_n^{\rm hel}(t)h_\perp(t),}
 \label{eq:v18-EH-from-E}
\end{equation}
\begin{equation}
 \boxed{
 \|Z_{n,t}^{\rm dip}\|_2^2
 =\|Z_{n,t}^{\rm EH}\|_2^2
  +\frac{|\langle Z_{n,t}^{\rm dip},h_\perp(t)\rangle|^2}
         {\|h_\perp(t)\|_2^2},}
 \label{eq:v18-helicity-debit}
\end{equation}
where the quotient is defined as zero at $h_\perp=0$.  The second term is the
exact positive action overestimate made by stopping after energy centring.
\end{theorem}

\begin{proof}
The vector $h_\perp$ is orthogonal to $e$ and
$\operatorname{span}\{e,h\}=\operatorname{span}\{e,h_\perp\}$.  Hence the
orthogonal projection onto the latter span is the sum of the two orthogonal
rank-one projections, with the second absent in the dependent case.  Since
$Z^{\rm dip}=(I-\mathsf P^E)\widehat Z^{\rm cal}$, applying the second short
gives \eqref{eq:v18-EH-from-E}.  Pythagoras gives
\eqref{eq:v18-helicity-debit}.
\end{proof}

\begin{assessment}[The helicity row is signed only when it meets the selected writer]
\label{ass:v18-helicity-sign}
The action debit in \eqref{eq:v18-helicity-debit} is positive.  Its effect on
a selected mixed writer is not: the relevant cross term is
$\gamma_n^{\rm hel}\langle h_\perp,W_n^{\rm nat}\rangle$, whose sign is not
fixed by energy, helicity, positivity of the Reynolds covariance, or the
separate norm windows.  The terminal-card comparison in
\Cref{sec:v18-terminal-card} therefore retains one signed helicity-incidence
row rather than replacing it by the positive action debit.
\end{assessment}

% =============================================================================
\section{Helical spectral regression of the caloric writer}
\label{sec:v18-helical-regression}
% =============================================================================

For $k\in\mathbb Z^3\setminus\{0\}$ let
$\mathsf C_k=\mathrm i k\times$ on the complex divergence-free plane
$k^\perp$.  It is Hermitian there and has eigenvalues $\sigma|k|$,
$\sigma\in\{+1,-1\}$.  Denote its spectral projections by
\begin{equation}
 \mathsf P_k^\sigma
 =\frac12\left(I+\sigma|k|^{-1}\mathsf C_k\right)
 \quad\text{on }k^\perp,
 \qquad
 \widehat e_\sigma(k)=\mathsf P_k^\sigma\widehat e(k).
 \label{eq:v18-helical-projections}
\end{equation}
Define the probability law
\begin{equation}
 \mu_t(k,\sigma)
 :=\frac{|\widehat e_\sigma(t,k)|^2}{\|e(t)\|_2^2},
 \qquad
 \sum_{k\ne0}\sum_{\sigma=\pm1}\mu_t(k,\sigma)=1,
 \label{eq:v18-helical-law}
\end{equation}
and the two real random variables
\begin{equation}
 X(k,\sigma):=\sigma|k|,
 \qquad
 F_r(k,\sigma):=|k|e^{-2\nu r|k|^2}
                =|X|e^{-2\nu rX^2}.
 \label{eq:v18-X-F}
\end{equation}
Expectations below are taken under $\mu_t$.

\begin{theorem}[Exact helical regression formula]
\label{thm:v18-helical-regression}
Let
\begin{equation}
 \mathsf M_t
 :=\begin{pmatrix}
  1&\mathbb EX\\
  \mathbb EX&\mathbb EX^2
 \end{pmatrix},
 \qquad
 \mathsf v_{t,r}
 :=\binom{\mathbb EF_r}{\mathbb E(XF_r)}.
 \label{eq:v18-regression-data}
\end{equation}
Then
\begin{equation}
 \boxed{
 \frac{\|Z_{n,t}^{\rm EH}\|_2^2}{\|e(t)\|_2^2}
 =\inf_{a,b\in\mathbb R}\mathbb E|F_{r_n(t)}-a-bX|^2
 =\mathbb EF_{r_n(t)}^2
  -\mathsf v_{t,r_n(t)}^{\mathsf T}
   \mathsf M_t^\dagger\mathsf v_{t,r_n(t)}.}
 \label{eq:v18-regression-residual}
\end{equation}
If $\operatorname{Var}X>0$, this becomes
\begin{equation}
 \boxed{
 \frac{\|Z_{n,t}^{\rm EH}\|_2^2}{\|e(t)\|_2^2}
 =\operatorname{Var}F_{r_n(t)}
  -\frac{\operatorname{Cov}(F_{r_n(t)},X)^2}
         {\operatorname{Var}X}.}
 \label{eq:v18-partial-variance}
\end{equation}
The first term on the right of \eqref{eq:v18-partial-variance} is exactly the
energy-only variance of \eqref{eq:v17-caloric-frequency-variance}; the
second is its positive helicity-explained debit.
\end{theorem}

\begin{proof}
On the $(k,\sigma)$ helical component,
\[
 h=\operatorname{curl}e=Xe,
 \qquad
 \widehat Z^{\rm cal}=F_re.
\]
Therefore
\[
 \frac{\|\widehat Z^{\rm cal}-ae-bh\|_2^2}{\|e\|_2^2}
 =\mathbb E|F_r-a-bX|^2.
\]
Theorem~\ref{thm:v18-EH-quotient} identifies the left-hand minimum with
$\|Z^{\rm EH}\|_2^2/\|e\|_2^2$.  Least squares with a possibly singular
Gram gives the Moore--Penrose formula in
\eqref{eq:v18-regression-residual}.  If $\operatorname{Var}X>0$, first
subtract the mean of $F_r$ and then regress on $X-\mathbb EX$; this gives
\eqref{eq:v18-partial-variance}.  The energy-only projection removes only
the constant regressor, hence its residual is $\operatorname{Var}F_r$.
\end{proof}

\begin{corollary}[Zero criterion and the signed-helical three-level gate]
\label{cor:v18-three-level-gate}
One has
\begin{equation}
 \boxed{
 Z_{n,t}^{\rm EH}=0
 \iff
 F_{r_n(t)}=a+bX
 \quad\mu_t\text{-almost everywhere}
 \text{ for some }a,b\in\mathbb R.}
 \label{eq:v18-affine-zero}
\end{equation}
In particular:
\begin{enumerate}[label=\textup{(\roman*)},leftmargin=3em]
\item if the occupied signed curl spectrum $\supp\mu_t\circ X^{-1}$ has at
      most two points, then $Z_{n,t}^{\rm EH}=0$;
\item a nonzero complete caloric contact requires at least three distinct
      occupied signed helical frequencies;
\item at zero age, $F_0=|X|$, so a nonzero instantaneous critical nonlinear
      work requires both helicity signs and at least three distinct signed
      curl levels;
\item at positive age, three or more levels of one helicity may have a
      nonzero residual because $x\mapsto xe^{-2\nu rx^2}$ is not affine on a
      general positive-frequency set.
\end{enumerate}
The level count here is not the spatial lattice rank of the occupied Fourier
subgroup.
\end{corollary}

\begin{proof}
The zero criterion is equality in a finite- or countable-dimensional least
squares problem.  Every function on a set of at most two real points is the
restriction of an affine function, proving item~\textup{(i)}.  Item
\textup{(ii)} follows from
\eqref{eq:v18-EH-contact}.  At zero age, $|X|=X$ on the positive half-line and
$|X|=-X$ on the negative half-line.  Thus one-helicity support is affine.  A
set of three distinct points meeting both half-lines cannot lie on one affine
branch of $|X|$, proving item~\textup{(iii)}.  The last statement is the
strict nonlinearity of the positive-age multiplier away from accidental
three-point collinearity.
\end{proof}

For an exact finite three-level packet, the residual has one scalar
coefficient.

\begin{theorem}[Exact three-level divided-difference action]
\label{thm:v18-three-level-formula}
Suppose the signed curl law is supported on three distinct values
$x_1,x_2,x_3$ with masses $p_i>0$, $\sum_i p_i=1$.  Put
\begin{equation}
 f_i:=|x_i|e^{-2\nu r x_i^2},
 \qquad
 \Delta_i:=\prod_{j\ne i}(x_i-x_j),
 \label{eq:v18-three-data}
\end{equation}
and define the second divided difference
\begin{equation}
 [x_1,x_2,x_3]F_r
 :=\sum_{i=1}^3\frac{f_i}{\Delta_i}.
 \label{eq:v18-divided-difference}
\end{equation}
Then
\begin{equation}
 \boxed{
 \frac{\|Z^{\rm EH}\|_2^2}{\|e\|_2^2}
 =\frac{|[x_1,x_2,x_3]F_r|^2}
        {\displaystyle\sum_{i=1}^3
          \frac1{p_i\Delta_i^2}}.}
 \label{eq:v18-three-level-action}
\end{equation}
The residual is nonzero exactly when the three points
$(x_i,f_i)$ are not collinear.  At $r=0$, it is nonzero exactly when the
three signed levels contain both signs.
\end{theorem}

\begin{proof}
On the weighted space $\mathbb R^3$ with inner product
$\langle a,b\rangle_p=\sum_i p_i a_ib_i$, the orthogonal complement of
$\operatorname{span}\{(1,1,1),(x_1,x_2,x_3)\}$ is generated by
\[
 q_i=\frac1{p_i\Delta_i}.
\]
The Lagrange identities
$\sum_i\Delta_i^{-1}=0$ and
$\sum_i x_i\Delta_i^{-1}=0$ show the required orthogonality.  Moreover,
\[
 \langle f,q\rangle_p
 =\sum_i\frac{f_i}{\Delta_i}
 =[x_1,x_2,x_3]F_r,
 \qquad
 \|q\|_p^2=\sum_i\frac1{p_i\Delta_i^2}.
\]
Projection onto the one-dimensional complement gives
\eqref{eq:v18-three-level-action}.  A second divided difference vanishes
exactly on affine data.  For $r=0$, strict convexity of $|x|$ across its kink
at zero gives the final criterion; on one half-line $|x|$ is affine.
\end{proof}

\begin{countertheorem}[Positive energy-only variance with zero invariant writer]
\label{cth:v18-homochiral-two-level}
There is a smooth real mean-zero divergence-free field with nonzero projected
quadratic source for which the V17 energy-only caloric variance is positive
at generic ages, but $Z^{\rm EH}=0$ at every age.

On $\Torus^3$ set
\begin{equation}
 U_1(x)=\bigl(0,\cos x_1,-\sin x_1\bigr),
 \qquad
 U_2(x)=\bigl(\cos2x_2,0,\sin2x_2\bigr),
 \qquad
 u=U_1+U_2.
 \label{eq:v18-explicit-homochiral}
\end{equation}
Then
\begin{equation}
 \operatorname{curl}U_1=U_1,
 \qquad
 \operatorname{curl}U_2=2U_2,
 \label{eq:v18-explicit-curl}
\end{equation}
and $\BNS(u)\ne0$.  Nevertheless, for every $r\ge0$,
\begin{equation}
 \boxed{Z_r^{\rm EH}=0.}
 \label{eq:v18-explicit-EH-zero}
\end{equation}
The energy-only writer has positive norm whenever
\begin{equation}
 e^{-2\nu r}\ne2e^{-8\nu r};
 \label{eq:v18-explicit-energy-positive}
\end{equation}
in particular it is positive at $r=0$.
\end{countertheorem}

\begin{proof}
The displayed curl identities follow by direct differentiation.  The raw
convective term is
\[
 (u\cdot\nabla)u
 =\bigl(-2\cos x_1\sin2x_2,
        -\sin x_1\cos2x_2,
         \cos x_1\cos2x_2\bigr).
\]
Its curl is
\[
 \bigl(-2\cos x_1\sin2x_2,
        \sin x_1\cos2x_2,
        3\cos x_1\cos2x_2\bigr),
\]
which is nonzero.  Hence the convective term is not a gradient and its Leray
projection is nonzero.

The two helical components have signed curl values $1$ and $2$.  For every
$r$, the multiplier $F_r$ has only two occupied values.  Corollary
\ref{cor:v18-three-level-gate} therefore gives
$Z_r^{\rm EH}=0$.  Both modes have nonzero critical energy.  The energy-only
variance of two values is positive exactly when those values differ, namely
when \eqref{eq:v18-explicit-energy-positive} holds.
\end{proof}

\begin{assessment}[Why the explicit field matters]
\label{ass:v18-explicit-scope}
The field in Countertheorem~\ref{cth:v18-homochiral-two-level} is not offered
as a terminal blow-up packet.  It is an exact smooth Navier--Stokes datum
with nonzero quadratic source.  It proves that positivity of the energy-only
frequency variance does not certify a physical critical writer after the
second universal invariant has been shorted.  The omission cannot be repaired
by excluding only single shells, rank-at-most-two spatial subgroups, or zero
nonlinearity.
\end{assessment}

% =============================================================================
\section{Insertion into the V17 terminal card}
\label{sec:v18-terminal-card}
% =============================================================================

Return to the selected interior interval $\supp\eta_n$.  Define
$h_\perp(t)$ and $\gamma_n^{\rm hel}(t)$ by
\eqref{eq:v18-helicity-orthogonal}--\eqref{eq:v18-helicity-coefficient}, and
put
\begin{equation}
 \mathcal E_n^{\rm EH}
 :=\int_{\supp\eta_n}\|Z_{n,t}^{\rm EH}\|_2^2\,\dd t,
 \label{eq:v18-EH-history-action}
\end{equation}
\begin{equation}
 \mathcal D_n^{\rm hel}
 :=\int_{\supp\eta_n}
   \frac{|\langle Z_{n,t}^{\rm dip},h_\perp(t)\rangle|^2}
        {\|h_\perp(t)\|_2^2}\,\dd t,
 \label{eq:v18-helicity-history-debit}
\end{equation}
with the integrand set to zero on $\{h_\perp=0\}$.

\begin{lemma}[History-level invariant action identity]
\label{lem:v18-history-action}
One has
\begin{equation}
 \boxed{
 \mathcal E_n^{\rm dip}
 =\mathcal E_n^{\rm EH}+\mathcal D_n^{\rm hel},
 \qquad
 0\le\mathcal E_n^{\rm EH}\le\mathcal E_n^{\rm dip}.}
 \label{eq:v18-history-action-split}
\end{equation}
Consequently,
\begin{equation}
 \boxed{
 \frac\nu{\mathcal R_n^2}\mathcal E_n^{\rm EH}
 \le e_{\rm nat}^*
 =\frac1e\log\frac{\tau_+}{\tau_-}.}
 \label{eq:v18-EH-action-window}
\end{equation}
\end{lemma}

\begin{proof}
Integrate \eqref{eq:v18-helicity-debit} and use the V17 bound
\eqref{eq:v17-dipole-window}.
\end{proof}

Keep the V17 source-first vector $G_{n,t}^{\rm nat}$ and selected writer
$W_{n,t}^{\rm nat}$.  Define
\begin{equation}
 N_n^{\rm EH}
 :=\int_{\supp\eta_n}
   \langle Z_{n,t}^{\rm EH},W_{n,t}^{\rm nat}\rangle\,\dd t,
 \qquad
 \mathfrak m_n^{\rm EH}:=\frac\nu{\mathcal R_n}N_n^{\rm EH},
 \label{eq:v18-EH-mixed}
\end{equation}
and
\begin{equation}
 \mathfrak e_n^{\rm EH}
 :=\frac\nu{\mathcal R_n^2}\mathcal E_n^{\rm EH}.
 \label{eq:v18-EH-e}
\end{equation}
The denominator remains the native selected-writer action
$\mathfrak d_n^{\rm nat}=\nu D_n^{\rm nat}$.

\begin{theorem}[Compact energy--helicity terminal Gram]
\label{thm:v18-EH-Gram}
Put
\begin{equation}
 \mathfrak c_n^{\rm EH}
 :=-\frac1{\mathcal R_n^2}
   \int_{\supp\eta_n}
   \langle G_{n,t}^{\rm nat},Z_{n,t}^{\rm EH}\rangle\,\dd t.
 \label{eq:v18-EH-source-contact}
\end{equation}
Then
\begin{equation}
 \boxed{
 \mathbb G_n^{\rm EH}
 :=\begin{pmatrix}
  \mathfrak a_n^{\rm nat}&\sigma_*A_n&-\mathfrak c_n^{\rm EH}\\
  \sigma_*A_n&\mathfrak d_n^{\rm nat}&\mathfrak m_n^{\rm EH}\\
  -\mathfrak c_n^{\rm EH}&\mathfrak m_n^{\rm EH}&\mathfrak e_n^{\rm EH}
 \end{pmatrix}\succeq0.}
 \label{eq:v18-EH-three-Gram}
\end{equation}
Its entries lie in the same compact region as the V17 Gram, with
$\mathfrak e_n^{\rm EH}\le\mathfrak e_n^{\rm nat}$ and
\begin{equation}
 \boxed{
 |\mathfrak m_n^{\rm EH}|
 \le\sqrt{d_{\rm nat}^*e_{\rm nat}^*}.}
 \label{eq:v18-EH-mixed-bound}
\end{equation}
The mixed row is acquired by the two positive actions
\begin{equation}
 \mathcal E_{n,\rm EH}^{\pm}
 :=\int_{\supp\eta_n}
  \|Z_{n,t}^{\rm EH}
    \pm\mathcal R_nW_{n,t}^{\rm nat}\|_2^2\,\dd t,
 \label{eq:v18-EH-positive-actions}
\end{equation}
through
\begin{equation}
 \boxed{
 \mathfrak m_n^{\rm EH}
 =\frac\nu{4\mathcal R_n^2}
  \left(\mathcal E_{n,\rm EH}^{+}
       -\mathcal E_{n,\rm EH}^{-}\right).}
 \label{eq:v18-EH-positive-read}
\end{equation}
\end{theorem}

\begin{proof}
The matrix is the Gram of
\[
 \frac{G_n^{\rm nat}}{\sqrt\nu\mathcal R_n},
 \qquad
 \sqrt\nu W_n^{\rm nat},
 \qquad
 \frac{\sqrt\nu}{\mathcal R_n}Z_n^{\rm EH}
\]
in the common history space.  Positivity and all inherited bounds follow as
in Theorem~\ref{thm:v17-native-three-Gram}, using
Lemma~\ref{lem:v18-history-action}.  Cauchy--Schwarz gives
\eqref{eq:v18-EH-mixed-bound}.  Expanding the positive actions and
subtracting gives \eqref{eq:v18-EH-positive-read}.
\end{proof}

Define the invariant mixed coefficient and selected contact by
\begin{equation}
 \kappa_n^{\rm EH}:=\frac{N_n^{\rm EH}}{D_n^{\rm nat}},
 \qquad
 \boxed{
 \mathcal C_n^{\rm EH,sel}
 :=\kappa_n^{\rm EH}\mathfrak s_n
 =-\sigma_*\mathcal R_n^2A_n
   \frac{\mathfrak m_n^{\rm EH}}{\mathfrak d_n^{\rm nat}}.}
 \label{eq:v18-EH-selected-contact}
\end{equation}
The source-first complete-contact ledger is
\begin{equation}
 \boxed{
 \mathcal C_n^{\rm cal}
 =\mathcal C_n^{\rm EH,comp}
  +\mathcal C_n^{\rm EH,sel}
  +\mathcal U_n^{\rm EH,wr},}
 \label{eq:v18-EH-source-ledger}
\end{equation}
where
\begin{equation}
 \mathcal C_n^{\rm EH,comp}
 :=-\int_{\beta_n}^{b_n}
   \langle c(t)-G_{n,t}^{\rm nat},Z_{n,t}^{\rm EH}\rangle\,\dd t,
 \label{eq:v18-EH-complement}
\end{equation}
and
\begin{equation}
 \mathcal U_n^{\rm EH,wr}
 :=-\int_{\supp\eta_n}
 \left\langle G_{n,t}^{\rm nat},
 Z_{n,t}^{\rm EH}-\kappa_n^{\rm EH}W_{n,t}^{\rm nat}
 \right\rangle\,\dd t.
 \label{eq:v18-EH-writer-residual}
\end{equation}
The proof is the same orthogonal projection and source-complement identity as
in Proposition~\ref{prop:v17-native-selected-contact}, now using
\eqref{eq:v18-EH-contact}.

% =============================================================================
\section{The signed helicity-transfer row and the complete incidence chain}
\label{sec:v18-helicity-transfer}
% =============================================================================

Define
\begin{equation}
 N_n^{\rm hel}
 :=\int_{\supp\eta_n}
   \gamma_n^{\rm hel}(t)
   \langle h_\perp(t),W_{n,t}^{\rm nat}\rangle\,\dd t.
 \label{eq:v18-helicity-transfer-row}
\end{equation}

\begin{theorem}[Exact transfer from the energy-only to the invariant card]
\label{thm:v18-helicity-transfer}
The V17 and V18 mixed numerators satisfy
\begin{equation}
 \boxed{
 N_n^{\rm nat}=N_n^{\rm EH}+N_n^{\rm hel}.}
 \label{eq:v18-numerator-transfer}
\end{equation}
Equivalently,
\begin{equation}
 \boxed{
 \mathfrak m_n^{\rm nat}-\mathfrak m_n^{\rm EH}
 =\frac\nu{\mathcal R_n}N_n^{\rm hel}.}
 \label{eq:v18-normalized-transfer}
\end{equation}
Define
\begin{equation}
 S_n^{\rm EH}
 :=-\sigma_*A_n
   \frac{\mathfrak m_n^{\rm EH}}{\mathfrak d_n^{\rm nat}},
 \qquad
 \boxed{
 \Delta_n^{\rm hel}:=S_n^{\rm nat}-S_n^{\rm EH}.}
 \label{eq:v18-helicity-defect}
\end{equation}
Then
\begin{equation}
 \boxed{
 \Delta_n^{\rm hel}
 =-\sigma_*A_n
  \frac{\nu N_n^{\rm hel}}
       {\mathcal R_n\mathfrak d_n^{\rm nat}}.}
 \label{eq:v18-helicity-defect-row}
\end{equation}
No sign follows from the positive action debit
$\mathcal D_n^{\rm hel}$.
\end{theorem}

\begin{proof}
Insert \eqref{eq:v18-EH-from-E} into the definition of
$N_n^{\rm EH}$.  The difference from $N_n^{\rm nat}$ is precisely
\eqref{eq:v18-helicity-transfer-row}.  Normalize and substitute into the
definition of the two selected readouts.
\end{proof}

Define the invariant source-to-head residual
\begin{equation}
 I_n^{\rm EH}
 :=\mathfrak c_n^{\rm cal}-\mathfrak c_n^{\rm rep}-S_n^{\rm EH}
 =\frac{\mathcal C_n^{\rm ho}-\mathcal C_n^{\rm EH,sel}}
        {\mathcal R_n^2}.
 \label{eq:v18-EH-incidence}
\end{equation}

\begin{corollary}[Complete smoothed-to-invariant incidence chain]
\label{cor:v18-complete-incidence}
The three source-to-head readouts satisfy
\begin{equation}
 \boxed{
 I_n^{\rm EH}
 =I_n^{\rm nat}+\Delta_n^{\rm hel}
 =I_n^{\rm sm}+\Delta_n^{\rm met}+\Delta_n^{\rm hel}.}
 \label{eq:v18-complete-incidence-chain}
\end{equation}
Thus a V16 smoothed clean card is a clean universal-invariant card only after
both the V17 heat-metric defect and the V18 helicity-incidence defect vanish.
\end{corollary}

\begin{proof}
Subtract the definitions of $I_n^{\rm EH}$ and $I_n^{\rm nat}$ and use
\eqref{eq:v18-helicity-defect}.  The second equality is
\eqref{eq:v17-incidence-transfer}.
\end{proof}

\begin{countertheorem}[Positive action debit does not determine the signed helicity transfer]
\label{cth:v18-debit-no-sign}
There are finite three-vector Gram cards with fixed positive helicity debit
and fixed norms of the selected writer for which $N^{\rm hel}$ is positive,
zero, or negative.
\end{countertheorem}

\begin{proof}
Take an orthonormal pair $(e_1,e_2)$, let
$Z^{\rm dip}=ae_1+be_2$ and let the helicity short remove $ae_1$.  The
positive debit is $a^2$.  For selected writers
$W_\theta=\cos\theta\,e_1+\sin\theta\,e_2$, the helicity-transfer row is
$a\cos\theta$, while $\|W_\theta\|=1$.  Varying $\theta$ gives all three
signs without changing the debit.  The example is a Gram-level logical
obstruction; it is not asserted to be a terminal Navier--Stokes history.
\end{proof}

\begin{firewall}[Helicity action and helicity incidence are different targets]
\label{fire:v18-action-not-incidence}
The positive quantity $\mathcal D_n^{\rm hel}$ says how much writer action is
removed by the universal helicity short.  The signed row $N_n^{\rm hel}$ says
how that removed direction meets the already selected terminal writer.  The
first does not determine the second.  Bounding the positive debit, the
helicity norm, or the V17 frequency variance separately cannot replace the
mixed same-history read in \eqref{eq:v18-helicity-transfer-row}.
\end{firewall}

% =============================================================================
\section{Corrected priority theorem and the helical donor export}
\label{sec:v18-branch}
% =============================================================================

\begin{theorem}[Invariant-first terminal-writer alternative]
\label{thm:v18-branch}
After passage to a priority subsequence, exactly one of the following occurs.
\begin{enumerate}[label=\textup{(V18.\arabic*)},leftmargin=6.1em]
\item \emph{Helicity-incidence survivor.}  There is $\delta>0$ such that
      \begin{equation}
       \boxed{
       |\Delta_n^{\rm hel}|
       \ge\delta\mathfrak c_n^{\rm cal}.}
       \label{eq:v18-helicity-survivor}
      \end{equation}
      The energy-only and energy--helicity allocations differ by a fixed
      record fraction.  The exact signed carrier is
      \eqref{eq:v18-helicity-transfer-row}.
\item \emph{Source-metric incidence survivor.}  The helicity defect tends to
      zero, but the V17 metric defect obeys
      \begin{equation}
       \boxed{
       |\Delta_n^{\rm met}|
       \ge\delta\mathfrak c_n^{\rm cal}}
       \label{eq:v18-metric-survivor}
      \end{equation}
      for some $\delta>0$.
\item \emph{Represented endpoint return.}  The first two defects vanish, but
      \begin{equation}
       \boxed{
       |\mathfrak c_n^{\rm rep}|
       \ge\delta\mathfrak c_n^{\rm cal}}
       \label{eq:v18-represented-survivor}
      \end{equation}
      on a subsequence.
\item \emph{Invariant aggregate incidence residual.}  The first three
      branches fail, but
      \begin{equation}
       \boxed{
       |I_n^{\rm EH}|
       \ge\delta\mathfrak c_n^{\rm cal}}
       \label{eq:v18-EH-incidence-survivor}
      \end{equation}
      for some $\delta>0$.  Only then is the signed aggregate
      $\mathcal C_n^{\rm EH,comp}+\mathcal U_n^{\rm EH,wr}$ opened.
\item \emph{Universal-invariant clean export.}  One has
      \begin{equation}
       \boxed{
       \Delta_n^{\rm hel}\to0,
       \quad
       \Delta_n^{\rm met}\to0,
       \quad
       \mathfrak c_n^{\rm rep}\to0,
       \quad
       I_n^{\rm EH}\to0.}
       \label{eq:v18-clean-hypotheses}
      \end{equation}
      Then
      \begin{equation}
       \boxed{
       S_n^{\rm EH}-\mathfrak c_n^{\rm cal}\longrightarrow0,}
       \label{eq:v18-clean-readout}
      \end{equation}
      and eventually
      \begin{equation}
       \boxed{
       -\sigma_*
       \frac{\mathfrak m_n^{\rm EH}}{\mathfrak d_n^{\rm nat}}
       \ge\frac1{\mathfrak V_*},
       \qquad
       \mathcal C_n^{\rm EH,sel}
       \ge\frac34\mathcal R_n^2.}
       \label{eq:v18-clean-lower}
      \end{equation}
      In particular,
      \begin{equation}
       \boxed{
       \sum_n\mathcal C_n^{\rm EH,sel}=+\infty.}
       \label{eq:v18-clean-divergence}
      \end{equation}
\end{enumerate}
On the last branch, every nonzero pointwise contact contributing to the
clean packet has nonzero helical regression residual and therefore requires
at least three occupied signed helical levels.
\end{theorem}

\begin{proof}
Apply successive subsequence selection to the nonnegative normalized
magnitudes in the displayed priority order.  The first four alternatives are
therefore exhaustive.  On the final branch, the definition
\eqref{eq:v18-EH-incidence} gives
$S_n^{\rm EH}-\mathfrak c_n^{\rm cal}\to0$.  Since
$\mathfrak c_n^{\rm cal}\ge3/2$ and
$A_n\le\mathfrak V_*$, the first inequality in
\eqref{eq:v18-clean-lower} follows exactly as in the V17 clean branch.  The
contact identity \eqref{eq:v18-EH-selected-contact} then gives the second
inequality for all sufficiently large $n$.  The record amplitudes grow
geometrically, so the series diverges.  The pointwise helical-level statement
is Corollary~\ref{cor:v18-three-level-gate}.
\end{proof}

\begin{definition}[Universal-invariant V18 export card]
\label{def:v18-export-card}
On branch \textup{(V18.5)} retain
\begin{equation}
 \boxed{
 \begin{aligned}
 \mathfrak C_n^{\rm V18}:=\bigl(&b_n,I_n^{\rm ann},\Gamma_n,
 \Theta_n,E_{n,\cdot},P_n^{\rm rep},\\
 &G_{n,\cdot}^{\rm nat},W_{n,\cdot}^{\rm nat},
 e(\cdot),h(\cdot),\mathsf P_{\cdot}^{\rm EH},
 Z_{n,\cdot}^{\rm EH},\\
 &\mathbb G_n^{\rm EH},\mathcal D_n^{\rm hel},N_n^{\rm hel},
 \Delta_n^{\rm hel},\Delta_n^{\rm met},\\
 &\mathcal C_n^{\rm cal},\mathcal C_n^{\rm rep},
 \mathfrak s_n,\kappa_n^{\rm EH},
 \mathcal C_n^{\rm EH,sel},I_n^{\rm EH}\bigr).
 \end{aligned}}
 \label{eq:v18-export-card}
\end{equation}
The card exports to NS--B the selected positive Reynolds source together with
its complete universal-invariant terminal writer.  In addition, every finite
helical restriction records the signed levels $x=\sigma|k|$, their actual
critical-energy weights, and the divided-difference coefficient
\eqref{eq:v18-divided-difference}.  This is the finite helical gate for the
rank-three donor audit.
\end{definition}

\begin{firewall}[Two notions of rank three]
\label{fire:v18-two-ranks}
The V18 gate counts distinct signed curl eigenvalues required by the
energy--helicity quotient.  The inherited NS--B donor programme uses spatial
lattice rank three, whose first motif consists of noncoplanar interacting
triads.  A field may pass one condition and fail the other.  V18 neither
renames a three-level homochiral packet as a spatial rank-three donor nor
uses spatial rank to bypass the helical invariant short.  A clean donor card
must retain both pieces of information on the same occurring history.
\end{firewall}

% =============================================================================
\section{Version V18 integrated theorem and proof-status ledger}
\label{sec:v18-integrated}
% =============================================================================

\begin{theorem}[Version V18 invariant-quotient and helical-rank theorem]
\label{thm:v18-integrated}
Preserving every theorem and stated limitation of Versions~V1--V17,
Version~V18 proves the following.
\begin{enumerate}[label=\textup{(T18.\arabic*)},leftmargin=5.7em]
\item The periodic quadratic source is simultaneously orthogonal to the
      critical energy vector and the critical helicity vector.
\item The possibly singular energy--helicity Gram has a canonical
      Moore--Penrose support projection, and the complete caloric contact is
      represented by the unique minimum writer modulo that invariant
      carrier.
\item The V17 writer action splits into the invariant writer action plus one
      exact positive helicity debit.  The corresponding selected mixed row is
      signed and independent of the debit.
\item The invariant writer action is the least-squares residual of
      $F_r=|X|e^{-2\nu rX^2}$ after regression on $1$ and the signed helical
      frequency $X$.
\item With nonzero signed-frequency variance, the exact residual is the V17
      energy variance minus the helicity-explained covariance square.
\item Every two-level signed-helical packet is caloric-flat in the invariant
      quotient.  For exactly three levels, the complete action is the
      explicit divided-difference formula
      \eqref{eq:v18-three-level-action}.
\item At zero age, nonzero critical nonlinear work requires both helicity
      signs and at least three signed levels.  An explicit two-frequency
      homochiral field has nonzero quadratic source and positive energy-only
      variance but zero invariant writer.
\item On the selected terminal history, source, selected writer, and
      invariant caloric writer form the compact positive Gram
      \eqref{eq:v18-EH-three-Gram}, whose mixed row is acquired by two
      positive actions.
\item The complete source-to-head verdict obeys the exact chain
      \eqref{eq:v18-complete-incidence-chain}.  The helicity incidence is
      tested before the V17 metric incidence.
\item The clean terminal branch exports the non-summable contact
      \eqref{eq:v18-clean-divergence} together with a distinct signed-helical
      three-level gate and the inherited spatial rank-three donor gate.
\end{enumerate}
No populated terminal helical card, finite source-stock upper bound, Dini
improvement, backward-uniqueness theorem, Liouville theorem, or global
three-dimensional regularity conclusion is supplied.
\end{theorem}

\begin{proof}
Items~\textup{(T18.1)--(T18.2)} are
\Cref{thm:v18-energy-helicity-null,thm:v18-EH-quotient}.  Item
\textup{(T18.3)} is
\Cref{thm:v18-helicity-short,thm:v18-helicity-transfer,cth:v18-debit-no-sign}.
Items~\textup{(T18.4)--(T18.6)} are
\Cref{thm:v18-helical-regression,cor:v18-three-level-gate,thm:v18-three-level-formula}.
Item~\textup{(T18.7)} is
\Cref{cor:v18-three-level-gate,cth:v18-homochiral-two-level}.  Item
\textup{(T18.8)} is \Cref{lem:v18-history-action,thm:v18-EH-Gram}.  The last
two items are
\Cref{cor:v18-complete-incidence,thm:v18-branch,def:v18-export-card}.
\end{proof}

\begin{firewall}[Exact stopping boundary after V18]
\label{fire:v18-stop}
V18 closes a structural omission upstream of the V17 metric bridge.  A later
positive Reynolds estimate, represented-head projection, compactness theorem,
or finite-stock argument cannot manufacture the missing helicity-incidence
row.  The first admissible continuation is therefore a populated evaluation
of $N_n^{\rm hel}$, a theorem controlling that exact signed same-history row,
or a clean export to NS--B.  Re-estimating the energy-only variance or
replacing the signed row by the positive debit would return to a card now
proved nonminimal.
\end{firewall}

\begingroup\small
\begin{longtable}{p{0.25\textwidth}p{0.19\textwidth}p{0.48\textwidth}}
\caption{NS--A Version V18 proof-status ledger.}\label{tab:v18-ledger}\\
\toprule Row & Status & Exact conclusion\\\midrule
\endfirsthead
\toprule Row & Status & Exact conclusion\\\midrule
\endhead
V1--V17 prefix & \PROVED{} preserved
 &The complete V17 preterminal source is retained byte-for-byte; no earlier
 theorem or limitation is rewritten.\\[0.3em]
Nonlinear energy cancellation & \INHERITED{} / re-used
 &$\langle\BNS(u),u\rangle=0$.\\[0.3em]
Nonlinear helicity cancellation & \PROVED{} exact
 &$\langle\BNS(u),\operatorname{curl}u\rangle=0$, hence the critical source
 has a second universal null direction.\\[0.3em]
Energy--helicity support projection & \PROVED{} exact
 &The Moore--Penrose formula \eqref{eq:v18-null-projection} handles both rank
 one and rank two without choosing a helicity basis.\\[0.3em]
Invariant caloric writer & \PROVED{} exact
 &$Z^{\rm EH}$ is the minimum writer modulo the universal quadratic-invariant
 carrier and represents the unchanged complete contact.\\[0.3em]
Helicity action debit & \PROVED{} positive
 &The V17 energy-only action equals the invariant action plus
 \eqref{eq:v18-helicity-history-debit}.\\[0.3em]
Helicity mixed incidence & \OPEN{} signed physical row
 &$N_n^{\rm hel}$ has no sign or upper cancellation estimate from the
 separate positive actions.\\[0.3em]
Helical regression formula & \PROVED{} exact
 &The invariant action is the residual after regression of the caloric
 multiplier on $1$ and signed curl frequency.\\[0.3em]
Three-level necessity & \PROVED{} pointwise
 &A nonzero caloric contact requires at least three occupied signed helical
 levels; at zero age both signs are necessary.\\[0.3em]
Three-level coefficient & \PROVED{} exact
 &The residual is the squared second divided difference divided by the
 explicit positive mass/separation denominator.\\[0.3em]
Energy-only variance sufficiency & \OBSTRUCTION{} explicit
 &A smooth nontrivial two-frequency homochiral datum has positive energy-only
 variance and zero energy--helicity writer.\\[0.3em]
Invariant terminal Gram & \PROVED{} positive compact
 &The selected source, writer, and $Z^{\rm EH}$ form
 \eqref{eq:v18-EH-three-Gram}; positive polarization acquires the mixed row.\\[0.3em]
Complete incidence chain & \PROVED{} exact
 &$I^{\rm EH}=I^{\rm sm}+\Delta^{\rm met}+\Delta^{\rm hel}$.\\[0.3em]
Populated helical terminal card & \OPEN{} finite acquisition
 &The attached hypothetical chronology contains no literal values of the
 signed helical law or $N_n^{\rm hel}$.\\[0.3em]
Clean NS--A export & \CONDITIONAL{} refined
 &Export requires helicity, metric, represented, and invariant aggregate
 residuals all to vanish in the stated order.\\[0.3em]
Three-dimensional regularity & \OPEN
 &No surviving helicity, metric, represented, aggregate, or replenishment
 branch is excluded for every datum.\\
\bottomrule
\end{longtable}
\endgroup

% =============================================================================
\section{Exact mandate after Version V18}
\label{sec:v18-next}
% =============================================================================

\begin{programme}[Evaluate the invariant row or export the corrected donor]
\label{prog:v18-next}
Preserve Versions~V1--V18 verbatim.  There is no automatic general-purpose
Version~V19.

A continuation is licensed only in the following order.
\begin{enumerate}[label=\textup{(V19.\arabic*)},leftmargin=5.9em]
\item Evaluate the signed helicity-transfer row
      \eqref{eq:v18-helicity-transfer-row} on the literal selected history.
      A nonzero normalized value is the first upstream survivor.  Do not
      replace it by the positive action debit.
\item On a vanishing-helicity branch, continue with the already existing V17
      metric, represented, and invariant aggregate residuals in the order of
      Theorem~\ref{thm:v18-branch}.
\item If a pointwise contact is retained, acquire the actual signed-helical
      frequency law.  A finite three-level packet is evaluated by
      \eqref{eq:v18-three-level-action}; diffuse helical support remains a
      typed spectral-spreading branch rather than a shell count.
\item Send the helical three-level packet to NS--B only together with the
      existing spatial rank-three, source-ancestry, Reynolds-covariance, and
      paid-stress incidences.  Neither rank notion substitutes for the other.
\item On the universal-invariant clean branch, terminate NS--A and continue
      only in NS--B with \eqref{eq:v18-export-card} and an independently
      occurring upper budget.
\item A new analytical continuation without populated data must control the
      actual signed row $N_n^{\rm hel}$ or the complete invariant residual
      $I_n^{\rm EH}$.  Another energy-only variance estimate, endpoint norm,
      source-action lower bound, path law, Dini modulus, or compactness
      compiler is inadmissible.
\end{enumerate}
\end{programme}

\begin{assessment}[Programme position after V18]
\label{ass:v18-position}
The native source quotient is now corrected through all universal quadratic
invariants used by the periodic Euler nonlinearity.  The caloric writer has a
coordinate-free Moore--Penrose short, an exact helical spectral formula, and
a literal three-level obstruction.  The next unknown is not another
concentration geometry.  It is the signed incidence of the removed helicity
direction with the already selected terminal writer, or the downstream
fixed-datum donor budget after that incidence passes.
\end{assessment}

\paragraph{Version V18 history.}
Preserved the complete V17 source.  Reopened only its claim of native
source-minimality after identifying the additional universal helicity null.
Proved simultaneous energy and helicity cancellation, constructed the
possibly singular energy--helicity support projection, and derived the exact
positive helicity debit.  Diagonalized the quotient in the helical Fourier
basis and obtained the regression, partial-variance, three-level necessity,
and exact divided-difference formulas.  Gave a smooth two-frequency
homochiral datum with nonzero quadratic source, positive energy-only
variance, and zero invariant writer.  Rebuilt only the third column of the
terminal Gram, derived the signed helicity-transfer row and the complete
smoothed-to-invariant incidence chain, and inserted the new row ahead of the
metric/represented/aggregate/clean priority.  Exported the resulting
signed-helical gate without identifying it with spatial lattice rank three.

\begin{assessment}[Append-only integrity]
\label{ass:v18-integrity}
The complete Version~V17 source preceding its sole final executable document
terminator is preserved verbatim.  Its preterminal SHA--256 digest is
\begin{center}\scriptsize\ttfamily
4bfaf8216f5bd51a10df334bb4c1e483f0650e613893138e40d776cfdc17b506
\end{center}
The present material begins at Section~\ref{sec:v18-scope}; the final command
below is again unique.
\end{assessment}

\addtocontents{toc}{\protect\endgroup}
% =============================================================================
% END VERSION V18 MATERIAL -- append later versions before final terminator
% =============================================================================

% APPEND ALL LATER VERSIONS IMMEDIATELY ABOVE THE SOLE FINAL LINE BELOW.

\clearpage
% =============================================================================
% VERSION V19: NEW MATERIAL.  The complete V1--V18 source above is preserved
% verbatim; only its sole active \end{document} line was removed before append.
% =============================================================================
\hypersetup{pdftitle={NS-A Terminal Source Concentration Programme: Cumulative V19},pdfauthor={A. Pelissier}}
\makeatletter
\addtocontents{toc}{\protect\begingroup\protect\makeatletter}
\addtocontents{toc}{\protect\def\protect\l@section{\protect\@dottedtocline{1}{0em}{4.7em}}}
\addtocontents{toc}{\protect\def\protect\l@subsection{\protect\@dottedtocline{2}{1.5em}{5.3em}}}
\addtocontents{toc}{\protect\makeatother}
\makeatother
\renewcommand{\thesection}{V19.\arabic{section}}
\renewcommand{\thesubsection}{V19.\arabic{section}.\arabic{subsection}}
\renewcommand{\theequation}{V19.\arabic{equation}}
\setcounter{section}{0}
\setcounter{equation}{0}
\renewcommand{\theHsection}{V19.\arabic{section}}
\renewcommand{\theHsubsection}{V19.\arabic{section}.\arabic{subsection}}
\renewcommand{\theHtheorem}{V19.\arabic{section}.\arabic{theorem}}
\renewcommand{\theHequation}{V19.\arabic{equation}}

\begin{center}
{\LARGE\bfseries NS--A: Version V19}\\[0.5em]
{\large The Actual Helical-Triad Incidence Matrix, Support-Specific Invariant\\
Quotient, and the Affine-Rigidity Export Stop}\\[0.5em]
{\normalsize 3 September 2026}
\end{center}
\addcontentsline{toc}{part}{Version V19: actual triad incidence and the affine-rigidity stop}

\begin{abstract}
Version~V18 removed the two universal diagonal quadratic nulls of the
three-dimensional Euler nonlinearity: energy and helicity.  That quotient is
canonical for the complete equation, but a selected finite donor bank may
possess additional \emph{support-specific} diagonal quadratic invariants when
its actually active helical triads do not connect the bank strongly enough.
Such invariants are not the accidental orthogonal complement of one observed
source vector.  They are fixed before the amplitudes are inspected and are
computed exactly from the same triadic donor table that the fixed-datum
programme is required to populate.

Version~V19 constructs the corresponding finite incidence matrix.  Its row
for an active helical triad with signed curl frequencies
$x_a,x_b,x_c$ is
\[
 (x_b-x_c,\ x_c-x_a,\ x_a-x_b).
\]
A helical-diagonal quadratic form is invariant under the complete finite
Galerkin nonlinearity if and only if its coefficient vector lies in the
kernel of this matrix.  The vectors $1$ and $x$ always lie in the kernel and
give energy and helicity.  The exact extra-invariant count is the rank of the
orthogonal projector
\[
 (I-\mathfrak A^\dagger\mathfrak A)
 -B(B^*B)^\dagger B^*,\qquad B=(\mathbf1\ \mathbf x).
\]
A two-overlap propagation criterion proves that the kernel is exactly
$\operatorname{span}\{1,x\}$ on an affinely connected active-triad bank.
This gives a finite termination test: after it passes there is no further
helical-diagonal invariant to discover.

For a nonrigid bank, V19 constructs the Moore--Penrose projection onto all
bank-invariant writers $M_qe$.  It proves the exact minimum-writer
Pythagoras and identifies the signed mixed transfer omitted by the
energy--helicity card.  For a full Navier--Stokes field the finite-bank null
is modified by one explicit outer-parent leakage
\[
 \ell_{\mathscr A}
 =A^{-1/4}P_{\mathscr A}
  \bigl(\mathcal B(u)-\mathcal B(P_{\mathscr A}u)\bigr).
\]
The complete contact equals the fully shorted finite-bank contact plus the
pairing of this leakage with the removed invariant writer.  Thus no finite
support invariant is silently promoted to a full-PDE invariant.

An explicit six-atom periodic bank has spatial lattice rank three and two
nonzero helical triads, but the triads are disconnected.  Its invariant
kernel has dimension four, while the global energy--helicity space has
dimension two.  A writer can therefore have positive energy--helicity
residual and zero physical Galerkin contact.  Spatial rank three, signed
helical three-level occupancy, and affine triad connectivity are distinct
requirements.

The resulting terminal gate is exact.  A candidate V18 export returns one
of: finite-head capture failure, outer-parent triad leakage, an additional
support-invariant mixed row, a nonrigid kernel whose selected mixed row is
explicitly null, all-shell triadic escape, or an affine-rigid closed donor
head on which the V18 energy--helicity quotient is genuinely maximal.  Only the last branch is exported to NS--B.  No Dini improvement,
finite-stock upper bound, backward-uniqueness theorem, Liouville theorem, or
global regularity conclusion is claimed.
\end{abstract}

% =============================================================================
\section{Why the actual triad bank is upstream of another invariant guess}
\label{sec:v19-scope}
% =============================================================================

Version~V18 proved that the complete periodic quadratic source is
orthogonal to the energy and helicity writers and constructed the resulting
Moore--Penrose quotient.  Its stopping statement correctly forbade replacing
the signed helicity-transfer row by a positive action debit.  There remains,
however, one earlier question whenever the selected terminal response is
handed to a finite Fourier donor card:
\[
 \boxed{
 \text{does the actually active donor bank have only the two universal
 energy--helicity invariants?}}
 \tag{V19.G0}
\]
The answer is not determined by spatial lattice rank or by the number of
occupied signed helical levels.  A finite Galerkin bank can split into
triadically disconnected components, or can be joined only through
one-vertex overlaps.  Each such failure can create additional diagonal
quadratic invariants of the finite internal dynamics.

The correct order is therefore
\begin{equation}
 \boxed{
 \begin{gathered}
 \text{actual helical donor table}
 \longrightarrow
 \text{active-triad incidence matrix}
 \longrightarrow
 \text{finite invariant kernel}\\
 \longrightarrow
 \text{outer-parent closure test}
 \longrightarrow
 \text{energy--helicity card or an additional invariant survivor}.
 \end{gathered}}
 \label{eq:v19-upstream-order}
\end{equation}
This is not a new ambient source norm.  It is a finite algebraic read of the
literal donor table already required at the NS--A/NS--B interface.

\begin{firewall}[Universal invariants versus support-specific invariants]
\label{fire:v19-universal-support}
The energy and helicity directions remain the universal diagonal quadratic
nulls used by the complete periodic equation.  V19 does not retract
Theorem~\ref{thm:v18-energy-helicity-null}.  It proves that a restricted
finite interaction bank may have additional exact invariants because some
triadic channels are absent.  Those extra invariants are valid for the
internal Galerkin source only.  Their promotion to the full Navier--Stokes
source requires the outer-parent leakage in
\Cref{sec:v19-tail-leakage} to vanish or be retained as a signed residual.
\end{firewall}

Version~V19 establishes the following chain.
\begin{enumerate}[label=\textup{(U19.\arabic*)},leftmargin=5.2em]
\item Every active helical triad contributes one exact linear row to the
      quadratic-invariant incidence matrix.
\item The kernel of the complete row bank is exactly the space of all
      helical-diagonal quadratic invariants of the finite Galerkin system.
\item Energy and helicity form a canonical two-dimensional subspace; the
      Moore--Penrose difference projector counts every additional invariant.
\item A finite two-overlap criterion implies exact affine rigidity and hence
      proves that energy and helicity exhaust the bank.
\item The actual state turns the coefficient kernel into a positive writer
      Gram and a source-minimal invariant quotient.
\item The full-PDE failure of a finite-bank null is exactly the outer-parent
      leakage, with no unnamed remainder.
\item A spatial-rank-three disconnected bank exhibits two additional
      invariants and proves that the V18 quotient is not automatically
      support minimal.
\item The terminal handoff is reduced to a finite PASS, typed obstruction,
      or all-shell escape decision.
\end{enumerate}

% =============================================================================
\section{Reality-paired helical atoms and the exact triad row}
\label{sec:v19-triad-row}
% =============================================================================

Let $\mathscr A$ be a finite family of \emph{reality-paired helical atoms}.
An atom
\[
 a=[k_a,\sigma_a]
\]
denotes the pair of Fourier modes $(k_a,\sigma_a)$ and
$(-k_a,\sigma_a)$, where $k_a\in\mathbb Z^3\setminus\{0\}$ and
$\sigma_a\in\{+1,-1\}$.  Its signed curl frequency is
\begin{equation}
 \boxed{x_a:=\sigma_a|k_a|.}
 \label{eq:v19-signed-frequency}
\end{equation}
Choose normalized helical vectors $h_a$ satisfying
\begin{equation}
 \mathrm i k_a\times h_a=x_ah_a,
 \qquad k_a\cdot h_a=0,
 \qquad h_a(-k_a)=\overline{h_a(k_a)}.
 \label{eq:v19-helical-basis}
\end{equation}
A real field on the bank is written
\begin{equation}
 u_{\mathscr A}
 =\sum_{a\in\mathscr A}
   \left(z_ah_ae^{\mathrm i k_a\cdot x}
   +\overline{z_ah_a}e^{-\mathrm i k_a\cdot x}\right).
 \label{eq:v19-bank-field}
\end{equation}
Let $\mathcal E_a=|z_a|^2$ denote the modal energy with the harmless common
Parseval normalization suppressed.

\begin{definition}[Active reality-paired helical triad]
\label{def:v19-active-triad}
An unordered triple $\mathfrak t=\{a,b,c\}\subset\mathscr A$ is
\emph{active} when representatives may be oriented so that
\begin{equation}
 k_a+k_b+k_c=0,
 \label{eq:v19-wavevector-triad}
\end{equation}
and the helical geometric coefficient
\begin{equation}
 \Gamma_{\mathfrak t}
 :=(\overline h_b\times\overline h_c)\cdot\overline h_a
 \label{eq:v19-geometric-coefficient}
\end{equation}
is nonzero, and the signed-frequency difference row
\[
 (x_b-x_c,\ x_c-x_a,\ x_a-x_b)
\]
is nonzero.  The scalar triple product is cyclic, so the same coefficient,
up to the fixed harmless phase convention for the three helical vectors,
occurs in all three modal equations.  The negative wavevector copy is the
same reality-paired triad and is not counted again.
\end{definition}

\begin{lemma}[Exact modal transfer vector of one active triad]
\label{lem:v19-triad-transfer}
For an isolated active triad $\mathfrak t=\{a,b,c\}$, there is a real cubic
scalar $\Xi_{\mathfrak t}(z)$, not identically zero, such that the nonlinear
modal-energy transfer is
\begin{equation}
 \boxed{
 \begin{pmatrix}
  \dot{\mathcal E}_a\\
  \dot{\mathcal E}_b\\
  \dot{\mathcal E}_c
 \end{pmatrix}_{\!\mathfrak t}
 =\Xi_{\mathfrak t}(z)
 \begin{pmatrix}
  x_b-x_c\\
  x_c-x_a\\
  x_a-x_b
 \end{pmatrix}.}
 \label{eq:v19-triad-transfer}
\end{equation}
The phase of the three amplitudes may be chosen so that
$\Xi_{\mathfrak t}(z)\ne0$.
\end{lemma}

\begin{proof}
The vector identity
\[
 \mathcal B(u)=\mathbb P(u\cdot\nabla u)
              =-\mathbb P(u\times\operatorname{curl}u)
\]
puts the Euler nonlinearity into a form symmetric in the two parent modes.
For the equation of the $k_a$ mode, the two parents are the reality
partners $-k_b,-k_c$.  Symmetrization gives the factor
\[
 \frac12(x_b-x_c)(\overline h_b\times\overline h_c).
\]
Projection onto $h_a$ therefore gives, after a fixed phase choice,
\[
 \dot z_a
 =\gamma_{\mathfrak t}(x_b-x_c)
   \overline z_b\,\overline z_c,
\]
with one nonzero complex number $\gamma_{\mathfrak t}$ proportional to
$\Gamma_{\mathfrak t}$.  Cyclicity of the scalar triple product gives the
same $\gamma_{\mathfrak t}$ in the two cyclic equations.  Consequently
\[
 \dot{\mathcal E}_a
 =2\operatorname{Re}
  \left(\gamma_{\mathfrak t}
  \overline z_a\overline z_b\overline z_c\right)(x_b-x_c),
\]
and similarly for $b,c$.  This is
\eqref{eq:v19-triad-transfer} with the common real factor in parentheses as
$\Xi_{\mathfrak t}$.  Since
$\gamma_{\mathfrak t}\ne0$, choose nonzero amplitudes and their total phase
so that its real part is nonzero.
\end{proof}

For a real coefficient vector $q=(q_a)_{a\in\mathscr A}$, define the
diagonal quadratic form
\begin{equation}
 \mathcal Q_q(u_{\mathscr A})
 :=\sum_{a\in\mathscr A}q_a\mathcal E_a.
 \label{eq:v19-quadratic-form}
\end{equation}
For an active triad, define its incidence row
\begin{equation}
 \boxed{
 r_{\mathfrak t}
 :=(x_b-x_c)e_a^*
   +(x_c-x_a)e_b^*
   +(x_a-x_b)e_c^*.}
 \label{eq:v19-incidence-row}
\end{equation}
Changing the cyclic orientation multiplies this row by $-1$ and therefore
does not change its kernel.

\begin{definition}[Actual helical-triad incidence matrix]
\label{def:v19-incidence-matrix}
Let $\mathscr T_{\mathscr A}$ be the complete set of active triads in the
predeclared bank.  The matrix
\begin{equation}
 \boxed{
 \mathfrak A_{\mathscr A}:
 \mathbb R^{\mathscr A}\longrightarrow
 \mathbb R^{\mathscr T_{\mathscr A}},
 \qquad
 (\mathfrak A_{\mathscr A}q)_{\mathfrak t}
 =r_{\mathfrak t}q}
 \label{eq:v19-incidence-matrix}
\end{equation}
is the \emph{actual quadratic-invariant incidence matrix} of the finite
helical donor table.
\end{definition}

\begin{theorem}[Exact finite classification of diagonal quadratic invariants]
\label{thm:v19-invariant-kernel}
Consider the orthogonal Galerkin Euler equation on the reality-paired bank
$\mathscr A$.  Then
\begin{equation}
 \boxed{
 \mathcal Q_q\text{ is conserved by every Galerkin trajectory}
 \quad\Longleftrightarrow\quad
 \mathfrak A_{\mathscr A}q=0.}
 \label{eq:v19-invariant-iff}
\end{equation}
Equivalently, the space of all translation-invariant, curl-diagonal
quadratic invariants on this bank is
\begin{equation}
 \boxed{\mathcal I_{\mathscr A}=\Ker\mathfrak A_{\mathscr A}.}
 \label{eq:v19-invariant-space}
\end{equation}
\end{theorem}

\begin{proof}
The derivative of \eqref{eq:v19-quadratic-form} is the sum over active
triads of their modal transfer contributions.  By
Lemma~\ref{lem:v19-triad-transfer}, the contribution of
$\mathfrak t=\{a,b,c\}$ is
\[
 \Xi_{\mathfrak t}(z)
 \bigl[q_a(x_b-x_c)+q_b(x_c-x_a)+q_c(x_a-x_b)\bigr]
 =\Xi_{\mathfrak t}(z)r_{\mathfrak t}q.
\]
If every row annihilates $q$, each contribution vanishes separately, proving
sufficiency.

Conversely, fix one active triad and set all other modal amplitudes to zero
at one instant.  Every other triad contribution to the derivative then
vanishes because at least one of its three amplitudes is zero.  The phase
choice in Lemma~\ref{lem:v19-triad-transfer} makes
$\Xi_{\mathfrak t}\ne0$.  Conservation for every initial datum therefore
forces $r_{\mathfrak t}q=0$.  Repeating this for every active triad proves
necessity.
\end{proof}

\begin{corollary}[Energy and helicity are always in the finite kernel]
\label{cor:v19-EH-in-kernel}
Let
\begin{equation}
 \mathbf1=(1)_{a\in\mathscr A},
 \qquad
 \mathbf x=(x_a)_{a\in\mathscr A}.
 \label{eq:v19-EH-vectors}
\end{equation}
Then
\begin{equation}
 \boxed{
 \mathbf1,\mathbf x\in\Ker\mathfrak A_{\mathscr A}.}
 \label{eq:v19-EH-kernel}
\end{equation}
They are the coefficient vectors of energy and helicity.
\end{corollary}

\begin{proof}
For every row,
\[
 (x_b-x_c)+(x_c-x_a)+(x_a-x_b)=0.
\]
Also
\[
 x_a(x_b-x_c)+x_b(x_c-x_a)+x_c(x_a-x_b)=0.
\]
Theorem~\ref{thm:v19-invariant-kernel} gives the invariant interpretation.
\end{proof}

\begin{firewall}[Class of invariants classified]
\label{fire:v19-invariant-class}
Theorem~\ref{thm:v19-invariant-kernel} classifies quadratic forms whose
self-adjoint Fourier multiplier is diagonal in the helical atom basis.  This
is the target-native class used by the Fourier donor table and contains
energy, helicity, shell weights, component energies, and every
support-specific diagonal invariant below.  V19 does not claim that an
arbitrary nontranslation-invariant or mode-mixing quadratic form has been
classified.
\end{firewall}

% =============================================================================
\section{Affine rigidity and the exact extra-invariant projector}
\label{sec:v19-affine-rigidity}
% =============================================================================

Put
\begin{equation}
 B_{\mathscr A}:=(\mathbf1\ \mathbf x):\mathbb R^2\to
 \mathbb R^{\mathscr A}.
 \label{eq:v19-EH-synthesis-coeff}
\end{equation}
The coefficient-space projections onto the complete invariant kernel and the
energy--helicity subspace are
\begin{equation}
 \boxed{
 K_{\mathscr A}:=I-\mathfrak A_{\mathscr A}^{\dagger}
                         \mathfrak A_{\mathscr A},
 \qquad
 K_{\mathscr A}^{\rm EH}
 :=B_{\mathscr A}(B_{\mathscr A}^*B_{\mathscr A})^\dagger
   B_{\mathscr A}^*.}
 \label{eq:v19-coefficient-projections}
\end{equation}
By \eqref{eq:v19-EH-kernel},
$\Ran K_{\mathscr A}^{\rm EH}\subseteq\Ran K_{\mathscr A}$.
Hence
\begin{equation}
 \boxed{
 K_{\mathscr A}^{\rm ex}
 :=K_{\mathscr A}-K_{\mathscr A}^{\rm EH}}
 \label{eq:v19-extra-projector}
\end{equation}
is itself an orthogonal projection.

\begin{definition}[Extra invariant nullity and affine rigidity]
\label{def:v19-extra-nullity}
Define
\begin{equation}
 \boxed{
 d_{\mathscr A}^{\rm ex}
 :=\operatorname{tr}K_{\mathscr A}^{\rm ex}
 =\dim\Ker\mathfrak A_{\mathscr A}
  -\operatorname{rank} B_{\mathscr A}.}
 \label{eq:v19-extra-nullity}
\end{equation}
The active triad bank is \emph{affine rigid} when
$d_{\mathscr A}^{\rm ex}=0$, equivalently
\begin{equation}
 \boxed{
 \Ker\mathfrak A_{\mathscr A}
 =\operatorname{span}\{\mathbf1,\mathbf x\}.}
 \label{eq:v19-affine-rigid}
\end{equation}
When $\mathbf x$ is nonconstant, the right side has dimension two.
\end{definition}

The row condition has a simple geometric meaning:
\begin{equation}
 r_{\mathfrak t}q=0
 \quad\Longleftrightarrow\quad
 \det\begin{pmatrix}
  1&x_a&q_a\\
  1&x_b&q_b\\
  1&x_c&q_c
 \end{pmatrix}=0.
 \label{eq:v19-collinearity}
\end{equation}
Thus every active triad requires the three points $(x_a,q_a)$,
$(x_b,q_b)$, $(x_c,q_c)$ to be collinear.

\begin{theorem}[Two-overlap propagation implies affine rigidity]
\label{thm:v19-two-overlap-rigidity}
Suppose the atoms admit an ordering $a_1,\ldots,a_N$ with the following
properties.
\begin{enumerate}[label=\textup{(R\arabic*)},leftmargin=3.5em]
\item The first three atoms form an active triad whose signed frequencies
      are not all equal.
\item For every $m\ge4$, there is an active triad containing $a_m$ and two
      earlier atoms $a_i,a_j$ with $x_{a_i}\ne x_{a_j}$.
\end{enumerate}
Then the bank is affine rigid.
\end{theorem}

\begin{proof}
Let $q\in\Ker\mathfrak A_{\mathscr A}$.  On the seed triad,
\eqref{eq:v19-collinearity} and the presence of at least two distinct
$x$-values give unique numbers $\alpha,\beta$ such that
$q_{a_i}=\alpha+\beta x_{a_i}$ for $i=1,2,3$.  This also covers the case in
which two seed frequencies coincide: the row condition first forces the two
corresponding $q$-values to agree.

Assume inductively that $q_a=\alpha+\beta x_a$ for every earlier atom.
For the triad containing $a_m,a_i,a_j$, the two earlier signed frequencies
are distinct.  Equation \eqref{eq:v19-collinearity} says that
$(x_{a_m},q_{a_m})$ lies on the unique line through
$(x_{a_i},q_{a_i})$ and $(x_{a_j},q_{a_j})$, which is the line
$q=\alpha+\beta x$.  Thus the same formula holds for $a_m$.  Induction gives
$q=\alpha\mathbf1+\beta\mathbf x$ on the whole bank.  The reverse inclusion
is Corollary~\ref{cor:v19-EH-in-kernel}.
\end{proof}

\begin{corollary}[Finite terminating rigidity audit]
\label{cor:v19-finite-rigidity-audit}
For a populated finite donor table, affine rigidity is decided exactly by
any one of the equivalent finite tests
\begin{equation}
 \boxed{
 d_{\mathscr A}^{\rm ex}=0
 \quad\Longleftrightarrow\quad
 K_{\mathscr A}=K_{\mathscr A}^{\rm EH}
 \quad\Longleftrightarrow\quad
 \operatorname{rank}\mathfrak A_{\mathscr A}
 =|\mathscr A|-\operatorname{rank} B_{\mathscr A}.}
 \label{eq:v19-rigidity-tests}
\end{equation}
The two-overlap condition of
Theorem~\ref{thm:v19-two-overlap-rigidity} is a sufficient combinatorial
certificate, not a necessary one.
\end{corollary}

\begin{proof}
The Moore--Penrose operators in
\eqref{eq:v19-coefficient-projections} are the orthogonal projections onto
the two nested subspaces.  Their difference vanishes exactly when the
subspaces agree.  Rank--nullity gives the last equality.
\end{proof}

\begin{proposition}[Disconnected triad components create extra invariants]
\label{prop:v19-components}
Partition the active-triad hypergraph into vertex-disjoint connected
components $\mathscr A_1,\ldots,\mathscr A_m$.  Then every coefficient vector
which is affine in $x$ separately on each component belongs to
$\Ker\mathfrak A_{\mathscr A}$.  Consequently,
\begin{equation}
 \boxed{
 \dim\Ker\mathfrak A_{\mathscr A}
 \ge\sum_{j=1}^m\operatorname{rank}(\mathbf1_{\mathscr A_j},
                            \mathbf x_{\mathscr A_j}).}
 \label{eq:v19-component-nullity}
\end{equation}
If every component contains at least two distinct signed frequencies, the
right side is $2m$.  Hence every bank with at least two such components has
extra invariant nullity at least two.
\end{proposition}

\begin{proof}
Every incidence row is supported inside one component.  On that component,
a coefficient of the form $\alpha_j+\beta_jx$ annihilates the row by the two
identities in Corollary~\ref{cor:v19-EH-in-kernel}.  The componentwise affine
spaces have disjoint supports, so their dimensions add.
\end{proof}

\begin{assessment}[What affine rigidity closes]
\label{ass:v19-rigidity-meaning}
Once \eqref{eq:v19-rigidity-tests} passes, no further
translation-invariant curl-diagonal quadratic null can be discovered on the
same finite active bank.  Energy and helicity then exhaust the target-native
invariant class.  A later failure must come from source/writer incidence,
outer-parent leakage, represented return, paid-stress transport, or a
non-diagonal structure outside the declared donor target---not from another
unseen diagonal invariant.
\end{assessment}

% =============================================================================
\section{The source-minimal writer modulo the complete bank invariant kernel}
\label{sec:v19-bank-quotient}
% =============================================================================

Let $\mathcal H_{\mathscr A}$ be the real finite Galerkin carrier of the
bank.  For $q\in\mathbb R^{\mathscr A}$, let $M_q$ be the self-adjoint
helical-diagonal multiplier
\begin{equation}
 M_qe_{\mathscr A}
 =\sum_{a\in\mathscr A}q_ae_a,
 \label{eq:v19-diagonal-multiplier}
\end{equation}
where $e_a$ denotes the reality-paired critical-energy component of
$e_{\mathscr A}=A^{1/4}u_{\mathscr A}$.
Define the invariant-writer synthesis
\begin{equation}
 \mathsf N_{e,\mathscr A}
 :\mathbb R^{\mathscr A}\to\mathcal H_{\mathscr A},
 \qquad
 \mathsf N_{e,\mathscr A}q
 :=M_{K_{\mathscr A}q}e_{\mathscr A}.
 \label{eq:v19-invariant-synthesis}
\end{equation}
Its positive Gram and support projection are
\begin{equation}
 \mathsf G_{e,\mathscr A}^{\rm inv}
 :=\mathsf N_{e,\mathscr A}^*\mathsf N_{e,\mathscr A},
 \qquad
 \boxed{
 \mathsf P_{e,\mathscr A}^{\rm inv}
 :=\mathsf N_{e,\mathscr A}
   (\mathsf G_{e,\mathscr A}^{\rm inv})^\dagger
   \mathsf N_{e,\mathscr A}^*.}
 \label{eq:v19-state-invariant-projection}
\end{equation}
The Moore--Penrose formula automatically removes invariant coefficient
vectors whose writer vanishes on the currently occupied amplitudes.

Let
\begin{equation}
 c_{\mathscr A}^{\rm int}
 :=A^{-1/4}P_{\mathscr A}
   \mathcal B(u_{\mathscr A})
 \label{eq:v19-internal-source}
\end{equation}
be the internal Galerkin critical source.

\begin{theorem}[Complete finite-bank source-null projection]
\label{thm:v19-bank-source-null}
For every $q\in\Ker\mathfrak A_{\mathscr A}$,
\begin{equation}
 \boxed{
 \langle c_{\mathscr A}^{\rm int},M_qe_{\mathscr A}\rangle=0.}
 \label{eq:v19-bank-null-pairing}
\end{equation}
Hence
\begin{equation}
 \boxed{
 c_{\mathscr A}^{\rm int}
 \perp\Ran\mathsf P_{e,\mathscr A}^{\rm inv}.}
 \label{eq:v19-bank-source-orthogonal}
\end{equation}
For any writer $Y\in\mathcal H_{\mathscr A}$, define
\begin{equation}
 Z_{\mathscr A}^{\rm inv}
 :=(I-\mathsf P_{e,\mathscr A}^{\rm inv})Y.
 \label{eq:v19-bank-minimal-writer}
\end{equation}
Then
\begin{equation}
 \boxed{
 \langle c_{\mathscr A}^{\rm int},Y\rangle
 =\langle c_{\mathscr A}^{\rm int},Z_{\mathscr A}^{\rm inv}\rangle,}
 \label{eq:v19-bank-contact-unchanged}
\end{equation}
while for every invariant writer
$V\in\Ran\mathsf P_{e,\mathscr A}^{\rm inv}$,
\begin{equation}
 \boxed{
 \|Y-V\|_2^2
 =\|Z_{\mathscr A}^{\rm inv}\|_2^2
  +\|\mathsf P_{e,\mathscr A}^{\rm inv}Y-V\|_2^2.}
 \label{eq:v19-bank-Pythagoras}
\end{equation}
Thus $Z_{\mathscr A}^{\rm inv}$ is the unique minimum-norm representative
of the finite-bank writer modulo every diagonal quadratic invariant of the
actual donor table.
\end{theorem}

\begin{proof}
For the Galerkin Euler equation,
\[
 \frac{\dd}{\dd t}\mathcal Q_q(u_{\mathscr A})
 =-\langle\mathcal B(u_{\mathscr A}),M_qu_{\mathscr A}\rangle.
\]
If $q\in\Ker\mathfrak A_{\mathscr A}$, Theorem
\ref{thm:v19-invariant-kernel} makes the left side zero for every state.
Because $M_q$ commutes with $A$ and the Galerkin projection,
\[
 \langle\mathcal B(u_{\mathscr A}),M_qu_{\mathscr A}\rangle
 =\langle A^{-1/4}P_{\mathscr A}\mathcal B(u_{\mathscr A}),
          M_qA^{1/4}u_{\mathscr A}\rangle,
\]
which is \eqref{eq:v19-bank-null-pairing}.  Linear combinations give
\eqref{eq:v19-bank-source-orthogonal}.  The Moore--Penrose operator in
\eqref{eq:v19-state-invariant-projection} is the orthogonal projection onto
the invariant-writer range.  Contact invariance and orthogonal Pythagoras
then prove the remaining statements.
\end{proof}

Let $\mathsf P_{e,\mathscr A}^{\rm EH}$ be the finite-bank version of the
V18 projection generated by $e_{\mathscr A}$ and
$h_{\mathscr A}=\operatorname{curl}e_{\mathscr A}$.  Its range is contained
in the full invariant-writer range.  Therefore
\begin{equation}
 \boxed{
 \mathsf P_{e,\mathscr A}^{\rm ex}
 :=\mathsf P_{e,\mathscr A}^{\rm inv}
  -\mathsf P_{e,\mathscr A}^{\rm EH}}
 \label{eq:v19-physical-extra-projection}
\end{equation}
is the orthogonal projection onto the physically nonvanishing additional
invariant writers.

\begin{corollary}[Exact extra-invariant debit and signed transfer]
\label{cor:v19-extra-invariant-transfer}
For
\begin{equation}
 Z_{\mathscr A}^{\rm EH}
 :=(I-\mathsf P_{e,\mathscr A}^{\rm EH})Y,
 \qquad
 R_{\mathscr A}^{\rm ex}
 :=\mathsf P_{e,\mathscr A}^{\rm ex}Y,
 \label{eq:v19-EH-extra-writers}
\end{equation}
one has
\begin{equation}
 \boxed{
 Z_{\mathscr A}^{\rm EH}
 =Z_{\mathscr A}^{\rm inv}+R_{\mathscr A}^{\rm ex},
 \qquad
 \|Z_{\mathscr A}^{\rm EH}\|_2^2
 =\|Z_{\mathscr A}^{\rm inv}\|_2^2
  +\|R_{\mathscr A}^{\rm ex}\|_2^2.}
 \label{eq:v19-extra-action-debit}
\end{equation}
For any selected writer $W$,
\begin{equation}
 \boxed{
 \langle Z_{\mathscr A}^{\rm EH},W\rangle
 =\langle Z_{\mathscr A}^{\rm inv},W\rangle
  +\langle R_{\mathscr A}^{\rm ex},W\rangle.}
 \label{eq:v19-extra-mixed-row}
\end{equation}
The action debit is positive, but the additional mixed row is signed.  If
the bank is affine rigid, then
\begin{equation}
 \boxed{
 \mathsf P_{e,\mathscr A}^{\rm inv}
 =\mathsf P_{e,\mathscr A}^{\rm EH},
 \qquad
 R_{\mathscr A}^{\rm ex}=0.}
 \label{eq:v19-rigid-projections-equal}
\end{equation}
\end{corollary}

\begin{proof}
The two orthogonal projections are nested, so their difference is an
orthogonal projection and
\[
 I-\mathsf P^{\rm EH}
 =(I-\mathsf P^{\rm inv})
  +(\mathsf P^{\rm inv}-\mathsf P^{\rm EH})
\]
is an orthogonal splitting.  This proves
\eqref{eq:v19-extra-action-debit}; pairing with $W$ gives
\eqref{eq:v19-extra-mixed-row}.  On an affine-rigid bank the invariant
coefficient space is generated by $1,x$, so its state-dependent writer range
is exactly $\operatorname{span}\{e_{\mathscr A},h_{\mathscr A}\}$, including
any rank collapse.  The projections therefore agree.
\end{proof}

\begin{firewall}[Extra coefficient nullity versus an occurring extra writer]
\label{fire:v19-nullity-not-incidence}
A positive integer $d_{\mathscr A}^{\rm ex}$ proves that the finite donor
bank has additional invariant coefficient directions.  It does not prove
that those directions are populated by the present state or meet the
selected writer.  The state Gram
\eqref{eq:v19-state-invariant-projection} removes unoccupied directions, and
the signed read in \eqref{eq:v19-extra-mixed-row} is still required before an
extra invariant is reported as a terminal obstruction.
\end{firewall}

% =============================================================================
\section{Outer-parent leakage is the exact full-PDE correction}
\label{sec:v19-tail-leakage}
% =============================================================================

Return to a smooth full periodic field $u$ and let $P_{\mathscr A}$ be the
orthogonal Fourier--helical projection onto a finite reality-paired bank.
Put
\begin{equation}
 u_{\mathscr A}=P_{\mathscr A}u,
 \qquad
 e_{\mathscr A}=P_{\mathscr A}e,
 \label{eq:v19-projected-state}
\end{equation}
and define the exact outer-parent leakage
\begin{equation}
 \boxed{
 \ell_{\mathscr A}[u]
 :=A^{-1/4}P_{\mathscr A}
   \bigl(\mathcal B(u)-\mathcal B(u_{\mathscr A})\bigr).}
 \label{eq:v19-outer-parent-leakage}
\end{equation}
Then
\begin{equation}
 P_{\mathscr A}c
 =c_{\mathscr A}^{\rm int}+\ell_{\mathscr A}[u].
 \label{eq:v19-source-split}
\end{equation}
The leakage contains every interaction with at least one parent outside the
finite bank, including high--high return into the bank.  It is one assembled
source term and is not split into a number of shell occurrences.

\begin{theorem}[Exact finite-bank invariant short inside the full equation]
\label{thm:v19-full-leakage-identity}
For every $q\in\Ker\mathfrak A_{\mathscr A}$,
\begin{equation}
 \boxed{
 \langle c,M_qe_{\mathscr A}\rangle
 =\langle\ell_{\mathscr A}[u],M_qe_{\mathscr A}\rangle.}
 \label{eq:v19-invariant-leakage-pairing}
\end{equation}
For every bank-supported writer $Y\in\mathcal H_{\mathscr A}$,
\begin{equation}
 \boxed{
 \langle c,Y\rangle
 =\langle c,(I-\mathsf P_{e,\mathscr A}^{\rm inv})Y\rangle
  +\langle\ell_{\mathscr A}[u],
           \mathsf P_{e,\mathscr A}^{\rm inv}Y\rangle.}
 \label{eq:v19-full-contact-short}
\end{equation}
In particular,
\begin{equation}
 \boxed{
 \left|
 \langle\ell_{\mathscr A}[u],
        \mathsf P_{e,\mathscr A}^{\rm inv}Y\rangle
 \right|
 \le
 \|\ell_{\mathscr A}[u]\|_2
 \|\mathsf P_{e,\mathscr A}^{\rm inv}Y\|_2.}
 \label{eq:v19-leakage-bound}
\end{equation}
Thus a support-specific invariant may be shorted in the full equation only
after its exact outer-parent leakage has been evaluated.
\end{theorem}

\begin{proof}
Since $M_qe_{\mathscr A}$ is bank supported,
\[
 \langle c,M_qe_{\mathscr A}\rangle
 =\langle P_{\mathscr A}c,M_qe_{\mathscr A}\rangle.
\]
Insert \eqref{eq:v19-source-split}.  The internal term vanishes by
Theorem~\ref{thm:v19-bank-source-null}, giving
\eqref{eq:v19-invariant-leakage-pairing}.  Decompose $Y$ with the orthogonal
projection \eqref{eq:v19-state-invariant-projection}.  Its projected part is
a linear combination of invariant writers, so
\eqref{eq:v19-invariant-leakage-pairing} applies and gives
\eqref{eq:v19-full-contact-short}.  Cauchy--Schwarz proves the last bound.
\end{proof}

\begin{corollary}[Time-integrated closure debit]
\label{cor:v19-time-leakage}
For a time-dependent finite bank and writer $Y(t)$, define
\begin{equation}
 \mathcal L_{\mathscr A}^{\rm triad}
 :=\int_I
 \left\langle
 \ell_{\mathscr A}[u(t)],
 \mathsf P_{e(t),\mathscr A}^{\rm inv}Y(t)
 \right\rangle\,\dd t.
 \label{eq:v19-time-leakage}
\end{equation}
Then the complete integrated source contact equals the fully invariant-
shorted contact plus $\mathcal L_{\mathscr A}^{\rm triad}$.  A nonzero
normalized value is a signed outer-parent or high--high-return survivor on
the same datum.
\end{corollary}

\begin{proof}
Integrate \eqref{eq:v19-full-contact-short}.  No positive part is taken
before the complete leakage pairing is assembled.
\end{proof}

\begin{assessment}[Why this is not shell counting]
\label{ass:v19-no-shell-count}
The norm of \eqref{eq:v19-outer-parent-leakage} may be estimated by fused
square functions or source-relative capacities, but the identity itself has
one source occurrence.  Decomposing the outer parents into shells and
summing their absolute contributions would generally duplicate cancellation.
The terminal verdict is the signed fused leakage in
\eqref{eq:v19-time-leakage}, or a proved upper bound for that same fused
quantity.
\end{assessment}

% =============================================================================
\section{A spatial-rank-three bank with two extra invariants}
\label{sec:v19-counterexample}
% =============================================================================

The following exact bank separates spatial rank, signed-helical occupancy,
and affine triad connectivity.

\begin{countertheorem}[Spatial rank three does not imply affine triad rigidity]
\label{cth:v19-rank-three-not-rigid}
Fix an integer $N\ge4$.  Let the first reality-paired triad have wavevectors
\begin{equation}
 a_1=(1,0,0),
 \qquad
 a_2=(0,1,0),
 \qquad
 a_3=(-1,-1,0),
 \label{eq:v19-small-triad}
\end{equation}
and the second have
\begin{equation}
 b_1=N(1,0,1),
 \qquad
 b_2=N(0,1,1),
 \qquad
 b_3=-N(1,1,2).
 \label{eq:v19-large-triad}
\end{equation}
Include their negative wavevector copies and choose positive helicity on all
six atoms.  Then:
\begin{enumerate}[label=\textup{(C\arabic*)},leftmargin=3.4em]
\item the spatial lattice span has rank three;
\item the only zero-sum wavevector triples in the bank are the two displayed
      triads and their negative copies;
\item both displayed helical triads are active;
\item the incidence matrix has rank two and invariant nullity four;
\item the energy--helicity coefficient space has dimension two, so
      $d_{\mathscr A}^{\rm ex}=2$.
\end{enumerate}
In particular, the component indicator
\begin{equation}
 q^{(1)}=(1,1,1,0,0,0)
 \label{eq:v19-component-indicator}
\end{equation}
is an invariant coefficient vector but is not of the form
$\alpha\mathbf1+\beta\mathbf x$.

If a Galerkin state has nonzero amplitude on all six atoms and
$Y=M_{q^{(1)}}e$, then
\begin{equation}
 \boxed{
 \langle c_{\mathscr A}^{\rm int},Y\rangle=0,
 \qquad
 (I-\mathsf P_{e,\mathscr A}^{\rm EH})Y\ne0,
 \qquad
 (I-\mathsf P_{e,\mathscr A}^{\rm inv})Y=0.}
 \label{eq:v19-counterexample-writer}
\end{equation}
Thus the energy--helicity quotient can assign positive writer action to an
exactly source-null support-specific invariant.
\end{countertheorem}

\begin{proof}
The vectors $a_1,a_2,b_1$ span $\mathbb R^3$, proving spatial rank three.
Every vector in the second triad has all coordinates divisible by $N$, while
every nonzero vector in the first triad has a coordinate equal to $\pm1$.
A mixed zero-sum triple with two large vectors would force a nonzero small
vector to be divisible by $N$; a mixed triple with one large vector has
large-vector norm at least $N\sqrt2>2\sqrt2$, exceeding the sum of the two
small-vector norms.  Hence there is no mixed triad.

For normalized helical vectors, direct evaluation of
\eqref{eq:v19-geometric-coefficient} gives
\[
 |\Gamma_{\mathfrak t_1}|^2=\frac{3+2\sqrt2}{8},
 \qquad
 |\Gamma_{\mathfrak t_2}|^2=\frac{7+4\sqrt3}{32}.
\]
Both coefficients are nonzero, so both triads are active.  Their signed frequency lists are
\[
 (1,1,\sqrt2),
 \qquad
 (N\sqrt2,N\sqrt2,N\sqrt6).
\]
Up to nonzero row factors, the incidence matrix is
\begin{equation}
 \mathfrak A_{\mathscr A}
 =\begin{pmatrix}
  1-\sqrt2&\sqrt2-1&0&0&0&0\\
  0&0&0&N(\sqrt2-\sqrt6)&N(\sqrt6-\sqrt2)&0
 \end{pmatrix}.
 \label{eq:v19-explicit-incidence}
\end{equation}
It has rank two, so its kernel has dimension four.  The vector
$q^{(1)}$ lies in that kernel.  If it were
$\alpha\mathbf1+\beta\mathbf x$, its zero values at the two distinct large
frequencies $N\sqrt2$ and $N\sqrt6$ would force $\alpha=\beta=0$, contrary
to its value one on the first component.

The first identity in \eqref{eq:v19-counterexample-writer} is
Theorem~\ref{thm:v19-bank-source-null}.  The full invariant projection fixes
$Y$ because $q^{(1)}$ is in the coefficient kernel, proving the third
identity.  If the energy--helicity residual vanished, then, on every occupied
atom, $q^{(1)}$ would agree with an affine function of $x$.  The preceding
argument rules this out, proving the middle identity.
\end{proof}

\begin{firewall}[Scope of the counterexample]
\label{fire:v19-counterexample-scope}
Countertheorem~\ref{cth:v19-rank-three-not-rigid} is an exact finite
Galerkin and donor-table obstruction.  It is not a singular Navier--Stokes
trajectory.  In the full equation, interactions with modes outside the two
components appear through the explicit leakage
\eqref{eq:v19-outer-parent-leakage}.  The example proves only the logically
necessary point: spatial lattice rank three and the universal
energy--helicity short do not certify that a selected finite bank has no
additional source-null writer.
\end{firewall}

% =============================================================================
\section{Insertion before the V18 terminal export}
\label{sec:v19-terminal-gate}
% =============================================================================

Let a candidate V18 terminal thread be given.  Before exporting it to the
finite donor programme, choose a predeclared finite reality-paired helical
head $\mathscr A_n$ and its orthogonal projection $\Pi_n$ such that the
selected source-first writer $W_{n,t}^{\rm nat}$ is head supported.  Put
\begin{equation}
 Y_{n,t}:=\Pi_n\widehat Z_{n,t}^{\rm cal},
 \qquad
 e_{n,t}:=\Pi_ne(t),
 \label{eq:v19-terminal-head-fields}
\end{equation}
and construct from the literal active donor table:
\[
 \mathfrak A_n,
 \quad
 K_n^{\rm ex},
 \quad
 \mathsf P_{n,t}^{\rm inv},
 \quad
 \mathsf P_{n,t}^{\rm EH,h}.
\]
The last two are the full bank-invariant and bank energy--helicity
projections at the occurring state $e_{n,t}$.

Define the two head writers
\begin{equation}
 Z_{n,t}^{\rm h,EH}
 :=(I-\mathsf P_{n,t}^{\rm EH,h})Y_{n,t},
 \qquad
 Z_{n,t}^{\rm h,inv}
 :=(I-\mathsf P_{n,t}^{\rm inv})Y_{n,t},
 \label{eq:v19-two-head-writers}
\end{equation}
and their mixed reads
\begin{equation}
 N_n^{\rm h,EH}
 :=\int_{\supp\eta_n}
   \langle Z_{n,t}^{\rm h,EH},W_{n,t}^{\rm nat}\rangle\,\dd t,
 \qquad
 N_n^{\rm h,inv}
 :=\int_{\supp\eta_n}
   \langle Z_{n,t}^{\rm h,inv},W_{n,t}^{\rm nat}\rangle\,\dd t.
 \label{eq:v19-head-mixed-reads}
\end{equation}
The additional invariant transfer and global-to-head capture residual are
\begin{equation}
 \boxed{
 N_n^{\rm ex}:=N_n^{\rm h,EH}-N_n^{\rm h,inv},
 \qquad
 N_n^{\rm cap}:=N_n^{\rm EH}-N_n^{\rm h,EH}.}
 \label{eq:v19-extra-capture-rows}
\end{equation}
They give the exact numerator chain
\begin{equation}
 \boxed{
 N_n^{\rm EH}
 =N_n^{\rm h,inv}+N_n^{\rm ex}+N_n^{\rm cap}.}
 \label{eq:v19-numerator-chain}
\end{equation}
No commutation of the terminal history projection with the Stokes operator
is assumed; every failure is contained in $N_n^{\rm cap}$.

Normalize the two signed transfer rows by
\begin{equation}
 \Delta_n^{\rm ex}
 :=-\sigma_*A_n
   \frac{\nu N_n^{\rm ex}}
        {\mathcal R_n\mathfrak d_n^{\rm nat}},
 \qquad
 \Delta_n^{\rm cap}
 :=-\sigma_*A_n
   \frac{\nu N_n^{\rm cap}}
        {\mathcal R_n\mathfrak d_n^{\rm nat}}.
 \label{eq:v19-normalized-head-defects}
\end{equation}
If
\begin{equation}
 S_n^{\rm h,inv}
 :=-\sigma_*A_n
   \frac{\nu N_n^{\rm h,inv}}
        {\mathcal R_n\mathfrak d_n^{\rm nat}},
 \label{eq:v19-head-selected-read}
\end{equation}
then
\begin{equation}
 \boxed{
 S_n^{\rm EH}
 =S_n^{\rm h,inv}+\Delta_n^{\rm ex}+\Delta_n^{\rm cap}.}
 \label{eq:v19-selected-chain}
\end{equation}
The corresponding source-to-head residual is
\begin{equation}
 \boxed{
 I_n^{\rm h,inv}
 :=\mathfrak c_n^{\rm cal}-\mathfrak c_n^{\rm rep}
   -S_n^{\rm h,inv}
 =I_n^{\rm EH}+\Delta_n^{\rm ex}+\Delta_n^{\rm cap}.}
 \label{eq:v19-incidence-chain}
\end{equation}

The full-source closure row on the same head is
\begin{equation}
 \boxed{
 \mathcal L_n^{\rm triad}
 :=\int_{\supp\eta_n}
 \left\langle
 \ell_{\mathscr A_n}[u(t)],
 \mathsf P_{n,t}^{\rm inv}Y_{n,t}
 \right\rangle\,\dd t.}
 \label{eq:v19-terminal-leakage}
\end{equation}
By Theorem~\ref{thm:v19-full-leakage-identity}, this is exactly the contact
lost by treating the finite bank invariants as full-PDE nulls.

\begin{theorem}[Triad-incidence gate before the NS--A/NS--B export]
\label{thm:v19-terminal-gate}
For any cofinal candidate V18 export and any predeclared exhausting sequence
of donor heads, pass to a priority subsequence.  At least one of the following
occurs.
\begin{enumerate}[label=\textup{(G19.\arabic*)},leftmargin=5.4em]
\item \emph{Finite-head capture failure:} for some $\delta>0$,
      \begin{equation}
       |\Delta_n^{\rm cap}|\ge\delta
       \label{eq:v19-capture-survivor}
      \end{equation}
      cofinally.  The global V18 writer is not represented by the declared
      finite donor head.
\item \emph{Outer-parent triad leakage:} for some $\delta>0$,
      \begin{equation}
       \frac{|\mathcal L_n^{\rm triad}|}{\mathcal R_n^2}
       \ge\delta
       \label{eq:v19-leakage-survivor}
      \end{equation}
      cofinally.  The internal finite-bank invariant is replenished by
      parents outside the bank.
\item \emph{Additional invariant transfer:} the preceding branches fail,
      $d_{\mathscr A_n}^{\rm ex}>0$, and for some $\delta>0$,
      \begin{equation}
       |\Delta_n^{\rm ex}|\ge\delta
       \label{eq:v19-extra-survivor}
      \end{equation}
      cofinally.  A support-specific invariant meets the selected terminal
      writer with record-size signed incidence.
\item \emph{Nonrigid selected-null kernel:} the capture and leakage rows
      vanish, $d_{\mathscr A_n}^{\rm ex}>0$, and
      \begin{equation}
       \Delta_n^{\rm ex}\longrightarrow0.
       \label{eq:v19-extra-null-selected}
      \end{equation}
      The additional invariant coefficient space is genuine, but its
      complete signed incidence with the selected writer is asymptotically
      null.  The selected scalar may pass this finite target-specific gate;
      no full-bank rigidity or coercivity conclusion is inferred.
\item \emph{All-shell triadic escape:} no finite donor-head sequence makes
      the capture and leakage rows vanish while retaining the selected
      scalar.  The output is a fused all-shell interaction survivor, not a
      finite invariant card.
\item \emph{Affine-rigid closed head:} there is a cofinal finite-head
      sequence for which
      \begin{equation}
       d_{\mathscr A_n}^{\rm ex}=0,
       \qquad
       \Delta_n^{\rm cap}\to0,
       \qquad
       \frac{\mathcal L_n^{\rm triad}}{\mathcal R_n^2}\to0.
       \label{eq:v19-rigid-clean-head}
      \end{equation}
      Then $\Delta_n^{\rm ex}=0$ identically, the V18 energy--helicity
      quotient is the complete diagonal invariant quotient of the actual
      donor head, and
      \begin{equation}
       \boxed{
       I_n^{\rm h,inv}-I_n^{\rm EH}\longrightarrow0.}
       \label{eq:v19-rigid-incidence-alignment}
      \end{equation}
      The corrected V18 card is eligible for the NS--B spatial-rank-three,
      Reynolds-covariance, and paid-stress audit.
\end{enumerate}
The nonrigid selected-null branch is a target-specific PASS only for the
predeclared scalar.  Its additional invariant kernel remains explicit, so
nullity alone is neither a general PASS nor an obstruction.
\end{theorem}

\begin{proof}
Apply successive subsequence selection to the nonnegative magnitudes of the
capture and normalized leakage rows.  If neither has a positive lower bound,
they tend to zero.  On every finite head, the integer
$d_{\mathscr A_n}^{\rm ex}$ and the signed transfer
$\Delta_n^{\rm ex}$ are then exact finite reads.  On a nonrigid head,
subsequence selection makes the transfer either retain a positive magnitude,
giving the third branch, or tend to zero, giving the fourth.  If no finite
sequence simultaneously captures the selected scalar and makes the leakage
small, the fifth branch is the definition of all-shell triadic escape.  On
an affine-rigid head, Corollary~\ref{cor:v19-extra-invariant-transfer} gives
$\Delta_n^{\rm ex}=0$, and
\eqref{eq:v19-incidence-chain} gives
\eqref{eq:v19-rigid-incidence-alignment}.  These alternatives exhaust the
priority order.
\end{proof}

\begin{corollary}[No further diagonal-invariant search after a rigid pass]
\label{cor:v19-no-more-invariants}
On branch \textup{(G19.6)}, every helical-diagonal quadratic writer which is
nonlinearly null on the actual finite donor head is a linear combination of
the energy and helicity writers.  Any subsequent clean-export failure is
therefore one of the already named V18 metric, represented, aggregate,
source-ancestry, Reynolds, paid-stress, or stock-budget rows.
\end{corollary}

\begin{proof}
This is the affine-rigidity identity
\eqref{eq:v19-affine-rigid} transported through the state synthesis in
Theorem~\ref{thm:v19-bank-source-null}.
\end{proof}

\begin{firewall}[Three different rank gates]
\label{fire:v19-three-ranks}
The following statements are distinct:
\[
 \boxed{
 \begin{gathered}
 \text{at least three occupied signed curl levels},\\
 \text{spatial Fourier-subgroup rank three},\\
 \text{affine rigidity of the active helical-triad incidence matrix}.
 \end{gathered}}
 \tag{V19.RK}
\]
The first makes the V18 caloric regression potentially nonzero.  The second
removes isolated rank-at-most-two $2$D$3$C carriers.  The third proves that
energy and helicity exhaust the finite diagonal invariant kernel.
Countertheorem~\ref{cth:v19-rank-three-not-rigid} shows that the second does
not imply the third.  A donor export must retain all three on the same
history when all three are part of the intended conclusion.
\end{firewall}

% =============================================================================
\section{Version V19 integrated theorem and proof-status ledger}
\label{sec:v19-integrated}
% =============================================================================

\begin{theorem}[Version V19 actual-triad and affine-rigidity theorem]
\label{thm:v19-integrated}
Preserving every theorem and limitation of Versions~V1--V18, Version~V19
proves the following.
\begin{enumerate}[label=\textup{(T19.\arabic*)},leftmargin=5.7em]
\item Every active reality-paired helical triad has the exact modal-transfer
      vector \eqref{eq:v19-triad-transfer}.
\item The kernel of the finite incidence matrix
      \eqref{eq:v19-incidence-matrix} is exactly the complete space of
      helical-diagonal quadratic invariants of the finite Galerkin donor
      system.
\item Energy and helicity always lie in this kernel.  The Moore--Penrose
      difference projector \eqref{eq:v19-extra-projector} gives the exact
      number and coefficient space of additional support-specific
      invariants.
\item A two-overlap active-triad chain is an explicit sufficient certificate
      of affine rigidity; the rank tests
      \eqref{eq:v19-rigidity-tests} are necessary and sufficient.
\item The actual state converts the invariant coefficient kernel into the
      positive writer Gram and projection
      \eqref{eq:v19-state-invariant-projection}.  The resulting writer is
      the unique minimum modulo every bank invariant.
\item The V18 energy--helicity writer equals the complete bank-invariant
      writer plus one orthogonal extra-invariant debit.  Its effect on the
      selected writer is the signed row
      \eqref{eq:v19-extra-mixed-row}.
\item In the full Navier--Stokes equation, failure of a finite-bank invariant
      is exactly the outer-parent leakage
      \eqref{eq:v19-outer-parent-leakage}; the contact identity is
      \eqref{eq:v19-full-contact-short}.
\item A literal spatial-rank-three six-atom bank has two disconnected active
      triads and two extra invariant directions.  It gives a writer with
      positive energy--helicity residual and zero internal source contact.
\item The V18 terminal export now passes through the exact finite gate of
      Theorem~\ref{thm:v19-terminal-gate}: capture failure, outer-parent
      leakage, additional invariant transfer, a nonrigid selected-null
      kernel, all-shell triadic escape, or an affine-rigid closed donor head.
\item On the last branch, the V18 quotient is maximal in the declared
      diagonal invariant class and NS--A terminates at the existing NS--B
      donor/paid-stress interface.
\end{enumerate}
No populated terminal donor table, upper paid-stress budget, critical
acceleration estimate, Dini improvement, terminal trace, backward uniqueness,
Liouville theorem, or global regularity conclusion is supplied.
\end{theorem}

\begin{proof}
Items~\textup{(T19.1)--(T19.3)} are
\Cref{lem:v19-triad-transfer,thm:v19-invariant-kernel,cor:v19-EH-in-kernel,def:v19-extra-nullity}.
Item~\textup{(T19.4)} is
\Cref{thm:v19-two-overlap-rigidity,cor:v19-finite-rigidity-audit}.
Items~\textup{(T19.5)--(T19.6)} are
\Cref{thm:v19-bank-source-null,cor:v19-extra-invariant-transfer}.
Item~\textup{(T19.7)} is
\Cref{thm:v19-full-leakage-identity,cor:v19-time-leakage}.
Item~\textup{(T19.8)} is
Countertheorem~\ref{cth:v19-rank-three-not-rigid}.  The last two items are
Theorem~\ref{thm:v19-terminal-gate} and
Corollary~\ref{cor:v19-no-more-invariants}.
\end{proof}

\begin{firewall}[Exact stopping boundary after V19]
\label{fire:v19-stop}
Version~V19 closes the possibility of an endless sequence of guessed
helical-diagonal invariant shorts.  The actual finite donor table has one
terminating rank computation.  If it is affine rigid, energy and helicity
are complete.  If it is not, the additional coefficient kernel and its
signed selected-writer incidence are explicit.  If the finite bank is not
closed under the full source, the exact outer-parent leakage is explicit.
A later compactness theorem, positive Reynolds estimate, or paid-stress
bound cannot retroactively populate any of these earlier rows.
\end{firewall}

\begingroup\small
\begin{longtable}{p{0.25\textwidth}p{0.19\textwidth}p{0.48\textwidth}}
\caption{NS--A Version V19 proof-status ledger.}\label{tab:v19-ledger}\\
\toprule Row & Status & Exact conclusion\\\midrule
\endfirsthead
\toprule Row & Status & Exact conclusion\\\midrule
\endhead
V1--V18 prefix & \PROVED{} preserved
 &The complete V18 preterminal source is retained byte-for-byte; no earlier
 theorem or limitation is rewritten.\\[0.3em]
One active helical triad & \PROVED{} exact
 &Its modal energy-transfer vector is the signed-frequency difference row
 \eqref{eq:v19-triad-transfer}.\\[0.3em]
Finite diagonal invariant space & \PROVED{} exact
 &$\mathcal I_{\mathscr A}=\Ker\mathfrak A_{\mathscr A}$.\\[0.3em]
Energy--helicity coefficient space & \INHERITED{} / embedded
 &$\operatorname{span}\{\mathbf1,\mathbf x\}$ is a canonical subspace of the
 finite invariant kernel.\\[0.3em]
Extra invariant nullity & \PROVED{} finite
 &$d_{\mathscr A}^{\rm ex}$ is the trace of the Moore--Penrose difference
 projector \eqref{eq:v19-extra-projector}.\\[0.3em]
Affine-rigidity criterion & \PROVED{} terminating
 &The rank test \eqref{eq:v19-rigidity-tests} is exact; the two-overlap chain
 is a sufficient combinatorial certificate.\\[0.3em]
State-level invariant writer & \PROVED{} positive
 &The Gram \eqref{eq:v19-state-invariant-projection} gives the unique minimum
 writer modulo all populated bank invariants.\\[0.3em]
Additional invariant action debit & \PROVED{} positive
 &The V18 action splits orthogonally as in
 \eqref{eq:v19-extra-action-debit}.\\[0.3em]
Additional invariant mixed incidence & \OPEN{} until populated
 &Its signed read is \eqref{eq:v19-extra-mixed-row}; nullity or positive debit
 does not determine its sign.\\[0.3em]
Nonrigid selected-null kernel & \PROVED{} target-specific branch
 &After the signed row is acquired, \eqref{eq:v19-extra-null-selected} may
 pass the predeclared scalar without proving full-bank rigidity.\\[0.3em]
Outer-parent finite-head leakage & \PROVED{} exact residual
 &Equation \eqref{eq:v19-full-contact-short} is the complete full-PDE
 correction to the finite invariant short.\\[0.3em]
Spatial rank three implies rigidity & \OBSTRUCTION{} explicit
 &Countertheorem~\ref{cth:v19-rank-three-not-rigid} has spatial rank three and
 two extra diagonal invariants.\\[0.3em]
V18-to-donor head capture & \OPEN{} physical incidence
 &$N_n^{\rm cap}$ must be populated on the literal selected history.\\[0.3em]
Affine-rigid closed donor head & \CONDITIONAL{} exact handoff
 &Once \eqref{eq:v19-rigid-clean-head} passes, the V18 quotient is maximal in
 the declared invariant class and exports to NS--B.\\[0.3em]
All-shell triadic escape & \OPEN{} typed alternative
 &Failure of every finite capture/closure head remains a fused all-shell
 interaction survivor, not a shell count.\\[0.3em]
Three-dimensional regularity & \OPEN
 &No surviving capture, leakage, extra-invariant, all-shell, donor-budget, or
 terminal PDE branch is excluded for every datum.\\
\bottomrule
\end{longtable}
\endgroup

% =============================================================================
\section{Exact mandate after Version V19}
\label{sec:v19-next}
% =============================================================================

\begin{programme}[Populate the actual donor head or stop NS--A]
\label{prog:v19-next}
Preserve Versions~V1--V19 verbatim.  There is no automatic generic
Version~V20.

A continuation is licensed only in the following order.
\begin{enumerate}[label=\textup{(V20.\arabic*)},leftmargin=5.9em]
\item On the literal selected terminal history, populate one finite
      reality-paired donor head, its complete active helical-triad list, and
      the exact matrix $\mathfrak A_n$.
\item Evaluate the global-to-head capture row $N_n^{\rm cap}$ and the fused
      outer-parent leakage $\mathcal L_n^{\rm triad}$ before using any finite
      invariant short.
\item If $d_{\mathscr A_n}^{\rm ex}>0$, acquire the signed extra-invariant
      mixed row $N_n^{\rm ex}$.  Do not replace it by the positive action
      debit or by the integer nullity.
\item If $d_{\mathscr A_n}^{\rm ex}>0$ but the acquired signed row
      $N_n^{\rm ex}$ vanishes, retain the additional invariant kernel and
      export only the predeclared scalar target.  Do not report full-bank
      rigidity or coercivity.
\item On an affine-rigid, capture-tight, leakage-null branch, terminate
      NS--A.  Export the V18 energy--helicity card, the exact rigidity
      certificate, and the existing spatial-rank-three and source-ancestry
      rows to NS--B.
\item If no finite head captures the scalar with vanishing leakage, retain
      all-shell triadic escape.  Continue only if a new fused square-function,
      source-relative capacity, or critical acceleration upper stock controls
      that exact escape on the same history.
\item Do not introduce another invariant by inspection, another shell-count
      argument, path law, Dini modulus, profile space, or generic response
      compiler.  The finite kernel calculation is the complete diagonal
      invariant audit.
\end{enumerate}
\end{programme}

\begin{assessment}[Programme position after V19]
\label{ass:v19-position}
The upstream invariant question is now finite and terminating.  V18 gives
the universal energy--helicity quotient.  V19 says exactly when that quotient
is also complete on the actual donor head, identifies every additional
support-specific invariant when it is not, and keeps the full-PDE correction
as one signed outer-parent leakage.  The first unpopulated object is now a
literal donor-table/head-incidence packet, not another theorem-level
normalization.  On the affine-rigid closed branch NS--A has reached its
intended single interface and must stop in favor of NS--B.
\end{assessment}

\paragraph{Version V19 history.}
Preserved the complete V18 source.  Moved upstream of its signed helicity
row to the actual finite helical donor table.  Derived the exact modal
transfer row of one active triad and assembled the finite incidence matrix
whose kernel is the complete helical-diagonal quadratic-invariant space.
Constructed the coefficient-space Moore--Penrose extra-invariant projector,
proved the two-overlap affine-rigidity theorem, and gave the exact finite rank
test.  Converted the coefficient kernel into the occurring state-level
invariant-writer Gram and proved the minimum-writer Pythagoras.  Identified
the positive extra-invariant debit and its independent signed mixed row.
For the full equation, proved that every finite-bank invariant failure is
exactly the fused outer-parent leakage.  Constructed an explicit spatial-
rank-three bank of two disconnected active helical triads with two extra
invariants.  Finally inserted head capture, leakage, and extra-invariant
rows before the V18 terminal export and stated the affine-rigid stop.

\begin{assessment}[Append-only integrity]
\label{ass:v19-integrity}
The complete Version~V18 source preceding its sole final executable document
terminator is preserved verbatim.  Its preterminal SHA--256 digest is
\begin{center}\scriptsize\ttfamily
6f15123a42dc9292ad16675f51985f761c122bc52382be35aa29bd822179050c
\end{center}
The present material begins at Section~\ref{sec:v19-scope}; the final command
below is again unique.
\end{assessment}

\addtocontents{toc}{\protect\endgroup}
% =============================================================================
% END VERSION V19 MATERIAL -- append later versions before final terminator
% =============================================================================

% APPEND ALL LATER VERSIONS IMMEDIATELY ABOVE THE SOLE FINAL LINE BELOW.

\clearpage
% =============================================================================
% VERSION V20: NEW MATERIAL.  The complete V1--V19 source above is preserved
% verbatim; only its sole active \end{document} line was removed before append.
% =============================================================================
\hypersetup{pdftitle={NS-A Terminal Source Concentration Programme: Cumulative V20},pdfauthor={A. Pelissier}}
\makeatletter
\addtocontents{toc}{\protect\begingroup\protect\makeatletter}
\addtocontents{toc}{\protect\def\protect\l@section{\protect\@dottedtocline{1}{0em}{4.7em}}}
\addtocontents{toc}{\protect\def\protect\l@subsection{\protect\@dottedtocline{2}{1.5em}{5.3em}}}
\addtocontents{toc}{\protect\makeatother}
\makeatother
\renewcommand{\thesection}{V20.\arabic{section}}
\renewcommand{\thesubsection}{V20.\arabic{section}.\arabic{subsection}}
\renewcommand{\theequation}{V20.\arabic{equation}}
\setcounter{section}{0}
\setcounter{equation}{0}
\renewcommand{\theHsection}{V20.\arabic{section}}
\renewcommand{\theHsubsection}{V20.\arabic{section}.\arabic{subsection}}
\renewcommand{\theHtheorem}{V20.\arabic{section}.\arabic{theorem}}
\renewcommand{\theHequation}{V20.\arabic{equation}}

\begin{center}
{\LARGE\bfseries NS--A: Version V20}\\[0.5em]
{\large The Internal-to-Exterior Daughter Source, the No-Finite-Reducing-Head\\
Theorem, and the Ancestry-Complete Open Donor Export}\\[0.5em]
{\normalsize 3 September 2026}
\end{center}
\addcontentsline{toc}{part}{Version V20: daughter-source boundary germ and open donor export}

\begin{abstract}
Version~V19 completed the helical-diagonal invariant audit of a declared
finite donor head and retained the incoming outer-parent leakage which
modifies those finite-bank invariants in the full equation.  That is the
correct correction for the contact of a head-supported invariant writer.
It is not, however, the complete source-ancestry boundary of the head.  The
same head modes can generate a nonlinear source at output frequencies
outside the head even when there are no outside parents and hence when the
V19 incoming leakage is exactly zero.

Version~V20 constructs this missing internal-to-exterior daughter source.
For a finite Fourier head $P$ and $v=Pu$, it is
\[
 d_P^{\rm out}=A^{-1/4}(I-P)\mathcal B(v).
\]
Together with the V19 incoming row it gives the minimum two-sided source
boundary packet.  The complete parent-output coefficient table splits
orthogonally into internal and exterior outputs, and its exterior Gram is a
finite positive matrix.  A finite polarization identity reconstructs this
matrix from the actual quadratic donor coefficients.

A sharp no-go theorem is proved.  A finite, reality-symmetric,
polarization-complete Fourier set is universally closed under the Euler
quadratic nonlinearity if and only if all of its wavevectors are collinear.
For every finite noncollinear set one can choose two divergence-free modes
whose interaction creates a nonzero exterior daughter.  Consequently no
finite polarization-complete spatial-rank-three head is a universal
nonlinear reducing carrier.  A restricted helical head may still have
coefficient cancellations, but those are decided by the displayed exterior
coefficient Gram, not by spatial rank or by the internal triad-incidence
matrix.

The theorem also separates source birth from stock transfer.  With no
outside field, the V19 incoming leakage, the head energy current, and the
head helicity current all vanish, while the exterior daughter source may be
nonzero.  Thus a newly generated exterior source cannot be erased by a zero
instantaneous energy or helicity current.  It must first be shorted against
a physically declared exterior source/return dictionary.  The resulting
Pythagoras separates represented daughter output from a genuinely new
source birth without charging either twice.

The V19 ``affine-rigid closed head'' is therefore corrected to an
\emph{affine-rigid ancestry-complete open head}.  The internal diagonal
invariant calculation remains exact, but export to NS--B additionally
requires the unrepresented daughter action to vanish at the record scale,
or an exact same-history ancestry map which identifies that daughter with
an already represented source.  Failure returns a daughter-source birth,
a source-germ boundary leak, or all-shell source escape.  No Dini
improvement, finite-stock upper bound, Liouville theorem, or global
regularity conclusion is claimed.
\end{abstract}

% =============================================================================
\section{The missing arrow is source birth, not another invariant}
\label{sec:v20-scope}
% =============================================================================

Version~V19 proved that, for every coefficient vector
$q\in\Ker\mathfrak A_{\mathscr A}$ and every head-supported invariant
writer $M_qe_{\mathscr A}$, the full-source failure of the finite Galerkin
null is exactly the incoming outer-parent leakage
\eqref{eq:v19-outer-parent-leakage}.  This statement remains correct.  It
computes the part of the full source which lands \emph{inside} the head from
parents not both contained in the head.

There is a different source row.  Two parents already inside the head may
write an output outside it.  This internal-to-exterior daughter is invisible
to every head-supported writer at the same instant, but it enlarges the
source atlas and can later return through the V19 incoming row.  Declaring
the head ``closed'' before this first-birth row is evaluated reverses the
source-ancestry order.

The corrected order is
\begin{equation}
 \boxed{
 \begin{gathered}
 \text{complete internal parent-output table}\\[-0.15em]
 \Downarrow\\[-0.15em]
 \text{internal versus exterior output Gram}\\[-0.15em]
 \Downarrow\\[-0.15em]
 \text{daughter first-birth short against the represented exterior atlas}\\[-0.15em]
 \Downarrow\\[-0.15em]
 \text{incoming return and finite invariant short}\\[-0.15em]
 \Downarrow\\[-0.15em]
 \text{ancestry-complete open donor export or all-shell source escape}
 \end{gathered}}
 \label{eq:v20-upstream-order}
\end{equation}
No additional helical-diagonal invariant is introduced.  The V19 incidence
matrix remains the complete diagonal invariant audit of the internal
Galerkin source.

\begin{firewall}[Three notions of closure]
\label{fire:v20-three-closures}
For a finite Fourier head the following statements are distinct.
\begin{enumerate}[label=\textup{(C\arabic*)},leftmargin=3.4em]
\item A selected head-supported writer is \emph{target transparent} when its
      pairing with the incoming leakage vanishes.
\item The occurring head state is \emph{source-ancestry closed relative to a
      declared exterior atlas} when every daughter source is represented by
      that atlas and every later return carries the same source ancestry.
\item The head is \emph{universally nonlinear reducing} when every
      head-supported divergence-free field generates no exterior source and
      no cross-boundary source for any complementary field.
\end{enumerate}
The first is one scalar, the second is a same-history source statement, and
the third is an operator identity.  None is inferred from another.  The
finite-rank theorem below shows that the third is unavailable for a
polarization-complete rank-three head.
\end{firewall}

Version~V20 proves the following chain.
\begin{enumerate}[label=\textup{(U20.\arabic*)},leftmargin=5.2em]
\item Every two nonparallel Fourier wavevectors admit divergence-free
      polarizations whose quadratic interaction has a nonzero output at
      their sum.
\item A finite polarization-complete Fourier head is universally
      self-closed exactly when all wavevectors are collinear.
\item The exterior parent-output synthesis and its positive Gram form a
      finite terminating daughter-source certificate.
\item The V19 incoming leakage and the new daughter output form the minimum
      two-sided source boundary packet; neither may be used as a surrogate
      for the other.
\item Energy and helicity have one signed boundary current, but that current
      may vanish while a new exterior source is born.
\item A Moore--Penrose short against the existing exterior source atlas
      separates represented daughter output from a genuine source birth.
\item The terminal export becomes an ancestry-complete open-head theorem,
      not a finite support-closure theorem.
\end{enumerate}

% =============================================================================
\section{A nonparallel pair creates an exterior daughter}
\label{sec:v20-pair-escape}
% =============================================================================

Use normalized Haar measure on $\Torus^3$.  For a finite set
$K=-K\subset\mathbb Z^3\setminus\{0\}$, let $P_K$ be the orthogonal Fourier
projection onto all divergence-free polarizations at the wavevectors in
$K$, and put $Q_K=I-P_K$.  Denote by $\mathscr H_K$ the real
finite-dimensional space of smooth divergence-free fields supported in
$K$.  The Leray-projected quadratic source is
\[
 \mathcal B(v)=\mathbb P(v\cdot\nabla v).
\]

\begin{lemma}[Finite additive escape]
\label{lem:v20-additive-escape}
If $K=-K$ is finite and is not contained in one rank-one subgroup of
$\mathbb Z^3$, then there exist nonparallel $p,q\in K$ such that
\begin{equation}
 p+q\ne0,
 \qquad
 p+q\notin K.
 \label{eq:v20-additive-escape}
\end{equation}
\end{lemma}

\begin{proof}
Choose a real linear functional $\xi\cdot k$ which is nonzero on every
$k\in K$ and has a unique maximizer $p\in K$.  Such $\xi$ exists because
one excludes only finitely many hyperplanes.  Since $K$ is not collinear,
choose $r\in K$ nonparallel to $p$.  Replacing $r$ by $-r$ if necessary,
put $q\in\{r,-r\}$ so that $\xi\cdot q>0$.  Then
\[
 \xi\cdot(p+q)>\xi\cdot p,
\]
so $p+q\notin K$.  The vectors are nonparallel, and the positive value of
$\xi\cdot(p+q)$ also gives $p+q\ne0$.
\end{proof}

\begin{lemma}[Explicit nonzero daughter coefficient]
\label{lem:v20-pair-daughter}
Let $p,q\in\mathbb Z^3\setminus\{0\}$ be nonparallel and set $k=p+q$.
There is a real field $v$ supported on $\{\pm p,\pm q\}$ for which
\begin{equation}
 \widehat{\mathcal B(v)}(k)\ne0.
 \label{eq:v20-nonzero-daughter}
\end{equation}
More precisely, put
\begin{equation}
 a=\frac{p\times q}{|p\times q|},
 \qquad
 b=\frac{P_{q^\perp}p}{|P_{q^\perp}p|},
 \label{eq:v20-explicit-polarizations}
\end{equation}
and prescribe the real Fourier amplitudes
$\widehat v(p)=a$, $\widehat v(q)=b$ and their conjugate negative copies.
Then
\begin{equation}
 \boxed{
 \widehat{\mathcal B(v)}(p+q)
 =\mathrm i\,|P_{q^\perp}p|\,a.}
 \label{eq:v20-explicit-daughter}
\end{equation}
The harmless common Fourier normalization depends only on the convention
used to write a real cosine pair.
\end{lemma}

\begin{proof}
The Fourier coefficient of the ordered quadratic term at $k=p+q$ is
\[
 \mathrm i\bigl((a\cdot q)b+(b\cdot p)a\bigr).
\]
The choice in \eqref{eq:v20-explicit-polarizations} gives
$a\cdot q=0$ and
$b\cdot p=|P_{q^\perp}p|>0$.  Moreover $a$ is perpendicular to both
$p$ and $q$, hence to $k$.  The Leray projection at $k$ therefore fixes
$a$, proving \eqref{eq:v20-explicit-daughter}.  The negative modes make the
field real and give the conjugate coefficient at $-k$.
\end{proof}

\begin{theorem}[No finite polarization-complete nonlinear head beyond one dimension]
\label{thm:v20-no-finite-reducing-head}
For a finite symmetric wavevector set $K=-K$, the following are equivalent.
\begin{enumerate}[label=\textup{(\roman*)},leftmargin=2.8em]
\item For every $v\in\mathscr H_K$,
      \begin{equation}
       Q_K\mathcal B(v)=0.
       \label{eq:v20-universal-self-closure}
      \end{equation}
\item The set $K$ is contained in one rank-one subgroup of
      $\mathbb Z^3$.
\end{enumerate}
On this branch $\mathcal B(v)=0$ for every $v\in\mathscr H_K$.
Consequently every finite polarization-complete Fourier head of spatial
lattice rank at least two, in particular every such rank-three head, has a
nonzero internal-to-exterior daughter operator.
\end{theorem}

\begin{proof}
Suppose first that $K$ is not collinear.  By
Lemma~\ref{lem:v20-additive-escape}, choose nonparallel $p,q\in K$ with
$p+q\notin K$.  Lemma~\ref{lem:v20-pair-daughter} gives a field
$v\in\mathscr H_K$ whose source has a nonzero Fourier coefficient at
$p+q$.  Hence \eqref{eq:v20-universal-self-closure} fails.

Conversely, suppose $K\subset\mathbb Z n$ for one nonzero primitive
$n\in\mathbb Z^3$.  Every $v\in\mathscr H_K$ has the form
$v(x)=V(n\cdot x)$ and incompressibility gives $n\cdot V=0$.  Therefore
\[
 (v\cdot\nabla)v=(v\cdot n)\,\partial_\theta V=0,
\]
so $\mathcal B(v)=0$.  This proves both implications and the final
statement.
\end{proof}

\begin{corollary}[The V19 closed-head phrase cannot mean universal support closure]
\label{cor:v20-v19-closed-qualification}
On a polarization-complete finite donor head of spatial rank three,
branch~\textup{(G19.6)} of Theorem~\ref{thm:v19-terminal-gate} cannot be
interpreted as universal invariance of the Fourier support under the
quadratic Navier--Stokes source.  Its valid content is the V19
capture/incoming-leakage/invariant statement for the selected history.
A further daughter-source row is required before source-ancestry closure is
reported.
\end{corollary}

\begin{proof}
Theorem~\ref{thm:v20-no-finite-reducing-head} excludes universal
self-closure.  The V19 identities involve the projection of the full source
back into the head and therefore remain valid.  The omitted component is
$Q_K\mathcal B(P_Ku)$, which is orthogonal to every head-supported writer
and hence does not appear in the V19 selected contact.
\end{proof}

\begin{firewall}[Restricted helical heads]
\label{fire:v20-restricted-helical}
Theorem~\ref{thm:v20-no-finite-reducing-head} concerns a physical Fourier
head which contains the complete divergence-free polarization fibre at each
retained wavevector.  A restricted helical subspace may have additional
geometric zeros.  No spatial-rank argument decides those zeros.  They are
computed by the exterior helical coefficient matrix of the next section.
Thus V20 returns an exact finite test rather than asserting that every
restricted helical bank leaks.
\end{firewall}

% =============================================================================
\section{The complete parent-output germ and its positive exterior Gram}
\label{sec:v20-parent-output}
% =============================================================================

Let $\mathscr A$ be a finite reality-paired helical bank as in
\Cref{sec:v19-triad-row}.  Let $E_{\mathscr A}$ be its real amplitude space
and let $\operatorname{Sym}^2E_{\mathscr A}$ be the symmetric parent-pair
space.  Write
\begin{equation}
 \mathcal B_{\rm s}(a,b)
 :=\frac12\mathbb P\bigl((a\cdot\nabla)b+(b\cdot\nabla)a\bigr),
 \qquad
 \mathcal B(a)=\mathcal B_{\rm s}(a,a).
 \label{eq:v20-symmetric-bilinear-source}
\end{equation}
Let $P_{\mathscr A}$ be the orthogonal helical projection and
$Q_{\mathscr A}=I-P_{\mathscr A}$.  Define the two finite synthesis maps
\begin{align}
 \mathscr C_{\mathscr A}^{\rm int}(a\odot b)
 &:=A^{-1/4}P_{\mathscr A}\mathcal B_{\rm s}(a,b),
 \label{eq:v20-internal-synthesis}\\
 \mathscr C_{\mathscr A}^{\rm out}(a\odot b)
 &:=A^{-1/4}Q_{\mathscr A}\mathcal B_{\rm s}(a,b).
 \label{eq:v20-exterior-synthesis}
\end{align}
Only the finite output set generated by sums of two retained wavevectors can
occur in \eqref{eq:v20-exterior-synthesis}; hence both maps are finite
matrices after helical output coordinates are chosen.

For oriented helical atoms
$\alpha=(p,\sigma)$, $\beta=(q,\tau)$ and an output atom
$\gamma=(k,\rho)$ with $k=p+q\ne0$, one compatible coefficient convention is
\begin{equation}
 \boxed{
 \mathscr C_{\gamma;\alpha\beta}
 =\frac{\mathrm i}{2}|k|^{-1/2}
  \overline h_{k,\rho}\cdot
  \left[
   (h_{p,\sigma}\cdot q)h_{q,\tau}
   +(h_{q,\tau}\cdot p)h_{p,\sigma}
  \right].}
 \label{eq:v20-helical-boundary-coefficient}
\end{equation}
The coefficient belongs to the internal or exterior matrix according as
$\gamma$ is or is not retained by $\mathscr A$.  Reality pairs are assembled
before the real matrix is formed.

\begin{theorem}[Exact internal/exterior source Pythagoras]
\label{thm:v20-source-Pythagoras}
For $z\in E_{\mathscr A}$ and its field $u_{\mathscr A}(z)$, put
\begin{equation}
 c_{\mathscr A}^{\rm int}(z)
 =\mathscr C_{\mathscr A}^{\rm int}(z\odot z),
 \qquad
 d_{\mathscr A}^{\rm out}(z)
 =\mathscr C_{\mathscr A}^{\rm out}(z\odot z).
 \label{eq:v20-state-sources}
\end{equation}
Then
\begin{equation}
 \boxed{
 A^{-1/4}\mathcal B(u_{\mathscr A}(z))
 =c_{\mathscr A}^{\rm int}(z)
  +d_{\mathscr A}^{\rm out}(z),}
 \label{eq:v20-source-split}
\end{equation}
and the two terms are orthogonal.  Consequently
\begin{equation}
 \boxed{
 \left\|A^{-1/4}\mathcal B(u_{\mathscr A}(z))\right\|_2^2
 =\left\|c_{\mathscr A}^{\rm int}(z)\right\|_2^2
  +\left\|d_{\mathscr A}^{\rm out}(z)\right\|_2^2.}
 \label{eq:v20-source-action-Pythagoras}
\end{equation}
The positive exterior parent Gram
\begin{equation}
 \boxed{
 \mathbb D_{\mathscr A}^{\rm out}
 :=(\mathscr C_{\mathscr A}^{\rm out})^*
   \mathscr C_{\mathscr A}^{\rm out}\succeq0}
 \label{eq:v20-exterior-Gram}
\end{equation}
satisfies
\begin{equation}
 \operatorname{rank}\mathbb D_{\mathscr A}^{\rm out}
 =\operatorname{rank}\mathscr C_{\mathscr A}^{\rm out},
 \qquad
 \mathbb D_{\mathscr A}^{\rm out}=0
 \iff
 \mathscr C_{\mathscr A}^{\rm out}=0.
 \label{eq:v20-exterior-rank}
\end{equation}
Thus its rank is the exact number of independent exterior source directions
visible to the complete symmetric parent-pair bank.
\end{theorem}

\begin{proof}
Equations \eqref{eq:v20-internal-synthesis} and
\eqref{eq:v20-exterior-synthesis} are the two orthogonal Fourier projections
of the same source.  Their sum gives \eqref{eq:v20-source-split}, and
orthogonality gives \eqref{eq:v20-source-action-Pythagoras}.  The Gram
claims are the standard identities
$\Ker(C^*C)=\Ker C$ and $\operatorname{rank}(C^*C)=\operatorname{rank} C$ in finite dimension.
\end{proof}

\begin{corollary}[Finite polarization acquisition]
\label{cor:v20-finite-polarization}
Let $e_1,\ldots,e_N$ be a real basis of $E_{\mathscr A}$ and write
$d(z)=d_{\mathscr A}^{\rm out}(z)$.  The complete exterior synthesis is
reconstructed by
\begin{align}
 \mathscr C_{\mathscr A}^{\rm out}(e_i\odot e_i)
 &=d(e_i),
 \label{eq:v20-polar-diagonal}\\
 2\mathscr C_{\mathscr A}^{\rm out}(e_i\odot e_j)
 &=d(e_i+e_j)-d(e_i)-d(e_j),
 \qquad i<j.
 \label{eq:v20-polar-offdiagonal}
\end{align}
Hence at most $N(N+1)/2$ source-output evaluations, or direct evaluation of
the coefficients \eqref{eq:v20-helical-boundary-coefficient}, populate the
complete daughter Gram.  No shell limit or compactness theorem enters this
finite acquisition.
\end{corollary}

\begin{proof}
Expand the homogeneous quadratic map
$d(z)=\mathscr C_{\mathscr A}^{\rm out}(z\odot z)$ at $e_i+e_j$ and use
symmetry of the parent synthesis.
\end{proof}

\begin{theorem}[Exceptional cancellation cone and source germ]
\label{thm:v20-daughter-cone}
Assume $\mathscr C_{\mathscr A}^{\rm out}\ne0$.  Then
\begin{equation}
 \mathcal Z_{\mathscr A}^{\rm out}
 :=\{z\in E_{\mathscr A}:d_{\mathscr A}^{\rm out}(z)=0\}
 \label{eq:v20-daughter-zero-cone}
\end{equation}
is a proper real algebraic cone with empty interior and Lebesgue measure
zero.  The linearized exterior source germ at $z$ is
\begin{equation}
 \mathscr L_{\mathscr A,z}^{\rm out}h
 :=2\mathscr C_{\mathscr A}^{\rm out}(z\odot h),
 \qquad
 \mathbb K_{\mathscr A,z}^{\rm out}
 :=(\mathscr L_{\mathscr A,z}^{\rm out})^*
   \mathscr L_{\mathscr A,z}^{\rm out}\succeq0.
 \label{eq:v20-tangent-daughter-Gram}
\end{equation}
It satisfies
\begin{equation}
 \boxed{
 \mathscr L_{\mathscr A,z}^{\rm out}z
 =2d_{\mathscr A}^{\rm out}(z).}
 \label{eq:v20-Euler-homogeneity}
\end{equation}
Thus $\mathbb K_{\mathscr A,z}^{\rm out}=0$ implies
$d_{\mathscr A}^{\rm out}(z)=0$, whereas a state-level cancellation
$d_{\mathscr A}^{\rm out}(z)=0$ need not be stable under donor variations.
\end{theorem}

\begin{proof}
Because the symmetric synthesis is nonzero, polarization shows that the
quadratic map $z\mapsto\mathscr C_{\mathscr A}^{\rm out}(z\odot z)$ is not
identically zero.  At least one real coordinate is therefore a nonzero
homogeneous quadratic polynomial.  Its zero set has empty interior and
Lebesgue measure zero, and the common zero set
\eqref{eq:v20-daughter-zero-cone} is contained in it.  Differentiating the
quadratic map gives \eqref{eq:v20-tangent-daughter-Gram}; substituting $h=z$
gives \eqref{eq:v20-Euler-homogeneity}.
\end{proof}

\begin{assessment}[What the finite daughter Gram decides]
\label{ass:v20-daughter-Gram-scope}
The operator identity
$\mathbb D_{\mathscr A}^{\rm out}=0$ is universal closure of the declared
helical bank.  The scalar identity
$d_{\mathscr A}^{\rm out}(z)=0$ is cancellation at one occurring state.  The
tangent identity
$\mathbb K_{\mathscr A,z}^{\rm out}=0$ is first-order closure for the donor
variation class.  A one-scalar NS--A target needs only the occurring source
and its ancestry.  Reuse as a finite NS--B donor bank additionally requires
the tangent row on the declared donor variations.  V20 keeps these targets
separate.
\end{assessment}

% =============================================================================
\section{The complete two-sided source boundary packet}
\label{sec:v20-two-sided-boundary}
% =============================================================================

Return to a smooth full periodic field $u$.  Put
\begin{equation}
 P=P_{\mathscr A},
 \qquad Q=I-P,
 \qquad v=Pu,
 \qquad w=Qu,
 \qquad c=A^{-1/4}\mathcal B(u).
 \label{eq:v20-head-complement}
\end{equation}
The V19 incoming row is
\begin{equation}
 \ell_{\mathscr A}^{\rm in}[u]
 :=A^{-1/4}P\bigl(\mathcal B(u)-\mathcal B(v)\bigr).
 \label{eq:v20-incoming-row}
\end{equation}
Define the complementary outgoing row
\begin{equation}
 \ell_{\mathscr A}^{\rm out}[u]
 :=A^{-1/4}Q\bigl(\mathcal B(u)-\mathcal B(w)\bigr).
 \label{eq:v20-outgoing-row}
\end{equation}
The pure head-head daughter is
\begin{equation}
 d_{\mathscr A}^{\rm out}[u]
 :=A^{-1/4}Q\mathcal B(v).
 \label{eq:v20-pure-daughter}
\end{equation}
Using \eqref{eq:v20-symmetric-bilinear-source},
\begin{equation}
 \ell_{\mathscr A}^{\rm out}[u]
 =d_{\mathscr A}^{\rm out}[u]
 +2A^{-1/4}Q\mathcal B_{\rm s}(v,w).
 \label{eq:v20-outgoing-components}
\end{equation}

\begin{theorem}[Exact two-sided source reduction]
\label{thm:v20-two-sided-source-reduction}
The full source has the two block identities
\begin{align}
 Pc
 &=A^{-1/4}P\mathcal B(v)
   +\ell_{\mathscr A}^{\rm in}[u],
 \label{eq:v20-head-source-block}\\
 Qc
 &=A^{-1/4}Q\mathcal B(w)
   +\ell_{\mathscr A}^{\rm out}[u].
 \label{eq:v20-exterior-source-block}
\end{align}
Equivalently,
\begin{equation}
 \boxed{
 c=A^{-1/4}P\mathcal B(v)
   +A^{-1/4}Q\mathcal B(w)
   +\ell_{\mathscr A}^{\rm in}[u]
   +\ell_{\mathscr A}^{\rm out}[u].}
 \label{eq:v20-complete-source-boundary}
\end{equation}
The decomposition into the head and complement is source reducing at the
occurring state $u$ exactly when
\begin{equation}
 \boxed{
 \ell_{\mathscr A}^{\rm in}[u]=0,
 \qquad
 \ell_{\mathscr A}^{\rm out}[u]=0.}
 \label{eq:v20-reducing-state}
\end{equation}
In particular, vanishing of the V19 incoming row alone does not imply source
reduction.
\end{theorem}

\begin{proof}
Apply $P$ and $Q$ to $c=A^{-1/4}\mathcal B(u)$ and add and subtract the two
internal sources $A^{-1/4}P\mathcal B(v)$ and
$A^{-1/4}Q\mathcal B(w)$.  Since $P,Q$ commute with $A$, this gives
\eqref{eq:v20-head-source-block}--\eqref{eq:v20-complete-source-boundary}.
The two blocks evolve without nonlinear source exchange at the occurring
state precisely when the two remainders vanish.
\end{proof}

\begin{corollary}[V19 invariant contact remains exact]
\label{cor:v20-v19-contact-retained}
For every internal Galerkin invariant coefficient
$q\in\Ker\mathfrak A_{\mathscr A}$,
\begin{equation}
 \boxed{
 \left\langle c,M_qe_{\mathscr A}\right\rangle
 =\left\langle
   \ell_{\mathscr A}^{\rm in}[u],M_qe_{\mathscr A}
  \right\rangle.}
 \label{eq:v20-v19-contact-identity}
\end{equation}
The outgoing row is absent from this scalar because it is Fourier
orthogonal to every head-supported writer.  Its omission from the scalar is
not permission to omit it from source ancestry.
\end{corollary}

\begin{proof}
The head-supported writer annihilates the $Q$ block.  The internal $P$
source is orthogonal to $M_qe_{\mathscr A}$ by
Theorem~\ref{thm:v19-bank-source-null}.  Equation
\eqref{eq:v20-head-source-block} gives the result.
\end{proof}

% =============================================================================
\section{One signed stock current does not detect daughter birth}
\label{sec:v20-stock-current}
% =============================================================================

Assume $u$ solves
\begin{equation}
 \partial_tu+\nu Au+\mathcal B(u)=0
 \label{eq:v20-NSE}
\end{equation}
on a smooth interval, and retain the Fourier projection $P$ of
\eqref{eq:v20-head-complement}.  Since $P$ commutes with $A$ and curl, put
\begin{align}
 E_P(t)&=\frac12\|Pu(t)\|_2^2,
 \label{eq:v20-head-energy}\\
 H_P(t)&=\frac12\langle Pu(t),\operatorname{curl}Pu(t)\rangle.
 \label{eq:v20-head-helicity}
\end{align}

\begin{theorem}[Exact energy and helicity boundary currents]
\label{thm:v20-energy-helicity-currents}
Define
\begin{align}
 J_P^E(t)
 &:=\langle\mathcal B(u(t)),Pu(t)\rangle,
 \label{eq:v20-energy-current}\\
 J_P^H(t)
 &:=\langle\mathcal B(u(t)),\operatorname{curl}Pu(t)\rangle.
 \label{eq:v20-helicity-current}
\end{align}
Then
\begin{align}
 \frac{\dd}{\dd t}E_P(t)
 +\nu\|A^{1/2}Pu(t)\|_2^2
 &=-J_P^E(t),
 \label{eq:v20-head-energy-balance}\\
 \frac{\dd}{\dd t}H_P(t)
 +\nu\langle APu(t),\operatorname{curl}Pu(t)\rangle
 &=-J_P^H(t).
 \label{eq:v20-head-helicity-balance}
\end{align}
The nonlinear cancellations on the complete field give the one-count
identities
\begin{equation}
 \boxed{
 J_P^E=-\langle\mathcal B(u),Qu\rangle,
 \qquad
 J_P^H=-\langle\mathcal B(u),\operatorname{curl}Qu\rangle.}
 \label{eq:v20-opposite-currents}
\end{equation}
Moreover,
\begin{equation}
 \boxed{
 J_P^E
 =\left\langle\ell_{\mathscr A}^{\rm in},A^{1/4}Pu\right\rangle,
 \qquad
 J_P^H
 =\left\langle\ell_{\mathscr A}^{\rm in},
                   A^{1/4}\operatorname{curl}Pu\right\rangle.}
 \label{eq:v20-current-incoming}
\end{equation}
Thus the instantaneous energy and helicity stock currents depend on the
incoming row, not on the pure daughter source
$d_{\mathscr A}^{\rm out}$.
\end{theorem}

\begin{proof}
Take the $L^2$ pairing of the projected equation with $Pu$ to obtain
\eqref{eq:v20-head-energy-balance}.  Pair instead with
$\operatorname{curl}Pu$ and use self-adjointness of curl on periodic
solenoidal fields to obtain \eqref{eq:v20-head-helicity-balance}.  The exact
Euler cancellations
\[
 \langle\mathcal B(u),u\rangle=0,
 \qquad
 \langle\mathcal B(u),\operatorname{curl}u\rangle=0
\]
give \eqref{eq:v20-opposite-currents}.  Finally,
$\langle\mathcal B(Pu),Pu\rangle=0$ and
$\langle\mathcal B(Pu),\operatorname{curl}Pu\rangle=0$.  Subtract these
internal cancellations and use \eqref{eq:v20-incoming-row} to obtain
\eqref{eq:v20-current-incoming}.
\end{proof}

\begin{countertheorem}[Zero incoming current with nonzero source birth]
\label{cth:v20-current-not-birth}
Let $K$ be a finite symmetric noncollinear wavevector set and choose the
field $v\in\mathscr H_K$ constructed in
Lemma~\ref{lem:v20-pair-daughter}.  Let $u$ be the smooth local
Navier--Stokes solution with initial datum $u(0)=v$.  At time zero,
\begin{equation}
 \boxed{
 \ell_K^{\rm in}[u(0)]=0,
 \qquad
 J_K^E(0)=J_K^H(0)=0,
 \qquad
 d_K^{\rm out}[u(0)]\ne0.}
 \label{eq:v20-zero-current-nonzero-birth}
\end{equation}
Thus neither a zero V19 incoming row nor zero energy/helicity boundary
current certifies source closure.
\end{countertheorem}

\begin{proof}
At $t=0$ the full field equals its head projection, so
$\mathcal B(u)-\mathcal B(P_Ku)=0$ and the incoming row vanishes.  Equation
\eqref{eq:v20-current-incoming} then gives both currents equal to zero.  The
exterior daughter is nonzero by
Lemma~\ref{lem:v20-pair-daughter}.  Smooth local existence for the smooth
periodic datum places the calculation on an actual Navier--Stokes
trajectory.
\end{proof}

\begin{assessment}[Source first birth precedes stock transfer]
\label{ass:v20-birth-before-stock}
In Countertheorem~\ref{cth:v20-current-not-birth}, the exterior mode has zero
amplitude at the initial time, so its quadratic energy begins only at higher
order even though its time derivative in amplitude is nonzero.  This is why
the source birth is invisible to the instantaneous energy current.  A
physical stock theorem may charge the later transfer, but it cannot be used
to erase the earlier source occurrence or assign it to an unrelated parent.
\end{assessment}

% =============================================================================
\section{Source-first short of the exterior daughter}
\label{sec:v20-daughter-short}
% =============================================================================

Let $\mathcal J_{\mathscr A,t}^{\rm out,rep}$ be a physically declared
closed subspace of the exterior source carrier at time $t$.  It contains
only sources whose first birth and common-history transport have already
been entered in the represented atlas.  Let
$R_{\mathscr A,t}^{\rm out}$ be its orthogonal projection.  Define
\begin{align}
 d_{\mathscr A,t}^{\rm rep}
 &:=R_{\mathscr A,t}^{\rm out}d_{\mathscr A}^{\rm out}[u(t)],
 \label{eq:v20-represented-daughter}\\
 d_{\mathscr A,t}^{\rm birth}
 &:=(I-R_{\mathscr A,t}^{\rm out})d_{\mathscr A}^{\rm out}[u(t)].
 \label{eq:v20-new-daughter}
\end{align}

\begin{theorem}[Exact daughter first-birth Pythagoras]
\label{thm:v20-daughter-Pythagoras}
For almost every smooth time,
\begin{equation}
 \boxed{
 \|d_{\mathscr A}^{\rm out}[u(t)]\|_2^2
 =\|d_{\mathscr A,t}^{\rm rep}\|_2^2
  +\|d_{\mathscr A,t}^{\rm birth}\|_2^2.}
 \label{eq:v20-daughter-Pythagoras}
\end{equation}
For every nonnegative time weight $\vartheta$,
\begin{equation}
 \boxed{
 \int\vartheta\|d_{\mathscr A}^{\rm out}\|_2^2
 =\int\vartheta\|d_{\mathscr A}^{\rm rep}\|_2^2
  +\int\vartheta\|d_{\mathscr A}^{\rm birth}\|_2^2.}
 \label{eq:v20-time-daughter-Pythagoras}
\end{equation}
The represented component receives no new source price.  The orthogonal
component has source-first-birth rank
\begin{equation}
 \dim\operatorname{span}
 \{d_{\mathscr A,t}^{\rm birth}:t\in I\},
 \label{eq:v20-daughter-birth-rank}
\end{equation}
which is the minimum additional exterior source dimension required by this
history.
\end{theorem}

\begin{proof}
Equations \eqref{eq:v20-represented-daughter} and
\eqref{eq:v20-new-daughter} are the orthogonal projection decomposition.
Pythagoras gives the pointwise identity; integration gives the weighted
identity.  The span of the orthogonal residual is, by the universal property
of orthogonal projection, the smallest subspace which must be added to the
represented atlas to contain every daughter source on the interval.
\end{proof}

\begin{firewall}[A represented daughter is not yet a represented return]
\label{fire:v20-daughter-return}
The condition $d^{\rm birth}=0$ identifies the exterior source with the
existing source atlas.  It does not by itself prove that a later incoming
outer-parent source is the response of that daughter.  That stronger
statement requires the already declared common-endpoint source-to-return
map.  If the map is absent, its incidence residual is retained.  If it is
present, the later return is charged to the original daughter first birth
and not again at the head boundary.
\end{firewall}

\begin{corollary}[Tangent daughter short for a reusable donor bank]
\label{cor:v20-tangent-daughter-short}
Let $E_{\mathscr A}^{\rm don}\subseteq E_{\mathscr A}$ be the predeclared
finite donor-variation space.  The positive tangent birth Gram is
\begin{equation}
 \boxed{
 \mathbb K_{\mathscr A,z}^{\rm birth}
 :=P_{E^{\rm don}}
  (\mathscr L_{\mathscr A,z}^{\rm out})^*
  (I-R_{\mathscr A}^{\rm out})
  \mathscr L_{\mathscr A,z}^{\rm out}
  P_{E^{\rm don}}
 \succeq0.}
 \label{eq:v20-tangent-birth-Gram}
\end{equation}
It vanishes exactly when every first-order exterior source generated by the
declared donor variations is represented.  Its rank is the minimum new
exterior source dimension visible to those variations.
\end{corollary}

\begin{proof}
The displayed matrix is the Gram of the synthesis
$(I-R_{\mathscr A}^{\rm out})
 \mathscr L_{\mathscr A,z}^{\rm out}P_{E^{\rm don}}$.  Its kernel, rank, and
vanishing statements follow from the ordinary Gram identities.
\end{proof}

% =============================================================================
\section{Insertion into the V19 terminal card}
\label{sec:v20-terminal-insertion}
% =============================================================================

Retain the candidate V19 terminal history, its predeclared finite helical
head $\mathscr A_n$, projection $\Pi_n$, record amplitude $\mathcal R_n$,
and time weight $\eta_n$.  Put
\begin{equation}
 v_n(t)=\Pi_nu(t),
 \qquad
 Q_n=I-\Pi_n,
 \label{eq:v20-terminal-head-state}
\end{equation}
and define the occurring exterior daughter source
\begin{equation}
 \boxed{
 d_{n,t}^{\rm out}
 :=A^{-1/4}Q_n\mathcal B(v_n(t)).}
 \label{eq:v20-terminal-daughter}
\end{equation}
Let $R_{n,t}^{\rm out}$ be the projection onto the physically declared
exterior source/return atlas transported to the same time and source metric.
Define
\begin{equation}
 d_{n,t}^{\rm birth}
 :=(I-R_{n,t}^{\rm out})d_{n,t}^{\rm out}
 \label{eq:v20-terminal-new-birth}
\end{equation}
and the scale-normalized daughter first-birth action
\begin{equation}
 \boxed{
 \mathfrak B_n^{\rm dau}
 :=\frac1{\nu\mathcal R_n^2}
   \int_{\supp\eta_n}
   \|d_{n,t}^{\rm birth}\|_2^2\,\dd t.}
 \label{eq:v20-normalized-daughter-action}
\end{equation}
The normalization is the same source-action scale used by the V10--V18
terminal card.  It does not turn the daughter action into a scalar contact or
a stock decrement.

If the NS--B target reuses a finite donor-variation space
$E_n^{\rm don}$, also define the integrated tangent birth Gram
\begin{equation}
 \boxed{
 \mathbb T_n^{\rm dau}
 :=\int_{\supp\eta_n}
 P_{E_n^{\rm don}}
 (\mathscr L_{\mathscr A_n,z_n(t)}^{\rm out})^*
 (I-R_{n,t}^{\rm out})
 \mathscr L_{\mathscr A_n,z_n(t)}^{\rm out}
 P_{E_n^{\rm don}}\,\dd t.}
 \label{eq:v20-terminal-tangent-Gram}
\end{equation}
This row is required only when the donor bank, rather than the single
occurring scalar, is intended to be reusable.

\begin{theorem}[Revised finite-head export gate]
\label{thm:v20-revised-export}
Apply the V19 priority theorem first.  On a candidate branch satisfying the
V19 affine-rigidity, capture, incoming-leakage, and additional-invariant
conditions, pass to a further priority subsequence.  Exactly one of the
following occurs.
\begin{enumerate}[label=\textup{(G20.\arabic*)},leftmargin=5.4em]
\item \emph{Daughter first-birth survivor.}  For some $\delta>0$,
      \begin{equation}
       \mathfrak B_n^{\rm dau}\ge\delta
       \label{eq:v20-daughter-survivor}
      \end{equation}
      cofinally.  The finite head emits a record-scale source component not
      generated by the represented exterior atlas.
\item \emph{Donor-germ boundary survivor.}  The occurring daughter action is
      negligible for the selected scalar, but on the declared reusable donor
      variations
      \begin{equation}
       \|\mathbb T_n^{\rm dau}\|_{\rm op}\not\longrightarrow0.
       \label{eq:v20-tangent-survivor}
      \end{equation}
      The selected state is a boundary-cancellation point, while the donor
      germ still creates unrepresented exterior sources.
\item \emph{Ancestry-incidence survivor.}  The unrepresented daughter action
      and tangent row vanish, but the physically declared map from the
      represented daughter atlas to the later incoming/return dictionary has
      a nonzero common-history residual.  The exterior source is represented
      as a vector but not as the claimed return ancestry.
\item \emph{Ancestry-complete open donor head.}  One has
      \begin{equation}
       \boxed{
       \mathfrak B_n^{\rm dau}\to0,
       \qquad
       \|\mathbb T_n^{\rm dau}\|_{\rm op}\to0
       \quad\text{when the reusable-bank target is active},}
       \label{eq:v20-open-head-pass}
      \end{equation}
      and the represented daughter-to-return incidence passes.  Then the
      finite head need not be support closed.  It is source-ancestry complete
      for the predeclared scalar and donor variation class, and the V18--V19
      card is eligible for the NS--B paid-stress and Reynolds audit.
\item \emph{All-shell daughter escape.}  No exhausting finite-head sequence
      simultaneously captures the selected scalar, controls the V19 incoming
      leakage, and makes the unrepresented daughter action negligible.  The
      output is one fused all-shell source-birth survivor.
\end{enumerate}
Branch~\textup{(G20.4)} replaces the phrase ``closed donor head'' by the
strictly weaker and physically correct phrase ``ancestry-complete open donor
head.''
\end{theorem}

\begin{proof}
The V19 gate is applied before the new row, so no daughter estimate can
repair a capture, incoming-leakage, or invariant-incidence failure.  On its
candidate passing branch, the nonnegative sequence
$\mathfrak B_n^{\rm dau}$ either has a positive lower bound on a subsequence
or tends to zero.  This gives the first split.  When the reusable donor-bank
target is active, the nonnegative operator norms of
\eqref{eq:v20-terminal-tangent-Gram} give the second split.  On the remaining
branch, the represented daughter vector must still be compared through the
actual common-history response map; failure is the third branch and passage
is the fourth.  If no finite sequence reaches these tests while preserving
the earlier V19 rows, the fifth branch is exactly the stated all-shell
escape.  Orthogonal Pythagoras in
Theorem~\ref{thm:v20-daughter-Pythagoras} guarantees that no represented and
new daughter action is counted twice.
\end{proof}

\begin{corollary}[Corrected NS--A to NS--B export packet]
\label{cor:v20-export-packet}
On branch~\textup{(G20.4)}, append to the V19/V18 export card the finite
boundary packet
\begin{equation}
 \boxed{
 \mathfrak D_n^{V20}
 =\bigl(
   \mathscr C_{\mathscr A_n}^{\rm out},
   \mathbb D_{\mathscr A_n}^{\rm out},
   d_{n}^{\rm rep},d_n^{\rm birth},
   \mathfrak B_n^{\rm dau},
   \mathbb T_n^{\rm dau},
   \mathcal R_n^{\rm dau\to ret}
  \bigr),}
 \label{eq:v20-export-packet}
\end{equation}
where $\mathcal R_n^{\rm dau\to ret}$ is the declared daughter-to-return
incidence residual.  The resulting card contains the internal invariant
kernel, incoming return, exterior first birth, and represented-return
ancestry on one fixed datum.  NS--A stops on that subsequence.
\end{corollary}

\begin{proof}
All entries are formed from the same finite head, occurring source history,
and exterior atlas.  The previous V19 packet supplies the internal invariant
and incoming rows; V20 supplies the exterior first-birth and tangent rows.
Branch~\textup{(G20.4)} makes every unrepresented boundary component
negligible at the declared target scale, leaving the existing NS--B
paid-stress and Reynolds estimates as the next rows.
\end{proof}

\begin{firewall}[Open does not mean unpaid]
\label{fire:v20-open-not-free}
An ancestry-complete open head may emit a nonzero represented daughter
source and may carry a nonzero signed energy or helicity boundary current.
``Open'' means these channels remain explicit.  ``Ancestry complete'' means
they are assigned to their first source occurrence and their later return is
not priced again.  No finite stock bound, sign, or smallness follows from
ancestry alone.
\end{firewall}

% =============================================================================
\section{Integrated V20 theorem and exact stopping line}
\label{sec:v20-integrated}
% =============================================================================

\begin{theorem}[Internal-to-exterior daughter theorem for the NS--A terminal card]
\label{thm:v20-integrated}
Preserving Versions~V1--V19, Version~V20 establishes the following.
\begin{enumerate}[label=\textup{(T20.\arabic*)},leftmargin=5.6em]
\item For every finite symmetric noncollinear wavevector set, two retained
      divergence-free modes can be polarized so that their quadratic
      interaction produces a nonzero exterior mode.
\item A finite polarization-complete Fourier head is universally
      self-closed under the Euler quadratic source if and only if it is
      collinear; on the collinear branch its nonlinearity vanishes.
\item Every finite helical bank has a complete internal/exterior
      parent-output synthesis and the positive exterior Gram
      \eqref{eq:v20-exterior-Gram}.  Its rank is the exact daughter-source
      output rank of the symmetric parent bank.
\item The actual daughter-null states form a proper algebraic cone whenever
      the exterior synthesis is nonzero.  The tangent Gram
      \eqref{eq:v20-tangent-daughter-Gram} distinguishes one-state
      cancellation from donor-germ closure.
\item The full source decomposes exactly as
      \eqref{eq:v20-complete-source-boundary}.  The V19 incoming leakage
      remains the exact correction to every head-supported invariant
      contact, but it is only one side of source reduction.
\item Energy and helicity cross the head boundary through one signed current
      each.  Countertheorem~\ref{cth:v20-current-not-birth} proves that both
      currents and the incoming leakage may vanish while a nonzero exterior
      source is born.
\item The exterior daughter has the exact represented/new-source
      Pythagoras \eqref{eq:v20-time-daughter-Pythagoras}.  Only the orthogonal
      residual enlarges the source atlas, and represented return still
      requires its common-history map.
\item The V19 affine-rigid branch is replaced by the revised alternatives of
      Theorem~\ref{thm:v20-revised-export}: daughter birth, tangent source
      germ, ancestry-incidence residual, all-shell source escape, or an
      ancestry-complete open donor head.
\item On the final branch, the enlarged packet
      \eqref{eq:v20-export-packet} reaches the intended NS--A/NS--B interface
      and NS--A terminates on that subsequence.
\end{enumerate}
No finite-stock upper budget, critical acceleration estimate, terminal Dini
improvement, terminal trace, backward uniqueness, Liouville theorem, or
global regularity conclusion is supplied.
\end{theorem}

\begin{proof}
Items~\textup{(T20.1)--(T20.2)} are
\Cref{lem:v20-additive-escape,lem:v20-pair-daughter,thm:v20-no-finite-reducing-head}.
Items~\textup{(T20.3)--(T20.4)} are
\Cref{thm:v20-source-Pythagoras,cor:v20-finite-polarization,thm:v20-daughter-cone}.
Item~\textup{(T20.5)} is
\Cref{thm:v20-two-sided-source-reduction,cor:v20-v19-contact-retained}.
Item~\textup{(T20.6)} is
\Cref{thm:v20-energy-helicity-currents,cth:v20-current-not-birth}.
Item~\textup{(T20.7)} is
\Cref{thm:v20-daughter-Pythagoras,cor:v20-tangent-daughter-short}.
The last two items are
Theorem~\ref{thm:v20-revised-export} and
Corollary~\ref{cor:v20-export-packet}.
\end{proof}

\begin{firewall}[Exact stopping boundary after V20]
\label{fire:v20-stop}
Version~V20 closes the interpretation of a finite rank-three donor head as
a universally reducing Fourier support.  Such a support is unavailable on a
polarization-complete finite head.  The admissible finite object is an open
source card with an internal invariant kernel, an exterior daughter germ,
an incoming return row, and a nonduplicating ancestry short.  A later
positive Reynolds estimate or paid-stress bound may be applied only after
these rows occur on the same history; it cannot manufacture their
incidence.
\end{firewall}

\begingroup\small
\begin{longtable}{p{0.25\textwidth}p{0.19\textwidth}p{0.48\textwidth}}
\caption{NS--A Version V20 proof-status ledger.}\label{tab:v20-ledger}\\
\toprule Row & Status & Exact conclusion\\\midrule
\endfirsthead
\toprule Row & Status & Exact conclusion\\\midrule
\endhead
V1--V19 prefix & \PROVED{} preserved
 &The complete V19 preterminal source is retained byte-for-byte; no earlier
 theorem or limitation is rewritten.\\[0.3em]
V19 diagonal invariant audit & \INHERITED{} exact
 &$\Ker\mathfrak A_{\mathscr A}$ remains the complete helical-diagonal
 invariant space of the internal Galerkin source.\\[0.3em]
Nonparallel pair daughter & \PROVED{} explicit
 &Equation \eqref{eq:v20-explicit-daughter} gives a nonzero exterior output
 for concrete divergence-free polarizations.\\[0.3em]
Finite universal self-closure & \OBSTRUCTION{} classified
 &A polarization-complete finite head is self-closed exactly on the collinear,
 nonlinear-null branch.\\[0.3em]
Exterior parent-output matrix & \PROVED{} finite
 &$\mathscr C_{\mathscr A}^{\rm out}$ and
 $\mathbb D_{\mathscr A}^{\rm out}$ are reconstructed by the literal
 parent-output coefficients or finite polarization.\\[0.3em]
Daughter-null state & \PROVED{} exceptional
 &For a nonzero exterior synthesis the daughter-null states form a proper
 algebraic cone.\\[0.3em]
Tangent daughter germ & \PROVED{} positive
 &$\mathbb K_{\mathscr A,z}^{\rm out}$ is the exact first-order exterior
 source Gram for the declared donor variations.\\[0.3em]
Incoming outer-parent row & \INHERITED{} / retained
 &It is the exact correction to a head-supported invariant contact, but not
 the complete source boundary.\\[0.3em]
Two-sided source reduction & \PROVED{} exact
 &Equations \eqref{eq:v20-head-source-block}--
 \eqref{eq:v20-complete-source-boundary} have no unnamed source remainder.\\[0.3em]
Energy/helicity boundary current & \PROVED{} one-count
 &The currents are opposite on the two blocks and depend on the incoming row.\\[0.3em]
Zero current implies no source birth & \OBSTRUCTION{} explicit
 &Countertheorem~\ref{cth:v20-current-not-birth} has zero incoming and stock
 currents but nonzero exterior daughter source.\\[0.3em]
Daughter first-birth short & \PROVED{} exact
 &Equation \eqref{eq:v20-time-daughter-Pythagoras} separates represented
 daughter output from the minimum new exterior source.\\[0.3em]
Daughter-to-return ancestry & \OPEN{} physical incidence
 &A represented exterior vector becomes a represented return only through the
 declared common-history map.\\[0.3em]
Affine-rigid closed head & \OBSTRUCTION{} terminology corrected
 &The valid finite PASS is an ancestry-complete open head, not a universal
 reducing support.\\[0.3em]
Ancestry-complete open export & \CONDITIONAL{} exact
 &When \eqref{eq:v20-open-head-pass} and the return-incidence row pass, the
 packet \eqref{eq:v20-export-packet} is eligible for NS--B.\\[0.3em]
All-shell daughter escape & \OPEN{} typed alternative
 &Failure of every finite ancestry-complete head is one fused source-birth
 escape, not a shell-count sum.\\[0.3em]
Three-dimensional regularity & \OPEN
 &No daughter, return, Reynolds, paid-stress, stock, trace, or rigidity branch
 is excluded for every datum.\\
\bottomrule
\end{longtable}
\endgroup

% =============================================================================
\section{Exact mandate after Version V20}
\label{sec:v20-next}
% =============================================================================

\begin{programme}[Populate the open boundary card or stop NS--A]
\label{prog:v20-next}
Preserve Versions~V1--V20 verbatim.  There is no automatic generic
Version~V21.

A continuation is licensed only in the following order.
\begin{enumerate}[label=\textup{(V21.\arabic*)},leftmargin=5.9em]
\item On the literal V19 donor head, populate the exterior parent-output
      synthesis $\mathscr C_{\mathscr A_n}^{\rm out}$ and its occurring
      daughter source \eqref{eq:v20-terminal-daughter}.  This is a finite
      coefficient and state read, not another invariant search.
\item Short that daughter against the already declared exterior source and
      return atlas.  Do not use energy or helicity current as a surrogate for
      source first birth.
\item If the unrepresented daughter action survives, export the resulting
      source-first-birth packet and follow only its actual common-endpoint
      return.  Do not price the later incoming row again.
\item If the occurring daughter is represented but the donor bank is to be
      reused under variations, evaluate the tangent birth Gram
      \eqref{eq:v20-terminal-tangent-Gram} on the declared variation space.
\item On a daughter-tight branch, evaluate the physical daughter-to-return
      incidence.  A failed map is the exact ancestry residual; do not repair
      it by enlarging the head after the result is inspected.
\item When the V19 internal rows and the V20 boundary rows all pass, stop
      NS--A and export the ancestry-complete open card to NS--B.  The next
      question is the paid-stress/Reynolds upper budget on that same card.
\item If no finite head reaches this pass, retain all-shell daughter escape.
      Continue only when a new source-relative square-function, capacity, or
      critical acceleration upper stock controls this exact fused escape.
\item Do not add another path law, Dini modulus, compactness space, shell
      count, guessed invariant, or generic response compiler.
\end{enumerate}
\end{programme}

\begin{assessment}[Programme position after V20]
\label{ass:v20-position}
The finite invariant question and the finite source-boundary question are
now separated.  V19 decides every helical-diagonal invariant of the internal
head.  V20 proves that a finite rank-three physical head is necessarily open
at the source level unless a restricted-polarization coefficient cancellation
occurs, and it gives the exact finite matrix deciding that cancellation.
The remaining unpopulated object is a literal exterior daughter/ancestry
packet on the selected terminal history.  Once it passes, NS--A has reached
its intended interface and further donor estimates belong in NS--B.
\end{assessment}

\paragraph{Version V20 history.}
Preserved the complete V19 source.  Moved upstream of its closed-head phrase
to the internal-to-exterior parent-output germ.  Proved the finite additive
escape lemma and constructed explicit divergence-free polarizations whose
interaction writes a nonzero exterior mode.  Classified finite
polarization-complete self-closed Fourier heads as exactly the collinear,
nonlinear-null heads.  Constructed the complete internal/exterior source
synthesis, its positive daughter Gram, finite polarization acquisition,
algebraic cancellation cone, and tangent source-germ Gram.  Added the exact
two-sided full-source decomposition while retaining the V19 invariant-contact
identity.  Proved the one-count energy/helicity boundary-current laws and an
actual smooth-trajectory counterexample with zero incoming current but
nonzero source birth.  Performed the source-first daughter short against a
represented exterior atlas and inserted its occurring and tangent rows into
the terminal gate.  Replaced universal finite support closure by the
ancestry-complete open donor export.

\begin{assessment}[Append-only integrity]
\label{ass:v20-integrity}
The complete Version~V19 source preceding its sole final executable document
terminator is preserved verbatim.  Its preterminal SHA--256 digest is
\begin{center}\scriptsize\ttfamily
2fd810e449af7f3815bd7d5a13fe4cbdd4db1f7fc70415bb5fd0af9da3f73517
\end{center}
The present material begins at Section~\ref{sec:v20-scope}; the final command
below is again unique.
\end{assessment}

\addtocontents{toc}{\protect\endgroup}
% =============================================================================
% END VERSION V20 MATERIAL -- append later versions before final terminator
% =============================================================================

% APPEND ALL LATER VERSIONS IMMEDIATELY ABOVE THE SOLE FINAL LINE BELOW.

\clearpage
% =============================================================================
% VERSION V21: NEW MATERIAL.  The complete V1--V20 source above is preserved
% verbatim; only its sole active \end{document} line was removed before append.
% =============================================================================
\hypersetup{pdftitle={NS-A Terminal Source Concentration Programme: Cumulative V21},pdfauthor={A. Pelissier}}
\makeatletter
\addtocontents{toc}{\protect\begingroup\protect\makeatletter}
\addtocontents{toc}{\protect\def\protect\l@section{\protect\@dottedtocline{1}{0em}{4.7em}}}
\addtocontents{toc}{\protect\def\protect\l@subsection{\protect\@dottedtocline{2}{1.5em}{5.3em}}}
\addtocontents{toc}{\protect\makeatother}
\makeatother
\renewcommand{\thesection}{V21.\arabic{section}}
\renewcommand{\thesubsection}{V21.\arabic{section}.\arabic{subsection}}
\renewcommand{\theequation}{V21.\arabic{equation}}
\setcounter{section}{0}
\setcounter{equation}{0}
\renewcommand{\theHsection}{V21.\arabic{section}}
\renewcommand{\theHsubsection}{V21.\arabic{section}.\arabic{subsection}}
\renewcommand{\theHtheorem}{V21.\arabic{section}.\arabic{theorem}}
\renewcommand{\theHequation}{V21.\arabic{equation}}

\begin{center}
{\LARGE\bfseries NS--A: Version V21}\\[0.5em]
{\large The Source-Native Exterior Response, the Cubic First-Return Germ,\\
and the Exact Daughter-to-Return Volterra Card}\\[0.5em]
{\normalsize 3 September 2026}
\end{center}
\addcontentsline{toc}{part}{Version V21: source-to-state response and daughter-to-return incidence}

\begin{abstract}
Version~V20 constructed the exterior daughter source of a finite Fourier
head and shorted that source against the represented exterior atlas.  Its
last open row was the common-history map from a represented daughter to the
later incoming return.  The first upstream correction is that source
representation and return representation are separated by an actual
source-to-state propagator.  A daughter residual which is small in source
action can still produce a record-sized return if that propagator amplifies
it.  Vanishing of the V20 birth action is therefore not by itself a
response-tight pass.

For one smooth Navier--Stokes history and one fixed Fourier split
$P+Q=I$, Version~V21 gives an exact factorization.  If
$v=Pu$, $w=Qu$, $y=A^{-1/4}w$, and
$d=A^{-1/4}Q\mathcal B(v)$ is the pure head--head daughter, then the exterior
state obeys a linear nonautonomous equation on the actual trajectory,
\[
 \partial_ty+\nu Ay+\mathscr L_Q^u(t)y=-d.
\]
The incoming source is another linear read of the same state,
$\ell^{\rm in}=\mathscr M_P^u(t)y$.  Its evolution family gives the literal
causal map
\[
 (\mathscr V_{P\leftarrow Q}^ud)(t)
 =-\int_a^t\mathscr M_P^u(t)\mathscr U_Q^u(t,s)d(s)\,\dd s.
\]
Thus the incoming row is exactly old exterior memory plus the return of the
occurring daughter.  The coefficients are carried by the same solution;
they are transport data and are not repriced as source births.

At a clean daughter birth, $Qu(a)=0$, the V20 source action has an exact
finite-stock interpretation:
\[
 \left.\frac{\dd^2}{\dd t^2}\frac12
 \|A^{-1/4}Qu(t)\|_2^2\right|_{t=a}
 =\|A^{-1/4}Q\mathcal B(Pu(a))\|_2^2.
\]
The first return is the cubic germ
\[
 \mathfrak R_P(v)=2P\mathcal B_{\rm s}(v,Q\mathcal B(v)),
\]
which satisfies
$\langle\mathfrak R_P(v),v\rangle=-\|Q\mathcal B(v)\|_2^2$.
Consequently every nonzero clean daughter has a nonzero immediate return
germ, although its energy and helicity currents vanish at the birth time.
The full head differs from the head-only Galerkin trajectory by
$\frac12h^2\mathfrak R_P(v)+O(h^3)$.

On the terminal card, Version~V21 transports the V20 represented and
unrepresented daughter components through the same causal operator before
shorting the output against the declared return atlas.  It constructs the
positive return Gram, the represented-return residual, and the
birth-to-return amplification action.  A finite-screen Volterra estimate
shows that a surviving return from a vanishing daughter forces divergence
of the dimensionless causal return gain.  The corrected passing branch
therefore requires source tightness, response tightness, and return-atlas
incidence.  No terminal Dini improvement, critical acceleration upper
budget, finite-stock contradiction, backward-uniqueness theorem, Liouville
theorem, or global regularity conclusion is claimed.
\end{abstract}

% =============================================================================
\section{The missing arrow is source to state before state to return}
\label{sec:v21-scope}
% =============================================================================

The V20 daughter short answers whether
\[
 d_P^{\rm out}(t)=A^{-1/4}Q\mathcal B(Pu(t))
\]
is already represented in the exterior source atlas.  It does not yet
answer what state that source generates or whether the generated state
writes the V19 incoming row.  The correct order is
\begin{equation}
 \boxed{
 \begin{gathered}
 \text{head--head daughter source}\\[-0.15em]
 \Downarrow\quad\text{actual exterior source-to-state propagation}\\[-0.15em]
 \text{exterior response state}\\[-0.15em]
 \Downarrow\quad\text{actual state-to-incoming read}\\[-0.15em]
 \text{head return source}\\[-0.15em]
 \Downarrow\quad\text{return-atlas short}\\[-0.15em]
 \text{represented follower or ancestry residual}.
 \end{gathered}}
 \label{eq:v21-upstream-order}
\end{equation}
A source residual which vanishes before the first arrow need not vanish
after the second unless the causal gain is controlled.

\begin{firewall}[Actual-trajectory transport is not a second source]
\label{fire:v21-transport-not-source}
The propagator below is evaluated from the already occurring smooth
trajectory and the predeclared Fourier split.  Its coefficients may contain
the complete exterior state.  They are not new forcing terms and are not
assigned a second source price.  The construction is a linear ancestry
factorization on one fixed history; it does not assert that deleting the
daughter would leave the coefficients unchanged.
\end{firewall}

Version~V21 proves the following chain.
\begin{enumerate}[label=\textup{(U21.\arabic*)},leftmargin=5.3em]
\item The exact exterior equation admits a radial-homogeneous linear
      factorization on the actual trajectory.
\item The V19 incoming source is a linear read of the same exterior state,
      giving one causal daughter-to-return Volterra operator.
\item At a clean birth the V20 daughter action is the curvature of a bounded
      source-native $H^{-1/2}$ exterior stock.
\item The first incoming return is a cubic head coefficient whose kinetic
      pairing is minus the daughter norm squared.
\item A source-first daughter split is propagated before any return-atlas
      positive short is taken.
\item Vanishing source action plus nonvanishing generated return is a
      dynamic-gain survivor, not an ancestry pass.
\item The V20 export is corrected by one finite return Gram and one actual
      return-incidence residual.
\end{enumerate}

% =============================================================================
\section{Radial-homogeneous factorization of the exact block equation}
\label{sec:v21-homogeneous-factorization}
% =============================================================================

Retain the smooth periodic equation
\begin{equation}
 \partial_tu+\nu Au+\mathcal B(u)=0
 \label{eq:v21-NSE}
\end{equation}
on $I=[a,b]\Subset[0,T_*)$.  Let $P$ be a finite orthogonal Fourier
projection commuting with $A$, the Leray projection, and curl.  Put
\begin{equation}
 Q=I-P,
 \qquad v=Pu,
 \qquad w=Qu,
 \qquad y=A^{-1/4}w.
 \label{eq:v21-block-variables}
\end{equation}
Write
\[
 \mathcal B(f,g)=\mathbb P((f\cdot\nabla)g),
 \qquad
 \mathcal B_{\rm s}(f,g)
 =\frac12\bigl(\mathcal B(f,g)+\mathcal B(g,f)\bigr).
\]
Define the pure daughter in the native critical-source metric by
\begin{equation}
 d(t)=A^{-1/4}Q\mathcal B(v(t)).
 \label{eq:v21-native-daughter}
\end{equation}

\begin{definition}[Actual radial exterior and return operators]
\label{def:v21-radial-operators}
For a smooth $Q$-state $h$, define
\begin{align}
 \mathscr L_Q^u(t)h
 &:={A^{-1/4}}Q\left[
       2\mathcal B_{\rm s}\bigl(v(t),A^{1/4}h\bigr)
       +\mathcal B_{\rm s}\bigl(w(t),A^{1/4}h\bigr)
     \right],
 \label{eq:v21-LQ}\\
 \mathscr M_P^u(t)h
 &:={A^{-1/4}}P\left[
       2\mathcal B_{\rm s}\bigl(v(t),A^{1/4}h\bigr)
       +\mathcal B_{\rm s}\bigl(w(t),A^{1/4}h\bigr)
     \right].
 \label{eq:v21-MP}
\end{align}
The coefficient one in the $w$-term is forced by Euler homogeneity:
$\mathcal B(w)=\mathcal B_{\rm s}(w,w)$.
\end{definition}

\begin{theorem}[Exact radial-homogeneous block factorization]
\label{thm:v21-exact-block-factorization}
The actual exterior state and the V20 incoming critical source satisfy
\begin{align}
 \boxed{
 \partial_ty+\nu Ay+\mathscr L_Q^u(t)y=-d,}
 \label{eq:v21-exterior-linear-equation}\\
 \boxed{
 \ell_P^{\rm in}[u](t)=\mathscr M_P^u(t)y(t),}
 \label{eq:v21-incoming-linear-read}
\end{align}
where
\[
 \ell_P^{\rm in}[u]
 =A^{-1/4}P\bigl(\mathcal B(u)-\mathcal B(v)\bigr).
\]
No source remainder is omitted.
\end{theorem}

\begin{proof}
Project \eqref{eq:v21-NSE} by $Q$ and use $u=v+w$:
\[
 \partial_tw+\nu Aw+Q\mathcal B(v)
 +2Q\mathcal B_{\rm s}(v,w)+Q\mathcal B(w)=0.
\]
Apply $A^{-1/4}$.  Since $A^{1/4}y=w$,
\[
 \mathscr L_Q^u(t)y
 =2A^{-1/4}Q\mathcal B_{\rm s}(v,w)
  +A^{-1/4}Q\mathcal B(w),
\]
which gives the first identity.  Similarly,
\[
 \mathscr M_P^u(t)y
 =2A^{-1/4}P\mathcal B_{\rm s}(v,w)
  +A^{-1/4}P\mathcal B(w)
 =\ell_P^{\rm in}[u].
\]
\end{proof}

\begin{lemma}[Evolution family on every smooth Sobolev screen]
\label{lem:v21-evolution-family}
For every integer $m\ge4$, the homogeneous equation
\begin{equation}
 \partial_tz+\nu Az+\mathscr L_Q^u(t)z=0,
 \qquad z(s)=z_s\in H^m_\sigma\cap\operatorname{Ran}Q,
 \label{eq:v21-homogeneous-evolution}
\end{equation}
has a unique solution for $s\le t\le b$.  It defines a strongly continuous
evolution family $\mathscr U_Q^u(t,s)$ with
\begin{equation}
 \mathscr U_Q^u(t,r)\mathscr U_Q^u(r,s)=\mathscr U_Q^u(t,s),
 \qquad
 \mathscr U_Q^u(s,s)=Q.
 \label{eq:v21-evolution-composition}
\end{equation}
Moreover,
\begin{equation}
 \boxed{
 \|\mathscr U_Q^u(t,s)z_s\|_{H^m}^2
 \le
 \exp\!\left(
  C_{m,\nu}\int_s^t
  \bigl(1+\|u(r)\|_{H^{m+2}}^2\bigr)\,\dd r
 \right)
 \|z_s\|_{H^m}^2.}
 \label{eq:v21-evolution-bound}
\end{equation}
\end{lemma}

\begin{proof}
Since $m\ge4$, $H^m(\Torus^3)$ is an algebra and embeds into
$W^{1,\infty}$.  Expanding
$\Lambda^m(a\cdot\nabla A^{1/4}z)$ and using the half-derivative gain of
$A^{-1/4}$ gives
\[
 \bigl|
 \langle\Lambda^m\mathscr L_Q^u(t)z,\Lambda^mz\rangle
 \bigr|
 \le
 C_m\|u(t)\|_{H^{m+2}}
       \|z\|_{H^{m+1}}\|z\|_{H^m}.
\]
The terms with all derivatives on $z$ use
$\|u\|_{W^{1,\infty}}$; every commutator term places a derivative on $u$
and is controlled by the $H^m$ algebra estimate.  Young's inequality yields
\[
 \frac12\frac{\dd}{\dd t}\|z\|_{H^m}^2
 +\frac\nu2\|z\|_{H^{m+1}}^2
 \le
 C_{m,\nu}\bigl(1+\|u(t)\|_{H^{m+2}}^2\bigr)
 \|z\|_{H^m}^2.
\]
Fourier--Galerkin approximation, weak compactness, and Gronwall give
existence and \eqref{eq:v21-evolution-bound}.  The same estimate for the
difference gives uniqueness; composition follows from uniqueness.
\end{proof}

\begin{theorem}[Exact daughter-to-return Volterra map]
\label{thm:v21-return-Volterra}
Define
\begin{align}
 \ell_P^{\rm old}(t)
 &:=\mathscr M_P^u(t)\mathscr U_Q^u(t,a)y(a),
 \label{eq:v21-old-return}\\
 (\mathscr V_{P\leftarrow Q}^uf)(t)
 &:=-\int_a^t
   \mathscr M_P^u(t)\mathscr U_Q^u(t,s)f(s)\,\dd s.
 \label{eq:v21-return-operator}
\end{align}
Then
\begin{equation}
 \boxed{
 \ell_P^{\rm in}[u]
 =\ell_P^{\rm old}+\mathscr V_{P\leftarrow Q}^ud.}
 \label{eq:v21-return-factorization}
\end{equation}
The operator is linear and causal.  If $d=d_1+d_2$, then
\begin{equation}
 \boxed{
 \ell_P^{\rm in}[u]
 =\ell_P^{\rm old}
  +\mathscr V_{P\leftarrow Q}^ud_1
  +\mathscr V_{P\leftarrow Q}^ud_2.}
 \label{eq:v21-return-source-split}
\end{equation}
The three output vectors need not be orthogonal; the signed identity is
assembled before positive action is taken.
\end{theorem}

\begin{proof}
Variation of constants in \eqref{eq:v21-exterior-linear-equation} gives
\[
 y(t)=\mathscr U_Q^u(t,a)y(a)
      -\int_a^t\mathscr U_Q^u(t,s)d(s)\,\dd s.
\]
Apply the exact read $\mathscr M_P^u(t)$ and use
\eqref{eq:v21-incoming-linear-read}.  Causality is the triangular integration
domain; linearity gives \eqref{eq:v21-return-source-split}.
\end{proof}

\begin{assessment}[What the Volterra row does and does not say]
\label{ass:v21-Volterra-scope}
Equation~\eqref{eq:v21-return-factorization} is the literal common-history
source-to-return map missing in V20.  It does not assert that the returned
source lies in a predeclared return dictionary, and it does not assert that
the old exterior term is represented.  Those are range-incidence tests.
\end{assessment}

% =============================================================================
\section{The source-native exterior stock and clean daughter birth}
\label{sec:v21-clean-birth}
% =============================================================================

Define the source-native exterior stock
\begin{equation}
 \mathscr S_Q^{-1/2}(t)
 :=\frac12\|A^{-1/4}w(t)\|_2^2
 =\frac12\|y(t)\|_2^2.
 \label{eq:v21-negative-stock}
\end{equation}
If $\lambda_1>0$ is the first mean-zero eigenvalue of $A$, then
\begin{equation}
 \boxed{
 0\le\mathscr S_Q^{-1/2}(t)
 \le\frac1{2\sqrt{\lambda_1}}\|u(t)\|_2^2.}
 \label{eq:v21-stock-global-bound}
\end{equation}

\begin{theorem}[Exact source-native exterior-stock balance]
\label{thm:v21-negative-stock-balance}
For every smooth time,
\begin{equation}
 \boxed{
 \frac{\dd}{\dd t}\mathscr S_Q^{-1/2}
 +\nu\|A^{1/2}y\|_2^2
 +\langle\mathscr L_Q^u(t)y,y\rangle
 =-\langle d,y\rangle.}
 \label{eq:v21-negative-stock-balance}
\end{equation}
The daughter source is paired with the exterior response state.  Its square
action is not the first derivative of this stock.
\end{theorem}

\begin{proof}
Pair \eqref{eq:v21-exterior-linear-equation} with $y$ and use
$\langle Ay,y\rangle=\|A^{1/2}y\|_2^2$.  The global bound follows from
$\|A^{-1/4}Qf\|_2\le\lambda_1^{-1/4}\|f\|_2$.
\end{proof}

\begin{definition}[Clean daughter birth]
\label{def:v21-clean-birth}
A time $a$ is a clean birth for the split $P+Q=I$ when
\begin{equation}
 Qu(a)=0.
 \label{eq:v21-clean-birth}
\end{equation}
Put
\begin{equation}
 v_0=Pu(a)=u(a),
 \qquad
 D_0=Q\mathcal B(v_0),
 \qquad
 d_0=A^{-1/4}D_0.
 \label{eq:v21-clean-birth-data}
\end{equation}
\end{definition}

\begin{theorem}[Clean-birth curvature in every spectral stock]
\label{thm:v21-clean-birth-curvature}
Assume \eqref{eq:v21-clean-birth}.  For every real $\theta$ for which the
smooth norm is read, set
\[
 \mathscr S_{Q,\theta}(t)=\frac12\|A^\theta Qu(t)\|_2^2.
\]
Then
\begin{equation}
 \boxed{
 \mathscr S_{Q,\theta}'(a)=0,
 \qquad
 \mathscr S_{Q,\theta}''(a)=\|A^\theta D_0\|_2^2.}
 \label{eq:v21-stock-curvature-hierarchy}
\end{equation}
In particular,
\begin{equation}
 \boxed{
 \left.\frac{\dd^2}{\dd t^2}
 \mathscr S_Q^{-1/2}(t)\right|_{t=a}
 =\|d_0\|_2^2.}
 \label{eq:v21-native-birth-curvature}
\end{equation}
\end{theorem}

\begin{proof}
At a clean birth, $w(a)=0$.  The $Q$-projected equation gives
$w_t(a)=-Q\mathcal B(v_0)=-D_0$.  For
$z_\theta=A^\theta w$,
\[
 \frac{\dd}{\dd t}\frac12\|z_\theta\|_2^2
 =\langle z_\theta,(z_\theta)_t\rangle.
\]
The first derivative vanishes at $a$.  A second differentiation gives
$\|(z_\theta)_t(a)\|_2^2=\|A^\theta D_0\|_2^2$.
\end{proof}

\begin{countertheorem}[A bounded stock does not bound integrated birth action]
\label{cth:v21-curvature-no-action-budget}
Already on the one-dimensional linear equation
\begin{equation}
 y_N'+\nu y_N=-d_N
 \label{eq:v21-scalar-model}
\end{equation}
on $[0,2\pi]$, the choice
\begin{equation}
 y_N(t)=N^{-1/2}\sin Nt,
 \qquad
 d_N(t)=-N^{1/2}\cos Nt-\nu N^{-1/2}\sin Nt
 \label{eq:v21-scalar-model-choice}
\end{equation}
satisfies
\begin{equation}
 \boxed{
 \sup_t\frac12|y_N(t)|^2\le\frac1{2N},
 \qquad
 \int_0^{2\pi}|d_N(t)|^2\,\dd t
 =\pi N+\frac{\pi\nu^2}{N}\longrightarrow\infty.}
 \label{eq:v21-stock-action-no-go}
\end{equation}
This is a linear control counterexample, not a Navier--Stokes trajectory.
\end{countertheorem}

\begin{proof}
Substitution gives \eqref{eq:v21-scalar-model}.  Orthogonality of
$\sin Nt$ and $\cos Nt$ gives the displayed action.
\end{proof}

\begin{assessment}[The remaining stock row]
\label{ass:v21-stock-scope}
The negative-order stock proves that the V20 action has a physical
infinitesimal meaning.  It does not make the action a monotone decrement.
An integrated upper budget still requires a one-sided turnover estimate, a
first-exit law, or another independently occurring action-to-stock
incidence.
\end{assessment}

% =============================================================================
\section{The cubic first-return germ}
\label{sec:v21-return-germ}
% =============================================================================

At a clean birth define
\begin{equation}
 \boxed{
 \mathfrak R_P(v_0)
 :=2P\mathcal B_{\rm s}\bigl(v_0,D_0\bigr),
 \qquad D_0=Q\mathcal B(v_0),}
 \label{eq:v21-cubic-return-germ}
\end{equation}
and
\begin{equation}
 \mathfrak r_P(v_0)=A^{-1/4}\mathfrak R_P(v_0).
 \label{eq:v21-critical-return-germ}
\end{equation}

\begin{theorem}[Exact first-return derivative]
\label{thm:v21-first-return-derivative}
At every clean birth,
\begin{equation}
 \boxed{
 \ell_P^{\rm in}[u](a)=0,
 \qquad
 \left.\frac{\dd}{\dd t}\ell_P^{\rm in}[u](t)\right|_{t=a}
 =-\mathfrak r_P(v_0).}
 \label{eq:v21-return-derivative}
\end{equation}
Equivalently, the Volterra kernel
\begin{equation}
 \mathscr K_{P\leftarrow Q}^u(t,s)
 :=-\mathscr M_P^u(t)\mathscr U_Q^u(t,s)
 \label{eq:v21-return-kernel}
\end{equation}
satisfies
\begin{equation}
 \boxed{
 \mathscr K_{P\leftarrow Q}^u(a,a)d_0
 =-\mathfrak r_P(v_0).}
 \label{eq:v21-kernel-diagonal-germ}
\end{equation}
\end{theorem}

\begin{proof}
At a clean birth, $y(a)=0$ and $y_t(a)=-d_0$.  Differentiate
$\ell_P^{\rm in}=\mathscr M_P^uy$:
\[
 (\ell_P^{\rm in})'(a)
 =\mathscr M_P^u(a)y_t(a)
 =-\mathscr M_P^u(a)d_0.
\]
Since $w(a)=0$ and $A^{1/4}d_0=D_0$, the last term is
$-2A^{-1/4}P\mathcal B_{\rm s}(v_0,D_0)$.
The kernel identity uses $\mathscr U_Q^u(a,a)=Q$.
\end{proof}

\begin{theorem}[Energy and helicity identities for the return germ]
\label{thm:v21-return-germ-invariants}
The cubic return germ obeys
\begin{align}
 \boxed{
 \langle\mathfrak R_P(v_0),v_0\rangle
 =-\|D_0\|_2^2,}
 \label{eq:v21-return-energy-identity}\\
 \boxed{
 \langle\mathfrak R_P(v_0),\operatorname{curl}v_0\rangle
 =-\langle D_0,\operatorname{curl}D_0\rangle.}
 \label{eq:v21-return-helicity-identity}
\end{align}
Hence
\begin{equation}
 D_0\ne0\quad\Longrightarrow\quad\mathfrak R_P(v_0)\ne0.
 \label{eq:v21-nonzero-daughter-nonzero-return}
\end{equation}
\end{theorem}

\begin{proof}
Since $v_0\in\operatorname{Ran}P$,
\[
 \langle\mathfrak R_P(v_0),v_0\rangle
 =\langle\mathcal B(v_0,D_0),v_0\rangle
  +\langle\mathcal B(D_0,v_0),v_0\rangle.
\]
The second term vanishes because $D_0$ is divergence free.  Integration by
parts in the first gives
\[
 \langle\mathcal B(v_0,D_0),v_0\rangle
 =-\langle\mathcal B(v_0),D_0\rangle
 =-\langle Q\mathcal B(v_0),D_0\rangle
 =-\|D_0\|_2^2.
\]
For helicity, differentiate at $\varepsilon=0$ the exact identity
\[
 \langle\mathcal B(v_0+\varepsilon D_0),
 \operatorname{curl}(v_0+\varepsilon D_0)\rangle=0.
\]
This gives
\[
 2\langle\mathcal B_{\rm s}(v_0,D_0),
          \operatorname{curl}v_0\rangle
 +\langle\mathcal B(v_0),\operatorname{curl}D_0\rangle=0.
\]
Since $\operatorname{curl}D_0\in\operatorname{Ran}Q$, the last term equals
$\langle D_0,\operatorname{curl}D_0\rangle$.
\end{proof}

\begin{corollary}[The V20 zero-current example turns on at first order]
\label{cor:v21-current-turnon}
Let $J_P^E,J_P^H$ be the V20 energy and helicity boundary currents.  At a
clean birth,
\begin{align}
 \boxed{(J_P^E)'(a)=\|D_0\|_2^2,}
 \label{eq:v21-energy-current-slope}\\
 \boxed{(J_P^H)'(a)=\langle D_0,\operatorname{curl}D_0\rangle.}
 \label{eq:v21-helicity-current-slope}
\end{align}
Thus Countertheorem~\ref{cth:v20-current-not-birth} has zero current at the
birth time but a strictly positive kinetic-energy current slope.
\end{corollary}

\begin{proof}
The internal cancellations give
\[
 J_P^E=\langle A^{1/4}\ell_P^{\rm in},v\rangle,
 \qquad
 J_P^H=\langle A^{1/4}\ell_P^{\rm in},\operatorname{curl}v\rangle.
\]
At a clean birth $\ell_P^{\rm in}(a)=0$.  Differentiate, use
\eqref{eq:v21-return-derivative}, and apply
\eqref{eq:v21-return-energy-identity}--
\eqref{eq:v21-return-helicity-identity}.
\end{proof}

\begin{theorem}[Second-order return to the head-only trajectory]
\label{thm:v21-head-only-secant}
Let $\bar v$ solve the finite Galerkin equation
\begin{equation}
 \partial_t\bar v+\nu A\bar v+P\mathcal B(\bar v)=0,
 \qquad
 \bar v(a)=v_0.
 \label{eq:v21-head-only-equation}
\end{equation}
For every Sobolev index $m$ and $h\downarrow0$,
\begin{align}
 w(a+h)
 &=-hD_0+O_{H^m}(h^2),
 \label{eq:v21-daughter-Taylor}\\
 \ell_P^{\rm in}(a+h)
 &=-h\mathfrak r_P(v_0)+O_{H^m}(h^2),
 \label{eq:v21-return-Taylor}\\
 v(a+h)-\bar v(a+h)
 &=\frac{h^2}{2}\mathfrak R_P(v_0)+O_{H^m}(h^3).
 \label{eq:v21-head-secant}
\end{align}
Moreover,
\begin{align}
 \frac12\|A^{-1/4}w(a+h)\|_2^2
 &=\frac{h^2}{2}\|d_0\|_2^2+O(h^3),
 \label{eq:v21-native-stock-Taylor}\\
 J_P^E(a+h)
 &=h\|D_0\|_2^2+O(h^2).
 \label{eq:v21-energy-current-Taylor}
\end{align}
\end{theorem}

\begin{proof}
The first two expansions are Taylor's theorem with
$w(a)=0$, $w_t(a)=-D_0$, and \eqref{eq:v21-return-derivative}.  The full
head equation is
\[
 \partial_tv+\nu Av+P\mathcal B(v)+A^{1/4}\ell_P^{\rm in}=0.
\]
At $a$, the full and head-only states and first derivatives agree.  Their
second-derivative difference is
\[
 \left.\frac{\dd^2}{\dd t^2}(v-\bar v)\right|_{t=a}
 =-A^{1/4}(\ell_P^{\rm in})'(a)
 =\mathfrak R_P(v_0).
\]
This proves the secant.  The last two expansions follow by squaring the
first and using Corollary~\ref{cor:v21-current-turnon}.
\end{proof}

% =============================================================================
\section{Finite return Gram and amplification of a small source residual}
\label{sec:v21-return-Gram}
% =============================================================================

Let $E^{\rm dau}\subseteq\operatorname{Ran}Q$ be a finite source screen
containing the range of the V20 exterior parent-output synthesis, and let
$F^{\rm ret}\subseteq\operatorname{Ran}P$ be a finite return screen.  Let
$P_E,P_F$ be their projections and define
\begin{equation}
 \mathscr V_{F\leftarrow E}^u
 :=P_F\mathscr V_{P\leftarrow Q}^uP_E.
 \label{eq:v21-finite-return-map}
\end{equation}

\begin{definition}[Positive daughter-to-return Gram]
\label{def:v21-return-Gram}
On $L^2(I;E^{\rm dau})$, put
\begin{equation}
 \boxed{
 \mathbb G_{E\to F}^{\rm ret}
 :=(\mathscr V_{F\leftarrow E}^u)^*
   \mathscr V_{F\leftarrow E}^u\succeq0.}
 \label{eq:v21-return-Gram}
\end{equation}
Thus
\begin{equation}
 \boxed{
 \langle f,\mathbb G_{E\to F}^{\rm ret}f\rangle
 =\|\mathscr V_{F\leftarrow E}^uf\|_{L^2(I;L^2)}^2.}
 \label{eq:v21-return-Gram-action}
\end{equation}
If $\mathsf D:\mathbb R^N\to L^2(I;E^{\rm dau})$ is the predeclared finite
synthesis of occurring daughter histories or donor variations, define
\begin{equation}
 \boxed{
 \mathbb G_{\mathsf D}^{\rm ret}
 :=(\mathscr V_{F\leftarrow E}^u\mathsf D)^*
   (\mathscr V_{F\leftarrow E}^u\mathsf D)\succeq0.}
 \label{eq:v21-finite-return-Gram}
\end{equation}
\end{definition}

\begin{proposition}[Return visibility and minimum response rank]
\label{prop:v21-return-rank}
The kernel of \eqref{eq:v21-return-Gram} is exactly the set of daughter
histories producing zero return on the declared output screen.  For the
finite synthesis $\mathsf D$,
\begin{equation}
 \operatorname{rank}\mathbb G_{\mathsf D}^{\rm ret}
 =\operatorname{rank}(\mathscr V_{F\leftarrow E}^u\mathsf D)
 \label{eq:v21-return-rank}
\end{equation}
is the minimum number of return-visible directions in that bank.
\end{proposition}

\begin{proof}
Both statements are ordinary Gram identities.  Equality of the response
range with an independently declared return atlas is not encoded by rank.
\end{proof}

\begin{theorem}[Finite-screen causal gain bound]
\label{thm:v21-finite-screen-gain}
Put $T_I=b-a$ and
\begin{equation}
 K_{E\to F}(I)
 :=\operatorname*{ess\,sup}_{a\le s\le t\le b}
 \left\|
 P_F\mathscr M_P^u(t)\mathscr U_Q^u(t,s)P_E
 \right\|_{\rm op}.
 \label{eq:v21-kernel-cap}
\end{equation}
This number is finite on every fixed smooth history and finite pair of
screens.  One has
\begin{equation}
 \boxed{
 \|\mathscr V_{F\leftarrow E}^u\|_{L^2\to L^2}
 \le T_IK_{E\to F}(I),}
 \label{eq:v21-Volterra-Schur}
\end{equation}
and hence
\begin{equation}
 \boxed{
 \langle f,\mathbb G_{E\to F}^{\rm ret}f\rangle
 \le T_I^2K_{E\to F}(I)^2\|f\|_{L^2(I;L^2)}^2.}
 \label{eq:v21-return-action-bound}
\end{equation}
\end{theorem}

\begin{proof}
Finite-dimensionality of the screens, smoothness of the coefficients, and
Lemma~\ref{lem:v21-evolution-family} imply finiteness.  If
$g=\mathscr V_{F\leftarrow E}^uf$, then
\[
 \|g(t)\|_2
 \le K_{E\to F}(I)\int_a^t\|f(s)\|_2\,\dd s.
\]
The Volterra integration operator has $L^2$ norm at most $T_I$ by
Cauchy--Schwarz and Fubini.  Squaring proves the action bound.
\end{proof}

\begin{corollary}[Vanishing source with surviving return forces gain escape]
\label{cor:v21-gain-escape}
Let $f_n$ be nonzero daughter histories on $I_n$.  If
\begin{equation}
 \|f_n\|_{L^2}^2=o(\nu\mathcal R_n^2),
 \qquad
 \|\mathscr V_n f_n\|_{L^2}^2
 \ge\delta\nu\mathcal R_n^2,
 \label{eq:v21-small-source-large-return}
\end{equation}
then
\begin{equation}
 \boxed{
 \frac{\|\mathscr V_nf_n\|_{L^2}}{\|f_n\|_{L^2}}
 \longrightarrow\infty,
 \qquad
 |I_n|K_n^{\rm ret}\longrightarrow\infty.}
 \label{eq:v21-return-gain-escape}
\end{equation}
\end{corollary}

\begin{proof}
The first conclusion is the quotient of the two assumptions; the second is
\eqref{eq:v21-Volterra-Schur}.
\end{proof}

\begin{firewall}[Return action is not additive across source labels]
\label{fire:v21-no-return-action-sum}
The source split \eqref{eq:v21-return-source-split} is linear and exact, but
its returned vectors may overlap.  Their squared norms are not summed as
independent physical payments.  The positive Gram is used on one
predeclared source bank or after the complete source-labelled panel has
been assembled.
\end{firewall}

% =============================================================================
\section{Insertion into the V20 terminal open-head card}
\label{sec:v21-terminal-insertion}
% =============================================================================

Retain a V20 candidate history with finite head $P_n$, source-to-return
interval $J_n=[a_n,b_n]$, record amplitude $\mathcal R_n$, and the V20
daughter split on its predeclared source window.  Extend the represented
and birth components by zero to $J_n$, and define
\begin{equation}
 d_n^{\rm out}=d_n^{\rm bg}+d_n^{\rm rep}+d_n^{\rm birth}.
 \label{eq:v21-terminal-daughter-split}
\end{equation}
Here $d_n^{\rm bg}$ contains all unselected daughter source.  Construct
$\mathscr U_n,\mathscr M_n,\mathscr V_n$ from the same actual trajectory and
fixed head.  Put
\begin{align}
 r_n^{\rm old}
 &:=\mathscr M_n(t)\mathscr U_n(t,a_n)A^{-1/4}Q_nu(a_n)
    +\mathscr V_nd_n^{\rm bg},
 \label{eq:v21-terminal-old-return}\\
 r_n^{\rm rep}
 &:=\mathscr V_nd_n^{\rm rep},
 \label{eq:v21-terminal-represented-return}\\
 r_n^{\rm birth}
 &:=\mathscr V_nd_n^{\rm birth}.
 \label{eq:v21-terminal-birth-return}
\end{align}
Then
\begin{equation}
 \boxed{
 \ell_n^{\rm in}
 =r_n^{\rm old}+r_n^{\rm rep}+r_n^{\rm birth}.}
 \label{eq:v21-terminal-return-assembly}
\end{equation}
No source is omitted.

Let $\mathbf S_n^{\rm ret}$ be the pointwise orthogonal projection onto the
physically declared incoming-return dictionary, and let
$\mathbf S_n^{\rm old}$ be the already declared old-tail dictionary.
Define
\begin{align}
 \mathfrak B_n^{\rm birth\to ret}
 &:=\frac{1}{\nu\mathcal R_n^2}
   \|r_n^{\rm birth}\|_{L^2(J_n;L^2)}^2,
 \label{eq:v21-birth-return-action}\\
 \mathfrak R_n^{\rm rep\to ret}
 &:=\frac{1}{\nu\mathcal R_n^2}
   \|(I-\mathbf S_n^{\rm ret})r_n^{\rm rep}\|_{L^2(J_n;L^2)}^2,
 \label{eq:v21-represented-return-residual}\\
 \mathfrak R_n^{\rm old}
 &:=\frac{1}{\nu\mathcal R_n^2}
   \|(I-\mathbf S_n^{\rm old})r_n^{\rm old}\|_{L^2(J_n;L^2)}^2.
 \label{eq:v21-old-return-residual}
\end{align}
The middle quantity is the previously symbolic V20 row
$\mathcal R_n^{\rm dau\to ret}$ in normalized square-action form.

\begin{theorem}[Exact return-atlas Pythagoras]
\label{thm:v21-return-atlas-Pythagoras}
The represented daughter return satisfies
\begin{equation}
 \boxed{
 \|r_n^{\rm rep}\|_{L^2}^2
 =\|\mathbf S_n^{\rm ret}r_n^{\rm rep}\|_{L^2}^2
  +\|(I-\mathbf S_n^{\rm ret})r_n^{\rm rep}\|_{L^2}^2.}
 \label{eq:v21-return-atlas-Pythagoras}
\end{equation}
The first term is a follower of the represented daughter source.  The
second is an ancestry-incidence residual.  The generated return
$r_n^{\rm birth}$ remains a follower of the V20 new source birth and is not
priced as another occurrence.
\end{theorem}

\begin{proof}
This is orthogonal Pythagoras for the return projection, integrated in time.
The ancestry statement follows from \eqref{eq:v21-terminal-return-assembly}.
\end{proof}

\begin{theorem}[V20 source tightness does not imply response tightness]
\label{thm:v21-source-response-tightness}
The V20 condition
\begin{equation}
 \mathfrak B_n^{\rm dau}
 =\frac{1}{\nu\mathcal R_n^2}
   \|d_n^{\rm birth}\|_{L^2}^2\longrightarrow0
 \label{eq:v21-V20-source-tight}
\end{equation}
does not by itself imply
$\mathfrak B_n^{\rm birth\to ret}\to0$.  When
$d_n^{\rm birth}\ne0$,
\begin{equation}
 \boxed{
 \mathfrak B_n^{\rm birth\to ret}
 =\bigl(\mathfrak g_n^{\rm birth}\bigr)^2
  \mathfrak B_n^{\rm dau},
 \qquad
 \mathfrak g_n^{\rm birth}
 =\frac{\|r_n^{\rm birth}\|_{L^2}}
        {\|d_n^{\rm birth}\|_{L^2}}.}
 \label{eq:v21-birth-gain-identity}
\end{equation}
Hence source tightness plus a positive lower return action forces
\begin{equation}
 \boxed{
 \mathfrak g_n^{\rm birth}\to\infty,
 \qquad
 |J_n|K_n^{\rm ret}\to\infty.}
 \label{eq:v21-birth-gain-divergence}
\end{equation}
\end{theorem}

\begin{proof}
The identity is the definition of the occurring gain after dividing by the
common normalization.  The conclusion is
Corollary~\ref{cor:v21-gain-escape}.
\end{proof}

% =============================================================================
\section{Corrected open-head alternative and export}
\label{sec:v21-export}
% =============================================================================

\begin{theorem}[Source-state-return gate for a finite terminal head]
\label{thm:v21-terminal-gate}
Apply the V19 internal gate and the V20 daughter-source gate first.  On a
candidate finite-head subsequence, exactly one of the following priority
branches occurs.
\begin{enumerate}[label=\textup{(G21.\arabic*)},leftmargin=5.5em]
\item \emph{Daughter first-birth survivor.}  The V20 normalized source action
      $\mathfrak B_n^{\rm dau}$ has a positive lower bound.
\item \emph{Birth-to-return amplification.}  One has
      $\mathfrak B_n^{\rm dau}\to0$, but for some $\delta>0$,
      \begin{equation}
       \mathfrak B_n^{\rm birth\to ret}\ge\delta.
       \label{eq:v21-amplified-birth-branch}
      \end{equation}
      The later packet is a follower of a small first birth with divergent
      causal gain; it is not a second source birth.
\item \emph{Old exterior-memory residual.}  The preceding rows vanish, but
      \begin{equation}
       \mathfrak R_n^{\rm old}\ge\delta
       \label{eq:v21-old-return-branch}
      \end{equation}
      cofinally.
\item \emph{Represented-daughter ancestry residual.}  The old row is
      represented, but
      \begin{equation}
       \mathfrak R_n^{\rm rep\to ret}\ge\delta.
       \label{eq:v21-return-incidence-branch}
      \end{equation}
      The daughter is represented at birth but its actual causal return is
      absent from the declared return atlas.
\item \emph{Reusable-bank return-gain survivor.}  The occurring scalar rows
      pass, but on the predeclared finite daughter-history synthesis
      \begin{equation}
       \|\mathbb G_{\mathsf D_n}^{\rm ret}\|_{\rm op}\to\infty.
       \label{eq:v21-return-Gram-survivor}
      \end{equation}
      This row is required only for reusable donor variations.
\item \emph{Ancestry-complete open return card.}  The V20 internal and source
      rows pass and
      \begin{equation}
       \boxed{
       \mathfrak B_n^{\rm birth\to ret}\to0,
       \qquad
       \mathfrak R_n^{\rm old}\to0,
       \qquad
       \mathfrak R_n^{\rm rep\to ret}\to0,}
       \label{eq:v21-return-pass}
      \end{equation}
      together with the bounded finite return-Gram row when reusable-bank
      semantics are active.  Then the V20 symbolic return incidence is
      populated by \eqref{eq:v21-return-operator}, and NS--A stops.
\item \emph{All-shell source-response escape.}  No exhausting finite-head
      sequence simultaneously controls the V19 internal rows, V20 daughter
      rows, and V21 causal response rows.  The output is one fused all-shell
      source-to-state-to-return survivor.
\end{enumerate}
\end{theorem}

\begin{proof}
The V19 and V20 priority order is retained, so no response estimate repairs
an earlier invariant, capture, incoming-leakage, or first-birth failure.  On
the V20 source-tight branch, the nonnegative birth-return action either has
a positive lower bound on a subsequence or tends to zero.  The two
nonnegative return-dictionary residuals give the next splits.  If a reusable
bank is intended, the top eigenvalue of the finite return Gram is either
bounded on a subsequence or escapes.  Passage of every row is the sixth
branch.  Failure of every exhausting finite head is exactly the final fused
escape.  The signed assembly and one-component Pythagoras prevent duplicate
source or return pricing.
\end{proof}

\begin{corollary}[Populated daughter-to-return export packet]
\label{cor:v21-export-packet}
On branch~\textup{(G21.6)}, append to the V20 packet
\begin{equation}
 \boxed{
 \begin{aligned}
 \mathfrak D_n^{V21}=\bigl(&
  \mathscr U_n,\mathscr V_n,\mathbb G_{\mathsf D_n}^{\rm ret},
  r_n^{\rm old},r_n^{\rm rep},r_n^{\rm birth},\\
 &\mathfrak B_n^{\rm birth\to ret},
  \mathfrak R_n^{\rm old},
  \mathfrak R_n^{\rm rep\to ret}
 \bigr).
 \end{aligned}}
 \label{eq:v21-export-packet}
\end{equation}
The formal V20 component $\mathcal R_n^{\rm dau\to ret}$ is thereby replaced
by the explicit residual \eqref{eq:v21-represented-return-residual}.  The
combined V19--V21 card is eligible for the fixed-datum NS--B paid-stress and
Reynolds-covariance audit.
\end{corollary}

\begin{proof}
Every entry is constructed from the same trajectory, finite head, source
metric, source atlas, and common endpoint.  The exact factorization gives
the occurrence incidence.  On the passing branch every unrepresented
response component is negligible at the declared scale, while represented
returns remain followers of their first source occurrences.
\end{proof}

\begin{firewall}[A bounded return Gram is not an upper budget]
\label{fire:v21-no-upper-budget}
A bounded causal return Gram is not a finite-stock upper budget, a
contraction coefficient, or a regularity estimate.  It says only that the
selected daughter-source atlas has a controlled response on the declared
history.  The NS--B paid-stress and Reynolds square-function rows, and any
critical acceleration or replenishment upper bound, remain independent.
\end{firewall}

% =============================================================================
\section{Integrated V21 theorem, proof-status ledger, and stopping line}
\label{sec:v21-integrated}
% =============================================================================

\begin{theorem}[Source-native daughter response and first-return theorem]
\label{thm:v21-integrated}
Preserving Versions~V1--V20, Version~V21 establishes the following.
\begin{enumerate}[label=\textup{(T21.\arabic*)},leftmargin=5.7em]
\item The exact $P/Q$ block equation has the radial-homogeneous
      source-to-state factorization \eqref{eq:v21-exterior-linear-equation}
      and the exact incoming read \eqref{eq:v21-incoming-linear-read}.
\item The actual-trajectory evolution family gives the causal
      daughter-to-return operator \eqref{eq:v21-return-operator} and the
      no-remainder assembly \eqref{eq:v21-return-factorization}.
\item The V20 daughter action is the clean-birth curvature of the bounded
      stock $\frac12\|A^{-1/4}Qu\|_2^2$.
\item The first return is the cubic coefficient
      \eqref{eq:v21-cubic-return-germ}, with the energy and helicity identities
      \eqref{eq:v21-return-energy-identity}--
      \eqref{eq:v21-return-helicity-identity}.
\item Every nonzero clean daughter has a nonzero return germ, positive
      kinetic-current slope, and the head secant \eqref{eq:v21-head-secant}.
\item Bounded negative-order stock does not control integrated daughter
      action; Countertheorem~\ref{cth:v21-curvature-no-action-budget}
      isolates the missing monotone or first-exit input.
\item The positive return Gram \eqref{eq:v21-return-Gram} is the exact action
      of the declared source-to-return map and satisfies the finite-screen
      bound \eqref{eq:v21-Volterra-Schur}.
\item A vanishing V20 source residual can have a nonvanishing return only
      through divergent causal gain.
\item The corrected terminal gate is Theorem~\ref{thm:v21-terminal-gate}.
      On its passing branch, the V20 daughter-to-return row is populated and
      NS--A reaches its intended NS--B interface.
\end{enumerate}
No branch is excluded for every smooth datum and no global regularity
conclusion is drawn.
\end{theorem}

\begin{proof}
The first two items are
Theorems~\ref{thm:v21-exact-block-factorization} and
\ref{thm:v21-return-Volterra}, with
Lemma~\ref{lem:v21-evolution-family}.  The clean-birth items are
Theorems~\ref{thm:v21-negative-stock-balance},
\ref{thm:v21-clean-birth-curvature},
\ref{thm:v21-return-germ-invariants}, and
\ref{thm:v21-head-only-secant}, together with
Countertheorem~\ref{cth:v21-curvature-no-action-budget}.  The return-action
items are Definition~\ref{def:v21-return-Gram},
Theorem~\ref{thm:v21-finite-screen-gain}, and
Corollary~\ref{cor:v21-gain-escape}.  The final item is
Theorem~\ref{thm:v21-terminal-gate} and
Corollary~\ref{cor:v21-export-packet}.
\end{proof}

\begingroup\small
\begin{longtable}{p{0.25\textwidth}p{0.19\textwidth}p{0.48\textwidth}}
\caption{NS--A Version V21 proof-status ledger.}\label{tab:v21-ledger}\\
\toprule Row & Status & Exact conclusion\\\midrule
\endfirsthead
\toprule Row & Status & Exact conclusion\\\midrule
\endhead
V1--V20 prefix & \PROVED{} preserved
 &The complete V20 preterminal source is retained byte-for-byte. No earlier
 theorem or limitation is rewritten.\\[0.3em]
V20 daughter first-birth short & \INHERITED{} exact
 &$d^{\rm out}=d^{\rm rep}+d^{\rm birth}$ remains the source-first split
 before any return is priced.\\[0.3em]
Exterior source-to-state equation & \PROVED{} exact
 &Equation \eqref{eq:v21-exterior-linear-equation} is the full exterior block
 equation on the actual trajectory.\\[0.3em]
State-to-incoming read & \PROVED{} exact
 &$\ell^{\rm in}=\mathscr M_P^uy$ contains the cross and exterior-self
 contributions with no source remainder.\\[0.3em]
Daughter-to-return map & \PROVED{} exact causal
 &Equation \eqref{eq:v21-return-factorization} separates old exterior memory
 from the causal response of the head daughter.\\[0.3em]
Source-native stock & \PROVED{} finite
 &$\frac12\|A^{-1/4}Qu\|_2^2$ is energy bounded and has clean-birth curvature
 $\|d\|_2^2$.\\[0.3em]
Integrated action from curvature & \OBSTRUCTION{} explicit
 &The linear model \eqref{eq:v21-stock-action-no-go} has vanishing stock
 height and divergent source action.\\[0.3em]
Cubic first-return germ & \PROVED{} finite
 &$\mathfrak R_P(v)=2P\mathcal B_{\rm s}(v,Q\mathcal B(v))$ is the diagonal
 return kernel at a clean birth.\\[0.3em]
Energy return identity & \PROVED{} positive
 &$\langle\mathfrak R_P(v),v\rangle=-\|Q\mathcal B(v)\|_2^2$; a nonzero
 clean daughter cannot have zero first return.\\[0.3em]
Head-only secant & \PROVED{} exact asymptotic
 &The full head differs from its Galerkin evolution by
 $\frac12h^2\mathfrak R_P(v)+O(h^3)$.\\[0.3em]
Finite return Gram & \PROVED{} positive
 &$\mathbb G^{\rm ret}=\mathscr V^*\mathscr V$ measures return-visible daughter
 directions on the declared source/output screens.\\[0.3em]
V20 source tightness implies response tightness & \OBSTRUCTION{} corrected
 &A vanishing birth source can generate a record-sized return only through
 divergent causal gain.\\[0.3em]
Old exterior return & \OPEN{} physical incidence
 &It must be shorted against the existing old-tail dictionary at the common
 endpoint.\\[0.3em]
Represented daughter return & \PROVED{} exact short / \OPEN{} value
 &Equation \eqref{eq:v21-return-atlas-Pythagoras} gives the exact residual;
 its terminal value is not populated by the theorem layer.\\[0.3em]
Reusable-bank return gain & \CONDITIONAL{} finite
 &The finite Gram \eqref{eq:v21-finite-return-Gram} is required only when
 donor variations are to be reused.\\[0.3em]
Ancestry-complete open export & \CONDITIONAL{} corrected
 &Pass requires source tightness, birth-response tightness, old-return
 representation, and daughter-to-return incidence.\\[0.3em]
All-shell source-response escape & \OPEN{} typed alternative
 &Failure of every finite card remains one fused source-to-state-to-return
 escape, not a shell-count sum.\\[0.3em]
Three-dimensional regularity & \OPEN
 &No critical upper action, stock, trace, rigidity, or global regularity
 theorem is inferred.\\
\bottomrule
\end{longtable}
\endgroup

% =============================================================================
\section{Exact mandate after Version V21}
\label{sec:v21-next}
% =============================================================================

\begin{programme}[Evaluate the actual return card or stop NS--A]
\label{prog:v21-next}
Preserve Versions~V1--V21 verbatim.  There is no automatic generic
Version~V22.

A continuation is licensed only in the following order.
\begin{enumerate}[label=\textup{(V22.\arabic*)},leftmargin=5.9em]
\item On the selected V20 finite head, evaluate the actual finite kernel
      moments of $\mathscr M_n(t)\mathscr U_n(t,s)$ needed by the declared
      return dictionary.  Do not enlarge the head after inspecting the
      return.
\item Evaluate $r_n^{\rm birth}=\mathscr V_nd_n^{\rm birth}$ before using
      $\mathfrak B_n^{\rm dau}\to0$ as a pass.  A surviving follower is the
      dynamic-gain branch and receives no second source price.
\item Short $r_n^{\rm old}$ against the existing old-tail dictionary and
      $r_n^{\rm rep}$ against the existing return atlas at the same endpoint.
      Assemble their signed sum before reading any total positive payment.
\item If reusable donor variations are intended, populate the finite return
      Gram on exactly the V20 exterior parent-output range.  A large top edge
      is the return-gain survivor; do not replace it by an ambient norm.
\item If the V19 internal, V20 source-boundary, and V21 response-boundary rows
      pass, terminate NS--A and export the combined card to NS--B.
\item If no finite head passes, retain all-shell source-response escape.
      Reopen it only with a new source-relative critical acceleration,
      capacity, or square-function upper stock on this exact fused carrier.
\item The clean-birth curvature theorem may be used as a local incidence
      check, but not as an integrated action budget without a new monotone or
      first-exit estimate.
\item Do not add another invariant family, path law, age law, Dini modulus,
      compactness space, shell count, or generic response compiler.
\end{enumerate}
\end{programme}

\begin{assessment}[Programme position after V21]
\label{ass:v21-position}
The source boundary and the response boundary are now separately complete.
V20 identifies which exterior source is born; V21 identifies the exact
causal state and return generated by that source on the actual trajectory.
At a clean birth the first return is a finite cubic coefficient, and at
finite lag it is one explicit Volterra matrix on the declared screens.  The
first unpopulated object is a physical value of this finite card, not
another theorem-level source or compactness construction.
\end{assessment}

\paragraph{Version V21 history.}
Preserved the complete V20 source.  Moved upstream of its symbolic
source-to-return incidence and inserted the missing source-to-state layer.
Used symmetric polarization and Euler homogeneity to rewrite the exact
exterior equation as a nonautonomous linear equation on the actual
trajectory.  Constructed its evolution family and exact causal
Daughter-to-return Volterra operator.  Identified the bounded $H^{-1/2}$
exterior stock whose clean-birth curvature equals the V20 daughter action,
and supplied a linear model showing that stock curvature does not yield an
integrated action budget.  Evaluated the first return as the cubic
coefficient $2P\mathcal B_{\rm s}(v,Q\mathcal B(v))$, proved its energy and
helicity identities, and derived the second-order difference from the
head-only Galerkin trajectory.  Constructed the finite positive return Gram
and causal kernel bound.  Inserted represented, birth-generated, and old
exterior responses into the terminal card and corrected the V20 pass by the
birth-to-return amplification row.

\begin{assessment}[Append-only integrity]
\label{ass:v21-integrity}
The complete Version~V20 source preceding its sole final executable document
terminator is preserved verbatim.  Its preterminal SHA--256 digest is
\begin{center}\scriptsize\ttfamily
85e6d39c493ab5e7972452ff3c941d67ff089a9bc7e47674eb9e0ec319701213
\end{center}
The present material begins at Section~\ref{sec:v21-scope}; the final
command below is again unique.
\end{assessment}

\addtocontents{toc}{\protect\endgroup}
% =============================================================================
% END VERSION V21 MATERIAL -- append later versions before final terminator
% =============================================================================

% APPEND ALL LATER VERSIONS IMMEDIATELY ABOVE THE SOLE FINAL LINE BELOW.
\clearpage
% =============================================================================
% VERSION V22: NEW MATERIAL.  The complete V1--V21 source above is preserved
% verbatim; only its sole active \end{document} line was removed before append.
% =============================================================================
\hypersetup{pdftitle={NS-A Terminal Source Concentration Programme: Cumulative V22},pdfauthor={A. Pelissier}}
\makeatletter
\addtocontents{toc}{\protect\begingroup\protect\makeatletter}
\addtocontents{toc}{\protect\def\protect\l@section{\protect\@dottedtocline{1}{0em}{4.7em}}}
\addtocontents{toc}{\protect\def\protect\l@subsection{\protect\@dottedtocline{2}{1.5em}{5.3em}}}
\addtocontents{toc}{\protect\makeatother}
\makeatother
\renewcommand{\thesection}{V22.\arabic{section}}
\renewcommand{\thesubsection}{V22.\arabic{section}.\arabic{subsection}}
\renewcommand{\theequation}{V22.\arabic{equation}}
\setcounter{section}{0}
\setcounter{equation}{0}
\renewcommand{\theHsection}{V22.\arabic{section}}
\renewcommand{\theHsubsection}{V22.\arabic{section}.\arabic{subsection}}
\renewcommand{\theHtheorem}{V22.\arabic{section}.\arabic{theorem}}
\renewcommand{\theHequation}{V22.\arabic{equation}}

\begin{center}
{\LARGE\bfseries NS--A: Version V22}\\[0.5em]
{\large The Target-Native Adjoint Return Read, Scalar Range Incidence,\\
and Removal of Volterra Target Inflation}\\[0.5em]
{\normalsize 3 September 2026}
\end{center}
\addcontentsline{toc}{part}{Version V22: target-native adjoint return and scalar ancestry}

\begin{abstract}
Version~V21 constructed the complete causal daughter-to-return Volterra map
for one fixed Fourier head.  That construction remains exact and is needed
when the target is the complete returned trajectory, a reusable donor bank,
or a return-atlas range statement.  The present continuation moves one arrow
upstream and distinguishes that strong target from the actual NS--A target on
one selected singular thread: a single signed return contribution to the
already frozen critical source--writer scalar.

The distinction is forced by the accompanying programmes.  A complete
stochastic coupling is stronger than one mixed Hilbert-source Gram; a
fixed-action scalar must be read before a phase or coefficient is refitted;
and a same-history mixed amplitude must be acquired before exterior or
Duhamel promotion.  Applied here, those observations say that the output
norm of the complete V21 return and its full range incidence are not necessary
to decide one predeclared critical contact.  What is necessary is one
backward adjoint writer on the same trajectory.

For a fixed head-supported target writer $\phi$, Version~V22 defines
\[
 \Psi_\phi(s)
 =\int_s^b
   \mathscr U_Q^u(t,s)^*\mathscr M_P^u(t)^*\phi(t)\,\dd t.
\]
It proves that $\Psi_\phi$ is the unique terminal-zero solution of the exact
backward adjoint equation and that
\[
 -\int_a^b\ip{\ell_P^{\rm in}(t)}{\phi(t)}\,\dd t
 =-\ip{y(a)}{\Psi_\phi(a)}
  +\int_a^b\ip{d(s)}{\Psi_\phi(s)}\,\dd s.
\]
Thus old exterior memory, represented daughter source, and new daughter
source are assembled directly on the one physical target without first
reconstructing the whole returned field.

The positive target gain is the single scalar
$\|\Psi_\phi\|_{L^2}^2/\|\phi\|_{L^2}^2$.  It gives the exact minimum source
action for a prescribed target return.  If the V20 daughter first-birth
action vanishes while its selected return stays of record size, this
one-dimensional gain diverges.  No orthogonal output channel is charged.
On finite screens, the unique Hilbert--Schmidt-minimal operator reproducing
the same target functional is rank one; the remaining part of the full
Volterra map is target null.  An explicit causal Volterra family has fixed
selected scalar and diverging complete return Gram, proving that the latter
is a target inflation for one scalar.

Version~V22 also replaces full return-atlas membership by the exact mixed
residual
\[
 \ip{(I-S)r}{(I-S)\phi}.
\]
It vanishes whenever the selected writer is represented, even if the returned
field has an arbitrarily large component outside the atlas.  The residual is
acquired by two positive polarization reads.  Consequently NS--A may stop and
export a target-native scalar ancestry card once the upstream material-cell
and boundary-current rows pass and this mixed scalar is populated.  The full
V21 response operator and return Gram remain mandatory only for the stronger
claim that a reusable bank or complete return trajectory has been acquired.
No terminal Dini improvement, critical upper action, rigidity theorem, or
global regularity conclusion is claimed.
\end{abstract}

% =============================================================================
\section{Scope: target semantics precede another return operator}
\label{sec:v22-scope}
% =============================================================================

The supplied fixed-datum Reynolds/paid-stress continuation places the
material--heat route mismatch, ultra-canonical frequency leverage, symmetric
strain, material-map curvature, and the viscous or Bernoulli material-cell
current before any captured interior Reynolds packet.  These rows are
inherited here and are not recomputed.  Version~V22 operates only on a branch
where those upstream tests have produced one actual, collar-free,
head-supported return writer.

The supplied boundary and quantum continuations make a second distinction.
A pathwise coupling, complete transition law, or complete exterior response
is a stronger target than one mixed scalar.  The NS--A terminal card has only
one return target at this stage: the contribution of the incoming source to
the predeclared critical writer.  The writer must be fixed before the return
is inspected.  If the target is later enlarged to a reusable donor bank or a
complete response trajectory, the full V21 Volterra map and return Gram are
retained without change.

\begin{firewall}[What is and is not being weakened]
\label{fire:v22-target-scope}
Version~V22 does not weaken the V21 source-to-state equation, causal evolution
family, or daughter-to-return factorization.  It changes only the amount of
that valid object required by the declared output.  Range membership of the
whole return, operator-norm control on every output direction, and a complete
return atlas remain necessary when those objects themselves are claimed.
They are not imposed merely to evaluate one same-history scalar.
\end{firewall}

% =============================================================================
\section{The exact backward-adjoint return writer}
\label{sec:v22-adjoint-writer}
% =============================================================================

Retain the V21 smooth interval $I=[a,b]$, the fixed Fourier split $P+Q=I$,
the exterior state $y=A^{-1/4}Qu$, the daughter source
$d=A^{-1/4}Q\mathcal B(Pu)$, and the actual-trajectory operators
$\mathscr L_Q^u(t)$ and $\mathscr M_P^u(t)$.  Write
\begin{equation}
 \mathscr A_Q^u(t):=\nu A+\mathscr L_Q^u(t).
 \label{eq:v22-AQ}
\end{equation}
Thus the homogeneous V21 evolution family solves
\begin{equation}
 \partial_t\mathscr U_Q^u(t,s)
 =-\mathscr A_Q^u(t)\mathscr U_Q^u(t,s),
 \qquad
 \mathscr U_Q^u(s,s)=Q.
 \label{eq:v22-evolution-equation}
\end{equation}

Fix a real smooth head-supported target writer
\begin{equation}
 \phi\in C^\infty(I;\operatorname{Ran}P),
 \qquad \phi\not\equiv0,
 \label{eq:v22-target-writer}
\end{equation}
chosen before the incoming source is inspected.  Its target return scalar is
\begin{equation}
 \mathfrak h_\phi^{\rm in}
 :=-\int_a^b\ip{\ell_P^{\rm in}[u](t)}{\phi(t)}\,\dd t.
 \label{eq:v22-target-return-scalar}
\end{equation}

\begin{definition}[Target-native adjoint return writer]
\label{def:v22-adjoint-writer}
Define, for $a\le s\le b$,
\begin{equation}
 \boxed{
 \Psi_\phi(s)
 :=\int_s^b
   \mathscr U_Q^u(t,s)^*
   \mathscr M_P^u(t)^*\phi(t)\,\dd t.}
 \label{eq:v22-adjoint-writer}
\end{equation}
All adjoints are taken in the physical $L^2$ metric on the already fixed
smooth Sobolev screen.  The definition is independent of a basis in the
head or exterior spaces.
\end{definition}

\begin{theorem}[Backward adjoint equation and exact scalar duality]
\label{thm:v22-adjoint-duality}
The writer $\Psi_\phi$ is the unique smooth terminal-zero solution of
\begin{equation}
 \boxed{
 -\partial_s\Psi_\phi(s)
 +\mathscr A_Q^u(s)^*\Psi_\phi(s)
 =\mathscr M_P^u(s)^*\phi(s),
 \qquad
 \Psi_\phi(b)=0.}
 \label{eq:v22-backward-adjoint}
\end{equation}
For every smooth exterior source $f$,
\begin{equation}
 \boxed{
 \int_a^b
 \ip{(\mathscr V_{P\leftarrow Q}^uf)(t)}{\phi(t)}\,\dd t
 =-\int_a^b\ip{f(s)}{\Psi_\phi(s)}\,\dd s.}
 \label{eq:v22-target-duality}
\end{equation}
The complete incoming target scalar has the no-remainder identity
\begin{equation}
 \boxed{
 \mathfrak h_\phi^{\rm in}
 =-\ip{y(a)}{\Psi_\phi(a)}
  +\int_a^b\ip{d(s)}{\Psi_\phi(s)}\,\dd s.}
 \label{eq:v22-complete-target-rectangle}
\end{equation}
\end{theorem}

\begin{proof}
For a nonautonomous evolution family generated by
$-\mathscr A_Q^u(t)$, differentiation in the entrance time gives
\[
 \partial_s\mathscr U_Q^u(t,s)
 =\mathscr U_Q^u(t,s)\mathscr A_Q^u(s).
\]
Taking adjoints and differentiating
\eqref{eq:v22-adjoint-writer} yields
\[
 \partial_s\Psi_\phi(s)
 =-\mathscr M_P^u(s)^*\phi(s)
  +\mathscr A_Q^u(s)^*\Psi_\phi(s),
\]
which is \eqref{eq:v22-backward-adjoint}; the terminal condition follows
from the empty integral.  Uniqueness follows by the backward version of the
V21 energy estimate, or equivalently by time reversal and uniqueness of the
forward homogeneous equation.

Using the V21 Volterra formula and Fubini's theorem,
\begin{align*}
 \int_a^b\ip{(\mathscr V_{P\leftarrow Q}^uf)(t)}{\phi(t)}\,\dd t
 &=-\int_a^b\int_a^t
   \ip{\mathscr M_P^u(t)\mathscr U_Q^u(t,s)f(s)}{\phi(t)}
   \,\dd s\,\dd t\\
 &=-\int_a^b
   \ip{f(s)}{
    \int_s^b\mathscr U_Q^u(t,s)^*
     \mathscr M_P^u(t)^*\phi(t)\,\dd t}
   \,\dd s,
\end{align*}
which is \eqref{eq:v22-target-duality}.  The old exterior term satisfies
\[
 \int_a^b
 \ip{\mathscr M_P^u(t)\mathscr U_Q^u(t,a)y(a)}{\phi(t)}\,\dd t
 =\ip{y(a)}{\Psi_\phi(a)}.
\]
Insert the V21 factorization
$\ell_P^{\rm in}=\ell_P^{\rm old}+\mathscr Vd$ into
\eqref{eq:v22-target-return-scalar} and use the two preceding identities.
\end{proof}

\begin{corollary}[Exact old, represented, and first-birth scalar assembly]
\label{cor:v22-source-scalar-split}
Suppose the occurring daughter has the V21 source-first decomposition
\begin{equation}
 d=d^{\rm bg}+d^{\rm rep}+d^{\rm birth}.
 \label{eq:v22-daughter-split}
\end{equation}
Define
\begin{align}
 \mathfrak h_\phi^{\rm old}
 &:=-\ip{y(a)}{\Psi_\phi(a)}
   +\int_a^b\ip{d^{\rm bg}(s)}{\Psi_\phi(s)}\,\dd s,
 \label{eq:v22-old-scalar}\\
 \mathfrak h_\phi^{\rm rep}
 &:=\int_a^b\ip{d^{\rm rep}(s)}{\Psi_\phi(s)}\,\dd s,
 \label{eq:v22-represented-scalar}\\
 \mathfrak h_\phi^{\rm birth}
 &:=\int_a^b\ip{d^{\rm birth}(s)}{\Psi_\phi(s)}\,\dd s.
 \label{eq:v22-birth-scalar}
\end{align}
Then
\begin{equation}
 \boxed{
 \mathfrak h_\phi^{\rm in}
 =\mathfrak h_\phi^{\rm old}
  +\mathfrak h_\phi^{\rm rep}
  +\mathfrak h_\phi^{\rm birth}.}
 \label{eq:v22-three-scalar-assembly}
\end{equation}
The represented and first-birth pieces are evaluated through the same
adjoint writer.  The latter is a follower of the already counted daughter
birth and receives no second source price.
\end{corollary}

\begin{proof}
Substitute \eqref{eq:v22-daughter-split} into
\eqref{eq:v22-complete-target-rectangle} and use linearity before taking any
positive part.
\end{proof}

\begin{assessment}[The first finite kernel moment]
\label{ass:v22-first-kernel-moment}
For one scalar target, the first physical return datum is not the complete
kernel $(t,s)\mapsto\mathscr M_P^u(t)\mathscr U_Q^u(t,s)$.  It is the single
backward moment $\Psi_\phi$ in \eqref{eq:v22-adjoint-writer}.  Computing the
full kernel is necessary only when the return trajectory, its complete range,
or several independently reusable output directions are themselves the
target.
\end{assessment}

% =============================================================================
\section{Positive target gain and the exact minimum source action}
\label{sec:v22-target-gain}
% =============================================================================

Put
\begin{equation}
 \mathcal X_I=L^2(I;\operatorname{Ran}Q),
 \qquad
 \mathcal Y_I=L^2(I;\operatorname{Ran}P),
 \label{eq:v22-source-output-spaces}
\end{equation}
and use the physical $L^2$ inner products.  On the fixed smooth card, the
V21 return map extends to every source/output screen for which the displayed
kernel norm below is finite.

\begin{definition}[Target-visible return gain]
\label{def:v22-target-gain}
For $\phi\ne0$, define
\begin{equation}
 \boxed{
 \mathfrak g_\phi^2
 :=\frac{\|\Psi_\phi\|_{\mathcal X_I}^2}
          {\|\phi\|_{\mathcal Y_I}^2}
 =\frac{\|(\mathscr V_{P\leftarrow Q}^u)^*\phi\|_{\mathcal X_I}^2}
        {\|\phi\|_{\mathcal Y_I}^2}.}
 \label{eq:v22-target-gain}
\end{equation}
This is one positive scalar.  It is the gain visible to the predeclared
output direction, not the operator norm on every return direction.
\end{definition}

\begin{theorem}[Target-native Thomson--Riesz principle]
\label{thm:v22-target-Riesz}
Let $\tau\in\R$.  If $\Psi_\phi\ne0$, then
\begin{equation}
 \boxed{
 \inf\left\{
  \|f\|_{\mathcal X_I}^2:
  -\int_a^b\ip{f(s)}{\Psi_\phi(s)}\,\dd s=\tau
 \right\}
 =\frac{\tau^2}{\|\Psi_\phi\|_{\mathcal X_I}^2}.}
 \label{eq:v22-minimum-source-action}
\end{equation}
The unique minimizer is
\begin{equation}
 \boxed{
 f_{\min}
 =-\frac{\tau}{\|\Psi_\phi\|_{\mathcal X_I}^2}
   \Psi_\phi.}
 \label{eq:v22-minimum-source}
\end{equation}
If $\Psi_\phi=0$, every daughter source has zero selected return, so a
nonzero prescribed $\tau$ is infeasible.
\end{theorem}

\begin{proof}
Cauchy--Schwarz gives
$|\tau|^2\le\|f\|^2\|\Psi_\phi\|^2$.  Equality holds exactly when $f$ is a
scalar multiple of $\Psi_\phi$, and the constraint fixes the multiple as in
\eqref{eq:v22-minimum-source}.  If $\Psi_\phi=0$, the functional is
identically zero.
\end{proof}

\begin{corollary}[Target gain is bounded by one causal kernel moment]
\label{cor:v22-kernel-gain-bound}
Suppose
\begin{equation}
 K_{P,Q}(I)
 :=\operatorname*{ess\,sup}_{a\le s\le t\le b}
 \|\mathscr M_P^u(t)\mathscr U_Q^u(t,s)\|_{\rm op}<\infty.
 \label{eq:v22-kernel-sup}
\end{equation}
Then
\begin{equation}
 \boxed{
 \|\Psi_\phi\|_{\mathcal X_I}
 \le |I|K_{P,Q}(I)\|\phi\|_{\mathcal Y_I},
 \qquad
 \mathfrak g_\phi\le |I|K_{P,Q}(I).}
 \label{eq:v22-target-gain-bound}
\end{equation}
\end{corollary}

\begin{proof}
For almost every $s$,
\[
 \|\Psi_\phi(s)\|_2
 \le K_{P,Q}(I)\int_s^b\|\phi(t)\|_2\,\dd t.
\]
By Cauchy--Schwarz,
\[
 \left(\int_s^b\|\phi(t)\|_2\,\dd t\right)^2
 \le(b-s)\int_s^b\|\phi(t)\|_2^2\,\dd t.
\]
Integrating in $s$ and interchanging the order gives an upper bound
$|I|^2\|\phi\|_{\mathcal Y_I}^2$.  Take square roots and divide by
$\|\phi\|_{\mathcal Y_I}$.
\end{proof}

Return to a terminal card with record amplitude $\mathcal R_n$.  Define the
dimensionless target gain
\begin{equation}
 \boxed{
 \mathfrak G_n^\phi
 :=\frac{\nu}{\mathcal R_n^2}
   \|\Psi_{n,\phi}\|_{L^2(J_n;L^2)}^2.}
 \label{eq:v22-normalized-target-gain}
\end{equation}
The inherited daughter first-birth action is
\begin{equation}
 \mathfrak B_n^{\rm dau}
 =\frac{\|d_n^{\rm birth}\|_{L^2}^2}
        {\nu\mathcal R_n^2}.
 \label{eq:v22-inherited-daughter-action}
\end{equation}

\begin{corollary}[Vanishing daughter action with a fixed selected return]
\label{cor:v22-target-gain-escape}
The normalized first-birth return scalar satisfies
\begin{equation}
 \boxed{
 \left|
  \frac{\mathfrak h_{n,\phi}^{\rm birth}}{\mathcal R_n^2}
 \right|^2
 \le
 \mathfrak B_n^{\rm dau}\,\mathfrak G_n^\phi.}
 \label{eq:v22-birth-target-product}
\end{equation}
Consequently, if
\begin{equation}
 \mathfrak B_n^{\rm dau}\longrightarrow0,
 \qquad
 \frac{|\mathfrak h_{n,\phi}^{\rm birth}|}{\mathcal R_n^2}
 \ge\delta>0,
 \label{eq:v22-small-birth-large-target}
\end{equation}
then
\begin{equation}
 \boxed{
 \mathfrak G_n^\phi
 \ge\frac{\delta^2}{\mathfrak B_n^{\rm dau}}
 \longrightarrow\infty.}
 \label{eq:v22-target-gain-divergence}
\end{equation}
This is a target-visible causal-gain survivor.  It is weaker than divergence
of the complete return norm and is exactly sufficient for the selected
scalar.
\end{corollary}

\begin{proof}
Apply Cauchy--Schwarz to
\eqref{eq:v22-birth-scalar}, divide by $\mathcal R_n^4$, and insert the two
normalizations.  The lower bound follows by rearrangement.
\end{proof}

% =============================================================================
\section{Finite-screen operator minimality and a causal target-inflation counterexample}
\label{sec:v22-operator-minimality}
% =============================================================================

The next theorem isolates the exact amount of a return operator visible to
one target.  It is stated on a finite source/output screen, which is the only
setting in which a reusable finite donor card is claimed.

\begin{theorem}[Unique Hilbert--Schmidt-minimal operator for one return target]
\label{thm:v22-minimal-return-operator}
Let $\mathcal X,\mathcal Y$ be real Hilbert spaces, let
$V\in\mathfrak S_2(\mathcal X,\mathcal Y)$, and fix
$0\ne\phi\in\mathcal Y$.  Put
\begin{equation}
 \psi:=V^*\phi
 \label{eq:v22-abstract-pullback}
\end{equation}
and define the rank-one operator
\begin{equation}
 \boxed{
 V_\phi^{\min}x
 :=\frac{\ip{x}{\psi}}{\|\phi\|_{\mathcal Y}^2}\,\phi.}
 \label{eq:v22-minimal-operator}
\end{equation}
Then
\begin{equation}
 (V_\phi^{\min})^*\phi=V^*\phi,
 \qquad
 \ip{V_\phi^{\min}x}{\phi}=\ip{Vx}{\phi}
 \quad(x\in\mathcal X).
 \label{eq:v22-same-target-functional}
\end{equation}
Every Hilbert--Schmidt operator $W$ with $W^*\phi=V^*\phi$ has the orthogonal
splitting
\begin{equation}
 \boxed{
 \|W\|_{\rm HS}^2
 =\frac{\|\psi\|_{\mathcal X}^2}{\|\phi\|_{\mathcal Y}^2}
  +\|W-V_\phi^{\min}\|_{\rm HS}^2.}
 \label{eq:v22-operator-Pythagoras}
\end{equation}
Consequently $V_\phi^{\min}$ is the unique minimum-Hilbert--Schmidt operator
which reproduces the complete linear target functional.
\end{theorem}

\begin{proof}
The adjoint of \eqref{eq:v22-minimal-operator} satisfies
\[
 (V_\phi^{\min})^*y
 =\frac{\ip{y}{\phi}}{\|\phi\|^2}\,\psi,
\]
so evaluation at $y=\phi$ proves
\eqref{eq:v22-same-target-functional}.  If
$K=W-V_\phi^{\min}$, then $K^*\phi=0$.  The Hilbert--Schmidt inner product of
the rank-one operator with $K$ is
\[
 \ip{V_\phi^{\min}}{K}_{\rm HS}
 =\frac{\ip{K\psi}{\phi}}{\|\phi\|^2}
 =\frac{\ip{\psi}{K^*\phi}}{\|\phi\|^2}=0.
\]
Also
$\|V_\phi^{\min}\|_{\rm HS}^2=\|\psi\|^2/\|\phi\|^2$.
Pythagoras gives \eqref{eq:v22-operator-Pythagoras} and uniqueness.
\end{proof}

\begin{corollary}[The target-positive return Gram is rank one]
\label{cor:v22-rank-one-target-Gram}
Let
\begin{equation}
 P_\phi
 =\frac{|\phi\rangle\langle\phi|}{\|\phi\|^2}
 \label{eq:v22-target-projection}
\end{equation}
be the projection onto the target line.  Then
\begin{equation}
 \boxed{
 V^*P_\phi V
 =\frac{|\psi\rangle\langle\psi|}{\|\phi\|^2}
 \succeq0,}
 \label{eq:v22-target-Gram}
\end{equation}
and
\begin{equation}
 \ip{x}{V^*P_\phi Vx}
 =\frac{|\ip{Vx}{\phi}|^2}{\|\phi\|^2}.
 \label{eq:v22-target-Gram-action}
\end{equation}
This positive scalar or rank-one Gram is the complete return-action object
for one selected writer.
\end{corollary}

\begin{proof}
Since $P_\phi Vx=\phi\ip{\psi}{x}/\|\phi\|^2$, applying $V^*$ gives the
rank-one formula.  Pair with $x$.
\end{proof}

\begin{countertheorem}[A causal Volterra map can have fixed target and divergent full Gram]
\label{cth:v22-Volterra-target-inflation}
Let $K:L^2(0,1)\to L^2(0,1)$ be
\begin{equation}
 (Kf)(t)=\int_0^t f(s)\,\dd s,
 \label{eq:v22-scalar-Volterra}
\end{equation}
and define
\begin{equation}
 V_Nf=(Kf,NKf)
 \in L^2(0,1;\R^2).
 \label{eq:v22-Volterra-family}
\end{equation}
For the fixed target $\phi(t)=(1,0)$,
\begin{equation}
 \boxed{
 V_N^*\phi=1-s
 \quad\text{for every }N,}
 \label{eq:v22-fixed-pullback}
\end{equation}
so the selected return functional is independent of $N$:
\begin{equation}
 \boxed{
 \ip{V_Nf}{\phi}_{L^2}
 =\int_0^1(1-s)f(s)\,\dd s.}
 \label{eq:v22-fixed-target-functional}
\end{equation}
Nevertheless,
\begin{equation}
 \boxed{
 \|V_N\|_{\rm HS}^2=\frac{1+N^2}{2},
 \qquad
 V_N^*V_N=(1+N^2)K^*K.}
 \label{eq:v22-divergent-full-Gram}
\end{equation}
Thus complete return gain can diverge entirely in a target-null output
channel while the selected physical scalar remains unchanged.
\end{countertheorem}

\begin{proof}
Fubini gives
\[
 \int_0^1(Kf)(t)\,\dd t
 =\int_0^1(1-s)f(s)\,\dd s,
\]
which proves the first two identities.  The kernel of $K$ is
$\one_{\{0<s<t<1\}}$, so
$\|K\|_{\rm HS}^2=1/2$.  Orthogonality of the two output coordinates gives
$\|V_N\|_{\rm HS}^2=(1+N^2)\|K\|_{\rm HS}^2$ and
$V_N^*V_N=(1+N^2)K^*K$.
\end{proof}

\begin{firewall}[A target-null return is not a physical obstruction]
\label{fire:v22-target-null-return}
Countertheorem~\ref{cth:v22-Volterra-target-inflation} does not say that a
large orthogonal return is physically irrelevant for every purpose.  It says
that it is irrelevant to the one target which annihilates it.  If NS--B asks
for a reusable response bank, complete range, or pathwise returned field, the
orthogonal output is part of that stronger target and the V21 full Gram must
be evaluated.
\end{firewall}

% =============================================================================
\section{Target-relative return-atlas incidence}
\label{sec:v22-target-atlas}
% =============================================================================

Let $\mathcal Y$ be a real output Hilbert space, let $S$ be the orthogonal
projection onto a physically declared return atlas, let $r\in\mathcal Y$ be
an occurring return, and let $\phi\in\mathcal Y$ be the frozen target writer.

\begin{theorem}[Exact mixed atlas residual for one scalar]
\label{thm:v22-mixed-atlas-residual}
The selected scalar has the orthogonal two-term decomposition
\begin{equation}
 \boxed{
 \ip{r}{\phi}
 =\ip{Sr}{S\phi}
  +\ip{(I-S)r}{(I-S)\phi}.}
 \label{eq:v22-atlas-scalar-split}
\end{equation}
Define the target-relative atlas residual
\begin{equation}
 \boxed{
 \Delta_S(r,\phi)
 :=\ip{(I-S)r}{(I-S)\phi}.}
 \label{eq:v22-target-atlas-residual}
\end{equation}
Then the atlas reproduces the selected scalar exactly if and only if
$\Delta_S(r,\phi)=0$.  Full range membership $(I-S)r=0$ is sufficient but
not necessary.  In particular, if $\phi\in\Ran S$, then
\begin{equation}
 \boxed{\Delta_S(r,\phi)=0\quad\text{for every }r.}
 \label{eq:v22-represented-target-automatic}
\end{equation}
\end{theorem}

\begin{proof}
Expand both vectors into their $S$ and $I-S$ components.  The two ranges are
orthogonal, so both cross terms vanish.  The remaining assertions follow
immediately.
\end{proof}

\begin{corollary}[Two positive reads acquire the signed mixed residual]
\label{cor:v22-positive-polarization}
Put
\begin{equation}
 x=(I-S)r,
 \qquad
 z=(I-S)\phi.
 \label{eq:v22-atlas-polar-vectors}
\end{equation}
Then
\begin{equation}
 \boxed{
 4\Delta_S(r,\phi)
 =\|x+z\|_2^2-\|x-z\|_2^2.}
 \label{eq:v22-atlas-polarization}
\end{equation}
Thus one signed target-incidence row is reconstructed by two positive norm
reads on the same output card.  Separate values of $\|x\|$ and $\|z\|$ do
not determine its sign.
\end{corollary}

\begin{proof}
Subtract the two parallelogram expansions.  The last statement follows, for
example, by taking unit vectors at angles $\theta$ and $\pi-\theta$.
\end{proof}

\begin{countertheorem}[Full return-atlas residual is strictly stronger]
\label{cth:v22-full-atlas-too-strong}
Let $\mathcal Y=\R^2$, let $S$ project onto $\R e_1$, and take
\begin{equation}
 \phi=e_1,
 \qquad
 r_N=e_1+Ne_2.
 \label{eq:v22-atlas-counterexample}
\end{equation}
Then
\begin{equation}
 \boxed{
 \ip{r_N}{\phi}=1,
 \qquad
 \Delta_S(r_N,\phi)=0,
 \qquad
 \|(I-S)r_N\|=N\longrightarrow\infty.}
 \label{eq:v22-atlas-counterexample-values}
\end{equation}
The declared atlas represents the complete selected scalar even though the
full return-range residual diverges.
\end{countertheorem}

\begin{proof}
All three identities are immediate from the orthogonal coordinates.
\end{proof}

\begin{assessment}[Same-history mixed incidence is the actual read]
\label{ass:v22-mixed-incidence}
The scalar target is not determined by a return norm and a writer norm on
separate histories.  It is the mixed same-history quantity
\eqref{eq:v22-target-atlas-residual}, or equivalently the positive
polarization pair \eqref{eq:v22-atlas-polarization}.  Choosing the writer or
atlas after observing the return would fit the residual and supplies no
physical incidence certificate.
\end{assessment}

% =============================================================================
\section{Insertion into the V19--V21 terminal open-head card}
\label{sec:v22-terminal-insertion}
% =============================================================================

Retain a candidate terminal head $P_n$, exterior projection $Q_n$, interval
$J_n=[a_n,b_n]$, record amplitude $\mathcal R_n$, and the V21 decomposition
\begin{equation}
 \ell_n^{\rm in}
 =r_n^{\rm old}+r_n^{\rm rep}+r_n^{\rm birth}.
 \label{eq:v22-inherited-return-assembly}
\end{equation}
Assume first that the upstream material-route, leverage, strain,
material-curvature, capture, and signed boundary-current tests have produced
a physical head-supported writer
\begin{equation}
 0\ne\phi_n\in L^2(J_n;\Ran P_n)
 \label{eq:v22-terminal-target-writer}
\end{equation}
which was frozen before the three returns were inspected.  If no such writer
has occurred, the output is the first surviving upstream NS--B/NS--A row and
the remainder of this section is not entered.

Let $\Psi_n$ be the adjoint writer
\eqref{eq:v22-adjoint-writer}.  Define the direct same-history scalar reads
\begin{align}
 \mathfrak h_n^{\rm old}
 &:=-\ip{r_n^{\rm old}}{\phi_n}_{L^2(J_n)},
 \label{eq:v22-terminal-old-read}\\
 \mathfrak h_n^{\rm rep}
 &:=-\ip{r_n^{\rm rep}}{\phi_n}_{L^2(J_n)},
 \label{eq:v22-terminal-rep-read}\\
 \mathfrak h_n^{\rm birth}
 &:=-\ip{r_n^{\rm birth}}{\phi_n}_{L^2(J_n)},
 \label{eq:v22-terminal-birth-read}\\
 \mathfrak h_n^{\rm in}
 &:=-\ip{\ell_n^{\rm in}}{\phi_n}_{L^2(J_n)}.
 \label{eq:v22-terminal-total-read}
\end{align}
By Theorem~\ref{thm:v22-adjoint-duality}, these same numbers are
\begin{align}
 \mathfrak h_n^{\rm old}
 &=-\ip{A^{-1/4}Q_nu(a_n)}{\Psi_n(a_n)}
   +\ip{d_n^{\rm bg}}{\Psi_n}_{L^2},
 \label{eq:v22-terminal-old-adjoint}\\
 \mathfrak h_n^{\rm rep}
 &=\ip{d_n^{\rm rep}}{\Psi_n}_{L^2},
 \label{eq:v22-terminal-rep-adjoint}\\
 \mathfrak h_n^{\rm birth}
 &=\ip{d_n^{\rm birth}}{\Psi_n}_{L^2}.
 \label{eq:v22-terminal-birth-adjoint}
\end{align}
In particular,
\begin{equation}
 \boxed{
 \mathfrak h_n^{\rm in}
 =\mathfrak h_n^{\rm old}
  +\mathfrak h_n^{\rm rep}
  +\mathfrak h_n^{\rm birth}.}
 \label{eq:v22-terminal-scalar-assembly}
\end{equation}

Let $\mathbf S_n^{\rm old}$ and $\mathbf S_n^{\rm ret}$ be the V21 old-tail
and represented-return projections.  Define the two target-relative atlas
residuals
\begin{align}
 \Delta_n^{\rm old,\phi}
 &:=-\ip{(I-\mathbf S_n^{\rm old})r_n^{\rm old}}
             {(I-\mathbf S_n^{\rm old})\phi_n}_{L^2},
 \label{eq:v22-old-target-residual}\\
 \Delta_n^{\rm rep,\phi}
 &:=-\ip{(I-\mathbf S_n^{\rm ret})r_n^{\rm rep}}
             {(I-\mathbf S_n^{\rm ret})\phi_n}_{L^2}.
 \label{eq:v22-rep-target-residual}
\end{align}
They replace the full norm residuals only for this one scalar target.  If the
full return field or reusable output bank is claimed, the V21 norm residuals
remain mandatory.

\begin{definition}[Target-native V22 return card]
\label{def:v22-terminal-card}
The minimum finite card for the selected return scalar is
\begin{equation}
 \boxed{
 \begin{aligned}
 \mathfrak D_n^{\rm V22,\phi}:=\bigl(&
  P_n,Q_n,J_n,\phi_n,\Psi_n,
  \|\phi_n\|_{L^2}^2,\|\Psi_n\|_{L^2}^2,\\
 &\mathfrak h_n^{\rm old},
  \mathfrak h_n^{\rm rep},
  \mathfrak h_n^{\rm birth},
  \Delta_n^{\rm old,\phi},
  \Delta_n^{\rm rep,\phi}
 \bigr).
 \end{aligned}}
 \label{eq:v22-terminal-card}
\end{equation}
Every signed mixed entry is acquired by positive polarization on one common
history.  No complete return trajectory is part of the card unless it is an
independent target.
\end{definition}

\begin{theorem}[Target-native finite-head alternative]
\label{thm:v22-terminal-gate}
Apply the V19 internal invariant gate, the V20 source-boundary gate, and the
upstream material/cell/boundary-current gate before the present theorem.  On
a branch carrying the fixed writer \eqref{eq:v22-terminal-target-writer},
pass to a priority subsequence.  Exactly one of the following applies.
\begin{enumerate}[label=\textup{(G22.\arabic*)},leftmargin=5.7em]
\item \emph{Target-scalar null return.}  One has
      \begin{equation}
       \frac{|\mathfrak h_n^{\rm in}|}{\mathcal R_n^2}
       \longrightarrow0.
       \label{eq:v22-scalar-null-return}
      \end{equation}
      The V21 response row passes for this selected scalar, irrespective of
      target-null output components.
\item \emph{Old-tail target-incidence survivor.}  The total selected return is
      nonnegligible and, for some $\delta>0$,
      \begin{equation}
       \frac{|\Delta_n^{\rm old,\phi}|}{\mathcal R_n^2}
       \ge\delta.
       \label{eq:v22-old-target-survivor}
      \end{equation}
      The old exterior response is not represented on the actual target.
\item \emph{Represented-daughter target-incidence survivor.}  The preceding
      residual vanishes, but
      \begin{equation}
       \frac{|\Delta_n^{\rm rep,\phi}|}{\mathcal R_n^2}
       \ge\delta.
       \label{eq:v22-rep-target-survivor}
      \end{equation}
      The daughter source is represented, while its return fails the declared
      target-relative atlas incidence.
\item \emph{Daughter first-birth survivor.}  The target-incidence residuals
      vanish, but the inherited source action satisfies
      \begin{equation}
       \mathfrak B_n^{\rm dau}\ge\delta.
       \label{eq:v22-daughter-source-survivor}
      \end{equation}
      This is the V20 source-first-birth branch and is not repriced by the
      return scalar.
\item \emph{Target-visible causal-gain survivor.}  The daughter action tends
      to zero, but
      \begin{equation}
       \frac{|\mathfrak h_n^{\rm birth}|}{\mathcal R_n^2}
       \ge\delta.
       \label{eq:v22-target-gain-survivor}
      \end{equation}
      Then $\mathfrak G_n^\phi\to\infty$ by
      \eqref{eq:v22-target-gain-divergence}.  The return remains a follower
      of the original first birth.
\item \emph{Target-native ancestry-complete return card.}  The target scalar
      is nonnegligible, the two mixed atlas residuals vanish, the daughter
      first-birth action tends to zero, and its target return tends to zero.
      Then
      \begin{equation}
       \boxed{
       \frac{\mathfrak h_n^{\rm in}
             -\mathfrak h_{n,\rm atlas}^{\rm old}
             -\mathfrak h_{n,\rm atlas}^{\rm rep}}
            {\mathcal R_n^2}
       \longrightarrow0,}
       \label{eq:v22-target-native-pass}
      \end{equation}
      where
      \begin{align*}
       \mathfrak h_{n,\rm atlas}^{\rm old}
       &:=-\ip{\mathbf S_n^{\rm old}r_n^{\rm old}}
                  {\mathbf S_n^{\rm old}\phi_n}_{L^2},\\
       \mathfrak h_{n,\rm atlas}^{\rm rep}
       &:=-\ip{\mathbf S_n^{\rm ret}r_n^{\rm rep}}
                  {\mathbf S_n^{\rm ret}\phi_n}_{L^2}.
      \end{align*}
      The selected return has complete ancestry at the scalar level and
      NS--A may export \eqref{eq:v22-terminal-card} without reconstructing the
      target-null response.
\item \emph{Stronger reusable-bank target.}  If the intended conclusion is
      range membership of the complete return, uniform response to donor
      variations, or a returned path law, the target has been enlarged.  The
      full V21 norm residuals and return Gram must then be populated; the
      scalar card does not substitute for them.
\end{enumerate}
If no exhausting finite head gives a physical writer and a convergent scalar
card, the output is all-shell target-return escape rather than a sum of shell
prices.
\end{theorem}

\begin{proof}
The upstream order is inherited and prevents a later return calculation from
repairing a material-route, leverage, curvature, capture, or boundary-current
failure.  On the remaining branch, first separate the scalar-null case.
For a nonnegligible scalar, apply successive subsequence selection to the
nonnegative normalized magnitudes of
$\Delta_n^{\rm old,\phi}$,
$\Delta_n^{\rm rep,\phi}$, and
$\mathfrak B_n^{\rm dau}$.  If the first two tend to zero and the daughter
action tends to zero, apply the same selection to
$|\mathfrak h_n^{\rm birth}|/\mathcal R_n^2$.
A positive lower bound gives branch~\textup{(G22.5)} and
Corollary~\ref{cor:v22-target-gain-escape}.  If it also tends to zero, use
Theorem~\ref{thm:v22-mixed-atlas-residual} separately on the old and
represented components and insert the exact scalar assembly
\eqref{eq:v22-terminal-scalar-assembly}; this gives
\eqref{eq:v22-target-native-pass}.  The final branch is a change of target,
not another subsequence of the scalar problem.  Failure of every finite head
to reach the stated card is exactly the fused all-shell alternative.
\end{proof}

\begin{corollary}[Corrected NS--A to NS--B target-native export]
\label{cor:v22-target-export}
On branch~\textup{(G22.6)}, append the card
$\mathfrak D_n^{\rm V22,\phi}$ to the V18--V20 source and invariant packet.
For the one predeclared critical scalar, this is a complete same-history
source-to-return ancestry certificate.  The following objects are not part of
that scalar export unless independently requested:
\begin{equation}
 \boxed{
 \begin{gathered}
 \text{the complete returned field and its full range,}\\
 \text{the full return Gram, or a pathwise exterior chronology}
 \end{gathered}}
 \label{eq:v22-not-required-objects}
\end{equation}
If NS--B reuses a donor variation bank, the V21 finite return Gram is appended
in addition; no contradiction arises because the targets are different.
\end{corollary}

\begin{proof}
The direct adjoint identities provide the complete scalar ancestry.  The two
mixed residuals certify its target-relative old and represented atlas
incidences, and the vanishing first-birth return excludes an unaccounted
selected follower.  The listed stronger objects contain target-null output
information by Theorem~\ref{thm:v22-minimal-return-operator} and are therefore
not forced by the scalar target.  Reusable-bank semantics invoke the stronger
V21 theorem by definition.
\end{proof}

% =============================================================================
\section{Integrated conclusion, proof ledger, and exact next acquisition}
\label{sec:v22-conclusion}
% =============================================================================

\begin{theorem}[Version V22 target-native return theorem]
\label{thm:v22-integrated}
Preserving Versions~V1--V21 verbatim, Version~V22 proves the following.
\begin{enumerate}[label=\textup{(V22.\arabic*)},leftmargin=5.5em]
\item A single frozen head-supported return writer has the exact backward
      adjoint pullback \eqref{eq:v22-adjoint-writer}, which solves
      \eqref{eq:v22-backward-adjoint} and gives the no-remainder source-to-
      target rectangle \eqref{eq:v22-complete-target-rectangle}.
\item Old exterior memory, represented daughter source, and new daughter
      source assemble on that one target through
      \eqref{eq:v22-three-scalar-assembly}; no output trajectory is repriced
      as a source birth.
\item The positive scalar $\mathfrak g_\phi^2$ is the exact target-visible
      return gain.  The minimum source action for a prescribed target scalar
      is \eqref{eq:v22-minimum-source-action}.
\item Vanishing daughter action with a record-sized selected return forces
      the target-specific gain divergence
      \eqref{eq:v22-target-gain-divergence}.
\item On finite screens, one target depends only on
      $(\mathscr V)^*\phi$.  The unique minimum operator realizing it is rank
      one, with Pythagoras \eqref{eq:v22-operator-Pythagoras}.
\item Countertheorem~\ref{cth:v22-Volterra-target-inflation} proves that the
      complete return Gram may diverge in a target-null channel while the
      selected scalar and adjoint writer remain fixed.
\item Full return-atlas membership is replaced, for one scalar only, by the
      exact mixed residual \eqref{eq:v22-target-atlas-residual}, acquired by
      the positive polarization \eqref{eq:v22-atlas-polarization}.
\item The target-native terminal alternative is
      Theorem~\ref{thm:v22-terminal-gate}.  Its passing branch exports one
      finite scalar ancestry card; the V21 full response card remains active
      only for reusable-bank, range, or path semantics.
\end{enumerate}
No branch is excluded for every smooth datum, and no global regularity
conclusion is drawn.
\end{theorem}

\begin{proof}
Items~\textup{(V22.1)--(V22.2)} are
Theorem~\ref{thm:v22-adjoint-duality} and
Corollary~\ref{cor:v22-source-scalar-split}.  Items
\textup{(V22.3)--(V22.4)} are
Theorem~\ref{thm:v22-target-Riesz},
Corollary~\ref{cor:v22-kernel-gain-bound}, and
Corollary~\ref{cor:v22-target-gain-escape}.  Items
\textup{(V22.5)--(V22.6)} are
Theorem~\ref{thm:v22-minimal-return-operator} and
Countertheorem~\ref{cth:v22-Volterra-target-inflation}.  Item
\textup{(V22.7)} is
Theorem~\ref{thm:v22-mixed-atlas-residual} and
Corollary~\ref{cor:v22-positive-polarization}.  The final item is
Theorem~\ref{thm:v22-terminal-gate} and
Corollary~\ref{cor:v22-target-export}.
\end{proof}

\begingroup\small
\begin{longtable}{p{0.25\textwidth}p{0.19\textwidth}p{0.48\textwidth}}
\caption{NS--A Version V22 proof-status ledger.}\label{tab:v22-ledger}\\
\toprule Row & Status & Exact conclusion\\\midrule
\endfirsthead
\toprule Row & Status & Exact conclusion\\\midrule
\endhead
V1--V21 prefix & \PROVED{} preserved
 &The complete V21 preterminal source is retained byte-for-byte.  No earlier
 theorem, branch, or limitation is rewritten.\\[0.3em]
NS--B material/cell ordering & \INHERITED{} upstream
 &Route mismatch, frequency leverage, strain, material curvature, capture,
 and signed boundary currents are tested before this return scalar.\\[0.3em]
V21 full Volterra map & \INHERITED{} exact / stronger target
 &Retained for complete return, range, reusable-bank, or path semantics; not
 required merely to evaluate one scalar.\\[0.3em]
Backward target writer & \PROVED{} exact
 &$\Psi_\phi$ solves the terminal-zero adjoint equation and is the complete
 pullback of the fixed output writer.\\[0.3em]
Source-to-target rectangle & \PROVED{} exact
 &Old state, background daughter, represented daughter, and first birth
 assemble with no source or response remainder.\\[0.3em]
Target-visible gain & \PROVED{} positive scalar
 &$\|\Psi_\phi\|^2/\|\phi\|^2$ is the complete gain seen by one writer and
 gives the sharp minimum source action.\\[0.3em]
Small birth / large target & \PROVED{} typed survivor
 &Vanishing daughter action with fixed target return forces divergence of the
 one-dimensional causal gain.\\[0.3em]
Full operator from one scalar & \OBSTRUCTION{} explicit
 &The causal Volterra family \eqref{eq:v22-Volterra-family} has fixed target
 functional and divergent full return Gram.\\[0.3em]
Return-atlas incidence & \PROVED{} target relative
 &The mixed residual $\Delta_S(r,\phi)$ is necessary and sufficient for the
 selected scalar; full response membership is stronger.\\[0.3em]
Positive finite acquisition & \PROVED{} exact
 &Every signed target/atlas mixed row is recovered from two same-history
 positive norm reads.\\[0.3em]
Target-native finite-head export & \CONDITIONAL{} exact
 &After the upstream physical writer occurs and the mixed scalar rows pass,
 NS--A may export the scalar card without target-null response data.\\[0.3em]
Reusable donor-bank export & \INHERITED{} V21
 &The complete finite return Gram remains mandatory when several donor
 variations or the response range are intended.\\[0.3em]
All-shell target-return escape & \OPEN{} typed alternative
 &Failure of every finite target card is retained as one fused scalar-return
 escape, not a shell-count sum.\\[0.3em]
Three-dimensional regularity & \OPEN
 &No terminal Dini gain, critical upper action, trace, rigidity, or global
 regularity theorem is inferred.\\
\bottomrule
\end{longtable}
\endgroup

% =============================================================================
\section{Exact mandate after Version V22}
\label{sec:v22-next}
% =============================================================================

\begin{programme}[Read one physical adjoint scalar or stop NS--A]
\label{prog:v22-next}
Preserve Versions~V1--V22 verbatim.  There is no automatic generic
Version~V23.

A continuation is licensed only in the following order.
\begin{enumerate}[label=\textup{(V23.\arabic*)},leftmargin=5.9em]
\item Apply the supplied NS--B material-route, leverage, strain,
      material-curvature, cell-capture, and signed boundary-current gate.
      Stop at its first surviving row.  Do not construct an adjoint writer
      before one physical target writer has occurred.
\item On a collar-free captured branch, freeze one head, one return writer
      $\phi_n$, and its intended scalar normalization before reading the
      return.
\item Solve the one backward adjoint equation
      \eqref{eq:v22-backward-adjoint} and acquire the direct positive reads
      for $\|\Psi_n\|^2$, the old/represented mixed atlas residuals, and the
      first-birth target scalar.
\item If the daughter source action vanishes while the target scalar survives,
      export the target-visible causal-gain packet
      \eqref{eq:v22-target-gain-divergence}.  Do not replace it by the norm of
      the complete response operator.
\item If the target-relative old and represented residuals vanish and the
      first-birth target is negligible, terminate NS--A and export
      \eqref{eq:v22-terminal-card} to NS--B.
\item Populate the V21 full return Gram only if the declared target is a
      reusable donor bank, complete response range, or pathwise returned
      trajectory.  State that stronger target explicitly before computing it.
\item If no finite head gives a convergent target scalar, retain all-shell
      target-return escape.  Reopen it only with a new source-relative
      critical acceleration, capacity, or square-function upper stock on the
      same fused carrier.
\item Do not add another invariant family, path law, age law, Dini modulus,
      concentration metric, shell count, or generic response compiler.
\end{enumerate}
\end{programme}

\begin{assessment}[Programme position after Version V22]
\label{ass:v22-position}
The V21 causal response remains valid, but the first unpopulated NS--A object
has become smaller and more physical.  It is one same-history mixed scalar
between the already occurring daughter history and one frozen return writer,
represented by the backward moment $\Psi_\phi$.  A full response kernel is no
longer a prerequisite for that scalar.  The next admissible work is therefore
an actual value of this finite target card, an upstream material or boundary
survivor, or the stronger reusable-bank packet after that stronger target is
explicitly declared.
\end{assessment}

\paragraph{Version V22 history.}
Preserved the complete V21 source.  Reviewed the new fixed-datum,
Yang--Mills, and Einstein--Standard--Model continuations and adopted their
common upstream correction: target semantics and same-history mixed incidence
precede construction of a stronger transport object.  Derived the exact
terminal-zero adjoint equation for one V21 return writer and the complete
old/source/return scalar rectangle.  Proved the target-native Riesz principle,
the one-dimensional causal-gain alternative, and the rank-one minimum return
operator on finite screens.  Supplied a causal Volterra counterexample with
fixed selected scalar and divergent full Gram.  Replaced full return-atlas
membership by the exact mixed target residual and its positive polarization
read.  Inserted these objects into the V19--V21 terminal head and stated the
corrected scalar-level export, while retaining the V21 full response card for
reusable-bank and path semantics.

\begin{assessment}[Append-only integrity]
\label{ass:v22-integrity}
The complete Version~V21 source preceding its sole final executable document
terminator is preserved verbatim.  Its preterminal SHA--256 digest is
\begin{center}\scriptsize\ttfamily
81e30128585e0da86d6e607d71fd0e0a485125b15a648342ad73595157dbec0a
\end{center}
The present material begins at Section~\ref{sec:v22-scope}; the final command
below is again unique.
\end{assessment}

\addtocontents{toc}{\protect\endgroup}
% =============================================================================
% END VERSION V22 MATERIAL -- append later versions before final terminator
% =============================================================================

% APPEND ALL LATER VERSIONS IMMEDIATELY ABOVE THE SOLE FINAL LINE BELOW.

\clearpage
% =============================================================================
% VERSION V23: NEW MATERIAL.  The complete V1--V22 source above is preserved
% verbatim; only its sole active \end{document} line was removed before append.
% =============================================================================
\hypersetup{pdftitle={NS-A Terminal Source Concentration Programme: Cumulative V23},pdfauthor={A. Pelissier}}
\makeatletter
\addtocontents{toc}{\protect\begingroup\protect\makeatletter}
\addtocontents{toc}{\protect\def\protect\l@section{\protect\@dottedtocline{1}{0em}{4.7em}}}
\addtocontents{toc}{\protect\def\protect\l@subsection{\protect\@dottedtocline{2}{1.5em}{5.3em}}}
\addtocontents{toc}{\protect\makeatother}
\makeatother
\renewcommand{\thesection}{V23.\arabic{section}}
\renewcommand{\thesubsection}{V23.\arabic{section}.\arabic{subsection}}
\renewcommand{\theequation}{V23.\arabic{equation}}
\setcounter{section}{0}
\setcounter{equation}{0}
\renewcommand{\theHsection}{V23.\arabic{section}}
\renewcommand{\theHsubsection}{V23.\arabic{section}.\arabic{subsection}}
\renewcommand{\theHtheorem}{V23.\arabic{section}.\arabic{theorem}}
\renewcommand{\theHequation}{V23.\arabic{equation}}

\begin{center}
{\LARGE\bfseries NS--A: Version V23}\\[0.5em]
{\large Source-Cyclic Pullback of the Adjoint Return, the Physical Daughter Gain,\\
and the Parent-Pair Visibility Alternative}\\[0.5em]
{\normalsize 3 September 2026}
\end{center}
\addcontentsline{toc}{part}{Version V23: physical daughter range and source-cyclic adjoint gain}

\begin{abstract}
Version~V22 correctly replaced the complete daughter-to-return operator by
one backward adjoint writer when the declared output is one frozen return
scalar.  Its Riesz theorem, however, minimized over the entire ambient
exterior source space.  The actual first-birth source on a finite
Navier--Stokes donor head is much smaller: it lies in the range of the
quadratic parent-to-exterior daughter synthesis, after the represented source
atlas has been removed.  An adjoint component orthogonal to that range cannot
be excited by any physical daughter of the declared head and must not improve
the physical source-action minimum.

The present continuation moves one arrow upstream and performs that
source-cyclic short.  If
\(
 C_t^{\rm birth}=(I-R_t^{\rm out})\mathscr C_{\mathscr A}^{\rm out}
\)
is the occurring first-birth synthesis, then
\[
 D_t^{\rm birth}=(C_t^{\rm birth})^*C_t^{\rm birth},
 \qquad
 P_t^{\rm birth}
 =C_t^{\rm birth}(D_t^{\rm birth})^\dagger
  (C_t^{\rm birth})^*
\]
is the canonical projection onto the instantaneous physical daughter range.
For the V22 adjoint writer \(\Psi_\phi\), only
\(\Psi_\phi^{\rm birth}=P_t^{\rm birth}\Psi_\phi\) contributes to the
first-birth scalar.  Writing
\(H_{\phi,t}^{\rm birth}=(C_t^{\rm birth})^*\Psi_\phi(t)\), one has the
exact finite identity
\[
 \|\Psi_\phi^{\rm birth}(t)\|_2^2
 =\left\langle H_{\phi,t}^{\rm birth},
 (D_t^{\rm birth})^\dagger H_{\phi,t}^{\rm birth}\right\rangle.
\]
This is the source-native target gain.  It gives the sharp minimum action
among sources which the actual daughter bank can produce.  The full V22
adjoint norm splits orthogonally into this physical gain and a daughter-null
component.  A causal Volterra example has divergent ambient adjoint norm and
constant physical daughter gain, so the ambient norm is not itself a
physical causal-amplification certificate.

The daughter scalar is also pulled back to one finite symmetric parent-pair
writer.  For the occurring head state \(v(t)\),
\[
 \mathfrak h_\phi^{\rm birth}
 =\int\left\langle v(t)\odot v(t),
 H_{\phi,t}^{\rm birth}\right\rangle\,dt.
\]
The same finite parent Gram supplies the exact alignment factor between the
occurring quadratic source and the target-visible adjoint.  If the normalized
daughter action tends to zero while the selected return remains of record
size, the \emph{projected} daughter gain diverges.  This divergence has an
exhaustive finite interpretation: either the target-weighted exterior source
metric softens or the parent-pair adjoint coefficient itself overdrives.
Ambient adjoint growth in the daughter-null complement is removed from the
scalar card.

Thus Version~V23 does not construct another response object.  It corrects the
source class of the V22 Riesz problem and reduces its surviving gain branch
to one finite parent-output Gram and one finite target vector on the already
selected head.  The upstream material-route, frequency-leverage, strain,
material-curvature, capture, and signed boundary-current gates remain prior.
No terminal Dini improvement, critical upper-action theorem, rigidity result,
or global three-dimensional regularity conclusion is claimed.
\end{abstract}

% =============================================================================
\section{Scope: source occurrence precedes an ambient adjoint norm}
\label{sec:v23-scope}
% =============================================================================

Retain the V22 terminal branch only after the inherited NS--B material--heat
route, ultra-canonical leverage, symmetric-strain, material-map-curvature,
cell-capture, and signed viscous/Bernoulli boundary-current gates have passed.
Thus one finite head \(\mathscr A\), one represented exterior source atlas,
and one physical return writer \(\phi\) have already been frozen on one
common history.

Version~V22 proves the exact first-birth scalar identity
\begin{equation}
 \mathfrak h_\phi^{\rm birth}
 =\int_I\left\langle d^{\rm birth}(t),
                  \Psi_\phi(t)\right\rangle\,\dd t,
 \label{eq:v23-inherited-birth-scalar}
\end{equation}
where \(\Psi_\phi\) is the terminal-zero adjoint writer of
\eqref{eq:v22-backward-adjoint}.  The identity is retained unchanged.
The correction concerns the positive quantity used to price it.

The ambient norm \(\|\Psi_\phi\|_{L^2(I;L^2)}\) is the Riesz norm of the
selected functional on every exterior source history.  But the occurring
source \(d^{\rm birth}\) is generated by the finite quadratic daughter map of
Version~V20.  The physical source class is therefore a proper subspace unless
that daughter map is onto the complete exterior screen.  The Riesz writer
must be shorted to this class before its norm is interpreted as physical
gain.

\begin{firewall}[The source projection is not a new source]
\label{fire:v23-projection-not-source}
The projection constructed below acts on the mathematical adjoint writer.
It does not create a new Navier--Stokes forcing, add a daughter occurrence,
replace the represented source atlas, or alter the exact V21--V22 causal
factorization.  It removes source directions which the declared quadratic
head cannot excite.  A stronger exterior-source target may retain the full
adjoint writer, but that stronger target is not charged to the one daughter
scalar.
\end{firewall}

% =============================================================================
\section{The physical first-birth daughter range}
\label{sec:v23-daughter-range}
% =============================================================================

Let \(E_{\mathscr A}\) be the real amplitude space of the finite
reality-paired helical head, identified with its synthesized head fields, and
put
\begin{equation}
 \mathcal Q_{\mathscr A}:=\operatorname{Sym}^2E_{\mathscr A}.
 \label{eq:v23-parent-space}
\end{equation}
Retain the exterior synthesis of Version~V20,
\begin{equation}
 \mathscr C_{\mathscr A}^{\rm out}:\mathcal Q_{\mathscr A}
 \longrightarrow\operatorname{Ran}Q_{\mathscr A},
 \qquad
 \mathscr C_{\mathscr A}^{\rm out}(a\odot b)
 =A^{-1/4}Q_{\mathscr A}\mathcal B_{\rm s}(a,b).
 \label{eq:v23-inherited-exterior-synthesis}
\end{equation}
Let \(R_t^{\rm out}\) be the orthogonal projection onto the represented
exterior source atlas at time \(t\), transported in the same physical source
metric.  Define the first-birth synthesis
\begin{equation}
 \boxed{
 C_t^{\rm b}:=(I-R_t^{\rm out})
              \mathscr C_{\mathscr A}^{\rm out}.}
 \label{eq:v23-birth-synthesis}
\end{equation}
Then the occurring first-birth daughter has the exact form
\begin{equation}
 \boxed{
 d^{\rm birth}(t)=C_t^{\rm b}\bigl(v(t)\odot v(t)\bigr).}
 \label{eq:v23-occurring-birth}
\end{equation}
Set
\begin{equation}
 D_t^{\rm b}:=(C_t^{\rm b})^*C_t^{\rm b}\succeq0,
 \qquad
 P_t^{\rm b}:=C_t^{\rm b}(D_t^{\rm b})^\dagger(C_t^{\rm b})^*.
 \label{eq:v23-birth-Gram-projection}
\end{equation}

\begin{theorem}[Canonical support of the physical daughter source]
\label{thm:v23-daughter-support}
For every time \(t\), \(P_t^{\rm b}\) is the orthogonal projection onto
\(\operatorname{Ran}C_t^{\rm b}\).  For every exterior writer \(\psi\), put
\begin{equation}
 H_t^{\rm b}[\psi]:=(C_t^{\rm b})^*\psi,
 \qquad
 \psi_t^{\rm b}:=P_t^{\rm b}\psi.
 \label{eq:v23-target-vector-projected-writer}
\end{equation}
Then
\begin{align}
 H_t^{\rm b}[\psi]&\in\operatorname{Ran}D_t^{\rm b},
 \label{eq:v23-target-in-Gram-range}\\
 \boxed{
 \langle C_t^{\rm b}q,\psi\rangle
 =\langle C_t^{\rm b}q,\psi_t^{\rm b}\rangle
 =\langle q,H_t^{\rm b}[\psi]\rangle}
 \qquad(q\in\mathcal Q_{\mathscr A}),
 \label{eq:v23-source-target-three-way}\\
 \boxed{
 \|\psi_t^{\rm b}\|_2^2
 =\left\langle H_t^{\rm b}[\psi],
  (D_t^{\rm b})^\dagger H_t^{\rm b}[\psi]\right\rangle,}
 \label{eq:v23-projected-gain-finite}\\
 \boxed{
 \|\psi\|_2^2
 =\|\psi_t^{\rm b}\|_2^2
  +\|(I-P_t^{\rm b})\psi\|_2^2.}
 \label{eq:v23-adjoint-source-Pythagoras}
\end{align}
The last term is invisible to every first-birth daughter generated by the
frozen head.
\end{theorem}

\begin{proof}
Write \(C=C_t^{\rm b}\) and \(D=C^*C\).  The Moore--Penrose equations give
\[
 (CD^\dagger C^*)^*=CD^\dagger C^*,
 \qquad
 (CD^\dagger C^*)^2
 =CD^\dagger DD^\dagger C^*
 =CD^\dagger C^*.
\]
Thus \(P=CD^\dagger C^*\) is an orthogonal projection with range contained
in \(\operatorname{Ran}C\).  If \(y=Cq\), then
\[
 Py=CD^\dagger Dq=Cq,
\]
because \(D^\dagger D\) projects onto
\((\Ker C)^\perp\) and \(C\) annihilates \(\Ker C\).  Hence
\(\operatorname{Ran}P=\operatorname{Ran}C\).

In finite dimension,
\(\operatorname{Ran}C^*=\operatorname{Ran}(C^*C)=\operatorname{Ran}D\),
which proves \eqref{eq:v23-target-in-Gram-range}.  Since \(PC=C\),
\[
 \langle Cq,\psi\rangle
 =\langle Cq,P\psi\rangle
 =\langle q,C^*\psi\rangle,
\]
giving \eqref{eq:v23-source-target-three-way}.  Finally,
\[
 \|P\psi\|^2
 =\langle\psi,P\psi\rangle
 =\langle C^*\psi,D^\dagger C^*\psi\rangle,
\]
and orthogonal Pythagoras gives
\eqref{eq:v23-adjoint-source-Pythagoras}.
\end{proof}

\begin{corollary}[Decomposable history projection]
\label{cor:v23-history-projection}
Assume the finite matrices \(R_t^{\rm out}\) are strongly measurable on
\(I\).  Then \(t\mapsto P_t^{\rm b}\) is measurable and
\begin{equation}
 (\mathbf P^{\rm b}f)(t):=P_t^{\rm b}f(t)
 \label{eq:v23-history-projection}
\end{equation}
is an orthogonal projection on \(L^2(I;\operatorname{Ran}Q_{\mathscr A})\).
Its range
\begin{equation}
 \mathcal D_I^{\rm b}
 :=\{f:P_t^{\rm b}f(t)=f(t)\text{ for a.e. }t\}
 \label{eq:v23-physical-history-class}
\end{equation}
is the closed physical first-birth source class.  For the V22 adjoint writer,
put
\begin{equation}
 \boxed{
 \Psi_\phi^{\rm b}(t):=P_t^{\rm b}\Psi_\phi(t),
 \qquad
 H_{\phi,t}^{\rm b}:=(C_t^{\rm b})^*\Psi_\phi(t).}
 \label{eq:v23-birth-adjoint-and-parent}
\end{equation}
Then
\begin{align}
 \boxed{
 \mathfrak h_\phi^{\rm birth}
 =\int_I\langle d^{\rm birth}(t),
                  \Psi_\phi^{\rm b}(t)\rangle\,\dd t,}
 \label{eq:v23-birth-scalar-projected}\\
 \boxed{
 \|\Psi_\phi^{\rm b}\|_{L^2}^2
 =\int_I\left\langle H_{\phi,t}^{\rm b},
 (D_t^{\rm b})^\dagger H_{\phi,t}^{\rm b}\right\rangle\,\dd t.}
 \label{eq:v23-history-gain-finite}
\end{align}
\end{corollary}

\begin{proof}
The Moore--Penrose inverse is a Borel function of a finite matrix, so the
pointwise projections are measurable.  Their operator norms are at most one;
hence multiplication by them defines a bounded self-adjoint idempotent on the
history space.  Its range is therefore closed.  Apply
Theorem~\ref{thm:v23-daughter-support} pointwise to
\eqref{eq:v23-inherited-birth-scalar} and integrate.
\end{proof}

% =============================================================================
\section{The source-native target gain and restricted Riesz principle}
\label{sec:v23-physical-gain}
% =============================================================================

\begin{definition}[Physical daughter gain]
\label{def:v23-physical-gain}
For \(\phi\ne0\), define
\begin{equation}
 \boxed{
 (\mathfrak g_\phi^{\rm b})^2
 :=\frac{\|\Psi_\phi^{\rm b}\|_{L^2}^2}
          {\|\phi\|_{L^2}^2}.}
 \label{eq:v23-physical-gain}
\end{equation}
This is the gain of the selected return functional restricted to sources
which the declared first-birth daughter bank can produce.
\end{definition}

\begin{theorem}[Source-native Thomson--Riesz principle]
\label{thm:v23-physical-Riesz}
Let \(\tau\in\mathbb R\).  If \(\Psi_\phi^{\rm b}\ne0\), then
\begin{equation}
 \boxed{
 \inf\left\{
  \|f\|_{L^2}^2:
  f\in\mathcal D_I^{\rm b},\ 
  \int_I\langle f,\Psi_\phi\rangle\,\dd t=\tau
 \right\}
 =\frac{\tau^2}{\|\Psi_\phi^{\rm b}\|_{L^2}^2}.}
 \label{eq:v23-physical-minimum-action}
\end{equation}
The unique minimizing source vector is
\begin{equation}
 \boxed{
 f_{\min}^{\rm b}
 =\frac{\tau}{\|\Psi_\phi^{\rm b}\|_{L^2}^2}
  \Psi_\phi^{\rm b}.}
 \label{eq:v23-physical-minimizer}
\end{equation}
In parent-pair coordinates a minimum representative, unique modulo
\(\Ker C_t^{\rm b}\), is
\begin{equation}
 q_{\min}(t)
 =\frac{\tau}{\|\Psi_\phi^{\rm b}\|_{L^2}^2}
  (D_t^{\rm b})^\dagger H_{\phi,t}^{\rm b}.
 \label{eq:v23-parent-minimizer}
\end{equation}
If \(\Psi_\phi^{\rm b}=0\), every physical first-birth source has zero
selected return and a nonzero \(\tau\) is infeasible.
\end{theorem}

\begin{proof}
For \(f\in\mathcal D_I^{\rm b}\),
\[
 \int_I\langle f,\Psi_\phi\rangle
 =\int_I\langle f,\Psi_\phi^{\rm b}\rangle.
\]
Cauchy--Schwarz therefore gives
\(
 |\tau|^2\le\|f\|_{L^2}^2\|\Psi_\phi^{\rm b}\|_{L^2}^2
\), with equality exactly when \(f\) is the scalar multiple in
\eqref{eq:v23-physical-minimizer}.  This vector belongs to
\(\mathcal D_I^{\rm b}\).

For \eqref{eq:v23-parent-minimizer}, use
\[
 C_t^{\rm b}(D_t^{\rm b})^\dagger H_{\phi,t}^{\rm b}
 =P_t^{\rm b}\Psi_\phi(t)
 =\Psi_\phi^{\rm b}(t).
\]
Moreover,
\[
 \int_I\langle q_{\min},D_t^{\rm b}q_{\min}\rangle\,\dd t
 =\frac{\tau^2}{\|\Psi_\phi^{\rm b}\|_{L^2}^4}
  \int_I\langle H_{\phi,t}^{\rm b},
     (D_t^{\rm b})^\dagger H_{\phi,t}^{\rm b}\rangle\,\dd t,
\]
which equals the minimum in
\eqref{eq:v23-physical-minimum-action}.  Adding a vector in
\(\Ker C_t^{\rm b}=\Ker D_t^{\rm b}\) changes neither the source nor its
action.  The zero case follows from
\eqref{eq:v23-source-target-three-way}.
\end{proof}

\begin{corollary}[Ambient gain Pythagoras]
\label{cor:v23-ambient-gain-Pythagoras}
The V22 gain has the exact split
\begin{equation}
 \boxed{
 \mathfrak g_\phi^2
 =(\mathfrak g_\phi^{\rm b})^2
  +(\mathfrak g_\phi^{\rm b,null})^2,}
 \label{eq:v23-gain-Pythagoras}
\end{equation}
where
\begin{equation}
 (\mathfrak g_\phi^{\rm b,null})^2
 :=\frac{
  \|(I-\mathbf P^{\rm b})\Psi_\phi\|_{L^2}^2
 }{\|\phi\|_{L^2}^2}.
 \label{eq:v23-null-gain}
\end{equation}
The second term cannot lower the action of any source in
\(\mathcal D_I^{\rm b}\).
\end{corollary}

\begin{proof}
Integrate \eqref{eq:v23-adjoint-source-Pythagoras} and divide by
\(\|\phi\|_{L^2}^2\).  The minimum-action statement is
Theorem~\ref{thm:v23-physical-Riesz}.
\end{proof}

\begin{countertheorem}[Ambient adjoint growth may be daughter null]
\label{cth:v23-ambient-gain-null}
Let \(I=[0,1]\),
\(\mathcal X=L^2(I;\mathbb R^2)\),
\(\mathcal Y=L^2(I;\mathbb R)\), and define the causal Volterra operators
\begin{equation}
 (V_Nf)(t)=\int_0^t\bigl(f_1(s)+Nf_2(s)\bigr)\,\dd s.
 \label{eq:v23-source-null-Volterra}
\end{equation}
For the fixed target writer \(\phi\equiv1\),
\begin{equation}
 V_N^*\phi(s)=(1-s)(1,N).
 \label{eq:v23-source-null-adjoint}
\end{equation}
Declare the physical daughter class
\begin{equation}
 \mathcal D^{\rm b}
 =\{(q,0):q\in L^2(I)\}.
 \label{eq:v23-example-source-class}
\end{equation}
Then
\begin{equation}
 \boxed{
 \|V_N^*\phi\|_{L^2}^2=\frac{1+N^2}{3}\longrightarrow\infty,
 \qquad
 \|P_{\mathcal D^{\rm b}}V_N^*\phi\|_{L^2}^2=\frac13.}
 \label{eq:v23-source-null-divergence}
\end{equation}
Every selected return generated by a physical daughter is independent of
\(N\).  Hence divergence of the full adjoint norm need not be a physical
daughter-gain survivor.
\end{countertheorem}

\begin{proof}
Fubini gives
\[
 \int_0^1(V_Nf)(t)\,\dd t
 =\int_0^1(1-s)\bigl(f_1(s)+Nf_2(s)\bigr)\,\dd s,
\]
which proves \eqref{eq:v23-source-null-adjoint}.  Projection onto
\eqref{eq:v23-example-source-class} deletes the second coordinate.  Since
\(\int_0^1(1-s)^2\,\dd s=1/3\), the two norms are as stated.  For
\(f=(q,0)\), the output contains no \(N\)-dependent term.
\end{proof}

\begin{countertheorem}[Raw parent coefficients do not measure source gain]
\label{cth:v23-raw-parent-no-gain}
At one time let \(C_\varepsilon:\mathbb R\to\mathbb R\) be
\(C_\varepsilon q=\varepsilon q\), with \(\varepsilon>0\), and let
\(\psi=1\).  Then
\begin{equation}
 D_\varepsilon=\varepsilon^2,
 \qquad
 H_\varepsilon=C_\varepsilon^*\psi=\varepsilon,
 \label{eq:v23-soft-example-data}
\end{equation}
so \(|H_\varepsilon|\to0\), while
\begin{equation}
 \boxed{
 \langle H_\varepsilon,D_\varepsilon^\dagger H_\varepsilon\rangle=1
 =\|P_{\operatorname{Ran}C_\varepsilon}\psi\|^2.}
 \label{eq:v23-soft-example-gain}
\end{equation}
Thus the parent target vector must be measured in the Moore--Penrose source
metric; its Euclidean norm alone is not a physical visibility coefficient.
\end{countertheorem}

\begin{proof}
All identities are direct substitutions.  Since
\(D_\varepsilon^\dagger=\varepsilon^{-2}\), the weighted value is one.
\end{proof}

% =============================================================================
\section{Exact alignment of the occurring quadratic source}
\label{sec:v23-alignment}
% =============================================================================

Put
\begin{equation}
 \mathcal A_\phi^{\rm b}:=
 \|\Psi_\phi^{\rm b}\|_{L^2}^2,
 \qquad
 \mathcal B^{\rm b}:=
 \|d^{\rm birth}\|_{L^2}^2,
 \qquad
 \tau_\phi^{\rm b}:=
 \mathfrak h_\phi^{\rm birth}.
 \label{eq:v23-three-physical-scalars}
\end{equation}
When both positive factors are nonzero, define the actual source--target
alignment
\begin{equation}
 \boxed{
 \eta_\phi^{\rm b}
 :=\frac{|\tau_\phi^{\rm b}|^2}
 {\mathcal B^{\rm b}\mathcal A_\phi^{\rm b}}.}
 \label{eq:v23-alignment-efficiency}
\end{equation}
Set \(\eta_\phi^{\rm b}=0\) when the numerator vanishes and one denominator
factor is zero.

\begin{theorem}[Exact action--gain--alignment factorization]
\label{thm:v23-alignment-factorization}
One has
\begin{equation}
 \boxed{0\le\eta_\phi^{\rm b}\le1,}
 \label{eq:v23-alignment-window}
\end{equation}
and, whenever the denominator is nonzero,
\begin{equation}
 \boxed{
 |\tau_\phi^{\rm b}|^2
 =\mathcal B^{\rm b}\mathcal A_\phi^{\rm b}
  \eta_\phi^{\rm b}.}
 \label{eq:v23-exact-alignment-factorization}
\end{equation}
Equality \(\eta_\phi^{\rm b}=1\) holds exactly when the occurring first-birth
history is a scalar multiple of \(\Psi_\phi^{\rm b}\) in
\(L^2(I;L^2)\).
\end{theorem}

\begin{proof}
Equation~\eqref{eq:v23-birth-scalar-projected} and Cauchy--Schwarz give
\(
 |\tau_\phi^{\rm b}|^2
 \le\mathcal B^{\rm b}\mathcal A_\phi^{\rm b}
\).  This proves the window and the displayed factorization by definition.
Equality in Cauchy--Schwarz gives the final assertion.
\end{proof}

Return to the terminal normalization of V22.  Define
\begin{equation}
 \boxed{
 \mathfrak G_n^{\rm b,phys}
 :=\frac{\nu}{\mathcal R_n^2}
   \|\Psi_n^{\rm b}\|_{L^2(J_n)}^2,}
 \label{eq:v23-normalized-physical-gain}
\end{equation}
where \(\Psi_n^{\rm b}=P_{n,t}^{\rm b}\Psi_n\).  Retain the V20 normalized
daughter action \(\mathfrak B_n^{\rm dau}\).

\begin{corollary}[Corrected first-birth gain law]
\label{cor:v23-corrected-gain-law}
The terminal first-birth scalar satisfies the exact identity
\begin{equation}
 \boxed{
 \left|
  \frac{\mathfrak h_n^{\rm birth}}{\mathcal R_n^2}
 \right|^2
 =\mathfrak B_n^{\rm dau}
  \mathfrak G_n^{\rm b,phys}
  \eta_n^{\rm b}.}
 \label{eq:v23-terminal-exact-product}
\end{equation}
In particular, if
\begin{equation}
 \mathfrak B_n^{\rm dau}\longrightarrow0,
 \qquad
 \frac{|\mathfrak h_n^{\rm birth}|}{\mathcal R_n^2}
 \ge\delta>0,
 \label{eq:v23-small-birth-fixed-target}
\end{equation}
then
\begin{equation}
 \boxed{
 \mathfrak G_n^{\rm b,phys}
 \ge\frac{\delta^2}{\mathfrak B_n^{\rm dau}}
 \longrightarrow\infty.}
 \label{eq:v23-physical-gain-divergence}
\end{equation}
The divergence occurs inside the physical first-birth source range.  The
ambient daughter-null term in \eqref{eq:v23-gain-Pythagoras} is irrelevant to
this conclusion.
\end{corollary}

\begin{proof}
Substitute
\(
 \mathcal B^{\rm b}=\nu\mathcal R_n^2\mathfrak B_n^{\rm dau}
\)
and
\(
 \mathcal A_\phi^{\rm b}
 =\mathcal R_n^2\mathfrak G_n^{\rm b,phys}/\nu
\)
into \eqref{eq:v23-exact-alignment-factorization}.  Since
\(\eta_n^{\rm b}\le1\), the lower bound follows.
\end{proof}

% =============================================================================
\section{The finite parent-pair target writer}
\label{sec:v23-parent-writer}
% =============================================================================

\begin{definition}[Parent-pair pullback of one return target]
\label{def:v23-parent-writer}
For almost every time define the symmetric bilinear form
\begin{equation}
 \mathcal H_{\phi,t}^{\rm b}(a,b)
 :=\left\langle
  C_t^{\rm b}(a\odot b),\Psi_\phi(t)
 \right\rangle.
 \label{eq:v23-parent-bilinear-form}
\end{equation}
Its coefficient vector in \(\mathcal Q_{\mathscr A}\) is exactly
\begin{equation}
 H_{\phi,t}^{\rm b}=(C_t^{\rm b})^*\Psi_\phi(t).
 \label{eq:v23-parent-vector-again}
\end{equation}
Let \(\mathsf R_{\phi,t}^{\rm b}\) be the self-adjoint operator on
\(E_{\mathscr A}\) representing this bilinear form:
\begin{equation}
 \mathcal H_{\phi,t}^{\rm b}(a,b)
 =\langle a,\mathsf R_{\phi,t}^{\rm b}b\rangle.
 \label{eq:v23-parent-selfadjoint}
\end{equation}
\end{definition}

\begin{theorem}[Exact quadratic parent read]
\label{thm:v23-quadratic-parent-read}
The occurring first-birth scalar is
\begin{equation}
 \boxed{
 \mathfrak h_\phi^{\rm birth}
 =\int_I
  \left\langle v(t),
   \mathsf R_{\phi,t}^{\rm b}v(t)\right\rangle\,\dd t
 =\int_I
  \left\langle v(t)\odot v(t),
   H_{\phi,t}^{\rm b}\right\rangle\,\dd t.}
 \label{eq:v23-parent-quadratic-read}
\end{equation}
The physical daughter gain is
\begin{equation}
 \boxed{
 \|\Psi_\phi^{\rm b}\|_{L^2}^2
 =\int_I\left\langle H_{\phi,t}^{\rm b},
 (D_t^{\rm b})^\dagger H_{\phi,t}^{\rm b}\right\rangle\,\dd t.}
 \label{eq:v23-parent-gain}
\end{equation}
For any \(a,b\in E_{\mathscr A}\), if
\begin{equation}
 T_{\phi,t}(z)
 :=\left\langle C_t^{\rm b}(z\odot z),
                  \Psi_\phi(t)\right\rangle,
 \label{eq:v23-diagonal-parent-query}
\end{equation}
then
\begin{equation}
 \boxed{
 4\mathcal H_{\phi,t}^{\rm b}(a,b)
 =T_{\phi,t}(a+b)-T_{\phi,t}(a-b).}
 \label{eq:v23-parent-polarization}
\end{equation}
Thus one occurring scalar needs only the diagonal value at \(v(t)\); the
complete parent matrix is required only when a reusable donor-variation bank
is the declared target.
\end{theorem}

\begin{proof}
Use \eqref{eq:v23-occurring-birth} in
\eqref{eq:v23-inherited-birth-scalar}, then apply
\eqref{eq:v23-source-target-three-way}.  This gives both expressions in
\eqref{eq:v23-parent-quadratic-read}.  Equation
\eqref{eq:v23-parent-gain} is
\eqref{eq:v23-history-gain-finite}.  Finally,
\[
 (a+b)\odot(a+b)-(a-b)\odot(a-b)=4a\odot b,
\]
and linearity of \(C_t^{\rm b}\) proves
\eqref{eq:v23-parent-polarization}.
\end{proof}

\begin{assessment}[Literal finite coefficients]
\label{ass:v23-literal-coefficients}
In the helical parent and exterior bases of Version~V20,
\(D_t^{\rm b}\) is the positive Gram of the represented-shorted coefficient
matrix and
\[
 H_{\phi,t}^{\rm b}=(C_t^{\rm b})^*\Psi_\phi(t).
\]
Hence the source-native gain is a finite Moore--Penrose quadratic form in the
actual parent-output coefficients.  No new shell theorem, response graph, or
ambient source norm is needed to populate it.
\end{assessment}

% =============================================================================
\section{Source softness versus parent-target coefficient overdrive}
\label{sec:v23-softness}
% =============================================================================

Define the normalized parent-target coefficient stock
\begin{equation}
 \mathfrak H_n^{\rm par}
 :=\frac{\nu}{\mathcal R_n^2}
   \int_{J_n}\|H_{n,t}^{\rm b}\|^2\,\dd t.
 \label{eq:v23-parent-coefficient-stock}
\end{equation}
When the denominator is nonzero, define the target-weighted inverse source
metric
\begin{equation}
 \boxed{
 \mathfrak S_n^{\rm src}
 :=\frac{
  \displaystyle\int_{J_n}
   \left\langle H_{n,t}^{\rm b},
    (D_{n,t}^{\rm b})^\dagger H_{n,t}^{\rm b}\right\rangle\,\dd t
 }{
  \displaystyle\int_{J_n}\|H_{n,t}^{\rm b}\|^2\,\dd t
 }.}
 \label{eq:v23-target-source-softness}
\end{equation}
Set \(\mathfrak S_n^{\rm src}=0\) when the denominator is zero.

\begin{theorem}[Exact finite gain factorization]
\label{thm:v23-finite-gain-factorization}
If \(\mathfrak H_n^{\rm par}>0\), then
\begin{equation}
 \boxed{
 \mathfrak G_n^{\rm b,phys}
 =\mathfrak H_n^{\rm par}\mathfrak S_n^{\rm src}.}
 \label{eq:v23-gain-softness-factorization}
\end{equation}
Consequently, if
\(\mathfrak G_n^{\rm b,phys}\to\infty\), then after passage to a
subsequence at least one of the following holds:
\begin{equation}
 \boxed{
 \mathfrak H_n^{\rm par}\longrightarrow\infty
 \quad\text{or}\quad
 \mathfrak S_n^{\rm src}\longrightarrow\infty.}
 \label{eq:v23-overdrive-or-softness}
\end{equation}
The first branch is parent-pair adjoint coefficient overdrive.  The second is
source-metric softness in the target direction.

If for some \(\mu>0\) the finite daughter Grams obey
\begin{equation}
 D_{n,t}^{\rm b}\succeq
 \mu\,\operatorname{supp}D_{n,t}^{\rm b}
 \qquad\text{for a.e. }t\in J_n,
 \label{eq:v23-uniform-source-floor}
\end{equation}
then
\begin{equation}
 \boxed{\mathfrak S_n^{\rm src}\le\mu^{-1}.}
 \label{eq:v23-softness-floor-bound}
\end{equation}
Thus the softness branch excludes every uniform source-visibility floor on
the target-carrying parent range.
\end{theorem}

\begin{proof}
The factorization is obtained by multiplying and dividing
\eqref{eq:v23-normalized-physical-gain}--
\eqref{eq:v23-history-gain-finite} by
\(\int\|H_{n,t}^{\rm b}\|^2\).  If both factors on the right of
\eqref{eq:v23-gain-softness-factorization} stayed bounded, their product
would stay bounded, proving the alternative after subsequence selection.
Under \eqref{eq:v23-uniform-source-floor}, the Moore--Penrose inverse obeys
\[
 (D_{n,t}^{\rm b})^\dagger
 \preceq\mu^{-1}\operatorname{supp}D_{n,t}^{\rm b}.
\]
Since \(H_{n,t}^{\rm b}\in\operatorname{Ran}D_{n,t}^{\rm b}\), integration
gives \eqref{eq:v23-softness-floor-bound}.
\end{proof}

\begin{firewall}[Softness is not a new invariant search]
\label{fire:v23-softness-not-invariant}
The source-softness branch refers only to the already constructed V20
exterior daughter Gram after the represented atlas short.  It does not reopen
the V19 invariant kernel, enlarge the helical head after inspecting the
target, or introduce a new ambient source metric.  The parent-coefficient
branch refers only to the finite vector
\((C_t^{\rm b})^*\Psi_\phi(t)\) on that same head.
\end{firewall}

% =============================================================================
\section{Nested source classes and the exact no-target-inflation order}
\label{sec:v23-source-hierarchy}
% =============================================================================

\begin{theorem}[Monotonicity under a declared source class]
\label{thm:v23-source-class-monotonicity}
Let \(\mathcal D_0\subseteq\mathcal D_1\subseteq\mathcal X\) be closed
source subspaces with orthogonal projections \(P_0,P_1\).  For any target
writer \(\psi\),
\begin{equation}
 \boxed{
 \|P_0\psi\|^2
 \le\|P_1\psi\|^2
 \le\|\psi\|^2.}
 \label{eq:v23-gain-monotonicity}
\end{equation}
For a prescribed nonzero scalar \(\tau\), the corresponding minimum actions
satisfy
\begin{equation}
 \boxed{
 \frac{\tau^2}{\|P_0\psi\|^2}
 \ge
 \frac{\tau^2}{\|P_1\psi\|^2}
 \ge
 \frac{\tau^2}{\|\psi\|^2},}
 \label{eq:v23-action-monotonicity}
\end{equation}
with the convention that division by zero means infeasibility.

For the occurring daughter history \(d^{\rm birth}\ne0\), take
\begin{equation}
 \mathcal D_0=\operatorname{span}\{d^{\rm birth}\},
 \qquad
 \mathcal D_1=\mathcal D_I^{\rm b}.
 \label{eq:v23-nested-physical-classes}
\end{equation}
Then
\begin{equation}
 \boxed{
 \frac{|\mathfrak h_\phi^{\rm birth}|^2}
      {\|d^{\rm birth}\|_{L^2}^2}
 \le
 \|\Psi_\phi^{\rm b}\|_{L^2}^2
 \le
 \|\Psi_\phi\|_{L^2}^2.}
 \label{eq:v23-three-gain-levels}
\end{equation}
The three levels correspond respectively to one occurring source history,
the complete physical daughter range of the frozen head, and arbitrary
exterior sources.
\end{theorem}

\begin{proof}
Nested orthogonal projections satisfy
\(P_1-P_0\succeq0\) and
\[
 \|P_1\psi\|^2
 =\|P_0\psi\|^2+
  \|(P_1-P_0)\psi\|^2.
\]
This proves \eqref{eq:v23-gain-monotonicity}; the action inequalities are
its reciprocal form and Theorem~\ref{thm:v23-physical-Riesz}.  For
\eqref{eq:v23-three-gain-levels}, the projection of \(\Psi_\phi\) onto the
line generated by \(d^{\rm birth}\) has squared norm
\(
 |\langle d^{\rm birth},\Psi_\phi\rangle|^2/
 \|d^{\rm birth}\|^2
\), and
\eqref{eq:v23-birth-scalar-projected} identifies the numerator.
\end{proof}

\begin{assessment}[The admissible target hierarchy]
\label{ass:v23-target-hierarchy}
For one already occurring scalar, the mixed read
\(\mathfrak h_\phi^{\rm birth}\) is sufficient.  A robustness statement over
all quadratic daughters of the frozen head uses
\(\Psi_\phi^{\rm b}\) and the V20 source Gram.  Arbitrary exterior-source
control uses the complete V22 adjoint writer.  A reusable output bank or full
returned path uses the V21 return Gram.  These are successively stronger
targets and are not interchangeable.
\end{assessment}

\begin{corollary}[Pruning gain outside the occurring source line]
\label{cor:v23-target-null-pruning}
Suppose
\begin{equation}
 \frac{|\mathfrak h_n^{\rm birth}|}{\mathcal R_n^2}\longrightarrow0.
 \label{eq:v23-occurring-line-null}
\end{equation}
Then no divergence of \(\mathfrak G_n^{\rm b,phys}\) or of the full V22 gain
can obstruct ancestry of this one scalar.  Such divergence lies outside the
one-dimensional source line actually tested by
\(d_n^{\rm birth}\), up to a component whose mixed scalar tends to zero.  It
is retained only when robustness over the full daughter range or ambient
exterior source class was declared before the data were inspected.
\end{corollary}

\begin{proof}
The leftmost term in \eqref{eq:v23-three-gain-levels} is precisely the
squared gain on the occurring source line.  Under
\eqref{eq:v23-occurring-line-null}, that line contributes no record-sized
selected scalar.  The two larger projections may still grow in orthogonal
source directions, but Theorem~\ref{thm:v23-source-class-monotonicity} shows
that those directions belong to stronger declared source classes.  They do
not change the exact scalar identity
\eqref{eq:v23-birth-scalar-projected}.
\end{proof}

% =============================================================================
\section{Corrected terminal alternative and export}
\label{sec:v23-terminal}
% =============================================================================

On the V22 terminal card define pointwise
\begin{equation}
 C_{n,t}^{\rm b}
 :=(I-R_{n,t}^{\rm out})
   \mathscr C_{\mathscr A_n}^{\rm out},
 \qquad
 D_{n,t}^{\rm b}:=(C_{n,t}^{\rm b})^*C_{n,t}^{\rm b},
 \label{eq:v23-terminal-birth-data}
\end{equation}
\begin{equation}
 P_{n,t}^{\rm b}
 :=C_{n,t}^{\rm b}(D_{n,t}^{\rm b})^\dagger
   (C_{n,t}^{\rm b})^*,
 \qquad
 H_{n,t}^{\rm b}:=(C_{n,t}^{\rm b})^*\Psi_n(t).
 \label{eq:v23-terminal-projection-vector}
\end{equation}
The corrected scalar card appends
\begin{equation}
 \boxed{
 \begin{aligned}
 \mathfrak D_n^{\rm V23,\phi}:=\bigl(&
  P_n,Q_n,J_n,\phi_n,\|\phi_n\|_{L^2}^2,
  C_{n,\cdot}^{\rm b},D_{n,\cdot}^{\rm b},
  P_{n,\cdot}^{\rm b},H_{n,\cdot}^{\rm b},\\
 &\mathfrak G_n^{\rm b,phys},\eta_n^{\rm b},
  \mathfrak h_n^{\rm old},\mathfrak h_n^{\rm rep},
  \mathfrak h_n^{\rm birth},\mathfrak h_n^{\rm in},
  \Delta_n^{\rm old,\phi},\Delta_n^{\rm rep,\phi}
 \bigr).
 \end{aligned}}
 \label{eq:v23-terminal-card}
\end{equation}
The complete ambient adjoint field and its norm are not retained in this
source-cyclic scalar card.  They are restored only for the stronger arbitrary
exterior-source target.  All direct signed scalar and atlas entries are the
same V22 reads.

\begin{theorem}[Version V23 source-cyclic terminal gate]
\label{thm:v23-terminal-gate}
Apply first the V19 invariant gate, the V20 daughter-source boundary gate,
the NS--B material-route, leverage, strain, curvature, capture, and boundary-
current gates, and the V22 old/represented target-incidence gate.  On the
remaining branch carrying one frozen writer \(\phi_n\), fix the intended
source use class before passing to a priority subsequence.  Exactly one of the
following applies.
\begin{enumerate}[label=\textup{(G23.\arabic*)},leftmargin=5.7em]
\item \emph{Daughter first-birth survivor.}  For some \(\delta>0\),
      \begin{equation}
       \mathfrak B_n^{\rm dau}\ge\delta.
       \label{eq:v23-source-birth-survivor}
      \end{equation}
      The source occurrence is already nonnegligible and is not repriced by
      a return gain.
\item \emph{Physical daughter-gain survivor.}  One has
      \(\mathfrak B_n^{\rm dau}\to0\), but for some \(\delta>0\),
      \begin{equation}
       \frac{|\mathfrak h_n^{\rm birth}|}{\mathcal R_n^2}
       \ge\delta.
       \label{eq:v23-physical-gain-survivor}
      \end{equation}
      Then \(\mathfrak G_n^{\rm b,phys}\to\infty\).  On a further
      subsequence this is either parent-pair adjoint coefficient overdrive or
      target-weighted daughter-source softness as in
      \eqref{eq:v23-overdrive-or-softness}.
\item \emph{Source-cyclic ancestry-complete scalar card.}  The declared
      target is the one occurring scalar,
      \begin{equation}
       \mathfrak B_n^{\rm dau}\to0,
       \qquad
       \frac{|\mathfrak h_n^{\rm birth}|}{\mathcal R_n^2}\to0,
       \label{eq:v23-scalar-birth-pass}
      \end{equation}
      the V22 old and represented mixed residuals vanish, and the total
      selected return is represented by the declared old/return atlas.  Then
      NS--A exports \eqref{eq:v23-terminal-card}.  Divergence of either larger
      gain outside the occurring source line is pruned by
      Corollary~\ref{cor:v23-target-null-pruning} and receives no scalar price.
\item \emph{Stronger physical source target.}  The target was declared in
      advance to be uniform over a source class larger than the occurring
      line.  For the complete quadratic daughter range, populate the finite
      vector \(H_{n,t}^{\rm b}\) and Gram \(D_{n,t}^{\rm b}\); target-null
      growth relative to the occurring source may then be a genuine
      robustness obstruction.  For arbitrary exterior sources, retain the
      full V22 adjoint writer.  For a reusable return bank or pathwise
      chronology, retain the full V21 return Gram.
\item \emph{All-shell physical daughter escape.}  No exhausting finite head
      captures the selected scalar while controlling the physical daughter
      range and the inherited upstream gates.  The output is one fused
      all-shell source-cyclic escape, not a sum of target-null adjoint norms.
\end{enumerate}
No later source projection repairs an earlier material, frequency, strain,
map-curvature, cell, boundary-current, invariant, or source-birth failure.
\end{theorem}

\begin{proof}
The inherited priority order fixes the physical head, source atlas, and return
writer before the present source short.  First select a subsequence on which
the nonnegative daughter action either has a positive lower bound or tends to
zero.  This gives branch~\textup{(G23.1)} or its complement.  On the latter,
select according to the normalized first-birth scalar.  A positive lower
bound gives branch~\textup{(G23.2)} by
Corollary~\ref{cor:v23-corrected-gain-law};
Theorem~\ref{thm:v23-finite-gain-factorization} gives its finite subdivision.

If the first-birth scalar also tends to zero, the result depends only on the
source use class fixed before subsequence selection.  For the one-scalar
target, Corollary~\ref{cor:v23-target-null-pruning} removes every gain outside
the occurring source line.  The passing V22 old/represented mixed residuals
and exact V22 scalar assembly then give branch~\textup{(G23.3)}.  If instead
a larger source class or output object was declared, its corresponding finite
parent, ambient adjoint, or full return rows remain mandatory and give
branch~\textup{(G23.4)}.  Failure of every finite head to reach the relevant
card is exactly branch~\textup{(G23.5)}.
\end{proof}

\begin{corollary}[Corrected source-native NS--A export]
\label{cor:v23-source-native-export}
On branch~\textup{(G23.3)}, the V22 scalar ancestry certificate remains valid
with the ambient gain replaced by the source-cyclic card
\eqref{eq:v23-terminal-card}.  On branch~\textup{(G23.2)}, export instead the
finite physical packet
\begin{equation}
 \boxed{
 \bigl(
  \mathscr A_n,
  C_{n,\cdot}^{\rm b},D_{n,\cdot}^{\rm b},
  H_{n,\cdot}^{\rm b},
  \mathfrak B_n^{\rm dau},
  \mathfrak G_n^{\rm b,phys},
  \mathfrak H_n^{\rm par},
  \mathfrak S_n^{\rm src}
 \bigr).}
 \label{eq:v23-gain-export}
\end{equation}
It is then the task of the fixed-datum donor programme to evaluate the actual
source-visibility floor or parent-pair coefficient overdrive on that same
history.  No new source norm is introduced.
\end{corollary}

\begin{proof}
Branch~\textup{(G23.3)} has zero selected first-birth residual and the V22
old/represented target incidences, so the V22 scalar assembly closes.  The
ambient-null adjoint term is orthogonal to every physical first-birth source
by Theorem~\ref{thm:v23-daughter-support}.  On branch~\textup{(G23.2)}, all
entries of \eqref{eq:v23-gain-export} are formed from the same finite head,
represented source atlas, physical adjoint writer, and first-birth history;
Theorem~\ref{thm:v23-finite-gain-factorization} gives the stated alternative.
\end{proof}

% =============================================================================
\section{Integrated conclusion, proof ledger, and exact next acquisition}
\label{sec:v23-conclusion}
% =============================================================================

\begin{theorem}[Version V23 physical daughter-gain theorem]
\label{thm:v23-integrated}
Preserving Versions~V1--V22 verbatim, Version~V23 proves the following.
\begin{enumerate}[label=\textup{(V23.\arabic*)},leftmargin=5.5em]
\item The represented-shorted quadratic daughter synthesis
      \(C_t^{\rm b}\) has canonical positive Gram \(D_t^{\rm b}\) and
      support projection \(P_t^{\rm b}\).
\item The V22 first-birth scalar depends only on the projected adjoint
      \(\Psi_\phi^{\rm b}=P_t^{\rm b}\Psi_\phi\), whose action is the finite
      Moore--Penrose form \eqref{eq:v23-history-gain-finite}.
\item The sharp minimum source action for the selected return, restricted to
      physical daughter sources, is \eqref{eq:v23-physical-minimum-action}.
      The ambient V22 minimum is a weaker lower bound for a stronger source
      class.
\item The ambient adjoint norm splits exactly into physical daughter gain and
      daughter-null gain.  Countertheorem~\ref{cth:v23-ambient-gain-null}
      shows that only the first has source meaning for this target.
\item The occurring scalar has the exact action--gain--alignment
      factorization \eqref{eq:v23-terminal-exact-product}; vanishing daughter
      action with a fixed return forces the projected physical gain to
      diverge.
\item The physical gain is determined by one finite symmetric parent-pair
      writer and the V20 exterior source Gram.  Its divergence gives the
      coefficient-overdrive/source-softness alternative
      \eqref{eq:v23-overdrive-or-softness}.
\item Source classes are nested target semantics.  One occurring scalar, the
      complete quadratic daughter range, arbitrary exterior sources, and a
      complete returned bank require successively stronger carriers.
\item The corrected terminal gate is
      Theorem~\ref{thm:v23-terminal-gate}.  Gain growth outside the occurring source line is pruned from the scalar
      obstruction list unless a larger source use class was declared.
\end{enumerate}
No terminal concentration branch is excluded unconditionally, and no global
regularity conclusion is asserted.
\end{theorem}

\begin{proof}
Items~\textup{(V23.1)--(V23.2)} are
Theorem~\ref{thm:v23-daughter-support} and
Corollary~\ref{cor:v23-history-projection}.  Item~\textup{(V23.3)} is
Theorem~\ref{thm:v23-physical-Riesz}.  Item~\textup{(V23.4)} is
Corollary~\ref{cor:v23-ambient-gain-Pythagoras} together with
Countertheorem~\ref{cth:v23-ambient-gain-null}.  Item~\textup{(V23.5)} is
Theorem~\ref{thm:v23-alignment-factorization} and
Corollary~\ref{cor:v23-corrected-gain-law}.  Item~\textup{(V23.6)} is
Theorems~\ref{thm:v23-quadratic-parent-read} and
\ref{thm:v23-finite-gain-factorization}.  Item~\textup{(V23.7)} is
Theorem~\ref{thm:v23-source-class-monotonicity} and
Corollary~\ref{cor:v23-target-null-pruning}, and the final item is
Theorem~\ref{thm:v23-terminal-gate}.
\end{proof}

\begingroup\small
\begin{longtable}{>{\raggedright\arraybackslash}p{0.27\textwidth}
                  >{\raggedright\arraybackslash}p{0.18\textwidth}
                  >{\raggedright\arraybackslash}p{0.47\textwidth}}
\caption{Version V23 proof-status ledger.}\label{tab:v23-ledger}\\
\toprule
\textbf{Row}&\textbf{Status}&\textbf{Exact advance or limitation}\\
\midrule
\endfirsthead
\toprule
\textbf{Row}&\textbf{Status}&\textbf{Exact advance or limitation}\\
\midrule
\endhead
V1--V22 prefix&\PROVED{} preserved
 &Complete V22 preterminal source retained byte-for-byte; no inherited theorem
 rewritten.\\[0.35em]
Upstream NS--B physical gate&\INHERITED{} prior
 &Material-route mismatch, leverage, strain, map curvature, cell capture, and
 signed boundary currents remain before the present source short.\\[0.35em]
First-birth synthesis&\PROVED{} exact
 &The represented-shorted map
 \(C_t^{\rm b}=(I-R_t^{\rm out})\mathscr C_{\mathscr A}^{\rm out}\)
 generates the occurring daughter source.\\[0.35em]
Physical source support&\PROVED{} exact
 &\(C_t^{\rm b}(D_t^{\rm b})^\dagger(C_t^{\rm b})^*\) is the canonical
 daughter-range projection.\\[0.35em]
Projected adjoint scalar&\PROVED{} exact
 &The first-birth return depends only on
 \(P_t^{\rm b}\Psi_\phi\); the daughter-null complement is invisible.\\[0.35em]
Restricted Riesz minimum&\PROVED{} sharp
 &The physical minimum action is
 \(|\tau|^2/\|P^{\rm b}\Psi_\phi\|^2\), with a unique source minimizer.\\[0.35em]
Ambient gain interpretation&\OBSTRUCTION{} explicit
 &A causal Volterra family has divergent full adjoint norm and constant
 physical daughter gain.\\[0.35em]
Parent source metric&\PROVED{} necessary
 &The gain is \(\langle H,D^\dagger H\rangle\); raw parent coefficients can
 collapse while physical gain remains fixed.\\[0.35em]
Actual-source alignment&\PROVED{} exact
 &Selected scalar equals daughter action times physical gain times one
 alignment efficiency.\\[0.35em]
Small birth / fixed return&\PROVED{} corrected survivor
 &The projected physical daughter gain, not merely the ambient gain, diverges.\\[0.35em]
Finite gain subdivision&\PROVED{} exact
 &Divergence is parent-pair coefficient overdrive or target-weighted source
 softness on the same finite head.\\[0.35em]
Source-class hierarchy&\PROVED{} target relative
 &One source line, the quadratic daughter range, arbitrary exterior sources,
 and a full response bank have monotone gains and reverse action minima.\\[0.35em]
Scalar ancestry export&\CONDITIONAL{} exact
 &After upstream and V22 mixed-incidence rows pass, gain outside the occurring source line is omitted from the scalar price.\\[0.35em]
Critical upper stock / rigidity&\OPEN{}
 &No source-visibility floor, parent-coefficient upper bound, terminal trace,
 epsilon regularity, backward uniqueness, or Liouville theorem is supplied.\\[0.35em]
Three-dimensional regularity&\OPEN
 &No surviving terminal branch is eliminated for every fixed datum.\\
\bottomrule
\end{longtable}
\endgroup

% =============================================================================
\section{Exact mandate after Version V23}
\label{sec:v23-next}
% =============================================================================

\begin{programme}[Populate one physical source-cyclic scalar or stop]
\label{prog:v23-next}
Preserve Versions~V1--V23 verbatim.  There is no automatic generic
Version~V24.

A continuation is licensed only in the following order.
\begin{enumerate}[label=\textup{(V24.\arabic*)},leftmargin=5.9em]
\item Apply the inherited material-route, ultra-canonical leverage, strain,
      material-curvature, capture, and signed boundary-current gate first.
      Stop at its first surviving row.
\item On a collar-free captured branch, freeze the finite helical head, the
      represented exterior source atlas, and one physical return writer
      before evaluating any target gain.
\item Populate the represented-shorted exterior coefficient matrix
      \(C_{n,t}^{\rm b}\), its positive Gram \(D_{n,t}^{\rm b}\), and the
      finite parent target vector
      \(H_{n,t}^{\rm b}=(C_{n,t}^{\rm b})^*\Psi_n(t)\).
\item If the physical gain diverges, evaluate first the exact factorization
      \eqref{eq:v23-gain-softness-factorization}.  Return either a target-
      weighted source-visibility collapse or a parent-pair adjoint
      coefficient overdrive on the same head.
\item Compare parent-coefficient overdrive with the already existing
      ultra-canonical frequency, strain, material-curvature, and critical-
      acceleration packets.  Do not introduce another ambient source norm.
\item If the first-birth scalar vanishes and the V22 old/represented mixed
      residuals pass, terminate NS--A and export
      \eqref{eq:v23-terminal-card} to NS--B.
\item Populate a complete parent matrix only for a declared reusable donor
      variation bank, the ambient adjoint only for arbitrary exterior-source
      control, and the V21 return Gram only for a complete return range or
      pathwise chronology.
\item If no finite head controls the physical daughter range, retain one
      fused all-shell source-cyclic escape.  Reopen it only with a new
      source-relative critical action, capacity, or square-function upper
      stock on the same carrier.
\item Do not add another invariant family, path law, age law, Dini modulus,
      concentration metric, shell count, source class, or generic response
      compiler.
\end{enumerate}
\end{programme}

\begin{assessment}[Programme position after Version V23]
\label{ass:v23-position}
The first unresolved scalar is no longer the ambient norm of a backward
adjoint writer.  It is the source-cyclic Moore--Penrose value
\[
 \int\left\langle
  (C_{n,t}^{\rm b})^*\Psi_n(t),
  (D_{n,t}^{\rm b})^\dagger
  (C_{n,t}^{\rm b})^*\Psi_n(t)
 \right\rangle\,\dd t
\]
on the actual finite daughter bank, together with its direct mixed return
scalar.  If this value overdrives, the exact remaining alternatives are a
soft physical daughter metric or growth of one finite parent-pair target
coefficient.  If the first-birth scalar vanishes, gain growth outside the occurring
source line is irrelevant to that scalar and the ancestry card may close.  The next admissible work is therefore one literal finite coefficient
read or the first inherited PDE survivor, not another transport construction.
\end{assessment}

\paragraph{Version V23 history.}
Preserved the complete V22 source.  Moved upstream from the ambient adjoint
norm to the represented-shorted quadratic daughter range of the already
selected finite head.  Constructed its canonical Moore--Penrose support,
projected the V22 adjoint writer to the physical source class, and proved the
sharp restricted Riesz principle.  Exhibited a causal target-null adjoint
counterexample and a source-softness counterexample to raw parent-coordinate
pricing.  Derived the exact action--gain--alignment factorization, pulled the
selected return to one finite symmetric parent-pair writer, and split physical
gain divergence into source softness or parent-coefficient overdrive.  Inserted
this correction into the V22 terminal gate and removed ambient daughter-null
adjoint growth from the one-scalar obstruction list.  No terminal PDE closure
or global regularity conclusion was claimed.

\begin{assessment}[Append-only integrity]
\label{ass:v23-integrity}
The complete Version~V22 source preceding its sole final executable document
terminator is preserved verbatim.  Its preterminal SHA--256 digest is
\begin{center}\scriptsize\ttfamily
95754e8cde05dd6b9430100f35850cd184a6eacddbadc6a1c597990797bd586f
\end{center}
The present material begins at Section~\ref{sec:v23-scope}; the final command
below is again unique.
\end{assessment}

\addtocontents{toc}{\protect\endgroup}
% =============================================================================
% END VERSION V23 MATERIAL -- append later versions before final terminator
% =============================================================================

% APPEND ALL LATER VERSIONS IMMEDIATELY ABOVE THE SOLE FINAL LINE BELOW.


\clearpage
% =============================================================================
% VERSION V24: NEW MATERIAL.  The complete V1--V23 source above is preserved
% verbatim; only its sole active \end{document} line was removed before append.
% =============================================================================
\hypersetup{pdftitle={NS-A Terminal Source Concentration Programme: Cumulative V24},pdfauthor={A. Pelissier}}
\makeatletter
\addtocontents{toc}{\protect\begingroup\protect\makeatletter}
\addtocontents{toc}{\protect\def\protect\l@section{\protect\@dottedtocline{1}{0em}{4.7em}}}
\addtocontents{toc}{\protect\def\protect\l@subsection{\protect\@dottedtocline{2}{1.5em}{5.3em}}}
\addtocontents{toc}{\protect\makeatother}
\makeatother
\renewcommand{\thesection}{V24.\arabic{section}}
\renewcommand{\thesubsection}{V24.\arabic{section}.\arabic{subsection}}
\renewcommand{\theequation}{V24.\arabic{equation}}
\setcounter{section}{0}
\setcounter{equation}{0}
\renewcommand{\theHsection}{V24.\arabic{section}}
\renewcommand{\theHsubsection}{V24.\arabic{section}.\arabic{subsection}}
\renewcommand{\theHtheorem}{V24.\arabic{section}.\arabic{theorem}}
\renewcommand{\theHequation}{V24.\arabic{equation}}

\begin{center}
{\LARGE\bfseries NS--A: Version V24}\\[0.5em]
{\large Range-Only Source Gain, the Actual Daughter Line,\\
and the Cofinal Support-Transport Stop}\\[0.5em]
{\normalsize 3 September 2026}
\end{center}
\addcontentsline{toc}{part}{Version V24: range-only source gain and actual-line susceptibility}

\begin{abstract}
Version~V23 moved the V22 backward adjoint writer into the represented-shorted
quadratic daughter range and expressed the resulting Moore--Penrose gain as a
finite parent-pair quadratic form.  It then factored that gain into a parent
coefficient norm and a target-weighted inverse daughter metric, producing a
``coefficient overdrive versus source softness'' alternative.  The factorization
is algebraically exact.  Its interpretation, however, still mixes two different
physical targets.

The present continuation goes one arrow upstream and separates them.  For a
finite daughter synthesis \(C\), daughter Gram \(D=C^*C\), target writer
\(\psi\), and parent target vector \(H=C^*\psi\), one has the exact singular-
value cancellation
\[
 \langle H,D^\dagger H\rangle
 =\|P_{\operatorname{Ran}C}\psi\|^2.
\]
Every nonzero singular value of \(C\) cancels.  Thus the minimum action among
\emph{physical daughter source vectors}, measured in their native exterior
\(L^2\) norm, depends only on the reachable source range and its angle with the
adjoint writer.  It does not grow merely because the parent-to-source synthesis
becomes soft.  Source softness matters for a different target: stability of the
source range across cutoffs, or payment in a separately declared parent-amplitude
metric.  A rescaling example keeps the physical source gain fixed while moving
arbitrarily between the two V23 factors, and a support-collapse example shows
that a uniform singular floor is instead the exact requirement for cofinal
transport of the Moore--Penrose source projection.

For the one scalar actually carried by NS--A, the target is smaller still.  If
\(d^{\rm birth}\) is the occurring daughter history and \(\Psi\) the frozen
adjoint writer, the exact susceptibility is
\[
 \mathcal G^{\rm occ}
 =\frac{|\langle d^{\rm birth},\Psi\rangle|^2}
        {\|d^{\rm birth}\|^2}
 =\|P_{\mathbb R d^{\rm birth}}\Psi\|^2.
\]
It is bounded by the V23 daughter-range gain.  Vanishing daughter action with a
fixed record-sized scalar therefore forces divergence already on the occurring
source line; no further parent-pair or inverse-metric split is needed.  For a
reusable fixed-datum first-variation bank, the inherited V20 tangent daughter
germ gives the intermediate target.  An explicit finite model has divergent
full quadratic-range gain but bounded tangent and occurring-line gains, proving
that the complete symmetric parent range is again a stronger target.

Version~V24 consequently removes ``source softness'' as an independent
finite-card causal-amplification mechanism for the one-scalar NS--A branch.  It
retains softness only as a cofinal support-transport obstruction or as part of a
separately declared parent-amplitude donor budget.  The surviving scalar branch
is a projected-adjoint writer amplification and must be evaluated by the already
existing material--heat, ultra-canonical frequency, strain, material-curvature,
and boundary-current gates.  No new source class, response graph, path law,
Dini modulus, or compactness compiler is introduced, and no global
three-dimensional regularity conclusion is claimed.
\end{abstract}

% =============================================================================
\section{Scope: distinguish source action from parent-amplitude action}
\label{sec:v24-scope}
% =============================================================================

Retain the complete V23 card only after the inherited material-route,
ultra-canonical leverage, symmetric-strain, material-map-curvature, cell-capture,
and signed boundary-current gates have passed.  At one fixed time, suppress the
record index and write
\begin{equation}
 C:\mathcal Q\longrightarrow\mathcal Y,
 \qquad D=C^*C,
 \qquad H=C^*\psi,
 \qquad P_C:=CD^\dagger C^*.
 \label{eq:v24-finite-card}
\end{equation}
Here \(\mathcal Q\) is the finite symmetric parent-pair space,
\(\mathcal Y\) is the exterior source space, and \(\psi\) is the already frozen
backward-adjoint return writer.  In the V23 application,
\[
 C=C_t^{\rm b},
 \qquad D=D_t^{\rm b},
 \qquad H=H_{\phi,t}^{\rm b},
 \qquad \psi=\Psi_\phi(t).
\]

There are three different costs which must not be identified:
\begin{enumerate}[label=\textup{(M\arabic*)},leftmargin=3.5em]
\item the native norm of an exterior source vector \(f\in\operatorname{Ran}C\);
\item the norm of a parent coefficient \(q\in\mathcal Q\) with \(Cq=f\);
\item the one actually occurring source line generated by the fixed datum.
\end{enumerate}
The V23 Moore--Penrose form answers the first question.  Its parent coefficient
norm and inverse-metric factor separately become physical only when the second
question is also part of the target.

\begin{firewall}[No change to the V23 identities]
\label{fire:v24-v23-identities}
Every V23 equality involving \(C^*C\), its Moore--Penrose inverse, the projected
adjoint writer, and the action--gain--alignment product remains valid.  Version
V24 changes only the interpretation of the factorization
\(\mathfrak G^{\rm b,phys}=\mathfrak H^{\rm par}\mathfrak S^{\rm src}\).
The product is a source-action gain; its two factors separately refer to a
chosen parent metric.  They are not two intrinsic causes of that source-action
gain.
\end{firewall}

% =============================================================================
\section{Exact singular-value cancellation}
\label{sec:v24-svd}
% =============================================================================

\begin{theorem}[Range-only formula for the physical daughter gain]
\label{thm:v24-range-only}
Let \(C,D,H,P_C\) be as in \eqref{eq:v24-finite-card}.  Then \(P_C\) is the
orthogonal projection onto \(\operatorname{Ran}C\), and
\begin{equation}
 \boxed{
 \langle H,D^\dagger H\rangle
 =\langle\psi,P_C\psi\rangle
 =\|P_C\psi\|^2.}
 \label{eq:v24-range-only}
\end{equation}
If
\begin{equation}
 C=\sum_{j=1}^r\sigma_j\,u_j\otimes v_j^*,
 \qquad \sigma_j>0,
 \label{eq:v24-svd}
\end{equation}
with orthonormal singular vectors, then
\begin{equation}
 \boxed{
 \begin{aligned}
 H&=\sum_{j=1}^r\sigma_j\langle u_j,\psi\rangle v_j,\\
 \langle H,D^\dagger H\rangle
 &=\sum_{j=1}^r|\langle u_j,\psi\rangle|^2.
 \end{aligned}}
 \label{eq:v24-svd-cancellation}
\end{equation}
Thus every nonzero singular value cancels from the native exterior-source
gain.
\end{theorem}

\begin{proof}
The Moore--Penrose equations imply that
\(CD^\dagger C^*=CC^\dagger\) is the orthogonal projection onto
\(\operatorname{Ran}C\).  Hence
\[
 \langle H,D^\dagger H\rangle
 =\langle C^*\psi,D^\dagger C^*\psi\rangle
 =\langle\psi,CD^\dagger C^*\psi\rangle
 =\langle\psi,P_C\psi\rangle.
\]
For \eqref{eq:v24-svd},
\[
 D=\sum_{j=1}^r\sigma_j^2v_j\otimes v_j^*,
 \qquad
 D^\dagger=\sum_{j=1}^r\sigma_j^{-2}v_j\otimes v_j^*.
\]
Substitution gives \eqref{eq:v24-svd-cancellation}.
\end{proof}

\begin{corollary}[Source-range invariance under parent reparameterization]
\label{cor:v24-reparameterization}
Let \(G:\mathcal Q\to\mathcal Q\) be invertible and put
\(\widetilde C=CG\).  Then
\begin{equation}
 \boxed{
 \operatorname{Ran}\widetilde C=\operatorname{Ran}C,
 \qquad
 \langle\widetilde C^*\psi,
  (\widetilde C^*\widetilde C)^\dagger
  \widetilde C^*\psi\rangle
 =\|P_C\psi\|^2.}
 \label{eq:v24-reparameterization}
\end{equation}
The native source-action gain is therefore invariant under every invertible
change of parent coordinates.  A nonunitary change does alter the Euclidean
parent-amplitude cost and must be declared if that cost is physical.
\end{corollary}

\begin{proof}
Invertibility of \(G\) gives equality of the two ranges.  Apply
Theorem~\ref{thm:v24-range-only} to \(\widetilde C\).
\end{proof}

\begin{corollary}[Exact rescaling law for the V23 factors]
\label{cor:v24-rescaling}
For \(\varepsilon>0\), put \(C_\varepsilon=\varepsilon C\).  Define
\begin{equation}
 \mathcal G_{\rm src}(C,\psi):=\|P_C\psi\|^2,
 \qquad
 \mathcal G_{\rm par}(C,\psi):=\|C^*\psi\|^2.
 \label{eq:v24-two-gains}
\end{equation}
Whenever \(\mathcal G_{\rm par}>0\), put
\(
 \mathcal S_{\rm par}:=\mathcal G_{\rm src}/\mathcal G_{\rm par}
\).
Then
\begin{equation}
 \boxed{
 \begin{aligned}
 \mathcal G_{\rm src}(C_\varepsilon,\psi)
 &=\mathcal G_{\rm src}(C,\psi),\\
 \mathcal G_{\rm par}(C_\varepsilon,\psi)
 &=\varepsilon^2\mathcal G_{\rm par}(C,\psi),\\
 \mathcal S_{\rm par}(C_\varepsilon,\psi)
 &=\varepsilon^{-2}\mathcal S_{\rm par}(C,\psi).
 \end{aligned}}
 \label{eq:v24-rescaling-law}
\end{equation}
Thus the coefficient/softness split can move arbitrarily while the physical
source-action gain is unchanged.
\end{corollary}

\begin{proof}
The range is unchanged for \(\varepsilon>0\).  The remaining identities follow
from \(C_\varepsilon^*\psi=\varepsilon C^*\psi\).
\end{proof}

\begin{countertheorem}[The two V23 factors are not intrinsic source-action branches]
\label{cth:v24-factor-label}
On \(\mathbb R^2\), let \(\psi_n=ne_1\).  Consider
\begin{equation}
 C_n^{(1)}=I,
 \qquad
 C_n^{(2)}=n^{-2}I.
 \label{eq:v24-two-scalings}
\end{equation}
Both families have the same source range and the same native source gain
\begin{equation}
 \boxed{
 \mathcal G_{\rm src}(C_n^{(1)},\psi_n)
 =\mathcal G_{\rm src}(C_n^{(2)},\psi_n)=n^2.}
 \label{eq:v24-same-source-gain}
\end{equation}
For the first family,
\begin{equation}
 \mathcal G_{\rm par}=n^2,
 \qquad \mathcal S_{\rm par}=1,
 \label{eq:v24-first-factorization}
\end{equation}
whereas for the second,
\begin{equation}
 \mathcal G_{\rm par}=n^{-2},
 \qquad \mathcal S_{\rm par}=n^4.
 \label{eq:v24-second-factorization}
\end{equation}
The same source-action amplification can therefore be labelled entirely as
parent-coefficient growth or entirely as inverse parent-metric growth.  The
labels acquire independent physical meaning only after a parent-amplitude
metric has itself been loaded as a target.
\end{countertheorem}

\begin{proof}
The ranges are both \(\mathbb R^2\), so
\(P_{C_n^{(i)}}=I\).  The remaining calculations are immediate from
\eqref{eq:v24-two-gains}.
\end{proof}

% =============================================================================
\section{Three target-native Riesz problems}
\label{sec:v24-three-costs}
% =============================================================================

\begin{definition}[Source, parent, and occurring-line gains]
\label{def:v24-three-gains}
For \(d=Cq\ne0\), define
\begin{align}
 \mathcal G_{\rm src}(C,\psi)
 &:=\|P_C\psi\|^2,
 \label{eq:v24-source-gain}\\
 \mathcal G_{\rm par}(C,\psi)
 &:=\|C^*\psi\|^2,
 \label{eq:v24-parent-gain}\\
 \mathcal G_{\rm occ}(d,\psi)
 &:=\frac{|\langle d,\psi\rangle|^2}{\|d\|^2}.
 \label{eq:v24-occurring-gain}
\end{align}
Set \(\mathcal G_{\rm occ}=0\) when \(d=0\).
\end{definition}

\begin{theorem}[Exact hierarchy and separate Thomson minima]
\label{thm:v24-three-minima}
The gains satisfy
\begin{equation}
 \boxed{
 0\le\mathcal G_{\rm occ}(d,\psi)
 \le\mathcal G_{\rm src}(C,\psi)
 \le\|\psi\|^2.}
 \label{eq:v24-gain-hierarchy}
\end{equation}
For every scalar \(\tau\), the three minimum-action problems are
\begin{align}
 \inf_{\substack{f\in\operatorname{Ran}C\\
                  \langle f,\psi\rangle=\tau}}
 \|f\|^2
 &=\frac{|\tau|^2}{\mathcal G_{\rm src}(C,\psi)},
 \label{eq:v24-source-minimum}\\
 \inf_{\substack{q\in\mathcal Q\\
                  \langle Cq,\psi\rangle=\tau}}
 \|q\|^2
 &=\frac{|\tau|^2}{\mathcal G_{\rm par}(C,\psi)},
 \label{eq:v24-parent-minimum}\\
 \inf_{\substack{a\in\mathbb R\\
                  \langle ad,\psi\rangle=\tau}}
 \|ad\|^2
 &=\frac{|\tau|^2}{\mathcal G_{\rm occ}(d,\psi)}.
 \label{eq:v24-line-minimum}
\end{align}
A zero denominator means that the corresponding nonzero target scalar is
unattainable.
\end{theorem}

\begin{proof}
Since \(d\in\operatorname{Ran}C\), the orthogonal projection onto
\(\mathbb Rd\) is dominated by the projection onto \(\operatorname{Ran}C\).
This proves \eqref{eq:v24-gain-hierarchy}.  Each minimum is the Riesz theorem
on the indicated Hilbert subspace.  For the parent problem, the representing
vector of \(q\mapsto\langle Cq,\psi\rangle\) is \(C^*\psi\).
\end{proof}

\begin{theorem}[Singular edges compare source and parent targets]
\label{thm:v24-singular-edge-comparison}
Let \(\sigma_-\) and \(\sigma_+\) be the smallest and largest positive
singular values of \(C\).  Then
\begin{equation}
 \boxed{
 \sigma_-^2\mathcal G_{\rm src}(C,\psi)
 \le\mathcal G_{\rm par}(C,\psi)
 \le\sigma_+^2\mathcal G_{\rm src}(C,\psi).}
 \label{eq:v24-singular-edge-comparison}
\end{equation}
Consequently, when \(\mathcal G_{\rm par}>0\),
\begin{equation}
 \boxed{
 \sigma_+^{-2}
 \le\frac{\mathcal G_{\rm src}}{\mathcal G_{\rm par}}
 \le\sigma_-^{-2}.}
 \label{eq:v24-softness-window}
\end{equation}
These inequalities compare two separately declared metrics.  They do not
alter the range-only identity \eqref{eq:v24-range-only}.
\end{theorem}

\begin{proof}
In the singular basis of \eqref{eq:v24-svd},
\[
 \mathcal G_{\rm par}
 =\sum_{j=1}^r\sigma_j^2|\langle u_j,\psi\rangle|^2,
 \qquad
 \mathcal G_{\rm src}
 =\sum_{j=1}^r|\langle u_j,\psi\rangle|^2.
\]
Bounding every \(\sigma_j^2\) by the two spectral edges proves the result.
\end{proof}

\begin{assessment}[What a soft daughter Gram actually means]
\label{ass:v24-softness-meaning}
A small positive singular value of \(C\) means that a unit parent coefficient
produces a small exterior source vector.  It raises the parent-amplitude cost of
producing a prescribed physical source.  It does not raise the minimum
\emph{exterior-source} action for a prescribed return scalar, because that
minimum is measured after the source has been produced and depends only on its
range.  The singular floor re-enters when the parent amplitudes themselves are
priced or when the source range must be transported stably across cutoffs.
\end{assessment}

% =============================================================================
\section{Cofinal support transport and the true role of a singular floor}
\label{sec:v24-support-transport}
% =============================================================================

\begin{theorem}[Quantitative transport of daughter support projections]
\label{thm:v24-projection-stability}
Let \(C_0,C_1:\mathcal Q\to\mathcal Y\) be finite matrices and let
\(P_i\) be the orthogonal projection onto \(\operatorname{Ran}C_i\).  If
\(C_i\ne0\), suppose its smallest positive singular value is at least
\(\sigma_i>0\); if \(C_i=0\), use the convention
\(\sigma_i^{-1}=0\).  Then
\begin{equation}
 \boxed{
 \|P_1-P_0\|_{\rm op}
 \le\|C_1-C_0\|_{\rm op}
       \left(\sigma_0^{-1}+\sigma_1^{-1}\right).}
 \label{eq:v24-projection-stability}
\end{equation}
In particular, operator-norm convergence of the daughter syntheses together
with one uniform positive singular edge transports the physical source range.
\end{theorem}

\begin{proof}
If \(C_0=0\), then \((I-P_1)P_0=0\).  Otherwise let
\(y\in\operatorname{Ran}C_0\) with \(\|y\|=1\), and choose its minimum-
norm preimage \(x=C_0^\dagger y\).  Then \(\|x\|\le\sigma_0^{-1}\), and
\[
 \operatorname{dist}(y,\operatorname{Ran}C_1)
 \le\|y-C_1x\|
 =\|(C_0-C_1)x\|
 \le\sigma_0^{-1}\|C_1-C_0\|_{\rm op}.
\]
Hence
\(
 \|(I-P_1)P_0\|
 \le\sigma_0^{-1}\|C_1-C_0\|
\).
Interchanging the indices gives
\(
 \|(I-P_0)P_1\|
 \le\sigma_1^{-1}\|C_1-C_0\|
\).
Finally,
\[
 P_1-P_0=(I-P_0)P_1-P_0(I-P_1),
\]
and the two terms have the preceding norms.
\end{proof}

\begin{corollary}[Continuity of the physical gain]
\label{cor:v24-gain-continuity}
For writers \(\psi_0,\psi_1\in\mathcal Y\),
\begin{equation}
 \boxed{
 \begin{aligned}
 &\left|
  \|P_1\psi_1\|^2-\|P_0\psi_0\|^2
 \right|\\
 &\quad\le
 (\|\psi_0\|+\|\psi_1\|)\|\psi_1-\psi_0\|
 +\|P_1-P_0\|\,\|\psi_0\|^2.
 \end{aligned}}
 \label{eq:v24-gain-continuity}
\end{equation}
Thus convergence of the writer and stable transport of the daughter support
imply convergence of the source-action gain.
\end{corollary}

\begin{proof}
Expand
\[
 \langle\psi_1,P_1\psi_1\rangle
 -\langle\psi_0,P_0\psi_0\rangle
\]
by adding and subtracting
\(\langle\psi_0,P_1\psi_0\rangle\), and use that orthogonal projections are
contractions.
\end{proof}

\begin{countertheorem}[Support projections need not survive a soft limit]
\label{cth:v24-support-collapse}
Let \(C_n=n^{-1}I\) on \(\mathbb R^m\) and let \(C_\infty=0\).  Then
\begin{equation}
 C_n\longrightarrow C_\infty
 \quad\text{in operator norm},
 \label{eq:v24-soft-convergence}
\end{equation}
but
\begin{equation}
 \boxed{
 P_{C_n}=I\quad\text{for every }n,
 \qquad
 P_{C_\infty}=0.}
 \label{eq:v24-projection-jump}
\end{equation}
For every nonzero fixed writer \(\psi\), the finite-card source gain remains
\(\|\psi\|^2\) and then drops to zero at the limit.  A positive singular floor
is therefore a cofinal support-transport condition, not a finite-card source-
action amplification mechanism.
\end{countertheorem}

\begin{proof}
Every \(C_n\) is invertible, whereas \(C_\infty\) has zero range.
\end{proof}

\begin{corollary}[Time-dependent cofinal gain transport]
\label{cor:v24-time-dependent-transport}
Let \(C_n(t)\) and \(C(t)\) be measurable finite matrices on a finite interval
\(I\), and let \(\Psi_n\to\Psi\) strongly in \(L^2(I;\mathcal Y)\).  Suppose
\begin{equation}
 \|C_n-C\|_{L^\infty_t\mathcal L}\longrightarrow0
 \label{eq:v24-C-cofinal}
\end{equation}
and all nonzero singular values of \(C_n(t)\) and \(C(t)\) have one uniform
positive lower bound.  Then
\begin{equation}
 \boxed{
 \int_I\|P_{C_n(t)}\Psi_n(t)\|^2\,\dd t
 \longrightarrow
 \int_I\|P_{C(t)}\Psi(t)\|^2\,\dd t.}
 \label{eq:v24-integrated-gain-transport}
\end{equation}
Without the singular floor, Countertheorem~\ref{cth:v24-support-collapse}
shows that no such conclusion follows from convergence of the synthesis alone.
\end{corollary}

\begin{proof}
Theorem~\ref{thm:v24-projection-stability} gives uniform convergence of the
support projections.  Apply Corollary~\ref{cor:v24-gain-continuity} pointwise
and integrate, using the strong \(L^2\) convergence and uniform boundedness of
orthogonal projections.
\end{proof}

% =============================================================================
\section{The actual first-birth line is the one-scalar target}
\label{sec:v24-actual-line}
% =============================================================================

Return to the V23 history interval \(J_n\).  Work in the history Hilbert space
\[
 \mathscr Y_n=L^2(J_n;\mathcal Y_n).
\]
Let
\begin{equation}
 d_n=d_n^{\rm birth},
 \qquad
 \Psi_n^{\rm b}(t)=P_{n,t}^{\rm b}\Psi_n(t),
 \qquad
 \tau_n=\mathfrak h_n^{\rm birth}
 =\langle d_n,\Psi_n^{\rm b}\rangle_{\mathscr Y_n}.
 \label{eq:v24-history-packet}
\end{equation}

\begin{definition}[Actual-line susceptibility]
\label{def:v24-actual-line}
When \(d_n\ne0\), let \(P_n^{\rm occ}\) be the orthogonal projection in
\(\mathscr Y_n\) onto \(\mathbb Rd_n\), and define
\begin{equation}
 \boxed{
 \mathcal G_n^{\rm occ}
 :=\|P_n^{\rm occ}\Psi_n^{\rm b}\|_{\mathscr Y_n}^2
 =\frac{|\tau_n|^2}{\|d_n\|_{\mathscr Y_n}^2}.}
 \label{eq:v24-actual-line-susceptibility}
\end{equation}
Set \(\mathcal G_n^{\rm occ}=0\) when \(d_n=0\).
\end{definition}

\begin{theorem}[Exact line/range hierarchy on the terminal history]
\label{thm:v24-history-hierarchy}
One has
\begin{equation}
 \boxed{
 0\le\mathcal G_n^{\rm occ}
 \le\|\Psi_n^{\rm b}\|_{\mathscr Y_n}^2
 \le\|\Psi_n\|_{\mathscr Y_n}^2.}
 \label{eq:v24-history-hierarchy}
\end{equation}
Moreover, the V23 alignment efficiency is exactly
\begin{equation}
 \boxed{
 \eta_n^{\rm b}
 =\frac{\mathcal G_n^{\rm occ}}
        {\|\Psi_n^{\rm b}\|_{\mathscr Y_n}^2}}
 \label{eq:v24-alignment-is-angle}
\end{equation}
when the denominator is nonzero.  Thus V23's action--gain--alignment identity
is the ordinary decomposition into the norm of the occurring source, the norm
of the projected writer, and their squared cosine.
\end{theorem}

\begin{proof}
The first inequality is the contraction property of the line projection.
The second follows from the pointwise orthogonal projection
\(P_{n,t}^{\rm b}\).  Finally,
\[
 \eta_n^{\rm b}
 =\frac{|\langle d_n,\Psi_n^{\rm b}\rangle|^2}
        {\|d_n\|^2\|\Psi_n^{\rm b}\|^2}
 =\frac{\mathcal G_n^{\rm occ}}
        {\|\Psi_n^{\rm b}\|^2}.
\]
\end{proof}

Define the normalized actual-line gain by
\begin{equation}
 \boxed{
 \mathfrak G_n^{\rm occ}
 :=\frac{\nu}{\mathcal R_n^2}\mathcal G_n^{\rm occ}.}
 \label{eq:v24-normalized-occurring-gain}
\end{equation}

\begin{corollary}[The one-scalar overdrive is already on the occurring line]
\label{cor:v24-line-overdrive}
Retain the V20 normalization
\[
 \mathfrak B_n^{\rm dau}
 =\frac{\|d_n\|_{\mathscr Y_n}^2}
       {\nu\mathcal R_n^2}.
\]
Then
\begin{equation}
 \boxed{
 \mathfrak G_n^{\rm occ}
 =\frac{
  |\mathfrak h_n^{\rm birth}/\mathcal R_n^2|^2
 }{\mathfrak B_n^{\rm dau}}.}
 \label{eq:v24-normalized-line-identity}
\end{equation}
If
\begin{equation}
 \mathfrak B_n^{\rm dau}\longrightarrow0,
 \qquad
 \frac{|\mathfrak h_n^{\rm birth}|}{\mathcal R_n^2}
 \ge\delta>0,
 \label{eq:v24-line-overdrive-hypotheses}
\end{equation}
then
\begin{equation}
 \boxed{
 \mathfrak G_n^{\rm occ}
 \ge\frac{\delta^2}{\mathfrak B_n^{\rm dau}}
 \longrightarrow\infty,}
 \label{eq:v24-line-overdrive}
\end{equation}
and therefore also
\begin{equation}
 \boxed{
 \frac{\nu}{\mathcal R_n^2}
 \|\Psi_n^{\rm b}\|_{\mathscr Y_n}^2
 \longrightarrow\infty.}
 \label{eq:v24-projected-adjoint-overdrive}
\end{equation}
No daughter singular value enters this conclusion.
\end{corollary}

\begin{proof}
Substitute the definitions into
\eqref{eq:v24-actual-line-susceptibility}.  The final implication is
Theorem~\ref{thm:v24-history-hierarchy}.
\end{proof}

\begin{assessment}[Termination for one scalar]
\label{ass:v24-one-scalar-stop}
For the one selected first-birth scalar, the data
\(
 \mathfrak h_n^{\rm birth}
\)
and
\(
 \mathfrak B_n^{\rm dau}
\)
already determine the sharp occurring-line susceptibility.  Computing the
complete daughter Gram, its smallest singular value, or every parent-pair
coefficient is a stronger robustness request.  Those quantities remain useful
for a reusable donor bank or cofinal source-range transport, but they cannot be
made prerequisites for closing this one scalar.
\end{assessment}

% =============================================================================
\section{The inherited fixed-datum tangent germ is the intermediate target}
\label{sec:v24-tangent}
% =============================================================================

Version~V20 already constructed the exterior source germ.  After the represented
source short, define at the occurring head state \(v(t)\)
\begin{equation}
 L_{v,t}^{\rm b}h
 :=2C_t^{\rm b}\bigl(v(t)\odot h\bigr),
 \qquad
 K_{v,t}^{\rm b}:=(L_{v,t}^{\rm b})^*L_{v,t}^{\rm b}.
 \label{eq:v24-tangent-synthesis}
\end{equation}
Let
\begin{equation}
 P_{v,t}^{\rm tan}
 :=L_{v,t}^{\rm b}(K_{v,t}^{\rm b})^\dagger
   (L_{v,t}^{\rm b})^*.
 \label{eq:v24-tangent-projection}
\end{equation}

\begin{theorem}[Actual line, tangent germ, and full quadratic range]
\label{thm:v24-tangent-hierarchy}
For almost every time,
\begin{equation}
 \boxed{
 d^{\rm birth}(t)
 =\frac12L_{v,t}^{\rm b}v(t),}
 \label{eq:v24-daughter-in-tangent}
\end{equation}
and
\begin{equation}
 \boxed{
 \operatorname{Ran}L_{v,t}^{\rm b}
 \subseteq\operatorname{Ran}C_t^{\rm b}.}
 \label{eq:v24-tangent-range-inclusion}
\end{equation}
The first-variation target gain is
\begin{equation}
 \boxed{
 \mathcal G_n^{\rm tan}
 :=\int_{J_n}
   \|P_{v,t}^{\rm tan}\Psi_n(t)\|^2\,\dd t
 =\int_{J_n}
   \left\langle J_{v,t},
    (K_{v,t}^{\rm b})^\dagger J_{v,t}\right\rangle\dd t,}
 \label{eq:v24-tangent-gain}
\end{equation}
where
\(
 J_{v,t}=(L_{v,t}^{\rm b})^*\Psi_n(t)
\).
It satisfies the exact hierarchy
\begin{equation}
 \boxed{
 \mathcal G_n^{\rm occ}
 \le\mathcal G_n^{\rm tan}
 \le\|\Psi_n^{\rm b}\|_{\mathscr Y_n}^2.}
 \label{eq:v24-three-range-hierarchy}
\end{equation}
Thus the occurring line is sufficient for one scalar, the tangent germ is the
correct fixed-datum first-variation carrier, and the complete symmetric parent
range is required only for arbitrary quadratic parent-pair reuse.
\end{theorem}

\begin{proof}
Quadratic homogeneity gives
\[
 L_{v,t}^{\rm b}v(t)
 =2C_t^{\rm b}(v(t)\odot v(t))
 =2d^{\rm birth}(t).
\]
Every value of \(L_{v,t}^{\rm b}\) is visibly in the range of \(C_t^{\rm b}\).
The Moore--Penrose projection formula and
Theorem~\ref{thm:v24-range-only} give
\eqref{eq:v24-tangent-gain}.  Since the complete history \(d_n\) belongs to
the pointwise tangent range, Cauchy--Schwarz gives
\[
 |\langle d_n,\Psi_n\rangle|^2
 =|\langle d_n,P^{\rm tan}_{v,\cdot}\Psi_n\rangle|^2
 \le\|d_n\|^2\mathcal G_n^{\rm tan}.
\]
This proves the first inequality in
\eqref{eq:v24-three-range-hierarchy}; the second follows from
\eqref{eq:v24-tangent-range-inclusion}.
\end{proof}

\begin{countertheorem}[Full quadratic-range gain can be tangent-null]
\label{cth:v24-full-range-tangent-null}
Let \(E=\mathbb R^2\) with basis \(e_1,e_2\), let
\(\mathcal Q=\operatorname{Sym}^2E\) have an orthonormal basis
\(q_{11},q_{12},q_{22}\), and let \(\mathcal Y=\mathbb R^2\) have basis
\(f_1,f_2\).  Define
\begin{equation}
 Cq_{11}=f_1,
 \qquad Cq_{12}=0,
 \qquad Cq_{22}=f_2,
 \qquad v=e_1.
 \label{eq:v24-tangent-counter-C}
\end{equation}
For \(\psi_N=f_1+Nf_2\), one has
\begin{equation}
 \boxed{
 \mathcal G_{\rm src}(C,\psi_N)=1+N^2,}
 \label{eq:v24-full-gain-diverges}
\end{equation}
whereas
\begin{equation}
 \boxed{
 \mathcal G_{\rm tan}(v;C,\psi_N)=1,
 \qquad
 \mathcal G_{\rm occ}(Cv^{\odot2},\psi_N)=1.}
 \label{eq:v24-tangent-gain-fixed}
\end{equation}
The divergent component lies in the transverse \(q_{22}\) daughter, which is
not generated by a first variation at the fixed datum \(v=e_1\).
\end{countertheorem}

\begin{proof}
The full range of \(C\) is \(\mathcal Y\), proving
\eqref{eq:v24-full-gain-diverges}.  The tangent synthesis is
\[
 L_vh=2C(v\odot h),
\]
whose range is \(\mathbb Rf_1\), because the only potentially active tangent
parent coordinates are \(q_{11}\) and \(q_{12}\).  Also
\(Cv^{\odot2}=f_1\).  Projection of \(\psi_N\) onto either relevant line gives
norm one.
\end{proof}

\begin{firewall}[The tangent germ is inherited, not a new source class]
\label{fire:v24-no-new-source-class}
Equation~\eqref{eq:v24-tangent-synthesis} is exactly the represented-shorted
version of the V20 tangent daughter germ.  Version~V24 does not introduce an
additional ambient class.  It identifies which of the already present targets
is appropriate: the occurring line for one scalar, the tangent germ for
fixed-datum first variations, and the full symmetric parent bank only for a
predeclared reusable two-parent target.
\end{firewall}

% =============================================================================
\section{Projected-adjoint amplification is the remaining physical mechanism}
\label{sec:v24-projected-adjoint}
% =============================================================================

Define the ambient and source-visible adjoint actions
\begin{equation}
 \mathcal A_n^{\rm amb}:=\|\Psi_n\|_{\mathscr Y_n}^2,
 \qquad
 \mathcal A_n^{\rm src}:=\|\Psi_n^{\rm b}\|_{\mathscr Y_n}^2.
 \label{eq:v24-two-adjoint-actions}
\end{equation}
When \(\mathcal A_n^{\rm amb}>0\), put
\begin{equation}
 \vartheta_n^{\rm ran}
 :=\frac{\mathcal A_n^{\rm src}}{\mathcal A_n^{\rm amb}}
 \in[0,1].
 \label{eq:v24-range-angle}
\end{equation}

\begin{theorem}[Range angle and the absence of a softness-driven finite gain]
\label{thm:v24-projected-adjoint-factorization}
One has the exact factorization
\begin{equation}
 \boxed{
 \mathcal A_n^{\rm src}
 =\mathcal A_n^{\rm amb}\vartheta_n^{\rm ran}.}
 \label{eq:v24-range-angle-factorization}
\end{equation}
If the normalized physical daughter gain diverges, then the normalized ambient
adjoint action diverges as well.  No collapse of a positive singular value of
\(C_{n,t}^{\rm b}\), with its range held fixed, can create this divergence.
Such a collapse can only obstruct stable transport of the source projection or
increase a separately priced parent-amplitude cost.
\end{theorem}

\begin{proof}
The factorization is the definition of \(\vartheta_n^{\rm ran}\).  Since
\(\mathcal A_n^{\rm src}\le\mathcal A_n^{\rm amb}\), divergence of the former
forces divergence of the latter.  Theorem~\ref{thm:v24-range-only} proves that
the finite source gain is independent of all nonzero singular values.
Theorem~\ref{thm:v24-projection-stability} and
Countertheorem~\ref{cth:v24-support-collapse} identify their separate cofinal
role.
\end{proof}

\begin{assessment}[Return to the upstream PDE gate]
\label{ass:v24-return-upstream}
Once \eqref{eq:v24-projected-adjoint-overdrive} occurs, no further
Moore--Penrose decomposition inside NS--A can explain it.  The object is an
actual backward writer, projected onto an actually reachable source range, and
already detected on the occurring source line.  Its growth must be evaluated
by the physical mechanisms which precede the donor short: material--heat route
mismatch, ultra-canonical frequency leverage, symmetric-strain amplification,
curvature of the material map, viscous or Bernoulli boundary current, or the
first named compactness/rigidity input.  A parent singular value is not added as
another causal mechanism for this one scalar.
\end{assessment}

% =============================================================================
\section{Corrected terminal alternative and export boundary}
\label{sec:v24-terminal}
% =============================================================================

\begin{theorem}[Version V24 range-native terminal alternative]
\label{thm:v24-terminal-alternative}
Retain one V23 terminal sequence after every inherited upstream physical gate
has passed, and normalize the first-birth scalar by
\(
 \mathcal R_n^2
\).
After passage to a priority subsequence, either branch~\textup{(G24.5)}
occurs, or one frozen finite head captures the occurring source line and exactly
one of branches~\textup{(G24.1)}--\textup{(G24.3)} occurs.  Row
\textup{(G24.4)} is appended only when its stronger target was declared
independently before the scalar was inspected.
\begin{enumerate}[label=\textup{(G24.\arabic*)},leftmargin=5.3em]
\item \emph{First-birth scalar null.}
      \begin{equation}
       \frac{|\mathfrak h_n^{\rm birth}|}{\mathcal R_n^2}
       \longrightarrow0.
       \label{eq:v24-scalar-null}
      \end{equation}
      Subject to the V22 old and represented mixed-incidence rows, the
      one-scalar ancestry card closes.  Gain outside the occurring source line
      is target-null.
\item \emph{Daughter-action survivor.}  There is \(b_*>0\) such that
      \begin{equation}
       \boxed{\mathfrak B_n^{\rm dau}\ge b_*.}
       \label{eq:v24-daughter-action-survivor}
      \end{equation}
      The occurring first-birth source already carries a nonvanishing physical
      action and is exported without an inverse-metric interpretation.
\item \emph{Actual-line projected-adjoint amplification.}  There is
      \(\delta>0\) such that
      \begin{equation}
       \mathfrak B_n^{\rm dau}\longrightarrow0,
       \qquad
       \frac{|\mathfrak h_n^{\rm birth}|}{\mathcal R_n^2}
       \ge\delta,
       \label{eq:v24-line-branch-data}
      \end{equation}
      and therefore
      \begin{equation}
       \boxed{
       \mathfrak G_n^{\rm occ}
       \le\mathfrak G_n^{\rm b,phys},
       \qquad
       \mathfrak G_n^{\rm occ}\longrightarrow\infty.}
       \label{eq:v24-line-branch-conclusion}
      \end{equation}
      This packet is returned directly to the inherited physical writer gate.
\item \emph{Explicitly stronger donor target.}  A reusable fixed-datum
      first-variation bank, an arbitrary quadratic parent bank, or cofinal
      source-range transport was declared independently.  Then the appropriate
      additional datum is respectively the tangent Gram
      \(K_{v,t}^{\rm b}\), the parent-amplitude gain
      \(\|C_t^{\rm b*}\Psi_t\|^2\), or a positive singular edge sufficient for
      Theorem~\ref{thm:v24-projection-stability}.  These stronger rows are not
      charged to the one-scalar branch.
\item \emph{All-shell occurring-source escape.}  No frozen finite head captures
      the actual first-birth line and its target scalar.  The fused all-shell
      source-line packet is retained until a source-relative critical action,
      capacity, or square-function upper stock is supplied on the same history.
\end{enumerate}
The V23 coefficient/softness factorization remains available on branch
\textup{(G24.4)}, but it is not an additional branch between
\textup{(G24.2)} and \textup{(G24.3)} for one selected scalar.
\end{theorem}

\begin{proof}
If no frozen finite head captures the occurring source line, one is in
branch~\textup{(G24.5)}.  Otherwise pass first to a subsequence on which the
normalized first-birth scalar tends to zero or has a positive lower bound.  The
first case is \textup{(G24.1)}.  In the second case, pass to a subsequence on
which \(\mathfrak B_n^{\rm dau}\) has a positive lower bound or tends to
zero.  These are \textup{(G24.2)} and \textup{(G24.3)}.  The conclusion in
the third branch is Corollary~\ref{cor:v24-line-overdrive}.  Row
\textup{(G24.4)} is not a competing subsequence branch; it records an
independently declared stronger target.  Theorems~\ref{thm:v24-three-minima},
\ref{thm:v24-tangent-hierarchy}, and
\ref{thm:v24-projection-stability} give its exact additional rows.
\end{proof}

\begin{corollary}[Corrected NS--A handoff]
\label{cor:v24-handoff}
On branch \textup{(G24.1)}, the V22 target-native scalar ancestry card may
close after its old and represented residuals pass.  On branch
\textup{(G24.2)}, export the actual daughter source and its action.  On branch
\textup{(G24.3)}, export
\begin{equation}
 \boxed{
 \left(
  d_n^{\rm birth},\Psi_n^{\rm b},
  \mathfrak h_n^{\rm birth},
  \mathfrak B_n^{\rm dau},
  \mathfrak G_n^{\rm occ}
 \right)}
 \label{eq:v24-line-export}
\end{equation}
together with the already inherited material-route, frequency, strain,
material-curvature, and boundary-current card.  No parent softness coordinate
is required unless branch \textup{(G24.4)} was independently declared.
\end{corollary}

\begin{proof}
The first statement is the V22 scalar assembly with target-null gains pruned.
The second retains the occurring source itself.  The third uses the exact
identity \eqref{eq:v24-normalized-line-identity}; every entry belongs to the
same history and no source is repriced.
\end{proof}

% =============================================================================
\section{Integrated conclusion, proof ledger, and exact next acquisition}
\label{sec:v24-conclusion}
% =============================================================================

\begin{theorem}[Version V24 range-only source-gain theorem]
\label{thm:v24-integrated}
Preserving Versions~V1--V23 verbatim, Version~V24 establishes the following.
\begin{enumerate}[label=\textup{(V24.\arabic*)},leftmargin=5.5em]
\item The V23 Moore--Penrose gain is exactly the squared projection of the
      adjoint writer onto the physical daughter range; all positive singular
      values cancel.
\item Invertible parent reparameterization leaves this source-action gain
      unchanged, while the V23 coefficient and inverse-metric factors change.
\item Native exterior-source action, parent-amplitude action, and the occurring
      source line have three different sharp Riesz minima.
\item A positive singular edge is needed for stable transport of the daughter
      support projection, not for the finite-card source gain itself.
\item The one-scalar susceptibility is exactly
      \(|\mathfrak h^{\rm birth}/\mathcal R^2|^2/
        \mathfrak B^{\rm dau}\); fixed scalar with vanishing daughter action
      forces projected-adjoint amplification on the occurring line.
\item The inherited V20 tangent source germ is the correct intermediate carrier
      for fixed-datum first variations.  Full quadratic-range gain may diverge
      entirely in a tangent-null direction.
\item The V23 source-softness coordinate is retained only for cofinal support
      transport or a separately loaded parent-amplitude donor metric.
\item The corrected terminal gate is
      Theorem~\ref{thm:v24-terminal-alternative}; the one-scalar amplification
      returns directly to the existing physical writer gate.
\end{enumerate}
No terminal PDE branch is excluded unconditionally and no global regularity
conclusion is asserted.
\end{theorem}

\begin{proof}
Items~\textup{(V24.1)--(V24.2)} are
Theorem~\ref{thm:v24-range-only},
Corollaries~\ref{cor:v24-reparameterization}--\ref{cor:v24-rescaling}, and
Countertheorem~\ref{cth:v24-factor-label}.  Item~\textup{(V24.3)} is
Theorems~\ref{thm:v24-three-minima} and
\ref{thm:v24-singular-edge-comparison}.  Item~\textup{(V24.4)} is
Theorem~\ref{thm:v24-projection-stability} and
Countertheorem~\ref{cth:v24-support-collapse}.  Item~\textup{(V24.5)} is
Theorem~\ref{thm:v24-history-hierarchy} and
Corollary~\ref{cor:v24-line-overdrive}.  Item~\textup{(V24.6)} is
Theorem~\ref{thm:v24-tangent-hierarchy} and
Countertheorem~\ref{cth:v24-full-range-tangent-null}.  The last two items are
Theorems~\ref{thm:v24-projected-adjoint-factorization} and
\ref{thm:v24-terminal-alternative}.
\end{proof}

\begingroup\small
\begin{longtable}{>{\raggedright\arraybackslash}p{0.27\textwidth}
                  >{\raggedright\arraybackslash}p{0.18\textwidth}
                  >{\raggedright\arraybackslash}p{0.47\textwidth}}
\caption{Version V24 proof-status ledger.}\label{tab:v24-ledger}\\
\toprule
\textbf{Row}&\textbf{Status}&\textbf{Exact advance or limitation}\\
\midrule
\endfirsthead
\toprule
\textbf{Row}&\textbf{Status}&\textbf{Exact advance or limitation}\\
\midrule
\endhead
V1--V23 prefix&\PROVED{} preserved
 &Complete V23 preterminal source retained byte-for-byte; no inherited theorem
 rewritten.\\[0.35em]
V23 Moore--Penrose gain&\PROVED{} reinterpreted
 &It is the squared source-range projection of the adjoint writer; every nonzero
 daughter singular value cancels.\\[0.35em]
Coefficient/softness split&\PROVED{} exact / \OBSTRUCTION{} as intrinsic source branch
 &The factors change under parent rescaling while their physical source-action
 product remains fixed.\\[0.35em]
Three action metrics&\PROVED{} sharp
 &Exterior-source, parent-amplitude, and occurring-line targets have separate
 Thomson minima and cannot be interchanged.\\[0.35em]
Source singular floor&\PROVED{} cofinal role
 &A uniform positive edge transports the support projection; without it the
 finite ranges can jump at the limit.\\[0.35em]
Actual-line susceptibility&\PROVED{} exact
 &It is \(|h/R^2|^2/\mathfrak B^{\rm dau}\) and is already sufficient for the
 one selected scalar.\\[0.35em]
Small daughter / fixed scalar&\PROVED{} corrected survivor
 &Projected-adjoint action diverges on the occurring source line; no softness
 attribution is needed.\\[0.35em]
Fixed-datum tangent germ&\PROVED{} inherited target identified
 &The V20 tangent range lies between the actual line and full quadratic daughter
 range.\\[0.35em]
Full-range robustness&\OBSTRUCTION{} explicit
 &A finite example has divergent full-range gain but fixed tangent and
 occurring-line gains.\\[0.35em]
Upstream physical explanation&\OPEN{} named
 &Frequency leverage, strain, material-map curvature, boundary current, critical
 action, compactness, or rigidity must control the projected writer.\\[0.35em]
Three-dimensional regularity&\OPEN
 &No surviving terminal branch is eliminated for every fixed datum.\\
\bottomrule
\end{longtable}
\endgroup

% =============================================================================
\section{Exact mandate after Version V24}
\label{sec:v24-next}
% =============================================================================

\begin{programme}[Evaluate the projected physical writer or terminate NS--A]
\label{prog:v24-next}
Preserve Versions~V1--V24 verbatim.  There is no automatic generic
Version~V25.

A continuation is licensed only in the following order.
\begin{enumerate}[label=\textup{(V25.\arabic*)},leftmargin=5.9em]
\item Apply the inherited material--heat route, ultra-canonical frequency,
      symmetric-strain, material-map-curvature, cell-capture, and signed
      viscous/Bernoulli boundary-current gate first.  Stop at its first survivor.
\item For one frozen return scalar, acquire only the actual first-birth scalar
      and daughter action.  Use
      \eqref{eq:v24-normalized-line-identity}; do not require the complete
      daughter Gram or a parent singular floor.
\item If the actual-line susceptibility diverges, evaluate the projected adjoint
      writer through the existing physical gate.  Return frequency leverage,
      deformation, map curvature, boundary current, or the first named PDE
      compactness/rigidity input.
\item Compute the tangent Gram only when a reusable fixed-datum first-variation
      bank was declared before inspection of the terminal scalar.
\item Compute the full parent metric and its singular edges only when arbitrary
      quadratic parent reuse or cofinal source-range transport is itself the
      target.  Treat loss of the singular floor as a support-transport
      obstruction, not as finite-card causal amplification.
\item If the first-birth scalar vanishes and the V22 old/represented residuals
      pass, terminate NS--A and export the target-native scalar ancestry card to
      NS--B.
\item If no finite head captures the occurring source line, retain one fused
      all-shell source-line escape and reopen it only with a source-relative
      critical action, capacity, or square-function upper stock on the same
      carrier.
\item Do not add another source class, response graph, path law, age law, Dini
      modulus, concentration metric, shell count, or generic compiler.
\end{enumerate}
\end{programme}

\begin{assessment}[Programme position after Version V24]
\label{ass:v24-position}
The first unresolved one-scalar object is no longer the factorization of a
Moore--Penrose gain.  It is the actual projected-adjoint writer amplification
\[
 \frac{\nu}{\mathcal R_n^2}
 \|P_{\mathbb R d_n^{\rm birth}}\Psi_n\|_{L^2}^2,
\]
which is already determined by the direct mixed scalar and daughter action.
Its explanation lies upstream in the physical material/frequency/deformation
and boundary-current packet.  A daughter singular floor remains necessary only
for a stronger cofinal source-range or parent-amplitude target.  Continuing to
split the one-scalar gain into further source metrics would therefore be
circular.
\end{assessment}

\paragraph{Version V24 history.}
Preserved the complete V23 source.  Evaluated the V23 Moore--Penrose form by
singular-value decomposition and proved that it is a range-only source-action
gain.  Separated exterior-source, parent-amplitude, occurring-line, and
fixed-datum tangent targets; exhibited rescaling, support-collapse, and
tangent-null counterexamples; proved quantitative transport of the source
support under a uniform singular edge; and reduced the selected scalar to its
exact actual-line susceptibility.  Reclassified source softness as a cofinal
support-transport or parent-metric row and returned the surviving one-scalar
overdrive to the inherited physical writer gate.  No terminal PDE closure or
global regularity conclusion was claimed.

\begin{assessment}[Append-only integrity]
\label{ass:v24-integrity}
The complete Version~V23 source preceding its sole final executable document
terminator is preserved verbatim.  Its preterminal SHA--256 digest is
\begin{center}\scriptsize\ttfamily
8566dd5a57766ccd4f6ef6598a3ae2e381b0d2fb41ae5a27120fee0415b72013
\end{center}
The present material begins at Section~\ref{sec:v24-scope}; the final command
below is again unique.
\end{assessment}

\addtocontents{toc}{\protect\endgroup}
% =============================================================================
% END VERSION V24 MATERIAL -- append later versions before final terminator
% =============================================================================

% APPEND ALL LATER VERSIONS IMMEDIATELY ABOVE THE SOLE FINAL LINE BELOW.


\clearpage
% =============================================================================
% VERSION V25: NEW MATERIAL.  The complete V1--V24 source above is preserved
% verbatim; only its sole active \end{document} line was removed before append.
% =============================================================================
\hypersetup{pdftitle={NS-A Terminal Source Concentration Programme: Cumulative V25},pdfauthor={A. Pelissier}}
\makeatletter
\addtocontents{toc}{\protect\begingroup\protect\makeatletter}
\addtocontents{toc}{\protect\def\protect\l@section{\protect\@dottedtocline{1}{0em}{4.7em}}}
\addtocontents{toc}{\protect\def\protect\l@subsection{\protect\@dottedtocline{2}{1.5em}{5.3em}}}
\addtocontents{toc}{\protect\makeatother}
\makeatother
\renewcommand{\thesection}{V25.\arabic{section}}
\renewcommand{\thesubsection}{V25.\arabic{section}.\arabic{subsection}}
\renewcommand{\theequation}{V25.\arabic{equation}}
\setcounter{section}{0}
\setcounter{equation}{0}
\renewcommand{\theHsection}{V25.\arabic{section}}
\renewcommand{\theHsubsection}{V25.\arabic{section}.\arabic{subsection}}
\renewcommand{\theHtheorem}{V25.\arabic{section}.\arabic{theorem}}
\renewcommand{\theHequation}{V25.\arabic{equation}}

\begin{center}
{\LARGE\bfseries NS--A: Version V25}\\[0.5em]
{\large Physical Daughter-Support Compression, the Exterior-Adjoint Energy Ledger,\\
and the Near-Shell Coercivity/Backscatter Gate}\\[0.5em]
{\normalsize 4 September 2026}
\end{center}
\addcontentsline{toc}{part}{Version V25: daughter-support compression and the exterior-adjoint gate}

\begin{abstract}
Version~V24 reduced one frozen return scalar to the actual-line susceptibility
\[
 \frac{\nu}{\mathcal R_n^2}
 \left\|P_{\mathbb R d_n^{\rm birth}}\Psi_n\right\|_{L^2}^2
 =\frac{|\mathfrak h_n^{\rm birth}/\mathcal R_n^2|^2}
        {\mathfrak B_n^{\rm dau}}.
\]
It then returned a divergent susceptibility to the inherited material--heat,
frequency, deformation, material-curvature, and boundary-current gate.  The
algebraic identity is exact.  Two upstream incidence rows must nevertheless be
resolved before that handoff is valid.

First, the V22 adjoint and every mixed scalar are linear in the chosen return
writer.  The displayed susceptibility is quadratic in its absolute
normalization.  A divergent value can therefore be produced by rescaling the
target writer unless its physical action has a fixed window.  Version~V25
constructs the scale-free unit-writer susceptibility and separates genuine
causal amplification from target-normalization overdrive.

Second, the V20 orthogonal short against an arbitrary represented exterior
atlas need not preserve the Fourier or source range of the quadratic daughter.
Its represented and residual vectors may acquire cancelling components outside
the support of the source being decomposed.  For one daughter target the
correct represented space is the projection of the declared atlas onto the
physical daughter range.  The resulting support-compatible residual is
source-minimal, remains inside the daughter range, and never has larger action
than the raw atlas residual.  The difference of the two mixed scalars is an
explicit source-support incidence defect.

After these corrections, the first-birth scalar sees only the projection of
the backward adjoint onto the fixed additive daughter support.  Version~V25
writes its exact terminal-zero equation and energy identity.  This adjoint is
not the material cotangent writer of the supplied NS--B deformation theorem:
the two solve different final-value equations, and a finite-dimensional
counterexample shows that bounded material-writer action does not control the
exterior adjoint without an explicit intertwiner.  The correct first gate is
instead the self-adjoint exterior secant form.  A positive shell-coercivity
margin bounds the normalized adjoint by the direct target forcing plus one
far-to-near adjoint current; failure gives an exact soft or antidissipative
source-shell direction.

For a predeclared canonical low-pass head, every pure head daughter lies in
\(K<\Lambda\le2K\).  The forcing of its projected adjoint splits exactly into
four physical pieces: direct head--daughter coupling, local exterior return
coupling involving only frequencies at most \(3K\), adjacent-octave adjoint
transport, and genuine exterior-parent high--high backscatter.  Thus an
amplified actual daughter line cannot be labelled ``ultra-canonical
frequency'' merely because the full adjoint has a remote tail.  The remote
part enters only through the displayed backscatter current.  No terminal
branch is excluded unconditionally and no global regularity conclusion is
claimed.
\end{abstract}

% =============================================================================
\section{Scope and the two omitted incidence rows}
\label{sec:v25-scope}
% =============================================================================

Retain the V20--V24 finite head, its exterior quadratic daughter synthesis,
the represented exterior source atlas, one frozen head-supported return writer
\(\phi\), and the terminal-zero adjoint writer
\begin{equation}
 -\partial_t\Psi_\phi
 +(\nu A+\mathscr L_Q^u(t)^*)\Psi_\phi
 =\mathscr M_P^u(t)^*\phi,
 \qquad
 \Psi_\phi(b)=0.
 \label{eq:v25-inherited-adjoint}
\end{equation}
Here and below the subscript in \(\Psi_\phi\) is read as \(\phi\); the
notation is kept visually separated from the following support projections.
All spaces are the real physical \(L^2\) spaces on the already fixed smooth
screen.

For an exterior first-birth history \(d^{\rm birth}\), the selected scalar is
\begin{equation}
 \mathfrak h_\phi^{\rm birth}
 :=\int_a^b\langle d^{\rm birth}(t),\Psi_\phi(t)\rangle\,\dd t.
 \label{eq:v25-inherited-birth-scalar}
\end{equation}
The V24 occurring-line identity is exact once both vectors in
\eqref{eq:v25-inherited-birth-scalar} have been physically fixed.  The first
question is whether their absolute target normalization and source support are
the ones intended by the terminal card.

\begin{firewall}[No retroactive weakening of the V20--V24 cards]
\label{fire:v25-no-retroactive-weakening}
The raw V20 atlas Pythagoras, the V21 Volterra equation, the V22 adjoint
identity, and the V23--V24 Moore--Penrose formulas remain exact for the objects
they declare.  Version~V25 proves a target-relative correction.  If the target
is the complete exterior source vector, the raw atlas projection remains
available.  If the target is one quadratic daughter and one scalar return, the
source must first be compressed to the range and Fourier support that the
quadratic head can actually produce.  Likewise, a physical absolute writer
normalization may be retained; its growth is then an explicit writer-action
survivor rather than being called causal amplification.
\end{firewall}

% =============================================================================
\section{Target normalization before causal gain}
\label{sec:v25-target-normalization}
% =============================================================================

\begin{theorem}[Homogeneity of the return writer and all selected scalars]
\label{thm:v25-target-homogeneity}
For every real scalar \(a\),
\begin{equation}
 \boxed{
 \Psi_{a\phi}=a\Psi_\phi,
 \qquad
 \mathfrak h_{a\phi}^{\rm old}=a\mathfrak h_\phi^{\rm old},
 \qquad
 \mathfrak h_{a\phi}^{\rm rep}=a\mathfrak h_\phi^{\rm rep},
 \qquad
 \mathfrak h_{a\phi}^{\rm birth}=a\mathfrak h_\phi^{\rm birth}.}
 \label{eq:v25-target-homogeneity}
\end{equation}
If \(d^{\rm birth}\ne0\) and \(\phi\ne0\), define
\begin{equation}
 \widehat\phi:=\frac{\phi}{\|\phi\|_{L^2(I)}},
 \qquad
 \widehat\Psi_\phi:=\frac{\Psi_\phi}{\|\phi\|_{L^2(I)}}.
 \label{eq:v25-unit-target}
\end{equation}
Then the unit-writer occurring-line susceptibility is
\begin{equation}
 \boxed{
 \mathfrak G_{\phi}^{\rm occ,unit}
 :=\frac{\nu}{\mathcal R^2}
   \left\|P_{\mathbb R d^{\rm birth}}
          \widehat\Psi_\phi\right\|_{L^2(I)}^2
 =\frac{\nu|\mathfrak h_\phi^{\rm birth}|^2}
        {\mathcal R^2
         \|d^{\rm birth}\|_{L^2(I)}^2
         \|\phi\|_{L^2(I)}^2}.}
 \label{eq:v25-unit-line-susceptibility}
\end{equation}
The unnormalized V24 line susceptibility equals
\begin{equation}
 \boxed{
 \mathfrak G_{\phi}^{\rm occ,raw}
 =\|\phi\|_{L^2(I)}^2
  \mathfrak G_{\phi}^{\rm occ,unit}.}
 \label{eq:v25-raw-unit-factorization}
\end{equation}
\end{theorem}

\begin{proof}
Equation~\eqref{eq:v25-inherited-adjoint} is linear in its forcing
\(\mathscr M_P^{u*}\phi\), and its terminal condition is zero.  Uniqueness
therefore gives \(\Psi_{a\phi}=a\Psi_\phi\).  Every displayed scalar is a
linear pairing with the same target writer or adjoint writer, which proves
\eqref{eq:v25-target-homogeneity}.  Orthogonal projection onto a nonzero line
gives
\[
 \left\|P_{\mathbb R d}\Psi\right\|^2
 =\frac{|\langle d,\Psi\rangle|^2}{\|d\|^2}.
\]
Apply this identity in the spacetime source Hilbert space and divide by
\(\|\phi\|^2\).  Equation~\eqref{eq:v25-raw-unit-factorization} follows
immediately.
\end{proof}

\begin{countertheorem}[Unnormalized amplification can be pure target rescaling]
\label{cth:v25-target-rescaling}
Let \(V:\mathcal X\to\mathcal Y\) be any bounded nonzero causal operator,
choose \(d\in\mathcal X\) and \(\phi\in\mathcal Y\) with
\(\langle Vd,\phi\rangle\ne0\), and set \(\phi_n=n\phi\).  Then
\begin{equation}
 \boxed{
 \|P_{\mathbb R d}V^*\phi_n\|^2
 =n^2\|P_{\mathbb R d}V^*\phi\|^2\longrightarrow\infty,}
 \label{eq:v25-rescaling-raw-diverges}
\end{equation}
whereas
\begin{equation}
 \boxed{
 \frac{\|P_{\mathbb R d}V^*\phi_n\|^2}{\|\phi_n\|^2}
 =\frac{\|P_{\mathbb R d}V^*\phi\|^2}{\|\phi\|^2}.}
 \label{eq:v25-rescaling-unit-fixed}
\end{equation}
Thus a frozen direction without a calibrated action scale does not define a
comparable cofinal gain.
\end{countertheorem}

\begin{proof}
The adjoint is linear, so \(V^*\phi_n=nV^*\phi\).  Projection and the norm are
respectively linear and quadratic.  Dividing by
\(\|\phi_n\|^2=n^2\|\phi\|^2\) proves the second identity.
\end{proof}

\begin{corollary}[Normalization-overdrive or calibrated causal amplification]
\label{cor:v25-normalization-fork}
If
\(\mathfrak G_{\phi_n}^{\rm occ,raw}\to\infty\), then after passage to a
subsequence either
\begin{equation}
 \boxed{\|\phi_n\|_{L^2(I_n)}\longrightarrow\infty}
 \label{eq:v25-target-action-overdrive}
\end{equation}
or
\begin{equation}
 \boxed{\mathfrak G_{\phi_n}^{\rm occ,unit}\longrightarrow\infty.}
 \label{eq:v25-unit-gain-overdrive}
\end{equation}
In particular, under a uniform physical target-action window
\begin{equation}
 0<c_\phi\le\|\phi_n\|_{L^2(I_n)}^2\le C_\phi<\infty,
 \label{eq:v25-target-action-window}
\end{equation}
the raw and unit-writer notions of amplification are equivalent up to fixed
constants.
\end{corollary}

\begin{proof}
Use the nonnegative product identity
\eqref{eq:v25-raw-unit-factorization}.  If the first factor is bounded on a
subsequence, divergence of the product forces divergence of the second.  The
last assertion is immediate from the two-sided window.
\end{proof}

% =============================================================================
\section{Support-compatible source ancestry}
\label{sec:v25-support-short}
% =============================================================================

The next correction is finite dimensional and precedes the adjoint estimate.
Let \(C:E\to\mathcal H\) be the complete quadratic daughter synthesis of one
frozen finite head and put
\begin{equation}
 P_C:=C(C^*C)^\dagger C^*.
 \label{eq:v25-complete-daughter-support}
\end{equation}
Thus \(P_C\) is the orthogonal projection onto \(\operatorname{Ran}C\).
Let \(\mathcal J\subseteq\mathcal H\) be the declared represented exterior
source atlas, with orthogonal projection \(R\).

\begin{definition}[Source-cyclic represented atlas]
\label{def:v25-source-cyclic-atlas}
Define
\begin{equation}
 \mathcal J_C:=P_C\mathcal J\subseteq\operatorname{Ran}C,
 \qquad
 R_C:=P_{\mathcal J_C}.
 \label{eq:v25-source-cyclic-atlas}
\end{equation}
The space \(\mathcal J_C\) is closed because \(\operatorname{Ran}C\) is
finite dimensional.  For \(d=Cq\), put
\begin{equation}
 d^{\rm rep,C}:=R_Cd,
 \qquad
 d^{\rm birth,C}:=(I-R_C)d.
 \label{eq:v25-support-compatible-split}
\end{equation}
\end{definition}

\begin{theorem}[Minimum support-compatible first-birth short]
\label{thm:v25-support-compatible-short}
For every \(d\in\operatorname{Ran}C\),
\begin{equation}
 \boxed{
 d=d^{\rm rep,C}+d^{\rm birth,C},
 \qquad
 d^{\rm rep,C}\perp d^{\rm birth,C},}
 \label{eq:v25-support-Pythagoras}
\end{equation}
and both terms belong to \(\operatorname{Ran}C\).  Moreover,
\begin{equation}
 \boxed{
 \|d^{\rm birth,C}\|
 =\operatorname{dist}(d,P_C\mathcal J)
 \le\operatorname{dist}(d,\mathcal J)
 =\|(I-R)d\|.}
 \label{eq:v25-support-residual-minimal}
\end{equation}
If \([R,P_C]=0\), then
\begin{equation}
 \boxed{R_Cd=Rd,
 \qquad d^{\rm birth,C}=(I-R)d.}
 \label{eq:v25-support-commuting-equality}
\end{equation}
Thus the raw V20 short and the source-cyclic short coincide exactly on a
represented atlas which reduces the physical daughter range.
\end{theorem}

\begin{proof}
The first identity and orthogonality are the projection theorem in the closed
subspace \(P_C\mathcal J\).  Since that subspace lies in
\(\operatorname{Ran}C\), and since \(d\in\operatorname{Ran}C\), both
components lie there.

For every \(j\in\mathcal J\),
\[
 \|d-P_Cj\|=\|P_C(d-j)\|\le\|d-j\|.
\]
Taking the infimum over \(j\) proves
\eqref{eq:v25-support-residual-minimal}.  If \(R\) commutes with \(P_C\),
then \(Rd\in\mathcal J\cap\operatorname{Ran}C=P_C\mathcal J\), while
\(d-Rd\) is orthogonal to \(P_C\mathcal J\).  Hence \(Rd\) is the
orthogonal projection of \(d\) onto that subspace.
\end{proof}

\begin{corollary}[Corrected daughter synthesis and action debit]
\label{cor:v25-corrected-birth-synthesis}
The target-native first-birth synthesis is
\begin{equation}
 \boxed{
 \widetilde C^{\rm b}:=(I-R_C)C.}
 \label{eq:v25-corrected-birth-synthesis}
\end{equation}
For every parent coefficient \(q\),
\begin{equation}
 \boxed{
 \|\widetilde C^{\rm b}q\|^2
 \le\|(I-R)Cq\|^2.}
 \label{eq:v25-corrected-birth-action}
\end{equation}
The nonnegative difference
\begin{equation}
 \boxed{
 \Delta_C^{\rm act}(q)
 :=\|(I-R)Cq\|^2-\|\widetilde C^{\rm b}q\|^2\ge0}
 \label{eq:v25-support-action-debit}
\end{equation}
is the action overestimate incurred by shorting before the represented atlas
is compressed to the physical daughter range.
\end{corollary}

\begin{proof}
Apply Theorem~\ref{thm:v25-support-compatible-short} to \(d=Cq\).
\end{proof}

\begin{theorem}[Exact scalar support-incidence residual]
\label{thm:v25-support-scalar-residual}
For an exterior adjoint writer \(\Psi\), define the raw and support-compatible
first-birth scalars
\begin{equation}
 h^{\rm raw}:=\langle(I-R)d,\Psi\rangle,
 \qquad
 h^{\rm C}:=\langle(I-R_C)d,\Psi\rangle.
 \label{eq:v25-two-birth-scalars}
\end{equation}
Then
\begin{equation}
 \boxed{
 h^{\rm raw}-h^{\rm C}
 =\langle(R_C-R)d,\Psi\rangle
 =:\Delta_C^{\rm mix}.}
 \label{eq:v25-support-mixed-residual}
\end{equation}
If \([R,P_C]=0\), then \(\Delta_C^{\rm mix}=0\).  In general, the raw and
source-cyclic cards are related only after this signed same-history residual
has been evaluated.
\end{theorem}

\begin{proof}
Subtract the two definitions in \eqref{eq:v25-two-birth-scalars}.  The
commuting case is \eqref{eq:v25-support-commuting-equality}.
\end{proof}

\begin{definition}[Support-corrected unit-writer susceptibility]
\label{def:v25-support-corrected-susceptibility}
For a nonzero support-compatible first-birth history and nonzero target writer,
put
\begin{equation}
 \boxed{
 \mathfrak G_{\phi,C}^{\rm occ,unit}
 :=\frac{\nu}{\mathcal R^2}
   \left\|P_{\mathbb R d^{\rm birth,C}}
          \widehat\Psi_\phi\right\|_{L^2(I)}^2
 =\frac{\nu|\mathfrak h_\phi^{\rm birth,C}|^2}
        {\mathcal R^2
         \|d^{\rm birth,C}\|_{L^2(I)}^2
         \|\phi\|_{L^2(I)}^2}.}
 \label{eq:v25-support-corrected-susceptibility}
\end{equation}
All later daughter-support adjoint estimates use this corrected line.  The raw
quantity of \eqref{eq:v25-unit-line-susceptibility} remains the appropriate
one only when the raw and source-cyclic shorts coincide, or when the raw
source vector itself is the declared target.
\end{definition}

\begin{remark}[Time-dependent finite cards]
\label{rem:v25-time-dependent-support}
The preceding projection statements are pointwise finite-dimensional facts.
On the smooth terminal history they may be applied at almost every time to
\(C_t\), \(P_{C_t}\), and the declared atlas \(\mathcal J_t\), then integrated.
The Moore--Penrose support projection is Borel measurable as a function of a
finite matrix, so no constant-rank hypothesis is required for the integrated
identities.  This target-relative quotient does not assert that
\(P_{C_t}j\) is a new physically occurring source history for every
\(j\in\mathcal J_t\); it asserts only equality in the daughter-supported
source quotient needed by the one declared scalar.  Full source-history
membership continues to use the raw atlas.
\end{remark}

\begin{countertheorem}[A raw atlas short can manufacture off-support pieces]
\label{cth:v25-off-support-short}
Let \(\mathcal H=\R^2\), let \(P_C\) project onto \(\R e_1\), take
\(d=e_1\), and let
\(\mathcal J_N=\operatorname{span}\{e_1+Ne_2\}\).  Then
\begin{equation}
 \boxed{
 \|(I-P_{\mathcal J_N})d\|^2
 =\frac{N^2}{1+N^2}\longrightarrow1,}
 \label{eq:v25-raw-atlas-false-birth}
\end{equation}
whereas
\begin{equation}
 \boxed{
 P_C\mathcal J_N=\R e_1,
 \qquad
 (I-P_{P_C\mathcal J_N})d=0.}
 \label{eq:v25-support-atlas-no-birth}
\end{equation}
The raw represented and residual vectors have cancelling \(e_2\)-components,
even though the occurring source has no \(e_2\)-component.
\end{countertheorem}

\begin{proof}
The orthogonal projection of \(e_1\) onto
\(e_1+Ne_2\) is \((e_1+Ne_2)/(1+N^2)\), giving
\eqref{eq:v25-raw-atlas-false-birth}.  Projection of the represented line by
\(P_C\) is exactly \(\R e_1\), proving the second statement.
\end{proof}

\begin{assessment}[Correction to the V23--V24 first-birth range]
\label{ass:v25-corrected-range}
For one quadratic daughter target, replace
\((I-R_t^{\rm out})\mathscr C_{\mathscr A}^{\rm out}\) by
\(\widetilde C_t^{\rm b}\) from
\eqref{eq:v25-corrected-birth-synthesis}, with
\(P_C\) the support of the complete exterior daughter synthesis.  Every
V23--V24 range, Riesz, and occurring-line identity then remains valid with
this corrected synthesis.  The old card passes unchanged when the represented
atlas already reduces the daughter range.  Otherwise
\eqref{eq:v25-support-mixed-residual} is the first upstream incidence row.
\end{assessment}

% =============================================================================
\section{The fixed additive daughter support}
\label{sec:v25-additive-support}
% =============================================================================

Let \(\mathscr K=-\mathscr K\subset\mathbb Z^3\setminus\{0\}\) be the finite
wavevector support of a predeclared head.  Define its exterior additive hull
\begin{equation}
 \mathscr D(\mathscr K)
 :=\{p+q\ne0:p,q\in\mathscr K\}
   \setminus\mathscr K,
 \label{eq:v25-additive-hull}
\end{equation}
and let \(\Pi_{\mathscr D}\) be the orthogonal Fourier projection onto all
exterior divergence-free polarizations at those wavevectors.  Restricted
helical output spaces are understood as subspaces of this fixed Fourier
support.

\begin{theorem}[Exact additive support of every pure head daughter]
\label{thm:v25-additive-support}
For every divergence-free \(v\) supported in \(\mathscr K\),
\begin{equation}
 \boxed{
 A^{-1/4}(I-P_{\mathscr K})\mathcal B(v)
 =\Pi_{\mathscr D}
  A^{-1/4}(I-P_{\mathscr K})\mathcal B(v).}
 \label{eq:v25-daughter-additive-support}
\end{equation}
Consequently,
\begin{equation}
 \boxed{
 \operatorname{Ran}\mathscr C_{\mathscr A}^{\rm out}
 \subseteq\operatorname{Ran}\Pi_{\mathscr D},}
 \label{eq:v25-daughter-range-in-hull}
\end{equation}
and the support-compatible first-birth source of
\eqref{eq:v25-corrected-birth-synthesis} also lies in this fixed hull.
For every writer \(\Psi\),
\begin{equation}
 \boxed{
 \langle d^{\rm birth,C},\Psi\rangle
 =\langle d^{\rm birth,C},\Pi_{\mathscr D}\Psi\rangle.}
 \label{eq:v25-scalar-sees-support}
\end{equation}
\end{theorem}

\begin{proof}
The Fourier coefficient of \((v\cdot\nabla)v\) at \(k\) is a sum over
\(p+q=k\) with \(p,q\in\mathscr K\).  Leray projection and the multiplier
\(A^{-1/4}\) preserve every nonzero Fourier wavevector.  The exterior
projection removes retained outputs, leaving only
\(\mathscr D(\mathscr K)\).  This proves
\eqref{eq:v25-daughter-additive-support} and the range inclusion.  The
corrected source lies in the complete daughter range by
Theorem~\ref{thm:v25-support-compatible-short}.  Orthogonality of the Fourier
projection proves the final identity.
\end{proof}

\begin{corollary}[Canonical low-pass daughters occupy one octave]
\label{cor:v25-one-octave-daughter}
Let
\begin{equation}
 P_K:=\one_{(0,K]}(\Lambda),
 \qquad
 Q_K:=I-P_K,
 \qquad
 \Pi_1:=\one_{(K,2K]}(\Lambda).
 \label{eq:v25-canonical-head-shell}
\end{equation}
If \(v=P_Ku\), then
\begin{equation}
 \boxed{
 A^{-1/4}Q_K\mathcal B(v)
 =\Pi_1A^{-1/4}Q_K\mathcal B(v).}
 \label{eq:v25-one-octave-source}
\end{equation}
Every support-compatible represented and first-birth component therefore
lies in the same octave \(K<\Lambda\le2K\).
\end{corollary}

\begin{proof}
If \(|p|,|q|\le K\), then \(|p+q|\le2K\).  The exterior projection imposes
\(|p+q|>K\).  Apply Theorem~\ref{thm:v25-additive-support}.
\end{proof}

% =============================================================================
\section{The exterior adjoint is not the material cotangent writer}
\label{sec:v25-adjoint-type}
% =============================================================================

The supplied NS--B deformation theorem controls a material cotangent writer
with prescribed terminal value.  In schematic form it solves
\begin{equation}
 (\partial_t+u\cdot\nabla)Z+(\nabla u)^{\mathsf T}Z=0,
 \qquad Z(b)=\zeta.
 \label{eq:v25-material-writer-equation}
\end{equation}
The V22 exterior adjoint instead solves
\eqref{eq:v25-inherited-adjoint}, has terminal value zero, contains viscosity
and the radial exterior secant, and is forced by the head-return block.  No
identity between these two writers follows from their separate definitions.

\begin{countertheorem}[Bounded material action does not control an exterior adjoint]
\label{cth:v25-material-does-not-control-adjoint}
On \(I=[0,1]\), let the material writer be the constant scalar
\(Z_N(t)\equiv1\).  Let \(\Psi_N\) solve
\begin{equation}
 -\Psi_N'(t)-N\Psi_N(t)=1,
 \qquad
 \Psi_N(1)=0.
 \label{eq:v25-ODE-adjoint}
\end{equation}
Then
\begin{equation}
 \boxed{
 \int_0^1|Z_N(t)|^2\,\dd t=1,}
 \label{eq:v25-material-action-fixed}
\end{equation}
whereas
\begin{equation}
 \boxed{
 \Psi_N(t)=\frac{e^{N(1-t)}-1}{N},
 \qquad
 \|\Psi_N\|_{L^2(0,1)}\longrightarrow\infty.}
 \label{eq:v25-adjoint-action-diverges}
\end{equation}
Thus no universal estimate can send the V24 exterior adjoint to a
material-writer action theorem without an additional intertwining equation.
The example is a finite linear model, not a Navier--Stokes trajectory.
\end{countertheorem}

\begin{proof}
Direct integration gives the displayed solution.  On \([0,1/2]\),
\[
 \Psi_N(t)\ge\frac{e^{N/2}-1}{N},
\]
so its squared \(L^2\) norm diverges.  The material norm is immediate.
\end{proof}

\begin{firewall}[Correct order of the physical gate]
\label{fire:v25-correct-gate-order}
The NS--B material-route, strain, material-curvature, and boundary-current
rows remain mandatory before a physical return writer \(\phi\) is declared.
Once \(\phi\) is frozen, growth of its exterior adjoint is a different
question.  It is governed first by the exterior adjoint equation below.
The material-writer trichotomy may be reused only after an explicit
same-history intertwiner identifies its writer with the relevant projection
of \(\Psi_\phi\).  No such intertwiner is assumed in V20--V24.
\end{firewall}

% =============================================================================
\section{Projected exterior-adjoint equation and exact energy ledger}
\label{sec:v25-projected-adjoint}
% =============================================================================

Fix the additive daughter projection \(\Pi=\Pi_{\mathscr D}\) of
Theorem~\ref{thm:v25-additive-support}.  It is independent of time and
commutes with \(A\).  Put
\begin{equation}
 \psi:=\Pi\Psi_\phi,
 \qquad
 \Psi^{\rm far}:=(I-\Pi)\Psi_\phi.
 \label{eq:v25-near-far-adjoint}
\end{equation}
Define the direct target forcing and the far-to-daughter current by
\begin{align}
 F^{\rm dir}
 &:=\Pi\mathscr M_P^u(t)^*\phi(t),
 \label{eq:v25-direct-forcing}\\
 F^{\rm tail}
 &:=-\Pi\mathscr L_Q^u(t)^*\Psi^{\rm far}(t).
 \label{eq:v25-tail-forcing}
\end{align}

\begin{theorem}[Exact daughter-support adjoint equation]
\label{thm:v25-projected-adjoint-equation}
The support-visible adjoint is the terminal-zero solution of
\begin{equation}
 \boxed{
 -\partial_t\psi
 +\nu A\psi
 +\Pi\mathscr L_Q^u(t)^*\Pi\psi
 =F^{\rm dir}+F^{\rm tail},
 \qquad
 \psi(b)=0.}
 \label{eq:v25-projected-adjoint-equation}
\end{equation}
Every support-compatible daughter scalar satisfies
\begin{equation}
 \boxed{
 \mathfrak h_\phi^{\rm birth,C}
 =\int_a^b
  \langle d^{\rm birth,C}(t),\psi(t)\rangle\,\dd t.}
 \label{eq:v25-projected-birth-scalar}
\end{equation}
The full remote adjoint enters this scalar only through the current
\(F^{\rm tail}\) in the evolution of \(\psi\).
\end{theorem}

\begin{proof}
Apply \(\Pi\) to \eqref{eq:v25-inherited-adjoint}.  Since \(\Pi\) is fixed
and commutes with \(A\),
\[
 -\partial_t\psi+\nu A\psi
 +\Pi\mathscr L_Q^{u*}\Pi\psi
 +\Pi\mathscr L_Q^{u*}(I-\Pi)\Psi_\phi
 =\Pi\mathscr M_P^{u*}\phi.
\]
Move the last adjoint term to the right.  The scalar identity is
\eqref{eq:v25-scalar-sees-support}.
\end{proof}

Define the self-adjoint exterior secant form on the daughter support by
\begin{equation}
 \mathscr S_\Pi(t)
 :=\frac12\Pi
   \bigl(\mathscr L_Q^u(t)+\mathscr L_Q^u(t)^*\bigr)
   \Pi.
 \label{eq:v25-shell-symmetric-part}
\end{equation}

\begin{theorem}[Exact backward energy ledger]
\label{thm:v25-adjoint-energy}
The projected adjoint obeys
\begin{equation}
 \boxed{
 \begin{aligned}
 \frac12\|\psi(a)\|_2^2
 &+\int_a^b
   \left[
    \nu\|A^{1/2}\psi(t)\|_2^2
    +\langle\mathscr S_\Pi(t)\psi(t),\psi(t)\rangle
   \right]\dd t\\
 &=\int_a^b
   \langle F^{\rm dir}(t)+F^{\rm tail}(t),\psi(t)\rangle\,\dd t.
 \end{aligned}}
 \label{eq:v25-adjoint-energy}
\end{equation}
The skew-adjoint part of the exterior secant contributes exactly zero.
\end{theorem}

\begin{proof}
Take the real \(L^2\) inner product of
\eqref{eq:v25-projected-adjoint-equation} with \(\psi\).  The time derivative
is \(-\frac12\partial_t\|\psi\|^2\), the Stokes term is
\(\nu\|A^{1/2}\psi\|^2\), and the quadratic form of
\(\Pi\mathscr L_Q^{u*}\Pi\) equals that of
\(\mathscr S_\Pi\).  Integrate and use \(\psi(b)=0\).
\end{proof}

\begin{definition}[Daughter-shell adjoint coercivity]
\label{def:v25-shell-coercivity}
Put
\begin{equation}
 \boxed{
 \kappa_\Pi(I)
 :=\operatorname*{ess\,inf}_{t\in I}
   \inf_{\substack{z\in\operatorname{Ran}\Pi\\\|z\|_2=1}}
   \left[
    \nu\|A^{1/2}z\|_2^2
    +\langle\mathscr S_\Pi(t)z,z\rangle
   \right].}
 \label{eq:v25-shell-coercivity}
\end{equation}
This is one finite physical coefficient on the frozen daughter support.  A
nonpositive value is a soft or antidissipative exterior-secant direction,
not a material-map conclusion.
\end{definition}

\begin{theorem}[Coercive control of the target-normalized adjoint]
\label{thm:v25-coercive-adjoint-bound}
If \(\kappa_\Pi(I)>0\), then
\begin{equation}
 \boxed{
 \kappa_\Pi(I)^2
 \|\psi\|_{L^2(I)}^2
 \le
 2\left(
   \|F^{\rm dir}\|_{L^2(I)}^2
   +\|F^{\rm tail}\|_{L^2(I)}^2
  \right).}
 \label{eq:v25-coercive-adjoint-bound}
\end{equation}
For \(\phi\ne0\), define
\begin{align}
 \mathfrak A_\Pi
 &:=\frac{\nu}{\mathcal R^2\|\phi\|_{L^2(I)}^2}
      \|\psi\|_{L^2(I)}^2,
 \label{eq:v25-normalized-shell-action}\\
 \mathfrak F^{\rm dir}_\Pi
 &:=\frac{\nu}{\mathcal R^2\|\phi\|_{L^2(I)}^2}
      \|F^{\rm dir}\|_{L^2(I)}^2,
 \label{eq:v25-normalized-direct-action}\\
 \mathfrak F^{\rm tail}_\Pi
 &:=\frac{\nu}{\mathcal R^2\|\phi\|_{L^2(I)}^2}
      \|F^{\rm tail}\|_{L^2(I)}^2.
 \label{eq:v25-normalized-tail-action}
\end{align}
Then
\begin{equation}
 \boxed{
 \mathfrak G_{\phi,C}^{\rm occ,unit}
 \le\mathfrak A_\Pi,
 \qquad
 \kappa_\Pi(I)^2\mathfrak A_\Pi
 \le2\bigl(
  \mathfrak F^{\rm dir}_\Pi
  +\mathfrak F^{\rm tail}_\Pi
 \bigr).}
 \label{eq:v25-normalized-adjoint-gate}
\end{equation}
\end{theorem}

\begin{proof}
The definition of \(\kappa_\Pi\) and
\eqref{eq:v25-adjoint-energy} give
\[
 \kappa_\Pi\|\psi\|_{L^2}^2
 \le
 \left|\int_I\langle F^{\rm dir}+F^{\rm tail},\psi\rangle\right|
 \le
 (\|F^{\rm dir}\|_{L^2}+\|F^{\rm tail}\|_{L^2})
 \|\psi\|_{L^2}.
\]
If \(\psi=0\) there is nothing to prove.  Otherwise divide by
\(\|\psi\|_{L^2}\), square, and use
\((x+y)^2\le2x^2+2y^2\).  The corrected daughter line lies in the support projection, so its
line projection is dominated by the full support projection.  This proves the
first inequality in
\eqref{eq:v25-normalized-adjoint-gate}.  Multiply the first bound by the
normalizing factor.
\end{proof}

\begin{corollary}[Exterior-adjoint softness or one actual forcing survivor]
\label{cor:v25-adjoint-trichotomy}
Suppose
\(\mathfrak G_{\phi_n,C_n}^{\rm occ,unit}\to\infty\).  After passage to a
subsequence at least one of the following occurs:
\begin{enumerate}[label=\textup{(A\arabic*)},leftmargin=3.2em]
\item \(\kappa_{\Pi_n}(I_n)\le0\) infinitely often, or its positive lower
      edge collapses to zero;
\item the normalized direct target forcing
      \(\mathfrak F^{\rm dir}_{\Pi_n}\) diverges;
\item the normalized far-to-daughter current
      \(\mathfrak F^{\rm tail}_{\Pi_n}\) diverges.
\end{enumerate}
More quantitatively, if
\begin{equation}
 \kappa_{\Pi_n}(I_n)\ge\kappa_*>0,
 \qquad
 \sup_n\left(
  \mathfrak F^{\rm dir}_{\Pi_n}
  +\mathfrak F^{\rm tail}_{\Pi_n}
 \right)<\infty,
 \label{eq:v25-adjoint-passing-window}
\end{equation}
then the unit-writer occurring-line susceptibility is uniformly bounded.
\end{corollary}

\begin{proof}
This is the contrapositive of
\eqref{eq:v25-normalized-adjoint-gate}, followed by a priority subsequence.
\end{proof}

% =============================================================================
\section{Canonical low-pass resolution: one octave or genuine backscatter}
\label{sec:v25-canonical-resolution}
% =============================================================================

Assume in this section that the head \(P_K\) of
\eqref{eq:v25-canonical-head-shell} was fixed before the terminal scalar was
read.  Put
\begin{equation}
 \Pi_2:=\one_{(2K,3K]}(\Lambda),
 \qquad
 \Pi_{>2}:=\one_{(2K,\infty)}(\Lambda),
 \qquad
 \Pi_{>3}:=\one_{(3K,\infty)}(\Lambda).
 \label{eq:v25-adjacent-shells}
\end{equation}
Write \(v=P_Ku\), \(w=Q_Ku\), and split the V21 secant blocks as
\begin{align}
 \mathscr L_vh
 &:=A^{-1/4}Q_K
    \bigl[2\mathcal B_{\rm s}(v,A^{1/4}h)\bigr],
 &
 \mathscr L_wh
 &:=A^{-1/4}Q_K
    \mathcal B_{\rm s}(w,A^{1/4}h),
 \label{eq:v25-L-split}\\
 \mathscr M_vh
 &:=A^{-1/4}P_K
    \bigl[2\mathcal B_{\rm s}(v,A^{1/4}h)\bigr],
 &
 \mathscr M_wh
 &:=A^{-1/4}P_K
    \mathcal B_{\rm s}(w,A^{1/4}h).
 \label{eq:v25-M-split}
\end{align}
Thus \(\mathscr L_Q^u=\mathscr L_v+\mathscr L_w\) and
\(\mathscr M_P^u=\mathscr M_v+\mathscr M_w\).

\begin{theorem}[Exact Fourier support of the four adjoint carriers]
\label{thm:v25-four-carrier-support}
For every head-supported \(\phi=P_K\phi\),
\begin{align}
 \boxed{\mathscr M_v^*\phi=\Pi_1\mathscr M_v^*\phi,}
 \label{eq:v25-Mv-one-octave}\\
 \boxed{
 \Pi_1\mathscr M_w^*\phi
 =\Pi_1\mathscr M_{\Pi_{(K,3K]}w}^*\phi,}
 \label{eq:v25-Mw-three-octave}\\
 \boxed{\Pi_1\mathscr L_v^*\Pi_{>3}=0.}
 \label{eq:v25-Lv-no-remote}
\end{align}
Consequently the complete forcing in the daughter-octave adjoint equation is
\begin{equation}
 \boxed{
 F_1
 =F^{\rm hv}+F^{\rm hw}+F^{\rm oct}+F^{\rm bs},}
 \label{eq:v25-four-force-sum}
\end{equation}
where
\begin{align}
 F^{\rm hv}
 &:=\mathscr M_v^*\phi,
 \label{eq:v25-force-hv}\\
 F^{\rm hw}
 &:=\Pi_1\mathscr M_{\Pi_{(K,3K]}w}^*\phi,
 \label{eq:v25-force-hw}\\
 F^{\rm oct}
 &:=-\Pi_1\mathscr L_v^*\Pi_2\Psi_\phi,
 \label{eq:v25-force-oct}\\
 F^{\rm bs}
 &:=-\Pi_1\mathscr L_w^*\Pi_{>2}\Psi_\phi.
 \label{eq:v25-force-bs}
\end{align}
The four terms are respectively direct head--daughter coupling, local
exterior return coupling, adjacent-octave adjoint transport, and genuine
exterior-parent backscatter.
\end{theorem}

\begin{proof}
It is enough to test each adjoint coefficient against a single exterior
Fourier mode \(h_r\).  In the \(v\)-term, every interaction has output
wavevector \(p+r\) with \(|p|\le K\).  If it pairs with
\(\phi\), then \(|p+r|\le K\), hence \(|r|\le2K\).  Since \(h_r\) is
exterior, \(|r|>K\), proving \eqref{eq:v25-Mv-one-octave}.

For \(h_r\) in the daughter octave and a head output, an exterior parent
wavevector \(q\) in the \(w\)-term must satisfy \(q+r=p\) with
\(|p|\le K\).  Therefore \(|q|\le|p|+|r|\le3K\), which proves
\eqref{eq:v25-Mw-three-octave}.

Finally, if \(h_r\) lies in the daughter octave, every output of its
interaction with \(v\) has modulus at most \(|r|+K\le3K\).  It is therefore
orthogonal to every adjoint input supported above \(3K\), proving
\eqref{eq:v25-Lv-no-remote}.  Substitute these three support identities into
\(F^{\rm dir}+F^{\rm tail}\) from
\eqref{eq:v25-direct-forcing}--\eqref{eq:v25-tail-forcing}.
\end{proof}

\begin{definition}[Dimensionless daughter-shell coercivity ratio]
\label{def:v25-dimensionless-coercivity}
On \(\operatorname{Ran}\Pi_1\), put
\begin{equation}
 \beta_K
 :=\operatorname*{ess\,sup}_{t\in I}
   \frac{\|\bigl(\mathscr S_{\Pi_1}(t)\bigr)_-\|_{\rm op}}
        {\nu K^2},
 \qquad
 \overline\kappa_K
 :=\frac{\kappa_{\Pi_1}(I)}{\nu K^2}.
 \label{eq:v25-beta-kappa}
\end{equation}
Here \((S)_-\) is the negative spectral part of the finite self-adjoint
shell form.
\end{definition}

\begin{lemma}[Viscous shell floor versus secant antidissipation]
\label{lem:v25-viscous-shell-floor}
One has
\begin{equation}
 \boxed{
 \overline\kappa_K\ge1-\beta_K.}
 \label{eq:v25-viscous-shell-floor}
\end{equation}
In particular, if \(\beta_K\le1-\epsilon\), then
\begin{equation}
 \boxed{
 \kappa_{\Pi_1}(I)\ge\epsilon\nu K^2.}
 \label{eq:v25-uniform-shell-floor}
\end{equation}
Thus collapse of the normalized coercive edge forces
\(\limsup\beta_K\ge1\): the symmetric exterior secant cancels at least
one viscous daughter-shell unit.
\end{lemma}

\begin{proof}
For \(z\in\operatorname{Ran}\Pi_1\),
\[
 \nu\|A^{1/2}z\|_2^2\ge\nu K^2\|z\|_2^2
\]
and
\[
 \langle\mathscr S_{\Pi_1}z,z\rangle
 \ge-\|\bigl(\mathscr S_{\Pi_1}\bigr)_-\|_{\rm op}\|z\|_2^2.
\]
Take the infimum over unit vectors and the essential infimum in time.
\end{proof}

\begin{corollary}[Four-way physical forcing alternative]
\label{cor:v25-four-force-alternative}
Assume the target-action window \eqref{eq:v25-target-action-window}, the
support-compatible first-birth short, and a uniform normalized shell floor
\begin{equation}
 \overline\kappa_{K_n}\ge\epsilon>0.
 \label{eq:v25-shell-floor-sequence}
\end{equation}
If the support-corrected unit-writer occurring-line susceptibility diverges, then at least one
of the four normalized actions
\begin{equation}
 \frac{\nu}{
  \mathcal R_n^2\|\phi_n\|_{L^2}^2(\nu K_n^2)^2}
 \|F_n^{\bullet}\|_{L^2}^2,
 \qquad
 \bullet\in\{\mathrm{hv},\mathrm{hw},\mathrm{oct},\mathrm{bs}\},
 \label{eq:v25-four-normalized-forces}
\end{equation}
diverges along a subsequence.
\end{corollary}

\begin{proof}
The shell floor gives
\(\kappa_{\Pi_1}\ge\epsilon\nu K^2\).  Apply
Theorem~\ref{thm:v25-coercive-adjoint-bound} and use
\[
 \|F^{\rm hv}+F^{\rm hw}+F^{\rm oct}+F^{\rm bs}\|^2
 \le4\sum_\bullet\|F^\bullet\|^2.
\]
If all four quantities in \eqref{eq:v25-four-normalized-forces} were bounded,
so would be the normalized shell adjoint action and hence the occurring-line
susceptibility.
\end{proof}

\begin{assessment}[Where remote frequency can actually enter]
\label{ass:v25-remote-frequency}
For a canonical low-pass head, the pure quadratic daughter source and every
support-compatible source-first-birth residual lie in
\(K<\Lambda\le2K\).  A remote adjoint tail is not itself a remote daughter
birth.  It affects the selected scalar only through
\(F^{\rm bs}\), or through the adjacent-octave term \(F^{\rm oct}\) before
reaching the daughter shell.  Thus the pre-existing far-frequency/backscatter
audit is entered through an explicit current, not by reading the support of
the full adjoint writer.
\end{assessment}

% =============================================================================
\section{Corrected V25 terminal alternative}
\label{sec:v25-terminal}
% =============================================================================

For a V24 terminal sequence, let \(d_n^{\rm raw}\) and
\(\mathfrak h_n^{\rm raw}\) denote the V20 raw-atlas first-birth history
and scalar.  Let \(d_n^{\rm C}\) and \(\mathfrak h_n^{\rm C}\) be their
source-cyclic counterparts from
\eqref{eq:v25-support-compatible-split}--\eqref{eq:v25-two-birth-scalars}, and
put
\begin{equation}
 \Delta_n^{\rm sup}
 :=\mathfrak h_n^{\rm raw}-\mathfrak h_n^{\rm C}.
 \label{eq:v25-terminal-support-residual}
\end{equation}

\begin{theorem}[Normalization, support, and exterior-adjoint terminal gate]
\label{thm:v25-terminal-gate}
Apply the inherited material-route, cell-capture, and signed boundary-current
rows first, and freeze the intended return writer before reading the return.
After passage to a priority subsequence, one of the following occurs.
\begin{enumerate}[label=\textup{(G25.\arabic*)},leftmargin=5.9em]
\item \emph{Target-normalization survivor.}  The writer-action scale has no
      uniform physical upper window, in particular
      \(\|\phi_n\|_{L^2(I_n)}\to\infty\) on a subsequence.  The raw V24 gain
      is not interpreted before this target action is retained.
\item \emph{Source-support incidence survivor.}  For some \(\delta>0\),
      \begin{equation}
       \boxed{
       \frac{|\Delta_n^{\rm sup}|}{\mathcal R_n^2}\ge\delta.}
       \label{eq:v25-support-survivor}
      \end{equation}
      The raw represented-source atlas does not reduce the physical daughter
      range on the selected scalar.
\item \emph{Daughter first-birth action.}  The support incidence passes, but
      the corrected action
      \begin{equation}
       \widetilde{\mathfrak B}_n^{\rm dau}
       :=\frac{\|d_n^{\rm C}\|_{L^2(I_n)}^2}
               {\nu\mathcal R_n^2}
       \label{eq:v25-corrected-daughter-action}
      \end{equation}
      has a positive lower bound.
\item \emph{Exterior-secant softness or antidissipation.}  The corrected
      daughter action tends to zero and the unit-writer scalar survives, but
      the shell coercivity in \eqref{eq:v25-shell-coercivity} is nonpositive
      or loses its declared normalized lower edge.
\item \emph{Direct target-forcing survivor.}  The coercive edge passes, but
      \(\mathfrak F^{\rm dir}_{\Pi_n}\) is unbounded.  On a canonical
      low-pass head this is one of the finite carriers
      \(F^{\rm hv}\) or \(F^{\rm hw}\).
\item \emph{Far-to-daughter current survivor.}  The coercive edge and direct
      forcing pass, but \(\mathfrak F^{\rm tail}_{\Pi_n}\) is unbounded.  On
      a canonical head this is adjacent-octave transport \(F^{\rm oct}\) or
      exterior-parent backscatter \(F^{\rm bs}\).
\item \emph{Target-native scalar ancestry pass.}  The target normalization,
      support incidence, shell coercivity, and direct/tail forcing windows all
      pass.  Then a fixed first-birth scalar with vanishing corrected daughter
      action is impossible.  If the first-birth scalar is negligible and the
      inherited V22 old and represented target residuals vanish, NS--A
      terminates on this scalar and exports its ancestry-complete card.
\item \emph{All-shell source-support escape.}  No predeclared finite head
      captures the occurring daughter line and its target scalar.  The fused
      all-shell packet is retained without fitting a later head.
\end{enumerate}
The supplied NS--B material-writer trichotomy may replace items
\textup{(G25.4)--(G25.6)} only after an explicit intertwiner identifies its
material cotangent writer with the daughter-support projection of the V22
exterior adjoint.
\end{theorem}

\begin{proof}
First pass to a subsequence on which the physical target-action norm is
bounded or diverges.  The latter is item~\textup{(G25.1)}.  On the bounded
branch apply Theorem~\ref{thm:v25-support-scalar-residual}; a nonnegligible
mixed residual gives item~\textup{(G25.2)}.  If it vanishes, the raw and
source-cyclic scalars agree at record scale, while
Theorem~\ref{thm:v25-support-compatible-short} shows that the corrected
source action does not exceed the raw one.

Pass next to a subsequence on which the corrected daughter action has a
positive lower bound or tends to zero.  These give item~\textup{(G25.3)} or
the small-action branch.  On the latter, a fixed scalar and the bounded target
norm force divergence of the support-corrected unit-writer susceptibility by
\eqref{eq:v25-support-corrected-susceptibility}.  Corollary
\ref{cor:v25-adjoint-trichotomy} then yields one of
items~\textup{(G25.4)--(G25.6)} unless all three physical windows pass.  If
they pass, Theorem~\ref{thm:v25-coercive-adjoint-bound} bounds the unit-line
susceptibility, contradicting a fixed scalar with vanishing action.  This is
item~\textup{(G25.7)}.  If no finite head enters the preceding argument, one
is in item~\textup{(G25.8)}.  The last qualification is
Countertheorem~\ref{cth:v25-material-does-not-control-adjoint} and
Firewall~\ref{fire:v25-correct-gate-order}.
\end{proof}

\begin{corollary}[Canonical low-pass refinement of the amplified branch]
\label{cor:v25-canonical-terminal-refinement}
On a predeclared canonical low-pass head satisfying the target-action and
source-support gates, a fixed scalar with vanishing daughter action yields at
least one of
\begin{equation}
 \boxed{
 \begin{gathered}
 \text{viscous daughter-shell cancellation by the symmetric secant},\\
 \text{direct head--daughter target coupling},\\
 \text{local exterior return coupling below }3K,\\
 \text{adjacent-octave adjoint transport},\\
 \text{or genuine exterior-parent high--high backscatter}.
 \end{gathered}}
 \label{eq:v25-canonical-five-way}
\end{equation}
No sixth branch called ``remote support of the daughter source'' exists.
\end{corollary}

\begin{proof}
Use Lemma~\ref{lem:v25-viscous-shell-floor} and Corollary
\ref{cor:v25-four-force-alternative}.  The daughter itself is one-octave by
Corollary~\ref{cor:v25-one-octave-daughter}.
\end{proof}

% =============================================================================
\section{Integrated conclusion, proof ledger, and exact next acquisition}
\label{sec:v25-conclusion}
% =============================================================================

\begin{theorem}[Version V25 support-native exterior-adjoint theorem]
\label{thm:v25-integrated}
Preserving Versions~V1--V24 verbatim, Version~V25 establishes the following.
\begin{enumerate}[label=\textup{(V25.\arabic*)},leftmargin=5.6em]
\item The V22 adjoint, every return scalar, and the V24 raw line gain have
      exact target-writer homogeneity.  Division by the target action gives
      the canonical unit-writer susceptibility.
\item Without a target-action window, raw gain divergence may be pure writer
      rescaling; with such a window, raw and unit gain are equivalent.
\item For the one daughter-supported scalar, the target-native represented
      quotient is the projection of the declared atlas onto the complete
      physical daughter range.  Its residual is support preserving and never
      larger than the raw V20 residual; full source-history membership remains
      a stronger raw-atlas statement.
\item Failure of the raw atlas to reduce the daughter range returns one
      positive action debit and one signed mixed support-incidence residual.
\item Every pure finite-head daughter lies in the additive Fourier hull of the
      head.  For a canonical low-pass head it lies exactly in the first
      exterior octave \(K<\Lambda\le2K\).
\item The support-visible exterior adjoint satisfies the exact terminal-zero
      equation \eqref{eq:v25-projected-adjoint-equation} and energy ledger
      \eqref{eq:v25-adjoint-energy}.
\item The exterior adjoint is not the NS--B material cotangent writer.  A
      finite counterexample excludes any automatic action bound from the
      latter to the former.
\item A positive daughter-shell secant margin bounds the normalized adjoint by
      direct target forcing plus one far-to-daughter current.  Divergence
      returns shell softness, direct forcing, or tail current.
\item On a canonical head the forcing splits exactly into head--daughter,
      local exterior, adjacent-octave, and high--high backscatter carriers.
\item The corrected terminal gate is Theorem~\ref{thm:v25-terminal-gate}.
\end{enumerate}
No terminal branch is excluded for every fixed datum and no global
three-dimensional regularity conclusion is asserted.
\end{theorem}

\begin{proof}
Items~\textup{(V25.1)--(V25.2)} are
Theorem~\ref{thm:v25-target-homogeneity}, Countertheorem
\ref{cth:v25-target-rescaling}, and Corollary
\ref{cor:v25-normalization-fork}.  Items~\textup{(V25.3)--(V25.4)} are
Theorems~\ref{thm:v25-support-compatible-short} and
\ref{thm:v25-support-scalar-residual}, together with Countertheorem
\ref{cth:v25-off-support-short}.  Item~\textup{(V25.5)} is
Theorem~\ref{thm:v25-additive-support} and Corollary
\ref{cor:v25-one-octave-daughter}.  Items~\textup{(V25.6)--(V25.8)} are
Theorems~\ref{thm:v25-projected-adjoint-equation},
\ref{thm:v25-adjoint-energy}, and
\ref{thm:v25-coercive-adjoint-bound}, with Countertheorem
\ref{cth:v25-material-does-not-control-adjoint}.  Item~\textup{(V25.9)} is
Theorem~\ref{thm:v25-four-carrier-support} and Corollary
\ref{cor:v25-four-force-alternative}.  The final item is
Theorem~\ref{thm:v25-terminal-gate}.
\end{proof}

\begingroup\small
\begin{longtable}{>{\raggedright\arraybackslash}p{0.27\textwidth}
                  >{\raggedright\arraybackslash}p{0.18\textwidth}
                  >{\raggedright\arraybackslash}p{0.47\textwidth}}
\caption{Version V25 proof-status ledger.}\label{tab:v25-ledger}\\
\toprule
\textbf{Row}&\textbf{Status}&\textbf{Exact advance or limitation}\\
\midrule
\endfirsthead
\toprule
\textbf{Row}&\textbf{Status}&\textbf{Exact advance or limitation}\\
\midrule
\endhead
V1--V24 prefix&\PROVED{} preserved
 &Complete V24 preterminal source retained byte-for-byte; no inherited theorem
 rewritten.\\[0.35em]
Target-writer normalization&\PROVED{} exact
 &Raw line gain is target-action times the unit-writer susceptibility; writer
 rescaling is separated before causal amplification is claimed.\\[0.35em]
Raw exterior atlas short&\PROVED{} exact / \OBSTRUCTION{} for supported target
 &It can create cancelling off-support components and overstate the one-daughter
 first-birth action.\\[0.35em]
Source-cyclic atlas&\PROVED{} minimum
 &Projection of the represented atlas into the complete daughter range gives
 the least support-compatible source residual.\\[0.35em]
Support mixed incidence&\PROVED{} exact signed row
 &Raw and support-compatible return scalars differ by
 \(\langle(R_C-R)d,\Psi\rangle\).\\[0.35em]
Additive daughter support&\PROVED{} exact
 &A finite head produces only its pairwise sumset; a canonical ball head
 produces daughters only in \((K,2K]\).\\[0.35em]
Material-writer handoff&\OBSTRUCTION{} without intertwiner
 &The material cotangent writer and exterior terminal-zero adjoint solve
 different equations; bounded action of one does not control the other.\\[0.35em]
Projected adjoint ledger&\PROVED{} exact
 &One fixed-support equation and energy identity retain direct target forcing,
 far-to-daughter current, viscosity, and the symmetric exterior secant.\\[0.35em]
Shell coercivity&\PROVED{} conditional upper bound
 &A positive edge bounds the unit-writer susceptibility; collapse returns a
 finite soft or antidissipative daughter-shell direction.\\[0.35em]
Canonical frequency resolution&\PROVED{} exact
 &The forcing is head--daughter, local exterior, adjacent-octave, or high--high
 backscatter; remote adjoint support is not a daughter birth.\\[0.35em]
Terminal concentration closure&\OPEN{} sharply reduced
 &Must populate the target norm, support residual, shell edge, and direct/tail
 forcing on the literal fixed-datum card.\\[0.35em]
Three-dimensional regularity&\OPEN
 &No surviving packet has been eliminated for every datum.\\
\bottomrule
\end{longtable}
\endgroup

% =============================================================================
\section{Exact mandate after Version V25}
\label{sec:v25-next}
% =============================================================================

\begin{programme}[Populate the support-native adjoint card or terminate]
\label{prog:v25-next}
Preserve Versions~V1--V25 verbatim.  Do not introduce another source class,
response graph, path law, age law, Dini modulus, concentration metric, shell
count, or Moore--Penrose factorization.

Proceed only in the following order.
\begin{enumerate}[label=\textup{(V26.\arabic*)},leftmargin=5.9em]
\item Retain the inherited material-route, cell-capture, and complete signed
      boundary-current gate which produces the physical return writer.  Record
      its absolute action scale before evaluating the return.
\item Replace the raw represented-source short by the source-cyclic short of
      \eqref{eq:v25-corrected-birth-synthesis}.  Evaluate
      \(\Delta_n^{\rm sup}\) before using the V24 line susceptibility.
\item On one frozen finite head, populate the additive-support projection,
      the shell secant edge \(\kappa_{\Pi_n}\), and the two actions
      \(\mathfrak F^{\rm dir}_{\Pi_n}\),
      \(\mathfrak F^{\rm tail}_{\Pi_n}\).
\item On a predeclared canonical low-pass head, replace the two forcing actions
      by the four literal reads
      \(F^{\rm hv},F^{\rm hw},F^{\rm oct},F^{\rm bs}\).  Stop at the first
      nonvanishing normalized carrier.  Send only \(F^{\rm bs}\) to the
      existing far-frequency/high--high backscatter audit.
\item Apply the supplied NS--B material-writer deformation theorem to this
      adjoint only if an explicit same-history intertwiner is acquired.  Do not
      identify the two writers from compatible norms or marginal supports.
\item If the target norm, support incidence, shell edge, and direct/tail
      forcing windows pass, the small-daughter/fixed-scalar branch is excluded
      by \eqref{eq:v25-normalized-adjoint-gate}.  Then use the V22 old and
      represented residuals to terminate or export the scalar ancestry card.
\item If no finite head captures the additive support and scalar, retain the
      fused all-shell source-support escape; do not fit a head after reading
      the remote adjoint.
\end{enumerate}
\end{programme}

\begin{assessment}[Programme position after Version V25]
\label{ass:v25-position}
The V24 actual-line ratio is no longer sent directly to a material cotangent
writer.  The first unpopulated physical card is
\[
 \left(
  \|\phi_n\|_{L^2}^2,
  \Delta_n^{\rm sup},
  \kappa_{\Pi_n}(I_n),
  \mathfrak F^{\rm dir}_{\Pi_n},
  \mathfrak F^{\rm tail}_{\Pi_n}
 \right),
\]
or, on a canonical head, the corresponding four force components of
\eqref{eq:v25-four-force-sum}.  These quantities are all attached to the
literal occurring daughter, the frozen return writer, and one fixed exterior
adjoint equation.  Their population is a physical PDE calculation.  Adding a
larger source metric or response carrier before these five reads would be
circular.
\end{assessment}

\paragraph{Version V25 history.}
Preserved the complete V24 source.  Separated target normalization from causal
gain; replaced the raw represented-source short by the minimum
support-compatible daughter short; proved the exact additive Fourier support
of a finite-head daughter; derived the fixed-support exterior-adjoint equation
and its backward energy ledger; showed by a finite counterexample that the
material cotangent writer does not control this adjoint without an
intertwiner; and resolved the canonical daughter-shell forcing into direct
head coupling, local exterior coupling, adjacent-octave transport, and genuine
high--high backscatter.  No terminal Dini gain, upper critical-action theorem,
or global regularity conclusion was claimed.

\begin{assessment}[Append-only integrity]
\label{ass:v25-integrity}
The complete Version~V24 source preceding its sole final executable document
terminator is preserved verbatim.  Its preterminal SHA--256 digest is
\begin{center}\scriptsize\ttfamily
9af746bbf4955a5a54db82581d9fd607d34cef3fdf25765f8336b268bdb04820
\end{center}
The present material begins at Section~\ref{sec:v25-scope}; the final command
below is again unique.
\end{assessment}

\addtocontents{toc}{\protect\endgroup}
% =============================================================================
% END VERSION V25 MATERIAL -- append later versions before final terminator
% =============================================================================

% APPEND ALL LATER VERSIONS IMMEDIATELY ABOVE THE SOLE FINAL LINE BELOW.


\clearpage
% =============================================================================
% VERSION V26: NEW MATERIAL.  The complete V1--V25 source above is preserved
% verbatim; only its sole final executable \end{document} line was removed.
% =============================================================================
\makeatletter
\addtocontents{toc}{\protect\begingroup\protect\makeatletter}
\addtocontents{toc}{\protect\def\protect\l@section{\protect\@dottedtocline{1}{0em}{4.7em}}}
\addtocontents{toc}{\protect\def\protect\l@subsection{\protect\@dottedtocline{2}{1.5em}{5.3em}}}
\addtocontents{toc}{\protect\makeatother}
\makeatother
\renewcommand{\thesection}{V26.\arabic{section}}
\renewcommand{\thesubsection}{V26.\arabic{section}.\arabic{subsection}}
\renewcommand{\theequation}{V26.\arabic{equation}}
\setcounter{section}{0}
\setcounter{equation}{0}
\renewcommand{\theHsection}{V26.\arabic{section}}
\renewcommand{\theHsubsection}{V26.\arabic{section}.\arabic{subsection}}
\renewcommand{\theHtheorem}{V26.\arabic{section}.\arabic{theorem}}
\renewcommand{\theHequation}{V26.\arabic{equation}}

\begin{center}
{\LARGE\bfseries NS--A: Version V26}\\[0.5em]
{\large The Support-Corrected Source Envelope, Signed Adjoint Work,\\
and the Scale-Local Secant-Amplification Gate}\\[0.5em]
{\normalsize 4 September 2026}
\end{center}
\addcontentsline{toc}{part}{Version V26: support-corrected source envelope and signed adjoint-work gate}

\begin{abstract}
Version~V25 placed the terminal-zero exterior adjoint on the fixed additive
Fourier support of a finite quadratic daughter and obtained a sufficient
coercive estimate from the minimum symmetric shell edge and the full
\(L^2\)-actions of its direct and tail forcings.  Those objects remain exact
and useful when the target is a reusable daughter bank or uniform control of
every direction in the exterior octave.  They are stronger than the declared
one-scalar target, however.  The minimum shell edge may fail in a direction
which the selected daughter residual does not use, and a forcing norm may diverge entirely in
a direction orthogonal to the actual adjoint path.  Neither event by itself
explains the selected return scalar.

The present continuation moves one arrow upstream.  From the already
selected support-corrected daughter residual \(d^{\rm C}(t)\), it constructs
the unique smallest \emph{fixed spatial source envelope} containing its
essential range.  This envelope is fixed before the return writer is read and
is contained in the first exterior octave on a canonical low-pass head.  The
selected scalar sees only the projection \(\chi\) of the V25 adjoint onto this
envelope.  Projecting the exact terminal-zero equation produces a no-remainder
ledger with seven signed carriers: the four V25 physical forcings, heat
mixing across the source envelope, symmetric secant mixing, and skew
secant/circulation mixing.  The broad shell edge and the seven full force
norms are not required.

The exact energy identity yields a sharp finite alternative.  Either the
negative symmetric-secant work carries at least half of the positive
endpoint-plus-viscous adjoint stock, or one of the seven signed carrier works
carries at least one fourteenth of that stock.  Every mixed work can be read
by two positive polarizations after the physical parabolic scale has been
fixed.  Under the inherited small-daughter/fixed-return hypothesis, the
normalized positive stock diverges; consequently one of these actual signed
carriers diverges.  Target-null shell softness and target-null force action
are removed from the one-scalar stopping rule.

On the antidissipative branch the quadratic exterior secant can be evaluated
further.  A periodic Lipschitz commutator identity gives, for an adjoint in
\(K<\Lambda\le2K\),
\[
 |\langle \mathscr S\chi,\chi\rangle|
 \le C\bigl(
  \|\nabla P_Ku\|_\infty
  +\|\nabla P_{(K,4K]}u\|_\infty
 \bigr)\|\chi\|_2^2.
\]
Thus actual secant antidissipation of viscous size forces a same-history
scale-local Lipschitz pulse of size \(c\nu K^2\).  Bernstein then gives a
finite enstrophy packet
\(
 \|\nabla P_{\le4K}u(t_*)\|_2^2\ge c\nu^2K
\)
at some time in the selected interval.  At the canonical energy cutoff
\(K=\mathcal R^2/E_*\), this is
\(c\nu^2\mathcal R^2/E_*\).  Remote frequencies do not enter this secant
quadratic form; they enter only through the already isolated high--high
backscatter work.  No terminal branch is excluded for every datum and no
global three-dimensional regularity conclusion is claimed.
\end{abstract}

% =============================================================================
\section{Why the V25 robust-bank gate is not the one-scalar gate}
\label{sec:v26-entry}
% =============================================================================

Retain the complete V25 card after the target normalization and
source-support incidence rows have passed.  On a predeclared canonical
low-pass head, write
\[
 P_K=\one_{(0,K]}(\Lambda),\qquad
 Q_K=I-P_K,\qquad
 \Pi_1=\one_{(K,2K]}(\Lambda),
\]
and retain the V25 terminal-zero daughter-octave adjoint
\begin{equation}
 -\partial_t\psi
 +\nu A\psi
 +\Pi_1\mathscr L_Q^u(t)^*\Pi_1\psi
 =F_1,
 \qquad
 \psi(b)=0,
 \label{eq:v26-inherited-octave-adjoint}
\end{equation}
where
\begin{equation}
 F_1
 =F^{\rm hv}+F^{\rm hw}+F^{\rm oct}+F^{\rm bs}
 \label{eq:v26-inherited-four-forces}
\end{equation}
is the exact four-carrier decomposition of
\eqref{eq:v25-four-force-sum}.  The support-corrected first-birth history is
denoted
\begin{equation}
 d^{\rm C}(t):=d^{\rm birth,C}(t),
 \qquad
 \mathfrak h^{\rm C}
 :=\int_I\langle d^{\rm C}(t),\psi(t)\rangle\,\dd t.
 \label{eq:v26-inherited-source-scalar}
\end{equation}
By V25, \(d^{\rm C}(t)\in\operatorname{Ran}\Pi_1\) almost everywhere.

\begin{firewall}[No new source class and no retroactive weakening]
\label{fire:v26-no-new-source}
The fixed envelope introduced below is generated by the essential range of
the already selected support-corrected daughter residual \(d^{\rm C}\).  It is not a new
Navier--Stokes forcing, a refitted Fourier head, or a larger daughter bank.
It also does not promote the target-relative residual \(d^{\rm C}\) to a
second physical occurrence: a full ancestry claim continues to use the raw
daughter and the original represented-source atlas.
The V25 shell edge and forcing-action estimates remain correct sufficient
conditions for their declared robust-bank target.  Version~V26 proves that
they are not necessary prerequisites for the weaker target consisting of one
selected support-corrected daughter residual and one frozen return scalar.
\end{firewall}

\begin{countertheorem}[A negative broad shell edge may be target null]
\label{cth:v26-broad-edge-null}
Let \(I=[0,1]\), \(\mathcal H=\R^2\), and let
\[
 T_N=\begin{pmatrix}1&0\\0&-N\end{pmatrix},
 \qquad f=e_1.
\]
The terminal-value problem
\begin{equation}
 -\psi_N'(t)+T_N\psi_N(t)=f,
 \qquad \psi_N(1)=0
 \label{eq:v26-edge-counter-ode}
\end{equation}
has the solution
\begin{equation}
 \boxed{
 \psi_N(t)=\bigl(1-e^{t-1}\bigr)e_1.}
 \label{eq:v26-edge-counter-solution}
\end{equation}
The minimum symmetric edge of \(T_N\) is \(-N\), whereas the complete
solution lies in \(\R e_1\) and sees the strictly positive edge one.  Thus a
nonpositive minimum edge on a larger shell need not obstruct a scalar whose
source envelope avoids the soft direction.
\end{countertheorem}

\begin{proof}
The second component of \eqref{eq:v26-edge-counter-ode} has zero forcing and
zero terminal value, so it vanishes.  The first component solves
\(-x'+x=1\), \(x(1)=0\), whose unique solution is
\(x(t)=1-e^{t-1}\).  The spectral statements are immediate.
\end{proof}

\begin{countertheorem}[A force action may diverge in a work-null direction]
\label{cth:v26-force-norm-null}
Let \(\chi=e_1\in\R^2\) and
\(G_N=e_1+Ne_2\).  Then
\begin{equation}
 \boxed{
 \langle G_N,\chi\rangle=1,
 \qquad
 \|G_N\|_2^2=1+N^2\longrightarrow\infty.}
 \label{eq:v26-force-norm-null}
\end{equation}
Hence the full action of one forcing component is a strictly stronger output
than its contribution to the selected scalar.  The same example applies in
spacetime by taking constant functions on a unit interval.
\end{countertheorem}

% =============================================================================
\section{The minimal fixed envelope of the selected daughter residual}
\label{sec:v26-occurring-envelope}
% =============================================================================

\begin{definition}[Essential fixed source envelope]
\label{def:v26-source-envelope}
Let \(d^{\rm C}\in L^2(I;\operatorname{Ran}\Pi_1)\).  Define its
source-null space and essential fixed spatial envelope by
\begin{equation}
 \mathcal N_d
 :=\left\{
  q\in\operatorname{Ran}\Pi_1:
  \langle d^{\rm C}(t),q\rangle=0
  \text{ for a.e. }t
 \right\},
 \qquad
 \boxed{\mathcal E_d:=\mathcal N_d^\perp.}
 \label{eq:v26-source-envelope}
\end{equation}
Let \(P_d\) be the orthogonal projection onto \(\mathcal E_d\).  Since
\(\operatorname{Ran}\Pi_1\) is finite dimensional, both spaces are closed.
Equivalently, \(\mathcal E_d\) is the essential linear span of the values of
\(d^{\rm C}\).
\end{definition}

\begin{theorem}[Canonical minimality of the fixed source envelope]
\label{thm:v26-source-envelope-minimal}
The envelope has the following properties.
\begin{enumerate}[label=\textup{(E\arabic*)},leftmargin=3.4em]
\item \(P_dd^{\rm C}(t)=d^{\rm C}(t)\) for almost every \(t\in I\).
\item If \(P\) is any time-independent orthogonal projection satisfying
      \(Pd^{\rm C}(t)=d^{\rm C}(t)\) almost everywhere, then
      \begin{equation}
       \boxed{P_dP=P_d,
       \qquad \operatorname{Ran}P_d\subseteq\operatorname{Ran}P.}
       \label{eq:v26-envelope-minimality}
      \end{equation}
\item The selected scalar depends only on
      \begin{equation}
       \boxed{\chi:=P_d\psi,}
       \label{eq:v26-envelope-adjoint}
      \end{equation}
      namely
      \begin{equation}
       \boxed{
       \mathfrak h^{\rm C}
       =\int_I\langle d^{\rm C}(t),\chi(t)\rangle\,\dd t.}
       \label{eq:v26-envelope-scalar}
      \end{equation}
\end{enumerate}
The construction uses the selected daughter residual but not the return writer
or the value of its scalar.
\end{theorem}

\begin{proof}
Choose an orthonormal basis \(q_1,\ldots,q_m\) of \(\mathcal N_d\).
For every \(j\), the scalar \(\langle d^{\rm C}(t),q_j\rangle\) vanishes
outside a null set.  The union of these finitely many null sets is null, and
outside it \(d^{\rm C}(t)\perp\mathcal N_d\), hence
\(d^{\rm C}(t)\in\mathcal E_d\).  This proves item~\textup{(E1)}.

If \(Pd^{\rm C}=d^{\rm C}\) almost everywhere and
\(q\in(\operatorname{Ran}P)^\perp\), then
\(\langle d^{\rm C}(t),q\rangle=0\) almost everywhere.  Thus
\((\operatorname{Ran}P)^\perp\subseteq\mathcal N_d\), and taking orthogonal
complements gives
\(\mathcal E_d\subseteq\operatorname{Ran}P\).  This is item~\textup{(E2)}.
Finally,
\[
 \langle d^{\rm C},\psi\rangle
 =\langle P_dd^{\rm C},\psi\rangle
 =\langle d^{\rm C},P_d\psi\rangle
\]
almost everywhere; integration gives item~\textup{(E3)}.  The definition of
\(\mathcal E_d\) contains no occurrence of \(\phi\) or
\(\mathfrak h^{\rm C}\).
\end{proof}

\begin{assessment}[Three target levels]
\label{ass:v26-three-target-levels}
The V23 pointwise daughter range is the correct source class for all
quadratic first-birth daughters of the frozen head.  The V25 additive octave
is the minimum fixed \(A\)-reducing Fourier envelope available without
tracking source-range motion.  The space \(\mathcal E_d\) is the minimum
fixed spatial envelope for the already selected support-corrected residual.
These are successively stronger use classes.  Version~V26 uses the last one
because the declared target is one return scalar; V25 remains available when
uniform reuse over the complete octave is intended.
\end{assessment}

% =============================================================================
\section{Exact source-envelope adjoint equation}
\label{sec:v26-envelope-equation}
% =============================================================================

Put
\begin{equation}
 \eta:=(I-P_d)\psi,
 \qquad
 \psi=\chi+\eta.
 \label{eq:v26-adjoint-envelope-split}
\end{equation}
Since \(P_d\le\Pi_1\), both \(\chi\) and \(\eta\) lie in the daughter octave.
On that octave define
\begin{equation}
 \mathscr S_1(t)
 :=\frac12\Pi_1
  \bigl(\mathscr L_Q^u(t)+\mathscr L_Q^u(t)^*\bigr)\Pi_1,
 \qquad
 \mathscr K_1(t)
 :=\frac12\Pi_1
  \bigl(\mathscr L_Q^u(t)-\mathscr L_Q^u(t)^*\bigr)\Pi_1.
 \label{eq:v26-symmetric-skew-secant}
\end{equation}
Thus \(\mathscr S_1^*=\mathscr S_1\),
\(\mathscr K_1^*=-\mathscr K_1\), and
\begin{equation}
 \Pi_1\mathscr L_Q^u(t)^*\Pi_1
 =\mathscr S_1(t)-\mathscr K_1(t).
 \label{eq:v26-adjoint-secant-split}
\end{equation}

For \(\bullet\in\{\mathrm{hv},\mathrm{hw},\mathrm{oct},\mathrm{bs}\}\),
put
\begin{equation}
 G^\bullet:=P_dF^\bullet.
 \label{eq:v26-projected-four-forces}
\end{equation}
Define the three source-envelope currents
\begin{align}
 G^A&:=-\nu P_dA\eta,
 \label{eq:v26-heat-envelope-current}\\
 G^S&:=-P_d\mathscr S_1\eta,
 \label{eq:v26-symmetric-envelope-current}\\
 G^K&:=P_d\mathscr K_1\eta.
 \label{eq:v26-skew-envelope-current}
\end{align}

\begin{theorem}[No-remainder source-envelope adjoint equation]
\label{thm:v26-envelope-adjoint-equation}
The target-visible adjoint \(\chi=P_d\psi\) is the terminal-zero solution of
\begin{equation}
 \boxed{
 \begin{aligned}
 -\partial_t\chi
 &+\nu P_dAP_d\chi
 +P_d\mathscr S_1P_d\chi
 -P_d\mathscr K_1P_d\chi\\
 &=G^{\rm hv}+G^{\rm hw}+G^{\rm oct}+G^{\rm bs}
   +G^A+G^S+G^K,
 \qquad \chi(b)=0.
 \end{aligned}}
 \label{eq:v26-envelope-adjoint-equation}
\end{equation}
The three additional terms have the following exact meanings.
\begin{enumerate}[label=\textup{(R\arabic*)},leftmargin=3.4em]
\item \(G^A\) is heat mixing between the support-corrected source envelope and the
      target-null part of the additive octave.
\item \(G^S\) is symmetric exterior-secant mixing across that envelope.
\item \(G^K\) is skew secant circulation or phase transport across that
      envelope.
\end{enumerate}
They are response currents of the same adjoint equation, not new nonlinear
source births.
\end{theorem}

\begin{proof}
Apply the fixed projection \(P_d\) to
\eqref{eq:v26-inherited-octave-adjoint}.  Since \(P_d\) is independent of
time,
\(-P_d\partial_t\psi=-\partial_t\chi\).  Substitute
\(\psi=\chi+\eta\), use \eqref{eq:v26-adjoint-secant-split}, and move every
term containing \(\eta\) to the right.  The four projected V25 forcings give
\eqref{eq:v26-projected-four-forces}; the remaining three terms are exactly
\eqref{eq:v26-heat-envelope-current}--\eqref{eq:v26-skew-envelope-current}.
The terminal condition follows from \(\psi(b)=0\).
\end{proof}

\begin{corollary}[Exact reducing criteria]
\label{cor:v26-envelope-reducing}
The heat current vanishes identically when \(\mathcal E_d\) reduces \(A\):
\begin{equation}
 [P_d,A]=0\quad\Longrightarrow\quad G^A=0.
 \label{eq:v26-heat-current-zero}
\end{equation}
The symmetric and skew currents vanish respectively when
\begin{equation}
 P_d\mathscr S_1(I-P_d)=0,
 \qquad
 P_d\mathscr K_1(I-P_d)=0.
 \label{eq:v26-secant-current-zero}
\end{equation}
If all three conditions hold, the support-corrected source envelope is reducing for
the complete daughter-octave adjoint generator.
\end{corollary}

\begin{proof}
Each statement is its defining off-diagonal block.  Since \(A\) and
\(\mathscr S_1\) are self-adjoint and \(\mathscr K_1\) is skew-adjoint,
vanishing of the displayed block also annihilates its adjoint block.
\end{proof}

\begin{assessment}[Why skew circulation is a separate row]
\label{ass:v26-skew-circulation}
The broad shell edge of V25 reads only the symmetric form
\(\nu A+\mathscr S_1\).  A skew transport can move adjoint mass from the
source-null complement into \(\mathcal E_d\) without changing the total
shell energy.  The current \(G^K\) is therefore invisible to a minimum
symmetric edge but visible to the selected scalar.  This is a range-incidence
or phase-transport row, not an additional source action.
\end{assessment}

% =============================================================================
\section{Signed work before force action}
\label{sec:v26-signed-work}
% =============================================================================

Define the positive endpoint-plus-viscous stock
\begin{equation}
 \boxed{
 \mathcal P_d[\chi]
 :=\frac12\|\chi(a)\|_2^2
   +\nu\int_I\|A^{1/2}\chi(t)\|_2^2\,\dd t.}
 \label{eq:v26-positive-adjoint-stock}
\end{equation}
The symmetric-secant work and its negative part are
\begin{align}
 \mathcal N_d[\chi]
 &:=\int_I\langle\mathscr S_1(t)\chi(t),\chi(t)\rangle\,\dd t,
 \label{eq:v26-symmetric-secant-work}\\
 \mathcal N_d^-[\chi]
 &:=\int_I
   \bigl(-\langle\mathscr S_1(t)\chi(t),\chi(t)\rangle\bigr)_+
   \,\dd t.
 \label{eq:v26-negative-secant-work}
\end{align}
For
\[
 \mathfrak C_{26}
 :=\{\mathrm{hv},\mathrm{hw},\mathrm{oct},\mathrm{bs},A,S,K\},
\]
define the seven signed carrier works
\begin{equation}
 \boxed{
 \mathcal W^\alpha_d
 :=\int_I\langle G^\alpha(t),\chi(t)\rangle\,\dd t,
 \qquad \alpha\in\mathfrak C_{26}.}
 \label{eq:v26-seven-signed-works}
\end{equation}

\begin{theorem}[Exact seven-carrier work ledger]
\label{thm:v26-signed-work-ledger}
One has the exact identity
\begin{equation}
 \boxed{
 \mathcal P_d[\chi]+\mathcal N_d[\chi]
 =\sum_{\alpha\in\mathfrak C_{26}}\mathcal W_d^\alpha.}
 \label{eq:v26-exact-work-ledger}
\end{equation}
The seven works have the equivalent common-history representations
\begin{align}
 \mathcal W_d^{\rm hv}
 &=\int_I\langle\phi,\mathscr M_v\chi\rangle\,\dd t,
 \label{eq:v26-work-hv}\\
 \mathcal W_d^{\rm hw}
 &=\int_I\left\langle
   \phi,\mathscr M_{\Pi_{(K,3K]}w}\chi
   \right\rangle\,\dd t,
 \label{eq:v26-work-hw}\\
 \mathcal W_d^{\rm oct}
 &=-\int_I\langle\Pi_2\Psi_\phi,\mathscr L_v\chi\rangle\,\dd t,
 \label{eq:v26-work-oct}\\
 \mathcal W_d^{\rm bs}
 &=-\int_I\langle\Pi_{>2}\Psi_\phi,\mathscr L_w\chi\rangle\,\dd t,
 \label{eq:v26-work-bs}\\
 \mathcal W_d^{A}
 &=-\nu\int_I\langle A\eta,\chi\rangle\,\dd t,
 \label{eq:v26-work-heat}\\
 \mathcal W_d^{S}
 &=-\int_I\langle\mathscr S_1\eta,\chi\rangle\,\dd t,
 \label{eq:v26-work-symmetric-leak}\\
 \mathcal W_d^{K}
 &=\int_I\langle\mathscr K_1\eta,\chi\rangle\,\dd t.
 \label{eq:v26-work-skew-leak}
\end{align}
Every term is assembled with its sign before a positive action is formed.
\end{theorem}

\begin{proof}
Take the real \(L^2\) inner product of
\eqref{eq:v26-envelope-adjoint-equation} with \(\chi\) and integrate.  The
time derivative gives \(\frac12\|\chi(a)\|^2\), because \(\chi(b)=0\).
Since \(P_d\chi=\chi\),
\[
 \langle P_dAP_d\chi,\chi\rangle
 =\langle A\chi,\chi\rangle
 =\|A^{1/2}\chi\|^2.
\]
The quadratic form of \(P_d\mathscr K_1P_d\) is zero, while the symmetric
part gives \(\mathcal N_d\).  This proves
\eqref{eq:v26-exact-work-ledger}.  The four first identities follow by
adjointness from the definitions in V25.  The final three are
\eqref{eq:v26-heat-envelope-current}--\eqref{eq:v26-skew-envelope-current}
and the facts \(P_d\chi=\chi\), \(\mathscr S_1^*=\mathscr S_1\), and
\(\mathscr K_1^*=-\mathscr K_1\).
\end{proof}

\begin{theorem}[Actual-path antidissipation or one signed carrier]
\label{thm:v26-actual-path-alternative}
If \(\mathcal P_d[\chi]>0\), then at least one of the following holds:
\begin{align}
 \boxed{
 \mathcal N_d^-[\chi]
 \ge\frac12\mathcal P_d[\chi],}
 \label{eq:v26-antidissipation-branch}
\end{align}
or there is \(\alpha\in\mathfrak C_{26}\) such that
\begin{equation}
 \boxed{
 |\mathcal W_d^\alpha|
 \ge\frac1{14}\mathcal P_d[\chi].}
 \label{eq:v26-one-signed-carrier}
\end{equation}
No minimum edge on a target-null shell direction and no norm of a work-null
force is used.
\end{theorem}

\begin{proof}
Since
\(
 -\mathcal N_d[\chi]\le\mathcal N_d^-[\chi]
\), the exact ledger gives
\[
 \mathcal P_d[\chi]
 =\sum_\alpha\mathcal W_d^\alpha-\mathcal N_d[\chi]
 \le
 \left|\sum_\alpha\mathcal W_d^\alpha\right|
 +\mathcal N_d^-[\chi].
\]
If \eqref{eq:v26-antidissipation-branch} fails, then
\(
 |\sum_\alpha\mathcal W_d^\alpha|
 >\mathcal P_d[\chi]/2
\).
There are seven terms, so the triangle inequality gives
\eqref{eq:v26-one-signed-carrier}.
\end{proof}

\begin{corollary}[Positive polarization of every signed carrier]
\label{cor:v26-work-polarization}
For every fixed \(\lambda>0\),
\begin{equation}
 \boxed{
 4\lambda\mathcal W_d^\alpha
 =\|G^\alpha+\lambda\chi\|_{L^2(I)}^2
  -\|G^\alpha-\lambda\chi\|_{L^2(I)}^2,
 \qquad \alpha\in\mathfrak C_{26}.}
 \label{eq:v26-work-polarization}
\end{equation}
In particular, \(\lambda=\nu K^2\) gives two positive reads at the physical
daughter-shell rate.  A full force norm is an optional stronger read, not a
prerequisite for these mixed scalars.
\end{corollary}

\begin{proof}
Expand the two squared norms in the real spacetime Hilbert space and subtract.
\end{proof}

\begin{firewall}[The seven works are one ledger, not seven source births]
\label{fire:v26-no-seven-births}
The four V25 terms are response mechanisms of one exterior adjoint.  The
three envelope terms are off-diagonal blocks created by compressing that same
response to the support-corrected source envelope.  They are not seven independent
nonlinear occurrences and are not assigned seven positive prices.  The
complete signed identity \eqref{eq:v26-exact-work-ledger} is assembled before
Theorem~\ref{thm:v26-actual-path-alternative} selects a surviving carrier.
\end{firewall}

% =============================================================================
\section{The symmetric secant sees only a scale-local Lipschitz packet}
\label{sec:v26-secant-locality}
% =============================================================================

We next evaluate branch \eqref{eq:v26-antidissipation-branch}.  Recall
\(\Lambda=A^{1/2}\).  All vector fields in this section are smooth,
mean zero, and divergence free.

\begin{lemma}[Periodic Lipschitz commutator at half energy]
\label{lem:v26-Lipschitz-commutator}
There is a constant \(C_{\Torus}\) such that, for every smooth
divergence-free vector field \(a\) and every smooth mean-zero vector field
\(z\),
\begin{equation}
 \boxed{
 \left|
  \langle[a\cdot\nabla,\Lambda]z,z\rangle
 \right|
 \le
 C_{\Torus}\|\nabla a\|_\infty
 \|\Lambda^{1/2}z\|_2^2.}
 \label{eq:v26-Lipschitz-commutator}
\end{equation}
\end{lemma}

\begin{proof}
The periodic half-Laplacian has the positive difference-kernel representation
\[
 \langle\Lambda z,z\rangle
 =\frac12\int_{\Torus^3}\int_{\Torus^3}
   |z(x)-z(y)|^2\mathcal K_\Lambda(x-y)\,\dd x\,\dd y,
\]
where \(\mathcal K_\Lambda\) is the periodization of a positive multiple of
\(|x|^{-4}\).  On a fixed fundamental cell it satisfies
\begin{equation}
 d_{\Torus}(0,h)|\nabla\mathcal K_\Lambda(h)|
 \le C\mathcal K_\Lambda(h).
 \label{eq:v26-periodic-kernel-derivative}
\end{equation}
For completeness, this follows termwise from
\(
 \mathcal K_\Lambda(h)=c\sum_{m\in\mathbb Z^3}|h+2\pi m|^{-4}
\): the singular term has the bound with equality up to a constant, and the
remaining uniformly convergent derivative series is controlled on the
compact fundamental cell by the positive kernel.

Put \(T=a\cdot\nabla\).  Since \(\operatorname{div}a=0\), \(T^*=-T\), and
therefore
\[
 \langle[T,\Lambda]z,z\rangle
 =-2\langle Tz,\Lambda z\rangle.
\]
Using the difference-kernel formula and integrating by parts in both
variables gives the exact identity
\[
 \langle[T,\Lambda]z,z\rangle
 =\frac12\int\!\!\int
 |z(x)-z(y)|^2
 (a(x)-a(y))\cdot\nabla\mathcal K_\Lambda(x-y)
 \,\dd x\,\dd y.
\]
The Lipschitz estimate
\(
 |a(x)-a(y)|\le\|\nabla a\|_\infty d_{\Torus}(x,y)
\), followed by \eqref{eq:v26-periodic-kernel-derivative}, bounds the last
integral by a fixed multiple of
\(
 \|\nabla a\|_\infty\langle\Lambda z,z\rangle
\), which is \eqref{eq:v26-Lipschitz-commutator}.
\end{proof}

Let
\begin{equation}
 v:=P_Ku,
 \qquad
 w_4:=\one_{(K,4K]}(\Lambda)u.
 \label{eq:v26-local-parent-fields}
\end{equation}

\begin{theorem}[Exact secant formulas and four-octave locality]
\label{thm:v26-secant-formula}
Let \(\chi=\Pi_1\chi\), and put
\begin{equation}
 z:=\Lambda^{-1/2}\chi.
 \label{eq:v26-z-from-chi}
\end{equation}
Then
\begin{align}
 \langle\mathscr L_v\chi,\chi\rangle
 &=\frac12
   \langle[v\cdot\nabla,\Lambda]z,z\rangle
   +\int_{\Torus^3}
     ((\Lambda z)\cdot\nabla)v\cdot z\,\dd x,
 \label{eq:v26-head-secant-formula}\\
 \langle\mathscr L_w\chi,\chi\rangle
 &=\frac14
   \langle[w_4\cdot\nabla,\Lambda]z,z\rangle
   +\frac12\int_{\Torus^3}
     ((\Lambda z)\cdot\nabla)w_4\cdot z\,\dd x.
 \label{eq:v26-exterior-secant-formula}
\end{align}
In particular, the quadratic exterior-secant form sees no parent frequency
above \(4K\).  There is a universal periodic constant \(C_0\) such that
\begin{equation}
 \boxed{
 \left|
  \langle\mathscr S_1(t)\chi,\chi\rangle
 \right|
 \le
 C_0\mathscr X_K(t)\|\chi\|_2^2,}
 \label{eq:v26-secant-Lipschitz-bound}
\end{equation}
where
\begin{equation}
 \mathscr X_K(t)
 :=\|\nabla v(t)\|_\infty
  +\|\nabla w_4(t)\|_\infty.
 \label{eq:v26-local-Lipschitz-packet}
\end{equation}
\end{theorem}

\begin{proof}
Since \(\chi=\Lambda^{1/2}z\), one has
\(A^{1/4}\chi=\Lambda z\).  The Leray projection and the exterior projection
may be removed when pairing with the divergence-free field \(z\).  Therefore
\[
 \langle\mathscr L_v\chi,\chi\rangle
 =\langle2\mathcal B_{\rm s}(v,\Lambda z),z\rangle
 =\langle(v\cdot\nabla)\Lambda z,z\rangle
  +\langle((\Lambda z)\cdot\nabla)v,z\rangle.
\]
Because \(v\cdot\nabla\) is skew-adjoint,
\[
 2\langle(v\cdot\nabla)\Lambda z,z\rangle
 =\langle[v\cdot\nabla,\Lambda]z,z\rangle,
\]
which proves \eqref{eq:v26-head-secant-formula}.  The definition of
\(\mathscr L_w\) contains \(\mathcal B_{\rm s}\), rather than
\(2\mathcal B_{\rm s}\), and gives the factors in
\eqref{eq:v26-exterior-secant-formula}.

To prove the support statement, expand the latter pairing in Fourier modes.
The two occurrences of \(z\) have wavevectors in \((K,2K]\).  If a parent
wavevector \(p\) of \(w\) contributes, it is the difference of those two
wavevectors and hence satisfies \(|p|\le4K\).  Since \(w=Q_Ku\), only
\(K<|p|\le4K\) remains, so \(w\) may be replaced by \(w_4\) exactly.

Apply Lemma~\ref{lem:v26-Lipschitz-commutator} to the commutator terms.  For
the remaining terms,
\[
 \left|
  \int ((\Lambda z)\cdot\nabla)a\cdot z
 \right|
 \le\|\nabla a\|_\infty\|\Lambda z\|_2\|z\|_2.
\]
On \(K<\Lambda\le2K\),
\[
 \|\Lambda z\|_2^2
 \le2K\|\Lambda^{1/2}z\|_2^2,
 \qquad
 \|z\|_2^2
 \le K^{-1}\|\Lambda^{1/2}z\|_2^2.
\]
Hence
\(
 \|\Lambda z\|_2\|z\|_2
 \le\sqrt2\|\Lambda^{1/2}z\|_2^2
 =\sqrt2\|\chi\|_2^2
\).
Combine the estimates for \(v\) and \(w_4\).  Since the quadratic form of
\(\mathscr S_1\) equals the real quadratic form of
\(\mathscr L_v+\mathscr L_w\), this proves
\eqref{eq:v26-secant-Lipschitz-bound}.
\end{proof}

\begin{corollary}[Actual antidissipation forces a scale-local gradient pulse]
\label{cor:v26-secant-pulse}
Assume \(K\ge1\).  If branch \eqref{eq:v26-antidissipation-branch} holds, then
\begin{equation}
 \boxed{
 \frac{
  \displaystyle\int_I
   \mathscr X_K(t)\|\chi(t)\|_2^2\,\dd t
 }{
  \displaystyle\int_I\|\chi(t)\|_2^2\,\dd t
 }
 \ge\frac{\nu K^2}{2C_0}.}
 \label{eq:v26-weighted-Lipschitz-lower}
\end{equation}
Consequently, there is a time \(t_*\in I\) such that
\begin{equation}
 \boxed{
 \mathscr X_K(t_*)\ge\frac{\nu K^2}{2C_0}.}
 \label{eq:v26-Lipschitz-pulse}
\end{equation}
Moreover, for another universal constant \(c_0>0\),
\begin{equation}
 \boxed{
 \|\nabla P_{\le4K}u(t_*)\|_2^2
 \ge c_0\nu^2K.}
 \label{eq:v26-enstrophy-pulse}
\end{equation}
At the canonical energy cutoff \(K=\mathcal R^2/E_*\),
\begin{equation}
 \boxed{
 \|\nabla P_{\le4K}u(t_*)\|_2^2
 \ge c_0\frac{\nu^2\mathcal R^2}{E_*}.}
 \label{eq:v26-canonical-enstrophy-pulse}
\end{equation}
\end{corollary}

\begin{proof}
The positive stock dominates the viscous action, and the octave support gives
\[
 \mathcal P_d[\chi]
 \ge\nu\int_I\|A^{1/2}\chi\|_2^2\,\dd t
 \ge\nu K^2\int_I\|\chi\|_2^2\,\dd t.
\]
Branch \eqref{eq:v26-antidissipation-branch} and
\eqref{eq:v26-secant-Lipschitz-bound} therefore imply
\[
 C_0\int_I\mathscr X_K\|\chi\|_2^2
 \ge\mathcal N_d^-[\chi]
 \ge\frac{\nu K^2}{2}
      \int_I\|\chi\|_2^2,
\]
which proves \eqref{eq:v26-weighted-Lipschitz-lower}.  The denominator is
positive on this branch, so the essential supremum gives
\eqref{eq:v26-Lipschitz-pulse}.

For a trigonometric polynomial supported in \(|k|\le4K\), Fourier
Cauchy--Schwarz gives the periodic Bernstein estimate
\[
 \|\nabla f\|_\infty
 \le C K^{3/2}\|\nabla f\|_2.
\]
Apply it separately to \(v\) and \(w_4\), then use orthogonality of their
Fourier supports to obtain
\[
 \mathscr X_K(t)
 \le C K^{3/2}\|\nabla P_{\le4K}u(t)\|_2.
\]
Insert \eqref{eq:v26-Lipschitz-pulse} and square.  This gives
\eqref{eq:v26-enstrophy-pulse}; substituting the canonical value of \(K\)
gives \eqref{eq:v26-canonical-enstrophy-pulse}.
\end{proof}

\begin{assessment}[No remote-frequency secant branch]
\label{ass:v26-no-remote-secant}
The negative quadratic secant work is generated only by the head and the
exterior state through frequency \(4K\).  The source-envelope symmetric and
skew mixing works also pair two daughter-octave adjoint vectors, so their
exterior parent in \(\mathscr L_w\) is restricted to the same range.  Remote
frequency can enter the selected ledger only through
\(\mathcal W_d^{\rm bs}\), the V25 high--high backscatter current.  It is not
created by a soft direction of the full remote adjoint.
\end{assessment}

% =============================================================================
\section{The inherited small-daughter branch forces one actual scalar carrier}
\label{sec:v26-small-daughter}
% =============================================================================

Retain the V25 physical target-action window
\begin{equation}
 0<c_\phi\le\|\phi_n\|_{L^2(I_n)}^2\le C_\phi<\infty,
 \label{eq:v26-target-window}
\end{equation}
and assume the support-incidence row has passed.  Put
\begin{equation}
 \widetilde{\mathfrak B}_n^{\rm dau}
 :=\frac{\|d_n^{\rm C}\|_{L^2(I_n)}^2}
          {\nu\mathcal R_n^2}.
 \label{eq:v26-daughter-action}
\end{equation}
Let \(P_{d,n}\), \(\chi_n\), and \(\mathcal P_{d,n}\) be the objects above
for the \(n\)-th card.

\begin{theorem}[Fixed scalar plus vanishing daughter action gives divergent positive stock]
\label{thm:v26-small-daughter-stock}
Suppose that, for some \(\delta>0\),
\begin{equation}
 \frac{|\mathfrak h_n^{\rm C}|}{\mathcal R_n^2}\ge\delta,
 \qquad
 \widetilde{\mathfrak B}_n^{\rm dau}\longrightarrow0.
 \label{eq:v26-small-daughter-fixed-scalar}
\end{equation}
Then
\begin{equation}
 \boxed{
 \mathfrak A_{d,n}
 :=\frac{\nu\|\chi_n\|_{L^2(I_n)}^2}
          {\mathcal R_n^2\|\phi_n\|_{L^2(I_n)}^2}
 \ge
 \frac{\delta^2}
      {C_\phi\widetilde{\mathfrak B}_n^{\rm dau}}
 \longrightarrow\infty.}
 \label{eq:v26-envelope-action-diverges}
\end{equation}
If \(K_n>0\) is the canonical-head radius, define
\begin{equation}
 \mathfrak P_{d,n}
 :=\frac{\mathcal P_{d,n}[\chi_n]}
          {K_n^2\mathcal R_n^2\|\phi_n\|_{L^2(I_n)}^2}.
 \label{eq:v26-normalized-positive-stock}
\end{equation}
Then
\begin{equation}
 \boxed{
 \mathfrak P_{d,n}\ge\mathfrak A_{d,n}\longrightarrow\infty.}
 \label{eq:v26-positive-stock-diverges}
\end{equation}
\end{theorem}

\begin{proof}
By \eqref{eq:v26-envelope-scalar} and Cauchy--Schwarz,
\[
 \delta^2\mathcal R_n^4
 \le|\mathfrak h_n^{\rm C}|^2
 \le\|d_n^{\rm C}\|_{L^2}^2\|\chi_n\|_{L^2}^2.
\]
Use
\(
 \|d_n^{\rm C}\|_{L^2}^2
 =\nu\mathcal R_n^2\widetilde{\mathfrak B}_n^{\rm dau}
\), multiply by
\(
 \nu/(\mathcal R_n^2\|\phi_n\|_{L^2}^2)
\), and apply the upper target window to prove
\eqref{eq:v26-envelope-action-diverges}.

Since \(\chi_n\) lies in \(K_n<\Lambda\le2K_n\),
\[
 \mathcal P_{d,n}[\chi_n]
 \ge\nu K_n^2\|\chi_n\|_{L^2}^2.
\]
Divide by the denominator in
\eqref{eq:v26-normalized-positive-stock}; the result is
\(
 \mathfrak P_{d,n}\ge\mathfrak A_{d,n}
\).
\end{proof}

For \(\alpha\in\mathfrak C_{26}\), put
\begin{align}
 \mathfrak W_{n}^{\alpha}
 &:=\frac{|\mathcal W_{d,n}^{\alpha}|}
 {K_n^2\mathcal R_n^2\|\phi_n\|_{L^2(I_n)}^2},
 \label{eq:v26-normalized-signed-work}\\
 \mathfrak N_n^-
 &:=\frac{\mathcal N_{d,n}^-[\chi_n]}
 {K_n^2\mathcal R_n^2\|\phi_n\|_{L^2(I_n)}^2}.
 \label{eq:v26-normalized-antidissipation}
\end{align}

\begin{corollary}[One actual signed carrier must diverge]
\label{cor:v26-one-carrier-diverges}
Under \eqref{eq:v26-small-daughter-fixed-scalar}, after passage to a
subsequence at least one of the following holds:
\begin{equation}
 \boxed{\mathfrak N_n^-\longrightarrow\infty,}
 \label{eq:v26-antidissipation-diverges}
\end{equation}
or, for one fixed
\(\alpha\in\mathfrak C_{26}\),
\begin{equation}
 \boxed{\mathfrak W_n^\alpha\longrightarrow\infty.}
 \label{eq:v26-work-diverges}
\end{equation}
On the first branch the scale-local pulse
\eqref{eq:v26-Lipschitz-pulse}--\eqref{eq:v26-enstrophy-pulse} occurs.  On the
second branch the divergence is a signed mixed read of one literal physical
carrier, not merely growth of its full \(L^2\)-norm.
\end{corollary}

\begin{proof}
Apply Theorem~\ref{thm:v26-actual-path-alternative} to each \(n\) and divide
by the positive normalizing denominator.  By
\eqref{eq:v26-positive-stock-diverges}, the right-hand lower fraction tends
to infinity.  There are finitely many branches, so a priority subsequence
fixes one of them.
\end{proof}

% =============================================================================
\section{Corrected terminal theorem after V26}
\label{sec:v26-terminal}
% =============================================================================

\begin{theorem}[Occurring-envelope and signed-work terminal alternative]
\label{thm:v26-terminal-alternative}
Apply the inherited material-route, cell-capture, complete signed
boundary-current, target-normalization, and source-support incidence rows in
their existing order.  Freeze a canonical low-pass head and its physical
return writer before reading the return.  After passage to a priority
subsequence, one of the following occurs.
\begin{enumerate}[label=\textup{(G26.\arabic*)},leftmargin=6.1em]
\item \emph{Target-normalization or source-support survivor.}  One of the
      first two rows of the V25 terminal gate survives.
\item \emph{Corrected daughter first-birth action.}  The quantity
      \(\widetilde{\mathfrak B}_n^{\rm dau}\) has a positive lower bound.
\item \emph{Source-envelope heat mixing.}  The normalized signed work
      \(\mathfrak W_n^A\) diverges.  The source-cyclic occurring envelope is
      not reducing for the Stokes operator on the selected scalar.
\item \emph{Source-envelope symmetric mixing.}  The normalized signed work
      \(\mathfrak W_n^S\) diverges.
\item \emph{Source-envelope skew circulation.}  The normalized signed work
      \(\mathfrak W_n^K\) diverges.  This is a phase/circulation transport of
      the existing adjoint, not a new source birth.
\item \emph{Actual secant antidissipation.}  The negative symmetric-secant
      work diverges.  Then the scale-local gradient and enstrophy packets
      \eqref{eq:v26-Lipschitz-pulse}--\eqref{eq:v26-enstrophy-pulse} occur on
      the same history.
\item \emph{Direct head--daughter work.}  The signed carrier
      \(\mathfrak W_n^{\rm hv}\) diverges.
\item \emph{Local exterior return work.}  The signed carrier
      \(\mathfrak W_n^{\rm hw}\) diverges and uses only exterior frequencies
      through \(3K_n\).
\item \emph{Adjacent-octave work.}  The signed carrier
      \(\mathfrak W_n^{\rm oct}\) diverges.
\item \emph{High--high backscatter work.}  The signed carrier
      \(\mathfrak W_n^{\rm bs}\) diverges.  This is the only branch sent to
      the inherited remote-frequency/high--high audit.
\item \emph{Target-native scalar ancestry pass.}  The first-birth scalar is
      negligible, the V22 old and represented target residuals vanish, and
      the scalar ancestry card terminates without reconstructing a robust
      all-direction shell response.
\item \emph{Fused all-shell selected-source escape.}  No predeclared finite
      head captures the selected support-corrected residual and its scalar.
\end{enumerate}
In particular, on a finite head the combination
\begin{equation}
 \widetilde{\mathfrak B}_n^{\rm dau}\to0,
 \qquad
 |\mathfrak h_n^{\rm C}|/\mathcal R_n^2\ge\delta>0
 \label{eq:v26-forbidden-silent-branch}
\end{equation}
cannot pass silently: one of items~\textup{(G26.3)--(G26.10)} must occur.
The minimum broad shell edge and the seven full force actions remain optional
stronger certificates for a reusable bank; they are not extra prerequisites
for item~\textup{(G26.11)}.
\end{theorem}

\begin{proof}
The first two and the final all-shell alternatives are inherited from V25.
On the finite-head branch, construct the minimum fixed support-corrected source
envelope by Theorem~\ref{thm:v26-source-envelope-minimal}.  If the
first-birth scalar is negligible, apply the V22 old and represented residual
gate; its passing branch is item~\textup{(G26.11)}.  Otherwise pass to a
subsequence on which it has a fixed record-normalized lower bound.  If the
corrected daughter action does not vanish, item~\textup{(G26.2)} occurs.  On
the vanishing branch, Corollary~\ref{cor:v26-one-carrier-diverges} selects
one of the negative secant or seven signed-work alternatives.  The meanings
and support of those alternatives are
Theorems~\ref{thm:v26-envelope-adjoint-equation},
\ref{thm:v26-signed-work-ledger}, and
\ref{thm:v26-secant-formula}.  The pulse conclusion is
Corollary~\ref{cor:v26-secant-pulse}.  Finiteness of the list gives one
priority subsequence.
\end{proof}

\begin{firewall}[What V26 does and does not close]
\label{fire:v26-no-global}
Version~V26 removes target-null shell softness and target-null forcing action
from the one-scalar terminal gate.  It does not show that the surviving
signed work, gradient pulse, adjacent-octave current, or backscatter current
is impossible.  It does not supply a terminal Dini improvement, a critical
unweighted action upper bound, terminal trace, epsilon regularity, backward
uniqueness, or a Liouville theorem.  No three-dimensional global regularity
conclusion is asserted.
\end{firewall}

% =============================================================================
\section{Version V26 integrated conclusion and proof-status ledger}
\label{sec:v26-integrated}
% =============================================================================

\begin{theorem}[Version V26 integrated conclusion]
\label{thm:v26-integrated}
Preserving every theorem and limitation of Versions~V1--V25, Version~V26
proves the following.
\begin{enumerate}[label=\textup{(V26.\arabic*)},leftmargin=5.8em]
\item The essential range of the occurring support-corrected daughter history
      has a unique minimum fixed spatial envelope, selected before the return
      writer and contained in the first exterior octave.
\item Projection of the V25 adjoint onto that envelope gives the exact
      terminal-zero equation \eqref{eq:v26-envelope-adjoint-equation}.  Heat,
      symmetric-secant, and skew-circulation mixing across the envelope are
      retained as signed response currents.
\item The endpoint, viscous, secant, and seven carrier terms form the exact
      signed ledger \eqref{eq:v26-exact-work-ledger}.
\item Either actual negative secant work carries half the positive adjoint
      stock, or one signed carrier carries one fourteenth of that stock.
      Full shell edges and force norms are not used.
\item Broad-edge and force-norm counterexamples prove that those stronger
      objects can fail or diverge entirely in target-null directions.
\item The symmetric secant has the exact fractional-commutator formulas
      \eqref{eq:v26-head-secant-formula}--
      \eqref{eq:v26-exterior-secant-formula} and sees exterior parents only
      through frequency \(4K\).
\item Actual secant antidissipation of viscous size forces the scale-local
      gradient pulse \eqref{eq:v26-Lipschitz-pulse} and enstrophy packet
      \eqref{eq:v26-enstrophy-pulse}.
\item On the inherited fixed-scalar/vanishing-daughter branch, the normalized
      occurring-envelope adjoint stock diverges and therefore one actual
      signed carrier diverges.
\item The corrected terminal alternatives are exactly those of
      Theorem~\ref{thm:v26-terminal-alternative}.
\end{enumerate}
No item excludes every terminal sequence for a fixed datum.
\end{theorem}

\begin{proof}
Items~\textup{(V26.1)--(V26.2)} are
Theorems~\ref{thm:v26-source-envelope-minimal} and
\ref{thm:v26-envelope-adjoint-equation}.  Items~\textup{(V26.3)--(V26.4)}
are Theorems~\ref{thm:v26-signed-work-ledger} and
\ref{thm:v26-actual-path-alternative}.  Item~\textup{(V26.5)} is
Countertheorems~\ref{cth:v26-broad-edge-null} and
\ref{cth:v26-force-norm-null}.  Items~\textup{(V26.6)--(V26.7)} are
Lemma~\ref{lem:v26-Lipschitz-commutator},
Theorem~\ref{thm:v26-secant-formula}, and
Corollary~\ref{cor:v26-secant-pulse}.  Item~\textup{(V26.8)} is
Theorem~\ref{thm:v26-small-daughter-stock} and
Corollary~\ref{cor:v26-one-carrier-diverges}.  The last item is
Theorem~\ref{thm:v26-terminal-alternative}.
\end{proof}

\begingroup\small
\begin{longtable}{>{\raggedright\arraybackslash}p{0.27\textwidth}
                  >{\raggedright\arraybackslash}p{0.18\textwidth}
                  >{\raggedright\arraybackslash}p{0.47\textwidth}}
\caption{Version V26 proof-status ledger.}\label{tab:v26-ledger}\\
\toprule
\textbf{Row}&\textbf{Status}&\textbf{Exact advance or limitation}\\
\midrule
\endfirsthead
\toprule
\textbf{Row}&\textbf{Status}&\textbf{Exact advance or limitation}\\
\midrule
\endhead
V1--V25 prefix&\PROVED{} preserved
 &Complete V25 preterminal source retained byte-for-byte; no inherited theorem
 rewritten.\\[0.35em]
Minimum fixed support-corrected source envelope&\PROVED{} exact
 &Unique smallest time-independent spatial subspace containing the essential
 source range; fixed before the return writer.\\[0.35em]
Source-envelope adjoint equation&\PROVED{} exact
 &Four V25 forces plus heat, symmetric, and skew envelope currents; no
 remainder and no new source birth.\\[0.35em]
Signed adjoint-work ledger&\PROVED{} exact
 &Endpoint, viscosity, symmetric secant, and seven mixed works are assembled
 before positive polarization.\\[0.35em]
Broad shell edge&\OBSTRUCTION{} as one-scalar prerequisite
 &A negative target-null direction can destroy the minimum edge while the
 support-corrected source envelope remains strictly stable.\\[0.35em]
Full forcing actions&\OBSTRUCTION{} as one-scalar prerequisite
 &They can diverge in work-null directions; the selected target uses the
 signed mixed work.\\[0.35em]
Actual-path alternative&\PROVED{} quantitative
 &Half-stock secant antidissipation or one signed carrier with at least one
 fourteenth of the positive stock.\\[0.35em]
Fractional secant identity&\PROVED{} exact / bounded
 &Periodic Lipschitz commutator plus strain term; no exterior parent above
 \(4K\) enters the quadratic form.\\[0.35em]
Scale-local gradient packet&\PROVED{} necessary branch
 &Secant antidissipation forces \(\mathscr X_K\gtrsim\nu K^2\) and
 \(\|\nabla P_{\le4K}u\|_2^2\gtrsim\nu^2K\) at one time.\\[0.35em]
Small-daughter/fixed-scalar branch&\PROVED{} cannot be silent
 &The target-visible positive stock diverges, hence one actual signed work or
 the secant packet diverges.\\[0.35em]
Exclusion of surviving carriers&\OPEN{} physical PDE
 &Requires evaluation of the selected signed work, a same-cell upper stock,
 or an existing compactness/rigidity theorem.\\[0.35em]
Three-dimensional regularity&\OPEN
 &No surviving terminal packet has been excluded for every datum.\\
\bottomrule
\end{longtable}
\endgroup

% =============================================================================
\section{Exact mandate after Version V26}
\label{sec:v26-next}
% =============================================================================

\begin{programme}[Evaluate one signed carrier; do not rebuild the broad shell]
\label{prog:v26-next}
Preserve Versions~V1--V26 verbatim.  Do not introduce another source class,
adjoint field, response graph, path law, shell count, Moore--Penrose
factorization, Dini modulus, or compactness space.

Proceed only in the following order.
\begin{enumerate}[label=\textup{(V27.\arabic*)},leftmargin=5.9em]
\item Retain the existing material-route, cell-capture, signed boundary-current,
      target-normalization, and support-incidence gates.  Construct the
      minimum fixed support-corrected source envelope before reading the return.
\item Evaluate the three off-envelope works
      \(\mathcal W^A_d,\mathcal W^S_d,\mathcal W^K_d\).  A nonzero normalized
      value is respectively a heat-range, symmetric-range, or phase/circulation
      incidence survivor.  Do not enlarge the source envelope after the read.
\item On the actual secant branch, use the same time and terminal cell as in
      \eqref{eq:v26-Lipschitz-pulse}.  Test whether the scale-local gradient or
      enstrophy packet is captured by the inherited material cell; otherwise
      return the literal spatial-incidence escape.  Do not introduce a new
      fitted ball.
\item Evaluate the signed direct works
      \(\mathcal W^{\rm hv}_d\) and \(\mathcal W^{\rm hw}_d\) from the finite
      head/daughter coefficients.  A full forcing norm is unnecessary unless
      reusable donor variations are explicitly declared.
\item Evaluate \(\mathcal W^{\rm oct}_d\) as one adjacent-octave current after
      represented returns have been fused.  Evaluate
      \(\mathcal W^{\rm bs}_d\) only through the inherited high--high
      backscatter audit.
\item If all seven signed works and the negative secant packet are
      subcritical relative to \(\mathcal P_d\), use
      \eqref{eq:v26-exact-work-ledger} directly: the fixed-scalar,
      vanishing-daughter branch is excluded without a broad shell edge.
\item If the first-birth scalar is negligible and the V22 old and represented
      residuals pass, terminate NS--A and export the target-native scalar
      ancestry card.  Open the full V25 shell edge or force actions only for
      an independently declared reusable-bank target.
\item If no finite head captures the support-corrected source envelope, retain fused
      all-shell selected-source escape.
\end{enumerate}
\end{programme}

\begin{assessment}[Programme position after Version V26]
\label{ass:v26-position}
The one-scalar problem no longer waits for the minimum edge and complete
forcing actions on every direction of the first exterior octave.  Its first
unpopulated finite card is
\[
 \boxed{
 \left(
  \mathcal W^A_d,\mathcal W^S_d,\mathcal W^K_d,
  \mathcal N_d^-,
  \mathcal W_d^{\rm hv},\mathcal W_d^{\rm hw},
  \mathcal W_d^{\rm oct},\mathcal W_d^{\rm bs}
 \right),}
\]
together with the already occurring positive stock \(\mathcal P_d\).  Every
entry is a signed same-history scalar in one exact adjoint ledger.  The
secant entry has already been reduced to a scale-local gradient/enstrophy
packet.  Continuing by another all-direction norm or source metric before
these reads are evaluated would be circular.
\end{assessment}

\paragraph{Version V26 history.}
Preserved the complete V25 source.  Replaced its robust all-direction shell
prerequisites, for the weaker one-scalar target, by the minimum fixed envelope
of the selected support-corrected daughter residual.  Derived the exact envelope-projected
adjoint equation with heat, symmetric-secant, and skew-circulation currents;
assembled the seven signed carrier works and proved the half-stock/fourteenth-
work alternative; exhibited finite counterexamples to target-native use of a
broad minimum edge and full forcing norm; proved a periodic Lipschitz
commutator formula for the exterior secant; localized its parent frequencies
to \(4K\); and derived the necessary scale-local gradient and enstrophy pulse.
No terminal Dini gain, upper critical-action estimate, compact-profile
rigidity, or global regularity conclusion was claimed.

\begin{assessment}[Append-only integrity]
\label{ass:v26-integrity}
The complete Version~V25 source preceding its sole final executable document
terminator is preserved verbatim.  Its preterminal SHA--256 digest is
\begin{center}\scriptsize\ttfamily
0809ff7938892195602c525094e3c01ccc95b8e7900344aa6df995bdfbdadea6
\end{center}
The present material begins at Section~\ref{sec:v26-entry}; the final command
below is again unique.
\end{assessment}

\addtocontents{toc}{\protect\endgroup}
% =============================================================================
% END VERSION V26 MATERIAL -- append later versions before final terminator
% =============================================================================

% APPEND ALL LATER VERSIONS IMMEDIATELY ABOVE THE SOLE FINAL LINE BELOW.


\clearpage
% =============================================================================
% VERSION V27: NEW MATERIAL.  The complete V1--V26 source above is preserved
% verbatim; only its sole final executable \end{document} line was removed.
% =============================================================================
\makeatletter
\addtocontents{toc}{\protect\begingroup\protect\makeatletter}
\addtocontents{toc}{\protect\def\protect\l@section{\protect\@dottedtocline{1}{0em}{4.7em}}}
\addtocontents{toc}{\protect\def\protect\l@subsection{\protect\@dottedtocline{2}{1.5em}{5.3em}}}
\addtocontents{toc}{\protect\makeatother}
\makeatother
\renewcommand{\thesection}{V27.\arabic{section}}
\renewcommand{\thesubsection}{V27.\arabic{section}.\arabic{subsection}}
\renewcommand{\theequation}{V27.\arabic{equation}}
\setcounter{section}{0}
\setcounter{equation}{0}
\renewcommand{\theHsection}{V27.\arabic{section}}
\renewcommand{\theHsubsection}{V27.\arabic{section}.\arabic{subsection}}
\renewcommand{\theHtheorem}{V27.\arabic{section}.\arabic{theorem}}
\renewcommand{\theHequation}{V27.\arabic{equation}}

\begin{center}
{\LARGE\bfseries NS--A: Version V27}\\[0.5em]
{\large The Heat-Reducing Source Carrier, Assembly-First Return Ledger,\\
and the Same-Cell Secant Pulse}\\[0.5em]
{\normalsize 4 September 2026}
\end{center}
\addcontentsline{toc}{part}{Version V27: heat-reducing source carrier and assembly-first return ledger}

\begin{abstract}
Version~V26 constructed the unique smallest fixed linear span of the selected
support-corrected daughter history and projected the terminal-zero exterior
adjoint onto it.  Its exact seven-carrier ledger is correct.  That linear
span need not reduce the fixed Stokes operator, however.  The resulting
heat-range work can therefore be created by compressing a known diagonal
heat propagation to a nonreducing source span.  Treating this compression
current as a terminal nonlinear obstruction would stop one arrow too late.

The present continuation moves upstream without changing the source, its
action, its represented-source short, or its selected scalar.  It dephases
the positive source Gram across the finite Stokes eigenspaces of the first
exterior octave and takes its support.  The result is the unique smallest
heat-reducing fixed carrier containing the occurring daughter history.  The
dephasing preserves the exact source action, may increase only the response-
carrier rank, and is selected before the return writer is read.  On this
carrier the V26 heat-range current vanishes identically.

The remaining adjoint equation has one assembled physical forcing and one
assembled dynamic cross-carrier current.  Its exact signed energy ledger has
only three primary outcomes: one-third negative secant work, one-third
dynamic range transfer, or one-third physical forcing work.  Only after one
aggregate survives is it split into symmetric/skew transfer or head--
daughter/local-exterior/adjacent-octave/backscatter forcing.  Finite
counterexamples show that opening those components before assembly can
manufacture arbitrarily large cancelling apparent survivors.

On the secant branch, the V26 frequency-local estimate remains valid on the
heat-reducing carrier.  The resulting finite-frequency enstrophy pulse is
then tested against the already inherited material cell at the same selected
time.  Either at least half of the pulse lies in that exact moving cell, with
an explicit density lower bound, or at least half lies outside it and gives a
literal spatial-incidence escape.  No fitted replacement ball is introduced.
Under the inherited fixed-scalar/vanishing-daughter hypothesis, the heat-
reduced positive adjoint stock diverges; hence one of the three assembled
primary outcomes diverges.  No terminal Dini gain, critical upper-action
estimate, compact-profile rigidity, or global three-dimensional regularity
conclusion is claimed.
\end{abstract}

% =============================================================================
\section{The Stokes-dephased source Gram is the upstream carrier}
\label{sec:v27-entry}
% =============================================================================

Retain the complete V26 finite-head packet.  Thus
\[
 P_K=\one_{(0,K]}(\Lambda),\qquad
 Q_K=I-P_K,\qquad
 \Pi_1=\one_{(K,2K]}(\Lambda),
\]
the support-corrected daughter history is
\(d^{\rm C}\in L^2(I;\operatorname{Ran}\Pi_1)\), and the terminal-zero
adjoint \(\psi\) satisfies \eqref{eq:v26-inherited-octave-adjoint}.  The V26
source envelope \(\mathcal E_d\) is the minimum fixed linear span of the
essential range of \(d^{\rm C}\).  It need not reduce \(A\).

\begin{firewall}[The V26 heat work remains an exact diagnostic]
\label{fire:v27-v26-heat-valid}
The identity \eqref{eq:v26-exact-work-ledger} and its heat work
\(\mathcal W_d^A\) remain exact on the V26 carrier.  Version~V27 does not
rewrite or negate that theorem.  It proves that this work diagnoses the
failure of the minimum linear source span to reduce the already known Stokes
operator.  Since a canonical source-generated heat-reducing completion
exists and preserves the source action, this diagnostic is resolved before
the remaining nonlinear carrier works are interpreted.
\end{firewall}

On the finite-dimensional real Hilbert space
\(\operatorname{Ran}\Pi_1\), write the distinct spectral values of \(A\) as
\begin{equation}
 \Sigma_K:=\operatorname{spec}\!\left(
 A\big|_{\operatorname{Ran}\Pi_1}\right),
 \qquad
 A\Pi_1=\sum_{\lambda\in\Sigma_K}\lambda E_\lambda,
 \label{eq:v27-shell-spectral-resolution}
\end{equation}
where the \(E_\lambda\) are mutually orthogonal spectral projections.  Every
\(\lambda\in\Sigma_K\) lies in \((K^2,4K^2]\).

\begin{definition}[Source Gram and Stokes dephasing]
\label{def:v27-heat-dephased-Gram}
Use the rank-one convention
\((x\otimes x)q=\langle q,x\rangle x\), and put
\begin{align}
 \mathbb Q_d
 &:=\int_I d^{\rm C}(t)\otimes d^{\rm C}(t)\,\dd t,
 \label{eq:v27-source-Gram}\\
 \boxed{
 \mathbb Q_d^{\rm heat}
 :=\sum_{\lambda\in\Sigma_K}
 E_\lambda\mathbb Q_dE_\lambda.}
 \label{eq:v27-heat-dephased-Gram}
\end{align}
Define the heat-reducing source carrier and its support projection by
\begin{equation}
 \boxed{
 \mathcal E_d^{\rm heat}:=\operatorname{Ran}\mathbb Q_d^{\rm heat},
 \qquad
 P_d^{\rm heat}
 :=\mathbb Q_d^{\rm heat}
   (\mathbb Q_d^{\rm heat})^\dagger.}
 \label{eq:v27-heat-carrier}
\end{equation}
The zero-source case is read with
\(\mathcal E_d^{\rm heat}=\{0\}\) and \(P_d^{\rm heat}=0\).
\end{definition}

\begin{theorem}[Canonical minimum heat-reducing completion]
\label{thm:v27-minimum-heat-carrier}
The carrier in \eqref{eq:v27-heat-carrier} has the following properties.
\begin{enumerate}[label=\textup{(H\arabic*)},leftmargin=3.5em]
\item It has the intrinsic block-range description
      \begin{equation}
       \boxed{
       \mathcal E_d^{\rm heat}
       =\bigoplus_{\lambda\in\Sigma_K}
        \operatorname*{ess\,span}_{t\in I}
        \{E_\lambda d^{\rm C}(t)\}.}
       \label{eq:v27-block-range}
      \end{equation}
\item It contains the occurring source and reduces the Stokes operator:
      \begin{equation}
       \boxed{
       P_d^{\rm heat}d^{\rm C}(t)=d^{\rm C}(t)
       \text{ a.e.},
       \qquad
       [P_d^{\rm heat},A]=0.}
       \label{eq:v27-source-contained-heat-reducing}
      \end{equation}
\item It is the unique smallest fixed reducing carrier with this property.
      If \(P\) is an orthogonal projection satisfying
      \([P,A]=0\) and \(Pd^{\rm C}=d^{\rm C}\) almost everywhere, then
      \begin{equation}
       \boxed{
       P_d^{\rm heat}P=P_d^{\rm heat},
       \qquad
       \operatorname{Ran}P_d^{\rm heat}
       \subseteq\operatorname{Ran}P.}
       \label{eq:v27-heat-minimality}
      \end{equation}
\item It lies between the V26 envelope and the complete daughter octave:
      \begin{equation}
       \boxed{
       P_dP_d^{\rm heat}=P_d,
       \qquad
       P_d^{\rm heat}\Pi_1=P_d^{\rm heat}.}
       \label{eq:v27-carrier-nesting}
      \end{equation}
\item Stokes dephasing preserves the complete daughter source action:
      \begin{equation}
       \boxed{
       \operatorname{tr}\mathbb Q_d^{\rm heat}
       =\operatorname{tr}\mathbb Q_d
       =\int_I\|d^{\rm C}(t)\|_2^2\,\dd t.}
       \label{eq:v27-action-preserved}
      \end{equation}
\item The same carrier is the finite heat-cyclic closure of the source:
      \begin{equation}
       \boxed{
       \mathcal E_d^{\rm heat}
       =\operatorname{span}
        \{e^{-sA}x:s\ge0,\ x\in\mathcal E_d\}.}
       \label{eq:v27-heat-cyclic-characterization}
      \end{equation}
\end{enumerate}
Every object in the theorem depends only on the occurring daughter history
and the fixed Stokes operator, not on the return writer or the value of the
selected scalar.
\end{theorem}

\begin{proof}
For every \(q\in\operatorname{Ran}\Pi_1\),
\begin{equation}
 \left\langle q,\mathbb Q_d^{\rm heat}q\right\rangle
 =\sum_{\lambda\in\Sigma_K}
   \int_I
   \left|\left\langle E_\lambda d^{\rm C}(t),q\right\rangle\right|^2
   \,\dd t.
 \label{eq:v27-dephased-quadratic-form}
\end{equation}
The kernel is therefore the orthogonal complement of the direct sum in
\eqref{eq:v27-block-range}; positivity and finite dimensionality give the
range formula.  For almost every \(t\), each spectral component
\(E_\lambda d^{\rm C}(t)\) belongs to the corresponding essential span.
Their finite sum is \(d^{\rm C}(t)\), proving source containment.  Each
summand in \eqref{eq:v27-block-range} lies in one eigenspace of \(A\), where
\(A\) acts as the scalar \(\lambda\); hence the support projection commutes
with \(A\).

Let \(P\) reduce \(A\) and contain the source.  On the finite spectral space
every \(E_\lambda\) is a polynomial in \(A\), so \(P\) commutes with
\(E_\lambda\).  Therefore
\[
 E_\lambda d^{\rm C}(t)
 =E_\lambda Pd^{\rm C}(t)
 =PE_\lambda d^{\rm C}(t)
\]
for almost every \(t\).  The complete block range lies in
\(\operatorname{Ran}P\), proving \eqref{eq:v27-heat-minimality}.  Since
\(\mathcal E_d\) is the essential span of the undecomposed source values,
source containment gives
\(\mathcal E_d\subseteq\mathcal E_d^{\rm heat}\); this and shell support
prove \eqref{eq:v27-carrier-nesting}.

Cyclicity of the finite-dimensional trace gives
\[
 \operatorname{tr}\mathbb Q_d^{\rm heat}
 =\sum_\lambda\operatorname{tr}(\mathbb Q_dE_\lambda)
 =\operatorname{tr}(\mathbb Q_d\Pi_1)
 =\operatorname{tr}\mathbb Q_d.
\]
The trace of \(d\otimes d\) is \(\|d\|^2\), proving
\eqref{eq:v27-action-preserved}.

The heat orbit of \(\mathcal E_d\) lies in every reducing carrier containing
\(\mathcal E_d\), hence in \(\mathcal E_d^{\rm heat}\).  Conversely,
Lagrange interpolation on the finite set \(\Sigma_K\) expresses every
\(E_\lambda\) as a polynomial in \(A\).  A finite collection of functions
\(s\mapsto e^{-s\lambda}\) has an invertible generalized Vandermonde matrix
at suitable distinct values of \(s\), so the same spectral projectors are
linear combinations of finitely many heat multipliers \(e^{-sA}\).  Thus
every \(E_\lambda x\), \(x\in\mathcal E_d\), belongs to the span of heat
translates of \(x\).  Apply \eqref{eq:v27-block-range} to obtain
\eqref{eq:v27-heat-cyclic-characterization}.
\end{proof}

\begin{countertheorem}[The minimum linear source span can manufacture heat work]
\label{cth:v27-linear-envelope-heat-current}
Let \(\mathcal H=\R^2\),
\[
 A=\begin{pmatrix}1&0\\0&4\end{pmatrix},
 \qquad
 f=\frac{e_1+e_2}{\sqrt2},
 \qquad
 g=\frac{e_1-e_2}{\sqrt2},
\]
and take the constant source history \(d^{\rm C}(t)=f\) on a unit interval.
Then the V26 source envelope is \(\R f\), whereas
\begin{equation}
 \boxed{
 \mathbb Q_d^{\rm heat}
 =\frac12e_1\otimes e_1+\frac12e_2\otimes e_2,
 \qquad
 \mathcal E_d^{\rm heat}=\R^2.}
 \label{eq:v27-two-mode-dephasing}
\end{equation}
For the decomposition \(\psi=f+g\), \(\chi=f\), \(\eta=g\), the V26 heat
work is
\begin{equation}
 \boxed{
 -\nu\langle A\eta,\chi\rangle
 =\frac{3\nu}{2}\ne0.}
 \label{eq:v27-compression-heat-work}
\end{equation}
On the heat-reducing carrier the complementary component is zero and the
heat current vanishes.  The selected source scalar and the source action are
unchanged.  This is a finite compression counterexample, not a
Navier--Stokes trajectory.
\end{countertheorem}

\begin{proof}
The ordinary source Gram is \(f\otimes f\).  Dephasing it across the two
eigenspaces of \(A\) gives \eqref{eq:v27-two-mode-dephasing}; its trace is
one, the same as the original rank-one Gram.  Moreover,
\[
 \langle Ag,f\rangle=\frac{1-4}{2}=-\frac32,
\]
which proves \eqref{eq:v27-compression-heat-work}.  Since
\(d^{\rm C}\perp g\), replacing the projection of \(\psi\) onto \(\R f\)
by its projection onto \(\R^2\) does not alter
\(\langle d^{\rm C},\psi\rangle\).
\end{proof}

\begin{assessment}[Rank increase is carrier completion, not source birth]
\label{ass:v27-rank-not-birth}
Stokes dephasing can increase rank, as
\eqref{eq:v27-two-mode-dephasing} shows.  Equation
\eqref{eq:v27-action-preserved} proves that it does not increase the source
action.  The additional directions are the minimum fixed response directions
needed to carry the heat orbit of the source already present.  They are not
orthogonal first births and receive no new source price.
\end{assessment}

% =============================================================================
\section{The heat-reduced exterior-adjoint equation}
\label{sec:v27-heat-reduced-adjoint}
% =============================================================================

Put
\begin{equation}
 P_{\rm h}:=P_d^{\rm heat},
 \qquad
 \chi_{\rm h}:=P_{\rm h}\psi,
 \qquad
 \eta_{\rm h}:=(\Pi_1-P_{\rm h})\psi.
 \label{eq:v27-heat-adjoint-split}
\end{equation}
Then \(\psi=\chi_{\rm h}+\eta_{\rm h}\), and
\begin{equation}
 \boxed{
 \mathfrak h^{\rm C}
 =\int_I\langle d^{\rm C}(t),\chi_{\rm h}(t)\rangle\,\dd t.}
 \label{eq:v27-heat-carrier-scalar}
\end{equation}
Define the daughter-octave secant
\begin{equation}
 \mathscr L_1(t):=\Pi_1\mathscr L_Q^u(t)\Pi_1,
 \qquad
 \mathscr L_1(t)^*=\mathscr S_1(t)-\mathscr K_1(t),
 \label{eq:v27-octave-secant}
\end{equation}
and put
\begin{align}
 F_{\rm h}&:=P_{\rm h}F_1,
 \label{eq:v27-heat-forcing}\\
 G_{\rm h}^{\rm tr}
 &:=-P_{\rm h}\mathscr L_1^*\eta_{\rm h}.
 \label{eq:v27-dynamic-transfer-current}
\end{align}

\begin{theorem}[Exact heat-reduced adjoint equation]
\label{thm:v27-heat-reduced-equation}
The target-visible adjoint on the minimum heat-reducing carrier is the
terminal-zero solution of
\begin{equation}
 \boxed{
 \begin{aligned}
 -\partial_t\chi_{\rm h}
 &+\nu A\chi_{\rm h}
 +P_{\rm h}\mathscr S_1P_{\rm h}\chi_{\rm h}
 -P_{\rm h}\mathscr K_1P_{\rm h}\chi_{\rm h}
 \\
 &=F_{\rm h}+G_{\rm h}^{\rm tr},
 \qquad \chi_{\rm h}(b)=0.
 \end{aligned}}
 \label{eq:v27-heat-reduced-equation}
\end{equation}
There is no heat-range current.  The only off-carrier term is the complete
dynamic secant transfer \(G_{\rm h}^{\rm tr}\).
\end{theorem}

\begin{proof}
Apply \(P_{\rm h}\) to \eqref{eq:v26-inherited-octave-adjoint}.  The
projection is fixed in time.  By
\eqref{eq:v27-source-contained-heat-reducing},
\(P_{\rm h}A=AP_{\rm h}\), so
\[
 P_{\rm h}A\eta_{\rm h}=AP_{\rm h}\eta_{\rm h}=0.
\]
Substitute \(\psi=\chi_{\rm h}+\eta_{\rm h}\), use
\eqref{eq:v27-octave-secant}, and move the single term containing
\(\eta_{\rm h}\) to the right.  This gives
\eqref{eq:v27-heat-reduced-equation}.  The terminal condition follows from
\(\psi(b)=0\).
\end{proof}

The dynamic transfer may be split, only after it is assembled, as
\begin{align}
 G_{\rm h}^{S}&:=-P_{\rm h}\mathscr S_1\eta_{\rm h},
 \label{eq:v27-symmetric-transfer-current}\\
 G_{\rm h}^{K}&:=P_{\rm h}\mathscr K_1\eta_{\rm h},
 \label{eq:v27-skew-transfer-current}
\end{align}
so that
\begin{equation}
 \boxed{G_{\rm h}^{\rm tr}=G_{\rm h}^{S}+G_{\rm h}^{K}.}
 \label{eq:v27-transfer-sum}
\end{equation}

\begin{firewall}[Do not close under the dynamic secant]
\label{fire:v27-no-dynamic-closure}
The heat completion is determined by the fixed operator \(A\) and the
occurring source before the return writer is read.  Completing the carrier
under \(\mathscr L_1(t)\) would be different: it would absorb the actual
state-dependent interaction which the return problem is meant to measure.
For example, a skew matrix can rotate a one-dimensional source carrier into
its orthogonal complement while preserving the total norm.  Declaring the
whole rotation orbit to be source support would erase a genuine circulation
current.  Version~V27 therefore removes only the universal heat compression
and retains the complete dynamic cross-carrier work.
\end{firewall}

% =============================================================================
\section{Assembly-first work ledger}
\label{sec:v27-assembly-first}
% =============================================================================

Define the positive endpoint-plus-viscous stock
\begin{equation}
 \boxed{
 \mathcal P_{\rm h}
 :=\frac12\|\chi_{\rm h}(a)\|_2^2
   +\nu\int_I\|A^{1/2}\chi_{\rm h}(t)\|_2^2\,\dd t,}
 \label{eq:v27-positive-stock}
\end{equation}
the symmetric-secant work and its negative part
\begin{align}
 \mathcal N_{\rm h}
 &:=\int_I
   \langle\mathscr S_1(t)\chi_{\rm h}(t),
          \chi_{\rm h}(t)\rangle\,\dd t,
 \label{eq:v27-secant-work}\\
 \mathcal N_{\rm h}^-
 &:=\int_I
   \bigl(-\langle\mathscr S_1(t)\chi_{\rm h}(t),
                  \chi_{\rm h}(t)\rangle\bigr)_+\,\dd t,
 \label{eq:v27-negative-secant-work}
\end{align}
and the two assembled signed works
\begin{align}
 \boxed{
 \mathcal W_{\rm h}^{\rm for}
 :=\int_I\langle F_1(t),\chi_{\rm h}(t)\rangle\,\dd t,}
 \label{eq:v27-assembled-forcing-work}\\
 \boxed{
 \mathcal W_{\rm h}^{\rm tr}
 :=\int_I\langle G_{\rm h}^{\rm tr}(t),
                    \chi_{\rm h}(t)\rangle\,\dd t
 =-\int_I\langle\eta_{\rm h}(t),
                    \mathscr L_1(t)\chi_{\rm h}(t)\rangle\,\dd t.}
 \label{eq:v27-assembled-transfer-work}
\end{align}
For \(\bullet\in\{\mathrm{hv},\mathrm{hw},\mathrm{oct},\mathrm{bs}\}\),
put
\begin{equation}
 \mathcal W_{\rm h}^{\bullet}
 :=\int_I\langle F^{\bullet},\chi_{\rm h}\rangle\,\dd t,
 \label{eq:v27-four-force-works}
\end{equation}
and define
\begin{align}
 \mathcal W_{\rm h}^{S}
 &:=-\int_I\langle\mathscr S_1\eta_{\rm h},
                         \chi_{\rm h}\rangle\,\dd t,
 \label{eq:v27-symmetric-transfer-work}\\
 \mathcal W_{\rm h}^{K}
 &:=\int_I\langle\mathscr K_1\eta_{\rm h},
                        \chi_{\rm h}\rangle\,\dd t.
 \label{eq:v27-skew-transfer-work}
\end{align}
Then
\begin{equation}
 \mathcal W_{\rm h}^{\rm for}
 =\mathcal W_{\rm h}^{\rm hv}
  +\mathcal W_{\rm h}^{\rm hw}
  +\mathcal W_{\rm h}^{\rm oct}
  +\mathcal W_{\rm h}^{\rm bs},
 \qquad
 \mathcal W_{\rm h}^{\rm tr}
 =\mathcal W_{\rm h}^{S}+\mathcal W_{\rm h}^{K}.
 \label{eq:v27-aggregate-splits}
\end{equation}

\begin{theorem}[Exact three-term assembly-first ledger]
\label{thm:v27-three-term-ledger}
The heat-reduced adjoint satisfies
\begin{equation}
 \boxed{
 \mathcal P_{\rm h}+\mathcal N_{\rm h}
 =\mathcal W_{\rm h}^{\rm for}
  +\mathcal W_{\rm h}^{\rm tr}.}
 \label{eq:v27-three-term-ledger}
\end{equation}
The V26 heat-range work is absent, and no physical forcing component or
dynamic transfer component is opened before its aggregate is formed.
\end{theorem}

\begin{proof}
Take the real \(L^2\) inner product of
\eqref{eq:v27-heat-reduced-equation} with \(\chi_{\rm h}\) and integrate.
The time derivative gives \(\frac12\|\chi_{\rm h}(a)\|^2\), the terminal
value being zero.  The Stokes term gives
\(\nu\int\|A^{1/2}\chi_{\rm h}\|^2\).  The internal skew form vanishes,
while the symmetric form gives \(\mathcal N_{\rm h}\).  The right side is
exactly \(\mathcal W_{\rm h}^{\rm for}+\mathcal W_{\rm h}^{\rm tr}\).
\end{proof}

\begin{theorem}[One-third secant, transfer, or forcing alternative]
\label{thm:v27-one-third-alternative}
If \(\mathcal P_{\rm h}>0\), then at least one of
\begin{align}
 \boxed{\mathcal N_{\rm h}^-\ge\frac13\mathcal P_{\rm h},}
 \label{eq:v27-one-third-secant}\\
 \boxed{|\mathcal W_{\rm h}^{\rm tr}|
        \ge\frac13\mathcal P_{\rm h},}
 \label{eq:v27-one-third-transfer}\\
 \boxed{|\mathcal W_{\rm h}^{\rm for}|
        \ge\frac13\mathcal P_{\rm h}}
 \label{eq:v27-one-third-forcing}
\end{align}
holds.
\end{theorem}

\begin{proof}
Since \(-\mathcal N_{\rm h}\le\mathcal N_{\rm h}^-\), the exact ledger
gives
\[
 \mathcal P_{\rm h}
 \le
 |\mathcal W_{\rm h}^{\rm for}
   +\mathcal W_{\rm h}^{\rm tr}|
 +\mathcal N_{\rm h}^-.
\]
If \eqref{eq:v27-one-third-secant} fails, then the absolute value of the
assembled right-hand work exceeds \(2\mathcal P_{\rm h}/3\).  If both
\eqref{eq:v27-one-third-transfer} and
\eqref{eq:v27-one-third-forcing} failed, the triangle inequality would make
that absolute value strictly smaller than \(2\mathcal P_{\rm h}/3\), a
contradiction.
\end{proof}

\begin{corollary}[Hierarchical opening of a surviving aggregate]
\label{cor:v27-hierarchical-opening}
Define
\begin{equation}
 \mathcal W_{\rm h}^{\rm loc}
 :=\mathcal W_{\rm h}^{\rm hv}+\mathcal W_{\rm h}^{\rm hw},
 \qquad
 \mathcal W_{\rm h}^{\rm tail}
 :=\mathcal W_{\rm h}^{\rm oct}+\mathcal W_{\rm h}^{\rm bs}.
 \label{eq:v27-local-tail-work}
\end{equation}
If \eqref{eq:v27-one-third-transfer} holds, then one of
\begin{equation}
 \boxed{
 |\mathcal W_{\rm h}^{S}|
 \ge\frac16\mathcal P_{\rm h},
 \qquad\text{or}\qquad
 |\mathcal W_{\rm h}^{K}|
 \ge\frac16\mathcal P_{\rm h}}
 \label{eq:v27-transfer-component}
\end{equation}
holds.  If \eqref{eq:v27-one-third-forcing} holds, then first one of
\begin{equation}
 |\mathcal W_{\rm h}^{\rm loc}|
 \ge\frac16\mathcal P_{\rm h},
 \qquad\text{or}\qquad
 |\mathcal W_{\rm h}^{\rm tail}|
 \ge\frac16\mathcal P_{\rm h}
 \label{eq:v27-local-tail-alternative}
\end{equation}
holds, and only then one of the two components of the selected pair has
magnitude at least \(\mathcal P_{\rm h}/12\).
\end{corollary}

\begin{proof}
Use \eqref{eq:v27-aggregate-splits} and the triangle inequality, first on
the two aggregate groups and then, only on the selected group, on its two
components.
\end{proof}

\begin{corollary}[Positive polarization of the assembled reads]
\label{cor:v27-aggregate-polarization}
For every \(\lambda>0\),
\begin{align}
 4\lambda\mathcal W_{\rm h}^{\rm for}
 &=\|F_{\rm h}+\lambda\chi_{\rm h}\|_{L^2(I)}^2
  -\|F_{\rm h}-\lambda\chi_{\rm h}\|_{L^2(I)}^2,
 \label{eq:v27-forcing-polarization}\\
 4\lambda\mathcal W_{\rm h}^{\rm tr}
 &=\|G_{\rm h}^{\rm tr}+\lambda\chi_{\rm h}\|_{L^2(I)}^2
  -\|G_{\rm h}^{\rm tr}-\lambda\chi_{\rm h}\|_{L^2(I)}^2.
 \label{eq:v27-transfer-polarization}
\end{align}
The same formula applies to a component only after the corresponding
aggregate branch has been selected.
\end{corollary}

\begin{proof}
Expand the two squared norms in the real spacetime Hilbert space and subtract.
\end{proof}

\begin{countertheorem}[Component growth before assembly can be pure cancellation]
\label{cth:v27-component-cancellation}
On a unit time interval take \(\chi=e_1\in\R^2\) and define
\[
 G_N^{\rm hv}=Ne_1,
 \qquad
 G_N^{\rm hw}=-Ne_1,
 \qquad
 G_N^{\rm oct}=G_N^{\rm bs}=0.
\]
Then
\begin{equation}
 \boxed{
 \mathcal W_N^{\rm hv}=N,
 \qquad
 \mathcal W_N^{\rm hw}=-N,
 \qquad
 \mathcal W_N^{\rm for}=0.}
 \label{eq:v27-force-cancellation}
\end{equation}
The same construction with \(G_N^S=Ne_1\) and
\(G_N^K=-Ne_1\) gives divergent symmetric and skew component works with
zero dynamic transfer work.  Hence a large component read is not a terminal
branch until the complete signed aggregate containing it has survived.  This
is a finite signed-work model, not a Navier--Stokes trajectory.
\end{countertheorem}

\begin{proof}
Each work is the scalar product with the constant vector \(e_1\), integrated
over an interval of length one.  The identities are immediate.
\end{proof}

% =============================================================================
\section{The secant pulse must use the inherited material cell}
\label{sec:v27-same-cell-pulse}
% =============================================================================

The V26 fractional secant estimate
\eqref{eq:v26-secant-Lipschitz-bound} applies to every vector in
\(\operatorname{Ran}\Pi_1\), hence to \(\chi_{\rm h}\).  Recall
\[
 \mathscr X_K(t)
 =\|\nabla P_Ku(t)\|_\infty
  +\|\nabla P_{(K,4K]}u(t)\|_\infty.
\]

\begin{theorem}[One-third secant work gives a scale-local pulse]
\label{thm:v27-sec-antidissipation-pulse}
Assume \eqref{eq:v27-one-third-secant}.  Then
\begin{equation}
 \boxed{
 \frac{
  \displaystyle\int_I
  \mathscr X_K(t)\|\chi_{\rm h}(t)\|_2^2\,\dd t
 }{
  \displaystyle\int_I\|\chi_{\rm h}(t)\|_2^2\,\dd t
 }
 \ge\frac{\nu K^2}{3C_0}.}
 \label{eq:v27-weighted-gradient-lower}
\end{equation}
Consequently there is \(t_*\in I\) such that
\begin{equation}
 \boxed{
 \mathscr X_K(t_*)\ge\frac{\nu K^2}{3C_0},}
 \label{eq:v27-gradient-pulse}
\end{equation}
and, for a universal \(c_1>0\),
\begin{equation}
 \boxed{
 \|\nabla P_{\le4K}u(t_*)\|_2^2
 \ge c_1\nu^2K.}
 \label{eq:v27-enstrophy-pulse}
\end{equation}
At the canonical energy cutoff \(K=\mathcal R^2/E_*\),
\begin{equation}
 \boxed{
 \|\nabla P_{\le4K}u(t_*)\|_2^2
 \ge c_1\frac{\nu^2\mathcal R^2}{E_*}.}
 \label{eq:v27-canonical-enstrophy-pulse}
\end{equation}
\end{theorem}

\begin{proof}
By octave support,
\[
 \mathcal P_{\rm h}
 \ge\nu\int_I\|A^{1/2}\chi_{\rm h}\|_2^2\,\dd t
 \ge\nu K^2\int_I\|\chi_{\rm h}\|_2^2\,\dd t.
\]
The V26 secant bound and \eqref{eq:v27-one-third-secant} give
\[
 C_0\int_I\mathscr X_K\|\chi_{\rm h}\|_2^2
 \ge\mathcal N_{\rm h}^-
 \ge\frac{\nu K^2}{3}
      \int_I\|\chi_{\rm h}\|_2^2,
\]
which proves \eqref{eq:v27-weighted-gradient-lower} and then
\eqref{eq:v27-gradient-pulse}.  The Bernstein argument in
Corollary~\ref{cor:v26-secant-pulse} applies verbatim with the changed
constant: since
\[
 \mathscr X_K(t)
 \le CK^{3/2}\|\nabla P_{\le4K}u(t)\|_2,
\]
substitution of \eqref{eq:v27-gradient-pulse} and squaring gives
\eqref{eq:v27-enstrophy-pulse}.  The canonical formula follows by replacing
\(K\) with \(\mathcal R^2/E_*\).
\end{proof}

Now retain the already constructed V8 material preimage
\begin{equation}
 \mathcal M_n(t)
 =\Phi_{t,b_n}^{-1}(B(x_n,3\rho_n)),
 \label{eq:v27-inherited-material-cell}
\end{equation}
with \(x_n\to x_*\in\Sigma_T\), \(\rho_n\to0\), and
\(|\mathcal M_n(t)|\le C\rho_n^3\) by
\eqref{eq:v8-material-volume}.  This cell is fixed by the inherited material
route; it is not selected after the pulse is located.

\begin{theorem}[Exact same-cell pulse or spatial-incidence escape]
\label{thm:v27-same-cell-pulse}
On a terminal sequence satisfying the inherited material-cell gate, let
\(t_n^*\) be the time supplied by
Theorem~\ref{thm:v27-sec-antidissipation-pulse} and put
\begin{align}
 \mathcal E_n^{\rm in}
 &:=\int_{\mathcal M_n(t_n^*)}
   |\nabla P_{\le4K_n}u(t_n^*,x)|^2\,\dd x,
 \label{eq:v27-inside-cell-enstrophy}\\
 \mathcal E_n^{\rm out}
 &:=\int_{\Torus^3\setminus\mathcal M_n(t_n^*)}
   |\nabla P_{\le4K_n}u(t_n^*,x)|^2\,\dd x.
 \label{eq:v27-outside-cell-enstrophy}
\end{align}
Then
\begin{equation}
 \boxed{
 \mathcal E_n^{\rm in}+\mathcal E_n^{\rm out}
 =\|\nabla P_{\le4K_n}u(t_n^*)\|_2^2.}
 \label{eq:v27-cell-enstrophy-split}
\end{equation}
Consequently, after passage to a subsequence, exactly one priority branch
occurs.
\begin{enumerate}[label=\textup{(C\arabic*)},leftmargin=3.4em]
\item \emph{Same-cell secant pulse.}
      \begin{equation}
       \boxed{
       \mathcal E_n^{\rm in}
       \ge\frac{c_1}{2}\nu^2K_n.}
       \label{eq:v27-same-cell-lower}
      \end{equation}
      In this case
      \begin{equation}
       \boxed{
       \operatorname*{ess\,sup}_{x\in\mathcal M_n(t_n^*)}
       |\nabla P_{\le4K_n}u(t_n^*,x)|^2
       \ge c_2\frac{\nu^2K_n}{\rho_n^3}.}
       \label{eq:v27-same-cell-density}
      \end{equation}
\item \emph{Spatial-incidence escape.}
      \begin{equation}
       \boxed{
       \mathcal E_n^{\rm out}
       \ge\frac{c_1}{2}\nu^2K_n.}
       \label{eq:v27-outside-cell-lower}
      \end{equation}
      At least half of the scale-local secant packet lies outside the exact
      inherited material cell.
\end{enumerate}
At the canonical cutoff, the right side of
\eqref{eq:v27-same-cell-lower} or
\eqref{eq:v27-outside-cell-lower} is
\(c\nu^2\mathcal R_n^2/E_*\).
\end{theorem}

\begin{proof}
Equation \eqref{eq:v27-cell-enstrophy-split} is the exact partition of the
torus into the inherited material cell and its complement.  By
\eqref{eq:v27-enstrophy-pulse}, one of the two nonnegative terms is at least
half of \(c_1\nu^2K_n\).  Pass to a priority subsequence to fix which one.
On the inside branch, divide by
\(|\mathcal M_n(t_n^*)|\le C\rho_n^3\) to obtain
\eqref{eq:v27-same-cell-density}.  The canonical formula follows from
\(K_n=\mathcal R_n^2/E_*\).
\end{proof}

\begin{assessment}[What same-cell capture does not prove]
\label{ass:v27-cell-qualification}
Theorem~\ref{thm:v27-same-cell-pulse} places the finite-frequency enstrophy
mechanism on, or away from, the already selected material cell.  It does not
identify the pointwise maximizer with the original source atom, does not
supply an epsilon-regularity contradiction, and does not turn the pulse into
an ancient profile.  The inherited diameter estimate
\eqref{eq:v8-material-diameter} only shows that the captured branch is
spatially tied to the terminal singular path.
\end{assessment}

% =============================================================================
\section{The fixed-scalar branch has only three assembled exits}
\label{sec:v27-small-daughter}
% =============================================================================

Retain the target window \eqref{eq:v26-target-window} and the corrected
daughter action
\[
 \widetilde{\mathfrak B}_n^{\rm dau}
 =\frac{\|d_n^{\rm C}\|_{L^2(I_n)}^2}
        {\nu\mathcal R_n^2}.
\]
Let \(P_{{\rm h},n}\), \(\chi_{{\rm h},n}\), and
\(\mathcal P_{{\rm h},n}\) be the heat-reduced objects on the \(n\)-th
card.  Define
\begin{align}
 \mathfrak A_{{\rm h},n}
 &:=\frac{\nu\|\chi_{{\rm h},n}\|_{L^2(I_n)}^2}
          {\mathcal R_n^2\|\phi_n\|_{L^2(I_n)}^2},
 \label{eq:v27-normalized-adjoint-action}\\
 \mathfrak P_{{\rm h},n}
 &:=\frac{\mathcal P_{{\rm h},n}}
 {K_n^2\mathcal R_n^2\|\phi_n\|_{L^2(I_n)}^2}.
 \label{eq:v27-normalized-positive-stock}
\end{align}

\begin{theorem}[Heat completion preserves the divergent target stock]
\label{thm:v27-heat-stock-diverges}
Suppose that, for some \(\delta>0\),
\begin{equation}
 \frac{|\mathfrak h_n^{\rm C}|}{\mathcal R_n^2}\ge\delta,
 \qquad
 \widetilde{\mathfrak B}_n^{\rm dau}\longrightarrow0.
 \label{eq:v27-fixed-scalar-small-daughter}
\end{equation}
Then
\begin{equation}
 \boxed{
 \mathfrak A_{{\rm h},n}
 \ge
 \frac{\delta^2}
 {C_\phi\widetilde{\mathfrak B}_n^{\rm dau}}
 \longrightarrow\infty,}
 \label{eq:v27-heat-action-diverges}
\end{equation}
and
\begin{equation}
 \boxed{
 \mathfrak P_{{\rm h},n}
 \ge\mathfrak A_{{\rm h},n}
 \longrightarrow\infty.}
 \label{eq:v27-heat-stock-diverges}
\end{equation}
\end{theorem}

\begin{proof}
By \eqref{eq:v27-heat-carrier-scalar} and Cauchy--Schwarz,
\[
 \delta^2\mathcal R_n^4
 \le\|d_n^{\rm C}\|_{L^2}^2
     \|\chi_{{\rm h},n}\|_{L^2}^2.
\]
Substitute
\(\|d_n^{\rm C}\|_{L^2}^2
 =\nu\mathcal R_n^2\widetilde{\mathfrak B}_n^{\rm dau}\),
multiply by
\(\nu/(\mathcal R_n^2\|\phi_n\|_{L^2}^2)\), and use the upper target
window.  This proves \eqref{eq:v27-heat-action-diverges}.  Since
\(\chi_{{\rm h},n}\) is supported in \(K_n<\Lambda\le2K_n\),
\[
 \mathcal P_{{\rm h},n}
 \ge\nu K_n^2\|\chi_{{\rm h},n}\|_{L^2}^2.
\]
Divide by the denominator in
\eqref{eq:v27-normalized-positive-stock}.
\end{proof}

Define the three normalized assembled outputs
\begin{align}
 \mathfrak N_{{\rm h},n}^-
 &:=\frac{\mathcal N_{{\rm h},n}^-}
 {K_n^2\mathcal R_n^2\|\phi_n\|_{L^2(I_n)}^2},
 \label{eq:v27-normalized-secant}\\
 \mathfrak W_{{\rm h},n}^{\rm tr}
 &:=\frac{|\mathcal W_{{\rm h},n}^{\rm tr}|}
 {K_n^2\mathcal R_n^2\|\phi_n\|_{L^2(I_n)}^2},
 \label{eq:v27-normalized-transfer}\\
 \mathfrak W_{{\rm h},n}^{\rm for}
 &:=\frac{|\mathcal W_{{\rm h},n}^{\rm for}|}
 {K_n^2\mathcal R_n^2\|\phi_n\|_{L^2(I_n)}^2}.
 \label{eq:v27-normalized-forcing}
\end{align}

\begin{corollary}[One assembled output must diverge]
\label{cor:v27-aggregate-diverges}
Under \eqref{eq:v27-fixed-scalar-small-daughter}, after passage to a
subsequence at least one of
\begin{equation}
 \boxed{
 \mathfrak N_{{\rm h},n}^-\to\infty,
 \qquad
 \mathfrak W_{{\rm h},n}^{\rm tr}\to\infty,
 \qquad
 \mathfrak W_{{\rm h},n}^{\rm for}\to\infty}
 \label{eq:v27-one-aggregate-diverges}
\end{equation}
holds.  On the first branch, Theorem~\ref{thm:v27-same-cell-pulse} applies.
On the second or third branch, components are opened only through
Corollary~\ref{cor:v27-hierarchical-opening}.
\end{corollary}

\begin{proof}
Apply Theorem~\ref{thm:v27-one-third-alternative} to each card and divide by
the positive normalizing denominator.  By
\eqref{eq:v27-heat-stock-diverges}, the selected lower bound diverges.
There are only three primary branches, so a priority subsequence fixes one.
\end{proof}

% =============================================================================
\section{Corrected terminal theorem after V27}
\label{sec:v27-terminal}
% =============================================================================

\begin{theorem}[Heat-reduced assembly-first terminal alternative]
\label{thm:v27-terminal-alternative}
Apply the inherited material-route, cell-capture, complete signed boundary-
current, target-normalization, represented-source, and support-incidence rows
in their existing order.  Freeze a canonical low-pass head and its physical
return writer before reading the return.  Construct the minimum heat-
reducing carrier from the occurring support-corrected daughter history before
using that writer.  After passage to a priority subsequence, one of the
following occurs.
\begin{enumerate}[label=\textup{(G27.\arabic*)},leftmargin=6.3em]
\item \emph{Inherited upstream survivor.}  One of the material, boundary,
      target-normalization, represented-source, or support-incidence gates
      fails.
\item \emph{Corrected daughter first-birth action.}
      \(\widetilde{\mathfrak B}_n^{\rm dau}\) has a positive lower bound.
\item \emph{Heat-reduced secant antidissipation.}  The normalized negative
      secant work survives.  Then the scale-local pulse is either captured in
      the exact inherited material cell through
      \eqref{eq:v27-same-cell-lower}, or it escapes that cell through
      \eqref{eq:v27-outside-cell-lower}.
\item \emph{Assembled dynamic range transfer.}
      \(\mathfrak W_{{\rm h},n}^{\rm tr}\) survives.  Only then is it split
      into symmetric-range and skew/circulation work.
\item \emph{Assembled physical forcing work.}
      \(\mathfrak W_{{\rm h},n}^{\rm for}\) survives.  Only then is it split
      first into local and tail forcing, and subsequently into
      head--daughter, local exterior, adjacent-octave, or high--high
      backscatter work.
\item \emph{Target-native scalar ancestry pass.}  The first-birth scalar is
      negligible and the V22 old and represented target residuals vanish.
      NS--A terminates without reconstructing a broad all-direction shell
      response.
\item \emph{Fused all-shell selected-source escape.}  No predeclared finite
      head captures the support-corrected source and its scalar.
\end{enumerate}
On every finite head, the combination
\begin{equation}
 \widetilde{\mathfrak B}_n^{\rm dau}\to0,
 \qquad
 |\mathfrak h_n^{\rm C}|/\mathcal R_n^2\ge\delta>0
 \label{eq:v27-no-silent-branch}
\end{equation}
cannot pass silently: one of items~\textup{(G27.3)--(G27.5)} occurs.  There
is no independent heat-range branch after the canonical completion.
\end{theorem}

\begin{proof}
The inherited upstream and all-shell branches precede the finite-card
argument.  On a finite head, Theorem~\ref{thm:v27-minimum-heat-carrier}
constructs the canonical carrier without changing the source action or
selected scalar.  If the first-birth scalar is negligible, apply the V22 old
and represented residual gate; its passing branch is item~\textup{(G27.6)}.
Otherwise pass to a subsequence on which the scalar has a fixed normalized
lower bound.  If the daughter action does not vanish, item~\textup{(G27.2)}
occurs.  On the vanishing branch,
Corollary~\ref{cor:v27-aggregate-diverges} selects one of the three assembled
outputs.  The secant output is refined by
Theorem~\ref{thm:v27-same-cell-pulse}; the other two are opened only through
Corollary~\ref{cor:v27-hierarchical-opening}.  This gives the complete list.
\end{proof}

\begin{firewall}[No global theorem in V27]
\label{fire:v27-no-global}
Version~V27 removes only the compression artifact created by using a
nonreducing source span for the known Stokes flow.  It does not show that the
remaining dynamic transfer, physical forcing, same-cell enstrophy pulse,
adjacent-octave current, or backscatter current is impossible.  It does not
supply a terminal Dini improvement, critical action upper bound, terminal
trace, epsilon regularity, backward uniqueness, or a Liouville theorem.  No
three-dimensional global regularity conclusion is asserted.
\end{firewall}

% =============================================================================
\section{Version V27 integrated conclusion and proof-status ledger}
\label{sec:v27-integrated}
% =============================================================================

\begin{theorem}[Version V27 integrated conclusion]
\label{thm:v27-integrated}
Preserving every theorem and limitation of Versions~V1--V26, Version~V27
proves the following.
\begin{enumerate}[label=\textup{(V27.\arabic*)},leftmargin=5.8em]
\item Stokes dephasing of the occurring daughter-source Gram gives the unique
      minimum fixed heat-reducing carrier containing that source.
\item The dephasing preserves total source action and is selected before the
      return writer; any rank increase is response-carrier completion, not a
      new source birth.
\item A two-mode exact counterexample shows that the V26 heat work can be
      nonzero solely because the minimum linear source span fails to reduce
      \(A\).
\item On the heat-reducing carrier, the adjoint equation has no heat current;
      its only off-carrier row is the complete dynamic secant transfer.
\item The exact signed ledger has three primary outputs: negative secant work,
      assembled dynamic transfer, and assembled physical forcing.
\item A one-third alternative and hierarchical component theorem prevent
      cancellation from being interpreted as multiple physical survivors.
\item On the secant branch, the finite-frequency enstrophy pulse is tested on
      the exact inherited material cell, producing same-cell capture or
      literal spatial-incidence escape.
\item On the fixed-scalar/vanishing-daughter branch, the heat-reduced positive
      stock diverges, so one of the three assembled outputs diverges.
\item The corrected terminal alternatives are exactly those of
      Theorem~\ref{thm:v27-terminal-alternative}.
\end{enumerate}
No item excludes every terminal sequence for a fixed datum.
\end{theorem}

\begin{proof}
Items~\textup{(V27.1)--(V27.3)} are
Theorem~\ref{thm:v27-minimum-heat-carrier} and
Countertheorem~\ref{cth:v27-linear-envelope-heat-current}.  Item
\textup{(V27.4)} is Theorem~\ref{thm:v27-heat-reduced-equation}.  Items
\textup{(V27.5)--(V27.6)} are
Theorem~\ref{thm:v27-three-term-ledger},
Theorem~\ref{thm:v27-one-third-alternative}, and
Corollary~\ref{cor:v27-hierarchical-opening}.  Item~\textup{(V27.7)} is
Theorems~\ref{thm:v27-sec-antidissipation-pulse} and
\ref{thm:v27-same-cell-pulse}.  Item~\textup{(V27.8)} is
Theorem~\ref{thm:v27-heat-stock-diverges} and
Corollary~\ref{cor:v27-aggregate-diverges}.  The final item is
Theorem~\ref{thm:v27-terminal-alternative}.
\end{proof}

\begingroup\small
\begin{longtable}{>{\raggedright\arraybackslash}p{0.27\textwidth}
                  >{\raggedright\arraybackslash}p{0.18\textwidth}
                  >{\raggedright\arraybackslash}p{0.47\textwidth}}
\caption{Version V27 proof-status ledger.}\label{tab:v27-ledger}\\
\toprule
\textbf{Row}&\textbf{Status}&\textbf{Exact advance or limitation}\\
\midrule
\endfirsthead
\toprule
\textbf{Row}&\textbf{Status}&\textbf{Exact advance or limitation}\\
\midrule
\endhead
V1--V26 prefix&\PROVED{} preserved
 &Complete V26 preterminal source retained byte-for-byte; no inherited theorem
 rewritten.\\[0.35em]
Stokes-dephased source Gram&\PROVED{} exact
 &Its support is the unique minimum heat-reducing carrier of the occurring
 daughter history.\\[0.35em]
Source action under completion&\PROVED{} invariant
 &Dephasing preserves the trace exactly; rank growth is not a second source
 occurrence.\\[0.35em]
V26 heat work&\PROVED{} refined
 &It is an exact nonreducing-compression diagnostic and vanishes on the
 canonical heat completion.\\[0.35em]
Heat-reduced adjoint equation&\PROVED{} exact
 &One physical forcing plus one dynamic cross-carrier current; no heat-range
 remainder.\\[0.35em]
Assembly-first work ledger&\PROVED{} exact
 &One-third secant, transfer, or forcing alternative; components are opened
 only after their aggregate survives.\\[0.35em]
Component-first interpretation&\OBSTRUCTION{} explicit
 &Divergent component works can cancel exactly inside a zero aggregate.\\[0.35em]
Secant pulse&\PROVED{} scale local
 &Negative secant work forces a finite-frequency gradient and enstrophy pulse.\\[0.35em]
Material-cell incidence&\PROVED{} exact alternative
 &At least half of the pulse is in the inherited moving cell or at least half
 is outside it.\\[0.35em]
Fixed scalar with vanishing daughter action&\PROVED{} cannot be silent
 &The heat-reduced positive stock diverges; one assembled primary output must
 diverge.\\[0.35em]
Exclusion of assembled outputs&\OPEN{} physical PDE
 &Requires a same-cell upper stock, signed dynamic-transfer estimate, or
 direct forcing/backscatter estimate.\\[0.35em]
Three-dimensional regularity&\OPEN
 &No surviving terminal packet has been excluded for every datum.\\
\bottomrule
\end{longtable}
\endgroup

% =============================================================================
\section{Exact mandate after Version V27}
\label{sec:v27-next}
% =============================================================================

\begin{programme}[Evaluate the three assembled reads; do not enlarge the carrier]
\label{prog:v27-next}
Preserve Versions~V1--V27 verbatim.  Do not introduce another source class,
adjoint field, response graph, path law, shell count, Moore--Penrose
factorization, Dini modulus, carrier completion, or compactness space.

Proceed only in the following order.
\begin{enumerate}[label=\textup{(V28.\arabic*)},leftmargin=5.9em]
\item Retain the inherited material-route, cell-capture, boundary-current,
      target-normalization, represented-source, and support-incidence gates.
      Construct the heat-reducing carrier from the source before reading the
      return writer.
\item Evaluate the three primary quantities
      \[
       \mathcal N_{\rm h}^-,\qquad
       \mathcal W_{\rm h}^{\rm tr},\qquad
       \mathcal W_{\rm h}^{\rm for}
      \]
      in the exact ledger \eqref{eq:v27-three-term-ledger}.  Do not inspect a
      component first.
\item On the secant branch, use the same time and the same inherited material
      cell in Theorem~\ref{thm:v27-same-cell-pulse}.  Return the inside packet
      or the literal outside-cell escape; no replacement ball is permitted.
\item On the transfer branch, evaluate the complete off-block current
      \(-\int\langle\eta_{\rm h},\mathscr L_1\chi_{\rm h}\rangle\) before
      separating symmetric and skew pieces.
\item On the forcing branch, assemble represented returns before separating
      local from tail forcing.  Open head--daughter/local-exterior or
      adjacent-octave/backscatter components only after the corresponding
      aggregate survives.
\item If all three primary reads are subcritical relative to
      \(\mathcal P_{\rm h}\), use the exact ledger directly to exclude the
      fixed-scalar/vanishing-daughter branch.
\item If the first-birth scalar is negligible and the V22 old and represented
      target residuals pass, terminate NS--A and export the target-native
      scalar ancestry card.
\item If no finite head captures the support-corrected source and its heat
      completion, retain fused all-shell selected-source escape.
\end{enumerate}
\end{programme}

\begin{assessment}[Programme position after Version V27]
\label{ass:v27-position}
The one-scalar card no longer contains a heat-range obstruction created by a
nonreducing compression, and it no longer opens six dynamic/forcing
components before their signed aggregates are known.  Its first unpopulated
finite card is
\[
 \boxed{
 \left(
  \mathcal N_{\rm h}^-,
  \mathcal W_{\rm h}^{\rm tr},
  \mathcal W_{\rm h}^{\rm for}
 \right),}
\]
together with the already occurring positive stock \(\mathcal P_{\rm h}\)
and, on the secant branch, the exact material-cell split
\((\mathcal E_n^{\rm in},\mathcal E_n^{\rm out})\).  Adding another carrier
or all-direction norm before these reads are evaluated would be circular.
\end{assessment}

\paragraph{Version V27 history.}
Preserved the complete V26 source.  Replaced its minimum linear source
carrier, for purposes of the fixed Stokes propagation, by the unique minimum
heat-reducing carrier obtained from the Stokes-dephased positive source Gram.
Proved exact source containment, minimality, heat-cyclic characterization,
and action preservation; exhibited a two-mode compression example producing
nonzero V26 heat work with no change in the source scalar; derived the
heat-reduced terminal-zero adjoint equation and its three-term signed ledger;
proved the one-third assembly-first alternative and hierarchical component
opening; retained the V26 secant estimate with the new carrier; and tested the
resulting enstrophy pulse against the exact inherited material cell.  No
terminal Dini gain, upper critical-action estimate, compact-profile rigidity,
or global regularity conclusion was claimed.

\begin{assessment}[Append-only integrity]
\label{ass:v27-integrity}
The complete Version~V26 source preceding its sole final executable document
terminator is preserved verbatim.  Its preterminal SHA--256 digest is
\begin{center}\scriptsize\ttfamily
67fe7cbe49d490c49096c44f3b9bec5eadaff479f8ebecb027d0532efbc333ec
\end{center}
The present material begins at Section~\ref{sec:v27-entry}; the final command
below is again unique.
\end{assessment}

\addtocontents{toc}{\protect\endgroup}
% =============================================================================
% END VERSION V27 MATERIAL -- append later versions before final terminator
% =============================================================================

% APPEND ALL LATER VERSIONS IMMEDIATELY ABOVE THE SOLE FINAL LINE BELOW.


% =============================================================================
% VERSION V28: NEW MATERIAL.  The complete V1--V27 source above is preserved
% verbatim; only its sole final executable document terminator was removed.
% =============================================================================

\clearpage
\makeatletter
\addtocontents{toc}{\protect\begingroup\protect\makeatletter}
\addtocontents{toc}{\protect\def\protect\l@section{\protect\@dottedtocline{1}{0em}{4.7em}}}
\addtocontents{toc}{\protect\def\protect\l@subsection{\protect\@dottedtocline{2}{1.5em}{5.3em}}}
\addtocontents{toc}{\protect\makeatother}
\makeatother
\renewcommand{\thesection}{V28.\arabic{section}}
\renewcommand{\thesubsection}{V28.\arabic{section}.\arabic{subsection}}
\renewcommand{\theequation}{V28.\arabic{equation}}
\setcounter{section}{0}
\setcounter{equation}{0}
\renewcommand{\theHsection}{V28.\arabic{section}}
\renewcommand{\theHsubsection}{V28.\arabic{section}.\arabic{subsection}}
\renewcommand{\theHtheorem}{V28.\arabic{section}.\arabic{theorem}}
\renewcommand{\theHequation}{V28.\arabic{equation}}

\begin{center}
{\LARGE\bfseries NS--A: Version V28}\\[0.5em]
{\large The Chronological Stokes Source Filtration, Target-Visible Birth Promotion,\\
and the Future-Born/Source-Null Return Split}\\[0.5em]
{\normalsize 4 September 2026}
\end{center}
\addcontentsline{toc}{part}{Version V28: chronological Stokes source filtration and birth-promotion ledger}

\begin{abstract}
Version~V27 replaced the minimum linear source span by the unique smallest
fixed Stokes-reducing carrier containing the complete support-corrected
daughter history.  That construction is exact and remains appropriate for a
hindsight or reusable-bank target.  It uses the whole terminal interval,
however.  A direction which first occurs late in the daughter history is
therefore placed in the carrier at every earlier time.  For source-first
chronology this is an anticipatory compression: it can attribute adjoint
forcing and transfer to a source direction before that direction has
occurred.

The present continuation moves one arrow upstream without changing the
source, its action, the support-compatible represented-source short, the
return writer, or the selected scalar.  It forms the cumulative daughter
Gram, dephases it in the Stokes eigenspaces, and takes the right-continuous
support filtration.  This is the unique smallest increasing Stokes-reducing
projection path containing the source as it occurs.  It has only finitely
many acquisition jumps on each finite octave.  The orthogonal jump
projections give a literal source-first decomposition: every source component
vanishes before its own acquisition time, the source action is a Pythagorean
sum, and the selected scalar is the corresponding signed sum.

Projecting the terminal-zero adjoint onto this chronological carrier gives a
piecewise smooth equation with an exact positive promotion debit at every
source-direction birth.  The endpoint-plus-promotion-plus-viscous stock obeys
one no-remainder signed ledger with negative symmetric-secant work, complete
physical forcing, and complete dynamic transfer.  The transfer then splits,
only after assembly, into a future-born source-direction current and a
permanently source-null circulation current.  A two-dimensional exact model
shows that the full-interval V27 carrier can absorb a large pre-birth forcing
packet which is absent from the chronological scalar card.  On the immediate
source-saturation branch the new ledger reduces exactly to V27.

Under the inherited fixed-scalar/vanishing-daughter hypothesis, the
chronological adjoint stock still diverges.  Hence one of only four assembled
outputs survives: negative secant work, future-born transfer, permanently
source-null transfer, or physical forcing.  The secant branch retains the
same inherited material-cell pulse alternative.  No new source class,
response graph, path law, Dini modulus, compactness space, critical upper
budget, or global three-dimensional regularity conclusion is introduced.
\end{abstract}

% =============================================================================
\section{The full-interval carrier is a hindsight carrier}
\label{sec:v28-entry}
% =============================================================================

Retain the complete V27 finite-head card.  Thus
\[
 P_K=\one_{(0,K]}(\Lambda),\qquad
 Q_K=I-P_K,\qquad
 \Pi_1=\one_{(K,2K]}(\Lambda),
\]
the selected support-corrected daughter history is
\(d^{\rm C}\in L^2(I;\operatorname{Ran}\Pi_1)\), where
\(I=[a,b]\), and the terminal-zero adjoint satisfies
\begin{equation}
 -\partial_t\psi+\nu A\psi+\mathscr L_1(t)^*\psi=F_1,
 \qquad \psi(b)=0,
 \label{eq:v28-inherited-adjoint}
\end{equation}
with \(\mathscr L_1=\Pi_1\mathscr L_Q^u\Pi_1\) and
\(F_1\) the already assembled physical forcing of
\eqref{eq:v25-four-force-sum}.  The selected scalar is
\begin{equation}
 \mathfrak h^{\rm C}
 =\int_a^b\langle d^{\rm C}(t),\psi(t)\rangle\,\dd t.
 \label{eq:v28-inherited-scalar}
\end{equation}

Let \(P_{\rm h}\) be the V27 projection onto the final heat-reducing
source carrier \(\mathcal E_d^{\rm heat}\).  It is generated by the complete
history on \([a,b]\).  Its use at times before every source direction has
appeared is therefore a hindsight compression.

\begin{firewall}[V27 remains exact for its declared fixed-carrier target]
\label{fire:v28-v27-valid}
No identity of Version~V27 is withdrawn.  The fixed projection \(P_{\rm h}\)
is the unique minimum Stokes-reducing carrier containing the complete
source history and its ledger is exact.  Version~V28 asks a stronger
chronological question: whether a direction is used before its first
occurrence in the selected daughter history.  The two cards coincide when
the final source carrier is already generated at the entrance.  They need
not coincide when source directions are acquired later.
\end{firewall}

% =============================================================================
\section{The cumulative Stokes-dephased source filtration}
\label{sec:v28-filtration}
% =============================================================================

Use the finite Stokes spectral resolution of \eqref{eq:v27-shell-spectral-resolution},
\[
 A\Pi_1=\sum_{\lambda\in\Sigma_K}\lambda E_\lambda.
\]
For \(a\le t\le b\), define the cumulative source Gram and its Stokes
dephasing by
\begin{align}
 \mathbb Q_d(t)
 &:=\int_a^t d^{\rm C}(s)\otimes d^{\rm C}(s)\,\dd s,
 \label{eq:v28-cumulative-Gram}\\
 \boxed{
 \mathbb Q_d^{\rm heat}(t)
 :=\sum_{\lambda\in\Sigma_K}
 E_\lambda\mathbb Q_d(t)E_\lambda.}
 \label{eq:v28-cumulative-heat-Gram}
\end{align}
For \(t<b\), put
\begin{equation}
 \boxed{
 P_{\rm chr}(t)
 :=\lim_{r\downarrow t,\ r>t}
 \operatorname{supp}\mathbb Q_d^{\rm heat}(r),}
 \qquad
 P_{\rm chr}(b):=
 \operatorname{supp}\mathbb Q_d^{\rm heat}(b),
 \label{eq:v28-chronological-projection}
\end{equation}
where the limit is in operator norm; it exists because the projections are
nested and the space is finite dimensional.  Set \(P_{\rm chr}(a-)=0\).

\begin{theorem}[Canonical chronological Stokes source filtration]
\label{thm:v28-chronological-filtration}
The projection path \(P_{\rm chr}\) has the following properties.
\begin{enumerate}[label=\textup{(F\arabic*)},leftmargin=3.5em]
\item It is right continuous and increasing:
      \begin{equation}
       \boxed{
       0\le P_{\rm chr}(s)\le P_{\rm chr}(t)\le\Pi_1
       \qquad(a\le s\le t\le b).}
       \label{eq:v28-filtration-monotone}
      \end{equation}
\item It reduces the Stokes operator:
      \begin{equation}
       \boxed{[P_{\rm chr}(t),A]=0.}
       \label{eq:v28-filtration-reduces-A}
      \end{equation}
\item It contains the occurring source at its time of occurrence:
      \begin{equation}
       \boxed{
       P_{\rm chr}(t)d^{\rm C}(t)=d^{\rm C}(t)
       \quad\text{for almost every }t\in I.}
       \label{eq:v28-source-contained-chronologically}
      \end{equation}
\item Its terminal value is the V27 carrier:
      \begin{equation}
       \boxed{P_{\rm chr}(b)=P_{\rm h}.}
       \label{eq:v28-final-carrier-equality}
      \end{equation}
\item It is minimal.  If \(R(t)\) is any right-continuous increasing path of
      orthogonal projections which commute with \(A\) and satisfy
      \(R(t)d^{\rm C}(t)=d^{\rm C}(t)\) almost everywhere, then
      \begin{equation}
       \boxed{P_{\rm chr}(t)\le R(t)\qquad(a\le t\le b).}
       \label{eq:v28-filtration-minimal}
      \end{equation}
\end{enumerate}
\end{theorem}

\begin{proof}
For \(s\le t\),
\[
 \mathbb Q_d^{\rm heat}(t)-\mathbb Q_d^{\rm heat}(s)
 =\sum_{\lambda}
 E_\lambda\left(\int_s^t d^{\rm C}(r)\otimes d^{\rm C}(r)\,\dd r\right)E_\lambda
 \succeq0.
\]
If \(0\preceq B\preceq C\), then \(\ker C\subseteq\ker B\), hence
\(\operatorname{Ran}B\subseteq\operatorname{Ran}C\).  The support
projections in \eqref{eq:v28-chronological-projection} are therefore nested,
which proves monotonicity and the existence and right continuity of the
finite-dimensional right limit.  Every dephased Gram commutes with all
\(E_\lambda\), hence its support commutes with \(A\), proving
\eqref{eq:v28-filtration-reduces-A}.  At \(t=b\), the definition is exactly
that of \(P_{\rm h}\).

Because an increasing finite-dimensional projection path can change only
when its integer rank increases, it is constant off a finite set.  On any
open interval on which it equals a fixed projection \(P\), choose a later
point \(t_1\) in the same interval.  The range of
\(\mathbb Q_d^{\rm heat}(t_1)\) lies in \(\operatorname{Ran}P\).  Hence,
for every \(x\in\operatorname{Ran}(I-P)\),
\[
 0
 =\langle x,\mathbb Q_d^{\rm heat}(t_1)x\rangle
 =\sum_\lambda\int_a^{t_1}
   |\langle E_\lambda x,d^{\rm C}(r)\rangle|^2\,\dd r.
\]
Thus every \(E_\lambda d^{\rm C}(r)\), and therefore
\(d^{\rm C}(r)\), lies in \(\operatorname{Ran}P\) for almost every
\(r\) in that interval.  The exceptional jump set is finite, proving
\eqref{eq:v28-source-contained-chronologically}.

Finally, let \(R\) be as in the statement.  By monotonicity, for almost
every \(s\le r\), the vector \(d^{\rm C}(s)\) belongs to
\(\operatorname{Ran}R(r)\).  Since \(R(r)\) commutes with every
\(E_\lambda\), it also contains every \(E_\lambda d^{\rm C}(s)\).
Therefore
\(\operatorname{Ran}\mathbb Q_d^{\rm heat}(r)\subseteq
\operatorname{Ran}R(r)\).  Letting \(r\downarrow t\) and using right
continuity gives \eqref{eq:v28-filtration-minimal}.
\end{proof}

\begin{corollary}[Finite acquisition times and orthogonal jump projections]
\label{cor:v28-acquisition-times}
There is a finite set \(\mathcal T_d\subset[a,b)\) such that
\(P_{\rm chr}\) is constant on each component of
\(I\setminus\mathcal T_d\).  For \(\tau\in\mathcal T_d\), define
\begin{equation}
 \boxed{
 \Delta_\tau
 :=P_{\rm chr}(\tau)-P_{\rm chr}(\tau-).}
 \label{eq:v28-jump-projection}
\end{equation}
Then every \(\Delta_\tau\) is an orthogonal projection commuting with
\(A\), the nonzero \(\Delta_\tau\) are mutually orthogonal, and
\begin{equation}
 \boxed{
 P_{\rm h}=\sum_{\tau\in\mathcal T_d}\Delta_\tau,
 \qquad
 \#\mathcal T_d\le\operatorname{rank}P_{\rm h}.}
 \label{eq:v28-jump-resolution}
\end{equation}
\end{corollary}

\begin{proof}
The rank of an increasing projection path is a nondecreasing integer bounded
by \(\operatorname{rank}\Pi_1\).  It can therefore increase only finitely
many times, and nested projections of equal rank are equal.  If
\(P_-\le P_+\), then \(P_-P_+=P_-=P_+P_-\), so
\(P_+-P_-\) is the orthogonal projection onto
\(\operatorname{Ran}P_+\ominus\operatorname{Ran}P_-\).  Successive
increments are orthogonal and telescope to the final projection.
\end{proof}

\begin{assessment}[No new source action is created]
\label{ass:v28-no-new-source-action}
The chronological filtration changes only when a new source direction first
appears.  By \eqref{eq:v28-source-contained-chronologically},
\[
 \int_I\|P_{\rm chr}(t)d^{\rm C}(t)\|_2^2\,\dd t
 =\int_I\|d^{\rm C}(t)\|_2^2\,\dd t.
\]
The projection path is generated from the already selected target-relative
source residual.  It is not a new Navier--Stokes source occurrence and does
not replace the raw daughter or the original represented-source atlas in a
full ancestry statement.
\end{assessment}

% =============================================================================
\section{Literal source-first atomization}
\label{sec:v28-source-atomization}
% =============================================================================

For each acquisition time, put
\begin{equation}
 d_\tau(t):=\Delta_\tau d^{\rm C}(t),
 \qquad
 \psi_\tau(t):=\Delta_\tau\psi(t).
 \label{eq:v28-birth-components}
\end{equation}

\begin{theorem}[Orthogonal chronological source decomposition]
\label{thm:v28-source-first-Pythagoras}
The acquisition components satisfy
\begin{equation}
 \boxed{
 d_\tau(t)=0
 \quad\text{for almost every }t<\tau.}
 \label{eq:v28-component-not-before-birth}
\end{equation}
Moreover,
\begin{align}
 \boxed{
 d^{\rm C}(t)=\sum_{\tau\in\mathcal T_d}d_\tau(t)
 \quad\text{for almost every }t,}
 \label{eq:v28-source-birth-sum}\\
 \boxed{
 \int_I\|d^{\rm C}(t)\|_2^2\,\dd t
 =\sum_{\tau\in\mathcal T_d}
   \int_\tau^b\|d_\tau(t)\|_2^2\,\dd t,}
 \label{eq:v28-source-action-Pythagoras}\\
 \boxed{
 \mathfrak h^{\rm C}
 =\sum_{\tau\in\mathcal T_d}
   \int_\tau^b
   \langle d_\tau(t),\psi_\tau(t)\rangle\,\dd t.}
 \label{eq:v28-scalar-birth-sum}
\end{align}
Thus every source component is charged at its first support acquisition and
is never used before that time.
\end{theorem}

\begin{proof}
If \(t<\tau\), then
\(P_{\rm chr}(t)\le P_{\rm chr}(\tau-)\).  The occurring source lies in
\(\operatorname{Ran}P_{\rm chr}(t)\) almost everywhere, whereas
\(\Delta_\tau\) is orthogonal to
\(\operatorname{Ran}P_{\rm chr}(\tau-)\).  This proves
\eqref{eq:v28-component-not-before-birth}.  The orthogonal resolution
\eqref{eq:v28-jump-resolution} and the final source containment give
\eqref{eq:v28-source-birth-sum}.  Orthogonality of the jump ranges gives the
action Pythagoras.  Finally,
\[
 \langle d_\tau,\psi\rangle
 =\langle\Delta_\tau d^{\rm C},\Delta_\tau\psi\rangle,
\]
and the integrand vanishes before \(\tau\); summing proves
\eqref{eq:v28-scalar-birth-sum}.
\end{proof}

\begin{corollary}[Per-birth Riesz action]
\label{cor:v28-per-birth-Riesz}
For
\begin{align}
 A_\tau^{\rm src}
 &:=\int_\tau^b\|d_\tau(t)\|_2^2\,\dd t,
 &
 A_\tau^{\rm wr}
 &:=\int_\tau^b\|\psi_\tau(t)\|_2^2\,\dd t,
 \label{eq:v28-birth-actions}\\
 h_\tau
 &:=\int_\tau^b
 \langle d_\tau(t),\psi_\tau(t)\rangle\,\dd t,
 \label{eq:v28-birth-scalar}
\end{align}
one has
\begin{equation}
 \boxed{|h_\tau|^2\le A_\tau^{\rm src}A_\tau^{\rm wr}.}
 \label{eq:v28-birth-Riesz-bound}
\end{equation}
The signed values \(h_\tau\) must be assembled before any positive part is
taken; the positive source actions already add without multiplicity by
\eqref{eq:v28-source-action-Pythagoras}.
\end{corollary}

\begin{proof}
This is Cauchy--Schwarz in
\(L^2((\tau,b);\operatorname{Ran}\Delta_\tau)\).
\end{proof}

% =============================================================================
\section{A full-interval carrier can import pre-birth target work}
\label{sec:v28-anticipation-counterexample}
% =============================================================================

\begin{countertheorem}[Exact anticipatory-compression model]
\label{cth:v28-anticipatory-carrier}
Let \(I=[0,1]\), \(\mathcal H=\mathbb R^2\),
\[
 A=\begin{pmatrix}1&0\\0&2\end{pmatrix},
 \qquad \nu=1,
 \qquad \mathscr L_1=0,
\]
and define
\begin{equation}
 d(t)=
 \begin{cases}
 e_1,&0\le t<\frac12,\\
 e_2,&\frac12\le t\le1,
 \end{cases}
 \qquad
 \psi(t)=(1-t)e_2.
 \label{eq:v28-counter-source-writer}
\end{equation}
Then \(\psi(1)=0\) and
\begin{equation}
 -\psi'(t)+A\psi(t)=F(t),
 \qquad
 F(t)=(3-2t)e_2.
 \label{eq:v28-counter-adjoint}
\end{equation}
The full-interval heat-reducing carrier is \(P_{\rm h}=I\), whereas
\begin{equation}
 P_{\rm chr}(t)=
 \begin{cases}
 e_1\otimes e_1,&0\le t<\frac12,\\
 I,&\frac12\le t\le1.
 \end{cases}
 \label{eq:v28-counter-chronological-carrier}
\end{equation}
The selected source scalar is
\begin{equation}
 \boxed{
 \int_0^1\langle d(t),\psi(t)\rangle\,\dd t
 =\frac18.}
 \label{eq:v28-counter-scalar}
\end{equation}
The fixed-carrier V27 stock and forcing work are both
\begin{equation}
 \boxed{
 \frac12\|\psi(0)\|^2
 +\int_0^1\|A^{1/2}\psi(t)\|^2\,\dd t
 =\int_0^1\langle F(t),\psi(t)\rangle\,\dd t
 =\frac76.}
 \label{eq:v28-counter-static-ledger}
\end{equation}
The chronological promotion-plus-viscous stock and forcing work are instead
\begin{equation}
 \boxed{
 \frac12\|e_2\psi(\tfrac12)\|^2
 +\int_{1/2}^1\|A^{1/2}\psi(t)\|^2\,\dd t
 =\int_{1/2}^1\langle F(t),\psi(t)\rangle\,\dd t
 =\frac5{24}.}
 \label{eq:v28-counter-chronological-ledger}
\end{equation}
Thus the full-interval carrier includes \(23/24\) units of forcing work in
the \(e_2\) direction before that source direction is acquired.  The
chronological card removes this anticipatory work while preserving the
selected scalar.
\end{countertheorem}

\begin{proof}
The adjoint equation is immediate.  The cumulative source span is
\(\mathbb Re_1\) before \(1/2\) and \(\mathbb R^2\) afterward, proving
\eqref{eq:v28-counter-chronological-carrier}.  Direct integration gives
\[
 \int_{1/2}^1(1-t)\,\dd t=\frac18,
\]
\[
 \frac12+2\int_0^1(1-t)^2\,\dd t
 =\frac12+\frac23=\frac76,
\]
and
\[
 \int_0^1(3-2t)(1-t)\,\dd t=\frac76.
\]
At the acquisition time, the promotion is
\(\frac12(1/2)^2=1/8\), while
\[
 2\int_{1/2}^1(1-t)^2\,\dd t=\frac1{12},
\]
so the chronological stock is \(1/8+1/12=5/24\).  The forcing integral on
\([1/2,1]\) has the same value.  The difference
\(7/6-5/24=23/24\) is the pre-birth forcing absorbed by the hindsight
carrier.  This is a finite linear counterexample to a carrier
interpretation, not a Navier--Stokes trajectory.
\end{proof}

% =============================================================================
\section{The chronological adjoint and its promotion ledger}
\label{sec:v28-chronological-ledger}
% =============================================================================

Define
\begin{equation}
 \chi_{\rm chr}(t):=P_{\rm chr}(t)\psi(t),
 \qquad
 \eta_{\rm chr}(t):=(\Pi_1-P_{\rm chr}(t))\psi(t).
 \label{eq:v28-chronological-split}
\end{equation}
The adjoint \(\psi\) is continuous on the finite smooth card.  Hence, at an
acquisition time \(\tau\),
\begin{equation}
 \boxed{
 \chi_{\rm chr}(\tau)-\chi_{\rm chr}(\tau-)
 =\Delta_\tau\psi(\tau).}
 \label{eq:v28-promotion-jump}
\end{equation}
Because \(\Delta_\tau\psi(\tau)\perp
P_{\rm chr}(\tau-)\psi(\tau)\),
\begin{equation}
 \boxed{
 \|\chi_{\rm chr}(\tau)\|_2^2
 -\|\chi_{\rm chr}(\tau-)\|_2^2
 =\|\Delta_\tau\psi(\tau)\|_2^2.}
 \label{eq:v28-promotion-Pythagoras}
\end{equation}

On each component of \(I\setminus\mathcal T_d\), the projection is constant.
Writing
\(\mathscr L_1^*=\mathscr S_1-\mathscr K_1\), one obtains
\begin{equation}
 \boxed{
 \begin{aligned}
 -\partial_t\chi_{\rm chr}
 &+\nu A\chi_{\rm chr}
 +P_{\rm chr}\mathscr S_1P_{\rm chr}\chi_{\rm chr}
 -P_{\rm chr}\mathscr K_1P_{\rm chr}\chi_{\rm chr}
 \\
 &=P_{\rm chr}F_1
   -P_{\rm chr}\mathscr L_1^*\eta_{\rm chr}.
 \end{aligned}}
 \label{eq:v28-piecewise-adjoint}
\end{equation}

Define the positive target-visible source-birth promotion,
viscous action, and total chronological stock by
\begin{align}
 \boxed{
 \mathcal B_{\rm chr}
 :=\frac12\sum_{\tau\in\mathcal T_d}
   \|\Delta_\tau\psi(\tau)\|_2^2,}
 \label{eq:v28-promotion-stock}\\
 \mathcal D_{\rm chr}
 &:=\nu\int_I\|A^{1/2}\chi_{\rm chr}(t)\|_2^2\,\dd t,
 \label{eq:v28-chronological-viscous-stock}\\
 \boxed{
 \mathcal P_{\rm chr}:=\mathcal B_{\rm chr}+\mathcal D_{\rm chr}.}
 \label{eq:v28-total-chronological-stock}
\end{align}
Define also
\begin{align}
 \mathcal N_{\rm chr}
 &:=\int_I
 \langle\mathscr S_1(t)\chi_{\rm chr}(t),
        \chi_{\rm chr}(t)\rangle\,\dd t,
 \label{eq:v28-chronological-secant}\\
 \mathcal N_{\rm chr}^-
 &:=\int_I
 \bigl(-\langle\mathscr S_1\chi_{\rm chr},
                    \chi_{\rm chr}\rangle\bigr)_+\,\dd t,
 \label{eq:v28-chronological-negative-secant}\\
 \mathcal W_{\rm chr}^{\rm for}
 &:=\int_I\langle F_1,\chi_{\rm chr}\rangle\,\dd t,
 \label{eq:v28-chronological-forcing}\\
 \mathcal W_{\rm chr}^{\rm tr}
 &:=-\int_I
 \langle\eta_{\rm chr},\mathscr L_1\chi_{\rm chr}\rangle\,\dd t.
 \label{eq:v28-chronological-transfer}
\end{align}

\begin{theorem}[Exact chronological promotion ledger]
\label{thm:v28-chronological-ledger}
The source-first projected adjoint satisfies
\begin{equation}
 \boxed{
 \mathcal P_{\rm chr}+\mathcal N_{\rm chr}
 =\mathcal W_{\rm chr}^{\rm for}
  +\mathcal W_{\rm chr}^{\rm tr}.}
 \label{eq:v28-chronological-ledger}
\end{equation}
The positive term \(\mathcal B_{\rm chr}\) is the exact writer energy
promoted when new source directions become physically available.  It is not
a second source action.
\end{theorem}

\begin{proof}
On a component \((s_0,s_1)\) of
\(I\setminus\mathcal T_d\), take the real \(L^2\) inner product of
\eqref{eq:v28-piecewise-adjoint} with \(\chi_{\rm chr}\) and integrate.
The internal skew form vanishes.  The time derivative contributes
\[
 \frac12\|\chi_{\rm chr}(s_0)\|_2^2
 -\frac12\|\chi_{\rm chr}(s_1-)\|_2^2.
\]
Summing over the finitely many intervals telescopes.  The terminal value is
zero because \(\psi(b)=0\).  At each internal acquisition time, the
uncancelled endpoint contribution is
\[
 \frac12\left(
 \|\chi_{\rm chr}(\tau)\|_2^2
 -\|\chi_{\rm chr}(\tau-)\|_2^2
 \right)
 =\frac12\|\Delta_\tau\psi(\tau)\|_2^2
\]
by \eqref{eq:v28-promotion-Pythagoras}; the same formula includes the
initial acquisition at \(a\) under the convention
\(P_{\rm chr}(a-)=0\).  The remaining terms are exactly those in
\eqref{eq:v28-chronological-ledger}.
\end{proof}

\begin{corollary}[Immediate source saturation recovers V27 exactly]
\label{cor:v28-v27-recovery}
Suppose
\begin{equation}
 P_{\rm chr}(t)=P_{\rm h}
 \quad\text{for almost every }t\in[a,b).
 \label{eq:v28-immediate-saturation}
\end{equation}
Then the only nonzero acquisition is the entrance acquisition,
\begin{equation}
 \mathcal B_{\rm chr}
 =\frac12\|P_{\rm h}\psi(a)\|_2^2,
 \label{eq:v28-entrance-promotion}
\end{equation}
and every term of \eqref{eq:v28-chronological-ledger} is exactly the
corresponding V27 term in \eqref{eq:v27-three-term-ledger}.
\end{corollary}

\begin{proof}
Under \eqref{eq:v28-immediate-saturation}, the chronological projection is
the fixed V27 projection off a null set.  Its jump from zero at the entrance
gives \eqref{eq:v28-entrance-promotion}; all spacetime integrals and the
complementary projection agree with V27.
\end{proof}

\begin{corollary}[Promotion-heavy birth or source-direction proliferation]
\label{cor:v28-promotion-rank}
Let \(r_{\rm h}=\operatorname{rank}P_{\rm h}>0\).  Then
\begin{equation}
 \boxed{
 \max_{\tau\in\mathcal T_d}
 \frac12\|\Delta_\tau\psi(\tau)\|_2^2
 \ge\frac{\mathcal B_{\rm chr}}{r_{\rm h}}.}
 \label{eq:v28-heavy-promotion}
\end{equation}
Consequently, on a sequence for which
\(\mathcal B_{{\rm chr},n}\ge\theta\mathcal P_{{\rm chr},n}\), either the
final carrier ranks diverge or a single source-direction acquisition carries
a fixed rank-normalized fraction of the chronological stock.
\end{corollary}

\begin{proof}
There are at most \(r_{\rm h}\) nonzero orthogonal jump projections and the
promotion stock is half the sum of their squared writer amplitudes.  The
largest summand is at least the average.
\end{proof}

% =============================================================================
\section{Future-born and permanently source-null transfer}
\label{sec:v28-transfer-split}
% =============================================================================

Since \(P_{\rm chr}(t)\le P_{\rm h}\le\Pi_1\), define
\begin{equation}
 P_{\rm fut}(t):=P_{\rm h}-P_{\rm chr}(t),
 \qquad
 P_{\rm null}:=\Pi_1-P_{\rm h}.
 \label{eq:v28-future-null-projections}
\end{equation}
These are orthogonal projections with orthogonal ranges, and
\begin{equation}
 \eta_{\rm chr}=P_{\rm fut}\psi+P_{\rm null}\psi.
 \label{eq:v28-eta-future-null}
\end{equation}
Define
\begin{align}
 \mathcal W_{\rm chr}^{\rm fut}
 &:=-\int_I
 \langle P_{\rm fut}(t)\psi(t),
        \mathscr L_1(t)\chi_{\rm chr}(t)\rangle\,\dd t,
 \label{eq:v28-future-born-transfer}\\
 \mathcal W_{\rm chr}^{\rm null}
 &:=-\int_I
 \langle P_{\rm null}\psi(t),
        \mathscr L_1(t)\chi_{\rm chr}(t)\rangle\,\dd t.
 \label{eq:v28-source-null-transfer}
\end{align}

\begin{theorem}[Exact chronology of the dynamic transfer]
\label{thm:v28-transfer-chronology}
The complete transfer has the assembly-first split
\begin{equation}
 \boxed{
 \mathcal W_{\rm chr}^{\rm tr}
 =\mathcal W_{\rm chr}^{\rm fut}
  +\mathcal W_{\rm chr}^{\rm null}.}
 \label{eq:v28-transfer-two-way}
\end{equation}
Moreover, almost everywhere,
\begin{equation}
 P_{\rm fut}(t)
 =\sum_{\substack{\tau\in\mathcal T_d\\\tau>t}}\Delta_\tau,
 \label{eq:v28-future-projection-sum}
\end{equation}
and hence
\begin{equation}
 \boxed{
 \mathcal W_{\rm chr}^{\rm fut}
 =-\sum_{\tau\in\mathcal T_d}
   \int_a^\tau
   \langle\Delta_\tau\psi(t),
          \mathscr L_1(t)\chi_{\rm chr}(t)\rangle\,\dd t.}
 \label{eq:v28-future-transfer-by-birth}
\end{equation}
Thus \(\mathcal W_{\rm chr}^{\rm fut}\) is precisely transfer into source
directions before their first occurrence.  The term
\(\mathcal W_{\rm chr}^{\rm null}\) is transfer into directions which never
belong to the selected source carrier on the card.
\end{theorem}

\begin{proof}
The orthogonal decomposition
\(\Pi_1-P_{\rm chr}=(P_{\rm h}-P_{\rm chr})+(\Pi_1-P_{\rm h})\)
gives \eqref{eq:v28-transfer-two-way}.  The jump resolution of the final
carrier implies \eqref{eq:v28-future-projection-sum}: at time \(t\), exactly
the increments acquired after \(t\) remain outside the chronological
carrier but inside the final carrier.  Substitute this finite sum into
\eqref{eq:v28-future-born-transfer} and interchange sum and integral.
\end{proof}

\begin{firewall}[Future-born transfer is not another source birth]
\label{fire:v28-no-double-birth}
The future-born current is generated by the already occurring dynamics acting
on a direction whose source incidence appears later.  At its acquisition
time the direction enters the carrier through the positive promotion term
\(\frac12\|\Delta_\tau\psi(\tau)\|^2\).  Neither the prior transfer nor the
later promotion is licensed as a second nonlinear source occurrence.  The
source birth remains the component \(d_\tau\) of
Theorem~\ref{thm:v28-source-first-Pythagoras}.
\end{firewall}

% =============================================================================
\section{The chronological four-way terminal work alternative}
\label{sec:v28-work-alternative}
% =============================================================================

\begin{theorem}[One-third chronological alternative]
\label{thm:v28-one-third}
If \(\mathcal P_{\rm chr}>0\), then at least one of
\begin{align}
 \boxed{
 \mathcal N_{\rm chr}^-
 \ge\frac13\mathcal P_{\rm chr},}
 \label{eq:v28-one-third-secant}\\
 \boxed{
 |\mathcal W_{\rm chr}^{\rm tr}|
 \ge\frac13\mathcal P_{\rm chr},}
 \label{eq:v28-one-third-transfer}\\
 \boxed{
 |\mathcal W_{\rm chr}^{\rm for}|
 \ge\frac13\mathcal P_{\rm chr}}
 \label{eq:v28-one-third-forcing}
\end{align}
holds.  On the transfer branch, only after the complete transfer survives,
one of
\begin{equation}
 \boxed{
 |\mathcal W_{\rm chr}^{\rm fut}|
 \ge\frac16\mathcal P_{\rm chr},
 \qquad\text{or}\qquad
 |\mathcal W_{\rm chr}^{\rm null}|
 \ge\frac16\mathcal P_{\rm chr}}
 \label{eq:v28-transfer-components}
\end{equation}
holds.
\end{theorem}

\begin{proof}
Since \(-\mathcal N_{\rm chr}\le\mathcal N_{\rm chr}^-\), the exact ledger
gives
\[
 \mathcal P_{\rm chr}
 \le
 \mathcal N_{\rm chr}^-
 +|\mathcal W_{\rm chr}^{\rm tr}|
 +|\mathcal W_{\rm chr}^{\rm for}|.
\]
At least one term is at least one third of the left side.  On the transfer
branch, use \eqref{eq:v28-transfer-two-way} and the triangle inequality only
after the aggregate has been selected.
\end{proof}

\begin{corollary}[The secant branch retains the exact inherited material-cell test]
\label{cor:v28-chronological-secant-cell}
On branch \eqref{eq:v28-one-third-secant}, there exists
\(t_*\in I\setminus\mathcal T_d\) such that
\begin{equation}
 \boxed{
 \mathscr X_K(t_*)\ge\frac{\nu K^2}{3C_0},}
 \label{eq:v28-chronological-gradient-pulse}
\end{equation}
and
\begin{equation}
 \boxed{
 \|\nabla P_{\le4K}u(t_*)\|_2^2\ge c_1\nu^2K.}
 \label{eq:v28-chronological-enstrophy-pulse}
\end{equation}
On the terminal sequence, the exact inherited material cell
\(\mathcal M_n(t_n^*)\) then satisfies the same inside/outside alternative as
Theorem~\ref{thm:v27-same-cell-pulse}.
\end{corollary}

\begin{proof}
The proof of Theorem~\ref{thm:v27-sec-antidissipation-pulse} uses only
\(\chi\in\operatorname{Ran}\Pi_1\), the V26 secant bound, and
\(\mathcal P\ge\nu K^2\int\|\chi\|^2\).  All three statements hold for
\(\chi_{\rm chr}\).  The finite acquisition set has zero measure, so the
maximizing time may be chosen outside it.  The Bernstein and exact
material-cell arguments of V27 are unchanged.
\end{proof}

% =============================================================================
\section{The fixed-scalar branch remains non-silent after chronology is restored}
\label{sec:v28-fixed-scalar}
% =============================================================================

On the terminal sequence define
\begin{align}
 \mathfrak A_{{\rm chr},n}
 &:=\frac{\nu\|\chi_{{\rm chr},n}\|_{L^2(I_n)}^2}
 {\mathcal R_n^2\|\phi_n\|_{L^2(I_n)}^2},
 \label{eq:v28-normalized-chronological-action}\\
 \mathfrak P_{{\rm chr},n}
 &:=\frac{\mathcal P_{{\rm chr},n}}
 {K_n^2\mathcal R_n^2\|\phi_n\|_{L^2(I_n)}^2}.
 \label{eq:v28-normalized-chronological-stock}
\end{align}

\begin{theorem}[Chronological stock divergence]
\label{thm:v28-chronological-stock-diverges}
Assume the inherited target window \eqref{eq:v25-target-action-window} and
\begin{equation}
 \frac{|\mathfrak h_n^{\rm C}|}{\mathcal R_n^2}\ge\delta>0,
 \qquad
 \widetilde{\mathfrak B}_n^{\rm dau}
 =\frac{\|d_n^{\rm C}\|_{L^2(I_n)}^2}
 {\nu\mathcal R_n^2}
 \longrightarrow0.
 \label{eq:v28-fixed-scalar-small-source}
\end{equation}
Then
\begin{equation}
 \boxed{
 \mathfrak A_{{\rm chr},n}
 \ge
 \frac{\delta^2}
 {C_\phi\widetilde{\mathfrak B}_n^{\rm dau}}
 \longrightarrow\infty,}
 \label{eq:v28-chronological-action-diverges}
\end{equation}
and
\begin{equation}
 \boxed{
 \mathfrak P_{{\rm chr},n}
 \ge\mathfrak A_{{\rm chr},n}
 \longrightarrow\infty.}
 \label{eq:v28-chronological-stock-diverges}
\end{equation}
\end{theorem}

\begin{proof}
By \eqref{eq:v28-source-contained-chronologically},
\[
 \mathfrak h_n^{\rm C}
 =\int_{I_n}\langle d_n^{\rm C},\chi_{{\rm chr},n}\rangle\,\dd t.
\]
Cauchy--Schwarz gives
\[
 \delta^2\mathcal R_n^4
 \le
 \|d_n^{\rm C}\|_{L^2}^2
 \|\chi_{{\rm chr},n}\|_{L^2}^2.
\]
Insert the definition of
\(\widetilde{\mathfrak B}_n^{\rm dau}\), multiply by
\(\nu/(\mathcal R_n^2\|\phi_n\|_{L^2}^2)\), and use the upper target
window.  Since \(\chi_{{\rm chr},n}\) lies in the first exterior octave,
\[
 \mathcal P_{{\rm chr},n}
 \ge\nu K_n^2\|\chi_{{\rm chr},n}\|_{L^2}^2,
\]
which proves the second assertion.
\end{proof}

Define the normalized assembled outputs by dividing
\(\mathcal N_{{\rm chr},n}^-\),
\(|\mathcal W_{{\rm chr},n}^{\rm fut}|\),
\(|\mathcal W_{{\rm chr},n}^{\rm null}|\), and
\(|\mathcal W_{{\rm chr},n}^{\rm for}|\) by the positive denominator in
\eqref{eq:v28-normalized-chronological-stock}.

\begin{corollary}[Only four chronological outputs can carry the divergence]
\label{cor:v28-four-outputs}
Under \eqref{eq:v28-fixed-scalar-small-source}, after passage to a priority
subsequence at least one of
\begin{equation}
 \boxed{
 \mathfrak N_{{\rm chr},n}^-\to\infty,
 \quad
 \mathfrak W_{{\rm chr},n}^{\rm fut}\to\infty,
 \quad
 \mathfrak W_{{\rm chr},n}^{\rm null}\to\infty,
 \quad
 \mathfrak W_{{\rm chr},n}^{\rm for}\to\infty}
 \label{eq:v28-four-output-divergence}
\end{equation}
holds.  The first branch gives the inherited same-cell secant pulse; the
second is a source-direction-before-birth transfer; the third is a
permanently source-null circulation; the fourth is the already assembled
physical forcing work.
\end{corollary}

\begin{proof}
Apply Theorem~\ref{thm:v28-one-third} and divide by the normalizing
denominator.  The chronological stock diverges by
Theorem~\ref{thm:v28-chronological-stock-diverges}.  If the transfer branch
is selected, apply \eqref{eq:v28-transfer-components}.  Finiteness of the
list gives a priority subsequence.
\end{proof}

% =============================================================================
\section{Corrected terminal theorem after V28}
\label{sec:v28-terminal}
% =============================================================================

\begin{theorem}[Source-first chronological terminal alternative]
\label{thm:v28-terminal-alternative}
Apply the inherited material-route, cell-capture, signed boundary-current,
target-normalization, represented-source, support-incidence, and finite-head
gates in their established order.  Freeze the physical return writer before
reading its scalar.  Construct the cumulative Stokes source filtration from
the occurring support-corrected daughter history.  After passage to a
priority subsequence, one of the following occurs.
\begin{enumerate}[label=\textup{(G28.\arabic*)},leftmargin=6.4em]
\item \emph{Inherited upstream survivor.}  One of the previously declared
      material, boundary, normalization, ancestry, support, or finite-head
      gates fails.
\item \emph{Corrected daughter first-birth action.}
      \(\widetilde{\mathfrak B}_n^{\rm dau}\) has a positive lower bound.
\item \emph{Immediate source saturation.}
      The final heat-reducing carrier is generated at the entrance.  The
      chronological ledger is exactly the V27 ledger by
      Corollary~\ref{cor:v28-v27-recovery}.
\item \emph{Target-visible source-direction promotion.}
      Delayed source directions occur and the positive promotion stock is
      nonnegligible.  It yields either a heavy acquisition through
      \eqref{eq:v28-heavy-promotion} or diverging final carrier rank.  The
      promotion is retained in the positive ledger and is not repriced as a
      second source.
\item \emph{Chronological secant antidissipation.}
      The scale-local gradient/enstrophy pulse is captured in the exact
      inherited material cell or escapes that cell.
\item \emph{Future-born transfer.}
      A fixed part of the dynamic return enters directions before their
      source acquisition, as in \eqref{eq:v28-future-transfer-by-birth}.
\item \emph{Permanently source-null transfer.}
      A fixed part of the dynamic return lies in the source-null complement
      of the complete selected daughter history.
\item \emph{Assembled physical forcing.}
      The complete forcing work survives.  Its local/tail and four physical
      components are opened only after this aggregate branch is selected.
\item \emph{Target-native ancestry pass.}
      The first-birth scalar is negligible and the V22 old and represented
      target residuals vanish.  NS--A exports the scalar ancestry card and
      stops on the subsequence.
\item \emph{Fused all-shell selected-source escape.}
      No finite head captures the occurring source together with its
      chronological heat filtration.
\end{enumerate}
On the fixed-scalar/vanishing-daughter branch, items~\textup{(G28.5)--(G28.8)}
are exhaustive by Corollary~\ref{cor:v28-four-outputs}.
\end{theorem}

\begin{proof}
The inherited gates and daughter-action branch are unchanged.  Immediate
saturation is Corollary~\ref{cor:v28-v27-recovery}.  The positive promotion
statement and its finite-rank alternative are
Corollary~\ref{cor:v28-promotion-rank}.  The exact chronological work
alternatives are Theorem~\ref{thm:v28-one-third},
Theorem~\ref{thm:v28-transfer-chronology}, and
Corollary~\ref{cor:v28-chronological-secant-cell}.  The fixed-scalar
exhaustion is Corollary~\ref{cor:v28-four-outputs}.  The final two rows are
the inherited scalar-export and all-shell alternatives with the fixed
carrier replaced by its source-first filtration.
\end{proof}

\begin{firewall}[What V28 does not close]
\label{fire:v28-no-global}
The chronological projection removes anticipatory carrier use; it does not
bound the surviving promotion, secant, transfer, or forcing packets.  A
large promotion is not a finite-stock contradiction, a future-born transfer
is not a second source birth, and a permanently source-null current is not
an ancient profile.  V28 supplies no terminal Dini improvement, critical
source-action upper bound, epsilon-regularity theorem, terminal trace,
backward uniqueness, Liouville theorem, or three-dimensional global
regularity conclusion.
\end{firewall}

% =============================================================================
\section{Version V28 integrated conclusion and proof-status ledger}
\label{sec:v28-integrated}
% =============================================================================

\begin{theorem}[Version V28 integrated conclusion]
\label{thm:v28-integrated}
Preserving every theorem and limitation of Versions~V1--V27, Version~V28
proves the following.
\begin{enumerate}[label=\textup{(V28.\arabic*)},leftmargin=5.9em]
\item The cumulative Stokes-dephased daughter Gram gives the unique minimum
      right-continuous increasing Stokes-reducing source filtration.
\item The filtration has finitely many orthogonal acquisition jumps and its
      final value is the V27 heat-reducing carrier.
\item The jump projections give an exact source-first decomposition of the
      source action and selected scalar; every component vanishes before its
      own acquisition time.
\item A two-dimensional exact model shows that the full-interval carrier can
      import forcing work before the corresponding source direction occurs.
\item The chronological adjoint ledger contains one positive source-direction
      promotion stock, viscous action, symmetric secant, assembled forcing,
      and assembled transfer, with no remainder.
\item On immediate source saturation, this ledger reduces exactly to V27.
\item The dynamic transfer splits after assembly into future-born and
      permanently source-null currents; the former has the exact per-birth
      formula \eqref{eq:v28-future-transfer-by-birth}.
\item Under fixed scalar and vanishing daughter action, the chronological
      stock diverges and one of four assembled physical outputs must diverge.
\item The negative-secant branch retains the exact inherited same-cell pulse
      or spatial-incidence escape.
\item The terminal alternatives are precisely those of
      Theorem~\ref{thm:v28-terminal-alternative}.
\end{enumerate}
No item excludes every terminal sequence for a fixed datum.
\end{theorem}

\begin{proof}
Items~\textup{(V28.1)--(V28.3)} are
Theorem~\ref{thm:v28-chronological-filtration},
Corollary~\ref{cor:v28-acquisition-times}, and
Theorem~\ref{thm:v28-source-first-Pythagoras}.  Item~\textup{(V28.4)} is
Countertheorem~\ref{cth:v28-anticipatory-carrier}.  Items
\textup{(V28.5)--(V28.6)} are
Theorem~\ref{thm:v28-chronological-ledger} and
Corollary~\ref{cor:v28-v27-recovery}.  Item~\textup{(V28.7)} is
Theorem~\ref{thm:v28-transfer-chronology}.  Items
\textup{(V28.8)--(V28.9)} are
Theorem~\ref{thm:v28-chronological-stock-diverges},
Corollary~\ref{cor:v28-four-outputs}, and
Corollary~\ref{cor:v28-chronological-secant-cell}.  The final item is
Theorem~\ref{thm:v28-terminal-alternative}.
\end{proof}

\begingroup\small
\begin{longtable}{>{\raggedright\arraybackslash}p{0.27\textwidth}
                  >{\raggedright\arraybackslash}p{0.18\textwidth}
                  >{\raggedright\arraybackslash}p{0.47\textwidth}}
\caption{Version V28 proof-status ledger.}\label{tab:v28-ledger}\\
\toprule
\textbf{Row}&\textbf{Status}&\textbf{Exact advance or limitation}\\
\midrule
\endfirsthead
\toprule
\textbf{Row}&\textbf{Status}&\textbf{Exact advance or limitation}\\
\midrule
\endhead
V1--V27 prefix&\PROVED{} preserved
 &Complete V27 preterminal source retained byte-for-byte; no inherited theorem
 rewritten.\\[0.35em]
Full-interval V27 carrier&\PROVED{} retyped
 &Exact hindsight/reusable-bank carrier; source-first only on the immediate
 saturation branch.\\[0.35em]
Cumulative source filtration&\PROVED{} canonical
 &Unique minimum increasing Stokes-reducing path containing the occurring
 daughter history.\\[0.35em]
Acquisition jumps&\PROVED{} finite
 &Orthogonal source-direction births, at most the final carrier rank.\\[0.35em]
Source-first Pythagoras&\PROVED{} exact
 &Action and selected scalar decompose without using any component before its
 acquisition time.\\[0.35em]
Anticipatory compression&\OBSTRUCTION{} explicit
 &A fixed full-interval carrier can absorb large target work before the source
 direction occurs.\\[0.35em]
Promotion ledger&\PROVED{} exact
 &Positive jump promotion plus viscosity balances signed secant, forcing, and
 dynamic transfer work.\\[0.35em]
Immediate-saturation recovery&\PROVED{} exact
 &The V28 and V27 ledgers coincide when all final source directions occur at
 the entrance.\\[0.35em]
Dynamic-transfer chronology&\PROVED{} exact
 &Future-born directions and permanently source-null directions form the two
 assembled transfer channels.\\[0.35em]
Promotion concentration&\PROVED{} finite alternative
 &One heavy target-visible acquisition or final source-direction rank
 proliferation.\\[0.35em]
Fixed scalar with vanishing daughter action&\PROVED{} non-silent
 &Chronological stock diverges; secant, future-born transfer, source-null
 transfer, or physical forcing must survive.\\[0.35em]
Material-cell pulse&\PROVED{} inherited exactly
 &The chronological secant packet is captured in the original moving cell or
 escapes it.\\[0.35em]
Upper control of surviving packets&\OPEN{} physical PDE
 &Requires a source-birth promotion budget, signed transfer estimate, same-cell
 upper stock, or assembled forcing/backscatter estimate.\\[0.35em]
Three-dimensional regularity&\OPEN
 &No surviving terminal packet has been excluded for every datum.\\
\bottomrule
\end{longtable}
\endgroup

% =============================================================================
\section{Exact mandate after Version V28}
\label{sec:v28-next}
% =============================================================================

\begin{programme}[Evaluate the chronological card; do not restore hindsight]
\label{prog:v28-next}
Preserve Versions~V1--V28 verbatim.  Do not introduce another source class,
adjoint field, response graph, path law, shell count, Moore--Penrose
factorization, carrier completion, Dini modulus, or compactness space.

Proceed only in the following order.
\begin{enumerate}[label=\textup{(V29.\arabic*)},leftmargin=5.9em]
\item Retain every inherited material, boundary, normalization, represented-
      source, support, and finite-head gate.  Build
      \(P_{{\rm chr},n}(t)\) from the source before reading the return writer.
\item Test immediate source saturation.  If it holds, use the V27 card
      exactly; do not duplicate the V28 notation.
\item On delayed source acquisition, acquire the positive promotion stock
      \(\mathcal B_{{\rm chr},n}\) and its actual acquisition times.  A
      promotion is charged to the same daughter source direction, not as a
      second birth.
\item Evaluate the three assembled chronological reads
      \[
       \mathcal N_{\rm chr}^-,\qquad
       \mathcal W_{\rm chr}^{\rm tr},\qquad
       \mathcal W_{\rm chr}^{\rm for}
      \]
      before opening components.
\item On the transfer branch, separate future-born and permanently source-null
      work only through \eqref{eq:v28-transfer-two-way}.  Pair a future-born
      current with its later promotion before any positive price is assigned.
\item On the secant branch, use the same selected time and exact inherited
      material cell.  Return the inside packet or literal outside-cell escape.
\item On the forcing branch, assemble represented returns before the existing
      local/tail and four-carrier decomposition is read.
\item If the first-birth scalar is negligible and the V22 old and represented
      residuals pass, terminate NS--A and export the target-native scalar
      ancestry card.
\item If no finite head captures the occurring source filtration, retain fused
      all-shell selected-source escape.
\end{enumerate}
No later compactness, terminal-trace, backward-uniqueness, or Liouville
argument is opened before this chronological card is resolved.
\end{programme}

\begin{assessment}[Programme position after Version V28]
\label{ass:v28-position}
The first unresolved finite object is no longer the hindsight triple of V27.
It is the source-first card
\[
 \boxed{
 \left(
  \mathcal B_{\rm chr},
  \mathcal N_{\rm chr}^-,
  \mathcal W_{\rm chr}^{\rm fut},
  \mathcal W_{\rm chr}^{\rm null},
  \mathcal W_{\rm chr}^{\rm for}
 \right),}
\]
together with the already occurring viscous stock and, on the secant branch,
the exact inherited material-cell split.  If the source filtration saturates
at the entrance, this card collapses back to V27.  If it does not, using the
full final carrier at earlier times would be anticipatory and therefore
circular with source-first ancestry.
\end{assessment}

\paragraph{Version V28 history.}
Preserved the complete V27 source.  Replaced its full-interval source carrier,
for source-first chronology only, by the right-continuous cumulative
Stokes-dephased support filtration.  Proved monotonicity, Stokes reduction,
minimality, finite acquisition, source containment, exact source-action and
scalar Pythagoras, and a finite anticipatory-compression counterexample.
Derived the piecewise adjoint equation, the positive target-visible
source-direction promotion, and the exact chronological signed ledger.
Separated complete transfer into future-born and permanently source-null
currents only after assembly.  Proved exact recovery of V27 on immediate
source saturation and retained the fixed-scalar divergence and inherited
same-cell secant pulse on the corrected chronological card.  No global
regularity conclusion was claimed.

\begin{assessment}[Append-only integrity]
\label{ass:v28-integrity}
The complete Version~V27 source preceding its sole final executable document
terminator is preserved verbatim.  Its preterminal SHA--256 digest is
\begin{center}\scriptsize\ttfamily
8f036c8bbd915f3868fec3618d639cc5bad2162c84e576f78228f11d2859c94e
\end{center}
The present material begins at Section~\ref{sec:v28-entry}; the final command
below is again unique.
\end{assessment}

\addtocontents{toc}{\protect\endgroup}
% =============================================================================
% END VERSION V28 MATERIAL -- append later versions before final terminator
% =============================================================================

% APPEND ALL LATER VERSIONS IMMEDIATELY ABOVE THE SOLE FINAL LINE BELOW.


% =============================================================================
% VERSION V29: NEW MATERIAL.  The complete V1--V28 source above is preserved
% verbatim; only its sole final executable document terminator was removed.
% =============================================================================

\clearpage
\makeatletter
\addtocontents{toc}{\protect\begingroup\protect\makeatletter}
\addtocontents{toc}{\protect\def\protect\l@section{\protect\@dottedtocline{1}{0em}{4.7em}}}
\addtocontents{toc}{\protect\def\protect\l@subsection{\protect\@dottedtocline{2}{1.5em}{5.3em}}}
\addtocontents{toc}{\protect\makeatother}
\makeatother
\renewcommand{\thesection}{V29.\arabic{section}}
\renewcommand{\thesubsection}{V29.\arabic{section}.\arabic{subsection}}
\renewcommand{\theequation}{V29.\arabic{equation}}
\setcounter{section}{0}
\setcounter{equation}{0}
\renewcommand{\theHsection}{V29.\arabic{section}}
\renewcommand{\theHsubsection}{V29.\arabic{section}.\arabic{subsection}}
\renewcommand{\theHtheorem}{V29.\arabic{section}.\arabic{theorem}}
\renewcommand{\theHequation}{V29.\arabic{equation}}

\begin{center}
{\LARGE\bfseries NS--A: Version V29}\\[0.5em]
{\large Positive-Time Analytic Saturation of the Raw Daughter,\\
Represented-Atlas Release, and the Corrected Fixed-Carrier Ledger}\\[0.5em]
{\normalsize 4 September 2026}
\end{center}
\addcontentsline{toc}{part}{Version V29: analytic raw-daughter saturation and represented-atlas release}

\begin{abstract}
Version~V28 constructed a canonical chronological Stokes filtration for the
support-corrected residual daughter $d^{\rm C}$.  Its Hilbert-space identities
are exact.  The present continuation moves one arrow upstream and separates
the raw Navier--Stokes daughter from the target-relative represented-source
short which produces $d^{\rm C}$.

On every positive-time smooth Navier--Stokes interval, the raw finite-head
daughter is real analytic in time.  A finite-dimensional analytic source has
the same Stokes-dephased support on every nonempty open subinterval.  Hence
the physical raw-daughter carrier saturates immediately at the entrance; it
has no interior source-direction births.  Its complete carrier is already
reconstructed by the time jets at any one interior time.

Delayed acquisitions of $d^{\rm C}=(I-R_{\rm C})d^{\rm raw}$ therefore do not
represent delayed births of the raw quadratic source.  They are caused by
the represented-atlas quotient: either the quotient creates a component
outside the actually occurring raw-source carrier, or it releases a raw
component which had previously been assigned to the represented part.  A
smooth finite counterexample shows that a moving represented projection can
produce an arbitrarily late residual-carrier jump from a constant analytic
raw source.  The V28 promotion and future-born transfer remain exact, but on
such a history they are retyped as residual-visibility promotion and
pre-release transfer, not as source-birth quantities.

Projecting the exterior adjoint onto the fixed raw-source carrier gives a
no-remainder ledger with entrance stock, viscosity, symmetric secant,
physical forcing, and source-null transfer.  On the source-supported branch,
a fixed residual scalar with vanishing residual action forces this fixed
ledger to diverge; no chronological promotion term is needed.  Thus the
first remaining obstruction is represented-atlas support/release or one of
the already existing fixed-carrier physical works.  No new source class,
response graph, compactness space, Dini modulus, or global regularity claim is
introduced.
\end{abstract}

% =============================================================================
\section{The raw daughter precedes the represented-source quotient}
\label{sec:v29-entry}
% =============================================================================

Retain a canonical finite head on one smooth terminal interval
$I=[a,b]\Subset(0,T_*)$:
\begin{equation}
 P_K=\one_{(0,K]}(\Lambda),\qquad
 Q_K=I-P_K,\qquad
 v(t)=P_Ku(t).
 \label{eq:v29-head}
\end{equation}
The occurring pure head-generated exterior daughter is
\begin{equation}
 \boxed{
 d^{\rm raw}(t)
 :=A^{-1/4}Q_K\mathcal B(v(t)).}
 \label{eq:v29-raw-daughter}
\end{equation}
By the inherited additive-support theorem, it is carried by the first
exterior octave
\begin{equation}
 \Pi_1=\one_{(K,2K]}(\Lambda).
 \label{eq:v29-daughter-octave}
\end{equation}

Let $R_{\rm C}(t)$ be the pointwise source-cyclic represented projection of
Version~V25.  The target-relative support-corrected residual used in
Versions~V26--V28 is
\begin{equation}
 \boxed{
 d^{\rm C}(t)=(I-R_{\rm C}(t))d^{\rm raw}(t).}
 \label{eq:v29-corrected-residual}
\end{equation}
The represented component is
\begin{equation}
 d^{\rm rep,C}(t)=R_{\rm C}(t)d^{\rm raw}(t).
 \label{eq:v29-represented-component}
\end{equation}
At almost every time,
\begin{equation}
 \boxed{
 d^{\rm raw}=d^{\rm rep,C}+d^{\rm C},\qquad
 d^{\rm rep,C}\perp d^{\rm C},\qquad
 \|d^{\rm raw}\|_2^2
 =\|d^{\rm rep,C}\|_2^2+\|d^{\rm C}\|_2^2.}
 \label{eq:v29-pointwise-source-Pythagoras}
\end{equation}

\begin{firewall}[The order of the two constructions]
\label{fire:v29-raw-before-short}
The field $d^{\rm raw}$ is an actual quadratic source of the smooth
Navier--Stokes trajectory.  The residual $d^{\rm C}$ is its orthogonal
innovation relative to a declared represented atlas.  The latter is an exact
target-relative source quotient, but its time-dependent projection may rotate
or release directions.  Source chronology must therefore be tested first on
$d^{\rm raw}$.  Chronology of $d^{\rm C}$ is a chronology of residual
visibility unless an independent theorem identifies the represented short
with a fixed physical source decomposition.
\end{firewall}

% =============================================================================
\section{Positive-time analyticity of the finite-head daughter}
\label{sec:v29-time-analyticity}
% =============================================================================

We first record the parabolic fact which was not used in Version~V28.
All Sobolev spaces below are the real mean-zero divergence-free periodic
spaces; their complexifications are used only inside the proof.

\begin{theorem}[Positive-time time analyticity of a smooth periodic trajectory]
\label{thm:v29-time-analyticity}
Let $u$ be a smooth solution of
\begin{equation}
 \partial_tu+\nu Au+\mathcal B(u)=0
 \label{eq:v29-NSE}
\end{equation}
on $(t_-,t_+)\times\Torus^3$, with $\nu>0$.  For every integer $m\ge4$,
the map
\begin{equation}
 t\longmapsto u(t)
 \label{eq:v29-time-map}
\end{equation}
is real analytic from $(t_-,t_+)$ into $H^m_\sigma$.
\end{theorem}

\begin{proof}
Fix $m\ge4$ and an interior compact interval $J\Subset(t_-,t_+)$.  Put
$X=H^{m-1}_\sigma$ and $X_{1/2}=H^m_\sigma$.  The Stokes operator is positive
self-adjoint on $X$ and, for every $0<\vartheta<\pi/2$, its complex analytic
semigroup satisfies
\begin{equation}
 \|e^{-\nu zA}\|_{X_{1/2}\to X_{1/2}}\le1,
 \qquad
 \|e^{-\nu zA}\|_{X\to X_{1/2}}
 \le C_\vartheta(\nu|z|)^{-1/2}
 \label{eq:v29-complex-semigroup}
\end{equation}
whenever $z\ne0$ and $|\arg z|\le\vartheta$.  The first estimate follows
mode by mode from $|e^{-\nu z|k|^2}|\le1$ when $\operatorname{Re}z\ge0$.
For the second, maximize
$r^{1/2}e^{-\nu(\operatorname{Re}z)r}$ over $r\ge0$ and use
$\operatorname{Re}z\ge|z|\cos\vartheta$.

Since $m\ge4$, the periodic product estimate gives the complex-bilinear bound
\begin{equation}
 \|\mathcal B(f,g)\|_{H^{m-1}}
 \le C_m\|f\|_{H^m}\|g\|_{H^m}.
 \label{eq:v29-bilinear-Hm}
\end{equation}
Choose $t_*\in J$ and, for $\rho>0$, let
\[
 \Omega_{\rho,\vartheta}(t_*)
 =\{t_*+z:0<|z|<\rho,\ |\arg z|<\vartheta\}.
\]
On bounded holomorphic $H^m_\sigma$-valued functions on this sector, continuous
at $t_*$, define
\begin{equation}
 \begin{aligned}
 (\Phi U)(t_*+z)
 :={}&e^{-\nu zA}u(t_*)\\
 &-z\int_0^1
 e^{-\nu(1-s)zA}
 \mathcal B\bigl(U(t_*+sz),U(t_*+sz)\bigr)\,\dd s.
 \end{aligned}
 \label{eq:v29-complex-mild-map}
\end{equation}
Equations \eqref{eq:v29-complex-semigroup}--\eqref{eq:v29-bilinear-Hm}
show that, on the ball $\|U\|_\infty\le M$,
\begin{equation}
 \|\Phi U-e^{-\nu(\cdot-t_*)A}u(t_*)\|_\infty
 \le C\rho^{1/2}M^2,
 \label{eq:v29-contraction-self}
\end{equation}
where the factor $\rho^{1/2}$ is
$\rho\int_0^1((1-s)\rho)^{-1/2}\dd s=2\rho^{1/2}$.
The same calculation gives
\begin{equation}
 \|\Phi U-\Phi V\|_\infty
 \le C\rho^{1/2}M\|U-V\|_\infty.
 \label{eq:v29-contraction-difference}
\end{equation}
Take $M=2\sup_{t\in J}\|u(t)\|_{H^m}+1$ and then choose $\rho_0>0$,
uniformly on $J$, so that the right sides of
\eqref{eq:v29-contraction-self}--\eqref{eq:v29-contraction-difference}
map the ball into itself and have contraction constant below one.  The fixed
point is holomorphic, and its restriction to the positive real ray is the
unique mild solution of \eqref{eq:v29-NSE} issuing from $u(t_*)$.

For any $t_0\in\operatorname{int}J$, choose $t_*<t_0$ with
$0<t_0-t_*<\rho_0/4$.  The sector based at $t_*$ contains a complex disk
centred at $t_0$ of positive radius.  Uniqueness identifies its real
restriction with the given smooth trajectory.  Thus $u$ is real analytic at
$t_0$ in $H^m_\sigma$.  Since $t_0$ was arbitrary, the theorem follows.
\end{proof}

\begin{corollary}[Analyticity of the raw finite-head daughter]
\label{cor:v29-raw-daughter-analytic}
For the fixed head \eqref{eq:v29-head},
\begin{equation}
 \boxed{
 t\longmapsto d^{\rm raw}(t)
 \text{ is real analytic on }(a,b)
 \text{ in the finite-dimensional space }\operatorname{Ran}\Pi_1.}
 \label{eq:v29-raw-analytic}
\end{equation}
\end{corollary}

\begin{proof}
The map $P_Ku(t)$ is analytic in every $H^m$ by
Theorem~\ref{thm:v29-time-analyticity}.  The quadratic map
$v\mapsto\mathcal B(v)$ is a continuous polynomial from $H^m$ to
$H^{m-1}$, and $A^{-1/4}Q_K$ is bounded on the finite output set.  Composition
preserves real analyticity.  The inherited additive-support theorem gives
\eqref{eq:v29-daughter-octave}.
\end{proof}

% =============================================================================
\section{Analytic sources saturate every open time interval}
\label{sec:v29-analytic-saturation}
% =============================================================================

The next statement is finite dimensional and independent of Navier--Stokes.
It is the exact reason the raw physical chronology has no interior direction
births.

\begin{theorem}[Analytic Stokes-support saturation]
\label{thm:v29-analytic-support-saturation}
Let $E$ be a finite-dimensional real Hilbert space, let
\begin{equation}
 A_E=\sum_{\lambda\in\Sigma}\lambda E_\lambda
 \label{eq:v29-finite-spectral-resolution}
\end{equation}
be self-adjoint, and let $f\in L^2(a,b;E)$ be real analytic on $(a,b)$.  For
$a<t\le b$, define
\begin{equation}
 \mathbb Q_f^{A}(t)
 :=\sum_{\lambda\in\Sigma}
 E_\lambda\left(\int_a^t f(s)\otimes f(s)\,\dd s\right)E_\lambda.
 \label{eq:v29-analytic-cumulative-Gram}
\end{equation}
Then
\begin{equation}
 \boxed{
 \ker\mathbb Q_f^{A}(t)=\ker\mathbb Q_f^{A}(b)
 \qquad(a<t\le b).}
 \label{eq:v29-kernel-saturation}
\end{equation}
Equivalently, the support projection of
$\mathbb Q_f^{A}(t)$ is independent of $t>a$.  Its right limit at $a$ is
already the final support.
\end{theorem}

\begin{proof}
Monotonicity gives
$\ker\mathbb Q_f^{A}(b)\subseteq\ker\mathbb Q_f^{A}(t)$.
Conversely, if $x\in\ker\mathbb Q_f^{A}(t)$, then positivity yields
\begin{equation}
 0
 =\langle x,\mathbb Q_f^{A}(t)x\rangle
 =\sum_{\lambda\in\Sigma}\int_a^t
   |\langle f(s),E_\lambda x\rangle|^2\,\dd s.
 \label{eq:v29-zero-analytic-Gram}
\end{equation}
Every scalar function
$s\mapsto\langle f(s),E_\lambda x\rangle$ is real analytic.  Equation
\eqref{eq:v29-zero-analytic-Gram} makes it zero almost everywhere on the
nonempty open interval $(a,t)$, hence zero there by continuity and then zero
on all of $(a,b)$ by the analytic identity theorem.  Therefore
$x\in\ker\mathbb Q_f^{A}(b)$, proving equality.
\end{proof}

\begin{corollary}[Immediate saturation of the physical raw daughter]
\label{cor:v29-raw-immediate-saturation}
Let
\begin{equation}
 P_{\rm raw}
 :=\operatorname{supp}
 \sum_{\lambda\in\Sigma_K}
 E_\lambda\left(
  \int_a^b d^{\rm raw}(s)\otimes d^{\rm raw}(s)\,\dd s
 \right)E_\lambda.
 \label{eq:v29-raw-carrier}
\end{equation}
Then, for every $t>a$, the cumulative raw-daughter support on $(a,t)$ is
exactly $P_{\rm raw}$.  In particular,
\begin{equation}
 \boxed{
 P_{\rm raw}(t)=P_{\rm raw}
 \quad(a\le t\le b)
 }
 \label{eq:v29-raw-saturates-entrance}
\end{equation}
for the right-continuous chronological convention, and the only possible
raw-source acquisition is the entrance acquisition.
\end{corollary}

\begin{proof}
Apply Theorem~\ref{thm:v29-analytic-support-saturation} to
$f=d^{\rm raw}$ and the finite Stokes spectral resolution on
$\operatorname{Ran}\Pi_1$.  The value at $a$ is the right limit of the
constant supports for $t>a$.
\end{proof}

\begin{theorem}[One-time jet reconstruction of the raw carrier]
\label{thm:v29-jet-carrier}
For every $t_0\in(a,b)$,
\begin{equation}
 \boxed{
 \operatorname{Ran}P_{\rm raw}
 =\bigoplus_{\lambda\in\Sigma_K}
   \operatorname{span}
   \left\{
    E_\lambda\partial_t^j d^{\rm raw}(t_0):j\ge0
   \right\}.}
 \label{eq:v29-jet-carrier}
\end{equation}
Because the target is finite dimensional, a finite set of derivatives
already spans the right side.  No bound on the largest required derivative
order follows from the carrier rank alone.
\end{theorem}

\begin{proof}
Fix $\lambda$ and let $M_\lambda$ be the span of
$E_\lambda d^{\rm raw}(t)$ over $t\in(a,b)$.  A vector $x\in E_\lambda E$
is orthogonal to every time jet at $t_0$ exactly when the analytic scalar
\[
 t\longmapsto\langle E_\lambda d^{\rm raw}(t),x\rangle
\]
has all derivatives zero at $t_0$.  It then vanishes near $t_0$, hence on the
complete connected interval.  Thus the orthogonal complement of the jet
span equals $M_\lambda^\perp$, proving equality.  Finite dimensionality makes
the increasing spans of finite jet sets stabilize.  The example
$f_N(t)=e_1+(t-t_0)^Ne_2$ shows that the first derivative exposing $e_2$ can
have arbitrarily large order even though the final rank is two.
\end{proof}

\begin{assessment}[No interior physical direction birth on a smooth card]
\label{ass:v29-no-interior-raw-birth}
A raw daughter mode may vanish at isolated times, and its amplitude may be
arbitrarily small near the entrance.  These facts do not create an interior
support birth.  Every raw Stokes direction which occurs anywhere on the
connected positive-time card is already present in the span of every
nonempty open initial subinterval.  What may remain difficult is its
quantitative conditioning, not its qualitative first occurrence.
\end{assessment}

% =============================================================================
\section{Delayed residual acquisition is represented-atlas release}
\label{sec:v29-atlas-release}
% =============================================================================

Let $P_{\rm chr}^{\rm C}(t)$ and $\Delta_\tau$ denote the V28 chronological
filtration and its jump projections, now emphasizing that they are formed
from $d^{\rm C}$ rather than from $d^{\rm raw}$.  For every delayed jump
$\tau>a$, put
\begin{equation}
 q_\tau(t)=\Delta_\tau d^{\rm C}(t),\qquad
 r_\tau(t)=\Delta_\tau d^{\rm rep,C}(t),\qquad
 f_\tau(t)=\Delta_\tau d^{\rm raw}(t).
 \label{eq:v29-three-direction-components}
\end{equation}

\begin{theorem}[Exact release identity at every delayed residual acquisition]
\label{thm:v29-release-identity}
For every delayed V28 acquisition time $\tau>a$,
\begin{equation}
 \boxed{
 f_\tau=q_\tau+r_\tau
 \quad\text{on }I,}
 \label{eq:v29-direction-source-sum}
\end{equation}
and
\begin{equation}
 \boxed{
 q_\tau(t)=0,
 \qquad
 f_\tau(t)=r_\tau(t)
 \quad\text{for almost every }t<\tau.}
 \label{eq:v29-pre-release-equality}
\end{equation}
For $\psi_\tau=\Delta_\tau\psi$, define
\begin{align}
 h_\tau^{\rm raw}
 &:=\int_a^b\langle f_\tau(t),\psi_\tau(t)\rangle\,\dd t,
 \label{eq:v29-raw-direction-scalar}\\
 h_\tau^{\rm rep}
 &:=\int_a^b\langle r_\tau(t),\psi_\tau(t)\rangle\,\dd t,
 \label{eq:v29-rep-direction-scalar}\\
 h_\tau^{\rm C}
 &:=\int_\tau^b\langle q_\tau(t),\psi_\tau(t)\rangle\,\dd t.
 \label{eq:v29-residual-direction-scalar}
\end{align}
Then
\begin{equation}
 \boxed{
 h_\tau^{\rm C}=h_\tau^{\rm raw}-h_\tau^{\rm rep}.}
 \label{eq:v29-release-scalar-ledger}
\end{equation}
Thus a delayed component of the support-corrected residual is a signed
release from the represented-source allocation, not a delayed occurrence of
the raw quadratic source.
\end{theorem}

\begin{proof}
Equation \eqref{eq:v29-direction-source-sum} is the projection of
\eqref{eq:v29-pointwise-source-Pythagoras}.  The V28 source-first theorem for
$d^{\rm C}$ gives $q_\tau=0$ before $\tau$, which yields
\eqref{eq:v29-pre-release-equality}.  Pair
\eqref{eq:v29-direction-source-sum} with $\psi_\tau$, integrate, and use the
pre-acquisition vanishing to replace the integral defining the residual term
by $[\tau,b]$.
\end{proof}

\begin{theorem}[Off-raw quotient support and on-raw represented release]
\label{thm:v29-on-off-raw}
The support-corrected residual has the orthogonal decomposition
\begin{equation}
 d^{\rm C}=d^{\parallel}+d^{\perp},
 \qquad
 d^{\parallel}:=P_{\rm raw}d^{\rm C},
 \qquad
 d^{\perp}:=(I-P_{\rm raw})d^{\rm C},
 \label{eq:v29-on-off-definition}
\end{equation}
with
\begin{align}
 \boxed{
 d^{\parallel}
 =d^{\rm raw}-P_{\rm raw}R_{\rm C}d^{\rm raw},}
 \label{eq:v29-on-raw-release}\\
 \boxed{
 d^{\perp}
 =-(I-P_{\rm raw})R_{\rm C}d^{\rm raw}.}
 \label{eq:v29-off-raw-leakage}
\end{align}
Consequently,
\begin{equation}
 \boxed{
 \|d^{\rm C}\|_2^2
 =\|d^{\parallel}\|_2^2+\|d^{\perp}\|_2^2.}
 \label{eq:v29-on-off-action}
\end{equation}
For the selected scalar,
\begin{equation}
 \boxed{
 \mathfrak h^{\rm C}
 =\mathfrak h^{\parallel}+\mathfrak h^{\perp},}
 \label{eq:v29-on-off-scalar}
\end{equation}
where
\begin{align}
 \mathfrak h^{\parallel}
 &:=\int_I
 \langle d^{\parallel},P_{\rm raw}\psi\rangle\,\dd t,
 \label{eq:v29-on-raw-scalar}\\
 \mathfrak h^{\perp}
 &:=\int_I
 \langle d^{\perp},(I-P_{\rm raw})\psi\rangle\,\dd t.
 \label{eq:v29-off-raw-scalar}
\end{align}
The term $d^{\perp}$ is entirely generated by the represented-source
projection; it has no counterpart in the actually occurring raw-source
carrier.  If $R_{\rm C}(t)$ reduces $P_{\rm raw}$ almost everywhere, then
$d^{\perp}=0$.
\end{theorem}

\begin{proof}
Since $P_{\rm raw}d^{\rm raw}=d^{\rm raw}$, applying $P_{\rm raw}$ and its
complement to \eqref{eq:v29-corrected-residual} gives
\eqref{eq:v29-on-raw-release}--\eqref{eq:v29-off-raw-leakage}.
Orthogonal Pythagoras gives \eqref{eq:v29-on-off-action}.  Insert
$I=P_{\rm raw}+(I-P_{\rm raw})$ in both arguments of the scalar pairing; the
cross terms vanish because the two ranges are orthogonal.  If
$[R_{\rm C},P_{\rm raw}]=0$, then $R_{\rm C}d^{\rm raw}\in\operatorname{Ran}P_{\rm raw}$,
which kills \eqref{eq:v29-off-raw-leakage}.
\end{proof}

\begin{corollary}[Analytic represented reads eliminate delayed residual acquisition]
\label{cor:v29-analytic-represented-short}
Suppose the represented source read
\begin{equation}
 t\longmapsto R_{\rm C}(t)d^{\rm raw}(t)
 \label{eq:v29-represented-read-path}
\end{equation}
is real analytic on $(a,b)$ in the finite daughter space.  Then $d^{\rm C}$
is real analytic, and its V28 chronological support saturates immediately:
\begin{equation}
 \boxed{
 P_{\rm chr}^{\rm C}(t)=P_{\rm h}^{\rm C}
 \quad(a\le t\le b).}
 \label{eq:v29-residual-immediate-saturation}
\end{equation}
In particular, this holds when $R_{\rm C}$ is a fixed projection, or when it
is a real-analytic operator-valued projection path.
\end{corollary}

\begin{proof}
Equation \eqref{eq:v29-corrected-residual} is the difference of two analytic
finite-dimensional paths.  Apply
Theorem~\ref{thm:v29-analytic-support-saturation}.
\end{proof}

\begin{corollary}[A delayed jump is a nonanalytic atlas-release witness]
\label{cor:v29-delayed-nonanalytic}
If the V28 residual filtration has a delayed nonzero jump at $\tau>a$, then
there exists $0\ne\xi\in\operatorname{Ran}\Delta_\tau$ such that
\begin{equation}
 \langle d^{\rm C}(t),\xi\rangle=0
 \quad\text{for almost every }t<\tau,
 \label{eq:v29-left-zero}
\end{equation}
but the same scalar is nonzero on a set of positive measure after $\tau$.
It cannot extend real analytically through $\tau$.  Since
$\langle d^{\rm raw}(t),\xi\rangle$ is analytic, the represented read
$\langle R_{\rm C}(t)d^{\rm raw}(t),\xi\rangle$ agrees with it before
$\tau$ and ceases to agree after $\tau$.  This is a literal represented-atlas
release.
\end{corollary}

\begin{proof}
A nonzero jump projection contributes positive trace to the post-$\tau$
cumulative residual Gram, so some $\xi$ in its range has nonzero residual
$L^2$ mass after $\tau$.  The pre-$\tau$ vanishing is
\eqref{eq:v28-component-not-before-birth}.  An analytic scalar which vanishes
on a nonempty open interval vanishes identically, proving the nonanalyticity
claim.  The final statement is
$d^{\rm C}=d^{\rm raw}-R_{\rm C}d^{\rm raw}$.
\end{proof}

% =============================================================================
\section{A smooth quotient can manufacture a delayed residual carrier}
\label{sec:v29-counterexample}
% =============================================================================

\begin{countertheorem}[Smooth represented-atlas motion creates a false delayed source birth]
\label{cth:v29-smooth-atlas-release}
Let $E=\mathbb R^2$, $A=I$, $I=[0,1]$, and let the raw source be the constant
analytic path
\begin{equation}
 d^{\rm raw}(t)=e_1.
 \label{eq:v29-counter-raw}
\end{equation}
Set $\tau=1/2$ and
\begin{equation}
 \theta(t)=
 \begin{cases}
 0,&t\le\tau,\\
 \exp\!\bigl(-(t-\tau)^{-2}\bigr),&t>\tau,
 \end{cases}
 \qquad
 v(t)=(\cos\theta(t),\sin\theta(t)),
 \label{eq:v29-flat-angle}
\end{equation}
and let $R_{\rm C}(t)=v(t)\otimes v(t)$.  Then $R_{\rm C}$ is a smooth
orthogonal-projection path, while
\begin{equation}
 \boxed{
 d^{\rm C}(t)
 =(I-R_{\rm C}(t))e_1
 =\sin^2\theta(t)e_1
  -\sin\theta(t)\cos\theta(t)e_2.}
 \label{eq:v29-counter-residual}
\end{equation}
The residual vanishes identically before $\tau$ and is nonzero immediately
after $\tau$.  Its right-continuous cumulative support jumps at $\tau$, even
though the raw source has been constant from the entrance.  For any
continuous terminal-zero writer with $\|\psi(\tau)\|=M$, the associated V28
promotion can be of order $M^2$ while the raw source action is fixed.
\end{countertheorem}

\begin{proof}
The flat function in \eqref{eq:v29-flat-angle} is $C^\infty$ and all of its
right derivatives vanish at $\tau$, so $R_{\rm C}$ is smooth.  Orthogonal
projection of $e_1$ onto $v$ is
$\cos\theta\,v=\cos^2\theta e_1+
\sin\theta\cos\theta e_2$, giving
\eqref{eq:v29-counter-residual}.  The path is zero for $t\le\tau$ and
nonzero for every $t>\tau$.  Since its direction
$\sin\theta e_1-\cos\theta e_2$ changes with $t$, every right neighborhood of
$\tau$ generates the final residual span.  Hence the right-continuous support
has a delayed jump.  Choosing, for example, a continuous $\psi_M$ with
$\psi_M(1)=0$ and $\|\psi_M(\tau)\|=M$ makes the jump contribution
$\frac12\|\Delta_\tau\psi_M(\tau)\|^2$ grow quadratically, while
$\int_0^1\|d^{\rm raw}\|^2\dd t=1$ is unchanged.
\end{proof}

\begin{assessment}[What the counterexample corrects]
\label{ass:v29-counterexample-scope}
The example does not invalidate the V28 projection or ledger identities.
It invalidates only the interpretation of every delayed residual-carrier
jump as a delayed physical source-direction birth.  The jump may be produced
entirely by a moving represented-source quotient, even when the raw source is
constant and analytic.
\end{assessment}

% =============================================================================
\section{The corrected fixed raw-carrier adjoint ledger}
\label{sec:v29-raw-ledger}
% =============================================================================

Return to the actual exterior adjoint
\begin{equation}
 -\partial_t\psi+\nu A\psi+\mathscr L_1(t)^*\psi=F_1,
 \qquad \psi(b)=0,
 \label{eq:v29-adjoint}
\end{equation}
where $F_1$ is the already assembled physical forcing.  Use the fixed raw
carrier \eqref{eq:v29-raw-carrier} and put
\begin{equation}
 \chi_{\rm raw}=P_{\rm raw}\psi,
 \qquad
 \eta_{\rm raw}=(\Pi_1-P_{\rm raw})\psi.
 \label{eq:v29-raw-adjoint-split}
\end{equation}
Because $P_{\rm raw}$ reduces $A$, no heat-range current and no chronological
promotion occur.

Write
\begin{equation}
 \mathscr L_1^*=\mathscr S_1-\mathscr K_1,
 \qquad
 \mathscr S_1^*=\mathscr S_1,
 \qquad
 \mathscr K_1^*=-\mathscr K_1.
 \label{eq:v29-secant-skew}
\end{equation}
Projection of \eqref{eq:v29-adjoint} gives
\begin{equation}
 \boxed{
 \begin{aligned}
 -\partial_t\chi_{\rm raw}
 &+\nu A\chi_{\rm raw}
 +P_{\rm raw}\mathscr S_1P_{\rm raw}\chi_{\rm raw}
 -P_{\rm raw}\mathscr K_1P_{\rm raw}\chi_{\rm raw}
 \\
 &=P_{\rm raw}F_1
   -P_{\rm raw}\mathscr L_1^*\eta_{\rm raw}.
 \end{aligned}}
 \label{eq:v29-raw-projected-adjoint}
\end{equation}

Define
\begin{align}
 \mathcal P_{\rm raw}
 &:=\frac12\|\chi_{\rm raw}(a)\|_2^2
   +\nu\int_I\|A^{1/2}\chi_{\rm raw}\|_2^2\,\dd t,
 \label{eq:v29-raw-positive-stock}\\
 \mathcal N_{\rm raw}
 &:=\int_I
   \langle\mathscr S_1\chi_{\rm raw},
          \chi_{\rm raw}\rangle\,\dd t,
 \label{eq:v29-raw-secant}\\
 \mathcal N_{\rm raw}^-
 &:=\int_I
   \bigl(-\langle\mathscr S_1\chi_{\rm raw},
                    \chi_{\rm raw}\rangle\bigr)_+\,\dd t,
 \label{eq:v29-raw-negative-secant}\\
 \mathcal W_{\rm raw}^{\rm for}
 &:=\int_I\langle F_1,\chi_{\rm raw}\rangle\,\dd t,
 \label{eq:v29-raw-forcing}\\
 \mathcal W_{\rm raw}^{\rm tr}
 &:=-\int_I
 \langle\eta_{\rm raw},\mathscr L_1\chi_{\rm raw}\rangle\,\dd t.
 \label{eq:v29-raw-transfer}
\end{align}

\begin{theorem}[Exact fixed raw-carrier ledger]
\label{thm:v29-raw-ledger}
The physical raw-source carrier satisfies
\begin{equation}
 \boxed{
 \mathcal P_{\rm raw}+\mathcal N_{\rm raw}
 =\mathcal W_{\rm raw}^{\rm for}
  +\mathcal W_{\rm raw}^{\rm tr}.}
 \label{eq:v29-raw-ledger}
\end{equation}
If $\mathcal P_{\rm raw}>0$, then at least one of
\begin{equation}
 \boxed{
 \mathcal N_{\rm raw}^-
 \ge\frac13\mathcal P_{\rm raw},
 \quad
 |\mathcal W_{\rm raw}^{\rm tr}|
 \ge\frac13\mathcal P_{\rm raw},
 \quad
 |\mathcal W_{\rm raw}^{\rm for}|
 \ge\frac13\mathcal P_{\rm raw}}
 \label{eq:v29-raw-one-third}
\end{equation}
holds.
\end{theorem}

\begin{proof}
Take the real inner product of
\eqref{eq:v29-raw-projected-adjoint} with $\chi_{\rm raw}$ and integrate from
$a$ to $b$.  The terminal value vanishes and the internal skew term is zero,
which gives \eqref{eq:v29-raw-ledger}.  If every inequality in
\eqref{eq:v29-raw-one-third} failed strictly, then
\[
 \mathcal P_{\rm raw}+\mathcal N_{\rm raw}
 \ge\mathcal P_{\rm raw}-\mathcal N_{\rm raw}^-
 >\frac23\mathcal P_{\rm raw},
\]
whereas
\[
 |\mathcal W_{\rm raw}^{\rm for}
   +\mathcal W_{\rm raw}^{\rm tr}|
 <\frac23\mathcal P_{\rm raw},
\]
contradicting the ledger.  Passing to non-strict inequalities proves the
stated alternative.
\end{proof}

\begin{theorem}[A fixed on-raw residual scalar cannot be hidden by vanishing residual action]
\label{thm:v29-on-raw-nonsilent}
Consider a terminal sequence of the preceding cards.  Put
\begin{equation}
 \mathcal A_n^{\parallel}
 :=\int_{I_n}\|d_n^{\parallel}(t)\|_2^2\,\dd t.
 \label{eq:v29-on-raw-action}
\end{equation}
If
\begin{equation}
 |\mathfrak h_n^{\parallel}|
 \ge\delta\mathcal R_n^2,
 \qquad
 \frac{\mathcal A_n^{\parallel}}
      {\nu\mathcal R_n^2}\longrightarrow0,
 \label{eq:v29-fixed-on-raw-vanishing-action}
\end{equation}
then
\begin{equation}
 \boxed{
 \int_{I_n}\|\chi_{{\rm raw},n}(t)\|_2^2\,\dd t
 \ge
 \frac{\delta^2\mathcal R_n^4}
      {\mathcal A_n^{\parallel}},}
 \label{eq:v29-raw-adjoint-lower}
\end{equation}
and, since $P_{{\rm raw},n}\le\Pi_{1,n}$,
\begin{equation}
 \boxed{
 \mathcal P_{{\rm raw},n}
 \ge
 \nu K_n^2
 \frac{\delta^2\mathcal R_n^4}
      {\mathcal A_n^{\parallel}}
 \longrightarrow\infty.}
 \label{eq:v29-raw-stock-diverges}
\end{equation}
Consequently, after passage to a subsequence, one of the three quantities in
\eqref{eq:v29-raw-one-third} diverges.  No promotion or future-born-source
term is needed.
\end{theorem}

\begin{proof}
Because $d_n^{\parallel}\in\operatorname{Ran}P_{{\rm raw},n}$,
\[
 \mathfrak h_n^{\parallel}
 =\int_{I_n}\langle d_n^{\parallel},
              \chi_{{\rm raw},n}\rangle\,\dd t.
\]
Cauchy--Schwarz gives \eqref{eq:v29-raw-adjoint-lower}.  Every vector in the
first exterior octave has Stokes eigenvalue at least $K_n^2$, so
\[
 \nu\int\|A^{1/2}\chi_{{\rm raw},n}\|_2^2
 \ge\nu K_n^2\int\|\chi_{{\rm raw},n}\|_2^2.
\]
Combine this with \eqref{eq:v29-raw-adjoint-lower}.  The hypothesis on
$\mathcal A_n^{\parallel}$ and $K_n,\mathcal R_n\to\infty$ gives divergence.
The final assertion is Theorem~\ref{thm:v29-raw-ledger}.
\end{proof}

\begin{corollary}[The secant branch keeps the inherited material cell]
\label{cor:v29-raw-secant-cell}
On the first branch of \eqref{eq:v29-raw-one-third}, the inherited
scale-local secant estimate and material-cell test apply with
$\chi_{\rm raw}$ in place of the residual chronological writer.  Hence the
same selected time carries either
\begin{equation}
 \int_{\mathcal M_n(t_n^*)}
 |\nabla P_{\le4K_n}u(t_n^*)|^2
 \gtrsim\nu^2K_n
 \label{eq:v29-inside-cell-pulse}
\end{equation}
or the identical lower packet lies outside the already selected material
cell.  No replacement cell is fitted.
\end{corollary}

\begin{proof}
The proof of the V26--V27 scale-local secant theorem uses only that the
writer is supported in the first exterior octave and the negative symmetric
secant carries a fixed fraction of its positive stock.  Both properties hold
for $\chi_{\rm raw}$.  The material-cell split is the same exact partition of
the resulting finite-frequency enstrophy packet.
\end{proof}

% =============================================================================
\section{Retyping the V28 card and the corrected terminal alternative}
\label{sec:v29-retyping}
% =============================================================================

\begin{theorem}[Exact status of the V28 delayed-acquisition quantities]
\label{thm:v29-v28-retyping}
The V28 chronological identities remain valid for the support-corrected
residual.  On an actual positive-time smooth Navier--Stokes card, however,
the following terminology is forced.
\begin{enumerate}[label=\textup{(R\arabic*)},leftmargin=3.5em]
\item The raw quadratic daughter has immediate Stokes-support saturation by
      \eqref{eq:v29-raw-saturates-entrance}; it has no delayed physical
      direction birth.
\item A delayed V28 jump is a \emph{residual-visibility acquisition}.  Its
      positive term
      \begin{equation}
       \frac12\|\Delta_\tau\psi(\tau)\|_2^2
       \label{eq:v29-residual-promotion}
      \end{equation}
      is target-writer promotion into the released quotient direction, not a
      second source payment.
\item The V28 future-born transfer is a \emph{pre-release transfer}: it acts
      in a direction before that direction becomes unrepresented, not before
      the raw daughter exists.
\item If the represented read is analytic, all delayed residual acquisitions
      disappear and V28 reduces exactly to the fixed-carrier card.
\item If a delayed acquisition remains, the first upstream output is the
      represented-atlas release of
      Theorem~\ref{thm:v29-release-identity}, together with any off-raw
      support term \eqref{eq:v29-off-raw-leakage}.
\end{enumerate}
\end{theorem}

\begin{proof}
The first item is Corollary~\ref{cor:v29-raw-immediate-saturation}.  The next
two items are Theorem~\ref{thm:v29-release-identity} and the definitions of
the V28 promotion and future transfer.  The fourth item is
Corollary~\ref{cor:v29-analytic-represented-short}.  The last is
Theorem~\ref{thm:v29-on-off-raw}.
\end{proof}

\begin{theorem}[Version V29 corrected terminal alternative]
\label{thm:v29-terminal-alternative}
Retain every inherited material, boundary, target-normalization, old-tail,
represented-return, and finite-head gate.  On a positive-time smooth
finite-head card, after passage to a subsequence, the first applicable output
is one of the following.
\begin{enumerate}[label=\textup{(V29.\arabic*)},leftmargin=5.9em]
\item \emph{Raw daughter null.}
      The occurring head generates no exterior daughter on the selected
      history.
\item \emph{Off-raw quotient support.}
      The selected residual scalar or action has a nonvanishing component in
      $(I-P_{\rm raw})d^{\rm C}$.  This component is generated by the
      represented-source quotient and is returned as a support-incidence
      residual.
\item \emph{Represented-atlas release.}
      The residual filtration has a delayed acquisition.  The raw source is
      already saturated, and the exact release scalar is
      \eqref{eq:v29-release-scalar-ledger}.  The V28 promotion and pre-release
      transfer are attached to this same represented-source release.
\item \emph{Nonvanishing on-raw residual action.}
      The target-relative innovation already carries a positive action on the
      actual raw-source carrier.
\item \emph{Fixed on-raw scalar with vanishing residual action.}
      The raw-carrier stock diverges and one of negative secant, assembled
      dynamic transfer, or assembled physical forcing survives by
      \eqref{eq:v29-raw-one-third}.  The secant branch lands in the original
      material cell or escapes it.
\item \emph{Target-native ancestry pass.}
      The first-birth scalar is negligible and the inherited old and
      represented target residuals vanish.  NS--A exports the scalar ancestry
      card without creating a delayed physical source birth.
\item \emph{All-shell raw-daughter escape.}
      No finite head captures the occurring raw source and its selected
      scalar.
\end{enumerate}
No branch is excluded for every fixed datum by the present note.
\end{theorem}

\begin{proof}
The raw daughter is either zero or has the immediate carrier
$P_{\rm raw}$.  Apply the orthogonal scalar and action split of
Theorem~\ref{thm:v29-on-off-raw}.  A nonzero off-raw component gives the
second branch.  On the complementary branch, delayed residual acquisition is
exactly the release alternative of
Theorem~\ref{thm:v29-release-identity}; if it is absent, the residual carrier
is fixed from the entrance.  Its action either remains nonvanishing or tends
to zero.  In the latter case, a fixed scalar gives
Theorem~\ref{thm:v29-on-raw-nonsilent} and the fixed-ledger alternatives.
The ancestry-pass and all-shell alternatives are the inherited terminal
exits once the finite card and scalar residuals are tested.
\end{proof}

\begin{firewall}[No source birth is inferred from a moving quotient]
\label{fire:v29-no-quotient-birth}
A support-corrected residual may be the correct source innovation for one
fixed target.  Its delayed appearance under a time-dependent represented
projection is not, without further incidence, a delayed occurrence of the
quadratic Navier--Stokes source.  The physical raw source, its represented
follower, and their signed difference remain one assembled history.  A
positive residual action or promotion is not priced as a second source birth.
\end{firewall}

% =============================================================================
\section{Version V29 integrated conclusion and proof-status ledger}
\label{sec:v29-conclusion}
% =============================================================================

\begin{theorem}[Version V29 integrated conclusion]
\label{thm:v29-integrated}
Preserving every theorem and limitation of Versions~V1--V28, Version~V29
establishes the following.
\begin{enumerate}[label=\textup{(V29.\arabic*)},leftmargin=5.9em]
\item Smooth periodic Navier--Stokes trajectories are real analytic in time on
      every positive-time Sobolev screen.
\item The raw finite-head daughter is analytic and its Stokes-dephased support
      saturates on every nonempty open initial interval.
\item One interior time-jet bank reconstructs the complete raw source carrier.
\item Delayed V28 acquisitions are represented-atlas releases or quotient-
      generated off-raw directions, not delayed raw-source births.
\item A smooth finite model produces a delayed residual carrier from a
      constant analytic raw source.
\item The fixed raw carrier has an exact no-promotion adjoint ledger.
\item A fixed on-raw residual scalar with vanishing residual action forces the
      raw-carrier stock and one physical ledger term to diverge.
\item The V28 promotion and future-born transfer are retyped as residual-
      visibility promotion and pre-release transfer.
\item The terminal alternatives are exactly those of
      Theorem~\ref{thm:v29-terminal-alternative}.
\end{enumerate}
No terminal Dini improvement, critical upper-action estimate, terminal trace,
backward uniqueness, Liouville theorem, or global regularity conclusion is
asserted.
\end{theorem}

\begin{proof}
Items~\textup{(V29.1)--(V29.3)} are
Theorem~\ref{thm:v29-time-analyticity},
Corollary~\ref{cor:v29-raw-daughter-analytic},
Theorem~\ref{thm:v29-analytic-support-saturation}, and
Theorem~\ref{thm:v29-jet-carrier}.  Item~\textup{(V29.4)} is
Theorems~\ref{thm:v29-release-identity} and
\ref{thm:v29-on-off-raw}.  Item~\textup{(V29.5)} is
Countertheorem~\ref{cth:v29-smooth-atlas-release}.  Items
\textup{(V29.6)--(V29.7)} are
Theorems~\ref{thm:v29-raw-ledger} and
\ref{thm:v29-on-raw-nonsilent}.  Item~\textup{(V29.8)} is
Theorem~\ref{thm:v29-v28-retyping}, and the final item is
Theorem~\ref{thm:v29-terminal-alternative}.
\end{proof}

\begingroup\small
\begin{longtable}{>{\raggedright\arraybackslash}p{0.27\textwidth}
                  >{\raggedright\arraybackslash}p{0.18\textwidth}
                  >{\raggedright\arraybackslash}p{0.47\textwidth}}
\caption{Version V29 proof-status ledger.}\label{tab:v29-ledger}\\
\toprule
\textbf{Row}&\textbf{Status}&\textbf{Exact advance or limitation}\\
\midrule
\endfirsthead
\toprule
\textbf{Row}&\textbf{Status}&\textbf{Exact advance or limitation}\\
\midrule
\endhead
V1--V28 prefix&\PROVED{} preserved
 &Complete V28 preterminal source retained byte-for-byte; no inherited theorem
 rewritten.\\[0.35em]
Positive-time time analyticity&\PROVED{} PDE
 &Holomorphic mild-solution construction on every interior Sobolev screen.\\[0.35em]
Raw daughter analyticity&\PROVED{} exact
 &Fixed-head quadratic daughter is a finite-dimensional analytic path.\\[0.35em]
Analytic support saturation&\PROVED{} exact
 &Every nonempty open initial interval generates the complete Stokes-dephased
 raw carrier.\\[0.35em]
One-time jet card&\PROVED{} finite
 &All time jets at any interior time span the raw carrier; a finite subset
 suffices, with no rank-only derivative-order bound.\\[0.35em]
V28 delayed physical births&\OBSTRUCTION{} removed
 &Delayed residual jumps are quotient release/support effects, not delayed
 occurrence of the raw daughter.\\[0.35em]
Smooth moving-atlas model&\OBSTRUCTION{} explicit
 &A constant analytic raw source acquires a delayed smooth residual carrier
 under a moving represented projection.\\[0.35em]
Off-raw quotient support&\PROVED{} exact residual
 &The component $(I-P_{\rm raw})d^{\rm C}$ is entirely projection generated.\\[0.35em]
Analytic represented short&\PROVED{} sufficient
 &Analytic represented read forces immediate residual saturation and exact
 recovery of the fixed-carrier branch.\\[0.35em]
Fixed raw-carrier ledger&\PROVED{} exact
 &Entrance plus viscous stock balances signed secant, assembled transfer, and
 assembled forcing, with no promotion term.\\[0.35em]
Fixed scalar / vanishing action&\PROVED{} non-silent
 &Raw-carrier adjoint stock diverges; secant, transfer, or forcing survives.\\[0.35em]
Material-cell landing&\PROVED{} inherited exactly
 &The secant pulse remains in the original material cell or gives literal
 outside-cell escape.\\[0.35em]
Represented-atlas release control&\OPEN{} physical incidence
 &Requires the actual same-history represented-source path or a signed release
 estimate; analyticity of the raw PDE source does not supply it.\\[0.35em]
Critical upper budget / regularity&\OPEN
 &No surviving physical ledger packet is excluded for every fixed datum.\\
\bottomrule
\end{longtable}
\endgroup

% =============================================================================
\section{Exact mandate after Version V29}
\label{sec:v29-next}
% =============================================================================

\begin{programme}[Evaluate the represented-atlas release before another source chronology]
\label{prog:v29-next}
Preserve Versions~V1--V29 verbatim.  Do not introduce another source
filtration, carrier completion, adjoint field, response graph, path law,
Moore--Penrose factorization, shell count, Dini modulus, or compactness space.

Proceed only in the following order.
\begin{enumerate}[label=\textup{(V30.\arabic*)},leftmargin=5.9em]
\item Build the raw daughter and its immediately saturated Stokes carrier
      before applying the represented-source short.
\item Evaluate the off-raw quotient component
      $(I-P_{\rm raw})d^{\rm C}$.  A surviving scalar or action is the exact
      source-support incidence residual.
\item On the on-raw branch, evaluate the direct signed release
      $d^{\rm raw}-P_{\rm raw}R_{\rm C}d^{\rm raw}$ and its old/represented
      scalar ledger.  Do not call its delayed visibility a source birth.
\item If the represented read is analytic or fixed, use the immediate
      fixed-carrier ledger; do not retain the V28 promotion notation.
\item If a nonanalytic release remains, pair its V28 promotion and pre-release
      transfer before taking any positive part, and export the assembled
      represented-atlas release as one ancestry obstruction.
\item On the raw fixed-carrier secant branch, keep the same selected time and
      inherited material cell.  On the transfer and forcing branches, retain
      complete assembly before opening components.
\item If the first-birth scalar is negligible and the inherited old and
      represented target residuals pass, terminate NS--A and export the
      target-native scalar ancestry card.
\item If no finite head captures the raw daughter, retain fused all-shell raw-
      source escape.
\end{enumerate}
No later compactness, terminal-trace, backward-uniqueness, or Liouville
argument is opened before this represented-atlas release card is resolved.
\end{programme}

\begin{assessment}[Programme position after Version V29]
\label{ass:v29-position}
The first unresolved finite object is no longer the five-entry chronological
birth card of Version~V28.  On an actual positive-time smooth trajectory the
raw source carrier is fixed at the entrance.  The first remaining card is
\begin{equation}
 \boxed{
 \left(
  \mathfrak h^{\perp},
  \|d^{\perp}\|_{L^2}^2,
  \mathfrak h^{\parallel},
  \|d^{\parallel}\|_{L^2}^2,
  \mathcal W_{\rm raw}^{\rm tr},
  \mathcal W_{\rm raw}^{\rm for},
  \mathcal N_{\rm raw}^-
 \right).}
 \label{eq:v29-first-unresolved-card}
\end{equation}
A delayed V28 residual jump is retained only as additional evidence that the
represented-source quotient releases or rotates a direction.  Repeating a
source-birth filtration after that point would be circular.
\end{assessment}

\paragraph{Version V29 history.}
Preserved the complete V28 source.  Returned upstream from the chronological
support-corrected residual to the raw finite-head quadratic daughter.  Proved
positive-time Sobolev time analyticity by a complex-sector mild contraction,
then proved that every analytic finite-dimensional source saturates its
Stokes-dephased support on every nonempty open interval.  Derived immediate
raw-source saturation and one-time jet reconstruction.  Separated the
support-corrected residual into an on-raw represented release and an off-raw
quotient-support component, and proved the exact per-acquisition raw minus
represented scalar ledger.  Constructed a smooth moving-projection
counterexample showing that a delayed residual carrier can be manufactured
from a constant analytic raw source.  Retyped the V28 promotion and
future-born transfer accordingly.  Finally derived the fixed raw-carrier
adjoint ledger and its non-silent fixed-scalar/vanishing-action alternative.
No global regularity conclusion was claimed.

\begin{assessment}[Append-only integrity]
\label{ass:v29-integrity}
The complete Version~V28 source preceding its sole final executable document
terminator is preserved verbatim.  Its preterminal SHA--256 digest is
\begin{center}\scriptsize\ttfamily
3d1c8f04d67d38c6795fc54c55d5ef40d245da8a45d6e4a28687f5eb321ab2e1
\end{center}
The present material begins at Section~\ref{sec:v29-entry}; the final command
below is again unique.
\end{assessment}

\addtocontents{toc}{\protect\endgroup}
% =============================================================================
% END VERSION V29 MATERIAL -- append later versions before final terminator
% =============================================================================

% APPEND ALL LATER VERSIONS IMMEDIATELY ABOVE THE SOLE FINAL LINE BELOW.

\end{document}
